\PassOptionsToPackage{hypertexnames=false}{hyperref}
% Height-source replacement: see work/remove-vkk-20261104/REPORT.md for scope and review status.
\documentclass[11pt]{article}
\usepackage[T1]{fontenc}
\usepackage[utf8]{inputenc}
\usepackage{lmodern,amsmath,amssymb,amsthm,mathtools}
\usepackage[margin=1in]{geometry}
\usepackage{booktabs,longtable,array,microtype,enumitem}
\usepackage[hidelinks]{hyperref}
\setlist{itemsep=3pt,topsep=5pt}
\setlength{\emergencystretch}{2em}
\newtheorem{theorem}{Theorem}[section]
\newtheorem{proposition}[theorem]{Proposition}
\newtheorem{lemma}[theorem]{Lemma}
\newtheorem{corollary}[theorem]{Corollary}
\theoremstyle{definition}
\newtheorem{definition}[theorem]{Definition}
\theoremstyle{remark}
\newtheorem{remark}[theorem]{Remark}
\newcommand{\Q}{\mathbb Q}
\newcommand{\Z}{\mathbb Z}
\newcommand{\N}{\mathbb N}
\newcommand{\F}{\mathbb F}
\newcommand{\Cset}{\mathcal C}
\providecommand{\mathscr}[1]{\mathcal{#1}}
\newcommand{\Sel}{\operatorname{Sel}}
\newcommand{\Sha}{\operatorname{Sha}}
\newcommand{\ord}{\operatorname{ord}}
\newcommand{\Norm}{\operatorname{N}}
\newcommand{\Tr}{\operatorname{Tr}}
\newcommand{\Trd}{\operatorname{Trd}}
\newcommand{\Close}{\operatorname{Close}}
\newcommand{\Test}{\operatorname{Test}}
\newcommand{\Val}{\operatorname{Val}}
\newcommand{\tw}[1]{2\mathbin{\uparrow\!\uparrow}#1}
\providecommand{\LogE}{\log_E}
\makeatletter
\renewcommand*\l@subsection{\@dottedtocline{2}{1.5em}{3.2em}}
\renewcommand*\l@section{\@dottedtocline{1}{0em}{3em}}
\makeatother
\title{$\mathbb{Z}$ is Diophantine in $\mathbb{Q}$\\[4pt]\large Version 17}
\author{}
\date{November 13, 2026}
\providecommand{\C}{\mathbb C}
\providecommand{\hF}{h_{\mathrm F}}
\providecommand{\OK}{\mathcal O_K}
\providecommand{\rad}{\operatorname{rad}}
\providecommand{\Jac}{\operatorname{Jac}}
\providecommand{\Ardeg}{\widehat{\deg}}
\providecommand{\A}{\mathbb A}
\providecommand{\Pj}{\mathbb P}
\DeclareMathOperator{\End}{End}
\DeclareMathOperator{\divisor}{div}
\providecommand{\vp}{v_p}
\providecommand{\HH}{\mathcal H}
\DeclareMathOperator{\Nm}{Nm}
\ifdefined\pdftrailerid\pdftrailerid{}\fi
\usepackage{graphicx}
\hypersetup{pdftitle={Z is Diophantine in Q, Version 17},pdfauthor={Elliot Glazer and attributed contributors},bookmarksnumbered=true}
\setcounter{tocdepth}{2}\setcounter{secnumdepth}{3}
\allowdisplaybreaks[1]
\makeatletter\g@addto@macro\UrlBreaks{\do\a\do\b\do\c\do\d\do\e\do\f\do\0\do\1\do\2\do\3\do\4\do\5\do\6\do\7\do\8\do\9}\makeatother
\begin{document}
\renewcommand{\thesection}{I.\arabic{section}}
\renewcommand{\theHsection}{q214.\arabic{section}}
\maketitle
\begin{abstract}
We present the proposed $214$-variable quartic family obtained by reciprocal
clearing on the shared-trace $224$-coordinate predecessor.  For the fixed
coefficient $A_N=149^{2N+1}$, the polynomial is the sum of squares of $188$
primitive integral quadratic residuals, all printed below.  The proposed
arithmetic conclusion is uniform for every $N$ beyond one fixed ineffective
threshold.  The complete new local, modularity, height and existence arguments
are followed by the full current Version~16 mathematical foundation, including
the inherited elliptic, scalar, prime-selection, pole and aggregate arguments.
The exact coefficient calculation gives $2919$ coordinate monomials, $2920$
parameter terms and height $149^{4N+2}+1$ for $N\ge53$; this is only a
coefficient threshold.  This assembled manuscript is an author proposal for
fresh review and does not itself confer Protocol~D or release acceptance.
\end{abstract}
\tableofcontents
\clearpage
\part{The proposed 214-variable construction}
\section{The theorem, conventions and reading order}\label{q214:main}
\begin{theorem}[The proposed $214$-variable family]\label{q214:theorem}
For an integer $N\ge1$ put $A_N=149^{2N+1}$.  In coordinates
$X_1,\ldots,X_{214}$, let $r_{i,N}$ be the $188$ rows specified in
Section~\ref{q214:rows}, and set
\[
 P_N=\sum_{i=1}^{188}r_{i,N}^2\in\Z[X_1,\ldots,X_{214}].
\]
Subject to the exact source theorems and applications developed in this
manuscript, there exists one fixed integer $N_0\ge1$ such that
\[
 \forall N\ge N_0\ \forall x\in\Q,\qquad
 x\in\Z\ \Longleftrightarrow\
 \exists (y_2,\ldots,y_{214})\in\Q^{213}:
 P_N(x,y_2,\ldots,y_{214})=0.
\]
Every $P_N$ has degree four and content one, and every one of its $214$
coordinates occurs.  For $N\ge53$, it has $2919$ distinct coordinate
monomials, $2920$ parameter terms, and largest absolute coefficient
$A_N^2+1=149^{4N+2}+1$.
\end{theorem}
Here $N$ and $A_N$ are fixed parameters of each polynomial, rather than
existential coordinates.  The proof gives no effective $N_0$ and no numerical
defining member.  In particular, $N=53$ is not asserted to define $\Z$.
The coefficient statements and formal row identities have finite exact
certificates; those checks do not prove the arithmetic equivalence.

\paragraph{Scope of this edition.}
Part~I contains the current $214$-variable argument, with all selected proof
and source fragments incorporated as mathematical text.  Part~II contains
the entire current upstream Version~16 body, including its proofs and bibliography.
Its $222$- and $314$-variable theorems and earlier constructions retain their
own scope; they are not the polynomial of Theorem~\ref{q214:theorem}.
The $214$ branch starts at the raw $224$-coordinate system, before the
incompatible paired-root and trace-quotient reductions.  The complete inverse
maps and arithmetic foundations retained in Part~II supply precisely the
interfaces identified in Part~I.
The transition in Section~\ref{q214:foundation-transition} identifies the
Version~16 corrections and their current consumers.  Statements in the retained
author packets about byte identity with Version~15 describe their earlier
freeze; they do not assert that no later analytic correction is incorporated
here.  Version~16 is the upstream foundation's edition, not a new numbered
release of the proposed $214$ research manuscript.

\paragraph{Navigation to current foundation passages.}
The retained source has direct links to the
\hyperref[q214foundation:gi-normalization]{GI horizontal normalization},
\hyperref[q214foundation:gi-cover]{fixed-lift finite-cover comparison},
\hyperref[q214foundation:tz-adaptation]{Tao--Ziegler proof adaptation}, and
\hyperref[q214foundation:cmh-surface]{surface-height supplier with corrected citations}.
The original raw-$224$ inverse maps remain in
Sections~\ref{v222:ratio} and~\ref{v222:reference}; the formal-parameter
supplier is Section~\ref{v15:logs}.  Historical Version~15 labels retain
their names inside the current Version~16 text.

\paragraph{Heights, fields and labels.}
For coprime integers $u,v$ with $v>0$, write
$h(u/v)=\log\max(|u|,v)$.  All logarithms are natural.  Use the original
stable Faltings height throughout.  The arithmetic field and curve are
\[
 K=\Q(\sqrt2),\qquad d=3-\sqrt2,\qquad
 E:y^2=x^3-d^2x,\qquad P=2(1+\sqrt2,2\sqrt2).
\]
The four nonzero integer labels are $(n,a,t,b)$, with $c=|a|$ and
$B=\max(|n|,|a|,|t|,|b|)$.  The input test is $\beta=51x^2+2504$.
The physical scale is the positive quarter-trace
$T=\frac14\Tr(\sum_{i=1}^4 x([n_i]P))$, and the two detector arguments
are $t_0=\lambda$, $t_1=T\lambda$.  The contact set is
\[
 \Theta_2=\Theta\cup(2-\Theta),\qquad
 \Theta=\{0,-\tfrac12,-1,\tfrac1{2399},-\tfrac{32}{407}\}.
\]
The exceptional set $\Sigma$ consists of the primes at most $76176$ and
$2029147$.  For each scale, $M_j$ is the squarefree product of primes
outside $\Sigma$ at a positive odd contact of $t_j$.
The normal surface field $L'$, the coefficient field $E_f$, and the auxiliary
field splitting a modular evaluation fibre have distinct roles and are never
identified.

\paragraph{How the proof joins.}
The assembly map in Section~\ref{q214:assembly} identifies each source,
constructed object and consumer.  Reciprocal clearing reconstructs the
physical raw-$224$ objects before logarithmic ratios are consumed.  The
changed factors give $M_j^3\mid\Delta$ and $M_j\mid J$ with no extra scalar.
The full-family height bound, applied to both transformed arguments, gives
$h(t_j)\le H_2M_j^{7/6}$.  The actual aggregate \emph{sum} theorem then
gives $B^2\le C|J_0|^{7/(6k)}+M_2$ when $M_j^k\mid J_0\ne0$ for both
scales.  Scalar growth and the Lucas bound force $\Delta=J=0$ for one
uniform ineffective threshold.  The completeness argument orders scalar
preparation, labels, prime selection, CRT and fresh pole witnesses to obtain
a single threshold valid for every integer input.

\paragraph{Operative averaging and retained historical text.}
The current averaging argument is the published product construction in
Section~\ref{q214:height-averaging}: it chooses $r$ separate occurrences of
$L'$-simple factors in the actual $A_a$-isogeny class, then uses a
dimension-two isogeny $C_i\to A_a$.  The quotient-filtration H3 retained
in Section~\ref{q214:height-supplier} is explicitly superseded for this
application.  The coarse support and conductor source fragments are retained
as provenance; the current character and conductor applications are M1 and
M5 of Section~\ref{q214:modularity}.

\paragraph{Attribution and review boundary.}
This edition serializes the public $214$ proposal and its attributed internal
model contributions under Elliot Glazer's direction.  The retained Version~16
body records its historical authorship and acknowledgements, including
Jos\'e Brox's starting construction and Gabriel Conant's contributions.
The present assembly and typesetting were prepared by a Sol~6.1 author agent.
Source statements concerning earlier inspection, pending status or historical
acceptance remain dated evidence.  They are not new acceptance records for
this exact TeX/PDF.  The independently authored public review overlay and
protected controller decision govern any later adoption.  Internal model
review is distinct from human peer review and formal verification.
The bibliography and source catalogue identify external works without
redistributing their full PDFs.  Exact source hashes, line coverage and
typesetting transformations are recorded in \nolinkurl{SOURCE_SPANS.json}.
\clearpage\section{Complete coordinate and row specification}\label{q214:rows}
Put $X_1=x$.  The following dictionary lists every coordinate, in the exact IR order.  Its names describe gates; the equations below, rather than a choice of field-component interpretation, define the polynomial.  The last column is the retained raw-$224$ coordinate identifier.  No coordinate is reserved for $A_N$.
{\small\begin{longtable}{@{}r p{0.77\linewidth} r@{}}\toprule
$i$ & Coordinate name & Parent\\\midrule\endhead
1 & \nolinkurl{input} & 1\\
2 & \nolinkurl{S.chart.j0.rescaled} & 2\\
3 & \nolinkurl{S.chart.j1.rescaled} & 3\\
4 & \nolinkurl{reciprocal.rho} & 4\\
5 & \nolinkurl{detector0.normalized.scaled_e} & 5\\
6 & \nolinkurl{detector0.normalized.s} & 6\\
7 & \nolinkurl{detector0.normalized.p} & 7\\
8 & \nolinkurl{detector0.factored.q} & 8\\
9 & \nolinkurl{detector1.normalized.scaled_e} & 9\\
10 & \nolinkurl{detector1.normalized.s} & 10\\
11 & \nolinkurl{detector1.normalized.p} & 11\\
12 & \nolinkurl{detector1.factored.q} & 12\\
13 & \nolinkurl{normalized_scalar.scaled_e} & 13\\
14 & \nolinkurl{normalized_scalar.s} & 14\\
15 & \nolinkurl{normalized_scalar.p} & 15\\
16 & \nolinkurl{normalized_scalar.factored.q} & 16\\
17 & \nolinkurl{U0.f.0} & 17\\
18 & \nolinkurl{U0.f.1} & 18\\
19 & \nolinkurl{primary.negativeReciprocalZ0.0} & 19\\
20 & \nolinkurl{primary.negativeReciprocalZ0.1} & 20\\
21 & \nolinkurl{logV0.fSquare.0.quadratic} & 21\\
22 & \nolinkurl{logV0.fSquare.1.quadratic} & 22\\
23 & \nolinkurl{logV1.fSquare.0.quadratic} & 25\\
24 & \nolinkurl{logV1.fSquare.1.quadratic} & 26\\
25 & \nolinkurl{logV2.fSquare.0.quadratic} & 27\\
26 & \nolinkurl{logV2.fSquare.1.quadratic} & 28\\
27 & \nolinkurl{logV3.fSquare.0.quadratic} & 31\\
28 & \nolinkurl{logV3.fSquare.1.quadratic} & 32\\
29 & \nolinkurl{logV4.fSquare.0.quadratic} & 35\\
30 & \nolinkurl{logV4.fSquare.1.quadratic} & 36\\
31 & \nolinkurl{L1.0} & 39\\
32 & \nolinkurl{L1.1} & 40\\
33 & \nolinkurl{LP1.0} & 41\\
34 & \nolinkurl{LP1.1} & 42\\
35 & \nolinkurl{L2.0.rescaled} & 43\\
36 & \nolinkurl{L2.1.rescaled} & 44\\
37 & \nolinkurl{LP2.0.rescaled} & 45\\
38 & \nolinkurl{LP2.1.rescaled} & 46\\
39 & \nolinkurl{L3.0.rescaled} & 47\\
40 & \nolinkurl{L3.1.rescaled} & 48\\
41 & \nolinkurl{LP3.0.rescaled} & 49\\
42 & \nolinkurl{LP3.1.rescaled} & 50\\
43 & \nolinkurl{L4.0.rescaled} & 51\\
44 & \nolinkurl{L4.1.rescaled} & 52\\
45 & \nolinkurl{LP4.0.rescaled} & 53\\
46 & \nolinkurl{LP4.1.rescaled} & 54\\
47 & \nolinkurl{inputSquare.quadratic} & 55\\
48 & \nolinkurl{Lucas.C.0.quadratic.rescaled} & 56\\
49 & \nolinkurl{Lucas.C.1.quadratic.rescaled} & 57\\
50 & \nolinkurl{Lucas.AT.0.quadratic.rescaled} & 58\\
51 & \nolinkurl{Lucas.AT.1.quadratic.rescaled} & 59\\
52 & \nolinkurl{Lucas.F0.quadratic} & 60\\
53 & \nolinkurl{Lucas.F1.quadratic} & 61\\
54 & \nolinkurl{testSquare.1.0.0.quadratic} & 62\\
55 & \nolinkurl{testSquare.1.0.1.quadratic} & 63\\
56 & \nolinkurl{testSquare.1.1.0.quadratic} & 64\\
57 & \nolinkurl{testSquare.1.1.1.quadratic} & 65\\
58 & \nolinkurl{testSquare.1.2.0.quadratic} & 66\\
59 & \nolinkurl{testSquare.1.2.1.quadratic} & 67\\
60 & \nolinkurl{testSquare.1.3.0.quadratic} & 68\\
61 & \nolinkurl{testSquare.1.3.1.quadratic} & 69\\
62 & \nolinkurl{testSquare.12.0.0.quadratic} & 70\\
63 & \nolinkurl{testSquare.12.0.1.quadratic} & 71\\
64 & \nolinkurl{testSquare.12.1.0.quadratic} & 72\\
65 & \nolinkurl{testSquare.12.1.1.quadratic} & 73\\
66 & \nolinkurl{testSquare.40.0.0.quadratic} & 74\\
67 & \nolinkurl{testSquare.40.0.1.quadratic} & 75\\
68 & \nolinkurl{testSquare.40.1.0.quadratic} & 76\\
69 & \nolinkurl{testSquare.40.1.1.quadratic} & 77\\
70 & \nolinkurl{testSquare.40.2.0.quadratic} & 78\\
71 & \nolinkurl{testSquare.40.2.1.quadratic} & 79\\
72 & \nolinkurl{testSquare.40.3.0.quadratic} & 80\\
73 & \nolinkurl{testSquare.40.3.1.quadratic} & 81\\
74 & \nolinkurl{testSquare.40.4.0.quadratic} & 82\\
75 & \nolinkurl{testSquare.40.4.1.quadratic} & 83\\
76 & \nolinkurl{reciprocal.rhoSquared} & 84\\
77 & \nolinkurl{reciprocal.rhoFourth} & 85\\
78 & \nolinkurl{reciprocal.contactPolynomial} & 86\\
79 & \nolinkurl{detector0.q1} & 87\\
80 & \nolinkurl{detector0.inverseWeight} & 88\\
81 & \nolinkurl{detector0.inverseWeightSquared} & 89\\
82 & \nolinkurl{detector0.inverseWeightFourthPower} & 90\\
83 & \nolinkurl{detector0.actualW.0} & 91\\
84 & \nolinkurl{detector0.actualW.1} & 92\\
85 & \nolinkurl{detector0.normOver2w0} & 93\\
86 & \nolinkurl{detector0.B} & 94\\
87 & \nolinkurl{detector0.normalized.p2.quadratic} & 95\\
88 & \nolinkurl{detector0.factored.v} & 96\\
89 & \nolinkurl{detector0.factored.q2} & 97\\
90 & \nolinkurl{detector1.t.quadratic.rescaled} & 98\\
91 & \nolinkurl{symmetric.detector1.deltaSquared.rescaled} & 99\\
92 & \nolinkurl{symmetric.detector1.deltaFourthPower.rescaled} & 100\\
93 & \nolinkurl{symmetric.detector1.contactPolynomial.rescaled} & 101\\
94 & \nolinkurl{detector1.q1} & 102\\
95 & \nolinkurl{detector1.inverseWeight} & 103\\
96 & \nolinkurl{detector1.inverseWeightSquared} & 104\\
97 & \nolinkurl{detector1.inverseWeightFourthPower} & 105\\
98 & \nolinkurl{detector1.actualW.0} & 106\\
99 & \nolinkurl{detector1.actualW.1} & 107\\
100 & \nolinkurl{detector1.normOver2w0} & 108\\
101 & \nolinkurl{detector1.B} & 109\\
102 & \nolinkurl{detector1.normalized.p2.quadratic} & 110\\
103 & \nolinkurl{detector1.factored.v} & 111\\
104 & \nolinkurl{detector1.factored.q2} & 112\\
105 & \nolinkurl{scalar.rPair1.quadratic} & 113\\
106 & \nolinkurl{scalar.contactProduct} & 114\\
107 & \nolinkurl{scalar.zRatio.0.rescaled} & 115\\
108 & \nolinkurl{scalar.zRatio.1.rescaled} & 116\\
109 & \nolinkurl{scalar.W2.quadratic.rescaled} & 117\\
110 & \nolinkurl{scalar.b.quadratic} & 118\\
111 & \nolinkurl{normalized_scalar.p2.quadratic} & 119\\
112 & \nolinkurl{normalized_scalar.factored.v} & 120\\
113 & \nolinkurl{normalized_scalar.factored.q2} & 121\\
114 & \nolinkurl{scalar.ratio1_unscaled_square.rescaled} & 122\\
115 & \nolinkurl{parity_scalar.C_Z0square_over_R} & 123\\
116 & \nolinkurl{reduced_scalar.root.rescaled} & 124\\
117 & \nolinkurl{scalar.r_power_2} & 125\\
118 & \nolinkurl{scalar.forced_positive_parameter} & 126\\
119 & \nolinkurl{reciprocal.rhoFifth} & 127\\
120 & \nolinkurl{reflected.detector1.deltaSixthPower.rescaled} & 128\\
121 & \nolinkurl{amplified.C1.power4.0} & 129\\
122 & \nolinkurl{amplified.C1.power4.1} & 130\\
123 & \nolinkurl{residue.I.power4.0} & 131\\
124 & \nolinkurl{residue.I.power4.1} & 132\\
125 & \nolinkurl{amplified.det0.factor0.0} & 135\\
126 & \nolinkurl{amplified.det0.factor0.1} & 136\\
127 & \nolinkurl{amplified.det0.factor1.0} & 137\\
128 & \nolinkurl{amplified.det0.factor1.1} & 138\\
129 & \nolinkurl{amplified.det0.factor2.0} & 139\\
130 & \nolinkurl{amplified.det0.factor2.1} & 140\\
131 & \nolinkurl{amplified.det0.factor3.0} & 141\\
132 & \nolinkurl{amplified.det0.factor3.1} & 142\\
133 & \nolinkurl{amplified.det0.factor4.0} & 143\\
134 & \nolinkurl{amplified.det0.factor4.1} & 144\\
135 & \nolinkurl{amplified.det0.product1.0} & 145\\
136 & \nolinkurl{amplified.det0.product1.1} & 146\\
137 & \nolinkurl{amplified.det0.product2.0} & 147\\
138 & \nolinkurl{amplified.det0.product2.1} & 148\\
139 & \nolinkurl{amplified.det0.product3.0} & 149\\
140 & \nolinkurl{amplified.det0.product3.1} & 150\\
141 & \nolinkurl{amplified.det1.factor0.0} & 151\\
142 & \nolinkurl{amplified.det1.factor0.1} & 152\\
143 & \nolinkurl{amplified.det1.factor1.0} & 153\\
144 & \nolinkurl{amplified.det1.factor1.1} & 154\\
145 & \nolinkurl{amplified.det1.factor2.0} & 155\\
146 & \nolinkurl{amplified.det1.factor2.1} & 156\\
147 & \nolinkurl{amplified.det1.factor3.0} & 157\\
148 & \nolinkurl{amplified.det1.factor3.1} & 158\\
149 & \nolinkurl{amplified.det1.factor4.0} & 159\\
150 & \nolinkurl{amplified.det1.factor4.1} & 160\\
151 & \nolinkurl{amplified.det1.product1.0} & 161\\
152 & \nolinkurl{amplified.det1.product1.1} & 162\\
153 & \nolinkurl{amplified.det1.product2.0} & 163\\
154 & \nolinkurl{amplified.det1.product2.1} & 164\\
155 & \nolinkurl{amplified.det1.product3.0} & 165\\
156 & \nolinkurl{amplified.det1.product3.1} & 166\\
157 & \nolinkurl{rootProduct.1.A.0.rescaled} & 167\\
158 & \nolinkurl{rootProduct.1.A.1.rescaled} & 168\\
159 & \nolinkurl{rootProduct.1.B.0.rescaled} & 169\\
160 & \nolinkurl{rootProduct.1.B.1.rescaled} & 170\\
161 & \nolinkurl{rootProduct.1.C.0.rescaled} & 171\\
162 & \nolinkurl{rootProduct.1.C.1.rescaled} & 172\\
163 & \nolinkurl{rootProduct.1.E.0.rescaled} & 173\\
164 & \nolinkurl{rootProduct.1.E.1.rescaled} & 174\\
165 & \nolinkurl{rootProduct.1.i.0.rescaled} & 175\\
166 & \nolinkurl{rootProduct.1.i.1.rescaled} & 176\\
167 & \nolinkurl{rootProduct.1.k.0.rescaled} & 177\\
168 & \nolinkurl{rootProduct.1.k.1.rescaled} & 178\\
169 & \nolinkurl{rootProduct.1.u.0} & 179\\
170 & \nolinkurl{rootProduct.1.u.1} & 180\\
171 & \nolinkurl{rootProduct.2.A.0.rescaled} & 181\\
172 & \nolinkurl{rootProduct.2.A.1.rescaled} & 182\\
173 & \nolinkurl{rootProduct.2.B.0.rescaled} & 183\\
174 & \nolinkurl{rootProduct.2.B.1.rescaled} & 184\\
175 & \nolinkurl{rootProduct.2.C.0.rescaled} & 185\\
176 & \nolinkurl{rootProduct.2.C.1.rescaled} & 186\\
177 & \nolinkurl{rootProduct.2.E.0.rescaled} & 187\\
178 & \nolinkurl{rootProduct.2.E.1.rescaled} & 188\\
179 & \nolinkurl{rootProduct.2.i.0} & 189\\
180 & \nolinkurl{rootProduct.2.i.1} & 190\\
181 & \nolinkurl{rootProduct.2.k.0} & 191\\
182 & \nolinkurl{rootProduct.2.k.1} & 192\\
183 & \nolinkurl{rootProduct.2.u.0} & 193\\
184 & \nolinkurl{rootProduct.2.u.1} & 194\\
185 & \nolinkurl{rootProduct.3.A.0.rescaled} & 195\\
186 & \nolinkurl{rootProduct.3.A.1.rescaled} & 196\\
187 & \nolinkurl{rootProduct.3.B.0.rescaled} & 197\\
188 & \nolinkurl{rootProduct.3.B.1.rescaled} & 198\\
189 & \nolinkurl{rootProduct.3.C.0.rescaled} & 199\\
190 & \nolinkurl{rootProduct.3.C.1.rescaled} & 200\\
191 & \nolinkurl{rootProduct.3.E.0.rescaled} & 201\\
192 & \nolinkurl{rootProduct.3.E.1.rescaled} & 202\\
193 & \nolinkurl{rootProduct.3.i.0} & 203\\
194 & \nolinkurl{rootProduct.3.i.1} & 204\\
195 & \nolinkurl{rootProduct.3.k.0} & 205\\
196 & \nolinkurl{rootProduct.3.k.1} & 206\\
197 & \nolinkurl{rootProduct.3.u.0} & 207\\
198 & \nolinkurl{rootProduct.3.u.1} & 208\\
199 & \nolinkurl{rootProduct.4.A.0.rescaled} & 209\\
200 & \nolinkurl{rootProduct.4.A.1.rescaled} & 210\\
201 & \nolinkurl{rootProduct.4.B.0.rescaled} & 211\\
202 & \nolinkurl{rootProduct.4.B.1.rescaled} & 212\\
203 & \nolinkurl{rootProduct.4.C.0.rescaled} & 213\\
204 & \nolinkurl{rootProduct.4.C.1.rescaled} & 214\\
205 & \nolinkurl{rootProduct.4.E.0.rescaled} & 215\\
206 & \nolinkurl{rootProduct.4.E.1.rescaled} & 216\\
207 & \nolinkurl{rootProduct.4.i.0} & 217\\
208 & \nolinkurl{rootProduct.4.i.1} & 218\\
209 & \nolinkurl{rootProduct.4.k.0} & 219\\
210 & \nolinkurl{rootProduct.4.k.1} & 220\\
211 & \nolinkurl{rootProduct.4.u.0} & 221\\
212 & \nolinkurl{rootProduct.4.u.1} & 222\\
213 & \nolinkurl{shared.alternate.negativeReciprocalZ0.0} & 223\\
214 & \nolinkurl{shared.alternate.negativeReciprocalZ0.1} & 224\\
\bottomrule\end{longtable}}
For $1\le i\le188$, the row below is $r_{i,N}$.  Each term is an exact integer coefficient times a power of $A_N$ and a coordinate monomial.  Long integer coefficients are grouped on separate display lines without abbreviation.  The definition is $P_N=\sum_i r_{i,N}^2$, with weight one on each of these primitive integral rows.  In particular the historical $222$-variable row weights are not imported.
\paragraph{Row 1.} \nolinkurl{inputSquare.definition}\nopagebreak
\begin{equation*}\begin{aligned}r_{1,N}={}&
 X_{1}^{2}\\ &
- X_{47}
\end{aligned}\end{equation*}
\paragraph{Row 2.} \nolinkurl{Lucas.C.0.definition}\nopagebreak
\begin{equation*}\begin{aligned}r_{2,N}={}&
 X_{35}\,X_{43}\\ &
+ 2\,X_{36}\,X_{44}\\ &
- X_{48}\,X_{118}
\end{aligned}\end{equation*}
\paragraph{Row 3.} \nolinkurl{Lucas.C.1.definition}\nopagebreak
\begin{equation*}\begin{aligned}r_{3,N}={}&
 X_{35}\,X_{44}\\ &
+ X_{36}\,X_{43}\\ &
- X_{49}\,X_{118}
\end{aligned}\end{equation*}
\paragraph{Row 4.} \nolinkurl{Lucas.AT.0.definition}\nopagebreak
\begin{equation*}\begin{aligned}r_{4,N}={}&
 X_{35}\,X_{39}\\ &
+ 2\,X_{36}\,X_{40}\\ &
- X_{50}\,X_{118}
\end{aligned}\end{equation*}
\paragraph{Row 5.} \nolinkurl{Lucas.AT.1.definition}\nopagebreak
\begin{equation*}\begin{aligned}r_{5,N}={}&
 X_{35}\,X_{40}\\ &
+ X_{36}\,X_{39}\\ &
- X_{51}\,X_{118}
\end{aligned}\end{equation*}
\paragraph{Row 6.} \nolinkurl{Lucas.F0.definition}\nopagebreak
\begin{equation*}\begin{aligned}r_{6,N}={}&
 X_{39}^{2}\\ &
+ 2\,X_{40}^{2}\\ &
+ 4\,X_{43}^{2}\\ &
- 4\,X_{43}\,X_{48}\\ &
+ 2\,X_{43}\,X_{50}\\ &
+ 8\,X_{44}^{2}\\ &
- 8\,X_{44}\,X_{49}\\ &
+ 4\,X_{44}\,X_{51}\\ &
+ X_{48}^{2}\\ &
- X_{48}\,X_{50}\\ &
+ 2\,X_{49}^{2}\\ &
- 2\,X_{49}\,X_{51}\\ &
- X_{52}\,X_{117}
\end{aligned}\end{equation*}
\paragraph{Row 7.} \nolinkurl{Lucas.F1.definition}\nopagebreak
\begin{equation*}\begin{aligned}r_{7,N}={}&
 2\,X_{39}\,X_{40}\\ &
+ 8\,X_{43}\,X_{44}\\ &
- 4\,X_{43}\,X_{49}\\ &
+ 2\,X_{43}\,X_{51}\\ &
- 4\,X_{44}\,X_{48}\\ &
+ 2\,X_{44}\,X_{50}\\ &
+ 2\,X_{48}\,X_{49}\\ &
- X_{48}\,X_{51}\\ &
- X_{49}\,X_{50}\\ &
- 2\,X_{53}\,X_{117}
\end{aligned}\end{equation*}
\paragraph{Row 8.} \nolinkurl{testSquare.12.0.0.definition}\nopagebreak
\begin{equation*}\begin{aligned}r_{8,N}={}&
 X_{31}^{2}\\ &
- 2\,X_{31}\,X_{33}\\ &
+ 2\,X_{32}^{2}\\ &
+ 4\,X_{32}\,X_{34}\\ &
+ X_{33}^{2}\\ &
+ 2\,X_{34}^{2}\\ &
- X_{62}
\end{aligned}\end{equation*}
\paragraph{Row 9.} \nolinkurl{testSquare.12.0.1.definition}\nopagebreak
\begin{equation*}\begin{aligned}r_{9,N}={}&
 X_{31}\,X_{32}\\ &
+ X_{31}\,X_{34}\\ &
- X_{32}\,X_{33}\\ &
- X_{33}\,X_{34}\\ &
- X_{63}
\end{aligned}\end{equation*}
\paragraph{Row 10.} \nolinkurl{testSquare.12.1.0.definition}\nopagebreak
\begin{equation*}\begin{aligned}r_{10,N}={}&
 6255001\\ &
- 5002\,X_{31}\\ &
+ X_{31}^{2}\\ &
- 102\,X_{31}\,X_{47}\\ &
+ 2\,X_{32}^{2}\\ &
+ 255102\,X_{47}\\ &
+ 2601\,X_{47}^{2}\\ &
- X_{64}
\end{aligned}\end{equation*}
\paragraph{Row 11.} \nolinkurl{testSquare.12.1.1.definition}\nopagebreak
\begin{equation*}\begin{aligned}r_{11,N}={}&
 X_{31}\,X_{32}\\ &
- 2501\,X_{32}\\ &
- 51\,X_{32}\,X_{47}\\ &
- X_{65}
\end{aligned}\end{equation*}
\paragraph{Row 12.} \nolinkurl{testSquare.40.0.0.definition}\nopagebreak
\begin{equation*}\begin{aligned}r_{12,N}={}&
 X_{17}^{2}\\ &
+ 2\,X_{18}^{2}\\ &
- X_{66}
\end{aligned}\end{equation*}
\paragraph{Row 13.} \nolinkurl{testSquare.40.0.1.definition}\nopagebreak
\begin{equation*}\begin{aligned}r_{13,N}={}&
 X_{17}\,X_{18}\\ &
- X_{67}
\end{aligned}\end{equation*}
\paragraph{Row 14.} \nolinkurl{testSquare.40.1.0.definition}\nopagebreak
\begin{equation*}\begin{aligned}r_{14,N}={}&
 X_{35}^{2}\\ &
- 2\,X_{35}\,X_{37}\\ &
+ 2\,X_{36}^{2}\\ &
+ 4\,X_{36}\,X_{38}\\ &
+ X_{37}^{2}\\ &
+ 2\,X_{38}^{2}\\ &
- X_{68}\,X_{117}
\end{aligned}\end{equation*}
\paragraph{Row 15.} \nolinkurl{testSquare.40.1.1.definition}\nopagebreak
\begin{equation*}\begin{aligned}r_{15,N}={}&
 X_{35}\,X_{36}\\ &
+ X_{35}\,X_{38}\\ &
- X_{36}\,X_{37}\\ &
- X_{37}\,X_{38}\\ &
- X_{69}\,X_{117}
\end{aligned}\end{equation*}
\paragraph{Row 16.} \nolinkurl{testSquare.40.2.0.definition}\nopagebreak
\begin{equation*}\begin{aligned}r_{16,N}={}&
 X_{39}^{2}\\ &
- 2\,X_{39}\,X_{41}\\ &
+ 2\,X_{40}^{2}\\ &
+ 4\,X_{40}\,X_{42}\\ &
+ X_{41}^{2}\\ &
+ 2\,X_{42}^{2}\\ &
- X_{70}\,X_{117}
\end{aligned}\end{equation*}
\paragraph{Row 17.} \nolinkurl{testSquare.40.2.1.definition}\nopagebreak
\begin{equation*}\begin{aligned}r_{17,N}={}&
 X_{39}\,X_{40}\\ &
+ X_{39}\,X_{42}\\ &
- X_{40}\,X_{41}\\ &
- X_{41}\,X_{42}\\ &
- X_{71}\,X_{117}
\end{aligned}\end{equation*}
\paragraph{Row 18.} \nolinkurl{testSquare.40.3.0.definition}\nopagebreak
\begin{equation*}\begin{aligned}r_{18,N}={}&
 X_{43}^{2}\\ &
- 2\,X_{43}\,X_{45}\\ &
+ 2\,X_{44}^{2}\\ &
+ 4\,X_{44}\,X_{46}\\ &
+ X_{45}^{2}\\ &
+ 2\,X_{46}^{2}\\ &
- X_{72}\,X_{117}
\end{aligned}\end{equation*}
\paragraph{Row 19.} \nolinkurl{testSquare.40.3.1.definition}\nopagebreak
\begin{equation*}\begin{aligned}r_{19,N}={}&
 X_{43}\,X_{44}\\ &
+ X_{43}\,X_{46}\\ &
- X_{44}\,X_{45}\\ &
- X_{45}\,X_{46}\\ &
- X_{73}\,X_{117}
\end{aligned}\end{equation*}
\paragraph{Row 20.} \nolinkurl{testSquare.40.4.0.definition}\nopagebreak
\begin{equation*}\begin{aligned}r_{20,N}={}&
 X_{52}^{2}\\ &
+ 8\,X_{53}^{2}\\ &
- X_{74}
\end{aligned}\end{equation*}
\paragraph{Row 21.} \nolinkurl{testSquare.40.4.1.definition}\nopagebreak
\begin{equation*}\begin{aligned}r_{21,N}={}&
 X_{52}\,X_{53}\\ &
- X_{75}
\end{aligned}\end{equation*}
\paragraph{Row 22.} \nolinkurl{reciprocal.symmetric.detector0.deltaSquared.definition}\nopagebreak
\begin{equation*}\begin{aligned}r_{22,N}={}&
 X_{4}^{2}\\ &
- X_{76}
\end{aligned}\end{equation*}
\paragraph{Row 23.} \nolinkurl{reciprocal.symmetric.detector0.deltaFourthPower.definition}\nopagebreak
\begin{equation*}\begin{aligned}r_{23,N}={}&
 X_{76}^{2}\\ &
- X_{77}
\end{aligned}\end{equation*}
\paragraph{Row 24.} \nolinkurl{reciprocal.reflected.detector0.contactPolynomial.definition}\nopagebreak
\begin{equation*}\begin{aligned}r_{24,N}={}&
 3813373161796\\ &
- 38133731617960\,X_{4}\\ &
- 170454839579608\,X_{4}\,X_{76}\\ &
- 68335335614837\,X_{4}\,X_{119}\\ &
+ 135708049801331\,X_{76}\\ &
- 126186643896348\,X_{76}\,X_{119}\\ &
- 81840384134614\,X_{77}\\ &
- 9293728180788\,X_{77}^{2}\\ &
+ 3895931520\,X_{77}\,X_{119}\\ &
- 3813373161796\,X_{78}\\ &
+ 314823294195284\,X_{119}
\end{aligned}\end{equation*}
\paragraph{Row 25.} \nolinkurl{reciprocal.reflected.detector0.weight.definition}\nopagebreak
\begin{equation*}\begin{aligned}r_{25,N}={}&
 1\\ &
- 10\,X_{4}\\ &
- 120\,X_{4}\,X_{76}\\ &
+ 210\,X_{4}\,X_{119}\\ &
+ 45\,X_{76}\\ &
- 120\,X_{76}\,X_{119}\\ &
+ 210\,X_{77}\\ &
+ 45\,X_{77}^{2}\\ &
- 10\,X_{77}\,X_{119}\\ &
- X_{78}\\ &
- X_{78}\,X_{79}\\ &
- 252\,X_{119}\\ &
+ X_{119}^{2}
\end{aligned}\end{equation*}
\paragraph{Row 26.} \nolinkurl{detector0.inverseWeight.definition}\nopagebreak
\begin{equation*}\begin{aligned}r_{26,N}={}&
 1\\ &
- X_{79}\,X_{80}
\end{aligned}\end{equation*}
\paragraph{Row 27.} \nolinkurl{detector0.inverseWeightSquared.definition}\nopagebreak
\begin{equation*}\begin{aligned}r_{27,N}={}&
 X_{80}^{2}\\ &
- X_{81}
\end{aligned}\end{equation*}
\paragraph{Row 28.} \nolinkurl{detector0.normOver2w0.definition}\nopagebreak
\begin{equation*}\begin{aligned}r_{28,N}={}&
 X_{83}^{2}\\ &
- 2\,X_{83}\,X_{85}\\ &
- 2\,X_{84}^{2}
\end{aligned}\end{equation*}
\paragraph{Row 29.} \nolinkurl{balanced.detector0.poleInput.definition}\nopagebreak
\begin{equation*}\begin{aligned}r_{29,N}={}&
 167281\,X_{79}\,X_{86}\\ &
- 166464\,X_{85}^{2}\\ &
+ 167281\,X_{86}
\end{aligned}\end{equation*}
\paragraph{Row 30.} \nolinkurl{detector0.p2.definition}\nopagebreak
\begin{equation*}\begin{aligned}r_{30,N}={}&
 X_{7}^{2}\\ &
- X_{87}
\end{aligned}\end{equation*}
\paragraph{Row 31.} \nolinkurl{detector0.firstConic}\nopagebreak
\begin{equation*}\begin{aligned}r_{31,N}={}&
 289\\ &
+ 34\,X_{5}\\ &
- 3\,X_{5}^{2}\\ &
- 17\,X_{86}\,X_{87}
\end{aligned}\end{equation*}
\paragraph{Row 32.} \nolinkurl{detector0.Bh.definition}\nopagebreak
\begin{equation*}\begin{aligned}r_{32,N}={}&
 17\,X_{5}\\ &
- X_{5}\,X_{86}\\ &
- 17\,X_{6}\,X_{88}
\end{aligned}\end{equation*}
\paragraph{Row 33.} \nolinkurl{detector0.q2.definition}\nopagebreak
\begin{equation*}\begin{aligned}r_{33,N}={}&
 X_{8}^{2}\\ &
- X_{89}
\end{aligned}\end{equation*}
\paragraph{Row 34.} \nolinkurl{detector0.secondConic}\nopagebreak
\begin{equation*}\begin{aligned}r_{34,N}={}&
 X_{6}^{2}\\ &
+ 34\,X_{6}\,X_{88}\\ &
- 17\,X_{86}\,X_{89}\\ &
- 867\,X_{88}^{2}
\end{aligned}\end{equation*}
\paragraph{Row 35.} \nolinkurl{symmetric.detector1.deltaSquared.definition}\nopagebreak
\begin{equation*}\begin{aligned}r_{35,N}={}&
 X_{90}^{2}\\ &
- 2\,X_{90}\,X_{117}\\ &
- X_{91}\,X_{117}\\ &
+ X_{117}^{2}
\end{aligned}\end{equation*}
\paragraph{Row 36.} \nolinkurl{symmetric.detector1.deltaFourthPower.definition}\nopagebreak
\begin{equation*}\begin{aligned}r_{36,N}={}&
 X_{91}^{2}\\ &
- X_{92}\,X_{117}
\end{aligned}\end{equation*}
\paragraph{Row 37.} \nolinkurl{reflected.detector1.contactPolynomial.definition}\nopagebreak
\begin{equation*}\begin{aligned}r_{37,N}={}&
 141822723674536\,X_{91}\,X_{117}\\ &
- 35893742479489\,X_{91}\,X_{120}\\ &
- 192222345736537\,X_{92}\,X_{117}\\ &
+ 3813373161796\,X_{92}\,X_{120}\\ &
- 3813373161796\,X_{93}\\ &
- 39896049934224\,X_{117}^{2}\\ &
+ 122376041313918\,X_{117}\,X_{120}
\end{aligned}\end{equation*}
\paragraph{Row 38.} \nolinkurl{reflected.detector1.weight.definition}\nopagebreak
\begin{equation*}\begin{aligned}r_{38,N}={}&
 X_{92}\,X_{120}\\ &
- X_{93}\\ &
- X_{93}\,X_{94}
\end{aligned}\end{equation*}
\paragraph{Row 39.} \nolinkurl{detector1.inverseWeight.definition}\nopagebreak
\begin{equation*}\begin{aligned}r_{39,N}={}&
 1\\ &
- X_{94}\,X_{95}
\end{aligned}\end{equation*}
\paragraph{Row 40.} \nolinkurl{detector1.inverseWeightSquared.definition}\nopagebreak
\begin{equation*}\begin{aligned}r_{40,N}={}&
 X_{95}^{2}\\ &
- X_{96}
\end{aligned}\end{equation*}
\paragraph{Row 41.} \nolinkurl{detector1.normOver2w0.definition}\nopagebreak
\begin{equation*}\begin{aligned}r_{41,N}={}&
 X_{98}^{2}\\ &
- 2\,X_{98}\,X_{100}\\ &
- 2\,X_{99}^{2}
\end{aligned}\end{equation*}
\paragraph{Row 42.} \nolinkurl{balanced.detector1.poleInput.definition}\nopagebreak
\begin{equation*}\begin{aligned}r_{42,N}={}&
 167281\,X_{94}\,X_{101}\\ &
- 166464\,X_{100}^{2}\\ &
+ 167281\,X_{101}
\end{aligned}\end{equation*}
\paragraph{Row 43.} \nolinkurl{detector1.p2.definition}\nopagebreak
\begin{equation*}\begin{aligned}r_{43,N}={}&
 X_{11}^{2}\\ &
- X_{102}
\end{aligned}\end{equation*}
\paragraph{Row 44.} \nolinkurl{detector1.firstConic}\nopagebreak
\begin{equation*}\begin{aligned}r_{44,N}={}&
 289\\ &
+ 34\,X_{9}\\ &
- 3\,X_{9}^{2}\\ &
- 17\,X_{101}\,X_{102}
\end{aligned}\end{equation*}
\paragraph{Row 45.} \nolinkurl{detector1.Bh.definition}\nopagebreak
\begin{equation*}\begin{aligned}r_{45,N}={}&
 17\,X_{9}\\ &
- X_{9}\,X_{101}\\ &
- 17\,X_{10}\,X_{103}
\end{aligned}\end{equation*}
\paragraph{Row 46.} \nolinkurl{detector1.q2.definition}\nopagebreak
\begin{equation*}\begin{aligned}r_{46,N}={}&
 X_{12}^{2}\\ &
- X_{104}
\end{aligned}\end{equation*}
\paragraph{Row 47.} \nolinkurl{detector1.secondConic}\nopagebreak
\begin{equation*}\begin{aligned}r_{47,N}={}&
 X_{10}^{2}\\ &
+ 34\,X_{10}\,X_{103}\\ &
- 17\,X_{101}\,X_{104}\\ &
- 867\,X_{103}^{2}
\end{aligned}\end{equation*}
\paragraph{Row 48.} \nolinkurl{scalar.rPair1.definition}\nopagebreak
\begin{equation*}\begin{aligned}r_{48,N}={}&
 6255001\\ &
+ 255102\,X_{47}\\ &
+ 2601\,X_{47}^{2}\\ &
- 4\,X_{105}
\end{aligned}\end{equation*}
\paragraph{Row 49.} \nolinkurl{scalar.four_contact_product.definition}\nopagebreak
\begin{equation*}\begin{aligned}r_{49,N}={}&
 14315517724\\ &
+ 291790992\,X_{47}\\ &
+ 15531948\,X_{47}\,X_{105}\\ &
+ 767410608\,X_{105}\\ &
- 512\,X_{105}^{2}\\ &
+ X_{106}
\end{aligned}\end{equation*}
\paragraph{Row 50.} \nolinkurl{scalar.W2.definition}\nopagebreak
\begin{equation*}\begin{aligned}r_{50,N}={}&
 \left(\begin{gathered}98194\cdot10^{100}\\{}+{}99495210116462795099322880135176040946699131400335\cdot10^{50}\\{}+{}89964224332595426527668691070368153600000000000000\end{gathered}\right)\,X_{109}\,X_{117}\\ &
- X_{116}^{2}
\end{aligned}\end{equation*}
\paragraph{Row 51.} \nolinkurl{scalar.b.definition}\nopagebreak
\begin{equation*}\begin{aligned}r_{51,N}={}&
 9792\,X_{109}\\ &
+ X_{109}\,X_{118}\\ &
+ 82479248\,X_{110}\,X_{118}
\end{aligned}\end{equation*}
\paragraph{Row 52.} \nolinkurl{scalar.p2.definition}\nopagebreak
\begin{equation*}\begin{aligned}r_{52,N}={}&
 X_{15}^{2}\\ &
- X_{111}
\end{aligned}\end{equation*}
\paragraph{Row 53.} \nolinkurl{scalar.firstConic}\nopagebreak
\begin{equation*}\begin{aligned}r_{53,N}={}&
 \left(\begin{gathered}6318946488646535718731759452\cdot10^{200}\\{}+{}29280662018444979247293017626849316401438015472961\cdot10^{150}\\{}+{}54172410366489483083614609845481185889826867345597\cdot10^{100}\\{}+{}46745916595796472604244988128011149397522493691010\cdot10^{50}\\{}+{}83632379691878234193920000000000000000000000000000\end{gathered}\right)\\ &
- \left(\begin{gathered}91789225864815\cdot10^{100}\\{}+{}13422229277396013244425859655544127655066108765484\cdot10^{50}\\{}+{}58365423573402134769787177776866918400000000000000\end{gathered}\right)\,X_{13}\\ &
- X_{13}^{2}\\ &
+ \left(\begin{gathered}45894612932407\cdot10^{100}\\{}+{}56711114638698006622212929827772063827533054382742\cdot10^{50}\\{}+{}29182711786701067384893588888433459200000000000000\end{gathered}\right)\,X_{110}\,X_{111}
\end{aligned}\end{equation*}
\paragraph{Row 54.} \nolinkurl{scalar.Bh.definition}\nopagebreak
\begin{equation*}\begin{aligned}r_{54,N}={}&
 17\,X_{13}\\ &
+ \left(\begin{gathered}8099049341013\cdot10^{100}\\{}+{}10007843759770236462743458204900952440152891949895\cdot10^{50}\\{}+{}69855772668241364832628280392076492800000000000000\end{gathered}\right)\,X_{13}\,X_{110}\\ &
- 17\,X_{14}\,X_{112}
\end{aligned}\end{equation*}
\paragraph{Row 55.} \nolinkurl{scalar.q2.definition}\nopagebreak
\begin{equation*}\begin{aligned}r_{55,N}={}&
 X_{16}^{2}\\ &
- X_{113}
\end{aligned}\end{equation*}
\paragraph{Row 56.} \nolinkurl{scalar.secondConic}\nopagebreak
\begin{equation*}\begin{aligned}r_{56,N}={}&
 X_{14}^{2}\\ &
- \left(\begin{gathered}275367677594445\cdot10^{100}\\{}+{}40266687832188039733277578966632382965198326296453\cdot10^{50}\\{}+{}75096270720206404309361533330600755200000000000000\end{gathered}\right)\,X_{14}\,X_{112}\\ &
+ \left(\begin{gathered}137683838797222\cdot10^{100}\\{}+{}70133343916094019866638789483316191482599163148226\cdot10^{50}\\{}+{}87548135360103202154680766665300377600000000000000\end{gathered}\right)\,X_{110}\,X_{113}\\ &
- \left(\begin{gathered}56870518397818821468585835070\cdot10^{200}\\{}+{}63525958166004813225637158641643847612942139256653\cdot10^{150}\\{}+{}87551693298405347752531488609330673008441806110377\cdot10^{100}\\{}+{}20713249362168253438204893152100344577702443219097\cdot10^{50}\\{}+{}52691417226904107745280000000000000000000000000000\end{gathered}\right)\,X_{112}^{2}
\end{aligned}\end{equation*}
\paragraph{Row 57.} \nolinkurl{parity_scalar.reduced_scalar.root.power4}\nopagebreak
\begin{equation*}\begin{aligned}r_{57,N}={}&
 X_{66}^{2}\\ &
- 8\,X_{67}^{2}\\ &
+ X_{114}\,X_{118}\\ &
+ X_{115}^{2}\\ &
- X_{116}\,X_{117}
\end{aligned}\end{equation*}
\paragraph{Row 58.} \nolinkurl{detector0.inverseWeightFourthPower.definition}\nopagebreak
\begin{equation*}\begin{aligned}r_{58,N}={}&
 X_{81}^{2}\\ &
- X_{82}
\end{aligned}\end{equation*}
\paragraph{Row 59.} \nolinkurl{detector1.inverseWeightFourthPower.definition}\nopagebreak
\begin{equation*}\begin{aligned}r_{59,N}={}&
 X_{96}^{2}\\ &
- X_{97}
\end{aligned}\end{equation*}
\paragraph{Row 60.} \nolinkurl{scalar.r_power_2.definition}\nopagebreak
\begin{equation*}\begin{aligned}r_{60,N}={}&
 X_{117}\\ &
- X_{118}^{2}
\end{aligned}\end{equation*}
\paragraph{Row 61.} \nolinkurl{scalar.ratio1_unscaled_square.definition}\nopagebreak
\begin{equation*}\begin{aligned}r_{61,N}={}&
 X_{108}^{2}\\ &
- X_{114}
\end{aligned}\end{equation*}
\paragraph{Row 62.} \nolinkurl{parity_scalar.scalar.ratio_node.consolidated_denominator.power4}\nopagebreak
\begin{equation*}\begin{aligned}r_{62,N}={}&
 X_{107}^{2}\\ &
- X_{115}\,X_{118}
\end{aligned}\end{equation*}
\paragraph{Row 63.} \nolinkurl{scalar.forced_prime.definition}\nopagebreak
\begin{equation*}\begin{aligned}r_{63,N}={}&
 22201\,X_{106}\,X_{118}\\ &
+ X_{118}\\ &
- A_N\,X_{106}
\end{aligned}\end{equation*}
\paragraph{Row 64.} \nolinkurl{reciprocal.reflected.detector0.deltaSixthPower.definition}\nopagebreak
\begin{equation*}\begin{aligned}r_{64,N}={}&
 X_{4}\,X_{77}\\ &
- X_{119}
\end{aligned}\end{equation*}
\paragraph{Row 65.} \nolinkurl{reflected.detector1.deltaSixthPower.definition}\nopagebreak
\begin{equation*}\begin{aligned}r_{65,N}={}&
 X_{91}\,X_{92}\\ &
- X_{117}\,X_{120}
\end{aligned}\end{equation*}
\paragraph{Row 66.} \nolinkurl{amplified.C1.power4.0.definition}\nopagebreak
\begin{equation*}\begin{aligned}r_{66,N}={}&
 X_{62}^{2}\\ &
+ 8\,X_{63}^{2}\\ &
- X_{121}
\end{aligned}\end{equation*}
\paragraph{Row 67.} \nolinkurl{amplified.C1.power4.1.definition}\nopagebreak
\begin{equation*}\begin{aligned}r_{67,N}={}&
 4\,X_{62}\,X_{63}\\ &
- X_{122}
\end{aligned}\end{equation*}
\paragraph{Row 68.} \nolinkurl{residue.I.power4.0.definition}\nopagebreak
\begin{equation*}\begin{aligned}r_{68,N}={}&
 225450225\\ &
- 180180\,X_{31}\\ &
+ 36\,X_{31}^{2}\\ &
- 3672\,X_{31}\,X_{47}\\ &
- 12\,X_{31}\,X_{64}\\ &
+ 72\,X_{32}^{2}\\ &
- 48\,X_{32}\,X_{65}\\ &
+ 9189180\,X_{47}\\ &
+ 93636\,X_{47}^{2}\\ &
+ 612\,X_{47}\,X_{64}\\ &
+ 30030\,X_{64}\\ &
+ X_{64}^{2}\\ &
+ 8\,X_{65}^{2}\\ &
- X_{123}
\end{aligned}\end{equation*}
\paragraph{Row 69.} \nolinkurl{residue.I.power4.1.definition}\nopagebreak
\begin{equation*}\begin{aligned}r_{69,N}={}&
 72\,X_{31}\,X_{32}\\ &
- 24\,X_{31}\,X_{65}\\ &
- 180180\,X_{32}\\ &
- 3672\,X_{32}\,X_{47}\\ &
- 12\,X_{32}\,X_{64}\\ &
+ 1224\,X_{47}\,X_{65}\\ &
+ 4\,X_{64}\,X_{65}\\ &
+ 60060\,X_{65}\\ &
- X_{124}
\end{aligned}\end{equation*}
\paragraph{Row 70.} \nolinkurl{residue.det0.factor0.0.reciprocal.definition}\nopagebreak
\begin{equation*}\begin{aligned}r_{70,N}={}&
 X_{23}^{2}\\ &
+ 8\,X_{24}^{2}\\ &
+ X_{66}\,X_{81}\\ &
+ X_{68}\,X_{79}\\ &
+ X_{82}\\ &
- X_{82}\,X_{125}
\end{aligned}\end{equation*}
\paragraph{Row 71.} \nolinkurl{residue.det0.factor0.1.reciprocal.definition}\nopagebreak
\begin{equation*}\begin{aligned}r_{71,N}={}&
 4\,X_{23}\,X_{24}\\ &
- 2\,X_{67}\,X_{81}\\ &
- 2\,X_{69}\,X_{79}\\ &
+ X_{82}\,X_{126}
\end{aligned}\end{equation*}
\paragraph{Row 72.} \nolinkurl{residue.det0.factor1.0.reciprocal.definition}\nopagebreak
\begin{equation*}\begin{aligned}r_{72,N}={}&
 X_{21}^{2}\\ &
+ 8\,X_{22}^{2}\\ &
- 2\,X_{52}\,X_{79}\\ &
+ X_{58}\,X_{81}\\ &
+ X_{74}\,X_{79}\\ &
+ X_{79}\\ &
+ X_{82}\\ &
- X_{82}\,X_{127}
\end{aligned}\end{equation*}
\paragraph{Row 73.} \nolinkurl{residue.det0.factor1.1.reciprocal.definition}\nopagebreak
\begin{equation*}\begin{aligned}r_{73,N}={}&
 4\,X_{21}\,X_{22}\\ &
+ 4\,X_{53}\,X_{79}\\ &
- 2\,X_{59}\,X_{81}\\ &
- 4\,X_{75}\,X_{79}\\ &
+ X_{82}\,X_{128}
\end{aligned}\end{equation*}
\paragraph{Row 74.} \nolinkurl{residue.det0.factor2.0.reciprocal.definition}\nopagebreak
\begin{equation*}\begin{aligned}r_{74,N}={}&
 X_{27}^{2}\\ &
+ 8\,X_{28}^{2}\\ &
+ X_{56}\,X_{80}\\ &
+ X_{60}\,X_{81}\\ &
+ X_{79}\,X_{123}\\ &
+ X_{82}\\ &
- X_{82}\,X_{129}
\end{aligned}\end{equation*}
\paragraph{Row 75.} \nolinkurl{residue.det0.factor2.1.reciprocal.definition}\nopagebreak
\begin{equation*}\begin{aligned}r_{75,N}={}&
 4\,X_{27}\,X_{28}\\ &
- 2\,X_{57}\,X_{80}\\ &
- 2\,X_{61}\,X_{81}\\ &
- X_{79}\,X_{124}\\ &
+ X_{82}\,X_{130}
\end{aligned}\end{equation*}
\paragraph{Row 76.} \nolinkurl{residue.det0.factor3.0.reciprocal.definition}\nopagebreak
\begin{equation*}\begin{aligned}r_{76,N}={}&
 X_{29}^{2}\\ &
+ 8\,X_{30}^{2}\\ &
+ X_{70}\,X_{79}\\ &
+ X_{81}\,X_{121}\\ &
+ X_{82}\\ &
- X_{82}\,X_{131}
\end{aligned}\end{equation*}
\paragraph{Row 77.} \nolinkurl{residue.det0.factor3.1.reciprocal.definition}\nopagebreak
\begin{equation*}\begin{aligned}r_{77,N}={}&
 4\,X_{29}\,X_{30}\\ &
- 2\,X_{71}\,X_{79}\\ &
- X_{81}\,X_{122}\\ &
+ X_{82}\,X_{132}
\end{aligned}\end{equation*}
\paragraph{Row 78.} \nolinkurl{residue.det0.factor4.0.reciprocal.definition}\nopagebreak
\begin{equation*}\begin{aligned}r_{78,N}={}&
 X_{25}^{2}\\ &
+ 8\,X_{26}^{2}\\ &
+ X_{54}\,X_{81}\\ &
+ X_{72}\,X_{79}\\ &
+ X_{82}\\ &
- X_{82}\,X_{133}
\end{aligned}\end{equation*}
\paragraph{Row 79.} \nolinkurl{residue.det0.factor4.1.reciprocal.definition}\nopagebreak
\begin{equation*}\begin{aligned}r_{79,N}={}&
 4\,X_{25}\,X_{26}\\ &
- 2\,X_{55}\,X_{81}\\ &
- 2\,X_{73}\,X_{79}\\ &
+ X_{82}\,X_{134}
\end{aligned}\end{equation*}
\paragraph{Row 80.} \nolinkurl{amplified.det0.product1.0.definition}\nopagebreak
\begin{equation*}\begin{aligned}r_{80,N}={}&
 X_{125}\,X_{127}\\ &
+ 2\,X_{126}\,X_{128}\\ &
- X_{135}
\end{aligned}\end{equation*}
\paragraph{Row 81.} \nolinkurl{amplified.det0.product1.1.definition}\nopagebreak
\begin{equation*}\begin{aligned}r_{81,N}={}&
 X_{125}\,X_{128}\\ &
+ X_{126}\,X_{127}\\ &
- X_{136}
\end{aligned}\end{equation*}
\paragraph{Row 82.} \nolinkurl{amplified.det0.product2.0.definition}\nopagebreak
\begin{equation*}\begin{aligned}r_{82,N}={}&
 X_{129}\,X_{135}\\ &
+ 2\,X_{130}\,X_{136}\\ &
- X_{137}
\end{aligned}\end{equation*}
\paragraph{Row 83.} \nolinkurl{amplified.det0.product2.1.definition}\nopagebreak
\begin{equation*}\begin{aligned}r_{83,N}={}&
 X_{129}\,X_{136}\\ &
+ X_{130}\,X_{135}\\ &
- X_{138}
\end{aligned}\end{equation*}
\paragraph{Row 84.} \nolinkurl{amplified.det0.product3.0.definition}\nopagebreak
\begin{equation*}\begin{aligned}r_{84,N}={}&
 X_{131}\,X_{137}\\ &
+ 2\,X_{132}\,X_{138}\\ &
- X_{139}
\end{aligned}\end{equation*}
\paragraph{Row 85.} \nolinkurl{amplified.det0.product3.1.definition}\nopagebreak
\begin{equation*}\begin{aligned}r_{85,N}={}&
 X_{131}\,X_{138}\\ &
+ X_{132}\,X_{137}\\ &
- X_{140}
\end{aligned}\end{equation*}
\paragraph{Row 86.} \nolinkurl{amplified.det0.product4.0.definition}\nopagebreak
\begin{equation*}\begin{aligned}r_{86,N}={}&
 X_{83}\\ &
- X_{133}\,X_{139}\\ &
- 2\,X_{134}\,X_{140}
\end{aligned}\end{equation*}
\paragraph{Row 87.} \nolinkurl{amplified.det0.product4.1.definition}\nopagebreak
\begin{equation*}\begin{aligned}r_{87,N}={}&
 X_{84}\\ &
- X_{133}\,X_{140}\\ &
- X_{134}\,X_{139}
\end{aligned}\end{equation*}
\paragraph{Row 88.} \nolinkurl{residue.det1.factor0.0.reciprocal.definition}\nopagebreak
\begin{equation*}\begin{aligned}r_{88,N}={}&
 X_{23}^{2}\\ &
+ 8\,X_{24}^{2}\\ &
+ X_{66}\,X_{96}\\ &
+ X_{68}\,X_{94}\\ &
+ X_{97}\\ &
- X_{97}\,X_{141}
\end{aligned}\end{equation*}
\paragraph{Row 89.} \nolinkurl{residue.det1.factor0.1.reciprocal.definition}\nopagebreak
\begin{equation*}\begin{aligned}r_{89,N}={}&
 4\,X_{23}\,X_{24}\\ &
- 2\,X_{67}\,X_{96}\\ &
- 2\,X_{69}\,X_{94}\\ &
+ X_{97}\,X_{142}
\end{aligned}\end{equation*}
\paragraph{Row 90.} \nolinkurl{residue.det1.factor1.0.reciprocal.definition}\nopagebreak
\begin{equation*}\begin{aligned}r_{90,N}={}&
 X_{21}^{2}\\ &
+ 8\,X_{22}^{2}\\ &
- 2\,X_{52}\,X_{94}\\ &
+ X_{58}\,X_{96}\\ &
+ X_{74}\,X_{94}\\ &
+ X_{94}\\ &
+ X_{97}\\ &
- X_{97}\,X_{143}
\end{aligned}\end{equation*}
\paragraph{Row 91.} \nolinkurl{residue.det1.factor1.1.reciprocal.definition}\nopagebreak
\begin{equation*}\begin{aligned}r_{91,N}={}&
 4\,X_{21}\,X_{22}\\ &
+ 4\,X_{53}\,X_{94}\\ &
- 2\,X_{59}\,X_{96}\\ &
- 4\,X_{75}\,X_{94}\\ &
+ X_{97}\,X_{144}
\end{aligned}\end{equation*}
\paragraph{Row 92.} \nolinkurl{residue.det1.factor2.0.reciprocal.definition}\nopagebreak
\begin{equation*}\begin{aligned}r_{92,N}={}&
 X_{27}^{2}\\ &
+ 8\,X_{28}^{2}\\ &
+ X_{56}\,X_{95}\\ &
+ X_{60}\,X_{96}\\ &
+ X_{94}\,X_{123}\\ &
+ X_{97}\\ &
- X_{97}\,X_{145}
\end{aligned}\end{equation*}
\paragraph{Row 93.} \nolinkurl{residue.det1.factor2.1.reciprocal.definition}\nopagebreak
\begin{equation*}\begin{aligned}r_{93,N}={}&
 4\,X_{27}\,X_{28}\\ &
- 2\,X_{57}\,X_{95}\\ &
- 2\,X_{61}\,X_{96}\\ &
- X_{94}\,X_{124}\\ &
+ X_{97}\,X_{146}
\end{aligned}\end{equation*}
\paragraph{Row 94.} \nolinkurl{residue.det1.factor3.0.reciprocal.definition}\nopagebreak
\begin{equation*}\begin{aligned}r_{94,N}={}&
 X_{29}^{2}\\ &
+ 8\,X_{30}^{2}\\ &
+ X_{70}\,X_{94}\\ &
+ X_{96}\,X_{121}\\ &
+ X_{97}\\ &
- X_{97}\,X_{147}
\end{aligned}\end{equation*}
\paragraph{Row 95.} \nolinkurl{residue.det1.factor3.1.reciprocal.definition}\nopagebreak
\begin{equation*}\begin{aligned}r_{95,N}={}&
 4\,X_{29}\,X_{30}\\ &
- 2\,X_{71}\,X_{94}\\ &
- X_{96}\,X_{122}\\ &
+ X_{97}\,X_{148}
\end{aligned}\end{equation*}
\paragraph{Row 96.} \nolinkurl{residue.det1.factor4.0.reciprocal.definition}\nopagebreak
\begin{equation*}\begin{aligned}r_{96,N}={}&
 X_{25}^{2}\\ &
+ 8\,X_{26}^{2}\\ &
+ X_{54}\,X_{96}\\ &
+ X_{72}\,X_{94}\\ &
+ X_{97}\\ &
- X_{97}\,X_{149}
\end{aligned}\end{equation*}
\paragraph{Row 97.} \nolinkurl{residue.det1.factor4.1.reciprocal.definition}\nopagebreak
\begin{equation*}\begin{aligned}r_{97,N}={}&
 4\,X_{25}\,X_{26}\\ &
- 2\,X_{55}\,X_{96}\\ &
- 2\,X_{73}\,X_{94}\\ &
+ X_{97}\,X_{150}
\end{aligned}\end{equation*}
\paragraph{Row 98.} \nolinkurl{amplified.det1.product1.0.definition}\nopagebreak
\begin{equation*}\begin{aligned}r_{98,N}={}&
 X_{141}\,X_{143}\\ &
+ 2\,X_{142}\,X_{144}\\ &
- X_{151}
\end{aligned}\end{equation*}
\paragraph{Row 99.} \nolinkurl{amplified.det1.product1.1.definition}\nopagebreak
\begin{equation*}\begin{aligned}r_{99,N}={}&
 X_{141}\,X_{144}\\ &
+ X_{142}\,X_{143}\\ &
- X_{152}
\end{aligned}\end{equation*}
\paragraph{Row 100.} \nolinkurl{amplified.det1.product2.0.definition}\nopagebreak
\begin{equation*}\begin{aligned}r_{100,N}={}&
 X_{145}\,X_{151}\\ &
+ 2\,X_{146}\,X_{152}\\ &
- X_{153}
\end{aligned}\end{equation*}
\paragraph{Row 101.} \nolinkurl{amplified.det1.product2.1.definition}\nopagebreak
\begin{equation*}\begin{aligned}r_{101,N}={}&
 X_{145}\,X_{152}\\ &
+ X_{146}\,X_{151}\\ &
- X_{154}
\end{aligned}\end{equation*}
\paragraph{Row 102.} \nolinkurl{amplified.det1.product3.0.definition}\nopagebreak
\begin{equation*}\begin{aligned}r_{102,N}={}&
 X_{147}\,X_{153}\\ &
+ 2\,X_{148}\,X_{154}\\ &
- X_{155}
\end{aligned}\end{equation*}
\paragraph{Row 103.} \nolinkurl{amplified.det1.product3.1.definition}\nopagebreak
\begin{equation*}\begin{aligned}r_{103,N}={}&
 X_{147}\,X_{154}\\ &
+ X_{148}\,X_{153}\\ &
- X_{156}
\end{aligned}\end{equation*}
\paragraph{Row 104.} \nolinkurl{amplified.det1.product4.0.definition}\nopagebreak
\begin{equation*}\begin{aligned}r_{104,N}={}&
 X_{98}\\ &
- X_{149}\,X_{155}\\ &
- 2\,X_{150}\,X_{156}
\end{aligned}\end{equation*}
\paragraph{Row 105.} \nolinkurl{amplified.det1.product4.1.definition}\nopagebreak
\begin{equation*}\begin{aligned}r_{105,N}={}&
 X_{99}\\ &
- X_{149}\,X_{156}\\ &
- X_{150}\,X_{155}
\end{aligned}\end{equation*}
\paragraph{Row 106.} \nolinkurl{rootProduct.1.equation1.0}\nopagebreak
\begin{equation*}\begin{aligned}r_{106,N}={}&
 X_{157}^{2}\\ &
+ 2\,X_{158}^{2}\\ &
+ X_{159}^{2}\\ &
+ 2\,X_{160}^{2}\\ &
+ 2\,X_{161}\,X_{163}\\ &
+ 4\,X_{162}\,X_{164}
\end{aligned}\end{equation*}
\paragraph{Row 107.} \nolinkurl{rootProduct.1.equation1.1}\nopagebreak
\begin{equation*}\begin{aligned}r_{107,N}={}&
 X_{157}\,X_{158}\\ &
+ X_{159}\,X_{160}\\ &
+ X_{161}\,X_{164}\\ &
+ X_{162}\,X_{163}
\end{aligned}\end{equation*}
\paragraph{Row 108.} \nolinkurl{rootProduct.1.equation2.0}\nopagebreak
\begin{equation*}\begin{aligned}r_{108,N}={}&
 3\,X_{118}\,X_{165}\\ &
- 2\,X_{118}\,X_{166}\\ &
- 3\,X_{118}\,X_{167}\\ &
+ 2\,X_{118}\,X_{168}\\ &
- X_{157}\,X_{161}\\ &
- 2\,X_{158}\,X_{162}
\end{aligned}\end{equation*}
\paragraph{Row 109.} \nolinkurl{rootProduct.1.equation2.1}\nopagebreak
\begin{equation*}\begin{aligned}r_{109,N}={}&
 X_{118}\,X_{165}\\ &
- 3\,X_{118}\,X_{166}\\ &
- X_{118}\,X_{167}\\ &
+ 3\,X_{118}\,X_{168}\\ &
+ X_{157}\,X_{162}\\ &
+ X_{158}\,X_{161}
\end{aligned}\end{equation*}
\paragraph{Row 110.} \nolinkurl{rootProduct.1.equation3.0}\nopagebreak
\begin{equation*}\begin{aligned}r_{110,N}={}&
 3\,X_{118}\,X_{165}\\ &
- 2\,X_{118}\,X_{166}\\ &
+ 3\,X_{118}\,X_{167}\\ &
- 2\,X_{118}\,X_{168}\\ &
- X_{159}\,X_{163}\\ &
- 2\,X_{160}\,X_{164}
\end{aligned}\end{equation*}
\paragraph{Row 111.} \nolinkurl{rootProduct.1.equation3.1}\nopagebreak
\begin{equation*}\begin{aligned}r_{111,N}={}&
 X_{118}\,X_{165}\\ &
- 3\,X_{118}\,X_{166}\\ &
+ X_{118}\,X_{167}\\ &
- 3\,X_{118}\,X_{168}\\ &
+ X_{159}\,X_{164}\\ &
+ X_{160}\,X_{163}
\end{aligned}\end{equation*}
\paragraph{Row 112.} \nolinkurl{rootProduct.1.equation4.0}\nopagebreak
\begin{equation*}\begin{aligned}r_{112,N}={}&
 44\,X_{117}\,X_{169}\\ &
- 48\,X_{117}\,X_{170}\\ &
- X_{157}^{2}\\ &
- 2\,X_{157}\,X_{159}\\ &
- 2\,X_{158}^{2}\\ &
- 4\,X_{158}\,X_{160}\\ &
- X_{159}^{2}\\ &
- 2\,X_{160}^{2}
\end{aligned}\end{equation*}
\paragraph{Row 113.} \nolinkurl{rootProduct.1.equation4.1}\nopagebreak
\begin{equation*}\begin{aligned}r_{113,N}={}&
 12\,X_{117}\,X_{169}\\ &
- 22\,X_{117}\,X_{170}\\ &
+ X_{157}\,X_{158}\\ &
+ X_{157}\,X_{160}\\ &
+ X_{158}\,X_{159}\\ &
+ X_{159}\,X_{160}
\end{aligned}\end{equation*}
\paragraph{Row 114.} \nolinkurl{rootProduct.1.equation5.0}\nopagebreak
\begin{equation*}\begin{aligned}r_{114,N}={}&
 X_{54}\\ &
- X_{169}^{2}\\ &
- 2\,X_{170}^{2}
\end{aligned}\end{equation*}
\paragraph{Row 115.} \nolinkurl{rootProduct.1.equation5.1}\nopagebreak
\begin{equation*}\begin{aligned}r_{115,N}={}&
 X_{55}\\ &
- X_{169}\,X_{170}
\end{aligned}\end{equation*}
\paragraph{Row 116.} \nolinkurl{rootProduct.1.equation6.0}\nopagebreak
\begin{equation*}\begin{aligned}r_{116,N}={}&
 11\,X_{54}\,X_{117}\\ &
- 24\,X_{55}\,X_{117}\\ &
- X_{117}\\ &
+ X_{165}^{2}\\ &
+ 2\,X_{166}^{2}
\end{aligned}\end{equation*}
\paragraph{Row 117.} \nolinkurl{rootProduct.1.equation6.1}\nopagebreak
\begin{equation*}\begin{aligned}r_{117,N}={}&
 3\,X_{54}\,X_{117}\\ &
- 11\,X_{55}\,X_{117}\\ &
- X_{165}\,X_{166}
\end{aligned}\end{equation*}
\paragraph{Row 118.} \nolinkurl{rootProduct.1.equation7.0}\nopagebreak
\begin{equation*}\begin{aligned}r_{118,N}={}&
 11\,X_{23}\,X_{117}\\ &
- 24\,X_{24}\,X_{117}\\ &
- X_{117}\\ &
+ X_{167}^{2}\\ &
+ 2\,X_{168}^{2}
\end{aligned}\end{equation*}
\paragraph{Row 119.} \nolinkurl{rootProduct.1.equation7.1}\nopagebreak
\begin{equation*}\begin{aligned}r_{119,N}={}&
 3\,X_{23}\,X_{117}\\ &
- 11\,X_{24}\,X_{117}\\ &
- X_{167}\,X_{168}
\end{aligned}\end{equation*}
\paragraph{Row 120.} \nolinkurl{alternate.direct.1.0}\nopagebreak
\begin{equation*}\begin{aligned}r_{120,N}={}&
 6\,X_{33}\,X_{167}\\ &
- 4\,X_{33}\,X_{168}\\ &
- 4\,X_{34}\,X_{167}\\ &
+ 12\,X_{34}\,X_{168}\\ &
- X_{157}\,X_{213}\\ &
- 2\,X_{158}\,X_{214}\\ &
+ X_{159}\,X_{213}\\ &
+ 2\,X_{160}\,X_{214}
\end{aligned}\end{equation*}
\paragraph{Row 121.} \nolinkurl{alternate.direct.1.1}\nopagebreak
\begin{equation*}\begin{aligned}r_{121,N}={}&
 2\,X_{33}\,X_{167}\\ &
- 6\,X_{33}\,X_{168}\\ &
- 6\,X_{34}\,X_{167}\\ &
+ 4\,X_{34}\,X_{168}\\ &
+ X_{157}\,X_{214}\\ &
+ X_{158}\,X_{213}\\ &
- X_{159}\,X_{214}\\ &
- X_{160}\,X_{213}
\end{aligned}\end{equation*}
\paragraph{Row 122.} \nolinkurl{rootProduct.1.equation9.0}\nopagebreak
\begin{equation*}\begin{aligned}r_{122,N}={}&
 X_{19}\,X_{157}\\ &
+ X_{19}\,X_{159}\\ &
+ 2\,X_{20}\,X_{158}\\ &
+ 2\,X_{20}\,X_{160}\\ &
- 6\,X_{31}\,X_{165}\\ &
+ 4\,X_{31}\,X_{166}\\ &
+ 4\,X_{32}\,X_{165}\\ &
- 12\,X_{32}\,X_{166}
\end{aligned}\end{equation*}
\paragraph{Row 123.} \nolinkurl{rootProduct.1.equation9.1}\nopagebreak
\begin{equation*}\begin{aligned}r_{123,N}={}&
 X_{19}\,X_{158}\\ &
+ X_{19}\,X_{160}\\ &
+ X_{20}\,X_{157}\\ &
+ X_{20}\,X_{159}\\ &
+ 2\,X_{31}\,X_{165}\\ &
- 6\,X_{31}\,X_{166}\\ &
- 6\,X_{32}\,X_{165}\\ &
+ 4\,X_{32}\,X_{166}
\end{aligned}\end{equation*}
\paragraph{Row 124.} \nolinkurl{rootProduct.2.equation1.0}\nopagebreak
\begin{equation*}\begin{aligned}r_{124,N}={}&
 X_{171}^{2}\\ &
+ 2\,X_{172}^{2}\\ &
+ X_{173}^{2}\\ &
+ 2\,X_{174}^{2}\\ &
+ 2\,X_{175}\,X_{177}\\ &
+ 4\,X_{176}\,X_{178}
\end{aligned}\end{equation*}
\paragraph{Row 125.} \nolinkurl{rootProduct.2.equation1.1}\nopagebreak
\begin{equation*}\begin{aligned}r_{125,N}={}&
 X_{171}\,X_{172}\\ &
+ X_{173}\,X_{174}\\ &
+ X_{175}\,X_{178}\\ &
+ X_{176}\,X_{177}
\end{aligned}\end{equation*}
\paragraph{Row 126.} \nolinkurl{rootProduct.2.equation2.0}\nopagebreak
\begin{equation*}\begin{aligned}r_{126,N}={}&
 3\,X_{117}\,X_{179}\\ &
- 2\,X_{117}\,X_{180}\\ &
- 3\,X_{117}\,X_{181}\\ &
+ 2\,X_{117}\,X_{182}\\ &
- X_{171}\,X_{175}\\ &
- 2\,X_{172}\,X_{176}
\end{aligned}\end{equation*}
\paragraph{Row 127.} \nolinkurl{rootProduct.2.equation2.1}\nopagebreak
\begin{equation*}\begin{aligned}r_{127,N}={}&
 X_{117}\,X_{179}\\ &
- 3\,X_{117}\,X_{180}\\ &
- X_{117}\,X_{181}\\ &
+ 3\,X_{117}\,X_{182}\\ &
+ X_{171}\,X_{176}\\ &
+ X_{172}\,X_{175}
\end{aligned}\end{equation*}
\paragraph{Row 128.} \nolinkurl{rootProduct.2.equation3.0}\nopagebreak
\begin{equation*}\begin{aligned}r_{128,N}={}&
 3\,X_{117}\,X_{179}\\ &
- 2\,X_{117}\,X_{180}\\ &
+ 3\,X_{117}\,X_{181}\\ &
- 2\,X_{117}\,X_{182}\\ &
- X_{173}\,X_{177}\\ &
- 2\,X_{174}\,X_{178}
\end{aligned}\end{equation*}
\paragraph{Row 129.} \nolinkurl{rootProduct.2.equation3.1}\nopagebreak
\begin{equation*}\begin{aligned}r_{129,N}={}&
 X_{117}\,X_{179}\\ &
- 3\,X_{117}\,X_{180}\\ &
+ X_{117}\,X_{181}\\ &
- 3\,X_{117}\,X_{182}\\ &
+ X_{173}\,X_{178}\\ &
+ X_{174}\,X_{177}
\end{aligned}\end{equation*}
\paragraph{Row 130.} \nolinkurl{rootProduct.2.equation4.0}\nopagebreak
\begin{equation*}\begin{aligned}r_{130,N}={}&
 44\,X_{117}\,X_{183}\\ &
- 48\,X_{117}\,X_{184}\\ &
- X_{171}^{2}\\ &
- 2\,X_{171}\,X_{173}\\ &
- 2\,X_{172}^{2}\\ &
- 4\,X_{172}\,X_{174}\\ &
- X_{173}^{2}\\ &
- 2\,X_{174}^{2}
\end{aligned}\end{equation*}
\paragraph{Row 131.} \nolinkurl{rootProduct.2.equation4.1}\nopagebreak
\begin{equation*}\begin{aligned}r_{131,N}={}&
 12\,X_{117}\,X_{183}\\ &
- 22\,X_{117}\,X_{184}\\ &
+ X_{171}\,X_{172}\\ &
+ X_{171}\,X_{174}\\ &
+ X_{172}\,X_{173}\\ &
+ X_{173}\,X_{174}
\end{aligned}\end{equation*}
\paragraph{Row 132.} \nolinkurl{rootProduct.2.equation5.0}\nopagebreak
\begin{equation*}\begin{aligned}r_{132,N}={}&
 X_{56}\\ &
- X_{183}^{2}\\ &
- 2\,X_{184}^{2}
\end{aligned}\end{equation*}
\paragraph{Row 133.} \nolinkurl{rootProduct.2.equation5.1}\nopagebreak
\begin{equation*}\begin{aligned}r_{133,N}={}&
 X_{57}\\ &
- X_{183}\,X_{184}
\end{aligned}\end{equation*}
\paragraph{Row 134.} \nolinkurl{rootProduct.2.equation6.0}\nopagebreak
\begin{equation*}\begin{aligned}r_{134,N}={}&
 1\\ &
- 11\,X_{56}\\ &
+ 24\,X_{57}\\ &
- X_{179}^{2}\\ &
- 2\,X_{180}^{2}
\end{aligned}\end{equation*}
\paragraph{Row 135.} \nolinkurl{rootProduct.2.equation6.1}\nopagebreak
\begin{equation*}\begin{aligned}r_{135,N}={}&
 3\,X_{56}\\ &
- 11\,X_{57}\\ &
- X_{179}\,X_{180}
\end{aligned}\end{equation*}
\paragraph{Row 136.} \nolinkurl{rootProduct.2.equation7.0}\nopagebreak
\begin{equation*}\begin{aligned}r_{136,N}={}&
 1\\ &
- 11\,X_{25}\\ &
+ 24\,X_{26}\\ &
- X_{181}^{2}\\ &
- 2\,X_{182}^{2}
\end{aligned}\end{equation*}
\paragraph{Row 137.} \nolinkurl{rootProduct.2.equation7.1}\nopagebreak
\begin{equation*}\begin{aligned}r_{137,N}={}&
 3\,X_{25}\\ &
- 11\,X_{26}\\ &
- X_{181}\,X_{182}
\end{aligned}\end{equation*}
\paragraph{Row 138.} \nolinkurl{alternate.direct.2.0}\nopagebreak
\begin{equation*}\begin{aligned}r_{138,N}={}&
 6\,X_{37}\,X_{181}\\ &
- 4\,X_{37}\,X_{182}\\ &
- 4\,X_{38}\,X_{181}\\ &
+ 12\,X_{38}\,X_{182}\\ &
- X_{171}\,X_{213}\\ &
- 2\,X_{172}\,X_{214}\\ &
+ X_{173}\,X_{213}\\ &
+ 2\,X_{174}\,X_{214}
\end{aligned}\end{equation*}
\paragraph{Row 139.} \nolinkurl{alternate.direct.2.1}\nopagebreak
\begin{equation*}\begin{aligned}r_{139,N}={}&
 2\,X_{37}\,X_{181}\\ &
- 6\,X_{37}\,X_{182}\\ &
- 6\,X_{38}\,X_{181}\\ &
+ 4\,X_{38}\,X_{182}\\ &
+ X_{171}\,X_{214}\\ &
+ X_{172}\,X_{213}\\ &
- X_{173}\,X_{214}\\ &
- X_{174}\,X_{213}
\end{aligned}\end{equation*}
\paragraph{Row 140.} \nolinkurl{rootProduct.2.equation9.0}\nopagebreak
\begin{equation*}\begin{aligned}r_{140,N}={}&
 X_{19}\,X_{171}\\ &
+ X_{19}\,X_{173}\\ &
+ 2\,X_{20}\,X_{172}\\ &
+ 2\,X_{20}\,X_{174}\\ &
- 6\,X_{35}\,X_{179}\\ &
+ 4\,X_{35}\,X_{180}\\ &
+ 4\,X_{36}\,X_{179}\\ &
- 12\,X_{36}\,X_{180}
\end{aligned}\end{equation*}
\paragraph{Row 141.} \nolinkurl{rootProduct.2.equation9.1}\nopagebreak
\begin{equation*}\begin{aligned}r_{141,N}={}&
 X_{19}\,X_{172}\\ &
+ X_{19}\,X_{174}\\ &
+ X_{20}\,X_{171}\\ &
+ X_{20}\,X_{173}\\ &
+ 2\,X_{35}\,X_{179}\\ &
- 6\,X_{35}\,X_{180}\\ &
- 6\,X_{36}\,X_{179}\\ &
+ 4\,X_{36}\,X_{180}
\end{aligned}\end{equation*}
\paragraph{Row 142.} \nolinkurl{rootProduct.3.equation1.0}\nopagebreak
\begin{equation*}\begin{aligned}r_{142,N}={}&
 X_{185}^{2}\\ &
+ 2\,X_{186}^{2}\\ &
+ X_{187}^{2}\\ &
+ 2\,X_{188}^{2}\\ &
+ 2\,X_{189}\,X_{191}\\ &
+ 4\,X_{190}\,X_{192}
\end{aligned}\end{equation*}
\paragraph{Row 143.} \nolinkurl{rootProduct.3.equation1.1}\nopagebreak
\begin{equation*}\begin{aligned}r_{143,N}={}&
 X_{185}\,X_{186}\\ &
+ X_{187}\,X_{188}\\ &
+ X_{189}\,X_{192}\\ &
+ X_{190}\,X_{191}
\end{aligned}\end{equation*}
\paragraph{Row 144.} \nolinkurl{rootProduct.3.equation2.0}\nopagebreak
\begin{equation*}\begin{aligned}r_{144,N}={}&
 3\,X_{117}\,X_{193}\\ &
- 2\,X_{117}\,X_{194}\\ &
- 3\,X_{117}\,X_{195}\\ &
+ 2\,X_{117}\,X_{196}\\ &
- X_{185}\,X_{189}\\ &
- 2\,X_{186}\,X_{190}
\end{aligned}\end{equation*}
\paragraph{Row 145.} \nolinkurl{rootProduct.3.equation2.1}\nopagebreak
\begin{equation*}\begin{aligned}r_{145,N}={}&
 X_{117}\,X_{193}\\ &
- 3\,X_{117}\,X_{194}\\ &
- X_{117}\,X_{195}\\ &
+ 3\,X_{117}\,X_{196}\\ &
+ X_{185}\,X_{190}\\ &
+ X_{186}\,X_{189}
\end{aligned}\end{equation*}
\paragraph{Row 146.} \nolinkurl{rootProduct.3.equation3.0}\nopagebreak
\begin{equation*}\begin{aligned}r_{146,N}={}&
 3\,X_{117}\,X_{193}\\ &
- 2\,X_{117}\,X_{194}\\ &
+ 3\,X_{117}\,X_{195}\\ &
- 2\,X_{117}\,X_{196}\\ &
- X_{187}\,X_{191}\\ &
- 2\,X_{188}\,X_{192}
\end{aligned}\end{equation*}
\paragraph{Row 147.} \nolinkurl{rootProduct.3.equation3.1}\nopagebreak
\begin{equation*}\begin{aligned}r_{147,N}={}&
 X_{117}\,X_{193}\\ &
- 3\,X_{117}\,X_{194}\\ &
+ X_{117}\,X_{195}\\ &
- 3\,X_{117}\,X_{196}\\ &
+ X_{187}\,X_{192}\\ &
+ X_{188}\,X_{191}
\end{aligned}\end{equation*}
\paragraph{Row 148.} \nolinkurl{rootProduct.3.equation4.0}\nopagebreak
\begin{equation*}\begin{aligned}r_{148,N}={}&
 44\,X_{117}\,X_{197}\\ &
- 48\,X_{117}\,X_{198}\\ &
- X_{185}^{2}\\ &
- 2\,X_{185}\,X_{187}\\ &
- 2\,X_{186}^{2}\\ &
- 4\,X_{186}\,X_{188}\\ &
- X_{187}^{2}\\ &
- 2\,X_{188}^{2}
\end{aligned}\end{equation*}
\paragraph{Row 149.} \nolinkurl{rootProduct.3.equation4.1}\nopagebreak
\begin{equation*}\begin{aligned}r_{149,N}={}&
 12\,X_{117}\,X_{197}\\ &
- 22\,X_{117}\,X_{198}\\ &
+ X_{185}\,X_{186}\\ &
+ X_{185}\,X_{188}\\ &
+ X_{186}\,X_{187}\\ &
+ X_{187}\,X_{188}
\end{aligned}\end{equation*}
\paragraph{Row 150.} \nolinkurl{rootProduct.3.equation5.0}\nopagebreak
\begin{equation*}\begin{aligned}r_{150,N}={}&
 X_{58}\\ &
- X_{197}^{2}\\ &
- 2\,X_{198}^{2}
\end{aligned}\end{equation*}
\paragraph{Row 151.} \nolinkurl{rootProduct.3.equation5.1}\nopagebreak
\begin{equation*}\begin{aligned}r_{151,N}={}&
 X_{59}\\ &
- X_{197}\,X_{198}
\end{aligned}\end{equation*}
\paragraph{Row 152.} \nolinkurl{rootProduct.3.equation6.0}\nopagebreak
\begin{equation*}\begin{aligned}r_{152,N}={}&
 1\\ &
- 11\,X_{58}\\ &
+ 24\,X_{59}\\ &
- X_{193}^{2}\\ &
- 2\,X_{194}^{2}
\end{aligned}\end{equation*}
\paragraph{Row 153.} \nolinkurl{rootProduct.3.equation6.1}\nopagebreak
\begin{equation*}\begin{aligned}r_{153,N}={}&
 3\,X_{58}\\ &
- 11\,X_{59}\\ &
- X_{193}\,X_{194}
\end{aligned}\end{equation*}
\paragraph{Row 154.} \nolinkurl{rootProduct.3.equation7.0}\nopagebreak
\begin{equation*}\begin{aligned}r_{154,N}={}&
 1\\ &
- 11\,X_{27}\\ &
+ 24\,X_{28}\\ &
- X_{195}^{2}\\ &
- 2\,X_{196}^{2}
\end{aligned}\end{equation*}
\paragraph{Row 155.} \nolinkurl{rootProduct.3.equation7.1}\nopagebreak
\begin{equation*}\begin{aligned}r_{155,N}={}&
 3\,X_{27}\\ &
- 11\,X_{28}\\ &
- X_{195}\,X_{196}
\end{aligned}\end{equation*}
\paragraph{Row 156.} \nolinkurl{alternate.direct.3.0}\nopagebreak
\begin{equation*}\begin{aligned}r_{156,N}={}&
 6\,X_{41}\,X_{195}\\ &
- 4\,X_{41}\,X_{196}\\ &
- 4\,X_{42}\,X_{195}\\ &
+ 12\,X_{42}\,X_{196}\\ &
- X_{185}\,X_{213}\\ &
- 2\,X_{186}\,X_{214}\\ &
+ X_{187}\,X_{213}\\ &
+ 2\,X_{188}\,X_{214}
\end{aligned}\end{equation*}
\paragraph{Row 157.} \nolinkurl{alternate.direct.3.1}\nopagebreak
\begin{equation*}\begin{aligned}r_{157,N}={}&
 2\,X_{41}\,X_{195}\\ &
- 6\,X_{41}\,X_{196}\\ &
- 6\,X_{42}\,X_{195}\\ &
+ 4\,X_{42}\,X_{196}\\ &
+ X_{185}\,X_{214}\\ &
+ X_{186}\,X_{213}\\ &
- X_{187}\,X_{214}\\ &
- X_{188}\,X_{213}
\end{aligned}\end{equation*}
\paragraph{Row 158.} \nolinkurl{rootProduct.3.equation9.0}\nopagebreak
\begin{equation*}\begin{aligned}r_{158,N}={}&
 X_{19}\,X_{185}\\ &
+ X_{19}\,X_{187}\\ &
+ 2\,X_{20}\,X_{186}\\ &
+ 2\,X_{20}\,X_{188}\\ &
- 6\,X_{39}\,X_{193}\\ &
+ 4\,X_{39}\,X_{194}\\ &
+ 4\,X_{40}\,X_{193}\\ &
- 12\,X_{40}\,X_{194}
\end{aligned}\end{equation*}
\paragraph{Row 159.} \nolinkurl{rootProduct.3.equation9.1}\nopagebreak
\begin{equation*}\begin{aligned}r_{159,N}={}&
 X_{19}\,X_{186}\\ &
+ X_{19}\,X_{188}\\ &
+ X_{20}\,X_{185}\\ &
+ X_{20}\,X_{187}\\ &
+ 2\,X_{39}\,X_{193}\\ &
- 6\,X_{39}\,X_{194}\\ &
- 6\,X_{40}\,X_{193}\\ &
+ 4\,X_{40}\,X_{194}
\end{aligned}\end{equation*}
\paragraph{Row 160.} \nolinkurl{rootProduct.4.equation1.0}\nopagebreak
\begin{equation*}\begin{aligned}r_{160,N}={}&
 X_{199}^{2}\\ &
+ 2\,X_{200}^{2}\\ &
+ X_{201}^{2}\\ &
+ 2\,X_{202}^{2}\\ &
+ 2\,X_{203}\,X_{205}\\ &
+ 4\,X_{204}\,X_{206}
\end{aligned}\end{equation*}
\paragraph{Row 161.} \nolinkurl{rootProduct.4.equation1.1}\nopagebreak
\begin{equation*}\begin{aligned}r_{161,N}={}&
 X_{199}\,X_{200}\\ &
+ X_{201}\,X_{202}\\ &
+ X_{203}\,X_{206}\\ &
+ X_{204}\,X_{205}
\end{aligned}\end{equation*}
\paragraph{Row 162.} \nolinkurl{rootProduct.4.equation2.0}\nopagebreak
\begin{equation*}\begin{aligned}r_{162,N}={}&
 3\,X_{117}\,X_{207}\\ &
- 2\,X_{117}\,X_{208}\\ &
- 3\,X_{117}\,X_{209}\\ &
+ 2\,X_{117}\,X_{210}\\ &
- X_{199}\,X_{203}\\ &
- 2\,X_{200}\,X_{204}
\end{aligned}\end{equation*}
\paragraph{Row 163.} \nolinkurl{rootProduct.4.equation2.1}\nopagebreak
\begin{equation*}\begin{aligned}r_{163,N}={}&
 X_{117}\,X_{207}\\ &
- 3\,X_{117}\,X_{208}\\ &
- X_{117}\,X_{209}\\ &
+ 3\,X_{117}\,X_{210}\\ &
+ X_{199}\,X_{204}\\ &
+ X_{200}\,X_{203}
\end{aligned}\end{equation*}
\paragraph{Row 164.} \nolinkurl{rootProduct.4.equation3.0}\nopagebreak
\begin{equation*}\begin{aligned}r_{164,N}={}&
 3\,X_{117}\,X_{207}\\ &
- 2\,X_{117}\,X_{208}\\ &
+ 3\,X_{117}\,X_{209}\\ &
- 2\,X_{117}\,X_{210}\\ &
- X_{201}\,X_{205}\\ &
- 2\,X_{202}\,X_{206}
\end{aligned}\end{equation*}
\paragraph{Row 165.} \nolinkurl{rootProduct.4.equation3.1}\nopagebreak
\begin{equation*}\begin{aligned}r_{165,N}={}&
 X_{117}\,X_{207}\\ &
- 3\,X_{117}\,X_{208}\\ &
+ X_{117}\,X_{209}\\ &
- 3\,X_{117}\,X_{210}\\ &
+ X_{201}\,X_{206}\\ &
+ X_{202}\,X_{205}
\end{aligned}\end{equation*}
\paragraph{Row 166.} \nolinkurl{rootProduct.4.equation4.0}\nopagebreak
\begin{equation*}\begin{aligned}r_{166,N}={}&
 44\,X_{117}\,X_{211}\\ &
- 48\,X_{117}\,X_{212}\\ &
- X_{199}^{2}\\ &
- 2\,X_{199}\,X_{201}\\ &
- 2\,X_{200}^{2}\\ &
- 4\,X_{200}\,X_{202}\\ &
- X_{201}^{2}\\ &
- 2\,X_{202}^{2}
\end{aligned}\end{equation*}
\paragraph{Row 167.} \nolinkurl{rootProduct.4.equation4.1}\nopagebreak
\begin{equation*}\begin{aligned}r_{167,N}={}&
 12\,X_{117}\,X_{211}\\ &
- 22\,X_{117}\,X_{212}\\ &
+ X_{199}\,X_{200}\\ &
+ X_{199}\,X_{202}\\ &
+ X_{200}\,X_{201}\\ &
+ X_{201}\,X_{202}
\end{aligned}\end{equation*}
\paragraph{Row 168.} \nolinkurl{rootProduct.4.equation5.0}\nopagebreak
\begin{equation*}\begin{aligned}r_{168,N}={}&
 X_{60}\\ &
- X_{211}^{2}\\ &
- 2\,X_{212}^{2}
\end{aligned}\end{equation*}
\paragraph{Row 169.} \nolinkurl{rootProduct.4.equation5.1}\nopagebreak
\begin{equation*}\begin{aligned}r_{169,N}={}&
 X_{61}\\ &
- X_{211}\,X_{212}
\end{aligned}\end{equation*}
\paragraph{Row 170.} \nolinkurl{rootProduct.4.equation6.0}\nopagebreak
\begin{equation*}\begin{aligned}r_{170,N}={}&
 1\\ &
- 11\,X_{60}\\ &
+ 24\,X_{61}\\ &
- X_{207}^{2}\\ &
- 2\,X_{208}^{2}
\end{aligned}\end{equation*}
\paragraph{Row 171.} \nolinkurl{rootProduct.4.equation6.1}\nopagebreak
\begin{equation*}\begin{aligned}r_{171,N}={}&
 3\,X_{60}\\ &
- 11\,X_{61}\\ &
- X_{207}\,X_{208}
\end{aligned}\end{equation*}
\paragraph{Row 172.} \nolinkurl{rootProduct.4.equation7.0}\nopagebreak
\begin{equation*}\begin{aligned}r_{172,N}={}&
 1\\ &
- 11\,X_{29}\\ &
+ 24\,X_{30}\\ &
- X_{209}^{2}\\ &
- 2\,X_{210}^{2}
\end{aligned}\end{equation*}
\paragraph{Row 173.} \nolinkurl{rootProduct.4.equation7.1}\nopagebreak
\begin{equation*}\begin{aligned}r_{173,N}={}&
 3\,X_{29}\\ &
- 11\,X_{30}\\ &
- X_{209}\,X_{210}
\end{aligned}\end{equation*}
\paragraph{Row 174.} \nolinkurl{alternate.direct.4.0}\nopagebreak
\begin{equation*}\begin{aligned}r_{174,N}={}&
 6\,X_{45}\,X_{209}\\ &
- 4\,X_{45}\,X_{210}\\ &
- 4\,X_{46}\,X_{209}\\ &
+ 12\,X_{46}\,X_{210}\\ &
- X_{199}\,X_{213}\\ &
- 2\,X_{200}\,X_{214}\\ &
+ X_{201}\,X_{213}\\ &
+ 2\,X_{202}\,X_{214}
\end{aligned}\end{equation*}
\paragraph{Row 175.} \nolinkurl{alternate.direct.4.1}\nopagebreak
\begin{equation*}\begin{aligned}r_{175,N}={}&
 2\,X_{45}\,X_{209}\\ &
- 6\,X_{45}\,X_{210}\\ &
- 6\,X_{46}\,X_{209}\\ &
+ 4\,X_{46}\,X_{210}\\ &
+ X_{199}\,X_{214}\\ &
+ X_{200}\,X_{213}\\ &
- X_{201}\,X_{214}\\ &
- X_{202}\,X_{213}
\end{aligned}\end{equation*}
\paragraph{Row 176.} \nolinkurl{rootProduct.4.equation9.0}\nopagebreak
\begin{equation*}\begin{aligned}r_{176,N}={}&
 X_{19}\,X_{199}\\ &
+ X_{19}\,X_{201}\\ &
+ 2\,X_{20}\,X_{200}\\ &
+ 2\,X_{20}\,X_{202}\\ &
- 6\,X_{43}\,X_{207}\\ &
+ 4\,X_{43}\,X_{208}\\ &
+ 4\,X_{44}\,X_{207}\\ &
- 12\,X_{44}\,X_{208}
\end{aligned}\end{equation*}
\paragraph{Row 177.} \nolinkurl{rootProduct.4.equation9.1}\nopagebreak
\begin{equation*}\begin{aligned}r_{177,N}={}&
 X_{19}\,X_{200}\\ &
+ X_{19}\,X_{202}\\ &
+ X_{20}\,X_{199}\\ &
+ X_{20}\,X_{201}\\ &
+ 2\,X_{43}\,X_{207}\\ &
- 6\,X_{43}\,X_{208}\\ &
- 6\,X_{44}\,X_{207}\\ &
+ 4\,X_{44}\,X_{208}
\end{aligned}\end{equation*}
\paragraph{Row 178.} \nolinkurl{rootProduct.S.retained_j.0}\nopagebreak
\begin{equation*}\begin{aligned}r_{178,N}={}&
 4\,X_{2}\\ &
+ X_{171}^{2}\\ &
+ 2\,X_{172}^{2}\\ &
- X_{173}^{2}\\ &
- 2\,X_{174}^{2}\\ &
- X_{175}^{2}\\ &
- 2\,X_{176}^{2}\\ &
+ X_{177}^{2}\\ &
+ 2\,X_{178}^{2}
\end{aligned}\end{equation*}
\paragraph{Row 179.} \nolinkurl{rootProduct.S.retained_j.1}\nopagebreak
\begin{equation*}\begin{aligned}r_{179,N}={}&
 2\,X_{3}\\ &
+ X_{171}\,X_{172}\\ &
- X_{173}\,X_{174}\\ &
- X_{175}\,X_{176}\\ &
+ X_{177}\,X_{178}
\end{aligned}\end{equation*}
\paragraph{Row 180.} \nolinkurl{scalar.S.direct_ratio.0}\nopagebreak
\begin{equation*}\begin{aligned}r_{180,N}={}&
 2\,X_{2}\,X_{107}\\ &
+ 4\,X_{3}\,X_{108}\\ &
- X_{19}\,X_{175}\\ &
+ X_{19}\,X_{177}\\ &
- 2\,X_{20}\,X_{176}\\ &
+ 2\,X_{20}\,X_{178}
\end{aligned}\end{equation*}
\paragraph{Row 181.} \nolinkurl{scalar.S.direct_ratio.1}\nopagebreak
\begin{equation*}\begin{aligned}r_{181,N}={}&
 2\,X_{2}\,X_{108}\\ &
+ 2\,X_{3}\,X_{107}\\ &
- X_{19}\,X_{176}\\ &
+ X_{19}\,X_{178}\\ &
- X_{20}\,X_{175}\\ &
+ X_{20}\,X_{177}
\end{aligned}\end{equation*}
\paragraph{Row 182.} \nolinkurl{reciprocal.aggregate.rootTrace.definition}\nopagebreak
\begin{equation*}\begin{aligned}r_{182,N}={}&
 8\,X_{4}\,X_{90}\\ &
- X_{161}^{2}\\ &
+ 2\,X_{161}\,X_{163}\\ &
- 2\,X_{162}^{2}\\ &
+ 4\,X_{162}\,X_{164}\\ &
- X_{163}^{2}\\ &
- 2\,X_{164}^{2}\\ &
- X_{175}^{2}\\ &
+ 2\,X_{175}\,X_{177}\\ &
- 2\,X_{176}^{2}\\ &
+ 4\,X_{176}\,X_{178}\\ &
- X_{177}^{2}\\ &
- 2\,X_{178}^{2}\\ &
- X_{189}^{2}\\ &
+ 2\,X_{189}\,X_{191}\\ &
- 2\,X_{190}^{2}\\ &
+ 4\,X_{190}\,X_{192}\\ &
- X_{191}^{2}\\ &
- 2\,X_{192}^{2}\\ &
- X_{203}^{2}\\ &
+ 2\,X_{203}\,X_{205}\\ &
- 2\,X_{204}^{2}\\ &
+ 4\,X_{204}\,X_{206}\\ &
- X_{205}^{2}\\ &
- 2\,X_{206}^{2}
\end{aligned}\end{equation*}
\paragraph{Row 183.} \nolinkurl{sharedT.reference.relation1.0}\nopagebreak
\begin{equation*}\begin{aligned}r_{183,N}={}&
 1752629952\\ &
+ 6299926668\,X_{17}\\ &
+ 2158699200\,X_{17}\,X_{213}\\ &
- 2185741488\,X_{17}\,X_{214}\\ &
- 8034039216\,X_{18}\\ &
- 2185741488\,X_{18}\,X_{213}\\ &
+ 4317398400\,X_{18}\,X_{214}\\ &
- 853624686\,X_{19}\\ &
+ 402282898\,X_{19}\,X_{213}\\ &
+ 632859681\,X_{20}\\ &
+ 804565796\,X_{20}\,X_{214}\\ &
+ 559299438\,X_{213}\\ &
- 615250593\,X_{214}
\end{aligned}\end{equation*}
\paragraph{Row 184.} \nolinkurl{sharedT.reference.relation1.1}\nopagebreak
\begin{equation*}\begin{aligned}r_{184,N}={}&
 760760576\\ &
+ 8034039216\,X_{17}\\ &
+ 2185741488\,X_{17}\,X_{213}\\ &
- 4317398400\,X_{17}\,X_{214}\\ &
- 12599853336\,X_{18}\\ &
- 4317398400\,X_{18}\,X_{213}\\ &
+ 4371482976\,X_{18}\,X_{214}\\ &
- 632859681\,X_{19}\\ &
- 804565796\,X_{19}\,X_{214}\\ &
+ 1707249372\,X_{20}\\ &
- 804565796\,X_{20}\,X_{213}\\ &
+ 615250593\,X_{213}\\ &
- 1118598876\,X_{214}
\end{aligned}\end{equation*}
\paragraph{Row 185.} \nolinkurl{sharedT.reference.relation2.0}\nopagebreak
\begin{equation*}\begin{aligned}r_{185,N}={}&
 1149539095796314755226869484\\ &
- 985061164914970179529857840\,X_{17}\\ &
- 1254821386791945912952706976\,X_{17}\,X_{19}\\ &
- 1153292377786262120636131146\,X_{17}\,X_{20}\\ &
- 4398006762181672234512370320\,X_{17}\,X_{21}\\ &
+ 387008901290036540531246720\,X_{17}\,X_{22}\\ &
+ 1532572190428498920661008288\,X_{17}\,X_{213}\\ &
+ 1294531802031560732990489418\,X_{17}\,X_{214}\\ &
- 1345259843509302298655015680\,X_{18}\\ &
- 1153292377786262120636131146\,X_{18}\,X_{19}\\ &
- 2509642773583891825905413952\,X_{18}\,X_{20}\\ &
+ 193504450645018270265623360\,X_{18}\,X_{21}\\ &
- 17592027048726688938049481280\,X_{18}\,X_{22}\\ &
+ 1294531802031560732990489418\,X_{18}\,X_{213}\\ &
+ 3065144380856997841322016576\,X_{18}\,X_{214}\\ &
+ 73337278472469394027868160\,X_{19}\\ &
- 686923528829960141567886528\,X_{19}\,X_{21}\\ &
+ 44414351330348966029310016\,X_{19}\,X_{22}\\ &
+ 96899299673172813900969984\,X_{20}\\ &
+ 22207175665174483014655008\,X_{20}\,X_{21}\\ &
- 2747694115319840566271546112\,X_{20}\,X_{22}\\ &
- 797514988115594446587154013\,X_{21}\\ &
- 1723051656304506319706496064\,X_{21}\,X_{66}\\ &
- 2606545027324788158562625392\,X_{22}\\ &
- 13784413250436050557651968512\,X_{22}\,X_{67}\\ &
+ 975462260105801738014285793\,X_{66}\\ &
+ 887540248202664802539362352\,X_{67}\\ &
- 588237978875945932189940160\,X_{213}\\ &
- 759606688303604527483518672\,X_{214}
\end{aligned}\end{equation*}
\paragraph{Row 186.} \nolinkurl{sharedT.reference.relation2.1}\nopagebreak
\begin{equation*}\begin{aligned}r_{186,N}={}&
 652461006517642156833128616\\ &
- 672629921754651149327507840\,X_{17}\\ &
- 576646188893131060318065573\,X_{17}\,X_{19}\\ &
- 1254821386791945912952706976\,X_{17}\,X_{20}\\ &
+ 96752225322509135132811680\,X_{17}\,X_{21}\\ &
- 8796013524363344469024740640\,X_{17}\,X_{22}\\ &
+ 647265901015780366495244709\,X_{17}\,X_{213}\\ &
+ 1532572190428498920661008288\,X_{17}\,X_{214}\\ &
- 985061164914970179529857840\,X_{18}\\ &
- 1254821386791945912952706976\,X_{18}\,X_{19}\\ &
- 1153292377786262120636131146\,X_{18}\,X_{20}\\ &
- 4398006762181672234512370320\,X_{18}\,X_{21}\\ &
+ 387008901290036540531246720\,X_{18}\,X_{22}\\ &
+ 1532572190428498920661008288\,X_{18}\,X_{213}\\ &
+ 1294531802031560732990489418\,X_{18}\,X_{214}\\ &
+ 48449649836586406950484992\,X_{19}\\ &
+ 11103587832587241507327504\,X_{19}\,X_{21}\\ &
- 1373847057659920283135773056\,X_{19}\,X_{22}\\ &
+ 73337278472469394027868160\,X_{20}\\ &
- 686923528829960141567886528\,X_{20}\,X_{21}\\ &
+ 44414351330348966029310016\,X_{20}\,X_{22}\\ &
- 651636256831197039640656348\,X_{21}\\ &
- 3446103312609012639412992128\,X_{21}\,X_{67}\\ &
- 1595029976231188893174308026\,X_{22}\\ &
- 3446103312609012639412992128\,X_{22}\,X_{66}\\ &
+ 221885062050666200634840588\,X_{66}\\ &
+ 1950924520211603476028571586\,X_{67}\\ &
- 379803344151802263741759336\,X_{213}\\ &
- 588237978875945932189940160\,X_{214}
\end{aligned}\end{equation*}
\paragraph{Row 187.} \nolinkurl{sharedT.reference.relation3.0}\nopagebreak
\begin{equation*}\begin{aligned}r_{187,N}={}&
 186960\\ &
- 667923\,X_{17}\,X_{19}\\ &
- 289924\,X_{17}\,X_{20}\\ &
+ 310419\,X_{17}\,X_{213}\\ &
- 61308\,X_{17}\,X_{214}\\ &
- 289924\,X_{18}\,X_{19}\\ &
- 1335846\,X_{18}\,X_{20}\\ &
- 61308\,X_{18}\,X_{213}\\ &
+ 620838\,X_{18}\,X_{214}\\ &
+ 1398264\,X_{66}\\ &
+ 658952\,X_{66}\,X_{213}\\ &
- 2073288\,X_{67}\\ &
+ 2635808\,X_{67}\,X_{214}\\ &
- 147928\,X_{213}\\ &
- 161376\,X_{214}
\end{aligned}\end{equation*}
\paragraph{Row 188.} \nolinkurl{sharedT.reference.relation3.1}\nopagebreak
\begin{equation*}\begin{aligned}r_{188,N}={}&
 54858\\ &
- 144962\,X_{17}\,X_{19}\\ &
- 667923\,X_{17}\,X_{20}\\ &
- 30654\,X_{17}\,X_{213}\\ &
+ 310419\,X_{17}\,X_{214}\\ &
- 667923\,X_{18}\,X_{19}\\ &
- 289924\,X_{18}\,X_{20}\\ &
+ 310419\,X_{18}\,X_{213}\\ &
- 61308\,X_{18}\,X_{214}\\ &
- 518322\,X_{66}\\ &
+ 658952\,X_{66}\,X_{214}\\ &
+ 2796528\,X_{67}\\ &
+ 1317904\,X_{67}\,X_{213}\\ &
- 80688\,X_{213}\\ &
- 147928\,X_{214}
\end{aligned}\end{equation*}
The rows are independently reconstructible from \nolinkurl{research/quartic214/construction/build.py}.  The separate Laurent implementation \nolinkurl{construction/check_exact.py} compares all twenty changed factor components, six inverse-weight relations and $162$ untouched rows, and checks degree, content, coordinate occurrence, coefficient profiles and the complete quartic table.  These are finite algebraic conclusions.  They do not certify the retained arithmetic theorems.
\clearpage
% Source: research/quartic214-manuscript-review/FOUNDATION_TRANSITION.md
\section{Current Version 16 foundation and the 214 consumers}\label{q214:foundation-transition}
\noindent\textit{Source: }\nolinkurl{research/quartic214-manuscript-review/FOUNDATION_TRANSITION.md}.\par

\begin{quote}\footnotesize Author: model task \nolinkurl{/root/q214_tex_coverage}, 12 November 2026. This is a new source-to-consumer comparison for independent review. It supplies no independent acceptance. The complete current foundation is \nolinkurl{manuscripts/quartic_v16.tex}, SHA-256 \nolinkurl{aeb2b6a42a03943d7f8c46682389c8da693d2903ba92ccc164bbbd236451de74}. The earlier \(V_{15}\) source, SHA-256 \nolinkurl{44ba3a8ec8521c29b81b26cfb9981cd6d5f3c76d54e7782438614842edb52c74}, remains immutable historical evidence.\end{quote}

\subsection{Current source rule}

The operative foundation in this edition is the complete Version 16 text reproduced in Part II. The \(V_{15}\) byte-identity statements in uniform-existence \(P_{02}\) and source-carry \(C_{01}\) describe the earlier frozen packet. They are superseded as statements of the \textbf{current receiving source}. Their hashes, earlier arguments and review limits remain historical evidence; no old approval is transferred automatically. The current source supplies the retained rank, restricted-converse, Kato, GI, Gaussian-prime, scalar and pole interfaces, with the three changes compared below. The 214 polynomial, its raw224 branch, the new local factors, the uniform ineffective threshold and the product averaging route are unchanged by this substitution.

The comparison used the actual complete normalized V15-to-current-V16 difference, the actual Fable audit passages concerning horizontal normalization and the four follow-ups, the current \(V_{16}\) passages themselves, and the actual 214 consumers. It follows both the editorial amendment and the later Tao--Ziegler precision amendment: the intermediate \nolinkurl{2f729...} and \nolinkurl{a9cb...} \(V_{16}\) hashes are not the final receiving source. Publication and pre-existing source-review scopes remain those recorded with the current foundation. This author does not claim a new reading of every external publication or a recursive reproof of the entire foundation.

\subsection{Exact horizontal identities before the GI finite cover}

The source change is \(V_{16}\) lines26819--26929, in ``Removing rational and small coordinates without losing unit norm''. The available regularized representation previously identified horizontal Taylor coefficients only modulo the horizontal lattice of \(\Gamma \). The later finite cover uses a smaller lattice Gamma-prime. Equality modulo \(\Gamma \) alone does not imply the required equality modulo Gamma-prime. The Fable audit gives an actual Heisenberg counterexample to that implication and identifies the required normalization; this is a substantive connecting-argument correction, even though the conclusion of the GI supplier remains the same.

\(V_{16}\) first fixes the actual internal real lifts: a small lift for each small coordinate, a rational lift for each rational coordinate, and \nolinkurl{{alpha h} beta} for each linear petal. For each horizontal weight \(i\), the discrepancy is in the image of \(\Gamma \) intersect \(G_{i,1}\), by the definition of that horizontal lattice. Choosing a preimage \(\gamma _{h,i}\), the orbit is right-multiplied by the ordered product of \(\gamma _{h,i}^{-\text{binom}(n,i)}\). This product is Gamma-valued, so the orbit modulo \(\Gamma \) is unchanged. The manuscript's Taylor criterion B.9 and product rule B.4 supply polynomiality, and the horizontal Taylor homomorphism subtracts exactly the discrepancy. Thus the identities hold as vector equalities for these fixed lifts \textbf{before} any coordinate removal.

Small and rational removal then subtract the actual fixed lifted terms. The new unit tests, lower-model circles, finite residue interpolation and fixed zero padding remain the constructions written in the source; the normalization changes none of their observed nilcharacters. When lattice generators are sent to their M-th powers, divide those same real lifts by \(M\). The identity \(\sum _{a} (x_{a}/M) M \phi (e_{a}) = \sum _{a} x_{a} \phi (e_{a})\) is an equality of vectors, hence also an equality modulo the smaller horizontal lattice. This checks the direction and scalar of the comparison map, rather than inferring it from equality modulo the larger lattice. Linear petals use \nolinkurl{{alpha h}(beta/M)}; core and independent coordinates use the indicated fixed/internal roots. Every horizontal weight, remaining letter and member of the retained dense internal family is included. A common standard denominator is retained, and an integer relation after division would multiply back to the original one.

The receiving 214 source-carry \(C_{08}\) requires the actual internal GI conclusion: fixed degree, a fixed finite family and uniform complexity/Lipschitz bounds, global unit-vector tests and exact full-limit vertical laws. The correction preserves the represented sequence and supplies the previously missing finite-cover identity in that construction. \(C_{09}\) then uses the same GI(k) with \(r=\operatorname{max}(s,4)\) and \(k=r-1\), retaining the normalized interval loss \(\alpha  \delta /(4 C_{\text{star}})\). Neither the finite family nor these thresholds depend on interval length or orbit parameters. \(C_{10}\) and uniform-existence \(P_{16}\) consume the resulting Gaussian theorem only after their fixed arithmetic data have been chosen. No stronger GI conclusion or fresh uniformity in the 214 input is introduced.

\subsection{Tao--Ziegler at the actual fixed modulus}

The source application is \(V_{16}\) lines19989--20169, especially the paragraph beginning ``We now adapt''. Its actual modulus is \(W=D \prod _{p\le v,p \text{ not }  \text{ dividing }  D} p\). After \(v\) exceeds the fixed prime divisors of \(D\) this is \(E \prod _{p\le v} p\), with \(E=D/\operatorname{rad}(D)\) fixed. The published modulus-dependent statements use a squarefree primorial. Thus Theorem11, Corollary3, Theorem13 and the polynomial-forms Theorem15 are used here through an \textbf{adaptation of their proofs}, not a literal substitution into their published hypotheses. The bounded-function Theorem9 has no such modulus parameter. This distinction supersedes any reading of the old source carry as a word-for-word application at the changed modulus.

For a unit residue \(b\), every divisibility condition at a prime dividing \(W\) is false, including at its extra prime powers. The normalization \(\phi (W)/W\) depends only on prime support and equals the primorial normalization once \(v\) contains the support of \(D\). At primes outside that support \(W\) is a unit. The divisor-sum factors there need not have numerically identical values for different \(W\) or residues: the operative assertion is that they satisfy the \textbf{same uniform estimates}. The fixed factor \(E\) is absorbed by the fixed polynomial coefficient bounds. The current source retains the resulting \nolinkurl{1+o(1)} polynomial-forms estimate, the bounded-model constant and the ordered choice of local box order, coefficient exponent, \(B_{0}\), global order and \(\rho \). It then fixes \(v\) before sending the analytic size to infinity and chooses \(v\) slowly by diagonalization. These are the estimates and quantifiers consumed downstream, not an identity of individual numerical local factors.

The identified earlier derivation in \nolinkurl{research/version16/review/current/REVIEW.md}, Section `Fixed-extra-factor Tao--Ziegler adaptation', was read as an argument, not used merely for its decision. Its Mobius divisor expansion restricts the divisors and their least common multiples to be squarefree. At primes dividing \(W\) the unit-residue condition removes the events independently of the exponent of \(E\). At other primes, invertibility preserves the available degree/root-count bounds; nonconstant polynomial differences persist and exceptional primes divide nonzero coefficient combinations controlled by coefficient height. Fixed \(E\) changes that bound by a fixed constant. The CRT/inradius error is governed by those squarefree divisor moduli. The same finitely many moment patterns feed the functional pseudorandomness, dense-model and generalized von Neumann steps; the separate independent-copy, cyclic-model and interval-boundary hypotheses remain required. This is reuse of an identified prior derivation at its exact fixed-data scope, followed by the new214 comparison below. This author does not claim to have repeated that reviewer's primary-source readings. Its precision addendum expressly rejects a numerical identity of local factors and preserves these estimates.

For 214, uniform-existence \(P_{15}\) first clears the two rational scales and merges equal contact coefficients across detectors while preserving every incidence. Within either detector the ten coefficients are distinct. \(P_{16}\) applies the fixed-data prime theorem to that single fixed integer list and to fixed avoidance, progression, slope and prime lower bounds. The field is the fixed ray class field of \(Q(i)\); no growing field or unrestricted Galois-prime theorem enters. The inequalities \(0\le \text{Lambda}'_{j}\le \text{Lambda}'\) permit the ordinary prime majorant. The resulting W-trick polynomials have the source's degree and coefficient bounds with all fixed data absorbed before the analytic limit. Repeated incidences use the same selected prime and CRT condition, rather than duplicating equal forms and violating the distinct-coefficient application.

The arithmetic data may depend on the particular (N,x), as P15--P16 explicitly allow. They are fixed before the auxiliary analytic parameter \(X_{\text{sel}}\) grows. No estimate uniform in (N,x), scale, cleared denominator or progression modulus is needed for completeness, and none becomes a lower bound on the family parameter \(N\). The separate obstruction-tail and slow-cutoff argument retained in \(C_{10}\) is unchanged. Thus the supplied fixed-data theorem has precisely the scope consumed at \(P_{16}\). This author comparison does not replace the internal majorant proof by a claim that a changed-modulus theorem applies literally.

The Fable follow-up did not reread every proof in that Tao--Ziegler chain and records that limitation. It also sketches an alternative progression-dilation route but expressly does not write it out completely. This edition \textbf{does not adopt that outline}. It retains the actual \(V_{16}\) proof adaptation and its source-review obligations. An auditor's unexamined portion is a limitation of that audit, not by itself a new mathematical objection; it neither supplies nor refutes a missing derivation.

\subsection{Height citations and the same-field product route}

\(V_{16}\) lines30658--30683 now cite Gaudron--Remond, \emph{Theoreme des periodes et degres minimaux d'isogenies}, Commentarii Mathematici Helvetici89 (2014),343--403, Section2.3 and Corollary8.4, DOI10.4171/CMH/322. The earlier Duke paper remains a different source for its own Theorem1.4/Section9 uses. The CMH source supplies the original-normalization lower bound and the one-sided half-log-degree height change at the corrected surface consumer.

The actual 214 primary-source-bridges already identifies that CMH article and its normalization conversion. The operative height-averaging proof applies the published GR2023 product construction over the actual L-prime, chooses \(r\) separate simple factors in the actual \(A_{a}-\text{isogeny}\) class, and uses a dimension-two isogeny from the selected \(C_{i}\) to \(A_{a}\) over that same field. All \(r\) occurrences, negative individual heights and a zero-dimensional complement remain covered. Neither this citation correction nor the \(V_{16}\) single-surface lemma substitutes an evaluation-splitting field for L-prime, replaces \(r\) by a set of distinct isogeny classes, or removes the published product supplier's indirect Remond dependencies. The exact received inequality is still H3.1, consumed with H2's multiplicity saving by \(H_{4}\). The historical quotient \(H_{3}\) remains superseded in the 214 application.

\subsection{Related uses and remaining acceptance}

The complete normalized source difference contains the GI insertion, fixed-lift division, Tao--Ziegler adaptation/precision wording, the two CMH citation changes, the added bibliography entry, publication metadata and the restored AI usage attribution. The other mathematical source text is unchanged after LF comparison. The bounded related-use comparison covers source-carry C01/C08--C10/C12--C13, uniform-existence P02/P08--P10/P15--P20, the current product averaging and primary-source height bridges, and the explicit scope notices before the historical packet. Old \(V_{15}\) line numbers in the retained packet are provenance locators; the source map and this transition identify the current receiving passages. The historical source texts and fourteen original public contracts are not rewritten.

The Fable Bogolyubov-numbering concern is expressly withdrawn in its follow-up after comparison with the published numbering; no manuscript repair is imported for that item. Its audit's other unexamined portions remain review-scope limitations, without inventing new objections or treating them as verification. The original audit and its \(V_{16}\) findings remain preserved in their versioned locations.

This transition is a changed foundation application. A fresh non-author must assess its exact source passages, the corrected connections, the actual 214 consumers and the assembled TeX/PDF. Old component or \(V_{16}\) acceptance is evidence at its own bound scope, not automatic acceptance of this new receiving application. Until that decision and the matching release binding exist, assembled-proof and record-release acceptance remain pending.
\clearpage
% Source: research/quartic214/PROOF.md
\section{Assembly of the selected 214 construction}\label{q214:assembly}
\noindent\textit{Source: }\nolinkurl{research/quartic214/PROOF.md}.\par

This is an author assembly for fresh public review. The construction and all source-to-consumer transitions below are proposed together. Historical source acceptance is separate from acceptance of this public adaptation. \href{https://github.com/elliotglazer/diophantine-z-q/blob/110cf72565be326fd567cd1a64a3b2e00213dcd1/research/quartic214/SOURCE_PATH_MAP.json}{\nolinkurl{SOURCE_PATH_MAP.json}} records exact original source hashes, selected ranges, adaptations, and reused public paths. It deliberately contains no generated reviewer approval.

\subsection{Exact family and domains}

Put \(A_{N}=149^{2N+1}\). Starting with the public \href{https://github.com/elliotglazer/diophantine-z-q/blob/110cf72565be326fd567cd1a64a3b2e00213dcd1/research/version15/quartics/quartic-222/provenance/shared-T-redundancy224-006/candidate224-ir.json}{raw224 predecessor}, \href{https://github.com/elliotglazer/diophantine-z-q/blob/110cf72565be326fd567cd1a64a3b2e00213dcd1/research/quartic214/construction/build.py}{the builder} removes ten fourth-power cache coordinates and their ten defining rows, reuses six existing power slots for inverse weights, and changes twenty factor component rows. The result has 214 coordinates and 188 primitive integral quadratic rows \(r_{i}\). Define \(P_{N}=\sum _{i=1}^{188}\) \(r_{i}^{2}\) after substituting the fixed coefficient \(A_{N}\). Over \(Q\), \(P_{N}=0\) is equivalent to every \(r_{i}=0\).

The exact \href{https://github.com/elliotglazer/diophantine-z-q/blob/110cf72565be326fd567cd1a64a3b2e00213dcd1/research/quartic214/construction/check_exact.py}{Laurent proof checker} reconstructs the deleted coordinates, checks the six inverse-power definitions, and compares all twenty factor components to the changed physical factor system up to a nonzero rational normalization and \(q^{4}\). It compares all 162 untouched rows. The factors have exponents (5,2,4), with exponent 3 in the special input term. They are \textbf{new arithmetic factors}; an old factor-exponent conclusion is not reused to prove their soundness. The equations \(q\) \(s_{1}=1\), \(s_{2}=s_{1}^{2}\), \(s_{4}=s_{2}^{2}\) require \(q\ne 0\). Completeness constructs witnesses in that domain. No bijection with the discarded \(q=0\) witness stratum is claimed.

The checker verifies all 214 variables occur, degree 4, content 1, and the full quartic table. Its coefficient proof identifies the unique \(A_{N}^{2}+1\) profile and bounds every other constant or linear profile strictly below it for \(A_{N}\ge 149^{107}\). Dividing those latter bounds by \(A_{N}^{2}\) shows they remain strict as \(N\) increases. This proves the coefficient-height statement for \(N\ge 53\) and makes no arithmetic statement about that value of \(N\).

\subsection{Source objects and changed applications}

The field is \(K=Q(\sqrt{2})\), the elliptic curve is \(E:y^{2}=x^{3}-(3-\sqrt{2})^{2}x\), and its generator is \(P=2(1+\sqrt{2},2\sqrt{2})\). The four nonzero labels are (n,a,t,b), \(c=|a|\), \(B=\operatorname{max}(|n|,|a|,|t|,|b|)\), and \(\beta =51x^{2}+2504\). The physical detector arguments are \(\lambda \) and \(T\lambda \), where \(T\) is the positive quarter-trace of the sum of the four abscissae. The ten contacts are \(\Theta \cup (2-\Theta )\), \(\Theta =\{0,-1/2,-1,1/2399,-32/407\}\). \(\Sigma \) consists of \(\text{primes}\le 76176\) and 2029147. These exact objects, rather than another contact center or scalar convention, are consumed throughout.

\begin{longtable}{@{}p{0.23\linewidth}p{0.73\linewidth}@{}}
\toprule
Application & Available result, constructed input, exact bridge and consumer\\
\midrule
\endhead
Physical recovery and reciprocal clearing & \href{https://github.com/elliotglazer/diophantine-z-q/blob/110cf72565be326fd567cd1a64a3b2e00213dcd1/research/quartic214/proofs/existence.md}{\nolinkurl{existence.md}}, ``Algebraic transport'', reconstructs the actual raw224 reference, root and scale objects using the predecessor's inverse maps. Forced-149 valuations recover \(w\) and the shared alternate ratio before division. Zero components and zero test numerators are retained. The full row identities are checked by the independent Laurent implementation.\\
Changed local extraction & \href{https://github.com/elliotglazer/diophantine-z-q/blob/110cf72565be326fd567cd1a64a3b2e00213dcd1/research/quartic214/proofs/uniform-existence.md}{\nolinkurl{uniform-existence.md}} P04--P07 applies the valuation-residue separation to the actual five changed factors at every consumed split or inert place, with arbitrary finite test valuations and zero tests. The trace/pole input is reconstructed literally. Positive reciprocal valuations precede logarithmic integrality. The outputs are the same integer residuals \(M_{j}^{3}\mid \Delta \) and \(M_{j}\mid J\), separately for each scale, with no extra scalar factor.\\
Full-family geometric support & \href{https://github.com/elliotglazer/diophantine-z-q/blob/110cf72565be326fd567cd1a64a3b2e00213dcd1/research/quartic214/proofs/geometry-carry.md}{\nolinkurl{geometry-carry.md}} G1--G4 and \href{https://github.com/elliotglazer/diophantine-z-q/blob/110cf72565be326fd567cd1a64a3b2e00213dcd1/research/quartic214/proofs/provider-scope-carry.md}{\nolinkurl{provider-scope-carry.md}} P1--P4 identify the actual public \(D_{210}\) cover, its field and Hom objects, original metric and CM branch. Its \(S_{10}\) bridge covers all primitive rational pairs, including the infinity contact, deep even contacts and \(M=1\). The ten determinant values only prove the finite support calculation; the public geometric proof supplies the actual cover and tame charts.\\
Original modular realization and conductor & \href{https://github.com/elliotglazer/diophantine-z-q/blob/110cf72565be326fd567cd1a64a3b2e00213dcd1/research/quartic214/proofs/modularity.md}{\nolinkurl{modularity.md}} M1--M6 constructs the controlled cochain and the actual original representations, checks the published modularity hypotheses at 5, absorbs the finite twist into the final primitive weight-two form, compares the same form at 3, and covers every residue prime by the pair \{5,3\}. \(M_{5}\) preserves inertia and monodromy in the level comparison and gives \(N_{\text{level}}\le C_{N}\) \(R^{2}\), \(\operatorname{rad}(N_{\text{level}})\mid R\). Its field \(L'\) is distinct from coefficient and evaluation-splitting fields.\\
Modular-Jacobian height & \href{https://github.com/elliotglazer/diophantine-z-q/blob/110cf72565be326fd567cd1a64a3b2e00213dcd1/research/quartic214/proofs/height-evaluation.md}{\nolinkurl{height-evaluation.md}}, \href{https://github.com/elliotglazer/diophantine-z-q/blob/110cf72565be326fd567cd1a64a3b2e00213dcd1/research/quartic214/proofs/height-different.md}{\nolinkurl{height-different.md}}, and \href{https://github.com/elliotglazer/diophantine-z-q/blob/110cf72565be326fd567cd1a64a3b2e00213dcd1/research/quartic214/proofs/height-metric.md}{\nolinkurl{height-metric.md}} prove the integral trace-dual, complete discriminant and metric links. \href{https://github.com/elliotglazer/diophantine-z-q/blob/110cf72565be326fd567cd1a64a3b2e00213dcd1/research/quartic214/proofs/height-assembly.md}{\nolinkurl{height-assembly.md}} A--D joins the same curve \(X_{1}(N_{\text{level}})/\operatorname{ker}\) \(\epsilon \), order, lattice, trace, metric and all fibre points to H0.3. The auxiliary evaluation field cancels in normalized height; it does not replace \(L'\).\\
Multiplicity, averaging and \(H_{4}\) & \href{https://github.com/elliotglazer/diophantine-z-q/blob/110cf72565be326fd567cd1a64a3b2e00213dcd1/research/quartic214/proofs/height-supplier.md}{\nolinkurl{height-supplier.md}} H1--H2 constructs the full coefficient-conjugate multiplicity over \(L'\). \href{https://github.com/elliotglazer/diophantine-z-q/blob/110cf72565be326fd567cd1a64a3b2e00213dcd1/research/quartic214/proofs/primary-source-bridges.md}{\nolinkurl{primary-source-bridges.md}} pins the Faltings, Poincaré, Fricke and lower-height applications. \textbf{The current \(H_{3}\) application is \href{https://github.com/elliotglazer/diophantine-z-q/blob/110cf72565be326fd567cd1a64a3b2e00213dcd1/research/quartic214/proofs/height-averaging.md}{\nolinkurl{height-averaging.md}}:} use the published product construction over \(L'\), select \(r\) separate simple factors in the actual \(A_{a}-\text{isogeny}\) class, and apply a dimension-two isogeny \(C_{i}\to A_{a}\). It supplies exactly H3.1 for \(H_{4}\). The historical quotient route retained in \nolinkurl{height-supplier.md} \(H_{3}\) is superseded in this application. \(H_{4}\) proves \(h(u)<C_{2}\) \(M_{S_{10}}(u)^{2103/1024}\) for every allowed rational \(u\), including CM and radical-one cases; its isolated rational certificate changes no row.\\
Source theorem carry & \href{https://github.com/elliotglazer/diophantine-z-q/blob/110cf72565be326fd567cd1a64a3b2e00213dcd1/research/quartic214/proofs/source-carry.md}{\nolinkurl{source-carry.md}} C01--C13 compares the exact public manuscript's rank, RC, Kato full-index/limiting-class, GI, Gaussian prime, scalar and pole statements with their actual 214 consumers. It does not import the final paired-root/trace-quotient reductions into this branch. Public review evidence is historical scope evidence, not a replacement for these statements and maps.\\
\bottomrule
\end{longtable}

The named source details are openly readable in the linked proofs and existing public manuscript. Selected older private proof fragments are in \href{https://github.com/elliotglazer/diophantine-z-q/blob/110cf72565be326fd567cd1a64a3b2e00213dcd1/research/quartic214/sources/}{sources/}. Published works are identified by edition, URL, hash and theorem/page in \href{https://github.com/elliotglazer/diophantine-z-q/blob/110cf72565be326fd567cd1a64a3b2e00213dcd1/research/quartic214/PRIMARY_SOURCES.json}{\nolinkurl{PRIMARY_SOURCES.json}}, rather than republished in full. The source metadata describes previous actual inspections; the present porter does not claim to have repeated all those primary-source readings.

\subsection{Soundness with one fixed threshold}

For every rational zero of the rows, the local argument supplies \(M_{j}^{3}\mid \Delta \) and \(M_{j}\mid J\). The two transformed source arguments \(2+1/t_{j}\) and \(2+1/(2-t_{j})\) cover exactly the excluded-contact domain, with inverse row norms 3 and 7. Discarding the fixed \(\Sigma \setminus S_{10}\) primes costs one fixed factor, since radicals are squarefree. The two blocks are disjoint outside \(\Sigma \). Thus \href{https://github.com/elliotglazer/diophantine-z-q/blob/110cf72565be326fd567cd1a64a3b2e00213dcd1/research/quartic214/proofs/uniform-existence.md}{\nolinkurl{uniform-existence.md}} P08--P10 gives \(h(t_{j})\le H_{2}\) \(M_{j}^{7/6}\), and the actual aggregate \textbf{sum} theorem gives

\begin{gather*}
B^{2} \le  C |J_{0}|^{7/(6k)} + M_{2}
\end{gather*}

whenever \(M_{j}^{k}\) divides the same nonzero integer \(J_{0}\) for each \(j\). No product divisibility between the scales is required. The graph-fibre argument checks the actual sum, all nonzero labels, and the published theorem's embedding condition; \(M_{2}\) is fixed and ineffective.

For \(\Delta \ne 0\), \(|\Delta |\le 14B^{4}\) and \(k=3\) bound \(B\) uniformly because \(14/9<2\). The unchanged scalar-S/four-contact preparation gives \(c\ge 149^{{L}_{N}}\), \(L_{N}=\lfloor 9(2N-1)/4\rfloor +1\), so all sufficiently large \(N\) exclude \(\Delta \ne 0\). The elementary Lucas lower bound and \(h(\beta )\le 2^{330770}\) \(149^{-L_{N}}c(1+\operatorname{log}(c)/3)\) then exclude nonzero reduced numerator \(J\) of \(n-\beta \). These are uniform bounds in \(x\) and all solutions, as made explicit in P11--P12. Hence \(\beta \) is integral, and squarefreeness of 51 forces \(x\) integral.

\subsection{Completeness and quantifiers}

\href{https://github.com/elliotglazer/diophantine-z-q/blob/110cf72565be326fd567cd1a64a3b2e00213dcd1/research/quartic214/proofs/uniform-existence.md}{\nolinkurl{uniform-existence.md}} P13--P20 fixes a single coefficient threshold \(N_{C}\) and constructs witnesses for every integer input and every \(N\ge N_{C}\). Scalar preparation precedes the nonzero Lucas labels and positive trace. The rational scale coefficients are cleared and merged with their incidence sets retained before the fixed-data prime theorem is applied. The CRT reference choice is compatible with all incidences. At selected primes the \textbf{changed factors} have the necessary positive depths; at opposite places their trace is a unit; at prepared finite primes and real embeddings their actual pole inputs satisfy the exact zero-preparation hypotheses. The RC theorem is applied to those fresh inputs' actual non-CM curves, with full Selmer corank before rank/finite-Sha consequences. Independent pole witnesses then supply the two detectors and all inverse weights.

The labels, fixed-data prime sets, moduli, slope and final witnesses may depend on (N,x). Their existence imposes no further lower bound on \(N\). Reconstructing the removed cache values and applying the exact row identities gives a zero of \(P_{N}\). With the fixed soundness threshold \(N_{S}\), take \(N_{0}=\operatorname{max}(N_{S},N_{C})\). This proves the asserted \(\exists N_{0}\forall N\ge N_{0}\forall x\) form relative to the exact retained source theorems and the applications proposed here. \(N_{0}\) remains ineffective.

\subsection{Review boundary}

The \href{https://github.com/elliotglazer/diophantine-z-q/blob/110cf72565be326fd567cd1a64a3b2e00213dcd1/research/quartic214/contracts/INDEX.json}{public application contracts} retain the original authors' nine obligations with new whole-file and verified public passage bindings. Original private contract hashes are provenance; old private line pins and approvals do not become current public pins or approvals. Required internal sources resolve to public files, and external primary sources retain edition, URL and hash metadata.

All current applications above require review of this exact public serialization and its dependencies. Prior scoped acceptance does not independently accept this assembly, establish a numerical \(N_{0}\), or certify a full example witness tuple. The isolated numeric and algebra checks establish only their declared finite conclusions. The public Astra/controller acceptance requirement remains in force.
\clearpage
% Source: research/quartic214/proofs/existence.md
\section{Reciprocal clearing on the 224-coordinate predecessor: 214 and 218}\label{q214:existence}
\noindent\textit{Source: }\nolinkurl{research/quartic214/proofs/existence.md}.\par

\begin{quote}\footnotesize Author: /root, integrating the attributed existence222 contribution. This is a new application on the literal public raw224 rows, not adoption of an earlier verdict. All source files and the author construction are bound by the author manifest for non-author review. This section is conditional on the retained elliptic, scalar, prime-selection, pole and aggregate-trace inputs of the frozen Version15 source. Its defining threshold remains ineffective. The 214 conclusion additionally uses the newly proved quadratic supplier in the sibling height packet; that theorem's inherited assumptions remain explicit.\end{quote}

\subsection{Statement and exact objects}

Let \(A_{N}=149^{2N+1}\), \(N\ge 1\). The portable builder starts with \nolinkurl{research/version15/quartics/quartic-222/provenance/shared-T-redundancy224-006/candidate224-ir.json}, the public shared-T-redundancy224-006 stage (224 variables,198 primitive quadratic rows), before the paired-root and trace-quotient steps. It produces a 214-variable family with188 rows and a 218-variable control with192 rows. Every polynomial is the sum of the squares of its displayed primitive integral rows. The complete tables and symbolic row identities are executable evidence, not a substitute for the arithmetic argument below.

For each family there is a fixed \(N_{0}\) such that all \(N\ge N_{0}\) define \(Z\) by rational projection of input coordinate1. This is the conditional conclusion established here. Its proof does not give an effective \(N_{0}\). The coefficient height \(\text{equals}149^{4N+2}+1\) for \(N\ge 53;53\) is solely a coefficient threshold. The parameter \(A_{N}\) is never a quantified variable.

Use exactly \(K=Q(\sqrt{2})\), \(d=3-\sqrt{2}\), \(D=d^{2}\), \(E:y^{2}=x^{3}-Dx\) and the generator \(P\) in the frozen source. The raw224 component bases, including the nonstandard \(W_{0}\) basis, are retained by literal rational component rows. The two actual weights are raw224 coordinates87 and102. The physical arguments are \(\lambda \) and \(T*\lambda \), with \(T=(1/4)\operatorname{Tr}(\sum _{i=1}^{4} x(Q_{i}))>0\). Reciprocal-scale and scalar-R rescalings are interpreted using their proved inverse maps, not by identifying a rescaled coordinate with an unscaled field value.

The exact ten-contact set is \(\Theta \) union(2-Theta), \(\Theta =\{0,-1/2,-1,1/2399,-32/407\}\); \(\Sigma \) consists of \(\text{primes}\le 76176\) and2029147. Write \(F_{2}\) for its monic polynomial, \(\delta =t-1\), \(R=F_{2}(t)/\delta ^{10}\) and \(q=\delta ^{10}/F_{2}(t)-1\) on \(\delta \ne 0\). The literal contact equation excludes every contact since \(F_{2}(1)\ne 0\). At \(\delta =0\) it forces \(q=-1\); the comparison radical is1 and no division by \(\delta \) is made. New inverse equations enforce \(q\ne 0\). They are allowed to remove some irrelevant old witnesses; no all-witness bijection with unrestricted raw224 is claimed.

\subsection{Algebraic transport, domains and source-interface recovery}

The removed coordinates are raw224 IDs23,24,29,30,33,34,37,38,133,134: the two rational components of each of five V-side fourth powers. Their defining equations solve them as explicit quadratic polynomials in retained reciprocal-abscissa-square coordinates. All other occurrences are in the twenty factor component rows; the symbolic checker verifies this occurrence statement and retains their actual component conventions.

In each detector reuse three power slots as s1,s2,s4 and impose \(q*s_{1}=1\), \(s_{2}=s_{1}^{2}\), \(s_{4}=s_{2}^{2}\). The214 factors use exponents(5,2,4), with special exponent3 below. Multiplying their defining equations by \(q^{-4}\) makes the old highest powers quadratic expressions in already stored squares. The218 control retains exponents(7,2,4), special exponent1; it also stores \(q_{3}\) and \(s_{3}\), with \(s_{3}=s_{1}*s_{2}\) and \(q_{3}*s_{3}=1\). These equations imply \(q_{3}=q^{3}\) and \(s_{3}=q^{-3}\) without further guards.

The symbolic Laurent-polynomial checker independently substitutes every eliminated coordinate and every weight, then compares all twenty factor component identities with the appropriate positive or negative constant row multipliers and \(q^{4}\). It compares all162 untouched rows exactly up to their primitive sign convention and verifies all power relations. This is an identity of formal polynomial expressions, not agreement at sample points. For214 the target factor system has the changed exponents; it is not the old order-four system. For218 it is literally the old system on \(q_{0}*q_{1}\ne 0\).

Every new rational solution therefore reconstructs the physical scalar/reference/root system by the raw224 inverse maps. The reconstructed deleted fourth-power outputs satisfy their old definitions identically, including \(W_{0}\); \(Q_{5}\) had already been removed in raw224 and has no hidden \(W_{0}\) consumer. Source locations: the public predecessor proof listed in \nolinkurl{SOURCE_PATH_MAP.json,} Version15 \nolinkurl{v222:reference} and its displayed polynomial F/E identities; \nolinkurl{v222:ratio}; \nolinkurl{v15:reductions}. The reference recovery uses Q0,Q1,Q3 to obtain the \(H\) curve and \(f\ne 0\), before interpreting \(w\). The unchanged forced149 scalar argument then gives \(v_{149}(f)=m>0\), \(v_{149}(S)=-m/2\), \(v_{149}(g)=m/2\). The unique lowest \(S\) term in the coefficient of \(w\) in \(Q_{2}\) gives \(w=g^{4}/f^{2}\). The former \(W_{0}\) is simply reconstructed as \(w^{2}\). No factor exponent is used in this chain.

The shared alternate ratio is recovered from \(D_{r}*T=N_{r}\) with both numerator and denominator having valuation -m/2 at149. Thus \(T\) is nonzero, all four alternate ratios regain their original meanings, and the \(k=0\) and zero-numerator branches are handled by the retained inverse maps. No division by a rational component or test numerator is added. The scalar four-contact normalizer and parity-scalar rows are untouched; the forced149 bound, real and finite scalar preparation, and nonzero integer root labels remain the same. These uses retain their earlier proofs at the exact source locations, rather than being inferred from finite row checks.

\subsection{New local factors and soundness for214}

At a consumed \(p\) outside \(\Sigma \) the fixed pairwise contact differences (equivalently the primitive-form determinants) and \(\delta \) are units; one contact factor has positive valuation. Thus \(\nu =v_{p}(R)>0\) is odd and \(v(q)=-\nu \). This statement concerns the physical argument \(t\) and does not assume its scale is a p-unit. At every place of \(K\) above \(p\), all reciprocal abscissa values have even finite valuation by the curve equation. The recovered H-P square has this same meaning; before reference recovery its parity is also supplied by \(f(H)*F_{0}=g^{2}\). Zeros have valuation+infinity.

Write \(D_{H}=f(H)\), \(A_{i}=f(U_{i})\), \(B_{i}=\sigma \) \(f(V_{i})\), \(C_{i}=L_{i}-L_{i}'\), \(I=L_{1}-\beta \) and \(G=F(L_{2},L_{3},L_{4})\), with the exact Lucas polynomial in the source. The five factors are

\begin{gather*}
T_{1}=1+q^{5} C_{2}^{2}+q^{2} D_{H}^{2}+q^{4} B_{1}^{4},\\
T_{2}=1+q^{5} G^{2}+q^{2} A_{3}^{2}+q^{4} B_{0}^{4},\\
T_{3}=1+q^{5} I^{4}+q^{3} A_{2}^{2}+q^{2} A_{4}^{2}+q^{4} B_{3}^{4},\\
T_{4}=1+q^{5} C_{3}^{2}+q^{2} C_{1}^{4}+q^{4} B_{4}^{4},\\
T_{5}=1+q^{5} C_{4}^{2}+q^{2} A_{1}^{2}+q^{4} B_{2}^{4}
\end{gather*}

Except in \(T_{3}\) the three nonconstant valuations lie respectively in odd integers,2+4Z,4+8Z. In \(T_{3}\) their residues are -nu,nu,2 modulo4 and4 modulo8. These classes are pairwise distinct and exclude0. Omitting zero-test terms preserves that property. Therefore every factor has a unique minimum, is nonzero, and has \(\text{valuation}\le 0\); at valuation0 its residue is1. This holds at both split places and at the unique inert place, for arbitrary negative test valuations as well as positive ones.

Let \(W=\prod \) \(T_{i}\) and \(\ell =\operatorname{max}_{v|p}v(W)\le 0\). The unchanged trace rows force \(w_{0}\ne 0\) and \(\tau =1/\operatorname{Tr}(W^{-1})\ne 0:\) \(w_{0}=0\) would force \(w_{1}=0\), contrary to \(W\ne 0\). The literal pole row gives \(b=\gamma ^{2}R/\operatorname{Tr}(W^{-1})^{2}\), \(\gamma =408/409\) a unit at \(p\). The trace inequality implies \(v_{p}(b)\le \nu +2\ell \). This valuation is odd. The retained pole obstruction forces an odd valuation of \(b\) to be positive, whence \(\ell >-\nu /2\). At a place attaining \(\ell \) every factor and every term has \(\text{valuation}>-\nu /2\). Thus

\begin{gather*}
v(C_{2}),v(C_{3}),v(C_{4}),v(G)>9\nu /4;\\
v(I)>9\nu /8; v(C_{1})>3\nu /8;\\
v(D_{H})>3\nu /4; v(A_{2})>5\nu /4;\\
v(A_{1}),v(A_{3}),v(A_{4})>3\nu /4; v(Bi)>7\nu /8
\end{gather*}

Integer rounding gives the required high order3, input orders at least1, and positive even reciprocal valuations (at least2). The source's factor lemma for this exact exponent pattern is Brox228 \nolinkurl{lem:residues} and \nolinkurl{prop:extract}; the argument has been reapplied to the actual recovered raw224 objects above.

Only after these positive reciprocal tests do all relevant nonidentity points lie in their respective formal groups. The source's constant-term logarithm estimate has relative precision4 at depth1. Write \(h=m/2\) for the reference index. The integral group identities yield logarithmic ratios \(u_{i}/h\) and \((u_{i}-n_{i})/(h-1)\), killing the possible two-torsion first. At least one of h,h-1 is a p-adic unit. That ratio is integral; the positive \(C_{i}\) test transfers integrality to the other plain ratio and then the full ratio. Relative errors may now be used as absolute errors modulo \(p^{4}\). For \(i>1\) the \(C_{i}\) test gives \(L_{i}=n_{i}\) modulo \(p^{3}\); for \(i=1\) the \(C_{1}\) and \(I\) tests give \(L_{1}=n=\beta \) modulo \(p\). Identity numerators, half-integral \(h\) and inert places are covered exactly as in Version15 \nolinkurl{Integrality before the formal congruences}; no invertibility of the tested numerator is needed. Consequently each comparison radical cubed divides the integer Lucas residual \(\Delta \), and each radical divides the reduced input residual \(J\). There is no added scalar factor \(c\) in these divisibilities.

\subsection{Height supplier and global consumers}

The sibling height theorem gives \(h(u)\le C_{2}\) \(M_{S_{10}}(u)^{\alpha }\) on the full five-contact family, \(\alpha =2103/1024<7/3\). To apply it to this unchanged public comparison set, let \(B_{\Sigma }\) be the product of the fixed primes in \(\Sigma \) outside \(S_{10}\). Then \(M_{S_{10}}(u)\le B_{\Sigma }\) \(M_{\Sigma }(u)\); all radicals are squarefree even when contact valuations are arbitrarily deep. Thus \(h(u)\le H\) \(M_{\Sigma }(u)^{7/3}\), with fixed \(H=C_{2}*B_{\Sigma }^{\alpha }\). No sharp numerical value of this \(H\) is needed or asserted.

Applying this bound to2+1/t and2+1/(2-t), whose inverse projective matrices have row norms3 and7, and multiplying the two inequalities after subtracting log7 gives

\begin{gather*}
h(t)\le H_{2} M_{2}^{\circ }(t)^{7/6}, H_{2}=\log 7+H
\end{gather*}

The two block radicals outside \(\Sigma \) are coprime and multiply to the same physical comparison radical. The inequalities also hold when \(h(t)\le \log 7\) and on the harmless \(t=1\) stratum (radical1). All contact domains were already excluded by the literal rows. This source-to-consumer passage uses the new height theorem on both transformed arguments, with index-set coverage and fixed discarded primes explicitly accounted for.

The retained aggregate trace theorem gives \(B^{2}\le M_{1}*h(T)+M_{2}\), \(B=\operatorname{max}(|n|,|a|,|t|,|b|)\), and \(M_{2}\) is ineffective. Since \(T=t_{1}/t_{0}\), \(h(T)\le h(t_{1})+h(t_{0})\). If each radical to power \(k\) divides a nonzero integer \(J_{0}\), then

\begin{gather*}
B^{2}\le C |J_{0}|^{7/(6k)}+M_{2}, C=2*M_{1}*H_{2}
\end{gather*}

The unchanged scalar conclusions are \(c=|a|>0\), \(B\ge c\), \(c\ge 149^{{L}_{N}}\), \(L_{N}=\text{floor}(9(2N-1)/4)+1\), and

\begin{gather*}
h(\beta )\le C_{3}*149^{-L_{N}}*c*(1+\operatorname{log}(c)/3), C_{3}=2^{330770}
\end{gather*}

They are the actual four-contact/parity-scalar conclusions in Version15 \nolinkurl{v222:fourcontact} and the scalar-S step, not the old \(c=|a-1|\) assertion.

If \(\Delta \ne 0\), \(|\Delta |\le 14B^{4}\). Take \(k=3\) to obtain \(B^{2}\le C*14^{7/18}*B^{14/9}+M_{2}\). The right exponent is strictly below2, so \(B\) is bounded by a fixed constant. Uniform \(c\ge 149^{{L}_{N}}\) contradicts this for all sufficiently large \(N\). Hence \(\Delta =0\). The retained integer Lucas argument then gives log \(B\ge (c-5)\operatorname{log}(c-2)\) for large \(c\) (the scalar-S variant is stronger).

If \(n-\beta \ne 0\), its reduced numerator \(J\) satisfies \(|J|\le 2B\) \(\operatorname{exp}(h(\beta ))\). Take \(k=1\); for \(B^{2}\ge 2M_{2}\) we have \(B^{5/6}\le 2C*2^{7/6}*\operatorname{exp}((7/6)h(\beta ))\). Thus

\begin{gather*}
\operatorname{log} B\le (6/5)\operatorname{log}(2C)+(7/5)\log 2+(7/5)h(\beta )
\end{gather*}

As \(N\) grows, the coefficient \((7/5)C_{3}*149^{-{L}_{N}}\) becomes at most1. The right side is then bounded by a fixed constant+c+(c log c)/3, contradicting the Lucas lower bound for sufficiently large \(c\). Therefore \(n=\beta =51X^{2}+2504\) is integral; for \(X=u/v\) reduced, \(v^{2}|51\) and squarefreeness of51 imply \(v=1\). Every threshold depends only on fixed constants and \(L_{N}\) grows monotonically, so a single ineffective \(N_{0}\) works for every larger \(N\).

\subsection{Fresh completeness and the218 control}

For an integer input use the unchanged scalar-S/four-contact/parity-scalar preparation, then the root labels, positive trace scale, fixed real neighborhoods and finite progression in the precise order of the frozen Version15 construction. Neither changed factors nor eliminated fourth-power caches occur in that preparation. The source chooses a common denominator dilation and the same ten contact forms per scale, merges repeated forms, and chooses the common reference index by the unchanged separated reduction-order congruences. No new prime pattern, scale or denominator is introduced here.

At a selected contact prime \(v(q)=-1\), all \(C_{i},I,G\) have \(\text{valuation}\ge 4\) and every reciprocal abscissa has \(\text{valuation}\ge 2\). The changed terms \(q^{5}\) \(C_{i}^{2}\) have \(\text{valuation}\ge 3\), \(q^{5}\) \(I^{4}\ge 11\) and \(q^{3}\) \(A_{2}^{2}\ge 1\); every other nonconstant term was already positive. Hence \(W=1\) modulo the chosen place. At the opposite place the newly proved noncancellation lemma still applies; trace residues are1 or2. The comparison pole input therefore has the same odd depth1 at selected primes and even valuations away from them.

At prepared finite primes \(q\) tends arbitrarily deeply to0, so every positive exponent used in the new factors makes its term small. The source's squareclasses and depths4,2,2 for the pole input at2,3,17 are unchanged. At both real embeddings \(q>0\) is small, all test powers are even, \(W\) tends to1, and the pole input tends to \(\gamma ^{2}/4\) in(0,17). The same simultaneous pole-realization theorem supplies independent witnesses for each detector. These are new complete witnesses for the changed factors, not reuse of old pole coordinates. The selected contact primes ensure \(q_{0},q_{1}\ne 0\) \((\text{also} \text{real} q>0)\), so all reciprocal weights exist. Reconstruct each deleted fourth power and apply the inherited exact raw224 coordinate maps; the symbolic identities give the new quadratic rows. Completeness is uniform for all sufficiently large \(N\) as in the parent scalar preparation.

For218 keep the parent's original factor exponents. The same exact clearing proof identifies its witnesses with the old ones on \(q_{0}*q_{1}\ne 0\). Soundness follows by reconstructing the unrestricted old system. Completeness uses the parent's selected witnesses with \(q>0\). Thus no claim that every old witness has \(q_{0}\ne 0\) is needed, and no new height theorem is needed for the218 control. This corrects the contribution's occasional unqualified phrase "solution-set transport": the all-solution bijection is restricted to the nonzero-weight domain; the rational projection agrees without that qualification by the two arguments just given.

\subsection{Scope and evidence}

The finite checker proves exact component identities, complete row occurrence coverage, inverse weights, unchanged rows, and full expansions, with deliberate corruptions rejected. These checks do not certify the retained arithmetic premises. The new214 argument has explicit local, completeness and height consumers above; their independent review and inventory coverage will be bound separately. No claim of public priority, a numerical defining parameter, or whole-proof completion is made. The construction deliberately branches at224; the final paired-root/trace-quotient reductions yielding222 are not counted again because they conflict with this clearing strategy.
\clearpage
% Source: research/quartic214/proofs/uniform-existence.md
\section{Existing 214: uniform existence and actual completeness interfaces}\label{q214:uniform-existence}
\noindent\textit{Source: }\nolinkurl{research/quartic214/proofs/uniform-existence.md}.\par
\begin{quote}\small \textbf{Current foundation overlay:} The Version~15 byte-identity and no-new-repair statements in this dated source describe its historical freeze.  This edition instead incorporates the complete current Version~16 foundation and the explicit source-to-consumer transition in Section~\ref{q214:foundation-transition}.  The historical text is retained as evidence; it does not negate or independently accept that changed application.\end{quote}

\begin{quote}\footnotesize Author: \nolinkurl{/root/d214_uniform_existence}, 2026-11-11. Author argument only; fresh non-author acceptance is required. No source, old packet, coefficient, witness table or public repository is changed. This is an additive clarification of the frozen existing 214 family, with its original hypotheses retained exactly. Source bindings and actual consumer passages are in \nolinkurl{SOURCE_USES.json} and \nolinkurl{CONTRACT.json.}\end{quote}

[P01: exact theorem] Let the 214 family be exactly the family with 188 primitive integral quadratic residuals in \nolinkurl{research/quadratic-ten-contact-integration/existence/artifacts/candidate214-brox/}, at private HEAD \nolinkurl{bfeb8400982c9a169e4fdd39fb68ec053621e8ed.} Its parameter is the fixed coefficient \(A_{N}=149^{2N+1}\), not an existential coordinate. Subject to the named full-family height and retained elliptic/rank interfaces below, the conclusion is

\begin{gather*}
\exists  N_{0} \in  Z_{\ge 1}, \text{ for }  \text{ every }  \text{ integer }  N\ge N_{0}, \text{ for }  \text{ every }  x \in  Q,\\
x \in  Z \Longleftrightarrow  \text{ there }  \exists  y \in  Q^{213} \text{ with }  P_{N}(x,y)=0
\end{gather*}

This is the same polynomial family. \(N_{0}\) is not computed. \(N\ge 53\) is a coefficient-height condition only; it need not be included in the arithmetic \(N_{0}\), and replacing \(N_{0}\) by \(\operatorname{max}(N_{0},53)\) is harmless. No \(N=53\) defining theorem or numerical defining member follows.

[P02: scope of source carry] The entire latest PR26 manuscript at public checkout HEAD \nolinkurl{1a49013b5683f58d2f88bb018a07d13770354157} and the old 214 frozen \nolinkurl{existence/sources/quartic_v15.tex} have the identical \textbf{raw byte} SHA256 \nolinkurl{44ba3a8ec8521c29b81b26cfb9981cd6d5f3c76d54e7782438614842edb52c74} and 1924438 bytes. This is stronger than normalized-line identity. Thus no analytic, restricted-converse, aggregate, scalar or prime-selection manuscript repair is newly imported from PR26 here. The later audit records can justify exact unchanged applications at their stated scope only after source/consumer comparison; PR status is not an arithmetic premise. This author read the actual source passages identified in \(\text{SOURCE}_{\text{USES}}\), including the sum-pole/graph argument, the prime theorem/obstruction-tail/fixed-field and diagonal statements, the CRT argument, scalar-S/four-contact preparation and zero-preparation proof. This author did not reconstruct every internal RC/analytic sublemma or final published source anew. Those unchanged internal foundations and their eligible source review remain a separately bound dependency; this packet closes the 214 application joins, not their entire recursive review.

[P03: exact retained inputs] Use \(K=Q(\sqrt{2})\), \(d=3-\sqrt{2}\), \(E:y^{2}=x^{3}-d^{2}\) \(x\), and the same \(P=2(1+\sqrt{2},2\sqrt{2})\). The four nonzero integer labels are (n,a,t,b), \(c=|a|\) and \(B=\operatorname{max}(|n|,|a|,|t|,|b|)\), so \(B\ge c\). The physical scale is \(T=(1/4)\operatorname{Tr}(\sum _{i} x([n_{i}]P))>0\). The physical arguments are \(t_{0}=\lambda \) and \(t_{1}=T\) \(\lambda \). The contact set is \(\Theta _{2}=\Theta \) union(2-Theta), \(\Theta =\{0,-1/2,-1,1/2399,-32/407\}\); \(\Sigma \) is the \(\text{primes}\le 76176\) together with 2029147. Each comparison radical \(M_{j}\) is the squarefree product of primes \(p\) outside \(\Sigma \) at a positive odd contact of \(t_{j}\). Do not substitute the sibling center-3/11 contact system or its input \(51(2X+1)^{2}+2504\). Here \(\beta =51x^{2}+2504\).

[P04: restored objects and branches] The old existence proof's exact Laurent clearing and raw224 inverse maps reconstruct the physical reference/scalar/root system before any ratio is consumed. Its recovered f,S,g and \(Q_{2}\) coefficient at149 imply \(w=g^{4}/f^{2}\); deleted fourth powers are their actual squares. The shared alternate ratio has numerator and denominator of valuation -v149(f)/2 and is nonzero. The unchanged scalar argument excludes reference torsion and \(H\) indices -1,0,1 before logarithms. No rational component or tested numerator is divided by. The new inverse rows enforce \(q_{j}\ne 0\), \(q_{j}=1/R_{j}-1\). They remove an irrelevant old \(q_{j}=0\) witness stratum, so this is not an all-old-witness bijection. The stratum \(\delta _{j}=t_{j}-1=0\) has \(q_{j}=-1\) and \(M_{j}=1\), and needs no division by \(\delta _{j}\). A contact \(t_{j}\) in \(\Theta _{2}\) is excluded by \(F_{2}(t_{j})(1+q_{j})=\delta _{j}^{10}\) and \(F_{2}(1)\ne 0\). A putative \(q_{j}=0\) has radical one, but is in any event excluded in 214. None of these statements confuses \(\delta _{j}\) with the integer Lucas residual \(\Delta \) below.

[P05: changed local factors] At a consumed \(p\) outside \(\Sigma \) and either place above \(p\), normalize \(v(p)=1\). Exactly one contact factor has positive odd valuation \(\nu \), \(\delta _{j}\) is a unit and \(v(q_{j})=-\nu \). The five actual 214 factors are

\begin{gather*}
1+q^{5} C_{2}^{2}+q^{2} D_{H}^{2}+q^{4} B_{1}^{4};\\
1+q^{5} G^{2}+q^{2} A_{3}^{2}+q^{4} B_{0}^{4};\\
1+q^{5} I^{4}+q^{3} A_{2}^{2}+q^{2} A_{4}^{2}+q^{4} B_{3}^{4};\\
1+q^{5} C_{3}^{2}+q^{2} C_{1}^{4}+q^{4} B_{4}^{4};\\
1+q^{5} C_{4}^{2}+q^{2} A_{1}^{2}+q^{4} B_{2}^{4}
\end{gather*}

Here \(D_{H}=f(H)\), \(A_{i}=f(U_{i})\), \(B_{i}=\sigma \) \(f(V_{i})\), \(C_{i}=L_{i}-L_{i}'\), \(I=L_{1}-\beta \) and \(G=L_{3}^{2}-L_{2}(L_{2}-2)L_{3}L_{4}+(L_{2}-2)^{2}L_{4}^{2}-1\), with \(f=1/x\) and \(f(O)=0\). Every nonzero reciprocal abscissa has even valuation, including \(f(H-P)\) restored before this step. In the generic factors the finite nonconstant valuations lie in odd integers, 2+4Z and 4+8Z. In the third factor the first three have residues \(-\nu \), \(\nu \), 2 modulo4 and the last is4 modulo8. They are pairwise distinct and distinct from zero. Zero tests omit leaves. Thus each factor is nonzero with \(\text{valuation}\le 0\), and has residue1 whenever its valuation is zero, at both split places and at the inert place. This argument treats arbitrary negative test valuations too.

[P06: trace/pole extraction] Write \(W\) for the product and \(\ell =\operatorname{max}_{v|p}v(W)\le 0\). The literal cleared trace equation \(w_{0}^{2}-2w_{0}*\tau -2w_{1}^{2}=0\), with \(W=w_{0}+\sqrt{2}*w_{1}\), forces \(w_{0}\ne 0\) since \(W\ne 0\), and \(\tau =1/\operatorname{Tr}(W^{-1})\ne 0\). The actual pole input is \(\mathfrak{b}=\gamma ^{2}\) \(R/\operatorname{Tr}(W^{-1})^{2}\), \(\gamma =408/409\), a p-unit. Trace gives \(v_{p}(\operatorname{Tr}(W^{-1}))\ge -\ell \), so \(v_{p}(\mathfrak{b})\le \nu +2\ell \). Its valuation is odd. The zero-preparation pole obstruction excludes negative odd valuations outside2,3,17 and hence forces it positive. Therefore \(\ell >-\nu /2\). At a place attaining \(\ell \) every factor valuation and every finite summand valuation exceeds \(-\nu /2\). Solving gives \(v(C_{2},C_{3},C_{4},G)>9\nu /4\), \(v(I)>9\nu /8\), \(v(C_{1})>3\nu /8\), \(v(D_{H})>3\nu /4\), \(v(A_{2})>5\nu /4\), \(v(A_{1},A_{3},A_{4})>3\nu /4\) and \(v(B_{i})>7\nu /8\). Integer rounding and even reciprocal valuations give cross/Lucas \(\text{orders}\ge 3\), input \(\text{order}\ge 1\), and every reciprocal \(\operatorname{depth}\ge 2\). No implication is extracted from a radical-one detector alone.

[P07: exact integer consumers] Positive reciprocal valuations put \(H,U_{i}\) in the primary formal group and \(H-P,V_{i}\) in its conjugate formal group, with the identity convention at zero numerator points. Write \(H=[m]R_{0}+T_{0}\), \(h=m/2\); all consumed \(p\) are odd. The integral identities \([m]U_{i}=[2u_{i}]H\) and \([m-2]V_{i}=[2(u_{i}-n_{i})](H-P)\) kill two-torsion. At least one of h,h-1 is a p-adic unit, so one full logarithmic ratio is integral. The positive \(C_{i}\) transfers integrality of the corresponding plain ratio to the other; formal logarithm units then give integrality of both full ratios. Only now is the relative error of order4 used as an absolute congruence. For \(i>1\), \(L_{i}=n_{i}\) mod \(p^{3}\); for \(i=1\), \(L_{1}=n=\beta \) mod \(p\). Consequently, for

\begin{gather*}
\Delta =t^{2}-a(a-2)t*b+(a-2)^{2} b^{2}-1,
\end{gather*}

one has \(M_{j}^{3}\) | \(\Delta \) for each \(j\), and \(M_{j}\) | \(J\) when \(J\) is the reduced numerator of n-beta. These are separate divisibilities, with no factor \(c\) and no requirement that \(M_{0}\) and \(M_{1}\) be coprime. Rational \(\beta \) is p-integral whenever an input congruence is consumed. If \(J\ne 0\), \(|J|\le 2B\) \(\operatorname{exp}(h(\beta ))\); if \(\Delta \ne 0\), \(|\Delta |\le 14B^{4}\). These are the exact objects used below.

[P08: aggregate source application] The actual aggregate input is the \textbf{sum} theorem, not the odd-label product theorem. On \(E_{0}:Y^{2}=X^{3}-X\) let P1,P2 be the pointed-isomorphism images of \(P\) and its conjugate. The independent squareclasses d,d' in \(K(i)\) provide an automorphism fixing \(P_{1}\) and negating \(P_{2}\) while fixing Z[i]; applying it to a torsion relation proves their Z[i]-independence modulo torsion. For arbitrary nonzero labels, at any common pole of summands in one block, the pole coefficient is \(\lambda _{i}^{-2}\) \(\sum _{j}\) \((n_{i}/n_{j})^{2}\), a nonzero positive rational multiple of \(\lambda _{i}^{-2}\). Cross-block common poles would contradict that independence. Restriction to a nonzero elliptic subgroup direction therefore excludes every positive-dimensional abelian translate through the actual tuple in a sum fiber. This covers negative, repeated and noncoprime labels with no parity condition.

[P09: family/normalization join] Apply final published Von Buhren2017 Theorem1.1, printed308, to \(A=E_{0}^{8}\), parameter \(P^{1}\), and the integral projective graph closure of \(F_{+}=(d/2)\sum \) \(X_{i}+(d'/2)\sum \) \(X_{i}'\). The fixed finite-rank group is generated by \(P_{1}\) in each first-block coordinate and \(P_{2}\) in each second-block coordinate. Its points lie in the theorem's \(C(\Gamma ,\epsilon )\) by taking error point0. Embed \(E_{0}\) by \(O(3O)\), and its product by Segre. Restriction of degree-n plane forms to a smooth plane cubic is surjective for all \(n\ge 1\) from \(0\to O(n-3)\to O(n)\to O_{E}(n)\to 0\) and \(H^{1}(P^{2},O(n-3))=0\). Tensoring and the Segre multiplication map provide the source's embedding condition(1.1) for the product. \(P^{1}\) uses \(O(1)\). The actual tuple avoids all coordinate poles since P1,P2 are nontorsion and labels nonzero. A translate in a projective graph fiber through this tuple would meet the function domain densely and make \(F_{+}\) constant there, contradicted by \(P_{08}\); graph boundary components do not change this. \(F_{+}(\text{tuple})=2T\). The fixed product canonical height is a positive multiple of the sum of all eight coordinate canonical heights, hence twice a multiple of the four original heights. Since \(h(2T)\le h(T)+\log 2\) and \(\widehat{h}([n_{i}]P)=n_{i}^{2}\) \(\widehat{h}(P)\), \(\widehat{h}(P)>0\), the theorem gives fixed \(M_{1}>0,M_{2}\ge 1\) with \(B^{2}\le M_{1}\) \(h(T)+M_{2}\) for every tuple. The constants are fixed before x,N or labels; the additive \(M_{2}\) is ineffective, exactly as in the published source. No full rank-one theorem is required for this aggregate step.

[P10: height domain and maps] Retain the sibling full-family supplier \(h(u)\le C_{2}\) \(M_{S_{10}}(u)^{\alpha }\), \(\alpha =2103/1024<7/3\), on every rational argument avoiding its four finite contacts; the full geometric/modular/height chain supplying it remains a separately reviewed interface. Put \(B_{\Sigma }=\prod \) of primes in \(\Sigma \) outside \(S_{10}\). Squarefreeness gives \(M_{S_{10}}\le B_{\Sigma }\) \(M_{\Sigma }\), even at deep contacts, so \(h(u)\le H\) \(M_{\Sigma }(u)^{7/3}\), \(H=C_{2}\) \(B_{\Sigma }^{\alpha }\) fixed. Apply this to \(u=2+1/t_{j}\) and \(u=2+1/(2-t_{j})\). Their excluded finite contacts correspond exactly to the ten excluded \(t_{j}\) contacts, and \(t_{j}=0,2\) are already excluded; the inverse projective row norms are3 and7. If \(h(t_{j})>\log 7\), multiply the two inequalities for \(h(t_{j})-\log 7\). Outside \(\Sigma \) the two block radicals are disjoint and their product is exactly \(M_{j}\), giving \(h(t_{j})\le H_{2}\) \(M_{j}^{7/6}\), \(H_{2}=\log 7+H\). If \(h(t_{j})\le \log 7\) it follows directly, including \(t_{j}=1\) with \(M_{j}=1\). Thus the argument covers the full rational solution domain, not a bounded example. Because \(t_{0}\ne 0\) and \(t_{1}=Tt_{0}\), \(h(T)\le h(t_{0})+h(t_{1})\). If \(M_{j}^{k}\) divides the same nonzero integer \(J_{0}\) separately for both \(j\), then

\begin{gather*}
B^{2} \le  C |J_{0}|^{7/(6k)}+M_{2}, C=2 M_{1} H_{2}
\end{gather*}

This is the source-to-physical-radical join; no product divisibility between scales is asserted.

[P11: a single soundness threshold] The actual scalar-S/four-contact argument gives \(c>0\), \(c\ge 149^{{L}_{N}}\) and \(h(\beta )\le C_{3}*149^{-{L}_{N}}*c(1+\operatorname{log}(c)/3)\), \(C_{3}=2^{330770}\), \(L_{N}=\text{floor}(9(2N-1)/4)+1\). This uses \(c=|a|\), not |a-1|. For \(\Delta \ne 0\), P07/P10 imply \(B^{2}\le D\) \(B^{14/9}+M_{2}\) where \(D=C*14^{7/18}\). Hence \(B\le K_{D}=\operatorname{max}(\sqrt{2M_{2}},(2D)^{9/4}):\) when \(B^{2}>2M_{2}\), move \(M_{2}\) and divide by \(B^{14/9}\), obtaining \(B^{4/9}<2D\). Choose \(N_{D}\) so \(149^{{L}_{N}}>K_{D}\) for all \(N\ge N_{D}\). This is possible by monotonicity and excludes \(\Delta \ne 0\) uniformly in every \(x\) and every solution.

[P12: Lucas lower bound and input threshold] With \(\Delta =0\) and a,t,b nonzero, set \(A_{0}=|a|=c\) and \(w=\text{sign}(a)(a-2)b\). Then \(t^{2}-c*t*w+w^{2}=1\). The elementary norm-one unit reduction in \(Z[(c+\sqrt{c^2-4})/2]\) gives \(|w|=U_{j}\) and |t| in \(\{U_{j-1},U_{j+1}\}\), where \(U_{0}=0,U_{1}=1,U_{j+1}=cU_{j}-U_{j-1}\). Nonzero \(w\) excludes \(j=0\). Reduction modulo c-2 \((if a>0)\) or c+2 \((if a<0)\), together with |a-2| dividing \(w\), gives \(j\ge c-2\). Induction gives \(U_{j}\ge (c-1)^{j-1}\); hence for \(c\ge 8\) the weaker bound log \(B\ge (c-5)\operatorname{log}(c-2)\) holds. If \(J\ne 0\) and \(B^{2}\ge 2M_{2}\), \(P_{10}\) at \(k=1\) gives

\begin{gather*}
\operatorname{log} B \le  K + \epsilon _{N} c(1+\operatorname{log}(c)/3),\\
K=(6/5)\operatorname{log}(2C)+(7/5)\log 2, \epsilon _{N}=(7/5)C_{3}*149^{-{L}_{N}}
\end{gather*}

Choose \(N_{e}\) so \(\epsilon _{N}\le 1\) for every \(N\ge N_{e}\). Put \(K_{c}=\operatorname{max}(\operatorname{exp}(12),2\operatorname{max}(K,0)+1,8,\sqrt{2M_{2}}+1)\). For \(c\ge K_{c}\),

\begin{gather*}
(c-5)\operatorname{log}(c-2)-c-(c/3)\operatorname{log} c\\
\ge  c[(\operatorname{log} c)/6-1-(\log 2)/2] > c/2 > K
\end{gather*}

Here \(c-5\ge c/2\) and \(c-2\ge c/2\) hold. Choose \(N_{c}\) so \(149^{{L}_{N}}\ge K_{c}\) for every \(N\ge N_{c}\). Then \(J\ne 0\) is impossible. Thus \(n=\beta \) is integral. Writing \(x=u/v\) reduced with \(v>0\), integrality forces \(v^{2}|51\); since51 is squarefree, \(v=1\). The same fixed \(N_{S}=\operatorname{max}(N_{D},N_{e},N_{c})\) works for every rational input and all solutions at every \(N\ge N_{S}\). This construction is existential, and \(K_{D},K_{c}\) depend on the ineffective \(M_{2}\).

[P13: uniform scalar completeness threshold] For integer \(x\), the retained four-factor scalar \(r_{4}=\beta (\beta -1)(\beta -2401)(32\beta +343)\) is an \(\text{ integer } \ge 1\). With \(R_{s}=A_{N}\) \(r_{4}/(1+149^{2} r_{4})\), one has \(R_{s}\ge A_{N}/(1+149^{2})\), independently of \(x\). Choose the fixed compact real neighborhoods from the scalar-S proof before finite progression moduli; on them Z0,Z1,F are bounded by constants \(M_{Z},M_{F}\) independent of x,N. The actual parity-normalized scalar is

\begin{gather*}
W_{s}=Z_{0}^{4}/R_{s}^{10}+Z_{1}^{2}/R_{s}^{5}+F^{2}/R_{s}^{4},\\
b_{s}=R_{s}(R_{s}+9792)W_{s}^{2}
\end{gather*}

For \(R_{s}\ge 1\), \(W_{s}\le K_{W}\) \(R_{s}^{-4}\), \(K_{W}=M_{Z}^{4}+M_{Z}^{2}+M_{F}^{2}\), so \(0<b_{s}\le 9793\) \(K_{W}^{2}\) \(R_{s}^{-6}\). Choose a fixed \(N_{C}\) so the right side \(\text{ is } <17\) for every \(N\ge N_{C}\). This is uniform in the integer input. \(F\ne 0\) and real positivity exclude \(b_{s}=0\). The neighborhoods and constants are those of the original scalar functions, before the later root-label-dependent comparison neighborhoods; the latter need no bound uniform in labels.

[P14: scalar finite preparation in correct order] For a particular x,N, put \(D_{s}=1+149^{2}\) \(r_{4}\). Both \(\text{gcd}(A_{N} r_{4},D_{s})\) and \(\text{gcd}(A_{N} r_{4}+9792D_{s},D_{s})\) equal1; thus \(R_{s}(R_{s}+9792)\) has square reduced denominator \(D_{s}^{2}\). Direct integer parity/congruences give \(v_{2}(R_{s})\) in \{2,3\} and \(v_{3}(R_{s})=v_{17}(R_{s})=0\). Hence \(1+9792/R_{s}\) is a square at2,3,17 \((\text{its} \text{difference} \text{from}1 \text{has} \operatorname{depth}\ge 3,\ge 1,\ge 1 \text{ respectively } )\), so \(R_{s}(R_{s}+9792)\) is a square there. At its finite odd-valuation numerator support choose an initial \(H\) formal enough; then choose the positive scalar index a still deeper and arbitrarily large, with [a]P in the fixed real neighborhood of \(O\). This makes the three terms of \(W_{s}\) integral at that support and deep enough at2,3,17 for the actual zero-preparation pole bounds. Compact-group recurrence in finite-index progressions gives the finite and real conditions together. After a is fixed, continuity preserves the scalar conditions for \(H\) in one progression \(h=h_{*}\) mod \(L\), \(\text{ with } 64|h_{*},64|L\) and the component-order divisibility included. This ordering permits support and depth to depend on x,N; the common \(N_{C}\) does not. The retained scalar/root chart guards hold on the chosen neighborhoods. No new condition \(a=1\) mod any modulus is imposed.

[P15: root labels, coefficients and incidences] Put \(n=\beta \). For the chosen positive \(a\ge 3\) choose \(j>0\) divisible by a-2; with the same Lucas recurrence put \(b=U_{j}/(a-2)\) and \(t=U_{j+1}\). Then n,a,t,b are positive nonzero integers and \(\Delta =0\). Fix their four root points and \(T>0\). Fix compact real neighborhoods for \(H\) and all comparison rational functions avoiding their denominators and chart guards. These neighborhoods may depend on the fixed labels. Enlarge the finite preparation set to contain \(\Sigma \), scale numerator/denominator primes, scalar-support primes and fixed guards. Choose \(D>0\) so all numbers \(D^{2}\) \(\theta /a_{*}\) are integers for \(a_{*}\) in \{1,T\}, \(\theta \) in \(\Theta _{2}\), and so \(v_{p}(a_{*}\lambda )<0\) with the needed depth at every prepared \(p\) when \(\lambda =u/(D^{2} w^{2})\). Use coefficients \(c_{a_{*},\theta }=-D^{2}\) \(\theta /a_{*}\); merge equal coefficients \textbf{across} detectors, retaining each incidence. Within a detector distinct \(\theta \) give distinct coefficients; at least ten distinct coefficients remain. The zero contact supplies coefficient0 and the prime \(u\). This is exactly the old center-two pattern.

[P16: prime theorem quantifiers and CRT] Apply the actual \(V_{15}\) simultaneous-prime theorem to this fixed list, prescribed divisibility q|w, fixed avoidance modulus \(M=\operatorname{rad}_{\text{odd}}(3*17*L)\) with7 omitted, and arbitrarily large slope/prime lower bounds. Choose these data before the analytic asymptotic parameter \((\text{call} it X_{\text{sel}}, \text{distinct} \text{from} \text{input} x)\). Its fixed-field proof first fixes the cutoff so positive cutoff mass exceeds the full obstruction tail, then fixes the ray field, then lets \(X_{\text{sel}}\) grow along sufficiently slow diagonal strip widths. No uniform estimate in x,N,T,D,L or number of contacts is needed. The selected Gaussian generators have residues9 mod16,1+3i mod7,i modM; adjacent norms \(d_{\text{minus}},d_{\text{plus}}\) have \(v_{2}=6,2\) and cross-sign gcd4, and all odd factors avoid \(L\). Define \(U=\text{lcm}(d_{\text{minus}}/4)\), \(V=\text{lcm}(d_{\text{plus}}/4)\). Full two-torsion makes d/4 kill the reduction of the double \(P\); \(\text{gcd}(U,V)=1\), \(\text{gcd}(U,L)=16\) and \(\text{gcd}(V,L)=1\). Thus \(h=0\) mod \(U\), \(h=1\) mod \(V\), \(h=h_{*}\) mod \(L\) are compatible. The real progression lemma gives arbitrarily large \(h\) in the original fixed compact neighborhoods. Set \(H=[h]P\), \(U_{i}=[hn_{i}]P\), \(V_{i}=[(h-1)n_{i}]P\). This gives natural formal \(\operatorname{depth}\ge 1\) at the correct primary/opposite places and preserves the scalar data. The optional strengthened separation \(p_{i}\) not dividing \(d_{j}\) is available from the weighted strip argument; natural depth1 alone suffices for this 214 construction. Every merged repeated coefficient carries one prime/generator and all its incidences; CRT is imposed once and supplies all those incidences.

[P17: changed factors at the actual fresh witnesses] At any selected contact \(p\), the slope and prime cutoff make \(p\) avoid D,w and both scale supports; selected primes are distinct and larger than \(w^{2}\). Thus exactly its indicated contact numerator has depth1 in each incident detector, \(\delta \) is a unit, \(v(q)=-1\), every reciprocal test has \(\operatorname{depth}\ge 2\), and every \(C_{i},I,G\) has \(\operatorname{depth}\ge 4\) by the full ratios \(n_{i}\) and relative formal error4. The changed factor terms have minimal depths \(q^{5}\) \(C_{i}^{2}:3\); \(q^{5}\) \(I^{4}:11\); \(q^{3}\) \(A_{2}^{2}:1\); \(q^{2}\) \(D_{H}^{2}\) or \(q^{2}\) \(A_{i}^{2}:2\); \(q^{2}\) \(C_{1}^{4}:14\); \(q^{4}\) \(B_{i}^{4}:4\). Therefore \(W=1\) mod the selected place. At the opposite place the \(P_{05}\) residue argument applies without a formalness assumption: \(W^{-1}\) reduces to0 or1, so \(\operatorname{Tr}(W^{-1})\) reduces to1 or2, a unit since \(p\) is odd. The new actual pole input \(\mathfrak{b}\) has valuation1 there. This is a fresh order-three-factor proof, not an appeal to the old order-four conclusion.

[P18: all other primes and completions] For each physical scale \(a_{*}\), the exact identity is

\begin{gather*}
R = [a_{*}^{10} \prod _{\theta  \in  \Theta _{2}}(u-(\theta /a_{*})D^{2} w^{2})] /\\
[(D^{2} w^{2})^{10} \delta ^{10}]
\end{gather*}

Every non-prime factor here is a rational square, including \(a_{*}^{10}\) for a nonsquare scale; each contact coefficient occurs once within a detector. Hence outside its selected prime numerators, \(v_{p}(R)\) is even. This includes primes dividing \(\delta \). Multiplication by \(\gamma ^{2}\) and division by the squared trace preserve that parity. At prepared \(p\), \(D\) and the fixed divisibility q|w make \(R=\prod _{\theta }(1+(1-\theta )/\delta )\) as close to1 as needed and the weight \(q=R^{-1}-1\) arbitrarily small. All positive exponents in the changed factors therefore make test terms small there uniformly on the fixed p-adic neighborhoods, and \(\operatorname{Tr}(W^{-1})\) is close to2. Choose closeness so \(R\) is a square at2,3,17. As \(v_{2}(\gamma )=3\) and \(v_{3}(\gamma )=v_{17}(\gamma )=1\), the actual pole input is a local square with depths4,2,2. At both real embeddings, as \(\lambda \) grows on the fixed neighborhoods, \(R=\prod _{r=1}^{5}(1-\alpha _{r}/\delta ^{2})\), \(\alpha _{r}>0\), so \(q>0\) tends to0; all factors are positive and \(W\) tends to1. Thus \(\mathfrak{b}\) tends to \(\gamma ^{2}/4\) in(0,17). The prime theorem permits an arbitrarily high fixed slope after these neighborhoods are chosen. This proves positivity and a nonzero trace globally, and \(q_{0},q_{1}>0\) so all new inverse weights exist.

[P19: actual pole theorem and its curve domain] The input just constructed has square reduced \(\text{denominator},0<\mathfrak{b}<17\), local squares at2,3,17 and depths4,2,2, meeting the exact preparation-zero predicate Psi0 (source requires3,1,2). In its proof set \(B_{\text{pole}}=17\mathfrak{b}\) and \(V_{\text{pole}}=B_{\text{pole}}^{2}(1-\mathfrak{b}/17)\). The positive squarefree representatives mu,theta have disjoint support avoiding2,3, \(\mu =17\) \(\text{mod}24,\theta =1\) mod24, mutual local squareness, and \(\theta \) may equal1. The simultaneous-descent lemma constructs positive squarefree \(n_{\text{pole}},n_{\text{pole}}\) \(\theta \), both11 mod12 and with full2-primary Selmer corank1; this uses full Selmer corank, including divisible Sha, before any finiteness claim. The retained RC is for \(E_{l}:y^{2}=x(x-l)(x+3l)\), \(l>0\) \(\text{ squarefree } ,l=11\) mod12. These curves have \(j=35152/9\), are non-CM, and are distinct from the CM index curve E/K. RC is applied separately to these exact two parameters; no constants uniform in them are required. The published forward analytic-rank-one input then gives rank1 and finite Sha, and Kummer gives points with the required nonzero abscissa class. The zero-preparation proof explicitly rescales those two points to produce e,p,q,s for this actual new input. Distinct detector certificates have independent structural coordinates and can be chosen separately after both inputs are fixed. Zero auxiliary witnesses remain allowed, and the finite conic/support charts include those branches. The later CRT or prime choices are not revisited.

[P20: final quantifiers and limitations] For every \(N\ge N_{C}\) and every integer \(x\), P14-P19 construct rational witnesses (existentially) in the original physical system, with the changed factors and nonzero inverse weights. Reconstruct deleted fourth powers and apply the exact raw224 coordinate maps and Laurent row identities; this gives a rational zero of the fixed 214 residuals. That argument allows the labels, selected primes, finite moduli, slope, final \(h\) and pole certificates to depend on the given (N,x); it imposes no further lower bound on \(N\). Taking \(N_{0}=\operatorname{max}(N_{S},N_{C})\) combines this with P11-P12 for all larger \(N\) and every rational input, proving \(P_{01}\) relative to the explicitly retained interfaces. This author did \textbf{not} produce a full numerical 214 witness tuple, even for a sample input; no sample could establish universal completeness. This packet gives no numerical defining threshold, no Busy Beaver coefficient replacement, no independent acceptance of itself. Full Protocol \(D\) closure additionally requires the fresh eligible review and exact assembly of unchanged RC/analytic and height foundations, not a green structural check.
\clearpage
% Source: research/quartic214/proofs/geometry-carry.md
\section{\(H_{01}:\) the actual \(D_{210}\) geometry and the small \(S_{10}\) application}\label{q214:geometry-carry}
\noindent\textit{Source: }\nolinkurl{research/quartic214/proofs/geometry-carry.md}.\par

Author supplement to \(\text{AUTHOR}_{{\text{FREEZE}}_{001}}\), \nolinkurl{/root/d214_modularity_closure}, 11 November 2026. Author proposal, pending nonauthor acceptance. This checks the actual earlier arguments and the changed application; it supplies no independent approval.

\subsection{\(G_{0}\). Locator correction and exact scope}

The suggested public SF08--10 locators concern modular curves, torsion differents and the modular-Jacobian estimate. They belong to \(H_{04}\), not \(H_{01}\). Public V15-HEIGHT-MODULAR accordingly supplies no \(D_{210}\) fine cover merely by accepting its modular curve. V15-FOUNDATION's third clause retains the original \textbf{six-contact} theorem; it is not by itself the five-contact \(D_{210}\) argument either. These are locator corrections, not new arithmetic objections.

There is a further distinction: H6-ORIGINAL-POLARIZED-QM-DATUM, H6-FINE-THREE-INTEGRAL-COVER and the \(B_{11}\) six-phase construction concern the discriminant-six surface, its auxiliary full-five model, and degree at most324. Reading their actual proofs exposed this difference. They must not be substituted for the discriminant210 \(H_{01}\) surface of degree576. A preliminary author message listing those \(H_{6}\) cards as if they were \(H_{01}\) providers was too broad. The actual providers are \nolinkurl{research/version14/proofs/h5-full-three-v1/PROOF.md}, its \nolinkurl{GEOMETRIC-ISOGENY.md} and parent \(D_{210}\) construction, together with the actual maximal-stack Yuan applications. Every assertion below concerns that \(D_{210}\) object.

The old private \(\text{FOUNDATION}_{\text{MAP}}\) F-H01-GEOMETRY already points to frozen \(V_{14}\) lines7605--8040 and8517--8655 and the cutoff proof. Its REPORT explicitly retains \(H_{01}\) source-closure limits; it does not identify the D6/B11 construction as proof of the \(D_{210}\) field. Thus the newly corrected locator suggestion is not an error in that old map. The old files remain untouched.

\subsection{\(G_{1}\). Exact text and retained public scopes}

\nolinkurl{GEOMETRY_TEXT_COMPARISON.json} binds three actual sources: private frozen \(V_{14}\), the private frozen \(V_{15}\) copy consumed elsewhere in the 214 packet, and current public \(V_{15}\) at accepted PR26 head1a49013b. The full \(D_{210}\) integral-geometry, fine-cover, field, Hom, and metric passage is \(V_{14}\) lines7605--8020, \(V_{15}\) lines7737--8152. Private and public \(V_{15}\) have \textbf{byte-identical} passages, raw SHA256 \nolinkurl{d8cea4a2552c8e5af575145b5db0e802bd6d96ff80da2e5d561a7e7da7c302d5}. The \(V_{14}\) passage differs only by LF/CRLF; its LF hash is \nolinkurl{c300ad111a16296fd05fa7c560a354222307124173de6790fa1c7390e1c7cac8}. No mathematical whitespace normalization was used.

The complete CM854R proof is \(V_{14}\) lines8519--8654 and \(V_{15}\) lines8618--8753. Private/public \(V_{15}\) are byte-identical, raw SHA256 \nolinkurl{42732844c948b234ba3820d2915b205aa0f909b1bdd0bd3d4fe8959fb0575412}; the LF content is the same as private \(V_{14}\). The CM supplier excerpt used by public inherited transport has SHA256 \nolinkurl{9abb3c65609362383a43ccbb463b2fc7dd4a6487c38be4c413f20bb1ed1fb02b}, equal to the private \(V_{14}\) raw excerpt when its terminal newline is included. This connects the actual proof, not merely its final inequality.

\nolinkurl{PUBLIC_GEOMETRY_SCOPES.json} extracts the precise final006 coverage decisions for ACTIVE-H19-D210-FIBER-FIELD, ACTIVE-H19-D210-GEOMETRIC-ISOGENIES, ACTIVE-H19-CM-ELLIPTIC-DISCRIMINANT, ACTIVE-H19-CM-HEIGHT, ACTIVE-H19-CM-GR2023 and CURRENT-YUAN-MODEL/INTEGRAL/METRIC. These decisions reuse exact retained derivations and explicitly preserve their prerequisites. They are evidence of public acceptance at the recorded scope, not a newly authored approval of the private \(S_{10}\) application.

The actual full \(D_{210}\) source proof was read in full. It constructs \(\text{Isom}_{O}(O/3O,A[3])\), proving the division-lattice assertion at ramified3. It distinguishes stack marking degree72, coarse degree36 and normalizer order576; proves trivial inertia on arbitrary bases; obtains projectivity by the component/genus/relative-ampleness sequence; and identifies the entire canonical Q-model, all multipliers and both half-planes. It checks Milne's full Hodge-orbit condition directly instead of using the corrected general8.17 implication. Its integral statement is only over Z[1/210]. The metric proof uses Yuan's maximal stack at bad primes, not a smooth fine model there or the disputed coarse Lemma2.1/Theorem2.2.

\nolinkurl{GEOMETRY_PROVIDER_COMPARISON.json} honestly records that these separately packaged supplier files were not found under the queried \nolinkurl{current/...} names in the old private handoff zip. No byte equality of an absent archive member is claimed. The \textbf{integrated proof text} equality above is established independently, and the actual publicly retained supplier is pinned and read. Its prior \(D_{210}\) index explicitly preserved five source-prerequisite bundles; public final006 acceptance preserves its exact source scope. This supplement does not infer resolution of any unexamined transitive prerequisite solely from that index or a PASS label.

\subsection{\(G_{2}\). Same cover, field, Hom and surface for every allowed a}

Let \(a=u/v\) be in lowest terms, \(v>0\), and a not in \{0,1,2401,-343/32\}. The omitted fifth contact is infinity. Choose the actual whole \(D_{210}\) full-three cover \(Y\) of the parent proof and its normalizer quotient line with precisely these five branch sections. The unbranched fiber at a is a torsor for the same constant group \(\Gamma \) of order576. Its commuting Galois action lies in the regular centralizer of \(\Gamma \), also of order576. The normal splitting field \(L\) therefore has degree \textbf{dividing}576. This argument depends on the same whole canonical cover and all-component rational action, not on the size of \(S\) or on a being integral.

All fiber3-torsion markings are defined over that \(L\). For any pair of the fiber surfaces, a homomorphism vanishing on3-torsion factors through[3]. Thus Hom/3Hom injects into \(\operatorname{Hom}(A[3],B[3])\). The finite Galois image on the integral Hom lattice is trivial modulo3, and the principal congruence subgroup mod3 is torsion-free; therefore the image is trivial. Every geometric Hom between actual conjugate surfaces is L-defined. The normalizer changes a commensurable integral lattice while retaining the rational Hodge object. Clearing denominators gives a geometric isogeny between the actual conjugate surfaces. The full Hodge-centralizer calculation gives \(\operatorname{End}^{0}(A)=B\) in the non-CM branch and the CM alternative otherwise. This is exactly the input of M0--M2, with no replacement surface or extra Hom field.

None of these statements uses \(B_{*}=\prod _{p\le 76176}p\), the old exponent19, the large-primorial inequality or the scalar character chosen later. Their complete proofs are the same in private frozen \(V_{15}\) and current public \(V_{15}\).

\subsection{\(G_{3}\). Fresh support extension to \(S_{10}\), with all contact cases}

Let \(S_{10}=\{2,3,5,7,11,13,17,37,41,2399\}\). Use the following primitive homogeneous forms in (u,v), in that exact order:

\begin{gather*}
v, u, u-v, u-2401v, 32u+343v
\end{gather*}

Their ten signed pair determinants are

\begin{gather*}
-1,-1,-1,-32,-1,-2401,343,-2400,375,77175
\end{gather*}

The nonunit factorizations are \(32=2^{5}\), \(2401=7^{4}\), \(343=7^{3}\), \(2400=2^{5}*3*5^{2}\), \(375=3*5^{3}\), \(77175=3^{2}*5^{2}*7^{3}\). Hence the \textbf{exact union} of their prime supports is \{2,3,5,7\}, contained in \(S_{10}\) and inverted in Z[1/210]. \nolinkurl{SMALL_SUPPORT_CERTIFICATE.json} gives every pair, coefficient vector and factorization. This finite calculation establishes distinct reductions and unit leading coefficients, not the arithmetic theorem by itself.

Define \(M\) to contain once every \(p\) outside \(S_{10}\) having positive odd valuation on any one of those five forms, and \(R=B_{S}\) \(M\), with \(B_{S}=\prod _{p \in  S_{10}}p\). For each \(p\) outside \(S_{10}\), primitiveness covers every case:

\begin{itemize}
\item If p|v, then \(p\) does not divide \(u\), and the other four forms are units because32 is a p-unit. The specialization is the infinity section. Its contact depth is \(v_{p}(v)\), since \(u\) is a unit.
\item If \(p\) does not divide \(v\), a is p-integral. The determinant calculation shows that at most one finite-contact form can have positive valuation. Its valuation is exactly the contact depth of a with that branch section: all \(v\) and32 denominators are p-units.
\item If no form has positive valuation, the point lies in the unbranched locus.
\end{itemize}

Thus if \(p\) does not divide \(R\), every contact depth that occurs is even, including infinity. The retained actual tame model has local equation \(w=z^{2}\) up to a unit at a branch section and is etale elsewhere. Over a finite unramified extension the unit has a square root (p is odd), and an even valuation has square uniformizer factor. Each such local fiber is therefore unramified. The entire normal fiber splitting field is consequently unramified: finite unramified local extensions and their normal closure remain unramified. This proves \(\operatorname{rad}(D_{L})|R\) for this \textbf{new} \(R\), with no cutoff76176 hypothesis.

The same projective fine cover over Z[1/210] and its universal abelian scheme give good reduction of each A/L at all places not over210, irrespective of whether L/Q itself ramifies there. Properness extends the L-point over the valuation ring; pullback of the universal abelian scheme supplies the reduction. In particular the support assertion needed by \(H_{0}\) is valid for this smaller \(S_{10}\). This does not infer the cover or its tame charts from the determinant table: both are the actual retained \(D_{210}\) constructions read above.

The changed application is precisely this support sharpening from the old \(R_{*}=B_{*}\) \(M_{5}(a)\) to \(R=B_{S}\) \(M_{S}(a)\). One cannot obtain it by substituting a smaller \(R\) into an inequality proved only with the old radical. The explicit local case argument supplies the needed extension. It applies to \textbf{all} rational a in the stated domain, including denominators, negative a, \(M=1\), \(p\) in \(S_{10}\) but larger than7, and arbitrarily deep contacts. Primes in \(S_{10}\) impose no asserted local unramified condition and are already in \(R\).

\subsection{\(G_{4}\). Metric and CM consumers of that exact field}

The metric comparison is independent of \(S\). On the same maximal \(D_{210}\) stack use \(t=32a\), branch values(0,32,76832,-343), sections \(q_{0}=(dt)^{2}/P(t),q_{1}=tq_{0}\), and \(J=2^{10}*3^{2}*5^{3}*7^{4}\). The actual proof extends \(J^{2}q_{i}\) vertically on the stack using semistability and controls their Petersson norms \(by4C_{\text{met}}\). Yuan's final isometry and the original stable normalization give

\begin{gather*}
h_{\omega ^{2}}(y)=2h_{F}(A_{y})+2\operatorname{log}(\pi )+\log 210
\end{gather*}

The positive log210 sign is retained. The same explicit integer estimate bounds the additive term \(by120<128\), so

\begin{gather*}
h(a)\le 2 \operatorname{max}(0,h_{F}(A_{a}))+128
\end{gather*}

The input point and universal \(A_{y}\) are the same as in \(G_{2}\). No coarse integral pullback or bad-prime smooth fine model enters this metric statement. The CURRENT-YUAN source scopes remain their exact maximal-stack/isometry applications; they do not by themselves prove the curve or the metric estimate. The actual integrated valuation/curvature/evaluation argument is what is carried here.

In the CM branch, retain that same A/L, all geometric endomorphisms over \(L\) and \([L:Q]\le 576\). \(G_{3}\) supplies \(\operatorname{rad}(D_{L})\le R\). The literal CM854R lemma assumes only these facts and \(R\ge 1\); it has no primorial or branch-coordinate hypothesis. Its dimension argument shows \(A\sim E^{2}\). L-defined idempotents give E/L and all its endomorphisms over \(L\). The derivative action embeds \(F=\operatorname{End}^{0}(E)\) into \(L\). Its same-field subgroup \(K=\sum (f \alpha _{i})(E[f])\) yields a quotient E'/L of maximal order, with \(A\sim (E')^{2}\) over \(L\). No coefficient-value field, \(L'\), or auxiliary good-reduction field replaces \(L\).

The discriminant tower gives \(|\Delta _{F}|\le 4\operatorname{rad}(D_{L})\le 4R\), including the dyadic multiplier4. The individual-embedding CM height formula bounds \(h_{F}(E')<3R\), hence \(h_{F}((E')^{2})<6R\). The same general dimension-two Gaudron--Rémond2023 Theorem1.9(1), over \(L\) and with no base extension, bounds an isogeny \((E')^{2}\to A\). The original stable half-log-degree inequality gives

\begin{gather*}
h_{F}(A)<6R+8(96+7+2+\operatorname{log} R)\le 854R
\end{gather*}

All objects and constants in that proof are identical to the actual public source. It uses the individual-embedding discriminant formula, not a class average requiring a \(Q(j)\) Weierstrass model; the temporary field evaluating stable height contributes no degree to the isogeny estimate. The final006 H19-CM records explicitly state this universally quantified \(R\ge 1\) interface. This is why the changed small-support application fits it once \(G_{3}\) is proved.

\subsection{\(G_{5}\). Dispositions and acceptance boundary}

Historical classification: prior substantive justification found and read for the actual \(D_{210}\) cover, same-L Hom, metric and CM854R arguments. The \(S_{10}\) cutoff/local-support bridge already appears in private cutoff \S{}1; this supplement verifies its full contact cases, determinant equality and exact primary-object joins. It is not claimed as a newly discovered missing arithmetic argument. Private heightH0 states these inputs but leaves their source qualifications inherited; incorporation of this additive comparison is pending.

Current author conclusion: H01's change of support introduces no additional mathematical hypothesis once the retained \(D_{210}\) source statements are supplied at their public accepted scopes. The source-to-private application and its consumer are now explicit. This is an author conclusion, awaiting the assigned nonauthor reviewer. Public final006's scoped approvals and byte identity do not alone certify every transitive primary-source prerequisite or the entire private 214 proof. Any such unresolved prerequisite must remain open in the root closure assessment; neither the mistaken SF08--10/D6 locators nor the ten-determinant check closes it.
\clearpage
% Source: research/quartic214/proofs/provider-scope-carry.md
\section{Exact retained \(D_{210}\) provider and conductor scope}\label{q214:provider-scope-carry}
\noindent\textit{Source: }\nolinkurl{research/quartic214/proofs/provider-scope-carry.md}.\par

Author supplement, \nolinkurl{/root/d214_modularity_closure}, 11 November 2026. Pending nonauthor review. Read together with \nolinkurl{GEOMETRY_CARRY.md} and \nolinkurl{PROOF_FINAL.md.} This supplies the detailed current-scope comparison requested after the first author freeze; it does not rewrite the older limited reviews.

\subsection{\(P_{1}\). Full \(D_{210}\) provider: sixteen actual applications}

The exact public final006 decisions are in \nolinkurl{PUBLIC_FULL_COVER_DECISIONS.json.} \nolinkurl{FULL_COVER_SOURCE_PROFILES.json} extracts the actual incoming assertions, hypotheses, source editions/hashes, constructed objects and connecting maps from \nolinkurl{protocol/current-coverage/INHERITED-TRANSPORT-PROPOSALS.json}. These are the sixteen \nolinkurl{ACTIVE-H19-FULL-COVER-*} applications, \textbf{not} a generic \(H_{19}\) verdict, D6/B11 proof or modular SF08--10 application. The actual entire \nolinkurl{h5-full-three-v1/PROOF.md} (SHA256054ec4a6...) and its \nolinkurl{GEOMETRIC-ISOGENY.md} (f889eab0...) were read. The composition below explains how the antecedents in each source profile are supplied by the previous construction.

\begin{longtable}{@{}p{0.23\linewidth}p{0.73\linewidth}@{}}
\toprule
Exact suffix of ACTIVE-H19-FULL-COVER-* & Supplier and actual construction-to-input match\\
\midrule
\endhead
MAXIMAL-STACK & Yuan final Chinese2024 \S{}2.1, pp230--231, aee74126...; use only \(N=1\) and the specified \(D_{210}\) division algebra/order/involution. It supplies the proper smooth DM maximal stack over Z[1/210] with universal surface. The separate Isom construction defines full3; no literal level-map statement is used at ramified3.\\
THREE-TORSION & Milne1986 Theorem8.2/Remark8.4, pp115--116, 86259344...; 3 is invertible on this base, A has dimension2, so A[3] is finite etale of rank81 and \(T_{3}\) has rank4. The projection T3/3 \(\to \) A[3] respects the O-action and multiplication-by3 transition maps.\\
DIVISION-LATTICE & Voight2021 \S{}13.3.4/13.3.11, pp206--207, exact final extract8f808986...; \(B_{3}\) is division, not split. Rank4 makes \(V_{3}\) rank1 over \(B_{3}\). A least \(w=\operatorname{ord}3(\operatorname{nrd})/2\) lattice element proves \(T_{3}=O_{3}\) \(x\). Thus A[3] is the regular O/3 module. Its unit \(\text{count}81-9=72\) is the stack-cover degree.\\
TATE-FAITHFULNESS & Milne1986 Lemma12.2, \(p_{122}\); inject the finite geometric stabilizer into GL4(Z3). Preservation of the full marking places it in1+3M4. The written binomial argument for prime-order elements proves torsion-freeness. It is not deduced merely from the Tate rank.\\
TRIVIAL-INERTIA & Stacks03YR, dated source82246ff7...; inertia of this proper DM cover is finite unramified. Its identity section is open and closed; its complement has empty geometric fibers. Hence inertia is trivial on every base and this same stack is an algebraic space with its universal scheme. This is stronger than geometric-point rigidity alone.\\
CURVE-FIBERS & Stacks0ADD, f067c166...; the smooth geometric fibers of that algebraic space have dimension1 and regular local rings, hence are schemes. The total space has not yet been assumed schematic.\\
GENUS-INPUT & Gonzalez--Rotger2006 pp937--941, e595d1a4..., for the selected \(D_{210}\) coarse curve. Every characteristic-zero marking-cover component maps finitely and nonconstantly to it. Its genus5 and Riemann--Hurwitz give \(\text{genus}\ge 5\) on each component. This is not inferred from degree72.\\
COMPONENT-BASE & Stacks0A18, e7df43c1...; the already proper smooth space is flat with geometrically reduced fibers. Its Stein component base is finite etale, with geometrically connected fibers. Every component of this finite-etale base meets the generic fiber.\\
GENUS-CONSTANCY & Stacks0E83, 3cac1b5d...; apply over that component base, using the schematic smooth proper connected curve fibers just proved. It carries \(\text{genus}\ge 2\) to every fiber component.\\
RELATIVE-AMPLENESS & Stacks0D3A,99f65961...; the relative canonical line has positive \(\text{degree}2g-2\) on each proper smooth fiber component and is ample there. Properness and a Noetherian base give relative ampleness locally near each base point.\\
SCHEME-FROM-AMPLE & Stacks0D32,a525102d...; that relative ample line makes the already constructed algebraic space a scheme. Trivial inertia or properness alone did not supply this.\\
PROJECTIVITY & Stacks0B45,d3057d39..., \(\text{implication}(6)\to (1)\); the base SpecZ[1/210] is affine and has an ample line, and \(Y\) is now a proper scheme with a relative ample line. Thus \(Y\) is projective.\\
PEL-CLASSIFICATION & Original MilneClay2005,6336f0a2..., Theorem7.4's full Hodge-orbit condition,8.6--8.7/8.16. Rank-one left \(B\) and the specified involution give one-dimensional compatible alternating forms and the right-centralizer \(\operatorname{GL}2\) type-C1. On the two-dimensional multiplicity factor, (v,Jv) is a basis and conjugates the full Hodge map into the same orbit. Both half-planes and all marking multipliers remain. The erroneous general8.17 implication is not used.\\
PRINCIPAL-POLARIZATION & Buzzard1997 p592,ed34a072..., before its level restrictions; recover the O-stable full Tate lattice, then its unique positive perfect polarization for the fixed involution. This changes no marking or unit multiplier. No auxiliary five-level \(D_{6}\) construction is imported.\\
ALGEBRAIC-FAMILY & Milne2005 Theorem14.8, \(p_{352}\); every arithmetic component has its integral lattice and polarization. A rational stabilizer preserving that lattice has positive multiplier a unit at every finite prime, hence1. Thus the variation is integrally polarized and the theorem supplies its abelian scheme/O maps algebraically. The marking's graph is a component union of a finite etale Isom cover, hence algebraic. The fine moduli map is bijective on geometric points; nonsingular characteristic-zero curve theory makes it an isomorphism.\\
CANONICAL-MODEL & Milne2005 pp354--355 type-A/C criterion and13.6--13.7,p349; the comparison transports the algebraic order, polarization, torsion quotients and marking under every automorphism of C/Q. Reflex field isQ. This is the required rational model, and uniqueness/Hecke-map compatibility identifies the same normalizer action. There is no fictitious Galois action on rational Betti homology or restriction to one component.\\
\bottomrule
\end{longtable}

For finite-etale Isom representability in this composition, trivialize the two finite etale group schemes over an etale cover. All group isomorphisms form a finite constant scheme, and the finitely many O-linearity identities select an open-and-closed subscheme. Faithfully flat descent glues these schemes and their tautological marking. Nonemptiness was independently supplied by DIVISION-LATTICE. Pullback of the existing universal abelian scheme carries all group maps identically. This establishes the specific construction needed here; it does not assert a universal family on an arbitrary coarse quotient.

The integral-to-rational lattice comparison also has actual maps: a full adelic lattice commensurable with the given one is realized after multiplying by an integer by quotienting by the corresponding finite torsion subgroup. O-stability descends the action. An isogeny preserving every integral Tate lattice has no nontrivial kernel/cokernel and is an isomorphism. The unique positive perfect representative restores the same polarization. This supplies the exact object classification consumed by ALGEBRAIC-FAMILY; it is not an equality of point counts alone.

\subsection{\(P_{2}\). The five old pending bundles, and what is retained}

The older H5-FULL-COVER-SOURCE-INDEX-REVIEW accepted a role index and explicitly left five bundles pending. That historical limited judgment is preserved. The later sixteen final006 source profiles above and the actual integrated proof support the following \textbf{new author comparison}, whose acceptance is still pending:

\begin{itemize}
\item 1. The fixed \(D_{210}\) datum and genus input are exactly MAXIMAL-STACK, PRINCIPAL-POLARIZATION and GENUS-INPUT, together with the selected order/involution in the actual \(D_{210}\) text. They are not supplied by the \(D_{6}\) original polarized datum.
\item 2. The finite-etale/Isom/stack-stability bundle is the explicit trivialization/equalities/descent construction above, with universal-scheme pullback and TRIVIAL-INERTIA. It applies to every geometric fiber over Z[1/210], not just characteristic zero.
\item 3. Riemann--Hurwitz and positive canonical degree are used in GENUS-INPUT and RELATIVE-AMPLENESS exactly at their proper smooth characteristic-zero/positive-genus curve scope; COMPONENT-BASE precedes GENUS-CONSTANCY. The current proof does not assume the total space was already a scheme.
\item 4. Finite torsion quotients, integral Tate detection and algebraic finite-etale markings are the actual constructions in \(P_{1}\) and \S{}6 of the full provider. The characteristic-zero nonsingular comparison is a morphism of curves with degree1, hence an isomorphism; no arbitrary complex isomorphism is descended.
\item 5. The central-idele and normalizer calculation is on that same entire canonical model: a rational scalar and its sign make any central idele trivial moduloU(3). The exact quotient \(\text{has}72/2=36\) marking classes and16 normalizer parities, hence576. Properness/good reduction is the valuative extension of points on that same projective \(Y\). Bad-prime metric statements use the separate Yuan maximal stack and its Hodge isometry, not a coarse-pullback lemma. These are respectively G2/G3/G4 in \(\text{GEOMETRY}_{\text{CARRY}}\).
\end{itemize}

Published canonical-model/PEL and CM reciprocity theorems retain their published statements as theorem inputs. The protocol does not require reproving their whole ancestry. The present author has not independently re-inspected every original final PDF/Stacks HTML byte in the sixteen old profiles; exact public accepted applications are carried at their existing source scope, while the actual entire internal \(D_{210}\) construction and its changed private application have been read and compared. This is the disclosed \S{}33 carry boundary for nonauthor review. If that exact carry is not accepted, those precise source applications remain open; a broad \(H_{19}\) or Yuan label cannot replace it.

\subsection{\(P_{3}\). Contact-line identity and tame branch composition}

The relevant actual predecessor passage, privateV14 lines7532--7604/currentV15 lines7664--7736, derives the \(D_{210}\) contact line from Gonzalez--Rotger's four published quartic quotients. It is distinct from both \(D_{6}\) and the modular j/1728 map. The quotient models have coefficient rows for \(m=30,42,70,105\) exactly as displayed there. Genus5, \(W=(C_{2})^{4}\) and Riemann--Hurwitz give the four known eight-point fixed sets, total fixed count40, genus-zero full quotient, and one further branch orbit. The four labels are independent and their product is210; therefore the five rational branch sections have labels30,42,70,105,210. Normalizing labels210,105,70 to infinity,0,1 fixes the coordinate.

\nolinkurl{CONTACT_LINE_CERTIFICATE.json} independently recomputes the four quartic I,J invariants and verifies their exact Legendre j-invariants for parameters1/2401,32/375,1/225,5/1029. It exhausts the36 pairs from the six cross-ratios of each first-two parameter. The last-two \(j\) equations leave uniquely \((\lambda ,\mu )=(2401,-343/32)\). This verifies the finite algebra consumed by the geometry proof; the published identification of the four quotient curves remains its identified published input.

Over Z[1/210], the level-change group36 and \(W\) of order16 have invertible orders. The generic level action is free. For a tame finite-order automorphism of a smooth curve, the fixed scheme is smooth along a fixed point (average a local tangent parameter to its eigencharacter); an isolated special fixed point lifts uniquely. Thus it cannot create a new special branch orbit or collision; intersecting distinct fixed loci would give a noncyclic stabilizer, impossible in a one-dimensional tame tangent action. The stable-curve Isom extension gives the same \(W\) action on the proper smooth model. The quotient remains genus zero with the same ordered three sections, hence that same normalized rational coordinate, and local anti-invariant parameter gives \(w=z^{2}\) up to a unit. This is the actual tame-model join used in \(G_{3}\), independent of the old primorial. It is not inferred from the determinant certificate.

\subsection{\(P_{4}\). Full11 semistability, torus rank and exact retained ramification source}

\nolinkurl{CONDUCTOR_INHERITED_PROFILES.json} records the actual public full11 and torus statements. The final Silverberg--Zarhin1995 AIF45,403--420 PDF was reacquired from Numdam; its hash is \textbf{\nolinkurl{fc62ba1c9d1315a0c441a3dd7de46a254f379289e630c407a3336c0230d91052}}, exactly the public source pin. Its actual Theorem3.4 and Theorem3.5, printed406, were read. The former gives a finite Galois semistable extension at every discrete valuation. The latter applies when n-torsion is rational over an unramified extension and \(n\ge 3\) is prime to residue characteristic. In our application \(n=11\), the ground field is \(E=L(A[11])\), and that unramified extension is the identity extension of \(E\). Thus every place over2,3,5,7 is covered; residue11 is excluded. No earlier \(n=3\) or5 occurrence is used as a substitute.

The same \(D_{210}\) B-action is defined over \(E\). Every actual endomorphism extends to the Neron model and preserves the semistable torus. If its geometric torus rank \(t>0\), restriction gives a unital homomorphism \(B\to M_{t}(Q),1\le t\le 2\). Its kernel is an ideal in simpleB and is not allB because1 maps to1. For \(t=1\) dimensions contradict injectivity; for \(t=2\) injectivity makes an isomorphism onto splitM2(Q), contradicting divisionB. Therefore \(t=0\) and semistable reduction is good. This proof is independent of chi,b and the size of \(R\).

E/Q is normal by the actual prime-to11 normalizer isogenies over \(L:\) at11 the algebra is split and a normalizer is Q11*GL2(Z11); a rational rescaling makes it integral invertible there, and denominators elsewhere are cleared by an integer prime to11. The original isogeny and every conjugate therefore identify the11-torsion kernels over \(L\). Every conjugate of \(E\) is \(E\). Lagrange strengthens the public inequality to divisibility

\begin{gather*}
[E:Q] | 576 |\operatorname{GL}4(F_{11})|
\end{gather*}

The public accepted source scope needed only the displayed larger \(\text{bound}<2^{66}\). The private sharper valuations(15,4,4,1) follow from this divisibility, not by substituting a new unproved finite-field estimate. Chi's new minimal-order choice does not change \(E\), A, \(B\) or these isogenies.

The retained ramification theorem is \textbf{Serre's1967 chapterVI, Local class field theory, in Cassels--Frohlich Algebraic Number Theory, \S{}\S{}4.1--4.3 pp155--159}, not Serre's1979 Local Fields. The public exact extract SHA is \nolinkurl{e0367b25a022eb27ec220e3c15b474a94fb96d70a6bb89b77f4a02c2231950fe;} its original retrieval URL was already recorded unknown. No new local copy or fresh final-page inspection is claimed. The exact source profile ACTIVE-H19-RAMIFICATION-NUMBERING and its final006 accepted decision are supplied in the carry records. Its content is general quotient compatibility of upper numbering, Hilbert's different formula and the full Artin tame/Swan conductor formula for finite local Galois extensions. It has no lower bound on degree or fixed list of ramification indices.

The new input to that unchanged theorem is the same kind of local finite Galois extension \(F'/Q_{p}\) with better wild-part valuations. \(Q(\chi )\) is the kernel field; it is tame at oddp and its dyadic inertia has \(\text{order}\le 2\). On \(F=E\) \(Q(\chi )\), \(b\) is a quadratic conjugation-invariant character and its kernel gives F'/F of \(\text{degree}\le 2\). At oddp this addition is tame, at2 it adds at most one factor2. Thus the actual constructed extensions satisfy all source hypotheses with \(v_{p}(e)\le (17,4,4,1)\). A/E is already good at2,3,5,7, and on \(G_{F'}\) all scalar/isogeny corrections are1, so the actual \(r_{\ell }\) is trivial on its local inertia for \(\ell \ne p\). The upper-numbering quotient map is therefore the one used to bound \textbf{this} representation's Swan conductor. The derivative/Hilbert/Herbrand calculation is explicitly derived in \(M_{5}\), retaining the tame term and allowing arbitrarily large tame index. The resulting conductor exponents(70,18,18,6) are a fresh numerical consequence; the \(\text{public}274<300\) conclusion itself is not carried as proof of them.

\nolinkurl{PUBLIC_CONDUCTOR_SCOPES.json} and \nolinkurl{PUBLIC_FULL_COVER_DECISIONS.json} point to the nonauthor public decisions for these exact inputs. They do not authorize an unreviewed changed application. This supplement, M1/M5's new lift and sharper arithmetic, and the exact same-form r3/r5 consumer remain subject to the assigned current independent review.
\clearpage
% Source: research/quartic214/proofs/modularity.md
\section{Original normalized modularity for the private 214 construction}\label{q214:modularity}
\noindent\textit{Source: }\nolinkurl{research/quartic214/proofs/modularity.md}.\par

\begin{quote}\footnotesize Final author revision after the first freeze. Author: \nolinkurl{/root/d214_modularity_closure}, internal Sol author lane, 11 November 2026. \textbf{Author proposal; no independent acceptance.} The frozen private repository is \nolinkurl{bfeb8400982c9a169e4fdd39fb68ec053621e8ed}. Sources and consumers are byte-bound in \nolinkurl{SOURCE_USES.json.} \nolinkurl{INITIAL_AUTHOR_ANALYSIS.json} records the actual order of exposure, including the scope summary read before reconstruction. This is an additive proof: no historical argument, verdict, count or polynomial is changed.\end{quote}

\subsection{\(M_{0}\). Exact input and conclusion}

For every rational \(a=u/v\) in the domain of private \nolinkurl{height/PROOF.md} \(H_{0}\), retain its original full-three-marked normalizer surface \(A=A_{a}\), normal field \(L\), and B-linear conjugate isogenies \(\mu _{\sigma }\). In the non-CM case the geometric endomorphism algebra is the division quaternion algebra \(B\) of discriminant 210. All conjugate Hom groups are L-defined, \(\mu _{1}=1\), and the choices are inflated from Gal(L/Q). These are geometric input assertions, not consequences of the modularity theorem. Put \(T=\operatorname{supp}(R)\), \(R=B_{S}\) \(M_{S}(a)\), and \(U=\operatorname{max}(\{2\} \cup  \{2^{v_{2}(p-1)}: \text{odd} p|R\})\).

The claims established below, conditional only on those exact geometric inputs and the source applications individually identified here, are: the new minimal-character lift gives genuine original representations \(r_{\ell }\); the same final normalized primitive weight-two \(f\) realizes \(r_{5}\) and \(r_{3}\); its actual level satisfies \(N\le 2^{68}\) \(3^{16}\) \(5^{16}\) \(7^{4}\) \(R^{2}\) and \(\operatorname{rad}(N)|R\); and the multiplicity and character estimates of H1--H2 hold over the actual \(L'\). The exact \(H_{01}\) geometric/metric/CM and retained conductor-source transport is supplied by \nolinkurl{GEOMETRY_CARRY.md} and \nolinkurl{PROVIDER_SCOPE_CARRY.md.} These additive comparisons require the same current nonauthor review; H3's Rémond inequality and the full214 assembly are separate.

Use covariant rational Tate modules, arithmetic Frobenius, and \(HT(Q_{\ell }(1))=-1\). Added algebraic coefficients have trivial Galois action. The names \(\chi \), \(b\), \(\epsilon _{\text{sign}}\), \(\epsilon _{f}\), \(\delta \) and the cyclotomic character \(\kappa \) are distinct.

\subsection{\(M_{1}\). The changed two-primary lift}

Let \(c(\sigma ,\tau )=\mu _{\sigma }\) \(\sigma (\mu _{\tau })\) \(\mu _{\sigma  \tau }^{-1}\) in Q*. Rational quasi-isogeny degree is defined by \(D(\mu )=\operatorname{deg}([n]\mu )/n^{4}\); multiplicativity, conjugation invariance and \(\operatorname{deg}([n])=n^{4}\) make this independent of \(n\). Thus \(D(c \mu )=c^{4}\) \(D(\mu )\). For \(a_{\sigma }=D(\mu _{\sigma })^{-1/4}\) with positive real fourth root,

\begin{gather*}
a_{\sigma } a_{\tau }/a_{\sigma  \tau }=1/|c(\sigma ,\tau )|
\end{gather*}

Consequently \(\nu _{\sigma }=a_{\sigma }\) \(\mu _{\sigma }\) has normalized sign discrepancy \(\epsilon _{\text{sign}}=c/|c|\), inflated from Gal(L/Q). This is the actual normalization, with no assertion that \(a_{\sigma }\) is a root of unity.

At an odd prime \(p\) in the prescribed nontrivial boundary set, write \(p-1=2^{v}\) \(d\), \(d\) odd, and choose a generator \(g\) modulo \(p\). The character \(g\) \(\to \) \(\zeta _{2^{v}}\) takes \(-1=g^{(p-1)/2}\) to -1. At all other odd primes choose 1; at 2 use the character modulo 4 only when required. Their product \(\chi \) has 2-power order \(m\) dividing \(U\), is tame at odd primes, and has dyadic conductor at most 4. This includes \(m=1\) and the empty boundary set.

Here is the complete local/global obstruction check. Use the constant-module squaring sequence 1 \(\to \) \(\{\pm 1\}\) \(\to \) \(\mu _{\infty }\) \(\to \) \(\mu _{\infty }\) \(\to \) 1. By local reciprocity, a finite character of \(G_{Q_{p}}\) has a character square root exactly when its character on \(Q_{p}^{*}\) does. The torsion factor of \(Q_{p}^{*}\) is cyclic of even order; its only obstruction is the value at -1. The valuation \(Z\) factor admits a finite-order square root. For odd \(p\), squaring on the pro-p principal units is an automorphism; for \(p=2\) the torsion-free factor \(1+4Z_{2}\) is procyclic and its generator's image may be given a square root of twice the order. Thus the criterion holds also at 2, with no omitted principal-unit obstruction. Local \(H^{2}(G_{Q_{p}},\{\pm 1\})\) has two elements, by Kummer and the local Brauer invariant. The boundary of \(\chi \) equals the sign class at every finite prime. The latter is zero outside \(T\) because its class inflates from the unramified local quotient, which has cohomological dimension one. The global Brauer injection and invariant-sum exact sequence determine the real invariant as well. See the final Milne ADT second edition, I.1.3/I.1.6(b,c), and the proof of II.2.1, printed 164, for the reciprocity and Brauer statements.

A continuous pointwise square root \(b_{0}\) of \(\chi \) has values in \(\mu _{2m}\). Equality of cohomology classes means \(\delta (b_{0})/\epsilon _{\text{sign}}\) is the coboundary of a continuous sign cochain; correcting \(b_{0}\) supplies \(b^{2}=\chi \) and \(\delta \) \(b=\epsilon _{\text{sign}}\). Normalize \(b(1)=1\). On \(G_{L}\), the cocycle is 1, so \(b\) is a character. For \(h\) in \(G_{L}\) and \(g\) in \(G_{Q}\), the identities \(\epsilon _{\text{sign}}(g,h)=\epsilon _{\text{sign}}(h,g)=1\) show \(b(gh)=b(g)b(h)=b(h)b(g)\), and then \(b(g h g^{-1})=b(h)\). Its kernel in \(G_{L}\) is therefore normal in \(G_{Q}\).

Outside \(T\), \(b|I_{p}\) is quadratic: \(L\) is unramified there, so \(I_{p}\) is contained in \(G_{L}\). Continuity of the finite cochain implies finitely many nontrivial inertia restrictions. Twist \(b\) by the quadratic character belonging to the product of the unwanted odd primes. Its local inertia character is the unique nontrivial quadratic one at each such odd prime; it introduces ramification only there and possibly at 2, already in \(T\). This preserves \(b^{2}\) and \(\delta \) \(b\). Whenever odd wild inertia lies in \(G_{L}\) \((\in  \text{particular} at p>7)\), \(b\) on that wild inertia is a character with square 1, hence is trivial because a pro-p group has no nontrivial image of order 2. No character assertion for \(b\) on all wild inertia at \(p=3\) is needed: \(M_{5}\) instead works over \(F\) and its quadratic extension \(F'\).

Let \(L'\) be the fixed field of \(\operatorname{ker}(b|G_{L})\). It is normal over \(Q\) and

\begin{gather*}
[L':Q] \le  2m[L:Q] \le  1152 U \le  2304 \phi (R)
\end{gather*}

This is the original controlled field in the private proof. Neither a coefficient-value field nor a later twisting field is substituted for it.

\subsection{\(M_{2}\). Modularity hypotheses for the actual original \(r_{5}\)}

Set \(t_{\sigma }=a_{\sigma }/b(\sigma )\) in one finite number field \(C\) containing these finite sets of values and splitting \(B\). The sign identity gives \(t_{\sigma }\) \(t_{\tau }\) \(c(\sigma ,\tau )=t_{\sigma  \tau }\), so \(\rho _{\ell }(\sigma )=t_{\sigma }\) \(\mu _{\sigma }\) \(\sigma \) is a genuine representation. Its B-commutation gives \(\rho _{\ell }=r_{\ell }\) direct-sum \(r_{\ell }\). On \(G_{L'}\), \(\mu =a=b=1\), hence \(\rho _{\ell }\) is the actual \(V_{\ell }(A)\). This identity also checks the actual object used in the later Faltings and conductor arguments.

Faltings' final Theorem 1(a,b), published \emph{Rational Points}, second edition (1986), printed 204, applies over \(L'\) to A and to products. A number field is finitely generated in characteristic zero. Its Tate representation is semisimple and its commutant is \(\operatorname{End}_{L'}(A)\) tensor \(Q_{\ell }\). All geometric endomorphisms are L-defined, so after extending to a splitting coefficient field this is \(M_{2}(E)\). The multiplicity commutant is \(E\), also after scalar extension; intertwiner equations are finite-dimensional linear equations and commute with extension. Semisimplicity persists under scalar extension in characteristic zero because the semisimple image algebra is separable. Thus \(r_{\ell }|G_{L'}\) and \(r_{\ell }\) are absolutely irreducible. Finite-index restriction and averaging also show semisimplicity upstairs. Continuity follows from continuity of the Tate action on this open subgroup and the finite coset decomposition. Away from \(T\) union \{ell\}, good reduction and the unramified field/cochain give unramified \(r_{\ell }\).

For \(\ell =5\), take a completion of \(L'\) above 5. It is a finite extension of \(Q_{5}\), hence complete discretely valued with perfect residue field. A is smooth proper over that field, so its analytification satisfies Scholze's final Corollary 1.8, pp.5--6. Theorems 9.1 and 9.3, pp.74--75, identify coherent/de Rham and etale algebraic/analytic realizations; the finite etale comparisons respect every \(Z/5^{n}\) transition and therefore their inverse limit. The final 2016 corrigendum, all four pages read, restricts the pro-etale covers and corrects the completed period-sheaf definition; it explicitly retains the comparison results. It does not impose good reduction at 5. Milne1986 Theorem15.1(a), p.129, identifies etale \(H^{1}\) with the Tate dual without another Tate twist. The surface's Hodge numbers 2,2 give Tate weights \{0,0,-1,-1\}; the two identical summands each have weights \{0,-1\}. Fontaine1994 Theorem3.8, p.142, supplies exact tensor/dual/subobject closure. Kato2004 \S{}9.2, p.187, explicitly supplies descent through finite local extensions and filtered base change. Hence this original \(r_{5}\) is de Rham over \(Q_{5}\) with those weights.

For oddness, on complex Betti homology conjugation interchanges the two Hodge summands, whereas \(\mu _{c}\) is holomorphic and preserves their types. Thus \(t_{c}\) \(\mu _{c}\) \(c\) has trace zero. Milne2005's complex uniformization and Theorem6.8/Corollary6.9, p.319, identify the lattice with \(H_{1}\). Multiplication by \(5^{n}\) identifies \((5^{-n} \text{Lambda})/\text{Lambda}\) with \(\text{Lambda}/5^{n}\) Lambda. Exponentiation commutes with holomorphic homomorphisms and conjugation to the conjugate variety. Clearing isogeny denominators and taking the inverse limit transports this operator to \(\rho _{5}(c)\). Thus \(\operatorname{tr}(r_{5}(c))=0\). Since \(r_{5}(c)^{2}=1\), it has eigenvalues 1,-1 and is odd.

Calegari's final EMS ICM2022 volume2 chapter, published 2023, definitions pp.611--612 and statement p.624, is the immediate published theorem: continuous irreducible geometric two-dimensional representations at \(p\ge 5\) with distinct weights and oddness are modular up to character twist. Every condition has now been checked at \(p=5\). No residual irreducibility assumption, local ordinary assumption, or missing \(p=3\) modularity theorem is being inserted. The historical Pan \(v_{2}\) Theorem1.0.4, p.5, was also read; its exceptional residual condition occurs only at 3. Final Pan JAMS body retrieval returned 403. We do not claim to have compared that body, or to have recursively re-audited its proof. The immediate final published Calegari statement is independently available and directly applicable under the protocol's published-theorem boundary.

\subsection{\(M_{3}\). Weight, finite twist and final primitive form}

The theorem gives \(r_{5}=U_{g}\) tensor \(\delta \) after finite coefficient enlargement, with a classical nonzero cuspidal eigenform \(g\) of positive integral weight \(k\). Select a normalized primitive constituent before determining \(k:\) Deligne--Serre1974 \S{}1.9, p.510, decomposes \(g\) into primitive oldforms; \S{}1.6, p.509, shows good-prime \(T_{q}\) commutes with each degeneracy map since its index is prime to \(q\). Uniqueness of the decomposition forces every nonzero constituent to have g's good-prime eigenvalues. \S{}1.7 normalizes \(a_{1}=1\). Its realization \(U_{f}\) therefore equals \(U_{g}\) by \S{}3.2, p.513, applied after semisimplification.

If \(k=1\), Deligne--Serre Theorem4.1, pp.513--514, gives a finite-image irreducible representation; \S{}3.1 uses the discrete complex topology, and \S{}\S{}8.2,8.5--8.6 realize it over a number field. It has a finite-image 5-adic realization with these same polynomials. After a finite local extension kills that image, \(r_{5}\) becomes two copies of \(\delta \). Kato \S{}9.2 preserves de Rhamness and weights under that extension; subobject closure gives two equal weights, contradicting \{0,-1\}. Thus \(k\ge 2\).

Use Kato's exact coefficient realization \(V=V_{E_{f},\lambda }(f)\), \S{}\S{}6.3 and9.4, pp.161--162,188. Its two-dimensional de Rham filtration has jumps 0,k-1; its middle step has dimension one. Small original levels are handled by \S{}\S{}5.4/8.3 through fine-level invariants; \(p\) may divide \(N\). Equations14.10.2--3, p.240, put the polynomial \(1-a_{q}\) \(X+\epsilon _{f}(q)q^{k-1}X^{2}\) on inverse arithmetic Frobenius in \(V\). Thus \(U_{f}=V^{\text{dual}}\), det \(U_{f}=\epsilon _{f}\) \(\kappa ^{k-1}\), and \(HT(U_{f})=\{0,1-k\}\). Equality follows first on semisimplifications; since \(U_{f}\) is an irreducible twist of \(r_{5}\), the matched two-dimensional object has composition length one and is actually isomorphic. There is no further Tate twist.

The delta-line in \(\operatorname{Hom}(U_{f},r_{5})\) is a de Rham subobject. If its weight is \(h\), \(\{h,h+1-k\}=\{0,-1\}\); \(k\ge 2\) gives \(h=0,k=2\). Put \(\alpha =\operatorname{det}(r_{5})/\kappa \). On \(G_{L'}\), the polarization determinant gives \(\alpha ^{2}=1\), hence \(\alpha \) has finite global image. Taking determinants gives \(\delta ^{2}=\alpha /\epsilon _{f}\), which is finite. Each nonzero scalar has at most two square roots; \(\delta \) is finite. Weight zero alone was not used to infer this.

To obtain a Dirichlet twist one may use the inherited Washington Theorem14.1; its final page was not reacquired here. The following replacement uses the final published Milne ADT already inspected: I.1.3 and1.6(c), pp.18--19, give global reciprocity for the idele class group, compatible with the local maps. Composing the finite character \(\delta \) with reciprocity gives a finite continuous idele-class character. The connected positive real component acts trivially. Every rational idele class has a representative whose finite components lie in \(\prod  Z_{p}^{*}\), unique after choosing positive real component; this follows by multiplying by the rational product of its finite uniformizer powers. A finite continuous character of that compact product factors through \((Z/rZ)^{*}\) for some \(r\). Take \(r\) to be the least such modulus; the Dirichlet character \(\eta \) defined below is then primitive of conductor \(r\). This choice changes neither \(\delta \) nor the controlled field \(L'\). Write \(\lambda \) for this idele-class character. In the decomposition \(C_{Q}=R_{>0}\) times \(\prod  Z_{p}^{*}\), the idele which is a uniformizer \(q\) at \(q\) and1 elsewhere is represented, after multiplication by \(q^{-1}\), by positive real component1/q and finite units \(q^{-1}\) at every \(p\ne q\) and1 at \(q\). Therefore, with arithmetic reciprocity, define \(\eta \) on \((Z/rZ)^{*}\) to be the inverse of \(\lambda \) restricted to the finite units; then \(\eta (q)=\delta (\text{Frob}_{q})\) for \(q\) prime to \(r\). If the published reciprocity convention is inverse arithmetic, invert that reciprocity map first. This retains the necessary inverse on the finite-unit character rather than silently reversing the twist. This derivation uses neither Washington's unreacquired page nor an unpublished theorem.

Shimura1973 \S{}4 equation4.1, p.531, sends \(f\) to a nonzero Gauss sum times \(h=\sum \) \(\eta (n)a_{n}(f)q^{n}\) at level \(Mr^{2}\). For \(\gamma =(A B;C D)\) in \(\Gamma _{0}(Mr^{2})\), the translation indices reindex by a \(\to \) \(aD^{2}\) mod \(r\); their weights transform by \(\eta (D)^{2}\), so \(h\) has nebentypus \(\epsilon _{f}\) \(\eta ^{2}\). The coefficient formula gives \(T_{q}\) \(h=\eta (q)a_{q}(f)h\) for \(q\) not dividing Mr, and \(h_{1}=1\). Select a primitive constituent using the same Deligne--Serre \S{}\S{}1.6--1.9 decomposition just proved. Its good-prime traces identify its Tate realization with \(U_{f}\) tensor \(\delta \), hence with the original \(r_{5}\). Call this \textbf{final} primitive form \(f\) and its actual level \(N\).

Precision repair: Li1975 Theorem5, p.298, proves uniqueness of normalized primitive forms from almost-all prime eigenvalues. It does not itself assert existence of a primitive constituent. The current \(V_{15}\) p294/Theorem5 citation should not bear that existence claim alone. The explicit Deligne--Serre extraction above repairs the application without changing its arithmetic conclusion. Li's actual Fricke statement is used separately below.

\subsection{\(M_{4}\). Original companion \(r_{3}\) and full prime coverage}

Fix a common coefficient number field containing the original \(a_{\sigma },b-\text{values}\), the \(B\) splitting field and final f's coefficients. Extend the original embeddings at 5 and3; do not alter either \(r_{\ell }\). For the comparison alone, spread the finitely many varieties and quasi-isogenies over a finite Galois auxiliary field \(K\) and clear denominators by an integer \(n\). Remove finitely many rational primes so all reductions and spread maps are smooth proper abelian schemes and \(n\) \(\mu _{\sigma }\) are homomorphisms; remove primes ramified in \(K\) as well. This comparison set is not \(T\) and its size enters no estimate.

At a remaining \(p\), choose v|p and an arithmetic Frobenius \(\sigma _{p}\). Reduction of \(\sigma _{p}\) A is the p-twist of \(A_{v}\). The source/target of \(\mu _{\sigma }\) is \(\sigma \) A \(\to \) A. Therefore

\begin{gather*}
e_{p} = \operatorname{reduction} (n \mu _{{\sigma }_{p}}) \circ  \text{ with }  \text{ relative }  Fr_{p,A_{v}}
\end{gather*}

is an actual endomorphism of \(A_{v}\). It is relative p-Frobenius, not the residue-cardinality Frobenius. Milne1986 Proposition12.9 and its definition, p.125, give a common rational degree polynomial on every \(V_{\ell }\) for \(\ell \ne p\). Put \(q_{p}=\operatorname{Tr}(e_{p})/n\) in \(Q\); retain division by \(n\). Finite etale \(\ell ^{j}\) torsion specializes compatibly with multiplication-by-ell transitions and \(n\) \(\mu _{\sigma }\), so

\begin{gather*}
\operatorname{tr} r_{\ell }(\sigma _{p}) = \iota _{\ell }((a_{{\sigma }_{p}}/b(\sigma _{p}))*q_{p}/2)
\end{gather*}

The factor 1/2 is from the two identical summands, not an omitted scalar. The algebraic number \(\tau _{p}\) in parentheses lies in the common coefficient field. The isomorphism at5 gives \(\iota _{5}(\tau _{p})=\iota _{5}(a_{p}(f))\), hence equality in that field. Applying \(\iota _{3}\) gives equality of traces with \(U_{f,3}\) outside finitely many primes. Serre's published 1968 text/1989 reprint, ChapterI \S{}2.2 Corollary2 and \S{}2.3 theorem and lemma, pp.I-8,I-10--11, gives density and semisimple trace uniqueness. \(r_{3}\) is absolutely irreducible by \(M_{2}\); thus \(U_{f,3}^{ss}=r_{3}\) has a single simple constituent, forcing the full \(U_{f,3}\) to be irreducible and isomorphic to \(r_{3}\). This is an isomorphism on \(G_{Q}\), including its restriction at 5. It is not only an equality of local semisimplified traces.

The same argument works for every coefficient prime, but only the pair 5 and3 is needed: at residue \(p\ne 5\) use \(r_{5}\), and at \(p=5\) use \(r_{3}\). This explicitly covers the full residue-prime index set without a \(p=\ell \) compatibility assertion.

\subsection{\(M_{5}\). Final level and changed-lift conductor bound}

Use Carayol1987 final Astérisque147--148, \S{}\S{}2.1/2.4/4.1, printed36,37,42. Its normalized new cuspidal form permits arbitrary Dirichlet nebentypus and \(\text{weight}\ge 2\). The full Theorem4 removes the unramified-principal-series condition of Theorem3; the following deduction affirms the conductor question of \S{}2.4. The separately stated weak theorem's extraordinary-cuspidal exclusion is not a restriction on Theorem4. \(Q\) has odd degree1. Thus all dyadic cases and all \(N\) are included. The cohomological dual and unramified norm normalization change no Artin conductor. M4's actual global isomorphism identifies the full local object, preserving inertia invariants and Swan terms; no inertia semisimplification drops monodromy. Consequently

\begin{gather*}
v_{p}(N)=a_{p}(r_{\ell }), \text{ for }  \ell \ne p,
\end{gather*}

with the pair \((\ell =5 \text{ for }  p\ne 5; \ell =3 \text{ for }  p=5)\). This replaces the old private parent9 Casselman/Carayol-TheoremA chain at this exact conclusion. It does not purport to accept that separately formulated automorphic chain.

For completeness the changed lift preserves every numerical hypothesis of parent9 \S{}\S{}1--3. \(E=L(A[11])\) is normal and [E:Q] divides \(576|\operatorname{GL}_{4}(F_{11})|\) by the retained prime-to11 normalizer isogenies. The order is 23842542097059840000 with valuations at (2,3,5,7) equal to (15,4,4,1). The full11-torsion semistability criterion is used only in those four residue characteristics, and the division-B torus argument makes reduction good. Put \(F=E\) \(Q(\chi )\), where \(Q(\chi )\) denotes its \textbf{kernel field}, and let F'/F kill the quadratic character \(b|G_{F}\). Both fields are normal. At odd \(p\) the added character and quadratic extension are tame; at2 the \(\chi \) inertia has \(\text{order}\le 2\) and the second extension \(\text{degree}\le 2\). Hence \(v_{p}(e(F'/Q))\le (17,4,4,1)\). This is a wild-part bound, with no upper bound on the tame part.

Here is the original derivative argument, independently restated. For the totally ramified part \(K/K_{0}\) of a finite Galois \(Q_{p}-\text{extension}\), an Eisenstein uniformizer polynomial of degree \(e\) has derivative terms of distinct valuations modulo \(e\). Its different exponent satisfies \(d_{K}\le e\) \(v_{p}(e)+e-1\). Hilbert's derivative product gives \(d_{K}=\sum _{i\ge 0}(|G_{i}|-1)\). The Herbrand change of variable therefore yields

\begin{gather*}
d_{K}/e = 1-1/e + \int _{0,\infty }(1-1/|G^{u}|)du
\end{gather*}

Before the last positive upper break \(b_{\text{last}}\) the integrand is at least1/2, so \(b_{\text{last}}\le 2\) \(v_{p}(e)\). Quotient compatibility of upper numbering then bounds the Swan conductor of the actual two-dimensional \(r_{\ell }\) by 2 \(b_{\text{last}}\), with tame \(\text{codimension}\le 2\). Thus exponents are \(\le (70,18,18,6)\). At \(p>7\), L/Q is tame since its degree divides576, A/L has good reduction, and \(b\) is trivial on wild inertia by \(M_{1}\). Thus \(a_{p}\le 2\) there and \(a_{p}=0\) outside \(T\). No11-torsion semistability argument is used at residue11. \nolinkurl{PROVIDER_SCOPE_CARRY.md} \(P_{4}\) supplies the exact current-source carry: the final Silverberg--Zarhin1995 Theorems3.4/3.5 have been directly read at the same public PDF hash; the exact previously accepted Serre1967 chapterVI ramification application is retained at its general finite-local-Galois-extension scope. It is not the1979 Local Fields book. The changed wild valuations and maps to the actual \(r_{\ell }\) are proved anew here.

Absorb baseline exponent2 at 2,3,5,7 into squarefree \(R^{2}\) to obtain

\begin{gather*}
N \le  C_{N} R^{2}, \operatorname{rad}(N)|R, C_{N}=2^{68} 3^{16} 5^{16} 7^{4}
\end{gather*}

\subsection{\(M_{6}\). Determinant, multiplicity and the exact modular quotient}

The top Tate exterior power of a surface has cyclotomic action \(\kappa ^{2}\). An isogeny acts there by its degree. Since \(a_{\sigma }^{4}\) \(\operatorname{deg}(\mu _{\sigma })=1\),

\begin{gather*}
\operatorname{det} \rho _{\ell } = b^{-4} \kappa ^{2} = \chi ^{-2} \kappa ^{2}
\end{gather*}

The final modular identity gives det \(\rho _{5}=\epsilon _{f}^{2}\) \(\kappa ^{2}\). Thus \((\epsilon _{f} \chi )^{2}=1\). Because \(U\) is an even power of2 and m|U, \(s=\operatorname{ord}(\epsilon _{f})\) divides \(U\). The weight-two -I condition gives \(\epsilon _{f}(-1)=1\). Also U|phi(R) and

\begin{gather*}
s \prod _{p|R}(1+1/p) \le  R \prod _{p|R}(1-p^{-2}) < R
\end{gather*}

The statement U|phi(R) includes \(R\) with no new prime outside \(S\). No uniform bound on \(U\) or \([E_{f}:Q]\) is asserted.

Let \(E_{f}=Q(a_{n}(f))\), \(e=[E_{f}:Q]\). Li's actual final \(p_{296}\) formula \(f|H_{N}=\gamma \) f|K, \(\gamma \ne 0\), shows that the span of \textbf{all} coefficient conjugates is Fricke-stable: for each embedding \(\tau \), conjugating that embedding produces another of the same field. Apply it at Li's \((M,1)=(N,1)\), with arbitrary \(\epsilon _{f}\), to the final form after absorption. Shimura1973 Theorem1, pp.526--527, then gives \(A_{f}/Q\), dimension \(e\), a quotient of \(J_{1}(N)\), differential span all \(f^{\tau }\) and the Euler factor identity. Its Tate module over \(\text{Qbar}_{5}\) is the sum of all \(U_{f,5}^{\tau }\); the cohomological statement is dualized consistently.

At a good place \(v\) of the actual \(L'\), with residue degree \(f_{v}\) over \(p\), tr \(r_{f}(\text{Frob}_{v})\) is an \(E_{f}\) element: it is the trace of the \(f_{v}-th\) power of a matrix with polynomial \(X^{2}-a_{p}\) X+p \(\epsilon _{f}(p)\). Its equality at the fixed embedding5 to half the rational Tate trace of A identifies it with a rational number in \(E_{f}\). Every embedding \(\tau \) fixes that number. Faltings semisimplicity and Serre's general-number-field density/trace criterion therefore identify every \(r_{f}^{\tau }|G_{L'}\) with the same absolutely irreducible \(W\). All \(e\) embeddings enter; no subset or bounded coefficient field is substituted.

Consequently \(((A_{f})_{L'})^{2}\) and \(A^{e}\) have isomorphic rational Tate modules. The rational isomorphism descends from coefficient extension because an invertible element exists in the finite-dimensional intertwiner space over \(Q_{5}:\) invertibility is a nonempty Zariski-open condition and \(Q_{5}\) is infinite. Faltings' End theorem applied to their product gives the Hom comparison. Rational combinations of actual Hom maps are dense in its \(Q_{5}\) extension, so an invertible Tate map may be chosen over \(Q\) and then denominators cleared; it is an L'-isogeny. Poincare reducibility, Milne1986 Proposition12.1 and following paragraph, p.122, applies over that same ground field. A is absolutely simple, so every simple factor of \((A_{f})_{L'}\) is L'-isogenous to A. Dimensions force \(e=2r\) and \((A_{f})_{L'}\sim A^{r}\) over \(L'\). This does not use invariance of Faltings height under isogeny.

The good-prime Hecke relation \(\epsilon _{f}(p)=(a_{p}^{2}-a_{p^{2}})/p\) and density in the finite cyclotomic quotient put the full image of \(\epsilon _{f}\) in \(E_{f}\). Hence \(Q(\zeta _{s})\) embeds into \(E_{f}\), \(r\ge \phi (s)/2\), and \(s/r\le 4:\) for \(s\ge 4\) a power of2 use \(\phi (s)=s/2\); for \(s=1,2\) use \(r\ge 1\).

Put \(H=\operatorname{ker}\) \(\epsilon _{f}\). Every coefficient conjugate has the same kernel \(H\), so the entire \(A_{f}\) differential span is H-invariant. Pullback under \(X_{1}(N)\to X_{H}\) identifies differentials with invariants; on Jacobians the pullback image is the invariant subvariety up to isogeny, since the norm composed with pullback is multiplication by deg and is invertible in Hom tensor \(Q\). The \(A_{f}\) quotient restricts surjectively to that image, as its cotangent map is injective on the same H-fixed differentials. Poincare reducibility over \(Q\) therefore supplies a Q-subvariety isogenous to \(A_{f}\). Base-changing to \(L'\) gives the actual \(B_{0}\) subset \(J_{L'}\) isogenous over \(L'\) to \(A^{r}\), so \(g\ge 2r\ge 2\). It supplies exactly the surfaces of \(H_{3}\). The modular evaluation-fiber field plays no role.

\subsection{\(M_{7}\). Historical and current dispositions}

Actual earlier arguments read include private frozen \(V_{14}\) lemma481--660, b-family8097--8347 (including original r5/r3 and CL1/CL2); cutoff \S{}1; parent9 \S{}\S{}1--3; heightH0--H4; the published-source Faltings/Poincare/Li bridge; public \(V_{15}\) lemma505--705, baseline1334--1505 and b-family; direct-p5 PROOF, FAMILY-CONNECTION, FIVE-TO-THREE and CLASSICAL-LOCAL; the later direct Carayol proof and source statement. These establish \textbf{prior justification found} for the general normalized modularity, same-form companion and numerical conductor derivations. The new minimal-character choice is an explicitly changed application requiring this fresh proof and fresh review. The old qualification rollup's final-Pan access limitation is preserved as history; it is not a reason to deny applicability of Calegari's actually read final published theorem. The Li existence locator needs the explicit repair \(M_{3}\). General source statements carry independently of \(\chi \); no approval carries to a revised constructed representation by itself.

Current incorporation is \textbf{pending} for this additive bridge. Current acceptance is \textbf{pending} a nonauthor reviewer on the final hashes. No final-Pan-body comparison or complete214 theorem is inferred. The actual \(D_{210}\) \(H_{01}\) provider/source composition and the exact retained Serre1967/full11 conductor applications are now supplied in \nolinkurl{GEOMETRY_CARRY.md} and \nolinkurl{PROVIDER_SCOPE_CARRY.md;} their current acceptance is pending the nonauthor reviewer. \nolinkurl{SOURCE_USES_FINAL.json} and \nolinkurl{FINDINGS.md} name the remaining scope limits precisely.
\clearpage
% Source: research/quartic214/proofs/height-evaluation.md
\section{Integral evaluation and the trace-dual lattice}\label{q214:height-evaluation}
\noindent\textit{Source: }\nolinkurl{research/quartic214/proofs/height-evaluation.md}.\par

\begin{quote}\footnotesize Author: internal model agent \nolinkurl{/root/d214_integral_evaluation}, 2026-11-10. This is a proposed source/application bridge, not independent acceptance. Frozen private repository: \nolinkurl{a81afd7996a20e6b7bb5e48670418541c036d7aa}. The exact target is Brox \(v_{17}\), \nolinkurl{research/quadratic-ten-contact-integration/height/sources/brox-v17.tex}, lines 87--99, 117--209 and its consumer at 218--225. Whole-file and passage bindings are in \nolinkurl{SOURCE_USES.json} and \nolinkurl{CONTRACT.json}.\end{quote}

The conclusion proposed here is exactly the integral map

\[
H^0(\mathcal X,\omega_{\mathcal X/R})\longrightarrow C^\vee,
\quad\alpha\longmapsto ((\alpha/dt)(P_i))_{i=1}^d,
\]

and, for every subset \(I\) of \(g\) coordinate positions, the bound
\(\widehat{\deg}\operatorname{pr}_I C^\vee\le\frac12\log N\mathfrak D(C)\).
Here \(R=O_{K}\) and \textbf{\(C\) is the unsaturated specialization of the actual normalization algebra}, not its integral closure \(R^{d}\). There is no change in constants or curve. The stable-height metric, analytic rank estimate, vertical different, torsion different and modular realization are outside this component's claimed conclusion.

\subsection{1. Objects and finite flatness}

Let \(F=K(t)\), \(L=K(X)\), \(A=R[t]\). The smooth geometrically connected curve X/K and its finite map to \(P^{1}\) have degree \(d\). Characteristic zero makes L/F separable. Let \(Y\) be the normalization of \(P^{1}_{R}\) in \(L\), and \(B\) its algebra over Spec A. The generic fibre of \(Y\) is the original \(X:\) a smooth curve is normal, and the finite map \(X\to P^{1}_{K}\) has the defining normalization property. The chosen \(t=2\) fibre is unramified and split into the distinct K-points \(P_{1},\ldots ,P_{d}\). These are the hypotheses actually supplied by v17:79--98.

For completeness, finiteness here can be proved directly, without an unchecked general appeal to excellence. Choose an F-basis \(e_{1},\ldots ,e_{d}\) of \(L\) integral over A; multiply arbitrary basis elements by suitable nonzero elements of A to arrange this. \(D=A[e_{1},\ldots ,e_{d}]\) is finite over A and spans \(L\). If \(b\) is integral over A, each \(\operatorname{Tr}(b e_{i})\) is in A: it is integral over A, belongs to \(F\), and A is integrally closed. The trace pairing is nondegenerate. Its Gram matrix in the \(e_{i}\) basis is invertible over \(F\). Therefore all such \(b\) lie in the finitely generated A-module obtained by applying that inverse matrix to \(A^{d}\). The integral closure \(B\) is an A-submodule of that module, hence finite because A is Noetherian. The same argument on R[1/t] glues the finite normalization \(Y\).

A is regular of dimension 2. At \(q\in \operatorname{Spec}\) \(B\), with \(p=q\cap A\), the normal local domain \(B_{q}\) has dimension at most 2 and is Cohen--Macaulay: Serre's criterion gives depth \(\ge \operatorname{min}(2,\operatorname{dim} B_{q})\), hence full depth in these dimensions (Stacks 031S). Going down for the integral inclusion \(A\subset B\) with A normal (00H8), together with incomparability, gives \(\operatorname{ht}(q)=\operatorname{ht}(p)\). Incomparability itself follows by localizing at \(p\) and reducing modulo \(p:\) an integral algebra over a field has only maximal primes, so two comparable primes over \(p\) coincide. Thus the local fibre dimension is zero and

\[
\dim B_q=\dim A_p+\dim(B_q/pB_q)=\dim A_p.
\]

All three hypotheses of miracle flatness 00R4 are now checked, including its dimension equality. Hence \(B_{q}\) is \(A_{p}-\text{flat}\) at every \(q\). Flatness is local, so \(B\) is finite flat, consequently finite projective, of rank \(d\). The same holds on the second target chart. \(Y\) is an integral normal proper flat arithmetic surface, finite over \(P^{1}_{R}\). No smoothness of its special fibres and no Gorenstein property of \(Y\) is asserted.

Consequently \(C=B\otimes _{A}\) \(R=B/(t-2)B\) is finite projective of rank \(d\) over \(R\). Its generic fibre is the \textbf{specified} \(K^{d}\) via evaluation at \(P_{i}\). Its map into \(K^{d}\) is injective because \(C\) is R-torsion-free. Each coordinate of an element of \(C\) is integral over \(R\) and belongs to \(K\), hence lies in \(R\). Thus \(C\subset O:=R^{d}\) with full rank. No claim that \(C\) is reduced modulo a finite prime, étale over \(R\), or normal is needed.

\subsection{2. A common regular model with the required direction}

Write \(M=\mathcal X\) for the given minimal regular semistable model. Take the reduced closure \(\Gamma \) of the graph of the identity \(X\to Y_{K}=X\) inside \(M\times _{R}\) \(Y\). Both projections are proper birational morphisms. Stacks 0BFP, parts (1)--(2), says that the maximal isomorphism locus of \(\Gamma \to M\) contains all codimension-zero and codimension-one points, since the codimension-one local rings of \(M\) are DVRs. Its complement \(T\) is a finite set of closed points: it is closed in a Noetherian surface and has dimension zero. If \(T\) is empty take \(Z=M\). Otherwise Stacks 0AHI applies to \(\Gamma \to M\) and supplies a finite sequence of point blowups

\[
Z=M_n\to M_{n-1}\to\cdots\to M_0=M
\]

dominating \(\Gamma \). Blowing up a closed regular point of a regular surface is regular: for regular parameters u,v its charts are T[v/u] and T[u/v], and at an exceptional point \(u\) and a lift of a uniformizer of the local ring of \(P^{1}_{\kappa (x)}\) form regular parameters. This also covers non-rational closed centres. The resulting \(Z\) is integral, proper over \(R\) and R-flat (its coordinate rings have no R-torsion); the generic fibre is \(X\), because every centre lies over a finite prime. Indeed \(\Gamma \to M\) was already an isomorphism on the entire generic fibre. The two maps are \(\alpha :Z\to M\) and \(\beta :Z\to Y\), both proper birational, with identity on \(X\).

In the published Javanpeykar paper, Lemma 2.3.1, proof on p.99, the point-blowup calculation is

\[
\omega_{M_i/R}=a_i^*\omega_{M_{i-1}/R}\otimes O_{M_i}(E_i).
\]

This is used from the \textbf{proof}, rather than confusing the lemma's point-height conclusion with the required sheaf statement. Since \(E_{i}\) is effective, pullback of a section of \(\omega _{M_{i-1}/R}\) is a section of \(\omega _{M_{i}/R}\); iterate to get \(\alpha ^{*}E\subset H^{0}(Z,\omega _{Z/R})\). The formula is identified by the generic differential, so this operation does not multiply a differential by a rational function. Equivalently the local calculation is the Cartier hypersurface \(uV-vU=0\) in \(P^{1}_{T}\); adjunction gives the exceptional twist +E.

Apply 0BFP again to \(\beta \). \(Y\) is normal, hence its codimension-one local rings are DVRs; \(Z\) is reduced. There is an open \(U\subset Y\) containing every codimension-one point with \(\beta ^{-1}(U)\cong U\). The complement has codimension at least two. The regular differential on \(Z\) therefore gives a section of \(\omega _{Y/R}|_{U}\), compatible with the same generic differential \(\alpha \). This direction descends sections from a regular model; it does \textbf{not} claim every canonical section on \(Y\) pulls back regularly to a resolution. Such a converse could require singularity hypotheses and is not used.

\subsection{3. The canonical module and its extension}

Use relative dualizing complexes normalized over \(R\), with \(\omega _{P^{1}_{R}/R}[1]\) the relative dualizing complex of \(P^{1}_{R}\). Transitivity (0ATX) and finite duality (0BUL, whose proof identifies the finite pushforward with Hom) identify the relative canonical module on the affine target chart with

\[
M_B\,dt:=\operatorname{Hom}_A(B,A)\,dt.
\]

\(B\) is projective, so higher \(\text{Ext}_{A}(B,A)\) vanish; there is only the indicated sheaf in relative degree \(-1\). The dt is the usual smooth differential generator on \(A^{1}_{R}\). The construction respects restrictions to affine opens, as explicitly required in 0BUL, so these local modules are one global \(\omega _{Y/R}\). One may also use the Cohen--Macaulay canonical-module statement in 0AWQ. The more general algebraic assertion 0C0I is \textbf{nonunique} and is insufficient by itself to choose the generic trace normalization; it is not used to choose that map.

Here is also a direct check of the reflexivity needed in \(v_{17}\). The trace isomorphism

\[
M_B\cong\{f\in L:\operatorname{Tr}_{L/F}(fB)\subset A\}
\tag{*}
\]

is Stacks 0BW2: \(A\to B\) is finite, A-nonzerodivisors remain nonzerodivisors in the domain \(B\), and L/F is finite separable. In particular \(M_{B}\) is B-torsion-free and has B-rank one. At a height-two \(q\) of \(B\), choose regular parameters x,y in \(A_{p}\) \((p=q\cap A)\). \(M_{B}\) is A-projective. Thus x,y is an \(M_{B}-\text{regular}\) sequence before localizing and remains so after localizing at \(q\). Both belong to \(q\); hence \(\operatorname{depth}_{B_{q}}(M_{B})_{q}\ge 2\). At a height-one prime torsion-freeness gives \(\operatorname{depth}\ge 1\). Thus \(M_{B}\) satisfies \(S_{2}\), and the normal-scheme criterion 0AVT Lemma 31.13.13 makes it reflexive. The same reasoning on the second chart gives global reflexivity. This argument does not require invertibility of \(M_{B}\) as a B-module.

The normal surface \(Y\) is \(S_{2}\); at points outside \(U\) its local depth is at least two. Lemma 31.13.12 in 0AVT then identifies \(\omega _{Y/R}\) with \(j_{*}(\omega _{Y/R}|_{U})\). Therefore the section from \S{}2 extends uniquely to all of \(Y\). We have a well-defined R-linear map \(E\to H^{0}(Y,\omega _{Y/R})\), taking each section to the same differential on \(X\).

\subsection{4. Trace normalization and specialization are the actual evaluation}

There are two separate comparisons; both are necessary.

First, on an open of the generic t-line avoiding branch points, the finite map is étale. The usual smooth identification of the canonical module with \(\Omega ^{1}\) is the one obtained from the conormal/adjunction construction of Stacks 0E9Z. Its construction uses the determinant of the conormal exact sequence, and therefore commutes with flat base change. Duality transitivity is 0ATX, the tensor comparison is 0B6U, and the flat base-change isomorphism is 0E9U. All schemes here are separated and finite type over the Noetherian base \(R\) (the category defined in 0F42); on the generic chart one can equally work over \(K\).

To compute the identification, use an étale neighbourhood splitting the finite étale cover (04HN). On that neighbourhood \(B\) becomes \(A^{d}\). Both the exterior/conormal construction and finite Hom duality are compatible with this flat base change. For a product of identity morphisms, the i-th differential \(f_{i}\) dt defines the coordinate functional \(b_{i}\mapsto f_{i}\) \(b_{i}\), and finite pushforward sums these functionals. Thus the map is

\[
(f_i\,dt)_i\longmapsto [\,(b_i)_i\mapsto\sum_i f_i b_i\,]\,dt.
\]

This formula is invariant under permuting the sheets and descends by faithful flatness. The sum is precisely the algebra trace. Consequently the trace map in (*) identifies the literal rational differential \(f\) dt with \(b\mapsto \operatorname{Tr}(fb)\). There is no degree factor \(d\), different factor, or chosen scalar. This argument also explains why a bare nonunique identification such as 0C0I would not settle the issue. At \(t=2\) all sheets are unramified in characteristic zero, so \(f=\alpha /dt\) is regular at every \(P_{i}\).

Second, because \(B\) is finite projective the natural map

\[
\operatorname{Hom}_A(B,A)\otimes_A R
\xrightarrow{\sim}\operatorname{Hom}_R(B\otimes_A R,R)
\tag{**}
\]

is an isomorphism. Explicitly \((\lambda \otimes 1)(b \text{mod}(t-2))=\lambda (b)\) \(\text{mod}(t-2)\); it is independent of the lift since \(\lambda ((t-2)b)=(t-2)\lambda (b)\). It is an isomorphism for free modules, and therefore for finite projective modules by expressing them as direct summands of finite free modules. This proves the required base change without assuming the map \(A\to R\) is flat.

After tensoring (**) with \(K\), \(C_{K}=K^{d}\). The preceding trace computation and compatibility of the trace of a finite projective algebra with base change (the trace of a matrix is unchanged by reducing its entries) give

\[
\bar\lambda(c)=\sum_i (\alpha/dt)(P_i)c(P_i).
\]

Since \(\bar\lambda (C)\subset R\), the vector \(((\alpha /dt)(P_{i}))\) belongs to exactly \(C^{\vee }\subset K^{d}\). Nondegeneracy of the dot-product trace on \(K^{d}\) identifies \(\operatorname{Hom}_{R}(C,R)\) with \(C^{\vee }\); it does not identify it with \(O\) unless \(C=O\). Thus the integral map asserted in v17:157--168 is obtained with its exact coordinate convention. No assumption about reduction of \(t=2\) at bad primes was needed.

\subsection{5. Trace-dual projection and arithmetic degree}

This part is an elementary local calculation, reconstructed here rather than imported from a review verdict. For every nonzero prime \(p\) of \(R\), choose an \(R_{p}-\text{basis}\) of \(C_{p}\), with matrix \(U\) in the coordinate basis of \(O_{p}\). \(U\) has nonzero determinant; trace on \(K^{d}\) is the dot product. The dual basis matrix is \((U^{T})^{-1}\), and the trace Gram matrix is \(U^{T}\) \(U\). Therefore

\[
\ell_{R_p}(C_p^\vee/O_p)=\ell_{R_p}(O_p/C_p)
=v_p(\det U)=\tfrac12v_p(\mathfrak D(C)).
\]

The inclusion \(O\subset C^{\vee }\) follows directly from \(C\subset O\). For every coordinate subset \(I\) (no special prime-dependent choice), the induced projection \(C^{\vee }/O\to pr_{I}(C^{\vee })/R^{I}\) is surjective: lift an element by its definition as an image and note \(pr_{I}(O)=R^{I}\). Its finite-prime lengths cannot increase. \(F_{I}=pr_{I}(C^{\vee })\) is finite and torsion-free, hence projective over the Dedekind ring \(R\), and contains \(R^{I}\) with finite quotient.

Equip \(F_{I}\) at every complex embedding with the standard Euclidean metric on \(K^{I}\). Its determinant basis \(e_{I}\) has norm 1. The determinant line is a fractional ideal containing \(R\) \(e_{I}\). Locally its length relative to \(R\) \(e_{I}\) equals \(\ell (F_{I}/R^{I})\). By the arithmetic-degree convention of v17:49--60 \((or by \text{multiplying} e_{I} by a \text{nonzero} \text{denominator} \text{ and }  \text{applying} \text{the} \prod  \text{formula})\),

\[
\widehat{\deg}F_I=\sum_p\ell(F_{I,p}/R_p^I)\log\#\kappa(p)
\le\tfrac12\log N\mathfrak D(C).
\]

This also checks the sign for an overlattice. The proof treats nonprincipal global modules by local bases, all finite primes including residue characteristic 2, and every \(I\) with \(|I|=g\le d\). The existence of an \(I\) for which evaluation has rank \(g\) belongs to the separate analytic/rank component; this component supplies the lattice conclusion for every such \(I\) once chosen.

\subsection{6. Exact consumer and bounded status}

The output is the same integral map and same trace discriminant used at v17:218--225. Composing with a chosen coordinate projection gives \(E\to F_{I}\); its generic determinant is nonzero only after the rank argument. Combining with the separate metric/determinant proof gives the unchanged Theorem 1.1 estimate. The modular-height consequence feeds \nolinkurl{height/PROOF.md} H0.3 and \(H_{4}\) and the retained existence/global bound, including the downstream 256 route. No replacement order, normalization, extension field, curve or metric has entered this component.

History: the earlier foundation report had already read a complete project proof of this route and specifically requested source/application closure. The retained quadratic integration review explicitly checked \(v_{17}\) at component scope and kept inherited foundational qualifications. Accordingly this is prior justification found plus an expanded source bridge, proposed for incorporation; it is not evidence that no earlier reviewer had supplied these details. No historical verdict is rewritten. A non-author must review this frozen bridge and its sources before the application can be called independently accepted. This component does not close the full arithmetic theorem or all of Protocol \(D\).
\clearpage
% Source: research/quartic214/proofs/height-different.md
\section{Different and modular-model bridge for the retained height bound}\label{q214:height-different}
\noindent\textit{Source: }\nolinkurl{research/quartic214/proofs/height-different.md}.\par

\begin{quote}\footnotesize Author: internal model \nolinkurl{/root/d214_different_bounds}, 2026-11-10. This is an author proposal for fresh non-author review. It neither writes nor supplies an independent acceptance. Frozen target: private repository commit \nolinkurl{a81afd7996a20e6b7bb5e48670418541c036d7aa}, \nolinkurl{research/quadratic-ten-contact-integration/height/sources/brox-v17.tex}, SHA-256 \nolinkurl{2ba5dd497276073d92db95f0119bede82d182997c66375e5be294157de6d9f7c}.\end{quote}

The proposed conclusion is that Brox v17:274--328, 354--451 and 461--488 have the required different bounds, including at every imperfect-residue vertical generic DVR. The resulting bound is unchanged:

\(\operatorname{log} \operatorname{Norm} \operatorname{disc}(C)/[K:Q] \le  d_{H} (\log 12 + \log N + 2\operatorname{log} d_{H})\).

Together with the separate integral and analytic evaluation interfaces, which require their own non-author review, this is exactly the input to \nolinkurl{height/PROOF.md} H0.3. No full-torsion splitting-field degree, extra factor two, or changed height constant is introduced. The source citation to Lemma49.12.6 is valid and now bound to its exact current tag0BWJ. The mathematical argument below supplies the expanded explicit source comparisons; incorporation and current independent acceptance remain pending.

\subsection{1. Fixed objects and the trace-discriminant identity}

Put \(R=O_{K}[t]\), where \(K\) is any finite extension used for semistability and splitting the \(t=2\) fibre. Let \(B\) be the normalization of \(R\) in \(K(X)\); let \(d=[K(X):K(t)]\). The integral-evaluation lane supplies finite projectivity of this same \(B\) of rank \(d\). Put \textbf{\(C=B/(t-2)B\)}, retaining this order rather than replacing it with its normalization \(O\). Its generic algebra is \(K^{d}\) after the chosen splitting extension. The determinant of \(B\to \operatorname{Hom}_{R}(B,R)\), \(b\mapsto (b'\mapsto \operatorname{Tr}(bb'))\), is a nonzero section of an invertible line. Base change \(t\mapsto 2\) gives exactly the trace determinant of \(C\), since trace of multiplication on a finite projective module commutes with base change. Thus the section's arithmetic intersection with \(t=2\) is log Norm disc(C). This is an equality of determinant ideals and does not assert that \(C\) is maximal.

At every height-one point \(q\) of Spec \(R\), \(V=R_{q}\) is a DVR and \(T=B\otimes _{R}\) \(V\) is its finite normal integral closure in the generic separable algebra. \(T\) can be semilocal; do not silently identify a single localization \(T_{m}\) with a finite V-algebra. Complete \(V\). Then \(T\otimes _{V}\) Vhat is the product of the complete DVRs \(U_{j}\) over every \(m_{j}\) above \(q\), with generic fields \(F_{j}\). Let \(e_{j}\) be the ramification index, \(f_{j}\) the full residue degree (including its inseparable degree), and \(r_{j}\) the Dedekind different exponent. Each \(U_{j}\) is finite free over Vhat. Filtering \(U_{j}/\pi _{V}\) \(U_{j}\) by powers of its uniformizer gives

\(\operatorname{rank}_{\text{Vhat}} U_{j} = \operatorname{dim}_{k}(U_{j}/\pi _{V} U_{j}) = e_{j} f_{j}\), hence \(\sum _{j} e_{j} f_{j}=d\).

This dimension calculation uses finiteness and flatness, not residue separability or a defectless-extension assertion for arbitrary valued fields. The completion decomposition and its whole degree \(d\) include every place; no preferred component is selected.

Under the generic trace, \(T^{\vee }=\{x:\operatorname{Tr}(xT)\subset V\}=\operatorname{Hom}_{V}(T,V)\). The valuation of the trace determinant is \(\text{length}_{V}(T^{\vee }/T):\) choose a V-basis, whose trace Gram matrix presents this cokernel. On each completed component \(U_{j}^{\vee }=\pi _{j}^{-r_{j}}U_{j}\). The length there is \(f_{j}\) \(r_{j}\). Consequently

\(v_{q}(\operatorname{disc}(B/R)) = \sum _{j} f_{j} r_{j}\).

This proves the local trace/discriminant formula used here directly. It removes dependence on an unchecked reading of Serre's book for this formula. Integrality \(r_{j}\ge 0\) follows because traces of integral elements lie in the integrally closed DVR \(V\), so \(U_{j}\subset U_{j}^{\vee }\). Completion preserves the trace pairing, its determinant valuation and the fractional duals by finite freeness.

\subsection{2. Lenstra's estimate really covers the vertical DVR}

The actual published source is Javanpeykar, ANT8(2014), \S{}4.1, printed pp112--113, Propositions4.1.1 and4.1.3 and Remark4.1.2. The retained PDF is the final publisher article, DOI10.2140/ant.2014.8.89; see \nolinkurl{SOURCE_USES.json}. The setup is a characteristic-zero DVR and its finite separable field extension. No perfect-residue assumption occurs. The remark explicitly distinguishes the separable-residue result in Serre and the imperfect-residue application. Proposition4.1.3 obtains a bound for each maximal ideal by completing it; it does not require the original integral closure to be a local ring.

Here is a direct reconstruction useful for checking conventions. For a finite extension of complete DVRs \(V\subset U\) with fraction-field degree \(n\), choose uniformizer \(\pi \) of \(V\). The element \(y=1/(n\pi )\) has field trace \(1/\pi \) and is therefore outside \(U^{\vee }\). As \(U^{\vee }\) is a fractional ideal of the DVR \(U\), it is strictly smaller than yU. Taking inverse fractional ideals gives \(D_{U/V}\) strictly larger than \(n\pi U\). With \(v_{U}(\pi )=e\) this yields

\(r < e(v_{V}(n)+1)\), equivalently \(r \le  e-1+e v_{V}(n)\).

This uses only characteristic zero, finite freeness, nondegeneracy of generic trace and total ordering of fractional ideals in a DVR. In particular it applies when the residue extension is inseparable. For a component of an algebra of total degree \(d\), its completed degree \(n\) is at most \(d\). In residue characteristic \(p\),

\(v_{V}(n) = v_{V}(p) v_{p}(n) \le  v_{V}(p) \text{floor}(\operatorname{log} d/\operatorname{log} p)\).

Now take the actual vertical point \(q=(\text{pfrak})\) of Spec \(O_{K}[t]\). Then \(V=O_{K}[t]_{\text{pfrak}}\), residue field \(k(\text{pfrak})(t)\), which generally is \textbf{imperfect}, fraction field \(K(t)\), and \(v_{V}(p)=e_{\text{pfrak}}\). \(V\) is excellent \((a \text{localization} of a \text{finite}-\text{type} \text{algebra} \text{ over }  O_{K})\), so the normalizations are finite. For \textbf{every} component \(j\) the displayed estimate gives

\(r_{j} \le  e_{j}-1 + e_{j} e_{\text{pfrak}} \text{floor}(\operatorname{log} d/\operatorname{log} p)\).

Multiplying by all \(f_{j}\) and using \(\sum _{j} e_{j} f_{j}=d\) bounds the whole discriminant coefficient by

\(d(1+e_{\text{pfrak}} \text{floor}(\operatorname{log} d/\operatorname{log} p))\).

There is no restriction \(p\le d\) needed: if \(p>d\) the floor is zero and the same estimate holds. At a horizontal branch point the residue characteristic is zero. In completed curve coordinates, after the unramified coefficient-field extension one has \(t-a=u^{e}\) \(v(u)\), \(v(0)\ne 0\). The derivative has order \(e-1\) because \(e\) is a unit. The local complete-intersection/Jacobian formula for the different therefore gives \(r=e-1\). Thus its discriminant coefficient is at most \(d\). This applies to the horizontal branch divisor at \(t=0\); \(t=1\) and infinity contribute no intersection with the section \(t=2\).

If normalization is finite étale over \(O_{K_{0}}[S^{-1},t,t^{-1},(t-1)^{-1}]\), then after any finite K/K0 it stays finite étale there. The base is normal and a finite étale algebra over a normal ring is normal; its generic algebra is the base-changed cover. It is therefore the normalization of that algebra on the open. This identifies the actual \(B\) there, and rules out additional vertical discriminant components outside \(S\) after the auxiliary extension. It does not assert that normalization commutes with arbitrary base change everywhere.

The horizontal intersection with \(t=0\) has normalized value log2. At pfrak above \(p\) the vertical intersection is \(f_{\text{pfrak}}\) logp. Summing the last bound, using \(\sum _{\text{pfrak}|p}e_{\text{pfrak}}\) \(f_{\text{pfrak}}=[K:Q]\) and \(\sum f_{\text{pfrak}}\le [K:Q]\), gives

\(\operatorname{log} \operatorname{Norm} \operatorname{disc}(C)/[K:Q] \le  d(\log 2 + \operatorname{log} R_{S} + |S|\operatorname{log} d)\).

The determinant Cartier divisor has no additional codimension-two divisor terms. Its proper pullback to \(t=2\) incorporates the intersections even when several divisors meet there. These statements concern the same nonmaximal fibre order \(C\) of \S{}1.

\subsection{3. Torsion different as an identity of fractional ideals}

Let \(V\) be the excellent characteristic-zero DVR occurring above, E/V an elliptic scheme, \(n\ge 1\), and \(E[n]=\operatorname{Spec}\) A. Its generic algebra \(L\) is finite étale; it need not be a field. The fact that \([n]\) is finite locally free of rank \(n^{2}\) can be checked using the field theorem Stacks0BFG, Lemma39.9.8, on every fibre: properness and finite fibres give finiteness. Stacks49.12.6/0BWJ applies with \(S=\operatorname{Spec}\) \(V\), \(X=Y=E\) smooth of relative dimension1, and \(f=[n]\) locally quasi-finite. It gives flatness, and its proof gives the LCI property. Thus the rank is \(n^{2}\) by its fibre ranks. The same application works over the regular base \(U\) used in \S{}4. Where \(n\) is invertible, the differential below is invertible, so \([n]\) is étale.

The graph of \([n]\) is the base change of the diagonal of the smooth relative curve E/V, hence a regular immersion into \(E\times _{V}\) \(E\). Projection is smooth. This gives an explicit local-complete-intersection factorization of \([n]\). Invariant differentials trivialize \(\Omega _{E/V}\) as a line pulled back from \(V\). The group law differentiates as addition at the identity; induction gives \([n]^{*}\omega =n\omega \) for positive \(n\). Therefore the relative differentials of \([n]\) have presentation \(O_{E}\) \(--n-\to \) \(O_{E}\) and \(\text{Fitt}_{0}\) \(\Omega _{[n]}=(n)\). This is \(n\), \textbf{not the degree \(n^{2}\)}.

Base change by the zero section \(E[n]\to \operatorname{Spec}\) \(V\) preserves finite flatness, the local-complete-intersection property, the displayed presentation and its Fitting ideal. For this finite flat LCI \(V\to A\), Stacks0BWE (Lemma49.10.1) checks the exact syntomic hypotheses, Stacks0BWF (49.10.4) gives invertibility of \(\omega _{A/V}\), and Stacks0BWG (49.12.3, on the archived section0BWB) identifies the scheme different with \(\text{Fitt}_{0}\) \(\Omega _{A/V}=(n)\). Stacks Definition49.9.1/0BW4 defines that scheme different as \(\text{Ann}(\omega _{A/V} / A\tau )\).

The generic trace identification is fixed, rather than an unspecified module isomorphism: Stacks0BW2 (49.8.2) maps \(x\in A^{\vee }\) to \(b\mapsto \operatorname{Tr}_{L/\text{Frac}(V)}(xb)\). Stacks0BW3 (49.8.3) maps \textbf{1 to \(\tau \)}. Under this map, the annihilator in the preceding paragraph is

\(\{a\in A: a A^{\vee }\subset A\}\).

Since \(A^{\vee }\) is an invertible fractional A-ideal (from0BWF) and contains1, this is exactly its inverse fractional ideal. Thus \((A^{\vee })^{-1}=nA\), and inverting again gives the literal equality inside \(L\)

\(A^{\vee }=n^{-1} A\).

This argument checks the scalar fixed by trace; an abstract isomorphism \(\omega \cong A\) by itself would not do so. It also explains how the Kähler different equals the Dedekind different here. In Stacks49.12.4/0BWH the local generator is mapped so that the Jacobian determinant maps to \(\tau \), consistently with the same equality. Brox's citation49.12.6 is the valid current tag0BWJ. It applies directly to \([n]\) as checked above and gives the same differential-determinant formula. The archived section0BWB includes Lemmas49.12.1--4, Example49.12.5 and Lemma49.12.6; a shortened web-text extraction had omitted the latter part. This author initially misclassified the citation as absent; fresh review caught that error before freeze. \nolinkurl{FINDINGS.md} preserves the correction.

Normalize A in \(L\), obtaining finite Atilde (excellence of V). Both A and Atilde are finite torsion-free V-modules with generic algebra \(L\). Trace duality reverses the actual inclusion \(A\subset \text{Atilde}:\)

\(\text{Atilde}^{\vee } \subset  A^{\vee } = n^{-1} A \subset  n^{-1} \text{Atilde}\).

Hence \(n\) \(\text{Atilde}^{\vee }\subset \text{Atilde}\). Componentwise, the Dedekind different contains nAtilde; equivalently \textbf{\(r_{j}\le v_{j}(n)\)}. No perfect-residue or monogenic-normalization hypothesis is needed. This is also valid for \(n=1\), when every exponent is zero.

For clarity, the intermediate-field direction can be proved rather than inferred from a slogan. Complete at each chosen place in a tower \(V\subset W\subset U\) of normal finite extensions. Write \(W^{\vee }_{V}=\alpha W\), possible since \(W\) is a DVR. Trace transitivity, tested also on every multiplication by \(b\in W\), says

\(U^{\vee }_{V} = \{x: \operatorname{Tr}_{U/W}(xU)\subset W^{\vee }_{V}\} = \alpha  U^{\vee }_{W}\).

Taking inverse fractional ideals yields \(D_{U/V}=D_{U/W}\) \(D_{W/V}\) \(U\), with exponent \(r_{U/V}=r_{U/W}+e_{U/W}\) \(r_{W/V}\). All differents have nonnegative exponents. Thus \(r_{U/V}\le v_{U}(n)=e_{U/W}\) \(v_{W}(n)\) implies \(r_{W/V}\le v_{W}(n)\). Passing to completions is legitimate by \S{}1; this proves the assertion at every place of every intermediate field. The estimate concerns a single-point field and its subfields; it requires no bound on a torsion splitting field.

\subsection{4. The actual primitive torsion quotient}

For \(N\ge 1\) and \(H\subset (Z/NZ)^{\times }\) \(\text{containing}-1\), let \(t_{H}=[(Z/NZ)^{\times }:H]\). Set

\(U=\operatorname{Spec} Z[1/(6N),t,t^{-1},(t-1)^{-1}]\),

\(E_{t}: y^{2}=x^{3}-3t/(t-1)x-2t/(t-1)\).

Direct substitution into \(\Delta =-16(4a^{3}+27b^{2})\), \(c_{4}=-48a\) and \(j=c_{4}^{3}/\Delta \) gives \(\Delta =1728t^{2}/(t-1)^{3}\) and \(j=1728t\). All denominators and \(\Delta \) are units on \(U\). It is an elliptic scheme there. As \(N\) is invertible, \(E_{t}[N]\) is finite étale. Its primitive locus is open and closed: remove \(E_{t}[m]\) for the finitely many proper divisors m|N; these are unions of components étale locally on \(U\). Thus the exact-order-N locus is finite étale as well, including \(N=1\) with the usual order-one point.

The constant \(H\) action is scalar multiplication. On primitive points \(aP=P\) implies \((a-1)P=0\), hence \(a=1\) modN. Étale locally the algebra is a product of copies of the base ring and \(H\) permutes these copies; the quotient is the product indexed by the H-orbits. This constructs a finite étale quotient even at residue characteristics dividing |H|. No averaging by |H| and no assumption |H| invertible is used. Freeness makes every orbit have |H| points.

Identify its generic fibre over \(Q(t)\) by the moduli definition of \(X_{1}(N)\). Over an algebraic closure, \(j\ne 0,1728\), an elliptic curve has automorphism group \(\{\pm 1\}\). Given another curve with that \(j\) and a primitive point, choose an isomorphism to \(E_{t}\); two choices differ \(by\pm 1\) and hence give the same H-orbit. This bijection is Galois compatible: changing the chosen isomorphism after a Galois action again differs only \(by\pm 1\). Thus it is an isomorphism of the finite étale \(Q(t)-\text{algebras}\), not just a count of their geometric points. Quadratic twisting is exactly this \(\pm 1\) ambiguity, already killed in \(H\). The generic cover is the actual \(X_{H}=X_{1}(N)/H\) with \(t=j/1728\). A normal finite étale cover over normal \(U\) is its normalization there.

Geometric primitive vectors are a single SL2(Z/NZ) orbit: a primitive pair (a,b) can be completed to a determinant-one matrix by Bézout modN. This is compatible with the standard analytic model of the connected modular curve; the construction introduces no full-level cyclotomic component. Counting primitive vectors gives \(N^{2}\prod _{p|N}(1-p^{-2})\); dividing by \(|H|=\phi (N)/t_{H}\) yields

\(d_{H}=N t_{H} \prod _{p|N}(1+p^{-1})\).

For \(N=1\) the products and \(\phi (1)\) are one, and \(d_{H}=1\); the positive-genus hypothesis later simply excludes irrelevant small-genus cases. Published von Känel2021 \S{}4, \(p_{8}\), (4.1) independently records the degree/index convention: the factor two occurs exactly \(\text{when}-I\) is absent, and the only possible branch values are the elliptic points and cusp. \(\text{Here}-I\) is present \(\text{because}-1\in H\), so no additional two divides the formula. The finite étale model already proves absence of branch away \(\text{from}0,1,\infty \) in characteristic zero.

A geometric H-orbit has Galois stabilizer containing the stabilizer of any chosen primitive point \(P\) in it. Consequently its field of definition is an intermediate field of \(Q(t)(P)/Q(t)\). This remains true after replacing \(Q\) by the auxiliary \(K\) and completing at a vertical generic point: each component field and each chosen place is covered by \S{}3. The quotient need not itself be a subalgebra factor of the integral nonnormal A; the generic field inclusion and its normalizations are the objects to which the intermediate-field conclusion applies.

\subsection{5. Same constant at every prime}

For every prime \(\text{pfrak}|p>3\) of \(K\), take \(V=O_{K}[t]_{\text{pfrak}}\). Both \(t\) and \(t-1\) are units there, and1728 is a unit, so \(E_{t}/V\) is an elliptic scheme even if p|N. Do \textbf{not} use the finite étale model \(U\) at such \(p\), since \(U\) inverted \(N\); instead use the finite flat torsion theorem of \S{}3. The quotient component field identified in \S{}4 is an intermediate field of a single N-torsion-point field. Its different exponent is at most \(v_{j}(N)=e_{j}\) \(v_{V}(N)\). By \S{}1 the full vertical discriminant coefficient is therefore

\(\sum _{j} f_{j} r_{j} \le  (\sum _{j} f_{j} e_{j}) v_{V}(N) = d_{H} v_{\text{pfrak}}(N)\).

After intersection with the section, summation over pfrak and division by [K:Q], this contributes at most \(d_{H}\) \(\sum _{p>3}v_{p}(N)\log p\le d_{H}\) logN. Every component is included; the algebra's rank is \(d_{H}\), not \(N^{2}\) and not the degree of a full torsion field.

At \(p=2\) and3 use \S{}2 rather than the elliptic model, which is not asserted smooth there. Their combined normalized bound is at most \(d_{H}(\log 6+2\operatorname{log} d_{H})\). This includes primes not dividing \(N\) and levels with \(d_{H}<p\). The horizontal \(t=0\) contribution is at most \(d_{H}\) log2. Altogether

\(\operatorname{log} \operatorname{Norm} \operatorname{disc}(C)/[K:Q] \le  d_{H}(\log 12+\log N+2\operatorname{log} d_{H})\).

This is Broxv17:461--488 with every local supplier matched to the actual objects. The independent evaluation conclusion, if accepted for this same \(C\), gives

\(\hF(J_{H}) \le  d_{H}/2 (\log 12+\log N+2\operatorname{log} d_{H}) + g_{H} \operatorname{log}(2/\sqrt{\pi})\).

The stable-height metric and the evaluation lattice are separate interfaces assigned to the other two author lanes; this document does not self-approve their conclusions. The prime summations show why the auxiliary \(K\) cancels, so no [K:Q] enters H0.3 or the later surface-isogeny field degree.

\subsection{6. History, incorporation and limits}

Actual earlier proof text read: frozen Broxv17 \S{}\S{}3--4 at274--488; \nolinkurl{height/PROOF.md} H0.3 and \(H_{4}\); \nolinkurl{research/v14-brox-order8/PROOF.md}45--56; its actual \nolinkurl{review/SOURCE_ASSESSMENT.md} sections ``The lattice and discriminant argument'' and ``The modular cover and surface''; its \nolinkurl{review/REPORT.md}19--33; \nolinkurl{research/v14-small-prime-conductor/review/REPORT.md} section ``Fresh height estimate''; and retained Protocol-D foundationREPORT/H04. The earlier source assessment already states the imperfect-residue Lenstra application, exact torsion derivative \(n\), reversed trace dual and intermediate-field conclusion. Its final review expressly retains published-source/whole-assembly limits. Therefore \textbf{prior justification found and read} is the historical disposition for these arguments. This packet is a source-pinned expansion and precise current-tag binding, not a claim to have first repaired a newly discovered arithmetic error.

Current incorporation: Broxv17 contains the substantive arguments but not this expanded bridge and current-tag expansion; the private height H0.3 retains their inherited status. This packet proposes the bridge for additive incorporation. Current acceptance: pending a reviewer other than this author after the packet is frozen, including the precise joins to the other two evaluation lanes. No global modularity, height-transfer, 214 existence or256 theorem is accepted by this bounded work. The source inventory records what was inspected; it does not claim exhaustive search of every historical archive, nor a new inspection of the complete Serre book. The relevant trace, different, tower and tame-horizontal assertions are instead derived explicitly above.
\clearpage
% Source: research/quartic214/proofs/height-metric.md
\section{Metric and archimedean evaluation bridge}\label{q214:height-metric}
\noindent\textit{Source: }\nolinkurl{research/quartic214/proofs/height-metric.md}.\par

\begin{quote}\footnotesize Author: internal model task \nolinkurl{/root/d214_metric_evaluation} (requested Sol61 lane), 2026-11-10. This is an author proposal awaiting a separately authored review. It supplies only Brox \(v_{17}\) lines62--78 and210--251, conditional on the exact integral map and trace-lattice bound specified below. It supplies neither acceptance of those algebraic inputs nor whole-214/256 closure. The immutable manuscript is private commit \nolinkurl{a81afd7996a20e6b7bb5e48670418541c036d7aa}; exact whole-file and passage bindings are in \nolinkurl{SOURCE_USES.json.}\end{quote}

\subsection{1. Published normalization and actual Jacobian}

Read the final published Javanpeykar2014, \S{}\S{}2.1/2.3, printed pp93--94/97--98, and final von Känel2021 \S{}2.1, printed \(p_{5}\). The former specifies the inner product \((i/2)\int _{X} \omega \wedge \overline{\eta }\), its determinant metric, and \(\text{hFal}(X)=\text{Ardeg}(\operatorname{det} p_{*}\omega )/[K:Q]\), equal to stable Faltings height of the Jacobian. The latter defines original relative height using the Néron top-differential line and all embeddings, then defines stable height by semistable extension and proves its extension invariance. Its paragraph after(2.2) explicitly distinguishes the alternative shifted height by \((g/2)\operatorname{log} \pi \). That shifted height is not used here. Sources were acquired from their publishers on2026-11-10; hashes and exact source locations are retained in \nolinkurl{SOURCE_USES.json.}

Apply Javanpeykar's identification to the \emph{same} smooth projective geometrically connected curve \nolinkurl{X/K}, genus \(g\ge 1\), and its minimal regular semistable model \(p:{\mathcal X}\to \operatorname{Spec} O_{K}\). Thus \(E=H^{0}({\mathcal X},\omega _{{\mathcal X}/O_{K}})\) with the displayed \(L^{2}\) metric has

\(\text{Ardeg}(E)=[K:Q] \hF(\operatorname{Jac} X)\).

This is the Jacobian itself, not an isogenous abelian variety, a modular factor, the singular normalization model, or an Arakelov metric on individual cotangent fibres. The relevant metric is on global differentials. Javanpeykar's published statement supplies the integral model identification; it is a theorem input, not a claim that a generic vector-space identification alone identifies lattices.

For an explicit check of the archimedean determinant scalar, choose symplectic integral cycles and normalized holomorphic forms \(\omega _{1},\ldots ,\omega _{g}\); write their period matrix \((1,\tau )\). Riemann's bilinear identity gives their Gram matrix \(Im \tau \). The analytic Jacobian is \(C^{g}/(Z^{g}+\tau Z^{g})\) and Abel--Jacobi pulls \(dz_{i}\) to \(\omega _{i}\). The positive volume of its fundamental parallelepiped is \(\operatorname{det} Im \tau \), so the squared norm of \(dz_{1}\wedge \ldots \wedge dz_{g}\) in the positive top-form convention is exactly \(\operatorname{det} Im \tau \), also the squared determinant norm of the \(\omega _{i}\). The top-form integral uses the usual positive ordering of real pairs; changing between block and interleaved complex wedges accounts for the orientation sign, not a positive metric scalar. No \nolinkurl{g!}, \(\pi ^{g}\), polarization degree, or isogeny degree occurs. A Kähler form raised to its gth power would require a factorial, but that is not this norm.

For any further number-field extension \nolinkurl{K'/K}, stable height is unchanged. Equivalently the normalized degree for the semistable model over \(K'\) again equals this same stable height. We need no assertion that the singular base change of a minimal regular model is already regular. In this argument choose one extension that simultaneously makes the curve semistable and splits the finite unramified fibre, then construct its model and \(E\) there. The degree in the height formula and the number of complex embeddings are both \nolinkurl{[K:Q]}. The auxiliary field is unrelated to the controlled field \(L'\) in the surface-isogeny argument.

\subsection{2. One global coordinate choice and its analytic bound}

Let \(\pi :X\to P^{1}_{t}\) be the same degree-d Belyi cover and let its split fibre over2 be \(P_{1},\ldots ,P_{d}\). Characteristic zero implies separability. Riemann--Hurwitz gives \(2g-2=-2d+\sum (e_{P}-1)\le -2d+3(d-1)=d-3\); over each of the three possible branch values the contribution is \(d\) minus the nonzero number of geometric points in that fibre. Hence \(d>2g-2\). A nonzero regular differential has zero divisor of degree \(2g-2\), including \(g=1\) where it has no zeros. Since every \(P_{j}\) is unramified, \(\omega /dt(P_{j})=0\) is precisely vanishing there. Thus

\(ev:E\otimes K\to K^{d}\), \(\omega \mapsto (\omega /dt(P_{j}))_{j}\)

has rank \(g\). Some \(g\times g\) minor is nonzero over \(K\). Fix its coordinate set \(I\subset \{1,\ldots ,d\}\) \textbf{once over \(K\)}. Each embedding sends that nonzero minor to a nonzero complex number, so the same \(I\) works at every embedding. An embedding permutes the labels of the geometric fibre but does not lose distinctness or unramifiedness. No separately chosen analytic minor is used to define the integral map.

The algebraic lane must supply the actual integral \(ev:E\to C^{\vee }\) for \(C=B/(t-2)B\), where \(B\) is the normalization algebra over \(O_{K}[t]\). Set \(F=pr_{I}(C^{\vee })\) with the standard Euclidean metric, and use the exact lattice bound \(\text{Ardeg}(F)\le (1/2)\operatorname{log} N \operatorname{disc}(C)\). Projection gives the generically invertible integral map \(\phi :E\to F\) used here. This packet does not replace \(F\) by the maximal order or its trace dual.

At every complex embedding the open disc \(D=\{|t-2|<1/2\}\) avoids \(0,1,\infty \). The inverse image consists of \(d\) disjoint sheets isomorphic to \(D\). On the sheet through \(P_{j}\), write \(\omega =f_{j}(t)dt\). The power-series identity, or the mean-value inequality, gives

\(|f_{j}(2)|^{2}\le 4/\pi  \int _{D} |f_{j}(t)|^{2} \text{dxdy}\le 4/\pi  ||\omega ||^{2}\),

because \((i/2)dt\wedge d\overline{t}=\text{dxdy}\). For completeness, expansion \(f_{j}(2+z)=\sum  a_{n} z^{n}\) on each smaller disc gives integral \(\pi \sum  |a_{n}|^{2} r^{2n+2}/(n+1)\ge \pi r^{2}|a_{0}|^{2}\); take \(r\uparrow 1/2\). The sheets are disjoint, so their positive integrals are parts of the integral on \(X\). This also rules out an unnoticed multiplicity in comparison with the global \(L^{2}\) norm.

In a Hermitian orthonormal basis of \(E_{\sigma }\), the row of the evaluation matrix at \(j\) represents the functional \(\omega \mapsto f_{j}(2)\). Its squared Euclidean norm is its dual Hermitian norm squared, bounded by \(4/\pi \). Hadamard therefore gives

\(||\operatorname{det} \phi ||_{\sigma }\le (2/\sqrt{\pi })^{g}\).

Here the target's standard coordinate basis is orthonormal; F's finite-place lattice does not alter its complex metric. One could obtain an operator-norm estimate using disjointness, but the displayed row/Hadamard proof already gives the exact constant consumed in \(v_{17}\). No strengthening is substituted downstream.

\subsection{3. Arithmetic degree, signs and consumer}

The determinant is a nonzero integral section \(s\) of \(L=\operatorname{det}(F)\otimes \operatorname{det}(E)^{-1}\). Its finite divisor is effective. The rank-one degree definition gives exactly

\(\text{Ardeg}(F)-\text{Ardeg}(E)=\operatorname{log} \#(L/O_{K} s)-\sum _{\sigma } \operatorname{log}||s||_{\sigma }\).

Thus \(\text{Ardeg}(E)\le \text{Ardeg}(F)+\sum _{\sigma } \operatorname{log}||\operatorname{det} \phi ||_{\sigma }\). This sign can also be checked by multiplication by \(m\) on the standard rank-one lattice: the finite term and the log norm both equal \nolinkurl{[K:Q]log|m|}. Inserting the two bounds and dividing by \nolinkurl{[K:Q]} proves precisely

\(\hF(\operatorname{Jac} X)\le \operatorname{log} N \operatorname{disc}(C)/(2[K:Q])+g \operatorname{log}(2/\sqrt{\pi })\).

The small positive constant is \(\operatorname{log}(2/\sqrt{\pi })=(1/2)\operatorname{log}(4/\pi )\); there is no extra \nolinkurl{g log g}, square-norm factor two, or half-weight for a conjugate embedding pair.

For the modular consumer, substitute \emph{the same} \(X=X_{H}=X_{1}(N)/H\), \(H\supset \{-1\}\), \(t=j/1728\), \(J_{H}=\operatorname{Jac} X_{H}\), \(d=d_{H}\), and \(g=g_{H}\) from Brox354--401. The modular-model and different lanes must establish that cover and the normalized discriminant bound \(d_{H}(\log 12+\operatorname{log} N+2\operatorname{log} d_{H})\). The result is exactly Brox453--488/H0.3, with no convention change when used in \nolinkurl{height/PROOF.md} H2--H4. In the actual non-CM realization \(H=\operatorname{ker} \epsilon \) with \(\epsilon \) even; the actual Jacobian contains an abelian surface over \(L'\), so \(g_{H}\ge 2\). That identity is not supplied by this metric lane.

There is no large-level, torsion-free-congruence-subgroup, \(\text{genus}\ge 2\), or absence-of-elliptic-points assumption in the evaluation proof. Every positive-genus modular curve for every positive \(N\) is covered, including genus1. Possible elliptic/cusp ramification is confined to the three omitted branch values; it cannot affect the disc over2. \(N=1\) and \(N=2\) give genus-zero curves (the covers have degrees1 and3); the modular section's explicit positive-genus hypothesis excludes them. Any other genus-zero level is excluded by that same hypothesis, and cannot be the actual surface-factor consumer since a zero Jacobian has no surface factor. No classification of small \(N\) is needed to cover all allowed positive-genus cases. If one extends the displayed modular inequality to genus0, its left side is \(\hF(0)=0\) and its right side is nonnegative, but that extension is not used in this contract.

\subsection{4. Scope and incorporation}

This is a source/application explanation of an already written argument; no numerical repair or manuscript edit is proposed. Historical justification was found in the actual earlier source assessment and report, including their withdrawn normalization concern, and in the current frozen \(v_{17}\) proof. We do not call this a newly repaired gap. The explanation can be incorporated by an exact pointer from an additive height-closure packet. Its author status, the still-separate algebraic/different inputs, and the need for fresh non-author acceptance must remain explicit. The 214 existence and the proposed256 work both consume the same full-family height supplier; the proof here is not restricted to a distinguished parameter \(\beta \).
\clearpage
% Source: research/quartic214/proofs/height-assembly.md
\section{Assembly at the actual modular-Jacobian consumer}\label{q214:height-assembly}
\noindent\textit{Source: }\nolinkurl{research/quartic214/proofs/height-assembly.md}.\par

\begin{quote}\footnotesize Author/integrator: \nolinkurl{/root}, 2026-11-10. This is an author proposal for a fresh non-author review. It supplements the two immutable packets at private commit \nolinkurl{a81afd7996a20e6b7bb5e48670418541c036d7aa}. The source/application target is Brox \(v_{17}\), SHA-256 \nolinkurl{2ba5dd497276073d92db95f0119bede82d182997c66375e5be294157de6d9f7c}, lines 47--488, and H0.3/H4 in \nolinkurl{research/quadratic-ten-contact-integration/height/PROOF.md}. No polynomial, parameter choice or constant is changed.\end{quote}

\subsection{A. The scoped theorem}

For a positive integer \(N\) and a subgroup \(H\) of \((Z/NZ)^{\times }\) containing \(-1\), put \(X_{H}=X_{1}(N)/H\), \(J_{H}=\operatorname{Jac}(X_{H})\), \(s=[(Z/NZ)^{\times }:H]\), and suppose its genus \(g\) is positive. For the map \(t=j/1728\) set

\(d=N s \prod _{p|N}(1+1/p)\).

The statement assembled here is

\(\hF(J_{H}) \le  d/2 (\log 12+\operatorname{log} N+2\operatorname{log} d) + g \operatorname{log}(2/\sqrt{\pi})\),

where \(\hF\) is the original stable Faltings height. This is the exact modular-Jacobian supplier already used by H0.3. It is a theorem about the actual quotient Jacobian, not a bound for an isogenous replacement or a full-level cover. Its proof is the following three linked project proofs, with their primary theorem applications recorded in their own contracts and source inventories.

\subsection{\(B\). Same curve, order, lattice and maps}

The modular-model construction in \nolinkurl{different/PROOF.md} \S{}4 identifies the generic primitive N-torsion points modulo \(H\) with the actual modular quotient, compatibly with Galois action. It provides the finite étale model away from 6N and the degree \(d\). The \(\pm 1\) ambiguity of a chosen elliptic-curve isomorphism is already killed in \(H\). At a prime dividing \(N\) we use the finite flat torsion argument, not the model with \(N\) inverted. Every component and every place is included in the discriminant sum. The positive-genus hypothesis removes the genus-zero cases without imposing an unrecorded large-level condition.

Choose one finite number field \(K\) defining the cover, making the curve semistable, and splitting its unramified fibre over \(t=2\). Existence of such an extension follows from semistable reduction and adjoining finitely many point coordinates. Its degree is not bounded. Let \(B\) be the normalization of \(O_{K}[t]\) in \(K(X_{H})\), and put \(C=B/(t-2)B\). The integral-evaluation proof uses this same \(B\) and this possibly nonmaximal order \(C\). Let \(E=H^{0}({\mathcal X},\omega _{{\mathcal X}/O_{K}})\) for the minimal regular semistable model \({\mathcal X}\) of \(X_{H}\). Then \nolinkurl{evaluation/PROOF.md} supplies the literal integral map

\(ev:E \to  C^{\vee }\), \(\omega  \to  (\omega /dt(P_{j}))_{j=1,...,d}\).

Here \(C^{\vee }\) is defined by the generic field trace, not by an arbitrary isomorphism of dualizing modules. The canonical finite-duality and specialization comparison fixes this normalization. The common regular model and codimension-one descent justify using the lattice of the regular semistable model even when the finite normalization model is singular. No claim that the latter is Gorenstein is used.

The metric proof fixes one coordinate set \(I\) of size \(g\) over \(K\) for which the evaluation minor is nonzero. It uses every fibre point to prove rank, then the same \(I\) at every embedding. Set \(F=pr_{I}(C^{\vee })\) with its standard Euclidean metric. The evaluation proof gives

\(\text{Ardeg}(F) \le  (1/2) \operatorname{log} \operatorname{Norm} \operatorname{disc}(C)\).

The metric proof supplies the original stable identification \(\text{Ardeg}(E)=[K:Q] \hF(J_{H})\) and the determinant bound \((2/\sqrt{\pi})^{g}\) at every embedding. Its arithmetic degree calculation therefore gives

\(\hF(J_{H}) \le  \operatorname{log} \operatorname{Norm} \operatorname{disc}(C)/(2[K:Q]) + g \operatorname{log}(2/\sqrt{\pi})\).

This chain keeps the same order, projected lattice, trace map, curve, metric and determinant convention throughout. It does not replace an integral lattice by its generic vector space.

\subsection{\(C\). Same discriminant and constant}

The different proof computes the trace determinant of the same \(B\) and its specialization to \(C\). Its local component formula is \(v(\operatorname{disc})=\sum _{j} f_{j} r_{j}\) with full residue degrees, including inseparable degrees. At \(p>3\) the torsion argument and intermediate-field comparison give the full vertical contribution at most \(d v_{p}(N)\). The possibly imperfect-residue estimate at 2 and 3 contributes at most \(d(\log 6+2\operatorname{log} d)\) after normalization. The horizontal \(t=0\) intersection contributes at most \nolinkurl{d log2}; \(t=1\) and infinity do not meet the section. Summing over all places of \(K\) cancels \nolinkurl{[K:Q]}, so

\(\operatorname{log} \operatorname{Norm} \operatorname{disc}(C)/[K:Q] \le  d(\log 12+\operatorname{log} N+2\operatorname{log} d)\).

Inserting this in \(B\) proves A with exactly its stated constant. In particular, the auxiliary \(K\) is never inserted into the surface-isogeny field degree \(d'=[L':Q]\) used later. \(K\) only occurs in a normalized stable-height calculation.

\subsection{\(D\). Actual H0.3 and the two downstream families}

The existing \(H_{1}\) constructs a final primitive weight-two form \(f\) with character \(\epsilon \), and \(H_{2}\) constructs an L'-defined surface-power subvariety of \(J=\operatorname{Jac}(X_{1}(N)/\operatorname{ker} \epsilon )\). Those constructions are retained inputs to the present application; this supplement does not re-prove normalized modularity. Weight two gives \(\epsilon (-1)=1\), so \(H=\operatorname{ker} \epsilon \) meets the subgroup hypothesis. Its index is \(\operatorname{ord}(\epsilon )=s\), and the actual surface factor gives \(g\ge 2\). Thus A applies to that exact \(J\) for every rational parameter in the full \(H_{0}\) domain. No restriction to the eventual distinguished \(\beta \) is introduced.

The genus bound \(g\le d\) used in \(H_{4}\) follows from Riemann--Hurwitz in the same Belyi cover: \(2g-2\le d-3\), hence \(g<d\) for positive \(g\). The remaining \(H_{4}\) relaxations, including \(\operatorname{log}(2/\sqrt{\pi})<\log 2\), are unchanged. They retain the same rational numerical certificate, conductor, character order and multiplicity assumptions. Consequently this supplement changes no value of \(C_{2}\), \(\alpha \), anchor exponent, 256-variable table or 214-variable parameter family. It supplies only the modular-Jacobian interface within their existing chains.

The 214 construction remains an existence family with an ineffective threshold \(N_{0}\); the separate condition \(N\ge 53\) only bounds coefficients. Neither this height supplier nor a passing finite checker makes \(N_{0}\) effective. No effective construction is asserted in this public package.

\subsection{\(E\). History and remaining closure}

The earlier source assessment and review already explained integral evaluation, trace duality, different bounds and the metric constant. This round compares those actual arguments to final primary sources and expands the joins. It does not claim a newly discovered arithmetic error or priority for the underlying proof. The old packets and verdicts stay immutable; the new reviewer decision applies only to the exact new supplement and its stated interfaces.

Acceptance of A does not accept \(H_{0}\) geometry/CM bounds, normalized modularity and companion trace assembly, the entire analytic/restricted-converse chain, or every remaining final-publication/source qualification in the prior foundation map. In particular the final 2022 Rémond body comparison remains a separate item. The current status overlay must carry those obligations forward rather than infer whole-proof closure from the success of this supplier.
\clearpage
% Source: research/quartic214/proofs/height-supplier.md
\section{Quadratic \(S_{10}\) contact supplier and its reflected effective consumers}\label{q214:height-supplier}
\noindent\textit{Source: }\nolinkurl{research/quartic214/proofs/height-supplier.md}.\par
\begin{quote}\small \textbf{Current application:} H0--H2 and H4 are consumed with the product averaging proof in Section~\ref{q214:height-averaging}.  The historical H3 quotient-filtration route is retained only as superseded evidence and is not used in this edition's 214 argument.\end{quote}

\begin{quote}\footnotesize Author: internal model task \nolinkurl{/root/sol61_height_foundation}, 2026-11-10. Author proposal, not independent acceptance. This selected H0--H4 text supplies the full-family height input for the 214 construction. Its scope is a complete connecting argument \textbf{relative to the inherited geometric, normalized modular-realization, modular-Jacobian height, root-height, scalar/anchor, rank and Lucas interfaces specified below}. It does not silently claim full published-source closure of those inherited proofs. Publication-version limitations are explicit in \nolinkurl{SOURCES.json} and \nolinkurl{FINDINGS.md.}\end{quote}

All heights \(h\) are absolute logarithmic Weil heights. All \(\hF\) are original stable Faltings heights. Tate modules use arithmetic Frobenius and \(\operatorname{det}(r_{f})=\epsilon *\chi _{\text{cyc}}\). Added algebraic scalar coefficients have trivial Galois action. Symbols for number fields, scalar characters and integer labels are kept distinct below.

\subsection{\(H_{0}\). Exact inherited interfaces and domain}

Let \(S=\{2,3,5,7,11,13,17,37,41,2399\}\), \(B_{S}=1857890364330\), and \(E(R)=\prod _{p|R}(1+1/p)\). For every finite rational \(a=u/v\) in lowest terms, \(v>0\), with \nolinkurl{a notin {0,1,2401,-343/32}}, define \(M=M_{S}(a)\) to count once each \(p\) outside \(S\) having positive odd valuation on one of

\nolinkurl{v,u,u-v,u-2401v,32u+343v}.

Put \(R=B_{S} M\) and \(\ell =\operatorname{log}\) \(R\). \(R\) is squarefree; \(M=1\) is allowed. The exact small-prime support argument is inherited from private\#8 cutoff proof \S{}1 and private\#9 conductor proof \S{}\S{}1--3. It uses the same full three-marked normalizer fiber as the parent manuscript, not a substituted surface: its branch sections are the five indicated contacts; outside210, even contact is unramified after an unramified extension, including infinity. The normal fiber field \(L\) has degree dividing576, is unramified outside supp(R), defines all geometric homomorphisms between the conjugate surfaces, and \(A_{a}\) has good reduction over \(L\) away from210. The same surface satisfies

\(h(a) \le  2 \operatorname{max}(0,\hF(A_{a}))+128\).                         (H0.1)

The CM branch gives \(\hF(A_{a})<854R\), over the same support and metric. These uniform source statements apply to every a above, not only the later constructed \(\beta \) or detector values. The non-CM surface has geometric endomorphism algebra the quaternion division algebra \(B\) of discriminant210, hence is absolutely simple. Its B-linear conjugate isogenies \(\mu _{\sigma }\) are L-defined, chosen with \(\mu _{1}=1\) and inflated from Gal(L/Q). The normalized sign cocycle and modularity-descent lemma are the actual parent \nolinkurl{quartic_v14.tex} passages481ff,8110ff; the latter lemma accepts any locally constant scalar splitting satisfying its cocycle equation. It identifies the \textbf{original} \(r_{5}\), with its finite twist absorbed into the final primitive weight-two form \(f\). The companion trace argument identifies the original \(r_{3}\) with that same \(f\); neither the absorbed finite twist nor its kernel field replaces the controlled lift or \(L'\).

The private\#9 conductor argument applies to every controlled lift that is tame at odd primes and has scalar Dirichlet conductor dividing4 at2. It proves for this final \(f\)

\(N \le  C_{N} R^{2}\), \(\operatorname{rad}(N)|R\), \(C_{N}=2^{68} 3^{16} 5^{16} 7^{4}\). (H0.2)

Here is why the freedom used below preserves that argument: the torsion field \(E_{\text{tor}}=L(A[11])\) is normal, \([E_{\text{tor}}:Q]\) divides \(576 |\operatorname{GL}_{4}(F_{11})|\), and the valuations of that integer at2,3,5,7 are15,4,4,1. Adjoining the kernel field of the scalar Dirichlet character and then the quadratic kernel of its lift adds at most two dyadic factors and only tame factors at odd primes. Thus \(v_{p}(e)\le 17,4,4,1\). The derivative/different calculation in private\#9 \S{}2 bounds the last positive upper break \(by2v_{p}(e)\), the Swan conductor of a two-dimensional representation \(by4v_{p}(e)\), and its tame codimension by2. Exponents70,18,18,6 follow. At \(p>7\), L/Q is tame, A/L has good reduction and the scalar lift is trivial on wild inertia; \(\text{exponent}\le 2\) follows. The original \(r_{5}\) is used at every \(p\ne 5\) and original \(r_{3}\) at5, so the level bound omits no prime. No bound on the tame part of \(e\) or the coefficient degree is substituted. This repeats the changed-lift connecting argument while retaining the parent published local-compatibility/modularity qualifications.

The retained modular-Jacobian project proof supplies, for \(H=\operatorname{ker} \epsilon \), \(J=\operatorname{Jac}(X_{1}(N)/H)\), genus \(g>0\) and actual j-map degree \(d\),

\(d=N s \prod _{p|N}(1+1/p)\), \(g\le d\), \(s=\operatorname{ord}(\epsilon )\),

\(\hF(J) \le  d/2(\log 12+\operatorname{log} N+2\operatorname{log} d)+g \operatorname{log}(2/\sqrt{\pi})\). (H0.3)

This is the differential-evaluation proof at Brox \(v_{17}\) Corollary4.3 (label cor:modularheight), not its large-primorial numerical theorem. It is retained as the precise versioned project height input already consumed by private\#9. It is valid after arbitrary finite extension; the auxiliary field splitting its evaluation fiber contributes no degree to the isogeny estimate. The local expanded proofs \nolinkurl{height-evaluation.md,} \nolinkurl{height-different.md,} \nolinkurl{height-metric.md} and \nolinkurl{height-assembly.md} supply this interface.

\subsection{\(H_{1}\). Minimal two-primary characters and the actual normalized field}

Let \(U(R)=\operatorname{max}(\{2\} \cup  \{2^{v_{2}(p-1)}: p \text{odd}, p|R\})\). At each odd prime where the sign cocycle has nontrivial local boundary, choose a generator \(g\) of the cyclic group modulo \(p\) and take the character sending \(g\) to a primitive \(2^{v_{2}(p-1)}\) root of unity. Since \nolinkurl{(p-1)/2} has exactly one fewer factor2, the character takes -1 to-1. Where the boundary is trivial choose the trivial local character. At2 choose the nontrivial character modulo4 only if its boundary is required. The product \(\chi \) has order \(m\) dividing \(U(R)\), is tame at every odd prime and has conductor dividing4 at2.

This is a choice inside the parent's instruction to choose an odd local character; it changes the representation and final \(f\) and therefore needs the fresh comparison here. The parent's local reciprocity criterion for the constant-coefficient squaring sequence says that the boundary is nontrivial exactly when the local character takes -1 to-1. Hence the new \(\chi \) has the same prescribed local boundaries as the original sign cocycle at every finite prime. The global Brauer injection/reciprocity argument used in the parent supplies the same sign-class equality, including the determined real invariant. A continuous pointwise square root, corrected by a sign cochain, gives

\(b^{2}=\chi \), \(\delta  b=\epsilon _{\text{sign}}\), \(b:G_{Q} \to  \mu _{2m}\).

Twist \(b\) by the parent's quadratic character to remove the finitely many unwanted inertia characters outside supp(R). This preserves \(b^{2}\) and \(\delta \) \(b\), stays in \(\mu _{2m}\), is tame at every odd added place, and adds possible support only at2, already in \(S\). On \(G_{L}\) the sign cocycle is trivial, so \(b\) is a character, invariant under conjugation by \(G_{Q}\). Its kernel field \(L'\) contains \(L\), is normal over \(Q\), and

\(d'=[L':Q] \le  2m[L:Q] \le 1152 U(R) \le 2304 \phi (R)\).    (H1.1)

The last deliberately conservative bound will be used numerically. Since the isogenies are inflated from Gal(L/Q) and \(\mu _{1}=1\), both \(\mu _{\sigma }\) and \(a_{\sigma }=(\operatorname{deg} \mu _{\sigma })^{-1/4}\) equal1 on \(G_{L}\). Consequently

\(\rho _{l}(\sigma )=b(\sigma )^{-1} a_{\sigma } \mu _{\sigma } \circ \sigma \)

restricts on \(G_{L'}\) to the \textbf{actual Tate representation of this \(A_{a}\)}, not an unspecified scalar twist. The normalized cocycle identity makes \(\rho _{l}\) a representation; the B-action gives \(\rho _{l}=r_{l} \oplus  r_{l}\). All geometric endomorphisms remain defined over \(L'\). Faltings semisimplicity and the commutant \(B \text{tensor} \text{Qbar}_{l}=M_{2}(\text{Qbar}_{l})\) imply that \(r_{l}|G_{L'}\) is absolutely irreducible. The parent modularity lemma applies with \(t_{\sigma }=a_{\sigma }/b(\sigma )\); changing \(\chi \) has preserved every hypothesis, including finite local constancy. It supplies a final primitive weight-two \(f\) with \(r_{5}=r_{f,5}\). The companion trace comparison supplies the same identity at3. For the multiplicity argument below one prime, say5, suffices; no assertion at a missing coefficient prime is needed.

On the fourth exterior power an isogeny acts by its degree, and Galois acts by \(\chi _{\text{cyc}}^{2}\). The normalization \(a_{\sigma }^{4} \operatorname{deg}(\mu _{\sigma })=1\) therefore gives \(\operatorname{det} \rho _{l}=b^{-4} \chi _{\text{cyc}}^{2}=\chi ^{-2} \chi _{\text{cyc}}^{2}\). On the other hand the modular decomposition gives \(\operatorname{det} \rho _{l}=\epsilon ^{2} \chi _{\text{cyc}}^{2}\). Thus

\((\epsilon  \chi )^{2}=1\) and \(s=\operatorname{ord}(\epsilon ) \mid  U(R)\). (H1.2)

Indeed \(U(R)\) is a power of two and is even; both \(\chi \) and \(\epsilon *\chi \) have orders dividing it. Weight two gives \(\epsilon (-1)=1\), so -1 belongs to \(H\). Finally \(U(R)|\phi (R)\) and

\begin{gather*}
s E(R) \le  U(R)E(R) \le  \phi (R)E(R)\\
= R \prod _{p|R}(1-p^{-2}) < R.\tag{H1.3}
\end{gather*}

No uniform bound on \(U(R)\) or on the coefficient field is asserted. The finite twist absorbed while choosing \(f\) is not identified with \(\chi \); determinant(H1.2) uses the original normalized representation itself.

\subsection{\(H_{2}\). Multiplicity over \(L'\), not only geometric multiplicity}

Let \(E_{f}=Q(a_{n}(f): n\ge 1)\), \(e=[E_{f}:Q]\). The modular abelian variety \(A_{f}\) is defined over \(Q\), has dimension \(e\), and over \(\text{Qbar}_{5}\) its Tate module is the sum of the two-dimensional coefficient conjugates \(r_{f}^{\tau }\). This is the standard modular-factor input: published Shimura Theorem1 gives the quotient, dimension and differentials; its Euler-factor assertion requires the coefficient-conjugate span to be Fricke-stable. For the final primitive weight-two newform this condition is supplied by the Li Fricke formula in the actual parent modularity argument. Thus the formula applies to this \(f\) with arbitrary nebentypus, not only to a \(\Gamma _{0}\) form. The Tate/cohomology dual convention is the same as that parent's modular identification. This is an explicit hypothesis match, not an appeal to a similarly named modular variety.

At every good place \(v\) of \(L'\) away from5 and the finitely many bad rational primes, \(\operatorname{tr}(r_{f}(\text{Frob}_{v}))=\operatorname{tr}(V_{5}(A_{a})(\text{Frob}_{v}))/2\) is rational. The trace on the left belongs to \(E_{f}\) before any coefficient embedding: if \(v\) has residue degree \(f_{v}\) over rational \(p\), it is the trace of the \(f_{v}-th\) power of a rational-prime Frobenius, a polynomial in \(a_{p}\) and \(p\) \(\epsilon (p)\). The equality in the chosen embedding at5 to a rational number therefore identifies that \(E_{f}\) element with the same rational number. Every coefficient embedding fixes it. Chebotarev and semisimple character uniqueness identify every \(r_{f}^{\tau }|G_{L'}\) with the same absolutely irreducible restricted representation \(W\). Its semisimplicity also follows directly from that of the abelian variety Tate modules; a finite-index restriction here introduces no unproved extension-class replacement.

Consequently the Tate modules of \((A_{f})_{L'}^{2}\) and \(A_{a}^{e}\) have identical semisimple characters, already over \(Q_{5}\). Faltings' isogeny theorem gives an L'-isogeny. In the semisimple L'-isogeny category, absolute simplicity of \(A_{a}\) forces every simple factor of \((A_{f})_{L'}\) to be isogenous to \(A_{a}\) over \(L'\). Comparing dimensions proves \(e\) is even and

\((A_{f})_{L'} \text{ is }  L'-\text{isogenous} \text{ to }  A_{a}^{r}\), \(r=e/2\).         (H2.1)

The Hecke relation \(a_{p}^{2}-a_{p^{2}}=p \epsilon (p)\) puts \(\epsilon (p)\) in \(E_{f}\) for every \(p\) prime to \(N\). Chebotarev for the finite cyclotomic quotient supplies primes in every invertible residue class modulo \(N\), hence the complete image of \(\epsilon \) lies in \(E_{f}\). Therefore \(Q(\zeta _{s}) \subset  E_{f}\), so \(r\ge \phi (s)/2\). If \(s\ge 4\) is a power of two, \(\phi (s)=s/2\); if \(s=1\) or2, use \(r\ge 1\). In all cases

\(s/r \le 4\).                                         (H2.2)

Every coefficient conjugate of \(f\) has the same diamond-character kernel \(H\), so all its differentials are fixed by \(H\). For the quotient map \(X_{1}(N)\to X_{H}\), pullback identifies the Jacobian of \(X_{H}\) up to isogeny with the H-invariant subvariety. The \(A_{f}\) quotient restricted to that invariant subvariety is surjective, as seen on those differentials; characteristic zero makes this a genuine surjective homomorphism, not merely an equality of dimensions. Poincare reducibility over \(Q\) supplies a subvariety of \(J\) isogenous over \(Q\) to \(A_{f}\). Base change and(H2.1) supply an \textbf{L'-defined} \(B_{0}\) inside \(J_{L'}\) isogenous over \(L'\) to \(A_{a}^{r}\). In particular \(g\ge 2r\ge 2\). The field splitting a modular evaluation fiber and the coefficient value field are not substituted for \(L'\).

\subsection{\(H_{3}\). Historical averaging route (superseded in this public application by height-averaging.md)}

Take an L'-isogeny \(A_{a}^{r} \to  B_{0}\), and let \(D_{0}=0 \subset  D_{1} \subset  ... \subset  D_{r}=B_{0}\) be the images of its first \(i\) factors. Each \(C_{i}=D_{i}/D_{i-1}\) is an L'-defined surface L'-isogenous to \(A_{a}\). Rémond's stable subvariety/quotient inequality and Bost's lower bound \(\hF(C)\ge -3 \operatorname{dim}(C)/2\) give

\begin{gather*}
\sum _{i} \hF(C_{i}) \le  \hF(B_{0})\\
\le  \hF(J)+3(g-2r)/2 \le  H_{J},
\end{gather*}

where \(H_{J}=\operatorname{max}(1,\hF(J)+3g/2)\). Thus one actual \(C_{i}\) has \(\hF(C_{i})\le H_{J}/r\). No invariance of \(\hF\) under isogeny is assumed.

Apply Gaudron--Rémond Theorem1.9(1) with reference variety \(C_{i}\) of dimension2, base field \(K=L'\), and extension equal to \(K\). The published bound gives an L'-isogeny \(C_{i}\to A_{a}\) of degree at most

\([14^{32} d' \operatorname{max}(1,\operatorname{log} d',\hF(C_{i}))]^{16}\).

The stable height increase along this direction is at most half the logarithm of that degree. Consequently, for any \(X\ge H_{J}/r\) with \(X\ge 1\),

\(\hF(A_{a}) \le  X+8[32\log 14+\operatorname{log} d'+\operatorname{log} \operatorname{max}(1,\operatorname{log} d',X)]\). (H3.1)

The small surface \(C_{i}\), its quotient maps and the final isogeny are all over \(L'\). The growing multiplicity \(r\) is never used as an isogeny dimension. The published GR theorem's original stable normalization and both isogeny directions were checked directly in its final publisher PDF, pp7--8. Rémond's authorJune2021 source explicitly states the stable version after Corollary1.3; its final2022 journal body remains unavailable through the subscription redirect. This is a source-version coverage limitation, not a claim that the known inequality is false or unpublished.

\subsection{\(H_{4}\). \(S_{10}\) quadratic constant, all \(M\ge 1\)}

Set \(k=9339/1000\), \(\delta =10^{-6}\), \(c_{0}=8k+\delta =74712001/1000000\), and \(T=C_{N} R^{2} E(R)\). By(H0.2),(H1.3), \(\log N\le \log C_{N}+2\ell \), \(\operatorname{log} d<\log C_{N}+3\ell \). Using(H0.3), \(g\le d\) and \(\operatorname{log}(2/\sqrt{\pi})<\log 2\), define

\(K_{0}=(\log 12+3\operatorname{log} C_{N})/2+\log 2+3/2\).

The rational interval certificate proves \(4+K_{0}/\operatorname{log} B_{S} < k\). It follows that \(H_{J}\le k d \ell \), including its max with1, because \(k\) \(d\) \(\ell >1\). Equations(H2.2),(H0.3) then give

\(H_{J}/r \le 4k T \ell  = X\).

The same certificate proves \(\ell \ge \log B_{S}>28\), \(\log 2304<\log B_{S}\), \(\operatorname{log}(4k)+\log C_{N}<4\log B_{S}\). Hence(H1.1) gives \(\log d'<2\ell \); and \(E(R)<R\), \(\log \ell \le \ell \) give \(\log X<8\ell \). Also \(X>\operatorname{max}(1,\log d')\). From(H3.1), \(\log 14<3\) \(\text{ and } 768<28\ell \),

\(\hF(A_{a}) < X+768+80\ell  < X+108\ell \).

Combining with(H0.1), \(216+128/\log B_{S} < \delta  C_{N} B_{S}^{2}\), gives

\(h(a) < c_{0} C_{N} R^{2} E(R) \log R\).                       (H4.1)

The CM bound also satisfies(H4.1): \(1708B_{S}+128<c_{0} C_{N} B_{S}^{2} \log B_{S}\), and the right side grows at least quadratically with \(R\) whereas1708R+128 is linear. Thus(H4.1) is uniform across both branches.

Define the exact positive real constant

\(C_{2}=c_{0} C_{N} B_{S}^{2} E(B_{S}) \log B_{S}\).

Since \(M\) is coprime to \(B_{S}\), \(E(R)=E(B_{S})E(M)\), and \(\log R=\log B_{S}+\log M\). Equation(H4.1) becomes

\(h(a) < C_{2} M^{2} E(M) [1+\log M/\log B_{S}]\).                (H4.2)

The portable rational certificate encloses each logarithm by an atanh series with a proved geometric tail. It proves strictly

\(C_{2} < 2^{7557/32}\).

Every \(p\) outside \(S\) is at least19. With \(\theta _{0}=\operatorname{log}(20/19)/\log 19\), the positive ratio \(\operatorname{log}(1+1/p)/\log p\) is at most \(\theta _{0}:\) its numerator decreases and denominator increases for \(p\ge 19\). Thus \(E(M)\le M^{\theta _{0}}\). Also \(1+\log M/\log B_{S} \le  M^{1/\log B_{S}}\) \(by1+x\le \operatorname{exp}(x)\). The certificate proves \(\theta _{0}+1/\log B_{S}<55/1024\). Consequently, including \(M=1\),

\(h(a) < C_{2} M^{\alpha }\), \(\alpha =2103/1024\).              (H4.3)

This is a fresh small-support calculation. The false large-primorial shortcut \(C_{N}^{2}<B_{S}\) is never used.

\clearpage
% Source: research/quartic214/proofs/height-averaging.md
\section{Published product route for the multiplicity saving}\label{q214:height-averaging}
\noindent\textit{Source: }\nolinkurl{research/quartic214/proofs/height-averaging.md}.\par

\begin{quote}\footnotesize Author: \nolinkurl{/root}, 2026-11-11. This is a new application proof for independent review, not an acceptance record. It changes the supplier and auxiliary surfaces in \(H_{3}\) of the frozen private quadratic-height proof. It preserves the inequality H3.1, H4's constants and all polynomial rows.\end{quote}

\subsection{Sources and historical disposition}

The original \(H_{3}\) applies the stable subvariety/quotient inequality in Rémond, \emph{Propriétés de la hauteur de Faltings}, Annali SNS Pisa 23 (2022), 1589--1596. Its June 2021 author version was read, including Corollary 1.2 and the stable-height statement after Corollary 1.3. On 2026-11-11 the final publisher body again redirected to subscriptions. That final-body comparison remains unavailable; no claim about its absence or incorrectness follows.

There is an existing alternative argument in the current public Version 15, \nolinkurl{v15:height:lem:surface}, especially lines 30617--30636. \(I\) read that argument and its recorded source-route disposition. It consumes the construction in the published proof of Gaudron--Rémond, \emph{Nouveaux théorèmes d'isogénie}, Mémoires SMF 176 (2023), Lemma 16.2(1), printed pp. 97--98. \(I\) independently downloaded the final publisher PDF, SHA256 \nolinkurl{b1e2e32bec626cf6d27fb7d6d8702b5c707e0c3dc002d979088e62d8cda514a4}, and read the text and image of printed \(p\). 98. The source constructs, over its given number field, an isogenous abelian variety of no larger original stable Faltings height which is a product of simple abelian varieties over that field. Its factors are explicitly allowed to repeat.

This published proof uses Rémond 2022, Corollary 1.3, and Rémond 2017, Theorem 1.6. They remain indirect dependencies. The present route does not assert that the inaccessible final 2022 text was read, and does not re-prove the published product supplier. It uses that construction at its exact stated scope. Public Version 15's accepted single-surface application is historical support for the route, not approval of the new multiplicity application below.

\subsection{Fixed-field product and multiplicity}

Retain the actual objects of H0--H2: the number field \(L'\), the absolutely simple abelian surface \(A=A_{a}\), the Jacobian \(J=\operatorname{Jac}(X_{1}(N)/\operatorname{ker} \epsilon )\) over \(L'\), its dimension \(g\), and an L'-defined subvariety \(B_{0}\) isogenous over \(L'\) to \(A^{r}\), with \(r=e/2\ge 1\). In particular \(g\ge 2r\). The auxiliary field used to split the evaluation fibre is not substituted for \(L'\). Geometry and the \(H_{2}\) realization are premises of this lemma and have their own source/application obligations.

Apply the construction in the proof of GR2023 Lemma 16.2(1) to \(J\) over \(L'\). It gives \(J_{0}=\prod \) \(D_{j}\), where every \(D_{j}\) is simple over \(L'\), \(J_{0}\) is L'-isogenous to \(J\), and \(\hF(J_{0})\le \hF(J)\). This is an isomorphism to a product, not merely an isogeny between the displayed factors and \(J_{0}\). No invariance of height under arbitrary isogeny is asserted.

Poincaré complete reducibility over \(L'\) and uniqueness of simple isogeny factors imply that the isogeny class of A occurs in this product with multiplicity at least \(r\). To see the multiplicity explicitly, \(B_{0}\sim A^{r}\) is a subobject of \(J\) in the semisimple category up to L'-isogeny. Passing to the isomorphic object \(J_{0}\) in that category preserves the multiplicity of every simple constituent. Since A is absolutely simple, it is simple over \(L'\); each constituent in its class has dimension two. Choose \(r\) distinct occurrences, name them \(C_{1},...,C_{r}\), and denote the product of all remaining factors by \(E\). Then

\begin{gather*}
J_{0} = C_{1} \times  ... \times  C_{r} \times  E,\\
\operatorname{dim}(E)=g-2r,\\
C_{i} \sim  A \text{ over }  L' \text{ for }  \text{ every }  i=1,...,r
\end{gather*}

Every index is covered, even when \(r\) grows with the nebentypus order. These \(C_{i}\) are new product factors, not the old filtration quotients \(D_{i}/D_{i-1}\). The only subsequent consumer of those old quotients was their dimension, L'-isogeny class and height bound. Those are supplied for the new factors here; no assertion identifies the two collections literally.

\subsection{Height average and transfer}

Use the original stable Faltings normalization throughout. Height is additive on products. The original-normalization lower bound is \(\hF(E)\ge -3\) \(\operatorname{dim}(E)/2\), from the actual GR2014 \S{}2.3 and Corollary 8.4 calculation recorded in \nolinkurl{quartic214-protocol-d/PRIMARY_SOURCE_BRIDGES.md}. It holds without a polarization hypothesis; if \(\operatorname{dim}(E)=0\) its height is zero. Consequently

\begin{gather*}
\sum _{i} \hF(C_{i}) = \hF(J_{0})-\hF(E)\\
\le  \hF(J)+3(g-2r)/2\\
\le  H_{J},\\
H_{J}=\operatorname{max}(1,\hF(J)+3g/2)
\end{gather*}

At least one of these actual \(r\) surfaces therefore has \(\hF(C_{i})\le H_{J}/r\). No assertion about a minimum on a varying field, or about average height over isogeny classes, is needed beyond the published product construction.

For this \(C_{i}\) apply GR2023 Theorem 1.9(1) over \(K=L'\), with extension \(K\) itself, source \(C_{i}\) and target A. Both are dimension-two varieties over the same field and are L'-isogenous. With \(d'=[L':Q]\), the degree bound in the final publisher statement gives an isogeny \(C_{i}\) \(\to \) A of degree at most

\begin{gather*}
[14^{32} d' \operatorname{max}(1,\operatorname{log} d',\hF(C_{i}))]^{16}
\end{gather*}

The one-sided stable-height change is at most half the logarithm of the degree in this direction. Thus for any \(X\ge \operatorname{max}(1,H_{J}/r)\),

\begin{gather*}
\hF(A) \le  X+8[32 \log 14+\operatorname{log} d'+\operatorname{log} \operatorname{max}(1,\operatorname{log} d',X)]
\end{gather*}

This is exactly H3.1. The later calculation in \(H_{4}\) uses only this displayed inequality, the \(H_{2}\) bound \(s/r\le 4\) and its separately supplied modular-Jacobian estimate. Neither \(H_{4}\) nor its rational certificate uses the old quotient presentation. The full-family domain, both CM/non-CM branches and the two transformed arguments used in 214 are unchanged. This proof does not itself accept those receiving applications or their other premises.

\subsection{Nine application obligations}

\begin{itemize}
\item 1. \textbf{Source hypotheses.} GR2023's construction is over a number field and begins with an abelian variety over that field. Here they are the actual \(L'\) and \(J\) of \(H_{2}\). The later isogeny theorem requires L'-isogenous dimension-two varieties: these are \(C_{i}\) and A, not \(J\) and A.
\item 2. \textbf{Constructed objects.} \(J_{0}\) and all simple product factors are defined over \(L'\). The realization of \(B_{0}\sim A^{r}\) inside \(J\) is retained as a separate \(H_{2}\) premise. The GR construction supplies the height comparison, not an informal change to an arbitrary isogenous product.
\item 3. \textbf{Index-set coverage.} Select \(r\) separate occurrences in the A-isogeny class. The entire family \(i=1,...,r\) participates in the sum; the repeated factors permitted by the source cannot be replaced by the set of distinct isogeny classes.
\item 4. \textbf{Maps.} Pass along the L'-isogeny \(J_{0}\sim J\) only in the isogeny category to count constituents. Use the actual product isomorphism for height additivity. Apply the final small isogeny in direction \(C_{i}\) \(\to \) A, with half-logarithmic cost retained.
\item 5. \textbf{Consumed object.} H3.1 bounds the same original \(A_{a}\). The new \(C_{i}\) replace auxiliary quotient surfaces; their identity with the old \(C_{i}\) is neither needed nor claimed. The received \(X\) and \(H_{J}/r\) are exactly those consumed in \(H_{4}\).
\item 6. \textbf{Conventions.} All heights are original stable \(\hF\); all logarithms are natural; \(d'\) is the degree of \(L'\), not the evaluation splitting field. No Tate/cohomology change is made.
\item 7. \textbf{Local conditions.} These published height/isogeny statements impose no new conductor, good-prime or integral-model constraint on the selected surfaces. Existing H0--H2 local hypotheses remain upstream; enlarging \(L'\) is not used.
\item 8. \textbf{Exceptions.} \(r\ge 1\) and \(g\ge 2r\); \(E\) of dimension zero is included. Absolute simplicity is required in the non-CM \(H_{2}\) branch. The original separate CM estimate does not use this lemma. Negative individual heights are allowed, and the maximum with 1 is retained.
\item 9. \textbf{Downstream dependencies.} Only the supplier for the \(H_{3}\) average and the auxiliary \(C_{i}\) change. H3.1 and all constants are identical, so downstream numerical arguments and rows need no modification, subject to review of this substitution. The indirect Rémond dependencies and final-body access limitation remain recorded.
\end{itemize}

Acceptance, if obtained, must bind this proof, these nine obligations, the actual \(H_{2}\) construction, H3/H4 consumer and published source version. The author makes no whole-214 Protocol \(D\) completion claim.
\clearpage
% Source: research/quartic214/proofs/primary-source-bridges.md
\section{Published-source bridges for the inherited multiplicity argument}\label{q214:primary-source-bridges}
\noindent\textit{Source: }\nolinkurl{research/quartic214/proofs/primary-source-bridges.md}.\par

\begin{quote}\footnotesize Author: \nolinkurl{/root}, 2026-11-10. This is a source/application supplement to the unchanged \nolinkurl{../quadratic-ten-contact-integration/height/PROOF.md}, H1--H3. It is author work awaiting the separately recorded review. It closes only the identified source-access and application records; it does not prove the inherited geometric or modular-realization premises.\end{quote}

\subsection{Faltings: the actual field and Tate objects}

\(I\) retrieved and read the published second edition of \emph{Rational Points} (1986), including the printed title/copyright pages and Faltings, ``Complements to Mordell'', Theorem 1(a)--(b), printed p.204 (PDF page 212). \(I\) inspected the page image as well as its extracted text. The exact downloaded PDF has SHA-256 \nolinkurl{3122011c0f12cc991fc0ff1601bb77cc1cd3e5ab2a0db023c0bf224ea576b6aa}; access URL: \url{https://yhzhumath.github.io/Mordell2024/RP.pdf}. This is a scan of the published edition, not an identification inferred from a search result. The publisher's edition identifier is \url{https://doi.org/10.1007/978-3-663-06812-9}.

The theorem concerns an abelian variety over a finitely generated field of characteristic zero and any prime \(\ell \). It gives semisimplicity of the rational Tate representation and the endomorphism comparison (indeed the printed statement gives the integral comparison). Apply it with the \textbf{same} number field \(L'\) constructed in \(H_{1}\) and \(\ell =5\). A number field is finitely generated over \(Q\); after inverting its discriminant, its ring of integers supplies the smooth finitely generated Z-algebra in the source's setup. This applies to \(A_{a}\), \((A_{f})\) over \(L'\), and their finite products. No field obtained by splitting an evaluation fiber is substituted for \(L'\).

\(H_{1}\) has shown that \(\rho _{5}\) restricted to \(G_{L'}\) is the actual covariant Tate action of \(A_{a}:\) the isogeny and scalar normalizations are 1 there. Under the retained assertion that all geometric endomorphisms are defined over \(L'\), the rational endomorphism algebra is \(B\). After extending scalars to \(\text{Qbar}_{5}\) it becomes \(M_{2}(\text{Qbar}_{5})\). The semisimple representation \(V_{5}(A_{a})\) therefore has one absolutely irreducible constituent \(W\) with multiplicity two: the endomorphism algebra of a semisimple representation is a product of matrix algebras, and here that product is the single \(M_{2}\). This is exactly the \(W\) consumed in \(H_{2}\), rather than a finite scalar twist. Semisimplicity remains valid after the coefficient extension because the finite-dimensional semisimple algebra generated by the image is separable in characteristic zero.

For completeness, the isogeny conclusion in \(H_{2}\) follows from the stated endomorphism theorem, without needing a differently normalized ``isogeny theorem''. Apply \((b)\) to \(X\) times \(Y\) and use the two factor idempotents. The off-diagonal block gives \(\operatorname{Hom}_{L'}(X,Y)\) tensor \(Q_{5}\) \(=\) \(\operatorname{Hom}_{{G}_{L'}}(V_{5}(X),V_{5}(Y))\). Take \(X=((A_{f})_{L'})^{2}\) and \(Y=A_{a}^{e}\). H2's trace and semisimple-character comparison identifies their rational Tate representations; that comparison remains a separately stated input. The invertible maps form a nonempty open subset of this Hom space. Rational combinations of a basis of \(\operatorname{Hom}_{L'}(X,Y)\) tensor \(Q\) are dense in its \(Q_{5}\) extension, so one such combination is invertible on Tate modules. Clear its denominators. Its image has full dimension, and its kernel has dimension zero, hence it is an L'-isogeny. This argument neither assumes integral Tate lattices agree nor introduces an extra field extension.

\subsection{Poincare reducibility over the stated ground field}

\(I\) read the author-hosted scan of Milne, ``Abelian varieties'', \emph{Arithmetic Geometry} (1986), Proposition 12.1 and the following decomposition paragraph, printed p.122 (PDF page 20). Its PDF SHA-256 is \nolinkurl{86259344a576e4ef75a9d81e9bf43143f872a72b14cf8fcc2703f3aea6349136}; source: \url{https://www.jmilne.org/math/articles/1986b.pdf}. This scan is image-only; the page image was inspected directly. The theorem constructs a complement using an ample sheaf and the reduced connected kernel of the associated dual morphism. Those data are over the original field, so the complement is over that field, not just over its algebraic closure.

For the surjection \(J\) over \(Q\) to \(A_{f}\) already supplied in \(H_{2}\), let \(K\) be the connected kernel. Apply Proposition 12.1 to \(K\) inside \(J\) over \(Q\). Its complement \(B'\) maps onto \(A_{f}\) with finite kernel; thus \(B'\) is a Q-defined subvariety Q-isogenous to \(A_{f}\). Base change to \(L'\) gives the \(B_{0}\) used in H2--H3. Separately, apply the same proposition over \(L'\) to decompose varieties into simple isogeny factors. Since \(A_{a}\) is absolutely simple under \(H_{0}\) and \(((A_{f})_{L'})^{2}\) is L'-isogenous to \(A_{a}^{e}\), every simple factor of \((A_{f})_{L'}\) is L'-isogenous to \(A_{a}\). Dimensions give \(e=2r\) and \((A_{f})_{L'}\) is L'-isogenous to \(A_{a}^{r}\). The division-algebra and absolute-simplicity premise remains H0's actual surface statement; Poincare reducibility does not supply that premise.

\subsection{Nine application obligations}

\begin{longtable}{@{}p{0.23\linewidth}p{0.73\linewidth}@{}}
\toprule
Obligation & Required input, constructed input, and connecting check\\
\midrule
\endhead
Source hypotheses & Faltings requires a characteristic-zero finitely generated field and an abelian variety. \(L'\) is the finite number field from \(H_{1}\) and all varieties above are defined over it. Milne applies to an abelian subvariety over its ground field; \(K\) is the connected kernel over \(Q\).\\
Constructed objects & The objects are \(A_{a}\), \((A_{f})_{L'}\), their products and the actual \(J=\operatorname{Jac}(X_{1}(N)/\operatorname{ker} \epsilon )\); the proof above never replaces \(A_{a}\) by a twist.\\
Domain and index coverage & Every rational a admitted by \(H_{0}\) is covered, separately in the non-CM branch. All \(e\) coefficient conjugates enter the unchanged trace comparison. The CM branch does not use the multiplicity argument.\\
Transition and comparison maps & Restriction is to \(G_{L'}\); factor idempotents recover Hom from End of a product. Rational approximation and denominator clearing yield an isogeny over \(L'\). The kernel complement is first constructed over \(Q\), then base changed.\\
Identity of the consumed object & The resulting \(W\), \(A_{a}^{r}\) and \(B_{0}\) are exactly those named in \(H_{2}\) and the filtration of \(H_{3}\). The dual cohomology convention is not substituted for the covariant Tate convention.\\
Coefficient and normalization conventions & The initial comparison is rational \(Q_{5}\), with \(\text{Qbar}_{5}\) used only to identify simple constituents. Integral lattice equality is not claimed. No Faltings-height invariance under isogeny is inferred.\\
Local conditions & The source theorems need no good-reduction or semistability hypothesis at every place. The good-place trace comparison is the separate \(H_{2}\) argument. \(\ell =5\) is a prime different from characteristic zero.\\
Exceptional cases & Products have matching positive dimension. H0's non-CM absolute simplicity is required; the CM branch remains separate. The complement may have a finite kernel, which is sufficient.\\
Downstream consumption & These identifications justify the L'-defined surfaces \(C_{i}\) in \(H_{3}\), conditional on its other inputs. The Rémond averaging inequality, modular-Jacobian height bound and normalized modularity are not proved by this supplement.\\
\bottomrule
\end{longtable}

\subsection{Li's Fricke condition for the modular quotient}

\(I\) read Li, ``Newforms and functional equations'', \emph{Mathematische Annalen} 212 (1975), pp.294--296, in Göttingen's digitized published volume \url{https://gdz.sub.uni-goettingen.de/id/PPN235181684_0212}. The library's IIIF manifest labels canvas 00000304 as printed p.296. \nolinkurl{PRIMARY_SOURCE_ACCESS.json} records the exact public page-image URLs and hashes. Page294 defines normalized primitive newforms; page296 defines the Fricke operator \(H_{M}\) and coefficient-conjugation operator \(K\) and, immediately before Lemma6, gives \(F|H_{M}=\gamma \) F|K for such a form. The character is arbitrary; no trivial-nebentypus assumption occurs there.

Use Li's \(M=N\) from \(H_{2}\) and Li's second level parameter 1. The final form \(f\) in \(H_{1}\) is the primitive weight-two form, after the permitted finite twist has been absorbed; it is that same final form to which the formula applies. For each coefficient embedding \(\tau \), \(f^{\tau }\) is a normalized primitive form at the same level with character \(\epsilon ^{\tau }\), as in the inherited modular-factor construction. Its Fricke image is a nonzero scalar times its coefficient conjugate. Conjugating the embedding gives another embedding of the coefficient field into \(C\), so the latter form is in the span of all coefficient conjugates. Fricke is invertible, hence this span is Fricke-stable. Any nonzero scalar difference between normalized slash conventions preserves that span; no numerical value of \(\gamma \) is consumed. This supplies the specific Fricke-span hypothesis of the already checked published Shimura Theorem1 used in \(H_{2}\). It does not prove the existence/modularity of \(f\), nor replace the independent normalized modularity application in H0--H1.

The nine obligations for this identification are: (1) a normalized primitive weight-two form with finite character, supplied by \(H_{1}\); (2) the same \(f\) and all its coefficient conjugates, not an untwisted predecessor; (3) every coefficient embedding \(\tau \), including conjugate embeddings; (4) the operator \(H_{N}\) and coefficient conjugation \(K\), with no discarded scalar affecting span; (5) the exact stable span consumed by Shimura; (6) Li's \((M,1)=(N,1)\) and complex coefficient conventions, with only span used; (7) no added good-prime condition or character-conductor restriction; (8) nonzero forms and invertible Fricke, hence nonzero \(\gamma \), including nontrivial nebentypus; and (9) the resulting modular quotient/Tate constituent in \(H_{2}\), with the retained Galois-conjugacy and modularity premises explicit. This is a source/application bridge awaiting the new independent review, not a new assertion that those premises were all audited here.

\subsection{Original-normalization lower bound used in \(H_{3}\)}

\(I\) also checked the publisher PDF of Gaudron--Rémond, ``Théorème des périodes et degrés minimaux d'isogénies'', \emph{Commentarii Mathematici Helvetici} 89 (2014), 343--403, DOI \url{https://doi.org/10.4171/CMH/322}, from \url{https://ems.press/content/serial-article-files/43351?nt=1}. The precise downloaded hash is in \nolinkurl{PRIMARY_SOURCE_ACCESS.json}. Definition2.1 and \S{}2.3, printed pp.351--352, define a stable height \(h\) and identify \(\hF=h-(\operatorname{dim} A/2)\operatorname{log}\) \(\pi \). Corollary8.4, printed p.400, proves \(h(A)\ge -(\operatorname{dim} A/2)\operatorname{log}(2\pi )\) without a polarization hypothesis. \(I\) checked both page images. Combining the two gives \(\hF(A)\ge -\operatorname{dim}\) A \(\operatorname{log}(\pi  \sqrt{2})\ge -3\operatorname{dim}\) A/2; the final numerical relaxation is also stated on p.352.

This applies over the same number field \(L'\) to every quotient \(C_{i}\) and \(J/B_{0}\) in \(H_{3}\). These are abelian varieties defined over \(L'\), with no principal polarization or new ramification condition required. All \(i=1,...,r\) are covered, including growing \(r\); for a zero-dimensional quotient the bound is 0. No isogeny-invariance assertion is used. Thus the lower bound consumed when isolating \(\hF(B_{0})\) has the required original stable normalization. This closes that particular Bost-normalization source/application record at author level, pending review; it does not supply Rémond's superadditivity inequality.

The nine obligations are supplied by: (1) an abelian variety over a number field, (2) the actual \(C_{i}\) and \(J/B_{0}\), (3) every \(i\) and both positive- and zero-dimensional quotients, (4) the displayed stable-height normalization conversion, (5) the identical \(\hF\) in \(H_{3}\), (6) the dimension-dependent correction exactly retained, (7) no extra local hypothesis, (8) no principal-polarization or positive-dimension exception, and (9) the bound sum \(\hF(C_{i})\le \hF(B_{0})\le \hF(J)+3(g-2r)/2\) when combined with the separately retained superadditivity input. The corollary proves its all-polarizations assertion, rather than leaving an omitted polarization choice to the application.

\subsection{Rémond publication-version check remains open}

On 2026-11-10 \(I\) rechecked the publisher record for \emph{Propriétés de la hauteur de Faltings}, Annali SNS Pisa 23 (2022), 1589--1596, DOI \url{https://doi.org/10.2422/2036-2145.202010_062}. The article-body link again redirected to subscriptions. The author page still links the June 2021 text \url{https://www-fourier.univ-grenoble-alpes.fr/~remond/2476.pdf}, already frozen in the old packet. Its Corollary 1.2 is the subvariety/quotient inequality, its post-Corollary 1.3 paragraph explicitly includes stable heights, and Proposition 2.1 gives the one-sided height variation used in \(H_{3}\). \(I\) compared those actual passages, not merely the abstract. The final 2022 body has \textbf{not} been compared; this record does not promote the author text to a checked final edition.

The new Faltings, Milne and Li source checks narrow the previous access qualification. They neither resolve the Rémond version comparison nor close the \(H_{0}\) geometry, modularity, modular-height or analytic prime-pattern chain. Historical files and verdicts remain unchanged. The new sources are recorded by URL, edition, hash and exact locator; the full third-party books are not added to this repository.
\clearpage
% Source: research/quartic214/proofs/source-carry.md
\section{Exact retained analytic and restricted-converse carry to existing214}\label{q214:source-carry}
\noindent\textit{Source: }\nolinkurl{research/quartic214/proofs/source-carry.md}.\par
\begin{quote}\small \textbf{Current foundation overlay:} The Version~15 byte-identity and no-new-repair statements in this dated source describe its historical freeze.  This edition instead incorporates the complete current Version~16 foundation and the explicit source-to-consumer transition in Section~\ref{q214:foundation-transition}.  The historical text is retained as evidence; it does not negate or independently accept that changed application.\end{quote}

Author \nolinkurl{/root/d214_uniform_existence}, 2026-11-11. This is an author comparison, not independent acceptance. The prior existence supplement remains frozen at manifest SHA256\nolinkurl{61509c3d17d642ad0b76fc1c6ffd435e68481e8e112143449faf5a188d368e05}. Public evidence is the root-frozen104-file subset at \nolinkurl{../../public-carry/}, manifest SHA256\nolinkurl{4137a55c240edb06d39f6332dc1145d8591f98045ff2ec62d817c97c0f8ab487}. Full read-only public checkout is HEAD\nolinkurl{1a49013b5683f58d2f88bb018a07d13770354157}. No acceptance is authored here.

[C01: carry rule and exact source] Both the latest public manuscript and the entire frozen214 \nolinkurl{existence/sources/quartic_v15.tex} are the same1924438 raw bytes, SHA256\nolinkurl{44ba3a8ec8521c29b81b26cfb9981cd6d5f3c76d54e7782438614842edb52c74}. Therefore the internal RC, Kato, GI and Gaussian proof objects below are unchanged source applications. Protocol D's permitted reuse of exact accepted derivations does not require pretending that this author re-proved every published theorem or every old proof. It does require the object/domain/map/consumer comparison made here. The source-first public final006 application and connection decisions, source-profile decisions and findings retain their original scopes and limitations. Their green result is not a premise of an arithmetic implication. Eligible review must decide whether these exact bindings carry; this author does not certify that decision.

[C02: restricted statements selected from broad targets] The relevant public targets are V15-FOUNDATION, V15-PROJECT-CLOSURE, V15-GAUSSIAN, KATO-1999-INDEX-AND-CLASS, GI-00-ASSEMBLY and their stated dependencies. The original huge polynomial's compiler, six-contact input and canonical code domain are not themselves needed to deduce a new214 solution; only the theorem components actually used here are selected. In particular the later short zero-preparation pole is not identified with the old eleven-witness degree75 predicate. The raw224 construction branches before the public223/222 paired-root/trace-quotient reductions. No approval of the latter is required to infer their unused equations. The retained scalar/root/reference/raw224 inverse maps are checked at their own current application targets and in the frozen214 algebra packet.

[C03: exact rank and normalization] Read the actual generator proof at V15:347--443: two-isogeny descent, local squareclasses, torsion reductions, duplication-height comparison and finite small-height saturation certificate give \(2E_{\text{old}}(K)=ZP0\) for \(d_{\text{old}}=75-53\sqrt{2}\). The normalized model is obtained with \(\epsilon =1+\sqrt{2}\), \(\lambda =\epsilon ^{2}=3+2\sqrt{2}\) and \((x_{\text{new}},y_{\text{new}})=(\epsilon ^{4}*x_{\text{old}},\epsilon ^{6}*y_{\text{old}})\). Explicitly \((75-53\sqrt{2})(17+12\sqrt{2})=3-\sqrt{2}\) and \((1641-1160\sqrt{2})(17+12\sqrt{2})=57-28\sqrt{2}\); the ordinate transformation gives -660+417sqrt2 over32. Thus the normalized curve and generator are exactly \(E:y^{2}=x^{3}-(3-\sqrt{2})^{2}*x\) and \(P=2(1+\sqrt{2},2\sqrt{2})\), not an unrelated rank-one curve. The pointed isomorphism sends doubles to doubles and torsion to torsion; the finite nonidentity root charts therefore produce actual nonzero integer labels (n,a,t,b). Current profileSF27 explicitly records this corrected transport and separately supplies the normalized \(n^{2}-6\) height bound. The aggregate step of the214 supplement needs only nontorsionness, but the scalar and formal-index steps need this actual subgroup identification.

[C04: exact RC input/output at the pole consumer] The retained restricted converse is

\begin{gather*}
\text{ for }  \text{ each }  l>0 \text{ squarefree } , l=11 \text{ modulo } 12,\\
\text{corank}_{Z_{2}} \operatorname{Sel}_{2^{\infty }}(E_{l}/Q)=1 \Longrightarrow  \operatorname{ord}_{s=1}L(E_{l},s)=1,\\
E_{l}:y^{2}=x(x-l)(x+3l)
\end{gather*}

The current proof computes \(j=35152/9\) and non-CM before applying its representation arguments. The CM index curve in \(C_{03}\) is a different curve. The simultaneous descent in V15:8870--8970 constructs exactly two such parameters \(n_{\text{pole}}\) and \(n_{\text{pole}}*\theta \), including \(\theta =1\), and establishes full Selmer corank1 before invoking RC; it does not equate Selmer corank and Mordell--Weil rank before proving finite Sha. The preparation-zero proof at11656--11739 then rescales the resulting actual points to rational e,p,q,s for the actual given \(b\). The new214 supplementP19 verifies every hypothesis for its fresh \(b\), so the same two applications remain within domain. The RC proof fixes each \(E_{l}\) and its conductor before constructing coefficient families; no bound uniform as \(l\) varies is claimed or needed. Each input-dependent \(b\) may give different \(l\), and the universal theorem applies separately to each. None of its proof constants becomes a lower bound on the family parameter \(N\).

[C05: Kato required full index set] Read the actual current proof at13347--13946 and the actual earlier \(v_{11}\) \nolinkurl{norms.md} and \nolinkurl{interface.md} arguments. For a fixed constituent with primitive conductor \(D=2^{a}*d\), the retained construction defines \(y_{m}=\text{cor}_{Q(\zeta _{Dm})/Q(\zeta _{m})}(z_{Dm} \text{tensor} e_{\chi })\) for every \(m\) with \((m,c*d_{s})=1\). \(K_{5}\) makes \(N_{K}=2Nd\) and \(\text{Nprime}=c*d_{s}\) coprime. Restricting to all \(m\) divisible by \(N_{K}\) gives every index required by Kato1999\S{}0.1, not merely the smaller2004\S{}13.1 set. The class has the integral \(j_{*}\) condition by the inertia/Leray and Mackey calculation before mapping to the open scheme with \(N_{K}\) inverted. The same fixed lattice-clearing \(\text{factor}2^{{c}_{E}}\) applies at every level. It is a unit at the height-one localizations away from2 used downstream.

[C06: Kato transitions and actual limiting class] For every prime \(q\) not dividing Nprime, if q|m then \(S_{\text{Dmq}}=S_{Dm}\) and Proposition8.12 contributes the empty product, factor1. If \(q\) does not divide \(m\), \(N_{K}|m\) forces \(q\) outside2Nd; it is good and unramified for \(\chi \). Twisted equivariance gives the exact factor

\begin{gather*}
1-a_{q}*\chi (q)*q^{-1}*\sigma _{q}^{-1}+q^{-1}*\sigma _{q}^{-2}\\
=P_{q}(q^{-1}*\sigma _{q}^{-1}), P_{q}(X)=1-a_{q}*\chi (q)*X+q*X^{2}
\end{gather*}

Arithmetic Frobenius acts on \(T*\chi \) and inverse deck action is in the substitution; the representation is not dualized again. For every \(n\ge 1\) the supports \(S_{D*N_{K}*2^{n}}\) and \(S_{D*2^{n}}\) equal prime(2Nd). Corestriction transitivity through the common upper field gives \(\text{cor}(y_{N_{K}*2^{n}})=y_{2^{n}}\), without an Euler factor or a field-degree multiplier. These equalities commute with transition maps; the cofinal dyadic inverse limit is exactly the Section0.7 divisor-parameter-one class \(xi_{K}\), distinguished from the modular symbol xi. Thus the localized coordinate ideal \(J\) in \(K_{15}\) is the ideal of this actual retained class. Rank-one proportionality would not suffice and is not used. The \(\text{full}/j_{*}\) local cone comparison preserves \(H^{1}\) and the same class/ideal; paying for the actual \(H^{2}\) cone and fixed local factors yields \(K_{2}\). This source repair is incorporated already and is unchanged in214's supplier. A narrow current connection locator13775--13778 summarizes the coordinate-ideal conclusion; it is not the whole index/class argument. This carry binds the actual13799--13899 argument as well as its complete application record, so that short locator is not mistaken for the proof.

[C07: RC assembly object preserved] The even bound uses one fixed-form, retained-smoothed class over one group ring, not separate varying-newform lattices. Smoothing has \(v_{2}(s_{F})=2e\), independent of the odd common conductor support; it is not divided out. The binary families have one fixed real sign and local squareclasses at2N. The common scalar/determinant construction acts on actual finite free cochain/local-condition/duality diagrams for fixed support, retaining maps and homotopies. The current terminal argument at6650--6721 compares the actual nonzero-address constant coefficients with the base derivative; it does not apply the unknown base estimate to itself. Its finite-degree/missing-address bound is chosen before the graph precision. The odd coefficient and conductor-one oriented-CM source replacements retain their current type, place and normalization restrictions via V15-PROJECT-CLOSURE. New214 only consumes the resulting universal RC implication in \(C_{04}\); it constructs no new coefficient family and alters none of these internal maps. Exact unchanged dependency carry therefore applies to the theorem already proved for every such fixed curve; this does not assert a new proof of all RC internals by this author.

[C08: internal GI has the required output] Read the actual current assembled target at23826ff and the earlier standalone00-assembly argument, including its replacement joins. For fixed \(k\ge 3\) and \(\delta >0\) the target gives one fixed finite family of connected simply connected \(\text{step}\le k\) nilmanifolds, uniform Lipschitz bounds and positive correlation with a linear nilsequence, independent of interval length, function and orbit parameters. The interval norm is zero extension divided by the interval indicator norm. The proof retains represented unit-vector families with full-limit vertical identities; it does not upgrade a weak standard-point identity by continuity. Degree/rank iteration reconstructs frequency data on each new family; projection removes padding from the entire LP04 matrix; rank zero supplies the degree drop. The local multifactor all-box proof, plain-quadruple source, same-step \(P_{1}\) increment linearization and local converse replace the affected old wrapper transitions. The base uses the published forward bracket theorem plus the local trivial-subgroup repair of the false general quadratic extension. No unpublished general GI wrapper or old false skew suspension is assumed as an external theorem. These exact internal derivations and their published input scopes are the existing source dependency, not a new theorem citation introduced for214.

[C09: GI to Gaussian relative input] Read the actual current Gaussian join at19660--19746. For a required norm order \(s\ge 2\) it sets \(r=\operatorname{max}(s,4)\), \(k=r-1\ge 3\), fixes pseudorandom complexity and majorant scaling \(\alpha \), then a positive norm threshold. The prime-dominated function is not asserted one-bounded. The relative inverse theorem consumes GI(k) as its absolute input. Truncation and the precise indicator ratios give normalized \(U^{r}[L]\) norm at least \(\alpha *\delta /(4C_{\text{star}})\); no monotonicity with constant1 for normalized interval norms is assumed. The returned finite family is independent of L,z,b and the ray-class union. Stagewise \(z_{j}\) and physical thresholds fix each modulus before its asymptotic nilsequence and majorant estimates, uniformly over the finite residues at that stage. Thus the small-norm contradiction uses a pseudorandom sequence along a slow cutoff, not a fixed-z or growing-field limit. This is exactly the input used by the unchanged Gaussian prime theorem; the214 coefficient list only changes fixed arithmetic data after which that theorem quantifies.

[C10: Gaussian finite-field and tail scope] The actual proof at17545--20359 uses fixed \(E=Q(i)\) ray fields, arithmetic Frobenius, Mitsui fixed-field Hecke estimates, ordinary prime majorants and the internal GI join; no unrestricted Galois-prime theorem or Frei--Tinkova preprint is silently assumed. For fixed coefficients the full obstruction tail is bounded separately from finite-cutoff local exclusions. First choose the cutoff \(L_{\text{cut}}\) so \(C/(L_{\text{cut}} \log L_{\text{cut}})<c/[32(\log L_{\text{cut}})^{4k}]\), then fix the field/ray class data and use the slow-width diagonal asymptotic. This retains positive mass after removing every forbidden odd common divisor, including the same-position obstruction. The rational214 coefficient list after clearing scale denominators and merging repeats is within that exact fixed-data theorem. The finite set and sizes may depend on (N,x), which is allowed; no estimate uniform in them is used to define \(N_{0}\). The separate actual reduction lemma12239--12268 applies because \(\pi =9\) mod16 and1+3i mod7 imply \(p=1\) mod8 and \((7/p)=-1\). The reductions of \(d=3-\sqrt{2}\) have opposite squareclasses; the two CM traces are opposite and the actual group orders are the unordered norms \(N(\pi -1),N(\pi +1)\). Full two-torsion and \(P\) being a double justify d/4 in CRT. This is the precise curve and generator already identified in \(C_{03}\).

[C11: source-profile arithmetic replacements selected exactly] The two final source-profile files have raw hashes\nolinkurl{3668753afd46bff7e3a1bc3bed63d2f102009d5bd3d379fa67411c55b2aaa2d2} and\nolinkurl{881213958ad12b1af53e3c6758361ea8e1c371833453bcac7d63ee48909ad971}, exactly the subject hashes in the public source-profile decision. Their actual mathematical statements were read, as well as selected current connection records; the review verdict is not substituted for those statements. SF22--26 give fixed-data counting, strip, rational coefficients, CRT and real recurrence. SF27 identifies the normalized index model. SF33's finite/real order of choices and SF40's specific scalar rates \(R^{-4},R^{-6}\) match the new supplementP13--P18. SF35 changes the baseline |a-1| criterion to |a| with its own soundness/completeness; SF39 gives the parity criterion's exponent9/4; the four-contact232 application retains that exponent and supplies \(L_{N}=\text{floor}(9(2N-1)/4)+1\). SF45 gives the exact zero-preparation domain and conic charts, including \(e=1\) and \(e=-1/3\). These are unchanged source applications in the literal raw224 system.

[C12: changes that are freshly checked, not carried by their conclusions] SF28's old reflected exponent7/4 is replaced in214 by the new full-family height transfer giving7/6; its arithmetic conclusion does not authorize the stronger bound. SF30/SF32 describe old order-four extraction and a shifted scalar at the baseline, so they cannot directly supply214's order-three factor proof or \(c=|a|\). The frozen supplementP05--P12 supplies that changed route and exact integer consumers anew. SF33's old comparison factors cannot supply new pole certificates without reevaluating actual \(b\); the frozen supplementP17--P19 performs that reevaluation and checks the fresh Psi0 hypotheses. These are the substantive new214 joins requiring fresh eligible acceptance. The source rank, RC theorem, prime-pattern theorem, scalar parity/four-contact theorem and exact raw maps themselves are unchanged; their previously checked statement/domain/map scopes can be carried only with the bindings recorded here.

[C13: exact disposition] Within the inspected source-to214 analytic/RC/rank/prime/scalar/pole joins, no unmatched mathematical required input remains after the frozen supplement is included. This is a bounded author conclusion: it does not newly validate every source theorem, the separate full-family height chain, contributor identities or final release. It gives concrete current application matching sufficient for an eligible reviewer to decide reuse at the established source scopes. The public final006 coverage decision expressly permits retained accepted source proofs while declining a fresh line-by-line reproof of all history; this author retains that distinction. Any source profile, theorem body, construction, consumer or inherited dependency change invalidates this carry. Historical classification is prior justification found and read, with current incorporation already present in the identical supplier; the fresh214 boundary supplement remains pending new independent acceptance. The defining threshold remains ineffective and no full numerical214 witness tuple has been produced.
\clearpage\section{Retained selected source proofs}\label{q214:selected}
These mathematical fragments are reproduced in full at the selected public ranges.  Their original coarse character/conductor route is retained as historical source evidence; M1 and M5 in Section~\ref{q214:modularity} supply the current character and conductor comparison.  The evaluation fragments are read together with the exact lattice, different and metric bridges above.  Historical source qualifications remain visible; no new acceptance is inferred from their inclusion.
\clearpage
% Source: research/quartic214/sources/support-section1.md
\subsection{1. The new radical and the actual modular objects}

Set \(S=\{2,3,5,7,11,13,17,37,41,2399\}\), \(B_{S}=\prod _{p \in  S}\) \(p\). For \(a=u/v\)
in lowest terms, \(v>0\), excluding infinity,0,1,2401,-343/32, let \(M_{S}(a)\)
be the squarefree product of \(p\) outside \(S\) at which one of

\begin{gather*}
v, u, u-v, u-2401v, 32u+343v
\end{gather*}

has positive odd valuation. Put \(R=B_{S}\) \(M_{S}(a)\). All prime products count
a prime once. \(R\) is squarefree, divisible by210, and \(R\ge B_{S}>e\).

The supplied geometric construction in M:7605--8040 is the same full
three-level normalizer fiber, not a different abelian surface. Its normal
fiber field \(L\) has degree dividing576; all geometric homomorphisms
between its conjugate fibers are defined over \(L\). Outside210 the quotient
has tame coordinates \(w=z^{2}\) at its five branch sections and is etale
elsewhere. Even contact is unramified after an unramified extension,
including contact with infinity. Thus its field discriminant and bad
reduction support lie in supp(R) for this \(S:\) the proof of M's \((R)\) uses
only the branch contacts and primes dividing210, not the cutoff76176.
The metric computation \((G)\), on this same surface, remains

\begin{gather*}
h(a)\le 2 \operatorname{max}(0,\hF(A_{a}))+128. (1)
\end{gather*}

Here \(h\) is absolute logarithmic height and \(\hF\) is original stable Faltings
height, with the parent's metric conversion and source qualifications.

In the non-CM case follow M's LR1--LR5 and \nolinkurl{quartic-height19}, with this
same support \(T=\operatorname{supp}(R)\). The scalar sign cocycle is inflated from Gal(L/Q).
At an odd ramified prime the chosen odd tame Dirichlet character has
order dividing p-1; at2 its conductor divides4. The combined character
has order \(m\le 2\) \(\prod _{p|R,p \text{odd}}(p-1)=2\) \(\phi (R)\). Its square-root lift
\(b\) and quadratic support-killing twist produce the same normal field \(L'\)
with degree \(\le 2m[L:Q]\le 2304\) \(\phi (R)\). The twist introduces no prime
outside \(T\). This does not identify the later nebentypus with that scalar
character: the actual compatible modular representations and determinant
comparison give \(\epsilon ^{2}\) trivial on \(G_{L'}\). The final primitive weight-two
newform, with its finite character absorbed as in LR3--LR5, has

\begin{gather*}
N\le 210^{300} R^{2}, \operatorname{rad}(N)|R. (2)
\end{gather*}

At \(p>7\) the field degree divides576, so the field inertia is tame;
the scalar lift is tame at odd primes and the representation conductor
exponent is at most2. The unchanged crude exponents at2,3,5,7 are below300.
No reduction of those exponents or new 2-primary character choice is
used. The same \(A_{a}\) over \(L'\) is an isogeny factor of
\(J=\operatorname{Jac}(X_{1}(N)/\operatorname{ker} \epsilon )\) over \(L'\). This is the actual factor, not a
coefficient-field replacement. These support arguments hold for all
rational a in the stated five-contact domain, not merely constructed
integer inputs.

The CM branch M:8517--8655 also has no large-primorial hypothesis. All
endomorphisms and the elliptic factor live over \(L\) with \([L:Q]\le 576\).
The CM quadratic field has discriminant \(\le 4\) \(\operatorname{rad}(D_{L})\le 4R\); the maximal
order elliptic comparison gives height \(<3R\), its square \(<6R\), and the
same-field isogeny comparison gives \(\hF(A_{a})<854R\). Its proof assumes
only \(R\ge 1\) and the support just checked. Both branches are retained.

\clearpage
% Source: research/quartic214/sources/conductor-sections1-3.md
\subsection{1. The same surface, support, and normal fields}

Keep \(S=\{2,3,5,7,11,13,17,37,41,2399\}\) and \(B_{S}=1857890364330\).
For \(a=u/v\) in lowest terms with \(v>0\), outside the five contacts
infinity,0,1,2401,-343/32, let \(M_{S}(a)\) be the squarefree product of
\(p\) outside \(S\) with positive odd valuation on one of

\begin{gather*}
v, u, u-v, u-2401v, 32u+343v
\end{gather*}

Set \(R=B_{S}\) \(M_{S}(a)\), \(T=\operatorname{supp}(R)\), and \(\ell =\operatorname{log}\) \(R\). The parent proof \S{}1
constructs the same surface \(A_{a}\) over a normal fiber field \(L\), with
[L:Q] dividing576, \(L\) unramified outside \(T\), and \(A_{a}\) having good
reduction over \(L\) away from210. All geometric homomorphisms between
conjugate fibers are L-defined. These statements apply to every rational
parameter in the stated domain; they are not restricted to integer
inputs. Keep the parent's original stable Faltings height convention and
geometric comparison

\begin{gather*}
h(a) \le  2 \operatorname{max}(0,\hF(A_{a})) +128. (1)
\end{gather*}

First consider the non-CM case. The actual argument in M:8042--8095
proves potential good reduction using the quaternion division action:
a positive torus rank \(t\le 2\) would give a unital injection \(B_{\text{quat}}\) into
\(M_{t}(Q)\), which is impossible. It then proves that \(E=L(A_{a}[11])\) is normal
over \(Q\). For every conjugate, the B-linear normalizer isogeny can be
scaled to have degree prime to11, since the maximal-order normalizer
at11 is \(Q_{11}^{*}\) \(\operatorname{GL}_{2}(Z_{11})\). The isogeny is L-defined and identifies the
11-torsion kernels of the original surface and its conjugate. Thus every
conjugate of \(E\) is \(E\). This normality is used, not inferred from a size
estimate. Its faithful action gives Gal(E/L) contained in \(\operatorname{GL}_{4}(F_{11})\), so

\begin{gather*}
[E:Q] \mid  576 |\operatorname{GL}_{4}(F_{11})|. (2)
\end{gather*}

This strengthens the printed inequality LR1 by retaining Lagrange
divisibility. The exact product of the four column choices is

\begin{gather*}
|\operatorname{GL}_{4}(F_{11})| = \prod _{i=0..3} (11^{4}-11^{i})\\
= 41393302251840000,\\
576 |\operatorname{GL}_{4}(F_{11})| =23842542097059840000
\end{gather*}

Its valuations at2,3,5,7 are respectively15,4,4,1. Full11-torsion makes
\(A_{a}\) semistable over \(E\) at every residue characteristic different from11;
the preceding torus argument makes this good reduction. Only2,3,5,7
use this torsion criterion here.

Retain the precise scalar construction M:8097--8179. The finite choice
\(\mu _{\sigma }\) is inflated from Gal(L/Q), with \(\mu _{1}=1\). Its normalized version
\(a_{\sigma }\) \(\mu _{\sigma }\), where \(a_{\sigma }=(\operatorname{deg} \mu _{\sigma })^{-1/4}\), has sign cocycle
\(\epsilon _{\text{sign}}\) inflated from Gal(L/Q). The controlled Dirichlet character
\(\chi \) is tame at every odd prime, has conductor dividing4 at2, and has
order \(m\le 2\) \(\phi (R)\). After the support-killing quadratic twist, the
cochain \(b\) still satisfies \(b^{2}=\chi \) and \(\delta \) \(b=\epsilon _{\text{sign}}\). Its restriction
to \(G_{L}\) is a conjugation-invariant character. Thus its kernel field \(L'\)
is normal and

\begin{gather*}
d'=[L':Q] \le  2m[L:Q] \le 2304 \phi (R). (3)
\end{gather*}

For the conductor computation only, set \(F=E\) \(Q(\chi )\), where \(Q(\chi )\)
means the fixed field of ker \(\chi \), not the field generated by character
values. On \(G_{F}\) the cochain \(b\) is a conjugation-invariant quadratic
character: this subgroup lies in \(G_{L}\) and \(\chi \) is trivial there. Its kernel
defines a normal field F'/Q, with \([F':F]\le 2\). This is M:8348--8354's
field; it is not a replacement for the surface factor field \(L'\).

At an odd \(p\), \(Q(\chi )/Q\) is locally tame and the quadratic extension
F'/F is tame. Base change of a tame local extension remains tame, so
these additions do not increase \(v_{p}\) of the ramification index. At2,
the inertia image of \(\chi \) has order at most2, and the second extension
also has ramification degree at most2. Unramified degrees of these
fields do not affect this assertion. For a completion of \(F'\) at a place
over \(p\) let \(e\) be its absolute ramification index. Equation(2), divisibility
of a local inertia order into the global Galois group, and these two
local additions give

\begin{gather*}
v_{2}(e)\le 17, v_{3}(e)\le 4, v_{5}(e)\le 4, v_{7}(e)\le 1. (4)
\end{gather*}

The tame part of \(e\) is not bounded by(4), and will not need to be.

\subsection{2. From the different to the last upper break}

Here is the ramification calculation with conventions explicit. Let
\(K/Q_{p}\) be a finite Galois extension with ramification index \(e\), maximal
unramified subfield \(K_{0}\), and uniformizer \(\pi \). Normalize \(v_{K}(\pi )=1\).
The minimal polynomial \(f(X)=X^{e}+\sum _{i=0..e-1}\) \(a_{i}\) \(X^{i}\) over \(K_{0}\)
is Eisenstein and \(O_{K}=O_{K_{0}}[\pi ]\). Therefore the different exponent is
\(\delta =v_{K}(f'(\pi ))\). This also follows from the monogenic trace-dual
description. Each nonzero term \(i\) \(a_{i}\) \(\pi ^{i-1}\) in the derivative has
valuation congruent to i-1 modulo \(e\); these residues, including the
leading derivative term \(e\) \(\pi ^{e-1}\), are distinct. The minimum valuation
cannot cancel, so

\begin{gather*}
\delta  \le  v_{K}(e \pi ^{e-1}) = e v_{p}(e)+e-1. (5)
\end{gather*}

The unramified base field is essential to \(v_{K}(e)=e\) \(v_{p}(e)\).
The assertion includes \(e=1\), with \(\delta =0\).

For completeness, the upper integral formula can be derived directly
from the same derivative. Put \(G=\text{Gal}(K/Q_{p})\), with \(G_{0}\) inertia and
\(G_{i}=\{\sigma  \in  G_{0}: v_{K}(\sigma  \pi -\pi )\ge i+1\}\) for integers \(i\ge 1\).
Since \(f'(\pi )=\prod _{\sigma  \in  G_{0}, \sigma \ne 1}(\pi -\sigma  \pi )\), counting
each positive valuation by its integer levels gives

\begin{gather*}
\delta =\sum _{i\ge 0}(|G_{i}|-1)
\end{gather*}

Use \(G_{t}=G_{\lceil t\rceil}\) for real \(t>0\) and define the Herbrand function
\(\phi (t)=\int _{0}^{t}\) \(|G_{s}|/e\) ds. Its inverse defines \(G^{u}=G_{\phi ^{-1}(u)}\).
Changing variables \(u=\phi (t)\) now gives, with no missing tame term,

\begin{gather*}
\delta /e = 1-1/e + \int _{0}^{\infty } (1-1/|G^{u}|) du. (6)
\end{gather*}

These definitions and quotient compatibility are the standard ones:
see \href{https://www.jmilne.org/math/CourseNotes/ANT.pdf}{Milne ANT v3.08, Proposition7.55, Remark7.56 and Theorem7.58}
and \href{https://www.jmilne.org/math/CourseNotes/CFTc.pdf}{Milne CFT v4.03, I.4, Proposition4.4}.
They are also the conventions invoked at M:8363--8378. The displayed
calculation supplies the particular integral used here, instead of
substituting into the manuscript's coarse numerical bound.

Let \(b_{\text{last}}\) be the last positive upper break, or0 when there is no
wild ramification. Until \(b_{\text{last}}\) the upper group is nontrivial, so the
integrand in(6) is at least1/2, except at irrelevant endpoints. Equations
(5)--(6) imply

\begin{gather*}
b_{\text{last}} \le 2 v_{p}(e). (7)
\end{gather*}

No abelian or Hasse--Arf assumption is used. Large tame inertia cancels
in the integral change of variables and causes no extra term in(7).

\subsection{3. The representation whose conductor is the final level}

For each coefficient prime \(\ell _{c}\), keep M:8179--8347's original

\begin{gather*}
\rho _{{\ell }_{c}}(\sigma )=b(\sigma )^{-1} a_{\sigma } \mu _{\sigma } \circ \sigma ,\\
\rho _{{\ell }_{c}} = r_{{\ell }_{c}} \oplus  r_{{\ell }_{c}}, \operatorname{dim} r_{{\ell }_{c}}=2
\end{gather*}

On \(G_{F'}\) all three scalar/isogeny corrections are trivial: \(F'\) contains
\(L\), \(\mu _{1}=a_{1}=1\), and \(b\) is killed. Thus \(\rho _{{\ell }_{c}}\) restricted to \(G_{F'}\) is
the Tate representation of this same \(A_{a}\). At \(p=2,3,5,7\), \(A_{a}\) has good
reduction over \(F'\) by \S{}1. For \(\ell _{c}\ne p\), the local inertia in \(G_{F'}\) acts
trivially. The original \(r_{{\ell }_{c}}\) therefore has finite inertia image,
and compatibility of upper numbering with quotients makes it trivial
on absolute upper groups above \(b_{\text{last}}\). Its Swan conductor is the
integral of \(\text{codim}(r_{{\ell }_{c}}^{I^{u}})\); it is at most2 \(b_{\text{last}}\). Its tame
codimension is at most2. Consequently(4) and(7) give

\begin{gather*}
a_{p}(r_{{\ell }_{c}})\le 2+4 v_{p}(e) \le  (70,18,18,6)\\
\text{ for }  p=(2,3,5,7), \text{ respectively } . (8)
\end{gather*}

For \(p>7\), [L:Q] dividing576 implies that L/Q is tame at \(p\).
\(A_{a}/L\) already has good reduction there, including \(p=11\). Wild inertia
lies in \(G_{L}\), so \(\mu =a=1\) on it. On this pro-p group \(b\) is a character,
its square \(\chi \) is tame, and its remaining image has order at most2;
hence \(b\) is trivial. This proves \(a_{p}(r_{{\ell }_{c}})\le 2\) for \(p>7\), and zero
outside \(T\). The prime11 argument does not apply the full11-torsion
criterion in residue characteristic11.

The object in(8) is the representation of the final primitive newform
\(f\), after its finite twist is absorbed. M:8191--8273 identifies \(r_{5}\)
exactly, then compares rational traces of the same quasi-isogenies
(including division by the denominator-clearing integer and the factor
1/2 from two equal summands). Equality in the common coefficient field,
Chebotarev and absolute irreducibility identify the original \(r_{3}\) with
the full three-adic realization of the same \(f\). The auxiliary comparison
field does not replace \(L'\), \(F'\), or the support \(T\).

Use CL1--CL2 at M:8274--8347: Casselman's primitive level is the local
conductor, with arbitrary nebentypus; Carayol's Theorem A supplies full
local Weil--Deligne compatibility over \(Q\), including the dyadic places.
Contragredient and unramified normalization twists preserve conductor.
The base degree is one, so the extra even-degree hypothesis does not
arise. No inertia semisimplification discards monodromy. Use \(r_{5}\) at
every \(p\ne 5\) and \(r_{3}\) at \(p=5\); this covers every prime. These established
identifications are retained source inputs, not new modularity proofs.

All four primes in(8) divide \(R\). Absorbing their baseline exponent2
into \(R^{2}\) gives

\begin{gather*}
N \le  C_{N} R^{2}, \operatorname{rad}(N)|R,\\
C_{N}=2^{68} 3^{16} 5^{16} 7^{4},\\
D_{0}=4608 C_{N}. (9)
\end{gather*}

This bound is the only conductor change in this packet.

\clearpage
% Source: research/quartic214/sources/aggregate-bridge.md
\section{Actual aggregate input for the 214 argument}\label{q214:source-aggregate-bridge}
\noindent\textit{Source: }\nolinkurl{research/quartic214/sources/aggregate-bridge.md}.\par

\begin{quote}\footnotesize Author: model task \nolinkurl{/root/d214_foundation_trace}, 2026-11-10. This is a proposed source/application bridge, not independent acceptance. All file and passage bindings are in \nolinkurl{FOUNDATION_MAP.json} (F-AGGREGATE). The target is the unchanged 214 proof at frozen private commit\nolinkurl{569a990bd3dfe141e2a1328f289fe62a82d13563}.\end{quote}

The conclusion needed at \nolinkurl{existence/PROOF.md:62} is a fixed inequality \(B^{2}\le M_{1}\) \(h(T)+M_{2}\) for all four nonzero integer labels, where \(T\) is the actual quarter-trace of the sum of their abscissae. No computable \(M_{2}\) is needed. The supplier is the internal proof in the frozen Version15 source, lines31238--31643, consumed explicitly at32449--32468. \(I\) read the sum-pole, cross-block independence, graph-family and canonical-height arguments, rather than inferring this conclusion from the product version.

The final published source was retrieved on2026-11-10: \(J\). Von Buhren, \emph{Borne de hauteur semi-effective pour le problème de Mordell--Lang dans une variété abélienne}, JTNB29(2017),307--320, DOI10.5802/jtnb.981, publisher-distributed PDF at \url{https://www.numdam.org/item/10.5802/jtnb.981.pdf}. Theorem1.1 and preceding definitions, printed308, require a projective integral parameter scheme, an integral closed family in an abelian variety times that scheme, suitable ample embeddings, a finite-rank subgroup, and avoidance of fiberwise positive-dimensional abelian translates. They bound the canonical height by a fixed geometric multiple of parameter height plus a subgroup-dependent constant. Points of the subgroup belong to its allowed thickening by taking the error point zero. The final body explicitly does not determine the additive constant. No theorem about an individual finite fiber is substituted for this family theorem.

\begin{itemize}
\item 1. \textbf{Source hypotheses.} Take \(A=E_{0}^{8}\), \(E_{0}:Y^{2}=X^{3}-X\), parameter space \(P^{1}\), and the reduced graph closure \(V\) of the rational function \(F_{+}=(d/2)\sum X_{i}+(d'/2)\sum X'_{i}\), where \(d=3-\sqrt{2}\) and \(d'=3+\sqrt{2}\). The graph over its nonempty domain is integral; its closure is integral and closed in the projective product. All geometry is fixed before the labels. The source's embedding condition(1.1) can be met concretely: embed \(E_{0}\) by \(O(3O)\) as a smooth plane cubic, and the product by its Segre embedding. For every \(n\ge 1\) the plane-cubic exact sequence and \(H^{1}(P^{2},O(n-3))=0\) give surjective restriction of degree-n forms; tensor products give the same surjectivity on the product. Thus this fixed symmetric ample embedding has the required surjectivity in every degree. \(P^{1}\) uses \(O(1)\). This spells out an implicit embedding choice in the retained project proof; it does not alter its rational function or inequality.
\item 2. \textbf{Constructed objects.} The eight coordinates are \(([n_{i}]P_{1})\) and \(([n_{i}]P_{2})\), \(i=1,\ldots ,4\), where P1,P2 are the two pointed-isomorphism images of the same \(P\) and its conjugate. The finite-rank subgroup is generated by \(P_{1}\) in each of the first four positions and \(P_{2}\) in each of the last four. It is fixed and finitely generated, and contains every required tuple. The extra equality of the paired integer labels is an allowed restriction of that domain.
\item 3. \textbf{Domain and index coverage.} Every \(n_{i}\) is a nonzero integer; negative, repeated and noncoprime labels are included. A nonzero multiple of either nontorsion base point is not \(O\), so the tuple lies in the domain of \(F_{+}\). The sum proof has no oddness hypothesis. The older product proof's odd-label domain is not used.
\item 4. \textbf{Maps.} Use \(x=dX\), \(y=d\sqrt{d}Y\) and their conjugates; these are pointed isomorphisms. The squareclasses of d,d' are independent in \(K(i):\) neither is a square because its norm is7, and \(dd'=7\) is not a square in \(K\). A positive square from \(K(i)\) would already be a square in the real field \(K\). The automorphism fixing \(\sqrt{d}\) and negating \(\sqrt{d'}\) fixes \(\operatorname{End}(E_{0})=Z[i]\), fixes \(P_{1}\) and negates \(P_{2}\). Applying it to a torsion relation proves independence over Z[i] modulo torsion. These are the actual two base points; no unrelated generic pair is substituted.
\item 5. \textbf{Consumed object.} \(F_{+}(q)=(1/2)\operatorname{Tr}(\sum x([n_{i}]P))=2T\). Canonical heights are preserved by the pointed isomorphisms and absolute conjugation. The eight-coordinate canonical-height sum is twice the four-coordinate sum. The fixed \(O(3O)\) product height is a fixed positive multiple of this sum. Since \(h(2T)\le h(T)+\log 2\), the factor2 is absorbed only into the fixed additive constant. The resulting \(T\) is the quarter-trace in the literal raw224 recovery and existence proof.
\item 6. \textbf{Conventions.} \(h\) is absolute logarithmic Weil height, and the elliptic canonical height is the normalization with \(2\widehat{h}(Q)=\text{lim}\) \(4^{-m}\) \(h(x([2^{m}]Q))\). Put \(hP=\widehat{h}(P)>0\). The family theorem gives \(\sum \widehat{h}([n_{i}]P)\le M\) \(h(2T)+M_{0}\) for fixed M,M0. Then \(hP\cdot B^{2}\le \sum \widehat{h}([n_{i}]P)\); division by hP and the preceding log2 loss give exactly the required inequality. This also avoids using a sharper rank-height estimate merely to obtain existence of M1,M2. Replace \(M_{1}\) by a larger positive value and \(M_{2}\) by \(\operatorname{max}(1,M_{2})\) if necessary.
\item 7. \textbf{Local conditions.} The source application is over algebraic number fields with absolute heights; no local integrality, splitting or coprimality condition is added. In the internal nonconstant-fiber proof, any nonzero abelian subgroup direction contains an elliptic direction \(E_{0}\to B\), obtained from a nonzero coordinate projection by duality and a polarization up to isogeny. Within one block, the common-pole leading coefficient is a nonzero tangent scalar times a sum of positive rational weights times \((n_{i}/n_{j})^{2}\). It cannot vanish. A pole cannot occur in both independent blocks. Thus \(F_{+}\) is nonconstant on every proposed positive-dimensional translate through \(q\).
\item 8. \textbf{Exceptional cases.} Nonzero labels exclude all coordinate poles. Their signs and coincidences do not affect the positive-square coefficient. If a translate through \(q\) lay in a projective graph-closure fiber, its nonempty open intersection with the rational-function domain would force \(F_{+}\) to be constant there, contrary to the pole argument. Boundary components introduce no exception through \(q\). P's infinite order is independently established in the same source by its nonzero formal parameter at2, lines31544--31568; full rank-one generation is unnecessary for this height step.
\item 9. \textbf{Downstream consumption.} The same recovered labels and \(T\) occur in \nolinkurl{existence/PROOF.md:9,62–78.} The nonzero physical arguments \(t_{0}=\lambda \) and \(t_{1}=T\lambda \) give \(h(T)\le h(t_{0})+h(t_{1})\). Separate divisibilities of each radical suffice; no product divisibility is claimed. This yields the displayed exponent7/(6k), \(\text{hence}14/9<2\) for the Lucas residual at \(k=3\). The additive constant remains fixed and ineffective, so it is absorbed only after choosing a sufficiently large uniform \(N_{0}\). No numerical defining threshold follows.
\end{itemize}

Historical disposition: prior justification found, including the independent effective-aggregate assessment in \nolinkurl{research/v14-small-prime-height/exploration/effective-aggregate/ASSESSMENT.md.} That assessment already distinguishes a valid qualitative bound from the unprovided effective additive constant; it is not an objection to the current ineffective theorem. Manuscript disposition: the sum/graph/quarter-trace argument is incorporated in the frozen Version15 supplier and referenced by the 214 proof. The explicit embedding choice and the \(\text{canonical}-\text{height}-\text{ to } -B^{2}\) derivation above are supplemental author explanations pending non-author comparison. Current review: the old component review retains the aggregate interface conditionally; this new final-source bridge has no independent acceptance yet. It does not close rank/descent, modularity, prime patterns or pole realization.
\section{Retained differential-evaluation proof}\label{q214:eval-source}
\subsection{Conventions and the evaluation estimate}

If $K$ is a number field, $\OK$ denotes its ring of integers.
For a nonzero ideal $\mathfrak a\subset\OK$, write
$\mathrm N\mathfrak a=\#(\OK/\mathfrak a)$.
If $E$ is a projective $\OK$-module with Hermitian metrics at
its complex embeddings, $\Ardeg(E)$ means the arithmetic degree
of its determinant. In rank one, for a nonzero integral section $s$,
\[
 \Ardeg(E)=\log\#(E/\OK s)
                  -\sum_{\sigma:K\hookrightarrow\C}\log\|s\|_\sigma.
\]
The definition is independent of the choice of $s$ by the product
formula. Both embeddings in a complex conjugate pair are included.

Throughout, $\hF$ is the original stable Faltings height.
For a curve $X/K$ with semistable reduction, its Jacobian height is
\begin{equation}\label{q214src:eval:eq:height-convention}
 \hF(\Jac X)=\frac{\Ardeg H^0(\mathcal X,
                                  \omega_{\mathcal X/\OK})}{[K:\Q]},
\end{equation}
where $\mathcal X$ is a minimal regular semistable model. The metric
on differentials at a complex embedding is
\[
 \langle\omega,\eta\rangle
       =\frac i2\int_{X(\C)}\omega\wedge\overline\eta.
\]
This is the determinant metric corresponding to the original Faltings
normalization; see \cite[Sections~2.1 and~2.3]{v15-height-Javanpeykar} and
\cite[Section~2.1]{v15-height-vonKaenel}. Passing to a finite extension on which
the curve is semistable does not change the stable height.

A \emph{Belyi map} is a finite morphism from a smooth projective
geometrically connected curve to $\mathbb P^1$, unramified outside
$0,1,\infty$. Let
\[
 \pi:X\longrightarrow\mathbb P^1
\]
be such a map, of degree $d$, where $X$ has genus $g\ge1$.
Write $t$ for the coordinate on the target.
Choose a number field $K$ over which $X$ and $\pi$ are defined,
$X$ is semistable, and every point of $\pi^{-1}(2)$ is rational.
Let $\mathcal Y$ be the normalization of $\mathbb P^1_{\OK}$ in
$K(X)$. It is finite and flat over $\mathbb P^1_{\OK}$.
On the affine $t$-chart write its algebra as $B$, a finite projective
$\OK[t]$-algebra of rank $d$, and put
\begin{equation}\label{q214src:eval:eq:order}
 C=B/(t-2)B\subset K^d,
 \qquad O=\OK^d.
\end{equation}
The inclusion in $K^d$ uses the $d$ distinct points of the split
fibre. Every element of $C$ is integral over $\OK$, so $C\subset O$.
Let $\mathfrak D(C)$ be the discriminant ideal of its trace pairing.

\begin{theorem}\label{q214src:eval:thm:evaluation}
With this notation,
\begin{equation}\label{q214src:eval:eq:evaluation-bound}
 \hF(\Jac X)
 \le\frac{\log\mathrm N\mathfrak D(C)}{2[K:\Q]}
       +g\log\frac2{\sqrt\pi}.
\end{equation}
\end{theorem}

The auxiliary field $K$ is used to split a finite fibre and to
compute a stable height. Its degree will cancel from the estimates.
No upper bound for $[K:\Q]$ is assumed in this theorem.

\subsection{Trace-dual lattices and regular differentials}

\begin{lemma}\label{q214src:eval:lem:trace-lattice}
Let $C\subset O=\OK^d$ be an order with $C\otimes K=K^d$, and
let
\[
 C^\vee=\{x\in K^d:\operatorname{Tr}_{K^d/K}(xC)\subset\OK\}.
\]
For any set $I$ of $g$ coordinate positions, set
$F=\operatorname{pr}_I(C^\vee)\subset K^g$.
Then $F$ is a projective $\OK$-module containing $\OK^g$.
With the standard Euclidean metric on $K^g$ at every embedding,
\[
 \Ardeg(F)\le\tfrac12\log\mathrm N\mathfrak D(C).
\]
\end{lemma}
\begin{proof}
Work locally at a finite prime $\mathfrak p$. The trace pairing
on $K^d$ is the ordinary dot product, and $O$ is self-dual.
Thus $O\subset C^\vee$, and dualizing the inclusion $C\subset O$
shows
\[
 \operatorname{length}(C^\vee/O)
 =\operatorname{length}(O/C)
 =\tfrac12v_{\mathfrak p}(\mathfrak D(C)).
\]
For example, if a basis matrix for $C$ relative to $O$ is $U$,
then the trace Gram matrix is $U^{\mathsf T}U$, which proves the
last equality. These equalities are valid locally even when
$C$ is not globally free.

Projection gives a surjection
$C^\vee/O\twoheadrightarrow F/\OK^g$. The latter quotient
therefore has no larger length at any prime.
A finite torsion-free module over the Dedekind ring $\OK$ is
projective. Finally, the standard basis of $\OK^g$ has determinant
norm one, so
\[
 \Ardeg(F)=\log\#(F/\OK^g)
 \le\tfrac12\log\mathrm N\mathfrak D(C).
\]
\end{proof}

\begin{lemma}\label{q214src:eval:lem:integral-evaluation}
Let $\mathcal X$ be a minimal regular semistable model of $X$.
The evaluations of a differential in
$E=H^0(\mathcal X,\omega_{\mathcal X/\OK})$ at the points of
$\pi^{-1}(2)$, expressed in the coordinate $t$, lie in $C^\vee$.
Equivalently, there is an integral map
\[
 E\longrightarrow C^\vee,
 \qquad \omega\longmapsto
       \left(\frac{\omega}{dt}(P)\right)_{P\in\pi^{-1}(2)}.
\]
\end{lemma}
\begin{proof}
First, normalization is finite because the base is excellent.
A normal surface is Cohen--Macaulay, and a finite morphism from
such a surface to a regular surface is flat by miracle flatness;
see \cite[Tag~00R4]{v15-height-Stacks}. Consequently $B$ is finite projective
over $\OK[t]$.

Take a common regular model $\mathcal Z$ dominating
$\mathcal X$ and $\mathcal Y$. Pullback under a blowup of regular
arithmetic surfaces takes regular relative differentials to regular
relative differentials. One can also read this directly from the
canonical-sheaf formula for a blowup; see
\cite[Lemma~2.3.1]{v15-height-Javanpeykar}.
Hence a differential in $E$ is regular on $\mathcal Z$.
The birational morphism $\mathcal Z\to\mathcal Y$ is an
isomorphism at every codimension-one point of the normal surface
$\mathcal Y$.

The dualizing module of a Cohen--Macaulay normal surface is
reflexive: it is torsion-free and satisfies Serre's condition
$(S_2)$. Thus sections defined at all codimension-one points extend
across the codimension-two complement; see
\cite[Tags~0AWQ and~0AVT, Lemmas~31.13.12--31.13.14]{v15-height-Stacks}.
The differential therefore defines a section of the relative
dualizing sheaf of $\mathcal Y$.

Finite duality identifies this sheaf on the affine chart with
\[
 \operatorname{Hom}_{\OK[t]}(B,\OK[t])\,dt;
\]
see \cite[Tags~0BUL and~0C0I]{v15-height-Stacks}.
This statement uses the dualizing \emph{module}; it does not require
that the normal surface have an invertible dualizing sheaf.
Because $B$ is finite projective, specializing the dual at $t=2$
gives $\operatorname{Hom}_{\OK}(C,\OK)\,dt$.
On the generic fibre the identification sends a differential
$f\,dt$ to the functional $b\mapsto\operatorname{Tr}(fb)$.
After splitting the fibre, its coordinates are exactly the values
$f(P)=\omega/dt(P)$. This proves the assertion.
\end{proof}

\begin{proof}[Proof of Theorem~\ref{q214src:eval:thm:evaluation}]
Riemann--Hurwitz gives
\[
 2g-2=-2d+\deg\operatorname{Ram}(\pi)\le d-3,
\]
since each of the three branch fibres contributes at most $d-1$.
In particular $d>2g-2$. A nonzero holomorphic differential cannot
vanish at all $d$ distinct points of $\pi^{-1}(2)$. Therefore the
evaluation map has rank $g$. Choose $g$ coordinate positions $I$
for which projection remains injective. By
Lemmas~\ref{q214src:eval:lem:trace-lattice} and~\ref{q214src:eval:lem:integral-evaluation},
we obtain an integral, generically invertible map
\[
 \phi:E\longrightarrow F=\operatorname{pr}_I(C^\vee),
 \qquad\Ardeg(F)\le\tfrac12\log\mathrm N\mathfrak D(C).
\]

Fix a complex embedding of $K$. The disc
$D=\{t:|t-2|<1/2\}$ is simply connected and avoids the branch
locus. Its inverse image is a disjoint union of $d$ analytic
sheets, each mapped isomorphically to $D$. On one such sheet write
$\omega=f(t)\,dt$. The mean-value inequality implies
\[
 |f(2)|^2\le\frac4\pi\int_D|f(t)|^2\,dx\,dy
       \le\frac4\pi\|\omega\|_{L^2}^2.
\]
Hence the norm of each evaluation functional is at most
$2/\sqrt\pi$. In an orthonormal basis of $E\otimes_\sigma\C$,
each row of the evaluation matrix has this same bound on its
Euclidean norm. Hadamard's determinant inequality gives
\[
 \|\det\phi\|_\sigma\le(2/\sqrt\pi)^g.
\]
The determinant is nonzero and integral. Taking its arithmetic
degree yields
\[
 \Ardeg(E)\le\Ardeg(F)
       +\sum_\sigma\log\|\det\phi\|_\sigma
 \le\tfrac12\log\mathrm N\mathfrak D(C)
       +[K:\Q]g\log(2/\sqrt\pi).
\]
Now apply \eqref{q214src:eval:eq:height-convention}.
\end{proof}

\subsection{The discriminant of the fibre}

Let $S$ be a finite set of rational primes, and put
\[
 R_S=\prod_{p\in S}p,\qquad \omega(S)=\#S.
\]
Assume that the normalization defining the Belyi cover is finite
\'etale over
\begin{equation}\label{q214src:eval:eq:good-open}
 \operatorname{Spec}\mathcal O_{K_0}
       [S^{-1},t,t^{-1},(t-1)^{-1}],
\end{equation}
where $K_0$ is a field of definition.
This property persists after any finite field extension:
finite \'etale base change is finite \'etale, and a finite \'etale
algebra over a normal base is normal. Thus it agrees with the
normalization of the new generic cover on this open set.
In particular the auxiliary splitting field in Section~1 introduces
no new vertical branch primes outside $S$.

\begin{lemma}\label{q214src:eval:lem:discriminant}
Under \eqref{q214src:eval:eq:good-open}, the order $C$ in \eqref{q214src:eval:eq:order}
satisfies
\[
 \frac{\log\mathrm N\mathfrak D(C)}{[K:\Q]}
 \le d\left(\log2+\log R_S+\omega(S)\log d\right).
\]
\end{lemma}
\begin{proof}
The determinant of the trace pairing of $B/\OK[t]$ defines an
effective discriminant divisor on the base. Formation of the trace
pairing commutes with specialization. Its intersection with the
section $t=2$ is therefore $\log\mathrm N\mathfrak D(C)$.
The coefficient of this divisor at a codimension-one point is
$\sum_j f_jr_j$, where $e_j,f_j,r_j$ are the ramification index,
residue degree, and different exponent of the corresponding
extensions of discrete valuation rings. They satisfy
$\sum_j e_jf_j=d$. These are the usual local
trace-discriminant and different identities; see
\cite[Chapter~III, Sections~3 and~6]{v15-height-Serre} and
\cite[Section~4.2]{v15-height-Javanpeykar}.

At a horizontal branch divisor the residue characteristic is zero.
The different exponent is $r_j=e_j-1$, so the discriminant
coefficient is at most $d$. The section $t=2$ meets $t=0$
with normalized intersection $\log2$, and meets neither $t=1$
nor $t=\infty$. Thus the horizontal contribution is at most
$d\log2$.

Let $\mathfrak p\mid p$ be a prime of $K$, and write
$e_{\mathfrak p}=v_{\mathfrak p}(p)$.
Lenstra's different bound, in the form
\cite[Proposition~4.1.3]{v15-height-Javanpeykar}, gives
\[
 r_j\le e_j-1+e_j e_{\mathfrak p}
                \left\lfloor\frac{\log d}{\log p}\right\rfloor.
\]
This bound also applies to the possibly imperfect residue field
at the generic point of a vertical fibre.
Its discriminant coefficient is consequently at most
\[
 d\left(1+e_{\mathfrak p}
                \left\lfloor\frac{\log d}{\log p}\right\rfloor\right).
\]
Only $p\in S$ can contribute. The intersection of the vertical
fibre over $\mathfrak p$ with the section has value
$f_{\mathfrak p}\log p$. Using
\[
 \sum_{\mathfrak p\mid p}e_{\mathfrak p}f_{\mathfrak p}=[K:\Q],
 \qquad \sum_{\mathfrak p\mid p}f_{\mathfrak p}\le[K:\Q],
\]
and
$\lfloor\log d/\log p\rfloor\log p\le\log d$
proves the estimate.
\end{proof}

\begin{theorem}\label{q214src:eval:thm:coverheight}
For a Belyi cover satisfying \eqref{q214src:eval:eq:good-open},
\begin{equation}\label{q214src:eval:eq:coverheight}
 \hF(\Jac X)
 \le\frac d2\left(\log2+\log R_S+\omega(S)\log d\right)
       +g\log\frac2{\sqrt\pi}.
\end{equation}
\end{theorem}
\begin{proof}
Combine Theorem~\ref{q214src:eval:thm:evaluation} and
Lemma~\ref{q214src:eval:lem:discriminant}.
\end{proof}

\subsection{The integral modular cover}

Let $N$ be a positive integer and let
$H\subset(\Z/N\Z)^\times$ contain $-1$. Write
\[
 X_H=X_1(N)/H,\qquad t_H=[(\Z/N\Z)^\times:H],
 \qquad J_H=\Jac X_H.
\]
The action is the diamond action. Assume that $X_H$ has positive
genus, and use the rational map $t=j/1728$ to the projective line.

\begin{lemma}\label{q214src:eval:lem:modular-model}
The map $X_H\to\mathbb P^1_t$ has degree
\[
 d_H=Nt_H\prod_{p\mid N}(1+p^{-1}).
\]
It is unramified outside $0,1,\infty$ and admits a finite \'etale
model over
\[
 U=\operatorname{Spec}\Z[1/(6N),t,t^{-1},(t-1)^{-1}].
\]
\end{lemma}
\begin{proof}
Over $U$ consider
\begin{equation}\label{q214src:eval:eq:universal-curve}
 E_t:\quad y^2=x^3-\frac{3t}{t-1}x-\frac{2t}{t-1}.
\end{equation}
Its discriminant and $j$-invariant are
\[
 \Delta(E_t)=\frac{1728t^2}{(t-1)^3},\qquad j(E_t)=1728t.
\]
Thus it is an elliptic curve over $U$. Since $N$ is invertible,
its scheme of points of exact order $N$ is finite \'etale.
The constant group $H$ acts freely on this scheme: if $aP=P$
and $P$ has exact order $N$, then $a=1$ in $(\Z/N\Z)^\times$.
The quotient is finite \'etale. This follows also from the
classification of finite \'etale covers by locally constant finite
sets; freeness is enough, even when a residue characteristic divides
$|H|$.

To identify its generic fibre, fix an elliptic curve with
$j$-invariant different from $0$ and $1728$ over an algebraic
closure. Its automorphism group is $\{1,-1\}$. The geometric
fibre of $X_1(N)$ is the set of primitive torsion points modulo
this automorphism group. Since $-1\in H$, the fibre of $X_H$
is precisely the set of primitive torsion points modulo $H$.
Different isomorphisms to \eqref{q214src:eval:eq:universal-curve} differ by
$\pm1$ and hence give the same $H$-orbit. In particular quadratic
twisting introduces no extra double cover. This identifies the
finite \'etale model with the normalization of $X_H$ over $U$.

There are $N^2\prod_{p\mid N}(1-p^{-2})$ primitive points.
Division by $|H|=\varphi(N)/t_H$ gives the asserted degree.
Equivalently, the congruence subgroup contains $-I$, so its
$\operatorname{SL}_2(\Z)$ index is already the degree of the
coarse map to $X(1)$; see \cite[Section~4]{v15-height-vonKaenel}.
The only possible branch points of that map are the two elliptic
points and the cusp, which become $0,1,\infty$ in the $t$ coordinate.
\end{proof}

\begin{lemma}\label{q214src:eval:lem:torsion-different}
Let $V$ be an excellent discrete valuation ring of characteristic zero,
let $E/V$ be an elliptic scheme, and let $n$ be a positive integer.
Write $E[n]=\operatorname{Spec} A$, and let $\widetilde A$ be its
normalization in its generic finite \'etale algebra. Under the generic
trace pairing, the complementary modules satisfy
\[
 A^\vee=n^{-1}A,\qquad
 \widetilde A^\vee\subseteq n^{-1}\widetilde A.
\]
The different of every field component of $\widetilde A/V$ therefore
contains $n$. The same assertion holds for every intermediate field
of such a component.
\end{lemma}
\begin{proof}
Multiplication $[n]:E\to E$ is finite flat and a local complete
intersection morphism. Its derivative in invariant differential bases
is multiplication by $n$. Consequently its relative differential
module has zeroth Fitting ideal $(n)$. For a finite flat local complete
intersection, this K\"ahler different is the different; see
\cite[Tag~0BWB, Lemmas~49.12.3 and~49.12.6]{v15-height-Stacks}.

Base change along the zero section gives $E[n]/V$. The relative
differential module and its zeroth Fitting ideal commute with base
change, while the morphism remains finite flat and a local complete
intersection. Its different is therefore $(n)$. Its relative
dualizing module is invertible. Under the trace pairing this is the
complementary fractional ideal $A^\vee$, whose inverse is the different;
see \cite[Tag~0BW0, Lemmas~49.8.2--49.8.3, and Tag~0BWB,
Lemma~49.12.4]{v15-height-Stacks}. Thus $A^\vee=n^{-1}A$ as fractional ideals,
not merely as abstract modules.

Normalization is finite, and inclusion reverses trace duals:
\[
 \widetilde A^\vee\subseteq A^\vee=n^{-1}A
                              \subseteq n^{-1}\widetilde A.
\]
Hence $n\widetilde A^\vee\subseteq\widetilde A$, which says exactly
that the different contains $n$. Finally, for normal discrete
valuation rings $V\subseteq W\subseteq U$, the different formula
\[
 \mathfrak D_{U/V}=\mathfrak D_{U/W}\mathfrak D_{W/V}U
\]
shows that $v_U(\mathfrak D_{U/V})\le v_U(n)$ implies
$v_W(\mathfrak D_{W/V})\le v_W(n)$, because the different exponent
of $U/W$ is nonnegative. Apply this in each field component and
at each place. These finite-duality and different arguments do not
require a perfect residue field.
\end{proof}

\begin{corollary}\label{q214src:eval:cor:modularheight}
Let $g_H$ be the genus of $X_H$. Then
\begin{equation}\label{q214src:eval:eq:modularheight}
 \hF(J_H)\le
 \frac{d_H}{2}\left(\log12+\log N+2\log d_H\right)
                  +g_H\log\frac2{\sqrt\pi}.
\end{equation}
\end{corollary}
\begin{proof}
Use the auxiliary field $K$ from Theorem~\ref{q214src:eval:thm:evaluation}.
At a prime $\mathfrak p$ of $K$ over $p>3$, localize the $t$-line
at the generic point of the special fibre. The family
\eqref{q214src:eval:eq:universal-curve} is an elliptic scheme over this discrete
valuation ring: its discriminant $1728t^2/(t-1)^3$ is a unit there.
By Lemma~\ref{q214src:eval:lem:modular-model}, the generic cover for $X_H$ is
the primitive $N$-torsion cover modulo $H$. Each of its field
components is an intermediate field of a field generated by a single
$N$-torsion point. Lemma~\ref{q214src:eval:lem:torsion-different} therefore bounds
every different exponent above $\mathfrak p$ by the valuation of $N$.
The coefficient of the vertical discriminant divisor is at most
\[
                    d_H v_{\mathfrak p}(N).
\]
After summing over these primes and dividing by $[K:\Q]$, the
contribution is at most $d_H\log N$.

At $p=2,3$, the estimate in Lemma~\ref{q214src:eval:lem:discriminant} gives at
most $d_H(\log6+2\log d_H)$. The horizontal intersection at
$t=2$ contributes at most $d_H\log2$. Thus the normalized
logarithmic discriminant of the fibre order is at most
\[
                 d_H(\log12+\log N+2\log d_H).
\]
Apply Theorem~\ref{q214src:eval:thm:evaluation}. In particular no degree of a
full torsion splitting field enters this estimate.
\end{proof}
\section{Retained changed-factor extraction proof}\label{q214:factor-source}
\subsection{Five factors and extraction at one place}
At a contact prime $p\notin\Sigma$, $\delta$ is a unit and
$\nu=v_p(R)=v_p(F_2(t_*))$ is positive and odd, and $v_p(q)=-\nu$. This follows
from the unique contact factor with positive odd valuation and
does not require any condition on the valuation of the scale. All valuations
of $K$ are unramified and normalized by $v(p)=1$.
For a point on $E$, every finite value of $v(f)$ is even: if
$v(x)<0$ then $2v(y)=3v(x)$, while if $v(x)>0$ then
$2v(y)=v(x)$; the remaining case is valuation zero.
This also holds for conjugated reciprocal abscissas. The common-point lemma identifies $F_0=f(H-P)$ on every solution,
including the exceptional leading-coefficient cases. Its even valuation
is therefore available before the exclusion of torsion classes and
before any formal-parameter quotient is used.

Use
\begin{align*}
 T_1&=1+q^5C_2^2+q^2D_H^2+q^4B_1^4,\\
 T_2&=1+q^5G^2+q^2A_3^2+q^4B_0^4,\\
 T_3&=1+q^5I^4+q^3A_2^2+q^2A_4^2+q^4B_3^4,\\
 T_4&=1+q^5C_3^2+q^2C_1^4+q^4B_4^4,\\
 T_5&=1+q^5C_4^2+q^2A_1^2+q^4B_2^4,
 \qquad W=T_1T_2T_3T_4T_5.
\end{align*}

\begin{lemma}\label{q214src:factor:lem:residues}
At a contact prime, every factor is nonzero and has valuation at most
zero. A factor of valuation zero is congruent to one in the residue field.
\end{lemma}
\begin{proof}
First suppose every displayed test in a factor is nonzero. In each
factor other than $T_3$, the three nonconstant valuations belong,
respectively, to the odd integers, to $2+4\Z$, and to $4+8\Z$.
For example $T_1$ has valuations
\[
 2v(C_2)-5\nu,\qquad 2v(D_H)-2\nu,\qquad4v(B_1)-4\nu.
\]
The second has exact factor two because $v(D_H)$ is even and $\nu$ is
odd. The third has exact factor four because $v(B_1)$ is even.
The term $q^2C_1^4$ in $T_4$ has valuation
$2(2v(C_1)-\nu)$, again twice an odd integer.

In $T_3$ the four nonconstant valuations are
\[
 4v(I)-5\nu,\quad 2v(A_2)-3\nu,\quad
 2v(A_4)-2\nu,\quad4v(B_3)-4\nu.
\]
The first three have residues $-\nu,\nu,2$ modulo four. These are
pairwise distinct and nonzero because $\nu$ is odd and the $A_i$
have even valuations. The fourth has residue four modulo eight;
it therefore differs from each of the first three and from zero.
Thus, in every factor, all finite term valuations, including zero
for the constant term, are distinct. Zero tests simply remove terms
from this argument. The minimum is attained exactly once. This proves
nonvanishing and the asserted bound. If the minimum is zero, every
nonconstant term has positive valuation, so the residue is one.
\end{proof}

\begin{proposition}\label{q214src:factor:prop:extract}
At each contact prime, the retained pole system and the cleared trace
equations imply $\nu>0$ and give all sixteen required test orders at
one place above that prime.
\end{proposition}
\begin{proof}
Write $W=w_0+w_1\rho$. The literal cleared trace equations are
\[
 w_0^2-2w_0\tau-2w_1^2=0,\qquad (1+q)\mathfrak b=\gamma^2\tau^2.
\]
By Lemma~\ref{q214src:factor:lem:residues}, $W\ne0$. The first equation therefore
gives $w_0\ne0$ and $\tau=1/\Tr(W^{-1})\ne0$; indeed $w_0=0$
would give $w_1=0$. Hence
$\mathfrak b=\gamma^2R/\Tr(W^{-1})^2$. The factor $\gamma$ is a unit at every consumed contact prime. If
$\ell=\max_{v\mid p}v(W)\le0$, the trace inequality gives
$v_p(\mathfrak b)\le\nu+2\ell$.
This valuation is odd; the retained pole obstruction forces it positive.
Consequently $\nu>0$ and $\ell>-\nu/2$.
At a place attaining $\ell$, the factor valuations are nonpositive
and add to $\ell$. Each factor valuation is therefore at least $\ell$,
and each of its term valuations is strictly greater than $-\nu/2$.

The resulting lower bounds, before integer rounding, are
\[
\begin{array}{c|c|c}
\text{test}&\text{strict lower bound}&\text{integer order}\\\hline
C_2,C_3,C_4,G&9\nu/4&3\\
I&9\nu/8&2\\
C_1&3\nu/8&1\\
D_H&3\nu/4&2\\
A_2&5\nu/4&1\\
A_1,A_3,A_4&3\nu/4&1\\
B_0,B_1,B_2,B_3,B_4&7\nu/8&1
\end{array}
\]
Here $\nu$ is a positive odd integer. The order two for $D_H$ uses
its even valuation. In fact each of the nine reciprocal-abscissa
tests has order at least two; order one suffices for its subsequent
formal-group use. Zero tests satisfy all the stated inequalities.
\end{proof}

\subsection{Integrality before the formal congruences}
At a contact place selected by Proposition~\ref{q214src:factor:prop:extract}, the
points $H,U_i$ are formal in the primary completion, and $H-P,V_i$
are formal in the conjugate completion. The reciprocal-abscissa
tests have positive even valuation, so every nonidentity point has
formal depth at least one. In the parameter $z=-x/y$, write
\[
 \log_E(Q)=z(Q)\mathcal A(Q),\qquad
 \mathcal A(Q)\equiv1\pmod{p^{4v(z(Q))}}.
\]
The value of the unit is taken to be one at the identity. This is
the constant-term formal logarithm estimate proved in~Section~\ref{v15:logs}.

The integral group identities
\[
                  [m]U_i=[2u_i]H,\qquad
                  [m-2]V_i=[2(u_i-n_i)](H-P)
\]
kill the torsion in $H=[m]R_0+T_0$. Their full logarithmic ratios
are $u_i/h$ and $(u_i-n_i)/(h-1)$, where
$(n_1,n_2,n_3,n_4)=(n,a,t,b)$. At least one of $h,h-1$ is a
$p$-adic unit, so at least one full ratio is integral. The formal
units preserve the valuation of its plain parameter ratio.
The positive-order test $C_i$ transfers integrality to the other
plain ratio, and then to the other full ratio.

Only after this transfer may the relative errors of order four
be used as absolute congruences. They give
\[
 L_i\equiv u_i/h\pmod{p^4},\qquad
 L_i'\equiv(u_i-n_i)/(h-1)\pmod{p^4}.
\]
Since $n_i=h(u_i/h)-(h-1)((u_i-n_i)/(h-1))$, the high cross tests
give $L_i\equiv n_i\pmod{p^3}$ for $i>1$. The first cross test
gives $L_1\equiv n\pmod p$, while $I$ gives $L_1\equiv\beta\pmod p$.
These arguments include identity numerator points, half-integral $h$,
and inert primes, where conjugation preserves the unique valuation.

\clearpage
% Source: research/version15/quartics/quartic-222/provenance/shared-T-redundancy224-006/PROOF.md
\section{Two redundant reference equations}\label{q214:raw224-source}
\noindent\textit{Source: }\nolinkurl{research/version15/quartics/quartic-222/provenance/shared-T-redundancy224-006/PROOF.md}.\par

This author-frozen successor removes Q5's two rational components from accepted raw224, retaining224 variables and198 residuals. Coordinates and parameter/input are unchanged. The parent manifest is \nolinkurl{a907e6dab18305d6b0f9443578c123d854eda6643d6e6c8b83db5c8dcaef0487} and its independent acceptance is7f79535dd88fcec653370ce495eaef4d4080e8791bdb40bc7cf205ac3710639f. The normative \(Q\) numbering and coefficients are the parent's \nolinkurl{FIELD-RELATIONS.json} first entry. No new identity is inferred from finite samples.

Write \(Q_{2}=M*w+L\). Its coefficient is
\(M=B+(23335979380411305/9142573010269636-r*256685364214285/4571286505134818)f+(44649296511110307/111996519375803041-r*2886885451827729/447986077503212164)S+3317607896141091533/7167777240051394624+r*677690583653850567/1791944310012848656\).

The coefficients are149-integral and the coefficient of \(S\) is a149-unit; \nolinkurl{COEFFICIENT-DOMAIN.json} gives exact norm fractions and residues. \(B=f^{2}\) is retained.

For any new solution, define \(g\) by the same polynomial as in the parent224 proof. The parent \nolinkurl{CONVERSE-RESULT.json} \(F\) and \(E\) certificates use only Q0,Q1,Q3, so these old reference equations already hold without \(Q_{2}\) orQ5. The same direct \(f=0\) contradiction also uses onlyQ0,Q1,Q3. Thus \(f\) is nonzero and the old \(H\) curve is recovered before interpreting \(w\). The unchanged scalar normalizer/core/pole equations give the inherited forced149 bound on Norm(f), independently of referencew and its outputW. Therefore \(\operatorname{ord}(f)=m>0\), \(\operatorname{ord}(S)=-m/2<0\), and \(\operatorname{ord}(g)=m/2\). Also \(A=a+cf\) is a unit, and \(D=\beta *S+\alpha \) is nonzero. No deleted relation has entered this argument.

The \(S\) term is the unique lowest-valuation term in \(M:\) \(B\) has order2m, the \(f\) term has orderm, and the constant has \(\text{order}\ge 0\). Hence \(M\ne 0\). Now form the expected value \(w*=g^{4}/f^{2}\). The parent's exact forward identity for \(Q_{2}\) gives \(Q_{2}(f,S,B,w*,T)=0\); it uses only the \(H\) curve, the reconstructed F/E, and \(D*T=N\), with f,A,D nonzero, not \(Q_{5}\). Since \(Q_{2}\) is linear in \(w\) with nonzero coefficientM, the retained \(Q_{2}=0\) forces \(w=w*\). Retained \(Q_{4}\) then forces \(W=w^{2}\). The parent's exact forward \(Q_{5}\) identity therefore gives \(Q_{5}=0\). All parent equations are recovered.

Old to new simply drops two satisfied equations. New to old is the identity on every coordinate and reconstructs those equations by the preceding argument. Thus there is no lost zero branch. The admissible \(A_{N}\) family and ineffective defining threshold are essential inherited qualifications; this is not asserted for arbitrary symbolicA.

Exact full expansion gives224 variables,198 residuals,degree4,content1,3014 coordinate monomials and3015 parameter terms. All coordinates occur. The coefficient height remains \(A_{N}^{2}+1\) for \(N\ge 53\), which is only a coefficient threshold. No manuscript adoption or independent acceptance is claimed by this author packet.
\clearpage\section{Finite certificates and exact computational sources}\label{q214:certificates}
The following files identify finite inputs, methods and conclusions used in Part~I.  Their byte hashes bind the exact public objects.  The links provide the full recoverable certificates or implementations.  Their finite conclusions are distinct from the arithmetic theorem and from review acceptance.
\medskip\noindent\textbf{\nolinkurl{sources/contact-line-certificate.json}}\par\noindent
The four quotient-quartic invariant checks and all $36$ cross-ratio candidate pairs used to recover the branch parameters $(2401,-343/32)$. The separate small-support certificate lists the five primitive forms and their ten pair determinants; the geometric proof supplies the cover and its tame charts.
\par\noindent SHA-256: \nolinkurl{64896dfbb8661e04c044a68aa8679ee5dc20c0da36cc9a57640a885b7ca46614}.\par
\href{https://github.com/elliotglazer/diophantine-z-q/blob/110cf72565be326fd567cd1a64a3b2e00213dcd1/research/quartic214/sources/contact-line-certificate.json}{Read the exact public source}.\par
\medskip\noindent\textbf{\nolinkurl{sources/small-support-certificate.json}}\par\noindent
The small exceptional support and exact finite support computations used by the full-family height argument.
\par\noindent SHA-256: \nolinkurl{0e7280479d5ace418b4383c814361812863cb097f424916b801c17ddbc8eb51e}.\par
\href{https://github.com/elliotglazer/diophantine-z-q/blob/110cf72565be326fd567cd1a64a3b2e00213dcd1/research/quartic214/sources/small-support-certificate.json}{Read the exact public source}.\par
\medskip\noindent\textbf{\nolinkurl{sources/public-provider-scopes.json}}\par\noindent
The exact public provider/source scopes used by the geometric and conductor carry arguments; this is source identification, not a fresh approval.
\par\noindent SHA-256: \nolinkurl{4b20650915af34d33413425c17131136be22698f09a0f125b011423b87df3654}.\par
\href{https://github.com/elliotglazer/diophantine-z-q/blob/110cf72565be326fd567cd1a64a3b2e00213dcd1/research/quartic214/sources/public-provider-scopes.json}{Read the exact public source}.\par
\medskip\noindent\textbf{\nolinkurl{height/NUMERIC_CERTIFICATE_H4.json}}\par\noindent
The rational logarithmic enclosures and strict \(H_{4}\) inequalities, including the exponent 2103/1024 and the bound on \(C_{2}\).
\par\noindent SHA-256: \nolinkurl{cd4afb85267aad1a22c0e6a765f86551d2c83a94fcece93939547fe7c06d0962}.\par
\href{https://github.com/elliotglazer/diophantine-z-q/blob/110cf72565be326fd567cd1a64a3b2e00213dcd1/research/quartic214/height/NUMERIC_CERTIFICATE_H4.json}{Read the exact public source}.\par
\medskip\noindent\textbf{\nolinkurl{height/numeric_certificate.py}}\par\noindent
The atanh-series rational enclosure method with its proved geometric tail and negative controls for the isolated \(H_{4}\) calculation.
\par\noindent SHA-256: \nolinkurl{2d6b1426b497b197782285280845ce4663936d4dee00f2ce3c5173bc58d39a1e}.\par
\href{https://github.com/elliotglazer/diophantine-z-q/blob/110cf72565be326fd567cd1a64a3b2e00213dcd1/research/quartic214/height/numeric_certificate.py}{Read the exact public source}.\par
\medskip\noindent\textbf{\nolinkurl{construction/EXACT-CHECKS.json}}\par\noindent
Recorded finite Laurent identities, row occurrence and polynomial-profile checks for the 214 construction.
\par\noindent SHA-256: \nolinkurl{0a160c8a7d0519621b4ebc27fb6e8cb2b8df2a9d25191b8daa2170db461a12b0}.\par
\href{https://github.com/elliotglazer/diophantine-z-q/blob/110cf72565be326fd567cd1a64a3b2e00213dcd1/research/quartic214/construction/EXACT-CHECKS.json}{Read the exact public source}.\par
\medskip\noindent\textbf{\nolinkurl{construction/check_exact.py}}\par\noindent
The separate implementation reconstructing deleted coordinates and comparing every changed and untouched row; it also verifies the full quartic and coefficient profiles.
\par\noindent SHA-256: \nolinkurl{07527831694fc6be9373e2dac92d292520fbb8c79128e4d51a1ae0128c96e214}.\par
\href{https://github.com/elliotglazer/diophantine-z-q/blob/110cf72565be326fd567cd1a64a3b2e00213dcd1/research/quartic214/construction/check_exact.py}{Read the exact public source}.\par
\medskip\noindent\textbf{\nolinkurl{construction/build.py}}\par\noindent
The deterministic raw224-to214 generator.  It fixes the parameter coefficient, prints the IR and expands the finite sum of squares; the manuscript assembler itself does not run this expansion.
\par\noindent SHA-256: \nolinkurl{09d958a38e5815a03b474033d439e84fec40b8a55f4109e95adb87b179fb8c17}.\par
\href{https://github.com/elliotglazer/diophantine-z-q/blob/110cf72565be326fd567cd1a64a3b2e00213dcd1/research/quartic214/construction/build.py}{Read the exact public source}.\par
\medskip\noindent\textbf{\nolinkurl{construction/artifacts/candidate214-brox/candidate214-ir.json}}\par\noindent
The exact 214-coordinate dictionary and 188 primitive integral rows printed in Section \ref{q214:rows}.
\par\noindent SHA-256: \nolinkurl{b15d6477439c3e4be1e98723b69053244debce59e8715683c28fee82aa847ca4}.\par
\href{https://github.com/elliotglazer/diophantine-z-q/blob/110cf72565be326fd567cd1a64a3b2e00213dcd1/research/quartic214/construction/artifacts/candidate214-brox/candidate214-ir.json}{Read the exact public source}.\par
\medskip\noindent\textbf{\nolinkurl{construction/artifacts/candidate214-brox/candidate214-polynomial.json}}\par\noindent
The complete symbolic quartic coefficient table: a term is \(c A_N^a\) times its coordinate multiset monomial, with \(A_N\) a coefficient and no additional coordinate.
\par\noindent SHA-256: \nolinkurl{9f461531df52ac73e02cf7a4ce270422715642c0d3112952cad4968d9a7e24ea}.\par
\href{https://github.com/elliotglazer/diophantine-z-q/blob/110cf72565be326fd567cd1a64a3b2e00213dcd1/research/quartic214/construction/artifacts/candidate214-brox/candidate214-polynomial.json}{Read the exact public source}.\par
\medskip\noindent\textbf{\nolinkurl{construction/artifacts/candidate214-brox/STATISTICS.json}}\par\noindent
The declared finite totals 214,188,2919,2920, degree four, content one and coefficient-only threshold53.
\par\noindent SHA-256: \nolinkurl{585635dab6320049320bffd5c1652c89db760d9f4e2ef19a10119011f70b5a68}.\par
\href{https://github.com/elliotglazer/diophantine-z-q/blob/110cf72565be326fd567cd1a64a3b2e00213dcd1/research/quartic214/construction/artifacts/candidate214-brox/STATISTICS.json}{Read the exact public source}.\par
\clearpage\section{Primary-source editions and inspection catalogue}\label{q214:catalogue}
\begingroup\raggedright
This catalogue retains the complete primary-source metadata accompanying the author packet.  An inspection statement describes its original named author and date.  It is not an assertion that the typesetter repeated that inspection, nor an independent acceptance of the present manuscript.  External books and PDFs are identified by their editions, locators, URLs and byte hashes; their complete contents are not republished here.
\paragraph{height}
\noindent\textbf{format}: height-sources-v1\par
\noindent\textbf{author}: /root/sol61\_height\_foundation\par
\noindent\textbf{parentPrivateCommit}: \nolinkurl{4acd40881c4abb6da7017c803a8f8d4b851d54cf}\par
\noindent\textbf{broxCommit}: \nolinkurl{0697019486451cdfdb71b3272d5f26c6270d2a6b}\par
\noindent\textbf{proofSha256}: \nolinkurl{25960e40419c7abea1f90ab4b35108169a5ca6ad04dccf3a3643e8d0e9bb7d7f}\par
\noindent\textbf{sourceUseInitialSha256}: \nolinkurl{7f027bcc37a8b4e888dee6f8890f0093c2e1d71176a75ec5e88f800c6ece9818}\par
\paragraph{files}
\medskip\noindent\textbf{files entry 1.}\par
\noindent\textbf{url}: \url{https://smf.emath.fr/system/files/filepdf/smf_mem_176_0.pdf}\par
\noindent\textbf{accessDate}: 2026-11-10\par
\noindent\textbf{sha256}: \nolinkurl{b1e2e32bec626cf6d27fb7d6d8702b5c707e0c3dc002d979088e62d8cda514a4}\par
\noindent\textbf{status}: final published PDF checked\par
\noindent\textbf{locator}: Theorem1.9(1),pp7--8; original stable normalization\par
\noindent\textbf{omitted\_local\_path}: GR2023-published.pdf\par
\medskip\noindent\textbf{files entry 2.}\par
\noindent\textbf{url}: \url{https://www-fourier.univ-grenoble-alpes.fr/~remond/2476.pdf}\par
\noindent\textbf{accessDate}: 2026-11-10\par
\noindent\textbf{sha256}: \nolinkurl{e23a019a7432e299014ead57a1ba51c3945b8351a64c96f339536786c8f79c5d}\par
\noindent\textbf{status}: authorJune2021PD\allowbreak{}F checked; final2022body unavailable via subscription redirect\par
\noindent\textbf{locator}: Corollary1.2,stable remark,Proposition2.1\par
\noindent\textbf{omitted\_local\_path}: Remond2021-author.pdf\par
\medskip\noindent\textbf{files entry 3.}\par
\noindent\textbf{url}: \url{https://www.jstage.jst.go.jp/article/jmath1948/25/3/25_3_523/_pdf/-char/ja}\par
\noindent\textbf{accessDate}: 2026-11-10\par
\noindent\textbf{sha256}: \nolinkurl{6a34cc46d06c0f5d4e7d380ddcbc45ec3143b19359bfebba5b1c8f00976f3308}\par
\noindent\textbf{status}: final published PDF checked\par
\noindent\textbf{locator}: Theorem1,pp526--527 including Fricke condition\par
\noindent\textbf{omitted\_local\_path}: Shimura1973-published.pdf\par
\medskip\noindent\textbf{files entry 4.}\par
\noindent\textbf{url}: \url{https://math.berkeley.edu/~ribet/Articles/korea.pdf}\par
\noindent\textbf{accessDate}: 2026-11-10\par
\noindent\textbf{sha256}: \nolinkurl{4c491a5294d1f4ec1b62855560aaea95cb64802d8fdd80fdd51ae2d2432f66ed}\par
\noindent\textbf{status}: archived author PDF checked; final publication identity not asserted\par
\noindent\textbf{locator}: \S{}3\par
\noindent\textbf{omitted\_local\_path}: Ribet-author.pdf\par
\paragraph{applications}
\paragraph{O1-supplier}
\noindent $\bullet$ U01\par
\noindent $\bullet$ U02\par
\noindent $\bullet$ U03\par
\noindent $\bullet$ U04\par
\noindent $\bullet$ U05\par
\noindent $\bullet$ U06\par
\noindent $\bullet$ U07\par
\noindent $\bullet$ U08\par
\noindent $\bullet$ U09\par
\noindent $\bullet$ U10\par
\noindent $\bullet$ U11\par
\noindent $\bullet$ U16\par
\paragraph{O1-reflection}
\noindent $\bullet$ U11\par
\noindent $\bullet$ U12\par
\paragraph{O9-consumers}
\noindent $\bullet$ U13\par
\noindent $\bullet$ U14\par
\noindent $\bullet$ U15\par
\paragraph{closureLimits}
\noindent $\bullet$ Final Rémond2022body not compared\par
\noindent $\bullet$ Faltings/Chebotarev/Poincare/Li and underlying geometry/modularity/Arakelov inputs retain parent source-version qualifications\par
\noindent $\bullet$ No independent inventory-coverage decision supplied\par
\noindent\textbf{copyright}: Local third-party PDFs are source evidence. Do not infer redistribution permission; release may retain URL/hash/locator without PDF copies.\par
\paragraph{bridges}
\paragraph{sources}
\medskip\noindent\textbf{sources entry 1.}\par
\noindent\textbf{id}: Faltings-Rational-Points.pdf\par
\noindent\textbf{url}: \url{https://yhzhumath.github.io/Mordell2024/RP.pdf}\par
\noindent\textbf{sha256}: \nolinkurl{3122011c0f12cc991fc0ff1601bb77cc1cd3e5ab2a0db023c0bf224ea576b6aa}\par
\noindent\textbf{bytes}: 11049110\par
\noindent\textbf{locator}: Second edition1986, Theorem1a-b, printed204/PDF212; title pages checked\par
\noindent\textbf{accessDate}: 2026-11-10\par
\noindent\textbf{inspectedBy}: /root\par
\noindent\textbf{status}: published page image directly read; current application authored separately, awaiting independent review\par
\noindent\textbf{repositoryCopy}: false\par
\medskip\noindent\textbf{sources entry 2.}\par
\noindent\textbf{id}: Milne-Abelian-Varieties-1986.pdf\par
\noindent\textbf{url}: \url{https://www.jmilne.org/math/articles/1986b.pdf}\par
\noindent\textbf{sha256}: \nolinkurl{86259344a576e4ef75a9d81e9bf43143f872a72b14cf8fcc2703f3aea6349136}\par
\noindent\textbf{bytes}: 3833583\par
\noindent\textbf{locator}: Published1986, Proposition12.1, printed122/PDF20\par
\noindent\textbf{accessDate}: 2026-11-10\par
\noindent\textbf{inspectedBy}: /root\par
\noindent\textbf{status}: published page image directly read; current application authored separately, awaiting independent review\par
\noindent\textbf{repositoryCopy}: false\par
\medskip\noindent\textbf{sources entry 3.}\par
\noindent\textbf{id}: Li-p294.jpg\par
\noindent\textbf{url}: \url{https://images.p.sub.uni-goettingen.de/iiif/image/gdz:PPN235181684_0212:00000302/full/1800,/0/default.jpg}\par
\noindent\textbf{sha256}: \nolinkurl{9b31b9fc6f0b176bff3b0ef85f8f07285c0d40cb282cc80d20401e24a6b60d49}\par
\noindent\textbf{bytes}: 991846\par
\noindent\textbf{locator}: Mathematische Annalen212(1975), printed294\par
\noindent\textbf{accessDate}: 2026-11-10\par
\noindent\textbf{inspectedBy}: /root\par
\noindent\textbf{status}: published page image directly read; current application authored separately, awaiting independent review\par
\noindent\textbf{repositoryCopy}: false\par
\medskip\noindent\textbf{sources entry 4.}\par
\noindent\textbf{id}: Li-p295.jpg\par
\noindent\textbf{url}: \url{https://images.p.sub.uni-goettingen.de/iiif/image/gdz:PPN235181684_0212:00000303/full/1800,/0/default.jpg}\par
\noindent\textbf{sha256}: \nolinkurl{3c4882694168d10693977176097644198fc312259ce7715872b063264cf78633}\par
\noindent\textbf{bytes}: 629475\par
\noindent\textbf{locator}: Mathematische Annalen212(1975), printed295\par
\noindent\textbf{accessDate}: 2026-11-10\par
\noindent\textbf{inspectedBy}: /root\par
\noindent\textbf{status}: published page image directly read; current application authored separately, awaiting independent review\par
\noindent\textbf{repositoryCopy}: false\par
\medskip\noindent\textbf{sources entry 5.}\par
\noindent\textbf{id}: Li-p296.jpg\par
\noindent\textbf{url}: \url{https://images.p.sub.uni-goettingen.de/iiif/image/gdz:PPN235181684_0212:00000304/full/1800,/0/default.jpg}\par
\noindent\textbf{sha256}: \nolinkurl{11f38136cd27108d507c7662d1ddb46d8c572874b314b69aa1ad8be16395aad3}\par
\noindent\textbf{bytes}: 758672\par
\noindent\textbf{locator}: Mathematische Annalen212(1975), printed296\par
\noindent\textbf{accessDate}: 2026-11-10\par
\noindent\textbf{inspectedBy}: /root\par
\noindent\textbf{status}: published page image directly read; current application authored separately, awaiting independent review\par
\noindent\textbf{repositoryCopy}: false\par
\medskip\noindent\textbf{sources entry 6.}\par
\noindent\textbf{id}: GR2014-published.pdf\par
\noindent\textbf{url}: \url{https://ems.press/content/serial-article-files/43351?nt=1}\par
\noindent\textbf{sha256}: \nolinkurl{5afec329530089fd62d686f5e75331c91fc97fbf775ad6c4b275fc514459f31e}\par
\noindent\textbf{bytes}: 576654\par
\noindent\textbf{locator}: Definition2.1 and\S{}2.3 pp351--352; Corollary8.4 p400; PDFpages9--10,58\par
\noindent\textbf{accessDate}: 2026-11-10\par
\noindent\textbf{inspectedBy}: /root\par
\noindent\textbf{status}: Published definition,corollary proof and pageimages read; source/application bridge awaiting independent review\par
\noindent\textbf{repositoryCopy}: false\par
\noindent\textbf{qualification}: URLs and exact byte hashes identify external published evidence. Full third-party books and images are local review copies, not redistributed in the packet. Source/application discussion is PRIMARY\_SOURCE\_BRIDGES.md.\par
\noindent\textbf{remondFinal}: Publisher article link redirects to subscription; final2022body not compared. Existing authorJune2021 version remains frozen.\par
\paragraph{modularity}
\noindent\textbf{role}: AUTHOR only\par
\noindent\textbf{date}: 2026-11-11\par
\paragraph{primary sources}
\medskip\noindent\textbf{primary\_sources entry 1.}\par
\noindent\textbf{id}: Calegari2023\par
\noindent\textbf{sha256}: \nolinkurl{d8e9a3ec3f2a4bedd674d95cf5f1e03fe26a67745c6f1cadf39b3fbf9e004cb5}\par
\noindent\textbf{bytes}: 834759\par
\noindent\textbf{locator}: definitions611--612; statement624\par
\noindent\textbf{url}: \url{https://ems.press/content/book-chapter-files/33140?nt=1}\par
\noindent\textbf{publication}: EMS ICM2022 vol2 (2023),610--651\par
\noindent\textbf{actual\_inspection}: actual statement text read; image-only Milne and Faltings pages additionally visually inspected\par
\noindent\textbf{consumer\_proof\_sections}: M2\par
\noindent\textbf{authoritative\_version}: published source; Milne author-hosted published chapter/edition\par
\medskip\noindent\textbf{primary\_sources entry 2.}\par
\noindent\textbf{id}: Carayol1987\par
\noindent\textbf{sha256}: \nolinkurl{cc2605dd390eae903712a650e925b322472bb8831a11d6282b40030c0753944b}\par
\noindent\textbf{bytes}: 1162324\par
\noindent\textbf{locator}: \S{}\S{}1.4,2.1,2.4,3.4.2,4.1/fullTheorem4; printed36,37,41,42\par
\noindent\textbf{url}: \url{https://www.numdam.org/item/AST_1987__147-148__33_0.pdf}\par
\noindent\textbf{publication}: Astérisque147--148(1987),33--47\par
\noindent\textbf{actual\_inspection}: actual statement text read; image-only Milne and Faltings pages additionally visually inspected\par
\noindent\textbf{consumer\_proof\_sections}: M5\par
\noindent\textbf{authoritative\_version}: published source; Milne author-hosted published chapter/edition\par
\medskip\noindent\textbf{primary\_sources entry 3.}\par
\noindent\textbf{id}: Scholze2013\par
\noindent\textbf{sha256}: \nolinkurl{f47a05bbcdd1ac69deb87ce4b6af5bbbc245cbbf0e99df74e8bf73b6d89227d7}\par
\noindent\textbf{bytes}: 688182\par
\noindent\textbf{locator}: Cor1.8pp5--6; Thms9.1,9.3pp74--75\par
\noindent\textbf{url}: \url{https://www.cambridge.org/core/services/aop-cambridge-core/content/view/9C5E8381BE3CD320ECADDDEAD2616B55/S2050508613000012a.pdf}\par
\noindent\textbf{publication}: Forum Math.Pi1(2013),e1\par
\noindent\textbf{actual\_inspection}: actual statement text read; image-only Milne and Faltings pages additionally visually inspected\par
\noindent\textbf{consumer\_proof\_sections}: M2\par
\noindent\textbf{authoritative\_version}: published source; Milne author-hosted published chapter/edition\par
\medskip\noindent\textbf{primary\_sources entry 4.}\par
\noindent\textbf{id}: Scholze2016\par
\noindent\textbf{sha256}: \nolinkurl{3caea8d89df1e7d047d4edc89a679cd1deb2a19ec350effb450305a9c8eef3d1}\par
\noindent\textbf{bytes}: 166017\par
\noindent\textbf{locator}: whole4-page correction, especiallyp2\par
\noindent\textbf{url}: \url{https://www.cambridge.org/core/services/aop-cambridge-core/content/view/423022854D448D337C999BCD505D8F75/S2050508616000044a.pdf}\par
\noindent\textbf{publication}: Forum Math.Pi4(2016),e6\par
\noindent\textbf{actual\_inspection}: actual statement text read; image-only Milne and Faltings pages additionally visually inspected\par
\noindent\textbf{consumer\_proof\_sections}: M2\par
\noindent\textbf{authoritative\_version}: published source; Milne author-hosted published chapter/edition\par
\medskip\noindent\textbf{primary\_sources entry 5.}\par
\noindent\textbf{id}: Fontaine1994\par
\noindent\textbf{sha256}: \nolinkurl{ace926ffda578dacfe8a333c38d450235506ebe18f558ee9bfcb96dc6f70a8ec}\par
\noindent\textbf{bytes}: 6497435\par
\noindent\textbf{locator}: \S{}\S{}3.4--3.8; Theorem3.8p142\par
\noindent\textbf{url}: \url{https://www.numdam.org/item/AST_1994__223__113_0.pdf}\par
\noindent\textbf{publication}: Astérisque223(1994),113--184\par
\noindent\textbf{actual\_inspection}: actual statement text read; image-only Milne and Faltings pages additionally visually inspected\par
\noindent\textbf{consumer\_proof\_sections}: M2--M3\par
\noindent\textbf{authoritative\_version}: published source; Milne author-hosted published chapter/edition\par
\medskip\noindent\textbf{primary\_sources entry 6.}\par
\noindent\textbf{id}: Kato2004\par
\noindent\textbf{sha256}: \nolinkurl{3c6e14b11fa60262db8aff782ce3cf4d83e9100c0be83621a7e4ce502cec605d}\par
\noindent\textbf{bytes}: 18010813\par
\noindent\textbf{locator}: \S{}\S{}5.4,6.3,8.3,9.2,9.4,14.10.2--3; pp155,161--162,181,187--188,240\par
\noindent\textbf{url}: \url{https://www.numdam.org/item/AST_2004__295__117_0.pdf}\par
\noindent\textbf{publication}: Astérisque295(2004),117--290\par
\noindent\textbf{actual\_inspection}: actual statement text read; image-only Milne and Faltings pages additionally visually inspected\par
\noindent\textbf{consumer\_proof\_sections}: M2--M3\par
\noindent\textbf{authoritative\_version}: published source; Milne author-hosted published chapter/edition\par
\medskip\noindent\textbf{primary\_sources entry 7.}\par
\noindent\textbf{id}: DeligneSerre1974\allowbreak{}\par
\noindent\textbf{sha256}: \nolinkurl{65b390f6d33e827e30c6c66bbc15421eca51db3180bdf5996dcee19047be97fc}\par
\noindent\textbf{bytes}: 2402609\par
\noindent\textbf{locator}: \S{}\S{}1.6--1.9,3.1--3.2,Theorem4.1,8.2,8.5--8.6; pp509--510,513--514,525--526\par
\noindent\textbf{url}: null\par
\noindent\textbf{publication}: Ann.Sci.ENS(4)7(1974),507--530\par
\noindent\textbf{actual\_inspection}: actual statement text read; image-only Milne and Faltings pages additionally visually inspected\par
\noindent\textbf{consumer\_proof\_sections}: M3--M4\par
\noindent\textbf{authoritative\_version}: published source; Milne author-hosted published chapter/edition\par
\medskip\noindent\textbf{primary\_sources entry 8.}\par
\noindent\textbf{id}: Milne1986\par
\noindent\textbf{sha256}: \nolinkurl{86259344a576e4ef75a9d81e9bf43143f872a72b14cf8fcc2703f3aea6349136}\par
\noindent\textbf{bytes}: 3833583\par
\noindent\textbf{locator}: Prop12.1p122; Prop12.9p125; Thm15.1(a)p129; scanpages20,23,27 visually read\par
\noindent\textbf{url}: null\par
\noindent\textbf{publication}: ArithmeticGeomet\allowbreak{}ry,Springer(1986),103--150\par
\noindent\textbf{actual\_inspection}: actual statement text read; image-only Milne and Faltings pages additionally visually inspected\par
\noindent\textbf{consumer\_proof\_sections}: M2,M4,M6\par
\noindent\textbf{authoritative\_version}: published source; Milne author-hosted published chapter/edition\par
\medskip\noindent\textbf{primary\_sources entry 9.}\par
\noindent\textbf{id}: Faltings1986\par
\noindent\textbf{sha256}: \nolinkurl{3122011c0f12cc991fc0ff1601bb77cc1cd3e5ab2a0db023c0bf224ea576b6aa}\par
\noindent\textbf{bytes}: 11049110\par
\noindent\textbf{locator}: Complements to Mordell,Theorem1(a,b)p204; scanp212 visually read\par
\noindent\textbf{url}: null\par
\noindent\textbf{publication}: RationalPoints,secondedition,Vieweg(1986)\par
\noindent\textbf{actual\_inspection}: actual statement text read; image-only Milne and Faltings pages additionally visually inspected\par
\noindent\textbf{consumer\_proof\_sections}: M2,M6\par
\noindent\textbf{authoritative\_version}: published source; Milne author-hosted published chapter/edition\par
\medskip\noindent\textbf{primary\_sources entry 10.}\par
\noindent\textbf{id}: MilneADT2006\par
\noindent\textbf{sha256}: \nolinkurl{2c6195ec76a974f3f336c77cb71cc3845b018aad2d43136716b3477a3bc5fb31}\par
\noindent\textbf{bytes}: 2136045\par
\noindent\textbf{locator}: I.1.3,1.5,1.6(b,c)pp18--19; II.2.1proofp164\par
\noindent\textbf{url}: \url{https://www.jmilne.org/math/Books/ADTnot.pdf}\par
\noindent\textbf{publication}: ArithmeticDualit\allowbreak{}yTheorems,2ndedition,BookSurge2006\par
\noindent\textbf{actual\_inspection}: actual statement text read; image-only Milne and Faltings pages additionally visually inspected\par
\noindent\textbf{consumer\_proof\_sections}: M1,M3\par
\noindent\textbf{authoritative\_version}: published source; Milne author-hosted published chapter/edition\par
\medskip\noindent\textbf{primary\_sources entry 11.}\par
\noindent\textbf{id}: MilneShimura2005\allowbreak{}\par
\noindent\textbf{sha256}: \nolinkurl{6336f0a23084dc0bf815a901c61a0f925dbc8dc1c395af8864c1d90bea69c4a1}\par
\noindent\textbf{bytes}: 1035957\par
\noindent\textbf{locator}: complex uniformization and Theorem6.8,Cor6.9p319\par
\noindent\textbf{url}: \url{https://www.jmilne.org/math/xnotes/svi2005.pdf}\par
\noindent\textbf{publication}: ClayMath.Proc4(2005),265--378\par
\noindent\textbf{actual\_inspection}: actual statement text read; image-only Milne and Faltings pages additionally visually inspected\par
\noindent\textbf{consumer\_proof\_sections}: M2\par
\noindent\textbf{authoritative\_version}: published source; Milne author-hosted published chapter/edition\par
\medskip\noindent\textbf{primary\_sources entry 12.}\par
\noindent\textbf{id}: Shimura1973\par
\noindent\textbf{sha256}: \nolinkurl{6a34cc46d06c0f5d4e7d380ddcbc45ec3143b19359bfebba5b1c8f00976f3308}\par
\noindent\textbf{bytes}: 1646652\par
\noindent\textbf{locator}: Theorem1pp526--527; \S{}4equation4.1p531\par
\noindent\textbf{url}: \url{https://www.jstage.jst.go.jp/article/jmath/25/3/25_3_523/_pdf}\par
\noindent\textbf{publication}: J.Math.Soc.Japan25(1973),523--544\par
\noindent\textbf{actual\_inspection}: actual statement text read; image-only Milne and Faltings pages additionally visually inspected\par
\noindent\textbf{consumer\_proof\_sections}: M3,M6\par
\noindent\textbf{authoritative\_version}: published source; Milne author-hosted published chapter/edition\par
\noindent\textbf{access}: Copied exact already-frozen published PDF from private height packet after current JStage500; source text directly read\par
\medskip\noindent\textbf{primary\_sources entry 13.}\par
\noindent\textbf{id}: Serre1968-reprint1989\par
\noindent\textbf{sha256}: \nolinkurl{971dbfdd4414d99ea21d7d0af63dbb790d7b1188e9aec37f422181a2b543dec4}\par
\noindent\textbf{bytes}: 5563667\par
\noindent\textbf{locator}: ChapterI\S{}2.2Corollary2pI-8;\S{}2.3theorem/lemmapI-10--11\par
\noindent\textbf{url}: \url{https://www.math.mcgill.ca/darmon/courses/18-19/gs/serre-mcgill.pdf}\par
\noindent\textbf{publication}: Abelianl-AdicRepresentati\allowbreak{}onsandEllipticCu\allowbreak{}rves,1968,Addison-Wesley1989reprin\allowbreak{}t\par
\noindent\textbf{actual\_inspection}: actual statement text read; image-only Milne and Faltings pages additionally visually inspected\par
\noindent\textbf{consumer\_proof\_sections}: M4,M6\par
\noindent\textbf{authoritative\_version}: published source; Milne author-hosted published chapter/edition\par
\medskip\noindent\textbf{primary\_sources entry 14.}\par
\noindent\textbf{id}: Li1975-p294\par
\noindent\textbf{sha256}: \nolinkurl{9b31b9fc6f0b176bff3b0ef85f8f07285c0d40cb282cc80d20401e24a6b60d49}\par
\noindent\textbf{bytes}: 991846\par
\noindent\textbf{locator}: Math.Ann.212(1975),printed294\par
\noindent\textbf{url}: \url{https://images.p.sub.uni-goettingen.de/iiif/image/gdz:PPN235181684_0212:00000302/full/1800,/0/default.jpg}\par
\noindent\textbf{actual\_inspection}: p296,p298 visually read; p294,p295 acquired, prior bridge read; no false claim of visual reading p294/p295\par
\noindent\textbf{consumer\_proof\_sections}: M3precision finding; M6Fricke\par
\medskip\noindent\textbf{primary\_sources entry 15.}\par
\noindent\textbf{id}: Li1975-p295\par
\noindent\textbf{sha256}: \nolinkurl{3c4882694168d10693977176097644198fc312259ce7715872b063264cf78633}\par
\noindent\textbf{bytes}: 629475\par
\noindent\textbf{locator}: Math.Ann.212(1975),printed295\par
\noindent\textbf{url}: \url{https://images.p.sub.uni-goettingen.de/iiif/image/gdz:PPN235181684_0212:00000303/full/1800,/0/default.jpg}\par
\noindent\textbf{actual\_inspection}: p296,p298 visually read; p294,p295 acquired, prior bridge read; no false claim of visual reading p294/p295\par
\noindent\textbf{consumer\_proof\_sections}: M3precision finding; M6Fricke\par
\medskip\noindent\textbf{primary\_sources entry 16.}\par
\noindent\textbf{id}: Li1975-p296\par
\noindent\textbf{sha256}: \nolinkurl{11f38136cd27108d507c7662d1ddb46d8c572874b314b69aa1ad8be16395aad3}\par
\noindent\textbf{bytes}: 758672\par
\noindent\textbf{locator}: Math.Ann.212(1975),printed296\par
\noindent\textbf{url}: \url{https://images.p.sub.uni-goettingen.de/iiif/image/gdz:PPN235181684_0212:00000304/full/1800,/0/default.jpg}\par
\noindent\textbf{actual\_inspection}: p296,p298 visually read; p294,p295 acquired, prior bridge read; no false claim of visual reading p294/p295\par
\noindent\textbf{consumer\_proof\_sections}: M3precision finding; M6Fricke\par
\medskip\noindent\textbf{primary\_sources entry 17.}\par
\noindent\textbf{id}: Li1975-p298\par
\noindent\textbf{sha256}: \nolinkurl{609bb6f3a234820789a6a07ed8219ac002e723e9fa46103e27a60d3b1dcdc386}\par
\noindent\textbf{bytes}: 868473\par
\noindent\textbf{locator}: Math.Ann.212(1975),printed298\par
\noindent\textbf{url}: \url{https://images.p.sub.uni-goettingen.de/iiif/image/gdz:PPN235181684_0212:00000306/full/1800,/0/default.jpg}\par
\noindent\textbf{actual\_inspection}: p296,p298 visually read; p294,p295 acquired, prior bridge read; no false claim of visual reading p294/p295\par
\noindent\textbf{consumer\_proof\_sections}: M3precision finding; M6Fricke\par
\medskip\noindent\textbf{primary\_sources entry 18.}\par
\noindent\textbf{id}: PanAuthorV2\par
\noindent\textbf{sha256}: \nolinkurl{590171129d28a65749cf92b1f6099dbacaa90cc6698828f0a4dfd3da943f4824}\par
\noindent\textbf{bytes}: 1552018\par
\noindent\textbf{locator}: Theorem1.0.4p5; Conjecture1.0.1p2\par
\noindent\textbf{url}: \url{https://arxiv.org/pdf/1901.07166v2}\par
\noindent\textbf{authoritative\_version}: historical author version, NOT final published body\par
\noindent\textbf{actual\_inspection}: actual theorem/definitions read\par
\noindent\textbf{use}: proof ancestry/history only; immediate live modularity input is Calegari2023\par
\medskip\noindent\textbf{primary\_sources entry 19.}\par
\noindent\textbf{id}: SilverbergZarhin\allowbreak{}1995\par
\noindent\textbf{sha256}: \nolinkurl{fc62ba1c9d1315a0c441a3dd7de46a254f379289e630c407a3336c0230d91052}\par
\noindent\textbf{bytes}: 1826653\par
\noindent\textbf{locator}: Theorems3.4/3.5,printed405--406\par
\noindent\textbf{url}: \url{https://www.numdam.org/article/AIF_1995__45_2_403_0.pdf}\par
\noindent\textbf{publication}: AnnalesInstitutF\allowbreak{}ourier45(1995),403--420\par
\noindent\textbf{actual\_inspection}: actualfinalTheor\allowbreak{}ems3.4/3.5andSection3text\allowbreak{}read;hashmatchesretai\allowbreak{}nedpublicsource\par
\noindent\textbf{consumer}: PROVIDER\_SCOPE\_CARRYP4;PROOF\_FINALM5\par
\medskip\noindent\textbf{primary\_sources entry 20.}\par
\noindent\textbf{id}: Serre1967-retained-ramification\par
\noindent\textbf{publication}: Cassels--FrohlichAlgebrai\allowbreak{}cNumberTheory(1967),ChapterVI,Localclassfieldt\allowbreak{}heory,128--161\par
\noindent\textbf{locator}: \S{}\S{}4.1--4.3pp155--159\par
\noindent\textbf{retained\_extract\_sha256}: \nolinkurl{e0367b25a022eb27ec220e3c15b474a94fb96d70a6bb89b77f4a02c2231950fe}\par
\noindent\textbf{retrieval\_url}: unknown in retainedsource record; no newlyinventedURL\allowbreak{}\par
\noindent\textbf{actual\_inspection}: retainedexactsou\allowbreak{}rceapplicationan\allowbreak{}dgeneraltheoremp\allowbreak{}rofilecheckedaga\allowbreak{}instcurrentobjec\allowbreak{}ts;notafreshprimary\allowbreak{}extractread\par
\paragraph{carry evidence}
\noindent $\bullet$ CONDUCTOR\_RAMIFICATION\_PROFILES.json\par
\noindent $\bullet$ PUBLIC\_FULL\_COVER\_DECISIONS.json\par
\noindent $\bullet$ PROVIDER\_SCOPE\_CARRY.mdP4\par
\paragraph{internal sources}
\medskip\noindent\textbf{internal\_sources entry 1.}\par
\noindent\textbf{id}: private-height\par
\noindent\textbf{sha256}: \nolinkurl{25960e40419c7abea1f90ab4b35108169a5ca6ad04dccf3a3643e8d0e9bb7d7f}\par
\noindent\textbf{bytes}: 23844\par
\noindent\textbf{locator}: H0--H4; constructionH0/H1; applicationsH0.2/H1/H2; consumersH2--H4\par
\medskip\noindent\textbf{internal\_sources entry 2.}\par
\noindent\textbf{id}: private-v14-source\par
\noindent\textbf{sha256}: \nolinkurl{b6f35c38c9e7370e8ee1b73948f9caa805ed194ca8551a9a61ab9f543ef06d5c}\par
\noindent\textbf{bytes}: 1164623\par
\noindent\textbf{locator}: 481--660;8097--8347 actual originalr5/r3/CL1/CL2\par
\medskip\noindent\textbf{internal\_sources entry 3.}\par
\noindent\textbf{id}: parent9\par
\noindent\textbf{sha256}: \nolinkurl{73ca3d625ba42160cd2a58b4479a9053ee487d95f99161bf62b28d71c0fe54da}\par
\noindent\textbf{bytes}: 18705\par
\noindent\textbf{locator}: \S{}\S{}1--3\par
\medskip\noindent\textbf{internal\_sources entry 4.}\par
\noindent\textbf{id}: cutoff\par
\noindent\textbf{sha256}: \nolinkurl{1a6ae253afd4e92fd2e06dcd3d1bf733464c0320d0861476fa88402b7440fd6f}\par
\noindent\textbf{bytes}: 16422\par
\noindent\textbf{locator}: \S{}1\par
\medskip\noindent\textbf{internal\_sources entry 5.}\par
\noindent\textbf{id}: public-v15\par
\noindent\textbf{sha256}: \nolinkurl{44ba3a8ec8521c29b81b26cfb9981cd6d5f3c76d54e7782438614842edb52c74}\par
\noindent\textbf{bytes}: 1924438\par
\noindent\textbf{locator}: lemma505--705; baseline1334--1505; b-family8252--8495\par
\medskip\noindent\textbf{internal\_sources entry 6.}\par
\noindent\textbf{id}: source-bridge\par
\noindent\textbf{sha256}: \nolinkurl{2c95d02a13ffa972eb6cb8d956b9063e970e33fd960e282c6d53ae8bb655c43c}\par
\noindent\textbf{bytes}: 12939\par
\noindent\textbf{locator}: Faltings/Poincare/Li actual source-application proofs\par
\medskip\noindent\textbf{internal\_sources entry 7.}\par
\noindent\textbf{id}: direct-p5-PROOF.md\par
\noindent\textbf{sha256}: \nolinkurl{995d60ad7b793c17096bbb55476f52104af45a84d634713e3c490fe9bbd6ffda}\par
\noindent\textbf{bytes}: 9670\par
\noindent\textbf{locator}: whole actual argument\par
\medskip\noindent\textbf{internal\_sources entry 8.}\par
\noindent\textbf{id}: direct-p5-FAMILY-CONNECTION.md\par
\noindent\textbf{sha256}: \nolinkurl{f8169651e1835464ef65927f92d0eb24f8af0ada1ab8b5da8b37cf4603ad34f9}\par
\noindent\textbf{bytes}: 1013\par
\noindent\textbf{locator}: whole actual argument\par
\medskip\noindent\textbf{internal\_sources entry 9.}\par
\noindent\textbf{id}: direct-p5-FIVE-TO-THREE-CONNECTION.md\par
\noindent\textbf{sha256}: \nolinkurl{56d613891398cb485f88aedd69c53f3afdb34a8cb167b89d4527e2538920fb35}\par
\noindent\textbf{bytes}: 6341\par
\noindent\textbf{locator}: whole actual argument\par
\medskip\noindent\textbf{internal\_sources entry 10.}\par
\noindent\textbf{id}: direct-p5-CLASSICAL-LOCAL-CONNECTION.md\par
\noindent\textbf{sha256}: \nolinkurl{d990c04c59fbf78ea085e018a6da82f269bb7735379b82255f2ffd4fac110c5e}\par
\noindent\textbf{bytes}: 8237\par
\noindent\textbf{locator}: whole actual argument\par
\medskip\noindent\textbf{internal\_sources entry 11.}\par
\noindent\textbf{id}: direct-p5-SOURCES.md\par
\noindent\textbf{sha256}: \nolinkurl{13cc8eed66c1e7714431c0cbef200a3a3c2e9853d5dcf2eab4aef7bc278ad5da}\par
\noindent\textbf{bytes}: 2932\par
\noindent\textbf{locator}: whole actual argument\par
\medskip\noindent\textbf{internal\_sources entry 12.}\par
\noindent\textbf{id}: direct-p5-SOURCE-STATUS.md\par
\noindent\textbf{sha256}: \nolinkurl{0658b7a9be745216f57d0ee71820c2e91f5c790132902bdc879d2dc44b833f8d}\par
\noindent\textbf{bytes}: 2180\par
\noindent\textbf{locator}: whole actual argument\par
\medskip\noindent\textbf{internal\_sources entry 13.}\par
\noindent\textbf{id}: D210-fullprovider\par
\noindent\textbf{sha256}: \nolinkurl{054ec4a656b82e1d1d17c7e9d3575afe722c42594f933256e7ea225f6a3f5e3b}\par
\noindent\textbf{bytes}: 19349\par
\noindent\textbf{locator}: wholeproof \S{}\S{}1--9,16sourceapplicat\allowbreak{}ions\par
\medskip\noindent\textbf{internal\_sources entry 14.}\par
\noindent\textbf{id}: D210-geometricisogeny\allowbreak{}\par
\noindent\textbf{sha256}: \nolinkurl{f889eab0b9f7880a647ebb34f10ac44a73aaf748f2041cb03dea756079d6ab9e}\par
\noindent\textbf{bytes}: 5819\par
\noindent\textbf{locator}: wholeactualargum\allowbreak{}ent\par
\medskip\noindent\textbf{internal\_sources entry 15.}\par
\noindent\textbf{id}: Yuan-finalbridge\par
\noindent\textbf{sha256}: \nolinkurl{cad3c3e1a78dc2f300a6150181184067ae5636a7339292c7fe74ef8298e89f69}\par
\noindent\textbf{bytes}: 4154\par
\noindent\textbf{locator}: wholeactualmaxim\allowbreak{}alstackandisomet\allowbreak{}rycomparison\par
\medskip\noindent\textbf{internal\_sources entry 16.}\par
\noindent\textbf{id}: H19-sourceprofiles\par
\noindent\textbf{sha256}: \nolinkurl{e9d209cbbecf3c0048cb7d5fb17ff4b428de23c8263659b84b2be8097816e9ec}\par
\noindent\textbf{bytes}: 4940417\par
\noindent\textbf{locator}: 16FULL-COVER+4conductor+H19RAMIFICATION+CM/D210profiles\par
\medskip\noindent\textbf{internal\_sources entry 17.}\par
\noindent\textbf{id}: public-final006coverage\allowbreak{}\par
\noindent\textbf{sha256}: \nolinkurl{9b7dc976e14eb57b940e6f478e921726f3158403b890589f3269ee6eec3f743c}\par
\noindent\textbf{bytes}: 1428503\par
\noindent\textbf{locator}: exactextractedsc\allowbreak{}opes\par
\paragraph{access findings}
\medskip\noindent\textbf{access\_findings entry 1.}\par
\noindent\textbf{source}: PanfinalJAMS2022\allowbreak{}\par
\paragraph{attempts}
\noindent $\bullet$ \url{https://www.ams.org/jams/2022-35-04/S0894-0347-2021-00991-1/S0894-0347-2021-00991-1.pdf}\par
\noindent $\bullet$ \url{https://www.ams.org/journals/jams/2022-35-04/S0894-0347-2021-00991-1/S0894-0347-2021-00991-1.pdf}\par
\noindent\textbf{result}: 403; final body not read; Calegari final statement read and used directly\par
\medskip\noindent\textbf{access\_findings entry 2.}\par
\noindent\textbf{source}: Washington1997Th\allowbreak{}eorem14.1\par
\noindent\textbf{result}: not reacquired; proofM3instead derives finiteDirichletc\allowbreak{}haracter from published MilneADTglobalre\allowbreak{}ciprocity and Qidele decomposition\par
\paragraph{uninspected live antecedents}
\noindent $\bullet$ The sixteen exactpublicfullD\allowbreak{}210sourceapplica\allowbreak{}tions are carried at their accepted scope, not freshly reaudited from every finalPDF/HTML; the wholeinternalpro\allowbreak{}of and currentprivateco\allowbreak{}nnection are directlyread. Independentcurre\allowbreak{}ntreviewmustacce\allowbreak{}ptthisexactcarry\allowbreak{}.\par
\noindent $\bullet$ No primaryfinalPanJ\allowbreak{}AMSbodycompariso\allowbreak{}n; immediatefinalpu\allowbreak{}blishedCalegarit\allowbreak{}heoremisdirectly\allowbreak{}checked.\par
\noindent $\bullet$ SeparateheightH3\allowbreak{}Remond/final214assembly\allowbreak{}remainrootlanes.\par
\noindent\textbf{supersedes\_scope\_statements\_only}: SOURCE\_USES.json remains frozen initialauthorhis\allowbreak{}tory; this finalinventory adds exactH01andcondu\allowbreak{}ctorcarry, without changing initialfiles.\par
\paragraph{final bindings}
\medskip\noindent\textbf{final\_bindings entry 1.}\par
\noindent\textbf{sha256}: \nolinkurl{814676050b389793d196346f45e2797f39e2c5be0f3fee8457bb709b2607606e}\par
\noindent\textbf{bytes}: 25777\par
\medskip\noindent\textbf{final\_bindings entry 2.}\par
\noindent\textbf{sha256}: \nolinkurl{b3cb8201de84f733b0518ad68a9c0d0bc1e78534ffd39aca5d8c70f1bae69e76}\par
\noindent\textbf{bytes}: 13337\par
\medskip\noindent\textbf{final\_bindings entry 3.}\par
\noindent\textbf{sha256}: \nolinkurl{368774d2b96a606df4dd5e81a3c4fecaf22f09f33f27e49e0a80e37b02157154}\par
\noindent\textbf{bytes}: 16324\par
\medskip\noindent\textbf{final\_bindings entry 4.}\par
\noindent\textbf{sha256}: \nolinkurl{48fd4208cabd031cabf9d71f4a9b9dc5e1f7176bed6489a9d7224f5d7dcfb7e5}\par
\noindent\textbf{bytes}: 227373\par
\medskip\noindent\textbf{final\_bindings entry 5.}\par
\noindent\textbf{sha256}: \nolinkurl{00643457f3c10e53fd31ab7b7d38755906c8dcc433939480f130d4339b2a3127}\par
\noindent\textbf{bytes}: 21154\par
\medskip\noindent\textbf{final\_bindings entry 6.}\par
\noindent\textbf{sha256}: \nolinkurl{64896dfbb8661e04c044a68aa8679ee5dc20c0da36cc9a57640a885b7ca46614}\par
\noindent\textbf{bytes}: 1501\par
\medskip\noindent\textbf{final\_bindings entry 7.}\par
\noindent\textbf{sha256}: \nolinkurl{0e7280479d5ace418b4383c814361812863cb097f424916b801c17ddbc8eb51e}\par
\noindent\textbf{bytes}: 1945\par
\medskip\noindent\textbf{final\_bindings entry 8.}\par
\noindent\textbf{sha256}: \nolinkurl{b97585cd239a9c7b3138dabb010eb959fc16f708bb1af4de2e35b182f5cc372a}\par
\noindent\textbf{bytes}: 3050\par
\paragraph{height evaluation}
\medskip\noindent\textbf{height\_evaluation entry 1.}\par
\noindent\textbf{sha256}: \nolinkurl{cc7f1947308a6c8db0d6d0d406bfecc4ef2363d710f694b1cc5fe9bb0ba8219f}\par
\noindent\textbf{bytes}: 16117\par
\noindent\textbf{id}: S-00R4\par
\noindent\textbf{title}: The Stacks Project\par
\noindent\textbf{version}: Continuously maintained online reference, exact raw HTML snapshot accessed 2026-11-10; no global release commit asserted\par
\noindent\textbf{publication\_status}: Published online mathematical reference\par
\noindent\textbf{url}: \url{https://stacks.math.columbia.edu/tag/00R4}\par
\noindent\textbf{access\_date}: 2026-11-10\par
\noindent\textbf{locator}: Lemma 10.128.1, statement and proof: all three miracle-flatness hypotheses\par
\noindent\textbf{access\_scope}: Statement, listed definitions and relevant proof read\par
\noindent\textbf{license}: HTML footer links GNU Free Documentation License\par
\medskip\noindent\textbf{height\_evaluation entry 2.}\par
\noindent\textbf{sha256}: \nolinkurl{05a6f4daf4e0ce0e22731a953dc355b411931313c9e549bbec1f7b9d0a1105c9}\par
\noindent\textbf{bytes}: 18424\par
\noindent\textbf{id}: S-031S\par
\noindent\textbf{title}: The Stacks Project\par
\noindent\textbf{version}: Continuously maintained online reference, exact raw HTML snapshot accessed 2026-11-10; no global release commit asserted\par
\noindent\textbf{publication\_status}: Published online mathematical reference\par
\noindent\textbf{url}: \url{https://stacks.math.columbia.edu/tag/031S}\par
\noindent\textbf{access\_date}: 2026-11-10\par
\noindent\textbf{locator}: Lemma 10.157.4, normal iff R1 and S2, including proof of normal \ensuremath{\Longrightarrow } depth at least min(2,dimension)\par
\noindent\textbf{access\_scope}: Statement, listed definitions and relevant proof read\par
\noindent\textbf{license}: HTML footer links GNU Free Documentation License\par
\medskip\noindent\textbf{height\_evaluation entry 3.}\par
\noindent\textbf{sha256}: \nolinkurl{f043e46feb56ec178ef2af81acf345fabc93d42e71a9f1e544e962edbe329af9}\par
\noindent\textbf{bytes}: 17764\par
\noindent\textbf{id}: S-00H8\par
\noindent\textbf{title}: The Stacks Project\par
\noindent\textbf{version}: Continuously maintained online reference, exact raw HTML snapshot accessed 2026-11-10; no global release commit asserted\par
\noindent\textbf{publication\_status}: Published online mathematical reference\par
\noindent\textbf{url}: \url{https://stacks.math.columbia.edu/tag/00H8}\par
\noindent\textbf{access\_date}: 2026-11-10\par
\noindent\textbf{locator}: Proposition 10.38.7, statement and proof: going down for integral extension of domains over normal domain\par
\noindent\textbf{access\_scope}: Statement, listed definitions and relevant proof read\par
\noindent\textbf{license}: HTML footer links GNU Free Documentation License\par
\medskip\noindent\textbf{height\_evaluation entry 4.}\par
\noindent\textbf{sha256}: \nolinkurl{1921d9e780f76f5ea003c57bb58b855748d4f9a333c4dad7d877c649ca1605f2}\par
\noindent\textbf{bytes}: 14961\par
\noindent\textbf{id}: S-0AHI\par
\noindent\textbf{title}: The Stacks Project\par
\noindent\textbf{version}: Continuously maintained online reference, exact raw HTML snapshot accessed 2026-11-10; no global release commit asserted\par
\noindent\textbf{publication\_status}: Published online mathematical reference\par
\noindent\textbf{url}: \url{https://stacks.math.columbia.edu/tag/0AHI}\par
\noindent\textbf{access\_date}: 2026-11-10\par
\noindent\textbf{locator}: Lemma 54.4.2, statement and proof: domination of proper modification by point blowups over finite regular dimension-two centres\par
\noindent\textbf{access\_scope}: Statement, listed definitions and relevant proof read\par
\noindent\textbf{license}: HTML footer links GNU Free Documentation License\par
\medskip\noindent\textbf{height\_evaluation entry 5.}\par
\noindent\textbf{sha256}: \nolinkurl{1715af2753bb039ae57ebd0d69d3b4efa1050bdd538629f8b17f4dae0751d4c1}\par
\noindent\textbf{bytes}: 17338\par
\noindent\textbf{id}: S-0BFP\par
\noindent\textbf{title}: The Stacks Project\par
\noindent\textbf{version}: Continuously maintained online reference, exact raw HTML snapshot accessed 2026-11-10; no global release commit asserted\par
\noindent\textbf{publication\_status}: Published online mathematical reference\par
\noindent\textbf{url}: \url{https://stacks.math.columbia.edu/tag/0BFP}\par
\noindent\textbf{access\_date}: 2026-11-10\par
\noindent\textbf{locator}: Lemma 33.17.3, statement parts (1)--(2), proof: codimension-one isomorphism locus for proper birational reduced source\par
\noindent\textbf{access\_scope}: Statement, listed definitions and relevant proof read\par
\noindent\textbf{license}: HTML footer links GNU Free Documentation License\par
\medskip\noindent\textbf{height\_evaluation entry 6.}\par
\noindent\textbf{sha256}: \nolinkurl{56877b9b7bb0e689cf3c5004fbe9eb4abd828d85652ff0306b99af928cdc2ad0}\par
\noindent\textbf{bytes}: 19975\par
\noindent\textbf{id}: S-0AWQ\par
\noindent\textbf{title}: The Stacks Project\par
\noindent\textbf{version}: Continuously maintained online reference, exact raw HTML snapshot accessed 2026-11-10; no global release commit asserted\par
\noindent\textbf{publication\_status}: Published online mathematical reference\par
\noindent\textbf{url}: \url{https://stacks.math.columbia.edu/tag/0AWQ}\par
\noindent\textbf{access\_date}: 2026-11-10\par
\noindent\textbf{locator}: Section 48.23, Lemma 48.23.1, statement and proof: canonical sheaf of connected CM scheme is a coherent CM module\par
\noindent\textbf{access\_scope}: Statement, listed definitions and relevant proof read\par
\noindent\textbf{license}: HTML footer links GNU Free Documentation License\par
\medskip\noindent\textbf{height\_evaluation entry 7.}\par
\noindent\textbf{sha256}: \nolinkurl{2d6f74cf8e80e4cab6b19656a6a6a50998fbb26805ead641658436109112ce65}\par
\noindent\textbf{bytes}: 32467\par
\noindent\textbf{id}: S-0AVT\par
\noindent\textbf{title}: The Stacks Project\par
\noindent\textbf{version}: Continuously maintained online reference, exact raw HTML snapshot accessed 2026-11-10; no global release commit asserted\par
\noindent\textbf{publication\_status}: Published online mathematical reference\par
\noindent\textbf{url}: \url{https://stacks.math.columbia.edu/tag/0AVT}\par
\noindent\textbf{access\_date}: 2026-11-10\par
\noindent\textbf{locator}: Section 31.13, Definition 31.13.1 and Lemmas 31.13.12--31.13.14: reflexivity, extension and normal S2 criterion\par
\noindent\textbf{access\_scope}: Statement, listed definitions and relevant proof read\par
\noindent\textbf{license}: HTML footer links GNU Free Documentation License\par
\medskip\noindent\textbf{height\_evaluation entry 8.}\par
\noindent\textbf{sha256}: \nolinkurl{d7e45af8f20d8af5017d2db20c2907ae5092031fc7650b7c561a4e398dfc13ed}\par
\noindent\textbf{bytes}: 18328\par
\noindent\textbf{id}: S-0BUL\par
\noindent\textbf{title}: The Stacks Project\par
\noindent\textbf{version}: Continuously maintained online reference, exact raw HTML snapshot accessed 2026-11-10; no global release commit asserted\par
\noindent\textbf{publication\_status}: Published online mathematical reference\par
\noindent\textbf{url}: \url{https://stacks.math.columbia.edu/tag/0BUL}\par
\noindent\textbf{access\_date}: 2026-11-10\par
\noindent\textbf{locator}: Lemma 49.15.1, statement and proof: finite Hom duality and compatibility with restrictions; 2024 corrected target of f shriek\par
\noindent\textbf{access\_scope}: Statement, listed definitions and relevant proof read\par
\noindent\textbf{license}: HTML footer links GNU Free Documentation License\par
\medskip\noindent\textbf{height\_evaluation entry 9.}\par
\noindent\textbf{sha256}: \nolinkurl{d36a3ecb38faa3db89bf72578c1b00c449ee985808e122c664366857a2c3a072}\par
\noindent\textbf{bytes}: 14286\par
\noindent\textbf{id}: S-0C0I\par
\noindent\textbf{title}: The Stacks Project\par
\noindent\textbf{version}: Continuously maintained online reference, exact raw HTML snapshot accessed 2026-11-10; no global release commit asserted\par
\noindent\textbf{publication\_status}: Published online mathematical reference\par
\noindent\textbf{url}: \url{https://stacks.math.columbia.edu/tag/0C0I}\par
\noindent\textbf{access\_date}: 2026-11-10\par
\noindent\textbf{locator}: Lemma 49.2.12, statement and proof: nonunique algebraic comparison; qualification recorded, not used as a canonical trace normalization\par
\noindent\textbf{access\_scope}: Statement, listed definitions and relevant proof read\par
\noindent\textbf{license}: HTML footer links GNU Free Documentation License\par
\medskip\noindent\textbf{height\_evaluation entry 10.}\par
\noindent\textbf{sha256}: \nolinkurl{8b93bd28adc3f32bb19dbf647502c3a140ee586997bc0c323cf25856a3447204}\par
\noindent\textbf{bytes}: 14246\par
\noindent\textbf{id}: S-0BW2\par
\noindent\textbf{title}: The Stacks Project\par
\noindent\textbf{version}: Continuously maintained online reference, exact raw HTML snapshot accessed 2026-11-10; no global release commit asserted\par
\noindent\textbf{publication\_status}: Published online mathematical reference\par
\noindent\textbf{url}: \url{https://stacks.math.columbia.edu/tag/0BW2}\par
\noindent\textbf{access\_date}: 2026-11-10\par
\noindent\textbf{locator}: Lemma 49.8.2, statement and proof: complementary module\ensuremath{\to }Hom via x\ensuremath{\mapsto }[b\ensuremath{\mapsto }Tr(bx)]\par
\noindent\textbf{access\_scope}: Statement, listed definitions and relevant proof read\par
\noindent\textbf{license}: HTML footer links GNU Free Documentation License\par
\medskip\noindent\textbf{height\_evaluation entry 11.}\par
\noindent\textbf{sha256}: \nolinkurl{69a7fc254581a7fbd20e41f7a56086d6184f6843be666e2dfe93658893bef4d7}\par
\noindent\textbf{bytes}: 14028\par
\noindent\textbf{id}: S-0BW3\par
\noindent\textbf{title}: The Stacks Project\par
\noindent\textbf{version}: Continuously maintained online reference, exact raw HTML snapshot accessed 2026-11-10; no global release commit asserted\par
\noindent\textbf{publication\_status}: Published online mathematical reference\par
\noindent\textbf{url}: \url{https://stacks.math.columbia.edu/tag/0BW3}\par
\noindent\textbf{access\_date}: 2026-11-10\par
\noindent\textbf{locator}: Lemma 49.8.3, statement (1) and proof: 1 maps to the trace element for flat finite generically etale extension\par
\noindent\textbf{access\_scope}: Statement, listed definitions and relevant proof read\par
\noindent\textbf{license}: HTML footer links GNU Free Documentation License\par
\medskip\noindent\textbf{height\_evaluation entry 12.}\par
\noindent\textbf{sha256}: \nolinkurl{cab9835bfd28b4c45507644b09cc1f56d78572a6a9bedb1fba0385fdb9b24536}\par
\noindent\textbf{bytes}: 13377\par
\noindent\textbf{id}: S-0F42\par
\noindent\textbf{title}: The Stacks Project\par
\noindent\textbf{version}: Continuously maintained online reference, exact raw HTML snapshot accessed 2026-11-10; no global release commit asserted\par
\noindent\textbf{publication\_status}: Published online mathematical reference\par
\noindent\textbf{url}: \url{https://stacks.math.columbia.edu/tag/0F42}\par
\noindent\textbf{access\_date}: 2026-11-10\par
\noindent\textbf{locator}: Situation 48.16.1: Noetherian base and finite-type separated scheme category\par
\noindent\textbf{access\_scope}: Statement, listed definitions and relevant proof read\par
\noindent\textbf{license}: HTML footer links GNU Free Documentation License\par
\medskip\noindent\textbf{height\_evaluation entry 13.}\par
\noindent\textbf{sha256}: \nolinkurl{793c7be1cd2a4f36fce5e49e10e1f5f04642924f7b2f6a6fae4117f4e718476e}\par
\noindent\textbf{bytes}: 21847\par
\noindent\textbf{id}: S-0ATX\par
\noindent\textbf{title}: The Stacks Project\par
\noindent\textbf{version}: Continuously maintained online reference, exact raw HTML snapshot accessed 2026-11-10; no global release commit asserted\par
\noindent\textbf{publication\_status}: Published online mathematical reference\par
\noindent\textbf{url}: \url{https://stacks.math.columbia.edu/tag/0ATX}\par
\noindent\textbf{access\_date}: 2026-11-10\par
\noindent\textbf{locator}: Lemma 48.16.3, statement and proof: canonical composition comparison for upper shriek\par
\noindent\textbf{access\_scope}: Statement, listed definitions and relevant proof read\par
\noindent\textbf{license}: HTML footer links GNU Free Documentation License\par
\medskip\noindent\textbf{height\_evaluation entry 14.}\par
\noindent\textbf{sha256}: \nolinkurl{9cc1165e76b572600e5ed912353529cea4eb290d0e7fcc86cb1a302d27e6eafb}\par
\noindent\textbf{bytes}: 14999\par
\noindent\textbf{id}: S-0AU0\par
\noindent\textbf{title}: The Stacks Project\par
\noindent\textbf{version}: Continuously maintained online reference, exact raw HTML snapshot accessed 2026-11-10; no global release commit asserted\par
\noindent\textbf{publication\_status}: Published online mathematical reference\par
\noindent\textbf{url}: \url{https://stacks.math.columbia.edu/tag/0AU0}\par
\noindent\textbf{access\_date}: 2026-11-10\par
\noindent\textbf{locator}: Lemma 48.17.1, statement and proof: restriction to opens\par
\noindent\textbf{access\_scope}: Statement, listed definitions and relevant proof read\par
\noindent\textbf{license}: HTML footer links GNU Free Documentation License\par
\medskip\noindent\textbf{height\_evaluation entry 15.}\par
\noindent\textbf{sha256}: \nolinkurl{9592cdfc190919da3ff6408cc8edb6b05fbe82fc546bafb2a5450997f773001c}\par
\noindent\textbf{bytes}: 18311\par
\noindent\textbf{id}: S-0B6U\par
\noindent\textbf{title}: The Stacks Project\par
\noindent\textbf{version}: Continuously maintained online reference, exact raw HTML snapshot accessed 2026-11-10; no global release commit asserted\par
\noindent\textbf{publication\_status}: Published online mathematical reference\par
\noindent\textbf{url}: \url{https://stacks.math.columbia.edu/tag/0B6U}\par
\noindent\textbf{access\_date}: 2026-11-10\par
\noindent\textbf{locator}: Lemma 48.17.9, statement (5) and proof: tensor comparison for perfect, in particular flat, morphism\par
\noindent\textbf{access\_scope}: Statement, listed definitions and relevant proof read\par
\noindent\textbf{license}: HTML footer links GNU Free Documentation License\par
\medskip\noindent\textbf{height\_evaluation entry 16.}\par
\noindent\textbf{sha256}: \nolinkurl{3b33f531f4fbed0748f632bc79ec2687c364cb7b0ea495916356e7f09e00b68a}\par
\noindent\textbf{bytes}: 15946\par
\noindent\textbf{id}: S-0E9U\par
\noindent\textbf{title}: The Stacks Project\par
\noindent\textbf{version}: Continuously maintained online reference, exact raw HTML snapshot accessed 2026-11-10; no global release commit asserted\par
\noindent\textbf{publication\_status}: Published online mathematical reference\par
\noindent\textbf{url}: \url{https://stacks.math.columbia.edu/tag/0E9U}\par
\noindent\textbf{access\_date}: 2026-11-10\par
\noindent\textbf{locator}: Lemma 48.18.1, statement and proof: flat base-change comparison\par
\noindent\textbf{access\_scope}: Statement, listed definitions and relevant proof read\par
\noindent\textbf{license}: HTML footer links GNU Free Documentation License\par
\medskip\noindent\textbf{height\_evaluation entry 17.}\par
\noindent\textbf{sha256}: \nolinkurl{73a846e8b92c1be1901fc8ceda12fcfc31a37f0f32cd0dad899be9ca66634df6}\par
\noindent\textbf{bytes}: 18222\par
\noindent\textbf{id}: S-0E9Z\par
\noindent\textbf{title}: The Stacks Project\par
\noindent\textbf{version}: Continuously maintained online reference, exact raw HTML snapshot accessed 2026-11-10; no global release commit asserted\par
\noindent\textbf{publication\_status}: Published online mathematical reference\par
\noindent\textbf{url}: \url{https://stacks.math.columbia.edu/tag/0E9Z}\par
\noindent\textbf{access\_date}: 2026-11-10\par
\noindent\textbf{locator}: Lemma 48.29.2, statement and conormal/adjunction proof: smooth differential identification; the conormal construction commutes with flat base change\par
\noindent\textbf{access\_scope}: Statement, listed definitions and relevant proof read\par
\noindent\textbf{license}: HTML footer links GNU Free Documentation License\par
\medskip\noindent\textbf{height\_evaluation entry 18.}\par
\noindent\textbf{sha256}: \nolinkurl{f2578b1976e269273994b8cb04dd07efd34006d887dcc7d077015ffa42338466}\par
\noindent\textbf{bytes}: 14213\par
\noindent\textbf{id}: S-04HN\par
\noindent\textbf{title}: The Stacks Project\par
\noindent\textbf{version}: Continuously maintained online reference, exact raw HTML snapshot accessed 2026-11-10; no global release commit asserted\par
\noindent\textbf{publication\_status}: Published online mathematical reference\par
\noindent\textbf{url}: \url{https://stacks.math.columbia.edu/tag/04HN}\par
\noindent\textbf{access\_date}: 2026-11-10\par
\noindent\textbf{locator}: Lemma 41.18.3, statement and proof: finite etale morphism splits into identity maps on etale neighbourhood\par
\noindent\textbf{access\_scope}: Statement, listed definitions and relevant proof read\par
\noindent\textbf{license}: HTML footer links GNU Free Documentation License\par
\medskip\noindent\textbf{height\_evaluation entry 19.}\par
\noindent\textbf{sha256}: \nolinkurl{ec890fb41c58dee038ee2e06b530954f3aab85035238cd54a4c0ef081fcb8ea0}\par
\noindent\textbf{bytes}: 1222544\par
\noindent\textbf{id}: S-JAV\par
\noindent\textbf{title}: Polynomial bounds for Arakelov invariants of Belyi curves\par
\noindent\textbf{author}: Ariyan Javanpeykar, with appendix by Peter Bruin\par
\noindent\textbf{version}: Final published Algebra \& Number Theory 8 (2014), 89--140\par
\noindent\textbf{publication\_status}: Final published paper\par
\noindent\textbf{doi}: 10.2140/ant.2014.8.89\par
\noindent\textbf{url}: \url{https://msp.org/ant/2014/8-1/ant-v8-n1-p04-s.pdf}\par
\noindent\textbf{access\_date}: 2026-11-10\par
\noindent\textbf{locator}: Arithmetic surface definition pp.95--96; Lemma 2.3.1 and its proof p.99\par
\noindent\textbf{access\_scope}: Definition and actual blowup proof read; point-height conclusion not substituted for formula\par
\noindent\textbf{distribution}: Local research evidence; full PDF excluded from default integration copy list\par
\paragraph{height different}
\medskip\noindent\textbf{height\_different entry 1.}\par
\noindent\textbf{id}: javanpeykar-published.pdf\par
\noindent\textbf{sha256}: \nolinkurl{ec890fb41c58dee038ee2e06b530954f3aab85035238cd54a4c0ef081fcb8ea0}\par
\noindent\textbf{url}: \url{https://msp.org/ant/2014/8-1/ant-v8-n1-p04-s.pdf}\par
\noindent\textbf{access\_date}: 2026-11-10\par
\noindent\textbf{publication\_status}: final publisher PDF\par
\noindent\textbf{version}: ANT8(2014),89--140, DOI10.2140/ant.2014.8.89\par
\noindent\textbf{publication\_scope}: Full third-party PDF retained only as local reviewer source; publication bundle may retain URL/hash and scoped excerpts instead.\par
\medskip\noindent\textbf{height\_different entry 2.}\par
\noindent\textbf{id}: stacks-0BFG.html\par
\noindent\textbf{sha256}: \nolinkurl{f2b650fba17a059317eb935e33d33572a84edb970a5a7fc954a2be3d8a6eb212}\par
\noindent\textbf{url}: \url{https://stacks.math.columbia.edu/tag/0BFG}\par
\noindent\textbf{access\_date}: 2026-11-10\par
\noindent\textbf{publication\_status}: Stacks live tag page archived as exact HTML snapshot; no Git commit asserted\par
\noindent\textbf{version}: content served at access date, byte-hash version\par
\noindent\textbf{publication\_scope}: Full third-party PDF retained only as local reviewer source; publication bundle may retain URL/hash and scoped excerpts instead.\par
\medskip\noindent\textbf{height\_different entry 3.}\par
\noindent\textbf{id}: stacks-0BVZ.html\par
\noindent\textbf{sha256}: \nolinkurl{40a236f53d34753c7301bd6c881018f086b149e6b482cf15d1fed7339b11c829}\par
\noindent\textbf{url}: \url{https://stacks.math.columbia.edu/tag/0BVZ}\par
\noindent\textbf{access\_date}: 2026-11-10\par
\noindent\textbf{publication\_status}: Stacks live tag page archived as exact HTML snapshot; no Git commit asserted\par
\noindent\textbf{version}: content served at access date, byte-hash version\par
\noindent\textbf{publication\_scope}: Full third-party PDF retained only as local reviewer source; publication bundle may retain URL/hash and scoped excerpts instead.\par
\medskip\noindent\textbf{height\_different entry 4.}\par
\noindent\textbf{id}: stacks-0BW0.html\par
\noindent\textbf{sha256}: \nolinkurl{0615dff37470db0bc6ed411b7e9b340bc8636ebe2ba73f92155eb41b02af85da}\par
\noindent\textbf{url}: \url{https://stacks.math.columbia.edu/tag/0BW0}\par
\noindent\textbf{access\_date}: 2026-11-10\par
\noindent\textbf{publication\_status}: Stacks live tag page archived as exact HTML snapshot; no Git commit asserted\par
\noindent\textbf{version}: content served at access date, byte-hash version\par
\noindent\textbf{publication\_scope}: Full third-party PDF retained only as local reviewer source; publication bundle may retain URL/hash and scoped excerpts instead.\par
\medskip\noindent\textbf{height\_different entry 5.}\par
\noindent\textbf{id}: stacks-0BW4.html\par
\noindent\textbf{sha256}: \nolinkurl{e00ca39df1423bf9650831740a12ba11073e26b97145501aed047e71e3b5eb29}\par
\noindent\textbf{url}: \url{https://stacks.math.columbia.edu/tag/0BW4}\par
\noindent\textbf{access\_date}: 2026-11-10\par
\noindent\textbf{publication\_status}: Stacks live tag page archived as exact HTML snapshot; no Git commit asserted\par
\noindent\textbf{version}: content served at access date, byte-hash version\par
\noindent\textbf{publication\_scope}: Full third-party PDF retained only as local reviewer source; publication bundle may retain URL/hash and scoped excerpts instead.\par
\medskip\noindent\textbf{height\_different entry 6.}\par
\noindent\textbf{id}: stacks-0BW6.html\par
\noindent\textbf{sha256}: \nolinkurl{1fd9f612c858f0111c257515bf42f527408d4fedf4e4351a01a5ab3abdfabe9f}\par
\noindent\textbf{url}: \url{https://stacks.math.columbia.edu/tag/0BW6}\par
\noindent\textbf{access\_date}: 2026-11-10\par
\noindent\textbf{publication\_status}: Stacks live tag page archived as exact HTML snapshot; no Git commit asserted\par
\noindent\textbf{version}: content served at access date, byte-hash version\par
\noindent\textbf{publication\_scope}: Full third-party PDF retained only as local reviewer source; publication bundle may retain URL/hash and scoped excerpts instead.\par
\medskip\noindent\textbf{height\_different entry 7.}\par
\noindent\textbf{id}: stacks-0BWB.html\par
\noindent\textbf{sha256}: \nolinkurl{3f08ca631d04dcd310e4831f2129c98fc5aad0e16a374bfcd2dffeb7a65338b0}\par
\noindent\textbf{url}: \url{https://stacks.math.columbia.edu/tag/0BWB}\par
\noindent\textbf{access\_date}: 2026-11-10\par
\noindent\textbf{publication\_status}: Stacks live tag page archived as exact HTML snapshot; no Git commit asserted\par
\noindent\textbf{version}: content served at access date, byte-hash version\par
\noindent\textbf{publication\_scope}: Full third-party PDF retained only as local reviewer source; publication bundle may retain URL/hash and scoped excerpts instead.\par
\medskip\noindent\textbf{height\_different entry 8.}\par
\noindent\textbf{id}: stacks-0BWE.html\par
\noindent\textbf{sha256}: \nolinkurl{a6ea73cec0946ca70e7c1bdfc99594f2a8e15ee232abab158d9a92a39cc9fbac}\par
\noindent\textbf{url}: \url{https://stacks.math.columbia.edu/tag/0BWE}\par
\noindent\textbf{access\_date}: 2026-11-10\par
\noindent\textbf{publication\_status}: Stacks live tag page archived as exact HTML snapshot; no Git commit asserted\par
\noindent\textbf{version}: content served at access date, byte-hash version\par
\noindent\textbf{publication\_scope}: Full third-party PDF retained only as local reviewer source; publication bundle may retain URL/hash and scoped excerpts instead.\par
\medskip\noindent\textbf{height\_different entry 9.}\par
\noindent\textbf{id}: stacks-0BWF.html\par
\noindent\textbf{sha256}: \nolinkurl{a1a842d28f6abc94ef2eadd9948f894988d02d1ea125b014f1ca441d702360b5}\par
\noindent\textbf{url}: \url{https://stacks.math.columbia.edu/tag/0BWF}\par
\noindent\textbf{access\_date}: 2026-11-10\par
\noindent\textbf{publication\_status}: Stacks live tag page archived as exact HTML snapshot; no Git commit asserted\par
\noindent\textbf{version}: content served at access date, byte-hash version\par
\noindent\textbf{publication\_scope}: Full third-party PDF retained only as local reviewer source; publication bundle may retain URL/hash and scoped excerpts instead.\par
\medskip\noindent\textbf{height\_different entry 10.}\par
\noindent\textbf{id}: von-kaenel-published.pdf\par
\noindent\textbf{sha256}: \nolinkurl{5749ce3a371c571655b44174e2a0d65c2e53dc179f518f0e52d585ff73da46f9}\par
\noindent\textbf{url}: \url{https://www.cambridge.org/core/services/aop-cambridge-core/content/view/5DA2DC5985591DFBC9A26C20BA22F49B/S2050509421000293a.pdf/the-effective-shafarevich-conjecture-for-abelian-varieties-of-dollartext-gl2dollar-type.pdf}\par
\noindent\textbf{access\_date}: 2026-11-10\par
\noindent\textbf{publication\_status}: final publisher PDF\par
\noindent\textbf{version}: FoM Sigma9(2021)e39, DOI10.1017/fms.2021.29\par
\noindent\textbf{publication\_scope}: Full third-party PDF retained only as local reviewer source; publication bundle may retain URL/hash and scoped excerpts instead.\par
\paragraph{height metric}
\medskip\noindent\textbf{height\_metric entry 1.}\par
\noindent\textbf{id}: J2014\par
\noindent\textbf{title}: Polynomial bounds for Arakelov invariants of Belyi curves\par
\noindent\textbf{author}: Ariyan Javanpeykar, appendix Peter Bruin\par
\noindent\textbf{publication}: Algebra \& Number Theory 8 (2014), 89--140\par
\noindent\textbf{doi}: 10.2140/ant.2014.8.89\par
\noindent\textbf{url}: \url{https://msp.org/ant/2014/8-1/ant-v8-n1-p04-s.pdf}\par
\noindent\textbf{accessed}: 2026-11-10\par
\noindent\textbf{version}: Final publisher PDF\par
\paragraph{evidence}
\noindent\textbf{path}: sources/javanpeykar2014.pdf\par
\noindent\textbf{sha256}: \nolinkurl{ec890fb41c58dee038ee2e06b530954f3aab85035238cd54a4c0ef081fcb8ea0}\par
\noindent\textbf{bytes}: 1222544\par
\paragraph{text}
\noindent\textbf{path}: sources/javanpeykar2014.txt\par
\noindent\textbf{sha256}: \nolinkurl{ae952c57b066d1cd0daaffb88a5407a84bc970a694bb8f1e1e27e2aa33935845}\par
\noindent\textbf{bytes}: 163512\par
\paragraph{locators}
\noindent $\bullet$ \S{}2.1 printed pp93--94: H0 inner product and determinant metric\par
\noindent $\bullet$ \S{}2.3 printed pp97--98: minimal regular semistable model, normalized arithmetic degree and stable Jacobian identification\par
\noindent\textbf{scope}: Direct final published statement and its definitions; Szpiro lemma cited there is not separately re-audited. The published identification is the permitted theorem input.\par
\medskip\noindent\textbf{height\_metric entry 2.}\par
\noindent\textbf{id}: VK2021\par
\noindent\textbf{title}: The effective Shafarevich conjecture for abelian varieties of GL2-type\par
\noindent\textbf{author}: Rafael von Känel\par
\noindent\textbf{publication}: Forum of Mathematics, Sigma 9 (2021), e39, 1--29\par
\noindent\textbf{doi}: 10.1017/fms.2021.29\par
\noindent\textbf{url}: \url{https://www.cambridge.org/core/services/aop-cambridge-core/content/view/5DA2DC5985591DFBC9A26C20BA22F49B/S2050509421000293a.pdf/the-effective-shafarevich-conjecture-for-abelian-varieties-of-dollartext-gl2dollar-type.pdf}\par
\noindent\textbf{accessed}: 2026-11-10\par
\noindent\textbf{version}: Final publisher PDF\par
\paragraph{evidence}
\noindent\textbf{path}: sources/vonkanel2021.pdf\par
\noindent\textbf{sha256}: \nolinkurl{5749ce3a371c571655b44174e2a0d65c2e53dc179f518f0e52d585ff73da46f9}\par
\noindent\textbf{bytes}: 667353\par
\paragraph{text}
\noindent\textbf{path}: sources/vonkanel2021.txt\par
\noindent\textbf{sha256}: \nolinkurl{6ddb8af13dc25fd1c5a2a8f8c822e8ab77ed65fde6a4456ab9609b6c50a671a2}\par
\noindent\textbf{bytes}: 161161\par
\paragraph{locators}
\noindent $\bullet$ \S{}2.1 printed p5: Néron differential-line norm, all-embedding degree, stable extension invariance\par
\noindent $\bullet$ \S{}2.1 printed p5 after(2.2): hF=hB-(g/2)log\ensuremath{\pi }\par
\noindent\textbf{scope}: Final definition/base-change paragraphs read, not a complete source audit of its cited semistable-reduction and Néron-model theorems.\par
\paragraph{uniform existence}
\medskip\noindent\textbf{uniform\_existence entry 1.}\par
\noindent\textbf{author}: Jérôme Von Buhren\par
\noindent\textbf{title}: Borne de hauteur semi-effective pour le problème de Mordell--Lang dans une variété abélienne\par
\noindent\textbf{edition}: JTNB29(2017),307--320, final published\par
\noindent\textbf{url}: \url{https://www.numdam.org/item/10.5802/jtnb.981.pdf}\par
\noindent\textbf{doi}: 10.5802/jtnb.981\par
\noindent\textbf{locator}: printed308 Theorem1.1, definition of C(Gamma,epsilon) and embedding(1.1)\par
\noindent\textbf{distribution}: local read-only source; no third-party PDF copied into this packet\par
\endgroup
\clearpage\section{Public-source resolution and reproducibility}\label{q214:resolution}
The historical path names in the retained author proofs are provenance locators.  The complete mathematical foundation they identify is present in Part~II, and the selected later proofs are present in Part~I.  The public reading aliases below resolve the retained historical locators; none requires private-repository access.  Exact file and passage bindings for the source packet remain in \nolinkurl{research/quartic214/SOURCE_PATH_MAP.json} and the application contracts.  They are not rewritten into current acceptance for this new manuscript.
\begin{description}[style=nextline]
\item[\nolinkurl{height/PROOF.md}] \nolinkurl{research/quartic214/proofs/height-supplier.md}
\item[\nolinkurl{existence/PROOF.md}] \nolinkurl{research/quartic214/proofs/existence.md}
\item[\nolinkurl{parent224-proof.md}] \nolinkurl{research/version15/quartics/quartic-222/provenance/shared-T-redundancy224-006/PROOF.md}
\item[\nolinkurl{quartic_v15.tex}] \nolinkurl{research/version15/snapshots/v15-222-004/manuscript/quartic_v15.tex}
\item[\nolinkurl{brox-v17.tex}] \nolinkurl{research/quartic214/sources/evaluation-sections1-4.tex}
\item[\nolinkurl{parent9-PROOF.md}] \nolinkurl{research/quartic214/sources/conductor-sections1-3.md}
\item[\nolinkurl{cutoff-PROOF.md}] \nolinkurl{research/quartic214/sources/support-section1.md}
\item[\nolinkurl{quartic214-protocol-d/PRIMARY_SOURCE_BRIDGES.md}] \nolinkurl{research/quartic214/proofs/primary-source-bridges.md}
\item[\nolinkurl{quartic214-height-closure/ASSEMBLY.md}] \nolinkurl{research/quartic214/proofs/height-assembly.md}
\item[\nolinkurl{HEIGHT_AVERAGING.md}] \nolinkurl{research/quartic214/proofs/height-averaging.md}
\item[\nolinkurl{GEOMETRY_CARRY.md}] \nolinkurl{research/quartic214/proofs/geometry-carry.md}
\item[\nolinkurl{AGGREGATE_BRIDGE.md}] \nolinkurl{research/quartic214/sources/aggregate-bridge.md}
\item[\nolinkurl{brox228.tex}] \nolinkurl{research/quartic214/sources/factor-extraction.tex}
\item[\nolinkurl{PUBLIC_FULL_COVER_DECISIONS.json}] \nolinkurl{research/quartic214/sources/public-provider-scopes.json}
\item[\nolinkurl{FULL_COVER_SOURCE_PROFILES.json}] \nolinkurl{research/quartic214/sources/public-provider-scopes.json}
\item[\nolinkurl{PUBLIC_GEOMETRY_SCOPES.json}] \nolinkurl{research/quartic214/sources/public-provider-scopes.json}
\item[\nolinkurl{PUBLIC_CONDUCTOR_SCOPES.json}] \nolinkurl{research/quartic214/sources/public-provider-scopes.json}
\item[\nolinkurl{CONDUCTOR_INHERITED_PROFILES.json}] \nolinkurl{research/quartic214/sources/public-provider-scopes.json}
\item[\nolinkurl{CONDUCTOR_RAMIFICATION_PROFILES.json}] \nolinkurl{research/quartic214/sources/public-provider-scopes.json}
\end{description}
\paragraph{Rebuilding this edition.} Run \nolinkurl{python3 -B research/quartic214-manuscript/assemble.py} from the repository root.  It reads bounded source files, emits this single TeX source and \nolinkurl{SOURCE_SPANS.json}, and performs no polynomial expansion or TeX build.  Its \nolinkurl{--check} mode verifies byte-for-byte deterministic output.  The separate PDF build and release metadata record the exact generated-source hash, toolchain and PDF hash.  Source mapping and successful compilation concern serialization, not arithmetic acceptance.
\clearpage\part{The complete current Version 16 mathematical foundation}
\begin{quote}\small\textbf{Scope of retained foundation.} The following is the entire mathematical body of the current upstream Version~16 source (SHA-256 \nolinkurl{aeb2b6a42a03943d7f8c46682389c8da693d2903ba92ccc164bbbd236451de74}), selected by \nolinkurl{docs/current-manuscript.json}.  Its original abstract, theorem statements, proofs, bibliography and historical provenance are retained, including the GI, Gaudron--R\'emond citation and Tao--Ziegler corrections identified in Section~\ref{q214:foundation-transition}.  Assertions about $222$, $314$ and earlier variable counts refer to those constructions.  They do not replace the $214$ row specification in Part~I.  Historical review and incorporation statements belong to that upstream edition; they do not accept the newly assembled 214 manuscript.  Section, theorem, equation and cross-reference numbering in this part is restored to its original convention.\end{quote}
\setcounter{section}{0}\setcounter{subsection}{0}\setcounter{equation}{0}\setcounter{theorem}{0}
\renewcommand{\thesection}{\arabic{section}}\renewcommand{\theHsection}{foundation.\arabic{section}}
\renewcommand{\theHequation}{foundation.\arabic{equation}}
\begin{abstract}
We construct an integer-coefficient polynomial in
$292\,890\,231\,905$ variables whose rational zero projection is exactly
$\Z$. Its degree is less than $2\mathbin{\uparrow\!\uparrow}7$, and the
absolute value of every coefficient is less than
$2\mathbin{\uparrow\!\uparrow}8$.

We show the minimum degree of a Diophantine definition for $\Z$ in $\Q$
is 4 and describe two such quartics. The first has $314$ variables,
$6\,672$ nonzero monomials, and absolute coefficients below $2^{6500}$.
Its complete decimal coefficient table is $264\,430$ bytes.
The second has $222$ variables and $2\,277$ nonzero monomials. Its coefficients are explicit polynomials in a positive integer $C$, for which we have no effective upper bound.

Well-known consequences of $\Z$ being Diophantine in $\Q$ include a refutation of Mazur's conjecture that the topological closure of the rational points of a variety over $\Q$ has finitely many components, and existence of a universal rational Diophantine polynomial. Moreover, the set $\operatorname{HTP}(\Q)$ of integer polynomials with a rational zero is many-one equivalent to the halting set $0'$.
\end{abstract}
% Original second table of contents omitted; unified navigation precedes Part I.

% Begin 01-specification-model.tex
\section{The main theorem and proof structure}\label{sec:main}
This construction builds on OpenAI's preprint
\emph{Hilbert's tenth problem over the rational numbers}~\cite{openaihtp}.

The strategy of using rank-one elliptic curves and existential
local-divisibility conditions to approach undecidability over $\Q$
goes back to Pheidas~\cite{pheidas2000}. His work also developed
partial valuation predicates anticipating the pole conditions used here.

Earlier definability results provide the logical context. Robinson
gave a first-order definition of $\Z$ in $\Q$
\cite[Theorem 3.1]{robinson1949}. Poonen gave a definition with
universal quantifiers followed by existential quantifiers
\cite[Theorems 3.1 and 4.1]{poonen2009}. Koenigsmann gave a universal
definition and proved that $\Q\setminus\Z$ is Diophantine in $\Q$
\cite[Theorem 1 and Corollary 2]{koenigsmann-defining}. These are
historical context, not additional premises of our construction.

We use $\N=\{0,1,2,\ldots\}$. For a reduced rational number $a/b$, with
$b>0$, put $h(a/b)=\log\max(|a|,b)$. Define $\tw1=2$ and
$\tw{(r+1)}=2^{\tw r}$; every tower is associated to the right.

\begin{theorem}[The fully decimal 314-variable quartic]\label{thm:v15-decimal}
There is an explicitly specified polynomial $P_{314}\in\Z[x_1,\ldots,x_{314}]$ of degree four and content one such that, for every $x\in\Q$,
\[x\in\Z\quad\Longleftrightarrow\quad\exists(x_2,\ldots,x_{314})\in\Q^{313}:P_{314}(x,x_2,\ldots,x_{314})=0.\]
It has 6672 nonzero monomials and coefficients of magnitude below $2^{6496}$. The complete ordinary decimal coefficient table has 264430 bytes. The construction and its proof are given in Section~\ref{sec:v14}, with the final constant-coordinate lift in Section~\ref{subsec:v14-ordinary}.
\end{theorem}

\begin{theorem}[The 222-variable family]\label{thm:v15-main}
For $N\in\N$ let $A_N=149^{2N+1}$. Let $Q_{1,N},\ldots,Q_{196,N}$ and the positive integer weights $w_i$ be the expressions in Section~\ref{v222:equations}, in222 coordinates. Put $P_N=\sum_{i=1}^{196}w_iQ_{i,N}^2$. There is an integer $N_0$ such that, for every $N\ge N_0$ and $x\in\Q$,
\[
x\in\Z\quad\Longleftrightarrow\quad\exists(z_2,\ldots,z_{222})\in\Q^{221}:P_N(x,z_2,\ldots,z_{222})=0.
\]
Every $P_N$ has degree four and content one. For $N\ge53$ it has2277 coordinate monomials and largest absolute coefficient $A_N^2+1=149^{4N+2}+1$. The defining threshold $N_0$ is ineffective; the coefficient calculation does not certify a particular defining member.
\end{theorem}
The additional222-variable reduction is proved in Sections~\ref{v222:construction}--\ref{v222:height}. The predecessor proof is given in the internal contact-height, aggregate-height, scalar, prime-selection and transformation sections culminating in Sections~\ref{v15x:reductions}--\ref{v15x:height}. The arithmetic foundation developed throughout this manuscript is reused there with its exact domains. The 248-variable starting construction is due to Jos\'e Brox; its new height argument, aggregate argument, and local construction are included here, rather than used as external black boxes. The 314-variable quartic is the other featured quartic. The large polynomial construction below supplies the original quantitative definition. Earlier quartic systems are retained as proof intermediates for the two featured quartics.

\begin{theorem}[A polynomial definition of the integers]\label{thm:target}
There exists one polynomial
\[
 P\in\Z[X_0,X_1,\ldots,X_{292890231904}]
\]
such that
\begin{align*}
 \deg P&<\tw7,\\
 |[X^\alpha]P|&<\tw8\qquad\text{for every multi-index }\alpha,\\
 q\in\Z&\ \Longleftrightarrow\
 \exists(w_1,\ldots,w_{292890231904})\in\Q^{292890231904},\quad
 P(q,w_1,\ldots,w_{292890231904})=0\quad(q\in\Q).
\end{align*}
All conclusions concern the same polynomial. Here $q\in\Z$ means the
ordinary embedded integers. The ambient variable count does not assert that
every variable occurs essentially.
\end{theorem}

The file \nolinkurl{ordinary-integer-definition.json} specifies the
polynomial $P_*$. The argument first represents integer indices by points
on an elliptic curve over $\Q(\sqrt2)$. A separate value relation identifies
the ordinary rational value of an index. The point coordinates and the
integer index have distinct roles throughout.

Theorem~\ref{thm:target} refutes Mazur's conjecture that the closure of
the rational points of a variety in its real locus has finitely many
connected components. Mazur proposed this in 1992
\cite{mazur1992}; his observation that it obstructs a Diophantine
definition of $\Z$ is recalled in \cite[p.~166]{mazur1995}.
It also gives a universal rational Diophantine
polynomial, with
$\operatorname{HTP}(\Q)$ many-one equivalent to $0'$.
Section~\ref{sec:four-corollaries} proves these three consequences.


The earlier construction gives a quartic in 299 total variables and an ordinary
decimal-coefficient quartic in 314 total variables. Section~\ref{sec:v14}
proves the reductions and specifies their exact coefficient tables.
The Version 13 constructions and their arithmetic inputs remain as proof
baselines. Version 11's index-set and retained-class identification in the
published Kato application is unchanged.

Version 10 clarified determinant specialization, the published
prime-two Euler-system input, the local valuation deduction, and the
identification of the saved templates, following the review in repository
Issue~2. It also corrects and expands the historical attribution.
The polynomial specifications and numerical bounds are unchanged.
We thank \texttt{manvirjaswal} for the detailed review
\footnote{\url{https://github.com/elliotglazer/diophantine-z-q/issues/2}}.

Version 9 incorporated Gabriel Conant's odd-index product-height argument
and asymmetric formal-depth refinement, developed with ChatGPT.
Section~\ref{sec:v9} specifies the resulting 523- and 540-variable quartics
and proves their connection to the arithmetic inputs established in this
manuscript. The earlier quartics below are retained as proof and source
baselines; they are not the current variable records.

% Begin quartic introduction addition
Section~\ref{sec:quartics} gives the explicit $4\,973$-variable quartic.
Its arithmetic source revises the base residual compiler of
Section~\ref{sec:degree8-original}, with a separate proof of soundness and
completeness. The complete coefficient table and exact algebraic conversion
certificate are included with this manuscript.
A third polynomial, in $962$ variables, is constructed in
Section~\ref{sec:v6-962}. Its complete monomial table records the large
integer coefficients using exact powers. The $4\,973$-variable polynomial
is retained with its ordinary, fully expanded integer coefficients.
Section~\ref{sec:degree-lower-bound} proves that a single polynomial
defining $\Z$ in $\Q$ must have total degree at least four, following
the agentic notes cited there. Thus the degree of a quartic definition is optimal.
% End quartic introduction addition

The final composition uses the following precise statements.
\begin{description}[style=nextline,leftmargin=1.8cm]
\item[$\mathsf G$] Canonical point codes form exactly the infinite cyclic
group $2E(\Q(\sqrt2))=\Z P_0$. The membership polynomial types arbitrary inputs;
raw addition describes the group law on canonical codes.
\item[$\mathsf M$] The actual guarded multiplication template has rational
witnesses exactly for triples with indices $a,b,ab$.
\item[$\mathsf H$] The six-contact height theorem holds with the constants
used by the precision computation and, separately, by multiplication.
\item[$\mathsf{PC}$] The explicit pole polynomial is soluble for every
input in Theorem~\ref{hyp:PC}.
\item[$\mathsf C$] The actual five-input natural polynomial and its
rational translation have the exact computation and substitution semantics
of Section~\ref{sec:compiler}, including the power template.
\item[$\mathsf B$] The templates, fixed constructors and private witness
substitutions specify the same saved expression $P_*$. The quantitative
syntax calculation applies to that expression.
\end{description}

\begin{lemma}[Assembly]\label{lem:assembly}
Assume $\mathsf G,\mathsf M,\mathsf H,\mathsf{PC},\mathsf C,\mathsf B$.
Then $P_*$ satisfies Theorem~\ref{thm:target}.
\end{lemma}
\begin{proof}
Section~\ref{sec:bridge} proves the exact ordinary-value relation using
one shared shift and disjoint private witness blocks. The literal
substitutions in $\mathsf B$ identify its existential polynomial with
$P_*$. Section~\ref{sec:size} proves the stated bounds for that same
expression.
\end{proof}

\begin{proof}[Proof of Theorem~\ref{thm:target}]
Section~\ref{sec:model}, including Proposition~\ref{prop:generator},
proves $\mathsf G$ with its canonical membership and addition formulas.
Theorem~\ref{thm:height} proves $\mathsf H$ for each of the two sets of
constants used below. Sections~\ref{sec:rc}--\ref{sec:limits}, with
Appendices~\ref{app:even-source}--\ref{app:endpoints}, prove the restricted
converse in Theorem~\ref{hyp:RC}. The simultaneous descent in
Section~\ref{sec:pole} then proves $\mathsf{PC}$. Using these results,
Section~\ref{sec:mult} proves $\mathsf M$ on all rational input codes,
including the finite false-triple guard.

Section~\ref{sec:compiler} proves $\mathsf C$ for the fixed five-input
natural polynomial and its rational model translation. Its witness
variables remain natural before that translation. Sections~\ref{sec:syntax}
and~\ref{sec:size} and Appendix~\ref{app:literal} prove $\mathsf B$, including
the fixed scalars and private witness substitutions. All six statements
therefore hold, and Lemma~\ref{lem:assembly} applies to the specified $P_*$. The versioned supporting arguments, their source identities and their stated hypotheses are indexed in \href{https://github.com/elliotglazer/diophantine-z-q/blob/main/research/version14/SUPPORTING-NOTES.md}{\texttt{research/version14/SUPPORTING-NOTES.md}}.
\end{proof}

\section{A finite specification of the polynomial}\label{sec:syntax}
The delivered specification is
\nolinkurl{construction/ordinary-integer-definition.json}, a
$23\,179\,783$-byte JSON expression with SHA256
\begin{center}\small\ttfamily
4d639d4c44e69ea360d36d4591b8fdc82269d68b9363af0157d1e5c1174c7af6.
\end{center}
It contains 23 library templates and seven named recipes. The root recipe
is \texttt{integers}; its only parameter is the rational variable at index
zero. Every other variable is existentially quantified over $\Q$.
Sources S01--S03 fix the construction and its semantics. Their historical
specification had $23\,180\,023$ bytes and SHA256
\begin{center}\small\ttfamily
a86790bbf1095d6b1cea9c66067ad8432dd35adc320254eecde21f9f98574714.
\end{center}
The delivered expression replaces two diagnostic \texttt{sourceFile}
strings and the path key in \texttt{sourceHashes} by portable relative
paths. Direct comparison shows that every polynomial-defining field is
unchanged. Since these three provenance fields contribute no equations,
both byte representations denote exactly the same integer polynomial,
with the same variables and bounds. The shipped
\nolinkurl{construction/SOURCE_MANIFEST.json} records both identities and
the normalization; the same manifest retains the historical source references.

\subsection{Literal operations and witness scope}
A scalar node is a variable, integer constant, sum, product, or a call to
an earlier template. A call substitutes its parameter expressions and a
private consecutive witness block into the callee's polynomial. Shared
nodes represent repeated subexpressions, not new equations. A recipe
sums squares of constraint polynomials or multiplies branch polynomials.
Over $\Q$,
\[
 \sum_i f_i^2=0\ \Longleftrightarrow\ \bigwedge_i f_i=0,
 \qquad\prod_i f_i=0\ \Longleftrightarrow\ \bigvee_i f_i=0.
\]
Private witness blocks can be moved to the front because the rational
domain is nonempty. A whole inactive product-call block may receive
arbitrary values. This does not discard fixed equations inside an active
callee. A rational constant $p/q$, $q\ne0$, is represented by a fresh
$c$ and the globally conjoined equation $qc=p$. In particular the cached
coordinates of $P_0$ satisfy
\[
 8g_0=1641,\qquad4g_1=-15465,\qquad32g_2=87483.
\]
These equations remain active outside internal disjunctions in the
multiplication, signed-encoding and exception templates. Inverses are
equations $uv=1$, never partial evaluation operations.

Strict positivity has the polynomial expression
\[
 s>0\ \Longleftrightarrow\
 \exists z_1,z_2,z_3,z_4\in\Q,\quad s\sum_i z_i^2=1.
\]
Indeed, if $1/s=a/b>0$, express $ab$ as four integer squares and divide
their coordinates by $b$. The reverse direction follows from the order
on $\Q$. Absolute value has the exact expression
\[
 a=|x|\quad\Longleftrightarrow\quad
 \exists z_1,\ldots,z_4\in\Q,\quad a=\sum_i z_i^2,\quad a^2=x^2.
\]

\subsection{Fixed finite constructors}
Three node types abbreviate finite polynomials. A
\texttt{fixed\_integer} inserts a fixed integer expression;
\texttt{psi\_component} inserts one coordinate of a division polynomial;
\texttt{finite\_guard\_product} inserts a product over a fixed box of
integers. Their bounds never depend on the free rational input.

For example, on $y^2=x^3+Ax$, write $\psi_n=f_n$ for odd $n$ and
$\psi_n=2yf_n$ for even $n$. The bases are
\begin{align*}
 f_0&=0,&f_1=f_2&=1,\\
 f_3&=3x^4+6Ax^2-A^2,&
 f_4&=2(x^6+5Ax^4-5A^2x^2-A^3).
\end{align*}
For $m\ge3$,
$f_{2m}=f_m(f_{m+2}f_{m-1}^2-f_{m-2}f_{m+1}^2)$.
For $m\ge2$, put $L=f_{m+2}f_m^3$, $R=f_{m-1}f_{m+1}^3$, and
$D_0=16(x^3+Ax)^2$. Then $f_{2m+1}=D_0L-R$ for even $m$ and
$L-D_0R$ for odd $m$. Negative indices satisfy $\psi_{-n}=-\psi_n$.
This is a terminating recursion. These are the short-Weierstrass
normalizations of \cite[\S\S5.5--5.6, Lemmas 5.20 and 5.24, Theorem 5.21]{sutherland}.
In its Jacobian-coordinate convention,
$x([n]P)=\phi_n(P)/\psi_n(P)^2$ and
$y([n]P)=\omega_n(P)/\psi_n(P)^3$.
The numerator formula has a minus sign:
\[
 \omega_n=\frac{\psi_{n+2}\psi_{n-1}^2-
                   \psi_{n-2}\psi_{n+1}^2}{4y},
 \qquad \psi_{2n}=2\psi_n\omega_n.
\]
Thus the fixed-scalar ordinate equation is
$2y([n]P)\psi_n(P)^4=\psi_{2n}(P)$.
For $n\ne0$ and a nonidentity point in the torsion-free carrier,
$\psi_n(P)\ne0$, so clearing these denominators is reversible.

The fixed prime catalogue consists of the first 22 primes above 324 in
each unit residue class modulo $W=4\prod_{p\le324}p$. Dirichlet's theorem
ensures termination. Put $N_s=(\prod_{p\le324}p)H_{\rm cat}$, where
$H_{\rm cat}$ is the catalogue product. Assign inverse abscissa zero
at the identity. For the normalized representation
$1/x([m]P_0)=(a+bs)/c$ with $a,b,c\in\Z$, $c>0$, and
$\gcd(a,b,c)=1$, define $M_0$ to be the least $m\ge1$ such that
$N_s^4\mid a,b$ and $\gcd(c,N_s)=1$ (S31).
This is a fixed integer, with a successful index bounded in
Section~\ref{sec:size}.
The fixed point coordinates in the guard have nonzero denominators by
the infinite order of $P_0$. Thus finite expansion gives one integer
polynomial without requiring a monomial expansion. Its arithmetic zero
set remains a separate issue. In particular, bounding the integer $T$
does not prove that it excludes all false multiplication triples.

\section{The elliptic model of integer indices}\label{sec:model}
The rank-one model and its generator calculation adapt
\cite[Proposition 2.1 and its proof]{openaihtp}.

Let
\[
 K_0=\Q(s),\quad s^2=2,\quad d=75-53s,\quad E:y^2=x^3-d^2x,
\]
and fix
\[
 P_0=\left(\frac{1641-1160s}{8},\frac{-123720+87483s}{32}\right).
\]
The identity has code $(1,0,0,0,0)$. An affine point has code
$(0,x_0,x_1,y_0,y_1)$ with $x=x_0+x_1s$, $y=y_0+y_1s$.
Linear independence of $1,s$ makes this unique. Write $\iota(n)$ for
the code of $[n]P_0$.

\begin{proposition}[Generator certificate]\label{prop:generator}
The point $P_0$ has infinite order and $2E(K_0)=\Z P_0$.
\end{proposition}
\begin{proof}
We give the argument of S05--S07, making its finite certificate explicit.
The ring of integers $\Z[s]$ is norm-Euclidean, hence has class number one.
Its units are $\{\pm(1+s)^k:k\in\Z\}$, so it has four unit squareclasses.
The sign vectors of $-1,1+s$ are independent, so every totally positive unit is a square.

The $x$-squareclass map in the two-isogeny descent on $E$ has image
contained in the eight classes generated by $-1,1+s,d$. Away from $(d)$,
odd $v(x)$ makes the minimum of $3v(x)$ and $v(x)+2v(d)$ odd, impossible
for $y^2$. The torsion points give $1,-1,d,-d$. Put $\epsilon=1+s$,
$\delta=3-s$. The change $(x,y)\mapsto(\epsilon^4x,\epsilon^6y)$ takes
$E$ to $C:y^2=x^3-\delta^2x$, and $R=(1+s,2s)$ provides the remaining
unit class. Thus the first image has order eight.

On $E':y^2=x^3+4d^2x$, a nonzero $x$ is positive in both real embeddings.
Its valuations are even. At $(s)$ the two terms have valuations $3m$
and $m+4$, with equality only at $m=2$. At $(d)$ the only odd case
would require $u^2=-4$ in $\F_7$. Thus the second image is trivial.
The exact two-isogeny sequence gives $[E(K_0):2E(K_0)]=8$; its possible
extra kernel term vanishes since a nonkernel rational two-torsion point
maps to the dual kernel.

The good reductions at $(17,s-6)$ and $(23,s-5)$ have 20 and 24 points.
Prime-to-residue-characteristic torsion injects, so the torsion group
has order at most four, hence is exactly the displayed full two-torsion. Mordell--Weil now gives rank one.
The $x$-class of $R$ has opposite real signs, whereas all torsion classes
have equal signs. Its free index is therefore odd.

The duplication forms on $C$ are
\[
 F_2(a,b)=(a^2+\delta^2b^2)^2,\qquad
 G_2(a,b)=4ab(a^2-\delta^2b^2).
\]
At either real embedding their maximum lies between
$\max(|a|,|b|)^4$ and $676\max(|a|,|b|)^4$. For primitive $a,b$ their
common ideal has valuations at most six at $(s)$ and four at $(\delta)$,
and none elsewhere. For the latter bound use that $-1$ is nonsquare
modulo seven. For the former, a local unit is $A+Bs$ with $A$ odd, and
$v_s((A+Bs)^2+1)$ is two or three according to the parity of $B$.
Hence
\[
 4h_x(T)-\log392\le h_x(2T)\le4h_x(T)+\log676.
\]
With $2\widehat h(T)=\lim_m4^{-m}h_x(2^mT)$ this gives
$h_x(T)\le2\widehat h(T)+\log392/3$. Direct doubling gives
$x(2R)=(57-28s)/8$, with height $\log41$. If $R=mS+T$ for torsion
$T$ and $|m|\ge3$, it follows that
\[
 h_x(S)\le\frac{\log(41^3\cdot676)}{108}+\frac{\log392}{3}<\log9.
\]
The strict inequality is the exact integer comparison
$41^3\cdot676\cdot392^{36}<9^{108}$.

A nonrational $x\in K_0$ in this height range has primitive polynomial
$aX^2+bX+c$ with $1\le a\le81$, $|b|\le162$, $|c|\le81$, and
discriminant $2k^2$, $1\le k\le162$. Both roots are included.
Rational possibilities have coprime numerator of absolute value at most
nine and denominator from one to nine. The exact finite square test of
$x^3-\delta^2x$ gives precisely
\[
 0,\quad\delta,\quad-\delta,\quad1+s,\quad23-17s,
 \quad2+4s,\quad(5-4s)/2.
\]
The certificate covers 7,004 primitive quadratics, their 14,008 roots,
and 111 rational candidates. To test whether $u+vs$ is a square, take
both signs of a rational square root $t$ of $u^2-2v^2$, test
$y_0^2=(u+t)/2$, then set $y_1=v/(2y_0)$ if $y_0\ne0$.
The $y_0=0$ case requires $v=0$ and $y_1^2=u/2$. Substitution verifies
each result. This is a necessary and sufficient test, including zero.
The last four listed coordinates are those of $R+T$ for $T\in C[2]$.
Every surviving nontorsion point therefore has height $\widehat h(R)$,
contradicting $\widehat h(R)=m^2\widehat h(S)$. This finite computation
is an explicit part of the proof, with a recorded independent replay;
it is not replaced by a rank oracle.

Thus $2C(K_0)=\Z(2R)$. The inverse coordinate change sends $2R$ to
$P_0$, as asserted.
\end{proof}

\subsection{Polynomial membership and addition}
Apart from the identity branch, membership in $2E(K_0)$ is expressed by
choosing $X,Y,H\in K_0$, $u\in\Q$, and imposing
\[
 Y^2=X^3-d^2X,\quad u\Norm(Y)=1,\quad2YH=3X^2-d^2,
 \quad x=H^2-2X,\quad y=H(X-x)-Y.
\]
Every nonidentity double has a preimage with nonzero ordinate.
Each field equation is two rational equations. Membership is an all-input
statement. Raw addition has the narrower exact contract
\[
 \operatorname{Add}_{\rm raw}(\iota(a),\iota(b),R)
             \quad\Longleftrightarrow\quad R=\iota(a+b).
\]
Its branches cover both identities, opposite points, distinct abscissas
and doubling with nonzero ordinate. Every later use supplies canonical
input codes. Multiplication first imposes membership of all three roots
and is sound on arbitrary rational tuples. In the natural predicate,
four initially untyped tuples enter multiplication calls, so are typed
before their squared values are added. Lagrange's theorem then gives
exactly $\iota(\N)$. This uses multiplication completeness; the generator
proposition alone does not supply the natural predicate.

% End 01-specification-model.tex


% Begin 02-height-pole-multiplication.tex
\section{The six-contact height estimate}\label{sec:height}
The contact and odd-support framework follows
\cite[Section 4]{openaihtp}. The effective geometric estimates below
give the height input used in this manuscript; the preprint's height
assertion is not taken as an additional premise.

Set
\[
 \Cset=\{0,-15,-20,-24,-30,-60\},\qquad
 F(u)=u(u+15)(u+20)(u+24)(u+30)(u+60),
\]
and $G(u)=7u^2+300u+3600$. For a finite prime set $S$ containing all
primes at most 324 and $u\notin\Cset$, define
\[
 M_S(u)=\prod_{\substack{p\notin S:\ v_p(u-c)>0\text{ odd}\\
                          \text{for some }c\in\Cset}}p.
\]
Negative valuations of $u$ do not enter the product. Put
\begin{align*}
 F_0&=324!,\qquad D=F_0\,3^{256F_0^4},\qquad G_0=116064000000,\\
 E_0&=10(DG_0)^2,\qquad Q_R=R^{RE_0},\\
 B_{\rm cm}&=D^2\,3^{21000D^2}+1,\\
 C_{\rm cm}&=2B_{\rm cm}+4096(1028+5D+84000D^2),\\
 H_6(S;R)&=2\cdot10^{13}\max(1,C_{\rm cm},Q_R^{17}),\\
 &\hspace{1em}R\ge\max(D,\max S,10^{12}).
\end{align*}
When the second argument is omitted, we use
$H_6(S)=H_6(S;\max(D,\max S,10^{12}))$.
The geometric construction is followed by the published prime-five
modularity theorem and a separate effective CM argument. The constants
in the displayed definition of $H_6(S;R)$ are unchanged.

\begin{lemma}[Modularity after rational scalar descent]
\label{lem:qm-modularity-descent}
Let $A$ be an abelian surface over a finite Galois extension of $\Q$,
with full geometric endomorphism algebra a quaternion division
algebra $B$ over $\Q$. Suppose that $B$-compatible quasi-isogenies
$\mu_\sigma:{}^\sigma A\to A$, with $\mu_1=1$, form a finite
locally constant choice on $G_\Q$, and put
\[
 c(\sigma,\tau)=\mu_\sigma{}^\sigma\mu_\tau\mu_{\sigma\tau}^{-1}
 \in\Q^\times.
\]
Suppose that a locally constant cochain $t:G_\Q\to C^\times$ in
a number field $C$ satisfies
\[
 t_\sigma t_\tau c(\sigma,\tau)=t_{\sigma\tau}.
\]
The added coefficients have trivial Galois action. Enlarge $C$ to
split $B$, and fix its embedding into $\overline{\Q}_5$.
If $\rho_5(\sigma)=t_\sigma\mu_\sigma\circ\sigma$
is written as $r_5\oplus r_5$ using the matrix units of $B$,
then $r_5$ is the Tate realization of a primitive weight-two
cuspidal newform, after a finite coefficient extension.
No bound on an auxiliary field or abelian variety is part of this assertion.
\end{lemma}

\begin{proof}
We apply the published prime-five modularity statement reported in
\cite[p.~624, with the definition on pp.~611--612]{calegari-icm2022}
to this same representation. We check its hypotheses and the finite
twist needed for the conclusion.

The cocycle equation gives $t_1=1$. Choose a finite Galois extension
$F/\Q$ over which $A$ and all its geometric endomorphisms are defined
and on whose Galois subgroup both locally constant choices are the
identity. Then
\[
 V_5(A)\otimes_{\Q_5}E\simeq
 r_5|_{G_F}\oplus r_5|_{G_F}
\]
for a finite coefficient extension $E/\Q_5$. The continuous Tate action
on this open subgroup and the finite coset decomposition prove
continuity of $r_5$ on $G_\Q$. Faltings' semisimplicity
and endomorphism theorem
\cite[Theorem 1(a)--(b), p.~204]{faltings1986-complements}
identifies the commutant on the left with $B\otimes E=M_2(E)$.
The multiplicity commutant on either copy is therefore $E$.
The same calculation after coefficient extension, together with
semisimplicity, shows that $r_5|_{G_F}$ is absolutely irreducible.
Thus $r_5$ is absolutely irreducible. Good reduction of $A$ outside
a finite set, enlarged by the primes ramified in $F/\Q$, proves
that $r_5$ is unramified outside a finite set.

At a place of $F$ above five, proper-smooth comparison
\cite[Corollary 1.8 and Theorems 9.1, 9.3]{scholze2013}, with
\cite{scholze-erratum}, makes $H^1_{\mathrm{et}}(A,\Q_5)$ de Rham.
The canonical dual is $V_5(A)$
\cite[Theorem 15.1(a), p.~129]{milne1986-abelian}.
In the convention that $\Q_5(1)$ has weight $-1$, its weights are
$0,0,-1,-1$. De Rham representations are closed under duals,
tensor products and subquotients
\cite[Theorem 3.8, p.~142]{fontaine1994-representations}.
The two identical summands therefore each have weights $0,-1$.
This conclusion descends to $\Q_5$: after a finite Galois local
extension $K'/\Q_5$, the full-rank filtered space
$(B_{\mathrm{dR}}\otimes r_5)^{G_{K'}}$ has its natural semilinear
Galois action. Descent of that space and each filtration step gives
the full-rank filtered space over $\Q_5$ with the same weights.
No good-reduction or crystalline hypothesis at five is used.

For oddness, fix an embedding into $\mathbb{C}$ and its complex conjugation
$c$. On Betti homology, $c$ interchanges the two Hodge summands,
whereas the holomorphic quasi-isogeny $\mu_c$ preserves their types.
Thus $t_c\mu_c c$ has trace zero. The complex lattice description
\cite[Theorem 6.8 and the preceding paragraph, p.~319]{milne-shimura2005}
identifies its $5$-power torsion functorially with the homology
lattice modulo $5^n$. This also respects conjugation; clearing
isogeny denominators and taking the inverse limit transfers the
trace identity to the Tate module under the chosen coefficient
embedding. Hence $2\operatorname{tr}(r_5(c))=0$.
Since $r_5(c)^2=1$, its eigenvalues are $1,-1$, so it is odd.

The cited published statement gives, after a finite coefficient extension,
$r_5\simeq U_g\otimes\delta$ for a nonzero classical cuspidal eigenform
$g$ of positive integral weight $k$. We use arithmetic Frobenius for the
associated semisimple representation. We first choose a normalized primitive
representative without assuming that $k=2$.
The decomposition in \cite[\S1.9, p.~510]{deligne-serre1974} writes $g$
as a sum of primitive oldforms $f_j(d_jz)$; its Eisenstein part is zero.
For a prime $\ell$ not dividing the original level, the coefficient formula
in \cite[\S1.6, p.~509]{deligne-serre1974} shows that $T_\ell$ commutes
with each such degeneracy map. Uniqueness of the decomposition implies
that every nonzero constituent has the same good-prime eigenvalues as $g$.
Choose one primitive constituent $f$, normalized by $a_1(f)=1$ as in
\cite[\S1.7, p.~509]{deligne-serre1974}. Its induced nebentypus and
good-prime characteristic polynomials agree with those of $g$.
Semisimple uniqueness \cite[Lemma~3.2, p.~513]{deligne-serre1974}
therefore gives $r_5\simeq U_f\otimes\delta$ with the same character.
All coefficient embeddings are extended compatibly with the one already fixed.

If $k=1$, the transformation under $-I$ makes the nebentypus odd.
Theorem~4.1 of \cite[pp.~513--514]{deligne-serre1974} supplies the
associated irreducible finite-image representation. Finiteness here uses
the discrete complex topology specified in its \S3.1; its Frobenius is
arithmetic. Moreover, \S\S8.2, 8.5--8.6, pp.~525--526, construct this
representation over a number field, so it has a finite-image $5$-adic
realization with the same good-prime polynomials. Lemma~3.2 identifies it
with $U_f$. After a finite local extension killing this finite image,
$r_5$ is the direct sum of two copies of the same character. De Rham
subobject closure and finite-extension compatibility then give two equal
Hodge--Tate weights, contrary to $\{0,-1\}$. These closure and
base-change properties are also stated in \cite[\S9.2, p.~187]{kato}.
Thus $k\geq2$.

For this range, let $E_f$ be the coefficient field of $f$ in Kato's
construction, with the place $\lambda$ above five fixed by the chosen
embedding. Write $V=V_{E_{f,\lambda}}(f)$ for its realization.
It has dimension two over $E_{f,\lambda}$ and $\dim_{E_f} S(f)=1$
\cite[\S6.3, pp.~161--162]{kato}. Section~9.4, p.~188, states that
$V$ is de Rham and gives the decreasing filtration
\[
 \operatorname{Fil}^i D_{\mathrm{dR}}(V)=
 \begin{cases}
 D_{\mathrm{dR}}(V),&i\leq0,\\
 S(f)\otimes_{E_f} E_{f,\lambda},&1\leq i\leq k-1,\\
 0,&i\geq k.
 \end{cases}
\]
This holds at every level, including levels divisible by five;
\S\S5.4 and 8.3, pp.~155 and 181, include the small-level objects
through fine-level invariants. No non-CM or crystalline hypothesis is used.
Writing $\varepsilon_f$ for the nebentypus and $\epsilon_5$ for the
cyclotomic character, \cite[(14.10.2)--(14.10.3), p.~240]{kato} gives
\[
 \det V=\epsilon_5^{1-k}\varepsilon_f^{-1}
\]
and puts the good-prime Hecke polynomial on inverse arithmetic Frobenius
in $V$. Hence arithmetic Frobenius on $V^\vee$ has the polynomial
defining $U_f$. Their semisimplifications agree by the same uniqueness
lemma. Since $U_f$ is irreducible, being a character twist of $r_5$,
the two-dimensional $V^\vee$ is itself irreducible and isomorphic to $U_f$.
Duality therefore yields
\[
 \operatorname{HT}(U_f)=\{0,1-k\},\qquad
 \det U_f=\varepsilon_f\epsilon_5^{k-1}.
\]
There is no further Tate twist in this identification.

The $\delta$-line in $\operatorname{Hom}(U_f,r_5)$ is Galois-stable;
tensor, dual and subobject closure make it de Rham. If its weight is $m$,
then
\[
 \{m,m+1-k\}=\{0,-1\}.
\]
Since $k\geq2$, comparison of the larger entries and then the smaller
entries gives $m=0$ and $k=2$.
To prove that the twist is finite, put $\alpha=\det(r_5)/\epsilon_5$.
The polarization pairing on $V_5(A)$ gives
$\det V_5(A)=\epsilon_5^2$, so $\alpha^2=1$ on $G_F$.
Its image on all of $G_\Q$ is consequently finite. The determinant
formula just proved now gives $\delta^2=\alpha/\varepsilon_f$.
The right side has finite image, and each nonzero scalar has at most two
square roots in characteristic zero. Thus $\delta$ has finite image;
this is not inferred from weight zero alone.

The fixed field of $\ker\delta$ is a finite abelian extension of
$\Q$. By \cite[Theorem 14.1, p.~321]{washington1997}, it lies in
a cyclotomic field. Let $\chi$ be the corresponding primitive
Dirichlet character, with
$\chi(q)=\delta(\operatorname{Frob}_{q,\mathrm{arith}})$.
If $M=\operatorname{level}(f)$ and $r=\operatorname{cond}(\chi)$,
Shimura's translate formula
\cite[Section 4, equation (4.1), p.~531]{shimura-factors}
gives a nonzero Gauss sum times
$h=\sum_{n\geq1}\chi(n)a_n(f)q^n$.
For completeness, conjugating the translation matrices by an
element $\left(\begin{smallmatrix}A&B\\C&D\end{smallmatrix}\right)
\in\Gamma_0(Mr^2)$ permutes their indices by
$a\mapsto aD^2\pmod r$. The weighted sum therefore transforms
with character $\varepsilon_f\chi^2$ and remains cuspidal.
The Hecke coefficient formula gives
$T_\ell h=\chi(\ell)a_\ell(f)h$ for $\ell\nmid rM$.
\paragraph{Operative primitive-form supplier for this receiving revision.}
For the current $214$ application, the existence assertion in the next
historical Li attribution is supplied instead by the explicit
Deligne--Serre extraction in Part~I, M3
(Section~\ref{q214:modularity}), using \cite[\S\S1.6--1.9, pp.~509--510]{deligne-serre1974}.
The nonzero cusp eigenform $h$ has a nonzero primitive constituent;
good-prime Hecke operators commute with its oldform degeneracy maps,
and uniqueness of that decomposition passes the same eigenvalues to it.
Normalize its first coefficient to $1$. Here the twisting character is
primitive of conductor $r$ as chosen above; in the Part~I Milne route,
take $r=\operatorname{cond}(\eta)$, the least modulus, with the same
Galois character and controlled field. The application is bound explicitly
in \nolinkurl{research/quartic214-protocol-d-progress/CURRENT_NATIVE_PRIMITIVE_APPLICATION.json}.
Li's uniqueness and Fricke uses remain separate. The following original
wording is retained as historical source evidence, qualified by this
operative supplier; no existence claim rests on Li's Theorem~5 alone.
Li's old/new decomposition
\cite[p.~294 and Theorem 5, p.~298]{li-newforms}
supplies a primitive weight-two newform with these good-prime
eigenvalues. Its Tate representation is $U_f\otimes\delta$ by
Chebotarev and semisimplicity, hence is the original $r_5$.
No conductor bound is inferred from this qualitative twisting step.
\end{proof}

\paragraph{Scope of the modularity input.}
The immediate supplier is the published statement in
\cite[p.~624]{calegari-icm2022}; Pan's work \cite{pan2022} is its
proof source. This use imposes no residual restriction at five.
It does not assert that every input to Pan's proof has been
reproved here. In particular, his ordinary inputs and his local
Hu--Tan ring calculations are distinct from the global Hu--Tan
support argument. The separate GL2-type and Serre weight-transfer
reconstructions are not inputs to the present lemma.

The exact source-to-object and historical-dependency map is given in
\cite{direct-p5-modularity-project}.

\begin{theorem}[Six-contact estimate]\label{thm:height}
For every such $S,R$ and every $u\in\Q\setminus\Cset$,
\[
 h(u)\le H_6(S;R)M_S(u)^{24}.
\]
\end{theorem}
\begin{proof}
Take $B=(2,-3)$, the maximal order $\mathcal O=\Z[i,(1+j)/2]$, and
$\mu=2j+ij$. With its positive sign choice, the form
$\Trd(\mu^{-1}x\overline y)$ has matrix
\[
 \begin{pmatrix}0&0&-1&1\\0&0&1&-2\\1&-1&0&0\\-1&2&0&0\end{pmatrix}
\]
and determinant one, in the ordered basis
$(1,i,(1+j)/2,(i+ij)/2)$. The Riemann form is the convention of
\cite[\S2, pp.~779--780, Lemma 3]{linyang}; its Rosati involution is
$a\mapsto\mu^{-1}\overline a\mu$.
The form printed in \cite[\S2.1]{babagranath},
$\Trd(\overline x\mu y)$, is six times this form, since $\mu^2=-6$.
It has determinant $6^4$ on this lattice. We use the determinant-one
form above; the positive scalar does not change the Rosati involution
or which isomorphisms preserve the form.
The corrected form has the prescribed positive Rosati involution, and the
principal polarization inducing it is unique \cite[p.~592]{buzzard}.
Thus it gives the abstract polarized datum in
\cite[\S2.1, Theorem~2.1]{babagranath}. For this order the two twisting
involutions are $w_2,w_3$, and they generate the full four-element
Atkin--Lehner group $W$. The theorem identifies the polarized forgetful
quotient with $V/W$, where $V=\mathfrak H/\mathcal O^1$.
The canonical model and the involutions are defined over $\Q$.
Lemma~3.1 and the fixed-point table of \cite[pp.~739--740]{babagranath}
give a genus-zero quotient with three distinct rational marked points.
The two elliptic stabilizer orders on $V$ are $2$ and $3$.
A finite subgroup fixing a point of a smooth characteristic-zero curve
acts faithfully on its tangent line, so it is cyclic. Hence a point fixed
by a nontrivial element of $W$ has $W$-stabilizer of order two, rather
than four. The three resulting branch indices are therefore $4,6,2$.
We take the unique coordinate $j$ with values $0,\infty,-16/27$ at
the CM points of discriminants $-4,-3,-24$, respectively.
This is the coordinate defined in \cite[Example~30]{linyang}.
The principal-form correction and the abstract coarse-moduli application
are detailed in \cite{a001-moduli-project}.

We also need its compatibility with the explicit Igusa invariants.
Put $s_2=\psi_4$, $s_3=\psi_6$,
$s_5=2^{12}3^5\chi_{10}$ and $s_6=2^{12}3^6\chi_{12}$.
The divisor calculations in \cite[Examples~29--30, pp.~795--797]{linyang}
give
\[
 \frac{s_3^2}{s_2^3}=c_1j,\qquad
 \frac{s_5}{s_2s_3}=c_2,\qquad
 \frac{s_6}{s_2^3}=c_3j+c_4.
\]
Its product-curve parameterization at elliptic invariant $1728$ gives
$c_4=2^{12}3^6/1728^2=1$. For the remaining constants use its
$Y'_5$ intersection at discriminant $-24$. The stated eta-coordinate
$x=-6\pm6\sqrt{-3}$ satisfies $x^2+12x+144=0$; substitution in the
published parameterization gives
\[
 s_2=-36x,\quad s_3=6912,\quad s_5=-248832x,
 \quad s_6=-32845824.
\]
The three ratios are $-16/27,1,11/27$, and their denominators are
nonzero. Since $j=-16/27$ at this point, $c_1=c_2=c_3=1$.
On the finite nonzero $j$-locus, a common weighted rescaling gives
$[s_2,s_3,s_5,s_6]=[j,j^2,j^3,j^3(j+1)]$ in weights $2,3,5,6$.
Equations~(13)--(14) of \cite[p.~793]{linyang}, followed by rescaling
by $-12$ in weights $1,2,3,5$, therefore give
\begin{equation}\label{eq:qm-normalized-igusa}
 [J_2,J_4,J_6,J_{10}]
 =[12(j+1),6(j^2+j+1),4(j^3-2j^2+1),j^3].
\end{equation}
This calculation uses the published intersection, divisor and eta-value
inputs at their stated scope; its source comparisons and weighted
conventions are in \cite{ly-d6-normalization-project}.
It does not use the product-isomorphism arguments of
\cite[Propositions~3.7--3.8]{babagranath}, the subsequent numerical
CM evaluation, or its explicit conic descent.
Let $Y_3$ be the generic full-three-marked moduli curve, with all its
determinant components. Its full adelic description and the same
positive-normalizer forgetful quotient are proved below, after the
integral construction. We first compute the arithmetic index which
that identification will permit us to apply.
Put $K_0=\widehat{\mathcal O}^{\,*}$,
$K_3=\ker(K_0\longrightarrow(\mathcal O/3\mathcal O)^*)$, and let
$N$ be the finite-adelic normalizer of $\widehat{\mathcal O}$.
It normalizes $K_3$, since it preserves the ideal $3\widehat{\mathcal O}$.

The completion $\mathcal O_p$ is maximal. Indeed a larger local order
inside $p^{-k}\mathcal O_p$ would intersect $\mathcal O[1/p]$ in a
larger global order inside $p^{-k}\mathcal O$; coordinate density in
$B_p$ makes that containment strict. At a split prime identify
$\mathcal O_p$ with $\operatorname{End}_{\Z_p}(L)$ for a rank-two
lattice $L$. If $g$ normalizes this order, the matrix units show that
$gL=I e_1\oplus I e_2=p^rL$. Thus the local normalizer is
$\Q_p^*\mathcal O_p^*$.
At each division prime the maximal order is the unique nonnegative
valuation ring, with $w=v_p(\operatorname{nrd})/2$.
Its unit group has valuation zero, while central elements have integral
valuation and a division uniformizer has half-integral valuation.
The quotient by center and units consequently has order two.
These local facts, including the prime two, are used with the precise
normalization of \cite[\S13.3]{voight2021}.
Since the division primes are exactly two and three, their independent
valuation parities give $|N/(\mathbb A_f^*K_0)|=4$.

A central finite idele acts trivially on this full level-three curve.
Multiply it by a rational scalar to cancel all finite valuations and
choose the sign so that its residue at three is one. The remaining
idele is in $K_3$; the rational central scalar is absorbed on the left
of the adelic quotient and acts trivially on the symmetric domain.
The given normalizer action thus factors through
$N/(\mathbb A_f^*K_3)$. Its quotient is the same $\mathbf P^1_j$,
and a geometric generic fiber is an orbit. Therefore
\[
 \deg(Y_3\longrightarrow\mathbf P^1_j)
 \le4[K_0:K_3]\le4|(\mathcal O/3\mathcal O)^*|
 \le4\cdot81=324.
\]
This uses the kernel definition of $K_3$, so no surjectivity of reduction
is needed. It also retains all components throughout.

We construct an integral model without changing this generic degree.
Write $R_0=\Z[1/30]$ in the following construction. The full proofs and
the exact source hypotheses for its six phases are recorded in
Section~\ref{v15:marking-descent}.

\emph{The initial model.}
Apply \cite[Theorem~2.1 and Corollary~2.3]{buzzard} to $V_1(5)$.
At five its determinant image is full because it contains the diagonal
matrices $\operatorname{diag}(a,1)$ for every unit $a$.
At the other split primes, maximal-order units have arbitrary determinant.
At a division prime, the maximal order contains the integers of the
unramified quadratic field, whose unit norms are all of $\Z_p^*$;
restriction of reduced norm gives the same surjectivity.
Thus the entire adelic determinant condition holds. The level is maximal
at the division primes, and $5\ge4$ is prime to six. The result is a
smooth projective fine model $X_{V_1}/R_0$ with its universal object.
Here Corollary~2.3 identifies the actual groupoid with the stack of a
scheme, so its objects have trivial automorphisms over arbitrary base
schemes, including nonreduced ones.

\emph{The full five-marking.}
The full-five marking sheaf and its $H$-orbit sheaf in
\cite[Definition~1.1, p.~593]{buzzard} apply at this prime-to-discriminant
level. Pulling the full-marking torsor back along the existing
$V_1(5)$-marking gives a finite \'{e}tale surjective $H$-torsor of
refinements. Its total space $X_5$ represents full five-markings and
retains every determinant component. It is smooth projective over $R_0$.
For a left-module marking $m$ and a unit $g$, set
$g\cdot m=m\circ r_g$, where $r_g(x)=xg$.
Since $r_h\circ r_g=r_{gh}$, this is the left action used below.
An automorphism preserving a full five-marking preserves its induced
$V_1(5)$-marking, so the full-five groupoid is also rigid.

\emph{The full three-marking.}
For an abelian scheme and an invertible integer $n$, multiplication by
$n$ is finite \'{e}tale of rank $n^4$ on its kernel in relative dimension
two. The fiberwise multiplication theorem, flatness and \'{e}tale-fiber
criteria give the assertion over arbitrary bases; properness and
quasi-finiteness give finiteness. Finite \'{e}tale schemes are finite
locally constant sheaves. Requiring a bijection to respect addition and
a finite generating set of $\mathcal O$ imposes finitely many equations
on a finite set. Hence
\[
 X_{35}=\operatorname{Isom}_{\mathcal O}
       (\mathcal O/3\mathcal O,A[3])\longrightarrow X_5
\]
is finite \'{e}tale. Its surjectivity requires a separate calculation.
On every geometric fiber, $T_3(A)$ is free of rank four over $\Z_3$.
The division algebra $B_3$ acts unitally on its four-dimensional rational
Tate module. Multiplication on any nonzero vector embeds $B_3$ into
that module, so it is a rank-one left $B_3$-module.
Its $\mathcal O_3$-stable lattice has a nonzero element $x$ of least
division valuation: the values are discrete and bounded below by the
minimum of the valuations of a lattice basis. For every $y$ in the
lattice, $w(yx^{-1})\ge0$, so $y\in\mathcal O_3x$; the reverse
inclusion follows from stability. Reduction modulo the central element
three gives $A[3]\simeq\mathcal O/3\mathcal O$.
Thus $X_{35}\to X_5$ is surjective in every allowed characteristic.
It is smooth projective over the affine, hence quasi-compact and
quasi-separated, base $R_0$.

\emph{Freeness and rigidity.}
Let $G_5=(\mathcal O/5\mathcal O)^*$ change only the five-marking.
A geometric stabilizer gives a polarized automorphism $\alpha$
preserving the three-marking. Polarized automorphisms form a finite
group, and the action on $T_3(A)$ is faithful, by the integral-lattice proof in Section~\ref{v15:tate-polarized}. The matrix of $\alpha$
lies in $1+3M_4(\Z_3)$, which is torsion-free.
Indeed, a nontrivial torsion element has a power of prime order $\ell$;
write it as $1+3^rM$ with $r\ge1$ and $M$ having a unit entry.
If $\ell\ne3$, the linear term in its $\ell$th power has entry
valuation $r$, while higher terms have valuation at least $2r$.
For $\ell=3$, these bounds are $r+1$ and at least $r+2$.
Neither power is the identity. Hence $\alpha=1$, and the five-marking
then forces the element of $G_5$ to be one.
For each nonidentity element of this constant group, its fixed subscheme
has no geometric points and is empty. This proves a free scheme action
without assuming that $|G_5|$ is invertible on $R_0$.

The same argument gives geometric rigidity of every full-three object.
To pass to an arbitrary base $T$, its automorphism acts identically on
the finite \'{e}tale $A[5]$, since morphisms of finite locally constant
sheaves are detected on geometric fibers. After adding a five-marking
it is the identity by the full-five rigidity above. The existence of
that faithfully flat marking cover for every object is proved below;
descent of equality of morphisms then gives rigidity over $T$.

\emph{The quotient and the exact descended object.}
The free finite-group quotient theorem gives a scheme
$Y=X_{35}/G_5$ and a $G_5$-torsor $q:X_{35}\to Y$.
Because $G_5$ is finite \'{e}tale, so is $q$.
Smoothness descends; the diagonal is closed after pullback, where it
is the union of the action graphs, and properness descends along this
finite surjection. An ample invertible sheaf descends in existence by
its norm, making $Y/R_0$ projective. This norm is used only for the
base quotient, not for its abelian-scheme polarization.
The regular-base projectivity theorem for abelian schemes makes the
universal abelian scheme projective over $X_{35}$ before any use of
duality. Projective schemes descend effectively along the finite
locally free surjection $q$. The relabeling isomorphisms obey the
cocycle identity, and faithful descent carries the group maps and their
identities to an abelian scheme $A/Y$.
Base-compatible duality identifies the pullback of $A^\vee$ with the
original dual. The particular polarization morphism therefore descends;
being an isomorphism is faithfully-flat local, and its geometric fibers
are the original polarizations. It remains principal, of degree one.
The $\mathcal O$-action and three-marking likewise descend as actual
maps. The projectivity, Picard and descent inputs, in this order, are
specified in Section~\ref{v15:marking-descent}.

\emph{The whole full-three functor.}
Let a full-three object be given over any $R_0$-scheme $T$.
On a geometric fiber, $A[5]$ has dimension four over $\mathbf F_5$ and
$\mathcal O/5\mathcal O\simeq M_2(\mathbf F_5)$ acts unitally.
Its matrix units split $A[5]=e_{11}A[5]\oplus e_{22}A[5]$; the two
summands are isomorphic through $e_{21},e_{12}$ and have dimension two.
For a basis $v_1,v_2$ of the first, $E_{ij}\mapsto e_{i1}v_j$ is a
left $M_2(\mathbf F_5)$-module isomorphism with the regular module.
Thus the finite \'{e}tale sheaf of full five-markings is surjective
for this arbitrary object and is a $G_5$-torsor $P\to T$.
The tautological extra marking gives an equivariant map $P\to X_{35}$,
whose invariant composite descends to $T\to Y$.
Conversely, a map to $Y$ pulls back the descended universal object.
These constructions are inverse after the faithfully flat cover $P$;
the rigidity just proved makes their comparison isomorphisms unique,
so they are inverse over $T$ and compatible with base change.
Consequently $Y$ represents the entire full-three problem.

\emph{The generic adelic curve and its rational model.}
We now identify this constructed full-three scheme with the precise
generic curve used in the degree calculation.
Use $B=(2,-3)$, $\mathcal O=\Z[i,(1+j)/2]$, $\mu=2j+ij$ and
\[
 E_0(x,y)=\operatorname{Trd}(\mu^{-1}x\bar y),\qquad
 a^\dagger=\mu^{-1}\bar a\mu,\qquad \mu^2=-6 .
\]
The principal-datum calculation above supplies the integral unimodular
positive form and its prescribed Rosati involution. In particular
we do not replace it by the six-times form in the printed source.
Put $K_n=\ker(K_0\to(\mathcal O/n\mathcal O)^*)$ and
$K_{35}=K_3\cap K_5$.
Everything in this paragraph is in characteristic zero. The integral
three-marking construction has already been proved separately.

Take $V=B$ as a left $B$-module. Its $B$-linear automorphisms are right
multiplications, and $R_b$ scales $E_0$ by $\operatorname{nrd}(b)$.
For an ordinary $B^*$ coordinate rather than its opposite group use
$b\mapsto R_{b^{-1}}$, whose multiplier is $\operatorname{nrd}(b)^{-1}$.
This inversion preserves every congruence subgroup $K_n$.
The identity $E_0(ax,y)=E_0(x,a^\dagger y)$ verifies the adjoint
condition of the PEL datum in
\cite[\S4.9]{Deligne1971}.
For $z=x+iy$ in either half-plane set
\[
 v_z=(z,1)^t,\qquad
 J_z=y^{-1}\begin{pmatrix}x&-|z|^2\\1&-x\end{pmatrix}.
\]
Then $J_z^2=-1$ and $J_zv_z=\mathrm i v_z$.
Right multiplication by $J_z$ is the complex structure for the
period map $u\mapsto uv_z$.
Deligne's homological convention has characters
$w^{-1},\bar w^{-1}$, so set
$h_z(a+b\mathrm i)=R_{(a+bJ_z)^{-1}}$.
In particular $h_z(\mathrm i)=-R_{J_z}$.
On the positive half-plane $E_0(u,h_z(\mathrm i)u)>0$;
on the other half-plane the effective positive form is $-E_0$.
Negative-norm elements exchange these choices, as allowed by the
homogeneous polarization in \cite[\S\S4.8--4.10]{Deligne1971}.
The left $B$-action on the complex Lie space has its reduced
characteristic polynomial. The associated conjugacy class of
cocharacters has the two standard weights and is rational.
Thus the reflex field is $\Q$.

Conversely, for a polarized $\mathcal O$-surface over $\mathbb C$ with this
Rosati action, $H_1(A,\Q)$ is a faithful four-dimensional left
$B$-module and hence has rank one, since $B$ is division.
In a rational frame its alternating form is a rational scalar
multiple of $E_0$. Indeed write $E'(x,y)=E_0(x,Ty)$.
The adjoint identities force $T$ to commute with left $B$, so
$T=R_b$. Alternation makes it self-adjoint for $E_0$, hence
$b=\bar b\in\Q$. The complex structure is likewise right multiplication
by an element squaring to $-1$, and gives one of the two displayed
half-planes. Its effective polarization chooses the positive sign.
These checks give the actual PEL and Lie-trace conditions of
\cite[\S4.10]{Deligne1971} on every object, without fixing a
determinant component.

In such a frame let $I=H_1(A,\Z)\subset B$. At a division prime
$p=2,3$, a least-valuation element $x$ of its full
$\mathcal O_p$-stable lattice gives $I_p=\mathcal O_px$, exactly by
$w(yx^{-1})\ge0$ as above. At a split prime, matrix idempotents give
$I_p=e_{11}I_p\oplus e_{22}I_p$ and $e_{21},e_{12}$ identify the
two rank-two $\Z_p$-summands. For a basis $v_1,v_2$ of the first,
$E_{ab}\mapsto e_{a1}v_b$ is a left $M_2(\Z_p)$-module isomorphism.
It extends to an invertible right multiplication on $B_p$.
Consequently $\widehat I=\widehat{\mathcal O}g$ for
$g\in B(\mathbb A_f)^*$.
Conversely $B\cap\widehat{\mathcal O}g$ is a full lattice with these
completions: clear the finitely many denominators and apply the
additive Chinese remainder theorem. This uses no multiplicative
strong approximation or ideal-class-number assertion.

Choose $c>0$ rational with
$v_p(c)=v_p(\operatorname{nrd}(g_p))$ at every prime.
The form $c^{-1}E_0$ on $I$ is integral and unimodular, since in its
local frame it is $E_0$ times the unit
$\operatorname{nrd}(g_p)/c$. Choose the effective sign with the
half-plane. Riemann's equivalence
\cite[Theorem 4.7]{Deligne1971} then gives the principally polarized
surface. The uniqueness of the principal polarization with this
Rosati involution in \cite[p.~592]{buzzard} identifies the same
polarization in our functor.

The actual analytic family is $\mathbb C^2/(Iv_z)$. On a simply connected
piece of a half-plane its homology local system is the constant
lattice $I$. Changing rational frame by $b\in B^*$ changes
$(I,z)$ to $(Ib^{-1},bz)$, with the actual torus isomorphism given by
$bv_z=(c_bz+d_b)v_{bz}$. The exponential sequence identifies
$A[n]=(n^{-1}I)/I$ with $I/nI$ as left $\mathcal O$-modules.
A generator of $\widehat I$ thus supplies the actual torsion marking;
its stabilizer is $K_n$. Every unit modulo $p^a$ lifts to a local
unit: lift it and a modular inverse and invert their product
$1+p^at$ by its convergent geometric series. This applies also at
division primes and proves the marking assertion at three directly.
All these frames are symplectic similitudes, since right multiplication
scales $E_0$ by reduced norm. No root of unity is trivialized.

Deligne's frames point out of the Tate module, whereas our markings
point into torsion, so invert the frames in
\cite[\S\S4.11--4.12]{Deligne1971}.
Explicitly put $a=g^{-1}$ when $\widehat I=\widehat{\mathcal O}g$.
Changing generator and rational frame gives
$g\mapsto ugb^{-1}$ and
$(z,a)\mapsto(bz,bau^{-1})$.
Thus left relabeling $u\cdot\alpha=\alpha\circ r_u$ is the canonical
right action by $u^{-1}$. We obtain the whole double quotient
\[
 B^*\backslash\bigl(\mathbb H^\pm\times B(\mathbb A_f)^*\bigr)/K
\]
for $K=K_{35}$ or $K_3$. This includes all rational frames, lattices
and determinant components. The explicit homology transitions and
\cite[\S4.15]{Deligne1971} identify the analytic families and their
markings, not just their geometric points.

To identify the same model over $\Q$, embed this PEL group in
$\operatorname{GSp}(V,E_0)$.
Proposition 1.15 of \cite{Deligne1971} gives a compact open $K'$
containing $K_{35}$ for which the PEL-to-Siegel map is a closed
immersion. Intersect $K'$ with the integral stabilizer of $\mathcal O$
and principal Siegel level three, both containing $K_{35}$.
This shrinking preserves that conclusion: the equivalence relations
in the proof decrease and the maps at the sufficiently small levels
are finite and unramified. This Siegel level is fine and has its
$\Q$-model from full principal level $N\ge3$ and a finite level
quotient \cite[\S\S4.16--4.17]{Deligne1971}.

There is an actual $\Q$-morphism from the constructed $X_{35}$ to
this Siegel model: forget the $\mathcal O$-action and take the
$K'$-orbit of its frame. Choose $N$ divisible by fifteen with
principal Siegel level $N$ contained in $K'$. Refine the
$\mathcal O$-marking to level $N$ on its finite etale surjective
marking cover. The lattice calculation supplies geometric nonemptiness.
The resulting $K'$-orbit is independent of refinement, since two
refinements agreeing at both three and five differ by $K_{35}\subset K'$.
It therefore descends over $\Q$. The primitive integral lattice and
the degree-one effective polarization in Deligne's construction are
the same ones used here, rather than an isogenous or scaled replacement.

Over $\mathbb C$ this morphism is the closed immersion just described.
Its PEL image is a canonical closed subscheme $Z$ over the reflex
field $\Q$
\cite[Theorem 4.21, Corollary 5.7 and Example 5.8]{Deligne1971};
the homological convention agrees with
\cite[2.2.5 and 2.3.1]{Deligne1979}.
The pullback of the ideal of $Z$ vanishes after extension to $\mathbb C$,
hence already over $\Q$. The morphism factors through $Z$ and becomes
an isomorphism over $\mathbb C$, so faithful flatness makes it an isomorphism
over $\Q$. This is the actual map to a fine moduli scheme, not descent
of an arbitrary complex isomorphism.

The projection at five gives
$K_3/K_{35}=(\mathcal O/5\mathcal O)^*$, by the local unit-lifting
calculation and independence of the factors at three and five.
The relabeling action and the canonical finite-adelic action agree
over $\mathbb C$ by the explicit $u^{-1}$ conversion. Both are $\Q$-defined,
so they agree over $\Q$. Their quotients agree, and the already
constructed $Y=X_{35}/G_5$ is the canonical level-$K_3$ curve.
The auxiliary five-marking contributes no factor to its degree.
The central $-1$ changes a full-three marking to its negative,
which is isomorphic by $[-1]$. On $X_{35}$ it changes both markings;
changing only the five-marking is not a stabilizer. This agrees with
the freeness and center calculation above.

Finally forget the three-marking. The unlevelled polarized curve is
connected: Buzzard's $V_1(5)$ curve is geometrically connected and
surjects onto this moduli problem, since every geometric object admits
a five-marking by the split matrix-module calculation.
Its distinguished analytic component $\mathbb H/\Gamma$ is therefore
the whole curve, with the same family $\mathbb C^2/(\mathcal Ov_z)$ and
principal form $E_0$. The fixed A001 application of
\cite[Theorem 2.1]{babagranath} identifies its polarized coarse
forgetful relation with the full positive-normalizer group $W$.
The four arithmetic parity classes above are exactly these classes:
$i,j$ normalize $\mathcal O$, $1+i$ is a unit of norm $-1$, and
$i(1+i)=2+i$ and $j$ have positive norms two and three.
They give the independent local parity classes, with no omitted sign.
The $\Q$-defined coarse model and its already fixed marked coordinate
therefore give the same map $Y\to\mathbf P^1_j$. This is a composition
of the actual $\Q$-defined forgetful maps, not an unspecified
fractional-linear change of $j$.
The normalizer index calculation now applies to this identified
$Y_3=Y_\Q$ and this same $j$, and gives the stated bound $324$.
The full generic connection and its source hypotheses are included
in Section~\ref{v15:marking-descent}.

Each geometric component of $Y_3$ is a torsion-free cover of the
norm-one orbifold of Euler characteristic $-1/3$, and hence has genus
at least two. The normalizer action extends to its smooth proper
model by uniqueness of the stable model, including permutations of
components. Fix $p>324$ and work after strict henselization.
The action is tame. On each component use its effective action.
A finite-order automorphism of order prime to $p$ cannot specialize
to the identity, by considering its first nonzero term modulo a
power of the uniformizer. At a fixed point the same power-series
argument shows that the tangent character of the stabilizer is
faithful. Thus the stabilizer is cyclic.

For a nonidentity generator the tangent eigenvalue minus one is a
unit. Its fixed-point equation therefore has a unique lift by
Hensel's lemma. All powers of the generator fix that lift. It follows
that no new branch orbit appears and that distinct branch orbits
cannot collide: their stabilizers at a purported common specialization
would belong to one cyclic group, whose nontrivial elements have the
same unique fixed lift. Averaging a parameter gives a local invariant
ring $R[[z^e]]$. Hence the quotient is smooth and retains the three
disjoint branch sections and their indices $4,6,2$.

Its genus-zero model with these three ordered sections is isomorphic
to $\mathbf P^1_{\mathbb Z_p}$ carrying the ordered triple
$(0,\infty,-16/27)$, which is disjoint since $2$ and $3$ are units.
The isomorphism agrees with $j$ on the generic fiber, by uniqueness
of a projective coordinate on three points. Thus it is the integral
model in the stated $j$-coordinate; no further denominator prime
has been discarded.

Pull back by $J(u)=-57600F(u)^2/G(u)^6$. If
\[
 H(u)=(u^2-60u-1800)(13u^2+480u+3600)(11u^2+600u+7200),
\]
direct calculation gives $G^6-97200F^2=H^2$. The pullback indices are
$2,6,2$. Tame normalization replaces $e$ by $e/\gcd(e,a)$ and leaves
$2,1,1$. The discriminants and resultants have prime support at most 13,
so this persists at $p>324$. The normalized cover of $\mathbf P^1_u$
has degree at most 324, not $12\cdot324$, and is branched only at the
six contacts. Locally it has equation $u-c=z^2$ after unramified
extension. Positive even contact valuation is unramified; negative
valuation specializes to the unbranched point at infinity. This is
exactly the support specified by $M_S(u)$.

Let $K$ be the splitting field of the fiber over $u$ of the normalized
cover. The fiber is finite \'{e}tale of degree at most $324$, so $K/\mathbb Q$
is normal with $[K:\mathbb Q]\le F_0=324!$. The preceding local
calculation makes it unramified outside
$T=S\cup\operatorname{supp}M_S(u)$. Choose a fiber $A/K$.
The proper model and its universal abelian scheme give good reduction
outside the places over $S$.

All conjugates of $A$ have the same geometric polarized moduli point.
Indeed the polarized equivalence relation in
\cite[Theorem~2.1]{babagranath}, with the corrected principal datum above,
gives $V_{\mathbb Z}=V/W$; its two twisting involutions generate all of
$W$ for discriminant six. On the finite nonzero $j$-locus,
\eqref{eq:qm-normalized-igusa} specifies the same point by its weighted
Igusa invariants. Here $j=J(u)$ is rational, finite and nonzero:
$F(u)\ne0$ and $\operatorname{disc}(G)=-10800<0$.
The exact coordinate connection is \cite{ly-d6-normalization-project}.

To define the Hom field without choosing a Galois action on a
labelled collection, form $W_0=\operatorname{Res}_{K/\mathbb Q}A$,
of dimension $g=2[K:\mathbb Q]$.
Write $W_0(\mathbb C)=\mathbb C^g/\Lambda$.
By \cite[Proposition~6.5, pp.~318--319]{milne-shimura2005},
every geometric endomorphism lifts uniquely to a complex-linear
map preserving $\Lambda$. Since $\Lambda$ spans $\mathbb C^g$
over $\mathbb R$, its action on $\Lambda$ is faithful.
Thus the geometric endomorphism group is a subgroup of
$\operatorname{End}_{\mathbb Z}(\Lambda)$, hence is free of rank
\[
 r\le(2g)^2=16[K:\mathbb Q]^2\le16F_0^2.
\]
Each member of a finite integral basis is defined over a finite
extension, so the Galois image on this lattice is finite.
Every finite subgroup of $\operatorname{GL}_r(\mathbb Z)$ injects
modulo three. Indeed, a nontrivial kernel would contain an element
$M$ of prime order $q$. Write $M=I+3^kB$, with $k\ge1$ and
$B\not\equiv0\pmod3$. If $q\ne3$, expansion modulo $3^{k+1}$
gives $M^q-I\equiv q3^kB\not\equiv0$.
If $q=3$, expansion modulo $3^{k+2}$ gives
$M^3-I\equiv3^{k+1}B\not\equiv0$, again a contradiction.
If $E$ is the kernel field of the Galois action and $L=EK$,
then $L/\mathbb Q$ is normal and
\[
 [L:\mathbb Q]\le F_0\,3^{r^2}
 \le F_0\,3^{256F_0^4}=D.
\]
At a prime $p$ outside $T$, choose $\ell\ne p$.
The field $K$ is unramified at $p$, and $A$ has good reduction
at every place above $p$. The residue fields are finite and hence
perfect, so \cite[Theorem~1, p.~493]{serre-tate-good-reduction}
makes inertia trivial on the Tate module at every such place.
The decomposition
$K\otimes_{\mathbb Q}\overline{\mathbb Q}
 \simeq\prod_{\sigma:K\hookrightarrow\overline{\mathbb Q}}
 \overline{\mathbb Q}$
identifies $W_{0,\overline{\mathbb Q}}$ with the product of the
conjugates of $A$, with Galois acting by transport of factors.
Thus
$V_\ell(W_0)=\operatorname{Ind}_{G_K}^{G_{\mathbb Q}}V_\ell(A)$.
Unramifiedness of $K$ makes inertia fix the coset indices,
so it acts trivially on this entire Tate representation.
The action of geometric endomorphisms on $V_\ell(W_0)$ is faithful.
Indeed, after fixing $\overline{\mathbb Q}\hookrightarrow\mathbb C$,
the finite torsion schemes have the same geometric points, and
\cite[equation~(41), p.~334]{milne-shimura2005} gives
$T_\ell(W_0)\simeq\Lambda\otimes\mathbb Z_\ell$
compatibly with endomorphisms.
An endomorphism with zero induced map therefore has zero
integral lattice matrix, and hence is zero.
This realization is Galois equivariant because it is defined on
torsion points. Inertia consequently fixes the full endomorphism
lattice by conjugation. Thus $E$, and hence $L$, is unramified
outside $T$.
Over $K$ the conjugate factors of $W_0$ are individually defined.
Their inclusions and projections identify its geometric
endomorphism lattice with the direct sum of all ordered Hom
groups between these factors. The group $G_E$ fixes every
endomorphism and $G_K$ fixes the factor labels.
Their intersection $G_L$ therefore fixes every such morphism.
Since its source and target are defined over $L$, Galois descent
for morphisms defines it over $L$. In particular $L$ defines
every geometric endomorphism and every ordered conjugacy
homomorphism of $A$.

By \cite[Theorems 6.7--6.8 and Corollary 6.9, p.~319]{milne-shimura2005}, polarizable complex
tori are abelian varieties and their analytic homomorphisms are algebraic.
After clearing lattice denominators, the same assertion identifies
rational Hodge endomorphisms with algebraic quasi-endomorphisms.
These endomorphisms are defined over $\overline{\Q}$: every automorphism
of $\mathbb C/\overline{\Q}$ fixes their action on all torsion points, and
torsion is Zariski dense.
On rational homology the complex structure is right multiplication
by an element $J\in B\otimes\mathbb R$ with $J^2=-1$. In
$\operatorname{End}_{\mathbb Q}(B)=B\otimes B^{\mathrm{op}}$,
commutation with $J$ requires each right coefficient to commute
with $J$. Its rational centralizer is either $\mathbb Q$ or an
imaginary quadratic field $k$. Hence the geometric endomorphism
algebra is either $B$ or $M_2(k)$.
We treat these two cases separately. Throughout, put
$M=M_S(u)$, $d=[L:\Q]\le D$, and $Q=Q_R$.

\paragraph{CM fibers.}
Suppose $\operatorname{End}^0(A_{\overline\Q})=M_2(k)$.
All its elements are defined over $L$. Multiples of matrix units
therefore give an elliptic curve $E/L$ with complex multiplication
and an $L$-isogeny $E^2\to A$. Let $\Delta<0$ be the discriminant
of the full geometric endomorphism order of $E$. The degree of
$j(E)$ is its order class number, so $h(\Delta)\le d\le D$.
The effective bound of \cite[Proposition 2.7]{fowler} gives
\[
 \sqrt{|\Delta|}\le2D^2\exp(21000D^2).
\]
This bound includes every imaginary quadratic order and has no
exceptional-field exclusion.

Here is the height comparison with its normalization explicit.
After a finite extension $E$ has everywhere good reduction, since
its $j$-invariant is integral. This extension does not change its
stable Faltings height. For a reduced CM period $\tau=x+iy$ we have
$\sqrt3/2\le y\le\sqrt{|\Delta|}/2$. The archimedean summand in
the original Faltings normalization is
\[
 -\log(2\pi)+\frac{\pi y}{6}-\frac12\log y
       -2\sum_{n\ge1}\log|1-e^{2\pi in\tau}|,
\]
by the discriminant formula \cite[Eq. (6.2)]{cmheight}.
For $\rho=e^{-2\pi y}<1/100$, the final sum is at most
$2\rho/(1-\rho)^2<200/9801$, and
$-\tfrac12\log y<1/10$. Averaging and discarding the negative
$-\log(2\pi)$ gives
\[
 h_F(E)\le\frac{\pi\sqrt{|\Delta|}}{12}+1
       \le D^2 3^{21000D^2}+1=B_{\rm cm}.
\]
Apply the published isogeny estimate
\cite[Theorem 1.4]{gr2014} over the original field $L$, with
source $E^2$. It gives an isogeny to $A$ of degree at most
\[
 \bigl(28^{256}d\max(h_F(E^2),\log d,1)^2\bigr)^{8192}.
\]
The maximum is at most $2B_{\rm cm}$. Height variation under
isogeny \cite{faltings1983} and additivity give
\begin{align*}
 h_F(A)&\le2B_{\rm cm}
   +4096\bigl(256\log28+\log d+2\log(2B_{\rm cm})\bigr)\\
 &\le2B_{\rm cm}+4096(1028+5D+84000D^2)=C_{\rm cm}.
\end{align*}
For the last line use $\log28<4$, $\log d\le D$, and
$\log(2B_{\rm cm})\le2+2D+42000D^2$.
No bound for the degree of the auxiliary good-reduction extension
enters the isogeny estimate.

\paragraph{Non-CM fibers and scalar descent.}
Now suppose $\operatorname{End}^0(A_{\overline\Q})=B$.
In particular $A$ is absolutely simple. For each
$\sigma\in\operatorname{Gal}(L/\Q)$ choose a geometric isomorphism
$\psi_\sigma:{}^\sigma A\to A$. The automorphism
$b\mapsto\psi_\sigma{}^\sigma b\psi_\sigma^{-1}$ of $B$ is inner.
Correcting it by Skolem--Noether and clearing denominators gives
a $B$-linear isogeny $\mu_\sigma:{}^\sigma A\to A$, defined over
$L$ because all the relevant Hom groups are defined there. Take
$\mu_1=1$ and inflate this finite choice to $G_\Q$.
The discrepancy
\[
 c(\sigma,\tau)=
 \mu_\sigma{}^\sigma\mu_\tau\mu_{\sigma\tau}^{-1}
\]
commutes with $B$, hence belongs to $\Q^\times$.
Set $a_\sigma=(\deg\mu_\sigma)^{-1/4}>0$ and
$\nu_\sigma=a_\sigma\mu_\sigma$ in the category with algebraic
coefficients. Throughout these coefficient extensions Galois acts on
the varieties and morphisms and acts trivially on the added scalars;
the induced homology maps are extended linearly. In particular, the
$a_\sigma$ are constant coefficients in the products below.
Here degree extends to rational quasi-isogenies by
\(D(u)=\deg([m]u)/m^4\), for any positive integer \(m\)
clearing denominators. This is independent of \(m\): comparison
after multiplication by a second denominator uses
\(\deg[n]=n^4\) and multiplicativity~\cite[Theorem~8.2, p.~115]{milne1986-abelian}. It is multiplicative and
invariant under field conjugation. Writing \(c=r/s\in\mathbf Q^\times\)
with \(s>0\), the equality \([s](cu)=[r]u\) gives
\(D(cu)=c^4D(u)\). Thus
\[
 D(\mu_\sigma)D(\mu_\tau)
 =c(\sigma,\tau)^4D(\mu_{\sigma\tau}),\qquad
 \frac{a_\sigma a_\tau}{a_{\sigma\tau}}
 =\frac1{|c(\sigma,\tau)|}.
\]
The positive fourth-root convention therefore gives
\(\varepsilon(\sigma,\tau)=c(\sigma,\tau)/|c(\sigma,\tau)|\).
Consequently
\[
 \nu_\sigma{}^\sigma\nu_\tau
       =\varepsilon(\sigma,\tau)\nu_{\sigma\tau},
 \qquad \varepsilon(\sigma,\tau)\in\{1,-1\}.
\]
Thus $\varepsilon$ is a normalized sign cocycle inflated from
$\operatorname{Gal}(L/\Q)$. Only finitely many algebraic
coefficients $a_\sigma$ occur.

We need a splitting with controlled ramification, not merely
existence of a projective lift. Identify
$H^2(G_\Q,\{1,-1\})$ with $\operatorname{Br}(\Q)[2]$.
At a finite prime outside $T$ the class of $\varepsilon$ is zero:
its restriction is unramified and the unramified local quotient
has cohomological dimension one. Let $T_-$ be the finite primes
where its invariant is $1/2$. For each odd $p\in T_-$ choose an
odd character of $(\Z/p\Z)^\times$; at $2\in T_-$ use the
nontrivial character modulo four. Their product is a Dirichlet
character $\chi$, unramified outside $T$, whose order $m$ satisfies
\[
 m\le2\prod_{\substack{p\in T\\p\ne2}}(p-1)
       \le2R^R M.
\]
Here the product over an empty set is one.

To check that $\chi$ has the required square-root obstruction,
use the exact sequence of constant discrete Galois modules
$1\to\{\pm1\}\to\mu_\infty\xrightarrow{z\mapsto z^2}
\mu_\infty\to1$. By local reciprocity, a finite-order character
of $G_{\Q_p}$ has a character square root exactly when its value
on $-1\in\Q_p^\times$ is $1$. Indeed this is the only obstruction
on the torsion part of $\Q_p^\times$; the valuation factor and
the procyclic principal-unit factors admit such roots.
Hence the boundary of $\chi$ and $[\varepsilon]$ agree at every finite
place: the local two-torsion Brauer group has order two, and the
square-root criterion distinguishes its two classes. The global
Brauer exact sequence determines the real invariant from the finite
ones and then identifies the two global classes
\cite[I.1.6(b) and proof of II.2.1, p.~164]{milne}.
Choosing a continuous pointwise square root of $\chi$ and
correcting it by a sign cochain consequently gives a continuous
cochain $\beta:G_\Q\to\mu_{2m}$ such that
\[
 \beta^2=\chi,\qquad
 \beta(\sigma)\beta(\tau)\beta(\sigma\tau)^{-1}
        =\varepsilon(\sigma,\tau).
\]
Its restriction to $G_L$ is a character invariant under $G_\Q$.
Outside $T$, its inertia restrictions are quadratic and nontrivial
at only finitely many primes. Twisting by the quadratic character
of the product of those odd primes kills all these restrictions:
at an odd prime there is a unique nontrivial quadratic character
on inertia. This introduces ramification only at those primes
and possibly at $2\in T$, and changes neither displayed identity.
We henceforth use this choice of $\beta$.

Let $L'$ be the kernel field of $\beta|_{G_L}$.
Conjugacy invariance shows that $L'/\Q$ is Galois, and
\[
 [L':\Q]\le2md\le4DR^R M\le QM.
\]
The last inequality follows from $R\ge\max(D,10^{12})$ and
$E_0=10(DG_0)^2$.

\paragraph{Modularity and a uniform conductor bound.}
Choose one number field $C$ containing all $a_\sigma$, the values of
$\beta$, and a splitting field of $B$. For each coefficient prime
$\ell$ choose an embedding of $C$ into $\overline{\mathbb Q}_\ell$
and extend scalars on $V_\ell(A)$ to a finite extension containing
its image. With the constant-coefficient convention above, define
\[
 \rho_\ell(\sigma)
   =\beta(\sigma)^{-1}a_\sigma\mu_\sigma\circ\sigma.
\]
The cocycle identities make this a representation of $G_\Q$.
It commutes with $B$, so matrix units give
$\rho_\ell\simeq r_\ell\oplus r_\ell$ with $\dim r_\ell=2$.
On $G_{L'}$ it is the actual Tate action of $A$.
Faltings' semisimplicity and endomorphism theorems
\cite{faltings1983} show that $r_\ell|_{G_{L'}}$ is absolutely
irreducible: after splitting $B$, its commutant is the scalars.
It is unramified outside $T\cup\{\ell\}$, since there $L$ is
unramified, $A$ has good reduction, and $\beta$ is unramified.

Apply Lemma~\ref{lem:qm-modularity-descent} with
$t_\sigma=a_\sigma/\beta(\sigma)$. It identifies this same $r_5$
with the Tate realization of a primitive weight-two newform $f$.
The qualitative modularity and finite-twist argument does not replace
$A$, $L'$, $\beta$, or $r_\ell$ in the ramification and height bounds below.
For this final newform, Li's Fricke formula
\cite[p.~296, before Lemma 6]{li-newforms} makes the span of its
coefficient conjugates Fricke-stable. Shimura's quotient construction
and Euler-factor identity \cite[Theorem 1]{shimura-factors} therefore
show that $r_5$ is a constituent of the Tate module of the modular abelian variety
$A_f$, or equivalently $r_5^\vee$ is its cohomological constituent.
This duality preserves the conductor.

The representations at all coefficient primes belong to this same system.
Enlarge $C$ to contain the coefficients of $f$, extending the original
embeddings and retaining the chosen embedding at five. Fix a coefficient
prime $\ell$. Work over a common finite extension of $\mathbb Q_\ell$ containing both coefficient images. For $p\notin T\cup\{5,\ell\}$ away also from the level of $f$,
choose a place $v$ of $L$ and an arithmetic Frobenius lift
$\sigma\in G_{\mathbb Q_p}$ whose residue action is $x\mapsto x^p$.
Since $L_v/\mathbb Q_p$ is unramified and $A$ has good reduction,
relative $p$-Frobenius maps $\overline A_v$ to the reduction of
${}^\sigma A$ at $v$. The actual isogeny $\mu_\sigma$ extends over
 the local abelian schemes. Composing its reduction with this Frobenius
 gives an endomorphism $u_p$ of $\overline A_v$.
Its rational Tate trace $t_p$ is independent of the coefficient prime
 different from $p$. Functorial specialization and the two-copy
 decomposition give
\[
 \operatorname{tr}r_\ell(\sigma)
 =\iota_\ell\bigl(a_\sigma\beta(\sigma)^{-1}t_p/2\bigr).
\]
At five the algebraic number in parentheses equals $a_p(f)$.
Injectivity of the embedding at five makes this an equality in $C$,
so transport to $\ell$ gives the trace of the corresponding classical
Tate realization $U_{f,\ell}$. Chebotarev and the characteristic-zero
semisimple trace criterion identify $r_\ell$ with $U_{f,\ell}^{\rm ss}$.
Since $r_\ell$ is absolutely irreducible, this semisimplification has
one simple composition factor; hence $U_{f,\ell}$ itself is irreducible
and isomorphic to $r_\ell$.
For each residue prime $p$ choose $\ell\ne p$. Restricting this
isomorphism to the decomposition group at $p$ identifies the local
Galois representations and their Artin conductors. The local conductor
comparison for the primitive newform then identifies this exponent
with $v_p(N_f)$~\cite[\S\S2.1, 2.4, 4.1]{carayol1987}.

Here are uniform local bounds, including the small primes.
For a finite extension of $\Q_p$ of degree at most $a$, its
normal closure has degree at most $a!$. For a finite Galois
local extension with ramification index $e$ and different
exponent $\delta$, let $f$ be the Eisenstein polynomial of a
uniformizer $\pi$ over the maximal unramified subfield. Then
$\delta\le e-1+e v_p(e)$. The derivative terms have distinct
valuations modulo $e$, so there is no cancellation at the
smallest valuation. The identity
$f'(\pi)=\prod_{\sigma\ne1}(\pi-\sigma\pi)$ over the inertia
conjugates gives $\delta=\sum_{i\ge0}(|G_i|-1)$. Changing variables
by the Herbrand function yields
$\delta/e=1-1/e+\int_0^\infty(1-|G^u|^{-1})\,du$.
Before the last positive upper break the integrand is at least
$1/2$, so that break is bounded by
$2\delta/e\le2(1+\log_2(a!))\le4a^2$
\cite[Sections 4.1--4.3]{serrelocal}.

At $p$, take $F=L_v$ and adjoin $A[3]$, using $A[5]$ if
$p=3$. This extension has degree at most $DG_0$ over $\Q_p$.
Raynaud's criterion \cite[Theorem 3.5]{silverbergzarhin}
gives semistable reduction there. Its wild inertia acts
trivially on each $V_\ell(A)$ with $\ell\ne p$.
The preceding break bound and the subgroup rule for upper
numbering therefore show that all absolute ramification
groups above $U_0=4(DG_0)^2$ act trivially on that Tate module.
On these groups $a_\sigma\mu_\sigma=1$ because they lie in $G_L$.
For odd $p$, the character $\beta|_{G_F}$ is trivial on wild
inertia: its square is the tame character $\chi|_{G_F}$, and a
pro-$p$ group has no quotient of order two. At $p=2$, local
reciprocity and the norm map give $\beta^4=1$ on the units of
$F$, since the $2$-primary local factor of $\chi$ has conductor
at most four. If $e=v_F(2)\le D$, the logarithm for $\widehat{\mathbf G}_m$ identifies
$U_F^{e+1}$ with $\mathfrak m_F^{e+1}$ since $e+1>e$
\cite[IV.6.4, p.~132]{silverman}. Multiplication by four gives
$(U_F^{e+1})^4=U_F^{3e+1}$. Thus the relative upper break of
$\beta$ is at most $3e$. To make the change of numbering explicit,
the same normal-closure estimate bounds the last upper break of
$F/\mathbb Q_2$ by $4D^2$. Above it, the absolute upper group
$G_{\mathbb Q_2}^{u}$ lies in $G_F$, and the subgroup formula is
\[
 G_{\mathbb Q_2}^{u}\cap G_F
       =G_F^{\psi_{F/\mathbb Q_2}(u)},
 \qquad \psi_{F/\mathbb Q_2}(u)\ge u\quad(u\ge0).
\]
Consequently $\beta$ is trivial on every absolute upper group with
$u>\max(4D^2,3e)$, which is at most $U_0$.
This proves the asserted absolute bound without a factor involving
$[L':\mathbb Q]$.
The Artin conductor exponent of $r_\ell$ is consequently
at most $2+2U_0\le E_0$.

For $p>R$, the extension $L_v/\Q_p$ is tame, since its
degree is at most $D<p$. The surface has good reduction,
and the same argument kills the wild part of $\beta$.
Hence the conductor exponent is at most two there.
Outside $T$ it is zero. Local--global compatibility for the
weight-two modular representation therefore bounds the
primitive level $N$ of $f$ by
\[
 N\le\Bigl(\prod_{p\le R}p^{E_0}\Bigr)M^2
      \le R^{RE_0}M^2=QM^2.
\]


\paragraph{Transfer from a modular Jacobian.}
Over $L'$ the Tate modules of $A$ and $J=J_1(N)$ have a common
constituent. Invariants in their finite-dimensional Hom space
commute with extension of the coefficient field. Faltings'
Hom theorem thus gives a nonzero $L'$-homomorphism $A\to J$.
Let $B_0\subseteq J$ be the image of this homomorphism.
Absolute simplicity implies that $A\to B_0$ is an $L'$-isogeny.
Apply \cite[Proposition~12.1, p.~122]{milne1986-abelian} over $L'$:
for an ample line bundle $\mathcal H$ on $J$, the reduced identity
component $Z$ of the kernel of
$J\xrightarrow{\phi_{\mathcal H}}J^\vee\to B_0^\vee$
is a complementary abelian subvariety over $L'$.
Indeed, its intersection with $B_0$ is finite because the
restriction is the polarization associated with
$\mathcal H|_{B_0}$, and $\dim Z\ge\dim J-\dim B_0$.
Thus $B_0\times Z\to J$, and hence $A\times Z\to J$,
are $L'$-isogenies. In particular $g=\dim J\ge2$.
The index of $\Gamma_1(N)$ is at most $N^2$, so both the
genus and the degree of its standard Belyi map are at most
$N^2$. By \cite[Theorem 1.1.1 and Section 2.3]{javanpeykar},
\[
 h_F(J)\le13000000N^{12}\le Q^{13}M^{24},\qquad
 g\le Q^3M^4,\qquad d'=[L':\Q]\le QM.
\]
The normalization in that theorem is the original stable
Faltings height of the Jacobian.

Apply \cite[Theorem 1.4]{gr2014} with source $J$ and target
$A\times Z$, both over $L'$. There is such an isogeny $\varphi$
with
\[
 \log\deg\varphi\le1024g^3
 \bigl(64g^2\log(14g)+\log d'
    +2\log\max(h_F(J),\log d',1)\bigr).
\]
For $Q\ge10^{12}$ and $M\ge1$, the three logarithms in order
are at most $5QM^4$, $2QM^4$, and $19QM^4$.
These follow from $\log Q\le Q$ and $\log M\le M^4/4$.
It follows that
\[
 \log\deg\varphi\le368640Q^{16}M^{24}
                       \le Q^{17}M^{24}.
\]
The general lower bound $h_F(Z)\ge-\tfrac32\dim Z\ge-2g$
is recorded in \cite[Section 9]{gr2014}.
Additivity and height variation now give
\[
 h_F(A)\le Q^{13}M^{24}+\tfrac12Q^{17}M^{24}+2Q^3M^4
          \le Q^{17}M^{24}.
\]
Thus both cases prove the exact estimate needed below:
\[
 \max(1,h_F(A))
      \le\max(1,C_{\rm cm},Q_R^{17})M_S(u)^{24}.
\]


Use the same Hauptmodul, normalized by the values $0,\infty,-16/27$
at the CM points of discriminants $-4,-3,-24$.
The calculation leading to \eqref{eq:qm-normalized-igusa}, using
\cite[equations~(13)--(14), Examples~29--30, pp.~793--797]{linyang}
and \cite{ly-d6-normalization-project}, gives in one common
weighted-projective representative
\[
 J_2=12(j+1),\quad J_4=6(j^2+j+1),\quad J_{10}=j^3,
 \qquad \psi_4=J_2^2-24J_4=144j,\quad \chi_{10}=-j^3/4.
\]
Write $z_m=\theta_m(\tau,0)$ for the same principally polarized surface,
where $m=(a/2,b/2)$ with $a,b\in\{0,1\}^2$; the characteristic is even
when $a^tb=0\pmod2$. Igusa's theta series and identities
\cite[pp.~833--834 and 848]{igusa1967} give
$\psi_4=2^{-2}\sum_{m\ {\rm even}}z_m^8$ and
$\chi_{10}=-2^{-14}\prod_{m\ {\rm even}}z_m^2$.
These are the half-characteristic coordinates used at $r=2$ in
\cite[\S2.3.4, pp.~30--31]{pazuki}; the exact convention and substitution
are recorded in \cite{igusa-theta-project}.
The quotient $\chi_{10}^2/\psi_4^5$ is unchanged by weighted
rescaling: multiplying a weight-$k$ coordinate by $s^k$ multiplies
both numerator and denominator by $s^{20}$. In particular, the
weight-dependent normalization in \cite[\S5, p.~792]{linyang}
cancels in this quotient. Since our $j$ is finite and nonzero, its denominator
is nonzero, and substitution gives
\[
 j=16\cdot144^5\frac{\chi_{10}^2}{\psi_4^5}
  =3779136\frac{(\prod_{m\ {\rm even}}z_m)^4}
                 {(\sum_{m\ {\rm even}}z_m^8)^5}.
\]
The two homogeneous polynomials have degree $40$ and coefficient
$\ell^1$ norms $3779136$ and $10^5$. Hence
$h(j)\le40h([z_m])+22$.
The product and sum use the ten even theta nulls of the same
principally polarized surface. Adjoining the six zero odd
coordinates does not change their projective height.

In \cite{pazuki}, the level $r=2$ theta coordinates are the sixteen
half-characteristic values used here; the six odd values vanish.
The permitted changes of theta structure permute coordinates and
multiply them by roots of unity. His Euclidean projective height
dominates our supremum projective height. Put
$X=\max(1,h_F(A))$ and
$X_D=\max(1,h_F(A)+\log\pi)\le3X$.
His Corollary~1.3(2), in the Deligne metric, gives
\[
 h_\Theta(A)\le\tfrac12 X_D+C_2\log(X_D+2),
 \qquad C_2=1000\cdot2^4(\log(2^4))^5\le1000\cdot2^{24}.
\]
Using $\log(X_D+2)\le2X_D$ proves
$h([z_m])\le10000\cdot2^{24}X$.

Write $j=a/b$ in lowest terms. The nonzero integer polynomial
$-57600bF(X)^2-aG(X)^6$ has $u$ as a root and coefficient norm at
most $2\cdot4000^6\max(|a|,|b|)$. Remove any initial power of $X$
and use the actual highest nonzero coefficient. The rational-root
theorem bounds both numerator and denominator of $u$, giving
$h(u)\le h(j)+\log(2\cdot4000^6)<h(j)+74$.

A field extension for the theta structure changes neither absolute
height. Since $40\cdot10000\cdot2^{24}+96<10^{13}$, the preceding
inequalities give
\[
 h(u)\le10^{13}\max(1,h_F(A))
 \le10^{13}\max(1,C_{\rm cm},Q_R^{17})M_S(u)^{24}
 =\tfrac12 H_6(S;R)M_S(u)^{24}.
\]
Every estimate used only the displayed lower bound on $R$, so the
conclusion holds for every permitted larger $R$ as well.
\end{proof}

The final value argument uses $S_{\rm ht}=\{p:p\le324\}$. Multiplication
uses $S_{\rm mul}=S_{\rm ht}\cup\mathcal H\cup\operatorname{supp}(c_*)$,
where $\mathcal H$ is the fixed catalogue. A common numerical $R$ may
bound both sets. It does not identify them.

\section{The pole predicate and its converse input}\label{sec:pole}
The pole construction and simultaneous-descent strategy adapt
\cite[Section 3]{openaihtp}.

Let $S_0=\{2,3,5,7,13\}$. With input $B$ and eleven witnesses
$b_0,\ldots,b_5,e,z,r_1,r_2,w$, define
\begin{align*}
 m&=(1-2b_0)(1+b_1)(1+2b_2)(1+4b_3)(1+6b_4)(1+12b_5),\\
 A&=mB,\qquad H=A(A-B^2),\\
 T_1&=A(A-e)(A+3e),\qquad T_2=A(Az^2-eH)(Az^2+3eH),\\
 Q_{\rm pole}&=B\Bigl(\sum_{i=0}^5(b_i^2-b_i)^2+(eHzw-1)^2\\
 &\hspace{4em}+(r_1^2-T_1)^2+(r_2^2-T_2)^2\Bigr).
\end{align*}
Write $\Phi(B)$ for rational solubility. For $B\ne0$, $m$ is a signed
divisor of 2730, and the equations give points $(A,r_1)\in E_e$ and
$(Az^2,zr_2)\in E_{eH}$ with $eHz\ne0$, where
$E_\ell:y^2=x(x-\ell)(x+3\ell)$. The polynomial has total degree 75.

\begin{lemma}[One-sided obstruction]\label{lem:pole}
If $B\ne0$ and $\Phi(B)$, then $v_p(B)$ is not negative and odd for
any $p\notin S_0$.
\end{lemma}
\begin{proof}
If $k=v_p(B)<0$ were odd, then $v_p(A)=k$, $v_p(H)=3k$.
Exactly one of $e,eH$ has even valuation. Square scaling makes its
parameter a unit, preserving the parity of the abscissa's valuation.
For a unit parameter the cubic has valuation $v_p(x)$ if $v_p(x)>0$
and $3v_p(x)$ if $v_p(x)<0$. Its being a square forces even valuation.
But both points have abscissa in the squareclass $A$, a contradiction.
\end{proof}

\begin{theorem}[Pole completeness, $\mathsf{PC}$]\label{hyp:PC}
For every rational $B$ with square reduced denominator,
$0<B<1/2$ and $v_p(B)\ge4$ for every $p\in S_0$ imply $\Phi(B)$.
\end{theorem}
The following statement is the $E=E_\ell$, $s_2(E)=1$ case of
\cite[Theorem 1.1]{openaipointwise}, and also follows as a special
case of \cite[Theorem 1.1]{openaigoldfeld}. We give the restricted
argument here; these preprints are cited for attribution, not assumed
as additional converse inputs.

\begin{theorem}[Restricted converse, $\mathsf{RC}$]\label{hyp:RC}
For every positive squarefree $\ell\equiv11\pmod{12}$,
\[
 \operatorname{corank}_{\Z_2}\Sel_{2^\infty}(E_\ell/\Q)=1
 \quad\Longrightarrow\quad\ord_{s=1}L(E_\ell,s)=1.
\]
\end{theorem}
Here $c_4=208\ell^2$, $\Delta=2304\ell^6$, and $j=35152/9$.
The calculation in \eqref{eq:rc-noncm-invariants} and the rational-integrality
argument following it show that these curves are non-CM, using the monic
integral CM class polynomial of \cite[\S2.1, pp.~612--613]{fowler}.
They are different from the CM generator curve. Theorem~\ref{hyp:RC} is proved in
Sections~\ref{sec:rc}--\ref{sec:limits}, with its coefficient and
normalization arguments in Appendices~\ref{app:even-source}--\ref{app:endpoints}.
We now prove Theorem~\ref{hyp:PC} from this restricted converse.

\subsection{The exact descent reduction}
The signed divisors of 2730 realize all real, $2$-adic and $3$-adic
squareclass signatures. After $-1,2,3$ fix sign and valuation parities,
the remaining unit vectors of $5,7,13$ are
$(0,1,1),(1,0,0),(0,1,0)$ and form a basis. Choose positive $m$ making
$A=mB$ square at two and a nonsquare unit squareclass at three.
Let $\beta,\theta$ be the positive squarefree representatives of $A,H$.
Then $\beta\equiv17\pmod{24}$, $\theta\equiv1\pmod{24}$, with disjoint
supports, $\beta>1$, and $\theta=1$ permitted. Each is a local square
at the support of the other. At a prime of $\beta$, the relevant odd
valuation of $B$ is positive, so $1-B/m$ is a square unit. At a prime
of $\theta$, $v(B)$ must be zero, since otherwise $v(H)$ is even;
there $B/m\equiv1$, so $A$ is a square.

The simultaneous construction in S11--S12 gives $n,n\theta$ positive
squarefree and congruent to eleven modulo twelve, both with two-Selmer
dimension three and with extra Kummer class $(\beta,1)$. Its matrix
conditions can be stated explicitly. For support primes $q_i$, set
$s_i=[-1/q_i]$, $\lambda_i=[3/q_i]$, and
\[
 (Kv)_i=\sum_{j\ne i}[q_j/q_i](v_j+v_i).
\]
One arranges $\ker K=\langle\mathbf1,\mathbf1_{\operatorname{supp}\beta}\rangle$
and $\lambda\notin\operatorname{im}K$ on both supports. Add $d-1$
prime slots of type $(s,\lambda)=(0,0)$ and one slot $q_0$ of type
$(1,1)$, where $d=\omega(\beta)$. The cross-block has columns
$e_i+e_d$. After the rank-one correction $K^\sharp=K+\lambda a^t$,
deleting $q_0$ leaves
\[
 \begin{pmatrix}D_1&C_1\\C_1^t&0\end{pmatrix},
 \qquad\mathbf1^tD_1\mathbf1=1.
\]
This is nonsingular: $C_1^tx=0$ makes $x$ constant, and the first
equation forces that constant and then $y$ to vanish. Adding the
support of $\theta$ preserves this block. Free $q_0$-symbols toggle
the new diagonal entries; the determinant's full-product coefficient
is one. Thus a single Boolean assignment works on both supports.
Reciprocity, CRT and Dirichlet realize it by distinct auxiliary primes.

In Kummer coordinates $(x+3\ell,x-\ell)$, a positive-second-coordinate
Selmer representative is $(3^\delta\prod q_i^{u_i},\prod q_i^{v_i})$.
The necessary local equations are
\[
 Ku+\operatorname{diag}(\lambda)v=\delta\lambda,\quad
 (K+\operatorname{diag}(s))v=0,\quad(\lambda+s)^tv=0.
\]
The final equation is the condition at three. The matrix properties
force $v=0$, $\delta=0$, $u\in\ker K$, so there are at most eight
classes in total. The extra class is actually local at every place,
including two, and independent of the two torsion classes, giving
dimension three. This sufficiency check for the particular class is
essential; arbitrary parity conditions at two would not suffice.

\begin{lemma}[Dimension three gives corank one]\label{lem:corank}
For an elliptic curve with full rational two-torsion,
$\dim_{\F_2}\Sel_2=3$ implies $\operatorname{corank}_{\Z_2}\Sel_{2^\infty}=1$.
\end{lemma}
\begin{proof}
The finite Kummer sequence gives $r+\dim\Sha[2]=1$.
Write $\Sha[2^\infty]=(\Q_2/\Z_2)^t\oplus F$ with $F$ finite.
The displayed divisible summand is the maximal divisible subgroup:
any divisible subgroup maps trivially to the finite quotient $F$.
The Cassels--Tate pairing on that quotient is nondegenerate and
alternating for an elliptic curve
\cite[Corollaries 4, 7 and Theorem 8]{poonenstoll}.
Thus no identification of divisible elements with infinitely
divisible elements in an arbitrary group is required.
Its finite primary cyclic factors are paired,
so $\dim F[2]$ is even. Thus $F=0$ and $r+t=1$. The restriction of
the pairing to $F[2]$ need not be nondegenerate and is not used.
\end{proof}

\begin{lemma}[Forward analytic rank zero and one]\label{lem:forward-rank}
Let $A/\Q$ be an elliptic curve. If $\ord_{s=1}L(A,s)=r\leq1$,
then $\operatorname{rank}A(\Q)=r$ and $\Sha(A/\Q)$ is finite.
Consequently $\operatorname{corank}_{\Z_2}\Sel_{2^\infty}(A/\Q)=r$.
\end{lemma}
\begin{proof}
Use the modularity theorem of Breuil, Conrad, Diamond and Taylor
\cite{bcdt} in the published formulation
\cite[Theorem C.13.6, p.~442]{silverman}.
Apply \cite[Theorem 1, pp.~430--431]{kolyvagin}, whose statement
includes both $r=0$ and $r=1$. Its Theorems 3--5 provide the
auxiliary imaginary quadratic field, the nonzero Heegner height,
and finiteness over that field; the restriction kernel for $\Sha$
is contained in the finite group $\Sha(A/\Q)[2]$.
The ordinary Kummer sequence
\[
0\longrightarrow A(\Q)\otimes\Q_2/\Z_2
\longrightarrow\Sel_{2^\infty}(A/\Q)
\longrightarrow\Sha(A/\Q)[2^\infty]\longrightarrow0
\]
gives the corank assertion.
For the modular quotient used below, let $N$ be the conductor of $A$.
Choose the finite surjective $\Q$-morphism
$\rho:X_0(N)\to A$ supplied by
\cite[Theorem C.13.6]{silverman}, and compose it with the finite
surjective forgetful morphism $u:X_1(N)\to X_0(N)$.
Put $f=\rho\circ u$. Pushforward of divisors preserves degree
and sends principal divisors to principal divisors by the function-field
norm \cite[Proposition II.3.6(c),(d)]{silverman}.
It therefore induces
\[
 q=f_*:J_1(N)=\operatorname{Pic}^0(X_1(N))
       \longrightarrow\operatorname{Pic}^0(A)\simeq A,
\]
where the last identification uses the identity point of $A$
\cite[Proposition III.3.4]{silverman}.
This is a homomorphism of abelian varieties over $\Q$: for the finite
flat map $f$, the norm on line bundles is
$\mathcal L\mapsto\det(f_*\mathcal L)\otimes
\det(f_*\mathcal O_{X_1(N)})^{-1}$;
this construction is defined over $\Q$, commutes with base change,
and induces the displayed map of Picard varieties. Section~\ref{v15:picard-norm} proves finite flatness and the relative-Picard construction used here.
To check surjectivity, work over $\overline{\Q}$ and fix any
$P_0\in X_1(N)(\overline{\Q})$. Given $a\in A(\overline{\Q})$,
choose $P$ with $f(P)=a+f(P_0)$. Then
$q([P]-[P_0])=a$. No rational point on $X_1(N)$ is required.
Independently, the $r=0$ finiteness
assertion follows directly from \cite[Corollary 14.3]{kato},
applied to the modular quotient $J_1(N)\twoheadrightarrow A$
over $\Q$ and the trivial character. This result includes the
prime two and imposes no residual irreducibility hypothesis.
\end{proof}

Applying Theorem~\ref{hyp:RC} to both $n$ and $n\theta$,
Lemma~\ref{lem:forward-rank} gives rank one and finite $\Sha$, hence
$\Sha[2]=0$.
The Kummer class is then represented by actual rational points. Their
abscissas have squareclass $\beta$, by the cubic equation. Write
$H=\theta c^2$ and rescale the second point to $E_{nH}$. If the first
abscissa is $Aq^2$, divide both parameters by $q^2$ and both coordinates
by $q^2,q^3$. The abscissas become $A,Az^2$, giving precisely the
eleven witnesses, with $w=(eHz)^{-1}$. If $\theta=1$, one can reuse
a point on the same curve. This proves Theorem~\ref{hyp:PC}, with both
restricted-converse applications in their stated domain.

\section{Multiplication and the finite false-triple guard}\label{sec:mult}
The local-test construction and prime-pattern strategy adapt
\cite[Section 2 and Lemma 6.1]{openaihtp}.

The actual multiplication template first forces all three input codes
to be members of $2E(K_0)$. These membership calls are conjuncts inside
the height certificate; they are not assumed caller preconditions.
Thus even untyped rational inputs, if accepted, have unique indices.
The template has 35,134 rational witnesses. Its construction is S03
and its mathematical specification is S13.

We use the following local conventions for the displayed integral
Weierstrass equation. Put $z=-x/y$, and put $w=1/x$ with $w(O)=0$.
This $w$ differs from the auxiliary variable $-1/y$ used in
\cite[IV.1]{silverman}. The formal expansions give
$w=z^2$ times an integral unit, so $v(w)=2v(z)$.
For a complete discrete valuation with $e=v(p)$, the formal
logarithm converges on $v(z)>0$ and is an isometry with inverse
exponential on $v(z)>e/(p-1)$
\cite[IV.5.5, IV.6.3--6.4, pp.~129--132]{silverman}.
Indeed its degree-$r$ coefficient has denominator dividing $r$,
so each relative error has valuation at least
$(r-1)v(z)-e v_p(r)>0$ in this range.
Thus $\log([n]Q)=n\log(Q)$ and
$v(z([n]Q))=v(z(Q))+v(n)$ for nonzero $Q$ and $n$ there.
These statements concern the formal neighborhood of our integral
equation. When identifying it with the full reduction kernel of
a minimal model, we use
\cite[VII.2.1--2.2, pp.~188--191]{silverman}.

For this particular equation the invariant differential is
\[
 \frac{dx}{2y}=\frac{dz}{\sqrt{1+4d^2z^4}},\qquad
 \log(z)=\sum_{j\ge0}
 \frac{(-1)^j\binom{2j}{j}d^{2j}}{4j+1}z^{4j+1}.
\]
This proves the displayed coefficients of $L$ below. At a contact
prime $p>324$, with $v(z)\ge1$ and $v(p)=1$, every omitted
term relative to $z$ has valuation at least
$r-1-v_p(r)\ge28$ for $r\ge29$.
Consequently $L(z)/16575=\log(z)(1+\epsilon_z)$ with
$v(\epsilon_z)\ge28$.

For indices $n_1,n_2,n_3$, each test group chooses $H=[h]P_0$,
$U_i=[z_i]P_0$, and sets $V_0=H-[c_*]P_0$,
$V_i=U_i-[c_*n_i]P_0$, where $c_*=4M_0$ is fixed. Use
\begin{align*}
 L(z)={}&16575z-6630d^2z^5+11050d^4z^9-25500d^6z^{13}\\
        &+68250d^8z^{17}-198900d^{10}z^{21}+612612d^{12}z^{25}.
\end{align*}
This is 16575 times the degree-25 truncation of the normalized formal
logarithm. The twelve tests are four inverse abscissas at $H,U_i$,
their four conjugate counterparts at $V_0,V_i$, the three differences
$A_i-\sigma(A'_i)$, and $A_1A_2-A_3$, where
\[
 A_i=\frac{L(z(U_i))}{L(z(H))},\qquad
 A'_i=\frac{L(z(V_i))}{L(z(V_0))}\quad(1\le i\le3),
\]
and $\sigma$ is the nontrivial automorphism of $K/\Q$. In particular
the primed ratio is conjugated before the difference is taken. Identity inverse abscissas are zero; all
denominators are guarded.

For a slope $t\notin\Cset\cup\{1\}$, put
\[
 R_0=\eta F(t)/(t-1)^8,\quad\eta^2=1,\quad
 a=1-R_0q^2\ne0,\quad R=\ell_0^2a^2R_0.
\]
Denote the twelve unscaled test values in group $g$ by $D_{gj}$.
Their normalized values are
\[
 C_{gj}=\frac{D_{gj}}{R_0^{k_j}(t-1)^{k_j+1}},\qquad
 k_j=\begin{cases}1,&1\le j\le8,\\25,&9\le j\le12.\end{cases}
\]
For six independent groups set
$W_g=\prod_j(1+RC_{gj}^2)$ and
$\Theta=\Tr\sum_gW_g^{-1}$, require them nonzero, and impose
$\Phi(R/\Theta^2)$. The pole input is nonzero, so the prefactor $B$
in $Q_{\rm pole}$ cannot make these constraints vacuous.

Here is the valuation deduction at a positive odd contact prime
$p\notin S_{\rm mul}$. Normalize every $v\mid p$ by $v(p)=1$ and
put $\delta=v_p(R_0)>0$. The contact roots are distinct modulo $p$,
and $t-1$ is a unit, so $\delta$ is the positive odd contact order.
The fixed $\ell_0$ is supported on $S_{\rm mul}$. Since
$\delta+2v_p(q)$ is odd, the two summands in $a=1-R_0q^2$ have
unequal valuations. Thus $v_p(a)\le0$ and
\[
 \alpha=v_p(R)=\delta+2v_p(a)\le\delta,
 \qquad \alpha\text{ is odd}.
\]
Every factor has valuation
$v(1+RC_{gj}^2)=\min\{0,\alpha+2v(C_{gj})\}\le0$.
If test $j$ fails its threshold $v(D_{gj})\ge k_j$ at $v$, then
\[
 v(1+RC_{gj}^2)
 \le\alpha+2(k_j-1)-2k_j\delta
 \le\alpha-2\delta<0.
\]
Consequently $v(W_g^{-1})\ge2\delta-\alpha$ for that group and
place. If every group fails at every place above $p$, the trace of
their sum satisfies
\[
 v_p(\Theta)\ge2\delta-\alpha,\qquad
 v_p(R/\Theta^2)\le3\alpha-4\delta\le-\delta<0.
\]
The latter valuation is odd, contradicting Lemma~\ref{lem:pole}.
Therefore one common group and place has its first eight tests of
valuation at least one and its last four of valuation at least 25.
This argument does not require all factors in that group to be units. The first tests put the points in the
formal group. At least one of $h,h-c_*$ is a unit at the prime, so its full-log
ratio is integral. The truncated logarithm is the full logarithm times
a unit congruent to one in valuation at least 28. Its ratio is therefore
integral. The difference test transfers integrality to the other truncated
ratio, then to the other full-log ratio. Only now do the relative errors
give absolute congruences modulo $p^{28}$. Write $A_i=L(z(U_i))/L(z(H))$. The difference test and the
full-log ratios give
\[
 z_i-hA_i\equiv0,\qquad
 z_i-c_*n_i-(h-c_*)A_i\equiv0\pmod{p^{25}}.
\]
Their difference is $c_*(n_i-A_i)\equiv0$.
Since $p\notin S_{\rm mul}$ implies $p\nmid c_*$, this yields
$n_i\equiv A_i\pmod{p^{25}}$.
The product test then implies $p^{25}\mid n_1n_2-n_3$.

To specify the height certificate, write the abscissa of each input as
$x_{i0}+x_{i1}s$, using $x_{i0}=x_{i1}=0$ at the identity. Choose
$a_{i0},a_{i1},b_i\in\Q$ with $b_i\ne0$ and
$b_i x_{ij}=a_{ij}$. Express each of these nine rational quantities
as a sum of three slopes. Impose the preceding slope predicate on
all 27 slopes, with private witnesses, and impose membership on the
three input codes. This conjunction is the height certificate.
The multiplication predicate conjoins it with the fixed scalar
graphs and the inverse of the finite guard below.

\begin{proposition}[Finite-box deduction]\label{prop:box}
Assuming these local test identities and the height estimate for
$S_{\rm mul}$ with constant $A_{\rm ht}$, every false triple accepted
by the unguarded certificate has $\max_i|n_i|<T=(48A_{\rm ht})^{13}$.
\end{proposition}
\begin{proof}
Put $f=n_1n_2-n_3\ne0$. For each of the 27 slopes,
$M_{S_{\rm mul}}(t)^{25}\mid f$, so $h(t)\le A_{\rm ht}|f|^{24/25}$.
The nine rational homogeneous quantities for the three abscissas
are sums of three slopes. The sum and quotient height bounds give
$h(x_j)\le6A_{\rm ht}|f|^{24/25}+2\log3$ for each rational component.
The fixed coordinate change and the duplication bounds yield
\[
 2n_i^2\widehat h(P_0)\le12A_{\rm ht}|f|^{24/25}+B_0,
\]
where $2\widehat h(P_0)>1$ and
$B_0=4\log3+\tfrac32\log2+2\log(1+\sqrt2)+\tfrac13\log676<10$.
For $N=\max|n_i|$, $|f|\le2N^2$, so
$N^2<24A_{\rm ht}N^{48/25}+10$.
If $N\ge5$, this implies $N^{2/25}<48A_{\rm ht}$, hence
$N<(48A_{\rm ht})^{25/2}\le T$. The case $N\le4$ is immediate.
\end{proof}

For each false triple $(a,b,c)\in[-T,T]^3$, write its fifteen code
coordinates as $A_j/B_j$ with $B_j>0$, put $D_{abc}=\prod B_j$, and
form
\[
 g_{abc}(X)=\sum_{j=1}^{15}
       (D_{abc}X_j-(D_{abc}/B_j)A_j)^2.
\]
The guard requires the product of all these factors to have an inverse.
A false accepted triple makes its own factor zero. At a true triple
every factor is positive by injectivity of codes, and the inverse exists.
Thus the sound direction uses the generator and corrected all-$S$
height theorem but does not use pole completeness.

% Replacement for the prime-pattern paragraph in manuscript Section 6,
% baseline lines 740--759. The existing application beginning
% "Apply this construction to the nine integral homogeneous quantities"
% follows unchanged. This fragment uses the existing greentao, moebius,
% and local-gi bibliography keys. No fixed polynomial constant is changed.

\begin{proposition}[The fixed twenty-two-prime pattern]\label{prop:fixed-prime-pattern}
Put
\[
 \mathcal B=\{0,-15,-20,-24,-30,-60,1\},\qquad
 P_s=\prod_{p\le324}p,\qquad W=4P_s.
\]
Let $\mathcal H$ consist of the first twenty-two primes greater than
$324$ in each unit residue class modulo $W$, as in the fixed constant
prescription, and set
\[
 \mathcal K=\{dh:d\mid P_s^4,\ d>0,\ h\in\mathcal H\}.
\]
For every $A\in\mathbb Z$ and every real $B_0$ there are integers
$u_1,u_2,v$, with $v\ne0$, such that, on putting $u_3=Av-u_1-u_2$,
each of the twenty-two values
\[
 v,\qquad u_j-bv\quad(1\le j\le3,\ b\in\mathcal B)
\]
equals $\epsilon k p$, where $\epsilon\in\{1,-1\}$,
$k\in\mathcal K$, and $p>B_0$ is prime. Every such $p$ splits in
$\mathbb Q(\sqrt2)$, the two reductions of
\[
 E:y^2=x^3-(75-53\sqrt2)^2x
\]
are good, and their orders $d_1,d_2$ satisfy
\[
 \gcd(d_1,d_2)=4,\qquad p\nmid d_1d_2.
\]
The prime values need not be distinct. The assertion concerns the two
orders at the same prime only.
\end{proposition}

\begin{proof}
Write $T=(u_1,u_2,v)$ and denote the twenty-two homogeneous forms by
$F_\nu$. Their coefficient vectors are
\[
 (0,0,1),\quad (1,0,-b),\quad(0,1,-b),\quad(-1,-1,A-b).
\]
Each form has a coefficient equal to $1$ or $-1$. Within each of the
three seven-form families, a nonzero two-by-two minor is a contact
difference, of absolute value at most $61$. Between two different
families the first two coordinates give a minor $1$ or $-1$; a pair
involving $v$ also has such a minor. Thus every pair is independent
over $\mathbb Q$ and over $\mathbb F_q$ for every prime $q>324$,
for every integer $A$. No condition on three-form dependencies is
required below.

Choose a bounded open rectangular box $\Omega\subset\mathbb R^3$
of positive volume whose closure avoids all the zero planes. For each
$\nu$ choose $\epsilon_\nu$ such that $\epsilon_\nu F_\nu>0$ on
$\overline\Omega$. These signs allow all integers $A$, including zero
and negative integers.

\paragraph{The fixed multipliers and initial progression.}
For each prime $q\le324$, the proportion of $T$ modulo $q^5$ for
which at least one $F_\nu(T)$ vanishes is at most $22/q^5<1$.
Choose a residue outside this union and lift it to a residue modulo
$q^8$. On this residue each $F_\nu$ has a fixed valuation
$e_{\nu,q}\in\{0,1,2,3,4\}$. Set
\[
 d_\nu=\prod_{q\le324}q^{e_{\nu,q}}.
\]
The unit $\epsilon_\nu F_\nu/d_\nu$ is now fixed modulo $8$ at
$2$ and modulo $q$ at each odd $q\le324$. Prescribe the prime
residues
\[
 c_2=1\pmod8,\qquad c_7=3\pmod7,\qquad
 c_q=1\pmod q\quad(q\le324,\ q\ne2,7).
\]
Choose $h_\nu\in\mathcal H$ in the unit class specified by
$h_\nu\equiv(\epsilon_\nu F_\nu/d_\nu)c_q^{-1}$ at these
moduli. There are at most twenty-two demands in any class, so the
$h_\nu$ can be chosen pairwise distinct. Put $k_\nu=d_\nu h_\nu$.
Dirichlet's theorem proves that this is the same finite catalogue as
the prescribed one; no extra catalogue prime is introduced.

At $h=h_\nu$, impose $F_\nu(T)=0$. This plane has $h^2$ points.
For every $\mu\ne\nu$, the restriction of $F_\mu$ is nonconstant,
so excluding
\[
 F_\mu(T)=0,\qquad \epsilon_\mu F_\mu(T)=-k_\mu
\]
removes at most $42h<h^2$ points. Choose a remaining point.
Among its lifts modulo $h^2$, the unit
$\epsilon_\nu F_\nu/(d_\nu h)$ runs through every residue
modulo $h$: changing a coordinate with coefficient $1$ or $-1$
by $h$ proves surjectivity. Choose this residue different from $0$
and $-1$. The restrictions for $\mu\ne\nu$ remain valid.

The Chinese remainder theorem now gives one vector residue $r_0$
modulo
\[
 W_0=\prod_{q\le324}q^8\prod_\nu h_\nu^2.
\]
On $r_0+W_0\mathbb Z^3$ the functions
$\Psi_\nu=\epsilon_\nu F_\nu/k_\nu$ are integral. At every prime
dividing $W_0$ they are nonzero, and at every odd such prime they
are different from $-1$. In the affine coordinates on this
progression every variable coefficient is divisible by each prime
dividing $W_0$. Indeed, $q^8/d_\nu$ is divisible by $q$ even when
$e_{\nu,q}=4$, and $h^2/k_\nu$ is divisible by $h$.

All these choices depend on $A$ and the fixed catalogue, but not on
$B_0$ or the cutoff $L$ used next. Fix also
\[
 Q_*=\max(324,\max\mathcal H).
\]
In both upper sieves below the excluded fixed primes are exactly
those at most $Q_*$; in particular this exclusion does not grow
with $L$. There are constants $0<\alpha_A<\beta_A$ such that
\[
 \alpha_A X\le\Psi_\nu(T)\le\beta_A X
 \quad(T\in X\Omega,\ 1\le\nu\le22).
\]

\paragraph{The lower count with its cutoff dependence.}
Take $L>\max(Q_*,B_0,3)$. At each prime $q\le L$ not dividing
$W_0$, exclude the forty-four affine planes
\[
 F_\nu(T)=0,\qquad \epsilon_\nu F_\nu(T)=-k_\nu.
\]
Each has $q^2$ points and all these primes exceed $324$. With
\[
 W_L=W_0\prod_{\substack{q\le L\\q\nmid W_0}}q,
\]
let $\mathcal R_L$ be the resulting vector residues modulo $W_L$
that reduce to $r_0$ modulo $W_0$. Then
\begin{equation}\label{eq:prime-density44}
 \frac{|\mathcal R_L|}{W_L^3}
 \ge W_0^{-3}\prod_{\substack{q\le L\\q\nmid W_0}}
       (1-44/q)
 \ge c_A(\log L)^{-44}.
\end{equation}
Here and below a positive constant with subscript $A$ may depend
on the choices already made, including the fixed catalogue, but
never on $L$, $X$, or a varying fiber prime. The last inequality
follows from $\log(1-44/q)=-44/q+O(q^{-2})$ and Mertens' theorem \cite[Theorem 1, with modulus $2$]{williams-mertens}.

Fix $L$. For $r\in\mathcal R_L$ apply the finite-complexity
prime asymptotic of \cite[Main Theorem, p.~1763]{greentao},
with its hypotheses supplied by \cite[Theorem 1.1]{moebius}
and Section~\ref{v15:gi-full}, to
\[
 n\longmapsto\frac{\epsilon_\nu F_\nu(r+W_Ln)}{k_\nu},
 \qquad r+W_Ln\in X\Omega.
\]
These are integral affine forms with pairwise independent linear
parts. Singleton classes in \cite[Definition 1.5]{greentao}
give complexity at most $20$, so the inputs $\mathrm{GI}(20)$
and $\mathrm{MN}(20)$ suffice. The local inverse-theorem proof supplies
the normalized finite-family statement with linear nilsequences of
step at most $20$; its constructions and published inputs are given
in the accompanying supplement.
All coefficient bounds and moduli are fixed before $X$ tends
to infinity. At $q\mid W_L$ the local factor is
$(1-1/q)^{-22}$. At $q\nmid W_L$, scaling and changing variables
preserve the twenty-two homogeneous zero planes, so
\[
 1-22/q\le
 \mathbb P(F_\nu\ne0\text{ for all }\nu)
 \le1-22/q+\binom{22}{2}q^{-2}.
\]
Consequently the local factor is $1+O(q^{-2})$, and
\[
 \log\beta_q\ge\log(1-22/q)-22\log(1-1/q)
 \ge-484/q^2.
\]
The singular product is therefore bounded below by an absolute
positive constant, uniformly in $L$ and $r$. The real factor is
$\operatorname{vol}(\Omega)X^3/W_L^3$.

There are finitely many $r$ for fixed $L$, so their asymptotic
errors can be made smaller than half their main terms with one
threshold for $X$, depending on $A,L$. This uses no uniform
prime-asymptotic error as $L$ varies. Summing over $r$ and using
\eqref{eq:prime-density44} gives a weighted lower bound
$c_A(\log L)^{-44}X^3$. Proper prime powers contribute at most
$O_A(X^{5/2}(\log X)^{23})$: there are
$O_A(X^{1/2}\log X)$ possible proper powers in each form's range,
each equation fixing a form value has $O_A(X^2)$ integer solutions
in the box, and each weight is $O_A((\log X)^{22})$.
Thus, for all sufficiently large $X$ depending on $A,L$,
\begin{equation}\label{eq:prime-lower}
 \#\{T\in X\Omega:\ T\bmod W_L\in\mathcal R_L,
                  \ \Psi_\nu(T)\text{ all prime}\}
 \ge a_A\frac{X^3}{(\log L)^{44}(\log X)^{22}},
 \qquad a_A>0.
\end{equation}

\paragraph{An upper sieve with uniform constants.}
We give the estimates needed for both upper bounds. Consider an
axis-parallel square of side $N\ge1$ in $\mathbb R^2$. For each sieving prime
$q$ let $\Omega_q\subset\mathbb F_q^2$ have size $\rho(q)$,
where $0<\rho(q)\le21q$ and $q>Q_*$. Put
\[
 g(q)=\rho(q)/q^2,\qquad
 h(q)=g(q)/(1-g(q)),
\]
and extend both multiplicatively on squarefree integers supported
on the sieving primes. For such $d$, CRT and the elementary count
of lattice points in each residue square give
\begin{equation}\label{eq:prime-residue-error}
 A_d=N^2g(d)+R_d,\qquad
 |R_d|\ll\rho(d)(N/d+1)\le Nd+d^2.
\end{equation}
The implied absolute constant is uniform in the square's location
and in the sets $\Omega_q$. The last inequality uses
$\rho(q)\le q^2$.

For $Z\ge2$ define
\[
 G(Z)=\sum_{d\le Z}h(d),\qquad
 \lambda_d=\frac{\mu(d)}{G(Z)\prod_{q\mid d}(1-g(q))}
       \sum_{\substack{m\le Z/d\\(m,d)=1}}h(m)
 \quad(d\le Z),
\]
and set $\lambda_d=0$ otherwise. All sums here are over the stated
squarefree support, with $g(1)=h(1)=1$. These are Selberg's weights.
They satisfy $\lambda_1=1$, $|\lambda_d|\le d$, and
\[
 \sum_{e,f\le Z}\lambda_e\lambda_f g([e,f])=G(Z)^{-1}.
\]
For completeness, set $y_r=\sum_{r\mid d}\lambda_dg(d)$.
The displayed quadratic form equals $\sum_r y_r^2/h(r)$,
because $1/g(m)=\sum_{r\mid m}1/h(r)$; the stated weights give
$y_r=\mu(r)h(r)/G(Z)$. The bound on the weights follows from
the inner sum being at most $G(Z)$ and
$(1-g(q))^{-1}\le q$ for $q>324$.

Squaring the divisor sum majorizes the indicator of avoiding
all $\Omega_q$ for $q\le Z$. Inserting
\eqref{eq:prime-residue-error}, and using $[e,f]\le Z^2$,
there are at most $Z^2$ pairs with $|\lambda_e\lambda_f|\le Z^2$.
It follows that
\begin{equation}\label{eq:prime-selberg}
 S(N,Z)\ll \frac{N^2}{G(Z)}+NZ^6+Z^8.
\end{equation}

Here are the two denominator estimates, including truncation.
For $z=Z^{1/100}$ put
\[
 H(z)=\prod_{\substack{q\le z\\q\text{ sieved}}}(1+h(q)).
\]
Give each squarefree divisor $d$ of this prime product probability
$h(d)/H(z)$. The events $q\mid d$ are independent and have
probability $g(q)$. Hence
\[
 \mathbb E\log d=\sum_{\substack{q\le z\\q\text{ sieved}}}
                  g(q)\log q
 \le21\sum_{q\le z}\frac{\log q}{q}
 \le42\log z=0.42\log Z
\]
for sufficiently large $z$, using
$\sum_{q\le z}(\log q)/q=\log z+O(\log\log z)$.
Indeed, \cite[Theorem 1, with modulus $2$]{williams-mertens}
gives $S(z):=\sum_{q\le z}q^{-1}
=\log\log z+C+O(1/\log z)$, including the prime $2$.
Partial summation in
$S(z)\log z-\int_2^z S(t)\,dt/t$ gives the asserted bound.
Markov's inequality gives probability at least $0.58$, and thus
at least $1/2$, that $d\le Z$. Every such divisor is included in
$G(Z)$, so
\begin{equation}\label{eq:prime-truncated-denominator}
 G(Z)\ge\tfrac12 H(Z^{1/100}).
\end{equation}
This proves a lower bound for the truncated sum, rather than
merely for an untruncated Euler product.

In the first application the sieved primes are
$q>Q_*$ with $q\equiv1\pmod4$, with at most one prime omitted,
and $g(q)=2/q-1/q^2$. Mertens' theorem in the fixed progression
$1\pmod4$ \cite[Theorem 1]{williams-mertens} gives
\[
 \log H(z)=\log\log z+O_{Q_*}(1).
\]
Omitting a prime $\ell$ changes $H(z)$ by $1-g(\ell)\ge1/2$
when $\ell\le z$, and has no effect otherwise. Therefore
$G(Z)\gg_{Q_*}\log Z$ uniformly in the omitted prime.
In the second application all primes $q>Q_*$ are sieved and
\[
 \frac{21}{q}-\frac{210}{q^2}\le g(q)\le\frac{21}{q}.
\]
Then $\log H(z)=21\log\log z+O_{Q_*}(1)$, with an error
independent of the individual sets $\Omega_q$, and hence
$G(Z)\gg_{Q_*}(\log Z)^{21}$. These statements and
\eqref{eq:prime-selberg} establish every uniform upper sieve
estimate used below.

\paragraph{The conjugate orders.}
Every prime counted by \eqref{eq:prime-lower} satisfies
$p\equiv1\pmod8$ and $p\equiv3\pmod7$. Thus it splits in
$\mathbb Q(\sqrt2)$ and $(7/p)=-1$. The two reductions of
$d=75-53\sqrt2$ have product $7$ and opposite square classes.
The two reduced curves are the quadratic twists of $y^2=x^3-x$
by these two classes. Their traces are opposite. The discriminant
$64d^6$ is a unit at these primes, so both reductions are good.

For $p\equiv1\pmod4$, choose $i\in\F_p$ with $i^2=-1$.
The automorphism $\iota:(x,y)\mapsto(-x,iy)$ satisfies
$\iota^2=[-1]$ and commutes with the arithmetic $p$-power
Frobenius $\varphi$, directly from its coordinate formula.
The injection
\[
 \operatorname{End}(\overline E)\otimes\Q_3
 \hookrightarrow\operatorname{End}_{\Q_3}(V_3\overline E)
\]
is \cite[III.7.4, pp.~89--91]{silverman}; here $p>324$.
Since $T^2+1$ is irreducible over $\Q_3$, the centralizer of
$\iota$ on this two-dimensional Tate space is $\Q_3(i)$.
The rational centralizer therefore has dimension at most two
and already contains $\Q(i)$, so it is exactly $\Q(i)$.
The Frobenius endomorphism is integral over $\Z$ by
\cite[V.2.3.1(b), p.~142]{silverman}. Hence it lies in $\Z[i]$;
write it as $a+bi$ with $a,b\in\Z$.
The trace and determinant formulas
\cite[III.8.6, p.~99]{silverman} give
\[
 a^2+b^2=p,\qquad d_1=p+1-2a,\qquad d_2=p+1+2a.
\]
Here $d_1,d_2$ are the cardinalities of the two reduced elliptic
curve groups; the opposite twist changes the Frobenius trace's
sign. In particular $a\ne0$, since $p$ is not a square.
Both orders are divisible by $4$ because all the two-torsion is
rational. Their sum is $2(p+1)$, of $2$-adic valuation exactly
$2$, so the $2$-part of their gcd is exactly $4$. Hasse's bound
\cite[V.1.1, p.~138]{silverman}
gives $0<d_i<2p$ for $p>324$, and $d_i\ne p$ since $4\mid d_i$.
Thus $p\nmid d_1d_2$.

If an odd prime $\ell$ divides both orders, it divides $a$ and
$p+1$, and $b^2\equiv-1\pmod\ell$. The progression conditions
already rule out every odd $\ell\le L$. It remains to discard
prime values with such an $\ell>L$.

\paragraph{The number of bad prime values.}
For a fixed form its prime value lies between $\alpha_A X$ and
$\beta_A X$. Representations $p=a^2+b^2$ lie in
$[-C_A\sqrt X,C_A\sqrt X]^2$. Overcount them to bound the
number of bad primes. If $L<\ell\le X^{1/8}$, there are at most
two residues $b_0$ with $b_0^2\equiv-1\pmod\ell$. Write
\[
 a=\ell U,\qquad b=b_0+\ell V.
\]
The required $(U,V)$ lie in a square of side
$N\asymp_A\sqrt X/\ell\gg_A X^{3/8}$.
For every $q>Q_*$ with $q\equiv1\pmod4$ and $q\ne\ell$,
the condition $a^2+b^2=0\pmod q$ is the union of two distinct
lines through the origin, of size $2q-1$. The affine substitution
is bijective modulo $q$. Take $Z=X^{1/100}$. Since the prime norm
is at least $\alpha_A X>Z$ for large $X$, every required pair
avoids these residues. Equations \eqref{eq:prime-selberg} and
\eqref{eq:prime-truncated-denominator} give
\[
 \#(U,V)\ll_A \frac{N^2}{\log X}+NZ^6+Z^8
 \ll_A\frac{X}{\ell^2\log X}.
\]
The absorption is uniform: $Z^6/N\ll_A X^{0.06-3/8}$ and
$Z^8/N^2\ll_A X^{0.08-3/4}$, both tending to zero faster
than $1/\log X$. Summing over $\ell>L$ gives
$O_A(X/(L\log X))$.

For $\ell>X^{1/8}$ use the lattice count directly. Since
$a\ne0$, there are $O_A(\sqrt X/\ell)$ possible multiples $a$,
with no additive constant. For each of the at most two root
classes there are $O_A(\sqrt X/\ell+1)$ choices of $b$.
Also $\ell\le|a|\le C_A\sqrt X$. Bounding the sum over primes
by the corresponding sum over positive integers yields
\[
 \sum_{X^{1/8}<\ell\le C_A\sqrt X}
       O_A\left(\frac{X}{\ell^2}+\frac{\sqrt X}{\ell}\right)
 \ll_A X^{7/8}+\sqrt X\log X.
\]
We have therefore proved the bound, uniformly in $L$,
\begin{equation}\label{eq:prime-bad-values}
 \#\{\text{bad prime values for a fixed form}\}
 \ll_A\frac{X}{L\log X}+X^{7/8}+\sqrt X\log X.
\end{equation}

\paragraph{The number of patterns on each nonzero fiber.}
Fix $\nu$ and $p\in[\alpha_A X,\beta_A X]$, and put
$c=\epsilon_\nu k_\nu p$. The integral solutions of $F_\nu(T)=c$
are parametrized by two original coordinates: solve for a
coordinate whose coefficient is $1$ or $-1$. For example, for
$F_\nu=-u_1-u_2+(A-b)v$, use $(u_1,v)$ and set
$u_2=(A-b)v-c-u_1$. This parametrization is valid over every
$\mathbb F_q$, and the points from $X\Omega$ lie in a square
of side $O_A(X)$ in these coordinates.

For $Q_*<q\le Z=X^{1/100}$, one has $q\nmid k_\nu$ and
$q<p$ when $X$ is large. Thus $c\ne0$ modulo $q$. Each remaining
form restricts to a nonconstant affine form. Their twenty-one
zero lines are distinct: if the restrictions of $F_j,F_k$ had
the same zero line, then for some $\lambda\ne0,\mu$,
\[
 F_j-\lambda F_k=\mu(F_\nu-c)
\]
as affine forms over $\mathbb F_q$. Comparing constants gives
$\mu c=0$, hence $\mu=0$, contradicting pairwise independence
of $F_j,F_k$. Parallel lines and multiple intersections cause
no difficulty. The union consequently has between
$21q-\binom{21}{2}$ and $21q$ points.

Every target pattern on the fiber avoids this union: its other
prime values are also at least $\alpha_A X>Z$, and every
$k_j$ is a unit at $q>Q_*$. Discarding the CRT progression
conditions only increases the count. The second denominator
estimate in \eqref{eq:prime-truncated-denominator} gives
\begin{equation}\label{eq:prime-fiber}
 \#\{\text{target patterns on }F_\nu=c\}
 \ll_A \frac{X^2}{(\log X)^{21}}+X^{1.06}+X^{0.08}
 \ll_A \frac{X^2}{(\log X)^{21}}.
\end{equation}
This is uniform in the fiber prime $p$ and independent of $L$.
In particular $q=p$ is outside the sieving range; no varying
exceptional-prime product has been suppressed.

\paragraph{The order of choices.}
Multiply \eqref{eq:prime-bad-values} by \eqref{eq:prime-fiber}
and sum over the twenty-two forms. The number of bad patterns
is at most
\[
 b_A\frac{X^3}{L(\log X)^{22}}
 +O_A\left(\frac{X^{23/8}}{(\log X)^{21}}
          +\frac{X^{5/2}}{(\log X)^{20}}\right),
\]
where $b_A$ is independent of $L$. First choose
$L>\max(Q_*,B_0,3)$ so large that
\[
 \frac{b_A}{L}<\frac{a_A}{2(\log L)^{44}}.
\]
Such an $L$ exists because $(\log L)^{44}/L\to0$.
Then keep $A,L$ and every preceding choice fixed, and take $X$
large enough for \eqref{eq:prime-lower}, the upper bounds, and
the remaining error to be less than
$a_A X^3/(2(\log L)^{44}(\log X)^{22})$.
A pattern survives. Its signed prime values tend to infinity
with $X$, $v\ne0$, and the preceding order calculation proves
all the assertions.
\end{proof}
Apply this construction to the nine integral homogeneous quantities
of a true triple. Each slope has at most six odd-contact primes.
At one such prime with reduced group cardinalities $d_1,d_2$,
which annihilate the reduced point, solve
\[
 h\equiv0\pmod{\operatorname{lcm}(M_0,d_1)},\qquad
 h\equiv c_*\pmod{d_2}.
\]
Compatibility follows from
\[
 \gcd(\operatorname{lcm}(M_0,d_1),d_2)
                  \mid M_0\gcd(d_1,d_2)=4M_0=c_*.
\]
Choose a solution larger than $c_*$. Each contact prime uses its own
private group within its own slope. No common shift across different
primes is imposed, and selected prime values need not be distinct.
The saved library allocates six private groups per slope and twenty-seven
private slope calls. Set $U_i=[n_i]H$ and assign one group to each prime. Fixed local
clearing at the old primes can be checked as follows.
Put $S_b=S_{\rm ht}\cup\mathcal H$, $B_p=4$ on $S_{\rm ht}$ and
$B_p=1$ on the remaining catalogue primes, $C_p=14B_p$, and
$D_p=51C_p$. The selected integral forms give
$|v_p(R_0)|,|v_p(t-1)|\le C_p$, while the small formal subgroup
in the definition of $M_0$ gives $v_p(D_{gj})\ge0$ at $S_b$.
Consequently $v_p(C_{gj})\ge-D_p$. The fixed normalization has
$v_p(\ell_0)=E_p=52C_p+v_p(12)+4$.

Weak approximation chooses $q$ small at $S_b$, integral at the
selected primes, and arbitrarily close to a real root of
$1-R_0q^2$, without equality; choose $\eta$ to make $R_0>0$.
Then $a$ is a unit at the specified finite primes. At $p\in S_b$,
\[
 v_p(RC_{gj}^2)\ge 2E_p-C_p-2D_p>v_p(12),\qquad
 v_p(\Theta-12)>v_p(12),
\]
so $v_p(R/\Theta^2)\ge103C_p+8\ge4$.
At a selected odd-contact prime, $v_p(R)=1$. Each of the twelve
summands of the trace of the six $W_g^{-1}$ has residue zero or one;
the assigned group and place contributes one. Since $p>324$, their
sum is a unit and the pole input has valuation one. At every other
prime outside $S_b$, its valuation is even. Thus its reduced
denominator is square. In both real embeddings each $W_g$ is positive;
as $a\to0$, $W_g\to1$, $\Theta\to12$, and $R/\Theta^2\to0$
through positive values. This verifies exactly $\mathsf{PC}$'s
hypotheses, including all nonzero guards. Negative indices, zero indices and
coincident points are handled by the group-law branches. Therefore
Theorem~\ref{hyp:PC} supplies the required pole witnesses. This proves
the full multiplication statement $\mathsf M$, including completeness
for every integer triple $(n_1,n_2,n_1n_2)$ and soundness on arbitrary
rational input codes.

% End 02-height-pole-multiplication.tex


% Begin 03-restricted-converse.tex
\section{The restricted converse and its dependency order}\label{sec:rc}
The organization of the restricted-converse argument follows
\cite{openaipointwise}, particularly its ring-class interpolation,
determinant transfer, coefficient systems and fixed-curve assembly.
The broader converse is stated in \cite{openaigoldfeld}. Our
specialization uses the same full Selmer-corank convention, including
any divisible part of $\Sha$, as those sources.

For a positive squarefree $\ell\equiv11\pmod {12}$ put
\[
 E_\ell:y^2=x(x-\ell)(x+3\ell).
\]
The distinct roots $0,\ell,-3\ell$ give full rational two-torsion.
Expanding the cubic gives $a_2=2\ell$, $a_4=-3\ell^2$ and
$a_1=a_3=a_6=0$. With the usual Weierstrass conventions,
$b_2=a_1^2+4a_2=8\ell$, $b_4=2a_4+a_1a_3=-6\ell^2$,
$b_6=0$, and $b_8=-a_4^2=-9\ell^4$. Hence
\begin{equation}\label{eq:rc-noncm-invariants}
\begin{split}
 c_4&=b_2^2-24b_4=208\ell^2,\\
 \Delta&=-b_2^2b_8-8b_4^3=2304\ell^6\ne0,\\
 j(E_\ell)&=c_4^3/\Delta=35152/9.
\end{split}
\end{equation}
For a CM elliptic curve, its $j$-invariant is a root of the monic
integral order class polynomial, including for nonmaximal orders
\cite[\S1, p.~607, and \S2.1, pp.~612--613]{fowler}.
A rational algebraic integer $a/b$ in lowest terms, $b>0$, is an integer:
a monic integral equation of degree $n$, multiplied by $b^n$, gives
$b\mid a^n$, so $b=1$. Since $35152/9$ is in lowest terms and has
denominator $9$, the curve $E_\ell$ is not CM.
The curve in this section is different from the fixed CM curve used
to represent integer indices.

The assertion needed in Theorem~\ref{hyp:RC} is
\[
 \operatorname{corank}_{\Z_2}\Sel_{2^\infty}(E_\ell/\Q)=1
 \quad\Longrightarrow\quad \ord_{s=1}L(E_\ell,s)=1.                 \tag{RC}
\]
We prove this assertion in this section and the next two sections.
Appendices~\ref{app:even-source}--\ref{app:endpoints} give the fixed-form,
theta, nonvanishing and normalization constructions used in this deduction.
The modularity theorem of Breuil, Conrad, Diamond and Taylor
\cite{bcdt}, in the published formulation
\cite[Theorems C.13.6 and C.16.4(a),(c), pp.~442, 451]{silverman},
makes $E=E_\ell$ modular and supplies a rational parametrization
from $X_0(N)$, where $N$ is the conductor of $E$.
Choose the rational weight-two newform $f\in S_2(\Gamma_0(N))$
with integral Fourier coefficients supplied by
\cite[Theorem 3.1]{amirhong}, using the normalized convention $a_1(f)=1$.
For every prime $q\nmid 2N$ its coefficient satisfies
$a_q(f)=q+1-\#E(\F_q)$.
Apply \cite[Theorem 1, pp.~526--527]{shimura-factors}
with $\Gamma=\Gamma_0(N)$. Since the coefficient field is $\Q$,
this gives an elliptic curve $A_f/\Q$ and a quotient
$q_f:J_0(N)\twoheadrightarrow A_f$ over $\Q$.
Here $T_q$ on the Jacobian is the Hecke correspondence whose
pullback on differentials is the weight-two operator, as in
\cite[equations (1.3)--(1.4), p.~525]{shimura-factors}.
The quotient satisfies $q_fT_q=[a_q(f)]q_f$ for $q\nmid N$.
Shimura's same theorem identifies the local Euler factors of $A_f$
and $f$ outside finitely many primes. Together with the displayed
coefficient equality, this makes the local Euler factors of $A_f$
and $E$ equal outside finitely many primes.
Thus \cite[\S5, Corollary 2, p.~361, with the 1984 erratum]{faltings1983}
places $A_f$ in the $\Q$-isogeny class of $E$.

For the present curve $E_\ell$, with $\ell$ odd, the displayed
discriminant is $2304\ell^6$, of valuation eight at two.
A change of Weierstrass model changes this valuation by a multiple
of twelve. Good reduction at two would give an integral model of
valuation zero, which is impossible. Thus $2\mid N$, and every
prime outside $N$ is already outside $2N$ in the trace comparison.

We spell out the operator normalization in this application.
The newform is a common eigenform for the Hecke operators, including
the level-prime operators. At $q\nmid N$ the determinant-$q$
correspondence acts on $f(z)\,dz$ by
\[
 \left(qf(qz)+\frac1q\sum_{j=0}^{q-1}f((z+j)/q)\right)dz.
\]
The coefficient of $e^{2\pi inz}$ is
$a_{qn}(f)+q a_{n/q}(f)$, with the second term zero unless $q\mid n$.
At a level prime the translation branches give $U_p$, with
coefficient $a_{pn}(f)$. Since $a_1(f)=1$, the eigenvalues are the
corresponding Fourier coefficients. The level-prime observation
does not assert a bad-prime trace formula absent from the cited
modularity statement. Shimura's equations (1.3)--(1.4) identify
this action on differentials with the action of the correspondence
on the Jacobian. The coefficient field is $\Q$, so its unital
coefficient action on $A_f$ sends $a_q(f)\in\Z$ to multiplication
by that integer. Hence $q_fT_q=[a_q(f)]q_f$ is an equality of the
actual homomorphisms. There is no torsion ambiguity: a homomorphism
of abelian varieties killed by an integer has connected image in
a finite group scheme and is zero in characteristic zero.

For clarity, the good-prime correspondence is also the unaveraged
geometric neighbor correspondence used below. An order-$q$ cyclic
subgroup of an elliptic curve with level-$N$ subgroup gives its
quotient with induced level structure. The dual $q$-isogeny returns
that level subgroup to itself, since multiplication by $q$ is
bijective on it. Thus the correspondence is symmetric under
exchanging its two projections. Its $q+1$ upper-triangular branches
fix the cusp at infinity, so $T_q[\infty]=(q+1)[\infty]$.
Consequently its action on $[x]-[\infty]$ is precisely the sum
of the $q+1$ corresponding degree-zero divisor classes. Composing
with a $\Q$-isogeny $A_f\to E$ preserves the Hecke identity and
introduces neither an averaging factor nor a cusp correction.
Choose a $\Q$-isogeny $\alpha:A_f\to E$ and put
$\varphi=\alpha q_f$ and $\pi(x)=\varphi([x]-[\infty])$.
Thus $\pi(\infty)=0$ and, for $q\nmid N$,
$\varphi T_q=[a_q(f)]\varphi$.
Fix the integral Tate lattice
and the local squareclasses allowed at the primes of $2N$.
Constants depending on this data are chosen before any conductor,
companion field or binary dimension varies. Quotienting by a rational
point of order two gives a rational two-isogeny. For the parity step,
apply \cite[Theorem 1.4 and the definition preceding Conjecture 1.2,
p.~568]{ddselmer} with $E/\Q$ and $p=2$. Its full $2^\infty$-Selmer
rank is the $\Z_2$-corank in (RC), so it gives root number $-1$
under the left side of (RC). This application does not replace Selmer
corank by Mordell--Weil rank and does not assume finite $\Sha$.
\begin{remark}[Why this family is the input to the converse]\label{ex:sec7-rc-scope}
The congruence $\ell\equiv11\pmod {12}$ comes from the pole-completeness
construction. In the proof of Theorem~\ref{hyp:PC}, the two parameters
$n,n\theta$ are positive squarefree numbers in this congruence class;
their two-Selmer groups have dimension three. Lemma~\ref{lem:corank}
gives full two-primary Selmer corank one, so these are exactly the
two applications required of Theorem~\ref{hyp:RC}.

In the converse argument we fix one such $\ell$. In particular $\ell$
is odd, is prime to three, and is greater than one, hence is not a
square. We use the displayed full rational two-torsion and non-CM
properties, the resulting rational two-isogeny, and the explicit roots
$0,\ell,-3\ell$. The last roots enter the $K$-split seed argument below.
The constants may depend on this fixed curve; there is no bound on
$\omega(\ell)$ and no assertion that they are uniform as $\ell$ varies.
The statements about coefficient graphs below are the rational,
full-two-torsion case used by $E_\ell$ and its quadratic twists.
\end{remark}

Here is the dependency order of the proof. A fixed-form Euler-system
calculation gives the even coefficient bound. A varying-companion
arithmetic estimate, proved by the two limits below, then gives the
genus-point lower bound. Only after that lower bound do we normalize
the odd coefficients. Their finite-stage forms and the even forms give
one common binary family. Finally a fixed-field determinant argument
excludes a zero derivative at the base. The even estimate does not use
that final determinant argument on Heegner classes.

\subsection{The even coefficient calculation}
Use normalized central values $L_*(h)$ with the real and imaginary
periods fixed by $\pi$ and the integral lattice. Only their valuations
are used; every subsequent comparison uses this same normalization.
For every twist whose two-primary Tate--Shafarevich group is finite, let
$s(h)=\operatorname{length}_{\Z_2}\Sha(E^{(h)}/\Q)[2^\infty]$.

\begin{proposition}[The fixed-form even scalar]\label{ex:sec7-even-scalar}
Use the fixed $E,\pi$, integral lattice and allowed old local types
specified above. Put $\Lambda=\Z_2[[t]]$. For a finite elementary
abelian two-group $G$, let $\chi_x=\chi_0(x\circ\varepsilon)$
$(x\in\widehat G)$ be a family of quadratic characters with a common
real sign and the same restrictions at every place dividing $2N$.
There is one scalar $u^{\rm ev}\in\Lambda[1/2][G]$ satisfying (3.2).
Its clearing exponent and the bounded error in that formula depend
only on the fixed data, not on $G$ or the good-prime conductor support.
The central identity is used only at corank-zero constituents with
nonzero central value; at a corank-zero constituent the same central
coordinate is zero exactly when the central value is zero.
\end{proposition}
\begin{proof}\leavevmode
The following construction is expanded in
Propositions~\ref{ex:app-kato-system} and~\ref{ex:app-even-scalar},
with its central calculation in Proposition~\ref{ex:app-central-calibration}.
The fixed-form construction uses one integral symbol combination with
nonzero projections to both real signs. Kato's integral maps, with
$j_*$ cohomology taken modulo each power of two before passage to the
limit, give classes whose inertia restrictions vanish integrally.
Indeed the low-degree Leray edge map gives this vanishing before
rationalization. Restriction of transfer is the Mackey sum of the
corresponding upper inertia restrictions. Hence unnormalized transfer
preserves the vanishing even when the intervening extension is ramified.
No character projector divided by a growing group order is used.

For a finite binary character group, let $F$ be the least common
multiple of the odd parts of all its character conductors, including
the base character. Write $Z_F$ for the compatible Iwasawa class
obtained by pushing the retained-smoothed Kato classes at levels
$F2^n$ through the integral coefficient map.
Here and throughout the common-scalar construction, the binary family
is required to have one real sign and fixed local squareclasses at
$2N$. Writing its characters as
$\chi_x=\chi_0(x\circ\varepsilon)$, this means
$\varepsilon(c)=1$ for complex conjugation $c$, and
$\varepsilon|_{G_{\Q_v}}=1$ at every $v\mid2N$.
Indeed every character of $G$ evaluates these elements to one, and
the characters separate $G$. The lower bound uses $G=1$; the
upper-transfer cubes lie in the fixed signed unit filter of their
coefficient source. The final cubes have the same property.
At a character $\chi$ its unused-prime factors are
\[
 E_q(\chi)=1-\frac{a_q\chi(q)}q\sigma_q^{-1}
                 +\frac1q\sigma_q^{-2}.                           \tag{3.1}
\]
They are the actual determinants of the unused good local terms.
Active ramified good primes have rational inertia $-1$, so have no such
rational local term. The fixed bad and dyadic corrections form a factor
$F_E$ with nonzero central value.

The rational divisibility input is Kato's Theorem 13.4(2),
together with the nonzero-class input in 13.5(2)
\cite{kato}. At two we justify the divisibility through
\cite[Theorem 0.8]{kato1999}; Appendix~\ref{app:even-source}
checks its stronger abelian-extension hypothesis and compares the
two cohomology conventions. We use the height-one primes not
containing two.
Rational irreducibility is preserved on a quadratic subgroup of the
non-CM open image. That subgroup contains a nontrivial unipotent,
which supplies the required one-dimensional fixed space. Oddness
supplies the two real eigenspaces, and purity supplies the common
Frobenius weight. This application does not use residual irreducibility
at two or the stronger hypotheses of Theorem 13.4(3).

The smoothing constants $c,d_s$ can be chosen to satisfy
\[
 c,d_s\equiv1\pmod {3N_{\rm odd}F},\qquad
 c,d_s\equiv1+2^e\pmod {2^{e+1}},
\]
where $e\ge\max\{3,v_2(N)\}$ is fixed by the original dyadic data.
In weight two and
cohomological degree one the two smoothing exponents are one.
Kato's explicit reciprocity formula therefore contributes valuation
$2e$, independent of $F$. For the central normalization, first descend
to the primitive character conductor and then use the invariant basis
$\sqrt h\,e_{\chi_h}$ of the quadratic character line. The resulting
varying scalar $1/h$ is a two-adic unit. The complete calculation in
Appendix~\ref{app:even-source} includes discriminant $4h$; a Gauss basis
there has the extra factor $\sqrt{\Delta_h/h}$, which must also be
included in its change of basis. The numerical sizes of $c,d_s$ may
grow, while their relevant valuation remains fixed.

For completeness, the integral calculation at two uses the positive
complex. Finite continuous cochains are filtered unions of finite free
modules and hence flat. Splitting the residual complex into cohomology
and disks, then using the nilpotent perturbation calculation of
Section~\ref{sec:limits}, gives finite free models with their maps.
In the fixed real cyclotomic tower the ordinary two-primary class group
vanishes. One proves this inductively: only the prime above two ramifies
in each quadratic step; the abelianization of a possible Hilbert
two-class tower is generated by that inertia of order two. A finite
two-group with cyclic abelianization is cyclic, and inertia excludes a
larger cyclic extension. The unit-signature calculation is separate.
A norm-minus-one unit in positive layers, and $-1$ at the base, has
signature vector a unit in the characteristic-two cyclic group algebra.
For a scalar coefficient the resulting residual dimensions are
\[
 \dim H^1=2^n+s_n-1,\qquad \dim H^2=s_n-1.
\]
For each fixed support $s_n$ is bounded. Universal coefficients give
generic residual rank one in degree one; the two trivial elliptic
factors give rank two. Unit cancellations thus identify the positive
complex over $\mathscr O[G]$ with $\mathscr O[G]^2[-1]$.
This last rank is independent of support, although preliminary models
need not have a uniform size.

A generator of the fixed real invariant lattice gives a connecting class $a$.
Its central image need not be saturated.
Twice an ordinary class lifts because the relevant real $H^1$ is killed
by two. Wedge $a$ with a lift of $2F_E Z_F$ in the original determinant
line. This produces a common scalar $u^{\rm ev}$.
The first assertion below holds for the full family. The central valuation
identity is asserted only at a constituent of Selmer corank zero with
$L_*(h)\ne0$:
\[
 2^{C_{\rm ev}}u^{\rm ev}\in\Lambda[G],\qquad
 v_2u^{\rm ev}_h(0)=v_2L_*(h)-2\omega(h)-s(h)+O_E(1).               \tag{3.2}
\]
Here $\Lambda=\Z_2[[t]]$ and
$\mathscr O=\widehat{\Lambda_{(2)}}$. The error and $C_{\rm ev}$ are
fixed before the binary group and support. To justify membership in
$\Lambda[G]$, first use rational characterwise divisibility and Fourier
inversion for membership in $\Lambda[1/2][G]$, without estimating the
Fourier denominator. Intersect coefficientwise with the integral
$\mathscr O[G]$ result. The explicit coefficient construction also covers field discriminant $4h$,
with its dyadic character retained before the cofinal norm limit, as in
Appendix~\ref{app:even-source}. At a corank-zero constituent the same coordinate
vanishes centrally if and only if the central value vanishes.
At such a nonzero corank-zero constituent, keeping the real saturation
and torsion terms gives the local formula
\[
 v_2u(0)=v_2\exp^*z-s-\sum_v g_v+2g+1-r_1-e_2.
\]
At an active good odd prime $g_v=2$, since $A^{I_v}=E[2]$.
Unused terms cancel (3.1); all other terms belong to the fixed local
list. This proves (3.2) with the indicated active-prime coefficient.
\end{proof}

\subsection{Nonvanishing with all support restrictions}\label{sec:nonvan}
\begin{proposition}[Nonvanishing in a prescribed unit filter]\label{ex:sec7-nonvanishing}
The curve $E$ is the fixed curve of this section. The following
assertion allows any nonempty signed unit filter and any additional
finite avoidance set; its bound on the number of prime factors does
not depend on either set.
A signed unit filter prescribes a real sign and unit squareclasses at
a finite set $S$ containing $2,3$ and the bad primes. Let $B$ be any
additional finite set of forbidden primes. If its root number is even,
there is a squarefree parameter in this filter, coprime to $B$, with
at most five prime factors and nonzero central value. In an odd filter
there is a simple-zero parameter, with no asserted prime-count bound.

\end{proposition}
\begin{proof}\leavevmode
Here is the reduction to the precise classical statements.
Choose a signed prime $q_0$ in the prescribed unit classes by
Dirichlet's theorem, avoiding $S\cup B$. Apply Hoffstein--Luo
\cite{hl} in the even case, or Bump--Friedberg--Hoffstein
\cite{bfh} in the odd case, to the actual curve $E^{(q_0)}$.
Require their positive discriminant $d$ to split at
$S\cup B\cup\{|q_0|\}$ and at the conductor primes. Then $q_0d$ has
the prescribed original classes. Hoffstein--Luo gives at most four
prime factors of $d$, so at most five in total.
The BFH clause used in the odd case is the simple-zero clause.

The root-number check includes ramified dyadic quadratic characters.
At the primes in $S$ equal squareclasses give isomorphic local curves.
At a good odd prime newly ramified in the parameter, the root factor is
$(-1/p)$ \cite[Proposition 3.3]{desjardins}. Its product outside $S$
is fixed by the real sign and the dyadic unit class. Thus unit classes
$3$ or $7$ modulo eight are allowed; the field discriminant may be four
times the parameter. We do not invoke an odd-discriminant root formula
on those parameters.
See also Proposition~\ref{ex:app-five-prime} for the same
five-prime bound and its subsequent use in the companion construction.
\end{proof}

\subsection{Coefficient symbols and the finite address lemmas}
We state the finite interface before using it for the even upper bound.
The graph language records quadratic residue data of the coefficient
indices. It does not replace a Selmer group by an abstract graph.
Throughout this subsection all vertices are in the rational
full-two-torsion case. Thus there is one independent residue bit per
pair; the split/simple distinction needed for other two-torsion fields
does not occur here.

\begin{definition}[Compatible typed graphs]\label{ex:sec7-typed-graph}
Fix the finite support $S$ of a coefficient source, containing two,
and the finite quadratic characters used to read its coefficients.
The \emph{type} of a prime $p\notin S$ records its values under these
characters and its unit squareclasses at $S$, including its class
modulo eight. We write the resulting finite elementary two-group
additively as $\mathcal L$. We retain all relations among these
coordinates; a tuple of arbitrary character values need not be a type.
Additional fixed splitting conditions restrict us to their actual
realizable types. For an odd theta stage, its reduction prime is part
of this fixed support during that stage.
Let $S_{\rm lab}$ also contain the ramification of any additional
stationary label tests. Enlarging $S_{\rm lab}$ records labels and
prime exclusions; it does not change the coefficient source or add
relaxed places to a Selmer complex.

A \emph{compatible typed graph} consists of a finite ordered vertex
set $V$, an individually realizable type $\ell_v\in\mathcal L$ for
each vertex, and quadratic residue bits. For prime representatives,
write $z_{vw}=0$ or $1$ according as $(p_v/p_w)=1$ or $-1$.
For each unordered pair just one direction is freely specified; the
other satisfies
\[
 z_{vw}+z_{wv}=\epsilon_v\epsilon_w,\qquad
 \epsilon_v=(p_v-1)/2\pmod2,
\]
where $\epsilon_v$ is part of the type. There are no self-residue bits.
Only these compatible, realizable data are used below. A fixed filter
prescribes the total type $\sum_v\ell_v$ needed for the coefficient
index, including the signed parameter convention of the source.
If that convention is $h=\sigma D$, then $D=\prod_v p_v>0$ is the
Fourier index and $\sigma$ is fixed by the source. The signed factors
$p_v^*=(-1/p_v)p_v$ will be used when the same graph defines a family
of quadratic characters.

Contracting two vertices replaces their type by its sum and every
incidence to an external vertex by the sum of the two incidences.
Equivalently its template is the product $p_vp_w$; it is represented
by a fresh prime with those data. Reverse incidences are determined
by reciprocity. No bit internal to the contracted cluster is retained.
A vertex is \emph{primary} if its generator is a square unit at the
fixed support and its type is neutral in every retained quadratic
and stationary coordinate. It is \emph{neutral relative to a graph}
if in addition all its incidences to the other vertices are zero.

Two symbols may read projections of one common type group. The
projections retain the relevant coordinates, preserve primary types,
and commute with the product and external-incidence rules just given.
We require each indicated projection to be surjective onto its own
realizable alphabet; we do not require the map to their product to
be surjective.
\end{definition}

\begin{definition}[Coefficient symbols]\label{ex:sec7-symbol}
A \emph{symbol} is the function $A$ with values in $\F_2$ obtained
from one fixed coefficient source and one fixed normalized depth
$\delta$. On an allowed positive squarefree index $D$ it is
\[
 A(D)=
 \begin{cases}
  1,&v_2(a(D))=\delta+\omega(D),\\
  0,&v_2(a(D))>\delta+\omega(D),
 \end{cases}
\]
and it is zero off that source's filter. For the even source, exact
phase removal makes the relevant coefficient rational. For the odd
source, (3.8) gives the required phase removal modulo the retained
power $2^{M_\theta}$; use its integer representative at that precision.
A coefficient zero modulo $2^{M_\theta}$ is invisible. Whenever a
unit is tested, the precision exceeds $\delta+\omega(D)$, so the test
is independent of the representative. No exact rationality of an
arbitrary lifted odd theta coefficient is asserted. The lower bound making these alternatives exhaustive
is part of the coefficient construction, not part of this definition.
For the even source it follows from (3.2) and the coefficient-square
formula of Proposition~\ref{ex:app-even-phase}.

For an odd source this definition at one finite stage is used only
through a specified finite squarefree-weight cutoff and at a stage
where the lower bound holds through that cutoff. The chosen cutoff
and all required divisions must be included in the operator budget.
The fixed minimum and the compatible finite tables are constructed
in Propositions~\ref{ex:app-odd-minimum} and~\ref{ex:app-odd-stabilization}.
The resulting limit symbol is a function on finite compatible typed
graphs. It is not a valuation rule at all weights of one finite form.
The same source and depth are kept when a vertex is omitted.
\end{definition}

\begin{lemma}[The coefficient rules]\label{ex:sec7-coefficient-rules}
After the stated coefficient lower bounds have been proved, the even
symbol and the stabilized odd symbol satisfy the following rules on
compatible typed graphs. At an odd stage a rule is asserted only when
all its inputs, outputs and coefficient divisions lie in the chosen
finite budget. Configurations in these rules, and outputs of
contractions, are nonempty.
\begin{enumerate}
\item \emph{Graph dependence and whole-prime omission.}
For a new vertex $q$ adjoining a list $U$, its reading factors through
the elementary quadratic extension unramified outside $S\cup U$.
It is determined by its type and its incidences with $U$.
The reading after omission of $p\in U$ uses the same original form,
phase characters and depth, with varying support
$U\setminus\{p\}$. That extension is unramified at the whole
rational prime $p$. Changing inertia at the omitted prime has no
effect on this reading.
\item \emph{Contraction.}
For the independent bit of a pair $v,w$, write
$\Delta_{vw}A=A|_{z_{vw}=0}+A|_{z_{vw}=1}$, changing the reverse bit
as required by reciprocity. Then
\[
                 \Delta_{vw}A(\Gamma)=A(\Gamma/(v,w)).
\]
The graph on the right has the summed type and \emph{every} summed
external incidence from Definition~\ref{ex:sec7-typed-graph}. Its
fresh prime representative is read using the omitted-prime source
in the first rule.
\item \emph{Primary isolation.}
If a configuration contains a primary vertex neutral relative to
all its other vertices, its symbol is zero.
\end{enumerate}
These are identities of the actual coefficient sources, rather than
axioms imposed on an arbitrary binary function.
\end{lemma}
\begin{proof}\leavevmode
Proposition~\ref{ex:app-coefficient-rules} proves the rules from the
integral shell recurrence, the unary subtraction and the weight-two
trace and inertia identities. Divisibility by the integer two precedes
reduction modulo the residue field. In the present full-two-torsion
case every new vertex is split in $\Q(E[2])=\Q$; the contraction
coefficient is one. The omitted form is constructed from the same
original source, so its unramifiedness removes the entire prime and
does not choose a new minimum. The final fresh-copy comparison in that
proof shows dependence on the whole typed graph, including overlapping
prime realizations. The same argument proves the rules for the
even source; no odd normalization is needed in that application.
\end{proof}

\begin{lemma}[Finite realization]\label{ex:sec7-realization}
Keep finitely many old prime vertices and all their mutual data fixed.
Given individually realizable types for finitely many new vertices,
every prescription of the independent forward bits involving a new
vertex can be realized by distinct fresh primes outside any prescribed
finite avoidance set. This includes primary vertices and extra fixed
splitting conditions. A product template can likewise be replaced by
a fresh prime with the same type and all external incidences whenever
that output type is one of the permitted realizable types.
\end{lemma}
\begin{proof}\leavevmode
For rational quadratic characters the fixed type is a reduced residue
condition at a modulus supported on $S_{\rm lab}$, including eight. Choose the
new primes successively. A desired residue bit at an old prime is a
nonzero residue class there. The Chinese remainder theorem combines
these classes with the fixed type, and Dirichlet's theorem supplies
infinitely many primes in the resulting reduced residue class.
At each step avoid all previously chosen primes and the finite
avoidance set. Reciprocity gives the required reverse bits. Apply
the same construction to the external residues and support type of
a product template. The graph dependence in
Lemma~\ref{ex:sec7-coefficient-rules} makes replacement by this fresh
representative legitimate for the symbol.
\end{proof}

\begin{definition}[Unit seeds and partitionable lists]\label{ex:sec7-seed-partition}
A \emph{unit seed} is a finite compatible typed graph with symbol
one. A \emph{minimum nonempty seed} has the least positive vertex
count among such seeds. This terminology does not exclude the
possibility that the empty configuration is also a unit.
A nonempty prime list is \emph{partitionable} if its vertices admit
a partition into nonempty clusters, with a contraction tree in each
cluster, whose contracted graph is a unit seed. These trees and the
resulting types and incidences constitute a contraction certificate.
This condition makes no assertion about a Selmer group.
\end{definition}

\begin{lemma}[Minimum seeds and contraction certificates]\label{ex:sec7-seed-certificate}
At the least positive vertex count of a unit seed, its value is
independent of all internal pair bits and depends only on its vertex
types. A contraction certificate for a nonempty list makes its
polynomial in the independent pair bits nonzero. In particular, a
list with the same total type as a minimum nonempty seed is
partitionable if it contains anchors for the first $r-1$ seed types
and a nonempty remainder, where $r$ is the seed's size.
\end{lemma}
\begin{proof}\leavevmode
If $r\ge2$, a nonzero derivative in an internal bit would, by
contraction, give a unit on $r-1$ nonempty vertices. This contradicts
minimality. If $r=1$ there is no internal bit. Iterating the contraction
identity along a forest gives its contracted symbol; an output equal
to one therefore proves that the original polynomial is nonzero.
For the last assertion keep the $r-1$ anchors and merge the whole
remainder into the last cluster. Equality of total types forces the
last seed type. The internal bits of the minimum seed do not matter.
All contractions have nonempty output.
\end{proof}

\begin{definition}[Addresses and terminal polynomials]\label{ex:sec7-address}
For an integer $b\ge1$, an \emph{address} is $x\in\F_2^b$.
A vertex with activation $\lambda\in(\F_2^b)^*\setminus\{0\}$ is
present at $x$ exactly when $\lambda(x)=1$. Stationary vertices are
present at every address. At a fixed address, retain all vertex
types and any fixed external incidences, and regard a symbol as a
function of the independent pair bits between the active terminals.
Its unique multilinear polynomial over $\F_2$ is the
\emph{terminal polynomial}. Its variables describe possible changes
of those bits; the initial prime realization already has fixed bit
values. Finite differences in these variables are formal derivatives.
\end{definition}

\begin{lemma}[One symbol with a stationary base]\label{ex:sec7-one-symbol-address}
Fix a coefficient symbol satisfying
Lemma~\ref{ex:sec7-coefficient-rules}, a nonempty stationary terminal
list, its types and any stationary external data. Suppose its terminal
polynomial is nonzero and its total data lie in the source's filter.
For every $b\ge1$ and every fixed assignment of the list's pair bits, finitely
many new primary vertices with nonzero linear activations can be
adjoined, with only their new incidences chosen, so that the symbol is
one at every $x\ne0$. The original vertices and all their mutual data
are unchanged. No new vertex is present at zero. Fixed extra splitting
and neutral external-incidence requirements may be included in the
primary types.
For a fixed even source the original prime representatives can be
kept throughout. For a stabilized odd symbol, the fixed data in this
statement are the abstract types and bits; its actual theta stage is
chosen before its prime representatives, as specified in the proof.
\end{lemma}

\begin{lemma}[Two symbols on one address family]\label{ex:sec7-two-symbol-address}
Fix one or two coefficient symbols satisfying
Lemma~\ref{ex:sec7-coefficient-rules}, reading compatible projections
of one common realizable alphabet. Fix $b\ge1$ and terminals with
nonzero linear activations, including at least one terminal for each
nonzero linear form. Stationary external data may also be fixed.
Assume that at every $x\ne0$ the active configuration is nonempty,
its total data satisfy every required filter, and each symbol's
terminal polynomial is separately nonzero.
Then, starting with any fixed assignment of these compatible types
and bits, finitely many new primary vertices with nonzero linear activations can be
adjoined so that all the symbols are one at every $x\ne0$.
Only incidences involving the new vertices are chosen. All original
data and the zero address are preserved, and the same fixed splitting
and neutral external-incidence conditions may be imposed as above.
If an odd limit symbol is present, this is first a statement about
the finite typed graphs. It does not fix rational prime representatives
before the eventual theta stage.
\end{lemma}
\begin{proof}[Proof of the two address lemmas]
The forest and network calculation in
Proposition~\ref{ex:app-simultaneous-address} proves both statements.
For the stationary case its grounded-network columns attach to a
fixed terminal, so no terminal activation anchors are required.
For variable terminals the stated anchors give the row-sum identity
on every active restriction. Singular networks isolate prescribed
derivatives at one nonzero address; varying networks supply the
corresponding difference when two symbols share a variable. The
two-symbol proof takes the product of the top forms in the ordinary
polynomial ring before Boolean reduction. Its nonzero coefficient
forces a single assignment that works at all nonzero addresses.
It does not require the two symbols to have a common nonzero
evaluation before the auxiliary vertices are added.

Each construction uses a finite number of vertices and coefficient
readings once $b$ and the terminal polynomials are fixed. For an odd
symbol, include all of them in one finite weight and operator budget,
then use Proposition~\ref{ex:app-odd-stabilization} to choose one actual
theta stage carrying those tables and divisions. Only afterward use
Lemma~\ref{ex:sec7-realization} to choose its primes. Any old
prime representatives already chosen at that stage may be retained
when they have its full prescribed types and bits, by the same
realization lemma. In the even stationary application all original
primes can be retained from the outset. The number of
auxiliaries and the required theta precision may depend on $b$ and
the terminal data. The source's depth and the valuation constants
attached to a unit do not change. No common finite stage for all
values of $b$ is asserted.
\end{proof}

\paragraph{The extra upper-bound argument.}
Nonvanishing alone does not bound the valuation of $L_*$, and (3.2)
by itself gives a lower bound. The needed upper bound uses a separate
even-symbol selection and transfer argument, which we now give.

\begin{proposition}[The partitionable even upper estimate]\label{ex:sec7-even-upper}
Fix one even coefficient source, its signed unit filter and its depth.
If a parameter $d$ in that filter has a nonempty partitionable prime
list and $E^{(d)}$ has full two-primary Selmer corank zero, then (3.3b)
holds. Its constant depends only on the fixed source, not on the
number or sizes of the primes of $d$ or on later finite exclusions.
\end{proposition}
\begin{proof}\leavevmode
For the even coefficient system described below, a unit symbol at a
parameter $d$ gives
\[
       L_*(d)\ne0,\qquad v_2L_*(d)\le2\omega(d)+C_0.              \tag{3.3a}
\]
This follows from its same-filter coefficient-square formula and its
fixed normalization. That even normalization uses the lower bound
from (3.2), not the as-yet undefined odd minimum.
Partitionability is the nonempty contraction-certificate condition in
Definition~\ref{ex:sec7-seed-partition}. By
Lemma~\ref{ex:sec7-seed-certificate}, its iterated finite difference
is nonzero and hence so is the terminal polynomial. An empty unit,
when present, satisfies its coefficient bound directly and does not
require the stationary address lemma.

For a partitionable parameter $d$ already known to have Selmer
corank zero, keep all its prime factors stationary.
Lemma~\ref{ex:sec7-one-symbol-address} gives arbitrarily large binary
cubes based at
$d$, with no old prime removed, whose other addresses satisfy (3.3a).
The fixed-form common scalar in (3.2) has one uniformly cleared
denominator on every such cube. Choose a two-adic precision larger
than the nonzero-address scalar valuation bound and its fixed
clearing exponent. Its sign-evaluation coefficient is a Boolean
function with the bounded bit degrees proved in
Section~\ref{sec:limits}. A cube dimension larger than their sum
therefore supplies a nonzero address congruent to the base at this
precision. If the base scalar were zero or had larger valuation than
the bound, so would that nonzero address, a contradiction.
The corank-zero central comparison in (3.2) consequently gives
\[
 L_*(d)\ne0,\qquad
             v_2L_*(d)-2\omega(d)-s(d)\le C_1.                  \tag{3.3b}
\]
The constants are fixed by the even system. The dimension needed
may depend on $d$; the valuation bound does not. This is the
cyclotomic missing-address argument on the fixed-form Kato scalar,
not the later Heegner transfer.
\end{proof}

\begin{proposition}[A bounded companion in the prepared filter]\label{ex:sec7-rough-companion}
For (3.3), restrict $h$ to positive squarefree parameters in the
prepared odd-sign unit filter containing one, prime to its fixed
support. Thus $h\equiv1\pmod8$ and $w(E^{(h)})=-1$.
The simple-zero parameters used in (3.4)--(3.6) lie in this filter.
There are constants $A,C_E$ fixed before $h$ varies such that,
outside any prescribed additional finite avoidance set, one can
choose a negative odd fundamental discriminant $k'$ which is a
square at every place dividing $2Nh$, avoids the extra-unit fields,
and has $L_*(hk')\ne0$. It satisfies (3.3), in particular
$\omega(k')\le A$. The valuation estimate uses the fixed even
sources of the original $E$, although the analytic nonvanishing
theorem will be applied to an auxiliary twist of $E$.
\end{proposition}
\begin{proof}\leavevmode

Here is why the actual rough companion is partitionable.
Prepare an even symbol for every possible total type at the finite
fixed support. There are finitely many. In our full rational
two-torsion case the alphabet is all split. A least nonempty unit
seed has no dependence on internal edge bits: an edge derivative
would be a smaller nonempty unit configuration. Its individual
labels determine it. The negative-discriminant parity at two makes
the total label nonzero, so this seed is nonempty.
A list with the same total label and enough copies of the seed
labels contracts to it: keep the first $r-1$ seed anchors and merge
all remaining vertices into one nonempty cluster. Its label is
forced to be the last seed label. Thus the list is partitionable,
regardless of the number or labels of the remaining vertices.

Before $h$ varies, fix a bounded multiset of dummy labels supplying
the required anchors and reserves for every prepared total type.
For each $h$, realize these labels by distinct primes avoiding the
fixed support, the support of $h$, and all already fixed exclusions.
Put in $B$ the entire fixed support, the support of $h$, the chosen
dummy primes, and every other already fixed forbidden prime.
Include any already fixed reduction prime in $B$ as well.
Prescribe the total $k'$ to be a negative odd fundamental
discriminant, square at every place of $2Nh$ and outside the
extra-unit cases. Absorb $h$ and the dummy product into the curve
and use the preceding nonvanishing theorem in its even filter.
It adds at most five new primes. For the actual resulting total
type, the already chosen dummy supply makes $hk'$ partitionable.
Its analytic nonvanishing gives Selmer corank zero by the forward
rank-zero theorem. Apply (3.3b), obtaining
\[
 \omega(k')\le A,\qquad
 v_2L_*(hk')\le 2\omega(h)+2\omega(k')+s(hk')+C_E.                \tag{3.3}
\]
Here $A$ is the fixed dummy count plus five. This is the argument
of the partitionable-even and rough-companion lemmas, with their
stationary primes and order of choices retained.
Neither prime sizes nor a uniform analytic threshold are asserted.
Excluding a later coefficient prime does not enlarge the old
Selmer support.
The finite source preparation and the distinction between fixed dummy
labels and their later prime realizations are also detailed in
Proposition~\ref{ex:app-rough-companion}.
\end{proof}

\subsection{The genus lower bound precedes normalization}
Fix a negative genus partner $h_*$ with $L_*(h_*)\ne0$.
Here $h_*$ is chosen odd fundamental with $|h_*|>4$ and in the
same squareclasses at $2N$ as the prepared positive filter containing
one. Throughout (3.4)--(3.6), $h$ belongs to that filter and is
prime to the fixed support, including $h_*$; the companion $k'$
is supplied by the preceding construction with the primes of $h_*$
also in its avoidance set.
For a known simple-zero $h$, let $P_h^{\rm gen}$ be the unaveraged
conductor-one genus sum over $\Q(\sqrt{hh_*})$.
Let $P_{h,k'}^{\rm rc}$ be the unaveraged conductor-$|h|$ ring-class
sum over $\Q(\sqrt{k'})$, weighted by $\chi_h$.
They are different points in different native lattices.
Write $j_{\rm gen}(h)$ and $j_{\rm rc}(h,k')$ for their index valuations.

\paragraph{The native genus lattice.}
For the genus point specifically, put
$K_h=\Q(\sqrt{hh_*})$ and $F_h=\Q(\sqrt h,\sqrt{h_*})$.
If $h>1$ and $\tau$ generates $\operatorname{Gal}(F_h/K_h)$, set
\[
 \Gamma_h=\{P\in E(F_h):\tau P=-P\},\qquad
 \Lambda_h=\Gamma_h/\Gamma_{h,\mathrm{tors}}.
\]
For $h=1$, set $\Gamma_1=E(K_1)$ and take its torsion-free quotient.
The native genus index is the divisibility index of the class of
$P_h^{\rm gen}$ in $\Lambda_h$, at the known simple-zero parameters
where this lattice has rank one. The CM reciprocity action and
reindexing the unaveraged genus-character sum give exact membership
$P_h^{\rm gen}\in\Gamma_h$, before multiplication by the later fixed
multiplier. The standard $h$-twist isomorphism identifies $\Gamma_h$
with $E^{(h)}(K_h)$. The remaining involution of $K_h/\Q$ fixes its
free quotient: the only other rational character component is the
$h_*$-component, whose rank is zero by the supplied forward theorem.
Thus $Q+\sigma Q\in E^{(h)}(\Q)$ represents twice the free class of
$Q\in E^{(h)}(K_h)$. The rational subgroup consequently contains
$2\Lambda_h$ and has index at most two. For $h=1$ this is ordinary
quadratic conjugation and the same norm argument. This calculation
uses the stated rank and CM inputs; it does not infer rank from
nonvanishing height, nor apply to the whole group $E(F_h)$.


\begin{proposition}[The genus bound before odd normalization]\label{ex:sec7-genus-lower}
Use $h_*,h,k'$ and the two native index valuations defined in the
preceding paragraph, with $h$ a known simple-zero parameter and $k'$
as in Proposition~\ref{ex:sec7-rough-companion}. Then (3.4)--(3.6)
hold. In particular their constant in (3.6) is independent of $h$.
The arithmetic lower estimate (3.5) holds uniformly for the admissible
companions with $\omega(k')\le A$; it is only asserted here at these
known simple-zero parameters.
\end{proposition}
\begin{proof}\leavevmode
The explicit height formula \cite[Theorem 1.1]{cst} applies twice.
All level primes split, conductors are prime to $N$, extra-unit fields
are excluded, and the parametrization is the same. Its conductor factor
is $c\sqrt{|\operatorname{disc}K|}$ and its sum is unaveraged.
The two Rankin products have the common factor $L(E^{(h)},s)$.
Each rational free line has index at most two in its native rank-one
lattice: conjugation fixes the free class, and the norm makes twice the
class rational. Thus
\[
 2j_{\rm rc}(h,k')-2j_{\rm gen}(h)
   =v_2L_*(hk')-v_2L_*(h_*)+\epsilon,\qquad |\epsilon|\le2,        \tag{3.4}
\]
up to the one fixed initial period convention.

The arithmetic estimate of Section~\ref{sec:limits}, for each fixed $A$,
is
\[
 2j_{\rm rc}(h,k')\ge
 4\omega(h)+s(h)+s(hk')-C_R(E,A),\qquad \omega(k')\le A.           \tag{3.5}
\]
Combining (3.3)--(3.5) cancels $s(hk')$ and gives
\[
 2j_{\rm gen}(h)\ge 2\omega(h)+s(h)-C.                            \tag{3.6}
\]
This calculation uses no odd normalization.
For the forward reference in (3.5), combine
Proposition~\ref{ex:limits-uniform-lower-bound} with the same-class
central identity in Proposition~\ref{ex:arith-central-calibration}.
Those arguments use the known nonzero point in the present setup;
they do not use an odd coefficient minimum or the final family.
\end{proof}

\subsection{A fixed nonzero witness and actual odd forms}
Choose one positive squarefree $D_w$ in the odd filter containing one,
after adjoining the primes of $h_*$ and their compatible unit tests.
Section~\ref{sec:nonvan} supplies $L'(E^{(D_w)},1)\ne0$.
The explicit height formula makes $P_{D_w}^{\rm gen}$ nontorsion.
Its unwanted $h_*$ component has rank zero. One fixed multiplier $L$
kills the bounded-degree torsion and the projection denominator, so
$Q=LP_{D_w}^{\rm gen}$ is a rational nontorsion point on $E^{(D_w)}$.

Let $C$ be complex conjugation on its Tate module and choose a nonzero
integral row $b$ with $bC=b$. The Kummer translation image of $Q$ is a
full lattice. Indeed, if it vanished, its rational cocycle would descend
to the Tate image. A nonidentity central scalar $z$ gives
\[
 (z-1)c(g)=(g-1)c(z),
\]
so that descended cocycle is a coboundary, contrary to Kummer
injectivity. Irreducibility makes a nonzero translation span the Tate
space, and intersection with any fixed finite kernel preserves full
rank. Twisting preserves open image by restriction to an index-at-most-two
subgroup.

Choose one such translation $v$ with $2b(v)\ne0$.
Multiplying conjugation by this translation gives square translation
$(1+C)v$, detected by $2b(v)$ with one fixed valuation $c_w$.
The element is a single actual Galois element on the compositum.
It is trivial on $\sqrt{D_w}$ and nontrivial on $\sqrt{h_*}$ because
$D_w>0>h_*$. Its restrictions therefore already agree on intersections.
Finite Chebotarev, with a place chosen above the resulting prime, gives
arbitrarily advanced stages $(M_\theta,\ell_0)$ such that
\[
 \rho(\operatorname{Fr}_{\ell_0})\equiv C\pmod {2^{M_\theta}},
 \quad (D_w/\ell_0)=1,\quad(h_*/\ell_0)=-1,
 \quad v_2\lambda_{M_\theta}(\overline Q)=c_w<M_\theta.            \tag{3.7}
\]
Here $\lambda_{M_\theta}=b\kappa_{M_\theta}$ is the actual finite-field
Kummer map on $E(\F_{\ell_0^2})$. It has Frobenius sign $+1$.
All reductions use that same place. This is a fixed-witness statement,
not a bound for arbitrary points or arbitrary reduction primes.

The weighted ternary theta construction gives actual holomorphic
forms at these stages. At original level primes the root weight uses
the summand selected by the cyclic level structure. At two, the
integral spectral projectors of $(1+x)/2$ specify that summand; the
operator itself is not in general idempotent. The nilpotent weight at
$h_*$ and the division-algebra residue-line weight at $\ell_0$ have the
prescribed conjugation and uniformizer signs. Their product is even.
The root-character field is fixed before $\ell_0$ and $M_\theta$.

Appendix~\ref{app:odd-source} gives these functions, their exact
transformation rules, and the direct extraction using the form $Q/m$.
Poisson summation for the resulting compactly supported local functions
gives the scalar transformation law and holomorphy at every cusp.
A fixed stabilizer multiple, for example one divisible by
$|\operatorname{GL}_4(\F_3)|$, clears the finite automorphism groups.
For the Hecke recurrence, a split or irreducible nonzero reduction has
$1+(-n/p)$ invariant lines, a nonzero nilpotent has one, and zero has
$p+1$. The unique integral primitive nilpotent neighbor supplies the
remaining square term. Tracking the scaling characters gives the
recurrence with the same character convention as the coefficient
comparison.

For an admissible $D$ with $(D/\ell_0)=1$, the oriented embedding
count from the conductor-one argument Section~\ref{v15:cm-oriented},
with the retained geometric and CM inputs stated there, gives
\[
 a_{M_\theta}(D)
   =c_0\,2^{\omega(N)}\zeta_{M_\theta}(D)
       \lambda_{M_\theta}(\overline{LP_D^{\rm gen}})
       \pmod {2^{M_\theta}},                                    \tag{3.8}
\]
where $c_0\ne0$ is fixed and $\zeta$ is a permitted root-of-unity phase.
For admissible $D$ with $(D/\ell_0)\ne1$, including $\ell_0\mid D$,
the local Schwartz support instead gives $a_{M_\theta}(D)=0$.
Formula (3.8) is asserted only on the first support.
There is no division by a class number. Formula (3.7) makes $D_w$
visible with fixed upper valuation in (3.8). It remains on every
stage's support, including the extra test $(D_w/\ell_0)=1$.

After phase removal, the coefficients used for valuation are integers.
Equation (3.6) gives a uniform lower bound for
$v_2a_{M_\theta}(D)-\omega(D)$ on visible bounded-weight sequences.
Take the finite integer values constant on a set belonging to the
chosen nonprincipal ultrafilter. This set of values is nonempty by
$D_w$, is bounded below, and hence has an attained minimum $\delta$.
If the resulting lower inequality failed for a fixed weight bound on
an ultrafilter-large set, choosing a violating index and then one of
finitely many smaller integer values would contradict minimality.
For any fixed finite operator budget, the coarse lower bound handles
the large weights and stabilization handles the bounded weights.
Consequently every required division is valid at one sufficiently
advanced actual stage. No infinite-precision modular form or
arbitrary-ring base-change assertion is needed.

\subsection{One common family and the final order of choices}
The even source is a fixed algebraic vector in the full good-Hecke
eigenspace. Waldspurger's Corollary 2, including its same-dyadic-squareclass
removal of the auxiliary condition, compares the actual squarefree
Fourier functionals \cite{waldspurger}. Local upper-unipotent projections
retain a nonzero initial functional; smoothness supplies lower-unipotent
invariance. The algebraic linear conditions choose a vector over a
number field. Galois conjugation followed by the matching finite
character twist shows that each normalized coefficient ratio is a
rational number times a two-power root of unity. This supplies the
valuation and nonvanishing comparison on the whole filter.
Only flat field extension of the algebraic model is used.
The comparison is proved at the original support $S_0$ before any extra
coprimality sieve $B$. The sieve is applied last, and every later good
recurrence prime avoids the full support $S_0\cup B$.

For both coefficient systems the integral operators have the required
omission and contraction rules. At a new prime $V_p$ raises the level
exponent only to one, with ramified quadratic character. The recurrence
makes the relevant differences divisible by the integer two before
residue reduction. Multiplying by ordinary theta, then taking a fresh
trace, removes the integral unary corrections. Two fresh traces vanish.
The matrix trace identity gives $t(gu^2)=t(g)$, hence elementary
quadratic labels. At a new odd level prime the weight-two local representation
has one unramified and one ramified line. The normalized-induction
and Whittaker calculation below identifies the Fourier $U_p$
eigenvalue with arithmetic Frobenius on the unramified line.
The diagonal-matrix calculation gives the two inertia differences.
These identities prove omission and primary isolation for the same
original form and depth, without choosing another minimum.

If the intrinsic odd limit symbol is a unit at the empty configuration,
choose an actual stage carrying that value with $M_\theta>\delta_o$.
The index $D=1$ lies in every moving support and then satisfies
$v_2a_{M_\theta}(1)=\delta_o<M_\theta$. The visibility implication
(O.8) gives $\ord_{s=1}L(E,s)=1$, which is exactly (RC).
No final companion or determinant bound is needed in this case.
For the rest of the final-family and missing-address argument, assume
that the empty configuration is not a unit. The bounded-weight witness
then supplies a least nonempty odd seed, as in
Appendix~\ref{app:endpoints}.

\paragraph{Fixing the final companion.}
Use a separate even symbol on negative discriminants square at $2N$,
excluding the fixed odd support, the ramification support of $E[8]$,
and the exceptional imaginary fields. Choose a unit $k$ of this
symbol. Then $L(E^{(k)},1)\ne0$. Fix $k$ and $K=\Q(\sqrt{k})$
before choosing the binary dimension. At least one prime ramified
in $K$ is outside all the old fixed label extensions and torsion
ramification. Choose every later moving theta prime $\ell_0$ away
from $k$. The intrinsic odd symbol continues to omit the primes of $k$.

At any such stage, the normal closure of the old odd label extension
is unramified outside its old support and $\ell_0$. The fresh ramified
prime of $K$ lies outside this set, so its intersection with $K$ is
$\Q$. Thus each old odd label still has a lift split in $K$.
The added stationary tests at the primes of $k$ may have relations;
include those relations in the common alphabet, rather than requiring
independent arbitrary prescriptions. Stabilize its separate projection
maps on the same ultrafilter. This preserves the old odd minimum.
Fresh ramification also makes $K$ independent of $\Q(E[8])$ and of
the old fields $\Q(i)$ and $\Q(\sqrt{\ell})$.

The forward theorem at this known nonzero central value gives
$s_2(E^{(k)}/\Q)=0$. The sum-and-difference isogeny over $\Q$ gives
\[
 s_2(E/K)=s_2(E/\Q)+s_2(E^{(k)}/\Q)=1.
\]
Indeed its kernel and cokernel on the relevant restriction and
corestriction maps are killed by a fixed power of two. Tensoring the
compact Selmer modules, equivalently the Pontryagin duals of the discrete
Selmer groups, with $\Q_2$ gives the direct sum and hence this identity
of coranks.
No integral direct-sum assertion is used.

\begin{lemma}[The nonempty seed for the present curve]\label{ex:sec7-ksplit-seed}
For $E=E_\ell$ and the $k,K$ just chosen, the positive even source
for $E^{(k)}$, normalized at index one, has a nonempty unit seed
consisting of $K$-split primes.
\end{lemma}
\begin{proof}\leavevmode
This is the specialization to the roots $0,\ell,-3\ell$ of the
coefficient calculation in Proposition~\ref{ex:app-nonempty-seeds}.
If no such nonempty unit existed, its integral $p^3$ coefficient
and a fresh trace would force
$a_p(E^{(k)})\equiv p+1\pmod8$ at every fresh good $K$-split
prime outside the support, whether or not it lies in the coefficient
filter. Indeed, off the filter $a(p)=0$ in (G.4a); the second fresh
prime $p'$ matches the old labels so that $pp'$ lies in the filter,
as used in (G.4d). The selected $p$ below need not lie in the filter.
The cited calculation uses positive Fourier indices and the positive
twist character convention; it derives this congruence inside the
assumed absence of nonempty units.

Choose such a prime with $p\equiv1\pmod4$, $(\ell/p)=-1$ and
$(k/p)=1$, avoiding the finite support. In fact this $p$ is outside
the square filter: quadratic reciprocity and $p\equiv1\pmod4$ give
$\prod_{r\mid\ell}(p/r)=(\ell/p)=-1$, so $p$ is a nonsquare at
some odd $r\mid\ell$, and $r\mid N$. This is why the off-filter
scope of the preceding calculation is needed. The fresh ramification of $K$
permits these simultaneous conditions. Because $k$ is a square modulo
$p$, the reductions of $E^{(k)}$ and $E_\ell$ are isomorphic over
$\F_p$. None of the three nonzero rational points of order two can
be doubled: their necessary square differences include, respectively,
$-\ell,\ell,-4\ell$. Each is a nonsquare under the chosen conditions.
The two-primary reduction group therefore has order exactly four,
contradicting the asserted congruence. This proof uses the particular
curve $E_\ell$; no nonempty-seed theorem for an arbitrary elliptic
curve is being asserted.
\end{proof}

We record the finite combinatorial realization used here. Each label
space is a finite elementary two-group; the common alphabet projects
onto each separately. Product surjectivity is unnecessary. Take
sufficient repeated copies of every label, with total label zero.
Keeping the seed anchors, contract the unused copies into the last
anchor. Minimal positive seed size makes the internal bits irrelevant,
so each of the two terminal polynomials is nonzero. Grounded Laplacian
minors express the constant and pair-contraction operators, including
singular matrices. A cyclic singular network isolates its target
radical and a path network isolates its selected pair coefficient.
In the product of the two top forms, taken in the ordinary polynomial
ring before Boolean reduction, a variable of exponent two uses one
singular program in each factor; exponent one uses one varying program.
The total degree kills the other contractions and leaves coefficient
one. Thus one choice works for both forms. CRT and Dirichlet realize
the prescribed forward residue bits by distinct primes; quadratic
reciprocity supplies the reverse bits.

\begin{proposition}[One family for the two coefficient tests]\label{ex:sec7-final-family}
Keep the fixed $E,h_*$, normalized odd symbol and final $k,K$ of the
preceding paragraphs. Assume, as in the remaining branch above,
that the intrinsic odd empty configuration is not a unit.
In this remaining branch, for every sufficiently large $b$ there is one
finite set $Q$ of distinct good odd primes split in $K$ and outside the
fixed support, and nonzero linear forms $\lambda_q:\F_2^b\to\F_2$,
such that the family
\[
 h_x=\prod_{q\in Q}(q^*)^{\lambda_q(x)}
 \quad(x\in\F_2^b),\qquad q^*=(-1)^{(q-1)/2}q,
\]
has $h_0=1$, all fixed local types unchanged, and, for $x\ne0$,
\[
 \ord_{s=1}L(E^{(h_x)},s)=1,\quad L_*(h_xk)\ne0,\quad
 j_{\rm gen}(h_x)\le\omega(h_x)+C_o,\quad
 v_2L_*(h_xk)\le2\omega(h_x)+C_e.                                \tag{3.9}
\]
The constants are independent of $b$. By (3.4), now with fixed $k$,
$2j_{\rm rc}(h_x,k)\le4\omega(h_x)+C_J$.
\end{proposition}
\begin{proof}\leavevmode
The packet argument just given supplies separate nonzero terminal
polynomials in one common alphabet. Apply
Lemma~\ref{ex:sec7-two-symbol-address}. The full proof in
Proposition~\ref{ex:app-simultaneous-address} and
Proposition~\ref{ex:app-final-family} preserves the projected total
labels, all stationary tests and the original two normalizations.

The order is fixed data and $D_w$, then the genus lower bound and odd
minimum, then the final $K=\Q(\sqrt k)$ and common alphabet, then
determinant precision $M$, then $b$ and its finite operator budget,
then a sufficiently advanced actual theta stage $M_\theta$ and its
prime realization. In particular $M_\theta$ is not $M$.
The finite graph construction, its local conditions and the coefficient
tests all occur at that same final stage. Appendix~\ref{app:endpoints}
spells out the absolute five-prime completion, normalization constants,
and the final use of the one corrected coordinate.
\end{proof}

% End 03-restricted-converse.tex


% Begin 03a-arithmetic.tex
\section{Actual classes, local conditions and central comparison}
\label{sec:arithmetic}
Let $K=\Q(\sqrt k)$ be imaginary, split at $2N$, and different from
$\Q(i)$ and $\Q(\sqrt{-3})$. Active primes are good odd primes split
in $K$. Derivative primes are good inert primes outside all current
supports. The actual finite coefficients are
\[
 A_i=(\Z/2^i\Z)[t]/t^i,\quad
 M_i=(T_2E\otimes A_i)(\chi_h\Psi_i).
\]
For a fixed binary group $G$, use $A_i[G]$ and the universal quadratic
character. The character $\Psi_i$ is finite, trivial at $2N$, ramified
at a moving conductor prime $r_i$, and has prescribed odd scalar
exponents at the finite active list. In a binary family let $Q$ be
the union of its prime supports. Apply
Proposition~\ref{ex:arith-moving-character} to an orientation above
every $q\in Q$, prescribing $\xi_i(\operatorname{Fr}_q)=1$ at each
finite stage. This prescription includes the primes unused by an
individual character, including every graph prime at augmentation.
The same $\Psi_i$ is used for all characters of the family. The finite
set $Q$ is fixed before choosing the sequence of moving primes $r_i$;
the trace bound $c_E$ does not depend on its size. The prime $r_i$
varies with precision. There is no assertion
that one fixed odd prime supports an infinitely ramified $\Z_2$ character.

\begin{definition}[Coefficient and local-condition conventions]
\label{ex:arith-data}
In Sections~\ref{sec:arithmetic} and~\ref{sec:limits}, retain the fixed
non-CM curve $E$, parametrization $\pi$, Tate lattice, and old local
types of Section~\ref{sec:rc}. Put $\Lambda=\Z_2[[t]]$ and
$\mathscr O=\widehat{\Lambda_{(2)}}$. A binary family uses one finite
elementary abelian two-group $G$, one real sign and common local
squareclasses at $2N$. Its universal varying quadratic character
is therefore trivial at the real place and at $2N$.
The active primes are good odd primes split in the indicated imaginary
quadratic field $K$; the moving character is the finite character
specified above, trivial at $2N$. The primitive and relaxed complexes
use the actual derived Kummer, unramified, full and orthogonal local
conditions specified below, with their maps and homotopies.
At derivative primes the four local conditions are those of
Lemma~\ref{ex:limits-derivative-switch}. Every determinant coordinate
uses the retained integral determinant basis.
\end{definition}
\begin{lemma}[Conjugate coefficients and local duality]\label{ex:arith-conjugate-duality}
For the finite coefficient module $M_i$ and coefficient ring
$R_i=A_i$ or $A_i[G]$ above, let $\tau$ be an actual complex
conjugation. Pullback of the Galois action gives an $R_i$-linear
identification $\tau^*M_i\simeq M_i^*(1)$ through
$(v,w)\mapsto e_2(v,\rho_T(\tau)w)$; it fixes $t$ and the group-ring
basis. With the conjugate places, actual orthogonal local conditions
and their isotropy homotopies, finite Selmer duality gives
\[
 C_U(M_i)\simeq
 R\operatorname{Hom}_{R_i}(C_{U^\perp}(M_i^*(1)),R_i)[-3].
\]
The local Kummer and derivative comparisons used here are proved
below and in Lemma~\ref{ex:limits-derivative-switch}.
\end{lemma}
\begin{proof}\leavevmode
\paragraph{The dual coefficient system.}
Fix an actual complex conjugation $\tau\in G_{\mathbf Q}$, so that
$\tau^2=1$. The moving character is constructed on the ring-class group:
\[
 \xi_i:\operatorname{Gal}(H_{r_i}/K)\longrightarrow
       \mathbf Z/2^{n_i}\mathbf Z,
 \qquad \Psi_i=(1+t)^{\xi_i},
\]
where $n_i$ is large enough that $(1+t)^{2^{n_i}}=1$ in $A_i$.
Conjugation acts on this ring-class group by inversion. Consequently
\[
 \Psi_i(\tau g\tau^{-1})=\Psi_i(g)^{-1}\qquad(g\in G_K).
\]
Let $R_i=A_i$ in the scalar construction and $R_i=A_i[G]$ in the
binary construction. Write $\alpha_i=\chi_h\Psi_i$, respectively
$\alpha_i=\chi_0\varepsilon\Psi_i$, where $\varepsilon:G_K\to G$
is the restriction of the universal rational quadratic character.
Then $\varepsilon^\tau=\varepsilon=\varepsilon^{-1}$ and
$\alpha_i^\tau=\alpha_i^{-1}$. Here $\tau^*M_i$ denotes pullback
of the Galois action; it does not act on $R_i$.

The Weil pairing gives the following $R_i$-linear perfect pairing:
\[
 M_i\otimes_{R_i}\tau^*M_i\longrightarrow R_i(1),\qquad
 (v,w)\longmapsto e_2(v,\rho_T(\tau)w).
\]
Indeed $\rho_T(\tau)^2=1$, the two scalar characters multiply to
one, and $e_2$ transforms by the cyclotomic character. Thus
$\tau^*M_i\simeq M_i^*(1)$. We do not identify $M_i$ with its
ordinary dual when $\Psi_i$ is nonquadratic. In this convention $t$
and every element of the group-ring basis are fixed. Inverting $1+t$
is already accounted for by pullback of the character.

At a split place $v$, the dual local condition is transported from
the conjugate place $\tau v$. At a fixed place it is transported
by the actual local representative of conjugation. Kummer conditions
at $2N$ are preserved by this transport. A full condition is paired
with the zero condition; the full local complex is replaced by zero
only at the localizations where its displayed contraction exists.
At a derivative prime the finite and singular planes are checked
with this transported pairing. In the chosen coordinates this is
the pairing $\beta(x,y)=\epsilon\langle x,Jy\rangle$.
The $J$ is therefore part of the duality convention, not an extra
identification of an arbitrary twist with its dual.

Use the coefficient-valued pairing in the finite Selmer diagrams,
with the conjugate local places and their isotropy homotopies. The
resulting duality is
\[
 C_U(M_i)\simeq
 R\operatorname{Hom}_{R_i}
       (C_{U^\perp}(M_i^*(1)),R_i)[-3].
\]
These are the maps and homotopies retained in both limits.
\end{proof}
% Source reconstruction of the finite arithmetic input I08/I09.
% This is a proof fragment, not an assertion about completion of RC.
\subsection{Finite Heegner descent and the derivative identities at two}

Fix the non-CM elliptic curve $E/\mathbf Q$, its conductor $N$, its
Tate module $T=T_2E$, and a modular parametrization
$\pi:X_0(N)\longrightarrow E$. Normalize $\pi(\infty)=0$.
Let $K$ be imaginary quadratic, split at every prime dividing $2N$,
and different from $\mathbf Q(i)$ and $\mathbf Q(\sqrt{-3})$.
The conductor $f$ below is prime to $2ND_K$; it may contain all the
active primes in a finite binary family and the moving prime.
Write $H_f$ for the ring-class field of the order
$\mathcal O_f=\mathbf Z+f\mathcal O_K$.

Let $I$ be any finite set of good inert primes outside $6fND_K$,
and put $n_I=\prod_{\ell\in I}\ell$ and $M_I=H_{fn_I}$.
At the end, a finite character $\theta$ of
$\Delta_f=\operatorname{Gal}(H_f/K)$ will take values in
$A^\times$, where $A=(\mathbf Z/2^m\mathbf Z)[t]/t^b$, or in
$A[G]^\times$ for a finite binary group $G$.
It has the form $\chi_h\Psi$, or the corresponding universal quadratic
character times $\Psi$. At $2N$, $\Psi$ is trivial and the universal
quadratic character has the fixed local type already specified in
the manuscript. These are hypotheses on the actual finite character.
Its construction is separate from the descent proved here.

\begin{proposition}[Finite Heegner descent with one multiplier]\label{ex:arith-finite-heegner}
With $E,\pi,K,f,I,M_I,A$ and the finite character $\theta$ as just
specified, there are integers $e\geq0$ and $L>0$ depending only
on the fixed curve, parametrization and allowed old local types,
and on any other prescribed fixed multiplier requirements, with
the following properties. Form the derivative over the complete
extension $H_{fn_I}/H_f$ and descend it before applying the character.
This gives finite classes $P_{I,m}$ with the old Kummer conditions,
the stated good unramified conditions and full active and moving
conditions. For $I=\varnothing$ the class is the Kummer class of
the $L$-multiplied unaveraged weighted point.

The unaveraged norms are (4.1). If $a\geq m+e$ and every new
derivative prime has $\rho_T(\operatorname{Fr}_\ell)\equiv
\rho_T(\tau)\pmod{2^{a+1}}$, the classes projected to precision
$m$ satisfy (4.2). These assertions hold over $A[G]$ before
character evaluation. They are realized by finite continuous
cocycles, local-condition representatives and boundary witnesses,
with the weighted coefficient map applied to all this data.
There is no multiplier depending on $K,f,I,m$, the truncation
length, the number of active primes or $|G|$.
\end{proposition}
\begin{proof}\leavevmode
\paragraph{1. One uniform torsion annihilator.}
By the non-CM open-image theorem \cite{serre}, choose $n\ge2$ such that
$1+2^nM_2(\mathbf Z_2)\subset\rho_T(G_{\mathbf Q})$, and choose
$g\in G_{\mathbf Q}$ with $\rho_T(g)=uI$, $u=1+2^n$.
For an actual complex conjugation $\tau$ set
\[
 z=g^2\tau g^2\tau^{-1},\qquad e=v_2(u^4-1)=n+2.       \tag{H1}
\]
The element $g^2$ fixes every quadratic field. Its image in the
abelian group $\operatorname{Gal}(H_d/K)$ is inverted by $\tau$,
for every $d$ and every allowed $K$. Thus $z$ fixes every such
$H_d$, while $\rho_T(z)=u^4I$. If $P\in E(H_d)[2^\infty]$,
then $(u^4-1)P=0$, so
\[
                       2^e E(H_d)[2^\infty]=0.          \tag{H2}
\]
This bounds the whole invariant torsion group, with no restriction
on the number of conductor primes. It also holds for a compositum
of the ring-class fields in a finite construction, since $z$ fixes
each of them.

For the matrix-span statement used elsewhere, square preimages of
$I+2^nE_{11}$, $I+2^nE_{12}$ and $I+2^nE_{21}$. The squares belong
to every quadratic subgroup, and their images, together with $I$,
are
\[
 I,\quad I+(2^{n+1}+2^{2n})E_{11},\quad
 I+2^{n+1}E_{12},\quad I+2^{n+1}E_{21}.
\]
Their $\mathbf Z_2$-span contains $2^{n+1}M_2(\mathbf Z_2)$:
subtract $I$, use that $1+2^{n-1}$ is odd, and obtain $E_{22}$
from $I-E_{11}$. This bound also depends only on $E$.

\paragraph{2. Point relations with the full old conductor.}
Choose the CM points $x_d\in X_0(N)(H_d)$ by compatible proper
ideal lattices and the fixed orientations at $N$, and put
$y_d=\pi(x_d)$. The split Heegner hypothesis identifies the curve
$X_K(N)$ in Cai--Chen--Liu, Theorem 2.15 \cite{ccl}, with $X_0(N)$:
$N_1=1$, and the condition $(*)$ reduces to $(d,N)=1$.
Take their old conductor to be any $d$ currently in use and the
new prime outside $6dND_K$. Their factor $u_m$ is one because
$\mathcal O_K^\times=\{\pm1\}$. Applying the elliptic quotient
to that theorem gives
\[
 \operatorname{Tr}_{H_{d\ell}/H_d}y_{d\ell}=a_\ell y_d
       \quad(\ell\text{ inert}),\qquad
 \operatorname{Tr}_{H_{dq}/H_d}y_{dq}
       =(a_q-\sigma_q-\sigma_q^{-1})y_d
       \quad(q\text{ split}).                            \tag{4.1}
\]
The inverse Frobenius pair at a split prime follows because their
product is the Artin symbol of the rational principal ideal $(q)$.
These are unaveraged sums. On degree-zero divisor classes,
$\varphi T_q=[a_q]\varphi$. Since $q\nmid N$,
$T_q[\infty]=(q+1)[\infty]$, so the displayed norm relations
for $y_d=\pi(x_d)$ have no cusp correction.
For other changes of CM convention, let $\mathcal C_N$ be the
subgroup of $J_0(N)(\overline\Q)$ generated by
$[c]-[\infty]$ for all cusps $c$. This is a finite group by
Manin--Drinfeld \cite[p.~59, Theorem]{elkik}.
Fix $C_\pi\geq1$ annihilating $\varphi(\mathcal C_N)$.
This single integer depends only on $N,\pi$.
For an Atkin--Lehner involution $W$ acting on $f$ by $w$, the exact
formula is $\pi(Wx)=[w]\pi(x)+\pi(W\infty)$; the final constant
lies in $\varphi(\mathcal C_N)$ and vanishes after multiplication
by $C_\pi$.

At an inert new prime $\ell$, a compatible place gives
\[
                      \overline y_{d\ell}
                         =F_\ell\overline y_d.           \tag{H4}
\]
Here $F_\ell$ is rational arithmetic Frobenius on the reduced
elliptic curve. This is the usual CM congruence, and it does not
require the old conductor to be an inert squarefree product.
Indeed $(d,\ell)=1$ makes the local CM order the maximal
unramified quadratic order. The CM elliptic curve has supersingular
reduction at this inert prime. Each of its $\ell+1$ cyclic raising
isogenies specializes to the quotient by the unique subgroup of
order $\ell$, namely the relative Frobenius kernel. The prime-to-$\ell$
cyclic $N$-subgroup specializes by the same Frobenius map.
The resulting equality of reduced modular points passes through
$\pi$, which extends over the good integral models at $\ell$.
This also proves (H4) after every old ideal-class translate, with
its place transported. Howard, Section 1.7, formula (13) \cite{howard2012}, records
this congruence in his conductor subfamily; his odd-prime descent
and residual-surjectivity theorem are not used here.

\paragraph{3. The relative group and invariant derivative class.}
The idelic description
$\operatorname{Pic}(\mathcal O_d)=
K^\times\backslash\mathbf A_{K,f}^\times/\widehat{\mathcal O}_d^\times$
identifies the kernel of
$\operatorname{Pic}(\mathcal O_{fn_I})\to
\operatorname{Pic}(\mathcal O_f)$ with
\[
 \prod_{\ell\in I}
 \frac{(\mathcal O_K/\ell\mathcal O_K)^\times}
      {(\mathbf Z/\ell\mathbf Z)^\times}
       \simeq\prod_{\ell\in I}C_{\ell+1}.               \tag{H5}
\]
There is no additional quotient by units: the only global units
are $\pm1$ and already belong to the denominator at every factor.
This proves (H5) without choosing a splitting of the full class
group. Let $\sigma_\ell$ be a generator of its $\ell$-factor,
with the local Artin orientation fixed. Set
\[
 D_\ell=\sum_{j=1}^{\ell}j\sigma_\ell^j,\qquad
 D_I=\prod_{\ell\in I}D_\ell,\qquad
 (\sigma_\ell-1)D_\ell=(\ell+1)-\operatorname{Tr}_\ell. \tag{H6}
\]
If $2^a$ divides $a_\ell$ and $\ell+1$ for all $\ell\in I$,
(4.1) and (H6) show that
\[
 c_{I,a}=\delta_a(C_\pi D_Iy_{fn_I})
 \in H^1(M_I,E[2^a])^{\operatorname{Gal}(M_I/H_f)}.     \tag{H7}
\]
For example its difference under $\sigma_\ell$ is the Kummer
class of
$C_\pi D_{I\setminus\{\ell\}}
 ((\ell+1)y_{fn_I}-a_\ell y_{fn_I/\ell})$, an actual
$2^a$-multiple of an $M_I$-point. The factors in (H5) commute.

\paragraph{4. Descend once, before twisting.}
Put $B=H_f$, $M=M_I$, $\Gamma=\operatorname{Gal}(M/B)$ and
$V_a=E[2^a]$. The five-term sequence for finite discrete coefficients
\cite[formula (4.1.4.1), p.~96]{nekovar} gives
\[
 0\longrightarrow H^1(\Gamma,V_a^{G_M})
 \longrightarrow H^1(B,V_a)
 \xrightarrow{\mathrm{res}} H^1(M,V_a)^\Gamma
 \longrightarrow H^2(\Gamma,V_a^{G_M}).                \tag{H8}
\]
By (H2), multiplication by $2^e$ is zero on the coefficient
module $V_a^{G_M}$ and hence on every cochain and cohomology
group of $\Gamma$ with those coefficients. For an invariant class
$c$, choose $b$ with $\mathrm{res}(b)=2^e c$ and define
\[
                         d_a(c)=2^e b.                  \tag{H9}
\]
Any two choices of $b$ differ by the first group in (H8), so (H9)
is independent of the choice. It is additive, natural for
coefficient reduction, and equivariant for transport of the whole
field diagram. Its restriction is exactly $2^{2e}c$.
This is one application to $M_I/H_f$, not one application for
each factor of (H5). The coefficient module in (H8) is untwisted
$E[2^a]$; the scalar and binary coefficient push comes only later.
Thus no invariant module containing a varying nonquadratic scalar
character has been silently substituted in this argument.
For $I=\varnothing$, the relative group is trivial and this same
definition gives $d_a(c)=2^{2e}c$. In particular it uses the same
normalizing multiplier for the primitive class and its derivatives.

\paragraph{5. Integral local descent at the old places.}
We record the local fact needed to put $d_a(c_{I,a})$ in the old
Kummer conditions. If $F'/F$ is a finite unramified extension of
nonarchimedean local fields and $B/F$ is an elliptic curve, then
\[
 H^1(\operatorname{Gal}(F'/F),B_0(F'))=0,\qquad
 c_B H^1(\operatorname{Gal}(F'/F),B(F'))=0,             \tag{H10}
\]
where $c_B=\#\Phi_B(\overline{k}_F)$ is the geometric component
group order. To prove the first assertion, reduce a cocycle to
the connected smooth special fiber. Lang's theorem \cite[Theorem~1 and its corollary]{lang1956} makes its
$H^1$ over the finite residue field zero; inflation--restriction
then makes its $H^1$ for the finite residue extension zero.
Lift the resulting coboundary through the smooth reduction map.
The remaining cocycle takes values in the formal group. For a
uniformizer of $F$, fixed by the unramified Galois group, each
successive formal quotient is an additive copy of the residue
field of $F'$ \cite[IV.3.2(a), p.~124]{silverman}. A normal basis makes this an induced module, so
its positive-degree cohomology vanishes. Correct the coboundary
successively in these quotients. The corrections converge in
$B_1(F')$, and the group is finite, so their limit solves all
the cocycle equations. The component exact sequence proves the
second assertion of (H10). No division by $[F':F]$ was used.

If a class $v\in H^1(F,B[2^a])$ becomes Kummer over $F'$, its
image in $H^1(F,B)[2^a]$ belongs to the kernel of restriction.
That kernel is the image of
$H^1(\operatorname{Gal}(F'/F),B(F'))$. Thus $c_Bv$ is Kummer
over $F$, by the Kummer exact sequence and (H10).

All ring-class extensions in use are unramified at $2N$, since
the conductors are prime to $2N$. The base fields there are
unramified extensions of $K_v=\mathbf Q_p$. Geometric component
groups are unchanged under unramified extension. Let $C_{\rm bad}$
be a common multiple of their orders for the finitely many
allowed local quadratic types at $2N$. The same integer works
for every conductor, every place above $2N$ and every $K$.
Choose a fixed positive integer $m_0$, allowing the other fixed
multiplier requirements in the argument to be included, and define
\[
 L=m_0C_\pi C_{\rm bad}2^{2e},\qquad
 b_{I,a}=m_0C_{\rm bad}d_a(c_{I,a})\in H^1(H_f,E[2^a]). \tag{H11}
\]
Thus any common fixed multiple of this descent multiplier and
the detection multiplier is permitted; every class is scaled by
the same $m_0$.
Then $b_{I,a}$ is Kummer at $2N$ and restricts to
$\delta_a(LD_Iy_{fn_I})$. At a good place away from $2fn_I$,
the same assertion has component order one. It gives the
unramified local condition, since at good odd places the finite
Kummer image equals the unramified image. At active and moving
primes the required condition is full. The integer $L$ is now
fixed, independently of $K,f,I,m,b$ and $G$.

\paragraph{6. Exact finite precision for one derivative prime.}
Let the retained torsion precision be $m$, choose $a\ge m+e$,
and require, for every retained derivative prime,
\[
 \rho_T(\operatorname{Fr}_\ell)\equiv J\pmod {2^{a+1}},
 \qquad J=\rho_T(\tau),\quad J^2=I.                    \tag{H12}
\]
Then $2^{a+1}$ divides $a_\ell$ and $\ell+1$, since
$\operatorname{tr}J=0$ and $\det J=-1$.
Choose $\ell\notin I$, and perform the construction for
$I\cup\{\ell\}$. Apply the other derivatives first, writing
\[
 X=D_Iy_{fn_I},\qquad Y=D_Iy_{fn_I\ell}.
\]
The $K$-prime $(\ell)$ splits completely in $H_{fn_I}$:
it is prime to the conductor and is a rational principal ideal
in the ring-class Artin map. Thus its residue field is
$\mathbf F_{\ell^2}$, and
\[
 F^2\overline X=\overline X,\qquad
 \operatorname{Tr}Y=a_\ell X,\qquad
 \overline Y=F\overline X,\qquad F^2-a_\ell F+\ell=0.  \tag{H13}
\]
The same splitting makes the old class $b_{I,a}$ at this place
exactly the Kummer class of $LX$, by its restriction identity
in (H11); there is no local descent ambiguity for this old class.
For $2^aQ=LX$ its finite coordinate is
\[
 f_{\rm old}=(F^2-1)\overline Q
             =(a_\ell F-(\ell+1))\overline Q.           \tag{H14}
\]
Although the two integer coefficients are highly divisible, $Q$
is a division point, so (H14) does not assert $f_{\rm old}=0$.

Choose a continuous cocycle representing $b_{I\cup\{\ell\},a}$
and choose $Z$ with $2^a Z=LD_\ell Y$.
Its restriction to $G_{H_{fn_I\ell}}$ is cohomologous to
$g\mapsto(g-1)Z$. Their difference is the coboundary of a single
element of $E[2^a]$; subtract the same global coboundary to make
the restriction exactly this Kummer cocycle. Denote the adjusted
representative by $b$. Restrict it to $G_{H_{fn_I}}$ and choose
a local inertia lift of $\sigma_\ell$. Put
\[
                         U=(\sigma_\ell-1)Z-b(\sigma_\ell).
\]
For $g\in G_{H_{fn_I\ell}}$, compare $b(g\sigma_\ell)$ with
$b(\sigma_\ell(\sigma_\ell^{-1}g\sigma_\ell))$.
Using the exact restricted Kummer representative gives $gU=U$.
Consequently $U$ is an actual point of $E(H_{fn_I\ell})$, and
\[
 2^aU=L((\ell+1)Y-a_\ell X).
\]
The point
\[
 U_0=L\frac{\ell+1}{2^a}Y-L\frac{a_\ell}{2^a}X        \tag{H15}
\]
also belongs to this field, with the same $2^a$-multiple. Its
coefficients are integers. Hence $U-U_0$ is upper-field
$2$-primary torsion and is killed by $2^e$, by (H2).
This comparison with an integer quotient point is necessary;
comparison with an arbitrary division point would give no bound.

Inertia acts trivially on the reduction of every algebraic point
at a good prime. Reduction is injective on $2$-primary torsion.
Thus the singular coordinate $b(\sigma_\ell)$ is
\[
 s_{\rm new}=(a_\ell-(\ell+1)F)\overline Q+\epsilon
             =(F-a_\ell)f_{\rm old}+\epsilon,
                    \qquad 2^e\epsilon=0.              \tag{H16}
\]
The second equality uses the exact endomorphism identity
\[
 (F-a_\ell)(a_\ell F-(\ell+1))
                         =a_\ell-(\ell+1)F.            \tag{H17}
\]
Reduction $T/2^aT\to T/2^mT$ is multiplication by
$2^{a-m}$ on geometric torsion points. It kills $\epsilon$.
Conditions (H12) then turn (H16) into $s_{\rm new}=Jf_{\rm old}$.

For the new finite coordinate, all inertia translates of $Y$
have the same reduction, so
\[
 \overline{LD_\ell Y}
       =L\frac{\ell(\ell+1)}2F\overline X.              \tag{H18}
\]
This is $2^a$ times an $F^2$-fixed reduced point, because
$2^{a+1}\mid\ell+1$. The local extension obtained by adjoining
the conductor $\ell$ is totally ramified, of degree $\ell+1$.
Its residue field is still $\mathbf F_{\ell^2}$. Lift the indicated
$2^a$-division point in the residue field; the discrepancy lies
in the formal group, where multiplication by $2^a$ has an
integral formal inverse because $\ell$ is odd
\cite[IV.2.3 and IV.3.2, pp.~122--124]{silverman}. This produces a
local $2^a$-division point of $LD_\ell Y$ in the upper local
field itself. The restricted Kummer class is therefore zero.
Choose a Frobenius lift in that upper local Galois group.
It acts trivially on $E[2^a]$ by (H12), so its cocycle value is
zero, independently of a coboundary adjustment. Hence
\[
                     f_{\rm new}=0,\qquad
                     s_{\rm new}=Jf_{\rm old}
                  \quad\hbox{at precision }2^m.         \tag{H19}
\]
The generator $\sigma_\ell$ fixes the tame orientation. A different
generator changes the corresponding singular coordinate by the
matching unit. Frobenius lifts and places in (H19) are part of
the retained local diagram. The extra bit in (H12) pays exactly
for the division by two in (H18). Repeat this calculation with
each prime of $I\cup\{\ell\}$ singled out: all other lower
conductors split at that prime by the same principal-ideal argument.
Thus the previously imposed singular conditions are retained.
The maximum of the required finite precisions suffices for the
whole finite set; neither (H11) nor its exponent is repeated.

\paragraph{7. The actual coefficient push and its local conditions.}
Reduce $b_{I,a}$ to $E[2^m]$, extend coefficients freely to $A$,
and only now use the character $\theta$ of $\Delta_f$.
Under Shapiro, followed by the equivariant coefficient map,
\[
 E[2^m]\otimes A[\Delta_f]\longrightarrow
 (E[2^m]\otimes A)(\theta),\qquad
                     v\otimes[\delta]\longmapsto\theta(\delta)v,
                                                               \tag{H20}
\]
obtain the class $P_{I,m}$ over $K$.
Here the induced module is written with the diagonal Tate action
and the left regular action on its class-group basis. Formula
(H20) is equivariant since left translation multiplies the value
by $\theta(g)$. It remains equivariant with values in $A[G]$,
before evaluation at any sign character.

For an explicit degree-one cochain formula, choose representatives
$r_\delta$ of $G_K/G_{H_f}$ using the normal quotient
$\Delta_f$, and write
$g r_\delta=r_{g\delta}h(g,\delta)$ with $h(g,\delta)\in G_{H_f}$.
The transferred cocycle is
\[
 (\mathcal T_\theta b)(g)=
 \sum_{\delta\in\Delta_f}
  \theta(r_{g\delta})\rho_T(r_{g\delta})b(h(g,\delta)). \tag{H21}
\]
The identity
$h(gk,\delta)=h(g,k\delta)h(k,\delta)$ proves the cocycle
identity with target action $\theta\rho_T$.
In degree zero the formula is the weighted sum
$\sum_\delta\theta(r_\delta)\rho_T(r_\delta)Q$; applying its
boundary gives (H21) applied to the boundary of $Q$.
More generally, put $R(g)=\theta(g)\rho_T(g)$ and, for
$g_1,\ldots,g_n$, set $\delta_i=g_{i+1}\cdots g_n\delta$.
The degree-$n$ map is
\[
 (\mathcal T_\theta c)(g_1,\ldots,g_n)
 =\sum_{\delta\in\Delta_f}R(r_{g_1\cdots g_n\delta})
       c\bigl(h(g_1,\delta_1),\ldots,h(g_n,\delta_n)\bigr).
\]
With the ordinary inhomogeneous differential, the relation for $h$
above gives $d\mathcal T_\theta=\mathcal T_\theta d$.
These are the integral Shapiro and transfer maps of
\cite[Sections 8.1.1--8.1.6, pp.~189--194]{nekovar}, written
with left cosets. Changing sections gives the integral homotopy
in Section 8.1.4; compatible sections give strictly transitive
transfer. No averaging factor occurs.

Localizing (H21) separates the cosets into decomposition-group
orbits: this is the local Mackey formula, a sum of the local
corestrictions from the transported places.
At $2N$, the deformation character $\Psi$ is trivial locally;
the fixed quadratic character need not be trivial. Under the
fixed quadratic-twist identification, the local map is the
ordinary integral norm on Kummer classes. One can see this
without a projector: the sum $\sum_\delta\chi(\delta)\delta Q$
is a point of the quadratic twist over the lower field, and
Kummer of this sum equals the cohomological transfer.
The extra $A$ coefficients take finite linear combinations of
such maps. Thus the old Kummer conditions proved in paragraph 5
survive (H20), also in the universal family with its common local
quadratic type. Local Kummer corestriction is natural even for
a ramified field extension; in the present old-place descent
the extensions used in (H10) are unramified.

At a derivative prime, $H_f/K$ splits completely and $\theta$
is trivial on the local Galois group. Each Mackey term uses a
transported place and a transported Frobenius. Transporting the
Tate basis transports $F$ as well, so (H19) is equivariant,
not an assertion that an arbitrary change-of-place matrix
commutes with a fixed untransported $J$. The same weights occur
on the upper and lower class. Summing gives, in the common target coordinates,
\[
 f(P_{I\cup\{\ell\},m})=0,\qquad
 s(P_{I\cup\{\ell\},m})=Jf(P_{I,m}).\tag{4.2}
\] Active and moving conditions are full; all other
good conditions are unramified and preserved by the same maps.

The split norm formula in (4.1) also survives this construction.
For a new conductor prime the class-field squares obtained by
retaining the derivative factors are Cartesian, by (H5) and the
same local-unit calculation at that prime. Restriction and
corestriction commute in these squares. If $b$ lifts $2^e c$
in (H9), its corestriction lifts $2^e$ times the corresponding
corestricted class; uniqueness after the second multiplication
by $2^e$ proves the required compatibility of $d_a$.
The weighted sums are transitive finite sums. Thus the two
horizontal Artin terms in (4.1) become $\theta(\sigma_q)$ and
$\theta(\sigma_q)^{-1}$, giving
$a_q-\theta(\sigma_q)-\theta(\sigma_q)^{-1}$ on the class.
For $I=\varnothing$, the class is precisely the Kummer image of
the $L$-multiplied weighted point, so the same statement gives
the central split-norm relation with that exact $L$.

\paragraph{The local diagrams underlying the coefficient push.}
At a retained derivative place put $q=\ell^2$ and choose tame
and unramified generators $\tau_v,\phi_v$ with
$\phi_v\tau_v\phi_v^{-1}=\tau_v^q$. The local coefficient
character is trivial. The two-generator tame resolution is
\[
 [M\xrightarrow{u\mapsto((\Phi_v-1)u,0)}M\oplus M
       \xrightarrow{(x,y)\mapsto(\Phi_v-q)y}M],
 \qquad \Phi_v=\rho(\phi_v).
\]
At the retained exponent $m\le a$, (H12) makes this the
zero-differential model $D^0\oplus(F_v\oplus S_v)[-1]
\oplus D^2[-2]$. Moreover $q\equiv1\pmod{2^{a+2}}$.
Represent the singular generator by a uniformizer adjusted by
a unit so its Kummer value on the chosen $\phi_v$ is zero.
Residue-unit Kummer characters allow this adjustment. The latter
congruence makes its self-cup zero: $-1$ is a $2^m$-th power,
so the adjusted uniformizer Kummer class has square zero. The finite self-cup is zero by unramified cohomological
dimension one. Retain the transported mixed pairing and the
$D^0,D^2$ pairing.

In the derivative diagram the old ring-class fields split
locally, and the new conductor extension is totally ramified.
For the latter choose the lower Frobenius in the upper subgroup
and take $\tau_{\rm upper}=\tau_{\rm lower}^e$; here $f=1$.
For split terms $e=f=1$. Transport these choices with the Tate
basis. With these compatible generators the local matrices, in the order
$D^0,F_v,S_v,D^2$, are
\[
 \operatorname{res}=\operatorname{diag}(1,f,e,ef),\qquad
 \operatorname{cor}=\operatorname{diag}(ef,e,f,1).
\]
The possible transfer shear is divisible by $(q-1)/2$ and
vanishes at retained precision. These matrices preserve each
of the four conditions retaining $D^0$, and satisfy the local
cup-product adjunctions. Split Mackey terms have $e=f=1$;
coefficient weights and reduction commute with these matrices.
An arbitrary new Frobenius lift is not used without transporting
the resulting finite/singular decomposition.

Compare the small free resolution with the continuous bar
resolution and choose augmentation-preserving comparison maps
and their composite homotopies. Continuous Hom into $M$ retains
these identities. Since the small model has zero differential,
the composite on it is the identity; normalize the other homotopy
to obtain a contraction. Here are its compatibility maps.
For this retained contraction $1-i_vp_v=dh_v+h_vd$ and a full
local cochain arrow $F$, take $F_L=p_wFi_v$ and restrict it to
$U_v$. The inclusion square has homotopy $h_wFi_v|_{U_v}$.
For composable arrows the transferred composition homotopy is
$p_zGh_wFi_v$, together with the transferred original bar
homotopy when present. Every such degree-minus-one map on
$U_v$ lands in $D^0$, which is retained in every condition;
thus these homotopies restrict too. For three composable maps,
write $K(G,F)=pGhFi$ and $L(H,G,F)=pHhGhFi$, with the
intermediate subscripts understood. Direct expansion gives
\[
 K(H,GF)-K(HG,F)+H_LK(G,F)-K(H,G)F_L=Ld-dL.
\]
The degree-minus-two map $L$ vanishes on every derivative
condition. Transfer the original bar comparison homotopies and
cubes as well; unrelated classwise boundaries are not substituted
for them. For the isotropy diagrams, $U_v\otimes U'_v$ is
free in degrees $0,1,2$, whereas the trace target is $A[-2]$.
Its negative Hom groups vanish, so the restricted zero pairing
has compatible null-homotopies and the requisite 2-homotopies.
This constructs the local data required in
\cite[8.6.2.1--8.6.2.4]{nekovar} before using 8.6.4.
At old Kummer places use the natural restriction/norm diagrams
of the multiplication-by-$2^m$ exact sequences, including
local torsion, and their derived scalar extensions; at good
unramified, full and zero places use the corresponding standard
subcomplexes. The witnesses in (H22) are chosen in these same
diagrams.

\paragraph{8. Actual finite Selmer-cone witnesses.}
Represent $P_{I,m}$ by a continuous global cocycle $p$.
At each constrained place, membership in the indicated local
condition gives a local cycle $u_v$ and a degree-zero cochain
$q_v$ with
\[
                   \operatorname{loc}_v p-\iota_vu_v=dq_v.
                                                               \tag{H22}
\]
These entries give a cycle in
$\operatorname{Cone}(C^\bullet(G_{K,S},M)\oplus\bigoplus U_v
\to\bigoplus C^\bullet(K_v,M))[-1]$.
For (H19), choose in the same way a degree-zero cochain whose
boundary is the difference of the two representatives being
compared. At the derivative primes use the already specified
$D^0$-retaining finite/singular local complexes, not full/zero
conditions. Their coordinate equality is the finite local
cohomology equality just proved. All coefficient groups and
coset sets at this stage are finite; the cocycles and the finite
list of witnesses factor through finite quotients.

Apply the coefficient push to these witnesses at the same time
as to the cocycle, including the natural Kummer comparison
homotopies. If $(i,p,h)$ is a retained contraction to a finite
target model, a relation $z-z'=dq$ becomes
$pz-pz'=d(pq)$. For a comparison $F$ to another model with
projection $p'$, the correction between $p'Fz$ and $p'Fi pz$
is the boundary of $p'Fhz$ (with the consistent contraction sign).
Thus the finite arithmetic provides actual cycles, boundaries
and comparison entries before compression. Adding a derivative
prime extends this same finite diagram. Existence and size
bounds for those target contractions, and the subsequent two
limits, are separate inputs; this proof makes no inference
from abstract isomorphism of cohomology groups to those bounds.

The resulting finite arithmetic uses the single multiplier (H11)
and the precision rule $a\ge m+e$, (H12). No factor depending on
$[H_f:K]$, $|G|$, $|I|$ or the number of active primes occurs.
\end{proof}

\subsection{The derived Kummer condition}
\begin{lemma}[The actual derived Kummer condition]\label{ex:arith-local-kummer}
For the actual local twist $B=E^{(h)}$ over $F=K_v$, the condition
$U_F=B(F)^\wedge_2[-1]\to R\Gamma(F,T_2B)$ includes the finite
two-primary torsion. Its reduction modulo $2^a$ has
$B(F)[2^a]$ in degree zero and $B(F)/2^aB(F)$ in degree one.
Its Kummer image is its exact local annihilator, and the comparison
has integral isotropy homotopies. At $v\mid2N$, triviality of the
moving character and the common old quadratic type identify the
scalar and binary versions by derived coefficient extension.
These are the local maps used in the Selmer-cone duality above.
\end{lemma}
\begin{proof}\leavevmode
For the actual local twist $B=E^{(h)}$ over $F=K_v$, put
\[
 U_F=B(F)^\wedge_2[-1]\longrightarrow R\Gamma(F,T_2B).
\]
The completion includes finite two-primary torsion. Derived reduction
modulo $2^a$ has $B(F)[2^a]$ in degree zero and
$B(F)/2^aB(F)$ in degree one. At $v\mid2N$, $\Psi_i$ is trivial,
so derived coefficient extension defines the scalar $A_i$ condition
for this actual twist. In the universal binary family every $h_x$
has the same local squareclasses at $2N$ as $h_0=1$. Hence every
sign character evaluates the universal character $\varepsilon$
trivially on $G_{K_v}$, and therefore $\varepsilon|_{G_{K_v}}=1$
in $G$, before reduction modulo two. The $A_i[G]$ Kummer condition
and its null-homotopy are consequently derived scalar extensions
of the fixed local data. Omitting the scalar twist, the universal
local-squareclass condition, or degree-zero torsion changes this argument.

The local finite dualities \cite[I.2.3, I.3.2--I.3.5]{milne} make the
Kummer map adjoint to $H^1(F,B[2^a])\to H^1(F,B)[2^a]$.
The latter has Kummer image as kernel, so the Kummer image is its exact
annihilator. To obtain the integral chain statement, write
$M=B(F)^\wedge_2$. The restricted cup map
$M[-1]\otimes^LM[-1]\to\Z_2[-2]$ has homotopy class determined by
$\operatorname{Hom}(M\otimes M,\Z_2)$. Finite isotropy makes every
value divisible by every power of two, hence zero.
A projective resolution therefore supplies a null-homotopy.
The possible ambiguity group in degree minus one vanishes.
Reduction modulo two, including the torsion in degree zero, gives the
perfect quotient-to-dual comparison; derived Nakayama lifts it.
The global Selmer duality theorem
\cite[Sections 5.1, 5.2.4, 5.4.3 and 6.3.4]{nekovar} is applied
to these actual local maps and homotopies.

For clarity the Selmer complex is the cone
\[
 C_U=\operatorname{Cone}\left(
 R\Gamma(G_{K,S},M_i)\oplus\bigoplus_{v\in S}U_v
 \longrightarrow\bigoplus_{v\in S}R\Gamma(K_v,M_i)\right)[-1].
                                                                    \tag{4.3}
\]
The compatible trace in (4.3) is coefficient-valued. A top-coefficient
functional used to test a finite Frobenius ring's perfection is not
used as a coefficient-reduction-compatible arithmetic trace.
\end{proof}

\subsection{Actual support comparisons}
At a tame split prime the local free model has
\[
 d^0(v)=(Av,Cv),\quad d^1(v,w)=-Cv+Dw,\quad
 A=\Phi-1,\ C=g-1,\ D=\Phi-N_q(g),
\]
\[
 N_q(g)=\frac{q+1}{2}+\frac{q-1}{2}g,\qquad DC=CA.                 \tag{4.4}
\]
At augmentation $g=1$, this is the direct sum of the unramified
complex $[\Phi-1]$ in degrees zero and one and the singular quotient
$[\Phi-q]$ in degrees one and two.

The support-change theorem \cite[Proposition 7.8.8(ii)]{nekovar}
retains arbitrary old local conditions and puts the unramified
condition at new places. Its finite-module proof applies at two.
Use it only when the actual coefficient representation is unramified
at the removed place. There are five uses here:
universal augmentation at $Q$; an unused scalar prime $q\nmid h$;
characteristic-two residual removal after $\chi_h=1$; the moving
prime after $t=0$; and a new good inert derivative prime outside
$h,r_i$ and the old support. In the last case the coefficient module
is unramified and the chosen finite local plane is the unramified
condition. It is not applied to an active ramified scalar prime
or to $r_i$ before central specialization. Replacing the unramified
condition by the full condition then gives the singular quotient
by a cone of Selmer cones. In particular the old complex is an actual
subcomplex of an upper triangular model, not just a complex with the
same Euler characteristic.

\begin{proposition}[One moving character for the complete active list]\label{ex:arith-moving-character}
Fix an allowed $K$ and any finite list of oriented active primes.
For every $n,B\geq1$ and finite avoidance set there are infinitely
many split rational primes $r\equiv1\pmod{2^B}$ and characters
$\xi:\operatorname{Gal}(H_r/K)\to\Z/2^n\Z$ that are surjective
on inertia at $r$, trivial on the local decomposition groups at
$2N$, and take value one at every prescribed active orientation.
They can be chosen with
\[
 v_2(2\pm a_r)\leq c_E,\qquad
 v_2\det(1\pm\operatorname{Fr}_r)\leq c_E,
\]
where $c_E$ depends only on the fixed Tate image, not on $K,n,B$
or the active list. One such character serves the entire binary
family before evaluation.
\end{proposition}
\begin{proof}\leavevmode
\paragraph{The simultaneous moving character.}
We give the construction for the non-CM curve used here. For every
$n,B\ge1$ and every finite avoidance set, there are infinitely many
split rational primes $r\equiv1\pmod{2^B}$ and characters
\[
 \xi:\operatorname{Gal}(H_r/K)\longrightarrow\Z/2^n\Z
\]
which are surjective on inertia at $r$, trivial at $2N$, and equal
to one at each prescribed oriented active prime. The trace bounds
below have one constant $c_E$, independent of $n,B,K$ and the active
list. The degrees of the finite fields and the sizes of the primes
chosen in the construction need not be uniformly bounded.

Choose prime ideals $\mathfrak p_1,\ldots,\mathfrak p_s$ in distinct
conjugate pairs, including the active orientations, one prime over
each prime dividing $2N$, and enough additional split prime ideals
to generate $\operatorname{Cl}(K)$. The latter exist outside any
finite set by Chebotarev in the Hilbert class field. Set $w_j=1$ at
the active primes and $w_j=0$ at the other chosen primes. Put
\[
 \mathcal R=\ker(\Z^s\longrightarrow\operatorname{Cl}(K)),
 \qquad (\beta_a)=\prod_j\mathfrak p_j^{m_j^{(a)}},
 \qquad \alpha_a=\beta_a/\overline\beta_a,
\]
where the $m^{(a)}$ form a $\Z$-basis of $\mathcal R$.
The generators $\beta_a$ may be fractional. Disjointness of the
conjugate pairs gives $v_{\mathfrak p_j}(\alpha_a)=m_j^{(a)}$.

Let $F_\infty=K(\mu_{2^\infty})$ and adjoin the $2^n$-th roots of
all the $\alpha_a$. Prescribe the multiplier
\[
 \zeta_{2^n}^{-\sum_jm_j^{(a)}w_j}
 \quad\hbox{on }\alpha_a^{1/2^n}.                              \tag{4.4a}
\]
These prescriptions respect every Kummer relation. Indeed, if
$\prod_a\alpha_a^{b_a}$ is a $2^n$-th power in $F_\infty$, its
root lies in a finite cyclotomic layer. At every nonzero-weight
prime that layer is unramified over $K$, since the prime is odd.
Consequently $2^n$ divides $\sum_ab_am_j^{(a)}$ at each such prime,
and hence divides $\sum_{a,j}b_am_j^{(a)}w_j$. The dyadic primes
have weight zero. This tests all multiplicative relations, including
those involving units; it does not assume independence of the
radicands. Kummer theory therefore gives the specified automorphism
of the finite abelian extension
\[
 L_\infty=F_\infty(\alpha_a^{1/2^n}:a)/F_\infty.
\]

Here is why these are precisely the extension conditions for $\xi$.
For a split $r$ avoiding the chosen primes, the ring-class sequence is
\[
 1\longrightarrow\mathbf F_r^\times
 \longrightarrow\operatorname{Pic}(\Z+r\mathcal O_K)
 \longrightarrow\operatorname{Cl}(K)\longrightarrow1.
\]
The kernel is the ratio of the two residue factors in
$(\mathcal O_K/r\mathcal O_K)^\times$ modulo diagonal rational
residues. There is no additional unit quotient because
$\mathcal O_K^\times=\{\pm1\}$. Write $q_j$ for the ideal-class
lift of $\mathfrak p_j$, and orient the connecting map $\delta$ so
that its principal relations are
\[
 \sum_jm_j^{(a)}q_j+\delta(\operatorname{res}\alpha_a)=0.
\]
The kernel and the $q_j$ generate the ring-class group. Any relation
between them projects to $\mathcal R$; subtracting its basis
relations leaves a relation within the kernel. Thus these are a
complete presentation, not merely necessary relations.
For $r\equiv1\pmod{2^n}$ the power-residue character
$\kappa_r:\mathbf F_r^\times\to\Z/2^n\Z$ is surjective.
The assignments $\xi(q_j)=w_j$ and $\xi\delta=\kappa_r$ extend
exactly when
\[
 \kappa_r(\operatorname{res}\alpha_a)
       =-\sum_jm_j^{(a)}w_j.                                  \tag{4.4b}
\]
Below we choose $r$ so that its radical Frobenius is (4.4a), which
gives (4.4b). Reversing the connecting-map convention reverses the
power-residue convention as well. No class number is divided out.

We next preserve all these conditions while controlling the Tate
trace. Let $J=\rho_T(G_{\Q(\mu_{2^\infty})})$. The determinant is
the cyclotomic character, so $J$ is the determinant-one part of the
fixed Tate image. Non-CM open image \cite{serre} gives a fixed
$d\ge2$ with
\[
 \Gamma_d:=\ker\bigl(\operatorname{SL}_2(\Z_2)
             \to\operatorname{SL}_2(\Z/2^d\Z)\bigr)\subset J.
\]
The subgroup $J_K=\rho_T(G_{F_\infty})$ has index at most two in
$J$, and therefore contains its squares. Write $U(a)$ and $V(b)$
for the upper and lower unipotent matrices. Then $J_K$ contains
$U(2^{d+1}\Z_2)$ and $V(2^{d+1}\Z_2)$. The identity
\[
 U(a)V(b)U\bigl(-a/(1+ab)\bigr)V\bigl(-b(1+ab)\bigr)
       =\operatorname{diag}(1+ab,(1+ab)^{-1})
\]
and Gaussian decomposition show that a subgroup containing
$U(2^k\Z_2),V(2^k\Z_2)$ contains $\Gamma_{2k}$.
Hence $\Gamma_m\subset J_K$ for $m=2d+2$.
Commuting $\operatorname{diag}(1+2^m,(1+2^m)^{-1})$ with the
upper and lower unipotents of level $m$ gives all upper and lower
unipotents of level $2m+1$, since
$v_2((1+2^m)^2-1)=m+1$. Applying the same decomposition gives
\[
 \Gamma_{4m+2}\subset\overline{[J_K,J_K]}.
\]
Every matrix commutator lifts to a commutator in $G_{F_\infty}$.
Such lifts fix $L_\infty$, whose Galois group over $F_\infty$ is
abelian. The image of the compact group $G_{L_\infty}$ is closed.
Thus, for the fixed integer $\nu=4m+2=8d+10$,
\[
                  \Gamma_\nu\subset\rho_T(G_{L_\infty}).       \tag{4.4c}
\]
This is a subgroup in the image of the actual radical-fixing
kernel. It requires no disjointness of the radical and Tate fields,
and is independent of the number of radicands and of $K$.

For the trace estimate one can take $c_E=2\nu+1=16d+21$.
We verify the elementary bound used here. Put $\delta=2^\nu$ and
$D(v)=\operatorname{diag}(v,v^{-1})$. For any
$A=\left(\begin{smallmatrix}a&b\\c&f\end{smallmatrix}\right)
\in\operatorname{SL}_2(\Z_2)$, one of
\[
       A,\quad AU(\delta),\quad AD(1+\delta),\quad AD(1+2\delta)
                                                                    \tag{4.4d}
\]
has both trace distances to $2$ and $-2$ of valuation at most $2\nu+1$.
If $A$ itself does not, its trace is within $2^{2\nu+2}\Z_2$ of
one of these two values, while its distance to the other has valuation
two. Its diagonal entries have the same parity. If both are even,
$c$ is odd by $af-bc=1$, and $AU(\delta)$ changes the trace by
$c\delta$, of valuation $\nu$. If both are odd, the two diagonal
multipliers change the trace by
\[
 \delta\frac{(a-f)+a\delta}{1+\delta},\qquad
 2\delta\frac{(a-f)+2a\delta}{1+2\delta}.
\]
If $v_2(a-f)<\nu$, the first change has valuation at most $2\nu-1$.
Otherwise write $a-f=\delta z$. For even $z$ the first change has
valuation $2\nu$; for odd $z$ the second has valuation $2\nu+1$.
All the chosen changes have valuation at least three. They therefore
leave the distance to the other forbidden value of valuation two,
and give the claimed bound for the nearer one.

Lift the prescribed Kummer automorphism to
$\sigma_0\in G_{F_\infty}$. By (4.4c), every right multiplier in
(4.4d) lifts to an element of $G_{L_\infty}$. Choose one giving the
trace bound and multiply $\sigma_0$ by that lift. The resulting
actual Galois element $\sigma$ fixes $F_\infty$, preserves every
prescribed radical multiplier, and satisfies
$v_2(2\pm\operatorname{tr}\rho_T(\sigma))\le c_E$.

Choose $B'\ge\max(B,n)$ with $B'>c_E$, and take a finite normal
extension of $\Q$ containing $K,\mu_{2^{B'}},E[2^{B'}]$ and all
the chosen radical roots. Restrict $\sigma$ to this extension.
Its restrictions on all intersections agree because it is already
an actual Galois element. Finite Chebotarev supplies infinitely many
rational primes in its Frobenius class. Choose a place realizing
this representative and avoid all the finitely many excluded primes.
Then $r$ splits in $K$ and $r\equiv1\pmod{2^{B'}}$. The Frobenius
formula on $\alpha_a^{1/2^n}$ gives (4.4b), and hence the character
$\xi$. Each local unit factor at $r$ maps onto the residue-ratio
kernel, so local reciprocity proves that $\xi$ is surjective on
inertia there. At $2N$ it is unramified with Frobenius value zero,
and is therefore trivial on the whole local decomposition group.
Conjugation negates its values because a prime ideal times its
conjugate is rational principal.

Finally $a_r\equiv\operatorname{tr}\rho_T(\sigma)\pmod{2^{B'}}$
and $\det(\operatorname{Fr}_r)=r$. Since $B'>c_E$, we obtain
\[
 v_2(2\pm a_r)\le c_E,\qquad
 v_2\det(1\pm\operatorname{Fr}_r)
       =v_2(1\pm a_r+r)\le c_E.
\]
The same bound controls $a_r-2\chi_h(r)$ for every constituent
unramified at $r$. Repeat with increasing $n,B$ and avoidance sets
to obtain distinct $r_i$. Taking $n_i\ge2i$ makes
$(1+t)^{2^{n_i}}=1$ in $A_i$: for $1\le j<i$,
$v_2\binom{2^{n_i}}j\ge n_i-v_2(j)\ge i$.
Thus $\Psi_i=(1+t)^{\xi_i}$ is the required finite character.
Its inertia generator can be oriented to have value $1+t$, and its
active oriented Frobenius values have odd exponent. One character
works for the entire fixed active list before any binary evaluation.
\end{proof}

\subsection{The same functional at every localization}
\begin{proposition}[The paired functional and its integral representatives]\label{ex:arith-paired-functional}
Use the retained inner-limit diagrams of
Lemma~\ref{ex:limits-inner-limit}, with the local conditions of
Definition~\ref{ex:arith-data}, and let $z$ be the actual retained
class. Transport $z$ by conjugation, lift it through the orthogonal
comparison at the generic point, and pair by the retained duality.
This defines one generic closed degree-two functional $\lambda_z$.
It has an integral representative after multiplication by $2^a$
over $\mathscr O[G]$, for a fixed old-support exponent $a$
independent of the active list and $G$.
For a scalar constituent with its primitive local conditions, its
paired functional has an integral representative at every height-one
prime $\mathfrak p\ne(2)$ of $\Lambda$, including $(t)$ and $(t+2)$.
These representatives have the same generic value for that primitive
constituent. The universal relaxed functional is asserted integral
after the fixed power of two over $\mathscr O[G]$; scalar regularity
of its corrected coordinate is obtained instead from the actual
primitive comparison of Lemma~\ref{ex:arith-common-coordinate}.
All these assertions concern the inner limit and the indicated
localizations, not inverses of $t$ in a finite Artin ring.
\end{proposition}
\begin{proof}\leavevmode
After the inner limit, put
$\mathscr O=\widehat{\Z_2[[t]]_{(2)}}$.
In (4.4), the residual active Frobenius is $(1+t)^aI$ with $a$ odd.
Thus $A,D$ are invertible over $\mathscr O[G]$.
The formulas
\[
 h^1(v,w)=A^{-1}v,\qquad h^2(u)=(0,D^{-1}u)
\]
contract the complex; the identities follow from $DC=CA$.
\paragraph{Which functional is paired with the class.}
Let $z$ denote the retained class and transport it by $\tau$ and the
preceding Weil identification to the conjugate coefficient system.
At a localization where the full local complexes at active and
moving primes are acyclic, the comparison from the orthogonal
condition to the transported condition is a quasi-isomorphism.
Apply its inverse in the derived category to the transported class,
and call the resulting class $z^\perp$. Define
\[
 \lambda_z(w)=\operatorname{PD}(w,z^\perp),\qquad
           w\in H^2(C_U).
\]
Thus $\lambda_z$ is a closed degree-two functional, with a cochain
representative obtained using the displayed local contractions and
the retained global duality map. The inverse is unique in the derived
category, so representatives built at distinct localizations define
the same generic functional. Integrality still requires the stated
integral contractions and fixed old-support correction; it does not
follow from abstract duality alone. The coordinate definition below uses this same $\lambda_z$.

Here the representatives must be constructed separately. Over
$\mathscr O[G]$, the active Frobenius operators $A,D$ are units,
and moving-prime inertia minus one is $t$ times a unit; the displayed
homotopies are integral. For the height-one assertion take the scalar
primitive constituent. Its active primes are precisely the primes
ramified in that constituent, so inertia minus one is $-2$, a unit
at $\mathfrak p\ne(2)$. At the moving prime inertia minus one is
invertible unless $\mathfrak p=(t)$; at $(t)$ the nonzero central
Frobenius determinants give integral adjugate contractions. In
particular $(t+2)$ is treated by the inertia contraction. Unused
good primes in the primitive constituent retain the unramified
condition and its orthogonal comparison; no inverse of a full local
Frobenius complex at such a prime is required. This proves the
scalar assertion only for the primitive functional. In the universal
relaxed family an unused prime can specialize to trivial inertia,
so the $-2$ contraction is not available there. Its correction is
the determinant identity of Lemma~\ref{ex:arith-common-coordinate}.
The common generic value alone would not prove any of these
integrality statements.

In the inertia contraction the maps are
\[
 h^1(v,w)=C^{-1}w,\qquad h^2(u)=(-C^{-1}u,0).
\]
\paragraph{The fixed old-support lattice and the exponent.}
Here is the integral content of the old-support comparison.  For each
$p\mid2N$ and each of the fixed allowed local quadratic types $\delta$,
choose once a bounded free $\mathbf Z_2$-model of the diagram
\[
 U_{p,\delta}\longrightarrow L_{p,\delta},\qquad
 U_{p,\delta}=(E^{(\delta)}(\mathbf Q_p)^\wedge_2)[-1],\qquad
 L_{p,\delta}=R\Gamma(\mathbf Q_p,T_2E^{(\delta)}),
\]
together with its Kummer map, cup pairing, isotropy homotopy, and
quotient-to-dual comparison.  Finite torsion is retained in $U_{p,\delta}$;
it is represented by a free resolution, not removed.  The local
duality and Kummer exact sequence identify that comparison integrally:
\[
 \operatorname{Cone}(U_{p,\delta}\to L_{p,\delta})
       \simeq R\operatorname{Hom}_{\mathbf Z_2}
                         (U_{p,\delta},\mathbf Z_2)[-2].
\]
This is the chain-level comparison of
Lemma~\ref{ex:arith-local-kummer}, with its specified isotropy
homotopy.  In particular it is not a rational isomorphism of free
parts of cohomology.  A quasi-isomorphism between these bounded free
models has an integral homotopy inverse: its bounded acyclic free
cone is contractible, by splitting from its last nonzero term
and proceeding downwards.  Consequently neither torsion lengths
nor inverse determinants of torsion differentials are denominators
in this comparison.

Every allowed $K$ splits at $2N$.  Its two completions above $p$ are
therefore $\mathbf Q_p$, with precisely two copies of this same
diagram, interchanged by conjugation.  The moving character and
the universal varying quadratic character are trivial on their
decomposition groups.  All old-place diagrams over $A_i[G]$ are thus
derived scalar extensions of this finite fixed list.  Tensoring the
chosen free models and their homotopy identities preserves all these
identities.  The remaining good unramified conditions use the
integral unramified comparison, and the derivative conditions use
their integral orthogonal planes.  Neither requires a new lattice
comparison when a prime is added.

For completeness, there is also no hidden group-order denominator in
duality.  The coefficient pairing is the $A_i[G]$-valued Weil pairing.
If one checks its perfection over $A_i$, the auxiliary pairing on the
group algebra is
\[
 (a,b)\longmapsto[1](ab),\qquad
 ([1](gh))_{g,h\in G}=(\mathbf1_{h=g^{-1}})_{g,h\in G}.
\]
Its matrix is a permutation matrix, with integral inverse.  It is not
the normalized trace $|G|^{-1}\operatorname{Tr}$.  The global
cup-product and trace maps are integral as well.  The hypotheses of
Selmer duality here are a perfect coefficient pairing and the actual
orthogonal local diagrams with their isotropy homotopies
\cite[Sections 6.2.5--6.2.8 and Theorem 6.3.4]{nekovar}; the old
Kummer diagrams just described supply those hypotheses.  Since $K$
is imaginary there is no real-place Tate correction in this argument.

Over $S=\mathscr O[G]$ the only remaining differences between the
transported conditions and their orthogonal conditions occur at the
active and moving places.  Their full local complexes have the
integral contractions displayed above.  Their direct sum, of any
finite size, is still an integral contraction.  The resulting
comparison $j:C_{U^\perp}\to\tau^*C_U$ is an integral homotopy
equivalence over $S$.  Write $j^{-1}$ for such a homotopy inverse and
$c$ for conjugate transport.  On the retained bounded free models,
\[
 z^\perp=j^{-1}c(z),\qquad
 \lambda_z=\operatorname{PD}(-,j^{-1}c(z))
\]
are integral cochains, and the second is closed.  Every map in this
formula is integral; no inverse on a torsion cohomology quotient is
taken.  The finite-model contractions and their retained limits
preserve these identities.
In detail, represent the comparison on the bounded free models and
contract its acyclic cone over $S$.  The cone contraction yields a
chain map $j^{-1}$ and homotopies $jj^{-1}\simeq1$ and
$j^{-1}j\simeq1$ over $S$.  Thus
$d(j^{-1}c(z))=j^{-1}c(dz)=0$ as a cochain identity, and its image
differs from $c(z)$ by an integral boundary.  Applying the integral
duality chain map gives a closed element of
$\operatorname{Hom}_S(C_U^2,S)$, with the claimed generic value.
This proves the existence of an actual integral closed lift, rather
than only a lift after rationalization.

With the conventions of this paper the actual class $z$ already
contains the fixed multiplier $L$.  Thus the preceding construction
permits $a=0$.  One may keep any larger fixed nonnegative $a$ in the
subsequent inequalities.  This does not change $z$, divide out $L$,
or change the determinant basis.  It proves the existence and
independence of the exponent from $G$, the active graph, and the
allowed varying $K$; a numerical computation of a new denominator
is unnecessary because no such denominator occurs in these integral
diagrams.  This conclusion concerns $S$.  The separate primitive
height-one argument above, and the corrected relaxed coordinate
comparison, retain exactly their stated scopes.
\end{proof}

\begin{definition}[The class-functional coordinate]
\label{ex:arith-coordinate}
For a constituent whose generic cohomology has dimensions $(1,1)$
in degrees one and two and vanishes otherwise, use the determinant
line
\[
 \mathcal D=\det\nolimits^{-1} C
       \simeq H^1(C)\otimes H^2(C)^*.
\]
The class-functional coordinate $u$ is the coordinate of
$[z]\otimes[\lambda]\in\mathcal D$ in the retained integral
determinant basis. Here $z$ is the actual class and $\lambda$ is
the closed functional just constructed. For a generically acyclic
constituent set $u=0$. The primitive and relaxed subscripts
specify the corresponding local conditions on $C$.
\end{definition}

\subsection{Central identity, torsion and the primitive coordinate}
\begin{lemma}[Central specialization preserves the same class]\label{ex:arith-central-specialization}
The finite classes of Proposition~\ref{ex:arith-finite-heegner}
satisfy the retained boundary (4.5), with the same fixed $L$.
For the functional of Proposition~\ref{ex:arith-paired-functional},
the actual central Kummer comparison $i$ satisfies
\[
 z_0=i(LmP),\qquad
 (\lambda_z)_0(i x)=\operatorname{PD}_K(x,c_K(LmP)).
\]
Here $P$ is the unaveraged ring-class point and $c_K$ is its
conjugate-dual transport. On a rank-one Mordell--Weil free lattice,
both free-line coefficients consequently have valuation
$j_{\rm rc}(h,k)+v_2(Lm)$. The statement identifies the same point
and functional; their integral torsion factors are still computed
in the integral Kummer complex and local quotient triangles.
\end{lemma}
\begin{proof}\leavevmode
The split norm relation gives an actual finite-stage boundary
\[
 z_i(0)-Lm_i\,j_i(b_h)=dv_i,\qquad
 m_i=a_{r_i}-2\chi_h(r_i).                                      \tag{4.5}
\]
The equality in global cohomology lifts to the Selmer cone because
the local degree-zero maps are onto and the degree-one maps inject.
Retain $v_i$. The uniform trace bound bounds $v_2(m_i)$ and the local
determinants by $c_E$; their adjugate identities multiplied by a fixed
power of two survive the limit.

% Insert after baseline line 1583, before the ordinary compact Kummer complex.
\paragraph{Specialization of the same paired functional.}
Put $A=\Lambda_{(t)}$, with residue field $\mathbf Q_2$.
The representative of $\lambda_z$ over $A$ is the one obtained from
the actual orthogonal comparison and the retained cup pairing.
At an active prime its contraction uses $C=-2\in A^\times$.
At the moving prime the central Frobenius determinants are nonzero;
their adjugate inverses therefore belong to $A$. Both contractions
commute with $A\to\mathbf Q_2$: their defining identities are matrix
identities with invertible denominators in $A$. In particular the
central full local complexes are rationally acyclic. The Kummer local
objects at these changed odd places are also rationally zero:
local points have finite two-primary completion, and the retained
central Frobenius matrices have nonzero determinants. Thus the
full/Kummer quotients are rationally acyclic as well. At the fixed old
places the Kummer maps and their duality homotopies are retained unchanged.

Let $j$ denote the resulting orthogonal comparison over $A$ and
let $c(z)$ be the conjugate-dual transport of the actual class.
The definition is $z^\perp=j^{-1}c(z)$ and
$\lambda_z=\operatorname{PD}(-,z^\perp)$.
Derived base change preserves the specified chain maps, cup products,
and homotopies. It also preserves $j^{-1}$, since the base change of
the two inverse identities is again the two inverse identities.
Consequently
\[
 (\lambda_z)_0
 =\operatorname{PD}_0\bigl(-,j_0^{-1}c(z_0)\bigr).
\]
To make the central identification explicit, write $R_0=C_U(M)_0$,
$R_0^c=C_{cU}(M^*(1))_0$, and $S_0=C_{U^\perp}(M^*(1))_0$.
Let $K_0$ be the ordinary rational Kummer complex and $K_0^c$ its
conjugate-dual transport. The actual comparison factors as
\[
 i:K_0\longrightarrow R_0,\qquad
 s:S_0\longrightarrow K_0^c,\qquad
 q:K_0^c\longrightarrow R_0^c,\qquad j_0=q s.
\]
All three maps are quasi-isomorphisms, by the preceding local
vanishing and the unchanged old conditions. The retained duality
homotopies give
\[
 \operatorname{PD}_{R}(i x,y)=\operatorname{PD}_{K}(x,s y),
 \qquad c(iP)=q c_K(P).
\]
The boundary (4.5) gives $z_0=i(LmP)$. Therefore
\[
 (\lambda_z)_0(i x)
 =\operatorname{PD}_{K}
       \bigl(x,s(qs)^{-1}q c_K(LmP)\bigr)
 =\operatorname{PD}_{K}\bigl(x,c_K(LmP)\bigr).
\]
Thus this is the Kummer-dual functional of the same transported point.
On the integral rank-one Mordell--Weil free lattice, the central
conjugation transport is an involutive automorphism (including the
quadratic sign), hence acts by $+1$ or $-1$. It preserves the point
index; it need not fix the point. The two free-line coefficients both have valuation
$j_{\rm rc}(h,k)+v_2(Lm)$. An orientation sign is a fixed unit;
there is no independently chosen scalar multiplying the functional.
The integral torsion correction is still computed from the ordinary
integral Kummer complex and the integral local quotient triangles,
not from their rationally acyclic reductions. This is the compatibility
needed to apply the arithmetic length calculation in (4.6).
\end{proof}
\begin{lemma}[Determinant specialization with no dimension jump]
\label{lem:v10-determinant-specialization}
Let $A$ be a discrete valuation ring with uniformizer $t$, fraction
field $F$, and residue field $k$. Let $C$ be a bounded finite free
complex over $A$. If
\[
 \dim_F H^i(C\otimes_A F)=\dim_k H^i(C\otimes_A k)
 \quad\text{for every }i,
\]
then its cohomology is free and commutes with reduction modulo $t$.
In particular, when the only nonzero dimensions are $(1,1)$ in
degrees one and two, the coordinate of an actual cycle and an actual
closed functional in $\det^{-1}C$ is regular at $t=0$ and specializes
to the coordinate of their central tensor in the specialized retained
determinant basis.
\end{lemma}
\begin{proof}
Smith reduction splits the bounded free complex into free summands with
zero differential and two-term summands $[A\xrightarrow{t^r u}A]$,
where $u$ is a unit and $r\ge0$. A summand with $r>0$ contributes
one extra residue-field cohomology dimension in each of two adjacent
degrees and contributes none over $F$. The dimension equalities
therefore exclude every such summand. The remaining two-term summands
are contractible over $A$. On the zero-differential summands both
cohomology and the class-functional tensor commute with base change.
The determinant identifications include the same contractible
summands before and after reduction, with their unit differentials;
there is no independent rescaling of the retained basis.
\end{proof}

\begin{proposition}[Arithmetic calibration at known simple zeros]\label{ex:arith-central-calibration}
Let $K=\Q(\sqrt{k})$ and $h$ satisfy the ring-class hypotheses
above, with every prime dividing $h$ split in $K$. Suppose
$L(E^{(h)},s)$ has a simple zero at $1$ and
$L(E^{(hk)},1)\ne0$. Then the specialization at $t=0$ of the generic coordinate in
Definition~\ref{ex:arith-coordinate} is regular and satisfies
\[
 v_2u_{\rm prim}(0)
 =2j_{\rm rc}(h,k)-4\omega(h)-s(h)-s(hk)
       -\sum_{p\mid k}t_p+O_{E,\pi}(1).
\]
The error depends on the fixed old local types and fixed multiplier,
but not on $h$ or $k$. Here $0\leq t_p\leq2$, and $t_p=2$ in the
present full rational two-torsion setting. Only after a bound on
$\omega(k)$ is imposed may the companion sum be absorbed.
The finite-Sha assertion in this calibration is supplied by the
two stated analytic hypotheses.
\end{proposition}
\begin{proof}\leavevmode
First we justify the specialization of the generic coordinate.
The forward theorem, Lemma~\ref{lem:forward-rank}, gives rank one
and finite $\Sha$ for the first twist and rank zero and finite $\Sha$
for its companion. Hence the central rational primitive complex has
cohomology dimensions $(1,1)$ in degrees one and two and zero
elsewhere. Over $A=\Lambda_{(t)}$, cancel the unit disks outside
these two degrees; Nakayama's lemma permits these cancellations
because the other central cohomology groups vanish. The retained
cycle has nonzero central class by (4.5), the height formula at the
known simple zero, and its specified nonzero multipliers. In the
remaining two-term free complex this forces a nonzero generic kernel.
Semicontinuity bounds its dimension by one. The Euler characteristic
is zero, so the generic dimensions are also $(1,1)$.
Lemma~\ref{lem:v10-determinant-specialization} now identifies the
specialization of the generic coordinate with the central tensor
computed in Lemma~\ref{ex:arith-central-specialization}, in the same
retained determinant basis. This step uses neither the numerical
calibration (4.6) nor the later uniform determinant bound.

For the ordinary compact Kummer complex $C_K$, universal coefficients
and the local degree-zero identifications give
\[
 H^1(C_K)_{\rm tors}=E^{(h)}(K)[2^\infty],\qquad
 0\to H^1(C_K)\otimes\Q_2/\Z_2\to\Sel_{2^\infty}
       \to H^2(C_K)_{\rm tors}\to0.
\]
The last finite group is Selmer modulo its maximal divisible subgroup.
Self-duality identifies the degree-three torsion with the degree-one
torsion and gives a unimodular pairing of the free degree-one and
degree-two parts. Thus their rational dimensions equal full Selmer
corank, including a possible divisible part of $\Sha$.

At known simple zeros, Lemma~\ref{lem:forward-rank}
gives rank one and finite $\Sha$.
For restriction of scalars and the product of the two twists, the
sum/difference isogenies are dual and compose to $[2]$.
Their free-lattice indices satisfy $IJ=2$, with regulator ratio
$2/I^2$. At odd primes the differential lattices agree; at two the
pullback determinants have bounded total valuation.
Milne's isogeny ratio \cite[proof of I.7.3, (7.3.1)]{milne} follows
from duality without assuming the BSD leading-term equality.
It applies here with finite $\Sha$, canonical heights and pullback
measures.

At every active $p\mid h$, $K$ splits, so the isogeny local factors
cancel exactly. At either split place the full/Kummer quotient
has length
$\dim_{\F_2}E[2]^{\operatorname{Fr}_p}=2$, since all of $E[2]$
is rational. The two places therefore contribute exactly $-4$. At a companion prime $p\mid k$, the rational component number
of the ramified twist is $\#E[2]^{\operatorname{Fr}_p}$, whereas
the restriction of scalars has connected special fiber.
Thus $t_p=\dim_{\F_2}E[2]^{\operatorname{Fr}_p}\le2$;
in the present full rational two-torsion case it is exactly two.
The result is
\[
 v_2u_{\rm prim}(0)
 =2j_{\rm rc}(h,k)-4\omega(h)-s(h)-s(hk)
                    -\sum_{p\mid k}t_p+O_{E,\pi}(1).             \tag{4.6}
\]
Only the companion sum can be absorbed when $\omega(k)$ is bounded.
No class number or unbounded degree-two torsion length is absorbed.

Here is the determinant and isogeny bookkeeping in more detail. This
calculation is used only at the known simple-zero constituents, where
the forward theorems give finite $\Sha$. Put
$\mathcal A=\operatorname{Res}_{K/\Q}(E^{(h)}_K)$ and
$\mathcal B=E^{(h)}\times E^{(hk)}$, and let $g_{\mathcal A}$ and
$s_{\mathcal A}$ be the lengths of their two-primary rational torsion
and Tate--Shafarevich group, respectively, with analogous notation
for $\mathcal B$. The ordinary Kummer complex has torsion lengths
$g_{\mathcal A},s_{\mathcal A},g_{\mathcal A}$ in degrees one, two
and three. A torsion disk with cokernel of length $e$ in degree $i$
contributes $(-1)^{i+1}e$ to the inverse determinant coordinate.
The central class $Lm_iP_{h,k}^{\rm rc}$ and its paired functional
each have free-coordinate valuation $j_{\rm rc}(h,k)+v_2(Lm_i)$.
Thus their ordinary Kummer coordinate has valuation
$2j_{\rm rc}+2v_2(Lm_i)+2g_{\mathcal A}-s_{\mathcal A}$.

At either place above the split moving prime $r_i$, the central
full/Kummer quotient has degree-two length
\[
 v_2(N_i),\qquad
 N_i=r_i+1-\chi_h(r_i)a_{r_i}
     =(r_i-1)-\chi_h(r_i)m_i.
\]
Use the moving-character construction with
$r_i\equiv1\pmod{2^{B_i}}$, $B_i>c_E$ tending to infinity.
Since $v_2(m_i)\le c_E$, both local lengths equal $v_2(m_i)$.
Their total negative contribution cancels the two occurrences of
$v_2(m_i)$ in the class-functional tensor. This equality persists
in the common limit. At each active split place the corresponding
quotient has only degree-two cohomology $E[2]$, of length two.
The two places above each active prime therefore contribute $-4$.

For clarity, orient the isogeny indices as
\[
 I=[\mathcal B(\Q)_{\rm free}:\phi\mathcal A(\Q)_{\rm free}],
 \qquad
 J=[\mathcal A(\Q)_{\rm free}:\psi\mathcal B(\Q)_{\rm free}].
\]
Then $IJ=2$ and
$\operatorname{Reg}_{\mathcal B}/\operatorname{Reg}_{\mathcal A}
=2/I^2$. Choose global generators $\omega_{\mathcal A}$ and
$\omega_{\mathcal B}$ of the top cotangent lattices of their
N\'eron models, and use their period measures over $\Q$ throughout.
The integral pullbacks
$\phi^*\omega_{\mathcal B}=a\omega_{\mathcal A}$ and
$\psi^*\omega_{\mathcal A}=b\omega_{\mathcal B}$ satisfy $ab=4$.
Thus $a,b\in\{\pm1,\pm2,\pm4\}$. Real change of variables gives
\[
 |a|\Omega_{\mathcal A}
 =\frac{\#\ker\phi_{\mathbf R}}{\#\operatorname{coker}\phi_{\mathbf R}}
   \Omega_{\mathcal B},
\]
where both finite groups have order at most four. Hence the period
ratio has uniformly bounded two-adic valuation, without a separate
discriminant normalization. With these differential choices the
finite measure factors are the products $C_{\mathcal A},C_{\mathcal B}$
of component numbers. The arithmetic isogeny identity gives
\[
 s_{\mathcal A}-s_{\mathcal B}
 =2g_{\mathcal A}-2g_{\mathcal B}
  +v_2\frac{\operatorname{Reg}_{\mathcal B}}{\operatorname{Reg}_{\mathcal A}}
  +v_2\frac{\Omega_{\mathcal B}}{\Omega_{\mathcal A}}
  +v_2\frac{C_{\mathcal B}}{C_{\mathcal A}}.
\]
The $2g_{\mathcal A}$ terms cancel on substitution into the
determinant formula. The remaining rational torsion is uniformly
bounded by the fixed local types at two. Component factors cancel
at every split prime. At $p\mid k$, put $E_0=E^{(h)}$; it has good reduction, and
$K_p/\Q_p$ is a tame ramified quadratic extension.
By \cite[Theorem~5.3 and Remark~5.4.1, pp.~299--300]{edixhoven},
the component group of the ramified twist is the fixed subgroup
of the good special fiber under its descent action $[-1]$:
\[
 \Phi(E_0^{(k)})\simeq\widetilde E_0[2].
\]
This is an isomorphism of finite \'etale residue-field group
schemes, so its rational component number is
$\#E[2]^{\operatorname{Fr}_p}$.
The Weil restriction of the good model has special fiber an
extension of $\widetilde E_0$ by a vector group, and is connected.
Thus its component number is one.
All other component factors cancel at split primes or equal one
at the remaining good primes. Consequently
\[
 v_2(C_{\mathcal B}/C_{\mathcal A})
   =\sum_{p\mid k}t_p,\qquad
 t_p=\dim_{\F_2}E[2]^{\operatorname{Fr}_p}.
\]
Substitution into the coordinate containing $-s_{\mathcal A}$
gives the negative sign in (4.6).
In the full rational two-torsion case each $t_p=2$.
\end{proof}
\begin{lemma}[Unused-prime comparison and one corrected coordinate]\label{ex:arith-common-coordinate}
At an unused good split prime, use the actual primitive-to-relaxed
maps and the retained class and functional above. Write $z$ for
the scalar Frobenius value. With $e_q=a_q-z-z^{-1}$ and $l_q$
the two-place quotient factor,
they satisfy (4.7) with $\varepsilon_q\in\Lambda^\times$.
For a finite binary family the universal correction $\eta$
lies in $\mathscr O[G]$, and its corrected coordinate
$\widetilde u=\eta u_{\rm rel}$ satisfies
$\widetilde u_x=\varepsilon_xu_{{\rm prim},x}$ at every character,
with $\varepsilon_x\in\Lambda^\times$.
At any character where the primitive coordinate is regular at
$t=0$, this identity preserves its central valuation and vanishing.
Only the product is specialized; $\eta_x$ itself may have a pole.
\end{lemma}
\begin{proof}\leavevmode
At an unused prime the exact point factor is $e_q=a_q-z-z^{-1}$.
The second factor comes from the same closed functional, as follows.
Let $i:C_P\to C_R$ be the primitive-to-relaxed map, and let
$j:C_S\to C_{P^\perp}$ be the paired strict-to-primitive map.
All old orthogonal conditions and their homotopies are retained.
At $q$ these are the maps from the unramified condition to the full
condition and from zero to the unramified condition. Naturality of
the cup pairing \cite[Sections 6.4.1--6.4.3]{nekovar} gives
\[
             \beta_R(v,iw)=\beta_P(jv,w).
\]
The actual norm identity is $z_R=e_q i(z_P)$, with its retained
boundary. The generic paired strict lift satisfies
$jz_S=e_q^c z_{P^\perp}$. Conjugation sends $z$ to $z^{-1}$, so
$e_q^c=e_q$. Consequently
\[
                 \lambda_R\circ i=e_q\lambda_P.
\]
Thus the class and its paired functional each contribute one
$e_q$; no integral self-duality of the primitive condition is assumed.
The two quotient determinants satisfy
\[
 \delta_q=\det(z\operatorname{Fr}_q-q)
              \det(z^{-1}\operatorname{Fr}_q-q)=q^4l_q.
\]
In the original integral determinant bases, the support triangle
therefore gives
\[
       l_q u_{\rm rel}=e_q^2\varepsilon_q u_{\rm prim},
             \qquad \varepsilon_q\in\Lambda^\times.             \tag{4.7}
\]
To see the unit assertion, retain the finite mapping-cone bases.
Their comparison is a sign; replacing them by the retained homotopy
equivalent free models contributes only determinants of unit disks.
These are $A_i$-units. Their limits and inverse limits are
$\Lambda$-units, and $q^{-4}$ is one also. An arbitrary
$\mathscr O$-unit would not suffice.

We check the integrality needed for the universal correction.
The graph-wide prescription at the start of this section applies
also to this unused prime. Modulo the maximal ideal
$(2,I_G)$ of $\mathscr O[G]$, every group element and fixed quadratic
sign becomes one, and the scalar $z$ becomes $1+t$. Since $E[2]$
is rational, $a_q$ is even at every good odd prime. Thus
\[
 \overline e_q=(1+t)+(1+t)^{-1}=\frac{t^2}{1+t}
       \ne0\quad\hbox{in }\mathbf F_2((t)).
\]
Consequently $e_q$ is a unit in $\mathscr O[G]$. This assertion uses
the prescribed exponent at every $q\in Q$, not merely at the primes
ramified in a particular constituent. With
\[
 l_q=\det(1-q^{-1}z\operatorname{Fr}_q)
       \det(1-q^{-1}z^{-1}\operatorname{Fr}_q),
\]
the trace and determinant of $\operatorname{Fr}_q$ are $a_q$ and $q$.
Expanding these two degree-two polynomials gives the exact identity
\[
 q^2l_q=e_q^2+(q-1)\bigl(q-1-(z+z^{-1})e_q\bigr).
\]
Because $q$ is odd and $e_q$ is a unit, this identity implies
$l_q/e_q^2-1\in2\mathscr O[G]$. Set
$v_q=(l_q/e_q^2-1)/2\in\mathscr O[G]$.
Then $1+v_q(1+g_q)$ selects $l_q/e_q^2$ at an inactive character
and one at an active character. Their finite product is therefore
one integral $\eta\in\mathscr O[G]$, without any division by the
group order. Put $\widetilde u=\eta u_{\rm rel}$.
It already is one common family; no interpolation of the individual
$\varepsilon_x$ is required. Formula (4.7) makes
$\widetilde u_x=\varepsilon_xu_{{\rm prim},x}$.
The central zero and valuations therefore transfer unchanged.
The multiplier $\eta_x$ may itself have a pole at $t=0$.
Only the product $\widetilde u_x$, identified with the regular
primitive coordinate times a $\Lambda$-unit, is specialized there.
\paragraph{Simultaneous prescriptions and local lifts.}
We spell out why the preceding calculation applies to the whole
graph. First choose its actual finite set $Q$ by
Proposition~\ref{ex:app-final-family}. Its primes are good, odd,
split in $K$, and outside $2ND_K$. The theta prime has already
been chosen and is an exclusion. Choose one prime of $K$ above
each $q\in Q$. In Proposition~\ref{ex:arith-moving-character},
``active list'' can be this entire list of orientations: the
construction uses only these splitting and avoidance properties,
not ramification in any particular quadratic constituent.

Here is the compatibility check in its class-group presentation.
Adjoin orientations above $2N$ and auxiliary split primes generating
$\operatorname{Cl}(K)$, all in distinct conjugate pairs. Assign weight
one to every member of $Q$ and weight zero to the other primes.
For the relation basis $m^{(a)}$ and
$\alpha_a=\beta_a/\overline\beta_a$ used there, a relation
$\prod_a\alpha_a^{b_a}\in
K(\mu_{2^\infty})^{\times 2^n}$ implies
\[
 2^n\mid\sum_a b_am_j^{(a)}\qquad(j\text{ corresponding to }Q).
\]
Indeed, a root belongs to a finite cyclotomic layer, which is
unramified at each of these odd primes, and its valuation is integral.
Multiplying by the weights and summing proves that the prescribed
Kummer multipliers respect every relation. At a resulting split
prime $r$, their power-residue values therefore satisfy every
relation in the presentation of $\operatorname{Pic}(\mathbf Z+r\mathcal
O_K)$. The resulting character has value one at every chosen
orientation in $Q$, value zero at $2N$, and is surjective on the
residue-ratio kernel, hence on inertia at $r$. The conjugate
orientation has the negative value because the product ideal is
rational principal. The radical-fixing subgroup (4.4c) preserves
all these prescriptions while imposing the two trace bounds.
Its level depends only on the Tate image. Enlarging $Q$ consequently
changes the finite Chebotarev field, but not $c_E$.

Perform this construction separately at each coefficient stage,
excluding $Q$, the theta prime, the old support and the previously
chosen moving primes. Taking $n_i\geq2i$ gives
$(1+t)^{2^{n_i}}=1$ in $A_i$, so these are actual finite characters.
They need not be restrictions of one character of $G_K$: the
retained inner limit is the limit of the finite diagrams of
Lemma~\ref{ex:limits-inner-limit}. At each fixed graph orientation
their values are already exactly $1+t$, so that identity survives
every coefficient reduction and that limit. The norm identities
(4.1), their coefficient maps, and their boundary witnesses hold
before character evaluation. This is the same construction for
all constituents. At $t=0$ the moving character is trivial, so
its added prescriptions do not alter the central quadratic family.

There is a local-lift issue only for the universal quadratic
character. Write $g_q\in G$ for its inertia generator at $q$;
thus $x(g_q)=(-1)^{\lambda_q(x)}$. Fix a lift of arithmetic
Frobenius at the chosen orientation and define $z_q$ using that
lift. The moving character is unramified there, and its value is
$1+t$. The fixed quadratic character is unramified there also.
Every other graph prime contributes an unramified quadratic
character at $q$. Changing the Frobenius lift can therefore only
replace $z_q$ by $g_qz_q$. For an inactive character $x$,
$x(g_q)=1$, and $x(z_q)$ is the actual unramified Frobenius scalar.
For an active character its value is immaterial to the correction,
which is one. Thus we never assign an unramified Frobenius
eigenvalue to a ramified constituent.

This independence also holds inside $R=\mathscr O[G]$, before
evaluation. Both choices of $z_q$ have reduction $1+t$ modulo
$(2,I_G)$, so both corresponding $e_q$ are units. Put
$r_q=l_q/e_q^2$ and $r'_q=l'_q/(e'_q)^2$. They agree in
$R/(g_q-1)$, hence $(1+g_q)(r'_q-r_q)=0$. Since $R$ is
two-torsion-free and $r_q=1+2v_q$, $r'_q=1+2v'_q$, this gives
\[
 (1+g_q)(v'_q-v_q)=0.
\]
The correction $1+(1+g_q)v_q$ is consequently independent of the
lift. Reversing the orientation replaces $z_q$ by $z_q^{-1}$
and leaves both $e_q$ and $l_q$ unchanged.

For clarity, write $s_q=z_q+z_q^{-1}$. Direct expansion gives
\[
 q^2l_q=(q-1)^2+a_q^2-a_q(q+1)s_q+qs_q^2,
 \qquad
 v_q=\frac{q-1}{2q^2e_q^2}
       \bigl(-(q+1)e_q^2+(q-1)-s_qe_q\bigr).
\]
Every term in the second expression is integral in $R$ because
$q$ is odd and $e_q$ is a unit. These inverses are taken only
after the inner limit and localization at $(2)$, never in the
finite ring $A_i[G]$. Reduction modulo $(2,I_G)$ gives
$\overline e_q=t^2/(1+t)$ and
$\overline l_q=t^4/(1+t)^2$. The latter identity also holds at
each sign evaluation modulo two. At augmentation the local
two-place determinant is $q^4l_q$, so its residual $t$-order is
exactly four. Multiplying the correction factors over the finite
set $Q$ proves the asserted universal integrality and all the
characterwise identities simultaneously.
\end{proof}

\begin{lemma}[The unknown base has zero coordinate if its derivative vanishes]\label{ex:arith-base-zero}
Use the base $h_0=1$ in the final fixed-$K$ construction, with
$L(E^{(k)},1)\ne0$ and full two-primary Selmer corank one for
$E/K$. If $L'(E/K,1)=0$, then
$u_{{\rm prim},0}(0)=\widetilde u_0(0)=0$ for the same coordinates
as above. The primitive coordinate is regular at $(t)$ by the local
argument below, and its corrected relaxed coordinate has the same
property by Lemma~\ref{ex:arith-common-coordinate}. No finite
$\Sha$ assertion at the base is required.
\end{lemma}
\begin{proof}\leavevmode
Put $A=\Lambda_{(t)}$, with residue field $\mathbf Q_2$ and
fraction field $F$. Use the actual finite free primitive complex
and its retained class and functional over $A$. The latter is
integral over $A$ by Proposition~\ref{ex:arith-paired-functional}.
The conductor-one height formula makes the specified point torsion
when $L'(E/K,1)=0$. Thus its rational Kummer class vanishes, and
the retained boundary (4.5) makes the specialization of the actual
cycle $z$ zero. No new normalization of $L$ is made here.

Full Selmer corank one gives central dimensions $(1,1)$ in degrees
one and two and zero in every other degree. This uses the full
Selmer group, including a possible divisible part of $\Sha$.
Over the DVR $A$, a bounded finite free complex decomposes, by
successive Smith reductions, into free zero-differential summands
and two-term disks $[A\xrightarrow{t^r}A]$, up to units. A disk
with $r>0$ contributes one additional residue-field cohomology
class in each of its two adjacent degrees. Semicontinuity and
Euler characteristic therefore leave only two generic possibilities:
the complex is acyclic, or its dimensions are $(1,1)$ in degrees
one and two.

In the acyclic case the specified coordinate is zero by definition.
In the second case the generic and central dimensions agree in
every degree. Lemma~\ref{lem:v10-determinant-specialization} applies:
no disk with $r>0$ occurs. All cancellable
disks have unit differential, so integral determinant base change
identifies the coordinate of $[z]\otimes[\lambda_z]$ with the
coordinate of its central tensor, in the retained determinant basis.
That tensor is zero because $[z]_0=0$. This proves regularity at
$(t)$ and $u_{{\rm prim},0}(0)=0$ directly. Finally
$\widetilde u_0=\varepsilon_0u_{{\rm prim},0}$ with
$\varepsilon_0\in\Lambda^\times$ by
Lemma~\ref{ex:arith-common-coordinate}, so the corrected coordinate
is also regular and zero at the center. This proof uses neither
the numerical lower bound of Proposition~\ref{ex:limits-uniform-lower-bound}
nor finite $\Sha$ at the unknown base.
\end{proof}

% End 03a-arithmetic.tex


% Begin 03b-limits.tex
\section{The two limits and determinant transfer}\label{sec:limits}
This section uses the actual diagrams of Section~\ref{sec:arithmetic}.
Their cohomology ranks alone would not suffice: maps to original
cochains, local comparisons, cup products, classes and boundary
witnesses are retained together. The same retained data include the
integral duality maps, their homotopy inverses, and the homotopies
expressing compatibility of duality with localization.

\subsection{Finite free models and the inner limit}
\begin{lemma}[Finite free models and retained inner-limit diagrams]\label{ex:limits-inner-limit}
Let $A$ be an Artin local ring with nilpotent maximal ideal, and
let a complex of flat $A$-modules have finite bounded residual
cohomology. This includes the actual continuous cochains with
finite free coefficient modules used here. It admits a bounded finite free model
with actual inclusion, projection and contraction homotopy.
Every fixed finite diagram of maps and homotopies, including the
duality and local-condition data used here, transports to these
models.

For one fixed active support and binary family, retain a tuple
whose model ranks are bounded over the finite coefficient stages.
One fixed nonprincipal ultrafilter gives a diagram over $\Lambda$
at all coefficient precisions, preserving its matrix identities
and retained homotopy inverses. This assertion imposes no rank
bound uniform over all active supports.
\end{lemma}
\begin{proof}\leavevmode
Let $A$ be an Artin local ring with nilpotent maximal ideal $\mathfrak m$.
A flat $A$-module is free, including in infinite rank: lift a basis
modulo $\mathfrak m$, use nilpotence for surjectivity, and flatness and
nilpotence for vanishing of the kernel. A continuous finite-coefficient
cochain module is a filtered union of modules of functions on finite
clopen partitions, hence is flat.

Suppose the residual cohomology of a complex of such modules is finite
and bounded. Split the residual complex into cohomology and disks,
and lift graded bases. Let $d_0$ be the lifted disk differential and
$d=d_0+\epsilon$, where $\epsilon$ has entries in $\mathfrak m$.
Use the graded decomposition into cohomology, disk sources and disk
targets. Let $i_0,p_0$ include and project the cohomology summand,
and let $h_0$ invert $d_0$ from disk targets to disk sources, zero
on the other summands. Then
\[
 p_0i_0=1,\quad h_0^2=h_0i_0=p_0h_0=0,\quad
 1-i_0p_0=d_0h_0+h_0d_0.
\]
With this normalized contraction, the finite geometric series give
\[
 \begin{aligned}
 d_C&=p_0\epsilon(1+h_0\epsilon)^{-1}i_0,&
 i&=(1+h_0\epsilon)^{-1}i_0,\\
 p&=p_0(1+\epsilon h_0)^{-1},&
 h&=h_0(1+\epsilon h_0)^{-1}.
 \end{aligned}
\]
Substitution using $d^2=0$ gives $pi=1$ and $1-ip=dh+hd$.
The resulting bounded finite free model carries actual maps and
homotopies to the original cochains. Transport any fixed finite
collection of arithmetic maps by $f\mapsto p'fi$. The identity
$1-ip=dh+hd$ transports their compatibility homotopies as well,
including the second homotopies in the cup-product diagrams.

Fix one nonprincipal ultrafilter on coefficient stages. Bounded ranks
can be fixed on a set in that ultrafilter. Each finite tuple of matrix
entries then has a compatible limit at every coefficient precision,
giving matrices over $\Lambda$. Every polynomial identity between
the matrices survives. Retained homotopy inverses remain homotopy
inverses. Mere stagewise quasi-isomorphisms without those maps would
not justify a nonflat coefficient change.
\end{proof}

\begin{lemma}[The reduced generic residual rank bound]\label{ex:limits-residual-bound}
Fix $E$, an integer $A\geq0$, and a number $j$ of derivative labels.
For the ring instances with $\omega(k)\leq A$, reduction to
characteristic two followed by the stated support comparison
and inversion of $t$ gives a generic residual model of rank at most
\[
 N(E,A,j)\leq56\omega(2N)+8A+56j+70.
\]
This bound is independent of the number and sizes of active primes
and of $v_2(r_i-1)$. It bounds the reduced generic model, not the
whole retained support-comparison tuple before $t$ is inverted.
\end{lemma}
\begin{proof}\leavevmode
All ranks in this proof are sums of the ranks in all degrees of one
complex. They are not the size of the diagram retaining all comparison
maps. Put $v=\omega(2N)$ and $m=v+1+j$. Let $S_0$ be the finite
places of $K$ above $2N$, the moving prime, and the derivative primes.
Let $T$ be the set of rational primes below $S_0$, together with the
primes ramified in $K/\Q$, and put $s=|T|$. Since $2\mid2N$,
\[
                    s\le v+A+1+j.                            \tag{5.1a}
\]
Primes added to $T$ solely because they ramify in $K$ are used for
the global estimate. They do not impose additional Selmer conditions.

We first give the global estimate at the residue field
$B_i/(t)=\F_2$, where $B_i=\F_2[t]/(t^i)$. The residual coefficient
module is the trivial module $W=E[2]=\F_2^2$. Write $S_T$ for all
finite places of $K$ above $T$, and let $r$ be the number of finite
rational primes ramified in $K$. Kummer theory gives the exact sequence
\[
0\longrightarrow\mathcal O_{K,S_T}^{\times}/
                 (\mathcal O_{K,S_T}^{\times})^2
\longrightarrow H^1(G_{K,S_T},\F_2)
\longrightarrow\operatorname{Cl}(\mathcal O_{K,S_T})[2]
\longrightarrow0.
\]
Here $2\in T$ and $\mu_2=\F_2$ as a Galois module. The unit term
has dimension $1+|S_T|$: the unit rank is $|S_T|$ and the roots of
unity contribute one dimension. Moreover
\[
 |S_T|\le2s-r,\qquad
 \dim_{\F_2}\operatorname{Cl}(\mathcal O_{K,S_T})[2]\le r-1.
                                                               \tag{5.1b}
\]
For completeness, the class-group inequality follows directly from
the imaginary quadratic genus argument. Conjugation acts by inversion
on the ideal class group. If $[\mathfrak a]$ has order dividing two,
write $\overline{\mathfrak a}=(\alpha)\mathfrak a$. Taking ideal
norms gives $N_{K/\Q}\alpha=1$, since this norm is positive. Hilbert
90 gives $\alpha=\beta/\bar\beta$ (take $\beta=1+\alpha$
unless $\alpha=-1$, and take $\beta=\sqrt{k}$ in that case), so
$\beta\mathfrak a$ is a
conjugation-invariant fractional ideal. Modulo rational principal
ideals such ideals are generated by the ramified prime ideals.
Their classes have exponent two, and the principal ideal of
$\sqrt{d}$, where $d$ is the squarefree radicand of $K$, supplies
a nonempty relation among them. Thus their
span has dimension at most $r-1$. The excluded field $\Q(i)$
would also satisfy the inequality, since its class group is trivial.
Finally an $S$-class group is a quotient of the ordinary class group,
and the two-rank of a quotient of a finite abelian group cannot
increase. This proves (5.1b), including the possible ramification
at two. The bound on $|S_T|$ counts one prime over each ramified
rational prime and at most two over each other prime.

Inflation from $G_{K,S_0}$ into $G_{K,S_T}$ is injective on $H^1$
with trivial coefficients. Consequently
\[
 h^0(G_{K,S_0},W)=2,\qquad
 h^1(G_{K,S_0},W)\le4s.
\]
The global Euler characteristic and cohomological-dimension statements
\cite[I.4.10 and I.5.1]{milne}, applied over the imaginary quadratic
field $K$, give
\[
 h^0-h^1+h^2=-2,\quad H^{>2}=0,\quad
 h^2=h^1-4.
\]
Thus the total residual cohomology dimension of the global complex
is at most $8s-2$, and in particular at most $8s+2$.
The latter weaker bound can also be obtained without genus theory:
the maximal pro-two group over $\Q$ unramified outside $T$ and
infinity has $s+1$ generators, by the squareclasses of $-1$ and
the primes in $T$; its subgroup for $K$ has at most $2s+1$
generators by the index-two Schreier bound. The maximal pro-two
group for $G_{K,S_0}$ is a quotient of this subgroup. Hence
$h^1(G_{K,S_0},W)\le4s+2$, giving total dimension at most $8s+2$.
In either computation the $H^2$ estimate is supplied by the global
Euler characteristic, not by a generator estimate alone.

We next bound the local terms in the Selmer cone (4.3). For a finite
place $w$, put $d_w=[K_w:\Q_2]$ if $w\mid2$ and $d_w=0$
otherwise. Local duality and the local Euler characteristic
\cite[I.2.3 and I.2.8]{milne} give, for the trivial module $W$,
\[
 (h^0,h^1,h^2)(K_w,W)=(2,4+2d_w,2).
\]
At an old Kummer place, the derived condition has residual dimensions
$(2,2+d_w)$ in degrees zero and one. Indeed
$E^{(h)}(K_w)^\wedge_2$ is the direct sum of a free
$\Z_2$-module of rank $d_w$ and finite two-primary torsion;
that torsion has two generators because $E^{(h)}[2]$ is rational.
Derived reduction retains its two-dimensional degree-zero torsion.
For the other conditions at the odd moving or derivative places,
the following table records the sum of the dimensions of the local
condition and of the full local complex. Cone shifts do not change
this sum.
\[
\begin{array}{c|c|c}
\text{condition }U_w&\sum h^a(U_w)&
       \sum h^a(U_w)+\sum h^a(K_w,W)\\ \hline
\text{old Kummer}&4+d_w&12+3d_w\\
\text{unramified, }w\nmid2&4&12\\
0,\ w\nmid2&0&8\\
\text{full, }w\nmid2&8&16\\
\text{derivative strict}&2&10\\
\text{derivative finite or singular}&4&12\\
\text{derivative relaxed}&6&14
\end{array}
\]
The last three rows use exactly the conditions retaining $D^0$
in Lemma~\ref{ex:limits-derivative-switch}; the relaxed condition
does not include $D^2$. There are at most two places over each
rational prime, and $\sum_{w\mid2}d_w=2$. The old prime two
therefore costs at most $30$, any other old rational prime at most
$24$, the moving prime at most $32$, and a derivative prime at
most $28$ (in fact it is inert and costs at most $14$).
Each is bounded by $48$. Thus all finite local terms cost at most
$48m$. No prime dividing $k$ alone is charged a local condition.

There are no real places of $K$. Its one complex place has full
cohomology $W$ in degree zero only. With the usual modified
archimedean convention its local complex is zero; even retaining
both an ordinary full term and a full local condition costs at
most four additional generators. Complex conjugation used to
transport duality is not an additional real local condition over
$K$. In particular no infinite two-periodic real resolution is
being included in this count.

At each finite coefficient stage, apply
Lemma~\ref{ex:limits-inner-limit} separately to these global and
local complexes and transport the actual local maps. The mapping
cone then has a finite free model whose total rank is bounded by
the sum just computed. Further cancellation can only decrease it.
The estimate applies before inverting $t$ to the smaller-support
model, uniformly in the moving prime and in the order of its
finite character: reduction modulo $t$ kills that character.
Allowing the weaker global estimate and the archimedean allowance
still gives
\[
 \begin{split}
 \operatorname{rank}_{\rm tot} C_{S_0}
 &\le(8s+2)+48m+4\\
 &\le8(s+1)+6+48(v+1+j)\\
 &\le56v+8A+56j+70.
 \end{split}\tag{5.1}
\]

It remains to transfer this estimate to the original support.
First consider a scalar ring instance. Reduction modulo two kills
$\chi_h$, so the coefficient system is unramified at every active
prime. The actual support comparison of
\cite[Proposition 7.8.8(ii)]{nekovar} therefore identifies the
smaller-support complex with the complex having unramified
conditions at those primes. At an oriented active place, reduction
of (4.4) has
\[
 C=0,\qquad A=D=\big((1+t)^e-1\big)I_2,
\]
where $e$ is odd; at the conjugate place the exponent is $-e$.
Both $((1+t)^e-1)/t$ and $((1+t)^{-e}-1)/t$ are units of
$\F_2[[t]]$. Thus the full, unramified and singular comparison
blocks all become contractible after inverting $t$. For each fixed
finite active support retain the comparison maps and homotopies,
take the inner limit, and only then invert $t$. Every active
comparison cone is then contractible, so the original complex has
the same generic residual cohomology as the bounded smaller-support
model. No count of active primes remains.

For clarity, a universal binary character is not trivial merely
after reduction modulo two. In the binary setting make the preceding
calculation after augmentation $\F_2[G]\to\F_2$. After the inner
limit and inversion of $t$, the ring $\F_2((t))[G]$ is local with
nilpotent augmentation ideal. A finite free model over this ring
has a minimal model obtained by cancelling unit entries; its ranks
are the dimensions of the cohomology after augmentation, and hence
obey (5.1). These are ranks over $\F_2((t))[G]$, not dimensions
over $\F_2((t))$. For scalar instances they are precisely the
generic residual ranks needed in the outer limit. The whole retained
comparison diagram can still grow with $G$ or with the active
support; no bound for that diagram is asserted. Finally the minimum
ranks over $\mathscr O$ are the dimensions after reduction to
$\F_2((t))$, so the same estimate bounds the scalar
$\mathscr O$-minimal models used below.
\end{proof}

\subsection{Evaluation and the outer valuation ring}
\begin{lemma}[Evaluation, coefficient change and the outer DVR]\label{ex:limits-outer-evaluation}
For the retained global models and for the primitive and derivative
Selmer models, evaluations at sequences of Galois elements satisfy
(5.2). A cycle defines an abstract crossed homomorphism of the
indicated product group. The induced map on $H^1$ is injective,
and the finite detecting-list comparison remains injective after
field extension. Direct evaluation of other Selmer models is asserted
only when their local maps induce isomorphisms on residual $H^0$
and injections on residual $H^1$. For orthogonal zero/full comparisons,
the assertion may instead be transported at a test ring where the
actual local comparison complexes are contractible.

For a sequence of ring instances with $\omega(k_j)\leq A$ and
a fixed number of derivative labels as in
Lemma~\ref{ex:limits-residual-bound}, the quotient
$\mathscr O_\infty$ in (5.3) is a DVR with uniformizer two.
The bounded minimal-model matrices, evaluations, duality maps and
their compatibility homotopies pass to it. Sequences whose
differentials tend to zero give actual crossed homomorphisms
there. No continuity assertion on a single original Galois
group, or completeness of $\mathscr O_\infty$, is required.
\end{lemma}
\begin{proof}\leavevmode
For fixed $K$ let $\Gamma_K=\prod_iG_K$ and include one simultaneous
complex conjugation. Evaluation of retained cochain maps at arbitrary
sequences $g=(g_i),h=(h_i)$ gives
\[
 e^1_gd=(\rho(g)-1)e^0,\qquad
 e^2_{g,h}d=\rho(g)e^1_h-e^1_{gh}+e^1_g.                          \tag{5.2}
\]
One fixed ultrafilter supplies these identities for all sequences.
A cycle therefore defines an abstract crossed homomorphism of the
product group. We do not assert that its restriction to one original
Galois group is a continuous pointwise limit.

Evaluation remains injective after field extension. At each finite
stage take representatives $V$ for $H^1$ and coefficient space $W$. Evaluate
$(c,w)$ by $b_c(g)-(\rho(g)-1)w$. The intersection of all kernels is
$(0,W^\Gamma)$. A finite list of at most $\dim V+\dim W$ evaluations
already has that intersection. Choose and pad these lists at the finite stages before taking their common limit. The comparison to $[W\to W^d]$ is an
isomorphism on $H^0$ and injective on $H^1$. Its cone has no residual
cohomology in degrees at most zero. Cancel unit disks there so that
the free cone starts in degree one; the same is true after every
coefficient extension. The long exact sequence gives the claimed
injection. For the primitive and derivative Selmer complexes the required
residual comparison with global cochains follows from their actual
local conditions. At an old Kummer place, derived reduction retains
$B(F)[2]$ in degree zero, isomorphically onto the local invariants,
and has the Kummer injection in degree one. The unramified condition
at a good place has the same degree-zero invariants as the full local
complex and injects in degree one. A full condition has the identity
map. Every derivative condition retains $D^0$ and includes its stated
degree-one subspace. The Selmer-cone exact sequence therefore gives
an isomorphism on global $H^0$ and an injection on global $H^1$.

This does not apply directly to the zero conditions in the retained
orthogonal complexes: at $(2,t)$ their local degree-zero invariants
need not vanish. Indeed, two zero local conditions with residual
invariant space $W$ can contribute
$(W\oplus W)/\operatorname{diag}(W)$ to Selmer $H^1$, and global
cochain evaluation kills this contribution. Such orthogonal complexes
are retained for duality, without asserting residual evaluation
injectivity. If an actual zero/full comparison is contractible at a
chosen test ring, its retained contraction identifies the two Selmer
models there and transports the evaluation injection. This use requires
that contraction at the chosen ring; it is not an assertion for arbitrary
unused relaxed primes at every scalar height-one localization.

Now allow $K_j,h_j$ to vary with $\omega(k_j)\le A$.
For copies $\mathscr O_j$ of $\mathscr O$, form
\[
 \mathscr O_\infty=
 \left(\prod_{\mathcal U}\mathscr O_j\right)/
 \{(x_j):v_2(x_j)\to_{\mathcal U}\infty\}.                         \tag{5.3}
\]
A nonzero element has valuation one fixed nonnegative integer on a
set in $\mathcal U$ and is that power of two times a unit.
Every nonzero ideal has an element of least valuation and is
principal. Thus (5.3) is a DVR with uniformizer two and characteristic-zero
fraction field. Completeness is unnecessary.

Integral minimal-model matrices and (5.2) pass to (5.3).
So do the retained duality maps, inverse homotopies and
duality-to-localization compatibility homotopies. After minimal
cancellation an integral homotopy equivalence is degreewise invertible:
its residual matrices are invertible, and Nakayama applies.
These inverse matrices and all compatibility identities pass to the
same outer quotient. Thus perfect duality and exact orthogonality
are preserved there; they are not inferred merely from isotropy.
Approximate cycles suffice: if $v_2(dw_j)\to\infty$, the degree-two
term in (5.2) tends to zero, so $w_\infty$ is an actual crossed
homomorphism. The finite detecting lists have bounded length by
(5.1); retain their cone contractions to preserve injection.
One must not replace approximate cycles by unexplained exact cycles.
\end{proof}

\subsection{One new prime detects two classes}
\begin{lemma}[The common kernel and finite realization of a rank-two test]\label{ex:limits-kernel-detection}
Use the exact kernel of sequences acting trivially on the full
Tate module, the relevant ring-class fields and quadratic
characters. On the characteristic-zero tests used below,
restriction of the evaluated cohomology to this kernel is injective.
When the moving scalar character is nonquadratic, two independent
classes have a rank-two localization at a single actual kernel
element. Finite Chebotarev realizes the required conjugation
coset and basis-cochain evaluations, with exact inner-limit
matrix identity (5.4) before any $\mathscr O$-minimal-model
denominators are applied.

At the residual test only scalar-inertia vanishing, bounded
cohomology and finite detecting lists are asserted. The quadratic
character exception at characteristic-zero height-one tests is
treated in Proposition~\ref{ex:limits-uniform-lower-bound}.
\end{lemma}
\begin{proof}\leavevmode
The exact kernel consists of sequences trivial on the full Tate
module, all relevant ring-class fields and quadratic characters.
The scalar element of Section~\ref{sec:arithmetic} is central modulo
this kernel, though not in the kernel itself.
For a cocycle zero on the kernel,
\[
 (u^4-1)b(g)=(g-1)b(z)
\]
makes it a coboundary over characteristic zero. The same fixed scalar
survives the outer quotient. To eliminate residual invariants, use moving-prime
inertia: its scalar difference is $t$, invertible in $\F_2((t))$.
These are two separate arguments.

On the characteristic-zero tests used for rank reduction, the fixed
finite matrix-span witness gives absolute irreducibility of each
rank-two summand. When the scalar character is not quadratic, moving
inertia has distinct characters on the two summands and simultaneous
conjugation exchanges them. The induced rank-four representation is
therefore absolutely irreducible with scalar endomorphisms. The outer
fraction field is one such characteristic-zero test; the quadratic
character exception at an inner test is treated separately below.

For two independent cohomology classes on these tests, the span of
their joint values on the exact kernel projects onto both copies of
this absolutely irreducible representation. A proper submodule with
both projections onto would be the graph of a scalar isomorphism,
contradicting independence. Thus the span is the full sum. If the
determinant of their two finite localizations vanished on every actual
kernel element, additivity and polarization would make it vanish on
the entire span, impossible. Hence one actual element supplies a
rank-two localization on these characteristic-zero tests.

The polarization identity itself is also valid in characteristic two,
but no residual irreducibility or residual rank-two localization is
asserted. Here $E[2]$ is trivial of rank two, so its scalar twist is
isotypic and the induced residual representation has multiplicity two.
At the residual test we use only scalar-inertia vanishing of $H^0$,
the bounded residual cohomology dimensions, and finite detecting lists
from intersections of linear evaluation kernels. Those lists need not
consist of one rank-two localization. The rank-two selection for the
uniform valuation bound occurs only after passage to the
characteristic-zero fraction field of the outer DVR.

Chebotarev is applied at finite stages, not to the product group.
Each of the finitely many original basis cochains in degrees zero
through two factors through a finite quotient. Include their quotient,
the required Tate precision, ring-class and quadratic fields and old
diagrams in one compositum. The element just found, multiplied by
conjugation, is a single actual automorphism of that compositum.
Its restrictions agree on intersections. Choose the new prime and
a place above it to realize that representative.

This also resolves the precision issue caused by localization.
Choose the new prime to match all original basis-cochain evaluations
at a common retained precision $n(i)\to\infty$. Taking the same inner
ultrafilter gives an exact equality of $\Lambda$-matrices
\[
                 L_j=E_{\sigma_j^2}.                            \tag{5.4}
\]
Only afterward compose with the inclusion $I_j$ from an
$\mathscr O$-minimal model:
$L_jI_j=E_{\sigma_j^2}I_j$.
The denominators in $I_j$ may depend on $j$; equality (5.4) makes this
irrelevant. Long $t$-jets are never treated as approximations in the
two-adic valuation of $\mathscr O$, where $t$ is a unit.
\end{proof}

\subsection{Switches and the varying-companion estimate}
\begin{lemma}[Derivative local conditions and determinant-preserving switch]\label{ex:limits-derivative-switch}
At a derivative prime use
$L_v=D^0\oplus(F_v\oplus S_v)[-1]\oplus D^2[-2]$ with the
transported pairing, and retain $D^0$ in all four conditions.
Take $U_{\rm str}=D^0$,
$U_{\rm rel}=D^0\oplus(F_v\oplus S_v)[-1]$,
$U_F=D^0\oplus F_v[-1]$, and $U_S=D^0\oplus S_v[-1]$.
Then strict and relaxed are orthogonal, the two plane conditions
are self-orthogonal, and their cone triangles are the ones below.

Over the stated fraction or test field, suppose the actual old
class spans generic $H^1$ and has nonzero finite localization $v$.
Its derivative from Proposition~\ref{ex:arith-finite-heegner}
has finite coordinate zero and singular coordinate $Jv$.
The strict generic $H^1$ is zero, the new generic $H^1$ is spanned
by this derivative, and the two class-functional determinant
coordinates have equal valuations, even when $v$ is not primitive
in the local integral lattice.
\end{lemma}
\begin{proof}\leavevmode
\paragraph{The strict and relaxed conditions at a derivative prime.}
The words strict and relaxed have a different local meaning at a
derivative prime from zero and full at an active or unused split
prime. At a derivative prime use the actual split local model
\[
 L_v=D^0\oplus(F_v\oplus S_v)[-1]\oplus D^2[-2]
\]
with zero differential at the retained precision, and retain its
pairing and isotropy homotopies. Here $D^0$ and $D^2$ are dual,
and $F_v,S_v$ are the finite and singular isotropic planes. Put
\[
 \begin{aligned}
 U_{\rm str}&=D^0,&
 U_{\rm rel}&=D^0\oplus(F_v\oplus S_v)[-1],\\
 U_F&=D^0\oplus F_v[-1],&
 U_S&=D^0\oplus S_v[-1].
 \end{aligned}
\]
Thus $U_{\rm str}^{\perp}=U_{\rm rel}$ and the two plane
conditions are self-orthogonal under the transported duality.
All four maps retain local degree-zero cohomology. In particular,
the relaxed condition here omits $D^2[-2]$ and is not the full
local complex. The actual cone triangles are consequently
\[
 C_{\rm str}\longrightarrow C_F\longrightarrow F_v[-1]
 \longrightarrow C_{\rm str}[1],\qquad
 C_{\rm str}\longrightarrow C_S\longrightarrow S_v[-1]
 \longrightarrow C_{\rm str}[1].
\]
These are the triangles used in the determinant switch. Zero/full
conditions at the other primes retain their previously stated
meaning. Exact annihilator and dimension assertions in the switch
argument are made after passage to the stated fraction or test
field, using these integral diagrams and their retained dualities.
At a new inert prime the finite and singular planes are free of rank
two, with cross pairing
$\beta(x,y)=\varepsilon\langle x,Jy\rangle$.
Suppose the old actual class $A$ spans generic $H^1$ and has
$f(A)=v\ne0$. Its derivative $B$ satisfies $f(B)=0$, $s(B)=Jv$.
The strict $H^1$ vanishes. The transported duality and localization
identities identify the relaxed localization image with the kernel
of its dual localization map, hence with its exact annihilator
\cite[Sections 6.3.3--6.3.4]{nekovar}. This image
contains $(v,0)$ and $(0,Jv)$, so these span
it. Hence the new generic $H^1$ is spanned by $B$.

The determinant comparison is exact even if $v$ is divisible.
In the fraction field choose $x$ with $v\wedge x$ of unit volume,
and similarly $y$ with $Jv\wedge y$ of unit volume.
The finite and singular determinant exact sequences have the same
strict factor, so their class-functional coordinates have equal
valuations. All strict torsion remains in that common factor.
Requiring $x$ integral would be wrong when $v$ is divisible.
\end{proof}

\begin{proposition}[Uniform ring lower bound and height-one regularity]\label{ex:limits-uniform-lower-bound}
Fix $E,\pi$, the old local types and a bound $A$.
For every ring instance satisfying the hypotheses of
Proposition~\ref{ex:arith-central-calibration} and
$\omega(k)\leq A$, the actual primitive coordinate satisfies
\[
 v_2u_{\rm prim}\geq-C(E,A),\qquad
 u_{\rm prim}\in\Lambda[1/2],
\]
where the first valuation is that of $\mathscr O$ and the
constant is independent of the active support and companion field.
Consequently its central specialization and (4.6) imply (3.5).

The regularity assertion also holds at the unknown full-corank-one
base used in Lemma~\ref{ex:arith-base-zero}. The uniform numerical
lower bound here is asserted only for the known simple-zero ring
instances; it is not asserted at that unknown base.
\end{proposition}
\begin{proof}\leavevmode
The starting generic kernel has rank one. Indeed, at a known simple
zero the ordinary class is nonzero. Equation (4.5), with its bounded
nonzero multiplier, makes the retained central class nonzero.
Over $\Lambda_{(t)}$ cancel the degree-zero and degree-three disks.
If the generic kernel of the resulting two-term free complex vanished,
the actual lifted cycle would vanish, contradicting its central value.
Semicontinuity bounds that kernel by the central dimension one.

Suppose no uniform lower bound exists. Choose instances with
$v_2u_j\to-\infty$, $\omega(k_j)\le A$, and write their minimal
completed models as $[\mathscr O_j^n\xrightarrow{D_j}\mathscr O_j^n]$,
fixing the bounded $n$ on an outer-large set. In each generic
one-dimensional kernel choose a primitive vector $v_j$.
Its outer image is nonzero, so the outer kernel dimension $s$ is
at least one. If $s>1$, one new derivative prime gives rank-two
localization including that vector. The relaxed localization image
is self-annihilating and contains the whole finite plane, hence equals
that plane. The strict kernel has dimension $s-2$ and the singular
switch adds no localized vector, so the new outer dimension is $s-2$.

The actual inner class is a nonzero scalar multiple of $v_j$.
Even if this scalar tends to zero, its generic localization is nonzero
on an outer-large set. Exact switch preservation therefore keeps
each inner determinant valuation. The new inner kernels again have
primitive vectors. Thus their outer dimension remains at least one;
a proposed drop to zero is already a contradiction.
There can be only finitely many drops. For each fixed number of them
(5.1) bounds the models; no bound uniform in an unbounded number of
switches is needed.

At terminal outer dimension one a maximal boundary minor $\Delta$
is nonzero. Its inner minors have one finite valuation $b$ on an
outer-large set. Determinant-one Schur operations in a square
two-term rank-one-kernel complex give
$u=\pm\alpha\beta/\Delta$, where $\alpha$ is an integral surviving
cycle coordinate and $\beta$ has valuation at least $-a$.
Consequently $v_2u_j\ge-a-b$, the required contradiction. We obtain
\[
                     v_2u_{\rm prim}\ge-C(E,A).                 \tag{5.5}
\]
Here (5.5) is asserted for the known simple-zero ring instances
used to obtain the varying-companion estimate. Their nonzero actual
central class supplied the rank-one start. It is not a uniform
valuation assertion at the unknown corank-one base.
Bounded matrix size alone would not prove it: $[2^j]$ is a
counterexample to that inference.

% Replace baseline lines 4112--4138.
\paragraph{Height-one regularity, including quadratic specializations.}
Fix a height-one prime $\mathfrak p\ne(2)$, put
$R=\Lambda_{\mathfrak p}$ and $k=\operatorname{Frac}(\Lambda/\mathfrak p)$,
and use the actual primitive scalar Selmer complex over $R$.
Here $R$ is a DVR and $k$ has characteristic zero. All subsequent
switches use the derivative conditions of
Lemma~\ref{ex:limits-derivative-switch}, retaining $D^0$.
The primitive old local conditions and these derivative conditions have
the evaluation injection of Lemma~\ref{ex:limits-outer-evaluation}.
Cancel the degree-zero and degree-three disks and then the unit disks.
Write $C=[R^s\xrightarrow{d}R^s]$ in degrees one and two, with
$d\otimes k=0$. Its generic kernel has dimension one. Choose a primitive
vector $v\in\ker d$. Such a vector exists because $\ker d$ is saturated
in $R^s$: divide a nonzero kernel vector by its largest common power
of a uniformizer. Thus $\bar v\ne0$ in
$V=H^1(C\otimes_R k)$ and $\dim_k V=s\ge1$.
The actual class $z$ is a nonzero fraction-field multiple of $v$.
We will only make switches whose finite localization is nonzero on
$\bar v$. Consequently their generic localization of $z$ is nonzero,
which is the hypothesis needed for the actual-class determinant switch.

If the scalar character specialized to $k$ is nonquadratic, the
rank-two detection of Lemma~\ref{ex:limits-kernel-detection} applies to
$\bar v$ and an independent vector whenever $s>1$. The finite
localization then has rank two. Exact duality makes the relaxed
localization image equal to the finite plane, so switching to the
singular plane lowers $s$ by two. We give the quadratic case in detail.

Let $\alpha$ be the specialized character on the index-two group.
Since the moving inertia generator has value $1+t$, quadraticity
implies $t(t+2)=0$ in $k$. The only possible primes are therefore
$(t)$ and $(t+2)$. The equality $\alpha^\tau=\alpha^{-1}=\alpha$
uses the actual conjugation $\tau$; it involves no automorphism of
$R$ or substitution for $t$. The action
\[
 c^\tau(g)=Jc(\tau^{-1}g\tau),\qquad J=\rho_T(\tau),
\]
on cocycles gives a $k$-linear involution on cohomology. Extend
$\alpha$ to the index-two overgroup by choosing its value on $\tau$
to be $+1$. The other extension differs by the character of $K$.
The induced rank-four representation is consequently
$W_+\oplus W_-$, with each $W_\pm$ of rank two. The matrix-span
witness in (H1)--(H2) and the paragraph following them makes their
restriction to the index-two group absolutely irreducible. Twisting
its matrices by nonzero scalar character values does not change their
linear span. The summands are nonisomorphic: an intertwiner restricts
to a scalar, whereas their actions of $\tau$ are $J$ and $-J$.
A scalar intertwining these two actions is zero since $2$ is invertible.
We write $V=V_+\oplus V_-$ for the resulting cohomology signs.

At a new inert derivative prime write $F,S$ for the finite and
singular planes. The action on $F$ is $J$, and that on $S$ is $-J$.
Indeed, the local rational Frobenius acts on the Tate module by $J$;
it fixes its square, the finite Frobenius generator, and conjugates
tame inertia by its $\ell$-th power. The retained congruences give
$\ell=-1$ in the limit, so pullback on the tame coordinate contributes
$\ell^{-1}=-1$. These assertions hold at every previously chosen
derivative prime as well. A common scalar sign would simply interchange
the names $+$ and $-$.
The transported cross pairing is exactly
\[
 \beta(x,y)=\varepsilon e_2(x,Jy),\qquad \varepsilon\in k^\times,
\]
from Lemma~\ref{ex:limits-derivative-switch}. The finite and singular
planes are isotropic. These are the retained local cup-product
identities, including the extra finite precision in (H12), not a
residual alternatingness assertion. Since $\det J=-1$ and $J^2=1$,
$e_2(x,Jy)=e_2(y,Jx)$. Choose nonzero $J$-eigenvectors $e_+,e_-$.
They form a basis and $e_2(e_+,e_-)\ne0$. The two local cohomology
signs are therefore
\[
 H_+=F(e_+)\oplus S(e_-),\qquad
 H_-=F(e_-)\oplus S(e_+).
\]
They are orthogonal nondegenerate symmetric hyperbolic planes.
In either plane the self-pairing of $aF(e_\pm)+bS(e_\mp)$ is
$2ab$ times a nonzero constant. Its only isotropic lines are thus
its finite and singular axes.

All old conditions are stable under this involution. Conjugate old
Kummer conditions are transported into each other; primitive ramified
full conditions occur at the conjugate pair; unused good conditions
are unramified; and both $F$ and $S$ at a previous derivative prime
are stable under the actions just computed. At the new prime let
$C_{\rm str}$ and $C_{\rm rel}$ have the actual truncated conditions
$D^0$ and $D^0\oplus(F\oplus S)[-1]$. Put
\[
 U=\operatorname{im}\bigl(H^1(C_{\rm rel}\otimes k)\longrightarrow
 F\oplus S\bigr),\qquad
 K_0=H^1(C_{\rm str}\otimes k).
\]
The retained localization-duality diagrams give $U=U^\perp$.
Specifically, localization's connecting map is its dual under the
strict/relaxed duality, so its kernel is both $U$ and $U^\perp$;
this is the exactness in \cite[6.3.3--6.3.4]{nekovar} used in
Lemma~\ref{ex:limits-derivative-switch}. Isotropy alone would not suffice.
The cone triangles, with $D^0$ retained, give exact sequences
\[
 0\longrightarrow K_0\longrightarrow V
   \longrightarrow U\cap F\longrightarrow0,
 \qquad
 0\longrightarrow K_0\longrightarrow V'
   \longrightarrow U\cap S\longrightarrow0,                 \tag{5.5a}
\]
where $V'$ is the switched fiber cohomology. There is no preceding
local degree-zero term, since the quotient conditions are
$F[-1]$ and $S[-1]$.
The space $U$ is stable under signs. Each $U_\pm\subset H_\pm$
is its own annihilator in $H_\pm$, hence is one of that plane's
two isotropic axes.

We also need one prime detecting prescribed classes of opposite
signs. Choose nonzero $c_+\in V_+$ and $c_-\in V_-$.
Their restrictions to the exact kernel $H$ of
Lemma~\ref{ex:limits-kernel-detection} are nonzero additive equivariant
maps. Their joint $k$-linear evaluation span in $W_+\oplus W_-$
surjects to each summand. Simplicity and nonisomorphism force this
span to be the full direct sum. At an element $h\in H$, evaluation
at $(\tau h)^2$ on a sign-$\delta$ cocycle is
\[
 c_\delta((\tau h)^2)=(1+\delta J)c_\delta(h)
 \qquad(\delta=+1,-1).                                      \tag{5.5b}
\]
This follows from additivity on $H$ and the sign relation; a
coboundary is zero on $H$. Thus the two finite-coordinate maps on
the joint span surject respectively onto $k e_+$ and $k e_-$.
There is an actual $h$ for which both coordinates are nonzero.
For completeness, if their product vanished at every actual $h$,
additivity and the identity
$Q(x+y)-Q(x)-Q(y)$ for that product would make its polar form
zero on all pairs of actual evaluations. It would then vanish on
their linear span, contradicting surjectivity onto the two axes.
Finite Chebotarev and the exact matrix identity (5.4) realize this
choice by a new inert prime, avoiding every previously used prime.
All character, Tate, cochain and old-diagram constraints are included
in the same finite quotient as in that lemma. No residual
irreducibility or infinite-level Chebotarev assertion is used.
The same argument for one nonzero sign class gives nonzero
localization in its finite axis.

Suppose first that both $V_+$ and $V_-$ are nonzero. Choose the
class in one sign to be a nonzero component of $\bar v$ and choose
any nonzero class in the other. Detection as above puts both finite
axes in $U$, so $U=F$ and $U\cap S=0$. Equation (5.5a) gives
$\dim V'=s-2$. The localization of $\bar v$ is nonzero: its two
sign components lie on different finite axes, so cannot cancel.

Suppose instead that only one sign occurs, say $V=V_+$, with
$s>1$. Detect $\bar v$ in its finite axis. Then
$U_+=F(e_+)$. Since $V_-=0$, (5.5a) gives
$U_-\cap F(e_-)=0$. The isotropic-axis calculation forces
$U_-=S(e_+)$. Consequently $K_0$ has sign dimensions $(s-1,0)$,
while $V'$ has sign dimensions $(s-1,1)$. This proves that the first
switch introduces the other sign and preserves total dimension.
Since $s-1\ge1$, the next switch uses the preceding mixed-sign
case and lowers the total dimension by two. After the first switch
choose the primitive vector in the new integral generic kernel;
include one of its nonzero sign components in the second detection.
Thus the actual-class nonvanishing requirement holds at both steps.
The case $V=V_-$ is identical with signs interchanged.

At each step the new class is the actual derivative
$P_{I\cup\{\ell\}}$ of Proposition~\ref{ex:arith-finite-heegner},
with the same fixed multiplier $L$. Equation (4.2) gives finite
coordinate zero and singular coordinate $Jf(P_I)$, which is nonzero
generically. The construction retains the old Kummer conditions and
all earlier singular conditions, by the final assertion following
(H19). Lemma~\ref{ex:limits-derivative-switch} therefore makes this
class span the new generic kernel and preserves the valuation of
its class-functional determinant coordinate. In particular we do
not replace the actual class by the chosen primitive vector.
That vector is used only to select the prime. The paired functional
is the one transported from the same actual class. The determinant
comparison uses the common strict determinant factor; fractional
complements to $f(P_I)$ and $Jf(P_I)$ have unit exterior volume.
Consequently a divisible localization creates no extra valuation,
and all strict torsion factors remain present on both sides.

Each new generic kernel has rank one and has a primitive integral
vector, so its fiber dimension is at least one. A proposed drop
from dimension two to zero is therefore a contradiction; that
configuration cannot occur. Whenever the dimension exceeds one,
one or two switches give a strict drop by two. Thus finitely many
switches leave dimension one. After cancellation the terminal
complex has one free generator in each of degrees one and two;
its differential is zero because its generic kernel has rank one.
The actual class is integral. To apply
Proposition~\ref{ex:arith-paired-functional} after these switches,
we check its orthogonal comparison on the added local conditions.
In the integral derivative model of
Lemma~\ref{ex:limits-derivative-switch}, $D^0$ pairs with $D^2$ and
$S$ pairs perfectly with $F$, while the pairing on $S$ is zero
with its retained isotropy homotopy. Hence
$(D^0\oplus S[-1])^\perp=D^0\oplus S[-1]$ as the specified local
condition, with the retained integral comparison maps. Conjugate
transport preserves this condition as well. At a derivative prime
the comparison from the orthogonal condition to the conjugate
condition is therefore the identity under those integral maps.
It requires no inverse of a localization vector or of a local
Frobenius operator. Using this identity in the Selmer-cone comparison,
together with the same old-place comparisons as before, constructs
the integral orthogonal lift of the new transported class. Composing
with the retained integral global duality gives a closed integral
degree-two functional. Generically this is exactly the functional
of Proposition~\ref{ex:arith-paired-functional}: both use the inverse
of the same orthogonal comparison, which is unique in the derived
category. Thus its primitive integrality proof extends to every
finite list of these derivative conditions. At the
active ramified primes this uses the unit $-2$; at the moving prime
it uses $t$ except at $(t)$, where the central Frobenius determinants
are nonzero and hence units in $R$. Unused good primes retain their
unramified condition. No inverse at an unused relaxed prime is used.
The terminal coordinate is the product of the two integral
coefficients, up to a determinant-basis unit, so belongs to $R$.
Valuation preservation transports this conclusion through every
switch. This proves regularity at every $\mathfrak p\ne(2)$,
including both quadratic specializations. Finally the normal-domain
intersection formula for $\Lambda[1/2]$ gives
$u_{\rm prim}\in\Lambda[1/2]$.

The same proof applies at the unknown full-corank-one base. If its
actual generic class is zero, its coordinate is zero. Otherwise its
nonzero generic cycle and the central dimension bound give the
rank-one start used above. This requires no uniform outer valuation
bound at that base. Combining (5.5), for the known simple-zero
instances, with (4.6) and $\sum_{p\mid k}t_p\le2A$ proves (3.5).
No fixed-field class number is absorbed into the varying-field bound.


\end{proof}

\subsection{Actual augmentation and the fixed-field count}
\begin{lemma}[Augmentation and the fixed-field rank budget]\label{ex:limits-fixed-field-budget}
Fix the final field $K$ and the family of
Section~\ref{sec:rc}, in the branch where the odd empty
configuration is not a unit. There are no derivative labels.
At augmentation all graph primes are unramified.
The actual old-support subcomplex gives (5.6). With
$d_K=\dim_{\F_2}\Sel_2(E/K)$,
$t_K=\operatorname{length}E(K)[2^\infty]$ and
$h_K=\operatorname{length}H^2(C_K)_{\rm tors}$, it has the rank,
Smith-disk and torsion bounds (5.7). These are fixed-$K$ bounds,
independent of the binary dimension and graph size.

At every sign character, the inverse new local blocks and the
resulting old-rank Schur complex have no Laurent terms below
$-4M$ modulo $2^M$, uniformly in the number of new blocks.
This is a bound on sign evaluations, not on coefficients in the
group basis. The finite number $h_K$ does not assume finiteness
of the whole $\Sha$ at the base.
\end{lemma}
\begin{proof}\leavevmode
We now use the final $K$ in the remaining branch of Section~\ref{sec:rc},
where the intrinsic odd empty configuration is not a unit.
There are no derivative labels in this
last binary family. At augmentation every graph prime is unramified,
and the actual support triangle gives
\[
 C_{\rm old}\longrightarrow(C_b)_0
       \longrightarrow\bigoplus_{q\in Q}Q_q\longrightarrow C_{\rm old}[1],
                                                                    \tag{5.6}
\]
with $C_{\rm old}$ an actual subcomplex of an upper triangular model.
The graph-wide prescription gives scalar $1+t$ at each chosen
orientation, even though the quadratic character is trivial at
augmentation. The Euler-factor calculation in
Lemma~\ref{ex:arith-common-coordinate} therefore gives
\[
 \overline{l_{q,0}}=\overline e_{q,0}^{\,2}
                  =\frac{t^4}{(1+t)^2}.
\]
The two-place local determinant differs only by the unit $q^4$.
Its residual $t$-order is consequently four, uniformly in $q$ and
the number of graph primes.

Put $T=T_2E$ and let $C_K$ be the ordinary compact Kummer
Selmer complex over the fixed imaginary quadratic field $K$,
with the derived Kummer condition, including its degree-zero
torsion, at every finite place. Lengths in this paragraph are
lengths over $\mathbb Z_2$. Write
\[
 d_K=\dim_{\mathbb F_2}\operatorname{Sel}_2(E/K),\qquad
 t_K=\operatorname{length}E(K)[2^\infty],\qquad
 h_K=\operatorname{length}H^2(C_K)_{\rm tors}.
\]
The finite Selmer group defines $d_K$, and perfection of $C_K$
makes its cohomology finitely generated, so $h_K$ is finite.
This does not assert that the full two-primary Tate--Shafarevich
group is finite. The already fixed companion has
$L(E^{(k)},1)\ne0$, so the sum-and-difference isogeny and the
forward corank-zero theorem give full Selmer corank one for $E/K$.
This includes any divisible part of the Tate--Shafarevich group.

Here are the cohomology calculations underlying the budget.
At finite coefficients, the local degree-zero maps are
isomorphisms onto local invariants and the local degree-one maps
are the Kummer injections. The Selmer cone therefore has
$H^0=E(K)[2]$ and $H^1=\operatorname{Sel}_2(E/K)$.
The Weil pairing and the exact orthogonality of the Kummer
conditions give the finite self-duality with shift three
\cite[5.1.1--5.1.6 and 6.3.4]{nekovar}. Since $K$ has no
real places and $E[2]$ is rational of rank two, this gives
\[
 \dim_{\mathbb F_2}H^i(C_K\otimes^{\mathbf L}\mathbb F_2)
       =(2,d_K,d_K,2)\quad(i=0,1,2,3),                         \tag{5.7a}
\]
with zero cohomology in other degrees. In particular a minimal
finite free $\mathbb Z_2$ model of $C_K$ has total rank $2d_K+4$.

For the integral torsion calculation, $H^0(C_K)=T^{G_K}=0$:
a compatible system of rational points of unbounded two-power
order cannot lie in the finite group $E(K)_{\rm tors}$.
The exact sequence for $0\to T\to T\otimes\mathbb Q_2
\to E[2^\infty]\to0$, and the local Kummer conditions, identify
\[
 H^1(C_K)_{\rm tors}=E(K)[2^\infty].
\]
Indeed this finite group injects as torsion global cohomology,
all of its elements satisfy the Kummer conditions, and a torsion
class maps to zero with rational coefficients, so is in this
image. Integral Kummer self-duality
$C_K\simeq R\operatorname{Hom}_{\mathbb Z_2}
(C_K,\mathbb Z_2)[-3]$ now gives
\[
 H^3(C_K)\simeq
 \operatorname{Ext}^1_{\mathbb Z_2}(H^1(C_K),\mathbb Z_2),
 \qquad \operatorname{length}H^3(C_K)=t_K.                     \tag{5.7b}
\]
It also identifies the free ranks in degrees one and two.
They both equal one: universal coefficients with
$\mathbb Q_2/\mathbb Z_2$ give
\[
0\longrightarrow H^1(C_K)\otimes\mathbb Q_2/\mathbb Z_2
 \longrightarrow\operatorname{Sel}_{2^\infty}(E/K)
 \longrightarrow H^2(C_K)_{\rm tors}\longrightarrow0,
\]
and the middle group has full corank one. Thus the total integral
torsion length of $C_K$ is $h_K+2t_K$.

At $t=0$ the moving character is trivial. At either place over
the good split moving prime $r_i$, the unramified condition is
the ordinary Kummer condition. Replacing it by the full condition
adds the singular complex
\[
 Q_{r_i,w}=[\,\mathbb Z_2^2
       \xrightarrow{\Phi_{r_i}-r_i}\mathbb Z_2^2\,]
       \quad\text{in degrees one and two}.
\]
Here $\Phi_{r_i}$ is arithmetic Frobenius on $T$, with trace
$a_{r_i}$ and determinant $r_i$. In particular
\[
 \det(\Phi_{r_i}-r_i)
   =r_i(r_i+1-a_{r_i})
   =r_i\det(1-\Phi_{r_i}),
\]
so this determinant is nonzero and has valuation at most $c_E$.
The same bound holds at the conjugate place. The singular
complexes have zero degree-one cohomology and finite degree-two
cokernels of total length at most $2c_E$.

This comparison is retained before the inner limit. The field
$K$ and $C_K$ are fixed, while the two moving Frobenius matrices
and the comparison maps may vary with $i$ and with the graph.
Pass to their limit at every coefficient precision using the
same retained finite diagram. A matrix whose determinant has
valuation at most $c_E$ cannot acquire a zero determinant in
this limit: the property is determined modulo $2^{c_E+1}$.
Consequently, if
$D_c=C_{\rm old}\otimes^{\mathbf L}_\Lambda\mathbb Z_2$,
the resulting triangle is
\[
 C_K\longrightarrow D_c\longrightarrow Q_r
      \longrightarrow C_K[1],                               \tag{5.7c}
\]
where $Q_r$ is the sum of two rank-two singular disks and
$H^2(Q_r)=J_r$ has length at most $2c_E$. The notation $r$
in this limiting complex does not designate a fixed rational
prime supporting characters of unbounded order.

Reducing (5.7c) modulo two, the two singular disks contribute at
most eight to total cohomology dimension. Equation (5.7a) thus
gives a finite free model of $C_{\rm old}$ over $\Lambda$ of
total rank at most $2d_K+12$, by cancelling unit entries at its
maximal ideal $(2,t)$. The integral long exact sequence of
(5.7c) gives
\[
 H^1(D_c)\simeq H^1(C_K),\qquad
 0\to H^2(C_K)\to H^2(D_c)\to J_r
       \to H^3(C_K)\to H^3(D_c)\to0.                         \tag{5.7d}
\]
In particular $H^3(D_c)$ is a quotient of $H^3(C_K)$.
The image of the torsion of $H^2(D_c)$ in $J_r$ has length at
most $\operatorname{length}J_r$, and its kernel is
$H^2(C_K)_{\rm tors}$. Hence, defining the symbol used below by
\[
             h_1:=\sum_i\operatorname{length}H^i(D_c)_{\rm tors},
\]
we have $h_1\le h_K+2c_E+2t_K$.

Finally, a finite free complex over $\mathbb Z_2$ decomposes
into free cohomology generators with zero differential and
Smith disks $[\mathbb Z_2\xrightarrow{2^e u}\mathbb Z_2]$,
$u\in\mathbb Z_2^\times$, $e\ge0$. Its only free cohomology
generators here are one in degree one and one in degree two.
If $m$ counts all its Smith disks, including those with $e=0$,
then $\operatorname{rank}_{\rm total}C_{\rm old}=2+2m$.
Every disk contributes $e$ to $h_1$, so the sum of the valuations
of all its nonzero Smith entries equals $h_1$. We have proved
\[
 \operatorname{rank}_{\rm total}C_{\rm old}\le2d_K+12,\qquad
 m\le d_K+5,\qquad h_1\le h_K+2c_E+2t_K.                     \tag{5.7}
\]
The model and its matrices may vary with the moving characters,
but these bounds depend only on the final fixed $K$ and fixed
old data. They do not depend on the binary dimension or graph
size. The varying-field estimate (5.1) is not used with $j=b$,
and $h_K$ is not treated as a varying-field constant.

The lift is made over $R=\Lambda[G]$ before passing to
$\mathscr O[G]$. Let $C$ be the retained finite free family model.
Its augmentation $C_0$ is the actual scalar cochain model, by the
retained coefficient-change maps. After augmentation the representation
is unramified at $Q$, so the support-change triangle (5.6) gives a
finite free model $D$ of $C_0$ whose underlying graded module is
$C_{\rm old}\oplus\bigoplus_q Q_q$ and whose differential is upper
triangular. This follows by representing the connecting map between
these bounded free complexes by a chain map and forming its cone.

Cancel unit disks in $C_0$. The same pivots and elementary operations
lift to $C$ over $R$, since the augmentation ideal is contained in
the Jacobson radical of $R$. Let $C'$ be the resulting family model,
so that $C'_0$ is minimal. Decompose $D$ over $\Lambda$ into a minimal
complex $D_{\min}$ and unit disks $E$. The actual comparison gives a
homotopy equivalence $C'_0\simeq D_{\min}$. A homotopy equivalence
between minimal bounded free complexes over a local ring is degreewise
invertible, by reduction modulo its maximal ideal and Nakayama.
Add to $C'$ the same unit disks $E_R$ and lift the resulting graded
isomorphism $C'_0\oplus E\simeq D$ to $R$. Its matrices are invertible
because their augmented determinants are units. Conjugating the
family differential by these lifted matrices makes its augmentation
exactly the differential of $D$. All matrices and their inverses at
this step lie in $\Lambda[G]$; thus their sign evaluations are power
series. The added disks are precisely those needed for $D$, so this
procedure does not enlarge its stated old-support rank.
The block containing the local quotients has
\[
 B_0=A=\operatorname{diag}(A_q),\qquad T_0=0
\]
in a differential $\left(\begin{smallmatrix}P&Q\\T&B\end{smallmatrix}\right)$.
Over $\mathscr O[G]$, cancel $B$. The surviving Schur differential
$P-QB^{-1}T$ augments to $P_0$.

Write $\det A_q=t^{d_q}v_q+2w_q$, $d_q\le4$, with $v_q$ a power-series
unit, and put
$C_q=\operatorname{adj}(A_q)/(t^{d_q}v_q)$.
For a character $x$, $B_xC=1+2F$ with both $C$ and $F$ having no
Laurent terms below degree $-4$. Therefore
\[
             B_x^{-1}=C\sum_{j=0}^{M-1}(-2F)^j\pmod {2^M}        \tag{5.8}
\]
has pole at most $4M$, independent of the number of blocks.
For completeness the adjacent-differential assertion is an exact
block identity. In the two degrees containing the new local block,
write the differential as
$d=\left(\begin{smallmatrix}P&Q\\T&B\end{smallmatrix}\right)$,
with $B$ invertible. Then
\[
 \begin{pmatrix}1&-QB^{-1}\\0&1\end{pmatrix}
 d\begin{pmatrix}1&0\\-B^{-1}T&1\end{pmatrix}
  =\begin{pmatrix}P-QB^{-1}T&0\\0&B\end{pmatrix}.
\]
If the preceding differential has components $(X,Y)^{\mathsf T}$,
then $d^2=0$ gives $TX+BY=0$. In the new source coordinates its
components are $(X,B^{-1}TX+Y)^{\mathsf T}=(X,0)^{\mathsf T}$.
If the following differential has components $(U,V)$, then
$UQ+VB=0$, and in the new target coordinates they are
$(U,UQB^{-1}+V)=(U,0)$. Thus the surviving entries of both adjacent
differentials are literally unchanged; no second inverse is
introduced there.

Here $A_q$ is the direct sum of the two rank-two singular
differentials at the conjugate graph places. Modulo two their
scalar entries are $t$ and $-t/(1+t)$, since their Frobenius
scalars are $1+t$ and $(1+t)^{-1}$ and $E[2]$ is trivial.
Consequently
\[
                  \det\overline A_q=\frac{t^4}{(1+t)^2}.
\]
This independently gives the residual order four in (5.6).
The factor $q^4$ in the Euler-factor convention is a unit and
does not affect it. Moreover every sign evaluation of a matrix
over $\Lambda[G]$ has the same reduction modulo two as its
augmentation, because each sign is congruent to one modulo two.
Thus $B_x=A+2V_x$ with $V_x$ a power-series matrix. For the
block-diagonal matrix $C$ in (5.8), the identities
\[
 A_q C_q=1+\frac{2w_q}{t^{d_q}v_q},\qquad
 B_xC=1+2F_x
\]
show that every entry of $C$ and $F_x$ has Laurent order at
least $-4$. Formula (5.8) therefore has Laurent order at least
$-4M$, because its $j$th summand has order at least $-4(j+1)$.
The estimate counts the largest pole of an entry, not the number
of summands or blocks. The remaining Schur differential has the
same bound since $P_x,Q_x,T_x$ are power-series matrices, and
the adjacent surviving entries have no poles at all. These
arguments take place after localization where $t$ is invertible;
they do not invert $t$ in a finite Artin coefficient ring.
\end{proof}

\subsection{Determinant clearing and the missing-address contradiction}
\begin{proposition}[A common denominator for the corrected coordinate]\label{ex:limits-fixed-field-clearing}
Under Lemma~\ref{ex:limits-fixed-field-budget}, use the paired
functional exponent $a$ of
Proposition~\ref{ex:arith-paired-functional} and the universal
correction $\eta$ of Lemma~\ref{ex:arith-common-coordinate}.
There are common $\Delta,Q\in\mathscr O[G]$ such that, wherever
$\Delta_x\ne0$,
$\widetilde u_x=2^{-a}\eta_xQ_x/\Delta_x^3$.
Putting $f=2^a\Delta^4$ and $d=\eta\Delta Q$ gives
$f_x\widetilde u_x=d_x$ at every character, including
$\Delta_x=0$, and $f_0\in\Lambda$. Moreover
\[
 v_2f_0(0)\leq a+4(h_K+2c_E+2t_K)=c_{\det},\qquad
 L_T(M)=16(d_K+5)M
\]
bounds the negative Laurent support of every $f_x$ modulo $2^M$,
uniformly in the sign character $x$. The elements $f,d$ themselves
belong to $\mathscr O[G]$ in the original group basis.
These bounds concern the same corrected coordinate throughout;
no pole bound for $d_x$ or for the group coefficients of $f$ is required.
\end{proposition}
\begin{proof}\leavevmode
Put $S=\mathscr O[G]$ and first perform the simultaneous new-local
cancellation of Lemma~\ref{ex:limits-fixed-field-budget}.  Its
determinant comparison factor $\kappa$ is an $S$-unit: the canceled
block is invertible over $S$, so its determinant and inverse are
integral over $S$.  Retain this factor, with the determinant-line
sign prescribed by the original basis.  Transport the actual cycle
$z$ and the closed functional $\mu=2^a\lambda_z$ to the remaining
complex $D$.  Both transported cochains are integral over $S$.
The augmentation of $D$ is the old complex over $\Lambda$, and
each sign-evaluated differential entry has pole at most $4M$
modulo $2^M$.  No bound on the entries of the transported $z$ or
$\mu$ is needed below.

Write $D_c=D_0\otimes_\Lambda\mathbf Z_2$.  Its integral Smith
decomposition gives graded decompositions
\[
 D^i=B_i\oplus H_i\oplus E_i,
\]
lifted by constant $\mathbf Z_2$ basis changes, where centrally the
only nonzero differentials are the Smith disks $E_i\to B_{i+1}$.
Here $H_1,H_2$ have rank one and all other $H_i$ vanish.  Let
$m_i=\operatorname{rank}E_i$, $m=\sum_i m_i$, and let $A_i$ be the
actual differential submatrix from $E_i$ to $B_{i+1}$.  Set
\[
 \delta_i=\det A_i,\qquad \Delta=\prod_i\delta_i.
\]
An empty determinant is one.  The entries in this definition are
the original entries of $D$, before any of these old-disk
cancellations.  Thus $\Delta$ has degree $m$ in those entries,
$\Delta_0\in\Lambda$, and
\[
 v_2\Delta_0(0)=\sum_{\text{central Smith disks}}v_2(\text{disk})
       \le h_1\le h_K+2c_E+2t_K.
\]

Here is the full denominator calculation.  Work temporarily in
$S[1/\Delta]$.  Group all sources and targets for a given degree,
and cancel the whole block $A_i$ once.  For that differential the
two elementary matrices are
\[
 \begin{pmatrix}1&0\\-CA_i^{-1}&1\end{pmatrix}
 \begin{pmatrix}A_i&B\\C&D'\end{pmatrix}
 \begin{pmatrix}1&-A_i^{-1}B\\0&1\end{pmatrix}
 =\begin{pmatrix}A_i&0\\0&D'-CA_i^{-1}B\end{pmatrix}.
\]
Both elementary matrices have determinant one.  In the adjacent
differentials, the equation $d^2=0$ makes the components into and
out of this canceled pair zero after these changes; the other
entries of those adjacent differentials are unchanged.  Hence a
cancellation does not modify any other surviving $A_j$.  Every
inverse used is $\operatorname{adj}(A_i)/\delta_i$, and each grouped
block is used only once.

More explicitly, in degree $k$ the projection associated to the
incoming cancellation is
\[
 (b,v)\longmapsto v-C A_{k-1}^{-1}b,
\]
and the inclusion associated to the outgoing cancellation is
\[
 v\longmapsto(-A_k^{-1}Bv,v).
\]
The other map in each of these two degrees is the corresponding
coordinate inclusion or projection.  Compose these maps in the
chosen cancellation order.  At each step each entry is linear in
the one inverse $A_i^{-1}$ at that step, and no $A_i$ is encountered
again.  Products of denominators therefore divide $\prod_i\delta_i$,
not a power depending on the number of individual Smith disks.
In particular, if $p:D\to H$ and $\iota:H\to D$ are the final
projection and inclusion, the two scalars
\[
 \alpha=(pz)^1\in H_1,
 \qquad \beta=(\mu\iota)^2\in H_2^*
\]
satisfy $Z:=\Delta\alpha\in S$ and $W:=\Delta\beta\in S$.
The final complex $H$ has just the map $H_1\to H_2$.

For the inverse determinant line, a disk in degrees $i,i+1$
contributes $\delta_i^{(-1)^i}$ to the coordinate of a cohomology
tensor.  This can be checked on the based disk
$e\mapsto\delta_i b$: replace the boundary basis $b$ by
$\delta_i b$ in $\bigotimes_k\det(D^k)^{(-1)^{k+1}}$.
Thus the remaining comparison factor is
\[
 T=\epsilon\prod_i\delta_i^{(-1)^i},\qquad
 P:=\Delta T=\epsilon\prod_{i\text{ even}}\delta_i^2\in S,
\]
where $\epsilon$ is the integral unit from the fixed basis ordering
and the constant Smith basis changes.  Including the previously
retained new-local factor, put
\[
 Q=\kappa ZWP\in S.
\]
This definition takes place over $S$, before sign evaluation.  It
uses no interpolation and no division by $|G|$.

If $\Delta_x\ne0$ and the constituent has the rank-one cohomology
pattern, its coordinate in the original determinant basis is exactly
\[
 \widetilde u_x=\eta_x\kappa_x\alpha_x(2^{-a}\beta_x)T_x
              =2^{-a}\eta_xQ_x/\Delta_x^3.
\]
This is the same $z$, functional, basis and multiplier $\eta$ as in
the statement.  If the constituent is generically acyclic, the
remaining map between the two lines is nonzero.  Its cycle
$\alpha_x$ is then zero over the fraction field, so $Q_x=0$ and the
formula agrees with the defined coordinate zero.  These are the
only generic patterns in the present family, by central
specialization and the central corank-one bound.

Now define $f=2^a\Delta^4$ and $d=\eta\Delta Q$ in $S$.
The preceding formula proves $f_x\widetilde u_x=d_x$ whenever
$\Delta_x\ne0$.  When $\Delta_x=0$, both $f_x$ and $d_x$ are zero;
the already defined coordinate $\widetilde u_x$ is finite in its
fraction field, and the same equality follows directly.  No rational
expression with a zero denominator is specialized.  In particular
$f_0\in\Lambda$ and
\[
 c=v_2f_0(0)\le c_{\det}:=a+4(h_K+2c_E+2t_K),\qquad
 L_T(M)=16(d_K+5)M.                                      \tag{5.9}
\]
Indeed $f$ is a polynomial of degree $4m$ in the differential
entries of $D$, with constant integral coefficients.  Each such
entry, at a sign character, has pole at most $4M$ modulo $2^M$.
A product of $4m$ entries has pole at most $16mM$, and sums cannot
increase this bound.  Finally $m\le d_K+5$.

Only the last pole estimate is characterwise.  The actual elements
$f,d$ belong to the original group ring $S$.  For $i<-L_T(M)$ all
sign values of $[t^i]f$ are divisible by $2^M$; the group-ring
coefficient itself need not be divisible by $2^M$.  No Laurent
support bound for $d$ or for the group coefficients of $f$ has been
used.  The constants $d_K,t_K,h_K$ are fixed only after the final
$K$ is fixed; $a$ and $c_E$ precede that choice.  None depends on
the graph or binary dimension.  A varying-$K$ estimate cannot be
deduced by treating $h_K$ as an absolute constant.
\end{proof}

\begin{proposition}[Finite coefficient transfer and the final contradiction]\label{ex:limits-coefficient-transfer}
Use the fixed field, common family, regular coordinates and
integral identity of Proposition~\ref{ex:limits-fixed-field-clearing}.
Let $B$ be the fixed upper bound on
$v_2\widetilde u_x(0)$ at every nonzero address supplied by
(3.9), (3.4), (4.6) and Lemma~\ref{ex:arith-common-coordinate}.
Choose $M,b$ as in (5.10), before the finite graph and its actual
theta stage are realized. For some nonzero address $x$,
\[
 v_2\bigl(\widetilde u_x(0)-\widetilde u_0(0)\bigr)
 \geq M-2c>B,\qquad c=v_2f_0(0)\leq c_{\det}.
\]
The convolution includes every negative Laurent coefficient
of $f_x$. If $L'(E/K,1)=0$, Lemma~\ref{ex:arith-base-zero}
makes $\widetilde u_0(0)=0$, contradicting the upper bound.
Since the fixed complementary value $L(E^{(k)},1)$ is nonzero,
this proves $\operatorname{ord}_{s=1}L(E,s)=1$.
\end{proposition}
\begin{proof}\leavevmode
For a coefficient of any element of $\mathscr O[G]$, its sign-value
function on $\F_2^b$ expands as
\[
 A(x)=a_\varnothing+
        \sum_{I\ne\varnothing}2^{|I|}a_I\prod_{i\in I}x_i .
\]
Its low bit is constant and its bit of index $j\ge1$ has Boolean
degree at most $2^{j-1}$. One can see this by extracting the
$Y^{2^j}$ coefficient of $(1+Y)^{A(x)}$ modulo two:
the first possible contribution of $x_I$ has degree at least
$2^{|I|}$ in $Y$, and $|I|\le2^{|I|-1}$.
Thus $q$ coefficient equalities modulo $2^M$ have sum of Boolean
degrees at most $q(2^{M-1}-1)$.

If that sum is less than $b$, multiply the complementary indicators
of the equations. Its degree is less than $b$, so its sum over the
Boolean cube is zero: every monomial misses a variable.
The support of this indicator is the simultaneous solution set.
It has even cardinality and contains zero, hence another address.

Apply this to the coefficients of $f$ from $-L_T(M)$ through zero
and the single constant coefficient of $d$. It suffices to choose
\[
 M>\max(c_{\det},B+2c_{\det}),\qquad
 b>(L_T(M)+2)(2^M-1),                                           \tag{5.10}
\]
where $B$ bounds the nonzero-address central valuations.
These choices precede the graph and its theta precision.
At the resulting nonzero address, negative coefficients of $f_x$
vanish modulo $2^M$, and its constant coefficient agrees with
$f_0(0)$ and has valuation $c$.
Positive coefficients may have smaller valuation. Nevertheless
$v_{\mathscr O}(f_x)\le c$, so integrality of
$d_x=f_x\widetilde u_x$ gives
$v_{\mathscr O}(\widetilde u_x)\ge-c$.
This estimate comes from the displayed product at the selected
nonzero address; it does not apply (5.5) to the unknown base.
Both $\widetilde u_x$ and $\widetilde u_0$ lie in $\Lambda[1/2]$
by the height-one argument above, so neither has negative powers
of $t$. In particular, since $f_0\in\Lambda$,
$[t^0]d_0=f_0(0)\widetilde u_0(0)$.
Put $F_0=[t^0]f_0$ and $F_x=[t^0]f_x$. Since $d_0$ is integral
and $[t^0]d_0=F_0\widetilde u_0(0)$, we also have
$v_2\widetilde u_0(0)\ge-c$, without using (5.5) at the base.
The convergent convolution has the form
\[
 [t^0]d_x=F_x\widetilde u_x(0)+E_x,\qquad
 E_x=\sum_{i<0}[t^i]f_x[t^{-i}]\widetilde u_x,
 \qquad v_2(E_x)\ge M-c.
\]
The untested coefficients below the Laurent cutoff have the same
divisibility by the pole bound, so this sum includes every negative
coefficient. Subtracting the two constant terms gives
\[
 F_x(\widetilde u_x(0)-\widetilde u_0(0))
 =([t^0]d_x-[t^0]d_0)-(F_x-F_0)\widetilde u_0(0)-E_x.
\]
Every term on the right has valuation at least $M-c$ and
$v_2(F_x)=c$. Consequently
\[
 v_2(\widetilde u_x(0)-\widetilde u_0(0))\ge M-2c>B.\tag{5.11}
\]

By (3.9), (3.4), (4.6) and the unit comparison (4.7), every nonzero address has
$v_2\widetilde u_x(0)\le B$ for one $B$ fixed before $b$; the finite
$\Sha$ terms have negative sign and may be discarded for this upper
bound. If the base derivative vanished, Section~\ref{sec:arithmetic}
would give $\widetilde u_0(0)=0$, contradicting (5.11).
Hence $L'(E/K,1)\ne0$. The complementary value $L(E^{(k)},1)$ is
nonzero, and
$L(E/K,s)=L(E,s)L(E^{(k)},s)$, so $E$ has analytic rank one.
This proves (RC), and hence Theorem~\ref{hyp:RC}, using the classical
applications and finite constructions specified above.
\end{proof}

% End 03b-limits.tex


% Begin 04-computation-value.tex
\section{Exact precision computation and polynomial compilation}\label{sec:compiler}
Encode a signed integer by
\[
 \operatorname{enc}(n)=\begin{cases}2n,&n\ge0,\\-2n-1,&n<0.\end{cases}
\]
The five natural inputs decode in the fixed order $(n,\kappa,a,l,j)$.
The decision predicate rejects unless
$n\notin\Cset$, $\kappa\ne0$, $a,l\ge1$, and $j=na$.
For nonzero multiples of $P_0$, $z=-x/y$ is well-defined. Set
\[
 z([l]P_0)=u_l+v_ls,\qquad\epsilon_l=|u_l|+|v_l|,
 \qquad r(a,j)=\tfrac12\Tr_{K_0/\Q}\frac{z([j]P_0)}{z([a]P_0)}.
\]
With the fixed constant $A_{\rm ht}^{(0)}=H_6(S_{\rm ht})$, its final
tests are
\[
 |r(a,j)-n|<\epsilon_l,\qquad
 4\epsilon_l<4^{-A_{\rm ht}^{(0)}|\kappa F(n)|^{24}}.
\]
These are strict tests on exactly computable rationals and integers.

The reference program S21 uses exact quadratic-field arithmetic and capped
integer operations. The ideal register intermediate representation S22
computes the same predicate on every five-natural tuple. Its independent
all-input argument covers normalized signed arithmetic, origin and
zero-ordinate cases, equal-abscissa addition, and every early rejection.
Nine exact polynomial identities verify the nonexceptional projective
formulas and preservation of $Y^2=X^3+AXZ^4$. Scalar loops decrease by
halving; the precision loop divides by four; factorial iteration is bounded.
This is a proof about the ideal intermediate representation, rather than
an assertion that a budgeted simulator always finishes.

For example, if $\epsilon=U/V>0$ with natural $U,V$ and integer $B\ge0$,
\[
 4\epsilon<4^{-B}\quad\Longleftrightarrow\quad
 4^B\le\left\lfloor\frac{V-1}{4U}\right\rfloor.
\]
This permits a finite cap. For positive $L$, $\sigma(x)=\min(x,L+1)$
commutes with saturated addition and multiplication. Powers and factorials
can be computed with this saturation. Capping an exponent is safe for
base at least two because an exponent above $L$ already makes the value
exceed $L$. Base zero, base one and exponent zero are separate cases.
This proves equality with the fixed-height expression without requiring
its full evaluation when the available precision is too small.

\subsection{The natural polynomial and the power interface}
The saved precision polynomial has $1\,305\,811$ natural variables,
the first five being inputs. Its SHA256 is
\begin{center}\small\ttfamily
f3e163f5c60b972c98b48647274f956dd7d78428e27ddc0543574a5f18120b8c.
\end{center}
Source S23 specializes a fixed universal natural halting polynomial at
the code of a wrapper. The wrapper inverts four Cantor pairs and then
runs the unchanged precision machine. Its pairing equations are
\[
 2t_i=(c_i+t_{i+1})(c_i+t_{i+1}+1)+2t_{i+1}\quad(0\le i<4),
 \qquad t_4=c_4.
\]
Their triangular intervals prove unique natural unpairing. The new
polynomial sums the square of the old halting polynomial and the four
squared differences above.

The universal compiler packs a finite run of a fixed $m$-state
increment/decrement/branch program into base-$B$ integers. Acceptance
means reaching the unique final halt with all registers zero; the fixed
interpreter includes a finite cleanup stage. Choose
\[
 c>0,\quad K=2^c,\quad B=4K>m,\quad T=B^q,\quad
 F_q=1+B+\cdots+B^q,\quad D=(K-1)F_q,
\]
and require every input to be below $K$. For natural integers, say
$X$ is masked by $Y$ when every set binary bit of $X$ occurs in $Y$.
Require each register history $R_j$ to be masked by $D$ and satisfy
$R_j<T$. Each state history $S_i$ is masked by $F_q$, and
\[
 \sum_{i=0}^{m-1}S_i=F_q,\qquad S_{\rm halt}=T.
\]
Thus register digits are below $K$, terminal register digits are zero,
and state digits are zero or one. Since $B>m$, the state sum has no
carries and selects exactly one state at each time.

The nonzero mask $N_j$ is determined by
\[
 K N_j=(R_j+D)\mathbin{\mathrm{AND}}(K F_q).
\]
Indeed each digit $r+K-1<2K<B$ has its $K$-bit set exactly when
$r>0$. For a decrement/branch instruction at state $i$ testing
register $j$, impose
\[
 V_i=S_i\mathbin{\mathrm{AND}}N_j,\qquad S_i=V_i+W_i.
\]
These are the nonzero and zero branch histories. Form the packed
increment indicator $A_j$ from states incrementing $j$, the decrement
indicator $C_j$ from the appropriate $V_i$, and each incoming-state
indicator $I_i$ from the states or branches directed to $i$.
They have zero-one digits, since there is just one active instruction
and branch at each time. The transition equations are
\[
 R_j+BC_j=\operatorname{initial}_j+BR_j+BA_j,\qquad
 S_i=\mathbf1_{i=0}+BI_i.
\]
Every register coefficient on either side is at most $K<B$;
uniqueness of base-$B$ digits gives the initial conditions and each
exact transition. The possible extra final slot vanishes because
halt has no outgoing instruction and the terminal registers are zero.
For $q=0$ the same equations require the initial state already to be
halt and the inputs to vanish. Conversely, any finite accepting run
fits after choosing $c$ sufficiently large, and supplies all these
masks and equations.

\subsection{The literal alpha and power formulas}
The alpha system below is \texttt{Exp\_Matrices.alpha\_equations},
with equivalence \texttt{alpha\_equivalence}; the power system is
\texttt{Exponential\_Relation.exp\_repr} and its polynomial expansion
\texttt{exp\_dioph} in the pinned AFP source \cite{afp}.
Their variables range over natural numbers. The alpha recurrence itself
is integer-valued; the bounds below justify its natural-valued uses.
Define $U_n=\alpha_b(n)$ by
\[
 U_0=0,\qquad U_1=1,\qquad U_{n+2}=bU_{n+1}-U_n\quad(n\ge0).
\]
For $b\ge3$ this recurrence gives
$0\le U_n<U_{n+1}$ and $U_n\ge n$. For $n\ge1$, the matrix
\[
 M=\begin{pmatrix}b&-1\\1&0\end{pmatrix},\qquad
 M^n=\begin{pmatrix}U_{n+1}&-U_n\\U_n&-U_{n-1}\end{pmatrix}
\]
has determinant one. Thus adjacent terms are coprime and
$U_n^2-bU_nU_{n-1}+U_{n-1}^2=1$.
Parameter congruence follows directly from the recurrence:
$b\equiv b'\pmod d$ implies
$\alpha_b(n)\equiv\alpha_{b'}(n)\pmod d$.
In particular $\alpha_2(n)=n$. This is parameter congruence,
not merely congruence between indices.

Every natural solution of $x^2+y^2=1+bxy$, $b>2$, is an unordered
adjacent pair. To prove this, order it as $X>Y$. If $Y=0$, $X=1$.
Otherwise $Z=bY-X$ satisfies $XZ=Y^2-1$, hence $0\le Z<Y$,
and $(Y,Z)$ satisfies the same equation. Descent reaches $(1,0)$;
reversing the steps generates exactly the adjacent pairs.
No ordering assumption on the original pair is needed.

For $k>0$, reduction of $M^k$ modulo $U_k$ is the invertible scalar
$U_{k+1}I$. Dividing an index by $k$ and using strict growth proves
$U_k\mid U_m$ if and only if $k\mid m$. More precisely, for $m>0$,
\[
                    U_k^2\mid U_m\quad\Longleftrightarrow\quad
                         kU_k\mid m.                           \tag{C.1}
\]
For $k=0$ both sides are false. For $k>0$ write $m=kq$ and expand
$M^{kq}=(U_kM-U_{k-1}I)^q$ modulo $U_k^2$. Its lower-left entry is
$qU_k(-U_{k-1})^{q-1}$, so coprimality proves (C.1), including $U_k=1$.

For $m>0$ put $v=U_{m+1}-U_{m-1}=\operatorname{tr}(M^m)$.
Cayley--Hamilton gives $M^{2m}\equiv-I\pmod v$.
Extend the recurrence by $U_{-j}=-U_j$. Every $n\ge0$ has
$n=2lm+\varepsilon j$ with $l\ge0$, $0\le j\le m$ and
$\varepsilon\in\{1,-1\}$; then
$U_n\equiv(-1)^l\varepsilon U_j\pmod v$.
The value sign need not equal the index sign. If $b\ge4$,
$v=bU_m-2U_{m-1}>2U_m$.

The actual alpha formula has natural witnesses $r,s,t,u,v,w,x,y$
and the following fifteen conditions:
\[
\begin{gathered}
 b>3,\quad u^2+t^2=1+but,\quad s^2+r^2=1+bsr,\quad r<s,\quad
 u^2\mid s,\quad v+2r=bs,\\
 w\equiv b\pmod v,\quad w\equiv2\pmod u,\quad w>2,\quad
 x^2+y^2=1+wxy,\\
 2X<u,\quad2X<v,\quad x\equiv X\pmod v,\quad
 2N<u,\quad x\equiv N\pmod u.
\end{gathered}                                                   \tag{C.2}
\]
A strict inequality $a<b$ is $b=a+1+d$; divisibility is $b=ad$;
congruence is $A+Mu=B+Mv$, with fresh natural witnesses.
At zero modulus the last equation says equality, so there is no
exception hidden in this replacement.

For soundness, the unordered Pell result gives $u=U_k$, $k>0$,
and $x=\alpha_w(n)$. The ordered pair gives $s=U_m,r=U_{m-1}$,
$m>0$. By (C.1), $u\mid m$. Parameter congruence gives
$x\equiv U_n\pmod v$ and $x\equiv n\pmod u$.
Reflect $n$ to $j\in[0,m]$. Since $2X<v$ and $2U_j<v$,
the congruence $X\equiv\pm U_j\pmod v$ forces $X=U_j$,
with the negative case forcing both zero.
Then $j\le X$, $2j<u$ and $2N<u$ force $N=j$ from
$N\equiv\pm j\pmod u$. Thus $X=U_N$.

Conversely choose an odd $u=U_k>2U_N$. Such terms are unbounded,
since adjacent terms are coprime and the sequence grows.
Set $t=U_{k+1}$, $m=ku$, $s=U_m$, $r=U_{m-1}$ and $v=bs-2r$.
Equation (C.1) gives $u^2\mid s$, and the required strict bounds hold.
The moduli $u,v$ are coprime: a common divisor is odd, divides $s$
and $2r$, hence divides both adjacent terms. CRT gives
$w\equiv2\pmod u$, $w\equiv b\pmod v$ with $w>2$.
Take $x=\alpha_w(N)$, $y=\alpha_w(N+1)$.
Parameter congruence verifies every remaining condition.
This includes $N=X=0$.

The literal power formula separates $Q=0,R=0,P=1$ and
$Q=0,R>0,P=0$. For $Q\ge1$ it imposes
\[
 a_0=\alpha_{Q+4}(R+1),\quad b=a_0+Q^2+2,\quad
 m+Q^2+1=bQ,\quad P<m,
\]
\[
 a_1=\alpha_b(R),\quad a_2=\alpha_b(R+1),\quad
 P+ba_1\equiv Qa_1+a_2\pmod m.                                  \tag{C.3}
\]
Induction gives $a_0\ge(Q+3)^R$, so $m>Q^R$ and $m>0$.
Let $E_R=Q U_R-bU_R+U_{R+1}$. The recurrence and
$bQ\equiv Q^2+1\pmod m$ give $E_0=1$ and
$E_{R+1}\equiv QE_R\pmod m$. Thus (C.3) says $P\equiv Q^R\pmod m$.
Both values lie in $[0,m)$, forcing equality. Conversely equality
supplies every witness. This includes $R=0$, $Q=1$ and $0^0=1$.

These formulas make the binomial and bit operations polynomial.
For $B=2^{n+1}$, the equation
\[
 (1+B)^n=B^k(Bq+C)+r,\quad C<B,\quad r<B^k
\]
selects exactly $\binom nk$, since all its base-$B$ coefficients
are below $B$, including zero for $k>n$.
The identity
$(1+t)^n=\prod_{i:\,\text{bit}_i(n)=1}(1+t^{2^i})$ over $\F_2$
makes odd binomial coefficients equivalent to bit inclusion.
Bit disjointness of $a,b$ is equivalent to inclusion of $a$ in $a+b$:
at the least shared bit there is no incoming carry and that bit
is cleared. Finally $x=z+a,y=z+b$ with $z,a,b$ pairwise bit-disjoint
says exactly $z=x\mathbin{\rm AND}y$.

\subsection{Interpreter, lowering and five-input specialization}
The universal program numbering is $e+1=2^w(2D+1)$.
At $w=0$ it halts. Otherwise the four base-$2^w$ digits at positions
$4p,\ldots,4p+3$ give the opcode, register and two targets at state $p$.
Opcode one increments, opcode two decrements or takes its zero branch;
every other opcode halts. Missing digits are zero, so decoding is total.

The interpreter stores a finite tape as
$0$ or $\operatorname{cons}(h,t)=2^h(2t+1)$.
Repeated halving decodes a nonempty head and tail; the tail is smaller
than the code. To update a register, save the finite reversed prefix,
replace the entry and restore the prefix. Absent entries are zero.
The reserved jump register remains zero. Copies restore their source;
aliased addition snapshots its operand before overwriting it.
Division maintains $x=qy+r$ and subtracts only when $r\ge y$;
the interpreter's divisors are at least two.
These invariants prove that one interpreter loop executes exactly one
decoded instruction, with finite work. A simulated halt enters cleanup
and clears all registers. An infinite simulation never enters cleanup.
Thus its all-zero halt is equivalent to virtual halting, without
requiring the virtual final tape itself to be zero.

The precision IR uses normalized signed pairs for integer arithmetic
and pairs over the basis $1,\sqrt2$ for field arithmetic.
Norm zero implies the field element is zero: after clearing
denominators, a primitive solution of $a^2=2b^2$ would have both
entries even. Every division is guarded by that exact test.
The Jacobian point formulas satisfy the curve equation and the affine
chord or tangent formulas after multiplying denominators.
Equal abscissas split into doubling and opposite-point branches;
zero ordinate doubles to the identity. The fixed projective point
is the same affine $P_0$. These are polynomial identities, independent
of the generator theorem.

The scalar loop halves its remaining count. The precision loop divides
its remainder by four. The capped arithmetic computes exactly
$\min(x,L+1)$; zero bases, unit bases and zero exponents are separate.
For bases at least two an exponent exceeding $L$ already exceeds the
cap. Thus the ideal IR terminates and agrees with the stated rational
predicate on all five natural inputs, including every early rejection.

Each primitive lowering uses fresh auxiliary registers, snapshots
overlapping sources, and preserves the live-register map at the end
of each finite block. Saturating decrement gives natural subtraction,
and repeated addition gives multiplication. All variable divisions
are guarded. Induction on completed loop iterations proves the same
claim for loops. Acceptance enters cleanup; rejection enters a
permanent self-loop. Hence the primitive machine halts exactly on
acceptance, although the ideal decision program is total.

For $\operatorname{pair}(a,b)=(a+b)(a+b+1)/2+b$, unpairing maintains
$\text{remainder}=\text{input}-s(s+1)/2$ and subtracts $s+1$
while $s<\text{remainder}$. At exit
$0\le\text{remainder}\le s$, so $b=\text{remainder}$, $a=s-b$.
Four repetitions put the five codes into registers zero through four.
Fresh scratch lies above the old register range and is cleared;
old work registers remain zero. Relocating every target by the prefix
length therefore preserves the precision machine.

The polynomial specialization fixes the universal program input to
this wrapper code, maps the old packed input to variable five, and
maps every old variable $i\ge2$ and private-block start by $i\mapsto i+7$.
The four pairing variables occupy five through eight, disjoint from
the five inputs. The appended four squared equations force their
unique Cantor values. Thus, for the exact saved polynomial $C$,
\[
 \forall\mathbf c\in\N^5,\qquad
 \exists\mathbf z\in\N^{1305806}\ C(\mathbf c,\mathbf z)=0
 \quad\Longleftrightarrow\quad
        \operatorname{accepts\_naturals}(\mathbf c)=\mathrm{True}. \tag{C.4}
\]
The natural compiler proof has no height, restricted-converse or
rational-model premise. Its use of the height expression concerns
only computation of that fixed integer expression.

\subsection{Translation to point codes}
The typing order is independent of the natural predicate. First,
\texttt{member} forces a canonical code in $D$, and raw addition on two
canonical elliptic inputs forces their canonical sum, including the
origin and two-torsion cases. The fixed scalar graph gives $[m]Q$ on
$D$. In \texttt{multiply}, the three root membership calls are forced
conjuncts inside \texttt{height\_certificate}. Thus the four scalar shifts
are in $D$ before they are passed, unchanged, to the 27 slope calls and
the six test-group calls within each slope. This supplies the hypotheses
of every raw addition in those groups. Only after globally typed
multiplication is established do the four square products and three
successive additions justify \texttt{natural}. Their outputs are canonical
by the raw-addition postcondition, so no separate output membership
assumption is needed. The signed-encoding and exceptional branches start
their addition chains at this natural code, the literal origin, or the
fixed signed generator. An inactive product branch is not assumed true.
The fixed generator equations are conjuncts outside the internal products
in all three templates that use them. These observations give the base
cases for the source-DAG induction below.
Apply (C.4) using the exact signed index model. Induction on the actual
source DAG gives each node the code of its signed integer value.
Natural guards constrain source variables, not intermediate sums or
products. Raw addition receives canonical codes; globally typed
multiplication justifies the natural guards. Fixed generator equations
and scalar templates supply constants.

The power-value recipe has three signed code inputs, seventy-nine
private natural witnesses, and an output code for the value of the
power-template polynomial. That value is not independently required
to vanish. Its saved root equality is $[17751579,15]$.
The precision parent has $54056$ explicit variables and $15845$
disjoint private blocks of seventy-nine variables. No parent variable
node accesses a private coordinate. Replacing each block by its
private power-value call therefore preserves the full existential
assignment. Only the final parent root is equated to the identity.

At each primitive formula, equality is a difference, conjunction is
a sum of squares, and disjunction is a product. A topological induction,
including the globally conjoined fixed equations, proves equality of
polynomial values under substitution, not merely equisatisfiability
of an unrelated formula. This is the required link to the saved AST.

\section{A variable-shift relation for ordinary values}\label{sec:bridge}
We prove the ordinary-value deduction from the exact index model,
the compiled precision relation, the height theorem and $\mathsf{PC}$.
These are supplied by Sections~\ref{sec:model}, \ref{sec:compiler},
\ref{sec:height} and Theorem~\ref{hyp:PC}, respectively. Capital letters below denote five-coordinate point codes;
their lower-case indices are ordinary integers. The external variable
$t$ is any rational number.

\subsection{Approximation for every fixed shift}
Define $\Close(N,\mathcal K,t)$ by choosing $A,L,J$, linking all five
signed codes to precision inputs in the order $(n,\kappa,a,l,j)$, and
imposing $|t-r(a,j)|<\epsilon_l$ on those same coordinates. Put
\[
 B(n,\kappa)=H_6(S_{\rm ht})|\kappa F(n)|^{24}.
\]
The precision contract gives
\begin{align}
 \Close(\iota(n),\iota(\kappa),t)&\Longrightarrow
                  |t-n|<e^{-B(n,\kappa)},\label{eq:close-sound}\\
 \Close(\iota(n),\iota(\kappa),n)&\quad\text{for every }
                  n\notin\Cset,\ \kappa\ne0.\label{eq:close-complete}
\end{align}
The first follows from $|t-n|<2\epsilon_l<4^{-B}/2<e^{-B}$.
For the second, positive multiples of $P_0$ approach the identity
simultaneously in the compact group $E(\mathbf R)\times E(\mathbf R)$.
This follows, for instance, by partitioning the compact closure of the
cyclic subgroup into arbitrarily small neighborhoods and subtracting
two sufficiently close multiples. If the resulting positive differences
were bounded, a fixed nonzero multiple would be the identity; infinite
order excludes this. Choose $l$ first to make the positive $\epsilon_l$
small enough. Since the derivative of $[n]$ at the identity is $n$
\cite[IV.2.3(a), p.~122]{silverman},
$r(a,na)\to n$ along positive multiples $a$ approaching it. Choose such
an $a$ and put $j=na$. This works for every already fixed nonzero shift
$\kappa$, regardless of its size.

\subsection{The six local tests}
For $N=\iota(n)$, $\mathcal K=\iota(\kappa)$, form the code for
$\kappa n$, choose $H=\iota(h)$, $U=\iota(s_0)$, and put
$V_0=H-\mathcal K$, $V_1=U-\iota(\kappa n)$.
Require $\kappa\ne0$, $H,V_0\ne O$. Define $w(Q)=1/x(Q)$ for
$Q\ne O$, $w(O)=0$. For $t\notin\Cset$ use
\begin{align*}
 D_1&=w(H),&D_2&=w(U),&D_3&=\sigma w(V_0),&D_4&=\sigma w(V_1),\\
 D_5&=z(U)/z(H)-t,&D_6&=\sigma(z(V_1)/z(V_0))-t.
\end{align*}
Here $\sigma$ is conjugation in $K_0$. Choose $\eta,q\in\Q$ and impose
\begin{align*}
 \eta^2&=1,&R_0&=\eta2730^6F(t),&a_0&=1-R_0q^2\ne0,\\
 R&=R_0a_0^2,&C_j&=D_j/F(t),&W&=\prod_{j=1}^6(1+RC_j^2)\ne0,\\
 \Theta&=\Tr_{K_0/\Q}(W^{-1})\ne0,&B_{\rm pole}&=R/\Theta^2,
 &\Phi(B_{\rm pole})&.
\end{align*}
Call this relation $\Test(N,\mathcal K,t)$. All quotients are guarded
polynomially. There is one common denominator $F(t)$, one product $W$,
one trace and one pole instance.

\begin{lemma}[Divisibility]\label{lem:test}
If $\Test(\iota(n),\iota(\kappa),t)$ holds, then
$M_{S_{\rm ht}}(t)\mid\kappa F(n)$.
\end{lemma}
\begin{proof}
Let $p\mid M_{S_{\rm ht}}(t)$. Exactly one $c\in\Cset$ has
$\delta=v_p(t-c)>0$ odd. The contact differences are units, so
$v_p(F(t))=\delta$; the field is unramified at $p$. Put $b=v_p(a_0)$.
If $q=0$, then $b=0$. Otherwise $\delta+2v_p(q)$ is odd and
nonzero, so $b=\min(0,\delta+2v_p(q))\le0$. In either case
$\alpha=v_p(R)=\delta+2b$ is odd. At each place above $p$,
\[
 v(1+RC_j^2)=\min(0,\alpha+2(v(D_j)-\delta)).
\]
Let $\lambda=\max_{v\mid p}v(W)\le0$. Then
$v_p(\Theta)\ge-\lambda\ge0$. The pole input has odd valuation
$\alpha-2v_p(\Theta)$, which is positive by Lemma~\ref{lem:pole}.
Therefore $\alpha>0$ and $\lambda> -\alpha/2$.

At a place attaining $\lambda$, each nonpositive factor valuation is
at least $\lambda$. Consequently
\[
 v(D_j)\ge\delta-3\alpha/4=\delta/4-3b/2>0
 \quad(1\le j\le6).
\]
Integral valuations are at least one, at the same place for all tests.
The first four put $H,U$ in one formal group and $V_0,V_1$ in the
conjugate formal group. At the identity we set both $z(O)=0$ and
$\log(O)=0$; thus the cases $U=O$ and $V_1=O$ obey the same ratio
identities below. Their denominators are nonzero by the guards on
$H$ and $V_0$. For a nonzero point in either formal group,
\[
 \log(Q)=z(Q)(1+e_Q),\qquad v(e_Q)\ge1.
\]
Here $p>324$ and $v(p)=1$. By
\cite[IV.5.5 and IV.6.4, pp.~129--132]{silverman},
the logarithm is an injective isometry on $v(z)\ge1$.
The integral invariant differential gives in relative degree $r-1$
the bound $(r-1)v(z)-v_p(r)\ge1$ for every $r\ge2$.
Thus the guarded nonzero points $H,V_0$ have nonzero logarithms.
The fifth test first makes $z(U)/z(H)$ integral and congruent to
$t$; multiplication by the units $1+e_Q$ preserves that congruence.
The relation $[s_0]H=[h]U$ then gives
$\log(U)/\log(H)=s_0/h$. Likewise the sixth gives
$(s_0-\kappa n)/(h-\kappa)\equiv t$.
Multiplying by the integer denominators and subtracting yields
$v_p(\kappa(n-t))\ge1$. Since $t\equiv c$, we get
$p\mid\kappa(n-c)$, hence $p\mid\kappa F(n)$. The squarefree
product divides as asserted. No unit assumption on $\kappa$ was used.
\end{proof}

\begin{lemma}[One suitable shift]\label{lem:test-complete}
For every $n\in\Z\setminus\Cset$ there is
$\kappa\ne0$ with $\Test(\iota(n),\iota(\kappa),n)$.
\end{lemma}
\begin{proof}
Choose the sign so $R_0=\eta2730^6F(n)>0$ and let $S'$ be its finite
prime support. Each sufficiently small local formal neighborhood has
finite index in the compact local group
\cite[VII.6.3, p.~200]{silverman}. A positive common multiple
of the quotient orders gives $\kappa$ with $[\kappa]P_0$ in every
chosen neighborhood. Take
\[
 H=[2\kappa]P_0,\quad U=[2\kappa n]P_0,\quad
 V_0=[\kappa]P_0,\quad V_1=[\kappa n]P_0.
\]
The formal parameter and the derivative of $[n]$ make each $C_j$
integral above $S'$ after sufficiently shrinking the neighborhoods.
Infinite order preserves the nonzero guards.

Choose rational $q$ integral at $S'$ and arbitrarily close to
$1/\sqrt{R_0}$ in $\mathbf R$, without equality. Denominators
$R_0m+1$ with nearest integer numerators give such choices, adjusting
by one if equality occurs. At $S'$, $a_0$ is a unit and
$W\equiv1$, $\Theta\equiv2$ modulo $p^{v_p(R_0)}$.
Integral trace justifies this also at the ramified place two.
The pole input has valuation at least four at two and at least six
at the other primes of $S_0$. Outside $S'$, its valuation is
$2v_p(a_0)-2v_p(\Theta)$ and is even, so its reduced denominator is
a square. In both real embeddings, $R$ tends to zero positively,
$W\to1$ and $\Theta\to2$. All guards hold and
$0<B_{\rm pole}<1/2$ for a sufficiently close choice.
Theorem~\ref{hyp:PC} supplies the pole witnesses.
\end{proof}

\subsection{Exact ordinary value}
Define $\Val(N,t)$ by membership of $N$ and the disjunction
\[
 \bigvee_{c\in\Cset}(N=\iota(c)\wedge t=c)
 \quad\vee\quad
 \exists\mathcal K\,
 (\Close(N,\mathcal K,t)\wedge\Test(N,\mathcal K,t)).
\]
The same five coordinates $\mathcal K$ are used in both calls.
In an accepting main branch $n\notin\Cset$, $\kappa\ne0$, and
Lemma~\ref{lem:test} and the height estimate imply
\[
 h(t)\le H_6(S_{\rm ht})|\kappa F(n)|^{24}=B(n,\kappa).
\]
For $t=a/b\ne n$ in reduced form, $|a-nb|\ge1$, hence
\[
 |t-n|\ge1/b\ge e^{-h(t)}\ge e^{-B(n,\kappa)},
\]
contradicting~\eqref{eq:close-sound}. Thus $t=n$.
Conversely choose $\kappa$ by Lemma~\ref{lem:test-complete}, then use
\eqref{eq:close-complete} for that same shift. Private witness blocks
are disjoint, so the two calls can be satisfied simultaneously.
The six remaining integers have exactly their displayed branches.

Therefore $\Val(\iota(n),t)$ is exactly the graph $t=n$.
Existentially quantifying $N$ gives ordinary $\Z$ in $\Q$.
S01 and S16 multiply the exception and main-branch polynomials and
conjoin membership by a sum of squares. The literal substitutions in
Section~\ref{sec:syntax} and Appendix~\ref{app:literal} identify this
polynomial with the saved root $P_*$. This proves the zero-set part of
Lemma~\ref{lem:assembly}.

% End 04-computation-value.tex


% Begin 05-bounds-appendices.tex
\section{The cutoff and polynomial size}\label{sec:size}
The literal constants in S03 define the catalogue product $H_{\rm cat}$,
its maximum $H_{\max}$, $N_s=P_sH_{\rm cat}$ with
$P_s=\prod_{p\le324}p$, the least formal multiple $M_0$, and $c_*=4M_0$.
They take
\[
 R=\max(D,H_{\max},c_*,10^{12}),\quad
 A_{\rm ht}=2\cdot10^{13}\max(1,C_{\rm cm},Q_R^{17}),\quad
 T=(48A_{\rm ht})^{13}.
\]
The pole clearing constant is $\ell_0=12P_s^{2916}H_{\rm cat}^{732}$.
The symbol $F_0=324!$ is distinct from the six-contact polynomial $F(u)$.

\begin{proposition}[Literal cutoff]\label{prop:cutoff}
The fixed prescription satisfies
\[
 R=D,\qquad T=(96\cdot10^{13})^{13}D^{2210G_0^2D^3}
            <2^{\,2^{\,2^{9013}}}.
\]
This is a statement about the prescribed integer, independently of its
semantic role in excluding false multiplication triples.
\end{proposition}
\begin{proof}
We give the bounds of S25. Exact small calculations give
$10^5<W<2^{512}$ and $2^{2000}<F_0<2^{2250}$.
For $q\ge10^5$, the unconditional progression estimate
\cite[Theorem 1.2, (1.10)--(1.11), Corollary 6.17 with $A=1$]{bmor}
applies for $x\ge\exp(0.03\sqrt q(\log q)^3)$, including the possible
exceptional-zero case. At $X=2^{W^3}$ its error is at most
$X/(160\log X)$ and the hypothesis follows from
$0.03\sqrt W(\log W)^3<(\log2)W^3$.
Hence $\vartheta(X;W,a)>X/(2W)$, giving more than
$X/(2W^4)>22$ primes in each unit class. No prime at most 324
belongs to such a class. Thus
\[
 H_{\max}<X,\quad H_{\rm cat}\le X^{22W},\quad
 \log_2N_s<2^{2054}.
\]
For every odd prime dividing $N_s$, take the actual reduction order of
$P_0$ at each place above it, and let their product be $m_0$.
All preceding positive multiples stay in the dyadic formal subgroup,
since $P_0$ starts there and the subgroup is closed under addition.
A nonidentity point there has negative abscissa valuation, so its
reciprocal is defined; the constructor accepts the identity directly.
Thus every earlier search step is defined without a nontorsion premise.
At the bad place $s\equiv3\pmod7$, the point reduces to $(2,6)$
on $y^2=x^3$ and is nonsingular because $2y\not\equiv0$.
The other place $s\equiv4\pmod7$ has good reduction.
Take the actual order of the reduced point in the nonsingular
reduction group, using
\cite[VII.2.1, p.~188]{silverman}.
At a good place the residue degree is at most two, so Hasse's
bound \cite[V.1.1, p.~138]{silverman} gives order at most
$(p+1)^2<p^3$ for odd $p$. The singular smooth group over
$\F_7$ has order seven. There are at most two places over each
rational prime, so $m_0\le N_s^6$.
If $[N_s^4m_0]P_0$ is the identity, the search has already succeeded.
Otherwise its reciprocal coordinate is defined and the following
estimates apply. At each odd place, $[m_0]P_0$ is formal, and multiplication by
$N_s^4$ gives
\[
 v(1/x([N_s^4m_0]P_0))\ge2+8v_p(N_s).
\]
At two normalize $v_s(s)=1$, so $v_s(2)=2$. Initially
$v_s(x)=-6$, $v_s(y)=-9$, and $v_s(z)=3>2$.
The log isometry therefore gives
$v_s(z([2]Q))=v_s(z(Q))+2$ on this subgroup, and hence
\[
 v_s(1/x([N_s^4m_0]P_0))\ge6+16v_2(N_s).
\]
Writing this reciprocal as $a+bs$, extraction of $a$ and $b$
uses division by $2$ and $2s$ and loses at most three units of
$v_s$. Thus both rational coefficients lie in
$N_s^4\Z_{(N_s)}$, where denominators in $\Z_{(N_s)}$ are
prime to $N_s$. This proves the required divisibility and guard. Therefore the literal search
succeeds by $N_s^4m_0$ and $M_0\le N_s^{10}$.
This uses actual finite orders, not numerical upper bounds incorrectly
treated as multiples of those orders.

Consequently $c_*<2^{2^{2059}}$, whereas
$2^{2^{8008}}<D<2^{2^{9010}}$, so $R=D$.
Next $C_{\rm cm}<2^{50000D^2}$ and
\[
 \log_2Q_D^{17}=170G_0^2D^3\log_2D>50000D^2.
\]
The maximum is therefore $Q_D^{17}$, giving the identity for $T$.
Finally $G_0<2^{37}$, $E_0<D^3$, $\log_2Q_D<D^5$,
$\log_2A_{\rm ht}<D^7$, and $\log_2T<D^8<2^{2^{9013}}$.
\end{proof}

\subsection{Quantitative syntax bounds}
Use coefficient $\ell^1$ norm for integer polynomials and the weighted
pair norm $\|a+bs\|=\|a\|_1+2\|b\|_1$. They are submultiplicative.
The division-polynomial recurrences give
\[
 \deg f_n\le(n+1)^3,\qquad\|f_n\|\le2^{100(n+1)^3}.
\]
For the even recurrence the exponent slack is
$4m^3-24m-27\ge9$ for $m\ge3$; for the odd recurrence it is at least
ten for $m\ge2$. This absorbs the fixed factors and addition.

The following facts have been checked for the actual saved expression:
the absolute values of all literal fixed constants and all occurring
division-polynomial indices are at most $T$; the specified syntax-height
traversal has root height at most 256; each model recipe has at most
$2^{20}$ constraints and each composition recipe has at most $2^{20}$
terms. The check accounts for all 23 templates, seven recipes, private
witness blocks, 173 scalar constant nodes and 50 finite constructors. Its exact
artifact identity and independent replay are recorded in the source map.
Together with the constructor bounds below, this discharges the
quantitative part of $\mathsf B$. The all-input arithmetic interpretation of the same constituents is the
separate substitution induction of Section~\ref{sec:compiler}.

Put $B_*=T+1$. Coordinates of $\iota(n)$ for
$|n|\le T$ have absolute numerator and positive denominator at most $2^{K_*}$, with
$K_*=2^{30}(T+1)^3$. Evaluate the division polynomials at $P_0$ and use
the fixed coordinate formulas; at most eight quadratic-field operations
are needed, each taking an available height bound $H$ to at most $3H^4$.
There are at most $8(T+1)^3$ guard factors. The resulting bounds for
each finite constructor are
\[
 \deg\le B_*^4,\qquad\log_2\max(1,\|\cdot\|_1)\le B_*^8.
\]

Give variables syntax height zero, constants height one, and add one
for each arithmetic operation or finite constructor. A template call
adds one plus the callee's intrinsic height and the maximum argument
height. The substitution inequalities
\[
 \deg P(Q)\le\deg P\max(1,\max_i\deg Q_i),\qquad
 \|P(Q)\|_1\le\|P\|_1\max(1,\max_i\|Q_i\|_1)^{\deg P}
\]
give degree and logarithmic norm at most $B_*^{8h}$ at positive height
$h$, with the bounded sums in the actual recipes treated as in S26.
This general call rule applies to polynomial arguments. In each outer
composition recipe the arguments are coordinate variables, so the call
is only a renaming, even when coordinates repeat. For any such map
$\sigma$,
\[
 \deg P(Y_{\sigma(i)})\le\deg P,\qquad
 \|P(Y_{\sigma(i)})\|_1
 =\sum_\beta\left|\sum_{\alpha\mapsto\beta}c_\alpha\right|
 \le\sum_\alpha|c_\alpha|=\|P\|_1.
\]
Thus these calls retain the callee's majorant height before charging
the sum-of-squares or product connective. The seven recipe heights are
$228,233,235,251,252,254,256$. A rule charging an extra unit for each
of these mere renamings would describe a looser syntax-depth bound,
not this proved degree-and-norm majorant.
For the model wrappers use the safe height bound four for the
generator equations and root differences; the callee contribution
dominates this bound. The literal height-two claim for those equations
in S26 is not used.
This height is a majorant for the compact constructor syntax, not for
an expanded monomial tree. The height-256 bound yields
\[
 \deg P_*\le(T+1)^{2048},\qquad
 |[X^\alpha]P_*|\le2^{(T+1)^{2048}}.
\]

Since $T$ is an integer, $T+1\le U=2^{2^{2^{9013}}}$. Hence
\[
 (T+1)^{2048}\le U^{2048}=2^{2^{2^{9013}+11}}
                     <2^{2^{2^{9014}}}.
\]
As $9014<65536=\tw4$, the degree is less than $\tw7$ and the
coefficient magnitude less than $\tw8$. The root recipe declares
$292\,890\,231\,904$ witnesses. Together with the preceding estimates,
these complete the bounds for the same $P_*$ in
Theorem~\ref{thm:target}.

\section{Consequences of the integer definition}
The proof of Theorem~\ref{thm:target} combines the following results.
\begin{center}\small
\begin{tabular}{p{.20\textwidth}p{.43\textwidth}p{.27\textwidth}}
\toprule
Result & Mathematical inputs & Proof location\\\midrule
Index group & Two-isogeny descent, reductions, finite height search &
Section~\ref{sec:model}\\
Six-contact bound & Quaternionic model, tame pullback, compatible-isogeny
height theorem, theta comparison & Section~\ref{sec:height}\\
Pole completeness & Simultaneous descent, restricted converse, forward
rank-one theorem & Section~\ref{sec:pole}\\
Multiplication & Index group, height bound, selected primes,
pole completeness, finite guard & Section~\ref{sec:mult}\\
Precision relation & Decision program, lowering, power and packed histories &
Section~\ref{sec:compiler}\\
Ordinary value & Shared-shift Close and Test, height and rational separation &
Section~\ref{sec:bridge}\\
Bounded $P_*$ & Literal substitutions and quantitative syntax &
Sections~\ref{sec:syntax}, \ref{sec:size}; Appendix~\ref{app:literal}\\\bottomrule
\end{tabular}
\end{center}

\subsection{Three corollaries}\label{sec:four-corollaries}
We apply the exact equality between the rational zero projection of
$P_*$ and ordinary $\Z$ proved in Theorem~\ref{thm:target}.

\begin{corollary}[Negation of Mazur's finite-component conjecture]
\label{cor:mazur}
Mazur's
finite-connected-components conjecture for the real-topological closure
of rational points, Conjecture~3 of \cite{mazur1992,mazur1995}, is
false. Its stronger Conjecture~2 is false as well.
\end{corollary}
\begin{proof}
We recall Mazur's projection argument, originating in
\cite{mazur1992} and repeated in \cite[p.~166]{mazur1995}.
Conjecture~3 in \cite[p.~165]{mazur1995} asserts that the closure
of the rational points of any variety over $\Q$ in its real locus has
finitely many connected components. In that source a variety means a
reduced scheme of finite type; irreducibility is not required.
Put $r=292\,890\,231\,904$ and
\[
 V=\bigl(\{P_*(t,w_1,\ldots,w_r)=0\}\subset\mathbf A^{r+1}_{\Q}\bigr)_{\rm red},
 \qquad C=\overline{V(\Q)}\subset V(\mathbf R).
\]
Reduction changes neither the rational nor the real points. Let
$\pi:V(\mathbf R)\to\mathbf R$ be projection to $t$.
Theorem~\ref{thm:target} gives $\pi(V(\Q))=\Z$. Since $\Z$ is closed in $\mathbf R$
and $\pi$ is continuous, $C\subset\pi^{-1}(\Z)$, so $\pi(C)\subset\Z$.
Every integer already occurs on $V(\Q)$, hence $\pi(C)=\Z$.
The image of any connected component of $C$ is a connected subset of
$\Z$, therefore a singleton. Points projecting to distinct integers
belong to distinct components. Thus $C$ has infinitely many connected
components, contradicting the stated conjecture. Conjecture~2 in
the same source would give the homotopy type of a finite simplicial
complex and hence only finitely many connected components, so it
implies Conjecture~3 and is contradicted too.
\end{proof}

\begin{corollary}[A universal rational Diophantine polynomial]
\label{cor:rational-universal}
There exist $m\ge1$ and a fixed polynomial
$U\in\Z[e,n,x_1,\ldots,x_m]$ such that, for every $e,n\in\N$,
\[
 \exists\mathbf x\in\Q^m\ U(e,n,\mathbf x)=0
 \quad\Longleftrightarrow\quad
 \text{machine $e$ halts on input $n$}.
\]
Here $e$ and $n$ are ordinary integer parameters, with $\N$ embedded
in $\Q$, and are not elliptic-curve coordinates. In particular, every
computably enumerable subset of $\N$ is represented by a suitable
fixed value of $e$ in this one polynomial.
\end{corollary}
\begin{proof}
Fix an effective enumeration of machines. The Davis--Putnam--Robinson--Matiyasevich theorem, \cite[Basic Theorem, p.~3]{matiyasevich1971},
applied to the universal halting relation gives a fixed
$D\in\Z[e,n,z_1,\ldots,z_s]$, with $s\ge1$, such that, for $e,n\in\N$,
\[
 \text{machine $e$ halts on $n$}
 \quad\Longleftrightarrow\quad
 \exists\mathbf z\in\Z^s\ D(e,n,\mathbf z)=0.
\]
To use integer witnesses in this statement, replace each nonnegative
witness by a sum of four integer squares. In the positive-integer
convention of the cited source, apply the theorem to the predicate
that machine $a-1$ halts on $b-1$, substitute $a=e+1$, $b=n+1$,
and use one plus that sum for each witness; see also
\cite[p.~2]{matiyasevich1971}. A dummy witness ensures $s\ge1$ if needed.
For each $z_i$ take a separate block $\mathbf w_i$ of $r$ rational
variables, where $r$ is the witness count of $P_*$. Define
\[
 U(e,n,\mathbf z,\mathbf w_1,\ldots,\mathbf w_s)
   =D(e,n,\mathbf z)^2+\sum_{i=1}^s P_*(z_i,\mathbf w_i)^2.
\]
A rational zero makes every summand zero. Theorem~\ref{thm:target}
then forces every $z_i$ to be an integer and yields the displayed
halting condition. Conversely, integer witnesses for $D$ have separate
rational witness blocks for $P_*$, giving a zero of $U$. All blocks
are private and the polynomial is fixed independently of $e,n$.
No variable, degree or coefficient bound proved for $P_*$ is asserted
for this new polynomial $U$.
\end{proof}

OpenAI already obtained Turing degree $0'$ for this set in
\cite[Corollary 7.1, pp.~56--57]{openaihtp}; that paper explicitly
does not assert many-one completeness. The stronger direction below
is $0'\le_m\operatorname{HTP}(\Q)$, obtained from the integer
definition through the universal polynomial. The reverse many-one
reduction is the standard enumeration of rational solutions.

\begin{corollary}[The many-one degree of $\operatorname{HTP}(\Q)$]
\label{cor:htp-degree}
\[
 \operatorname{HTP}(\Q)\equiv_m 0',
 \qquad 0'=\{e\in\N:\text{machine $e$ halts on input $e$}\}.
\]
Here $\operatorname{HTP}(\Q)$ is the set of codes of integer-coefficient
polynomials, of arbitrary finite arity, having a rational zero, under a
fixed effective coding.
\end{corollary}
\begin{proof}
Given $e$, specialize the fixed polynomial of
Corollary~\ref{cor:rational-universal} at its ordinary integer parameters
to obtain $U(e,e,\mathbf x)$. Computing its coefficient list and code is
effective, and
\[
 e\in0'\quad\Longleftrightarrow\quad
 \ulcorner U(e,e,\mathbf x)\urcorner\in\operatorname{HTP}(\Q).
\]
This proves $0'\le_m\operatorname{HTP}(\Q)$.

Conversely, from the code of a polynomial $f$ construct a machine that
ignores its input and enumerates all rational tuples of the appropriate
arity. It evaluates $f$ using exact rational arithmetic and halts as soon
as it finds a zero. Its index $j(f)$ is computable from the finite code
of $f$, and
\[
 f\in\operatorname{HTP}(\Q)\quad\Longleftrightarrow\quad j(f)\in0',
\]
where the left side uses the code of $f$. Thus
$\operatorname{HTP}(\Q)\le_m0'$. Invalid polynomial codes, if present in
the coding convention, are sent to a fixed nonhalting index. Using
rational-coefficient polynomial codes gives the same result, since
clearing their finitely many denominators is effective and preserves
rational zeros. The two reductions prove the assertion.
\end{proof}

More generally, an integer Diophantine problem translates to a rational
one by adjoining $P_*$ for every integer witness, with private witness
blocks. Natural witnesses can first be expressed as sums of four integer
squares. Bounds for the resulting polynomial require their own composition
calculation; the bounds for $P_*$ do not automatically transfer.

For example, a machine that enumerates ZFC proofs and halts upon finding
a contradiction has a rational Diophantine halting condition by
Corollary~\ref{cor:rational-universal}. This equivalence does not assert
that a contradiction exists, and no additional explicit polynomial is
specified here.

% Begin degree-eight construction
\section{The base residual compiler}
\label{sec:degree8-original}

This section records the base residual compiler used as a construction
stage for the quartic in Section~\ref{sec:quartics}. Its degree-eight
sum-of-squares encoding fixes the arithmetic interfaces and conventions;
it is not a third final polynomial specification. The construction uses
a joint comparison of four elliptic indices, a Lucas equation, and the
five-contact height estimate proved below. The archived counts describe
this base compiler before its reciprocal revision and quadratic lift.

\begin{proposition}[Finite specification of the base compiler]\label{prop:o8-finite}
The prescription below specifies a polynomial
\[
 F_8\in\Z[X,V_1,\ldots,V_{15671}]
\]
of degree exactly eight, with exactly $129497$ nonzero monomials and
maximum absolute coefficient
\[
 685783949153203534233601<2^{80}.
\]
These statements are finite identities about the displayed construction
and its integer coefficient table.
\end{proposition}

\begin{theorem}[Arithmetic interpretation]\label{thm:o8-conditional}
For every
$x\in\Q$,
\[
 x\in\Z\quad\Longleftrightarrow\quad
 \exists(v_1,\ldots,v_{15671})\in\Q^{15671},\qquad
 F_8(x,v_1,\ldots,v_{15671})=0.
\]
\end{theorem}

The required height estimate is proved below; Theorem~\ref{thm:height}
uses a different contact set. The seven-variable pole lemma below is proved from
Theorem~\ref{hyp:RC}; it is not inferred from Theorem~\ref{hyp:PC},
which has different variables and a restricted completeness domain.
The finite coefficient calculation does not establish the height bound.
The supplied unpublished Freitag archive is the provenance of this
construction and of the original coefficient table. It is not cited here
as an arithmetic theorem. The published inputs already used in
Sections~\ref{sec:model} and \ref{sec:mult} remain
available, together with the quaternion trace theorem specified below.

\subsection{The height input and the exact pole predicate}

Put
\[
 S_* =\{p:p\le76176\},\qquad
 \mathcal D=\{0,1,2401,-343/32\},
\]
\[
 H_*=2(4608^{165888}+6^{100})
          \left(\prod_{p\le76176}p\right)^{24}+128.
\]
For $a\in\Q\setminus\mathcal D$, let $M_5(a)$ be the product of
the primes outside $S_*$ at which the point $(a:1)$ has positive odd
contact order with one of $\infty,0,1,2401,-343/32$. Explicitly, choose
local homogeneous coordinates $(u:v)$ and subtract
$\min(v_p(u),v_p(v))$ from the valuations of $v$ and of
$u-\delta v$ for $\delta\in\mathcal D$.

\begin{theorem}[The five-contact height estimate]\label{hyp:o8-height}
For every $a\in\Q\setminus\mathcal D$,
\[
 h(a)\le H_*M_5(a)^{24}.
\]
The displayed constant satisfies $\log_2H_*<4783227$.
\end{theorem}
\begin{proof}
Write $s=a$, $M(s)=M_5(a)$, $H=H_*$,
$B_*=\prod_{p\le76176}p$, and
$\mathcal C=\{\infty,0,1,2401,-343/32\}$.
We prove the geometric comparison, the non-CM estimate through the
modular level, and the CM estimate over the same fiber field.
All heights use the original stable Faltings normalization.
\paragraph{Recovering the contact line from published quotients}
Let $B/\mathbb Q$ be the indefinite quaternion algebra of discriminant $210$,
$\mathcal O$ a maximal order, $X$ its maximal-level Shimura curve, and $W$
its Atkin--Lehner group. The standard genus formula gives
\[
 g(X)=1+\frac{\varphi(210)}{12}=5,
\]
because the elliptic-point terms vanish: $5$ splits in $\mathbb Q(i)$ and
$7$ splits in $\mathbb Q(\sqrt{-3})$, so neither quadratic field embeds in
$B$. The group $W$ is $(\mathbb Z/2\mathbb Z)^4$ with involutions $w_m$
indexed by square classes of divisors of $210$.

The published models of \cite[Theorem 4.2 and Table 2, pp.~938--939]{GonzalezRotger2006} have the following coefficient rows
$(a,b,c,d,e)$ for $y^2=ax^4+bx^3+cx^2+dx+e$:
\[
\begin{array}{c|rrrrr}
 m&a&b&c&d&e\\\hline
30&-43&-686&-2915&-1372&-172\\
42&-43&600&-986&-600&-43\\
70&-43&256&-130&-768&-387\\
105&-43&2&505&12&-1548
\end{array}
\]
The cited theorem identifies each with $X/\langle w_m\rangle$.
The genus formula and the Atkin--Lehner group law used here also appear
in the same published paper, Proposition 5.2 and the paragraph following
Proposition 5.3, pp.~940--941.

Here is a deduction of all the required branch information, rather than
an additional assumption about it. Each of these four genus-one quotients
has eight fixed points upstairs, by Riemann--Hurwitz. A stabilizer for the
$W$ action on a characteristic-zero smooth curve is cyclic, and therefore
has order at most two. Each nonempty fixed set of an involution is a union
of $W$-orbits of size eight. Riemann--Hurwitz for $X\to X/W$ now gives
\[
 8=16(2g(X/W)-2)+\sum_{1\ne w\in W}\#\operatorname{Fix}(w).
\]
The four known fixed sets already contribute $32$. Thus $g(X/W)=0$,
the total is $40$, and there is exactly one further orbit of eight fixed
points. Inertia generators around the branch points have product one.
The four classes $30,42,70,105$ are independent in $W$, and their product
is the class $210$. Hence there are exactly five branch points, each of
index two, with distinct labels $30,42,70,105,210$. Since the involutions
are defined over $\mathbb Q$ and each label occurs only once, each branch
point is rational. Three of them trivialize the genus-zero quotient.

Normalize $b_{210}=\infty$, $b_{105}=0$, $b_{70}=1$, and write
$b_{42}=\lambda$, $b_{30}=\mu$. On a genus-one quotient the residual
group $(\mathbb Z/2)^3$ has translation subgroup of order four, necessarily
the full two-torsion group after choosing an origin. Its quotient has the
same Jacobian $j$-invariant: it is multiplication by two over the algebraic
closure. The remaining double cover of $X/W$ branches at the four points
other than $b_m$. Therefore the four quartic invariants determine the four
Legendre invariants of these four-point subsets.

For a quartic put
\[
 I=12ae-3bd+c^2,\quad
 J_q=72ace+9bcd-27ad^2-27b^2e-2c^3,\quad
 j=\frac{6912I^3}{4I^3-J_q^2}.
\]
The exact computation gives the Legendre parameters
$1/2401,32/375,1/225,5/1029$, in row order. Enumerating the six cross-ratio
values of each of the first two parameters and testing the last two
invariants leaves exactly
\[
 \lambda=2401,\qquad\mu=-343/32.
\]
The script \texttt{check\_height.py} reproduces this rational computation.
Every nonzero determinant of the five primitive contact forms has prime
support in $\{2,3,5,7\}$.

\paragraph{Integral geometry and the degree 576 field}
Choose $\mu_0\in\mathcal O$ with $\mu_0^2=-210$. The fixed-maximal-order
statement of \cite[Section 1, pp.591--592]{buzzard} supplies such an element
in this chosen $\mathcal O$, and states that the associated involution
preserves $\mathcal O$.
Use the maximal moduli stack $\mathscr X$ for this order and involution.
Yuan's maximal-level theorem \cite[Section 2.1, pp.~230--231]{Yuan2024}
makes it regular, proper and semistable over $\mathbb Z$, with a universal
abelian surface, and smooth over $R_0=\mathbb Z[1/210]$.
Its generic stabilizer is $\{\pm1\}$.
We now construct the entire fine principal-level-three model over $R_0$.
This construction uses the maximal-level statement only: an
$\mathcal O$-linear level map, without an isomorphism requirement, would
not define the required full marking.

\paragraph{The full marking at the ramified prime}
\label{sec:h5-full-three-construction}
Over $\mathscr X_{R_0}$ consider the functor of triples
\[
 (A,i,\eta),\qquad
 \eta:(\mathcal O/3\mathcal O)_S\xrightarrow{\ \sim\ }A[3],
\]
where $\eta$ is an isomorphism of finite etale group schemes respecting
the left $\mathcal O$-action. Morphisms preserve $i$ and $\eta$.
Since $3$ is invertible, $A[3]$ is finite etale of rank $3^4$
\cite[Theorem 8.2 and Remark 8.4, pp.~115--116]{milne1986-abelian}.
At a geometric point, $V_3(A)$ is a faithful module of dimension four
over $\mathbb Q_3$ for the division algebra $B_3$ of dimension four;
it therefore has rank one over $B_3$.
Every full left $\mathcal O_3$-stable lattice $L$ in $B_3$ is principal.
Indeed, choose $x\in L\setminus\{0\}$ with least valuation
\[
 w(x)=\tfrac12\operatorname{ord}_3\operatorname{nrd}(x).
\]
Such a least value exists since the lattice is bounded and the value
group is $\tfrac12\mathbb Z$. For $y\in L$, one has
$w(yx^{-1})\ge0$, hence $y\in\mathcal O_3x$; stability proves the
reverse inclusion. The valuation and maximal-order assertions are
\cite[13.3.4 and 13.3.11, pp.~206--207]{voight-quaternion}; the
factor $1/2$ is retained. Consequently $T_3(A)/3$, and therefore $A[3]$,
is the regular left $\mathcal O/3\mathcal O$-module. This argument uses
the division algebra at $3$, not a split matrix order.

Etale-locally on the base both group schemes are constant. Their
isomorphisms form a finite etale scheme, and the equalities expressing
$\mathcal O$-linearity on a finite integral basis select an open and
closed subscheme. Descent gives the representable finite etale
$\operatorname{Isom}_{\mathcal O}$ cover $\mathscr Y\to\mathscr X_{R_0}$.
The preceding lattice calculation proves that every geometric fiber
is nonempty. The radical of $\mathcal O_3$ has residue field
$\mathbb F_9$ and square $3\mathcal O_3$ up to a unit, so
$\mathcal O/3\mathcal O$ has $81$ elements, radical of size $9$, and
$72$ units. The automorphism group of its regular left module is the
unit group of the opposite ring. Thus this cover has degree $72$
over the stack. It is proper and smooth over $R_0$ and carries the
pullback of the universal abelian scheme.

\paragraph{Rigidity and projectivity of this same cover}
A geometric stabilizer is finite, since $\mathscr X$ is proper
Deligne--Mumford. Its action on $T_3(A)$ is faithful
\cite[Lemma 12.2, p.~122]{milne1986-abelian}. A stabilizer preserving
$\eta$ acts in $1+3M_4(\mathbb Z_3)$, which has no nonidentity torsion.
To see this, an element of prime order $l$ would have the form
$1+3^aM$, with $a\ge1$ and $M\not\equiv0\pmod3$. If $l\ne3$, its
$l$th power has nonzero first term $l3^aM$ modulo $3^{a+1}$.
If $l=3$, its cube minus $1$, divided by $3^{a+1}$, is
\[
 M+3^aM^2+3^{2a-1}M^3\not\equiv0\pmod3.
\]
Both contradict the assumed order. All geometric stabilizers of
$\mathscr Y$ are trivial. This also gives rigidity on arbitrary bases:
its inertia is finite and unramified, the identity section is both open
and closed, and its complement has empty geometric fibers. Hence the
entire inertia is the identity section. The stack is therefore an
algebraic space $Y$, by \cite[Tag 03YR]{stacks-h5-fine}, and the same
universal abelian scheme descends to it.

Here projectivity needs an argument; proper algebraic spaces need not
be schemes. Each geometric fiber of $Y/R_0$ is a proper smooth
algebraic-space curve, and is a scheme by
\cite[Tag 0ADD]{stacks-h5-fine}. Every generic component maps finitely
and nonconstantly to the genus-five coarse curve already identified
above from \cite{GonzalezRotger2006}. Riemann--Hurwitz gives genus at
least five on each such component. The algebraic-space Stein
factorization, with the finite-etale component assertion
\cite[Theorem 76.36.5 and Lemma 76.36.9, Tag 0A18]{stacks-h5-fine},
gives a finite etale component base $R'_0/R_0$, over which the fibers
are geometrically connected. Smooth proper genus is locally constant
\cite[Tag 0E83]{stacks-h5-fine}. Every component of this finite etale
base meets the generic fiber, so every geometric fiber component of
$Y$ has genus at least two. The line $\Omega^1_{Y/R_0}$ is consequently
ample on every fiber. The fiberwise ampleness theorem for proper
algebraic spaces \cite[Tag 0D3A]{stacks-h5-fine} makes it relatively
ample near every base point. Tag 0D32 then makes $Y$ a scheme with a
relatively ample line. Finally Tag 0B45, implication (6) to (1), gives
projectivity, since the affine base itself has an ample line bundle.
This proves smooth proper projective representability of the actual
full-marking functor over $R_0$.

\paragraph{The complete generic moduli problem}
It remains to identify the whole generic fiber, rather than an
unspecified form of its complex curve. Let $V=B$ with its regular left
$B$-action and the alternating form $\psi$ of the selected positive
polarization. It satisfies
\[
 \psi(bx,y)=\psi(x,b^*y),\qquad b^*=\mu_0^{-1}\bar b\mu_0.
\]
The space of compatible alternating forms is one-dimensional over
$\mathbb Q$. After splitting $B$, write $V=W_1\otimes W_0$ with both
factors of dimension two. The involution on $\operatorname{End}(W_1)$
is adjoint to a symmetric form, and a compatible alternating form is
that symmetric form tensored with an alternating form on $W_0$.
The latter space is one-dimensional. Its similitude group is the
right-multiplication copy of $(B^{\mathrm{op}})^\times$; after splitting
it is $\operatorname{GSp}(W_0)$, hence the connected type-$C_1$ case
of \cite[Example 8.6 and Proposition 8.7, p.~325]{milne-shimura2005}.

We use the full Hodge-orbit condition of
\cite[Theorem 7.4 and its preceding definition, pp.~321--322]{milne-shimura2005}.
The general implication printed in Proposition 8.14 and Theorem 8.17
requires the author's correction; the polarization condition alone
does not suffice in general. For this exact datum the additional
condition follows directly. Over $\mathbb R$, the centralizer of the
left algebra acts on $W_0$ as $\operatorname{GL}_2(\mathbb R)$.
An admissible complex structure is $J$ on $W_0$ with $J^2=-1$.
For $v\ne0$, $(v,Jv)$ is a basis, and its matrix $g$ satisfies
$J=gJ_0g^{-1}$. Thus $1\otimes g$ lies in the actual centralizer
$G(\mathbb R)$, commutes with the fixed left $B$-action, and conjugates
the entire Hodge map $a+bi\mapsto1\otimes(a+bJ)$ to the fixed map.
Both half-planes occur. Positive rational multiples refer here only
to representatives that are actual polarizations on the same
abelian variety; the rational comparison retains the full
$\mathbb Q^\times$-class of the alternating form and allows either
sign of similitude multiplier. No positive-multiplier subgroup
replaces $G(\mathbb R)$.

The rational homology of a surface in this moduli problem has rank
one over $B$, so a $B$-linear identification with $V$ exists. The
one-dimensional calculation of forms makes it a similitude, and the
preceding calculation verifies the missing Hodge-orbit condition
for this identification. A basis $b_i$ of $B$ can be encoded by the
tensors $\psi(x,b_i y)$, as in loc. cit., 8.16. The determinant
condition is the corresponding trace condition. These checks give
all the conditions of 7.4 without using the erroneous general
implication from 8.17. The reflex field is $\mathbb Q$: the
cocharacter class is the usual $(1,0)$ class for this quaternionic
inner form of $\operatorname{GL}_2$. The standard involution identifies
the opposite group with the conventional $B^\times$ description and
preserves the order, center, normalizer and principal congruence
subgroup, compatibly with the maximal-level datum.

An integral $\mathcal O$-linear Tate isomorphism is automatically a
similitude, by the same one-dimensional calculation. At $3$ its
multiplier is a unit, since both integral forms are perfect; it is
not required to be $1$ modulo $3$. Every isomorphism modulo $3$
lifts, since $\mathcal O_3^\times\to(\mathcal O_3/3\mathcal O_3)^\times$
is surjective. Thus the full marking has stabilizer exactly $U(3)$,
with all its multipliers and geometric components. Fixing one
Weil-pairing value would give a different functor.

\paragraph{Integral lattices and the canonical rational model}
For a complex marked surface choose integral $\mathcal O$-linear
Tate isomorphisms at every finite prime, matching $\eta$ at $3$.
At split primes these exist by the matrix-unit decomposition of a
stable lattice; at division primes the valuation argument above
applies. Their $U(3)$-orbit is independent of the choices. Together
with the rational polarization class and $B$-action this gives the
quadruple in Milne's moduli description. Conversely the finite-adelic
orbit specifies a full lattice in the rational Tate space. An
isogenous abelian variety realizes this lattice: after multiplying
by an integer, take the quotient by the corresponding finite torsion
subgroup. Stability gives the integral $\mathcal O$-action, and the
same orbit gives the marking at $3$.

The compatible rational polarization class has a unique positive
representative perfect on this lattice. Locally its dual-lattice
discrepancy is scalar; a positive rational scalar cancels the
finitely many valuations, and two such scalars have ratio a positive
rational unit at every prime, hence ratio $1$. Equivalently, the
canonical principal-polarization assertion
\cite[p.~592]{buzzard} gives this polarization from the integral
order action and fixed involution. That assertion precedes the
prime-to-discriminant level restrictions in Buzzard's Section 2;
we do not apply that section's level-representability theorem at $3$.
Restoring this positive representative does not rescale the marking
or fix its unit multiplier. A rational quasi-isogeny carrying the
entire integral Tate lattice to itself is an actual isomorphism,
since any nontrivial finite kernel or cokernel is detected on some
prime-power torsion lattice. These constructions are inverse.

The comparison is algebraic. On each arithmetic component use the
natural integral lattice local system with its Hodge structures.
An arithmetic stabilizer preserving the lattice has positive
polarization multiplier and is a unit at every finite prime, so its
multiplier is $1$. The polarization is therefore integral and flat
on that component. The family equivalence
\cite[Theorem 14.8, p.~352]{milne-shimura2005} gives the abelian scheme
and its $\mathcal O$-action algebraically. The flat full marking is
algebraic as well: its graph is a union of connected components in
the finite etale isomorphism cover. The fine functor yields a
morphism to $Y_{\mathbb C}$. The preceding object classification
makes it bijective on geometric points, hence an isomorphism of
smooth complex curves. Indeed it has degree one in characteristic
zero, and the local discrete valuation rings agree. This uses every
component and commutes with forgetting level.

Finally the comparison commutes with every automorphism of
$\mathbb C/\mathbb Q$. The order action, polarization, torsion
subgroup quotients and markings are algebraic, and the adelic
comparison is transported through torsion, not through a supposed
Galois action on rational Betti homology. The type-$A/C$ canonical
model criterion in \cite[pp.~354--355]{milne-shimura2005} therefore
identifies $Y_{\mathbb Q}$ with the entire canonical $U(3)$ curve.
Uniqueness and the Hecke-map statement, loc. cit., 13.6--13.7, identify
its maximal-level forgetful map and canonical right normalizer
action. The CM and reciprocity results underlying that published
criterion remain its named inputs. The author's corrections to
8.14/8.17 and to proof intermediates in 5.22 and 11.2 are recorded
in the full-cover source packet; no corrected general implication
or Hausdorff argument is silently substituted here.

Thus $Y$ is the required entire principal-level-three scheme with
its universal surface. Its degree over the stack is $72$, while
over the generic coarse curve the signs $\{\pm1\}$ identify opposite
markings, giving degree $36$. No field extension or extra level
has entered this construction. At primes dividing $210$, the height
argument below continues to use the maximal stack and its Hodge line,
not an asserted smooth fine model at those primes.

Work first over $R_0=\mathbb Z[1/210]$. The effective level-change group
has order $36$, prime to every residue characteristic of $R_0$. Its generic
action is free. For a finite tame group acting on a smooth relative curve,
the fixed locus of an element is smooth locally along its fixed points;
an isolated special fixed point consequently lifts. Thus no new fixed
point occurs here, and $Y_{R_0}\to X_{R_0}$ is finite etale. The $W$ action
extends to the smooth stable model of $X$. Its stabilizers remain cyclic,
of order at most two. Distinct involution fixed loci cannot intersect;
each fixed locus is etale over the base. Consequently $X_{R_0}/W$ is a
smooth genus-zero curve with five disjoint branch sections of index two.
The three rational sections already chosen identify it with
$\mathbb P^1_{R_0}$. Their values are the five elements
of $\mathcal C$.

Set $K=\widehat{\mathcal O}^{\times}$, $U=U(3)$,
$Z=\mathbb A_f^\times$, and $N=N_{B(\mathbb A_f)^\times}
(\widehat{\mathcal O})$. Normalization preserves $3\widehat{\mathcal O}$,
so $U$ is normal in $N$. The canonical adelic action is defined over the
reflex field $\mathbb Q$ \cite[2.2.5 and 2.3.1, pp.~269--270]{Deligne1979}; this refers to the entire canonical model, not
an individually selected complex component. Every central finite idele
is a rational scalar times a unit idele. At three its scalar residue is
$1$ or $-1$, which a rational sign removes. Thus $Z\subset\mathbb Q^\times U$
acts trivially on $Y$.

At three the quaternion algebra is ramified. Its maximal order modulo
three has $81$ elements, radical of size $9$, and hence $72$ units.
It follows that
\[
 1\longrightarrow (\mathcal O/3\mathcal O)^\times/\{\pm1\}
 \longrightarrow\Gamma=N/(ZU)\longrightarrow W\longrightarrow1,
 \qquad |\Gamma|=36\cdot16=576.
\]
Above an unbranched point the actions of $W$ and of the effective
level-change kernel are free. The fiber therefore is a torsor of size
$576$ for this constant group. The Galois permutations commute with its
regular action, so belong to its regular centralizer, also of order
$576$. The normal splitting field $L$ consequently satisfies
$[L:\mathbb Q]\mid576$: the Galois group is a subgroup of that
regular centralizer. In particular $[L:\mathbb Q]\le576$.

All fiber surfaces and their three-torsion markings are defined over $L$.
For two such surfaces, a homomorphism annihilating the three-torsion
factors through $[3]$. Hence reduction embeds
\[
 \operatorname{Hom}(A,B)/3\operatorname{Hom}(A,B)
 \hookrightarrow\operatorname{Hom}(A[3],B[3]).
\]
The finite Galois image on the integral Hom lattice is trivial modulo
three. The principal congruence subgroup modulo three in an integral
general linear group is torsion-free, so this image is trivial. Thus
every geometric homomorphism between fiber surfaces is defined over $L$.

At a prime outside $210$, tame local coordinates for the quotient are
$w=z^2$ up to a unit at a branch section, and etale elsewhere. A base
point with even contact gives no ramification after an unramified
extension. The universal abelian scheme extends at the resulting
integral point of $Y_{R_0}$. Therefore
\[
 \operatorname{rad}(D_LN_{\rm bad})\mid B_*M(s).
\tag{R}\label{eq:h5-R}
\]
This assertion includes contact with infinity. It is not a statement
obtained merely from the numerical prime support of pairwise determinants.

For the geometric assertion, use the rational PEL description of the
same full-three curve. A finite-adelic normalizer move fixes the rational
Hodge object and replaces its integral lattice by a commensurable one.
One integer clears the finitely many local denominators; the finite
torsion quotient in the full-cover construction then gives an isogeny
$f:A\longrightarrow A_\sigma$ between the actual conjugate fiber surfaces.
This complex isogeny is geometric over $\overline{\mathbb Q}$. Indeed,
torsion points of both algebraic surfaces are algebraic, so $f$ and each
of its $\operatorname{Aut}(\mathbb C/\overline{\mathbb Q})$-conjugates
agree on all source torsion. Torsion is Zariski dense and the target is
separated, so these morphisms agree. The invariant graph descends by the
affine-ideal argument: on affine charts its unique reduced Gr\"obner
basis has invariant coefficients, hence coefficients in
$\overline{\mathbb Q}$.

In the non-CM case the full geometric endomorphism algebra is $B$.
More explicitly, the rank-one left-$B$ Hodge calculation gives
$B\otimes K_0^{\mathrm{op}}$, where $K_0=B\cap\mathbb R[J]$ is either
$\mathbb Q$ or an imaginary quadratic field. The latter is the CM
alternative, so $K_0=\mathbb Q$ here. The Hodge equivalence and the
geometric descent just used identify this with
$\operatorname{End}^0_{\overline{\mathbb Q}}(A)$; transport through $f$
gives the same conclusion for $A_\sigma$.

Let $F$ be the chosen abstract quadratic field, with specified embeddings
$i:F\hookrightarrow\operatorname{End}^0(A)$ and
$i_\sigma:F\hookrightarrow\operatorname{End}^0(A_\sigma)$. Put
$j(x)=f^{-1}i_\sigma(x)f$. Skolem--Noether, in the form
of \cite[Corollary 7.1.5, p.~97]{voight2021}, gives $b\in B^\times$
such that $j(x)=b^{-1}i(x)b$ for every $x\in F$. Consequently
$g=f b^{-1}$ satisfies $g i(x)=i_\sigma(x)g$. Choose a positive integer
$n$ clearing its denominators in
$\operatorname{Hom}(A,A_\sigma)\otimes\mathbb Q$. Then $h=ng$ is an
actual homomorphism with the same intertwining identity. Clearing the
denominators of its rational inverse shows that $h$ is an isogeny.
The simultaneous Hom argument above places $h$, both rational $F$-actions,
and their intertwining identity over the same $L$.

Thus each non-CM conjugate surface is of GL$_2$-type over $L$, and these
surfaces admit pairwise $F$-linear isogenies over $L$. This assertion
does not choose coherent isogenies or construct a compatible Galois
system. The separate scalar normalizer maps below supply the
$B$-compatible maps used in the paragraph ``The scalar lift with
controlled support''; the subsequent paragraph ``The compatible modular
representations'' constructs the representations and checks their
cross-prime traces. Those arguments, including their support conditions,
are separate inputs. The bounded descent to a GL$_2$-type variety over
$\mathbb Q$ is also a separate conclusion.

\paragraph{The metric comparison, with its exact additive budget}
Put $t=32s$, $(a_1,a_2,a_3,a_4)=(0,32,76832,-343)$ and
$P(t)=\prod_i(t-a_i)$. The sections
\[
 q_0=(dt)^2/P(t),\qquad q_1=tq_0
\]
are holomorphic bicanonical sections on $X$, with no common zero. At a
finite branch point use $t-a_i=z^2$; at infinity use $t=z^{-2}$.
At infinity $q_0$ vanishes to order two and $q_1$ does not vanish.

For each prime let $\delta_p=\max_{i<j}v_p(a_i-a_j)$, so
\[
 J=\prod_p p^{\delta_p}=2^{10}3^2 5^3 7^4=2765952000.
\]
Both $J^2q_i$ extend across the maximal regular stack. Indeed, at the
generic point of a vertical component semistability makes $p$ a
uniformizer, and the relative differential satisfies $v(du)\ge v(u)$.
Writing $r_i=v(t-a_i)$ gives $v(dt)\ge\max r_i$. If $v(t)<0$, all $r_i$
equal $v(t)$ and both sections have nonnegative valuation. Otherwise all
$r_i\ge0$ and at most one exceeds $\delta_p$. In either case
$2\max r_i-\sum r_i\ge-2\delta_p$. This proves the required vertical
integrality; regularity handles codimension two.
This extension takes place on the maximal stack itself. It uses no
integral coarse-model pullback from Yuan's Lemma 2.1 or the coarse
arithmetic-intersection formula of his Theorem 2.2.

For the archimedean estimate write $r_i=|t-a_i|$, $U_0=1+\sum_i r_i$ and
\[
 ds_\rho^2=\frac{U_0}{\prod_i r_i}|dt|^2,\qquad
 C_{\rm met}=16(12\cdot76832+9)^3=12540153541944538512.
\]
Its pullback is smooth and positive at every branch point and infinity.
The curvature formula, $\Delta U_0=\sum_i r_i^{-1}$ and
$|\nabla U_0|^2\le16$ give
\[
 -K_\rho\ge\frac{\sum_i\prod_{j\ne i}r_j}{2U_0^3}.
\]
When $|t|\le2\cdot76832+2$, at most one $r_i<1/2$, and
$U_0\le12\cdot76832+9$. The displayed bound is at least
$C_{\rm met}^{-1}$. Outside this disk, $|t|/2\le r_i\le3|t|/2$
and $U_0\le7|t|$, giving the stronger bound $1/1372$.
The maximum principle against the curvature $-1$ metric implies
$ds_\rho^2/C_{\rm met}\le ds_{\rm hyp}^2$. Yuan's Petersson norm is
$\|d\tau\|=2\operatorname{Im}\tau$. Consequently
\[
 \|q_0\|_{\rm Pet}\le4C_{\rm met}/U_0,\qquad
 \|q_1\|_{\rm Pet}\le4C_{\rm met}|t|/U_0\le4C_{\rm met}.
\]

Yuan's isometry \cite[Theorem 1.1, p.~229, and metric definitions, pp.~228--229]{Yuan2024} identifies the square of the Hodge bundle with
$210\omega_{\mathscr X}^{\otimes2}$ with the induced metric. The Hodge
height in his normalization equals the original stable Faltings height
plus $\log\pi$ in dimension two. Thus
\[
 h_{\bar\omega^{\otimes2}}(y)=2h_F(A_y)+2\log\pi+\log210.
\]
The sign of $\log210$ is positive: the sublattice $210\omega^{\otimes2}$
has height $h_{\bar\omega^{\otimes2}}-\log210$. Evaluate on an extension
where the point extends to the proper maximal stack; its universal
abelian scheme computes stable height. The integral sections and the
supremum estimate yield
\[
 h(s)\le2h_F(A_y)+2\log\pi+\log210+2\log J
                       +\log(4C_{\rm met})+\log32.
\]
The integer comparison
$210J^2(4C_{\rm met})32\cdot16<2^{160}$, together with $\pi<4$ and
$\log2<3/4$, bounds the additive term by $120<128$. Therefore
\[
 h(s)\le2\max(0,h_F(A_y))+128.
\tag{G}\label{eq:h5-G}
\]


% Level-based height argument; the stated source qualifications remain in force.
% Uses the existing baseline bibliography keys, plus GaudronRemond2023 and
% the published 2021 edition of Voight, chapter 43.
\paragraph{A level bound for the non-CM fibers}
\label{sec:h5-level-repair}

This argument allows the scalar character order and coefficient field to
vary. It does not assert bounded same-field descent to a rational
GL$_2$-type variety. Let $s$ be a rational point outside the five contacts,
let $L/\mathbb Q$ be the normal fiber field constructed above, and let
$A/L$ be a non-CM fiber surface. Thus $[L:\mathbb Q]\mid576$, all
geometric homomorphisms between conjugate fibers are defined over $L$,
and $\operatorname{End}^0(A_{\overline{\mathbb Q}})=B$, the quaternion
division algebra of discriminant $210$. Put
\[
 R=B_*M(s),\qquad T=\{p:p\mid R\}.
\]
The geometric argument gives that $L$ is unramified outside $T$ and
$A$ has good reduction outside the primes above $210$. All conjugate
QM structures lie in the same normalizer orbit.

\paragraph{Good reduction after full torsion.}
First a QM surface with the indicated division action has potentially
good reduction at every finite place. Indeed, pass to an extension with
semistable reduction, using \cite[Theorem 3.4, p.~406]{silverbergzarhin}.
Every endomorphism extends uniquely to the N\'eron model by its mapping
property. On the identity component of the geometric special fiber it
preserves the maximal torus $\mathcal T$: this torus is the maximal
connected affine subgroup, and every homomorphism from a torus to an
abelian variety is zero. If its rank $t$ is positive, restriction gives
a unital homomorphism
\[
 B\longrightarrow\operatorname{End}^0(\mathcal T)=M_t(\mathbb Q),
 \qquad 1\le t\le2.
\]
The kernel is a two-sided ideal of the simple algebra $B$, and the map
is unital, so it is injective. For $t=1$ dimensions prohibit this;
for $t=2$ it would be an isomorphism from a division algebra to
$M_2(\mathbb Q)$, also impossible. Hence $t=0$. The semistable identity
component is abelian, which is good reduction. These N\'eron-model and
reduction conventions are those of \cite[\S3, pp.~405--406]{silverbergzarhin}.

Set $E=L(A[11])$. We next justify that $E/\mathbb Q$ is Galois.
For each $\sigma$, the normalizer relation supplies a compatible
isogeny ${}^{\sigma}A\to A$ of degree prime to $11$. Here is the
prime-to-$11$ detail. In the analytic description
$A_z=\mathbb C^2/(\mathcal O(z,1)^t)$ of the published 2021 edition
of Voight \cite[equations (43.1.2)--(43.1.3)]{voight2021}, a positive normalizer element
$b=\left(\begin{smallmatrix}a&b'\\c&d\end{smallmatrix}\right)$
gives the scalar compatible map $(cz+d)I:A_{bz}\to A_z$.
When $\mathcal O b\subset\mathcal O$, its degree is
$[\mathcal O:\mathcal O b]=\operatorname{Nrd}(b)^2$.
At $11$ the algebra is split and the maximal-order normalizer is
$\mathbb Q_{11}^*\operatorname{GL}_2(\mathbb Z_{11})$.
Rescale $b$ rationally to make it an $11$-adic unit, then clear the
remaining denominators by an integer prime to $11$. This produces
the claimed isogeny. Changes of level marking use isomorphisms.
All these geometric homomorphisms are $L$-defined by construction.
Consequently the $11$-torsion representations of $A$ and its conjugates
over $L$ have the same kernel. Thus every conjugate of $E$ equals $E$.

The faithful action on $A[11]$ now gives
\[
 [E:\mathbb Q]\le576|\operatorname{GL}_4(\mathbb F_{11})|
 <2^{10}11^{16}<2^{66}.
\tag{LR1}
\]
For residue characteristic different from $11$, full rational
$11$-torsion gives semistable reduction over $E$ by
\cite[Theorem 3.5, p.~406]{silverbergzarhin}. The preceding torus
argument makes this good reduction. In particular it applies at
each of $2,3,5,7$.

\paragraph{The scalar lift with controlled support.}
Choose the isogenies $\mu_\sigma:{}^\sigma A\to A$ over $L$ from
the normalizer construction above, retaining their full $B$-linearity.
This can be checked on every component in the rational PEL description
of Section~\ref{sec:h5-full-three-construction}. For $n$ normalizing
$\widehat{\mathcal O}$, the points $(h,a)$ and $(h,an)$ have the same
rational left $B$-module $V$ and Hodge structure,
with commensurable lattices $\Lambda$ and $\Lambda_n$. The identity on
$V$ commutes with $B$. At $11$ the normalizer has the form
$n_{11}=q u$, with $q\in\mathbb Q_{11}^{\times}$ and
$u\in\mathcal O_{11}^{\times}$, so $\Lambda_{n,11}=q\Lambda_{11}$.
Multiplication by $11^{-v_{11}(q)}$, followed by an integer prime to
$11$ clearing the remaining denominators, gives $r\in\mathbb Q^*$
with $r\Lambda_n\subset\Lambda$ and equality at $11$.
The Hodge-lattice equivalence turns $r\,\mathrm{id}_V$ into a
$B$-linear isogeny from the translated
surface to the original one, of degree prime to $11$. The accepted
integral-lattice identification retains the actual $\mathcal O$-labels;
changes of the full-three marking are $\mathcal O$-linear isomorphisms.
The two-surface Hom argument places these maps over the same $L$.
In the analytic chart of \cite[equations (43.1.2)--(43.1.3)]{voight2021},
the action on every $A_z$ is $i_z=i_\infty$. The displayed map
$\lambda I:A_{bz}\to A_z$, with $\lambda=cz+d$, therefore satisfies
$\lambda I\,i_{bz}(x)=i_z(x)\lambda I$ for $x\in B$;
$\lambda(bz,1)^t=b(z,1)^t$ verifies its lattice direction.
Take $\mu_1=1$ and inflate this finite choice to $G_{\mathbb Q}$.
Their discrepancies $c(\sigma,\tau)$ commute with $B$, so lie in
$\mathbb Q^*$. Put $a_\sigma=(\deg\mu_\sigma)^{-1/4}>0$ and
$\nu_\sigma=a_\sigma\mu_\sigma$, with trivial Galois action on the
added algebraic coefficients. The rational degree used above satisfies
$D(cu)=c^4D(u)$ for $c\in\mathbb Q^*$. Its multiplicativity and the
positive fourth-root convention give normalized discrepancy
$c(\sigma,\tau)/|c(\sigma,\tau)|$. Thus the discrepancies of the
$\nu_\sigma$ form a sign cocycle
$\varepsilon\in Z^2(G_{\mathbb Q},\{\pm1\})$, inflated from
$\operatorname{Gal}(L/\mathbb Q)$.

The sign class is unramified, hence zero, locally outside $T$.
At each odd prime where its Brauer invariant is $1/2$, choose an odd
Dirichlet character modulo $p$; at $2$ use the odd character modulo
$4$ if required. Their product $\chi$ is tame at every odd prime,
has local conductor at most $4$ at $2$, and has order
\[
 m\le2\prod_{\substack{p\in T\\p\ne2}}(p-1)\le2R.
\]
For constant coefficient groups, the boundary in
$1\to\{\pm1\}\to\mu_\infty\xrightarrow{2}\mu_\infty\to1$
is the obstruction to a character square root. Local reciprocity
identifies this obstruction with the value at $-1$: the cyclic
torsion factor has precisely this obstruction, while the valuation
and principal-unit factors admit character square roots.
Thus the boundary of $\chi$ equals $[\varepsilon]$ at every finite
place. Brauer reciprocity and the global-to-local injection give
equality globally. A pointwise square root, corrected by a sign
cochain, gives
\[
 b:G_{\mathbb Q}\longrightarrow\mu_{2m},\qquad
 b^2=\chi,\qquad\delta b=\varepsilon.
\tag{LR2}
\]
Only this \emph{sign} cochain is asserted to take root-of-unity values;
the full scalar splitting also involves the $a_\sigma$.

On $G_L$, $b$ is a conjugation-invariant character. Outside $T$ its
inertia characters are quadratic. Only finitely many are nontrivial;
twist by the quadratic character of the product of those odd primes
to kill them. At each odd prime there is a unique nontrivial quadratic
inertia character. This changes no square or coboundary in (LR2),
and introduces possible ramification only at $2\in T$.
The kernel field $L'$ of $b|G_L$ is normal over $\mathbb Q$ and
\[
 [L':\mathbb Q]\le2m[L:\mathbb Q]\le2304R.
\tag{LR3}
\]

\paragraph{The compatible modular representations.}
Choose a number field containing all normalization scalars, the values
of $b$, and a splitting field of $B$. For each coefficient prime $\ell$
and an embedding into $\overline{\mathbb Q}_\ell$, set
\[
 \rho_\ell(\sigma)=b(\sigma)^{-1}a_\sigma\mu_\sigma\circ\sigma.
\]
The cocycle identity makes this a representation commuting with $B$;
after splitting $B$, matrix units give $\rho_\ell=r_\ell\oplus r_\ell$
with $\dim r_\ell=2$. On $G_{L'}$, $\rho_\ell$ is the Tate action of
$A$. Every geometric endomorphism is defined over $L\subset L'$, and
$\operatorname{End}^0(A_{\overline{\mathbb Q}})=B$. Faltings'
semisimplicity and endomorphism theorems therefore apply to the actual
Tate action on $G_{L'}$. After splitting $B$, its commutant is $M_2$;
the original matrix-unit factor has scalar commutant and is semisimple.
The same calculation after every finite coefficient extension proves
that this original $r_\ell$ is absolutely irreducible for every $\ell$.
The representation is unramified outside $T\cup\{\ell\}$,
by (LR2) and the reduction and field properties.
Apply Lemma~\ref{lem:qm-modularity-descent} with
$t_\sigma=a_\sigma/b(\sigma)$. It gives a primitive weight-two
newform $f$ whose Tate realization is exactly the original $r_5$.
Write $N$ for the level of this final form, after the finite character
in that lemma has been absorbed. The finite character $\delta$ absorbed in the lemma is not identified
with the controlled character $\chi=b^2$; the original $L'$, $b$ and
$r_\ell$ remain those in (LR2)--(LR3) and the estimates below.
The comparison at every coefficient prime uses the original family,
not a replacement of the controlled field or its ramification set.
Choose an auxiliary finite Galois field $K_{\rm aux}$ defining the
finitely many isogenies $\mu_\sigma$, their twists, and the scalar
choices. Choose a positive integer $m$ clearing all quasi-isogeny
denominators. After excluding a finite set $S$ of rational primes,
including those ramified in $K_{\rm aux}$,
these data spread to abelian schemes and homomorphisms $m\mu_\sigma$.
The auxiliary field $K_{\rm aux}$ and comparison set $S$ do not replace
$L'$ or $T$.

Let $p\notin S$ and choose a place of $K_{\rm aux}$ above $p$.
An arithmetic Frobenius lift $\sigma_p$ acts on the residue field by
$x\mapsto x^p$. Reduction identifies ${}^{\sigma_p}A$ with the
$p$-Frobenius twist of the reduced surface $A_v$, so
\[
 e_p=\overline{m\mu_{\sigma_p}}\circ\operatorname{Fr}_{p,A_v}
 \in\operatorname{End}(A_v).
\]
Here $\operatorname{Fr}_{p,A_v}$ is relative $p$-Frobenius. Replacing it
by Frobenius for the full residue-field cardinality would give a
different power of $\sigma_p$. The rational degree polynomial of an
endomorphism is its characteristic polynomial on every prime-to-$p$
rational Tate module, by
\cite[Proposition 12.9 and the preceding definition, p.~125]{milne1986-abelian}.
Thus $q_p=\operatorname{Tr}(e_p)/m\in\mathbb Q$ is independent of the
Tate prime. This division by $m$ retains the original quasi-isogeny.

For $\ell\ne p$, reduction on $\ell$-power torsion is compatible with
these spread homomorphisms. Indeed, the torsion is finite \'etale and
becomes a disjoint union of sections over the strict henselization.
It therefore identifies $\mu_{\sigma_p}\circ\sigma_p$ on rational
Tate modules with $e_p/m$. Since $\rho_\ell=r_\ell\oplus r_\ell$, this
gives
\[
 \operatorname{Tr}r_\ell(\sigma_p)
   =\frac12\,\iota_\ell\!\left(a_{\sigma_p}/b(\sigma_p)\right)q_p.
\]
The factor $1/2$ comes from the two equal summands of the full Tate
representation.

Fix a coefficient prime $\ell$. Choose a number field $F_{\rm coeff}$
containing the original scalar and splitting field and the coefficients
of the primitive form $f$. Choose an embedding into
$\overline{\mathbb Q}_5$ extending the fixed one used to identify the
original $r_5$ with the realization of $f$, and an embedding into
$\overline{\mathbb Q}_\ell$ extending the fixed embedding for the
original $r_\ell$. At $\ell$, work over a common finite coefficient
extension containing both coefficient images. Before either embedding,
\[
 \tau_p=\tfrac12\bigl(a_{\sigma_p}/b(\sigma_p)\bigr)q_p
 \in F_{\rm coeff}.
\]
Outside the finite set also containing $\ell,5$ and the level of $f$,
the identification at five gives
$\iota_5(\tau_p)=\iota_5(a_p(f))$. Injectivity gives equality in
$F_{\rm coeff}$, and the embedding at $\ell$ gives equality of the
traces of the original $r_\ell$ and the corresponding classical
Tate realization $U_{f,\ell}$.

Use the established absolute irreducibility of the original $r_\ell$.
Chebotarev and Brauer--Nesbitt identify $r_\ell$ with
$U_{f,\ell}^{\rm ss}$ over the common coefficient extension.
No independent semisimplicity assertion about $U_{f,\ell}$ is needed:
its semisimplification is one simple two-dimensional factor, so
$U_{f,\ell}$ has composition length one and is itself irreducible.
Hence $r_\ell\simeq U_{f,\ell}$ as representations of the full
$G_{\mathbb Q}$. The finite comparison exclusions neither restrict
this isomorphism to a subgroup nor enlarge $T$.

For each rational prime $p$ one may choose $\ell\ne p$ and restrict
this isomorphism to its decomposition group. In particular, the same
form $f$ may still be compared with $r_3$ at the prime $5$ and with
$r_5$ at the other places. The comparison asserts no new numerical
conductor bound. The modular level is the conductor of these realizations by the
published result recalled below. Li's Fricke relation and Shimura's
Jacobian construction, as in the baseline modularity argument, also
place $r_5$ in the Tate module of $J_1(N)$.

\paragraph{The primitive level equals the local conductor.}
Let $f\in S_2(\Gamma_0(N),\eta)$ be the primitive newform obtained
after absorbing the finite character twist, and let $U_{f,\ell}$
be its two-dimensional Tate realization for the chosen coefficient
embedding. Here $\eta$ is the nebentypus of $f$; it is not identified
with the controlled character $\chi=b^2$ in (LR2).
Carayol's published conductor result for normalized primitive
newforms gives
\[
 a_p(U_{f,\ell})=v_p(N),\qquad p\ne\ell.
\tag{CL1}
\]
Indeed, \cite[\S2.1, p.~36]{carayol1987} permits arbitrary
nebentypus and coefficient number fields, and weight two.
Section~2.4 on p.~37 states the conductor assertion also at primes
dividing $N$; the conclusion following the full Theorem~4 in
\S4.1 on p.~42 proves that assertion. We use the full theorem,
not its weaker version with an extraordinary-cuspidal exclusion.
Thus the assertion includes dyadic primes and nonsquarefree level.
The cohomological realization in that source is dual to the Tate
realization; duality preserves the Artin conductor. The unramified
norm twist in its Hecke-correspondence convention likewise changes
no conductor exponent.

For each rational prime $p$, choose an auxiliary prime $\ell\ne p$.
The preceding compatible-family comparison identifies the actual
$r_\ell$ constructed from $A$, $a_\sigma$, and $b$ with
$U_{f,\ell}$ for this same final newform and coefficient embedding.
Restriction to $G_{\mathbb Q_p}$ preserves this isomorphism, hence
preserves inertia-invariant dimensions and the Swan conductor.
Consequently
\[
 a_p(r_\ell)=a_p(U_{f,\ell})=v_p(N),\qquad p\ne\ell.
\tag{CL2}
\]
The conductor includes the tame term and Swan term, or equivalently
the full Weil--Deligne conductor with its monodromy contribution.
No semisimplification on inertia is taken. Applying (CL2) with an
auxiliary prime different from $p$ controls every prime divisor of
$N$; no local compatibility theorem at $p=\ell$ is required.

\paragraph{The conductor bound.}
Let $F=E\mathbb Q(\chi)$, where $\mathbb Q(\chi)$ means the kernel
field of the character, not its value field. On $G_F$, $b$ is a
conjugation-invariant quadratic character. Its kernel gives a Galois
extension $F'/\mathbb Q$ of degree at most two over $F$.
For $p=2,3,5,7$, let $e_E,e_{F'}$ denote local ramification indices
over $\mathbb Q_p$. From (LR1), $v_p(e_E)\le65$.
At odd $p$ both added extensions are tame, whereas at $2$ each
adds ramification index at most two. Consequently
\[
 v_p(e_{F'})\le65\ (p=3,5,7),\qquad v_2(e_{F'})\le67.
\tag{LR4}
\]
The possibly unbounded tame ramification index causes no problem.
For a finite Galois extension of $\mathbb Q_p$ with ramification
index $e$ and different exponent $\delta$, an Eisenstein polynomial
over the maximal unramified subfield gives
\[
 \delta\le e-1+ev_p(e).
\]
Indeed the derivative terms have distinct valuations modulo $e$,
and its leading derivative term has the stated valuation. In upper
numbering the different formula is
\[
 \delta/e=1-1/e+\int_0^\infty(1-|G^u|^{-1})\,du.
\]
Before the last positive break the integrand is at least $1/2$.
Thus (LR4) bounds this break for $F'/\mathbb Q_p$ by $136$.
These ramification and conductor conventions are those of
\cite[\S\S4.1--4.3, pp.~155--159]{serrelocal}.
Over $F'$ the scalar correction is trivial and $A$ has good reduction.
Above that break, therefore, the two-dimensional representation is
trivial on upper wild groups. Its Swan conductor is at most $272$,
and its tame conductor is at most two. Hence
\[
 a_p(r)\le274<300\qquad(p=2,3,5,7).
\]
For $p>7$, $L/\mathbb Q$ is tame because its degree divides $576$,
$A/L$ has good reduction, and $b$ is trivial on wild inertia:
on that pro-$p$ group its square is tame and its remaining image has
order at most two. Thus $a_p(r)\le2$ there; it is zero outside $T$.
Using the coefficient-prime comparison above proves
\[
 N\le210^{300}R^2.
\tag{LR5}
\]

\paragraph{The original constant is preserved.}
Apply the accompanying modular-Jacobian transfer calculation with
$C=210^{300}$ and (LR3)--(LR5). It gives
\[
 h_F(A)<2^{24}210^{3600}R^{24}
 <2^{28824}R^{24}<4608^{165888}R^{24}.
\tag{LR6}
\]
The last comparison uses $210<2^8$, $4608>2^{12}$, and
$28824<1990656=12\cdot165888$. Together with the separately proved
CM bound and geometric height comparison, (LR6) yields exactly
\[
 h(s)\le
 \bigl(2(4608^{165888}+6^{100})B_*^{24}+128\bigr)M(s)^{24}.
\]
No contact, coefficient constant, or radical exponent has been changed.
% Height consequence of the preceding modularity and level argument.

% Actual level-based repair: no uniform coefficient-field bound is used.
% The modular representation and level estimate are proved in the companion
% height-level-repair fragment. This file proves its height consequence.
\paragraph{Transfer from the level, allowing unbounded coefficient degree.}
Let $A/L$ be the non-CM quaternionic surface. Write $R\geq1$ for the
squarefree product of the prescribed rational bad primes, including $2$,
and assume $[L:\mathbb Q]\leq576$. The modularity construction provides
a weight-two newform of level $N$ and a field $L'$ over which the Tate
modules of $A$ and $J=J_1(N)$ share an irreducible constituent. In the
notation of the sign-splitting construction, the auxiliary Dirichlet
character has order
\[
 m\leq 2\prod_{\substack{p\mid R\\p\ne2}}(p-1)\leq2R.
\]
The sign-splitting cochain has values in $\mu_{2m}$. If one instead
uses the full scalar splitting map, it also includes the square root
of the degree map; that factor equals one on $G_L$. Thus in either
notation the restriction to $G_L$ is a character valued in
$\mu_{2m}$, and its kernel field satisfies
\[
 d'=[L':\mathbb Q]\leq2m[L:\mathbb Q]\leq2304R.
\tag{LT1}
\]
The quadratic twist used to remove extra ramification does not increase
this bound, since $2$ divides $2m$. There is no bound here on $m$
independent of $R$.

Faltings' Hom theorem, followed by scalar extension in the Hom space,
gives a nonzero $L'$-homomorphism $A\to J$. Absolute simplicity of $A$
and Poincar\'e reducibility therefore give an $L'$-isogeny
$A\times Z\to J$ for an abelian variety $Z/L'$. Put $g=\dim J$;
then $g\geq2$. The modular genus estimate
\cite[(5.9), p.~13]{vonKaenel2021} gives $g\leq N^2/24$.
The standard Belyi map has degree at most $N^2$ and
\cite[Theorem~1.1.1 and \S2.3]{javanpeykar} gives
\[
 h_F(J)\leq13000000N^{12}.
\tag{LT2}
\]
All heights in this argument use the original stable Faltings metric.

Suppose the companion local argument has proved $N\leq C R^2$,
where $C\geq1$. Then
\[
 g\leq C^2R^4/24,\qquad
 h_F(J)\leq13000000C^{12}R^{24}.
\]
Set $l=\log C$ and $r=\log R$. The published
Gaudron--R\'emond Theorem~1.9(1) \cite[pp.~7--8]{GaudronRemond2023}, combined with half-logarithmic isogeny-height variation,
gives
\[
 h_F(A)+h_F(Z)\leq h_F(J)+V,
 \quad
 V=2g^2\left(8g^2\log(7g)+\log d'
       +\log\max(1,\log d',h_F(J))\right).
\tag{LT3}
\]
Here the source variety for the isogeny estimate is $J_{L'}$.
The three logarithms in (LT3) satisfy respectively
\[
 \log(7g)\leq2l+4r,\quad
 \log d'<8+r,\quad
 \log\max(1,\log d',h_F(J))<17+12l+24r.
\]
For the last assertion, the bound in (LT2) also dominates $1$ and
$\log d'$; use $8+r\leq9R^{24}$.
It follows that
\[
 V\leq
 \frac{C^8R^{16}}{24^4}(32l+64r)
 +\frac{C^4R^8}{288}(25+12l+25r).
\tag{LT4}
\]
Using $l\leq C^4/4$, $r\leq R^8/8$, and also
$r\leq R^{16}/16$ in the second summand gives
\[
 \frac{V}{C^{12}R^{24}}
 \leq\frac{16}{24^4}+\frac{25+3+25/16}{288}<1.
\tag{LT5}
\]
The Bost lower bound in \cite[(2.2), p.~5]{vonKaenel2021}
gives $h_F(Z)\geq-\tfrac12\dim Z\log(2\pi^2)>-\tfrac32g$.
Consequently
\[
 h_F(A)<(13000000+1+1)C^{12}R^{24}
       <2^{24}C^{12}R^{24}.
\tag{LT6}
\]
This calculation needs no uniform bound on the degree of the newform's
coefficient field and no assertion that all its geometric factors are
isogenous to $A$.

For the companion conductor estimate one may take $C=210^{300}$.
Since $210<2^8$ and $4608>2^{12}$,
\[
 2^{24}C^{12}=2^{24}210^{3600}
 <2^{28824}<2^{1990656}<4608^{165888}.
\tag{LT7}
\]
Thus the non-CM height inequality is strictly within the original
coefficient budget, once the stated modularity and level estimate have
been established. In the five-contact application $R=B_*M(s)$, so
the radical exponent is exactly $24$, as required.

% Original 15671 architecture: independent CM branch, version 1.
% This lemma is conditional on the separately audited field/ramification lemma.
\paragraph{The CM height estimate over the same field}
Let $A/L$ be an abelian surface of CM type, let $B$ be a quaternion
division algebra over $\mathbb Q$ embedded in
$\operatorname{End}^0_{\overline L}(A)$, and suppose that all geometric
endomorphisms of $A$ are defined over $L$. Assume
\[
 d=[L:\mathbb Q]\le576,\qquad
 \operatorname{rad}(|D_L|)\le R,\qquad R\ge1.
\]
All heights below are stable Faltings heights in Faltings's original
normalization. We prove
\begin{equation}\label{eq:original-cm-bound}
 h_{\rm F}(A)<854R<6^{100}R^{24}.
\end{equation}
No assertion about the existence or degree of $L$ is included in this lemma.

First, $A_{\overline L}$ is isogenous to $E^2$ for a CM elliptic curve
$E$. Here is a dimension argument which makes the alternative cases
explicit. If $A$ were geometrically simple, its division endomorphism
algebra would act faithfully on the four-dimensional rational homology.
Its dimension therefore divides four. As it contains $B$, it would equal
$B$; but $B$ has no commutative semisimple subalgebra of dimension four,
contrary to CM type. If the two elliptic factors were not isogenous,
their endomorphism algebra would be a product of commutative fields and
could not contain $B$. Finally, if $E$ had no CM, the endomorphism algebra
would be $M_2(\mathbb Q)$, and the embedding of the four-dimensional
division algebra $B$ would be an impossible isomorphism.

The elliptic factors and the isogeny may be taken over $L$. Indeed,
the rational idempotents in $\operatorname{End}^0(A)$ are defined over
$L$, and the image of an integral multiple of such an idempotent is an
elliptic subvariety over $L$. Endomorphisms and homomorphisms of the
resulting factors extend to rational endomorphisms of $A$ by the
inclusions and rational projections; these are fixed by
$\operatorname{Gal}(\overline L/L)$. Thus the factor isogenies, and all
endomorphisms of $E$, are defined over $L$.

Put $F=\operatorname{End}^0(E)$, an imaginary quadratic field, and
$\mathcal O=\operatorname{End}(E)$. Differentiation gives an injective
unital map $F\longrightarrow\operatorname{End}_L(\operatorname{Lie}E)=L$:
it is nonzero, and every nonzero elliptic endomorphism in characteristic
zero is separable. Consequently $F\subset L$.

We next pass to the maximal order without extending $L$. Choose
$f\ge1$ with $f\mathcal O_F\subset\mathcal O$, choose an integral
basis $\alpha_1,\alpha_2$ of $\mathcal O_F$, and form the finite
subgroup scheme
\[
 K=\sum_{i=1}^2(f\alpha_i)(E[f])\subset E.
\]
Each $f\alpha_i$ is defined over $L$, so $K$ and the quotient
$E'=E/K$ are defined over $L$. Under a complex uniformization
$E(\mathbb C)=\mathbb C/\Lambda$, this subgroup is
$\mathcal O_F\Lambda/\Lambda$; thus $E'(\mathbb C)$ has lattice
$\mathcal O_F\Lambda$ and endomorphism ring $\mathcal O_F$.
The ring cannot be larger, since the rational endomorphism algebra
remains $F$. This is also a direct special case of the same-field
maximal-order theorem of R\'emond, but the construction just given is
the proof used here. It follows that $A'=(E')^2$ is $L$-isogenous to $A$.

Write $\Delta_F$ for the fundamental discriminant of $F$. The
discriminant tower formula implies that every rational prime dividing
$\Delta_F$ divides $D_L$. A quadratic fundamental discriminant has
odd prime exponents one and exponent at two at most three. Therefore
\[
 |\Delta_F|\le4\operatorname{rad}(|D_L|)\le4R.
\]
The multiplier four includes the possible factor $2^3/2$.

For completeness we give an elementary analytic bound on $h_{\rm F}(E')$.
CM elliptic curves have potentially good reduction. After a finite
extension with everywhere good reduction, the elliptic discriminant
formula of Silverman, reproduced as equation (6.2) in
\cite{cmheight}, gives
\[
 h_{\rm F}(E')=\frac1{[L':\mathbb Q]}\sum_{\sigma:L'\hookrightarrow\mathbb C}
 \left(-\log(2\pi)-\log|\eta(\tau_\sigma)|^2
                     -\tfrac12\log\operatorname{Im}\tau_\sigma\right).
\]
Here $L'$ is used only to evaluate the stable height; it is not the
field used in the isogeny estimate below. The discriminant term is zero
because reduction is good. We use the individual-embedding formula
(6.2), not the class-average assertion of Proposition 6.1, so no
hypothesis on Weierstrass coefficients lying in $\mathbb Q(j)$ is needed.

For each embedding choose a reduced quadratic parameter of discriminant
$\Delta_F$. If $y=\operatorname{Im}\tau$, then
\[
 \frac{\sqrt3}{2}\le y=\frac{\sqrt{|\Delta_F|}}{2a}
                   \le\frac{\sqrt{|\Delta_F|}}2,
 \qquad a\in\mathbb Z_{\ge1}.
\]
Put $r=|e^{2\pi i\tau}|<1/16$. The product defining $\eta$ gives
\[
 -\log|\eta(\tau)|^2
 \le\frac{\pi y}{6}+2\sum_{n\ge1}\frac{r^n}{1-r^n}
 \le\frac{\pi y}{6}+\frac{2r}{(1-r)^2}
 <\frac{\pi\sqrt{|\Delta_F|}}{12}+\frac{32}{225}.
\]
Also $-\tfrac12\log y\le\tfrac14\log(4/3)<1/12$.
Discarding the negative term $-\log(2\pi)$ and using $\pi<4$ shows
\[
 h_{\rm F}(E')<\sqrt{|\Delta_F|}+1\le2\sqrt R+1\le3R.
\]
The height is additive on products, hence $h_{\rm F}(A')<6R$.

Apply \cite[Theorem 1.9(1)]{GaudronRemond2023} over $L$, with source
$A'$ and target $A$ and with no extension of the base field.
For a surface, its quantity $\Xi(A')$ is
\[
 \Xi(A')=\left(14^{32}d\max(1,\log d,h_{\rm F}(A'))\right)^8.
\]
The general assertion supplies an isogeny $A'\to A$ of degree at most
$\Xi(A')^2$; no identification of CM type with that paper's
separate ``MM'' condition is made or needed. The isogeny-height
inequality then gives
\[
 h_{\rm F}(A)\le6R+
 8\log\left(14^{32}\,576\max(1,\log576,6R)\right).
\]
One can see the last inequality directly from invariant differentials:
after semistable base change, pullback on determinant differentials has
integral cokernel, while its archimedean squared norm is multiplied by
the isogeny degree. Thus the height difference is half the logarithm
of that degree minus a nonnegative normalized cokernel length.

Finally, $\log14<3$, $\log576<7$, and
$\max(1,\log576,6R)\le7R$. Since $\log7<2$ and $\log R\le R$,
\[
 h_{\rm F}(A)<6R+8(96+7+2+\log R)
             \le854R.
\]
As $6^4=1296>854$ and $R\ge1$, this proves
\eqref{eq:original-cm-bound}. In particular the original
$6^{100}R^{24}$ allowance absorbs this CM estimate with the field
degree still exactly bounded by 576.

\paragraph{The unconditional numerical certificate and scalar detector}
The specified integer $H$ satisfies
\[
 \log_2H<4783227.
\]
Indeed the exact sieve enumerates $7499$ primes in $S_*$, with largest
prime $76163$, and their product has bit length $109445$. Hence
$B_*<2^{109445}$. With $A_0=4608^{165888}$, we have
$6^{100}<A_0$ and $128<A_0$, whence $H<8A_0B_*^{24}$. Since $4608<2^{13}$,
\[
 \log_2H<3+165888\cdot13+24\cdot109445=4783227.
\]
This is a check of the number H, not a proof that H is a height constant.

For completeness, the scalar product used by the original polynomial is
\[
 r(\beta)=\beta(\beta-1)(\beta-2401)(32\beta+343)(2\beta-1).
\]
At a prime $p>76176$, if $v_p(\beta)<0$, its valuation is
$5v_p(\beta)$ and so detects odd contact with infinity. If $\beta$ is
$p$-integral and has odd contact with a finite member of $\mathcal C$,
only that factor has positive valuation; every other factor is a unit.
The extra root $1/2$ introduces only exceptional collision primes
(notably factors of $4801$ and $718$). Thus every prime of $M(\beta)$ has
odd valuation in $r(\beta)$. If a separate comparison proves
$M(\beta)^{48}\mid s\ne0$, the five-contact estimate just proved gives
$h(\beta)\le H\sqrt{|s|}$.

Combining the non-CM and CM estimates with \eqref{eq:h5-R} and
\eqref{eq:h5-G} gives
\[
 h(s)\le 2(4608^{165888}+6^{100})(B_*M(s))^{24}+128
       \le H_*M(s)^{24},
\]
because $M(s)\ge1$. The numerical certificate above proves the
remaining assertion about $H_*$.
\end{proof}


Define $\Psi(b)$ by seven rational witnesses $t,U,e,p,q,s,h$ and four
equations, with $B=17b$:
\begin{align}
 U(1+bt^2)&=b(1-bt^2),&
 p^2&=B(1-e)(1+3e),\label{eq:o8-pole}\\
 h&=e(B-U^2),&
 Bq^2&=(s^2-Bh)(s^2+3Bh).\notag
\end{align}

\begin{lemma}[Exact seven-variable pole statement]\label{hyp:o8-pole}
If $b\ne0$ and $\Psi(b)$, then $v_p(b)$ is not a negative odd integer
for every $p\notin\{2,3,17\}$. Conversely, $\Psi(b)$ holds whenever
$b>0$, the denominator of $b$ in lowest terms is a square, and $b$ is
a square in each of $\Q_2,\Q_3,\Q_{17}$.
\end{lemma}
\begin{proof}
These are Lemmas~\ref{lem:original-seven-parity-soundness} and
\ref{lem:original-seven-parity-completeness}. Their proof gives the
simultaneous two-twist construction explicitly, and applies the existing
Theorem~\ref{hyp:RC} and Lemma~\ref{lem:corank}. It introduces no new
simultaneous-realization assumption.
\end{proof}

The lemma includes the degenerate values of the seven witnesses
permitted by \eqref{eq:o8-pole}. No inverse condition on $e$, $h$, $s$,
or $U$ may be inserted into this particular polynomial. For $b\ne0$,
the first equation itself implies $1+bt^2\ne0$. The predicate is soluble
at $b=0$; the norm argument below accounts for that case where necessary.

% Additive proof for the original 15671-variable polynomial. No circuit changes.
\subsubsection{The seven-witness parity predicate and its finite correction}
\label{subsec:original-seven-parity}
Put $B=17b$. The predicate used here is exactly the existence of
$t,U,e,p,q,s,h\in\mathbb Q$ satisfying
\begin{align}
 U(1+bt^2)&=b(1-bt^2),\label{eq:seven-param}\\
 p^2&=B(1-e)(1+3e),\label{eq:seven-first}\\
 h&=e(B-U^2),\label{eq:seven-h}\\
 Bq^2&=(s^2-Bh)(s^2+3Bh).\label{eq:seven-second}
\end{align}
Denote it by $\Psi(b)$. No nonzero conditions are added to this formula.

\begin{lemma}[The exact one-sided assertion]
\label{lem:original-seven-parity-soundness}
If $b\ne0$ and $\Psi(b)$, then $v_\ell(b)$ is not a negative odd
integer for any prime $\ell\notin\{2,3,17\}$.
\end{lemma}
\begin{proof}
Suppose $k=v_\ell(b)<0$ is odd. Equation~\eqref{eq:seven-param}
forces $1+bt^2\ne0$: if it vanished, its right side would be $2b\ne0$.
If $t=0$, then $U=b$. Otherwise $v_\ell(bt^2)$ is odd and nonzero,
so $(1-bt^2)/(1+bt^2)$ is a unit. Thus $v_\ell(U)=k$ in all cases,
and $v_\ell(B-U^2)=2k$.

If $e=0$, equation~\eqref{eq:seven-first} gives $p^2=B$, impossible.
Write $j=v_\ell(e)$. If $j<0$, the right side of
\eqref{eq:seven-first} has valuation $k+2j$, which is odd.
If $j>0$, it has valuation $k$, again odd. Consequently $j=0$.
This includes the possibilities $p=0$, $e=1$, and $e=-1/3$.
Now $h\ne0$ and $v_\ell(Bh)=3k$ is odd. If $s\ne0$, its square
has even valuation, different from $3k$, so the two factors on the
right of~\eqref{eq:seven-second} have the same finite valuation.
Their product has even valuation and is nonzero. The same conclusion
holds when $s=0$, since the product is $-3(Bh)^2$.
The left side is either zero or has odd valuation, a contradiction.
\end{proof}

The omitted guards therefore cause no loss of the asserted soundness.
At $b=0$, the equations require $U=p=h=s=0$ and permit arbitrary
$t,e,q$. For $b\ne0$, the six-witness source formula is recovered
exactly by $e_{\rm old}=Be$, $p_{\rm old}=Bp$, $q_{\rm old}=Bq$,
with $s,t,U$ unchanged. Its first curve residual is $B^2$ times
\eqref{eq:seven-first}'s residual and its second is $B$ times
\eqref{eq:seven-second}'s residual after~\eqref{eq:seven-h}.
Division by $B$ reverses this transformation.

\begin{lemma}[Simultaneous realization from the existing restricted converse]
\label{lem:original-seven-simultaneous-realization}
Use Theorem~\ref{hyp:RC}, the restricted converse $\mathsf{RC}$ for
$E_l:y^2=x(x-l)(x+3l)$ with $l>0$ squarefree and $l\equiv11\pmod{12}$.
Let $\beta,\theta>0$ be squarefree, with disjoint supports avoiding
$2,3$, such that $\beta\equiv17\pmod{24}$,
$\theta\equiv1\pmod{24}$, and each is a local square at every prime
dividing the other. Then there is a positive squarefree $n$, prime
to $6\theta$, for which both $E_n$ and $E_{n\theta}$ have rational
points with nonzero abscissa in $\beta\mathbb Q^{\times2}$.
The case $\theta=1$ is allowed.
\end{lemma}
\begin{proof}
Here are the descent steps, so that no separate realization theorem
is being assumed. Write $B_0=\operatorname{supp}(\beta)$,
$J=\operatorname{supp}(\theta)$, and $d=|B_0|\ge1$.
For a support prime $q_i$, put $s_i=[-1/q_i]$, $\lambda_i=[3/q_i]$,
and define over $\mathbb F_2$
\[
 (Kv)_i=\sum_{j\ne i}[q_j/q_i](v_j+v_i).
\]
Add $d-1$ prime slots $I$ of type $(s,\lambda)=(0,0)$ and one slot
$q_0$ of type $(1,1)$. Let the $B_0$-by-$I$ symbol matrix have columns
$e_i+e_d$ for $1\le i<d$. Set all other symbols incident with $I$
to zero, prescribe $[q_0/q_b]$ as the corresponding row sum, and
obtain reverse symbols by reciprocity. The $q_0$--$J$ symbols remain
free. Mutual squareness gives zero row sums for both directed
$B_0$--$J$ blocks; hence $K\mathbf1_{B_0}=0$ on both supports.

Choose $a$ with ones at one fixed member of $B_0$ and at $q_0$.
The matrix $K^\sharp=K+\lambda a^t$ kills $\mathbf1$ on both sides.
After deleting $q_0$, its smaller-support matrix has the form
\[
 M_0=\begin{pmatrix}D&C\\C^t&0\end{pmatrix},
 \qquad \mathbf1^tD\mathbf1=1.
\]
The columns of $C$ span the even-sum hyperplane. If $M_0(x,y)^t=0$,
then $x=c\mathbf1$; multiplication of the first equation by
$\mathbf1^t$ gives $c=0$, and injectivity of $C$ gives $y=0$.
This also proves the case $d=1$, where $M_0=(1)$.
Adding $J$ preserves $M_0$. Each free $q_0$--$J$ bit toggles just
one new diagonal entry. The enlarged determinant is multilinear in
these bits, with full-product coefficient $\det M_0=1$, so some
Boolean choice makes it nonzero. CRT and Dirichlet's theorem realize
the prescribed symbols by distinct fresh primes, with types imposed
modulo $12$. Thus $n=\prod_{B_0\cup I\cup\{q_0\}}q$ and $n\theta$
are squarefree and congruent to $11$ modulo $12$.

The projection $T=1-\mathbf1_{B_0}a^t$ satisfies $K=K^\sharp T$.
Nonsingularity of the deleted matrix gives
$\ker K^\sharp=\langle\mathbf1\rangle$, and therefore
\[
 \ker K=\langle\mathbf1,\mathbf1_{B_0}\rangle,
 \qquad \lambda\notin\operatorname{im}K
\]
on both supports. For the latter assertion, $Kx=\lambda$ would
make $Tx-\mathbf1_{B_0}$ constant, contradicted by applying $a^t$.

For clarity, the local descent needed here is the manuscript's full
two-descent in coordinates $(x+3l,x-l)$. The two torsion classes are
$(l,1)$ and $(3,-l)$. At an odd good prime, the local image consists
of the four pairs with even valuations; at a prime dividing $l$,
the two torsion classes span it by their independent valuation
vectors. At $3$ it is the set of all first-coordinate classes with
square second coordinate. At $2$ both coordinate valuations are
even: the only problematic case $x-l=2u$, $u$ odd, would give
$(y/2)^2\equiv l\pmod8$, impossible for $l\equiv3\pmod4$.
When both differences are divisible by $4$, their difference $4l$
and the cubic equation give even valuations for both. At infinity
the first coordinate is positive, and the positive-second-coordinate
subgroup has index two. Its possible representatives consequently are
$(3^\delta\prod q_i^{u_i},\prod q_i^{v_i})$ with necessary equations
\[
 Ku+\operatorname{diag}(\lambda)v=\delta\lambda,
 \quad (K+\operatorname{diag}(s))v=0,
 \quad (\lambda+s)^tv=0.
\]
Reciprocity gives $K+K^t=ss^t+\operatorname{diag}(s)$.
Hence $\ker(K+\operatorname{diag}(s))=\ker K^t\cap s^\perp$
is a line. The functional $\lambda^t$ is nonzero on this line,
because $\lambda\notin\operatorname{im}K$ and
$\lambda^t\mathbf1=0$, whereas $s^t\mathbf1=1$.
The displayed equations force $v=0$, $\delta=0$, and $u\in\ker K$.
There are at most eight Selmer classes.
The class $(\beta,1)$ is actually local everywhere: it is the identity
at $2$ and infinity, the first torsion class at $3$, satisfies the
torsion comparison at support primes by $K\mathbf1_{B_0}=0$, and
is unramified at the other odd primes. It is independent of the two
torsion classes since $B_0$ is nonempty and $q_0\notin B_0$.
Thus each two-Selmer group has dimension exactly three.

Lemma~\ref{lem:corank} gives full two-primary Selmer
corank one. Explicitly, the Kummer sequence gives
$r+\dim\Sha[2]=1$; writing
$\Sha[2^\infty]=(\mathbb Q_2/\mathbb Z_2)^t\oplus F$, the
alternating Cassels--Tate pairing on the finite quotient makes
$\dim F[2]$ even, so $F=0$ and $r+t=1$.
Apply $\mathsf{RC}$ to $n$ and $n\theta$, and the manuscript's
published forward analytic-rank-one theorem. Both curves have rank
one and finite $\Sha$, hence $\Sha[2]=0$. Kummer exactness now
represents $(\beta,1)$ by actual rational points. These are not
two-torsion, so their abscissas are nonzero of class $\beta$.
When $\theta=1$ the same point can be used twice.
\end{proof}

\begin{lemma}[The completeness scope actually used]
\label{lem:original-seven-parity-completeness}
Assume only the existing $\mathsf{RC}$ and the published descent and
forward-rank facts used above. If $b>0$ has square reduced denominator
and is a square in $\mathbb Q_2,\mathbb Q_3,\mathbb Q_{17}$, then
$\Psi(b)$ holds.
\end{lemma}
\begin{proof}
Choose $t\in\mathbb Q$ close to a solution of $bt^2=1$ at infinity
and at $2,3,17$. In~\eqref{eq:seven-param} this makes $U$ arbitrarily
small at those places. We can and do require
$V=B(B-U^2)>0$ and $V$ locally square at all three primes.
Explicit sufficient conditions are $U^2<B$ in $\mathbb R$,
$v_2(U^2/B)\ge3$, and $v_3(U^2/B),v_{17}(U^2/B)\ge1$.
The denominator tends locally to $2\ne0$ and is positive globally.

Let $\beta,\theta$ be the positive squarefree representatives of
$B,V$. The number $B=17b$ is square at $2$, a nonsquare unit class
at $3$, and has odd valuation at $17$. Thus
$\beta\equiv17\pmod{24}$ and $\beta>1$;
$\theta\equiv1\pmod{24}$ and $17\nmid\theta$.
At any prime $\ell\mid\beta$ outside $2,3,17$, the valuation of
$b$ is odd and therefore positive, since its reduced denominator is
square. The same quotient calculation as in soundness gives
$v_\ell(U)=v_\ell(b)$; hence $1-U^2/B$ is a square unit.
This proves that $V$, and hence $\theta$, is square at every prime
of $\beta$, including $17$ by construction. The supports are disjoint.
Conversely, if $\ell\mid\theta$, then $\ell\notin\{2,3,17\}$
and $v_\ell(B)$ is even. Odd valuation of $B(B-U^2)$ forces
$v_\ell(B)=2v_\ell(U)$ and strict cancellation, so
$B/U^2\in1+\ell\mathbb Z_\ell$ is square. This proves the other
mutual-square condition. If $U=0$, then $\theta=1$ and this last
argument is unnecessary.

Apply the preceding lemma. Write $V=\theta c^2$ and rescale the
second point by $(x,y)\mapsto(c^2x,c^3y)$ to put it on $E_{nV}$.
If the first point has abscissa $Ba^2$, divide both curve parameters
by $a^2$ and the point coordinates by $a^2,a^3$. The resulting
points have abscissas $B,Bs^2$ with $s\ne0$, on $E_{e_0}$ and
$E_{e_0V}$ respectively. Write their ordinates as $y_1,y_2$.
Set
\[
 e=e_0/B,\qquad p=y_1/B,\qquad
 q=y_2/(B^2s),\qquad h=e(B-U^2).
\]
Direct substitution verifies all four equations. Every division
made in choosing these witnesses is by a nonzero number; none is
an additional equation or guard in the original polynomial.
\end{proof}

\paragraph{Finite correction and the unguarded norm root.}
At $2$ use valuation parity followed by unit coordinates $[-1],[5]$;
at $3$ and $17$ use valuation parity followed by the nonsquare-unit bit.
The columns for $2,3,5,7,11,13,17$ are exactly
\[
\begin{pmatrix}
1&0&0&0&0&0&0\\0&1&0&1&1&0&0\\0&1&1&0&1&1&0\\
0&1&0&0&0&0&0\\1&0&1&0&1&0&1\\0&0&0&0&0&0&1\\
0&1&1&1&1&0&0
\end{pmatrix}.
\]
The accompanying exact computation verifies rank seven and all 128
distinct subset signatures. Rational solutions of $e^2-e=0$ are
exactly $0,1$. Therefore the original eight bits and two products
\[
 A=(1+e_2)(1+2e_3)(1+12e_{13})(1+16e_{17}),\quad
 D=(1+4e_5)(1+6e_7)(1+10e_{11})(1-2\eta)
\]
give exactly $AD=(-1)^\eta\prod p^{e_p}$, always nonzero.
For nonzero integral $r$, choose $\delta=\operatorname{sign}(r)$
and the unique subset product $\mu$ having the same local signature
as $\delta r$. Taking $AD=\delta\mu$ makes $ADr$ positive and
locally square at $2,3,17$.

Choose $J_0$ sufficiently deep at all places over primes dividing
$\mu r$, then $S$ in a sufficiently divisible subgroup satisfying
the anchor. Since $f=z^2$ times a formal unit, the coefficients of
$f(J_0),f(S),z(S)/z(J_0)$ can attain any prescribed finite lower
valuation bounds. This remains true at ramified and bad primes by
working in sufficiently deep formal subgroups. Thus at $\ell\mid r$
all normalized leaves are integral, and at $\ell\mid\mu$ with
$\ell\nmid r$ all nonconstant leaves can be made divisible by $\ell$.
The norm tree $N_r(a,b)=a^2+rb^2$ has its constant leaf $1$ on
the all-left path; its root $W$ is a unit at these primes.
In the actual application $r=\beta(\beta-1)(\beta-2401)
(32\beta+343)(2\beta-1)$ with $\beta=X^2+2$ and $X\in\mathbb Z$.
It is nonzero and $|r|>1$ since $X^2\ne2399$ and $\beta\ge2$.
Consequently $W\ne0$, and $b=ADrW^2$ is positive, locally square
at $2,3,17$, and has square reduced denominator. Completeness applies.

For soundness no global guard $W\ne0$ is needed. At any prime where
$v_\ell(r)$ is odd, the two nonzero summands of $a^2+rb^2$ have
different valuation parity. Such a norm vanishes only if both inputs
vanish. Recursing to the constant leaf $1$ proves $W\ne0$ at every
prime relevant to an odd-contact argument. Outside the fixed
exceptional set, $AD$ is a unit and
$v_\ell(ADrW^2)=v_\ell(r)+2v_\ell(W)$.
The exact original norm argument therefore applies unchanged. If
$r$ has no relevant odd valuation, there is no local conclusion to
prove. No guard or witness has been inserted, and the original
polynomial and coefficient bound are unchanged.

\subsection{Elliptic charts and the anchor}

Use the normalized curve and point
\[
 K=\Q(\sqrt2),\quad d=3-\sqrt2,\quad D=d^2=11-6\sqrt2,
 \quad E:y^2=x^3-Dx,
\]
\[
 P_*=(57/8-7\sqrt2/2,-165/8+417\sqrt2/32),
 \qquad G=2E(K)=\Z P_*.
\]
The identification of $G$ follows from Proposition~\ref{prop:generator}
by the unit change of variables $\lambda=3+2\sqrt2$.
Write $f=1/x$ and $z=-x/y$. A nonzero point of $G$ is represented by
eight field variables $u,v,w,z,f,h,x,y$ satisfying
\begin{equation}\label{eq:o8-chart}
 f=u^2,\quad v^2+df=1,\quad w^2+v^2=2,\quad h=vw,
 \quad zh+u=0,\quad xf=1,\quad yz+x=0.
\end{equation}
Each field equality means its two rational coefficient equalities.
\begin{lemma}[The sparse point chart]\label{lem:o8-chart}
The equations \eqref{eq:o8-chart} parametrize precisely the nonzero points
of $G$, and their reciprocal functions are defined.
\end{lemma}
\begin{proof}
It gives $x=u^{-2}$, $y=vw/u^3$ and the three square root differences
$x=(1/u)^2$, $x-d=(v/u)^2$, $x+d=(w/u)^2$.
Here is a direct verification that these squares give precisely the
required doubles. Write $e_1=0,e_2=d,e_3=-d$ and choose
$r_i^2=x-e_i$ with $r_1r_2r_3=-y$. Put
\[
 B=r_1r_2+r_1r_3+r_2r_3,\qquad
 X'=x+B,\qquad
 Y'=-(r_1+r_2)(r_1+r_3)(r_2+r_3).
\]
Since $X'-e_i=(r_i+r_j)(r_i+r_k)$, the point $(X',Y')$ lies on
$E$. None of the pair sums is zero, because the $e_i$ are distinct.
Its tangent slope is $-(r_1+r_2+r_3)$, and the duplication equations
give $[2](X',Y')=(x,y)$.
Conversely, if $(x,y)=[2](a,b)$ is nonzero, then $b\ne0$ and direct
duplication gives
\[
 x-e_i=\left(
 \frac{(a-e_i)^2-(e_i-e_j)(e_i-e_k)}{2b}\right)^2.
\]
Choose the signs of these square roots to have product $y$, and recover
$u,v,w,h,z,f$ from \eqref{eq:o8-chart}. The group $G$ is torsion-free, so its nonzero points have
$x,y\ne0$. Thus the sparse chart and the membership equations in
Section~\ref{sec:model} have the same nonzero image.
The inverse relation $xf=1$ excludes zero denominators. In particular
$u,h,z,y$ are nonzero.
\end{proof}

Subtraction uses the equations
\[
 \ell x-y=j,\quad \ell a+b=j,\quad
 x'=\ell^2-x-a,\quad \ell x'+y'=j,\quad
 f'x'=1,\quad z'y'+x'=0
\]
for $(x,y)-(a,b)$, with an inverse for the rational norm of $x-a$.
Doubling uses the usual tangent formula with its nonzero denominators.

The baseline duplication estimates give
\[
 4h(x(Q))-\log392\le h(x(2Q))
       \le4h(x(Q))+\log676.
\]
Since $h(x(P_*))=\log41$, passage to the canonical height gives
$2\widehat h(P_*)\ge\log41-\frac13\log392>1$, and hence
\begin{equation}\label{eq:o8-indexheight}
 h(x(nP_*))\ge n^2-6\qquad(n\ne0).
\end{equation}
The deliberately weaker constant six suffices throughout this section.

Let $\mathcal T(a)$ be the seven-variable formula
\begin{equation}\label{eq:o8-T}
 u_0^2-2u_1^2-3u_2^2+6u_3^2=1,\qquad
 (a-2u_0)^2-8v_1^2-12v_2^2+24v_3^2=4.
\end{equation}
Then $\mathcal T(\Q)=\Z_{(2)}\cap\Z_{(3)}$. Indeed the quaternion
algebra $(2,3)$ is ramified exactly at $2,3$ and split at infinity.
Koenigsmann's trace-sum theorem, Proposition~6 of
\cite{koenigsmann-defining}, identifies the sum of two norm-one trace
sets with the intersection of these two local rings. Eliminating the
scalar coordinate of the second norm-one element gives exactly
\eqref{eq:o8-T}, with no translation of $a$.

Introduce twenty-five rational variables
\begin{equation}\label{eq:o8-chain}
 c_0=3,\qquad c_{i+1}=c_i^2\quad(0\le i<24),\qquad C=c_{24}.
\end{equation}
Thus $C=3^{2^{24}}$, but this large integer is not a coefficient in
$F_8$. For $S=sP_*\ne O$ and $f(S)=f_0+f_1\sqrt2$, impose
$Ca_i=f_i$ and $\mathcal T(a_i)$, for $i=0,1$.
At the prime above $3$, $P_*$ reduces to the nonidentity point
$(\sqrt2,0)$ of order two, and $v_3(z(2P_*))=1$.
The formal group therefore gives
$v_3(f(sP_*))=2+2v_3(s)$ whenever this valuation is positive.
The anchor forces
\begin{equation}\label{eq:o8-anchor-bound}
 |s|\ge3^{2^{23}-1}>(20H_*)^2.
\end{equation}
For the last inequality use Theorem~\ref{hyp:o8-height},
$\log_2 3>3/2$, and
$3(2^{23}-1)=25165821>4(5+4783227)=19132928$.
The allowed indices contain a nonzero subgroup: sufficiently divisible
indices make $f(sP_*)/C$ integral at both $2$ and $3$. The anchor can
therefore be imposed together with any finite list of formal divisibilities.

% Original 15,671-witness anchor. No polynomial or constant is changed.
\subsubsection{The literal seven-witness semilocal predicate}\label{subsec:o8-anchor-proof}

For $t\in\mathbb Q$, let $\mathcal T(t)$ mean that there are
$u_0,u_1,u_2,u_3,v_1,v_2,v_3\in\mathbb Q$ with
\[
 u_0^2-2u_1^2-3u_2^2+6u_3^2=1,
 \qquad
 (t-2u_0)^2-8v_1^2-12v_2^2+24v_3^2=4.
\]
Then
\[
 \mathcal T(t)\quad\Longleftrightarrow\quad
 t\in\mathbb Z_{(2)}\cap\mathbb Z_{(3)}.
\]
Indeed, use the quaternion convention $i^2=2$, $j^2=3$ and $ij=-ji$.
Its reduced norm is $x_0^2-2x_1^2-3x_2^2+6x_3^2$ and its reduced trace
is $2x_0$. Put $v_0=(t-2u_0)/2$. The two displayed equations say
exactly that $t$ is a sum of two reduced traces of norm-one elements.
No eighth variable is needed, since $v_0$ is recovered by this rational
expression with fixed nonzero denominator~$2$.

In the notation of \cite[Definition~4 and Proposition~6]{koenigsmann-defining},
this is $T_{2,3}=S_{2,3}+S_{2,3}$. That proposition identifies
$T_{a,b}$ with the intersection of the local integer rings at the
ramified places, with a separate real interval only when infinity is
ramified. Here the Hilbert symbols are $-1$ at $2$ and $3$, and $1$
elsewhere. At $3$ this follows from $(2/3)=-1$; at every other odd
prime both parameters are units. At $2$, writing $2=2^1\cdot1$ and
$3=2^0\cdot3$, the usual exponent is
$0+1\cdot(3^2-1)/8+0=1$. The algebra splits over $\mathbb R$ because
both parameters are positive. Thus the ramification set is precisely
$\{2,3\}$, which proves the claimed equivalence in both directions.
This uses Koenigsmann's sum of two trace sets, not a predicate with
an additional finite offset.

\subsubsection{The original anchor and its exact index subgroup}\label{subsec:o8-anchor-subgroup}

Write $K=\mathbb Q(\rho)$, $\rho^2=2$, and use the literal normalized
curve and generator
\[
 E:y^2=x^3-(11-6\rho)x,
 \qquad
 P=\left(\frac{57}{8}-\frac72\rho,
          -\frac{165}{8}+\frac{417}{32}\rho\right).
\]
As in the original construction, use Proposition~\ref{prop:generator}, transported by the unit change
of variables, to obtain $2E(K)=\mathbb ZP$. For a nonzero $S=nP$, put
$f(S)=1/x(S)=f_0+f_1\rho$. Retain the exact constant chain
\[
 c_0=3,\qquad c_{j+1}=c_j^2\quad(0\le j<24),\qquad C=c_{24}.
\]
It forces $C=3^{2^{24}}$, with $25$ shared rational variables and
$25$ quadratic-or-linear equations. The anchor introduces two
variables $a_i$, imposes $Ca_i=f_i$, and imposes $\mathcal T(a_i)$
for each $i\in\{0,1\}$. It has $16$ private witnesses and six
quadratic equations, exactly as in the saved module.

Set $N=2^{24}$ and $m=N/2=2^{23}$. We claim that, for nonzero $n$,
\[
 \text{the anchor holds on }nP
 \quad\Longleftrightarrow\quad
 n\in 2\cdot3^{m-1}\mathbb Z.
\]
First consider the place above~$3$. It is unramified quadratic, with
residue field $\mathbb F_9$ and normalized valuation $v(3)=1$.
For rational $a,b$ one has
$v(a+b\rho)=\min(v_3(a),v_3(b))$: after dividing by the smaller
power of $3$, linear independence of $1,\bar\rho$ over $\mathbb F_3$
precludes cancellation.
The curve has good reduction there. Direct calculation gives
\[
 \bar P=(\bar\rho,0),\qquad
 z(2P)=
 \frac{3396356159088-4135204317756\rho}{3287811625151},
 \qquad z=-x/y.
\]
The reduction of $P$ has order two. The denominator in this expression
is a $3$-unit and both numerator coefficients have valuation one.
Consequently $nP$ is in the formal kernel exactly when $n$ is even,
and $v(z(2P))=1$.

For completeness, the valuation of every such multiple can be computed
without an unproved uniformity assertion. The integral formal group
in the parameter $z$ has inverse $-z$, so its multiplication series
are odd. Their linear coefficient is the corresponding integer.
Thus, for a $3$-unit integer $u$, $[u]z=uz+O(z^3)$ preserves positive
valuation, whereas $[3]z=3z+O(z^3)$ increases valuation by exactly one
at depth at least one. The higher terms have valuation at least
$3v(z)>v(z)+1$. Apply this repeatedly to $[n/2](2P)$, allowing
negative integers by negation. It gives
\[
 v(z(nP))=1+v_3(n),\qquad
 v(f(nP))=2+2v_3(n)\quad(2\mid n).
\]
The second equality follows from
$f=z^2(1-(11-6\rho)f^2)$ on the formal branch.
The integral formal-group construction and its linear multiplication
coefficient are given in \cite[IV\S1, Example~IV.2.2.3 and
Proposition~IV.2.3(a)]{silverman}. No minimal-model assertion at
the bad prime~$2$ is used below.

At the place above~$2$, normalize $v(2)=1$, so $v(\rho)=1/2$.
The two summands of a rational pair have valuations of opposite
integer parity after multiplication by~$2$, and hence cannot cancel.
The displayed coordinates give
\[
 v(x(P))=-3,\qquad v(y(P))=-9/2,\qquad v(z(P))=3/2.
\]
Thus $P$ lies on the formal branch at infinity, even though reduction
at~$2$ is bad. Every integral multiple lies in this formal subgroup.
One can see the integral branch directly by setting $w=-1/y$:
the equation becomes $w=z^3-(11-6\rho)zw^2$, which has a unique
solution $w\in z^3\mathcal O[[z]]$; its group law is the integral
formal group just cited. In particular $f(nP)$ is integral at~$2$.
Since
\[
 v(a+b\rho)=\min(v_2(a),v_2(b)+1/2),
\]
its rational coefficients are both $2$-integral. As $C$ is a $2$-unit,
the two anchor conditions at~$2$ impose no further restriction on $n$.

At~$3$, the two anchor conditions are equivalent to
$v(f(nP))\ge N$. Positive valuation of $f$ implies negative valuation
of $x$, hence membership in the formal kernel, so $n$ is even.
The preceding exact formula makes the remaining condition
$v_3(n)\ge m-1$. Conversely these two conditions give both required
coordinate divisibilities at~$3$, and the condition at~$2$ already
holds. The proved characterization of $\mathcal T$ supplies its
rational witnesses, proving the subgroup claim.

In particular, anchor indices contain a nonzero finite-index subgroup
and have positive representatives. They can be combined with any
finite collection of requirements that indices be divisible by fixed
integers: take a common multiple with $2\cdot3^{m-1}$. More generally,
any finitely many formal-neighborhood conditions whose preimages are
finite-index subgroups can be imposed in the same way. This statement
does not claim compatibility with arbitrary nonzero residue conditions.

The old and new reciprocal coordinates are equivalent for this anchor.
With $\lambda=3+2\rho$, one has $f_{\rm old}=\lambda^2f$ and
\[
 \begin{pmatrix}f_{{\rm old},0}\\f_{{\rm old},1}\end{pmatrix}
 =\begin{pmatrix}17&24\\12&17\end{pmatrix}
 \begin{pmatrix}f_0\\f_1\end{pmatrix}.
\]
This matrix has determinant one and integral inverse, so it preserves
coordinatewise membership in
$C(\mathbb Z_{(2)}\cap\mathbb Z_{(3)})^2$ in both directions.

\subsubsection{The unchanged numerical cutoff}\label{subsec:o8-anchor-cutoff}
In this calculation $H=H_*$ is the constant in Theorem~\ref{hyp:o8-height}.

The anchor gives $|n|\ge3^{2^{23}-1}$. Conditional only on the separately
audited original height bound $\log_2H<4783227$, it follows that
$|n|>(20H)^2$. Indeed $\log_2 3>3/2$ and $\log_2 20<5$, and
\[
 3(2^{23}-1)=25165821
 >19132928=4(5+4783227).
\]
Thus the original $24$ squarings suffice whenever that exact height
input is established. This calculation does not replace the required
five-contact height theorem by a six-contact theorem or by a larger
constant. No part of this proof changes the saved polynomial.

\subsection{The exact norm tests}

For a nonzero rational contact $r$ and rational tests $X_i$ of positive
orders $L_i$, apply the binary operation
$N_r(a,b)=a^2+rb^2$ to the ordered leaves
$1,X_1/r^{L_1},\ldots,X_k/r^{L_k}$ in the following fixed order.
First pair $1$ with $X_1/r$, where $L_1=1$. Group the remaining
tests by increasing order, retaining their original order within each
group. Pair successive tests within each group and carry any last
unpaired test. Then pair successive resulting nodes at each level and
carry any last node until only a root $W$ remains. The polynomial
equations for the bottom pairs are
\[
 rw-r-X_1^2=0,\qquad
 r^{2L}w-X_i^2-rX_j^2=0;
\]
an unpaired normalized leaf satisfies $r^Lw-X_i=0$.
Higher nodes satisfy $w-a^2-rb^2=0$. Powers of $r$ are computed
by binary squaring. Add an inverse for $r$, a variable for $W^2$,
and the predicate $\Psi(rW^2)$.

At a prime with $R=v_p(r)$ odd, the two summands in every norm node
have distinct valuation parity. There is no cancellation. The constant
leaf is on the all-left path, so $v_p(W)\le0$ if $R<0$.
Consequently $v_p(rW^2)>0$ first implies $R>0$.
It also gives $v_p(W)>-R/2$. Descending any norm node, a bound
$v_p(N_r(a,b))>-cR$ with $c<1$ gives
$v_p(a)>-cR/2$ and $v_p(b)>-(c+1)R/2$.
Every normalized leaf therefore has valuation greater than $-R$.
Thus $v_p(X_i)>(L_i-1)R\ge L_i-1$, or $v_p(X_i)\ge L_i$.
The root cannot vanish at such a prime. Indeed $N_r(a,b)=0$ forces
$a=b=0$ over $\Q_p$, and descent would make its constant leaf zero.
Thus Lemma~\ref{hyp:o8-pole} applies even though the polynomial
has no separately imposed inverse for $W$.

For completeness at a positive prime $r=p$, tests of orders $L_i$
can be made so deep that every nonconstant normalized leaf is
$p$-integral and the root $W$ is a $p$-adic unit. Then $rW^2$ has
square reduced denominator. If $r$ is a local square at $2,3,17$,
the completeness half of Lemma~\ref{hyp:o8-pole} applies.

\subsection{Forty-one contacts and one common comparison point}

Set
\[
 R=47027=31\cdot37\cdot41,\quad W_0=25878772920,
 \quad c=18484837801,\quad V_0=32W_0.
\]
For nine rational targets $a_i$, put $A_i=Ra_i+3$ and
$w_i=A_iv-u$. Use the forty-one nonzero contacts
\begin{equation}\label{eq:o8-contacts}
 v/V_0,\qquad u-\delta v\ (\delta\in\mathcal D),\qquad
 -(w_i-\delta v)\ (\delta\in\mathcal D,\ 1\le i\le9).
\end{equation}
For integer $a_i$, the prime-pattern argument of
Section~\ref{sec:mult} applies on
$v=32W_0(c+W_0y)$, $u=c+W_0x$ and gives positive ordinary prime
values for every distinct contact, all congruent to $c$ modulo $W_0$.
Repeated targets use the same prime values. These are local squares
at $2,3,17$. The precise application uses at most forty-one distinct
affine forms and their successor exclusions. At $31,37,41$ there are
at most eight slopes since every $A_i$ reduces to $3$; at $p\ge43$
there are at most forty. Choosing two slopes outside that set proves
local admissibility. The small-prime conditions are built into $c$.
The full uniform fiber sieve and same-prime order argument are given
in Subsection~\ref{subsec:o8-pattern-proof} below. No conditions relating orders at
different selected primes are needed.

An odd normalized contact of $u/v$ or $w_i/v$ outside $S_*$ forces
some raw contact in \eqref{eq:o8-contacts} to have odd valuation.
To see this, let $m$ be the minimum of the two coordinate valuations.
For a finite contact of odd depth, if $m$ is even its distinguished
linear form has odd valuation; if $m$ is odd another finite contact
form has valuation $m$. At an odd contact with infinity the same
argument uses the denominator if $m$ is even and a finite form if
$m$ is odd. All differences of distinct finite contact points are units
outside $S_*$. The signs and $V_0$ do not affect parity there.

For nonzero points $Q_i=n_iP_*$, $1\le i\le4$, take $Q_0=P_*$,
one common $J=jP_*\ne O$, and four rational proxies $b_i$.
At each raw contact introduce $X_0=hP_*$ and $X_i=m_iP_*$ and
the subtraction charts $Y_i=X_i-4Q_i$, $0\le i\le4$.
The points $4Q_i$ are computed by two doublings. Put
\[
 L(Z)=\sum_{j=0}^{15}
       \frac{(-1)^j\binom{2j}{j}D^j}{4j+1}Z^{4j+1}.
\]
This polynomial has degree sixty-one; the first omitted logarithm term
has degree sixty-five. Test order one for
\[
 f(X_i)\ (0\le i\le4),\quad f(J^\sigma),\quad
 f(Y_i^\sigma)\ (0\le i\le4),
\]
and order sixty-four for
\[
 L(z(X_i))/L(z(X_0))-b_i\ (1\le i\le4),\qquad
 z(Y_i^\sigma)/z(J^\sigma)\ (0\le i\le4).
\]
The inverses in these formulas are imposed explicitly. For each field
test $a+b\sqrt2$ of order $L_i$, use the rational tests
$(\theta^2-2)b$ and $a+\theta b$ of that order, together with
$\theta^2-2$ of order one. There is one common rational $\theta$
for all forty-one comparisons. This gives twenty-three rational tests
of order one and eighteen of order sixty-four.

At an odd raw contact the norm tests select a place above $p$, and
the conjugate parameter ratios imply
$h\equiv4$ and $m_i\equiv4n_i\pmod {p^{64}}$.
In particular $h$ is a unit. The logarithm expansion and its degree
sixty-five remainder give
\begin{equation}\label{eq:o8-proxies}
 b_i\in\Z_p,\qquad b_i\equiv n_i\pmod {p^{64}}.
\end{equation}
This uses the same formal-group estimates as Section~\ref{sec:mult}.

For completeness choose a consistent orientation at each distinct
selected prime. Let $d_{1,p},d_{2,p}$ be the two reduction orders.
The ordinary-prime conditions give $\gcd(d_{1,p},d_{2,p})=4$ and
$p\nmid d_{1,p}d_{2,p}$. Choose
$j=\operatorname{lcm}_p d_{2,p}$. For each occurrence at $p$ choose
\[
 h\equiv0\pmod {d_{1,p}},\qquad
 h\equiv4\pmod {d_{2,p}},\qquad
 h\equiv4\pmod {p^{64+v_p(j)}}.
\]
These congruences are compatible. Avoid $h=0,4,-4$, and set
$X_i=hn_iP_*$. They give all the required formal conditions and
parameter ratios relative to the single point $J$. Increase the last
precision if necessary and choose $\theta$ last by the Chinese
remainder theorem. Repeated occurrences of a prime use the same
orientation. Every subtraction denominator is nonzero, and the
formal-logarithm unit calculation at $p$ proves $L(z(X_0))\ne0$.

% Independent arithmetic justification for the unchanged original 15671-witness circuit.
\subsubsection{The original forty-one-prime pattern}\label{subsec:o8-pattern-proof}
Let
\[
 \mathcal D=\{0,1,2401,-343/32\},\quad R=47027=31\cdot37\cdot41,
 \quad W=25878772920=4\prod_{q\le29}q,
\]
\[
 c=18484837801,\qquad V_0=32W=828120733440.
\]
For integers $a_1,\ldots,a_k$, $k\le9$, put $A_i=Ra_i+3$,
$w_i=A_iv-u$. There are $5+4k$ displayed positions
\begin{equation}\label{eq:old41-list}
 v/V_0,\qquad u-\delta v,\qquad
 u-(A_i-\delta)v\quad(\delta\in\mathcal D,\ 1\le i\le k).
\end{equation}
For every lower bound these positions can simultaneously be positive
primes beyond that bound, congruent to $c$ modulo $W$, with the
following property. For the curve
$E:y^2=x^3-(3-\sqrt2)^2x$, both reductions above each prime are good,
and their cardinalities $d_1,d_2$ satisfy
\[
 \gcd(d_1,d_2)=4,\qquad p\nmid d_1d_2.
\]
This is a same-prime statement and asserts no cross-prime gcd condition.

\begin{proof}
Use $u=c+Wx$, $v=V_0(c+Wy)$. Put
\[
 \mathcal S=\mathcal D\cup\bigcup_i(A_i-\mathcal D),\quad
 K_s=V_0s\in W\mathbb Z,
 \quad\Psi_0=c+Wy,\quad\Psi_s=c+Wx-K_s(c+Wy).
\]
Remove duplicate targets. The remaining slopes are all distinct.
Indeed, the nonzero integral differences in $\mathcal D-\mathcal D$
are $\pm1,\pm2400,\pm2401$, none divisible by $31$.
The integral sums in $\mathcal D+\mathcal D$ are
$0,1,2,2401,2402,4802$, with residues $0,1,2,14,15,28$ modulo $31$,
none equal to $A_i\equiv3$. These facts exclude intersections between
secondary families and between primary and secondary families.
The $m=1+|\mathcal S|\le41$ forms have pairwise independent linear
parts. Repeated targets may be restored afterwards with repeated values.

At $q\mid W$, the forms equal $c$; direct reduction gives
$c\equiv1\pmod8$, $c\equiv3\pmod7$, and $c\equiv1\pmod q$ at
all other odd primes $q\le29$. Thus the forms and their successors
are units at every odd $q\mid W$. For $q>29$, the affine map from
$(x,y)$ to $(u,v)$ is invertible. At $q=31,37,41$, all $A_i=3$
modulo $q$, so $|\mathcal S\bmod q|\le8$. For $q\ge43$ this
cardinality is at most $40$. In either case choose distinct residues
$b,e$ outside the slopes. Put $\gamma=V_0^{-1}$ modulo $q$ and,
reversing $b,e$ if necessary, require $e-b\ne-\gamma$. Then
\[
 v=(e-b)^{-1},\quad u=bv
\]
satisfies
\[
 u-sv=v(b-s)\ne0,\quad u-sv+1=v(e-s)\ne0,
 \quad\gamma v\ne0,\quad\gamma v+1\ne0.
\]
This proves all the value and successor local conditions. No condition
on a successor is imposed at $2$.

Add an unused integer coordinate $z$, and choose a fixed rectangular
box $\Omega$ inside the cone $v>0$, $u/v>\max(0,|A_1|,\ldots,|A_k|)+2413$ in the
homogeneous leading coordinates. All forms on $X\Omega$ lie between
$\alpha X$ and $\beta X$ for positive fixed constants. Fix $Q>82$
containing every prime divisor of $W$ and every nonzero $K_s-K_t$.
Impose one good residue at every prime up to $Q$. For a cutoff $L>Q$,
exclude $\Psi_\nu=0,-1$ at every additional prime up to $L$.
The surviving proportion is
\[
 \gg_{\mathbf a}\prod_{Q<q\le L}(1-2m/q)
 \gg_{\mathbf a}(\log L)^{-2m},
\]
by \cite[Theorem~1, modulus~2]{williams-mertens}.

For fixed $L$, apply \cite[Main Theorem, p.~1763]{greentao} to every
surviving lattice residue in the convex box. The singleton partition
in \cite[Definition~1.5 and Lemma~1.6]{greentao} bounds complexity by
$m-2\le39$. The inputs $\mathrm{MN}(39)$ and $\mathrm{GI}(39)$ follow
from \cite[Theorem~1.1]{moebius} and Section~\ref{v15:gi-full}.
At imposed primes the local factor is $(1-1/q)^{-m}$; at the other
primes it is $1+O(m^2/q^2)$. The singular product therefore has a
positive lower bound independent of $L$. All coefficient and modulus
bounds are fixed before $X$ tends to infinity. Summing the finitely
many asymptotics and removing proper prime powers gives
\begin{equation}\label{eq:old41-lower}
 \#\{\text{prime patterns}\}\ge
 c_{\mathbf a}\frac{X^3}{(\log L)^{2m}(\log X)^m}.
\end{equation}
The constant $c_{\mathbf a}>0$ is independent of $L$; the threshold
for $X$ may depend on it. Proper powers contribute at most
$O_{\mathbf a}(X^{5/2}(\log X)^{m+1})$ before removing von Mangoldt
weights, since each fixed value has $O_{\mathbf a}(X^2)$ fibers.

We give the uniform upper estimates needed to impose the order condition.
For forbidden sets of cardinality $\rho(q)$ in a two-dimensional box
of side $N$, put $g(q)=\rho(q)/q^2$, $h(q)=g(q)/(1-g(q))$, and extend
multiplicatively on squarefree support. CRT gives
\[
 A_d=N^2g(d)+R_d,\qquad |R_d|\ll Nd+d^2.
\]
With $G(Z)=\sum_{d\le Z}h(d)$, take
\[
 \lambda_d=\frac{\mu(d)}{G(Z)\prod_{q\mid d}(1-g(q))}
 \sum_{\substack{e\le Z/d\\(e,d)=1}}h(e)\quad(d\le Z).
\]
Then $\lambda_1=1$, $|\lambda_d|\le d$, and the Selberg quadratic
form equals $G(Z)^{-1}$. To verify the last assertion, use
$y_r=\sum_{r\mid d}\lambda_dg(d)=\mu(r)h(r)/G(Z)$ and
$1/g(n)=\sum_{r\mid n}1/h(r)$. The number of sifted points is thus
\begin{equation}\label{eq:old41-upper}
 \ll N^2/G(Z)+NZ^6+Z^8.
\end{equation}
For $\rho(q)=2q-1$, supported on $q\equiv1\pmod4$ with at most one
prime omitted, $G(Z)\gg_Q\log Z$, uniformly in that omission.
For $\rho(q)=(m-1)q$ at all $q>Q$,
$G(Z)\gg_{Q,m}(\log Z)^{m-1}$. These follow from
\cite[Theorem~1, moduli~4 and~2]{williams-mertens}. Truncation is
justified by using $z=Z^{1/512}$ and the probability distribution
$h(d)/\prod_{q\le z}(1+h(q))$: its expected $\log d$ is at most
$40\sum_{q\le z}(\log q)/q<\tfrac12\log Z$ for large $Z$.
Thus at least half the mass lies below $Z$. Removing one norm-sieve
prime changes its product by a factor at least $1/2$.

For a selected prime, $(7/p)=-1$ and $p\equiv1\pmod8$. The two
reductions of $3-\sqrt2$ have product $7$ and opposite square classes,
so their traces are opposite. Their discriminants $64d^6$ are units.
The automorphism $(x,y)\mapsto(-x,iy)$ is rational over $\mathbb F_p$.
By the Tate-module injection \cite[III.7.4]{silverman}, its rational
centralizer is $\mathbb Q(i)$: at the auxiliary prime $3$ the
centralizer has dimension two because $T^2+1$ is irreducible.
Frobenius is integral by \cite[V.2.3.1(b)]{silverman}; hence it is
$a+bi\in\mathbb Z[i]$. The trace and determinant formulas
\cite[III.8.6]{silverman} give
\[
 a^2+b^2=p,\qquad d_1=p+1-2a,\qquad d_2=p+1+2a.
\]
Full rational two-torsion gives $4\mid d_1,d_2$, and
$v_2(d_1+d_2)=2$, so their gcd has $2$-part exactly $4$.
Also $0<d_i<2p$ and $d_i\ne p$, so $p\nmid d_1d_2$.
An odd common divisor $\ell$ divides $a$ and $p+1$, and forces
$b^2\equiv-1\pmod\ell$. The successor conditions exclude $\ell\le L$.

The number of bad prime values of any fixed form is
\begin{equation}\label{eq:old41-bad}
 \ll_{\mathbf a}X/(L\log X)+X^{7/8}+\sqrt X\log X.
\end{equation}
Indeed, for $L<\ell\le X^{1/8}$, write
$a=\ell r$, $b=b_0+\ell s$, with at most two $b_0$ modulo $\ell$.
The box has side $N\asymp\sqrt X/\ell\gg X^{3/8}$.
Use the norm case of \eqref{eq:old41-upper} with $Z=X^{1/100}$;
its error is uniformly $o(N^2/\log X)$. Sum
$O(X/(\ell^2\log X))$. For $\ell>X^{1/8}$, $a\ne0$ gives
$O(\sqrt X/\ell)$ choices with no additive constant, and the
$b$ count is $O(\sqrt X/\ell+1)$. Summation proves the remaining terms.

On a fiber $\Psi_0=p$, use $(x,z)$; the other forms are
$c+Wx-K_s p$. On a fiber $\Psi_{s_0}=p$, use $(y,z)$ and
\[
 x=K_{s_0}y+\frac{p-c(1-K_{s_0})}{W},\qquad
 \Psi_s=p+(K_{s_0}-K_s)\Psi_0.
\]
The quotient is integral for an actual fiber. For
$Q<q\le X^{1/100}<p$, the remaining zero lines are distinct and
parallel, with total cardinality $(m-1)q$. The second case of
\eqref{eq:old41-upper} bounds each fiber by
$O_{\mathbf a}(X^2/(\log X)^{m-1})$, independently of $p,L$.
Multiply by \eqref{eq:old41-bad} and sum over $m$ positions.
The bad patterns are at most
\[
 C_{\mathbf a}\frac{X^3}{L(\log X)^m}
 +O_{\mathbf a}\left(
 X^{23/8}/(\log X)^{m-1}+X^{5/2}/(\log X)^{m-2}\right).
\]
Choose $L$ so large that $C_{\mathbf a}/L$ is less than half
$c_{\mathbf a}/(\log L)^{2m}$, then choose $X$ large.
A good pattern survives, proving the assertion.
\end{proof}

\subsubsection{A shared auxiliary point with private shifts}\label{subsec:o8-common-J-proof}
Let $\mathcal P$ be the finite set of distinct selected primes, all
larger than $76176$. At each $p$ fix an orientation of its two places,
and let $d_{p,1},d_{p,2}$ be the corresponding group cardinalities.
Set
\[
 j=\operatorname{lcm}_{p\in\mathcal P}d_{p,2}>0,
 \qquad J=[j]P_*.
\]
For each $p$, independently choose an integer $h_p$ satisfying
\begin{equation}\label{eq:old41-private-crt}
 h_p\equiv0\pmod{d_{p,1}},\qquad
 h_p\equiv4\pmod{d_{p,2}},\qquad
 h_p\equiv4\pmod{p^{64+v_p(j)}}.
\end{equation}
The first two congruences are compatible because their modulus gcd
is $4$. The last is compatible because $p\nmid d_{p,1}d_{p,2}$.
The resulting arithmetic progression contains an $h_p$ outside
$\{0,4,-4\}$. Put
\[
 X_0=[h_p]P_*,\qquad X_i=[h_p n_i]P_*,\qquad
 Y_i=[(h_p-4)n_i]P_*,\quad n_0=1.
\]
These are private for each contact occurrence, whereas $J$ is common.
Use the same orientation at repeated occurrences of $p$.

At the primary place all $X_i$ are formal; at the opposite place
$J$ and all $Y_i$ are formal. The formal logarithm is a homomorphism
and an isometry on $p\mathbb Z_p$ by
\cite[IV.6.4 and its proof; VII.2.1--2.2]{silverman}. The hypotheses hold
because $p$ is odd, split, and of good reduction. For any integers
$k,j$ whose multiples of $P_*$ lie in the formal group,
\[
 \frac{\log([k]P_*)}{\log([j]P_*)}=\frac{k}{j}.
\]
To see this without defining a logarithm of $P_*$ itself, apply the
homomorphism law to $[j]([k]P_*)=[k]([j]P_*)$. Nonzero formal
multiples have nonzero logarithm. Therefore
\[
 v_p\left(\frac{z(Y_i)}{z(J)}\right)
 =v_p\left(\frac{(h_p-4)n_i}{j}\right)\ge64,
 \qquad
 \frac{\log(X_i)}{\log(X_0)}=n_i.
\]
This proves the required local comparison conditions. Although other
$d_{q,2}$ may be divisible by $p$, their effect is exactly the term
$v_p(j)$ in \eqref{eq:old41-private-crt}; no cross-prime gcd assertion
is needed.

For the original truncation
\[
 L(z)=z\sum_{r=0}^{15}\frac{(-1)^r\binom{2r}{r}}{4r+1}(Dz^4)^r,
\]
the full logarithm has the same series extended over all $r\ge0$.
This follows by integrating
$(1+4Dz^4)^{-1/2}$, the invariant differential in the parameter
$z=-x/y$. If $v_p(z)\ge1$ and $p>76176$, every omitted term has
valuation at least $65v_p(z)$: its integer numerator can only increase
valuation, and for $n=4r+1\ge65$,
$n v_p(z)-v_p(n)\ge65v_p(z)$.
For $n<p$ this is immediate; for $n\ge p$ use $n-65\ge v_p(n)$.
Thus $L(z)/\log(z)\in1+p^{64}\mathbb Z_p$. The required primary
truncated-log ratios are congruent to the integers $n_i$ modulo $p^{64}$.

Conversely, in an accepted comparison at any $p>76176$, its conjugate
parameter-ratio tests imply
\[
 m_i-4n_i\equiv0\pmod{p^{64}},\qquad h-4\equiv0\pmod{p^{64}},
\]
since $j$ is a nonzero integer. Hence $h$ is a $p$-unit. The same
logarithm identity and truncation estimate give
$L(X_i)/L(X_0)\equiv m_i/h\equiv n_i\pmod{p^{64}}$ and show
that this ratio is integral. The rational proxy test then gives the
claimed integrality and congruence for every $b_i$.

\subsubsection{The common root approximation and parity boundary}\label{subsec:o8-root-proof}
For each distinct selected $p$, retain the chosen square root
$\rho_p$ of $2$. After constructing the private points, each field
test $A+B\sqrt2$ of order $k$ has $v_p(A+B\rho_p)\ge k$.
Choose a common integer $\theta$ sufficiently close to $\rho_p$ at
every selected prime that
$v_p(\theta-\rho_p)\ge k-v_p(B)$ for every nonzero coefficient $B$,
and at least one. There are finitely many requirements; CRT supplies
such an integer. Since $\theta+\rho_p$ is a unit,
\[
 v_p((\theta^2-2)B)\ge k,\qquad v_p(A+\theta B)\ge k.
\]
Repeated primes require the same root and impose compatible precision
bounds. Conversely, the root test and these two rational tests imply
the field test at the selected root. This verifies the circuit's one
common $\theta$, not a separate root variable per contact.

Every raw contact parameter in the complete assignment is a positive
prime congruent to $c$ modulo $W$. It is therefore a square in
$\mathbb Q_2,\mathbb Q_3,\mathbb Q_{17}$. The binary norm batch has
unit root at its contact prime, because normalized leaves are integral
and its leftmost constant leaf is $1$. Its root is positive and nonzero
in the real embedding. The argument $r\mathcal W^2$ consequently is
positive, has square reduced denominator, and is locally square at
$2,3,17$. These are the exact hypotheses of the original specialized
parity-completeness input. This verifies its consumption but does not
constitute a new proof of that input, which is supplied here by
Lemma~\ref{lem:original-seven-parity-completeness}. In particular no integrality of
the comparison points at these three small primes is required here.

\subsection{Scalar height and the defining relation}

Put $\beta=X^2+2$ and
\[
 r(\beta)=\beta(\beta-1)(\beta-2401)(32\beta+343)(2\beta-1).
\]
This is nonzero for every rational $X$, since $X^2=2399$ has no
rational solution. Introduce $J_0\ne O$ and apply a scalar norm
batch to the two coefficients of each of
$f(J_0),f(S),z(S)/z(J_0)$, of orders $1,1,1,1,48,48$.
In this batch replace the pole argument by $\delta\mu rW^2$, where
$\delta\in\{1,-1\}$ and $\mu$ is a squarefree product of
$2,3,5,7,11,13,17$. Precisely, eight Boolean variables produce
\[
 A=(1+e_2)(1+2e_3)(1+12e_{13})(1+16e_{17}),\quad
 B=(1+4e_5)(1+6e_7)(1+10e_{11})(1-2\eta),
\]
and the pole argument is $ABrW^2$. These factors have no valuation
outside $S_*$. Every odd five-contact of $\beta$ gives an odd
valuation of $r(\beta)$ there. At infinity its valuation is five
times that of $\beta$; at a finite contact all other factors are units.
Thus the scalar tests force $M_5(\beta)^{48}\mid s$, and
Theorem~\ref{hyp:o8-height} gives
\begin{equation}\label{eq:o8-scalarheight}
 h(\beta)\le H_*\sqrt{|s|}.
\end{equation}

For integral $X$, the scalar $r$ is a nonzero integer. The square
classes of $2,3,5,7,11,13,17$ form a basis of
$\Q_2^*/\Q_2^{*2}\times\Q_3^*/\Q_3^{*2}
\times\Q_{17}^*/\Q_{17}^{*2}$.
Choose $\delta$ and $\mu$ so that $\delta\mu r$ is positive and
a square at these three places. Choose $J_0$ in the required formal
subgroups and then choose $S$ sufficiently divisible relative to $J_0$
at all primes dividing $r\mu$, together with the anchor conditions.
All normalized nonconstant leaves are then integral at primes dividing
$r$, and vanish at primes dividing $\mu$ but not $r$. The root is
a unit at these primes. Consequently $\delta\mu rW^2$ has square
reduced denominator, and Lemma~\ref{hyp:o8-pole} supplies the
scalar pole witnesses.

Take $A=nP_*,S=sP_*,T=tP_*,Z=zP_*\ne O$, and use
$(\beta,\xi,\tau,\zeta)$ as the four proxies. Impose
\begin{equation}\label{eq:o8-Lucas}
 \tau^2-\xi(\xi-2)\tau\zeta+(\xi-2)^2\zeta^2=1.
\end{equation}
The nine rational targets in \eqref{eq:o8-contacts} are a nonzero
common denominator $e$ and the eight coefficients of
$ex(A),ex(S),ex(T),ex(Z)$. Require its forty-one comparison batches,
the scalar batch, and the anchor just described.

If every associated five-contact radical has its $K$th power dividing
a nonzero integer $N$, Theorem~\ref{hyp:o8-height} and
$Ra_i+3=u/v+w_i/v$ give
\[
 h(a_i)<2H_*|N|^{24/K}+14,
 \qquad h(x(Q))<7H_*|N|^{24/K}
\]
for each of the four points. The second bound follows by applying the
first to $e$ and both coefficients of $ex(Q)$ and absorbing the
constant smaller than fifty into $H_*$.

\begin{lemma}[Lucas growth]\label{lem:o8-lucas}
If $s,t,z\in\Z$, $|s|=a\ge7$, $z\ne0$, and
\[
 t^2-s(s-2)tz+(s-2)^2z^2=1,
\]
then $|t|>\exp(a-3)$. Conversely each such $s$ admits positive
integers $t,z$ satisfying this equation.
\end{lemma}
\begin{proof}
Put $w=\operatorname{sgn}(s)(s-2)z$ and
$\rho=(a+\sqrt{a^2-4})/2$. The equation becomes
$t^2-atw+w^2=1$. Every positive norm-one unit of $\Z[\rho]$ is a
power of $\rho$: multiply it by a power into $[1,\rho)$, write the
result as $v=A+B\rho$, and use
$B=(v-v^{-1})/(\rho-\rho^{-1})\in[0,1)\cap\Z$.
Thus $v=1$. If $U_0=0$, $U_1=1$ and
$U_{k+1}=aU_k-U_{k-1}$, the identities for $\rho^{\pm k}$ give
$|w|=U_k$ and $|t|\in\{U_{k-1},U_{k+1}\}$ for some $k\ge1$.
Modulo $a-2$ one has $U_k\equiv k$, and modulo $a+2$ one has
$U_k\equiv(-1)^{k-1}k$. The divisibility of $w$ by $|s-2|$ therefore
gives $k\ge a-2$ when $s>0$ and $k\ge a+2$ when $s<0$.
Induction gives $U_k\ge(a-1)^{k-1}$, whence
$|t|\ge(a-1)^{a-4}>e^{a-3}$ for $a\ge7$.
For existence take $k=a-2$, $z=U_k/(a-2)$ if $s>0$, and
$k=a+2$, $z=U_k/(a+2)$ if $s<0$, in both cases with $t=U_{k+1}$.
\end{proof}

\begin{proof}[Proof of Theorem~\ref{thm:o8-conditional}]
Suppose the relation holds, and let
$B=\max(|n|,|s|,|t|,|z|)$ and
$\Delta=t^2-s(s-2)tz+(s-2)^2z^2-1$.
If $\Delta\ne0$, the congruences \eqref{eq:o8-proxies} imply that
every coordinate contact radical to the sixty-fourth power divides
$\Delta$. Since $|\Delta|<5B^4$, the index-height and shared-height
bounds give
\[
 B^2-6<7H_*|\Delta|^{3/8}<14H_*B^{3/2}.
\]
Hence $B<(14H_*+6)^2\le(20H_*)^2$, contradicting the anchor.
Thus $\Delta=0$. Lemma~\ref{lem:o8-lucas} gives
$|t|>\exp(|s|-3)$ because $z\ne0$.

If $\beta\ne n$, put $N=|\operatorname{num}(\beta-n)|$ in lowest
terms. The same congruences give
\[
 B^2<7H_*N^{3/8}+6,\qquad N\le2e^{h(\beta)}B,
\]
whence $\log B<\frac8{13}\log(16H_*)+\frac3{13}h(\beta)$.
Writing $a=|s|>(20H_*)^2$, \eqref{eq:o8-scalarheight} implies
$h(\beta)<a/20$ and $\log(16H_*)<a/20$. Therefore
$\log B<11a/260<a-3$, another contradiction.
We obtain $\beta=n\in\Z$, and $X\in\Q$, $X^2=\beta-2\in\Z$
imply $X\in\Z$ by reducing its denominator.

Conversely let $X\in\Z$ and set $A=\beta P_*$. Choose the scalar
height data and then an anchored $S$ with all required divisibilities.
Lemma~\ref{lem:o8-lucas} gives positive integers $t,z$ satisfying the
index equation; set $T=tP_*$ and $Z=zP_*$. Choose an integer $e$
clearing their abscissa coefficients. The ordinary prime pattern
provides \eqref{eq:o8-contacts}. Set the proxies equal to their integer
indices, choose $J$ and the private points as above, and finally choose
the common $\theta$. The norm and pole completeness statements supply
all remaining witnesses. No subsequent choice changes the target points.
\end{proof}

\subsection{The base compiler prescription and finite verification}

Every equation above is implemented with named arithmetic intermediates
so that each rational residual has degree at most four. A field product
uses $(a+b\sqrt2)(c+d\sqrt2)=(ac+2bd)+(ad+bc)\sqrt2$.
Scalar inverses are equations $uv=1$. Field inverses use the two
coefficient equations of the field product being one; subtraction has
the separate scalar inverse of its abscissa-difference norm.
The logarithm uses Horner evaluation in $(Dz^4)^2$ with block length
two. The two doublings of the fixed point $P_*$ also use arithmetic
intermediates; inserting its large precomputed fourth multiple as a
coefficient would give a different polynomial. The chart is exactly
\eqref{eq:o8-chart}, including the intermediate $h=vw$.

For each nonzero rational residual, multiply by the least common multiple
of its coefficient denominators, then divide all coefficients by their
positive greatest common divisor. Denote the resulting integral rows
by $R_1,\ldots,R_{15107}$. In the variable order specified by
\texttt{degree8-original/build.py} and its residual table, define
\begin{equation}\label{eq:o8-polynomial}
 F_8=\sum_{i=1}^{15107}R_i^2.
\end{equation}
This equation is an exact finite prescription for every coefficient.
The residual table, its full collected expansion, and the construction
certificate accompany the manuscript. No huge constant is treated as
an additional free parameter. In particular, all twenty-five variables
in \eqref{eq:o8-chain} belong to the stated witness count.

The recursive module counts are
\[
\begin{array}{c|rr}
 \text{module}&\text{new variables}&\text{residuals}\\\hline
 \text{pole}&7&4\\
 \text{anchor}&16&6\\
 \text{scalar batch, including pole}&34&31\\
 \text{forty-one-test batch, including pole}&58&55\\
 \text{one joint comparison, including batch}&373&360\\
 \text{root, including all its calls}&15646&15082\\
 \text{shared anchor chain}&25&25
\end{array}
\]
The root has $303$ direct variables and $285$ direct residuals; its
calls are one anchor, one scalar batch, and forty-one comparisons.
Thus $303+16+34+41\cdot373=15646$ and
$285+6+31+41\cdot360=15082$. Expansion of \eqref{eq:o8-polynomial}
gives the finite assertions of Proposition~\ref{prop:o8-finite}.
Over $\Q$, a sum of squares vanishes exactly when every summand
vanishes, so the residual interpretation proves the theorem.
This computation checks the polynomial and its compilation. The proof
of the height estimate is the arithmetic argument given above.

% End degree-eight construction
% Begin explicit quartic constructions
\section{An explicit quartic defining the integers}
\label{sec:quartics}
This is the second final polynomial of the manuscript. The original huge
polynomial remains unchanged. The base residual compiler in
Section~\ref{sec:degree8-original} supplies construction machinery and
inherited arithmetic inputs; it is not a third final polynomial.

\begin{theorem}[An explicit quartic definition]\label{thm:quartic}
With precisely the inherited arithmetic hypotheses and qualifications of
Section~\ref{sec:degree8-original} and Appendix~\ref{app:sources}, the
coefficient file below defines one integer polynomial $\Phi$ of total
degree $4$, with $72\,925$ nonzero monomials in $4\,973$ variables.
For every $q\in\Q$,
\[
 q\in\Z\quad\Longleftrightarrow\quad
 \exists(w_1,\ldots,w_{4972})\in\Q^{4972},\qquad
 \Phi(q,w_1,\ldots,w_{4972})=0.
\]
There are $4\,972$ existential witnesses besides the original input.
The coefficient bound is
\[
 \|\Phi\|_\infty=685783949153203534233890<2^{80}.
\]
\end{theorem}

\subsection{Literal specification and coefficient file}
The complete coefficient table is \nolinkurl{quartics/quartic-4973/polynomial.expanded.txt}.
After one comment line, each line gives a nonzero signed integer coefficient
followed by a nondecreasing list of variable indices. Repeated indices give
their multiplicities as exponents; a coefficient alone is the constant.
Thus \texttt{-6 0 2 2} denotes $-6X_0X_2^2$. The original input is $X_0$.
Like terms are collected, and the rows are in lexicographic index-list order.
The occurrence-frequency relabeling is recorded in
\nolinkurl{quartics/quartic-4973/variable-map.json} and fixes index zero.
The gzip companion has exactly the same uncompressed bytes.
The coefficient table has $1\,622\,053$ bytes, including
its comment and line endings, and SHA-256
\begin{center}\footnotesize\ttfamily
094f6a7ed17e7cc8d5d17a2f809127bbaffd0fe343cb2a32d18ac97e61279876
\end{center}

This specifies every coefficient without an auxiliary choice of witnesses.

\subsection{Arithmetic source and retained foundations}
\label{subsec:quartic-source}
The selected source uses the following retained height, contact and
reciprocal-logarithm proofs before the new source equivalence is established.
An element of $K$ is represented by two rational coefficients, with
\[
 (a_0,a_1)(b_0,b_1)=(a_0b_0+2a_1b_1,a_0b_1+a_1b_0),\qquad
 a^{-1}=\frac{(a_0,-a_1)}{a_0^2-2a_1^2}\quad(a\ne0).
\]
The inverse denominator is nonzero by the irrationality of $\sqrt2$.

% BLOCK HEIGHT19.
\begin{theorem}[A sharper five-contact height estimate]\label{thm:quartic-height19}
Keep $S_*$, $\mathcal D$ and $M_5$ from Theorem~\ref{hyp:o8-height}, and put
\[
 B_*=\prod_{p\le76176}p,\qquad
 A_{19}=2^{109}210^{1800},\qquad H_{19}=2A_{19}B_*^{19}+128.
\]
For every $a\in\mathbb Q\setminus\mathcal D$,
\[
 h(a)\le H_{19}M_5(a)^{19},
 \qquad \log_2H_{19}<2093966.
\]
Here $h$ is the absolute logarithmic Weil height.
\end{theorem}
\begin{proof}
We use the geometric construction and compatible representations in the
proof of Theorem~\ref{hyp:o8-height}, not merely its final exponent-$24$
inequality. Write $R=B_*M_5(a)$ and $C=210^{300}$.
The geometric comparison there is
\[
 h(a)\le2\max(0,h_F(A))+128.
\]
All $h_F$ in this proof use the original stable Faltings metric.
The CM estimate in that proof is $h_F(A)<854R$.

In the non-CM case the same construction gives an absolutely simple
quaternionic surface $A/L$, a normal field $L'/\mathbb Q$, and the final
primitive weight-two newform $f$, after the finite twist has been
absorbed, with
\[
 [L:\mathbb Q]\le576,\quad d'=[L':\mathbb Q]\le2304R,\quad
 N\le CR^2,\quad \operatorname{rad}(N)\mid R.
\]
Its two-dimensional compatible representation $r_\ell$ satisfies
\[
 V_\ell(A)|_{G_{L'}}\simeq
       (r_\ell\oplus r_\ell)|_{G_{L'}},\qquad
 (\det r_\ell)^2|_{G_{L'}}=\chi_{\rm cyc}^2|_{G_{L'}}.
\]
These are (LR3)--(LR5) and the determinant calculation in the existing
modularity argument. There is no bounded-degree hypothesis on the
coefficient field of $f$.

Let $\epsilon$ be the nebentypus of this final form. In the Tate-module
convention, with arithmetic Frobenius, $\det r_\ell=\epsilon\chi_{\rm cyc}$;
see \cite[\S12.5, (12.5.1)--(12.5.2) and Remark 12.5.1,
pp.~119--120]{diamond-im}. Equivalently, at a good prime its
characteristic polynomial is $X^2-a_pX+p\epsilon(p)$.
A consistent contragredient convention replaces the finite character by
its inverse, with the same kernel and order. The character $\epsilon^2$
is trivial on $G_{L'}$, so its image order divides $d'$. Consequently
\[
 t=\operatorname{ord}(\epsilon)\le2d'\le4608R.
\]
This bounds the character order, not the coefficient-field degree.

Put $H=\ker\epsilon\subset(\mathbb Z/N\mathbb Z)^*$ and
\[
 \Gamma_H=\left\{
 \begin{pmatrix}a&b\\c&d\end{pmatrix}\in\Gamma_0(N):
 d\bmod N\in H\right\}.
\]
Since $\epsilon(-1)=1$, this group contains $-I$.
Let $X_H=X_1(N)/H$ for the diamond action and $J_H=\operatorname{Jac}(X_H)$.
The curve, diamond action and quotient are defined over $\mathbb Q$.
The rational model here parametrizes a point of exact order $N$, up to
$H$, and is the model specified by \cite[\S2.1, (2.1.8)]{mascot-moduli};
it introduces no cyclotomic field-degree factor.

Write $q:J_1(N)\to A_f$ for the newform quotient. Its pulled-back
differentials are spanned by all Galois conjugates of $f$; see
\cite[Theorem~1(i),(iv), p.~526]{shimura-factors} and
\cite[\S3.1]{baker-et-al-finiteness}. Every $h\in H$ fixes those differentials,
since its eigenvalue on each conjugate is a conjugate of $\epsilon(h)=1$;
the diamond convention is \cite[\S3.2.1]{baker-et-al-finiteness}.
Homomorphisms of abelian varieties in characteristic zero are determined
by their differentials. Hence $q\circ\langle h\rangle=q$.
For
\[
 e_H=\frac1{|H|}\sum_{h\in H}\langle h\rangle
       \in\operatorname{End}_{\mathbb Q}^0(J_1(N))
\]
we therefore have $qe_H=q$.

Let $H_{\rm eff}$ be the image of $H$ in
$\operatorname{Aut}(X_1(N))$, and let $\pi:X_1(N)\to X_H$ be the quotient.
On Jacobians,
\[
 \pi^*\pi_*=\sum_{h\in H_{\rm eff}}h,\qquad
 \pi_*\pi^*=[|H_{\rm eff}|].
\]
The averages over $H$ and $H_{\rm eff}$ agree. In particular the trivial
action of $-1$ is not counted twice in the first identity.
The image of $e_H$ is thus $\mathbb Q$-isogenous to $J_H$, and its
restriction of $q$ is surjective. Poincare reducibility makes $A_f$
a $\mathbb Q$-isogeny factor of $J_H$.

The accepted realization of $r_\ell$ in $A_f$ puts it in $V_\ell(J_H)$.
Over $L'$, Faltings' Hom theorem and extension of scalars give a nonzero
homomorphism $A\to J_H$. Absolute simplicity and Poincare reducibility
give an $L'$-isogeny $A\times Z\to J_H$. In particular
$g=\dim J_H\ge2$. The Jacobian has its canonical principal polarization;
no bound on a polarization degree of $A$ or $Z$ is being assumed.

The map $X_H\to X(1)$ has degree
\[
 d_H=[\operatorname{SL}_2(\mathbb Z):\Gamma_H]
     =Nt\prod_{p\mid N}(1+1/p).
\]
There is no factor of two here, because $-I\in\Gamma_H$.
It is branched only over the cusp and the two elliptic points, so
$\deg_B(X_H)\le d_H$. The congruence-curve genus bound gives
$g\le1+d_H/12\le d_H$, using $g\ge2$.
These degree and genus conventions, and the resulting height inequality,
are recorded in \cite[\S4, (4.1)--(4.2), p.~8]{vonKaenel2021}.
For every prime $p\ge5$, $(1+1/p)^6\le p$, while the possible factors
at two and three have product at most two. It follows that
\[
 d_H\le9216CR^{19/6}<T,\qquad
 T=2^{14}CR^{19/6},\qquad g<T.
\]
By \cite[Theorem~1.1.1 and \S2.3]{javanpeykar},
\[
 h_F(J_H)\le13000000\,d_H^5g<2^{24}T^6.
\]
This is the generic Belyi-curve estimate; no squarefree-level restriction
or trivial-nebentypus assumption enters it.

Apply \cite[Theorem~1.9(1), pp.~7--8]{GaudronRemond2023} with source
$(J_H)_{L'}$ and target $A\times Z$. Combining its degree bound with
half-logarithmic isogeny-height variation gives
\[
 h_F(A)+h_F(Z)\le h_F(J_H)+V,\quad
 V=2g^2\left(8g^2\log(7g)+\log d'
     +\log\max(1,\log d',h_F(J_H))\right).
\]
The Bost lower bound in \cite[(2.2), p.~5]{vonKaenel2021} gives
$h_F(Z)>-3g/2$. Put $u=\log T$. Then
\[
 V\le16T^4(2+u)+2T^2(25+7u)\le T^6.
\]
Indeed $\log d'<8+\log R\le8+u$, and the last logarithm in $V$ is
less than $17+6u$. For $T\ge16$ the inequalities
$2+\log T\le T^2/32$ and $25+7\log T\le T^4/4$ follow at $16$
and remain true by differentiation. Thus
\[
 h_F(A)<(2^{24}+1)T^6+\tfrac32T
       <2^{25}T^6=A_{19}R^{19}.
\]
The earlier CM estimate is also less than this bound.
The geometric comparison and $M_5(a)\ge1$ prove the assertion.
Finally $B_*<2^{109445}$, $210<2^8$ and $A_{19}B_*^{19}>64$ give
\[
 H_{19}<4A_{19}B_*^{19}
       <2^{\,2+109+1800\cdot8+19\cdot109445}
       =2^{2093966}.
\]
\end{proof}

% BLOCK ROOT37: no denominator-clearing variable is added to this source.
\begin{lemma}[Thirty-seven contacts for rational coefficients]
\label{lem:quartic-rational37}
Keep $K=\Q(\sqrt2)$, $d=3-\sqrt2$, $D=d^2$, and
$E:y^2=x^3-Dx$ from Section~\ref{sec:degree8-original}.
Let $A,S,T,Z$ be nontorsion points of $E(K)$.
Order the eight rational coefficients of
$x(A),x(S),x(T),x(Z)$ as $a_1,\ldots,a_8$. Put
\[
 R_{\rm c}=47027,\quad W_0=25878772920,\quad c=18484837801,
 \quad V_0=32W_0,
\]
\[
 A_i=R_{\rm c}a_i+3,\quad w_i=A_iv-u.
\]
Use the following thirty-seven nonzero raw contacts:
\begin{equation}\label{eq:quartic-contacts37}
 r_0=v/V_0,\qquad u-\delta v\quad(\delta\in\mathcal D),
 \qquad -(w_i-\delta v)\quad(1\le i\le8,\ \delta\in\mathcal D).
\end{equation}
Here and below the order within $\mathcal D$ is
$(0,1,2401,-343/32)$. The common five positions are numbered $0,\ldots,4$;
the other positions are numbered $5,\ldots,36$ in the displayed order.

Every normalized odd contact of $u/v$ or $w_i/v$ outside $S_*$ is
detected by an odd valuation at a raw position belonging to that slope.
If the radicals for both slopes $u/v,w_i/v$ have their $e$th powers
dividing an integer $N\ne0$, then, with $H=H_{19}$,
\[
 h(a_i)<2H|N|^{19/e}+14.
\]
Consequently the two coefficients of any one of the four abscissas give
\begin{equation}\label{eq:quartic-direct-height}
 h(x(Q))<4H|N|^{19/e}+30<7H|N|^{19/e}
 \quad(Q\in\{A,S,T,Z\})
\end{equation}
whenever all its required positions have those divisibilities.

Conversely, for any fixed eight rational coefficients, the inherited
ordinary-prime construction supplies these raw contacts with every finite
contact a positive selected prime and $r_0$ a positive rational square
times a selected prime. Repeated positions may have equal values.
These contacts have the positivity and local squareclasses used by Lemmas~\ref{lem:quartic-conic-completeness} and \ref{lem:quartic-linear-pole}.
\end{lemma}
\begin{proof}
Nonzero raw contacts imply $v\ne0$ and exclude the four finite branch
values for every slope. At a prime outside $S_*$ all differences of
distinct members of $\mathcal D$ are units. Write $m$ for the minimum
of the valuations of the two coordinates of one slope. If an odd finite
normalized contact has $m$ even, its distinguished raw linear form has
odd valuation. If $m$ is odd, another finite form has valuation $m$.
For an odd contact at infinity use the denominator in the first case
and a finite form in the second. The unit $V_0$ and the signs do not
change parity. This proof uses rational coordinates and also covers
negative valuations; it does not presuppose that a contact is prime.

Apply Theorem~\ref{thm:quartic-height19} to the two slopes. The identity
$R_{\rm c}a_i+3=u/v+w_i/v$ and the elementary logarithmic height inequalities
give the first bound, since $\log(6R_{\rm c})<14$. For
$x(Q)=a+b\sqrt2$ use
$h(x(Q))\le h(a)+h(b)+h(\sqrt2)+\log2$.
This proves \eqref{eq:quartic-direct-height}, since $H>128$ and $|N|\ge1$.

For the converse choose a positive integer $D_0$ with $D_0a_i\in\Z$
for every $i$, after fixing the eight coefficients. This integer is a
choice in the existence proof, not a variable of the source. Set
\[
 T_0=V_0D_0^2,\qquad p_0=c+W_0y,\qquad
 u=c+W_0x,\qquad v=T_0p_0.
\]
For each distinct slope
$b\in\mathcal D\cup\bigcup_i(A_i-\mathcal D)$ put
$K_b=T_0b\in W_0\Z$ and $\Psi_b=u-K_bp_0$.
Together with $p_0$, these are at most thirty-seven integer affine forms,
each congruent to $c$ modulo $W_0$. Distinct slopes give independent
linear parts. Exact duplicate slopes are removed only while applying
the prime theorem and restored at all their original positions afterward.

The value and successor admissibility in
Section~\ref{subsec:o8-pattern-proof} remains valid. At primes dividing
$W_0$ use the same residue $c$. If a prime $\ell>29$ divides $D_0$,
all $K_b$ vanish modulo $\ell$, and $u=p_0=1$ avoids zero and minus one.
If $\ell>29$ does not divide $D_0$, the change from $(u,p_0)$ to $(u,v)$
is invertible. For $\ell=31,37,41$ there are at most eight slope residues,
since $R_{\rm c}a_i\equiv0$; for $\ell\ge43$ there are at most thirty-six.
Choose distinct missing residues $b_0,e_0$ and order them so that
$e_0-b_0\ne-T_0^{-1}$. Then
\[
 v=(e_0-b_0)^{-1},\qquad u=b_0v
\]
make every $\Psi_b$ different from zero and minus one, and make
$p_0=v/T_0$ different from zero and minus one. Interchanging the two
missing residues permits the required ordering, because $\ell$ is odd.

Choose a fixed positive cone
$u/v>\max_i|A_i|+2413$, $v>0$. Enlarge the fixed initial cutoff past
$82$, every prime dividing $D_0$, and the nonzero determinants of the
distinct forms. The inherited prime-pattern proof then applies with its
same unused third coordinate and with complexity at most $35$.
For its fiber upper estimates, fixing $\Psi_{b_0}=p$ gives the same
one-variable forms
$\Psi_b=p+(K_{b_0}-K_b)p_0$.
All coefficients, the cone, the denominator $D_0$ and the cutoff are
fixed before the box grows. Thus the proof uses precisely the inherited
fixed-system prime and successor estimates, not a uniform estimate in a
growing denominator. It gives the positive ordinary primes and the
same-prime elliptic-order conditions used there.

Now $r_0=D_0^2p_0$. At a prime $\ell\mid D_0$ outside $S_*$, both $u$
and $p_0$ are units, $v_\ell(v)=2v_\ell(D_0)$, and
$v_\ell(A_i)\ge-v_\ell(D_0)$. Hence $w_i=A_iv-u$ is a unit.
Both slopes have negative valuation $-2v_\ell(D_0)$, and their normalized
infinite depths are even. No new normalized odd contact is created at
such a denominator prime. At $p_0$, the odd valuation of $r_0$ is one.

For each occurrence choose its private comparison witnesses sufficiently
deep that its norm root is a unit at its assigned prime. The remaining
denominator valuations in the pole argument occur in the square of that
root. The raw contacts are positive, their reduced denominators are
squares, and they are squares in $\Q_2,\Q_3,\Q_{17}$ by the retained
residue $c$ and the square factor $D_0^2$.
Lemmas~\ref{lem:quartic-conic-completeness} and
\ref{lem:quartic-linear-pole} supply the selected ordinary-pole witnesses. Equal prime occurrences use the
same orientation and may use separate private points. No condition on
reduction orders at different primes has been added.
\end{proof}

% BLOCK RECIPROCAL: the released complete proofs with the stated uniform range.
\begin{lemma}[Reciprocal logarithm and its error]
\label{lem:quartic-reciprocal-log}
\label{lem:quartic-uniform-reciprocal}
Put $f=1/x$, $z=-x/y$, and define
\[
 a_0=1,\qquad a_1=-\frac25,\qquad
 a_j=\frac{4j-5}{4j+1}a_{j-1}\quad(j\ge2),\qquad
 G(w)=\sum_{j\ge0}a_jw^j.
\]
On the formal branch at the identity,
\[
 \log_E(z)=zG(Df^2).
\]
For every $J\ge0$ put $G_J(w)=\sum_{j=0}^{J}a_jw^j$.
For every integer $10\le e\le64$, set $m=\lceil e/4\rceil$ and
\[
 G_{m-1}(w)=\sum_{j=0}^{m-1}a_jw^j,\qquad
 L_e(z,f)=zG_{m-1}(Df^2).
\]
Let $p>76176$ be split and choose either place above it. If the point is
nonidentity and lies in the formal group at that place, so
$r=v_p(z)\ge1$ and $v_p(f)>0$, then
\[
 v_p\bigl(\log_E(z)-L_e(z,f)\bigr)\ge(4m+1)r,
 \qquad \frac{L_e(z,f)}{\log_E(z)}\in1+p^e\Z_p.
\]
\end{lemma}
\begin{proof}
The curve equation gives
\[
 f=z^2(1-Df^2),\qquad w=Df^2,\qquad Dz^4=\frac{w}{(1-w)^2}.
\]
Coefficient comparison in the defining recurrence gives
\[
 4w(1-w)G'(w)+(1+w)G(w)=1-w.
\]
Indeed the coefficient of $w$ is $5a_1+a_0=-1$, and each later
coefficient gives $(4j+1)a_j=(4j-5)a_{j-1}$. Differentiating the
curve identity yields $z\,dw/dz=4w(1-w)/(1+w)$. Consequently
\[
 \frac{d}{dz}\{zG(w)\}=\frac{1-w}{1+w}
                   =(1+4Dz^4)^{-1/2}.
\]
The square root has constant term one. This is the normalized invariant
differential for $z$, and $zG(w)$ has zero constant term, which proves
the formal identity.

At the stated place, $D$ is a unit and the formal branch gives
$v_p(f)=2r$, hence $v_p(w)=4r$. The denominator of $a_j$, for $j\ge1$,
divides $5\prod_{k=2}^j(4k+1)$ and therefore divides $(4j+1)!$.
Thus
\[
 v_p(a_j)\ge-v_p((4j+1)!)\ge-\frac{4j+1}{p-1}.
\]
For $j\ge m$ with $4j+1<p$, the coefficient is $p$-integral and
$v_p(a_jw^j)\ge4mr$. Otherwise $j\ge19044$, so $3j\ge4m$ and
$(4j+1)/(p-1)\le j$. Therefore
\[
 v_p(a_jw^j)\ge4jr-j\ge4mr.
\]
These lower bounds tend to infinity. They control the entire infinite
tail, including terms whose denominators are divisible by $p$.
Multiplication by $z$ proves the first estimate. Formal logarithm isometry
gives $v_p(\log_E(z))=r$, and proves the relative estimate.
\end{proof}

\begin{lemma}[Global nonvanishing]
\label{lem:quartic-reciprocal-nonzero}
For every $f\in K$ and every $J\ge2$,
\[
                 G_J(Df^2)\ne0.
\]
In particular $L_e(z,f)\ne0$ whenever $z\ne0$ and $10\le e\le64$.
\end{lemma}
\begin{proof}
Let $\nu_2$ be the rational exponent-of-two valuation and put
\[
 v(a+b\sqrt2)=\tfrac12\nu_2(a^2-2b^2)
             =\min\{\nu_2(a),\nu_2(b)+\tfrac12\}.
\]
The two terms in the norm have valuations of opposite parity when finite.
Norm multiplicativity and the displayed minimum prove the valuation
identities. Its finite values lie in $\tfrac12\Z$, and $v(D)=0$.
Every $a_j$ with $j\ge1$ has valuation one, because $a_1=-2/5$ and
all subsequent recurrence factors have odd numerator and denominator.
If $f=0$, the assertion is immediate. If $t=v(f)\ge0$, the constant
term of $G_J(Df^2)$ is uniquely least in valuation. If $t<0$, then
$t\le-1/2$ and the nonconstant valuations $1+2jt$ strictly decrease;
the last is at most $1-J<0$, below the constant term as well. In either
case a uniquely least term proves nonvanishing. No good-reduction
assumption at two is used.
\end{proof}

\subsection{Conventions and inherited hypotheses}
\label{sec:v4-inherited}

Retain the field and curve of Section~\ref{sec:degree8-original}:
\[
 K=\Q(\sqrt2),\qquad d=3-\sqrt2,\qquad D=d^2,
 \qquad E:\ y^2=x^3-Dx.
\]
For an affine point with nonzero abscissa write $f(P)=1/x(P)$ and
$p(P)=y(P)/x(P)$. The rational function $z=-x/y$ extends with value zero
at $O$ and $T_0=(0,0)$; its other poles are $(d,0)$ and $(-d,0)$.
The letter $J$ denotes a common comparison point. We write $\jmath$ for
$1/z(J)$, to distinguish that field element from a point or an integer index.
The implementation conjugates the complete alternate-comparison identity:
its boundary value is $1/\overline{z(J)}$, and its primary and target
coefficients in that identity are conjugated as well.

The following arguments retain all inherited qualifications of
Section~\ref{sec:degree8-original} and Appendix~\ref{app:sources}. In particular, the restricted converse, rank and original
arithmetic assertions retain their stated proof status. The historical
prime-pattern source used arXiv:1009.3998v5 with its erratum. Here its
inverse-theorem premise is supplied by
Section~\ref{v15:gi-full}, with the published
transference inputs \cite{greentao,moebius}; their stated hypotheses
remain in force. The modularity step uses
Lemma~\ref{lem:qm-modularity-descent} and the published prime-five
statement in \cite[p.~624]{calegari-icm2022}. Its hypotheses, weight-two
conclusion and finite twist are checked for the actual representation
in that lemma. The source-status map in
\cite{direct-p5-modularity-project} separates this input from its
proof ancestry and the preserved GL2-type/Serre reconstruction.
The Yuan inputs used here are checked against Section 2.1 and
Theorem 1.1 of the published Chinese article, including its metric
definitions. The stack argument above does not use the integral
coarse-model Lemma 2.1 or Theorem 2.2; no equivalence of the complete
Chinese and English proofs is asserted. The stable Faltings metric,
arithmetic-Frobenius, diamond-quotient and final-twist conventions are
unchanged. None of the finite source or polynomial checks strengthens these
assumptions.

In particular we use, in exactly their previous scopes, the generator and
torsion statements $2E(K)=\Z P_*$ and $E(K)_{\mathrm{tors}}=E[2](K)$;
the formal-group logarithm and valuation statements; the fixed-system
prime-pattern theorem with arbitrary prescribed finite exclusions; the
height, scalar, anchor and Lucas lemmas; and the restricted-converse and
simultaneous-realization inputs used for the pole predicate. The one
additional explicit trace application below is identified in
Lemma~\ref{lem:quartic-anchor-four}.

Put
\[
 B_* =\prod_{p\le76176}p,\qquad
 H=2(2^{109}210^{1800})B_*^{19}+128,
 \qquad \log_2 H<2093966.
\]
This is $H_{19}$ of Theorem~\ref{thm:quartic-height19}. The
source parameters remain high order $40$, beta order $12$, scalar orders
$(1,1,1,1,32,32)$ and anchor constant $C=3^{2^{25}}$.

We also retain Lemma~\ref{lem:quartic-reciprocal-log}. Explicitly,
write
\[
 \sqrt{1-w}\sum_{j\ge0}\frac{\binom{2j}{j}}{4^j(4j+1)}w^j
       =\sum_{j\ge0}a_jw^j,\qquad
 G_m(w)=\sum_{j=0}^m a_jw^j.
\]
The polynomials used here are $G_9$ and $G_2$. They are nonzero at every
argument $Df^2$ with $f\in K$, by the previously proved dyadic argument.
For a nonidentity formal-group point at either place above a split prime
$p>76176$, the approximation $-G_9(Df^2)/p(P)$ has relative order at least
$40$, and $-G_2(Df^2)/p(P)$ has relative order at least $12$. Actual formal
membership is required: small $z$ alone is not the premise of this lemma.

\subsubsection{The retained scalar, anchor and beta interfaces}
The scalar input is still $\beta=X^2+2$. The scalar batch tests the two
rational coefficients each of $f(J_0),f(S),z(S)/z(J_0)$, with orders
$(1,1,1,1,32,32)$ in that order. Its original pole predicate is retained.
The anchor constant is computed by
\[
 c_0=3,\qquad c_{i+1}=c_i^2\quad(0\le i<25),\qquad C=c_{25}.
\]
Thus the integer $C=3^{2^{25}}$ is a witness value, not a polynomial
coefficient. The four-witness anchor below consumes this same terminal
value. The height targets are the eight rational abscissa coefficients
of Lemma~\ref{lem:quartic-rational37}.

For every beta-bearing comparison, the common and beta quotient equations
are explicitly
\[
 \Lambda p_H=G_H,\qquad B_A=C_Ap_A,\qquad B_A\Lambda=G_A.
\]
The global nonvanishing in Lemma~\ref{lem:quartic-reciprocal-nonzero}
forces $\Lambda,p_H,B_A,p_A\ne0$. Hence $C_A$ is the ratio of the beta
and common approximate logarithms, with their two minus signs cancelled.
The beta residual is $C_A-2\beta$. In full comparisons $G_H$ has order
$40$ and $G_A$ order $12$; in beta-only comparisons both have order $12$.
The Lucas numerators use the normalized equations below, whose permitted
zero-ordinate branches must not be excluded by this beta argument.

\subsection{Selective comparisons and the local elliptic chart}

\begin{lemma}[Three-mode coverage]\label{lem:quartic-three-mode-coverage}
Number the unchanged raw contacts from $0$ to $36$. Retain full comparisons
at $0,\ldots,4$ and $21,\ldots,28$, beta-only comparisons at
$5,\ldots,12$, and Lucas-only comparisons at $13,\ldots,20$ and
$29,\ldots,36$. The contacts carrying Lucas information cover the three
points $S,T,Z$, and those carrying beta information cover $A,T$.
\end{lemma}
\begin{proof}
The first five positions are shared. Each of $A,S,T,Z$ has its original
block of eight finite positions, in that order. Thus the Lucas set contains
the five shared positions and all three finite blocks of $S,T,Z$; the beta
set contains the five shared positions and the two finite blocks of $A,T$.
The unchanged two-slope rational-height argument applies separately to
these sets. No position is deleted when two slopes, raw values, or assigned
primes coincide. The two maxima used below are consequently
$B_L=\max(|s|,|t|,|z|)$ and $B_{AT}=\max(|n|,|t|)$.
\end{proof}

\begin{lemma}[Doubled formal comparisons]
\label{lem:quartic-doubled-private-comparison}
At a split good prime above the cutoff, if $v(z(R))>0$, then either $R$ or
$R+T_0$ is in the formal group. Whenever the relevant rational function is
defined,
\[
 v(\LogE([2]R))=v(z(R)),
\]
with infinite valuation for $R=O$ or $T_0$.
\end{lemma}
\begin{proof}
For an affine nontorsion point,
$z(R)^2=x(R)/(x(R)^2-D)$. If $v(x)=0$, the right side has nonpositive
valuation. If $v(x)<0$, the point is formal. If $v(x)>0$, translation by
$T_0$ gives $x'=-D/x$ and $z'=-z$, so the translated point is formal with
the same parameter valuation. Doubling kills $T_0$ and is an isometry on
the formal group at an odd good prime. The asserted equality follows from
the retained logarithm isometry. The two other nonzero two-torsion points
are poles of $z$ and are not in the domain.
\end{proof}

Private primary points may therefore lie in $E(K)$ rather than $2E(K)$.
If $[2]P=[2\alpha(P)]P_*$, their coefficients $\alpha(P)$ may be
half-integral. They are integral at each consumed odd prime. One compares
logarithms only after writing an actual integer-multiple identity between
formal points; torsion is not discarded by equating these half-indices.
The primary $f$ tests still put the primary points themselves in the
formal group. An alternate test of
$z(P-Q)/z(J)$, together with the opposite-place $f(J)$ test, implies
$v(z(P-Q))>0$, and the lemma supplies the needed doubled logarithm.

\begin{lemma}[Joint primary and difference graph]
\label{lem:quartic-joint-chart}
Let $Q=(a,aq)$ be nontorsion and $\jmath\ne0$. Set
\[
 F_Q=1/a,\quad c=a+DF_Q,\quad
 r=\jmath/q,\quad s_Q=c/q,\quad t_Q=1/q^2,\quad u=rs_Q.
\]
The three field equations
\begin{align}
 k&=(A-r)p,\\
 kp+up-2Drf-4Dt_QA&=0,\\
 2D(A+r)f+s_Qk+4Dt_QA+\jmath q&=0
\end{align}
give the exact graph $f=1/x(P)$, $p=y(P)/x(P)$,
$A=\jmath z(P-Q)$ for affine primary points $P$ with $x(P)\ne0$
and finite $z(P-Q)$. In particular this primary chart excludes $O$ and
$T_0$. The equations permit $A=0$ and do not
require inverses of $A,k$ or of a difference parameter.
\end{lemma}
\begin{proof}
Set $x=p^2+Df$. The target identity is $c^2-q^4=4D$.
Reversible linear combinations of the displayed equations give
\[
 A U=\jmath H_2,\qquad A V=\jmath H_1,
\]
where
\[
 U=q(x+Df)+cp,\quad V=q^2(x-Df)-4D,\quad
 H_1=q(x+Df)-cp,\quad H_2=x-Df-q^2.
\]
Direct expansion yields
\[
 H_2V-H_1U=-4Dq^2(fp^2+Df^2-1).
\]
Thus $xf=1$ and $P=(x,xp)$ is on $E$. If $U=V=0$, then for
$v\in\{d,-d\}$ one has
$p=-2v/q$ and $x=v(a+v)/(a-v)$, precisely $P=Q+(v,0)$.
At either point $H_2=(4D-q^4)/q^2\ne0$: equality $q^4=4D$ would imply
$N_{K/\Q}(q)^2=28$, impossible for a rational norm. These genuine poles
cannot satisfy the graph equations. Else $U$ or $V$ is nonzero and
$A/\jmath$ is uniquely determined. The chord identities
\[
 H U=N H_2,\qquad H V=N H_1,
\quad H=(p-q)(1+DfF_Q),\quad
 N=pq(1-DfF_Q)+2D(f+F_Q)
\]
identify it with $z(P-Q)$ on the chord domain and hence at every regular
point. At $P=Q$ and $Q+T_0$ the value is zero. Opposition is covered by
$V\ne0$. At $p=0$, the curve has $x=\pm d$, $U=2qx\ne0$, and the graph
still gives its unique regular value. Conversely substituting any point
and its regular difference value reverses these identities and recovers
all three equations with $k=(A-r)p$.
\end{proof}

The last $p=0$ cases are needed below. They are excluded at consumed places
by the positive primary-$f$ valuation, but need not be excluded globally
for every normalized numerator.

\subsection{Private conics, packed field tests and the ordinary pole}

Each contact $r$ has its own rational variables $\theta,s_c,S_c$ with
\[
 S_c=s_c^2,\qquad \theta^2-rS_c=2.
\]
In particular $r,s_c\ne0$. For a field value $F_i=a_i+b_i\sqrt2$ of
required order $e_i$, define
\[
 Y_i=a_i+\theta b_i,\qquad
 T_i=2a_iY_i-a_i^2+2b_i^2
     =Y_i^2-(\theta^2-2)b_i^2,\qquad
 r^{2e_i}W_i=T_i.
\]
The upper tree uses successive pairs $N_r(u,v)=u^2+rv^2$, starting with
the literal leaf $1$ followed by the $W_i$ in stable order, carrying an
unpaired final child unchanged. Denote its root by $W$.

\begin{lemma}[Local conic and tree deductions]
\label{lem:quartic-private-conic-packing}
Assume the ordinary pole gives no negative odd valuation to nonzero
$rW^2$ outside its fixed exceptional primes. At any $p>76176$ with odd
$R=v_p(r)$, the equations above imply $R>0$ and choose one root
$\rho\in\Q_p$ of $2$ such that
$v_p(a_i+\rho b_i)\ge e_i$ for every retained field.
\end{lemma}
\begin{proof}
The conic makes $\delta=v_p(\theta^2-2)=R+2v_p(s_c)$ odd.
It cannot be negative, since then $v_p(\theta^2)$ would be odd; it cannot
be zero. Thus $\delta>0$, $\theta$ is a unit, and one root $\rho$ satisfies
$v_p(\theta-\rho)=\delta$ while the other factor is a unit.
The summands in $N_r(u,v)$ have opposite valuation parity and cannot
cancel. The literal leaf proves $W\ne0$. If $R<0$, ascent along that
leaf gives $v_p(W)\le0$, contradicting the pole conclusion. Hence $R>0$.
The pole conclusion gives $v_p(W)>-R/2$. Descent through the tree uses
\[
 v_p(N_r(u,v))>-cR\ \Longrightarrow
 v_p(u)>-cR/2,\quad v_p(v)>-(c+1)R/2,
\]
with $0\le c<1$, and gives $v_p(W_i)>-R$. Consequently
$v_p(T_i)>(2e_i-1)R\ge2e_i-1$, so $v_p(T_i)\ge2e_i$.
The two summands of $T_i$ also have opposite parity, yielding
$v_p(Y_i)\ge e_i$ and $\delta+2v_p(b_i)\ge2e_i$.
It follows that $v_p((\theta-\rho)b_i)\ge e_i$ and hence the claimed
field order. Zero summands are treated with infinite valuation. No
integrality of $a_i,b_i$ and no global squareclass of $F_i$ was assumed.
\end{proof}

The proof uses $v_p(T_i)\ge2e_i$, not the generally unjustified stronger
bound $2e_iR$. It is independent of the elliptic meaning of the fields.

\begin{lemma}[Conic completeness]\label{lem:quartic-conic-completeness}
In the retained integral-input construction, each contact
$r=d_0^2p_0>0$ admits the needed private conic and arbitrarily deep
approximation to its chosen split branch, independently of the other
contact occurrences.
\end{lemma}
\begin{proof}
Here $p_0\equiv1\pmod8$ and $d_0\ne0$ is a $p_0$-unit. The ring
$\Z[\sqrt2]$ is norm Euclidean: subtracting nearest integers from two
coefficients leaves $|a|,|b|\le1/2$ and $|a^2-2b^2|\le1/2<1$.
The split prime ideal therefore gives integers $a,b$ with $a^2-2b^2=p_0$,
after multiplying by the norm-minus-one unit if needed. Since $b\ne0$,
$(\theta_0,s_0)=(a/b,1/(d_0b))$ is a rational conic point. The parametrization
\[
 \Delta=\frac{2(rs_0-\theta_0 t)}{t^2-r},\qquad
 s_c=s_0+\Delta,\quad \theta=\theta_0+t\Delta
\]
and rational approximation give simultaneous approximation at the finite
set of required places. There is no rational pole $t^2=r$.
Near the selected root, $\delta=1+2v_{p_0}(s_c)$ can be made arbitrarily
large. After all field coefficients are fixed, choose it so that
$\delta+2v_{p_0}(b_i)\ge2e_i$ for every $i$. The original field orders
then make all $W_i$ integral; the constant path makes $W$ a $p_0$-unit.
At the real place the tree is positive, even if some $W_i$ are negative.
The extra approximations used for the following pole may be imposed at
the same time. No conic variable is shared between occurrences.
\end{proof}

\begin{proposition}[Linear-unit ordinary pole]
\label{lem:quartic-linear-pole}
Set $g=3-2\theta$. Replace only the ordinary pole module by
\begin{gather*}
 \ell=gW,\quad z_\ell=\ell^2,\quad c_r=re,\qquad
 p_r^2-17r-34c_r+51ec_r=0,\\
 w=17e-c_rz_\ell,\quad j_s=s_r^2,\quad k_r=rq_r,\\
 j_s^2+34j_sw-867w^2-17q_rk_r=0.
\end{gather*}
It has the one-sided pole conclusion used in
Lemma~\ref{lem:quartic-private-conic-packing}. It has integral-input
completeness under the same inherited realization and rank hypotheses.
The scalar-height pole module is unchanged.
\end{proposition}
\begin{proof}
At an odd prime with odd $v_p(r)$, the conic argument makes $\theta$
integral, and
$(3-2\theta)(3+2\theta)=1-4rs_c^2$ is a unit. Thus $g$ is a unit.
Suppose $b=rW^2\ne0$ has negative odd valuation outside $2,3,17$.
The first conic $p_r^2=17r(1-e)(1+3e)$ forces $e$ to be a unit: positive,
negative or infinite valuation of $e$ would make the right side have odd
valuation. Since $v_p(\ell)=v_p(W)$, one has
$v_p(w)=v_p(b)<0$, an odd integer. In
\[
 (s_r^2-17w)(s_r^2+51w)=17rq_r^2
\]
the left side is nonzero with even valuation, since the valuations of
$s_r^2$ and $w$ differ in parity. The right side is zero or has odd
valuation. This contradiction proves the asserted exclusion, including
zero ordinate branches of the first conic. If $W=0$, the conclusion has
no nonzero pole input; explicit witnesses are
$\ell=z_\ell=p_r=q_r=k_r=0$, $e=1$, $c_r=r$, $w=17$, $s_r=17$, $j_s=289$.

For completeness retain the generalized fixed-$U$ realization argument
under its original hypotheses. For $b>0$ with square reduced denominator
and square in $\Q_2,\Q_3,\Q_{17}$, it applies to a prescribed $U_0$ when
\[
 U_0^2/(17b)<1\text{ at infinity},\quad
 v_2(U_0^2/(17b))\ge3,\quad
 v_3(U_0^2/(17b)),v_{17}(U_0^2/(17b))\ge1,
\]
and $v_p(U_0)=v_p(b)$ at every other odd-support prime of $b$.
Indeed, put $B=17b$ and $V=B(B-U_0^2)$, and let
$\beta_{\rm sf},\eta_{\rm sf}$ be their positive squarefree representatives.
The displayed conditions make $V$ positive and a square in
$\Q_2,\Q_3,\Q_{17}$. Consequently
\[
 \beta_{\rm sf}>1,\quad \beta_{\rm sf}\equiv17\pmod{24},\quad
 \eta_{\rm sf}\equiv1\pmod{24},\quad 17\nmid\eta_{\rm sf}.
\]
At an odd-support prime of $B$ outside $2,3,17$, the square reduced
denominator makes $v_p(b)$ positive, and the hypothesis on $U_0$ gives
$v_p(U_0^2/B)>0$. Thus $1-U_0^2/B$ is a square unit, so
$\eta_{\rm sf}$ is a square in $\Q_p$ and $p\nmid\eta_{\rm sf}$.
The same conclusion at $17$ follows from the displayed smallness condition.
Conversely, at a prime dividing $\eta_{\rm sf}$ the necessary strict
cancellation in $B-U_0^2$ gives $v_p(B)=2v_p(U_0)$ and makes
$B/U_0^2$ a square unit. Thus $\beta_{\rm sf}$ is a square in $\Q_p$.
In particular $\gcd(\beta_{\rm sf},\eta_{\rm sf})=1$, and both
mutual-square conditions hold. If $U_0=0$, then $V=B^2$ and
$\eta_{\rm sf}=1$, so the converse condition is empty. These are precisely
the hypotheses consumed by the inherited simultaneous-realization proof;
no new rank or converse statement is substituted.

Now set $b=rW^2$ and $U_0=rW^2g$. In the complete construction
$r=d_0^2p_0$ and $W$ is a $p_0$-unit, so $v_{p_0}(b)=1$ and all
other valuations of $b$ are even. Hence its reduced denominator is a
square. The retained local squares of $r$ at $2,3,17$ also make $b$ a
square at those places. After the fields are fixed, approximate
the private conic also near $(\theta,s_c)=(3/2,1/(2\sqrt r))$ at infinity,
$2,3,17$. These local points exist by the original contact construction.
The continuous function $(r/17)W(\theta)^2g(\theta)^2$ vanishes there,
so it satisfies the displayed smallness conditions while the selected
prime still makes $W$ a unit. The local-unit argument gives the required
odd-support valuations of $U_0$. The resulting old three-conic witnesses satisfy
\[
 p_o^2=17b(1-e)(1+3e),\quad h_o=e(17b-U_0^2),\quad
 (s_o^2-17bh_o)(s_o^2+51bh_o)=17bq_o^2.
\]
Here $r,W,b$ are nonzero. Set
\[
 p_r=p_o/W,\quad w=h_o/b,\quad
 s_r=s_o/b,\quad q_r=q_o/(r^2W^3).
\]
These substitutions give $w=e(17-r\ell^2)$ and the final conic in the
displayed normalized system; the other coordinates are its definitions.
Thus the witness conversion uses only stated nonzero quantities.
Equivalently the unnormalized form
has $L=r\ell$, $z_L=r^2z_\ell$, $w_L=r^2w$, $s_L=rs_r$,
$q_L=r^2q_r$, $j_L=r^2j_s$, $k_L=r^2k_r$, with $e,p_r,c_r$ unchanged.
All divisions in the reverse scaling are powers of the conic-proved
nonzero $r$. This also proves that normalization loses no zero branch.
\end{proof}

\begin{lemma}[Four-witness anchor]\label{lem:quartic-anchor-four}
For rational $a,b$, both belong to $\Z_{(2)}\cap\Z_{(3)}$ if and only if
$9+3(a^2+b^2)$ is a norm-one trace for the quaternion algebra $(2,3)$.
Consequently the original two-coefficient anchor is equivalent to
\[
 CH-9C^2-3f_0^2-3f_1^2=0,\qquad
 H^2-8U^2-12V^2+24W_a^2-4C^2=0.
\]
\end{lemma}
\begin{proof}
By Definition 4, Observation 5 and Facts (a), (b), (c), (e) in the proof of
Proposition 6 of \cite{koenigsmann-defining}, the algebra is ramified
exactly at $2,3$, the trace set is unrestricted at the other places, and it
is integral at its two ramified places
$2,3$, contains integral traces with residues $1$ at $2$ and $0$ at $3$,
and is determined by these local trace conditions. These are the explicit
trace facts used here, in addition to the older sum-of-traces criterion.
If $m=\min(v_p(a),v_p(b))<0$, then at $3$ the sum of squares has valuation
$2m$, while at $2$ it has valuation $2m$ for unequal valuations and
$2m+1$ for equal finite valuations. Thus $9+3(a^2+b^2)$ would be
nonintegral, proving one implication, with zeros covered by infinite
valuation. Conversely integral $a,b$ give a trace congruent to zero at
$3$. At $2$ it is odd or belongs to $4\Z_2$. For the latter case put
$x_0=h/2\in2\Z_2$, where $h$ is this trace. If $x_0\equiv2\pmod4$, then
$(x_0^2-1)/3\equiv1\pmod8$ is a square, giving
$x_0^2-3v^2=1$. If $x_0\equiv0\pmod4$, then $1-x_0^2/6$ is a square,
and $u=v=1$, $w^2=1-x_0^2/6$ gives
$x_0^2-2u^2-3v^2+6w^2=1$. The local trace criterion supplies rational
witnesses. Finally divide the two homogeneous equations by $C^2$,
which is allowed because the retained constant $C$ is nonzero. This is
an equality of existential projections, not a rational formula choosing
quaternion witnesses.
\end{proof}

\subsection{Local and normalized Lucas tests}

Write
\[
 \mathcal F(s,t,z)=t^2-s(s-2)tz+(s-2)^2z^2-1.
\]
\begin{lemma}[Local Lucas residual]\label{lem:quartic-local-lucas}
The primary and alternate logarithm comparisons, before any proxy equality
is used, give local integral quotients $C_s,C_t,C_z$ congruent to the
integer root indices modulo $p^{40}$. An order-$40$ test of
$\mathcal F(C_s,C_t,C_z)$ implies $p^{40}\mid\mathcal F(s,t,z)$.
\end{lemma}
\begin{proof}
The implication follows by substitution into the polynomial with integer
coefficients. One may evaluate the field residual using
\[
 V=(C_s-2)C_z,\quad D_L=C_sC_t-V,\quad
 E_L=C_t^2-VD_L-1.
\]
These are quadratic equations, and $E_L=\mathcal F(C_s,C_t,C_z)$.
They replace the three proxy comparisons and the global proxy equation.
They do not impose $E_L=0$ as a field equality. Removal of the proxy
coordinates requires their exclusive incidence in that equation and those
comparison arguments. Integer completeness and the reverse input
implication use the two height contradictions below; no conversion of
arbitrary old proxy witnesses is asserted.
\end{proof}

\begin{lemma}[Normalized local Lucas evaluation]
\label{lem:quartic-normalized-lucas}
Let $\Lambda p_H=G_H$ be the common negative-logarithm equation. For the
three Lucas numerators replace their quotient blocks and the preceding
evaluation by
\begin{gather*}
 B_i\Lambda=G_i\quad(i=s,t,z),\qquad P B_s=p_s,\\
 (U+V)p_t=P B_t,\qquad K_Lp_t=(1-2P)B_t,\\
 Vp_z=(1-2P)B_z,\qquad E_L=U^2-K_LV-P^2.
\end{gather*}
An order-$40$ test of $E_L$ has the same required integer consequence.
Every integer input still admits the changed witnesses.
\end{lemma}
\begin{proof}
The global nonvanishing of the retained $G_i$ implies
$\Lambda,p_H,B_s,B_t,B_z\ne0$. Also $p_t\ne0$: otherwise the two
equations with left factor $p_t$ would force both $P=0$ and $1-2P=0$.
At a consumed place all primary points are formal, so their $p_i$ are
nonzero. Put $C_i=B_i/p_i$. Then $P=1/C_s$ and direct substitution gives
\[
 E_L=P^2\mathcal F(C_s,C_t,C_z).
\]
Since $C_s$ is locally integral, the implication
$v(E_L)\ge40\Rightarrow v(\mathcal F(C_s,C_t,C_z))\ge40$ is valid even
when $C_s$ is not a unit. No local denominator has been discarded.

For completeness the integer $s\ne0$ and all root points are fixed before
the prime-pattern theorem is used. Its prescribed finite-exclusion version
allows every chosen prime to avoid the divisors of $s$. Then $C_s$ is a
unit at the consumed place. The old order-$40$ residual implies the new
one, and the displayed equations assign the new witnesses. This imposes
no compatibility between group orders at different primes.

Globally $p_s$ or $p_z$ may vanish at an occurrence not being consumed in
the local argument. Such points have $x=\pm d$ and are covered by the
regular extension in Lemma~\ref{lem:quartic-joint-chart}; they must not be
excluded by reusing an earlier blanket nonzero assertion. At a consumed
place $C_s=2$ gives $P=1/2$ and $V=K_L=0$, since both $p_t$ and $p_z$
are nonzero there. No division by $1-2P$ is needed. Globally the branch
$P=1/2$, $p_z=0$ forces $K_L=0$ but need not force $V=0$; the latter
value may be chosen in the completeness construction.
\end{proof}

The retained field lists can now be stated without proxy coordinates.
In a full occurrence there are the five primary $f$ values, the common
$\overline{f(J)}$, the beta residual, $E_L$, and five alternates.
In a beta-only occurrence there are two primary $f$ values, the common
field, the beta residual and two alternates. In a Lucas-only occurrence
there are four primary $f$ values, the common field, $E_L$ and four
alternates. These lists have lengths $13,6,10$. The formal fields all have
order one. The beta residual and beta alternate have order $12$; the
Lucas residual and Lucas alternates have order $40$. The denominator
alternate has order $e_0=40$ in full and Lucas-only occurrences, and
$e_0=12$ in beta-only occurrences. Its $G$ precision is the same $e_0$.

\begin{theorem}[The two height contradictions]
\label{thm:quartic-selective-arithmetic}
Retain the inherited scalar and contact-height conclusions, and the integer
root indices $n,s,t,z$ of the chosen carrier, with $s,z\ne0$. Let
$\beta=X^2+2$ and $\beta'=2\beta$. Suppose the
local tests above give $C_i\equiv n_i\pmod{p^{e_i}}$ at all consumed
contacts, with beta residual $C_A-\beta'$. Put $H_*=5H$. Suppose
\[
 n_Q^2-6<7H_*N^{19/e},\qquad
 h(\beta')<a/40,\qquad \log(40H_*)<a/40,
 \qquad a=|s|>(24H_*)^{10},
\]
where the first inequality is the contact-height bound whenever the
appropriate radical powers divide $N$. Then the changed arithmetic source
has the same integer-input conclusion as Version 3.
\end{theorem}
\begin{proof}
If $\Delta=\mathcal F(s,t,z)\ne0$, every relevant order-$40$ radical
power divides $|\Delta|$. By three-mode coverage, all three heights in
$B_L=\max(|s|,|t|,|z|)$ satisfy the assumed inequality. Since
$|\Delta|<5B_L^4$, the same elementary estimate as in Version 3 gives
\[
 B_L^{2-76/40}<24H_*.
\]
This contradicts $B_L\ge a>(24H_*)^{10}$. Therefore $\Delta=0$.
The displayed lower bound on $a$ gives $|s|\ge7$, and $z\ne0$ was
retained. The Lucas lemma therefore gives the stronger individual conclusion
$|t|>\exp(a-3)$, not merely a bound on $\max(|t|,|z|)$.

If $\beta'\ne n$, put $N=|\operatorname{num}(\beta'-n)|$.
The beta contacts cover $A,T$, so with $B=\max(|n|,|t|)$ one has
\[
 B^2<7H_*N^{19/12}+6,\qquad
 N\le2e^{h(\beta')}B.
\]
Thus $B^{2-19/12}<40H_*e^{(19/12)h(\beta')}$ and
\[
 \log B<\frac{12+19}{40(24-19)}a=\frac{31}{200}a<a-3,
\]
contrary to $B\ge|t|$. Hence $\beta'=n$ is an integer. $\beta'=2\beta=2(X^2+2)$;
writing $X=u/v$ in lowest terms gives $v^2\mid2$, hence again $v=1$.
\end{proof}


\subsection{General root carrier}
\label{sec:quartic-general-root-variant}

\begin{lemma}[Integer doubled root indices]
\label{lem:quartic-doubled-indices}
Let $R=(1+\sqrt2,2\sqrt2)$, so direct duplication gives $2R=P_*$.
Every $Q\in E(K)$ has a unique integer $n_Q$ with $2Q=n_QP_*$.
It is nonzero exactly when $Q$ is nontorsion. For such a point,
\[
 h(x(Q))>n_Q^2/4-4.
\]
\end{lemma}
\begin{proof}
The inherited generator assertion gives $2Q=n_QP_*$. The torsion
assertion makes its kernel exactly $E[2](K)$, so equivalently
$Q=n_QR+\tau$, $\tau\in E[2](K)$. Uniqueness follows from the infinite
order of $P_*$. The unchanged height bounds give
\[
 n_Q^2-6\le h(x(2Q))\le4h(x(Q))+\log676.
\]
Since $(6+\log676)/4<4$, the stated inequality follows.
\end{proof}

The root points $A,S,T,Z,J_0,J$ may now be any nontorsion points of $E(K)$.
Their indices always mean the integer doubled indices of the lemma. They
are distinct from the half-integer coefficients used for private points
in the earlier carrier convention.

\begin{lemma}[Small general-root charts]
\label{lem:quartic-general-root-charts}
For $A,S,T,Z$ retain $x,y$ and the actual first-doubling square $q=x^2$,
and impose $y^2=xq-Dx$. For $S$ also retain $xf=1$, $yz+x=0$.
Keep the nonzero tangent denominator, or the direct tangent
equation of Lemma~\ref{lem:quartic-tangent-fusion}. For each of $J_0,J$
use $f,z,k$ with
\[
 k=fz,\qquad f=z^2-Dk^2,
\]
with $z\ne0$ enforced by its retained inverse. These charts reconstruct exactly
the required nontorsion root points.
\end{lemma}
\begin{proof}
The first equations are the curve equation with a reused square. The
nonzero tangent denominator excludes all points with $y=0$, hence all
torsion. The two extra S equations preserve its reciprocal and parameter.
In the other chart, $z\ne0$ implies $f\ne0$: $f=0$ would force $k=0$
and then $z=0$. Set $x=1/f$, $p=-1/z$, $y=xp$. Dividing the second
equation by $z^2$ gives $fp^2=1-Df^2$, so $P=(x,y)$ lies on the curve
and $xy\ne0$. Conversely every nontorsion point gives these witnesses.
These are rational maps on the stated nonzero domains. The new equations
cannot justify deleting an unrelated guard without its incidence proof.
\end{proof}

For clarity, the inherited scalar argument still applies to the new
indices. At each odd contact prime of $\beta$ outside $S_*$, the
constant-left scalar norm tree excludes negative odd raw valuation and
gives all six assigned coefficient orders. The first four put $S,J_0$
in the formal group at every place above that prime. Writing
$2S=sP_*$ and $2J_0=j_0P_*$, with nonzero integer $j_0$, the last two
give $v_p(s/j_0)\ge32$, hence $v_p(s)\ge32$. The logarithm comparison
uses the exact identity $[2j_0]S=[2s]J_0$. It is valid also at inert
primes: both rational coefficients are tested and the extension is
unramified outside two. The unchanged scalar raw polynomial and its
small-prime correction detect every prime of $M_5(\beta)$, including
an odd infinite contact. Consequently
\[
 M_5(\beta)^{32}\mid s,\qquad
 h(\beta)\le H M_5(\beta)^{19}\le H|s|^{19/32}.
\]
The height theorem's domain condition holds. The value $\beta=X^2+2$
is none of $0,1,-343/32$, and $\beta=2401$ would require the nonsquare
integer $2399$ to be a rational square. Doubling both indices therefore
changes neither the scalar ratio nor its divisibility argument.

Take $H_*=5H$ and $\beta'=2\beta$. The unchanged contact-height bound
$h(x(Q))<7HN^{19/e}$ and the preceding lemma give
\[
 n_Q^2-6<28HN^{19/e}+7<35HN^{19/e}=7H_*N^{19/e}.
\]
The scalar ratio is still $s/j_0$, since doubling both indices cancels
the factor two. Its order-$32$ tests give
$h(\beta)\le H|s|^{19/32}$. The full anchor at both $2$ and $3$ is
retained. Its $3$-adic consequence, using $2S=sP_*$, is
\[
 s\in2\cdot3^{2^{24}-1}\Z\setminus\{0\},\qquad
 \log_2|s|>25165822.5.
\]
Here formal doubling at $3$ is an isometry and the retained calculation
is $v_3(f(sP_*))=2+2v_3(s)$. The $2$-adic part of the anchor is not
replaced by this consequence. Since $\log_2H_*<2093969$, one has
\[
 |s|>(24H_*)^{10},\qquad |s|>(40H_*)^{32/13},\qquad
 h(2\beta)<|s|/40,\quad\log(40H_*)<|s|/40.
\]
For example the first two required binary cutoffs are respectively
$10\cdot2093974$ and $(32/13)\cdot2093975$, both below the displayed
anchor lower bound. Thus Theorem~\ref{thm:quartic-selective-arithmetic}
applies with unchanged orders. Both beta-bearing comparators must now
test $C_A-2\beta$; the scalar definition $\beta=X^2+2$ and its raw scalar
polynomial remain unchanged.


\subsubsection{Doubled targets}
\label{sec:quartic-target-two-variant}

\begin{lemma}[Targets $2Q_i$ and $P_*$]
\label{lem:quartic-target-two}
The general-root construction may use targets $2Q_i=n_iP_*$ and fixed
denominator target $P_*=2R$, requiring only one doubling of each variable
root. It preserves the local integer congruences and integer completeness.
\end{lemma}
\begin{proof}
For soundness write $2X_0=hP_*$, $2X_i=m_iP_*$ and $2J=j_dP_*$.
The root $J$ is nontorsion, so $j_d\in\Z\setminus\{0\}$.
The doubled differences have indices $h-2$ and $m_i-2n_i$.
Lemma~\ref{lem:quartic-doubled-private-comparison} and the alternate
tests give
\[
 v_p((h-2)/j_d)\ge e_0,\qquad
 v_p((m_i-2n_i)/j_d)\ge e_i,
\]
where $e_0=40$ in full/Lucas mode and $e_0=12$ in beta-only mode.
Since $j_d$ is integral, $h\equiv2$ and $m_i\equiv2n_i$ at those orders.
In particular $h$ is a unit at the odd prime. The exact identity
$[2h]X_i=[2m_i]X_0$ permits logarithm comparison on the primary formal
group, giving ratio $m_i/h$. The retained relative errors, with
$e_i\le e_0$, give $C_i\equiv n_i\pmod{p^{e_i}}$.

For completeness let $N_{p,1},N_{p,2}$ be the two full finite group
cardinalities supplied by the inherited ordinary-prime pattern. It gives
$\gcd(N_{p,1},N_{p,2})=4$ and $p\nmid N_{p,1}N_{p,2}$. Each reduction
has exactly four rational two-torsion points, so its doubling image has
order $M_{p,i}=N_{p,i}/4$. Consequently $M_{p,1}$ and $M_{p,2}$ are
coprime and prime to $p$. Every reduction of a doubled global point is
killed by the corresponding $M_{p,i}$; no global reduction-surjectivity
or cyclicity claim is needed.
Put $j_c=\operatorname{lcm}_p M_{p,2}$ and $J=j_cP_*$, whose doubled
index is $j_d=2j_c$. Separately for each contact choose
\[
 k\equiv0\pmod{M_{p,1}},\quad
 k\equiv1\pmod{M_{p,2}},\quad
 k\equiv1\pmod{p^{40+v_p(j_c)}},\qquad k\notin\{0,1,-1\}.
\]
The three moduli are compatible. Set $X_0=kP_*$ and $X_i=kn_iP_*$.
The primary points are formal at the first place. Their differences from
the chosen targets, and $J$, are formal at the second. The alternate
logarithm ratios are $(k-1)/j_c$ and $(k-1)n_i/j_c$, while the primary
ratios are $n_i$. This supplies every required order. The finite exclusions
avoid zero points and the old open-chart exceptional cases. There is no
simultaneous CRT condition on different contact occurrences.
\end{proof}

The fixed target can be specified directly by
\[
 \begin{array}{c|cc}
 &1&\sqrt2\\\hline
 x(P_*)&57/8&-7/2\\
 q=y(P_*)/x(P_*)&-87/41&129/164\\
 F_Q=1/x(P_*)&456/1681&224/1681\\
 s_Q=(x(P_*)+DF_Q)/q&-55123/14637&19301/58548\\
 t_Q=1/q^2&7624/14161&39904/127449
 \end{array}
\]
Direct substitution verifies the curve and target identities. The source
stores $x,q,F_Q$ and the conjugates of $s_Q,t_Q$. Its remaining coefficients
are affine in $\jmath$:
\[
 r=\jmath\left(-\frac{232}{357}+\frac{86}{357}\sqrt2\right),\qquad
 u=\jmath\left(\frac{291673}{127449}-\frac{29440}{42483}\sqrt2\right),
\]
when $\jmath=1/\overline{z(J)}$. The ten constant bindings and four
affine bindings replace the fixed derivation, subject to full incidence
checks. Their isolated coefficient calculation does not bound the complete
sum of squares.



\subsection{Exact implementation refinements}


\begin{lemma}[Direct tangent equation]\label{lem:quartic-tangent-fusion}
For a curve point with $q=x^2$, replace the inverse and slope equations
$(2y)h=1$, $\lambda=(3q-D)h$ by
$2y\lambda=3q-D$. This preserves all retained rational coordinates.
\end{lemma}
\begin{proof}
The new equation and the curve force $y\ne0$. If $y=0$, then $x=0$ or
$x^2=D$, and $3q-D$ is respectively $-D$ or $2D$, a contradiction.
The unique inverse is restored by $h=(2y)^{-1}$. There is also the
polynomial restoration
\[
 h=\frac{2\lambda(3q-2D)-9xy}{4D^2},
\]
because $(3q-D)(3q-2D)-9x(x^3-Dx)=2D^2$.
For a retained second doubling its input is on the curve: with
$t=X-x$, $Y_0=-y-\lambda t$, the identity
\begin{align*}
 Y^2-X^3+DX={}&y^2-x^3+Dx+t(2y\lambda-3x^2+D)\\
 &+t^2(\lambda^2-X-2x)+(Y-Y_0)(Y+Y_0)
\end{align*}
proves curve preservation from the first doubling equations. Target-two
sources have only the first four variable-root doublings; no second set
of savings or equations is present in that branch.
\end{proof}





\begin{lemma}[Shared field-square evaluation]
\label{lem:quartic-shared-field-square}
For the common field $F=a+b\sqrt2=\overline{f(J)}$, bind once
$c=a^2+2b^2$, $d_s=2ab$. In each designated order-one packed field replace
\[
 Y=a+\theta b,\qquad PW=2aY-a^2+2b^2
\]
by $PW=c+\theta d_s$, with the same $P=r^2$. This is an exact projection
preserving $W$, every higher norm node and the pole witnesses.
\end{lemma}
\begin{proof}
Substitution gives $2a(a+\theta b)-a^2+2b^2=c+\theta d_s$.
The inverse map restores $Y=a+\theta b$, without division. All zero
cases are included. The two coefficients describe $F^2$, but their
packed use still enforces the original order-one condition on $F$;
they must not be relabelled as a new order-one test of $F^2$. The same
identity permits reuse of an already available primary-$f$ square in
its own packed numerator, after checking that square's defining rows.
\end{proof}



\begin{lemma}[First packed denominator and norm fusion]
\label{lem:quartic-first-norm-fusion}
For the first order-one primary field only, replace its virtual denominator
$W_0=T_0/r^2$ by $W_0=T_0/r$, and its first norm $N=1+rW_0^2$ by
\[
 T_0=2a_0Y_0-a_0^2+2b_0^2,\qquad rN-r-T_0^2=0.
\]
Keep every other denominator, the conic, all higher tree nodes and the
ordinary pole. The arithmetic input projection is unchanged.
\end{lemma}
\begin{proof}
The earlier tree proof still gives $v_p(W_0)>-R$, so $v_p(T_0)>0$.
More generally $v_p(T)\ge2e-1$ suffices for an order-$e$ field test:
opposite parity gives $v_p(Y)\ge e$ and
$\delta+2v_p(b)\ge2e-1$, hence
$\delta+v_p(b)\ge e+(\delta-1)/2\ge e$.
Thus the first field still has order one, and every other deduction is
unchanged. For completeness the old depth-two choice has $v_{p_0}(T_0)\ge2$
and $R=1$, so the new virtual leaf is integral and the root remains a unit.
For fixed nonzero $r$, all new tree values remain rational polynomials in
$\theta$, allowing the same exceptional-place pole approximation.

After that arithmetic change, the displayed fusion is exact: set
$T_0=rW_0$ forward and $W_0=T_0/r$ backward. The retained conic proves
$r\ne0$. These maps preserve $N$ and all later nodes, including $N=0$
or $T_0=0$ and negative rational $r$. The unguarded case $r=0$ would be
invalid and is not admitted. The old $r^2$ coordinate remains for other
leaves. The changed field is primary $f(H)$, not the separately shared
common field $\overline{f(J)}$.
\end{proof}


\subsection{Composition, completeness and finite realization}
\label{sec:quartic-composition}

For integral input first set $\beta=X^2+2$ and
take $A=\beta P_*$, whose doubled index is $2\beta$. Choose $J_0,S$ together in the old subgroup with the
required scalar and full anchor divisibilities. Their doubled-index ratio is unchanged. The Lucas existence lemma
then gives $t,z$ for the now fixed $s$, and take $T=tR$,
$Z=zR$. Fix the eight abscissa coefficients and denominator data before
using the prime-pattern theorem. Exclude the divisors of the fixed nonzero
$s$ as well as every earlier finite exceptional set.

Next choose one common comparison point $J$ using the stated least common
multiple, then choose the private points separately for each occurrence,
using precisely the target-two CRT rule.
Construct all joint-graph, reciprocal-$G$, normalized-ratio and
shared-square fields. Only then choose each private conic point and its
norm/pole witnesses by simultaneous approximation. Repeated primes retain
their common orientation but separate conic and private-point witnesses.
The constructions above give every required equation. In the reverse
direction the conic and pole decode field orders at every odd-contact
prime, without an ordinariness assumption. The graph and doubled logarithm
arguments give integer-index congruences, and the two height contradictions
give the integer input. The exact fusions preserve the specified
retained coordinates; arithmetic source replacements use the full
soundness and completeness argument rather than a claim about arbitrary
old private witness tuples.

Finally each selected source must be compiled afresh into rational
quadratic residual equations with the original input coordinate retained.
For real, hence rational, values the all-weights-one sum of their squares
vanishes exactly when every residual vanishes. Clearing nonzero rational
row denominators preserves those zero sets and gives an integral polynomial
of degree at most four.
This elementary last step does not prove the inherited arithmetic inputs.
The exact source, complete flattening map, allowed algebraic transformations,
and full coefficient collection must be checked for the single selected
endpoint. None of the local identities, allocation counts or isolated
coefficient bounds in this appendix is an assertion of that final check.
\begin{theorem}[The selected arithmetic source]
\label{thm:quartic-selected-source}
With precisely the inherited arithmetic hypotheses and source-version
qualifications stated above, the source consisting of the selected
equations has rational input projection $\Z$. Its input coordinate is the
original rational number $X$.
\end{theorem}
\begin{proof}
First suppose that rational witnesses satisfy the source. At every prime
above the cutoff with an odd raw-contact valuation, the private conic,
linear-unit pole and constant-left norm tree give the required field
orders, by Lemmas~\ref{lem:quartic-private-conic-packing} and
\ref{lem:quartic-first-norm-fusion}. This argument starts at arbitrary
odd-contact primes and does not assume that they are ordinary. The joint
graph supplies actual primary points and their regular difference
parameters. Positive primary-$f$ valuations put the primary points in the
formal group; the common conjugate-$f$ test and each alternate test allow
Lemma~\ref{lem:quartic-doubled-private-comparison} to be applied to the
doubled difference. Lemma~\ref{lem:quartic-target-two}, with denominator
order $12$ in beta-only mode and $40$ otherwise, gives the integer-index
congruences for each retained numerator. The beta quotient equations and
Lemma~\ref{lem:quartic-normalized-lucas} therefore give the precise local
tests used in Theorem~\ref{thm:quartic-selective-arithmetic}.

The retained scalar system gives its order-$32$ divisibility, while the
four-witness anchor is equivalent to the original two-place condition.
The general-root height and anchor estimates establish every cutoff of
that theorem with $H_*=5H$. The three-mode position sets cover exactly
the two maxima needed in its proof. It follows first that
$\mathcal F(s,t,z)=0$, then that $2\beta=n\in\Z$. Write $X=u/v$ in
lowest terms with $v>0$. Since $2X^2=2\beta-4$ is integral, $v^2$ divides
$2$; hence $v=1$ and $X\in\Z$.

Conversely fix any $X\in\Z$. Set $\beta=X^2+2$ and $A=\beta P_*$.
Choose the scalar pair $S,J_0$ in the retained subgroup using the
inherited scalar and full anchor construction, and choose $t,z$ by the
inherited Lucas existence lemma for the resulting nonzero integer doubled
index $s$. Set $T=tR$, $Z=zR$. All four roots and the two common roots
are nontorsion and admit the general-root charts. Fix their rational
abscissa coefficients and denominator data before applying
Lemma~\ref{lem:quartic-rational37}. Its fixed-system prime theorem permits
the finite exclusion of divisors of $s$, as required by the normalized
Lucas construction.

Choose one common $J$ by the least common multiple in
Lemma~\ref{lem:quartic-target-two}; at each occurrence choose its private
points using that lemma's three compatible congruences. No congruences
between different occurrences are imposed. The joint graph, exact
reciprocal evaluators, beta quotients and normalized Lucas equations then
have the stated witnesses and local orders. The explicit order of choices
in Section~\ref{sec:quartic-composition} now chooses each private conic
after its field coefficients are fixed. Conic approximation supplies the
integral packed leaves and a unit norm root; the fixed-parameter pole
realization supplies its remaining witnesses. This remains true with the
first denominator changed to $r$. The shared square and direct tangent
equations are exact identities with the reconstruction maps proved above.
Thus every equation of the selected source has a rational solution at
this input. This proves both implications without identifying arbitrary
old and new private witness tuples.
\end{proof}

The finite evaluations used here can themselves be written quadratically.
For precision $40$ bind four field pairs
\[
 w=Df^2,\qquad w_2=w^2,\qquad w_3=w_2w,\qquad U_6,
\]
with the fourth pair defined by the quadratic equation
\[
 U_6=\frac{a_3+a_4w+a_5w_2+
                  w_3(a_6+a_7w+a_8w_2+a_9w_3)}{a_6}.
\]
The quadratic output expression is
\[
 G_9(w)=a_0+a_1w+a_2w_2+a_6w_3U_6,
 \qquad a_6=-\frac{418}{16575}\ne0.
\]
For precision $12$ bind only $w=Df^2$ and use
$G_2(w)=1+a_1w+a_2w^2$. Substitute these expressions directly into
$\Lambda p_H=G_H$ and the beta or normalized-Lucas equations $B_i\Lambda=G_i$. Every product is quadratic in
the bound rational coordinates. Expansion gives each required $a_j$
coefficient exactly. The scale is one; no finite-error estimate is used
as a substitute for these polynomial identities.
\subsection{Exact source specification and audit boundary}
The source just proved has $4\,935$ rational coordinates,
including input zero, and $4\,513$ residual equations.
Its distinct name in \nolinkurl{quartics/quartic-4973/source-entry.json} is
\begin{center}\footnotesize\ttfamily
target2-tangent-first-norm-normalized-linear-pole-k40-q12-a25-h19-s32-4a82d9a8dc1b
\end{center}
The compressed source \nolinkurl{quartics/quartic-4973/source-residuals.jsonl.gz} has SHA-256
\begin{center}\footnotesize\ttfamily
94f8a1972cbfc0cf6493d418e628a18482307dec5c427db72e86d3eacb1a6f5d
\end{center}
Its uncompressed bytes have SHA-256
\begin{center}\footnotesize\ttfamily
4a82d9a8dc1b5c92eb40a49751c0a6aae073c8f765ca5e276a74802fd56b84be
\end{center}
The source contains $11$ modules. Its complete construction map is \nolinkurl{quartics/quartic-4973/construction-map.json}.
The independent source audit checks the full module equations, exact
ordered expression lists and contact positions, root and scalar interfaces,
all retained guards, the anchor terminal coordinate and complete flattening.
It is recorded in \nolinkurl{quartics/quartic-4973/source-audit.json}. The source-auditor SHA-256 is
\begin{center}\footnotesize\ttfamily
58bf1da4007247bc4558eacce6bdbbc83d2394ac2a47f128390f20f26567d50a
\end{center}
Those checks establish finite source identities. The source's rational
input projection is proved in Theorem~\ref{thm:quartic-selected-source},
with the inherited qualifications unchanged. A source digest is not an
arithmetic proof, and the historical base compiler remains construction
machinery rather than another final polynomial.

\subsection{Quadratic reductions preserving the rational input projection}
\label{subsec:quartic-reductions}

We describe the algebraic reductions applied to the new, fixed source.
All scalar coordinates range over $\Q$, and coordinate zero is the original
input $X$. Each field equation denotes its two rational coefficient
equations. These reductions prove equality of rational input projections
to that source; its arithmetic interpretation is
Theorem~\ref{thm:quartic-selected-source}.

\paragraph{Definitions, row operations, and interfaces.}
An equation $cv-F=0$, with $c\in\Q^*$ and $F$ independent of $v$,
permits substitution $v=F/c$ in every other equation and deletion of $v$.
Restriction is the forward solution map, and adjoining $v=F/c$ is its
inverse. Each use below retains degree at most two after substitution.
The initial quadratic lift uses the reverse operation: introduce targets
for products of two available coordinates until the source residuals are
quadratic. Evaluating these product definitions and forgetting their
targets give the two solution maps. A previously defined product can be
reused only with its retained defining equation.

An operation replacing $R_i$ by $aR_i+bR_j$, while retaining $R_j$,
has inverse $R_i=(R_i'-bR_j)/a$ when $a\ne0$ and $i\ne j$.
Thus every recorded linear row operation preserves the simultaneous zero
set. A shear replaces an old witness by $v=v'+cw$, with fixed rational
$c$ and another witness $w$ retained. Its forward and reverse coordinate
maps are $v'=v-cw$ and $v=v'+cw$. This is an invertible rational coordinate
change and preserves quadratic degree.

Every removed internal coordinate must have no occurrence outside the
replaced rows, unless the stated substitution is made in \emph{every}
such row. Retained coordinates keep their identities and values. All new
coordinates are fresh, and input coordinate zero is never a target of
deletion or redefinition. These conditions apply to the current system
after all preceding transformations.
\paragraph{Elimination modulo retained equations.}
For a recorded relation elimination, a current defining row has the form
$x-N(y)=0$ up to nonzero rational scale, with $N$ quadratic and independent
of the deleted witness $x$. For each other affected row $F_j(x,y)$ the
certificate gives an exact polynomial identity
\[
 F_j(N(y),y)=G_j(y)+\sum_k M_{jk}(y)H_k(y),
\]
up to the recorded nonzero row scales. Every $H_k$ is an unchanged,
$x$-free row, and every new row $G_j$ is at most quadratic. Forgetting
$x$ takes an old solution to a new one, since the $H_k$ vanish.
Conversely adjoin $x=N(y)$: the same identities recover every affected
$F_j=0$. Complete occurrence checks ensure that no other row uses $x$.
This is an identity over $\Q$, not a division by a possibly zero field
value. The retained support rows are never replaced in this operation.

\paragraph{The recorded contact-power circuits.}
The circuit replacements preserve the boundary values $r^2,r^{24},r^{80}$.
The replacement definitions are
\[
 r_2=r^2,\quad r_4=r_2^2,\quad r_8=r_4^2,\quad
 r_{16}=r_8^2,\quad r_{24}=r_{16}r_8,\quad
 r_{40}=r_{24}r_{16},\quad r_{80}=r_{40}^2.
\]
Their exact outputs agree with the old acyclic product definitions.
For either direction retain $r$ and the boundary outputs, and assign
each internal target its displayed power of $r$. These maps work also
for $r=0$ and use no inverse. Each deleted internal target has no use
outside the replaced definitions; every new internal target is fresh.
The certificate checks those conditions and the exact output identities
in the current system. Only circuit events with these retained powers
are included in this description.

\paragraph{A quadratic presentation of the parity block.}
This replacement concerns only its rational equations, with unchanged
input $b$. Its old private variables $t,U,e,p,q,s,h$ satisfy
\begin{align*}
 U(1+bt^2)&=b(1-bt^2),&
 p^2&=17b(1+2e-3e^2),\\
 h&=e(17b-U^2),&
 17bq^2&=(s^2-17bh)(s^2+51bh).
\end{align*}
Use instead twelve private rational variables
$u,v,a,z,e,p,c,w,s,q,j,k$ and nine quadratic equations
\begin{gather*}
 z=u^2,\quad a=bv,\quad z+va=b^2,\quad c=be,\\
 p^2-17b-34c+51ec=0,\qquad w-17bc+cz=0,\\
 j=s^2,\qquad k=bq,\qquad
 j^2+34jw-867w^2-17qk=0.
\end{gather*}
For the forward map put $\delta=1+bt^2$. It is nonzero: otherwise
$bt^2=-1$ and the first old equation gives $0=2b$, a contradiction.
Keep $e,p,s,q$ and set
\[
 u=U,\quad v=2bt/\delta,\quad a=bv,\quad z=U^2,\quad
 c=be,\quad w=bh,\quad j=s^2,\quad k=bq.
\]
The identity
$[b(1-bt^2)]^2+b(2bt)^2=b^2(1+bt^2)^2$ proves the first conic;
substitution proves the other equations.

Conversely $u^2+bv^2=b^2$. Set $h=e(17b-z)$ and keep $e,p,s,q$.
If $b=0$, take $U=t=0$. If $b\ne0$ and $u+b\ne0$, take
$U=u$ and $t=v/(u+b)$. In this case the first old residual, multiplied
by $(u+b)^2$, is
\[
 (u-b)(u+b)^2+b(u+b)v^2
       =(u+b)(u^2+bv^2-b^2)=0.
\]
If $b\ne0$ and $u=-b$, then $v=0$; take $U=b$ and $t=0$.
In every case $U^2=z$, and the remaining old equations follow from
$c=be$, $w=bh$, $j=s^2$, and $k=bq$. The possible sign change of $U$
is harmless because all old witnesses and deleted lifting coordinates
are private. These maps prove equality of existential projections;
they are not asserted to be inverse on the private witness tuples.
This rule is applied only at the certificate's height-parity blocks;
the ordinary comparison pole is the linear-unit source predicate already
proved above. The actual multiplicity is given in the following census.

\subsection{Finite realization and proof of the quartic theorem}
The fresh quadratic lift introduces only defining product equations.
Their evaluation and forgetting maps preserve rational solutions while
fixing the original input. The certificate then applies exactly the
operations below, under the current-row and complete-incidence hypotheses
just proved. The resulting file \nolinkurl{quartics/quartic-4973/residuals.json.gz}
contains $4\,552$ nonzero integer residuals of degree at most two.
Every weight in \nolinkurl{quartics/quartic-4973/weights.json} is one, and
\[
 \Phi(X_0,\ldots,X_{4972})=\sum_{i=1}^{4552}R_i^2.
\]
At rational values each square is nonnegative, so this sum vanishes if and
only if all residuals vanish. Nonzero rational row normalization preserves
each row's zero set before this integral sum is formed.

The certificate \nolinkurl{quartics/quartic-4973/certificate.jsonl.gz} records $551$ transformations:
\[
\begin{array}{lr}
\text{height-parity replacements}&1\\
\text{contact-power circuit replacements}&13\\
\text{invertible linear row operations}&266\\
\text{rational definition eliminations}&121\\
\text{invertible rational shears}&2\\
\text{eliminations modulo retained equations}&148\\
\end{array}
\]
Its SHA-256 is
\begin{center}\footnotesize\ttfamily
81ca768459a13b76dbdf4bfe81e94fbab9ab99f14a5731f6b1a1b03e28c1df2b
\end{center}
The independent verifier \nolinkurl{quartics/quartic-4973/verify_certificate.py} has SHA-256
\begin{center}\footnotesize\ttfamily
b7180b01029b3d9364d94914401a4f1f66b190762f10b99a38e9fd4f49437ffe
\end{center}
It checks the named source binding, fresh lift, every recorded
polynomial identity and reconstruction map, defining order, current-row
supports, every occurrence of each removed coordinate, final quadratic
degrees, and the bijective variable map fixing input zero. It then compares
the complete collected integer expansion, including every cross term.
The endpoint verification receipt is \nolinkurl{quartics/quartic-4973/evidence/final-verification.json}.

\begin{proof}[Proof of Theorem~\ref{thm:quartic}]
Theorem~\ref{thm:quartic-selected-source} proves both rational-input
implications for the selected arithmetic source. For any fixed rational
input, every defining lift and every recorded transformation has the
forward and reverse solution maps proved above. Their composition preserves
the existence of rational witnesses and keeps that input fixed. Hence the
final residual system has rational input projection exactly $\Z$.
The all-one sum of its squares vanishes precisely on its simultaneous
zero set, so the same equivalence holds for $\Phi$.

The literal coefficient table is the verified complete integral expansion
of that sum. Its nonzero degree-four term makes its degree exactly four;
the collected coefficients give the stated monomial count and maximum.
The checked variable map gives the stated ambient coordinates, one of
which is the input. Thus all numerical
claims and the input-projection equivalence concern this one polynomial.
\end{proof}

The finite checks do not replace the inherited rank, height, prime-pattern,
formal-group, norm or pole arguments. Their qualifications remain those
stated in the source proof and Appendix~\ref{app:sources}.

% Begin Version 5 additions
\section{Arithmetic interfaces for the smaller quartic}
\label{sec:v5-interfaces}
% Additive interfaces for the 1272-coordinate source.
% These use the original exponent-24 constant H_*, not H_19.
% They do not alter any existing source equation or polynomial artifact.

\begin{lemma}[The original scalar certificate at order thirty-two]
\label{v5:scalar32-bridge}
Let $K=\Q(\sqrt2)$, $d=3-\sqrt2$, $D=d^2$, and
$E:y^2=x^3-Dx$, with the point $P_*$ of
Section~\ref{sec:degree8-original}. Put
\[
 \beta=X^2+2,\qquad
 r=\beta(\beta-1)(\beta-2401)(32\beta+343)(2\beta-1),
 \qquad X\in\Q.
\]
For $h_0,a\in\Z\setminus\{0\}$ and points $J_0=h_0P_*$ and $S=aP_*$,
take the
two rational coefficients, in the basis $1,\sqrt2$, of
$f(J_0),f(S),z(S)/z(J_0)$ in that order. Denote them by
$X_1,\ldots,X_6$, and put
\[
 (m_1,\ldots,m_6)=(1,1,1,1,32,32),\qquad l_i=X_i/r^{m_i}.
\]
Use the same ordered norm tree as the original scalar construction:
\[
 \begin{split}
 N_r(u,v)&=u^2+rv^2,\\
 n_0&=N_r(1,l_1),\quad n_1=N_r(l_2,l_3),\quad
 n_3=N_r(l_5,l_6),\\
 W&=N_r\bigl(N_r(n_0,n_1),N_r(l_4,n_3)\bigr).
 \end{split}
\]
Let
\[
 c\in\left\{\pm\prod_{p\in I}p:
                   I\subseteq\{2,3,5,7,11,13,17\}\right\}.
\]
If $b=crW^2\ne0$ satisfies the original pole predicate $\Psi(b)$,
then
\[
 M_5(\beta)^{32}\mid a,\qquad
 h(\beta)\le H_*|a|^{3/4}.
\]
Here $M_5$ and the exact constant $H_*$ are those of
Theorem~\ref{hyp:o8-height}; in particular
$1\le H_*<2^{4783227}$.

For every integral $X$, the scalar certificate is complete with
successively chosen positive indices $h_0,a$. One may require $h_0$
to be divisible by any prescribed positive integer and $a$ to be
arbitrarily divisible and arbitrarily large. After $S$ has been chosen,
the normalized-leaf conditions and the nonvanishing of $W$ persist when
$h_0$ is replaced by every integer in one suitable congruence class
$h\equiv h_0\pmod L$. The modulus $L$ may also be required to be
divisible by any prescribed positive integer.
\end{lemma}

\begin{proof}
The value $r$ is nonzero: $\beta\ge2$, and $\beta=2401$ would
give $X^2=2399$, which has no rational solution. At a prime
$p>76176$, the multiplier $c$ is a unit. Every prime in $M_5(\beta)$
has odd valuation in $r$, by the contact computation in the proof of
Theorem~\ref{hyp:o8-height}. One can also observe that a negative
valuation of $\beta=X^2+2$ is even, so only its finite contacts
contribute here.

Write $R=v_p(r)$. The pole obstruction in Lemma~\ref{hyp:o8-pole}
and $b\ne0$ make the odd valuation $R+2v_p(W)$ positive. In a
norm node $u^2+rv^2$ the two nonzero summands have opposite valuation
parity, so there is no cancellation. The constant leaf on the
all-left path first excludes $R<0$. Then $v_p(W)>-R/2$.
Descending the tree, the implication
\[
 v_p(N_r(u,v))>-tR,\quad t<1
 \quad\Longrightarrow\quad
 v_p(u)>-tR/2,\quad v_p(v)>-(t+1)R/2
\]
shows that every leaf $l_i$ has valuation greater than $-R$.
Thus $v_p(X_i)>(m_i-1)R\ge m_i-1$, and integrality of the
valuation gives $v_p(X_i)\ge m_i$. This is the original norm
argument with the two assigned orders changed to thirty-two.

The order-one tests put $J_0$ and $S$ in the formal group at each
place above $p$. Such a completion is unramified, and the formal
logarithm is an isometry on that group. Since
$[h_0]S=[a]J_0$, it follows that
\[
 v_{\mathfrak p}\!\left(\frac{z(S)}{z(J_0)}\right)
       =v_p(a/h_0).
\]
Both rational coefficients of this ratio have valuation at least
thirty-two, so $v_p(a/h_0)\ge32$. As $h_0$ is a nonzero integer,
$v_p(a)\ge32$. This proves $M_5(\beta)^{32}\mid a$.
Theorem~\ref{hyp:o8-height} now gives
$h(\beta)\le H_*M_5(\beta)^{24}\le H_*|a|^{3/4}$.
No exponent-nineteen height estimate is used.

For completeness, let $X\in\Z$. The integer $r$ is nonzero and
$|r|>1$. The original seven-prime squareclass basis, proved in the
finite-correction argument following
Lemma~\ref{lem:original-seven-parity-completeness}, permits a choice
of $c$ such that $cr>0$ and $cr$ is square in
$\Q_2,\Q_3,\Q_{17}$. Let $\mathcal S$ be the finite prime support
of the integer $cr$.

At the places above $\mathcal S$, choose $h_0$ sufficiently
divisible that $J_0$ is in arbitrarily deep formal subgroups.
This remains possible after requiring any fixed divisibility of
$h_0$, because these subgroups have finite index in the local
elliptic groups. Having fixed this nonzero $J_0$, choose $a$
sufficiently divisible that $S=aP_*$ is still deeper. The
coefficients of $f(J_0),f(S)$ and $z(S)/z(J_0)$ then attain every
finite lower valuation bound needed here. This includes ramified
and bad primes, by working in sufficiently deep formal subgroups.
We arrange that all normalized leaves are integral when $p\mid r$
and divisible by $p$ when $p\mid c$ but $p\nmid r$.
The index $a$ may simultaneously exceed any fixed bound.

At $p\mid r$, the root of the displayed tree is congruent to one
modulo $p$, because $r\equiv0$ and the all-left leaf is one.
At $p\mid c$, $p\nmid r$, all other leaves vanish modulo $p$,
so the root is again congruent to one. Thus $W$ is a unit on
$\mathcal S$, and in particular is nonzero. At these primes
$v_p(crW^2)=v_p(cr)\ge0$; outside $\mathcal S$ its valuation is
even. Hence $b$ has square reduced denominator, is positive, and
has the same local squareclasses as $cr$ at two, three, and
seventeen. Lemma~\ref{lem:original-seven-parity-completeness}
supplies the scalar pole witnesses.

Finally, with $S$ fixed, these scalar tests are rational functions
of $J_0$ defined at the chosen point. Their coefficient maps are
continuous in the product of the completions above each prime in
$\mathcal S$. The finite lower valuation bounds persist in
sufficiently small neighborhoods. A common finite-index subgroup
of the integer orbit gives a congruence $h\equiv h_0\pmod L$
contained in all those neighborhoods. Enlarging $L$ preserves
that property. The same unit calculation proves the asserted
stability for the final shared point.
\end{proof}

\begin{lemma}[The fixed multiplier seventeen with preparation zero]
\label{v5:preparation-zero-bridge}
Let $b\in\Q^\times$, put $B=17b$, and define $\Psi_0(b)$ by
the existence of $e,p,q,s\in\Q$ satisfying
\[
 p^2=B(1-e)(1+3e),\qquad h=e(B-b^2),\qquad
 Bq^2=(s^2-Bh)(s^2+3Bh),
\]
where $h$ is the displayed expression, not a further free choice.
Then $\Psi_0(b)$ excludes negative odd valuations of $b$ at every
prime outside $\{2,3,17\}$. Conversely, it holds whenever $b$ has
square reduced denominator, is positive and square in
$\Q_2,\Q_3,\Q_{17}$, and
\[
 b<17,\qquad v_2(b)\ge3,\qquad v_3(b)\ge1,
 \qquad v_{17}(b)\ge2.
\]
No additional nonzero condition on $e,p,q,s,h$ is imposed.
\end{lemma}

\begin{proof}
Taking the preparation witnesses $t=0$ and $U=b$ in
\eqref{eq:o8-pole} gives exactly the displayed equations.
The obstruction therefore follows from
Lemma~\ref{lem:original-seven-parity-soundness}, including its
degenerate witness cases.

For completeness put
\[
 V=B(B-b^2)=B^2(1-b/17).
\]
The real inequality makes $V$ positive. The valuation bounds make
$1-b/17$ a square at two, three, and seventeen: its difference
from one has valuation at least three at two and at least one
at the two odd primes. Thus $V$ is square at these three places.

Let $\mu$ and $\theta$ be the positive squarefree representatives
of $B$ and $V$. The prescribed squareclasses of $b$ give
$\mu\equiv17\pmod{24}$; in particular $\mu>1$ and its support
contains seventeen but neither two nor three. Also
$\theta\equiv1\pmod{24}$, and seventeen is absent from its
support. At a prime $\ell\mid\mu$ outside the three exceptions,
$v_\ell(b)$ is positive and odd, since the reduced denominator
of $b$ is square. Hence $1-b/17$ is a square unit there, so $V$
and $\theta$ are local squares. The same assertion at seventeen
was already arranged. These facts make the supports disjoint.

Conversely, if $\ell\mid\theta$, then
$\ell\notin\{2,3,17\}$. If $v_\ell(b)>0$, the valuation
of $V$ is $2v_\ell(b)$; if $v_\ell(b)<0$, it is
$3v_\ell(b)$, again even because the denominator is square.
Thus an odd valuation of $V$ forces $v_\ell(b)=0$ and
$1-b/17\equiv0\pmod\ell$. It follows that
$B=17b\equiv17^2\pmod\ell$ is a nonzero square unit.
This proves the other mutual-squareness condition.

Lemma~\ref{lem:original-seven-simultaneous-realization} now gives
the two rational points in the squareclass of $B$. To make the
conversion explicit, write $V=\theta c^2$ and scale the second
point to the curve with parameter $nV$. If the first abscissa is
$Ba^2$, scale both points and parameters by $a$ so that their
abscissas are $B$ and $Bs^2$, with $s\ne0$, and their curve
parameters are $e_0$ and $e_0V$. For their ordinates $y_1,y_2$, set
\[
 e=e_0/B,\qquad p=y_1/B,\qquad q=y_2/(B^2s),\qquad
 h=e(B-b^2).
\]
The first curve equation gives $p^2=B(1-e)(1+3e)$.
Since $e_0V=B^2h$, the second gives
$Bq^2=(s^2-Bh)(s^2+3Bh)$. These are the required witnesses.
The only arithmetic existence theorem used in this specialization
is the existing simultaneous-realization lemma, with its stated
restricted-converse and descent dependencies.
\end{proof}

\begin{lemma}[Real bounds independent of later progression moduli]
\label{v5:real-progression-bridge}
Let $G$ be the closure of the integer orbit
\[
 \{([n]P_*,[n]P_*^\sigma):n\in\Z\}
       \subset E(\mathbb R)\times E^\sigma(\mathbb R),
\]
and let $G^0$ be its identity component. Fix a finite collection of
rational functions formed from nonzero integral multiples and fixed
translates of either elliptic coordinate, with all denominators
nonzero as rational functions on the relevant curve. Suppose their
pole set is contained in a finite union of inverse images of finite
sets under these nonzero multiplication-and-translation maps on
either factor. There is a bound $C$ such that every nonconstant
integer progression contains arbitrarily large positive indices at
which all these functions are defined and their real absolute values
are at most $C$. The bound is chosen from the functions before the
progression modulus is specified.

If $m=|G/G^0|$, $m\mid L$, and $m\mid h_0$, the progression
$h\equiv h_0\pmod L$ is dense in $G^0$, and its positive tail
meets every nonempty open subset of $G^0$ infinitely often.
\end{lemma}

\begin{proof}
The real elliptic groups are compact Lie groups. Their closed
subgroup $G$ is a compact abelian Lie group with finitely many
components, and $G^0$ is a torus. The projection of $G^0$ onto
each elliptic identity circle is surjective. Indeed the projected
integer orbit is infinite, since $P_*$ has infinite order and each
real embedding is injective on the group of algebraic points.
If the connected projection were trivial, the projected $G$ would
be finite; a nontrivial connected closed subgroup of a circle is
the whole circle.

Each nonzero multiplication map on a circle has finite fibers.
Consequently each of the prescribed inverse images of a finite
set has Haar measure zero in every component of $G$. The finite
union of pole sets cannot cover a component. Choose in each
component a nonempty open neighborhood whose compact closure
avoids the poles. Continuity bounds all the fixed rational
functions on these finitely many compact sets by one constant $C$.

For a progression $h=r+kn$, its closure is
$[r]P_*+[k]G$. Multiplication by the positive integer $k$ maps
$G^0$ onto itself, so this closure contains an entire component
of $G$. The positive tail of the progression is dense in its
closure: compactness gives unbounded positive multiples tending
to the identity, so each negative multiple is also a limit of
positive multiples. It therefore meets the chosen neighborhood
in that component at arbitrarily large positive indices.

If $m\mid k$, then $[k]G=G^0$. If also $m\mid r$, the coset
is $G^0$ itself. Apply this with $k=L$ and $r=h_0$ for the final
assertion.
\end{proof}

\section{The shared-point quartic in 962 variables}
\label{sec:v6-962}

This section gives a new complete quadratic compilation of the shared-point
construction in the public agentic notes~\cite{agentic-improvements}.
The shared comparison points, reciprocal-trace detectors, fixed height
scales, direct Lucas test, and affine anchor change are arithmetic
and algebraic improvements recorded in those notes. The simultaneous prime theorem used to choose
the shared points is Sun's Theorem~4.1, in the precise form of
Theorem~\ref{thm:v5-simultaneous-primes} below. Our contribution here is an
independently specified gate program, its source verification, and
further changes of coordinates and circuit identities leading to a
962-variable polynomial after 32 certified eliminations. The
definition does not depend on an unavailable coefficient table.
The relevant passages in the public notes are
\href{https://github.com/JoseBrox/diophantine-z-q_improvements/blob/f4d1b5e6c23d4adeb3ae1f986ab0bf87884bdff6/paper/diophantine_optimization_20261106_v3.tex#L2354-L3044}{shared-point construction},
\href{https://github.com/JoseBrox/diophantine-z-q_improvements/blob/f4d1b5e6c23d4adeb3ae1f986ab0bf87884bdff6/paper/diophantine_optimization_20261106_v3.tex#L3352-L3405}{fixed height scales},
and
\href{https://github.com/JoseBrox/diophantine-z-q_improvements/blob/f4d1b5e6c23d4adeb3ae1f986ab0bf87884bdff6/paper/diophantine_optimization_20261106_v3.tex#L4135-L4758}{quadratic compilation and anchor translation}.

We retain the arithmetic conclusions and conventions of
Section~\ref{sec:degree8-original}, including the restricted-converse
input to Lemma~\ref{hyp:o8-pole}. In this section the height estimate is
Theorem~\ref{hyp:o8-height}, with exponent 24 and its original constant
$H_*$; it is not the exponent-19 estimate used for the preceding quartic.
Thus
\[
  1\le H_*<2^{4783227}.
\]
The exact finite compilation and coefficient assertions below require no
arithmetic conjecture. Identifying the rational input projection with
$\Z$ uses the stated arithmetic results and the simultaneous prime theorem.

The point, pole, and correction variables below specify the
arithmetic source predicate. The final compilation subsection
gives explicit maps to its implementation coordinates.

\subsection{Points, local parameters, and quotient domains}
\label{subsec:v5-point-source}

Use the normalized curve and generator already established in
Section~\ref{sec:degree8-original}:
\[
 K=\Q(\sqrt2),\quad d=3-\sqrt2,\quad D=d^2,
 \quad E:y^2=x^3-Dx,
\]
\[
 P_*=(57/8-7\sqrt2/2,-165/8+417\sqrt2/32),
 \qquad 2E(K)=\Z P_*.
\]
The last assertion follows from Proposition~\ref{prop:generator} by the
unit change of variables recorded there. Conjugation of $K/\Q$ is
denoted by $\sigma$. All field equations below mean both rational
component equations in the basis $1,\sqrt2$.

The source uses nine nonidentity points
\[
 A,S,T,Z,H,U_A,U_S,U_T,U_Z\in2E(K).
\]
Each is represented as the double of an affine halfpoint
$(\alpha,\eta)\in E(K)$, using four rational halfpoint coordinates.
The equations for a double are
\begin{equation}\label{eq:v5-double}
 \eta^2=\alpha^3-D\alpha,\qquad
 2\eta m=3\alpha^2-D,\qquad
 x=m^2-2\alpha,\qquad y=m(\alpha-x)-\eta.
\end{equation}
Products of degree greater than two in this display are evaluated by
the quadratic gates specified below. Set
\[
 Q_0=P_*,\quad(Q_1,Q_2,Q_3,Q_4)=(A,S,T,Z),
 \quad(U_0,U_1,U_2,U_3,U_4)=(H,U_A,U_S,U_T,U_Z).
\]
Compute $[4]Q_i$ by two doublings, or by fixed rational arithmetic for
$Q_0$, and compute
\[
 V_i=U_i-[4]Q_i\quad(0\le i\le4).
\]
For the subtraction of $(x',y')$ from $(x,y)$ we retain an inverse of
$x-x'$ in $K$ and use
\[
 m=(y+y')/(x-x'),\quad x''=m^2-x-x',\quad
 y''=m(x-x'')-y.
\]

\begin{lemma}[Domains of the point quotients]\label{lem:v5-point-domains}
The equations \eqref{eq:v5-double} force $\eta\ne0$ and describe a
finite nonidentity point of $2E(K)$. Every such point has nonzero
abscissa and ordinate. The guarded differences $V_i$ have the same
properties. Consequently $f(Q)=1/x(Q)$ and $z(Q)=-x(Q)/y(Q)$ are
defined and nonzero at every point where they are used here.
\end{lemma}
\begin{proof}
If $\eta=0$, the curve equation gives $\alpha=0$ or $\alpha^2=D$.
The right side of the slope equation is then $-D$ or $2D$, respectively,
a contradiction. The ordinary duplication formulas therefore apply.
They give
\[
 x([2](\alpha,\eta))=
 \left(\frac{\alpha^2+D}{2\eta}\right)^2\ne0,
\]
because $D=d^2$ and the numerator is positive at either real embedding.
An ordinate-zero point has abscissa $0,d$, or $-d$.
The latter two are not squares in $K$: their rational norm is 7,
whereas the norm of a square is a rational square. Thus the doubled
ordinate is nonzero. A guarded affine difference is finite and is
again in $2E(K)$, so the same argument applies. Conversely any
nonidentity point of $2E(K)$ has an affine halfpoint of nonzero
ordinate and is represented by these formulas.
\end{proof}

Write
\[
 a_0=1,\quad a_1=-2/5,\quad
 a_j=\frac{4j-5}{4j+1}a_{j-1}\ (j\ge2),\qquad
 G_J(w)=\sum_{j=0}^J a_jw^j,
\]
and, for $e=16,52$, put
\begin{equation}\label{eq:v5-log}
 \ell_e(Q)=z(Q)G_{e/4-1}(Df(Q)^2).
\end{equation}
Lemma~\ref{lem:quartic-reciprocal-nonzero} applies to both truncation
degrees 3 and 12. Thus all these values are nonzero globally.

\begin{lemma}[The logarithm estimate at unramified places]
\label{lem:v5-unramified-log}
Let $p>76176$ be a rational prime and $\mathfrak p\mid p$ a place
of $K$. If a nonidentity point $Q$ is in the formal group at
$\mathfrak p$, normalize $v(p)=1$ and put $r=v(z(Q))$. Then $r\ge1$,
and for $e\in\{16,52,64\}$,
\[
 v(\log_E Q)=r,\qquad
 v(\log_E Q-\ell_e(Q))\ge(e+1)r,\qquad
 \ell_e(Q)/\log_E Q\in1+p^e\mathcal O_{K_{\mathfrak p}}.
\]
Here $\ell_{64}$ is defined by the same formula with $G_{15}$.
Ratios with numerator precision $e$ and denominator precision $e'$
in this set have the valuation of their full logarithm ratio and
relative error of valuation at least $\min(e,e')$.
\end{lemma}
\begin{proof}
The local extension is unramified, including when $p$ is inert.
The curve has good reduction and $D$ is a unit. On the formal branch,
\[
 f=z^2(1-Df^2),\qquad w=Df^2,\qquad v(w)=4r.
\]
The recurrence gives
$4w(1-w)G'(w)+(1+w)G(w)=1-w$ and hence the same identity
$\log_E Q=zG(w)$ as in Lemma~\ref{lem:quartic-reciprocal-log}.
For $j\ge1$ the denominator of $a_j$ divides $(4j+1)!$, so
\[
 v(a_j)\ge-\frac{4j+1}{p-1}.
\]
Set $m=e/4\le16$. If $j\ge m$ and $4j+1<p$, then
$v(a_jw^j)\ge4mr$. Otherwise $j\ge19044$ and
\[
 v(a_jw^j)\ge4jr-j\ge3jr\ge4mr.
\]
The lower bounds tend to infinity, proving the estimate for the
entire tail. Every nonconstant term of $G(w)$ has positive valuation,
so the leading term gives $v(\log_E Q)=r$. Division proves the
relative estimate, and division of two relative estimates proves
the ratio assertion. No step uses residue degree one.
\end{proof}

\subsection{The complete rational source}
\label{subsec:v5-complete-source}

The ordinary input is $X$, and $\beta=X^2+2$. Define
\[
 (k_1,k_2,k_3,k_4)=(16,52,52,52),\qquad
 L_i=\frac{\ell_{k_i}(U_i)}{\ell_{52}(H)},\qquad
 L'_i=\sigma\!\left(\frac{\ell_{k_i}(V_i)}{\ell_{52}(V_0)}\right).
\]
The sixteen field tests $D_j$, with their orders $m_j$, are exactly
\begin{equation}\label{eq:v5-test-list}
\begin{array}{c|c}
 f(U_0),\ldots,f(U_4),\ \sigma f(V_0),\ldots,\sigma f(V_4)&1\quad(10)\\
 L_i-L'_i\ (1\le i\le4)&16,52,52,52\\
 L_1-\beta&16\\
 F(L_2,L_3,L_4)&52,
\end{array}
\end{equation}
where
\begin{equation}\label{eq:v5-lucas-polynomial}
 F(U,V,W)=V^2-U(U-2)VW+(U-2)^2W^2-1.
\end{equation}
There is no rational proxy for an integer point label.

For $j=1,\ldots,4$ write $x(Q_j)=x_{j,0}+x_{j,1}\sqrt2$ and set
\[
 \mathcal A_0=1,\quad
 \mathcal A_{2j-1}=x_{j,0}^2,\quad
 \mathcal A_{2j}=(1+x_{j,1}^2)^2.
\]
These are positive nonzero rational squares. In fact $x(Q_j)$ is
positive at both real embeddings by Lemma~\ref{lem:v5-point-domains},
so $x_{j,0}>0$. Retain one rational coordinate $\lambda$ and put
\[
 \Theta=\{0,-1/2,-1,1/2399,-32/407\},\qquad
 F_5(T)=\prod_{\theta\in\Theta}(T-\theta),\qquad
 t_r=\mathcal A_r\lambda,\quad\delta_r=t_r-1.
\]
For every $0\le r\le8$, require $t_r\notin\Theta\cup\{1\}$ and define
\begin{align}
 R_r&=\mathcal A_r^{-5}F_5(t_r)/\delta_r^6,
 &\widetilde D_{r,j}&=D_j/\delta_r^{m_j+1},\label{eq:v5-detector-data}\\
 W_r&=\prod_{j=1}^{16}
 \left(1+R_r(\widetilde D_{r,j}/R_r^{m_j})^2\right),
 &T_r&=\Tr_{K/\Q}(W_r^{-1}),\qquad B_r=R_r/T_r^2.
 \label{eq:v5-detector-product}
\end{align}
Require $W_r,T_r\ne0$. For $B=B_r$ retain four rational witnesses
$e_r,s_r,p_r,q_r$ and impose
\begin{equation}\label{eq:v5-zero-pole}
 \begin{aligned}
 C_r&=17B_r,\qquad h_r=e_r(C_r-B_r^2),\\
 p_r^2&=C_r(1-e_r)(1+3e_r),\\
 C_rq_r^2&=(s_r^2-C_rh_r)(s_r^2+3C_rh_r).
 \end{aligned}
\end{equation}
The first two expressions are arithmetic definitions. The last two
are the preparation-zero specialization of the original pole predicate.
There is no nonzero guard on $e_r,s_r,p_r,q_r$, or $h_r$.

The scalar block shares the point $H$. Define
\[
 r_\beta=\beta(\beta-1)(\beta-2401)(32\beta+343)(2\beta-1).
\]
It is nonzero for rational $X$: the only positive possible root not
excluded by $\beta\ge2$ would require $X^2=2399$.
Let $Y_1,\ldots,Y_6$ be, in order, the two rational components of
$f(H),f(S),z(S)/z(H)$, and set
\[
 (e_1,\ldots,e_6)=(1,1,1,1,32,32),\qquad
 y_i=Y_i/r_\beta^{e_i},\qquad N_r(u,v)=u^2+rv^2.
\]
The scalar norm tree has no padding leaf. It is
\begin{equation}\label{eq:v5-scalar-tree}
 \begin{aligned}
 n_0&=N_{r_\beta}(1,y_1),\quad n_1=N_{r_\beta}(y_2,y_3),\quad
 n_3=N_{r_\beta}(y_5,y_6),\\
 W_\beta&=N_{r_\beta}\bigl(N_{r_\beta}(n_0,n_1),
                          N_{r_\beta}(y_4,n_3)\bigr).
 \end{aligned}
\end{equation}
Use eight Boolean coordinates $\epsilon_-,\epsilon_p$, for
$p\in\{2,3,5,7,11,13,17\}$, impose $\epsilon^2-\epsilon=0$, and set
\[
 c=(1-2\epsilon_-)\prod_p(1+(p-1)\epsilon_p),\qquad
 b_\beta=c r_\beta W_\beta^2,\qquad C_\beta=17b_\beta.
\]
Require $b_\beta\ne0$. Retain five rational witnesses
$t_\beta,e_\beta,s_\beta,p_\beta,q_\beta$ and impose
\begin{align}
 (1+b_\beta t_\beta^2)U_\beta
    &=b_\beta(1-b_\beta t_\beta^2),
 &h_\beta&=e_\beta(C_\beta-U_\beta^2),\label{eq:v5-scalar-preparation}\\
 p_\beta^2&=C_\beta(1-e_\beta)(1+3e_\beta),
 &C_\beta q_\beta^2&=(s_\beta^2-C_\beta h_\beta)
                     (s_\beta^2+3C_\beta h_\beta).
 \label{eq:v5-scalar-pole}
\end{align}
Here $U_\beta,h_\beta$ are arithmetic values. The preparation equation
forces $1+b_\beta t_\beta^2\ne0$: otherwise $b_\beta t_\beta^2=-1$
but its right side is $2b_\beta\ne0$.

Finally fix the integers
\begin{equation}\label{eq:v5-fixed-constant}
 C=3^{2^{26}},\qquad \mathsf S=C^2=3^{2^{27}},\qquad
 \kappa=\mathsf S^2=3^{2^{28}}.
\end{equation}
The symbol \texttt{S} in the coefficient files means $\mathsf S$;
it is not the point $S$ and is not a quantified coordinate.
Write $f(S)=f_0+f_1\sqrt2$, retain three private rational witnesses
$u,v,w$, and impose
\begin{equation}\label{eq:v5-physical-anchor}
 u^2+6v^2-2w^2-3(f_0^2+f_0f_1-f_1^2)^2+2\mathsf S^2=0.
\end{equation}
This completes the rational source. Every rational expression has the
displayed guard or a domain implication already proved. Equations
\eqref{eq:v5-zero-pole} and \eqref{eq:v5-scalar-pole} retain their
degenerate witness values.

\begin{lemma}[The three-witness anchor]\label{lem:v5-three-witness-anchor}
Equation~\eqref{eq:v5-physical-anchor} is equivalent to
$f_0/C,f_1/C\in\Z_{(2)}\cap\Z_{(3)}$. For the nonzero integer label
$a$ of $S=[a]P_*$ it implies
\[
 0\ne a\in2\,3^{2^{25}-1}\Z,\qquad
 |a|>(21H_*)^{13/2}>(20H_*)^4.
\]
Its permitted labels contain arbitrarily divisible positive integers.
\end{lemma}
\begin{proof}
Put $R=\Z_{(2)}\cap\Z_{(3)}$. The quaternion algebra $(2,3)_\Q$
ramifies exactly at 2 and 3. The quadratic embedding criterion
\cite[Proposition 14.6.7]{voight2021} shows, for rational $t\ne0$,
that $1-2/(3t^2)$ is represented by $2x^2+3y^2-6z^2$ precisely
when it is nonsquare in both $\Q_2$ and $\Q_3$. This last condition
is exactly $t\in R$. Indeed at 2 a positive valuation of $t$ gives
odd valuation of the expression, valuation zero gives residue 3
modulo 8, and negative valuation gives a square in $1+8\Z_2$.
At 3, nonnegative valuation of $t$ gives odd negative valuation,
while negative valuation gives a square in $1+3\Z_3$.
The expression is nonzero because $t^2=2/3$ has no rational solution.
Multiplying by $3t^2$ gives
\[
 t\in R\quad\Longleftrightarrow\quad
 \exists u,v,w\in\Q:\ u^2+6v^2-2w^2-3t^2+2=0.
\]
The case $t=0$ has witnesses $(0,0,1)$. A rational-square value of
$1-2/(3t^2)$ is excluded by the division-algebra criterion as well:
a nonzero trace-zero quaternion cannot square to a nonzero rational
square in a division algebra.

The form $A^2+AB-B^2$ has no nontrivial zero over $\F_2$ or $\F_3$.
Scaling a pair therefore proves
$v_p(a^2+ab-b^2)=2\min(v_p(a),v_p(b))$ for $p=2,3$.
Apply the preceding three-witness formula to
$(f_0^2+f_0f_1-f_1^2)/C^2$ and multiply its witnesses by $C^2$.
This gives exactly \eqref{eq:v5-physical-anchor}.

The local calculation in Section~\ref{subsec:o8-anchor-subgroup}
applies with exponent $2^{26}$: the reduction of $P_*$ at 3 has
order two and $v_3(z(2P_*))=1$. For even nonzero $a$ it gives
$v_3(f([a]P_*))=2+2v_3(a)$. At 2 the rational components of
$f([a]P_*)$ are integral. The anchor is consequently equivalent to
$0\ne a\in2\,3^{2^{25}-1}\Z$. Its minimum possible absolute value
has binary logarithm greater than
\[
 1+\tfrac32(2^{25}-1)=50331647.5,
\]
whereas
$\log_2((21H_*)^{13/2})<\tfrac{13}{2}(5+4783227)=31091008$.
The last comparison with $(20H_*)^4$ follows from
$21^{13}>20^8$ and $H_*\ge1$. The subgroup description proves the
last assertion.
\end{proof}

\subsection{Local tests and the height implication}
\label{subsec:v5-soundness}

\begin{lemma}[One place for every test]\label{lem:v5-one-place}
Let $R\in\Q^*$, $D_j\in K$, and positive integers $m_j$ define
\[
 W=\prod_j(1+R(D_j/R^{m_j})^2),\quad
 T=\Tr(W^{-1}),\quad B=R/T^2,
\]
with $W,T\ne0$ and a preparation-zero pole certificate for $B$.
If $p>76176$ and $v_p(R)$ is odd, then $v_p(R)>0$ and some
$\mathfrak p\mid p$ has $v_{\mathfrak p}(D_j)\ge m_j$ for every $j$.
\end{lemma}
\begin{proof}
Set $\alpha=v_p(R)$. The nonconstant term in each nontrivial factor
has odd valuation $\alpha+2(v(D_j)-m_j\alpha)$; unramifiedness
makes these valuations integral. The factor valuation is its minimum
with zero. Thus $\lambda_{\mathfrak p}=v_{\mathfrak p}(W)\le0$.
Put $\lambda=\max_{\mathfrak p\mid p}\lambda_{\mathfrak p}$.
Integrality of trace gives $v_p(T)\ge-\lambda$. Hence
\[
 v_p(B)\le\alpha+2\lambda.
\]
The left side is odd and, by Lemma~\ref{hyp:o8-pole}, not negative.
It is positive, so $\alpha>0$ and $\lambda>-\alpha/2$.
At a place attaining $\lambda$, each factor has valuation at least
$\lambda$, because all factor valuations are nonpositive. For
nonzero $D_j$ this implies
\[
 v(D_j)>(m_j-3/4)\alpha\ge m_j-3/4.
\]
Integrality gives $v(D_j)\ge m_j$; a zero $D_j$ satisfies it as well.
\end{proof}

\begin{lemma}[The two-anchor comparison]\label{lem:v5-two-anchors}
Write the integer labels of $(A,S,T,Z)$ as $(n,a,t,b)$.
If the first fourteen tests in \eqref{eq:v5-test-list} have their
specified orders at a place $\mathfrak p\mid p$, $p>76176$, then
\[
 L_i\in\mathcal O_{K_{\mathfrak p}},\qquad
 L_i\equiv n_i\pmod{p^{k_i}},\qquad
 (n_1,n_2,n_3,n_4)=(n,a,t,b).
\]
This holds at split and inert primes.
\end{lemma}
\begin{proof}
Write $H=[h]P_*$ and $U_i=[u_i]P_*$. The order-one tests put
$H,U_i$ in the formal group at $\mathfrak p$ and $V_0,V_i$ in the
formal group at the conjugate place. Positive valuation of $f=1/x$
gives formal membership on this good model. The guarded points give
$h\ne0,4$. Their full logarithm ratios, after conjugating the second,
are
\[
 r_i=u_i/h,\qquad r'_i=(u_i-4n_i)/(h-4).
\]
These identities follow from the homomorphism property on equal
formal multiples, and do not require $P_*$ itself to be formal.
At least one of $h,h-4$ is a $p$-unit. Its ratio is integral, and
Lemma~\ref{lem:v5-unramified-log} gives the same valuation to its
truncated ratio. The cross test makes the other truncated ratio
integral too. Both relative errors are therefore absolute errors,
and $L_i\equiv r_i\equiv r'_i$ modulo $p^{k_i}$. Multiplication
by $h,h-4$ and subtraction give $4L_i\equiv4n_i$. Divide by four.
At an inert prime conjugation preserves the single local valuation,
so the same argument applies.
\end{proof}

\begin{lemma}[Height transfer]\label{lem:v5-height-transfer}
Let $M_r$ be the odd-contact radical of $t_r$ for the five contacts
$\Theta$, omitting primes at most 76176. If $M_r^k$ divides a
nonzero integer $N$ for every $r$, then
\[
 h(x(Q_i))<7H_*|N|^{24/k}\qquad(1\le i\le4).
\]
At every prime counted in $M_r$, the source imposes all sixteen
field tests with their specified orders at one place above that prime.
\end{lemma}
\begin{proof}
The fractional linear map $\phi(t)=2+1/t$ carries $\Theta$ to
$\infty,0,1,2401,-343/32$. Its integral matrix has determinant $-1$.
All primes dividing the relevant fixed coefficients and contact
differences, including differences from 1, lie in
$\{2,3,5,7,11,37,109,439,2399\}$. Thus outside the cutoff
$M_r=M_5(\phi(t_r))$, and
\[
 h(t_r)\le h(\phi(t_r))+\log3
       \le(H_*+\log3)M_r^{24}.
\]
At a prime of $M_r$, $\delta_r$ is a unit. Exactly one contact
factor has odd positive valuation, while $\mathcal A_r^{-5}$ is
a rational square. Therefore $v_p(R_r)$ is odd.
Lemma~\ref{lem:v5-one-place} applied to $\widetilde D_{r,j}$ gives
the asserted orders for $D_j$.

Set $T_0=(H_*+\log3)|N|^{24/k}$. The case $r=0$ gives
$h(\lambda)\le T_0$, and the other cases give
$h(\mathcal A_r\lambda)\le T_0$. Hence
$h(\mathcal A_r)\le2T_0$. For the first component,
$h(x_{i,0})\le T_0$. For $x_{i,1}=a/b$ in lowest terms,
\[
 h(1+x_{i,1}^2)=\log(a^2+b^2)\ge2h(x_{i,1}),
\]
so $h(x_{i,1})\le T_0/2$. Finally
\[
 h(x(Q_i))\le h(x_{i,0})+h(x_{i,1})+\tfrac32\log2
 <7H_*|N|^{24/k}.
\]
\end{proof}

\begin{lemma}[The two reduced group cardinalities]
\label{lem:v5-two-reductions}
Let $p>76176$ be prime and let $\pi\in\Z[i]$ have norm $p$, with
$\pi\equiv9\pmod{16}$ and $\pi\equiv1+3i\pmod7$.
The two reductions of $E/K$ at the places above $p$ have group
cardinalities $N(\pi-1)$ and $N(\pi+1)$, in some order.
\end{lemma}
\begin{proof}
The congruences give $p\equiv1\pmod8$ and $p\equiv3\pmod7$.
Thus $p$ splits in $K$ and $(7/p)=-1$. The reductions of
$d=3-\sqrt2$ have product 7 and opposite square classes. The two
curves $y^2=x^3-d^2x$ are consequently opposite quadratic twists
of $y^2=x^3-x$ and have opposite traces.

The Frobenius calculation already given in
Section~\ref{subsec:o8-pattern-proof} applies: the order-four
automorphism is defined over $\F_p$, its rational centralizer is
$\Q(i)$, and integrality puts Frobenius in $\Z[i]$.
Write it as $a+bi$, so $a^2+b^2=p$ and the two cardinalities are
$p+1-2a$ and $p+1+2a$. Full rational two-torsion makes both
cardinalities divisible by four. Since $p\equiv1\pmod4$, this
forces $a$ odd and $b$ even. The real part of $\pi$ is odd and its
imaginary part is even. Unique factorization in $\Z[i]$ therefore
makes $a+bi$ one of $\pm\pi,\pm\overline\pi$; the associates
with factor $i$ have the wrong parity. Hence the unordered
cardinalities are exactly $N(\pi-1),N(\pi+1)$.
\end{proof}

\begin{theorem}[Arithmetic interpretation of the source]
\label{thm:v5-source-projection}
With the inherited arithmetic results and
Theorem~\ref{thm:v5-simultaneous-primes}, the rational source of
Section~\ref{subsec:v5-complete-source} has input projection exactly $\Z$.
\end{theorem}
\begin{proof}[Soundness]
Let $\mathcal B=\max(|n|,|a|,|t|,|b|)$ and
$\Delta=F(a,t,b)\in\Z$. Polynomial substitution preserves
congruences between integral local elements. Lemmas
\ref{lem:v5-two-anchors} and \ref{lem:v5-height-transfer} therefore
give $M_r^{52}\mid\Delta$ whenever $\Delta\ne0$.
The index-height estimate \eqref{eq:o8-indexheight} yields
\[
 \mathcal B^2-6<7H_*|\Delta|^{24/52},\qquad
 |\Delta|<5\mathcal B^4,
\]
the latter because $\mathcal B\ge|a|>7$. Consequently
\[
 \mathcal B^{2/13}<(7\cdot5^{6/13}+6)H_*<21H_*.
\]
For the last strict inequality, $5^6 7^{13}<15^{13}$.
This contradicts Lemma~\ref{lem:v5-three-witness-anchor}. Thus
$\Delta=0$, and Lemma~\ref{lem:o8-lucas}, with $b\ne0$, gives
$\log\mathcal B>|a|-3$.

Lemma~\ref{v5:scalar32-bridge} applies to the exact tree
\eqref{eq:v5-scalar-tree} and gives
\[
 h(\beta)\le H_*|a|^{3/4}<|a|/20.
\]
If $\beta\ne n$, put $N=|\operatorname{num}(\beta-n)|$ in lowest
terms. The order-16 test gives $M_r^{16}\mid N$. Therefore
\[
 \mathcal B^2<7H_*N^{3/2}+6,\qquad
 N\le2e^{h(\beta)}\mathcal B,
\]
and $\log\mathcal B<2\log(32H_*)+3h(\beta)$.
The anchor also gives $\log(32H_*)<|a|/20$, so the right side
is less than $|a|/4<|a|-3$, a contradiction. Hence
$\beta=n\in\Z$. Since $X\in\Q$ and $X^2=n-2\in\Z$,
the reduced denominator of $X$ is one.
\end{proof}

\begin{proof}[Completeness]
Fix $X\in\Z$, set $n=X^2+2$, and choose the Boolean correction
$c$ so that $cr_\beta>0$ is square in $\Q_2,\Q_3,\Q_{17}$,
as in the finite squareclass computation accompanying
Lemma~\ref{hyp:o8-pole}. Let $\mathcal S$ be the prime support of
$cr_\beta$. The scalar completeness construction in
Lemma~\ref{v5:scalar32-bridge} first chooses $h_0>0$ divisible by
64 with $[h_0]P_*$ sufficiently deep at these primes. It then
chooses a positive label $a$ sufficiently divisible relative to
$h_0$ and satisfying the anchor. The normalized scalar leaves are
integral at primes dividing $r_\beta$ and divisible at primes
dividing $c$ but not $r_\beta$. Choose positive $t,b$ with
$F(a,t,b)=0$ by Lemma~\ref{lem:o8-lucas}.
This fixes the four root points and all scales $\mathcal A_r$.

Temporarily put $H=[h]P_*$ and $U_i=[hn_i]P_*$.
All comparison tests are rational functions of $[h]P_*$ with
nonzero denominators at $h=h_0$. Their poles lie above finite
sets under fixed nonzero multiplication and translation maps
on the two real factors, as required by the real progression lemma.
Continuity preserves the scalar
conditions at every place over $\mathcal S$ and bounds below the
valuations of all sixteen comparison tests at the places over
$2,3,17$, throughout a progression
\[
 h\equiv h_0\pmod L,\qquad64\mid L.
\]
One obtains $L$ by putting sufficiently small open formal subgroups
inside the finitely many translated neighborhoods and taking a
common multiple of their indices. Lemma~\ref{v5:real-progression-bridge}
gives, independently of future progression moduli, real bounds for
the tests attained by arbitrarily large positive $h$ in each
nonconstant progression. The direct polynomial $F$ introduces
no denominator and has the same boundedness property.

Choose a positive integer $D_0$ such that every
$D_0^2\theta/\mathcal A_r$ is integral, and put
\[
 v=D_0^2w^2,\qquad\lambda=u/v.
\]
The contact positions are the integer forms
$u-(D_0^2\theta/\mathcal A_r)w^2$. Delete repeated forms, but retain
their incidences in the nine detectors. The five forms within one
detector are distinct, and the zero contact is the common form $u$.
Apply Theorem~\ref{thm:v5-simultaneous-primes} to these forms.
Prescribe arbitrarily high divisibility of $w$, a sufficiently
large lower bound for $u/w^2$, and primes larger than 76176 and
all prime factors of the fixed scales and contact differences.
In its odd squarefree avoidance modulus include every odd prime
factor of $3\cdot17L$ except 7.
Writing this modulus as $M$, the prescribed Gaussian classes are
\[
 \pi_p\equiv9\pmod{16},\qquad
 \pi_p\equiv1+3i\pmod7,\qquad
 \pi_p\equiv i\pmod M.
\]

Let the resulting Gaussian generators be $\pi_p$, and write
$d_-(p)=N(\pi_p-1)$ and $d_+(p)=N(\pi_p+1)$ for the cardinalities
of the two reduced elliptic groups, orienting the places by
Lemma~\ref{lem:v5-two-reductions}. These cardinalities annihilate
the reductions of $P_*$; their exact point orders are not needed.
They satisfy
\[
 v_2(d_-)=6,\quad v_2(d_+)=2,\quad
 \gcd(d_-(p),d_+(q))=4
\]
for every pair of selected primes. No odd prime factor of $L$
divides either order. For 7 this follows separately from
$d_+\equiv6$ and $d_-\equiv2\pmod7$.
Set $M_-=\operatorname{lcm}_p d_-(p)$ and
$M_+=\operatorname{lcm}_p d_+(p)$. The congruences
\[
 h\equiv0\pmod{M_-},\quad h\equiv4\pmod{M_+},\quad
 h\equiv h_0\pmod L
\]
are compatible: their pairwise gcds are $4,64,4$, respectively.
Choose a sufficiently large positive $h$ in this progression
attaining the real bounds. It avoids $0,4,-4$, so all point and
subtraction guards hold. At a selected prime, $H,U_i$ are formal
at the place of order $d_-$, while $V_0,V_i$ are formal at its
conjugate place. Their exact ratios are $n_i$. Lemma
\ref{lem:v5-unramified-log} gives all sixteen required orders,
including the order-52 polynomial test because $F(a,t,b)=0$.

We check that the nine preparation-zero poles can be filled.
For each detector,
\[
 R_r=\frac{\prod_{\theta\in\Theta}
                   (u-(\theta/\mathcal A_r)v)}
                 {v^5\delta_r^6}.
\]
The denominator is a rational square. At the five distinct contact
primes of this detector $R_r$ has valuation one, and elsewhere
its valuation is even. The primes do not divide $v$: otherwise
two distinct forms in the same detector would give the same
prime value, impossible since $w>0$ and their coefficients are
distinct. The lower bound for the primes excludes the fixed scale
and difference denominators. At either place over a
contact prime a factor of $W_r$ has valuation zero or negative,
and its inverse has residue one or zero, respectively. At the
successful place every inverse factor has residue one.
Consequently $T_r$ has residue one or two and is a unit.
It follows that $B_r$ has positive odd valuation at its contact
primes and even valuation elsewhere. Its reduced denominator is
a square.

The prescribed Gaussian classes make every contact prime square
in $\Q_2,\Q_3,\Q_{17}$. The last product formula therefore makes
$B_r$ square in each of those fields. The uniform real bounds
and $\lambda\longrightarrow+\infty$ give
\[
 R_r\sim\mathcal A_r^{-6}\lambda^{-1},\quad
 W_r=1+O(\lambda^{-3}),\quad T_r=2+O(\lambda^{-3}),\quad
 B_r\sim(4\mathcal A_r^6\lambda)^{-1}.
\]
In particular $0<B_r<17$. At a small prime $p=2,3,17$, put
$\nu_p=v_p(v)$ and $\eta_{p,r}=v_p(\mathcal A_r)$.
The common form $u$ is a selected prime greater than these small
primes, so $v_p(u)=0$ and $v_p(\lambda)=-\nu_p$.
The fixed neighborhoods give $v(D_j)\ge-c_p$. For sufficiently
large prescribed $\nu_p$ one has
\[
 v_p(R_r)=\nu_p-6\eta_{p,r},\qquad
 v(\widetilde D_{r,j})\ge-c_p+(m_j+1)(\nu_p-\eta_{p,r}).
\]
The nonconstant factor term has valuation at least
$3\nu_p+(10m_j-8)\eta_{p,r}-2c_p$, tending to infinity.
Thus $T_r$ is as close to 2 as needed, including at the ramified
prime 2, and
\[
 v_p(B_r)=\nu_p-6\eta_{p,r}-2v_p(2).
\]
Prescribing the divisibility of $w$ before applying the prime theorem
therefore ensures $v_2(B_r)\ge3$, $v_3(B_r)\ge1$,
$v_{17}(B_r)\ge2$. Lemma~\ref{v5:preparation-zero-bridge} supplies
the four witnesses in every coordinate pole.

Finally the progression preserved all scalar leaf conditions.
The constant leaf on the all-left path makes $W_\beta$ a unit at
every prime dividing $r_\beta$; the divisibility of the other
leaves makes it a unit at primes dividing $c$ but not $r_\beta$.
Since $|r_\beta|>1$, it is nonzero. Therefore
$b_\beta=cr_\beta W_\beta^2$ is positive, has square reduced
denominator, and is locally square at $2,3,17$.
Lemma~\ref{hyp:o8-pole} supplies the scalar preparation and pole
witnesses for this final $H$. The anchor already holds for $S$.
All witnesses have now been supplied without changing the root points.
\end{proof}

% Replace the former exact-compilation subsection with this fragment.
% Retain the preceding shared-point arithmetic source and its hypotheses.
% No claim here removes or proves its simultaneous-prime input.
\subsection{The quadratic compilation in 962 variables}
\label{subsec:v6-962-compilation}

We compile the shared-point source above using different point charts,
pole coordinates, and arithmetic circuits. Every change has explicit
rational witness maps. The arithmetic source, including its simultaneous
prime input, is unchanged. Throughout this subsection
\[
 K=\Q(\sqrt2),\qquad D=(3-\sqrt2)^2=11-6\sqrt2,
 \qquad \mathsf S=3^{2^{27}}.
\]
An equation in $K$ means its two rational component equations. The fixed
integer $\mathsf S$ is a coefficient, not a coordinate. An assignment of
a quadratic expression to one rational coordinate adds one quadratic
equation. Affine expressions require no coordinate.

\begin{theorem}[The 962-variable quartic]
\label{thm:v6-962-polynomial}
The complete table specified below defines
\[
 P_{962}\in\Z[X_0,\ldots,X_{961}]
\]
with degree four, exactly $55409$ nonzero coordinate monomials, integer
coefficient content one, and coefficient height
\[
 H(P_{962})=64\mathsf S^2=64\cdot3^{2^{28}}<2^{429496736}.
\]
It is the sum of squares of $919$ quadratic residuals and has the same
rational input projection as the shared-point source of
Theorem~\ref{thm:v5-source-projection}. Consequently, under that theorem's
arithmetic hypotheses, for every $x\in\Q$,
\[
 x\in\Z\quad\Longleftrightarrow\quad
 \exists X_1,\ldots,X_{961}\in\Q,\quad
 P_{962}(x,X_1,\ldots,X_{961})=0.
\]
The $962$ variables include the input and $961$ existential witnesses.
\end{theorem}

\paragraph{A quadratic chart for the nine points.}
For each point introduce $t,h\in K$ and set
\begin{equation}\label{eq:v6-point-chart}
 u=t^2,\qquad X=h^2,\qquad
 u(X-u/4)=D,\qquad Y=h(u/2-X).
\end{equation}
The three assignments use six rational coordinates. The displayed
constraint contributes two quadratic residuals. The four rational
components of $t,h$ replace the four halfpoint components.

\begin{lemma}[Exact point-chart maps]\label{lem:v6-point-chart}
The chart \eqref{eq:v6-point-chart} describes precisely the finite doubles
used in the shared-point source, with the same rational witnesses up to
the following invertible changes. It preserves all required nonzero
point-coordinate domains. The rational component $X_0$ of $X$ is positive.
\end{lemma}
\begin{proof}
For an old halfpoint $(a,b)$ with $b^2=a^3-Da$ and its tangent equation,
both $a,b$ are nonzero. Indeed $b=0$ would give $a=0$ or $a^2=D$,
whereas the tangent numerator $3a^2-D$ would be $-D$ or $2D$. Then
$a=0$ is excluded by the curve equation. Define
\[
 t=b/a,\qquad h=(a^2+D)/(2b)=a/t-t/2.
\]
The curve equation gives $t^2=a-D/a$. Substitution gives
$t^2(h^2-t^2/4)=D$. The doubling formulas give exactly $X=h^2$
and $Y=h(t^2/2-h^2)$.

Conversely, a chart solution has $t\ne0$ because $D\ne0$. Set
\[
 a=th+t^2/2,\qquad b=at,\qquad m=h+t.
\]
If $a=0$, then $h=-t/2$, contradicting the chart equation. Thus $b\ne0$.
Direct expansion gives
\[
 a^2-t^2a=D,\quad b^2=a^3-Da,\quad
 2bm=3a^2-D,\quad m^2-2a=h^2,
\]
and $m(a-h^2)-b=h(t^2/2-h^2)$. The substitutions are inverse, so they
recover exactly the old finite double. Lemma~\ref{lem:v5-point-domains}
therefore supplies its nonzero abscissa and ordinate, and the same remains
true for the unchanged guarded differences and subsequent doublings.

Both real embeddings of $D$ are positive. Since $t\ne0$, the chart
equation implies $h^2=t^2/4+D/t^2>0$ under both embeddings. Hence
$X_0$, the average of these two positive values, is positive.
\end{proof}

The five affine subtraction formulas keep their original inverses of
the differences of abscissas. Neither a subtraction guard nor a point
occurrence is deleted. Relative to the direct halfpoint doubling circuit,
the chart saves two rational arithmetic coordinates for each of the nine
points.

\paragraph{The logarithm polynomial.}
Retain the source's coefficients
\[
 a_0=1,\quad a_1=-2/5,\quad
 a_j=a_{j-1}(4j-5)/(4j+1).
\]
For $P(w)=\sum_{j=0}^{12}a_jw^j$ write
\begin{equation}\label{eq:v6-log-square-completion}
 P(w)=a_{12}Q(w)^2+R(w),\qquad
 \deg Q=6,\quad Q\text{ monic},\quad\deg R\le5.
\end{equation}
This identity is constructed without any witness division. Starting
with the leading coefficient of $Q$ equal to one, match degrees eleven
down to six successively. At each step the new coefficient is determined
by division by $2a_{12}\ne0$. The remainder then has degree at most five.
After computing $w^2,w^3$, every term through $w^6$ is quadratic in
$1,w,w^2,w^3$. Assigning $Q$ and then $a_{12}Q^2+R$ therefore uses
four field assignments, rather than five. The eight order-$52$
logarithms save sixteen rational coordinates. The order-$16$ logarithms
keep their exact degree-three polynomials. No truncation order changes.

Usually the final logarithm is $\ell=zP(Df^2)$, with
$f=1/x$ and $z=-x/y$. For $U_0$ and $V_0$ only, omit the separate $z$
coordinates and impose the direct quotient equations
\begin{equation}\label{eq:v6-log-quotient}
 y\ell=-xP(Df^2).
\end{equation}
The polynomial output is an affine expression in assigned coordinates,
so these are quadratic component equations. The ordinate is nonzero by
Lemma~\ref{lem:v5-point-domains}; hence \eqref{eq:v6-log-quotient}
is exactly the old logarithm. The source's logarithm-nonvanishing result
and every logarithm-ratio domain remain unchanged. In the scalar block
replace $z(H)\rho=z(S)$ by
\begin{equation}\label{eq:v6-scalar-ratio}
 x(H)\rho=-y(H)z(S).
\end{equation}
Both $x(H),y(H)$ are nonzero, so this has the same unique value
$\rho=z(S)/z(H)$.

\paragraph{Products, tests, and detector evaluation.}
For each of the three orders, compute the elementary symmetric
coefficients of $\prod_j(1+qD_j^2)$ by a balanced product tree.
At an internal node with coefficient vectors $a_i,b_j$, assign each
component of $\sum_{i+j=k}a_ib_j$ as one quadratic expression.
This is the ordinary coefficient-convolution identity. It preserves
every coefficient and every test, including its assigned order.

The degree-ten detector polynomial is evaluated in two blocks:
\begin{equation}\label{eq:v6-detector-block-evaluation}
 \sum_{j=0}^{10}c_jq^j
 =\sum_{j=0}^{4}c_jq^j
   +q^5\left(\sum_{j=0}^{5}c_{j+5}q^j\right).
\end{equation}
Compute $q^2,q^3,q^4,q^5$ using four scalar assignments, assign the
two components of the parenthesized block, and assign the two final
components directly. The low block is not assigned separately.
Every final summand is quadratic in the available coordinates. This
uses eight rational coordinates. The degree-two and degree-four
polynomials use their direct power sums, with their two component
sums assigned once. These formulas are identities over $K$.

The Lucas expression is regrouped as
\[
 t^2-a(a-2)tb+((a-2)b)^2-1
 =t^2-(at)C+C^2-1,\qquad C=(a-2)b.
\]
Thus one field multiplication is removed without changing the tested
polynomial. The fixed quintic is evaluated by its exact coefficients
using $t^2$; it remains
$F_5(t)=t(t+1/2)(t+1)(t-1/2399)(t+32/407)$.

\paragraph{Weight quotients and their domains.}
Fix a detector and write $A$ for its scale, $\delta=A\lambda-1$,
and $e$ for its weight. The scales are
\[
 1,\quad x(Q)_0^2,\quad (1+x(Q)_1^2)^2
 \qquad(Q=A,S,T,Z).
\]
They are all positive by Lemma~\ref{lem:v6-point-chart}. Retain the
inverse $\delta i=1$, and replace the separate fifth-power product
and inverse of $F_5$ by
\begin{equation}\label{eq:v6-weight-domain}
 F_5(A\lambda)e=(A\delta)^4(A\delta).
\end{equation}
Here $(A\delta)^4$ denotes the output of its retained power-assignment
gates; the displayed product is quadratic in the available coordinates.
The right side is nonzero, so this equation forces both $F_5(A\lambda)$
and $e$ nonzero. In particular the removed inverse of $F_5$ is recovered
uniquely. Define the low-order weight by
\[
 \delta q_1=e i^2,
\]
where $i^2$ is likewise a retained square-assignment output. This gives
the unchanged $q_1=e/\delta^3$. Retain the repeated-squaring
values $e^6,e^{24},e^{96}$ and set
\[
 c=q_1e^6,\qquad q_{16}=ce^{24},\qquad q_{52}=ce^{96}.
\]
These are the old $q_1e^{30}$ and $q_1e^{102}$, with one shared product.

If the detector product is $W=w_0+w_1\sqrt2$, retain $w_0j=1$ and use
\begin{equation}\label{eq:v6-norm-final-weight}
 w_0n=(w_0^2-2w_1^2)/2,\qquad ev=in,\qquad B=nv.
\end{equation}
The trace guard is essential and is retained. It gives $W\ne0$ in $K$,
so $\operatorname{Norm}(W)\ne0$ and
$n=\operatorname{Norm}(W)/(2w_0)\ne0$. Equations
\eqref{eq:v6-weight-domain}--\eqref{eq:v6-norm-final-weight} therefore
give
\[
 B=\frac{i}{e}\left(\frac{\operatorname{Norm}(W)}{2w_0}\right)^2\ne0,
\]
the original detector input. The reverse map restores $R=i/e$, its
old inverse intermediates, and $n^2$ by their unique rational values.
No inverse is inferred from an unguarded zero-over-zero quotient.

\paragraph{The scalar norm and finite corrections.}
The scalar contact polynomial $r$ and its inverse $ri=1$ remain.
Writing $f=f(H)_0$, $u=f(H)_1$, and $g=f(S)_0$, retain the leaf
$\ell=ui$ and replace two norm-tree
nodes by
\[
 rn_0=r+f^2,\qquad rn_1=u\ell+g^2.
\]
with the existing leaf $ui$ retained. Since $r\ne0$, these equal
$n_0=1+r(fi)^2$ and $n_1=(ui)^2+r(gi)^2$. All other nodes and all
orders in the scalar norm tree are unchanged.

Replace the eight Boolean correction choices and their seven product
coordinates by $c_0,\ldots,c_7\in\Q$ satisfying
\begin{equation}\label{eq:v6-correction-chain}
 c_0^2=1,\qquad
 (c_j-c_{j-1})(c_j-p_jc_{j-1})=0\quad(1\le j\le7),
\end{equation}
where $(p_1,\ldots,p_7)=(2,3,5,7,11,13,17)$. Induction shows
$c_j\ne0$ and $c_j/c_{j-1}\in\{1,p_j\}$. Thus $c_7$ is exactly an
allowed signed squarefree correction. Conversely every original
choice determines such a chain. Recover the original Booleans by
\[
 \epsilon_-=(1-c_0)/2,\qquad
 \epsilon_j=\frac{c_j/c_{j-1}-1}{p_j-1}.
\]
All denominators are nonzero. The scalar preparation coordinate,
its defining quotient, and the explicit inverse of its pole input
are retained.

\paragraph{The ten pole blocks.}
Write $b$ for a pole input, $C=17b$, and $U$ for its prepared value.
For the nine coordinate poles $U=b$. In the scalar pole $U$ retains
the source's preparation formula and its denominator domain.
Every $C$ is nonzero: this was proved above for coordinate detectors,
and the scalar input has its retained inverse.

Replace the old witness $e_0$ by $z=Ce_0$. Introduce the five
intermediates
\[
 P=p^2,\quad F=U^2,\quad H=z(C-F),\quad V=s^2,\quad Q=q^2,
\]
and impose
\begin{equation}\label{eq:v6-scaled-poles}
 CP=(C-z)(C+3z),\qquad CQ=(V-H)(V+3H).
\end{equation}
Every equation is quadratic. Under $z=Ce_0$, the first residual is
$C$ times the old residual $p^2-C(1-e_0)(1+3e_0)$, and the second
is unchanged because $H=Ce_0(C-U^2)$. The inverse map $e_0=z/C$
is defined. This saves one arithmetic coordinate at each pole,
without changing the pole predicate or any preparation condition.

\paragraph{Counts and the final polynomial.}
The initial allocation is
\[
\begin{array}{l|r|r}
 &\text{coordinates}&\text{structural residuals}\\\hline
 \text{input}&1&0\\
 \text{nine point charts}&36&18\\
 \text{common slope}&1&0\\
 \text{nine coordinate poles}&36&18\\
 \text{scalar pole with preparation}&5&2\\
 \text{private anchor}&3&1\\
 \text{signed correction chain}&8&8\\\hline
 \text{total}&90&47
\end{array}
\]
The complete program has $904$ arithmetic coordinates and their
$904$ defining equations. Hence the raw counts are $994$ coordinates
and $951$ quadratic residuals.

The private anchor uses exactly the old affine change
\[
 u=\mathsf S+(U-3V+3W)/4,\quad
 v=\mathsf S/2+V/4,\quad w=3\mathsf S/2+W/4.
\]
Its determinant is $1/64$. If $T=f_0^2+f_0f_1-f_1^2$, sixteen times
the old anchor after substitution is
\[
 -48T^2+U^2-6UV+6UW+15V^2-18VW+7W^2+8\mathsf S U.
\]
No anchor condition or witness is removed.

Each certified compiler elimination solves a row $aY+q=0$ with
$a\in\Q^*$ and $q$ independent of $Y$, substitutes $Y=-q/a$
everywhere, and removes that row. It is accepted only if every
remaining residual is still quadratic. The reverse map reinserts
the unique displayed value of $Y$. There are exactly $32$ such
eliminations, no row additions, and no omission of an unused coordinate.
The input is retained. Consequently
\[
 90+904-32=962,\qquad47+904-32=919.
\]
Normalize the remaining rows to primitive integer polynomials
$q_1,\ldots,q_{919}$ and define
\begin{equation}\label{eq:v6-final-polynomial}
 P_{962}(X_0,\ldots,X_{961})=\sum_{\nu=1}^{919}q_\nu^2.
\end{equation}
The complete coefficient table is
\nolinkurl{quartics/quartic-962/polynomial.json}, with SHA-256
\begin{center}\small
\nolinkurl{6b798de7d00ac3c7d725ca0a69d39f1eac1ddb02fee2756cb6323770b952af32}.
\end{center}
An entry with coefficient $c$, exponent \texttt{sPower} $k$, and
coordinate list $[i_1,\ldots,i_d]$ means
$c\mathsf S^kX_{i_1}\cdots X_{i_d}$. Every entry is included and repeated
indices give powers. This defines an ordinary integer polynomial.

\begin{proof}[Proof of Theorem~\ref{thm:v6-962-polynomial}]
For a source solution apply the point-chart, correction-chain and
pole substitutions, and evaluate all new assignment and quotient
coordinates in dependency order. The identities and domain proofs
above show that every structural residual vanishes. Conversely apply
the inverse substitutions, restore every deleted quotient by its
proved nonzero denominator, and evaluate the source assignments.
Every source point, logarithm, test, detector input, scalar value and
anchor value is recovered. These maps fix $X_0$. The affine anchor
change, certified eliminations and nonzero row normalizations preserve
the same projection. Over $\Q$, \eqref{eq:v6-final-polynomial} vanishes
exactly when all $919$ residuals vanish. This proves the projection
statement, and Theorem~\ref{thm:v5-source-projection} gives its conditional
arithmetic interpretation.

Independent exact convolution gives $55409$ distinct monomials.
Of their coefficients, $55401$ are $\mathsf S$-independent integers,
seven are $c\mathsf S$ with $0<|c|\le768$, and one is
$64\mathsf S^2$. The largest absolute $\mathsf S$-independent coefficient
has bit length $1526$. Thus no coefficient vanishes on specialization
and $64\mathsf S^2$ is the exact height. It occurs at $X_{79}^2$,
where $X_{79}=U$. The coefficient of $X_0^4$ is one, proving degree
four and integer content one. Finally
$\log_2(64\mathsf S^2)<6+(8/5)2^{28}<429496736$ because $3^5<2^8$.
\end{proof}

The frozen package includes the builder, raw systems, all $904$ gate
definitions, the exact elimination and normalization certificate,
compact coordinate map, final residuals, complete coefficient table,
source-equivalence proofs and independent reviews. Its manifest has
SHA-256
\begin{center}\small
\nolinkurl{36fca72cb0eb126193ffa43d890ff09d5e29cbf1274561d2632ed1e227163787}.
\end{center}
The independent checker verifies every gate, all $894$ certificate
operations, the complete sum-of-squares expansion, support,
integrality, content and height. A fresh isolated reproduction gives
the same seven deterministic artifacts byte for byte. These finite
checks establish the stated algebraic facts; they do not prove any
inherited arithmetic input.


% End Version 5 main additions
\section{Why a single defining polynomial must have degree at least four}
\label{sec:degree-lower-bound}

The argument in this section comes from the agentic notes
\cite[section on unconditional lower bounds]{agentic-improvements}.
We give its specialization to $\Q$, including the treatment of singular
cubics, cones and geometric reducibility.  It is independent of the
arithmetic inputs to our constructions.

For $F\in\Q[T,Y_1,\ldots,Y_r]$, write
\[
 S_F=\{t\in\Q:\exists\,y\in\Q^r,\ F(t,y)=0\}.
\]
All degrees below are total degrees, including the input variable $T$.

\begin{theorem}[Degree obstruction]
\label{thm:degree-four-necessary}
If $S_F=\Z$, then $\deg F\geq4$.
\end{theorem}

We use two published results.  The integral form of Hilbert's
irreducibility theorem says that an irreducible polynomial
$g(T,Y)\in\Q(T)[Y]$ admits infinitely many specializations $T=t\in\Z$
which preserve its degree in $Y$ and are irreducible over $\Q$.
Deleting finitely many values also avoids its coefficient poles and
any prescribed nonzero polynomial in $t$.
This integral formulation is recalled in the introduction of
M\"uller~\cite{Muller2002HIT}; merely obtaining specializations in
$\Q$ would not suffice here.
The second result is the characteristic-zero cubic theorem of
Colliot-Th\'el\`ene and Salberger
\cite[\S1, Proposition~1.3, pp.~520--521]{CTS1989Cubic}:
a geometrically integral cubic hypersurface of dimension at least two
over a characteristic-zero field is unirational over that field if it
has a rational point and is not a cone. This includes singular cubics;
the proof treats the codimension-one singular case and uses
\cite[Remark~2.3.1, p.~56]{CTSSD1987} for the regular-in-codimension-one case.
We apply it only over $\Q$, after removing every vertex direction.
The exact source and quotient-hypothesis connection is recorded in
\cite{cubic-source-project}. In particular, no unirationality assertion
is made for a cone over a plane cubic.

\begin{lemma}\label{lem:nonconstant-rational-not-integral}
A nonconstant rational function $R\in\Q(U_1,\ldots,U_m)$ cannot take
only integer values at the rational points of a nonempty Zariski-open
subset of $\mathbb A^m_{\Q}$ on which it is defined.
\end{lemma}
\begin{proof}
Shrink the open subset to a nonempty principal open subset and choose
a connected real open ball in it, avoiding the denominator of $R$.
Rational points are dense in this ball.  If their images belong to the
closed set $\Z$, continuity puts the image of the entire ball in $\Z$.
Connectedness then makes $R$ constant on the ball.  Its numerator minus
that constant times its denominator vanishes on a real open set, hence
is the zero polynomial.  This contradicts nonconstancy.
\end{proof}

Consequently, if an integral variety $X/\Q$ is unirational and
$f\in\Q(X)$ is nonconstant, the rational values of $f$ cannot all be
integers.  Indeed, compose $f$ with a dominant rational map from
projective space to $X$, and restrict to a nonempty affine open subset
of its domain.  Dominance keeps this composite nonconstant, so
Lemma~\ref{lem:nonconstant-rational-not-integral} applies.

\begin{lemma}\label{lem:quadratic-projection-finite}
If $\deg F\leq2$ and $S_F\subseteq\Z$, then $|S_F|\leq2$.
\end{lemma}
\begin{proof}
Otherwise choose rational zeros $a,b,c$ with distinct input
coordinates.  If they lie on an affine line, the restriction of $F$
to that line has three zeros and degree at most two, so vanishes
identically.  The input is nonconstant on the line and takes every
rational value, a contradiction.

If the points are not collinear, let $L$ be their rational affine
plane.  If $F|_L=0$, a line in $L$ with nonconstant input gives the same
contradiction.  If $F|_L$ factors over $\Q$ into affine linear
polynomials, its zero set consists of at most two rational lines.
One contains two of the three selected points and again has
nonconstant input.  A nonzero polynomial of degree at most one
cannot vanish at three noncollinear points.

It remains to consider an irreducible quadratic on $L$.  Its
projective conic has a smooth rational point: a reduced irreducible
quadratic in characteristic zero has at most one geometric singular
point, while the three chosen affine points are distinct.  A conic
with a smooth rational point is rational.  Explicitly, intersecting
it with the pencil of lines through that point gives a rational
parametrization by taking the second intersection point.  The input
function is nonconstant because it takes three different values.
Composing it with this parametrization contradicts
Lemma~\ref{lem:nonconstant-rational-not-integral}.
\end{proof}

\begin{proof}[Proof of Theorem~\ref{thm:degree-four-necessary}]
Suppose $S_F=\Z$ and $\deg F\leq3$.  The polynomial is neither zero
nor a nonzero constant.  Lemma~\ref{lem:quadratic-projection-finite}
excludes degree at most two.  It also excludes a reducible cubic:
every nonconstant irreducible factor then has degree at most two,
and the finite union of their projections cannot equal $\Z$.
Thus $F$ is irreducible of degree three.

Homogenize $F$ to $H$ on a rational vector space $E$, using $X_0$ as
the homogenizing coordinate and $X_1$ for the input coordinate.
Thus on $X_0\ne0$ the input is $X_1/X_0$.
The homogeneous polynomial $H$ is irreducible over $\Q$.
If it is geometrically reducible, its geometric irreducible factors
form a single Galois orbit with equal degrees.  In characteristic
zero they occur without repetition.  Since the total degree is
three, there are exactly three conjugate linear factors.
Every rational zero lies on each of their hyperplanes.
Take rational zeros normalized by $X_0=1$ with inputs zero and one.
Their whole affine joining line lies on all three hyperplanes and
its input takes every value in $\Q$, a contradiction.
We may therefore assume $H$ is geometrically integral.

Define the rational vertex subspace
\[
 W=\{v\in E:D_vH=0\text{ identically}\},\qquad
 D_vH(x)=\left.\frac{d}{ds}H(x+sv)\right|_{s=0}.
\]
These are linear conditions on $v$.  In characteristic zero,
$D_vH=0$ is equivalent to $H(x+sv)=H(x)$ for all $x,s$.
Thus $H$ descends to a cubic $\overline H$ on $E/W$.
After a basis of $W$ is extended to one of $E$, this simply deletes
the coordinates on which $H$ does not depend.
The quotient has no remaining vertex direction.  This assertion
holds over the algebraic closure as well, since the kernel of a
matrix over $\Q$ commutes with extension of scalars.  Hence its
projective hypersurface is not a cone.

For $v\in W$ and a rational zero $x$ with $x_0=1$, translation by
$sv$ gives rational zeros whose inputs, whenever defined, are
\[
 \frac{x_1+s v_1}{1+s v_0}.
\]
Unless $v_1=x_1v_0$, this is a nonconstant fractional linear function
of $s$, with cofinite image in $\Q$.  Its image cannot be contained
in $\Z$.  Applying the resulting equality first at input zero and
then at input one gives $v_0=v_1=0$ for every $v\in W$.
Therefore $X_0$ and $X_1$ descend to independent linear forms on
$E/W$.

Choose quotient coordinates beginning with these forms.
The affine chart $X_0=1$ of $\overline H=0$ has exactly the same
input projection as $F=0$.  To see both inclusions, project rational
points from $E$ to $E/W$, and in the reverse direction use a rational
linear section of that projection.  This section preserves $X_0$
and $X_1$, and $H$ is constant along each fiber.

If the quotient projective hypersurface has dimension at least two,
\cite[Proposition~1.3]{CTS1989Cubic} applies. It is geometrically integral,
has a rational point in $X_0=1$, is not a cone, and its ground field $\Q$
has characteristic zero. It is therefore unirational. Its input function is
nonconstant, since its rational values include all integers.
The consequence of Lemma~\ref{lem:nonconstant-rational-not-integral}
then contradicts that its input projection is contained in $\Z$.

If the quotient has dimension one, its affine chart is defined by
an irreducible polynomial $G(T,Y)\in\Q[T,Y]$.  If its degree in $Y$
is zero, its input projection is finite.  If that degree is one,
solving the linear equation for $Y$ gives a rational point above
every rational $T$ outside a finite set, also impossible.
If its degree in $Y$ is at least two, Gauss's lemma makes it
irreducible in $\Q(T)[Y]$.  Integral Hilbert irreducibility supplies
an integer $t$ for which $G(t,Y)$ remains irreducible of that same
degree.  It has no rational root, contradicting $S_G=\Z$.
Finally, a zero-dimensional projective quotient has only finitely
many rational points and cannot have input projection $\Z$.
These exhaust the possibilities and prove the theorem.
\end{proof}

\begin{remark}
This lower bound concerns a single polynomial equation and counts
degree in every variable.  It allows systems of quadratic equations;
over $\Q$, their sum of squares is quartic.  It puts no lower bound
here on the number of variables in a degree-four definition.
\end{remark}


\appendix
\renewcommand{\thesection}{A\arabic{section}}
\section{Dependency order in the arithmetic argument}\label{app:obligations}
The even coefficient bound uses stationary-base selection as well as
integrality. The odd minimum uses a fixed nonzero witness and the prior
genus-index lower bound. Support removal is applied only to the actual
unramified specialization. The functional bound and determinant clearing
are proved separately, and the same corrected coordinate is used at every
address. These constructions are given in Sections~\ref{sec:rc}--\ref{sec:limits}
and Appendices~\ref{app:even-source}--\ref{app:endpoints}.

The conversion to ordinary integers uses the canonical model and total
multiplication, followed by the exact computation relation. Only the
shared-shift separation argument in Section~\ref{sec:bridge} identifies
the external rational variable with the ordinary value of its integer
index. The literal substitution and size calculation then apply to the
same expression $P_*$.

\section{Source identities and scope}\label{app:sources}
The following catalogue records the base collection of 83 source
identifiers; parity, construction and uniformity supplements appear
separately. The released construction files and their exact identities
are recorded in \nolinkurl{construction/SOURCE_MANIFEST.json}.
Current proof-source selection and versioned supporting arguments are
indexed in \nolinkurl{research/version14/SUPPORTING-NOTES.md}.
\begin{longtable}{p{.07\textwidth}p{.58\textwidth}p{.25\textwidth}}
\toprule ID & Source basename & Use\\\midrule\endhead
S01 & \texttt{ordinary-integer-polynomial.js.txt} & Outer construction\\
S02 & \texttt{rational-circuit-translator.js.txt} & Substitution semantics\\
S03 & \texttt{finite-polynomial-expressions.js.txt} & Fixed constructors\\
S04 & \texttt{variable-shift-polynomial-library.js.txt} & Close and Test\\
S05 & \texttt{elliptic-generator-proof.txt} & Generator certificate\\
S06 & \texttt{explicit-elliptic-data.txt} & Rank, torsion, codes\\
S07 & \texttt{verify-elliptic-generator.js.txt} & Exact finite algorithm\\
S08 & \texttt{six-contact-height-proof.txt} & Contacts and constants\\
S09 & \texttt{z-in-q-height-independent-audit.txt} & Height proof\\
S10 & \texttt{explicit-pole-predicate.txt} & Pole polynomial\\
S11 & \texttt{z-in-q-pole-completeness-audit.txt} & Simultaneous descent\\
S12 & \texttt{parity-frozen-4eaee875.tex} & Raw local matrices\\
S13 & \nolinkurl{multiplication-certificate-correctness-proof.txt} & Multiplication proof\\
S14 & \texttt{twenty-two-prime-pattern-proof.txt} & Selected primes\\
S15 & \texttt{z-in-q-variable-shift-bridge.txt} & Ordinary-value graph\\
S16 & \texttt{z-in-q-ordinary-integer-assembly-proof.txt} & Outer composition\\
S17 & \texttt{z-in-q-restricted-converse-assembly.txt} & RC assembly inputs\\
S18 & \texttt{ring-limits-frozen-4eaee875.tex} & Finite models and limits\\
S19 & \texttt{paired-functional-and-uniformity-proof.txt} & Paired functional\\
S20 & \texttt{determinant-transfer-audit.txt} & Clearing and transfer\\
S21 & \texttt{variable-shift-precision.py.txt} & Exact decision program\\
S22 & \texttt{precision-register.js.txt} & Ideal IR and lowering\\
S23 & \texttt{specialize-precision.js.txt} & Wrapper specialization\\
S24 & \texttt{integer-halting-correctness-proof.txt} & Natural compiler\\
S25 & \texttt{z-in-q-explicit-cutoff-bound.txt} & Literal cutoff\\
S26 & \texttt{z-in-q-polynomial-size-bound-proof.txt} & Syntax/norm argument\\
S27 & \texttt{z-in-q-numerical-polynomial-bounds.txt} & Tower comparison\\
S28 & \texttt{determinant-switch-audit.txt} & Switches and outer rank\\
S29 & \texttt{ring-determinants-frozen-4eaee875.tex} & Ring data and derivatives\\
S30 & \texttt{preliminaries-frozen-4eaee875.tex} & Selmer coranks and lengths\\
S31 & \texttt{fixed-certificate-constants.js.txt} & Exact least-multiple search\\
S32 & \texttt{pointwise-frozen-4eaee875.tex} & Distinct point indices\\
\bottomrule
\end{longtable}
In particular S30, S18, S29 and S32 are, respectively, Sections
2, 4, 5 and 9 of OpenAI's \emph{A pointwise 2-converse for elliptic
curves with rational two-torsion}~\cite{openaipointwise}, frozen at
\texttt{openai/math} commit
\nolinkurl{4eaee87548250b2682affda7d05d947c05554bae}.
The source filenames are provenance labels for those sections, not
independent publications or independent arithmetic assumptions.

The supplement manifest also records the revised fixed-degree CM proof,
the complete natural compiler proof, actual arithmetic and analytic
producer arguments, finite-Artin and determinant-count clarifications,
the fixed BFH witness application, exact substitution checks, and the
new fixed-form, weighted-theta, operator and normalization derivations.
For the height estimate, the derivation in Section~\ref{sec:height}
supersedes the historical arithmetic-height arguments in S08--S09 and
the earlier CM supplements. Those historical alternatives are not
additional premises. Primary theorem applications remain explicit in
the mathematical text.
The version records identify the source files used in the construction.

% End 05-bounds-appendices.tex


% Begin 06-literal-construction.tex
\section{Literal finite-constructor specification}\label{app:literal}
The supplementary file \nolinkurl{ordinary-integer-definition.json}
is part of this specification. It is a finite directed acyclic description,
not an expanded monomial list. The following rules determine one polynomial
from its bytes. Diagnostic labels and intended-theorem strings add no equations.

The root is \texttt{integers}. Variable zero is the external input and
variables one through $292890231904$ are its rational witnesses.
A nonnegative point reference $r$ means $X_r,\ldots,X_{r+4}$.
References $-1,-2,-3$ respectively mean
\[
 (1,0,0,0,0),\quad(0,g_0,-145,g_1,g_2),\quad
                      (0,g_0,-145,-g_1,-g_2).
\]
Every model recipe conjoins
$8g_0-1641=0$, $4g_1+15465=0$, $32g_2-87483=0$.

\subsection{Scalar nodes and fixed expressions}
Read nodes in their stored topological order.
A \texttt{const} is its decimal integer; \texttt{var} is its local variable;
\texttt{add} and \texttt{mul} are ordinary polynomial operations.
For a \texttt{call} to a template with $p$ parameters, substitute the
specified node polynomials for variables $0,\ldots,p-1$, and $X_{w+j}$
for variable $p+j$, where $w$ is the witness start. The value is the
callee's entire root polynomial.

A \texttt{fixed\_integer} evaluates a finite expression using literals,
references to preceding definitions, addition, multiplication,
nonnegative powers, maximum, factorial, the catalogue product or maximum,
and the least formal multiple. The catalogue uses the first 22 primes
in increasing order in each unit class. The searches are deterministic
and terminate by Section~\ref{sec:size}. They produce constants, not
quantified witnesses.

A \texttt{psi\_component} uses the division-polynomial recurrence in
Section~\ref{sec:syntax}, pair multiplication
$(a,b)(c,d)=(ac+2bd,ad+bc)$ and $A=(-11243,7950)$.
It returns component zero or one; every index is fixed.
The positive fixed point codes use
\[
 x_n=x(P_0)-\frac{\psi_{n+1}(P_0)\psi_{n-1}(P_0)}{\psi_n(P_0)^2},
 \qquad y_n=\frac{\psi_{2n}(P_0)}{2\psi_n(P_0)^4}.
\]
Nontorsion makes these fixed denominators nonzero.
Negative indices negate both ordinate components.

A \texttt{finite\_guard\_product} ranges over
$(a,b,c)\in[-T,T]^3\cap\Z^3$ with $ab\ne c$, with fixed $T$.
Write the fifteen coordinates of the three codes as reduced
$A_j/B_j$, $B_j>0$, and put $D=\prod_jB_j$. Its factor is
\[
                 \sum_{j=1}^{15}(DX_j-(D/B_j)A_j)^2.
\]
Multiply these factors, with empty product one. This is an integer
polynomial with precisely the finite false-triple zero set used above.

\subsection{Recipe constraints and witness scope}
A model-evaluation recipe sums squares of the generator equations,
listed natural calls, scalar-constant calls, add/multiply/power-value calls,
and the five coordinate differences of \texttt{rootEquality}.
A source-variable row with a non-null natural witness start contributes
exactly that natural call. Operation kinds zero, one and two select
addition, multiplication and power value. Source-node and provenance
fields add no equations.

A sum-of-squares composition sums squares of its callee polynomials;
a product composition multiplies them. Parameter \texttt{[var,i]} means
$X_i$. Each private block uses its recorded start and length.
The checked ranges are disjoint from parameters and other private blocks.
These rules include globally conjoined fixed equations in active callees.

\subsection{The saved templates and the mathematical predicates}
\label{subsec:v10-template-dictionary}
The following dictionary identifies the relations encoded by the saved
templates, in addition to the decoding rules above. Point arguments are
canonical five-coordinate codes. Field values have their two rational
components in the basis $(1,\sqrt2)$; conjugation negates the second
component. An inverse equation guards each displayed quotient.
Conjunction is encoded by a sum of squares, and disjunction by a
product, with the private witness blocks specified in the file.
The identities are identities of zero sets over $\Q$; squaring a
callee polynomial can change the literal polynomial without changing
the relation it defines.

\begin{longtable}{p{.29\textwidth}p{.65\textwidth}}
\toprule Saved template & Relation and constraints\\\midrule\endhead
\nolinkurl{member} & Membership in the canonical code domain $D=2E(K)$
of Section~\ref{sec:model}.\\
\nolinkurl{raw_add}, \nolinkurl{raw_z}, \nolinkurl{raw_inverse_x} &
The group-law relation on typed inputs, $z=-x/y$, and $w=1/x$,
respectively, with $z(O)=w(O)=0$.\\
\nolinkurl{log25} & The two-component equality to the displayed
scaled degree-25 polynomial $L$ in Section~\ref{sec:mult}.\\
\nolinkurl{pole} & $Q_{\rm pole}(B)^2$, where $Q_{\rm pole}$ is the
eleven-witness polynomial of Section~\ref{sec:pole}. The outer square
comes from the singleton conjunction added by the rational builder.\\
\nolinkurl{test_group} & The twelve field tests of
Section~\ref{sec:mult}, including $A_i-\sigma(A'_i)$ for $i=1,2,3$.
Its four input points are $[c_*]P_0,[c_*]P_1,[c_*]P_2,[c_*]P_3$.
It types the first four auxiliary points and obtains the second four
by subtraction; its two denominator points and logarithms are nonzero.\\
\nolinkurl{slope} & The normalized twelve-test products in six private
groups, their traced inverses, and $\Phi(R/\Theta^2)$ from
Section~\ref{sec:mult}, with all contact and denominator guards.\\
\nolinkurl{height_certificate} & The 27 slopes summing in triples to
the nine homogeneous coordinate quantities in that section, together
with membership of the three input points.\\
\nolinkurl{multiply.scalar}, \nolinkurl{multiply.guard},
\nolinkurl{multiply} & The fixed $[c_*]$ graph; the inverse of the
finite false-triple product; and their conjunction with the height
certificate at the fixed $\ell_0$.\\
\nolinkurl{natural} & Four model squares, formed by multiplication,
whose sum is the input code, as in Section~\ref{sec:compiler}.\\
\nolinkurl{value.signed_encode} & $N=\iota(n)$ and a natural code
$C=\iota(c)$ with $n=c/2$ for even $c$, and $n=-(c+1)/2$ for odd
$c$. The two branches use $c=2q$ and $c=2q+1$ with $q\ge0$.\\
\nolinkurl{value.abs} & $a=|x|$: $a$ is a sum of four rational
squares and $a^2=x^2$.\\
\nolinkurl{value.close_coordinates} & The strict coordinate window
of Section~\ref{sec:bridge}, using the typed points $A,L,J$,
$r=\operatorname{Tr}(z(J)/z(A))/2$ and
$\epsilon=|(z(L))_0|+|(z(L))_1|>0$, with $|t-r|<\epsilon$.
This template alone does not assert the precision relation.\\
\nolinkurl{value.test} & Exactly $\Test(N,\kappa,t)$ of
Section~\ref{sec:bridge}: $C_1=[\kappa n]P_0$,
$V_0=H-\kappa P_0$, $V_1=U-C_1$, the six displayed raw-$z$ tests,
their normalization by $F(t)$, and the single guarded pole input.
These ratios use $z$, not $L$.\\
\nolinkurl{value.exceptions} & The six branches
$N=\iota(c)$ and $t=c$, $c\in\Cset$, with the constant points
computed by the fixed group-law chains.\\
\nolinkurl{scalar.integer_constant_}\,$i$ & The fixed integer
scalar graph requested by the translated precision computation.\\
\bottomrule
\end{longtable}

The composition recipes make the remaining identifications explicit.
\texttt{power\_value} translates the source power template with its
three input codes and output code; \texttt{precision} translates
the five-natural-input precision polynomial, in the order
$(n,\kappa,a,\ell,j)$ after signed encoding. The recipe
\texttt{close} introduces $A,L,J$, forms those five signed encodings,
conjoins \texttt{precision}, and calls
\nolinkurl{value.close_coordinates}. Thus it is the full $\Close$.
The recipe \texttt{main} introduces one shared $\kappa$ and conjoins
\texttt{close} and \nolinkurl{value.test} at exactly that same code.
The product recipe \texttt{value\_union} joins this branch with
\nolinkurl{value.exceptions}; \texttt{value} also imposes
\texttt{member} on $N$. Finally \texttt{integers} existentially
introduces $N$ and calls \texttt{value} with external coordinate
$t=X_0$. These are precisely $\Val$ and its input projection.
The ordered calls and private ranges are checked by
\nolinkurl{auditOrdinaryIntegerPolynomial} in
\nolinkurl{construction/ordinary-integer-polynomial.js}.
\nolinkurl{scripts/check_construction.py} regenerates the saved
expression byte for byte. The dictionary supplies the mathematical
identification; those finite checks verify the stored syntax and
wiring, not the arithmetic theorems used to prove the predicates.

The following parameter and witness counts are extracted from the exact
specification, without expanding its polynomial.
\begin{longtable}{p{.47\textwidth}rrr}
\toprule Template & Parameters & Witnesses & Root node\\\midrule\endhead
\nolinkurl{member} & 5 & 7 & 149\\
\nolinkurl{raw_add} & 15 & 6 & 308\\
\nolinkurl{raw_z} & 7 & 1 & 68\\
\nolinkurl{raw_inverse_x} & 7 & 0 & 54\\
\nolinkurl{log25} & 4 & 0 & 229\\
\nolinkurl{pole} & 1 & 11 & 106\\
\nolinkurl{test_group} & 44 & 162 & 515\\
\nolinkurl{slope} & 22 & 1,298 & 2681\\
\nolinkurl{height_certificate} & 36 & 35,106 & 260\\
\nolinkurl{multiply.scalar} & 10 & 1 & 136\\
\nolinkurl{multiply.guard} & 15 & 1 & 22\\
\nolinkurl{multiply} & 15 & 35,134 & 83\\
\nolinkurl{natural} & 5 & 140,604 & 79\\
\nolinkurl{value.signed_encode} & 10 & 140,654 & 138\\
\nolinkurl{value.abs} & 2 & 4 & 23\\
\nolinkurl{value.close_coordinates} & 16 & 53 & 271\\
\nolinkurl{value.test} & 11 & 35,282 & 775\\
\nolinkurl{value.exceptions} & 6 & 406 & 487\\
\nolinkurl{scalar.integer_constant_0} & 10 & 1 & 136\\
\nolinkurl{scalar.integer_constant_1} & 10 & 1 & 136\\
\nolinkurl{scalar.integer_constant_2} & 10 & 1 & 136\\
\nolinkurl{scalar.integer_constant_3} & 10 & 1 & 136\\
\nolinkurl{scalar.integer_constant_4} & 10 & 1 & 136\\
\bottomrule\end{longtable}

\begin{longtable}{lrrr}
\toprule Recipe & Parameters & Witnesses & Total variables\\\midrule\endhead
\nolinkurl{power_value} & 20 & 17,786,698 & 17,786,718\\
\nolinkurl{precision} & 25 & 292,889,492,836 & 292,889,492,861\\
\nolinkurl{close} & 11 & 292,890,196,199 & 292,890,196,210\\
\nolinkurl{main} & 6 & 292,890,231,486 & 292,890,231,492\\
\nolinkurl{value_union} & 6 & 292,890,231,892 & 292,890,231,898\\
\nolinkurl{value} & 6 & 292,890,231,899 & 292,890,231,905\\
\nolinkurl{integers} & 1 & 292,890,231,904 & 292,890,231,905\\
\bottomrule\end{longtable}


The five scalar-constant templates use the values in
\texttt{library.fixedConstants}. Large program literals are retained
byte-for-byte in the supplementary specification. Replacing them by
unspecified sufficiently large constants would change this definition.

\section{Primary theorem applications}
\begin{longtable}{p{.26\textwidth}p{.65\textwidth}}
\toprule Input & Exact application boundary\\\midrule\endhead
Selmer parity \cite[Theorem 1.4]{ddselmer} & $E/\Q$, $p=2$, full
$2^\infty$-Selmer corank; no finite-$\Sha$ hypothesis.\\
Kato \cite{kato} & Fixed weight-two form and symbol, integral transfer and
fixed-valuation smoothing; rational divisibility is Theorem 13.4(2),
with oddness, rational irreducibility and the unipotent condition.\\
BFH and Hoffstein--Luo \cite{bfh,hl} & Actual signed-prime twist with
trivial character and arbitrary finite splitting set. Only the
nonvanishing clause has the four-prime bound before the extra prime.\\
Waldspurger and Ramsey \cite{waldspurger,ramsey} & Full good-Hecke
eigenspace, same dyadic squareclass, and algebraic vectors by flat
field extension. No arbitrary-ring base change in weight $3/2$.\\
CM norms and heights \cite{ccl,cst} & Split level,
conductors coprime to level and the new prime, no extra units,
unaveraged sums and one fixed parametrization.\\
Conductor-one reduction Section~\ref{v15:cm-oriented} & Maximal CM order,
one inert prime, fixed orientations and same-place ideal action.
The geometric, CM and good-reduction inputs remain separate.\\
Duality \cite{milne,nekovar} & Actual derived Kummer conditions and their
homotopies. Support change only after the coefficient representation
is unramified at the removed place.\\
Forward finiteness \cite{gz,kolyvagin} & Only at known analytic
nonvanishing/simple zeros or a known nontorsion Heegner trace.\\
Isogeny ratio \cite[proof of I.7.3]{milne} & Finite $\Sha$, dual
sum/difference maps, canonical heights and pullback measures;
no BSD leading-term assumption.\\
Height \cite{calegari-icm2022,javanpeykar,gr2014,fowler,cmheight,pazuki} &
The same rationally descended $r_5$: absolute irreducibility, finite
ramification, de Rham weights and oddness; then the published prime-five
input, finite twist, explicit conductor and modular-Jacobian transfer.
CM order class number is bounded using the original Hom field.\\
Primes \cite{greentao,moebius,bmor} and Section~\ref{v15:gi-full} & Fixed finite-complexity system
after fixing sieve cutoff, and the unconditional explicit progression
estimate.\\\bottomrule
\end{longtable}
The released specification, natural precision polynomial and construction
files are identified in \nolinkurl{construction/SOURCE_MANIFEST.json}.
For the selected supporting proof arguments, see
\nolinkurl{research/version14/SUPPORTING-NOTES.md}. A source identity is not a proof verdict.

% End 06-literal-construction.tex


% Begin 07-even-source.tex
\section{Fixed-form arithmetic and the even coefficient source}\label{app:even-source}
This appendix supplies the constructions used in (3.1)--(3.3).
Throughout, $E=E_\ell$ is fixed, with the lattice, parametrization and
period convention of Section~\ref{sec:rc}. Constants may depend on these
data and on the prescribed finite list of local types. They do not depend
on the number or sizes of the later good conductor primes.

\subsection{One retained class, including discriminant \texorpdfstring{$4h$}{4h}}
All numbered primary references in this subsection are to Kato~\cite{kato}.

Fix a non-CM elliptic curve $E/\mathbf Q$ of conductor $N$, its weight-two newform $f$,
$T=T_2E$, and $V=T\otimes\mathbf Q_2$. In the application
$E=E_\ell$ is fixed first. Put
$\Lambda=\mathbf Z_2[[\Gamma]]=\mathbf Z_2[[t]]$, where $\Gamma$
is the real cyclotomic Galois group. Let
$\kappa:G_{\mathbf Q}\to\Gamma\subset\Lambda^\times$ be the quotient map,
and set $A=\Lambda[1/2]$.
We use arithmetic Frobenius throughout. The action on the cyclotomic
group-ring coefficient is the inverse deck action.

\paragraph{The precise output.}
\begin{proposition}[One retained Kato system and its two comparisons]
\label{ex:app-kato-system}
Use the fixed curve, lattice and cyclotomic convention just specified.
Fix $e\ge\max\{3,v_2(N)\}$ large enough for the fixed initial local data;
the smoothing choice is (K5).
Let $G$ be a finite elementary abelian two-group, let
$\psi:G_{\mathbf Q}\to G$ be a continuous homomorphism, and let
$\chi_x=\chi_0(x\circ\psi)$, for $x\in\widehat G$, be a finite
family of quadratic characters with one fixed real sign and fixed
local squareclasses at the primes dividing $2N$.
Since the characters of $G$ separate
its elements, $\psi(c)=1$ for complex conjugation $c$, and
$\psi|_{G_{\mathbf Q_v}}=1$ at every $v\mid2N$. All varying conductor primes
are good odd primes for $E$. Let $F$ be the squarefree least common
multiple of the odd parts of the primitive conductors. There is a
single retained-smoothed class
\[
 Z_F\in 2^{-c_E}H^1\bigl(G_{\mathbf Q,S},
        T\otimes\Lambda[G](\chi_0\psi\kappa^{-1})\bigr),
 \qquad S=\operatorname{prime}(2NF),                         \tag{K1}
\]
where $c_E$ is fixed before $F$ and $G$.
There is also one $F_E\in\Lambda$ with $F_E(0)\ne0$, fixed before
$F,G$, such that at every height-one prime $\mathfrak p$ of $A$,
\[
 \ell_{A_{\mathfrak p}}H^2(C_{S,x})_{\mathfrak p}
 \leq
 \ell_{A_{\mathfrak p}}
 \left(H^1(C_{S,x})/A F_E(Z_F)_x\right)_{\mathfrak p}.       \tag{K2}
\]
Here $C_{S,x}$ is the full rational global Iwasawa complex of the
constituent. For an odd squarefree signed parameter $h$ giving
$\chi_x=\chi_h$, central reciprocity for this same class gives
\[
 v_2\exp^*\bigl(F_E(Z_F)_x(0)\bigr)
 =v_2L_*(h)+\sum_{q\mid F,\ q\nmid h}
        v_2\#E^h(\mathbf F_q)+2e+v_2F_E(0)+\epsilon_h,       \tag{K3}
\]
when the value is nonzero. The numbers $\epsilon_h$ belong to a
bounded set depending only on the fixed curve, period conventions
and local types at two. A zero central $L$-value gives zero dual
exponential. Formula (K3) does not assert that its remaining unit is
one. No part of (K1)--(K3) compares integral lattices of varying newforms.
\end{proposition}
\begin{proof}\leavevmode
\paragraph{The retained route and its comparison.}
All classes in this subsection retain the smoothing integers of (K5).
The comparison used below is the exact identity (E.1), obtained by
integral weighted corestriction and the norm relation of
\cite[Proposition 8.12]{kato}. We do not divide by the smoothing
operators, or compare a chosen generator with a separately normalized
zeta class for the varying newform. In particular the central factor
is $s_F$ of (K6), whose valuation is $2e$ for every $F$.
The use of power-series zero finiteness in paragraph~7 establishes
nonvanishing of the supplied Euler system in both generic sign
components; it is not a comparison of two Euler systems.

\paragraph{1. One original lattice and one symbol combination.}
Kato (8.1.3), p.~180, constructs retained-smoothed classes
\[
 {}_{c,d_s}z_m(f,1,1,\xi,S_m)
 \in H^1\bigl(\mathbf Z[\zeta_m,1/2],L_f(1)\bigr),
 \qquad S_m=\operatorname{prime}(2mN),                      \tag{K4}
\]
where $L_f=V_{\mathcal O_\lambda}(f)$ is the image lattice of the
fixed form. For $\xi\in\mathrm{SL}_2(\mathbf Z)$, section~8.9,
p.~185, explicitly constructs the trace homomorphism into the
integral induced lattice
\[
 \mathbf Z_2[\operatorname{Gal}(\mathbf Q(\zeta_m)/\mathbf Q)]
       \otimes V_{2,\mathbf Z_2}(Y_1(N)).
\]
Section~8.11 projects this lattice onto $L_f$. Auxiliary modular
curve levels used in the trace do not change its target lattice.
Choose a rational identification $L_f(1)\otimes\mathbf Q_2=V$.
One fixed $2^{c_E}$ sends its lattice into $T$.

By Kato~13.6--13.7, p.~227, the symbol span has a nonzero projection
to each Betti sign line. The union of the two rational kernels does
not contain the whole span. Choose an integral finite linear
combination $\xi$ with both projections nonzero, clearing its fixed
denominators once. Norm relations and reciprocity are linear, so
the same combination works at every $m$. Although (K4) is written
with one symbol, it henceforth denotes this fixed combination.

\paragraph{2. Retain one smoothing choice for the whole family.}
Fix $e\geq\max\{3,v_2(N)\}$, enlarged only for fixed initial data,
and take
\[
                 c=d_s=1+2^e\,3N_{\rm odd}F.              \tag{K5}
\]
These integers exceed one, are prime to $6NF$, are one modulo $N$
and $F$, and are one modulo every dyadic quadratic conductor under
consideration. Distinctness or primality of $c,d_s$ is not a
hypothesis of Kato's construction. The exponents $(u,v)$ in (4.2.4),
p.~143, are $(1,1)$ when $k=2,r=r'=1$. The retained-smoothed
version of Theorem~6.6, p.~163, therefore has central scalar
\[
 s_F=(c^2-c)(d_s^2-d_s)=c(c-1)d_s(d_s-1),
                   \qquad v_2(s_F)=2e.                   \tag{K6}
\]
No smoothing factor is divided out. The prime factors of $cd_s$
exclude some Euler-system indices in Kato~13.1; they are not added
to the inverted arithmetic support $S$, as is explicit in (8.1.3)
and in the definition of $H^q$ in section~8.2. In particular their
number creates no local $H^2$ correction.

\paragraph{3. The integral $j_*$ condition survives ramified transfer.}
For a number field $K$, let $X_K=\operatorname{Spec}\mathcal O_K[1/2]$
and $j_K:\operatorname{Spec}K\to X_K$. At finite coefficients
$M=T/2^aT$, the low-degree Leray sequence gives an injection
\[
 H^1(X_K,j_{K,*}M)\longrightarrow H^1(K,M)
\]
whose image is the kernel of all inertia restrictions outside two.
Indeed the higher-direct-image stalk formula and Leray sequence
\cite[Tags 03Q9 and 03QA]{stackscohom} identify the stalks of
$R^1j_{K,*}M$ at a finite place with $H^1(I_v,M)$;
the edge map is restriction to these stalks. The generic-point map
$j_K$ is affine, hence quasi-compact and quasi-separated, as required
by the stalk formula. Kato's section~8.2,
pp.~180--181, defines the integral groups by the inverse limit of
these finite-coefficient groups, so this assertion applies to the
classes (K4) before rationalization.

Let $L/K$ be finite, with absolute Galois groups $H\subset G_K$.
For a class $z\in H^1(H,M)$ and an inertia subgroup $I\subset G_K$,
the Mackey formula is
\[
 \operatorname{res}_{I}^{G_K}\operatorname{cor}_{H}^{G_K}z
 =\sum_{g\in I\backslash G_K/H}
   \operatorname{cor}_{I\cap gHg^{-1}}^{I}
   \operatorname{res}_{I\cap gHg^{-1}}^{gHg^{-1}}(g z).    \tag{K7}
\]
The intersection in each summand is the upper inertia subgroup
for the corresponding place of $L$. If all upper inertia
restrictions vanish, every summand is zero. This does not require
$L/K$ to be unramified. It is a finite double-coset identity at each
finite coefficient precision. Uniqueness under the Leray injection
makes the extensions across $X_K$ compatible in $a$. Thus the
integral $j_*$ condition survives corestriction, including at primes
ramified in a quadratic twisting field.

\paragraph{4. Primitive twisting is an actual Euler system.}
Fix a constituent $\chi=\chi_h$ of primitive conductor
$D=|\Delta_h|=2^a d$, $d$ odd, and a basis $e_\chi$ of
$\mathbf Z_2(\chi)$. Put
$\Sigma=\operatorname{prime}(2Ndcd_s)$.
For every positive integer $m$ with $(m,cd_s)=1$, set
\[
 K_m=\mathbf Q(\zeta_m),\quad
 L_m=\mathbf Q(\zeta_{Dm}),\quad
 y_m=\operatorname{cor}_{L_m/K_m}
             (z_{Dm}\otimes e_\chi).                      \tag{K8}
\]
The character is trivial over $L_m$, so tensoring is an integral
map there. Formula (K7) places (K8) in the precise $j_*$ group
required by Kato~13.1 on its index set
$\operatorname{prime}(m)\cap\Sigma=\{2\}$. The larger definition is
permitted by (8.1.3), Sections~8.9--8.11 and Proposition~8.12:
$S_{Dm}=\operatorname{prime}(2NDm)$ is disjoint from
$\operatorname{prime}(cd_s)$, since $d\mid F$ and (K5) holds.
Each class uses the original $f,\xi,c,d_s$. We will also use this
same construction on the different index set required by
\cite[Section 0.1]{kato1999}, as checked below.

For $m\mid m'$, transitivity of corestriction reduces the norm
of (K8) to Proposition~8.12, p.~186. For a deck automorphism $g$,
equivariance gives
\[
 \operatorname{cor}((gz)\otimes e_\chi)
       =\chi(g)^{-1}g\operatorname{cor}(z\otimes e_\chi).
\]
Consequently the factor at a newly introduced good prime $q$ is
\[
 E_q(\chi)=1-\frac{a_q\chi(q)}q\sigma_q^{-1}
                         +\frac1q\sigma_q^{-2}.           \tag{K9}
\]
In Kato~13.1, pp.~224--225, the polynomial is
$P_q(X)=\det(1-\operatorname{Fr}_qX\mid T\chi)$ and the norm
factor is $P_q(q^{-1}\sigma_q^{-1})$, exactly (K9).
Repeated prime powers introduce no factor because Proposition~8.12
has a product over newly added support primes. Thus the subfamily
on Kato~13.1's index set is an Euler system of the stated lattice,
not a class inferred from proportionality of rank-one Iwasawa
modules. The required norm relations on the 1999 index set are
verified separately below.

For $m=2^n$ with $n\geq a$, the field $L_m$ can be replaced by
$\mathbf Q(\zeta_{d2^n})$: apply repeated-two norm compatibility
from exponent $n+a$ to $n$. This introduces no factor $2^a$.
The construction therefore covers discriminant $4h$ on the same
cofinal tail. At the coefficient level it retains the full character
$\chi_2\chi_{\rm odd}$. A dyadic character nontrivial on $\Gamma$
changes its action by the integral automorphism
$1+t\mapsto-(1+t)$, or $t\mapsto-2-t$. A dyadic sign character
instead changes the finite real-sign component. Neither operation
is a division by a sign projector.

\paragraph{5. One coefficient map produces every constituent.}
For $n$ large enough to contain the dyadic characters, put
$A_{F,n}=\operatorname{Gal}(\mathbf Q(\zeta_{F2^n})/\mathbf Q)$.
Apply Shapiro to $z_{F2^n}$, then the map on group bases
\[
 \Phi_n:\mathbf Z_2[A_{F,n}]\longrightarrow
                  \mathbf Z_2[\Gamma_n\times G],\qquad
 [s]\longmapsto\chi_0(s)[\bar s][\psi(s)].                 \tag{K10}
\]
With inverse deck actions, both sides transform under $g\in G_{\mathbf Q}$
by the inverse of its group image. On the target the extra scalar
is $\chi_0(g)$; since $\chi_0^{-1}=\chi_0$, (K10) is equivariant.
This is an integral coefficient homomorphism. Shapiro restriction
to an inertia subgroup is the direct sum of the upper restrictions;
the same zero restrictions used in (K7) therefore remain zero after
(K10). In particular the map acts on actual integral classes.
The map in (K10) is defined on this finite group for the following
reason. The constituent indexed by the trivial character of $G$ is
$\chi_0$, and
$x\circ\psi=\chi_x\chi_0$ for every $x\in\widehat G$.
All these quadratic Dirichlet characters have conductor dividing
$F2^n$ once $n$ is large enough. If an element of $G_{\mathbf Q}$
acts trivially on $\mathbf Q(\zeta_{F2^n})$, every character of $G$
is therefore one on its image under $\psi$. Character separation
makes that image the identity. Thus both $\chi_0$ and $\psi$
factor through $A_{F,n}$. Explicitly, for $g\in G_{\mathbf Q}$,
\[
 \Phi_n([g^{-1}s])
 =\chi_0(g)\,[\bar g]^{-1}[\psi(g)]^{-1}\Phi_n([s]).
\]
This verifies equivariance on each integral group basis before
evaluation at any character. The quotient to the real tower sends
each basis element to its image; its matrix has entries zero and
one. The transition maps and character evaluations are also maps
of integral coefficient modules. None of these operations is the
idempotent $|G|^{-1}\sum_g x(g)[g]$.


The maps $\Phi_n$ commute with the norm transition maps.
Proposition~8.12 has no new prime when $n$ increases, so their
images define the single class $Z_F$ in (K1).
The finite-coefficient degree-zero groups are finite and satisfy
Mittag--Leffler. The inverse-limit exact sequence
\cite[Corollary 4.1.3, p.~96]{nekovar} therefore has zero
$\varprojlim^1 H^0$ term, so it gives this degree-one class
in completed cochains.
The quotient from the full cyclotomic tower to $\Gamma$ is the
integral quotient of group bases, equivalently corestriction to
the real subfield. No factor $(1+\tau)/2$ occurs.

After evaluation at $x$, (K10) is the coefficient form of the
unnormalized corestriction weighted by $\chi_x$. Factor that
corestriction through the primitive conductor field in (K8).
Transitivity and Proposition~8.12 give the exact identity
\[
              (Z_F)_x=
       \prod_{q\mid F/d}E_q(\chi_x)\,y_{2^\infty}.        \tag{E.1}
\]
It is an identity of the given classes. The same smoothing choice
appears on both sides. There is no group-order denominator and no
comparison with a canonical zeta class of $f\otimes\chi_x$.

For clarity, at a constituent write $\chi_h=\chi_2\chi_{\rm odd}$ and
$|\Delta_h|=2^a d$. Before taking the cofinal limit, the corresponding
scalar coefficient map is exactly
\[
 [s]\longmapsto\chi_2(s)\chi_{\rm odd}(s)[s|_{\Gamma_n}],
                    \qquad n\ge\max\{a,2\}.\tag{E.2}
\]
Thus the dyadic character is part of the module action, including
when the field discriminant is $4h$.

\paragraph{6. Central trace and period normalization.}
First use (E.1) to remove the unused tame primes. The central
specialization of (K8) is the corestriction of $z_D\otimes e_\chi$
from $\mathbf Q(\zeta_D)$ to $\mathbf Q$, with omitted support
$\operatorname{prime}(2Nd)$. For $D=1$ use the trivial character
and the same omitted support. This follows from the same repeated-two
norm identity even when $D=4|h|$. One must do this before using a
Gauss basis; the Gauss sum of an imprimitive character modulo a
larger cyclotomic level can be zero.

Theorem~9.7, p.~189, identifies the dual exponential of $z_D$ with
the retained-smoothed de Rham zeta element. Dual exponential commutes
with corestriction, whose de Rham map is field trace. Thus the
coefficient of $e_\chi$ is the exact unnormalized sum
\[
                    \sum_{g\in\operatorname{Gal}(\mathbf Q(\zeta_D)/\mathbf Q)}
                       \chi(g)\,g(\exp^*z_D).             \tag{K12}
\]
Theorem~6.6, p.~163, evaluates precisely this sum. In weight two,
$r=r'=1$, its Betti sign is $\epsilon=\chi(-1)$ and its power of
$2\pi i$ is zero. If $e_B^\epsilon$ is a fixed rational generator
of that Betti line and
$\delta(f,1,\xi)^\epsilon=b_\epsilon e_B^\epsilon$, then
$b_\epsilon\in\mathbf Q^\times$ is fixed and
\[
     \operatorname{per}_f\bigl((\text{K12})\bigr)^\epsilon
          =s_F L_{\operatorname{prime}(2Nd)}(f,\chi,1)\,
                        b_\epsilon e_B^\epsilon.          \tag{K13}
\]
Choose the original de Rham generator so that its Betti period
on this line is $\Omega_E^\epsilon e_B^\epsilon$, absorbing a
fixed rational comparison in $b_\epsilon$. Formula (K13) says
that the scalar multiplying $e_\chi$ in (K12) is
$s_F b_\epsilon L_{\operatorname{prime}(2Nd)}(f,\chi,1)/
\Omega_E^\epsilon$. Its Galois character is $\chi$.

To descend it, take the invariant basis $\sqrt h\,e_\chi$ of the
quadratic character line. With the usual complex convention
$\Omega_E^-=i|\Omega_E^-|$ and $\sqrt h=i\sqrt{|h|}$ for $h<0$,
division by $\sqrt h$ gives, up to the already fixed rational
period choice, exactly
\[
             s_F b_\epsilon\,\frac{L_*(h)}{h}
       \prod_{p\mid 2N}\text{(the omitted local factor at $p$)}.
                                                               \tag{K14}
\]
The signed integer $h$ is odd, so $h^{-1}$ has valuation zero.
This computation also explains the usual Gauss-basis description:
for the primitive character, $\tau(\chi)^2=\Delta_h$ and
$\tau(\chi)e_\chi$ is invariant; replacing it by $\sqrt h e_\chi$
costs the fixed dyadic factor $\sqrt{\Delta_h/h}$. It does not cost
$\varphi(D)$ or any power of the number of prime divisors of $h$.
The discriminant-$4h$ case has ratio two in this last basis change,
and the direct basis calculation (K14) already includes it.
One can verify the normalization without choosing an idempotent.
Write $\beta=\exp^*z_D$ in the original de Rham line and put
$\beta_\chi=\sum_g\chi(g)g\beta$. Then
$u\beta_\chi=\chi(u)\beta_\chi$ for every $u$ in the finite
cyclotomic group. Thus $\beta_\chi\otimes e_\chi$ is invariant,
and its coordinate in the invariant basis
$\sqrt h\,e_\chi$ is $\beta_\chi/\sqrt h$.
Theorem~6.6 uses precisely $\beta_\chi$, rather than
$\beta_\chi/\varphi(D)$. Theorem~9.7 applies here since
$1\le r=r'=1\le k-1=1$; it has no restriction that $p$ be odd.
The dual form $f^*$ in Theorem~6.6 equals $f$, because the elliptic
newform has rational coefficients.

In the real-period convention,
$L_*(h)=\sqrt{|h|}L(E^h,1)/|\Omega_E^\epsilon|$, up to the
fixed rational period comparison already included in $b_\epsilon$.
For $h>0$,
$\Omega_E^+\sqrt h=|\Omega_E^+|\sqrt{|h|}$; for $h<0$,
$\Omega_E^-\sqrt h=-|\Omega_E^-|\sqrt{|h|}$.
Consequently
\[
 \frac{L(E^h,1)}{\Omega_E^\epsilon\sqrt h}
     =\frac{L_*(h)}{h}
\]
with the same fixed period comparison. This proves the signed
factor in (K14), not only its valuation. For a primitive quadratic
character, $\tau(\chi)^2=\Delta_h$ and
$\tau(\chi)/\sqrt h$ equals a fixed choice of
$\sqrt{\Delta_h/h}$. Changing from its Gauss basis to the displayed
invariant basis multiplies the coordinate by this ratio, which
is one or two up to sign. In particular the discriminant-$4h$
case adds no factor depending on the number of odd conductor primes.


Identifying this generator with a minimal invariant differential (or its dual, according to the fixed convention)
for $E^h$ at two has bounded valuation: odd squarefree $h$ belongs to
one of finitely many classes in $\mathbf Q_2^\times/\mathbf Q_2^{\times2}$,
and changing a representative by a unit square scales the quadratic
twist differential by a two-adic unit. The local factors at $2N$
also form a fixed finite nonzero list. At a good active prime the
local $L$-factor is one; at an unused good prime (K9) at $t=0$ is
$\#E^h(\mathbf F_q)/q$. This proves (K3), including all its
uniformity claims and the assertion about zero values.

\paragraph{7. Every hypothesis of Kato~13.4(2).}
Apply the theorem to the integral system (K8), after multiplying
by the fixed lattice comparison if needed. Its hypotheses are:
\begin{itemize}
\item Each generic sign component is nonzero. For a finite-order
cyclotomic character $\rho$, corestriction and the $\rho$-sum combine
into the single $(\chi\rho)$-weighted sum on the original class.
Theorems~6.6 and~9.7 apply with the same $\xi,c,d_s$. Their smoothing
factors are nonzero, since $c,d_s>1$ cannot be roots of unity.
Theorem~13.5(2), p.~227, gives only finitely many zero $L$-values
among characters supported on the fixed finite set
$\operatorname{prime}(2Nd)$. There are infinitely many $\rho$ of
each parity. A torsion Iwasawa class can have a nonzero specialization
at only finitely many such characters, since a nonzero annihilating
power series has only finitely many cyclotomic zeros. Hence both
generic sign components are nonzero.
\item The two real eigenspaces of $T\chi$ have rank one.
\item The theorem uses inverse Frobenius in its purity condition.
On $T\chi$ its good eigenvalues have complex absolute value
$q^{-1/2}$, so its integer $w$ is $-1$.
\item $V\chi$ is irreducible; twisting preserves irreducibility of
the non-CM rational Tate representation.
\item The image over $\mathbf Q(\mu_{2^\infty})$ contains an open
subgroup of $\mathrm{SL}_2(\mathbf Z_2)$. Intersecting with
$\ker\chi$ keeps an open subgroup, which contains a nonidentity
unipotent. It fixes a one-dimensional subspace and gives condition~(v).
\end{itemize}
These are all rational hypotheses. The residual irreducibility,
integral cokernel and $p\ne2$ conditions of part~(3) are not used.
Theorem~13.4(1) gives torsion $\mathbf H^2(T\chi)$, and part~(2)
gives, at every height-one prime away from two,
\[
 \ell\mathbf H^2(T\chi)_{0,\mathfrak p}
       \leq\ell(A_{\mathfrak p}/J_{\mathfrak p}),\qquad
 \mathbf H^2_0=\ker(\mathbf H^2\to\mathbf H^2_{\rm loc}). \tag{K15}
\]

The finiteness and Euler characteristic (12.2.1)--(12.2.2), p.~220,
give generic rank one in $\mathbf H^1$. Its localization away
from two is torsion-free. Indeed a torsion element for a scalar
prime would come from invariants of the corresponding finite-dimensional
rational coefficient algebra at that prime; its coefficient action is
abelian, whereas the open commutator subgroup on $V\chi$ has no
fixed vectors. The coefficient exact sequence excludes these
invariants. The scalar specialization need not be a finite-order
character. Thus $\mathbf H^1_{\mathfrak p}$ is free of rank one.
Localization commutes with $\operatorname{Hom}$ for this finitely
presented module. Writing $y_{2^\infty}=a e$ in a local basis gives
$J_{\mathfrak p}=(a)$. Formula (K15) is exactly the index bound for
the supplied class, not for an arbitrarily rescaled generator.

\paragraph{The published prime-two Euler-system input.}
For $p=2$ we use the rational, away-from-two case of
\cite[Theorem 0.8, pp.~317--318]{kato1999}, with the stronger
hypothesis in its condition 0.6(ii), to justify the application of
\cite[Theorem 13.4(2)]{kato}. For our non-CM elliptic representation
the image of $G_{\Q^{\rm ab}}$ contains an open subgroup of
$\mathrm{SL}_2(\Z_2)$, as stated in
\cite[Example 0.6.2, pp.~317--318]{kato1999}, using Serre's open-image
theorem. A quadratic character $\chi$ is trivial on $G_{\Q^{\rm ab}}$.
Thus the same subgroup acts on $V\chi$ and contains a matrix
$\left(\begin{smallmatrix}1&b\\0&1\end{smallmatrix}\right)$ with
$b\ne0$. Its fixed space over $\Q_2$ has dimension one. This supplies
the unipotent in $G_{\Q^{\rm ab}}$, not merely in the cyclotomic
subgroup. The latter subgroup also contains this open
$\mathrm{SL}_2$ subgroup, so $V\chi$ remains absolutely irreducible
over $\Q(\mu_{2^\infty})$. Indeed, a line invariant under both upper
and lower nonidentity unipotents cannot exist, even over an algebraic
closure. No inference from irreducibility over $G_\Q$ alone is needed.

We check the index set and the supplied class in that theorem
explicitly. Put $N_K=2Nd$ and $N'_K=cd_s$. These integers are
coprime by (K5), since $d\mid F$. Section~0.1 of \cite{kato1999}
requires every index $m$ with $N_K\mid m$ and $(m,N'_K)=1$.
The enlarged definition (K8) supplies all these classes. They
originate in the integral $j_*$ group over
$\mathbf Z[\zeta_m,1/2]$ and map to full cohomology over
$\mathbf Z[\zeta_m,1/N_K]$. The varying construction support
$S_{Dm}$ is not an instruction to invert all its primes in this
last scheme. Apply the fixed lattice map $2^{c_E}$ from (K4) if
needed, identically at every level.

For a prime $\ell\nmid N'_K$, if $\ell\mid m$ then
$S_{Dm\ell}=S_{Dm}$, so Proposition~8.12 gives norm factor one.
If $\ell\nmid m$, the condition $N_K\mid m$ implies
$\ell\nmid2Nd$, and hence
$S_{Dm\ell}\setminus S_{Dm}=\{\ell\}$. The equivariance identity
before (K9) gives the factor
$P_\ell(\ell^{-1}\sigma_\ell^{-1})=E_\ell(\chi)$, with
$P_\ell(X)=1-a_\ell\chi(\ell)X+\ell X^2$. These are exactly the
two cases of (0.1.1) in \cite{kato1999}. Smoothing primes are
excluded by $N'_K$ and need not be inverted.

For the class in Section~0.7 of that source take its parameter
$d=1$; this parameter is distinct from our odd conductor part $d$.
For every $n\geq1$,
\[
 S_{DN_K2^n}=S_{D2^n}=\operatorname{prime}(2Nd).
\]
Proposition~8.12 therefore gives
\[
 \operatorname{cor}_{\mathbf Q(\zeta_{DN_K2^n})/
                          \mathbf Q(\zeta_{D2^n})}
             z_{DN_K2^n}=z_{D2^n}.
\]
Transitivity of corestriction, through the two intermediate fields
in the common upper field, now yields
\[
 \operatorname{cor}_{\mathbf Q(\zeta_{N_K2^n})/
                          \mathbf Q(\zeta_{2^n})}y_{N_K2^n}
 =\operatorname{cor}_{\mathbf Q(\zeta_{D2^n})/
                          \mathbf Q(\zeta_{2^n})}
                         (z_{D2^n}\otimes e_\chi)
 =y_{2^n}.
\]
No Euler factor occurs because the supports are equal. No field
degree occurs because this is corestriction of the actual upper
class, not corestriction of a restriction. These equalities commute
with transition maps, and $n\geq1$ is a cofinal tail. Write
$\xi_K$ for the Section~0.7 class, to distinguish it from the
modular-symbol combination $\xi$. It is exactly the full-cohomology
image of $y_{2^\infty}$, multiplied by $2^{c_E}$ if the fixed
lattice map was needed. That factor is a unit at every height-one
prime under consideration, so its coordinate ideal is precisely
$J_{\mathfrak p}$. The preceding
nonvanishing argument supplies a nontorsion class on each sign
component, purity gives $w=-1$, and oddness gives condition
0.6(iv$_{\mathfrak p}$). At the height-one primes used here, $2$ is
invertible and the localized Iwasawa algebra is a DVR. Theorem 0.8
then bounds the coordinate ideal of this actual retained class. Its
strong residual and integral-cokernel assumptions concern primes
containing $2$, which are not used here.

We also account for the difference between the two published
cohomology conventions. Write $C_j$ for the $j_*$ Iwasawa complex of
\cite[Sections 8.2 and 12.2]{kato}, and $C_S$ for \emph{étale}
cochains on the open arithmetic scheme obtained by inverting $2$
and every odd prime in the actual ramification support $S$,
including the fixed twist, as in \cite[Sections 0.1 and 0.5]{kato1999}.
The support modulus called $N$ in that source is taken to contain
these primes; it is not merely the conductor of the untwisted curve.
Work in one rational sign component and localize at $\mathfrak p$.
Their quotient is the direct sum of the odd local complexes (K16).
Each differential has nonzero determinant, so its degree-one
cohomology vanishes. Hence $H^1(C_j)\simeq H^1(C_S)$ and
$H^2(C_j)\hookrightarrow H^2(C_S)$. At an odd prime $q$, the
unramified local term has cohomological dimension one after inverting
$2$. The localization triangle identifies the degree-two cohomology
of (K16) with full local $H^2(\Q_q,T\chi)$ in Iwasawa cohomology.
The map from $H^2(C_S)$ is the usual local restriction. Exactness
therefore identifies $H^2(C_j)$ with the kernel of all those odd
local restrictions. Intersecting with the kernel of restriction at
$2$ gives
\[
 \ker\bigl(H^2(C_j)\to H^2_{\rm loc,2}\bigr)
 \simeq
 \ker\bigl(H^2(C_S)\to\!\bigoplus_{q\in S}H^2_{\rm loc,q}\bigr).
\]
This is precisely the comparison of the subscript-zero groups in
the two sources. The $H^1$ identification preserves the actual
Euler-system class and its coordinate ideal $J$. Thus Theorem 0.8
gives (K15) in the convention used here. The argument for (K16)
below is local and does not use (K15). We do not invoke the
odd-prime-only Perrin--Riou argument or an unchecked version of
Rubin's theorem at $2$.

\paragraph{8. Paying for the full global complex.}
At an odd prime put $W_q=H^1(I_q,V\chi)$. The rational cone from
$j_*$ to full cochains is
\[
 [\,W_q\otimes A\xrightarrow{1-\sigma_q^{-1}\operatorname{Fr}_{W_q}}
          W_q\otimes A\,],\qquad\deg=(1,2).              \tag{K16}
\]
Its determinant $D_q$ is nonzero, since the cyclotomic Frobenius
exponent is nonzero. The map is injective even after localization
at a prime dividing $D_q$; it is not being reduced modulo that
prime. Thus the cone contributes only degree-two length
$v_{\mathfrak p}(D_q)$. The long exact sequence shows that full
$H^2$ increases by at most their sum and that $H^1$ is unchanged.
This inequality does not require vanishing of the next $j_*$ group.

At a good inactive prime, $W_q=V\chi(-1)$, so $D_q=E_q(\chi)$
with exactly the convention of (K9). At a good active prime,
ramified quadratic inertia acts by $-1$; its rational cohomology
vanishes and $D_q=1$. This includes every active base prime,
regardless of how many there are. At an old bad prime retain
$H^1(I_q,V\chi)$ itself; replacing it by $V^{I_q}\chi(-1)$
would in general be false.

Choose a fixed integral $F_2$ divisible over $A$ by the
characteristic elements of $\mathbf H^2_{\rm loc}(T\chi)$
for the finitely many allowed local types at two and finite
sign components. Its central value can be chosen nonzero:
derived local control, and the absence of local Iwasawa $H^3$,
identify its rational central $H^2$ with $H^2(\mathbf Q_2,V\chi)=0$.
The last equality follows from local duality and finiteness of
local elliptic torsion. Hence the characteristic elements have
no factor $t$. Kato (12.2.3), p.~220, gives the required finiteness
and torsion of this local Iwasawa module.

At the finitely many old bad primes, choose stable lattices in
the $W_q$ and multiply all the determinants $D_q$ for the finitely
many local twisting types. Clear their fixed scalar denominators.
Their central values are nonzero, since the cokernel of
$1-\operatorname{Fr}_{W_q}$ is $H^2(\mathbf Q_q,V\chi)=0$
by the same local duality. Call this product $F_{\rm bad}$ and set
$F_E=F_2F_{\rm bad}$.

The quotient $\mathbf H^2/\mathbf H^2_0$ injects into the local
$\mathbf H^2_{\rm loc}$, so $F_2$ pays for it in (K15).
Multiplying a class in a free rank-one module by a nonzero scalar
adds that scalar's valuation to its index. Formula (E.1) already
pays for every unused good determinant in (K16); $F_{\rm bad}$
pays for the remaining old ones. Adding these length inequalities
proves (K2).
Here is the length calculation with all contributions separated.
Fix a height-one prime $\mathfrak p$ of $A$ and write
$R=A_{\mathfrak p}$. Let $M$ be the localized degree-one module
for the primitive $j_*$ system, let $y=y_{2^\infty}$, and write
$\operatorname{ind}_R(y)=\ell_R(M/Ry)$.
The preceding freeness argument gives $M\simeq R$.
Put $Q_\chi=\prod_{q\mid F/d}E_q(\chi)$.
All primes in $F/d$ are good: the family has identical local
characters at primes dividing $N$, so a bad prime cannot ramify
in one constituent and be unramified in another.
The cone (K16) has zero degree-one cohomology, hence identifies
$M$ with the localized full degree-one module. Its long exact
sequence and (K15) give
\[
\begin{split}
 \ell_R H^2(C_{S,x})_{\mathfrak p}
 &\le \ell_R\mathbf H^2(T\chi)_{0,\mathfrak p}
       +\ell_R\mathbf H^2_{\rm loc}(T\chi)_{\mathfrak p}
       +\sum_{q\in S,\ q\ne2}v_{\mathfrak p}(D_q)\\
 &\le \operatorname{ind}_R(y)
       +v_{\mathfrak p}(F_2)+v_{\mathfrak p}(F_{\rm bad})
       +v_{\mathfrak p}(Q_\chi)\\
 &=\operatorname{ind}_R(F_EQ_\chi y)
  =\operatorname{ind}_R\bigl(F_E(Z_F)_x\bigr).
\end{split}
\]
The second line uses $D_q=1$ at active good primes, and
$D_q=E_q(\chi)$ at inactive good primes. It does not charge an
extra fixed constant for each active prime.

For completeness, nonvanishing at the center of the fixed
correction factors follows before multiplying over local types.
Local derived control and $H^3_{\rm Iw}=0$ identify
$\mathbf H^2_{\rm loc}/t\mathbf H^2_{\rm loc}$ rationally with
$H^2(\mathbf Q_2,V\chi)$. Local duality identifies the latter
with the dual of $H^0(\mathbf Q_2,V\chi)$, which is zero because
local elliptic torsion is finite. Nakayama at $(t)$ therefore
excludes a factor $t$ from its characteristic element.
At an old odd prime the cokernel of
$1-\operatorname{Fr}_{W_q}$ is $H^2(\mathbf Q_q,V\chi)=0$,
by the same duality. Its determinant at $t=0$ is nonzero.
Taking products over the finite list of local types, followed by
one fixed power of two to clear scalar denominators, gives
$F_E\in\Lambda$ with $F_E(0)\ne0$. This operation is independent
of the binary dimension and all later good conductor primes.


\paragraph{9. Uniformity and the boundary of this proposition.}
The only lattice denominator is fixed before the family. The
smoothing integers may grow, but their central valuation is $2e$.
The multiplier $F_E$ depends only on the original local types.
Transfer, coefficient push, real quotient and norm compatibility
are integral. The proof never divides by a conductor degree,
$|G|$, a class number, or a Fourier projector. It keeps one class
through (K1), (K2), and (K3).
\end{proof}

We now use this same class in the positive complex. The preceding
argument supplies its integral construction, its rational height-one
divisibility and its exact central reciprocity normalization.

\begin{proposition}[One integral even scalar for the family]\label{ex:app-even-scalar}
Retain Proposition~\ref{ex:app-kato-system} and assume, as for the
present curve, that $E[2](\Q)\simeq(\Z/2\Z)^2$. Let $C^+$ be the
fiber of the ordinary global cochains to the ordinary real cochains,
with real degree zero retained. It has a compatible finite free model
under the coefficient changes used here. For
$B=\mathscr O[G]$, it satisfies $C^+\otimes B\simeq B^2[-1]$.
The fixed real invariant generator, the positive lift of the retained
class, and the original determinant basis define the single scalar
$u^{\rm ev}$ by (E.4), with
\[
 u^{\rm ev}\in\Lambda[1/2][G],\qquad
 2^{c_E}u^{\rm ev}\in\Lambda[G].
\]
Here $c_E$ may be enlarged for the fixed data but is independent of
$G$ and the later good conductor primes. The assertion requires the
common real sign and local squareclasses at $2N$ in
Proposition~\ref{ex:app-kato-system}.
\end{proposition}
\begin{proof}\leavevmode
Let $C^+$ be the fiber of ordinary global cochains to ordinary real
cochains. Real $H^0$ is included. Over a finite Artin quotient,
continuous cochains are filtered unions of finite free function
modules, hence flat and therefore free. A residual splitting into
finite cohomology and disks lifts by the nilpotent perturbation
series to a finite free model with actual homotopies. Minimal models
at successive quotients are related by invertible matrices:
a homotopy equivalence reduces to an isomorphism of minimal residual
complexes, and Nakayama applies. Lifting these matrices and comparison
homotopies makes the models compatible. Their acyclic, termwise
surjective comparison cones have acyclic inverse limit.
Here is the compatibility argument. At successive finite coefficient
quotients, reduction of a minimal finite free model is a minimal model
of the reduced cochains. Two such models are homotopy equivalent;
reducing a homotopy inverse modulo the maximal ideal shows that the
comparison is degreewise invertible. Lift these invertible matrices
to the next quotient and change bases. The model differentials then
reduce to the preceding differentials exactly. The comparison maps
to the original cochains agree after reduction up to chain homotopy.
Since reduction of continuous cochains is termwise surjective, lift
that homotopy and adjust the upper comparison by its boundary. The
comparison maps are now exactly compatible and remain
quasi-isomorphisms.

Let $\mathcal K_n$ be their acyclic comparison cones. Their transition
maps are termwise surjective. If $s$ denotes the transition shift on
the product, successive lifting gives the exact sequence of complexes
\[
 0\longrightarrow\varprojlim_n\mathcal K_n
 \longrightarrow\prod_n\mathcal K_n
 \xrightarrow{\,1-s\,}\prod_n\mathcal K_n
 \longrightarrow0.
\]
Products of complexes of modules preserve acyclicity, so the long
exact cohomology sequence makes $\varprojlim_n\mathcal K_n$ acyclic.
The inverse-limit comparison is therefore a quasi-isomorphism from
the bounded finite free limit model to the continuous coefficient
complex. This also identifies its derived coefficient changes with
the finite-quotient constructions.
This proves the needed perfection and coefficient change.

% Replace baseline lines 3310--3335; retains equation tag (E.3).
\paragraph{The residual positive complex.}
Write $K_n=\mathbf Q(\zeta_{2^{n+2}})^+$, and let $S_n$ consist of
the places above two and a fixed finite set of odd rational primes.
Put $r_n=2^n$ and $s_n=\#S_n$. We first compute with the scalar
coefficient $\mathbf F_2=\mu_2$.

The ordinary two-class group of $K_n$ is trivial. Indeed this holds
for $K_0=\mathbf Q$. Assuming it for $K_{n-1}$, let $H$ be the
Hilbert two-class field of $K_n$, with the real places split.
The canonical extension $H/K_{n-1}$ is Galois with a finite
two-group $G_H$. Only the unique dyadic prime can ramify, and its
inertia $I$ has order two. The quotient of $G_H^{\rm ab}$ by the image
of $I$ gives an everywhere unramified abelian two-extension of
$K_{n-1}$ with real places split, hence is trivial by induction.
Thus $G_H^{\rm ab}$ is cyclic, generated by an element of order at most
two. The Frattini quotient of $G_H$ has one generator, so $G_H$ is cyclic and
has order at most two. Since it contains $K_n$, we obtain $H=K_n$.
The $S_n$-class group is a quotient, so its two-part also vanishes.

The signatures of units span all $r_n$ real signs. For $n\ge1$ put
$m=n+2$, $\zeta=\zeta_{2^m}$, and
\[
 u=1+\zeta+\zeta^{-1}
   =\zeta^{-1}\frac{\zeta^3-1}{\zeta-1}.
\]
This is a unit: the numerator and denominator generate the same
dyadic prime ideal in the cyclotomic field. Its real conjugates are
$1+2\cos(2\pi a/2^m)$ for odd $0<a<2^{m-1}$.
The negative ones have $2^m/3<a<2^{m-1}$, and their number is
\[
 2^{m-2}-\left\lfloor\frac{2^{m-1}}3+\frac12\right\rfloor,
\]
an odd integer. Hence $N_{K_n/\mathbf Q}(u)=-1$.
The signature vector of $u$ has augmentation one in the local ring
$\mathbf F_2[\operatorname{Gal}(K_n/\mathbf Q)]$ and is therefore
a unit. Its Galois translates span all signatures. At $n=0$ use $-1$.

Kummer theory now gives
$\dim H^1(G_{K_n,S_n},\mathbf F_2)=r_n+s_n$, and its signature map
to $\mathbf F_2^{r_n}$ is onto. The map in degree zero is the diagonal
injection. The positive complex is the fiber of ordinary global
cochains to the sum of ordinary real cochains, so its long exact
sequence gives $\dim H^1(C_n^+)=r_n+s_n-1$.
The vanishing two-part of the $S_n$-class group identifies global
$H^2$ with $\operatorname{Br}(\mathcal O_{K_n,S_n})[2]$.
The local invariant description has $r_n+s_n$ coordinates with one
sum relation. Projection to the $r_n$ real coordinates is onto,
because $s_n\ge1$, and has kernel of dimension $s_n-1$.
Together with signature surjectivity this gives
\[
 \dim H^1(C_n^+)=2^n+s_n-1,\qquad
 \dim H^2(C_n^+)=s_n-1,
 \qquad H^i(C_n^+)=0\quad(i\ne1,2).
\]
The last assertion uses the usual global-to-real comparison in
degrees greater than two; degree zero vanishes by the diagonal injection.

For each fixed support, $s_n$ is bounded. The dyadic prime is unique;
each fixed odd prime has open decomposition group in the real
cyclotomic $\mathbf Z_2$-extension, since its image in
$\mathbf Z_2^\times/\{\pm1\}$ has nonzero two-adic logarithm.
Use now the finite free model and coefficient change already constructed.
Its scalar reduction is a bounded finite free $\mathbf F_2[[t]]$-complex,
and its finite cyclotomic specializations are derived quotients by
$t^{2^n}$. Smith decomposition shows that torsion disks and their Tor
terms contribute bounded dimensions, while a free cohomology generator
contributes $2^n$. The displayed dimensions therefore give generic
rank one in degree one and zero in all other degrees.

For the present curve, $E[2]$ is the direct sum of two trivial
$\mathbf F_2$-representations. Thus the residual positive complex has
rank two in degree one only. The residue field of the local ring
$B=\mathscr O[G]$ is $\mathbf F_2((t))$; the augmentation ideal of
the finite two-group lies in its maximal ideal. A minimal finite free
$B$-model has zero residual differential, so the just-computed dimensions
force exactly two generators in degree one and no others. Consequently
\[
                    C^+\otimes B\simeq B^2[-1].\tag{E.3}
\]
This uses the local ring $B$, not semisimplicity of $\mathbf F_2[G]$.

The common real sign is needed here. The real cyclotomic character
and $\varepsilon$ are trivial on complex conjugation, so the real
action is the fixed action on $T_2E\otimes\chi_0$, extended to the
group ring. Its invariant module is the fixed free rank-one
$\Z_2$-lattice extended to $\Lambda[G]$, and then to
$\mathscr O[G]$. Thus a generator can be chosen before $G$ varies.
Real $H^1$ is killed by two over these coefficient rings, so the
positive lift below costs the same single factor two for every
binary dimension.
Choose a generator of the fixed real invariant lattice and its
connecting class $a$. Since real $H^1$ is killed by two,
$2^{c_E+1}F_EZ_F$ has an integral positive lift $b$.
For a basis $\theta$ of the original inverse determinant line, define
\[
                         u^{\rm ev}\theta=2^{-c_E}a\wedge b.
 \tag{E.4}
\]
The wedge does not depend on the choice of the positive lift.
Indeed, the difference of two lifts is in the image of real
degree-zero cohomology by the cone exact sequence. That module
is generated by the chosen real invariant generator, so the
difference is a multiple of $a$ and wedges to zero with $a$.
After rationalization set $z=F_EZ_F$ and choose any lift
$\widetilde z$ of $z$. Formula (E.4) then reads
\[
              u^{\rm ev}\theta=2a\wedge\widetilde z.
\]
Thus the factor $2^{c_E}$ used to make the lift integral cancels
exactly. Increasing the fixed clearing exponent does not change
the rational scalar. Character evaluation and derived central
specialization preserve this equality and its cone maps. The
class whose dual exponential occurs in (E.7) is consequently
the same $F_E(Z_F)_h$ as in (K3), and the single remaining factor
two is exactly the $+1$ in (E.7).

\paragraph{From the full cohomology bound to the same determinant scalar.}
Here is the rational divisibility step explicitly. Fix a character
$x$ of $G$, a height-one prime $\mathfrak p$ of $A$, and
$R=A_{\mathfrak p}$. Put $z=F_E(Z_F)_x$, without changing its
normalization. The full global degree-one module is free of rank one
over $R$ by the argument following (K15) and the injectivity of the
maps in (K16). Write it as $Re$, and write $z=\alpha e$.
The higher real cohomology vanishes over $R$, since two is invertible.
Global degree-zero cohomology is zero: an invariant would restrict
to an invariant of the open nonabelian subgroup on the elliptic
representation, on which the scalar coefficient character is trivial.
Thus the positive-cone exact sequence gives
\[
 0\longrightarrow R a\longrightarrow
 H^1(C^+_x\otimes R)\longrightarrow Re\longrightarrow0,
 \qquad
 H^2(C^+_x\otimes R)=H^2(C_{S,x}\otimes R).
\]
The element $a$ here is the image of the fixed real invariant generator
in (E.4), and it is primitive in this positive degree-one module,
because the quotient $Re$ is free. Choose a lift $\widetilde e$ of $e$;
then $(a,\widetilde e)$ is a basis. The equality already proved from
(E.4) becomes
\[
                 u_x^{\rm ev}\theta_x
                  =2\alpha\,a\wedge\widetilde e.
\]
Let $d_1,\ldots,d_r$ be the nonzero Smith entries of a finite free
model of $C^+_x\otimes R$. All its cohomology outside degrees one and
two is zero, its degree-one cohomology is free, and its degree-two
cohomology is torsion. Its inverse determinant lattice, under the
rational identification with $\bigwedge^2H^1$, is therefore generated,
up to an $R$-unit, by
$(\prod_i d_i)a\wedge\widetilde e$. It follows that
\[
 \begin{split}
 v_{\mathfrak p}(u_x^{\rm ev})
 &=v_{\mathfrak p}(\alpha)-\sum_i v_{\mathfrak p}(d_i)\\
 &=\ell_R(H^1(C_{S,x}\otimes R)/Rz)
       -\ell_R H^2(C_{S,x}\otimes R)\ \geq\ 0,
 \end{split}                                                   \tag{E.4a}
\]
where two is a unit and the inequality is exactly (K2).
This uses the full $H^2$ corrected in (K16), rather than only Kato's
$\mathbf H^2_0$.

The intersection of the height-one valuation rings of the normal
domain $A$ gives $u_x^{\rm ev}\in A$ for every $x$. Since $|G|$ is
a power of two, the character idempotents belong to $A[G]$ and these
memberships give $u^{\rm ev}\in A[G]$. This establishes membership
only; it is not a bound for the denominator from character inversion.
The uniform bound comes separately from (E.3): in its free rank-two
$\mathscr O[G]$-module the integral vectors $a,b$ of (E.4) have integral
wedge, so $2^{c_E}u^{\rm ev}\in\mathscr O[G]$ with the fixed exponent
already used for the original lattice. Coefficientwise intersection
in the original group basis gives


\[
              A\cap\mathscr O=\Lambda,\qquad
                  2^{c_E}u^{\rm ev}\in\Lambda[G].                 \tag{E.5}
\]
Indeed an element of $A$ has bounded powers of two as denominators
in its power-series coefficients; membership in $\mathscr O$ makes
their minimum valuation nonnegative. The coefficientwise intersection
retains $c_E$, without a term $v_2|G|$.
\end{proof}

\subsection{Central saturation and the numerical lower bound}
\begin{proposition}[Central calibration of the even scalar]\label{ex:app-central-calibration}
Use the scalar and determinant basis of
Proposition~\ref{ex:app-even-scalar}. At a corank-zero constituent
$E^h$ with finite $\Sha[2^\infty]$, its nonzero central value satisfies
\[
 v_2u_h^{\rm ev}(0)=v_2L_*(h)-2\omega(h)-s(h)+\epsilon_h,
 \qquad |\epsilon_h|\le C_{\rm loc},
\]
with one constant for the fixed local types. At such a constituent
$u_h^{\rm ev}(0)=0$ if and only if $L_*(h)=0$.
Consequently every admissible $h$ with $L_*(h)\ne0$ satisfies
$v_2L_*(h)\ge2\omega(h)+s(h)-c_E-C_{\rm loc}$: the forward
rank-zero theorem in Lemma~\ref{lem:forward-rank} supplies the
preceding corank and finiteness hypotheses in this case. The central saturation corrections are those
in (E.6)--(E.7); no primitivity after inverting $t$ is assumed to
survive central specialization.
\end{proposition}
\begin{proof}\leavevmode
At a corank-zero constituent put
\[
 g=\ell_{\Z_2}E^h(\Q)[2^\infty],\quad
 g_v=\ell_{\Z_2}E^h(\Q_v)[2^\infty],\quad
 r_i=\ell_{\Z_2}H^i(\mathbf R,T_h)\quad(i=1,2).
\]
Let $i$ be the index length of the image of global $H^1$ in the
free rank-one quotient $L_2=H^1(\Q_2,T_h)/H_f^1(\Q_2,T_h)$.
Orthogonal Kummer duality, with its real terms, gives
\[
 \ell H^2_{\rm glob}=s(h)-i+\sum_v g_v+r_2-g.                     \tag{E.6}
\]
Write $K=\ker(H^1_{\rm glob}\to H^1_{\mathbf R})$, let $k$ be its
torsion length and $m$ its free index. If $P=H^1(C_h^+)$, $p_0$ is
its torsion length and $a_s$ the saturation index of the real image
in $P/P_{\rm tors}$, exactness gives
\[
 \ell\operatorname{coker}(H^1_{\rm glob}\to H^1_{\mathbf R})
       =r_1-g+k-m,\qquad k=p_0+a_s.
\]
For the numerical calculation suppose also $L_*(h)\ne0$.
The wedge in (E.4) adds $a_s+1+v_{\rm free}(z)-m$.
Conversion to the inverse determinant adds $p_0-\ell H^2(C_h^+)$.
Substitution of these equalities cancels $m$ and $a_s$ exactly:
\[
 v_2u_h^{\rm ev}(0)
 =v_2\exp^*z-s(h)-\sum_v g_v+2g+1-r_1-e_2,                      \tag{E.7}
\]
where $e_2$ is the valuation of $\exp^*$ on a generator of $L_2$.
Primitivity of the real image after inverting $t$ has not been
specialized to $t=0$.

Here are the real local convention and the exact length sequence used
in (E.6). Put $T=T_2(E^h)$ and $D=E^h[2^\infty]$. Since $E^h[2]$ is
rational, complex conjugation $J$ satisfies $J\equiv1\pmod2$.
The integral projectors $(1\pm J)/2$ give $T=T^+\oplus T^-$.
At the real place use the bounded Kummer complex
\[
 U_{\mathbf R}=\tau_{\le1}R\Gamma(\mathbf R,T)
   =[\,T^+\oplus T^-\xrightarrow{\,J-1\,}T^-\,]
\]
in degrees zero and one, with its map to ordinary real cochains.
Modulo $2^n$ its degree-zero cohomology is
$T^+/2^n\oplus(T^-/2^n)[2]=E^h(\mathbf R)[2^n]$, and its degree-one
image is the real Kummer line $T^-/2$. The latter can also be checked
using a rectangular period lattice: dividing the nonidentity real
component, represented by half the imaginary period, produces the
minus-line cocycle. Thus this line is its exact annihilator under
finite real Tate duality. The ordinary completion
$E^h(\mathbf R)^\wedge_2[-1]$ alone would lose the first summand of
degree-zero cohomology.

At a finite place take the derived Kummer map from
$E^h(\Q_v)^\wedge_2[-1]$. Its reduction has full local degree-zero
cohomology and injective Kummer degree-one map. If $C_{\rm Sel}$
denotes the resulting ordinary global Selmer cone, these assertions
identify
\[
 H^1(C_{\rm Sel}\otimes^{\mathbf L}\Z/2^n)
       =\operatorname{Sel}_{2^n}(E^h/\Q).
\]
At corank zero the discrete Selmer group is finite, the Mordell--Weil
rank is zero, and its order is that of $\Sha[2^\infty]$. The
orthogonal-Kummer specialization of finite Poitou--Tate duality
\cite[5.7.1.4--5.7.1.7]{nekovar}, followed by the compact/discrete
limits, gives
\[
\begin{split}
0\longrightarrow L_2/J_2
 \longrightarrow\operatorname{Sel}_{2^\infty}(E^h/\Q)^\vee
 \longrightarrow H^2_{\rm glob}
 \longrightarrow \bigoplus_{v\in S_f}H^2(\Q_v,T)\oplus H^2(\mathbf R,T)
 \longrightarrow H^0(\Q,D)^\vee\longrightarrow0 ,
\end{split}
\]
where $J_2$ is the image of global $H^1$. There is no real
degree-one quotient: the compact real Kummer image is all of
$H^1(\mathbf R,T)$, and its discrete orthogonal Kummer group is zero.
At finite $v\ne2$ the compact Kummer image is likewise all of local
$H^1$, while at two its quotient is $L_2$. The real degree-two term
is retained. The finite-level images satisfy Mittag--Leffler, and
high-degree localization at the real place terminates the sequence.

This also makes $H^2_{\rm glob}$ finite. The cone sequence then
makes $H^2(C_{\rm Sel})$ finite, and universal coefficients give
\[
 0\to E^h(\Q)/2^n
 \to\operatorname{Sel}_{2^n}(E^h/\Q)
 \to H^2(C_{\rm Sel})[2^n]\to0.
\]
Comparison with the Kummer sequence gives
$\ell H^2(C_{\rm Sel})=s(h)$, with no extra real torsion. Equivalently
the cone sequence is
\[
0\to L_2/J_2\to H^2(C_{\rm Sel})\to H^2_{\rm glob}
 \to \bigoplus_{v\in S_f}H^2(\Q_v,T)\oplus H^2(\mathbf R,T)
 \to H^0(\Q,D)^\vee\to0.
\]
Its alternating lengths are exactly (E.6).

For the positive cone, $H^2_{\rm glob}\to H^2(\mathbf R,T)$ is onto:
any real component's image in $H^0(\Q,D)^\vee$ can be canceled by
a component at two, since restriction
$H^0(\Q,D)\to H^0(\Q_2,D)$ is injective. Thus its exact sequences are
\[
0\to T^+\xrightarrow{a}P\to K\to0,\qquad
0\to\operatorname{coker}(H^1_{\rm glob}\to H^1_{\mathbf R})
 \to H^2(C_h^+)\to H^2_{\rm glob}\to H^2_{\mathbf R}\to0.
\]
If $k$ and $p_0$ are the torsion lengths of $K$ and $P$, $m$ is the
free index of $K$ in global $H^1$, and $a_s$ is the saturation index
of the real line in $P/P_{\rm tors}$, these sequences give
\[
\ell\operatorname{coker}(H^1_{\rm glob}\to H^1_{\mathbf R})
 =r_1-g+k-m,\qquad k=p_0+a_s.
\]
The wedge $2^{-c_E}a\wedge b$ contributes
$a_s+1+v_{\rm free}(z)-m$, and the inverse determinant comparison
contributes $p_0-\ell H^2(C_h^+)$. Substitution leaves
$1+v_{\rm free}(z)+i-s(h)-\sum_vg_v+2g-r_1$.
Finally $v_2\exp^*z=v_{\rm free}(z)+i+e_2$, proving (E.7)
without specializing a primitivity assertion from the punctured
family.
Apply (K3) to the retained class $z=F_E(Z_F)_h$ used in the wedge
above. The primitive-conductor trace and the invariant
$\sqrt h\,e_{\chi_h}$ basis have already been computed in
paragraph~6, including discriminant $4h$. Absorb the fixed
smoothing and old-support factors into $\epsilon_{\rm rec}$.
Then
\[
 v_2\exp^*z
 =v_2L_*(h)+\sum_{q\ {\rm unused}}v_2\#E^h(\mathbf F_q)
                         +\epsilon_{\rm rec},\qquad
 |\epsilon_{\rm rec}|\le C_E.
\]
At unused good primes $g_q=v_2\#E^h(\mathbf F_q)$, so these
terms cancel. At an active good prime the inertia invariants in
$E^h[2^\infty]$ are $E[2]$, of length two for the present curve.
The remaining local types have fixed bounds, including global
torsion by its injection at two.
Thus
\[
 v_2u_h^{\rm ev}(0)=v_2L_*(h)-2\omega(h)-s(h)+\epsilon_h,\qquad
 |\epsilon_h|\le C_{\rm loc}.                                  \tag{E.8}
\]
The central rational positive complex has only rank-two $H^1$, so
this is specialization of (E.4). Corank-zero duality also shows
$u_h^{\rm ev}(0)=0$ if and only if $L_*(h)=0$.
At $L_*(h)\ne0$, Lemma~\ref{lem:forward-rank} supplies the
corank-zero and finite-$\Sha$ hypotheses just used.
Equations (E.5) and (E.8), already with $G=1$, yield
\[
          v_2L_*(h)\ge2\omega(h)+s(h)-c_E-C_{\rm loc}.             \tag{E.9}
\]
No binary-family dimension is needed for this lower bound.
\end{proof}

% Replaces baseline manuscript.tex lines 3416--3510 inclusive.
% Baseline SHA256 525c8c05e7a5acc83d6b1f89c1056fd1321bd7490fd227a63d1a520e941b5197.
\subsection{The actual even form and its phase}
\begin{proposition}[An actual even form with rational normalized phase]\label{ex:app-even-phase}
Fix a sign $\sigma$, unit squareclasses at a finite support $S_0$
containing $2N$, and a positive squarefree reference $D_0$ in that
filter with $L_*(\sigma D_0)\ne0$. For every fixed finite extra
exclusion set $B$ prime to $D_0$, there is an algebraic holomorphic
cuspidal form of weight $3/2$, normalized by $a(D_0)=1$, with the
scalar characters $\psi,\eta$ constructed below and
$\psi/\eta=\chi_{-\sigma}$. It has the prescribed coefficient
support. For every retained positive squarefree $D$ in the original
$S_0$ filter,
\[
 a(D)=z_Dr_D,\qquad z_D\in\mu_{2^\infty},\quad r_D\in\Q,
 \qquad r_D^2=L_*(\sigma D)/L_*(\sigma D_0).
\]
It satisfies the all-index good-prime recurrence (E.12) outside
$S_0\cup B$. Nonzero retained squarefree coefficients have integral
$v_2a(D)$, and $v_2a(D)-\omega(D)$ has a finite attained minimum.
The form and the initial sieve are fixed before this minimum is
chosen. Additional exclusions do not authorize a new minimum.
\end{proposition}
\begin{proof}\leavevmode
Fix a sign $\sigma$ and unit squareclasses at a finite support $S_0$
containing two and every prime dividing $N$. Let $D_0>0$ be squarefree
in these prescribed classes, with $L_*(\sigma D_0)\ne0$.
In (E.10)--(E.11), $D$ ranges over positive squarefree parameters in
these same unit squareclasses at every prime of $S_0$.
Fix also a finite extra exclusion set $B$ with
$\gcd(D_0,\prod_{p\in B}p)=1$.

\paragraph{A nonzero reference functional and its scalar local type.}
Use the global Shimura--Waldspurger packet with underlying
$\operatorname{PGL}_2$ representation $E^{(\sigma)}$ in the
normalization of Waldspurger's Corollary~2. The nonzero central value
at $\sigma D_0$ supplies the member with a nonzero $D_0$-Whittaker
functional. Equivalently, take the nonzero theta lift associated with
$E^{(\sigma D_0)}$ and the additive character of frequency $D_0$;
the change of additive character identifies its underlying packet
with the preceding one. The nonvanishing and the local identification
are the content of Waldspurger's global theta theorem and its
Fourier-functional calculation
\cite[Theorem~1 and Proposition~1]{waldspurgerICM}
and \cite[Proposition~18 and its proof, pp.~466--467]{waldspurger}. The positive real frequency selects the
holomorphic member of lowest weight $3/2$. The lift is in the
unary-orthogonal cuspidal spectrum.

Here is a local construction of a scalar type that retains the
reference functional. At $p\in S_0$, write $\ell=\ell_{D_0}$ and
choose $v$ with $\ell(v)\ne0$. Normalize Haar measure in
\[
 P_Rv=\frac1{\operatorname{vol}(p^{-R}\Z_p)}
       \int_{p^{-R}\Z_p}
          e_p(-D_0b)\pi(n(b))v\,db .
\]
Then $\ell(P_Rv)=\ell(v)$. Let $M(u)$ be the unit action normalized
by the Weil multiplier. Explicitly, if the chosen torus lifts have
cocycle $(u,v)_p$, their theta-cubed multiplier $\gamma$ has
two-power root-of-unity values and satisfies
$\gamma(u)\gamma(v)=(u,v)_p\gamma(uv)$; put
$M(u)=\gamma(u)^{-1}\pi(\widetilde m(u))$.
It is therefore an action of
$\Z_p^\times$ and
$M(u)\pi(n(b))M(u)^{-1}=\pi(n(u^2b))$.
Changing $u$ to its inverse changes none of the following support
conditions. Direct substitution gives
\[
 \ell(M(u)P_Rv)=
 {\bf1}_{\,u^2\equiv1\;(\mathrm{mod}\;p^R)}
                    \ell(M(u)v).
\]
The integral is a finite sum on this smooth vector.

The normalized $M(-1)$ acts by a scalar $\delta\in\{1,-1\}$.
Choose a continuous unit character $\alpha$ of two-power order
with $\alpha(-1)=\delta$. Such a character exists: for odd $p$,
$-1$ lies in the two-primary factor of the finite cyclic group
of residue-field units; for $p=2$ use the factor
$\{1,-1\}$ of $\Z_2^\times$.
Take $R$ large enough that both $v$ and $\alpha$ are constant on
the small neighborhoods appearing above. The projection
\[
 w=\int_{\Z_p^\times}\alpha(u)^{-1}M(u)P_Rv\,d^\times u
\]
satisfies $M(u)w=\alpha(u)w$. Its $\ell$-value is nonzero:
the neighborhoods of $1$ and $-1$ contribute the same nonzero
value. For odd $p$ the solution set is
$1+p^R\Z_p$ together with $-1+p^R\Z_p$.
For $p=2$ and $R\ge3$ it is
\[
 (1+2^{R-1}\Z_2)\ \cup\
 (-1+2^{R-1}\Z_2).
\]
Thus the four residue classes modulo $2^R$ form two neighborhoods;
there are not four independently assignable phases.
For completeness, the noncancellation in this projection is an exact
calculation. Normalize multiplicative Haar measure to give
$\Z_p^\times$ volume one, and put
$K_R=1+p^R\Z_p$ for odd $p$, and $K_R=1+2^{R-1}\Z_2$
for $p=2$. Increase $R$ until $M(k)v=v$ and $\alpha(k)=1$
for $k\in K_R$. On $K_R$ the integrand after applying $\ell$
is $\ell(v)$. On $-K_R$ it is
$\alpha(-1)^{-1}\ell(M(-1)v)=\delta^{-1}\delta\ell(v)$.
It vanishes elsewhere by the displayed additive projection formula.
Consequently
\[
                 \ell(w)=2\operatorname{vol}(K_R)\ell(v)\ne0.
\]
This calculation takes place in characteristic zero; reduction at two
is not involved. For the later support projection choose $m$ so that
the required unit squareclass is a union $A$ of residue classes modulo
$p^m$. Its operator is
\[
 \Pi_A=\frac1{p^m}\sum_{b\bmod p^m}
       \left(\sum_{a\in A}e_p(-ab/p^m)\right)\pi(n(b/p^m)).
\]
On the Fourier functional of integral frequency $D$ it acts as
$\mathbf1_A(D)$. It therefore preserves the reference value and
annihilates every excluded frequency. Conjugation by $M(u)$ permutes
the summands and replaces $A$ by $u^2A=A$, proving the asserted
commutation. Its denominators are fixed before the form is normalized.


The vector $P_Rv$, and then $w$, is fixed by upper unipotents
$n(\Z_p)$, because $D_0$ is a unit. Smoothness makes $w$
fixed by lower unipotents at sufficiently deep level. Together with
the unit eigencharacter, the elementary decomposition of
$K_0(p^m)$ gives a scalar $K_0$ type for sufficiently large $m$.
The Weil normalization contributes only two-power roots of unity,
so the resulting classical character has two-power order.
The forced values at $-1$ give the global parity required for a
weight-$3/2$ scalar character: the product of the finite central
actions is the inverse of the holomorphic archimedean action,
and the theta-cubed multiplier has the same product relation
on rational $-I$. Thus the quotient character satisfies
$\psi(-1)=1$. At $2$ one may, if desired, choose
$\alpha$ of arbitrarily large two-power conductor before increasing
$R$; this does not enlarge the prime support.

Apply the finite additive Fourier idempotent for the prescribed
unit-squareclass subset modulo a sufficiently high power of $p$.
This idempotent keeps the $D_0$ coefficient. Its support condition
is stable under multiplication by unit squares, so it commutes
with the unit action and retains a scalar type, at a possibly
deeper level. At every prime outside $S_0$ take the spherical vector.
The packet there is unramified; its standard squarefree coefficient
$c_p^0$ is $1$ at both valuations zero and one
\cite[Section~I.3, pp.~377--378]{waldspurger}.
This includes the primes dividing $D_0$. Local Whittaker uniqueness
and the tensor-product factorization show that the resulting
global reference coefficient is still nonzero.

\paragraph{The exact full-space comparison.}
Let the resulting scalar character be $\psi$ and put
$\eta=\psi\chi_{-\sigma}$. Both have two-power order.
Let $V$ be the full good-Hecke eigenspace of this scalar character,
at the level just constructed. Keep all old multiplicities. In
Waldspurger's notation its integral-weight eigenform and character are
\[
 \phi_0=f_E\otimes\eta,\qquad
 \chi_0=\psi\chi_{-1},\qquad
 \phi_0\otimes\chi_0^{-1}=f_E\otimes\chi_\sigma .
\]
This also fixes the signs in the packet convention: Waldspurger's
positive index $D$ twists the last representation by $\chi_D$.
The Hecke eigenvalues in $V$ are $\eta(p)a_p(E)$.

The nonzero space implies hypothesis (H1) by Waldspurger's
Proposition~2 \cite[Section~I.2, p.~377]{waldspurger}. In explicit terms, at a principal-series
place its two inducing characters take value $1$ at $-1$.
Corollary~2 \cite[pp.~483--484]{waldspurger} applies to every $f\in V$ and any
positive squarefree pair whose ratio is a square in every
$\Q_p$ at a prime dividing the level. Our $D,D_0$ satisfy
this condition because the level has prime support $S_0$.
The final paragraph of that proof removes (H2) for this
same-$\Q_2$-squareclass comparison. No assumption of
non-supercuspidality at two is required.

\paragraph{The printed coefficient identity and its dyadic hypothesis.}
For clarity, write $\mathcal L_D=L(E^{(\sigma D)},1)$ and let
$N_{\rm W}$ denote the level in Waldspurger's notation, not the
conductor of $E$. In \cite[Corollary~2, pp.~483--484]{waldspurger}
take
\[
 k=3,\quad \chi=\psi,\quad \chi_0=\psi\chi_{-1},\quad
 \rho_0=\rho_E\otimes\eta,\quad n_1=D_0,\quad n_2=D.
\]
Then $\rho_0\otimes\chi_0^{-1}=\rho_E\otimes\chi_\sigma$.
The corollary, in its printed order, is
\[
 a_f(D_0)^2\mathcal L_D\,\psi(D/D_0)D^{1/2}
   =a_f(D)^2\mathcal L_{D_0}D_0^{1/2}.
\]
The passage from the automorphic central argument $1/2$ to the
classical argument $1$ is the same for both terms. Dividing only
by the nonzero reference value gives (E.10). In particular this
is a literal square identity, with neither complex conjugation
nor an unspecified absolute-value normalization.

Every prime of $N_{\rm W}$ belongs to $S_0$ at this point of the
construction. The prescribed local classes make $D/D_0$ a square
in $\mathbb Q_p^\times$ at each such prime, including $p=2$.
In the proof of the corollary, the local coefficient functions
at these primes have equal values at $D,D_0$. Outside the level,
the squarefree arguments have valuation zero or one and the
unramified coefficient function has value one in both cases.
The real coefficient contributes exactly the power
$(D/D_0)^{(k-2)/4}$ before squaring. The quotient of epsilon
factors is computed on pp.~483--484: its Hilbert-symbol product
disappears by reciprocity and its remaining character is
$\chi(D/D_0)=\psi(D/D_0)$. Thus these local terms do not leave
an additional prime-dependent factor in the displayed identity.

The final paragraph on p.~484 explicitly explains the use of
(H2): it is needed in that argument only to compare Fourier
indices in different squareclasses of $\mathbb Q_2^\times$.
Our pair is in one squareclass. Hence the corollary as stated
requires (H1) here, but does not impose (H2), a nonsupercuspidal
condition at two, or a new dyadic conductor restriction. The
nonzero scalar eigenspace supplies (H1) by
\cite[Proposition~2, p.~377]{waldspurger}. Moreover the proof on
p.~483 extends its linear relation to every vector in the full
eigenspace before squaring. The resulting identity therefore
applies simultaneously to old multiplicities, and later to both
Galois-conjugated and character-twisted vectors at a common level.
At that later level the same squareclass test is repeated; its
prime support is still $S_0$.
Putting $n_1=D_0,n_2=D$ in the corollary, with weight $3/2$, gives
\[
 a_f(D)^2=
 \psi(D/D_0)\,
 \frac{\sqrt D\,L(E^{(\sigma D)},1)}
      {\sqrt{D_0}\,L(E^{(\sigma D_0)},1)}
 a_f(D_0)^2.                                      \tag{E.10}
\]
The period cancels, so the quotient in this formula is the quotient
of the specified $L_*$-values. No product over the primes dividing
$D$ or $D_0$ is missing: the unramified factors at squarefree
valuations are exactly those checked above.

For clarity, (E.10) gives proportional coefficient functionals on
the whole space. If $\ell_0(f)=a_f(D_0)$ and $\ell_D(f)=a_f(D)$,
then $\ell_D^2=C\ell_0^2$ as a polynomial identity on $V$.
If $C=0$, then $\ell_D=0$. Otherwise, choose a square root $c$ of
$C$. The polynomial ring on the finite-dimensional vector space
$V$ is an integral domain, so
$(\ell_D-c\ell_0)(\ell_D+c\ell_0)=0$ implies
$\ell_D=c\ell_0$ or $\ell_D=-c\ell_0$.
The scalar is common to every vector. This argument does not
assume that $V$ is one-dimensional.

\paragraph{An algebraic vector and the Galois comparison.}
Ramsey's section model \cite[Definition~3.1 and Theorem~4.2]{ramsey} realizes the half-integral-weight space as
sections of the indicated divisor sheaf on $X_1(4N_{\rm aux})$.
Here $N_{\rm aux}$ is supported on $S_0$ and includes the chosen
local depths; it need not equal the conductor $N$ of $E$.
Over characteristic-zero fields, flat field extension commutes
with these sections. Cuspidality, scalar character, support
conditions and prescribed good-Hecke eigenvalues are algebraic
linear conditions. Although infinitely many coefficients or
primes can occur in their description, a finite subset cuts out
the same subspace, by finite dimensionality.

One need not impose Petersson orthogonality as an additional
transcendental linear condition. The elliptic good-Hecke eigenspace
is automatically orthogonal to unary theta series: at a good prime,
the unary eigenvalues have absolute value $p+1$, whereas
$|\eta(p)a_p(E)|\le2\sqrt p<p+1$. The unary space and its
orthogonal complement are stable under the Hecke operators and
their adjoints. Thus their elliptic eigenspaces are disjoint.
The nonzero reference functional on the complex space is consequently
nonzero on its algebraic model. Choose such an algebraic $g$ and
scale it to $a(D_0)=1$.
This uses flat field extension only, not Ramsey's arbitrary-ring
base-change assertion in weight $3/2$; see \cite[Remark~4.4]{ramsey}.
Here the field-extension assertion can be checked without a
cohomology-vanishing theorem. Work over a number field $K$ containing
the character values and the roots of unity specifying the cusps and
the support idempotents. A finite affine cover of the modular curve
computes global sections as the kernel of the difference of restriction
maps to pairwise intersections. Tensoring with a field extension is
flat and preserves that kernel. Thus the sections in Ramsey's
Definition~3.1 commute with $K\subset\overline\Q\subset\mathbf C$,
also in weight $3/2$. Impose vanishing at the cusps, the diamond
character, the support idempotents, and the kernels of
$T(p^2)-\eta(p)a_p(E)$ inside this finite-dimensional space.
The span in its dual of all the defining linear equations has a finite
basis. These conditions therefore commute with the same field
extensions. Since the reference functional is nonzero after extension
to $\mathbf C$, it is nonzero before extension; a vector in a $K$-basis has
nonzero reference value and may be divided by that value.
This also places every coefficient of the chosen form in one number
field. No assertion about base change to a residue field, nor about
vanishing of $H^1$ in weight $3/2$, is used.


Let $\tau\in\operatorname{Aut}(\overline{\Q}/\Q)$.
Choose an odd integer $e$ describing its action on all values of
$\eta$ and $\psi$, and put $\nu=\eta^{(e-1)/2}$.
Coefficientwise conjugation gives a form $g^\tau$ with character
$\psi^e$ and good eigenvalues $\eta(p)^e a_p(E)$.
One can reduce this conjugation assertion to integral weight by
multiplying by the ordinary theta series: it has rational
coefficients, and its fourth power has trivial scalar character.
The algebraic divisibility condition at the cusps is precisely the
one in Ramsey's section model, so division after conjugation is
again a half-integral-weight form.
The coefficientwise twist $g\otimes\nu$ has character
$\psi\nu^2=\psi^e$ and exactly these good eigenvalues.
The identities hold because $\psi/\eta$ is quadratic.
The twist is modular at a level still supported on $S_0$; this
follows as well by writing it as the finite additive translation
sum using the Gauss sum of the primitive character inducing $\nu$. The full multiplier, cusp and conjugation calculation is in Section~\ref{v15:theta-transport}.

Pass to a common such level. Both forms lie in its full
good-Hecke eigenspace, so the proportionality of its $D$ and $D_0$
functionals gives
\[
 \tau(a(D))
 =\frac{\nu(D)a(D)}{\nu(D_0)a(D_0)}
 =\nu(D/D_0)a(D).                                  
\]
There is no vanishing reference denominator: $D_0$ is prime to
$S_0$ and $a(D_0)=1$. No rationality of the central-value quotient
has been assumed in obtaining this equality.

Write the Dirichlet character as a product of its finite
prime-power factors $\eta_p$. For every $p$ dividing its conductor
choose $u_p\in\Z_p^\times$ with $u_p^2=D/D_0$, and set
\[
                 z_D=\prod_{p\mid\operatorname{cond}(\eta)}
                                      \eta_p(u_p).
\]
Then $z_D$ is a two-power root of unity,
$z_D^2=\eta(D/D_0)$, and
\[
 \tau(z_D)=z_D^e,\qquad
 \nu(D/D_0)=z_D^{e-1}.
\]
Changing the selected odd representative of $e$ has no effect on
the last identity, because $D/D_0$ is a local square at the
conductor. The preceding Galois identity now shows that $a(D)/z_D$ is fixed by
every $\tau$, hence lies in $\Q$.
Different choices of the local square roots change $z_D$ by a
sign, which changes only the sign of this rational number.

Finally, $\psi/\eta=\chi_{-\sigma}$ is trivial on $D/D_0$:
its conductor is supported at two and the ratio is a square there.
Thus $\psi(D/D_0)=z_D^2$. Dividing (E.10) by $z_D^2$ gives
\[
 a(D)=z_Dr_D,\qquad z_D\in\mu_{2^\infty},\qquad r_D\in\Q,
 \qquad r_D^2=\frac{L_*(\sigma D)}{L_*(\sigma D_0)}.       \tag{E.11}
\]
The argument also covers the zero coefficient case. In particular,
every nonzero coefficient in (E.11) has integral normalized
two-adic valuation.

\paragraph{Recurrence, later exclusions, and valuations.}
Only now impose the extra finite coprimality sieve $B$.
The zero-or-retain Fourier sieve is a finite additive translation
operator. It preserves cuspidality, keeps every coefficient whose
index is prime to $B$, and commutes with good $T(p^2)$ for
$p\notin S_1=S_0\cup B$. In particular it keeps $a(D_0)=1$.
The squarefree comparison (E.11) was proved before this enlargement,
so it requires no extra squareclass hypotheses at the primes of $B$.
After the sieve it is asserted at the retained indices.
The standard good $T(p^2)$ formula, with $\psi^2=\eta^2$, gives
\[
 a(p^2n)=\{\eta(p)a_p(E)-\psi(p)(-n/p)\}a(n)
                       -p\eta(p)^2a(n/p^2),\qquad p\notin S_1.
                                                        \tag{E.12}
\]
For squarefree $D$, write $a(Dp^{2r})=\eta(p)^r b_r$.
The first multiplier is
$a_p(E)-(\psi/\eta)(p)(-D/p)$; subsequent multipliers satisfy
$b_{r+1}=a_p(E)b_r-pb_{r-1}$. All are integers, including when
$p\mid D$. Iteration over the good primes gives the same assertion
for every retained square shell after removal of the $\eta$ phase.
More explicitly, for $p\notin S_1$ set
$A_{p,0}(D)=1$, $A_{p,1}(D)=a_p(E)-(\psi/\eta)(p)(-D/p)$,
and
$A_{p,r+1}(D)=a_p(E)A_{p,r}(D)-pA_{p,r-1}(D)$ for $r\ge1$.
Then, for every retained index $Dm^2$ with $D$ squarefree,
\[
 a(Dm^2)=z_D\eta(m)r_D
              \prod_{p^r\Vert m}A_{p,r}(D).
\]
The formula includes $p\mid D$, when the initial Legendre symbol
is zero. When applying the recurrence at one prime, the other
prime-square factors do not alter that Legendre symbol. This proves
the formula by induction on the prime factors of $m$, with
$a(n/p^2)=0$ when $p^2\nmid n$. All factors $A_{p,r}(D)$ are
integers. Thus every coefficient used by the extraction and shell
operators of Proposition~\ref{ex:app-coefficient-rules} is a rational
number times a two-power root of unity, including coefficients with
non-squarefree index. In the chosen residue field such a root reduces
to one. A rational coefficient of valuation zero reduces to the
element $1$ of its prime field $\F_2$. Consequently the fresh
coefficient-one reading of that proposition detects precisely equality
in the fixed minimum-depth bound. It uses the same $g$, $D_0$ and
normalization throughout.


Combining (E.9) and (E.11) makes the values
$v_2a(D)-\omega(D)$ at nonzero retained squarefree coefficients a
bounded-below subset of $\Z$. It is nonempty because $D_0$ is retained,
and therefore has an attained finite minimum. The form and the initial
finite sieve are fixed before this minimum is chosen. Restricting
later fresh primes by further exclusions does not authorize choosing
another minimum.
\end{proof}

% End 07-even-source.tex


% Begin 08-odd-source.tex
\section{Actual odd theta forms and their operators}\label{app:odd-source}
We now construct the forms in (3.8), with the same $h_*$, $D_w$, $L$
and actual Kummer stages as in Section~\ref{sec:rc}.
Write $m=|h_*|$. An admissible $D$ is positive, squarefree, prime to
the fixed support and in the prescribed unit classes. Then
$K_D=\Q(\sqrt{Dh_*})$ is imaginary, split at $2N$, has no extra
units, and $Dh_*$ is an odd fundamental discriminant.
The conductor-one point is the unaveraged sum
\[
 P_D=\sum_{[\mathfrak a]\in\operatorname{Pic}(\mathcal O_{K_D})}
       \chi_{h_*}(\Norm\mathfrak a)\,
                         \pi(x_D^{[\mathfrak a]}).               \tag{O.1}
\]
It is distinct from the conductor-$|h|$ ring-class point used in
Section~\ref{sec:limits}.

% Replace immutable baseline lines 3532--3591 inclusive.
% Keeps (O.2). This expands the existing non-CM argument.
\subsection{A fixed witness and one compatible sequence}
\begin{proposition}[One witness and a compatible detecting sequence]\label{ex:app-detector}
For the fixed $E=E_\ell$ under the full-corank-one hypothesis of
Theorem~\ref{hyp:RC}, use the prepared positive unit filter containing
$1$. There exist a negative odd fundamental discriminant $h_*$ with
$L(E^{h_*},1)\ne0$, a positive fundamental discriminant $D_w$ in
that filter with $L'(E^{D_w},1)\ne0$, and a fixed integer $L$ such
that $LP_D$ belongs to the $D$-twist rational line for every
admissible $D$ and $Q=LP_{D_w}$ is nontorsion.
There is one sequence of actual stages with $M\to\infty$, good
primes $\ell_0$, and additive maps
\[
 \lambda_M:E(\F_{\ell_0^2})\longrightarrow\Z/2^M\Z,
 \qquad v_2\lambda_M(\overline Q)=c_w<M,
\]
where $c_w$ is fixed. At these stages Frobenius squared is trivial
on $E[2^M]$, $(D_w/\ell_0)=1$, $(h_*/\ell_0)=-1$, and
$\lambda_M$ has Frobenius sign $+1$.
Any specified finite prime exclusion is permitted. Additional
field prescriptions require a compatible common lift; no preservation
of $c_w$ under arbitrary later field prescriptions is asserted.
\end{proposition}
\begin{proof}\leavevmode
We use the positive unit-squareclass filter containing one, with its
fixed support enlarged before choosing the witness. The rational
Selmer parity theorem \cite[Theorem 1.4, p.~568]{ddselmer}, applied
over $\Q$ with $p=2$, gives root number $-1$ for $E$ under the
full-corank-one assumption. Apply clause (ii)
of \cite[pp.~543--544]{bfh} to its weight-two newform, with every
old prime in the splitting set, to obtain a negative fundamental
discriminant $h_*$ with $L(E^{(h_*)},1)\ne0$. Including two makes
$h_*$ odd and excludes the extra-unit fields. Now add the primes
of $h_*$ to that set and apply clause (i) to obtain a positive
fundamental discriminant $D_w$ with $L'(E^{(D_w)},1)\ne0$.
Splitting at every old prime puts $D_w$ in the filter containing one
and makes it prime to $h_*$. These choices precede every detecting
prime and every theta precision.

The conductor-one genus character for
$K_D=\Q(\sqrt{Dh_*})$ cuts out
$F_D=\Q(\sqrt D,\sqrt{h_*})$. Thus the unaveraged sum $P_D$ in (O.1)
belongs to $E(F_D)$. Its rational character components are the
$D$- and $h_*$-components. In the degenerate case $D=1$ these are
the trivial and $h_*$ characters of the actual quadratic group.
The $h_*$-component has rank zero by the forward rank-zero theorem.
For $D_w$, the unaveraged height formula
\cite[Theorem 1.1]{cst} has Rankin derivative
$L'(E^{(D_w)},1)L(E^{(h_*)},1)\ne0$, so $P_{D_w}$ is nontorsion.
Here the conductor is one, all primes of $N$ split, and the
extra-unit fields have been excluded.

We make the fixed multiplier explicit. Choose distinct good rational
primes $p_1,p_2$ for $E$ and put
\[
 B_j=\prod_{f=1}^{4}\#E(\F_{p_j^f}),\qquad B=B_1B_2.
\]
For every number field of degree at most four, prime-to-$p_j$
torsion injects into a reduction group with residue degree at most
four. The two primes therefore show that $B$ kills its entire
torsion subgroup. Put $B_{\rm cusp}=C_\pi$, where $C_\pi$
annihilates the image of the cuspidal subgroup defined after (4.1).
This includes every Atkin--Lehner translation $\pi(W\infty)$,
which need not vanish merely because $\pi(\infty)=0$. For $\Delta_D=\operatorname{Gal}(F_D/\Q)$,
the integral character operator
\[
 4e_D=\frac4{|\Delta_D|}
                  \sum_{\gamma\in\Delta_D}\chi_D(\gamma)\gamma
\]
is defined also when $D=1$. The difference
$4P_D-4e_DP_D$ is torsion, since its only possible rational
component is the rank-zero $h_*$-component. Hence
$4BP_D=B(4e_DP_D)$ transforms exactly by $\chi_D$.
Choose $L$ to be a common multiple of the previously required fixed
multiplier and $4BB_{\rm cusp}$, for example their product. This
retains every earlier descent and local-condition factor.
Then $LP_D$ is a rational point of
$E^{(D)}$ under the standard twist identification, for every
admissible $D$. In particular
\[
                         Q=LP_{D_w}\in A(\Q),
                 \qquad A=E^{(D_w)},
\]
is nontorsion. No assertion that the original $P_D$ is
$E(K_D)$-rational is required.

Choose compatible division points $Q_n$ with
$2^nQ_n=Q$ and $2Q_{n+1}=Q_n$, and let
$c_Q(g)=(gQ_n-Q_n)_n\in T_2A$ be the Kummer cocycle.
Write $\rho_A$ for the Tate representation and $H=\ker\rho_A$.
On $H$, this cocycle is a homomorphism. Its image is compact,
and hence a closed $\Z_2$-submodule of $T_2A$. Moreover
\[
 c_Q(ghg^{-1})=\rho_A(g)c_Q(h)\qquad(h\in H),
\]
so its rational span is invariant under the Tate image.

The curve $A$ is non-CM. Open image \cite{serre} supplies both
irreducibility of its rational Tate representation and a central
scalar $z\ne1$ in its Tate image. If $c_Q(H)=0$, the rational
cocycle descends to that image. Commuting $z$ with $g$ gives
\[
 (z-1)c_Q(g)=(g-1)c_Q(z).
\]
Since $z-1$ is invertible over $\Q_2$, this is a coboundary.
It contradicts the Kummer injection
$A(\Q)\otimes\Q_2\hookrightarrow H^1(\Q,V_2A)$ and
nontorsion of $Q$. Thus the translation image is nonzero.
Irreducibility makes its rational span all of $V_2A$.
It consequently contains $2^{a_0}T_2A$ for some $a_0$.

Let $F_0/\Q$ be any fixed finite Galois auxiliary extension,
including $\Q(\sqrt{D_w},\sqrt{h_*})$ and any fixed ring-class
data to be retained with the action of complex conjugation.
Put $H_0=H\cap G_{F_0}$. It has finite index $d$ in $H$.
For $h\in H$, $h^d\in H_0$ and
$c_Q(h^d)=d\,c_Q(h)$. Hence $c_Q(H_0)$ still contains
$2^aT_2A$ for some fixed $a$.
This argument uses no independence between $F_0$ and the Tate
or division-point fields.

Fix actual complex conjugation $\tau$. Since $D_w>0$, the real
twist identification identifies its matrix on $T_2A$ with its
matrix $C$ on $T=T_2E$. We have $C^2=1$ and $\det C=-1$.
Its dual $+1$-eigenspace therefore contains a primitive integral
row $b:T\to\Z_2$, with $bC=b$ and $b(T)=\Z_2$.
Choose $t_0\in T$ with $b(t_0)=1$, and choose $h\in H_0$ with
$c_Q(h)=v=2^at_0$. Set $\sigma=h\tau$. If $d_\tau=c_Q(\tau)$,
then $(1+C)d_\tau=0$, because $\tau^2=1$. Consequently
\[
 \rho_A(\sigma)=C,\qquad
 c_Q(\sigma^2)=(1+C)(v+d_\tau)=(1+C)v,\qquad
 b\,c_Q(\sigma^2)=2^{a+1}.
\]
Put $c_w=a+1$. The element $\sigma$ and this finite valuation
are fixed before precision varies. Its restriction to $F_0$
is that of $\tau$: it fixes $\sqrt{D_w}$ and negates
$\sqrt{h_*}$. Since $h$ fixes $\sqrt{D_w}$, the matrix of
$\sigma$ on $T_2E$ is also $C$.

For each $M>c_w$, take the finite Galois extension generated by
$F_0$, $E[2^M]$, $A[2^M]$, and all $2^M$-division points of $Q$.
Restrict this same actual $\sigma$ to the extension. Finite
Chebotarev supplies infinitely many good rational primes $\ell_0$
in its conjugacy class. Choose a place realizing this representative,
and exclude two, all bad primes, the support of $D_wh_*$ and any
other finite set. Thus, with the torsion basis transported by this
place,
\[
 \operatorname{Fr}_{\ell_0}\equiv C\pmod{2^M},\qquad
 (D_w/\ell_0)=1,\qquad(h_*/\ell_0)=-1.
\]
The determinant also gives $\ell_0\equiv-1\pmod{2^M}$.
In particular $\operatorname{Fr}_{\ell_0}^2$ acts trivially on
$E[2^M]$.

Here is the finite-field map without a normalization ambiguity.
Put $n=2^M$ and write $F$ for Frobenius on the reduction of $E$.
For $P\in E(\F_{\ell_0^2})$, choose $R$ with $nR=P$ and set
\[
                    \kappa_M(P)=F^2R-R\in E[n]=T/nT.
\]
Changing $R$ by $n$-torsion does not change this value.
The map is additive, with kernel $nE(\F_{\ell_0^2})$.
It is onto: $F^2-1$ is a nonzero separable isogeny, so for
$S\in E[n]$ one can solve $(F^2-1)R=S$ and take $P=nR$.
Also $\kappa_M(FP)=F\kappa_M(P)$, by choosing $FR$ as a
division point. Therefore $\lambda_M=b\kappa_M$ has
Frobenius sign $+1$.

Good reduction identifies the reduction of the global division
cocycle with this exact map. The chosen place makes the twist
identification defined over $\F_{\ell_0}$, since
$(D_w/\ell_0)=1$. Thus the fixed witness satisfies
\[
 \lambda_M:E(\F_{\ell_0^2})\longrightarrow\Z/2^M\Z,\qquad
 v_2\lambda_M(\overline Q)=c_w<M.                         \tag{O.2}
\]
Choosing increasing $M$ and excluding previous primes gives one
sequence of actual stages. Subsequent finite prime exclusions
do not alter $c_w$. New field prescriptions require a compatible
common lift and, if added after this choice, are not asserted to
preserve this particular constant.
\end{proof}

% Completion audit replacement for baseline lines 3593--3748.
\subsection{Local functions and the theta sum}
\begin{proposition}[The actual odd theta form and its recurrence]\label{ex:app-theta-form}
Fix an actual stage of Proposition~\ref{ex:app-detector} with $M\ge2$.
Use $m=|h_*|$, the same $L$, the Eichler order of level $N$ in the
quaternion algebra ramified at $\ell_0,\infty$, and the local
Schwartz functions specified below, including (O.3).
With $c_0=|\operatorname{GL}_4(\F_3)|=24\,261\,120$, the sum
$H_M$ in (O.4) is well defined modulo $2^M$ in the character ring.
It gives an actual holomorphic form
\[
 g_M=\sum_n a_M(n)q^n,\qquad a_M(n)=H_M(mn),
 \qquad g_M\in M_{3/2}(\Gamma_0(K),\psi),
\]
with level supported only at the specified coefficient primes.
The characters are exactly
$\psi=\eta_B(m/\cdot)$ and $\eta=\eta_B\chi_{h_*}$, so
$\psi/\eta=\chi_{-1}$. For every prime outside that support,
the recurrence (O.6) holds at every integer index, with nonintegral
indices interpreted as zero. Changing the lifts of $P_M$ loses no
precision. A specified further local coefficient sieve retains this
conclusion outside the enlarged support.
\end{proposition}
\begin{proof}\leavevmode
In this appendix $M=M_\theta$; take $M\geq2$.  Retain the
same $m=|h_*|$, $L$, and detecting prime $\ell_0$ as in (O.2).
Since $\operatorname{Fr}_{\ell_0}^2=1$ on $E[2^M]$, evaluation at
arithmetic Frobenius squared identifies the finite-field Kummer
quotient with $T_2E/2^M T_2E$.  Equivariance and $bC=b$ imply
$\lambda_M(\operatorname{Fr}_{\ell_0} P)=\lambda_M(P)$.
The determinant gives $\ell_0\equiv-1\pmod4$; hence
$(m/\ell_0)=(-1/\ell_0)(h_*/\ell_0)=1$.
\paragraph{The finite-field evaluation map.}
Write $k_0=\mathbb F_{\ell_0^2}$ and let $F$ denote its arithmetic
Frobenius. For $P\in E(k_0)$, choose $Q_P\in E(\overline{k_0})$
with $2^M Q_P=P$. Then
\[
 \kappa_M(P)=FQ_P-Q_P\in E[2^M]
\]
is independent of the choice of $Q_P$, since $F$ is the identity
on $E[2^M]$. This is an additive map. It kills $2^ME(k_0)$,
and its kernel is exactly that subgroup: a zero value means that
$Q_P$ itself lies in $E(k_0)$. The finite group
$E(k_0)/2^ME(k_0)$ has order $|E(k_0)[2^M]|=2^{2M}$, as does
$E[2^M]$, so the induced injection is an isomorphism.
Composing $\kappa_M$ with the fixed row $b$ modulo $2^M$ gives
$\lambda_M$. The arithmetic Frobenius over $\mathbb F_{\ell_0}$
commutes with $F$ and acts on the torsion by $C$. Hence
$\kappa_M(\operatorname{Fr}_{\ell_0}P)=C\kappa_M(P)$ and
$bC=b$ proves the asserted Frobenius sign.
For the fixed witness, specialization of a compatible division
point identifies this evaluation with the cocycle value at
$\sigma^2$, whose $b$-value is $2^{a+1}$. Thus the visibility
valuation is fixed as the precision grows.

By \cite[Main Theorem 14.6.1]{voight-quaternion}, there is a quaternion
algebra $B/\mathbb Q$ ramified exactly at $\ell_0$ and infinity.
Let $R$ be an Eichler order of level $N$.  Choose an enhanced
supersingular reference curve $\mathcal E_0$ with endomorphism order $R$, as in
\cite[Proposition 3.6, pp.~441--443]{ribet-geometry-1990}.
For an enhanced supersingular curve $\mathcal E$, use the right $R$-module
$\operatorname{Hom}(\mathcal E_0,\mathcal E)$ convention of
\cite[Proposition 3.3 and Remark 3.5(a), pp.~440--441]{ribet-geometry-1990}.
Writing its completed lattice as $z\widehat R$ gives
$X=B^\times\backslash\widehat B^\times/\widehat R^\times$,
with endomorphism order $R_z=B\cap z\widehat Rz^{-1}$ and automorphism
group $R_z^\times$.  Here the enhancement is the cyclic subgroup of
order $N$; arbitrary prime-power exponents, including at two, are allowed.
In the source's Section 3 setup one may choose its auxiliary prime
outside $N\ell_0$.
By \cite[Remark 3.5(b), p.~441]{ribet-geometry-1990}, Frobenius acts
on this set by right multiplication by a division uniformizer at
$\ell_0$ and by $1$ elsewhere, and this action is an involution.
Consequently every corresponding geometric coarse point is fixed by
Frobenius squared.  Its closed-point residue degree therefore divides
two, so it is a point of $X_0(N)(\mathbb F_{\ell_0^2})$.
This asserts coarse-point rationality; no chosen scalar-Frobenius
elliptic model is needed.

Define $P_M(z)=\lambda_M(L\pi_{\ell_0}(z))$.  Here the specialization is
legitimate because $\ell_0\nmid N$: the modular curve has good
reduction and the map to the good-reduction elliptic curve extends
over the local base.  The good Hecke correspondences also extend.
Writing $\pi$ as the Abel--Jacobi map based at infinity followed by
the elliptic quotient gives, for $p\nmid N\ell_0$,
\[
 \sum_{z'\text{ a }p\text{-neighbor of }z}\pi_{\ell_0}(z')
       =a_p(E)\pi_{\ell_0}(z).
\]
Indeed the quotient intertwines $T_p$ with $[a_p(E)]$ and
$T_p(\infty)=(p+1)\infty$.  For an Atkin--Lehner involution
$W_{p^e}$ the analogous identity has eigenvalue $w_{p^e}$ and the
constant $\pi(W_{p^e}\infty)$, which $L$ kills.  Frobenius acts on
the supersingular set as the right division-uniformizer operation.
Consequently $P_M$ has the asserted good Hecke and Atkin--Lehner
eigenvalues modulo $2^M$, and uniformizer eigenvalue $+1$.
Choose arbitrary integral lifts of this finite function.

\paragraph{The finite Schwartz function.}
All local functions below are zero off their specified supports.
Set $V=B^0$ and $Q_B(x)=\operatorname{Nrd}(x)=-x^2$.
For $p\nmid2Nm\ell_0$ use $1_{R_p^0}$.
For odd $p^e\Vert N$, use the split integral unit support.  The
cyclic level quotient selects one of the two eigensummands and its
root $r\in\mathbb Z_p^\times$.  The projectors
$(x+r)/(2r)$ and $(r-x)/(2r)$ are integral, so reduction and Hensel
lifting show that the selected root modulo any fixed power of $p$
is locally constant.  Weight it by a unit character $\alpha_p(r)$
of finite $2$-power order with
\[
 \alpha_p(-1)=w_{p^e}\chi_{h_*}(p)^e. \tag{O.3}
\]
Such a character exists: at an odd prime the cyclic $2$-primary
part of $\mathbb F_p^\times$ admits a character taking $-1$ to
$-1$, and the trivial character supplies the other parity.
Conjugation by $R_p^\times$ preserves the cyclic quotient and its
selected summand.  Multiplication of $x$ by a unit $u$ preserves
the summands and replaces $r$ by $ur$, giving multiplier
$\alpha_p(u)$.

At two require $e_x=(1+x)/2\in R_2$ and
$x^2=r^2\in\mathbb Z_2^{\times2}$.  The eigenvalues
$a=(1+r)/2$ and $a'=(1-r)/2$ of $e_x$ differ by the unit $r$.
Its projectors $(e_x-a')/r$ and $(a-e_x)/r$ are integral and remain
distinct modulo two.  If $2^e\Vert N$, the cyclic quotient is a
free rank-one $\mathbb Z/2^e$-module.  Of these complementary
projectors exactly one acts as the identity on it, since the only
idempotents of $\mathbb Z/2^e$ are zero and one.  This defines the
selected root without a choice of its parity.  For every odd $u$,
\[
 e_{ux}=u e_x+(1-u)/2.
\]
This affine transformation preserves the two projectors
individually, so the selected root changes to $ur$, even when
the two eigenvalue parities switch.  The same root character rule
therefore applies at two.  A character with either prescribed
parity exists, using the trivial character or the character modulo
four.  If $2\nmid N$ there is no root character at two.

For $p\mid m$, require $x\in M_2(\mathbb Z_p)^0$ and $v_pQ_B(x)=1$.
The reduction of $x$ is a nonzero nilpotent: zero reduction would
force $p^2\mid Q_B(x)$.  All nonzero nilpotents are conjugate to
$X_0=\left(\begin{smallmatrix}0&1\\0&0\end{smallmatrix}\right)$.
The centralizer of $X_0$ in $\operatorname{GL}_2(\mathbb F_p)$ is
$\{\left(\begin{smallmatrix}a&b\\0&a\end{smallmatrix}\right):a\ne0\}$
and has square determinant.  Thus
\[
 \epsilon_p(gX_0g^{-1})=(\det g/p)
\]
is well-defined.  Conjugation by $u$ multiplies it by $(\det u/p)$,
and scalar multiplication by $u\in\mathbb Z_p^\times$ multiplies
it by $(u/p)$, as follows from conjugation of $X_0$ by
$\operatorname{diag}(u,1)$.

At $\ell_0$ require $x\in R_{\ell_0}^0$ with nonzero reduction
$\bar x=t\beta$ in the trace-zero line of $\mathbb F_{\ell_0^2}$,
and use $(t/\ell_0)$.  Order units act trivially on the residue
field.  A division uniformizer acts by Frobenius, which sends
$\beta$ to $-\beta$, giving multiplier $-1$.
Unit scaling gives $(u/\ell_0)$.  Moreover $\beta^2$ is a nonsquare
in $\mathbb F_{\ell_0}$: otherwise $\beta=\pm a$ for
$a\in\mathbb F_{\ell_0}$, contrary to its nonzero trace-zero condition.
Both $-1$ and $\beta^2$ are nonsquares, so $Q_B(x)=-t^2\beta^2$
is a nonzero square.  Since $(m/\ell_0)=1$, extracted indices have
$(D/\ell_0)=1$.

These supports are compact and open, and all weights are locally
constant.  Additional fixed unit-squareclass conditions multiply
them by locally constant characteristic functions.  They are
preserved by $x\mapsto ux$, which multiplies the norm by $u^2$.
Let $\mathcal O$ be the valuation ring at a place above two of the
field generated by the fixed root-character values at primes
dividing $N$.  All other weights are quadratic, so this field is
independent of $M$ and $\ell_0$.
With $f=\prod f_p$, the preceding calculations give
\[
 f(u^{-1}xu)=\chi_{h_*}(\operatorname{Nrd}u)^{-1}f(x)
       \quad (u\in\widehat R^\times),
 \qquad f(vx)=\alpha(v)f(x)\quad(v\in\widehat{\mathbb Z}^{\times}).
\]
The ramified local quadratic character on units at $p\mid m$ is
$(\cdot/p)$; all other order-unit character factors are trivial.
The scaling character has
\[
 \alpha(-1)=
 \Bigl(\prod_{p^e\Vert N}w_{p^e}\chi_{h_*}(p)^e\Bigr)
             \Bigl(\prod_{p\mid m}(-1/p)\Bigr)(-1)=1.
\]
For the last equality, $m\equiv3\pmod4$ and the first product is
$w(E^{h_*})=+1$.  More explicitly,
$w(E)=-\prod w_{p^e}$ and, since $(h_*,N)=1$,
$w(E^{h_*})=w(E)\chi_{h_*}(-N)$.
The latter twist formula also follows by multiplying the local
unramified-twist factors $\chi_{h_*}(p)^{v_pN}$ at $p\mid N$
and the factors $(-1/p)$ at the newly ramified good primes.

Let $\eta_B$ denote the classical character associated with $\alpha$
in the inverse-idele convention.  The formula
$\eta_B(d)=\prod_p\alpha_p(d)$ for positive integers $d$ prime to
the support is verified in the transformation calculation below.

Set $c_0=|\operatorname{GL}_4(\mathbb F_3)|=24\,261\,120$,
so $v_2(c_0)=9$, and define
\[
 H_M(n)=c_0\sum_{z\in X}
 \frac{P_M(z)\chi_{h_*}(\operatorname{Nrd}z)}{|R_z^\times|}
       \sum_{x\in V(\mathbb Q),\ Q_B(x)=n}f(z^{-1}xz). \tag{O.4}
\]
Left multiplication embeds each $R_z^\times$ into
$\operatorname{GL}_4(\mathbb Z)$; reduction modulo three is
injective on finite subgroups.  Indeed, a prime-order element
$1+3^aA$ in its kernel, with $a\geq1$ and $A$ nonzero modulo three,
has nonzero first term in its prime-th power modulo $3^{a+1}$,
or modulo $3^{a+2}$ when its order is three, a contradiction.
Thus $|R_z^\times|\mid c_0$.
The right-conjugation identity proves independence of the right
representative.  On the left, rational reduced norms are positive
and the global quadratic idele character is trivial on them.
The transformation law below shows that each summand is an actual
holomorphic theta series with the same level and character;
$c_0/|R_z^\times|$ is an integer.
Consequently the coefficients lie in the fixed character ring
$\mathcal O$.  Changing any lift of $P_M$ by $2^M$ changes every
coefficient by an element of $2^M\mathcal O$.  In particular
there is no loss of precision from an automorphism denominator.

\subsection{Transformation law and recurrence}
Use the standard global additive character, self-dual measures,
and the Gaussian $e^{-2\pi Q_B}$ at infinity.  Weil's operator
formulas \cite[I.13]{weil} give multiplication by $\mathbf e_p(bQ_B)$ for
$n(b)$, scaled pullback $f_p(x)\mapsto\gamma_{Q_B}(u)f_p(ux)$ for
$m(u)$ with $u$ a unit, and a constant times Fourier transform
for the Weyl element.  The upper unipotents $n(\mathbb Z_p)$ fix
$f_p$, because $Q_B$ is integral on its support.  If $L'_p$ is any
lattice of translations fixing $f_p$, its Fourier transform has
support in $(L'_p)^\vee$.  Choose $A_p$ with
$p^{A_p}Q_B((L'_p)^\vee)\subseteq\mathbb Z_p$.  Fourier conjugation
then proves that $n^-(p^{A_p}\mathbb Z_p)$ fixes $f_p$.
This includes $p=2$. At two take $A_2\ge2$ and use the comparison
with the ordinary theta multiplier given below; no splitting of
the metaplectic cover over all diagonal units is asserted.
For clarity, the local normalizations and the lower-unipotent
argument are as follows. Take
$\mathbf e_\infty(t)=e^{2\pi it}$ and
$\mathbf e_p(t)=e^{-2\pi i\{t\}_p}$, so that
$\prod_v\mathbf e_v(r)=1$ for $r\in\mathbb Q$.
For $Q(x)=x^{\mathsf T}Ax$ put
$B_Q(x,y)=Q(x+y)-Q(x)-Q(y)=2x^{\mathsf T}Ay$ and set
\[
 \mathcal F_Q\phi(y)=\int\phi(x)\mathbf e_p(B_Q(x,y))\,dx_Q,
 \qquad dx_Q=|\det(2A)|_p^{1/2}\,dx.
\]
Here $dx$ is the coordinate product measure with
$\operatorname{vol}(\mathbb Z_p^3)=1$.
At the real place use the same formula
$dx_Q=|\det(2A)|_\infty^{1/2}dx$ with coordinate Lebesgue measure.
Thus $\mathcal F_Q^2\phi(x)=\phi(-x)$.
If $\phi(x+\lambda)=\phi(x)$ for $\lambda\in L'$, translation
in this integral shows that its Fourier transform vanishes unless
$\mathbf e_p(B_Q(\lambda,y))=1$ for every $\lambda\in L'$.
This is precisely the asserted support in $(L')^\vee$.
Since this dual lattice is compact, an integer $A_p$ with
$p^{A_p}Q((L')^\vee)\subset\mathbb Z_p$ exists. The operator
for $n^-(c)$ is a Fourier conjugate of multiplication by
$\mathbf e_p(-cQ)$; consequently it fixes $\phi$ whenever
$c\in p^{A_p}\mathbb Z_p$. A scalar used to normalize the Weyl
operator cancels with its inverse in this conjugation.

All the specified local functions have such a translation lattice.
At an odd level prime the unit-root condition and its character
depend on a finite residue quotient. At two the square-unit
condition is a congruence modulo eight; the integral projectors
selecting the cyclic quotient, and the finite root character,
again depend on a finite residue quotient. At a prime dividing
$m$, valuation one and the nonzero-nilpotent sign depend on finite
residue data. The division-prime function is already a residue
function. Each support is contained in a lattice and is a finite
union of cosets of a sufficiently small lattice. This proves the
Schwartz--Bruhat and translation assertions for the actual weights.

For the reference $Q_0=\sum X_i^2$, the dyadic dual of
$\mathbb Z_2^3$ is $\tfrac12\mathbb Z_2^3$ and
$\operatorname{vol}_{dx_{Q_0}}(\mathbb Z_2^3)=2^{-3/2}$.
In particular its Fourier transform is
$2^{-3/2}1_{\frac12\mathbb Z_2^3}$, not $1_{\mathbb Z_2^3}$.
The lower-unipotent conclusion nevertheless holds at depth two,
because $4Q_0(\frac12\mathbb Z_2^3)\subset\mathbb Z_2$.
These positive measure factors affect neither the diagonal
eigenvalue quotient nor the counting coefficients of the theta sum.


The diagonal factor can be computed directly.  For
$Q_B=\sum_{i=1}^3 a_iX_i^2$, the one-dimensional Weil-index identity
\[
 \frac{\gamma(ua_i)}{\gamma(a_i)}
       =\frac{\gamma(u)}{\gamma(1)}(u,a_i)_p
\]
and multiplication under orthogonal sums give, relative to the
ordinary theta multiplier cubed, the factor
$(u,a_1a_2a_3)_p$.  These identities are \cite[II.25 and II.28]{weil};
the latter obtains the Hilbert symbol from the four-dimensional
quaternion norm form.  For $B=(-a,-b)$ the norm on $B^0$ is
$aX^2+bY^2+abZ^2$, whose determinant is $(ab)^2$.
There is therefore no additional determinant character for $Q_B$.
\paragraph{The normalized diagonal factor.}
The Weil-index identity used here can be derived without choosing
Gauss-sum signs. The norm form of the quaternion algebra $(u,a)$
is $\langle1,-u,-a,ua\rangle$. By multiplicativity under orthogonal
sums and $\gamma(-q)=\gamma(q)^{-1}$, its index is
\[
 \frac{\gamma(1)\gamma(ua)}{\gamma(u)\gamma(a)}.
\]
For a split quaternion algebra the norm form is hyperbolic and its
index is one; for the division algebra its index is minus one by
\cite[II.28, Proposition 4]{weil}. These are exactly the two values
of $(u,a)_p$. Solving this equality gives
\[
 \frac{\gamma(ua)}{\gamma(a)}
   =\frac{\gamma(u)}{\gamma(1)}(u,a)_p.
\]
Applying it to three diagonal entries gives the ordinary
one-dimensional theta factor cubed times $(u,\det Q_B)_p$.
This derivation includes $p=2$. A change of diagonalization changes
the determinant only by a square, so it changes no symbol.
For $Q_B/m$ the determinant becomes $m^{-3}\det Q_B$, whose
squareclass is $m$. Thus the extra factor is precisely $(u,m)_p$.
This establishes the determinant-character contribution in (O.5)
with the same additive character as the original theta multiplier.

\paragraph{An operator construction of the common double cover.}
Here is a direct justification that the same lifts may be used in
the comparison. It also fixes the signs of the diagonal operators.
Work over $k=\mathbb Q_p$, including $p=2$, or $k=\mathbb R$,
with the additive character just specified. For a nondegenerate
ternary $Q$ write $\gamma_Q=\gamma(Q)$ and define
\[
 \begin{split}
 U_Q(b)\phi(x)&=\mathbf e_k(bQ(x))\phi(x),\\
 D(a)\phi(x)&=|a|_k^{3/2}\phi(ax),\\
 W_Q&=\gamma_Q\mathcal F_Q,\qquad
 M_Q(a)=r_Q(a)D(a),\qquad
 r_Q(a)=\frac{\gamma(aQ)}{\gamma(Q)}.
 \end{split}
\]
The Weyl matrix in this convention is
$w=\left(\begin{smallmatrix}0&1\\-1&0\end{smallmatrix}\right)$.
The opposite Fourier-sign convention changes both compared
representations together.

The Fourier formula of \cite[I.14, Theorem~2]{weil}, or its
completion-of-the-square calculation, says for $b\ne0$ that
\[
 \int_{k^3}
 \mathbf e_k\bigl(bQ(z)+B_Q(z,v)\bigr)\,dz_Q
   =|b|_k^{-3/2}\gamma(bQ)\mathbf e_k(-b^{-1}Q(v)).
\]
The integral has the oscillatory meaning of that theorem.
Applying it to a Schwartz function gives the operator identity
\[
 \mathcal F_Q U_Q(b)\mathcal F_Q
 =\gamma(bQ)U_Q(-b^{-1})D(b^{-1})
                   \mathcal F_Q^{-1}U_Q(-b^{-1}).
\]
Thus no choice of an individual square root of a Gauss sum is
needed. The index identities proved above also give
\[
 \gamma(Q)^2=(\det Q,-1)_k\gamma(1)^6,
 \qquad \gamma(1)^4=(-1,-1)_k,
 \qquad r_Q(a)r_Q(b)=(a,b)_k r_Q(ab).
\]
For the middle equality apply the quaternion-norm formula to
$\langle1,1,1,1\rangle$; the first follows by applying the same
formula to $\langle1,-a_i,1,-a_i\rangle$ for each diagonal
entry $a_i$ of $Q$.

Consequently the displayed operators satisfy
\[
 \begin{gathered}
 U_Q(b)U_Q(c)=U_Q(b+c),\qquad
 M_Q(a)U_Q(b)M_Q(a)^{-1}=U_Q(a^2b),\\
 M_Q(a)M_Q(b)=(a,b)_kM_Q(ab),\qquad
 W_QM_Q(a)W_Q^{-1}=M_Q(a^{-1}),\\
 W_Q^2=(-1,-1)_kM_Q(-1),\\
 W_QU_Q(b)W_Q
 =(-b,-1)_kU_Q(-b^{-1})M_Q(-b^{-1})W_QU_Q(-b^{-1})
 \quad(b\ne0).
 \end{gathered}
\]
For the last equality, the preceding Fourier identity leaves
the scalar $\gamma(Q)^2\gamma(bQ)^2$; the index-square formula
reduces it to $(-b,-1)_k$. Every scalar in these relations is
independent of the ternary form. Interpreting a scalar $-1$ as
the central element $\varepsilon$ gives one common double
cover of $\operatorname{SL}_2(k)$. Indeed these relations reduce
every word to one of
$\varepsilon^eU(x)M(a)$ or
$\varepsilon^eU(x)M(a)WU(y)$, with $e\in\{0,1\}$.
The ordinary Bruhat decomposition makes their projections
unique, apart from $\varepsilon^e$; the operators send
$\varepsilon$ to $-\operatorname{id}$, so the central kernel
has exactly two elements. This realizes the rank-three pullback
of Weil's metaplectic operators, rather than assuming a section
over the diagonal units.

For $Q_0=\langle1,1,1\rangle$ these operators are the tensor
products of the three ordinary one-dimensional theta operators.
For either $Q$ or $Q_0$, changing a lift multiplies its operator
by the same $-1$. Hence the quotient of their eigenvalues is
independent of the lift. At two the relation
$M_Q(-1)^2=(-1,-1)_2\operatorname{id}=-\operatorname{id}$
explicitly retains the order-four lift of $-I$.

On a unit $a$ the quotient of the two diagonal index factors is
\[
 \frac{r_Q(a)}{r_{Q_0}(a)}=(a,\det Q)_k.
\]
Together with $f_p(ax)=\alpha_p(a)f_p(x)$, this is exactly
the eigenvalue quotient $\alpha_p(a)(a,\det Q)_p$ used below.
There is no remaining eighth-root factor at two.

\paragraph{The metaplectic comparison at two and the extracted character.}
We specify how the preceding index calculation fixes the multiplier,
including its dyadic part. Let
\[
 Q_0=X_1^2+X_2^2+X_3^2,\qquad
 \phi_{0,p}=1_{\mathbb Z_p^3},\qquad
 \theta(\tau)=\sum_{n\in\mathbb Z}e^{2\pi i n^2\tau}.
\]
Use the same additive character and metaplectic lifts for the
rank-three Weil representations of $Q_B$ and $Q_0$. The latter
theta series is literally $\theta(\tau)^3$; this fixes the meaning
of the ordinary theta multiplier cubed, including its eighth roots
of unity. For either representation the nontrivial central element
of the metaplectic cover acts by $-1$. Consequently the ratio of
two eigenvalues computed using the same lift is independent of
that lift. In particular a central factor in multiplying the
three lifts in the displayed triangular factorization cancels
in this ratio. We compare eigenvalues, and do not choose a group
section of the cover over $\Gamma_0(2^{A_2})$.

Here is the comparison on the three factors. Both upper unipotents
act trivially on the indicated vectors. Choose $A_p$ large enough
that both lower unipotents do also. This is the Fourier-support
argument above for $f_p$; for $\phi_{0,p}$ it requires no positive
depth at an odd prime and depth two at two, since the bilinear dual
of $\mathbb Z_2^3$ for $Q_0$ is $\tfrac12\mathbb Z_2^3$ and
$4Q_0(\tfrac12\mathbb Z_2^3)\subset\mathbb Z_2$.
The Weyl scalar cancels in conjugating an upper unipotent to a
lower one, so it introduces no additional eigenvalue. On a
diagonal unit $u$, the scaling part acts by $\alpha_p(u)$ on
$f_p$ and by one on $\phi_{0,p}$. The quotient of the Weil-index
parts, by the preceding calculation, is
\[
 \frac{\gamma(uQ_B)/\gamma(Q_B)}
      {(\gamma(u)/\gamma(1))^3}
       =(u,\det Q_B)_p.                                      \tag{O.5a}
\]
If the inverse Weil-index convention is used for the diagonal
operator, both numerator and denominator are inverted; the
quotient is unchanged because the Hilbert symbol has order two.
Thus this computation does not leave an undetermined inverse
or a dyadic root of unity. These operator comparisons are precisely
the scaled pullback, Fourier conjugation and scalar relations of
\cite[I.13--I.14 and II.25--II.28]{weil}.

Here $\det Q$ means the determinant of the matrix in
$Q(x)=x^{\mathsf T}Ax$. The associated bilinear Gram matrix is
$2A$. If that convention is used instead, both $Q$ and $Q_0$
acquire the same factor $2^3$ in their determinants, which cancels
in the comparison. In particular no extra squareclass of two
is introduced at the dyadic place.

For a general ternary form $Q$ in this argument, at each prescribed
prime take $A_p\ge1$ and increase $A_p$
also to the conductor of $u\mapsto(u,\det Q)_p$. The eigenvalue
relative to the ordinary theta multiplier, on a matrix whose
upper-left entry is $a$, is then
\[
                    \alpha_p(a)(a,\det Q)_p.                  \tag{O.5b}
\]
Indeed the triangular factorization has just the two trivial
unipotent factors and this diagonal factor, and its lift ambiguity
cancels as explained above. This also verifies multiplication:
the upper-left entry of a product is congruent to the product
of the upper-left entries modulo $p^{A_p}$, so both local
characters in (O.5b) multiply. The argument therefore applies
to every matrix in the indicated congruence subgroup at a prescribed
prime, not only to a choice of generators. At an unramified prime
outside the support, the full spherical invariance established
below treats matrices with nonunit upper-left entry as well.

For $Q_B$ the determinant squareclass is one at every place,
including two. Thus (O.5b) has no determinant factor. For the
extracted form $Q_B'=Q_B/m$ it is $m^{-3}\equiv m$.
At an odd prime $p\mid m$ and a unit $u$, and at two, respectively,
the additional factors are explicitly
\[
 (u,m)_p=(u/p)^{v_p(m)},\qquad
 (u,m)_2=(-1)^{((u-1)/2)((m-1)/2)}.
\]
Here $m$ is odd and $m\equiv3\pmod4$, so the second factor is
$(-1)^{(u-1)/2}$. It is already trivial on $1+4\mathbb Z_2$.
At every odd prime outside $m$ the factor on units is one. There
is consequently no new level prime, and there is no omitted
dyadic character.

Finally take a positive integer $d$ prime to the complete level.
The finite operator used in the classical transformation is the
one at $\gamma^{-1}$, with upper-left entry $d$. The root-character
factor is therefore $\prod_p\alpha_p(d)=\eta_B(d)$.
Hilbert reciprocity gives
\[
 \prod_{p\mid 2m}(d,m)_p
   =\prod_{q\mid d}(m/q)^{v_q(d)}=(m/d).                       \tag{O.5c}
\]
Indeed the infinite symbol is one because $m,d>0$, all other
finite symbols are one, and at an odd $q\mid d$ the symbol is
$(m/q)^{v_q(d)}$. This proves the displayed product without
discarding its dyadic factor. Both sides extend to the same
even Dirichlet character. Since $h_*=-m$, division by the genus
character now gives
\[
 \frac{\eta_B(d)(m/d)}{\eta_B(d)\chi_{h_*}(d)}=(-1/d).
\]
Thus the characters in (O.5), with the stated inverse-idele
convention for $\eta_B$, are exactly the characters supplied
by these local Weil operators.

Take $A_p$ also at least the conductor exponent of $\alpha_p$.
For $\gamma=\left(\begin{smallmatrix}a&b\\c&d\end{smallmatrix}\right)$
with $c\in p^{A_p}\mathbb Z_p$ and $a$ a unit, use
\[
 \gamma=n^-(c/a)m(a)n(b/a).
\]
The scalar is the diagonal theta factor times $\alpha_p(a)$.
In passing to the classical transformation law one applies the
finite operator at $\gamma^{-1}$, whose upper-left entry is $d$.
Thus the classical Dirichlet character is
$\eta_B(d)=\prod_p\alpha_p(d)$, not its inverse.  Its value at
$-1$ is one, as required by the weight-$3/2$ theta multiplier.

There are no additional level primes hidden in this argument.
Outside the prescribed support, $p$ is odd, $B_p$ is split, and
$R_p^0=\mathfrak{sl}_2(\mathbb Z_p)$ is self-dual for the bilinear
form of $Q_B=-a^2-bc$.  Its determinant is a unit.  Hence the
Fourier transform of $1_{R_p^0}$ equals itself; the unramified
Weil constant is one, as also follows from \cite[II.27]{weil} by integrating
over this self-dual lattice.  Upper unipotents, diagonal units,
and the Weyl element fix it.  These generate
$\operatorname{SL}_2(\mathbb Z_p)$, so the exponent there is zero.
At the finitely many prescribed primes the preceding $A_p$ suffice.
Conjugation by any $z\in\widehat B^\times$ is an isometry of $Q_B$,
commutes with the Weil representation, and preserves the
self-dual measure.  The same level and character therefore work
for every summand indexed by $z$.

\cite[III, Theorem 6]{weil} applies to this Schwartz--Bruhat function on
$V(\mathbb A)$ and proves invariance of its rational theta sum.
The normalized archimedean expression is
$\sum_x f(x)e^{2\pi i\tau Q_B(x)}$: the factor $y^{3/4}$ in the
Weil action is removed in passing to the classical form.
It has holomorphic weight $3/2$.  At a rational cusp, the slash
operation replaces the finite Schwartz function by its finite
Weil transform.  Its support meets $V(\mathbb Q)$ in a finite
union of rational lattice cosets, with bounded denominators.
Positive definiteness gives nonnegative rational exponents of
bounded denominator and normal convergence for $\Im\tau>0$.
This proves holomorphy at every cusp.  Theorem 6 alone asserts
automorphic invariance; this last argument supplies the
holomorphy conclusion used here.

\paragraph{Extraction.}
The support at every $p\mid m$ has $v_pQ_B=1$.  All nonzero
indices are integral and divisible by $m$.  Repeat the preceding
construction for $Q_B'=Q_B/m$ with the same finite weights and the
Gaussian for $Q_B'$.  This quadratic form is integral on the
support.  Away from $m$ it differs by a unit, so the good-prime
self-duality argument remains valid.  Its determinant squareclass
is $m^{-3}\equiv m$, giving the extra Dirichlet character
$(m/\cdot)$.  Hence
\[
 g_M(\tau)=\sum_n H_M(mn)q^n\in
      M_{3/2}(\Gamma_0(K),\psi).
\]
Here $K$ has only the stated prime support, and
\[
 \psi=\eta_B(m/\cdot),\qquad
 \eta=\eta_B\chi_{h_*},\qquad
 \psi/\eta=\chi_{-1}. \tag{O.5}
\]
This directly constructs the extracted form without a
half-integral-weight extraction operator.
The inverse in this classical passage can also be checked before
any character notation is introduced. Put
$g_\tau=n(x)m(\sqrt y)$ for $\tau=x+iy$, so $g_\tau i=\tau$.
For $\gamma\in\operatorname{SL}_2(\mathbb Q)$ write
$\gamma g_\tau=g_{\gamma\tau}k_\infty$ with
$k_\infty\in\operatorname{SO}_2(\mathbb R)$. In the adelic group,
\[
 (g_{\gamma\tau},1_f)
   =\gamma\,(g_\tau,\gamma_f^{-1})(k_\infty^{-1},1_f).
\]
This equality also holds with consistent metaplectic lifts.
The global rational lifts are compatible for the two forms:
\cite[II.30, Proposition~5]{weil} gives
$\prod_v\gamma_v(Q)=\prod_v\gamma_v(aQ)=1$ for rational $Q,a$,
and hence the local normalizing factors in $W_Q$ and $M_Q(a)$
have product one. The product formula also gives
$\prod_v|\det(2A)|_v^{1/2}=1$, so the product of the stated
self-dual measures is the measure for which the rational lattice
has covolume one. Apply the rational invariance of
\cite[III.41, Theorem~6]{weil} to the displayed adelic identity.
The archimedean Gaussian for either positive ternary form has
the same weight $3/2$; therefore its factor from $k_\infty^{-1}$
is the same as that for $\theta^3$. The relative finite factor
is evaluated at $\gamma_f^{-1}$, whose upper-left entry is $d$.
Thus the relative classical character is
$\prod_p\alpha_p(d)(d,\det Q)_p$, over the prescribed primes.
This proves the inverse-idele convention directly. It compares
the transformation identities and does not divide by a value
of a theta function.

For the extracted form there is likewise no measure or degree
normalization in its coefficients. With exactly the same finite
weights, replace only the quadratic form and archimedean Gaussian
by $Q_B/m$ and $e^{-2\pi Q_B/m}$. The summand indexed by $x$ now
has exponent $Q_B(x)/m$. The original support makes $Q_B(x)$ an
integer with valuation one at every prime dividing $m$, so these
exponents are nonnegative integers. Coefficient $n$ is exactly
the original sum over $Q_B(x)=mn$, namely $H_M(mn)$.
The self-dual measure changes by $|m|_v^{-3/2}$ locally, but a
theta coefficient is a discrete sum and contains no Haar
measure; this change is used only in the Fourier operators.
Outside the original support it is a unit rescaling of the same
unimodular local form, and at primes dividing $m$ the
lower-unipotent depth can be increased within the existing
support. The determinant changes by $m^{-3}$, whose squareclass
is $m$, so (O.5c) supplies exactly $(m/\cdot)$.
In particular no new level prime, factor of $m$, or dyadic
root of unity is introduced by extraction.


\paragraph{All-index recurrence and its phases.}
Let $p$ avoid the full support and put $s=\eta_B(p)$.
Locally use $R_p=M_2(\mathbb Z_p)$ and $V_p=\mathfrak{sl}_2$.
The $p+1$ neighbors are indexed by lines of $\mathbb F_p^2$.
For integral $x$ with nonzero reduction and $Q_B(x)=n$, the number
of preserved lines is $1+(-n/p)$, including the nilpotent case.
If $x\in pR_p^0$, all $p+1$ lines are preserved; the excess
over the preceding formula is $p$.

A nonintegral $x$ in a neighbor lies in $p^{-1}R_p^0$.
Write $x=y/p$ with $y$ primitive and $Q_B(y)=p^2n$.
Then $\bar y$ is nonzero nilpotent.  In a basis with
$\bar y=X_0$, write
$y=\left(\begin{smallmatrix}a&b\\c&-a\end{smallmatrix}\right)$.
Here $p\mid a,c$, $b$ is a unit, and
$p^2\mid-a^2-bc$ forces $p^2\mid c$.  Exactly the neighbor
belonging to the kernel line makes $y/p$ integral.  The integral
and nonintegral parts therefore exhaust the sum with no genericity
assumption on $n$.

Scaling by $p$ at all other places multiplies the weight by $s$.
Thus the primitive nonintegral contribution is
$s^{-1}H_M(p^2n)-H_M(n)$, and the integral contribution is
$(1+(-n/p))H_M(n)+psH_M(n/p^2)$.  Their sum is
\[
 s^{-1}H_M(p^2n)+(-n/p)H_M(n)+psH_M(n/p^2).
\]
Here and below a coefficient at a nonintegral index is zero.
For completeness, if $N_{ij}$ counts neighbor kernels from
$i$ to $j$ and $w_i=|\operatorname{Aut}(i)|$, counting actual
isogenies and dualizing gives
$N_{ij}w_j=N_{ji}w_i$, or $N_{ij}/w_i=N_{ji}/w_j$.
This proves the self-adjointness with the mass $1/w_i$ appearing
in $H_M$.  A neighbor changes the norm character by
$\chi_{h_*}(p)$.  Moving the neighbor operator onto $P_M$ and
using its actual modular-parametrization identity consequently
gives
\[
 H_M(p^2n)+s(-n/p)H_M(n)+ps^2H_M(n/p^2)
     =s\chi_{h_*}(p)a_p(E)H_M(n)\pmod {2^M\mathcal O}.
\]
This computation explains both powers of $s$ even when $s$ is
not quadratic.

Put $a_M(n)=H_M(mn)$ and
$\eta=\eta_B\chi_{h_*}$.  Since
$(-mn/p)=(m/p)(-n/p)$ and $\eta(p)^2=s^2$, one obtains
\[
 a_M(p^2n)=
 \{\eta(p)a_p(E)-\psi(p)(-n/p)\}a_M(n)
        -p\eta(p)^2a_M(n/p^2)\pmod {2^M\mathcal O}. \tag{O.6}
\]
The proof covers $p\mid n$ and $p^2\mid n$ explicitly.
At zero the coefficients vanish by the local support, and negative
indices vanish by positive definiteness.
An added finite coprimality or unit-squareclass sieve is realized
by the corresponding local characteristic function for $Q_B'$;
the same modularity proof applies.  At primes outside the enlarged
support the three indices $p^2n,n,n/p^2$ have the same retain-or-zero
status.  Hence these sieves preserve the displayed recurrence.
\end{proof}

\subsection{The unaveraged point comparison}
\paragraph{CM and specialization inputs.}
For the conductor-one comparison, the fixed-type ideal action, its CM
reciprocity convention and forward inert reduction are recorded in
Section~\ref{v15:cm-reduction}.
Common good models and the same-place specialization of the labelled
prime-to-$\ell_0N$ ideal quotients, including the selected level and
residue orientations, are proved in
Section~\ref{v15:cm-reduction}.
The enhanced supersingular class-set and Frobenius inputs remain separate.
\paragraph{The actual weighted coefficient comparison.}
\begin{proposition}[The unaveraged weighted theta comparison]
\label{ex:app-theta-comparison}
Put $m=|h_*|$ and $F=\mathbf Q(\sqrt{Dh_*})$, where $D$ is a
positive squarefree parameter in the fixed odd filter. Assume
$(D/\ell_0)=1$ at the chosen actual stage, so that $\ell_0$ is
inert in $F$. Keep the local Schwartz functions, the actual function
$P_M(z)=\lambda_M(L\pi_{\ell_0}(z))$, and the integer $c_0$
of (O.4). Then a $2$-power root of unity $\zeta_M(D)$ satisfies
\[
 H_M(mD)=c_0\,2^{\omega(N)}\zeta_M(D)
                   \lambda_M(\overline{LP_D})\pmod {2^M}.\tag{O.7}
\]
The root of unity belongs to the field generated by the fixed
characters at the primes of $N$. No class number is divided out.
\end{proposition}


\begin{proof}
The filter makes $Dh_*$ an odd fundamental discriminant, with
$F$ split at every prime dividing $2N$. Its unit group is
$\{\pm1\}$. The condition $Q_B(x)=mD$ is equivalent to
$x^2=Dh_*$. Together with $(1+x)/2\in R_{z,2}$ and the other
integrality conditions, this says exactly that
\[
 \mathcal O_F\longrightarrow R_z,
 \qquad \sqrt{Dh_*}\longmapsto x
\]
is an optimal embedding. Optimality follows because the source
order is already maximal. All nonzero local weights on such an
embedding are roots of unity. The valuation-one condition at each
prime of $h_*$ holds because $(D,h_*)=1$; at $\ell_0$ the residue
of $x$ is nonzero. The fixed unit filters hold by admissibility.

First classify the local embeddings. Away from $N\ell_0$ there
is one local orbit: an $\mathcal O_{F,p}$-stable lattice is locally
a fractional ideal, hence principal, and this also covers ramified
primes dividing $Dh_*$. At $p^e\Vert N$, the algebra is split.
The two integral idempotents decompose both lattices of the
Eichler lattice pair. Its quotient is cyclic of order $p^e$,
so its elementary exponents lie wholly in one summand. The two
possibilities give exactly two orientations. This applies at two
using the integral spectral projectors specified in the construction.
At $\ell_0$, the maximal order of the division algebra has residue
field $\mathbf F_{\ell_0^2}$. The two residue-field embeddings
of the unramified quadratic order give the two orientations.
Consequently there are $2^{\omega(N)+1}$ local orientation tuples.

Fix one tuple. Choose a rational embedding
$\iota:F\hookrightarrow B$ and write $x_0=\iota(\sqrt{Dh_*})$.
Such an embedding exists because $F$ is a field at the two ramified
places of $B$, namely $\ell_0$ and infinity
\cite[Proposition 14.6.7, $(v)\Rightarrow(iii)$]{voight-quaternion}.
Choose $z_0$ realizing
the selected integral local embeddings. The adelic representatives
with this tuple are exactly
\[
       \widehat F^{\times}z_0\widehat R^{\times}.
\]
Their quotient by the rational centralizer $F^{\times}$ is
\[
 F^{\times}\backslash\widehat F^{\times}/
      (\widehat F^{\times}\cap z_0\widehat R^{\times}z_0^{-1})
 =F^{\times}\backslash\widehat F^{\times}/\widehat{\mathcal O}_F^{\times}
 =\operatorname{Pic}(\mathcal O_F).
\]
Thus this is an actual class-group torsor, not an average or a
quotient of a class-group orbit.

For an ideal-class representative $a\in\widehat F^{\times}$,
the element $a$ centralizes $x_0$, and therefore
\[
 f((az_0)^{-1}x_0(az_0))=f(z_0^{-1}x_0z_0).
\]
On the other hand,
\[
 \chi_{h_*}(\operatorname{Nrd}(az_0))
 =\chi_{h_*}(\operatorname{N}_{F/\mathbf Q}a)
                    \chi_{h_*}(\operatorname{Nrd}z_0).
\]
The first factor is exactly the genus character in (O.1).
To see that it is well-defined on the class group, for $D>1$
the biquadratic field $\mathbf Q(\sqrt D,\sqrt{h_*})$ has
discriminant $D\,h_*\,(Dh_*)=(Dh_*)^2$. Its relative
discriminant over $F$ is therefore one. The character obtained
by restricting $\chi_{h_*}$ to $G_F$ is unramified and, under
class field theory, is $a\mapsto
\chi_{h_*}(\operatorname{N}_{F/\mathbf Q}a)$. For $D=1$ it
is the trivial character. Ideal representatives can be chosen
prime to $h_*$, giving exactly the elementary norm formula in
(O.1). Reversing the Artin convention has no effect because this
character is quadratic.

We use the direct conductor-one reduction argument in
Section~\ref{v15:cm-oriented}. Its retained inputs are
CM ideal-class reciprocity, potentially good CM reduction and inert-prime
supersingularity, extension and specialization of endomorphisms and
prime-to-$\ell_0$ subgroup quotients, and the geometric and Tate-lattice
description of supersingular pairs used above. The latter input, including
its Frobenius compatibility, remains separate from the reduction argument.
The proof computes the labelled target torsor using the pointwise
centralizer $F^\times$, independently of reduction surjectivity. Ideal
representatives prime to $N\ell_0$ then give actual subgroup quotients
whose reduction, at the one fixed place, agrees with the same ideal action.
Thus reduction is an equivariant bijection onto the full orientation tuple
whose $\ell_0$ sign is induced by the fixed CM type. Its multiplicity is
one. Restricting to any underlying supersingular pair gives the same
bijection on its fiber. This uses maximal $\mathcal O_F$, every split
prime-power level (including powers of two), and $\ell_0\nmid N\operatorname{disc}(F)$.
The other level orientations have geometric Atkin--Lehner anchors; the
opposite $\ell_0$ sign is obtained by Frobenius, not by claiming it is in
the image of the same fixed-type reduction map.

Choose the reduction place compatibly with (O.2), and choose $z_0$ as the
reduction of the actual reference Heegner point $x_D$ in (O.1). Writing
\[
 \chi_F([a])=\chi_{h_*}(\operatorname{N}_{F/\Q}a),\qquad
 z_a=[az_0],
\]
equivariance gives the exact finite-sum identity
\[
 \sum_{[a]\in\operatorname{Pic}(\mathcal O_F)}
       \chi_F([a])P_M(z_a)
 =\lambda_M\!\left(
       \overline{L\sum_{[a]\in\operatorname{Pic}(\mathcal O_F)}
                         \chi_F([a])\pi(x_D^{[a]})}\right)
 =\lambda_M(\overline{LP_D})\quad\text{in }\Z/2^M\Z.
\]
Here the supersingular coarse points and their parametrization values are
defined over $\mathbb F_{\ell_0^2}$, so every term is in the domain
of the same $\lambda_M$. Additivity, specialization, and multiplication
by the fixed integer $L$ commute with this finite sum. If the Artin
action is written with inverses, reindexing $[a]\mapsto[a]^{-1}$
leaves it unchanged since $\chi_F$ is quadratic. Changing the
reference point within its torsor multiplies the sum only by
$\chi_F$ of that reference class.

The weights are constant on each embedding orbit used in the
following count. Indeed a unit of the definite order $R_z$ has
positive integral reduced norm with integral inverse, so its norm
is one. The conjugation law for $f$ therefore gives weight one
for its action. Along the class-group torsor the Schwartz factor
is constant because $a$ centralizes $x_0$, and the remaining
norm factor is exactly $\chi_F([a])$. Hence the finite-sum identity
just proved applies with the single phase
$\chi_{h_*}(\operatorname{Nrd}z_0)f(z_0^{-1}x_0z_0)$.


For each fixed $z$, the stabilizer of an embedding in $R_z^\times$
is $\mathcal O_F^\times=\{\pm1\}$. Each embedding orbit thus
has $|R_z^\times|/2$ elements. The denominator $|R_z^\times|$
in (O.4) leaves the factor $1/2$ per embedding class. The
vectors $x$ and $-x$ are not identified by this stabilizer:
at $\ell_0$ their residues define the two different embeddings
of the residue field, whereas conjugation by an order unit acts
trivially on that field. They already lie in the two different
$\ell_0$ orientation tuples counted above. The
contribution of the reference orientation tuple is
\[
 \frac{c_0}{2}\,\zeta_M(D)
                         \lambda_M(\overline{LP_D}),\qquad
 \zeta_M(D)=\chi_{h_*}(\operatorname{Nrd}z_0)
                               f(z_0^{-1}x_0z_0),
\]
up to the harmless quadratic reference-class factor just described.
All nonzero factors in this displayed phase are fixed-character
values or signs. Thus it is a $2$-power root of unity; its field
does not grow with $D$, $M$, or $\ell_0$.

It remains to account for the other orientations; they are not
all in the image of the one sign component of the reduction map.
At $p^e\Vert N$, use the local Atkin--Lehner normalizer. It
reverses the selected root, so multiplies the Schwartz weight by
\[
              \alpha_p(-1)=w_{p^e}\chi_{h_*}(p)^e.
\]
It multiplies the norm-weighted parametrization value by the same
sign: its reduced norm has valuation $e$, and the actual
parametrization transforms by $w_{p^e}$ after $L$ kills the fixed
cuspidal translation. The two signs multiply to one. This
normalizer has trivial components at the other primes, so it
changes no other orientation or local weight.

At $\ell_0$, a division-algebra uniformizer exchanges the two
residue embeddings. It sends the trace-zero residue $\beta$ to
$-\beta$, multiplying its weight by $(-1/\ell_0)=-1$. Its
reduced norm has valuation one and hence contributes
$\chi_{h_*}(\ell_0)=-1$. Its action on supersingular pairs is
Frobenius. Since $\lambda_M$ has Frobenius sign $+1$ and the
parametrization is defined over $\mathbf Q$, this introduces no
further sign. Again the two factors multiply to one.
The uniformizer operation used here may be described with arithmetic
or geometric Frobenius: they have the same action on the supersingular
coarse points and their parametrization values over $\mathbb F_{\ell_0^2}$.
The identity $\lambda_M(FP)=\lambda_M(P)$ from (O.2) consequently
applies to the very operation that exchanges the residue embeddings.
Its norm has valuation one, and $\chi_{h_*}$ is unramified at
$\ell_0$, so its norm-character value is exactly
$\chi_{h_*}(\ell_0)=-1$, independently of the uniformizer's unit part.
Likewise, at $p^e\Vert N$, the norm of the Atkin--Lehner normalizer
has valuation $e$ and $\chi_{h_*}$ is unramified at $p$. Its
norm-character value is exactly $\chi_{h_*}(p)^e$. These observations
fix the norm signs in both orientation changes.


It follows that every one of the $2^{\omega(N)+1}$ tuples has
the same contribution as the reference tuple. Multiplying by
their number gives precisely $c_0 2^{\omega(N)}$, as asserted.
This counts individual embeddings with their stabilizer; a
conjugate-pair formulation of the geometric correspondence must
not be used to insert or remove another factor of two.

Finally, $c_0/|R_z^\times|$ is an integer for every $z$. Thus a
change of the chosen integer lift of $P_M(z)$ by $2^M$ changes
the actual coefficient by an element of $2^M\mathcal O$.
Every preceding equality is consequently valid at the original
retained precision. No denominator depending on the class number,
the discriminant, or the stage has entered.
\paragraph{Precision in the orientation count.}
The factor $1/2$ in the orbit count does not require dividing a
congruence modulo $2^M$ by two. Indeed, $-1\in R_z^\times$ implies
$2\mid |R_z^\times|$, and the chosen integer $c_0$ is divisible by
$|R_z^\times|$. Hence $c_0/2$ is an integer. Choose lifts of all
$P_M(z)$ before summing. Each orbit of embeddings contributes
$(c_0/|R_z^\times|)(|R_z^\times|/2)P_M(z)=(c_0/2)P_M(z)$,
with its prescribed unit weight. Replacing any lift by a multiple
of $2^M$ changes this expression by a multiple of $2^M$.
The comparison with the reduced point may therefore be made at
precision $M$, separately for each orientation. Summing the
$2^{\omega(N)+1}$ orientations gives the stated multiplier
$c_0 2^{\omega(N)}$ without any loss of precision.
In particular, write $c_\theta=9+\omega(N)$, the valuation of
the fixed integer $c_0 2^{\omega(N)}$. At the fixed witness $D_w$,
(O.2) and (O.7) imply
\[
 v_2a_M(D_w)=c_\theta+c_w
       \qquad\text{whenever }M>c_\theta+c_w.
\]
The phase is a unit and the congruence error has strictly larger
valuation, so this is an equality, not merely a lower bound.
Discarding the finitely many earlier stages is sufficient. Thus the
visibility used in choosing the odd minimum retains a fixed witness
and a fixed valuation; neither the class number nor the changing
prime $\ell_0$ introduces a new loss of precision.

\end{proof}

\begin{proposition}[Visibility and the native genus index]\label{ex:app-visibility}
Use the actual odd form and unaveraged point of
Propositions~\ref{ex:app-theta-form} and~\ref{ex:app-theta-comparison}.
There is one constant $C_{\rm det}$, independent of the actual stage
and admissible $D$, for which (O.8) holds. Explicitly, for a known
simple-zero $D$,
$v_2a_M(D)\ge\min\{M,j_{\rm gen}(D)-C_{\rm det}\}$.
For every admissible $D$, the strict inequality $v_2a_M(D)<M$
implies $\operatorname{ord}_{s=1}L(E^D,s)=1$ and
$j_{\rm gen}(D)\le v_2a_M(D)+C_{\rm det}$.
If an admissible odd-sign $D$ is not a simple zero, its coefficient
vanishes modulo $2^M$. No genus index of a torsion point is defined
or used in this assertion.
\end{proposition}
\begin{proof}\leavevmode
For a known simple-zero $D$, write $P_D=m_DG+T$ in its native
rank-one free lattice. The same $L$ kills $T$.
An integral finite-field Kummer map cannot decrease
$v_2(m_D)$, so (O.7) gives the lower bound by the native index.
The rational $D$-twist subgroup has index at most two in the native
genus lattice by the calculation preceding
Proposition~\ref{ex:sec7-genus-lower}. Absorb this and the fixed
initial convention in $C_{\rm det}$. The first inequality below
holds for a known simple-zero $D$; the visibility implication holds
for every admissible $D$:
\[
 \begin{split}
 v_2a_M(D)&\ge\min\{M,j_{\rm gen}(D)-C_{\rm det}\},\\
 v_2a_M(D)<M&\ \Longrightarrow\
  \ord_{s=1}L(E^D,s)=1,\quad
       j_{\rm gen}(D)\le v_2a_M(D)+C_{\rm det}.
 \end{split}                                                    \tag{O.8}
\]
If an odd-sign $D$ is not a simple zero, the height formula gives a
torsion genus point, killed by $L$, and its coefficient is zero
modulo $2^M$. Combined with the first inequality on simple-zero
parameters, this gives the all-$D$ truncated bound used below,
without assigning a rank-one lattice index to a torsion point.
\end{proof}

\subsection{Integer divisions and actual trace operators}
\begin{proposition}[Arithmetic realization of the coefficient rules]\label{ex:app-coefficient-rules}
Use the rational full-two-torsion graph conventions of
Definitions~\ref{ex:sec7-typed-graph} and~\ref{ex:sec7-symbol}.
For an even source of Proposition~\ref{ex:app-even-phase}, use its
fixed minimum depth. For an odd source, first fix a finite operator
budget $|U|\le W$, take $R\ge2$, and choose an actual stage satisfying
the depth inequalities of Proposition~\ref{ex:app-odd-minimum}.
Then the integer divisions (O.9)--(O.11) are valid at every Fourier
index, with the stated truncation for the odd source.
The fresh coefficient-one reading is
$\overline{(t(\operatorname{Fr}_q)\theta H_U)[1]}=\mathcal A(Uq)$.
On all tested nonempty configurations, these actual symbols depend
only on their typed graph and satisfy the rules of
Lemma~\ref{ex:sec7-coefficient-rules}: a pair difference
contracts the endpoints and adds their types and every external
incidence; whole-prime omission uses the same source, phase, depth
and stage; and an isolated primary vertex of neutral type and
neutral external incidences makes the symbol zero.
Every contraction has nonempty output. Stationary character tests
are carried along in the same product rule. This statement asserts
neither values at untested odd-stage weights nor an identity at an
empty contraction output.
\end{proposition}
\begin{proof}\leavevmode
All new primes below avoid the full coefficient support $S$,
including every late sieve and the current $\ell_0$.
Here $S$ is the actual coefficient support.  The coefficients of
$g$ vanish at all indices not prime to $S$.
For the odd source this is forced at $2,N,h_*,\ell_0$ by the
local Schwartz support and the extraction at $|h_*|$.
At every additional coefficient-sieve prime it is imposed by an
exact zero-or-retain Fourier sieve, which commutes with the good
recurrence outside the enlarged support.  This is the vanishing
used in the condition $(m,S)=1$ in the unary series below.
Later stationary prime-selection tests may enrich the common
label alphabet without resieving the intrinsic odd source;
such tests do not enlarge this coefficient-support condition.
Use the minimum-depth inequalities proved in Appendix~\ref{app:endpoints}.
The general bookkeeping assigns weight one to a split prime of
$\Q(E[2])$ and one-half to a prime with transposition Frobenius, called
simple. In our actual full-two-torsion family every prime is split,
$w(U)=|U|$, and the depth $\delta$ is an integer. Thus only matching
parity, meaning $e_U=0$, is needed here; the displayed general formula
also records where the unused simple-prime term would enter.
For a finite set $U$, put
\[
 n_U=\lceil\delta+w(U)\rceil,\quad e_U=n_U-\delta-w(U),\qquad
 B_U=2^{-n_U}\prod_{p\in U}(U_p-\eta(p)V_p)g,
\]
where $U_p(\sum a(n)q^n)=\sum a(pn)q^n$ and $V_pg(\tau)=g(p\tau)$.
These are operations on the actual form $g$. We verify the level
at a new odd prime using the classical theta multiplier. Write the
old level as $4A$, with $p\nmid4A$, and put
$\beta=\left(\begin{smallmatrix}p&0\\0&1\end{smallmatrix}\right)$.
For $\gamma=\left(\begin{smallmatrix}a&b\\c&d\end{smallmatrix}\right)
\in\Gamma_0(4Ap)$ one has
\[
 \beta\gamma\beta^{-1}=
       \begin{pmatrix}a&pb\\c/p&d\end{pmatrix}\in\Gamma_0(4A).
\]
The square-root factor in the theta multiplier is unchanged on
evaluating the latter matrix at $p\tau$. The quotient of its
arithmetic factor by that of $\gamma$ is
$(c/p\,/d)/(c/d)=(p/d)$. Cubing the multiplier leaves this
quadratic factor unchanged. Hence $V_pg$ has level $4Ap$ and
character $\psi(p/\cdot)$.

At this level the normalized sum
\[
 U_{p^2}h(\tau)=\frac1{p^2}\sum_{r\bmod p^2}
                             h((\tau+r)/p^2)
\]
preserves the same weight and character. Indeed, put
$\alpha_r=\left(\begin{smallmatrix}1&r\\0&p^2\end{smallmatrix}\right)$.
For $\gamma\in\Gamma_0(4Ap)$ choose
$r'\equiv(b+rd)/(a+rc)\pmod{p^2}$. Then
\[
 \alpha_r\gamma=\gamma'_r\alpha_{r'},\qquad
 \gamma'_r=
 \begin{pmatrix}
  a+rc&(b+rd-(a+rc)r')/p^2\\
  p^2c&d-cr'
 \end{pmatrix}\in\Gamma_0(4Ap).
\]
The map $r\mapsto r'$ is a permutation, since $a+rc$ is a unit
modulo $p$ and the fractional linear map has determinant one.
Writing $d'=d-cr'$, its nebentypus factor is unchanged since
$d'\equiv d\pmod{4Ap}$. Its theta arithmetic factor is unchanged
as well: $d'\equiv d\pmod c$, $d'\equiv d\pmod4$, and
$(p^2c/d')=(c/d')=(c/d)$. These are the usual Kronecker-symbol
identities for the theta multiplier on $\Gamma_0(4)$.
Its square-root factor, evaluated at $(\tau+r')/p^2$, is
$(c\tau+d)^{1/2}$. Thus the sum transforms with the same
theta multiplier cubed and the same character. The formulas also
hold for the triangular matrices, and holomorphy at the cusps is
preserved by the finite rational translates.
On Fourier expansions $U_{p^2}V_p=U_p$. Consequently both terms
of $U_p-\eta(p)V_p$ have level exponent at most one at $p$ and
character $\psi(p/\cdot)$. Iteration introduces no higher exponent
at previously added primes.

The same computation applies to the unary corrections. Start with
the odd-character unary theta series at the old support, including
its fixed coprimality sieve. Its $J$-series is its transform by
$V_J$, for squarefree $J\subset U$. A prime $p\mid J$ therefore
has exponent at most one and the extra character $(p/\cdot)$;
a prime $p\nmid J$ is absent from its original level. Multiplication
by the ordinary theta series introduces no odd level prime. This
proves the precise level and oldform alternatives used in (O.12).


Write an index as $Jrm^2$ with $J\subset U$ and $r$ squarefree,
prime to $U$. Products of prime sets are understood here.
The originating squarefree parameter in the coefficient of $g$ is
$D=(U\setminus J)r$, not the squarefree part $Jr$ of the tested index.
Remove the phase $\eta(p)^s$ in (O.6) and put
\[
 b_{-1}=0,\quad b_0=1,\quad b_1=a_p-\kappa,\quad
 b_{s+1}=a_pb_s-pb_{s-1},\quad
 \kappa=(\psi/\eta)(p)(-D/p).
\]
For $p\in U\setminus J$ the factor $b_s-b_{s-1}$ is odd;
for $p\in J$ the factor $b_{s+1}-b_s$ is even.
Consequently, for the even source,
\[
 v_2B_U[Jrm^2]\ge w(r)
              +\tfrac12\#\{p\in J:p\text{ is simple}\}-e_U.      \tag{O.9}
\]
The right side is an integer on the allowed filter and is
nonnegative. In our full-two-torsion application all primes are
split and $w=\omega$.
For the odd source, fix an integer $R\ge1$ and a weight bound $W$
first. At one sufficiently advanced actual stage, simultaneously for
all $U$ with $w(U)\le W$ and every Fourier index, (O.9) holds with
its right side replaced by its minimum with $R$. Thus the untruncated
bound on a prescribed bounded range of squarefree weights follows by
choosing $R$ and then the stage. For the unary-removal and two-trace
tests below, choose $R\ge2$ after fixing the operator budget.
For a split prime $a_p\equiv p+1\pmod4$. The recurrence gives the
same residue on each divided shell,
\[
        (b_{s+1}-b_s)/2\equiv
     c_p(D):=\{a_p-(\psi/\eta)(p)(-D/p)-1\}/2\pmod2.              \tag{O.10}
\]

\paragraph{The two integral shell calculations.}
Here every prime is split in the rational two-torsion field, so
$a_p\equiv p+1\pmod4$. For $p\mid D$ one has $\kappa=0$.
Reducing the recurrence modulo two gives $b_{s+1}=b_{s-1}$,
with $b_0=1$ and $b_1=0$. Thus $b_s-b_{s-1}$ is odd for
all $s\ge0$, including $s=0$ with $b_{-1}=0$.
For $p\nmid D$, one has $\kappa\in\{1,-1\}$, so every $b_s$
is odd. Write $d_s=b_{s+1}-b_s$. The recurrence modulo four gives
$d_s\equiv p d_{s-1}\pmod4$ for $s\ge1$.
Each $d_s$ is even, and multiplication of an even residue modulo
four by an odd number fixes that residue. Consequently
\[
 d_s\equiv d_0=a_p-\kappa-1\pmod4\qquad(s\ge0).
\]
This proves (O.10) without dividing a residue-field identity by two.
For an index $Jrm^2$, each $p\in J$ therefore supplies an actual
factor two, whereas the factors for $p\in U\setminus J$ are integral
units modulo two. The squarefree input has
$\omega(D)=|U|-|J|+\omega(r)$. Its depth bound plus these $|J|$
factors gives
$\delta+\omega(D)+|J|=n_U+\omega(r)$ before division by
$2^{n_U}$. All removed character factors are units.
For the odd source the same argument is a congruence modulo
$2^M\mathcal O_{\rm coef}$ before that division. It therefore
retains precision $M-n_U$ afterward; the subsequent division
in (O.11) retains precision $M-n_U-1$.
Choose $i\in\{0,1\}$ with $\eta\chi_{-4}^i$ odd and let
\[
 \Theta_J=\sum_{\substack{m>0\\(m,S)=1}}
                      \eta(m)\chi_{-4}(m)^i m q^{Jm^2}.
\]
The Poisson calculation in Section~\ref{v15:unary-poisson} gives this genuine unary weight-$3/2$ form.
Its scalar character is $\eta\chi_{-4}^i(-J/\cdot)$ and every
$p\mid J$ enters the level at exponent one.
In matching parity the shell recurrence gives
$B_U[Jm^2]-B_U[J]\eta(m)\in2\mathcal O_{\rm coef}$.
Thus
\[
                 H_U=\frac{B_U-\sum_JB_U[J]\Theta_J}{2}          \tag{O.11}
\]
is integral; in opposite parity take $H_U=B_U$.
This is division by the integer two, not reduction modulo the
maximal ideal of a ramified coefficient ring. For odd-stage
forms, the coarse bound and bounded-weight stabilization proved
below pay every fixed denominator before choosing the stage.

Write $\mathfrak m_{\rm coef}$ for the maximal ideal of the chosen
coefficient valuation ring $\mathcal O_{\rm coef}$. In the following
trace argument a reduced expansion means coefficientwise reduction
modulo $\mathfrak m_{\rm coef}$. The preceding divisions by integer
powers of two have already been justified in $2^a\mathcal O_{\rm coef}$;
reduction modulo the maximal ideal alone would not justify them.
No saturation of the integral Fourier module is assumed, and no
identification of its abstract quotient with all reduced expansions
is needed.
Multiply (O.11) and the undivided corrections by
$\theta=1+2\sum_{m>0}q^{m^2}$. They are actual weight-two forms.
For a fixed stage and $U$, choose a finite complete coefficient-field
extension containing all constituents and retain their old
multiplicities and diamond translates in the finite-dimensional
characteristic-zero space. Use modular Galois representations on
cuspidal constituents and the two character summands on Eisenstein
constituents. Holomorphy alone does not exclude the latter.
Diamond translates remain integral:
their correction coefficients have the form
\[
 \tfrac12 B_U[J]\{\psi_U(u)-\eta(u)\chi_{-4}(u)^i(-J/u)\},
\]
whose character ratio is quadratic. Good Hecke operations preserve
integrality by
$(T_qf)[n]=f[qn]+q(\langle q\rangle f)[n/q]$.
Finitely many Fourier coordinates inject this space into a finite
coordinate space. The module generated by these translates is
therefore finite over the coefficient DVR, and closed.
Chebotarev and continuity make scalar trace and determinant
operators limits of good Hecke and diamond operations, so they
preserve this module. No integral eigenspace projector is used.

\paragraph{Reading an added vertex.}
For the actual full-two-torsion family, all primes are split,
$\delta\in\mathbb Z$, and $n_U=\delta+|U|$.  Put $F_U=\theta H_U$
and let $q\notin S\cup U$ be a fresh prime, where $S$ includes
every late sieve and the current moving theta prime.
Write $u=\prod_{p\in U}p$ and let $\mathcal A(D)$ denote the
normalized valuation symbol on the tested range: it is one exactly when
$v_2a(D)=\delta+\omega(D)$ on the filter. For an odd source, that
range has a fixed weight bound chosen before the valid actual stage,
as specified before (F.7). No finite-stage value at larger weights is
asserted; the compatible limit tables used later are a separate object.
Expanding the operators, every term containing a $V_p$ has zero
coefficient at $q$, and no unary index $Jm^2$ with $J\subset U$
equals $q$.  Therefore
\[
 B_U[q]=2^{-n_U}a(uq),\qquad
 H_U[q]=2^{-n_U-1}a(uq).
\]
At coefficient one the diamond term in the good Hecke formula
vanishes.  Since $F_U\equiv H_U\pmod{2\mathcal O_{\rm coef}}$,
\[
 \overline{(t(\operatorname{Fr}_q)F_U)[1]}
   =\overline{H_U[q]}
   =\overline{\,2^{-\delta-|U|-1}a(uq)\,}
   =\mathcal A(uq)\in\mathbb F_2.
\]
Here residues lie in the coefficient residue field, with
$\mathbb F_2$ its prime subfield.  Every permitted phase is a
two-power root of unity and reduces to one.  For an odd source,
fix the finite operator budget $|U|\le W$ and $R\ge2$ first;
the stage choice in (O.9) may be taken with
$M_\theta\ge\max(0,\delta+W)+R$.
Thus $M_\theta-n_U-1\ge1$, so the normalized finite-stage error
is still divisible by the integer two.  The same reading holds
with value zero off the fixed or moving coefficient support.

The diamond correction above is a sum of integral multiples of
$\theta\Theta_J$; its fresh coefficient $q$ is zero after reduction.
It introduces neither an extra phase nor a denominator in this
reading.  All readings and subsequent omissions use this same
$g,\eta,\psi,\delta$ and actual stage; only the set $U$ changes.

For any fixed Fourier index $n$, Chebotarev approximation lets one
test a finite product of trace operators using distinct fresh primes
$q_1,\ldots,q_a$ avoiding $n$ and the level. The good Hecke formula
then reads the coefficient at $nq_1\cdots q_a$ of the original vector.
Different indices may use different fresh primes; the conclusion is
an identity of coefficientwise reduced expansions.

The determinant identity uses the actual diamond formula, not merely
the reduction of its eigenvalues. For a good prime $r$,
$d(\operatorname{Fr}_r)=r\langle r\rangle$. The formula above shows
that $(d(\operatorname{Fr}_r)-1)F_U$ reduces to a sum of integral
multiples of the undivided corrections $\theta B_U[J]\Theta_J$:
the scalar $r\psi_U(r)-1$ lies in $\mathfrak m_{\rm coef}$.
A fresh trace kills every such correction. Indeed, to test at $n$
choose $q\nmid nJ$; then $nq$ cannot have squarefree part $J$, and
$\theta\equiv1\pmod{\mathfrak m_{\rm coef}}$. This remains true
after further fresh traces. Continuity and the closed integral
module therefore give, on the reduced trace orbit of $F_U$,
\[
 d(u)t(v)=t(v),\qquad
 t(u)t(v)=t(uv)+t(uv^{-1}),\qquad t(1)=0.
\]
The middle equality follows from the characteristic-zero matrix
identity
$\operatorname{tr}(A)\operatorname{tr}(B)
=\operatorname{tr}(AB)+\det(B)\operatorname{tr}(AB^{-1})$,
using the determinant identity just proved; the last follows from
$\operatorname{tr}(1)=2$.

Two fresh primes outside $S\cup U$ contribute two to the squarefree
weight in (O.9). Thus $B_U[nq_1q_2]$ is divisible by four, no unary
correction occurs there, and $H_U[nq_1q_2]$ remains divisible by two.
Here $R\ge2$ and the actual-stage budget were fixed before the test.
It follows that $t(g)t(h)F_U=0$ after reduction. Applying the middle
identity to $t(gh)t(h)$ gives $t(gh^2)F_U=t(g)F_U$.
Invariance under all squares, and continuity, give factorization
through the elementary quadratic quotient unramified outside
$S\cup U$. Over $\Q$ this quotient is finite: its squareclasses
are generated by $-1$, $2$, and the primes in that support.
This does not assert abelianity of the characteristic-zero
Galois representations.

At $p\in U$, first keep the summands of $F_U$ separately.
The block $\theta B_U$ and each unary block with $p\in J$ have
level exponent at most one and primitive quadratic character at $p$.
The unary blocks with $p\notin J$ are degeneracies of forms at a
level omitting $p$. This distinction is preserved when retaining all
old multiplicities. In particular, the latter blocks cannot acquire a
Steinberg constituent merely by being viewed at a larger level.

Enlarge the characteristic-zero coefficient field to realize all the
cuspidal newform constituents and the Eisenstein character pairs of
these blocks. On a constituent and all its degeneracies, define $t(g)$
and $d(g)$ to be the scalar trace and determinant of its two-dimensional
Galois representation, acting as the identity on its multiplicity
space. Distinct characteristic-zero constituents with the same good
eigenvalues have the same trace and determinant by Chebotarev, so
these definitions agree on coinciding summands. In particular they
give well-defined operators on their sum. No decomposition modulo
two, and no residual irreducibility assertion, enters this construction.

For a cuspidal eigenform $f$, let $\Pi_f$ denote its unitary
adelic normalization: at a good prime $q$ its Satake parameters
$\alpha_q,\beta_q$ satisfy
$a_q=\sqrt q(\alpha_q+\beta_q)$ and
$\alpha_q\beta_q=\chi(q)$.
Apply \cite[\S\S0.3--0.7, Theorem~A, pp.~409--411]{carayol}
to $\Pi_f^\vee$, and set $\rho_f=\sigma(\Pi_f^\vee)^\vee$.
Here $F=\Q$, the real weight is $2$, and the coefficient place
lies above $2$. The degree of $F$ is odd, so the additional
finite discrete-series-place hypothesis for even degree does
not apply. The local prime $p$ is odd.
Carayol's reciprocity sends geometric Frobenius to a uniformizer,
and his $\omega_p$ is $|\cdot|_p$.
His normalized-induction formula is
\[
 \sigma(\operatorname{Ind}(\xi_1,\xi_2))
 =\xi_1^{-1}\omega_p^{-1/2}\oplus
   \xi_2^{-1}\omega_p^{-1/2}.
\]
Thus if $\Pi_{f,p}=\operatorname{Ind}(\mu_1,\mu_2)$,
$\rho_f$ on arithmetic Frobenius has eigenvalues
$\sqrt p\,\mu_1(p),\sqrt p\,\mu_2(p)$ whenever these
characters are unramified. In particular its good polynomial
is $X^2-a_qX+q\chi(q)$. Both the automorphic contragredient and
the dual on the Galois side are part of this convention.

For a block with ramified character, its local conductor is
exactly one: it is at most the level exponent and at least the
conductor of its determinant.
Casselman's newvector theorem and conductor calculation
\cite[Theorem~1, pp.~302--306]{casselman} show that a special
representation has conductor one only for an unramified twist;
a ramified twist has conductor at least two. A supercuspidal
representation has conductor at least two. Consequently
$\Pi_{f,p}=\operatorname{Ind}(\mu_1,\mu_2)$ with $\mu_1$
unramified and $\mu_2$ of conductor one, after interchanging
the characters if necessary. Its newvector space is a line.
On inertia the ramified character is the prescribed quadratic
character.

We check the Fourier normalization of the operator on this line.
Use normalized induction, so that
\[
 v\left(\begin{pmatrix}a&x\\0&d\end{pmatrix}g\right)
 =\mu_1(a)\mu_2(d)|a/d|_p^{1/2}v(g).
\]
Put $K_0(p)=\{g\in\mathrm{GL}_2(\Z_p):g_{21}\in p\Z_p\}$.
Use its character $k\mapsto\mu_2(k_{22})$.
This is the lower-right-character convention. Casselman's
upper-left-character convention is carried to it by
$h_p=\left(\begin{smallmatrix}0&1\\-p&0\end{smallmatrix}\right)$,
which normalizes $K_0(p)$ and interchanges its diagonal entries
under conjugation. We label the inducing characters so that
$\mu_1$ is unramified, using the interchange isomorphism for
irreducible principal series.
Normalize Haar measure by $\operatorname{vol}(K_0(p))=1$.
The character-weighted double-coset operator for
$K_0(p)\operatorname{diag}(p,1)K_0(p)$, with value one at
$\operatorname{diag}(p,1)$, is
\[
 H_p=\sum_{b\in\Z/p\Z}
       \Pi_{f,p}\left(\begin{pmatrix}p&b\\0&1\end{pmatrix}\right).
\]
Take the newvector supported on $B K_0(p)$ with $v(1)=1$.
Evaluating at $1$ gives
\[
 (H_pv)(1)=p\mu_1(p)p^{-1/2}=\sqrt p\,\mu_1(p).
\]
This is the Fourier operator $U_p^{(2)}$. Indeed, in the
Whittaker model use an additive character of conductor $\Z_p$.
Writing $a(n)=\operatorname{diag}(n,1)$ and
$n(x)=\left(\begin{smallmatrix}1&x\\0&1\end{smallmatrix}\right)$,
\[
 a(n)\begin{pmatrix}p&b\\0&1\end{pmatrix}=n(nb)a(np)
 \quad(n\in\Z_p),\qquad
 (H_pW)(a(n))=pW(a(np)).
\]
Normalize the archimedean Whittaker function by
$W_\infty(a(y))=y\exp(-2\pi y)$ for $y>0$.
Comparison of its positive-$n$ term with $y f(iy)$ gives the
finite Whittaker coefficient $a_n/n$. The displayed equality
therefore sends $a_n$ to $a_{pn}$. Hence
$U_p^{(2)}=\sqrt p\,\mu_1(p)$ on the new line, exactly the
arithmetic-Frobenius eigenvalue on the unramified Galois line.
This scalar is algebraic; the square roots only express the
unitary normalization. Old multiplicities at primes other than
$p$ carry the same scalar. Casselman's oldspace dimension formula
\cite[Corollary, p.~306, and Theorem~4, p.~313]{casselman}
does not introduce an additional multiplicity at this conductor-one
prime. For Eisenstein constituents, the two Galois characters and
the divisor-sum Fourier coefficients give the same equality directly;
Carayol's cuspidal theorem is not applied to them.

Choose an inertia element $\tau$ on which the quadratic character
is $-1$, and a Frobenius lift $\phi$. On a ramified block, a basis
of the two character lines gives
\[
 \rho(\tau)=\begin{pmatrix}1&0\\0&-1\end{pmatrix},\qquad
 \rho(\phi)=\begin{pmatrix}\lambda&0\\0&\mu\end{pmatrix},
 \qquad U_p^{(2)}=\lambda,\quad d(\phi)=\lambda\mu.
\]
For $\rho(g)=\left(\begin{smallmatrix}a&b\\c&e\end{smallmatrix}\right)$,
the two inertia differences are
\[
 t(g\tau)-t(g)=-2e,\qquad
 t(g\tau\phi)-t(g\phi)=-2e\mu.
\]
It follows in characteristic zero that
\[
 U_p^{(2)}\{t(g\tau\phi)-t(g\phi)\}
       =d(\phi)\{t(g\tau)-t(g)\}.                       \tag{O.12}
\]
On a block coming from level prime to $p$, the representation is
unramified and both differences vanish. This proves the identity
on its entire old multiplicity space even if $U_p^{(2)}$ is not
scalar or semisimple on that space. It therefore proves (O.12)
on the actual sum containing $F_U$.

These are identities on characteristic-zero forms before any
reduction. To relate them to the integral trace module above, choose
finitely many Fourier coordinates that inject its ambient space.
Every generator of that module has integral coordinates; hence the
module is a submodule of a finite free module over the complete
coefficient DVR. It is finite and closed. Simultaneous Chebotarev
approximation in the finitely many Galois representations makes
$t(g)$ and $d(g)$ limits of good Hecke and diamond operators on this
space. They preserve this integral module, so the identity can be
evaluated on the integral $F_U$ and then reduced. This does not
require integral projectors onto constituents. In particular no
division by the two in the displayed inertia differences is used.


Apply (O.12) at coefficient one.  Put $V=U\setminus\{p\}$,
$v=\prod_{r\in V}r$, so $n_U=n_V+1$.  For a fresh prime $r$ set
\[
 C_p(vr)=\frac{a_p-(\psi/\eta)(p)(-vr/p)-1}{2}\in\Z,
 \qquad c_p(vr)=C_p(vr)\bmod2.
\]
The two surviving operator terms at $pr$ and the good recurrence give
\[
 H_U[pr]\equiv
   \eta(p)C_p(vr)\,2^{-n_V-1}a(vr)\pmod{2\mathcal O_{\rm coef}}.
\]
For the even source this is an equality before reduction.  For the
odd source the error is in $2^{M_\theta-n_U-1}\mathcal O_{\rm coef}$,
which is contained in $2\mathcal O_{\rm coef}$ at our chosen stage.
Choose fresh $q,q_\tau,r_0,r_1$ representing respectively
$g,g\tau,g\phi,g\tau\phi$ in one sufficiently precise joint quotient
for the $U$ and $V$ forms and all retained characters.
The omitted form uses the same $g,\eta,\psi,\delta$ and stage,
and its level omits the entire rational prime $p$.
Thus the two $V$ symbols at $vr_0,vr_1$ agree.  Their Legendre
signs at $p$ are opposite, so $c_p(vr_0)+c_p(vr_1)=1$ in $\F_2$.
The coefficient-one reading and the reduced trace identity now give
\[
             \mathcal A_U(Uq_\tau)+\mathcal A_U(Uq)=\mathcal A_V(Vr_0)\qquad\text{in }\F_2.
\]
Here the subscript records the operator set, not a new coefficient
source or minimum. The factor $\eta(p)$ reduces to one; no factor
two is cancelled in (O.12). For every retained quadratic character
$\chi$ unramified at $p$, the representative of the omitted vertex
satisfies
\[
 \chi(g\phi)=\chi(g)\chi(\phi)
            =\chi(\operatorname{Fr}_q)\chi(\operatorname{Fr}_p).
\]
For a fixed-support character this is the product of the two vertex
labels. For the character of a third vertex $r\in V$, it adds the
two bits $(r/q)$ and $(r/p)$ in $\F_2$. Reciprocity then gives the
reverse bits from the same endpoint types. All additional stationary
tests were included in the joint quotient, so they obey this formula
as well. This proves the full product-label omission rule, including
every external incidence.

A fresh vertex with trivial label in the elementary extension for
$S\cup U$ has the same reduced reading as $t(1)$, hence gives zero.
An isolated primary vertex has precisely this fixed label and these
zero incidence bits. A contracted cluster of primary vertices has
the same property when its external incidences are neutral. The
result is still a nonempty configuration, including the one-vertex
case; no contraction deletes the last vertex.

Finally the fixed-$U$ factorization gives dependence on the entire
finite typed graph. For two disjoint prime realizations of the same
graph, choose a fresh bridge copy with the prescribed types and
matching incidences to both realizations. CRT and Dirichlet provide
the primes successively, and reciprocity determines their reverse
bits. Replace the vertices of the first copy by bridge vertices one
at a time, then the bridge vertices by those of the second copy.
At each step the other vertices form the same fixed $U$, and the
changed vertex has the same elementary Frobenius data for $S\cup U$;
its symbol therefore agrees. If the original realizations overlap,
first replace one by a fresh copy avoiding both original prime sets,
with incidences matching that first realization. A vertex and its
own replacement are never tested together, so their mutual bit can
be chosen freely. The resulting two disjoint realizations admit the
preceding bridge argument without conflicting prescriptions.
Only graphs of the original bounded size are evaluated; extra bridge
incidences are prime-selection conditions. Thus these are rules for
the actual arithmetic graph symbols, at the already chosen valid
stage. The simultaneous finite graph budget and stabilization are
treated below.
\end{proof}

% End 08-odd-source.tex


% Begin 09-normalization-endpoints.tex
\section{Nonvanishing, normalization and the final family}\label{app:endpoints}
The uses of the coefficient and address propositions have the following
order. First apply only their even-source case, whose depth is already
fixed by Proposition~\ref{ex:app-even-phase}. This gives the rough
companion and the genus lower bound. These in turn establish the odd
minimum and its finite-stage inequalities. Only then apply the odd
case of Proposition~\ref{ex:app-coefficient-rules} and the two-symbol
address construction. In particular the proof of the rough companion
does not use the later odd minimum.
We give the choice order underlying (3.3)--(3.9). This also separates
the two valuation arguments: the even upper transfer uses the common
Kato scalar, while the final contradiction uses the common Heegner
coordinate of Section~\ref{sec:limits}.

\subsection{An absolute bound of five new primes}
The theorem of Hoffstein--Luo \cite[p.~436]{hl} supplies fundamental
discriminants with nonzero central value and at most four prime
factors. Its finite splitting set is arbitrary. The bound four is
absolute; no uniform size threshold for those discriminants is needed.

\begin{proposition}[Nonvanishing with five primes and finite local prescriptions]\label{ex:app-five-prime}
Let $A/\Q$ be a fixed modular elliptic curve for this application,
let $S$ contain $2,3$ and all its bad primes, and let $B$ be a finite
prime exclusion set. In every nonempty signed unit-squareclass filter
at $S$ of even root number, there is a signed odd squarefree parameter
$u$, prime to $S\cup B$, with $L(A^u,1)\ne0$ and $\omega(u)\le5$.
Only the intermediate parameter supplied by Hoffstein--Luo is required
to be fundamental; the final field discriminant may be $4u$.
The number five is independent of the prescribed finite set. No
uniform bound on the sizes of the selected primes is asserted.
\end{proposition}
\begin{proof}\leavevmode
Let $A/\Q$ be any fixed modular elliptic curve for this application,
$S$ contain two, three and all its bad primes, and $B$ be any finite
exclusion set. Fix a nonempty signed unit filter of even root number.
By CRT and Dirichlet choose one signed prime $q_0$ in its classes,
outside $S\cup B$. Apply the theorem to the actual curve $A^{q_0}$,
putting
\[
                 S'=S\cup B\cup\{|q_0|\}\cup
                               \{p:p\mid N_{A^{q_0}}\}
\]
in its splitting set. The theorem gives a fundamental $d$ with
$\chi_d(p)=1$ for all $p\in S'$ and at most four prime factors.
The functional equation forces $d>0$, since $A^{q_0}$ has even
root number and every conductor prime splits. Splitting at two
gives $d\equiv1\pmod8$. Thus $q_0d$ has all original unit
classes, is prime to $B$, is squarefree, and has at most five
prime factors. If $q_0d\equiv3$ or $7\pmod8$, its field
discriminant is $4q_0d$; only $d$ was required to be fundamental
in the theorem. This is compatible with the coefficient
construction (E.2).
\end{proof}

Here $h$ is restricted to the prepared positive odd-sign unit
filter containing one, as in (3.3), so $h\equiv1\pmod8$.
The multiset of dummy labels and its size are fixed before $h$;
its realizing primes are chosen after $h$ and avoid its support.

For the rough companion the curve itself varies:
write $a$ for the bounded dummy product and take $A=E^{ha}$.
Choose $q_0$ so $k_0=aq_0$ is negative, $1$ modulo eight and
a square at every prime of $2Nh$, with the required exclusions.
Then $k_0$ is odd fundamental, prime to $N_{E^h}$, and
\[
          w(E^{hk_0})=w(E^h)\chi_{k_0}(-N_{E^h})=+1.
\]
This proves the sign directly for the actual application.
Apply Hoffstein--Luo to $E^{hk_0}$ with every conductor,
old, $h$, dummy and forbidden prime in its splitting set.
The resulting $k'=k_0d$ adds at most five new primes.
The constants and sizes in the analytic theorem may depend on
this varying curve and the whole finite set; the bound five does not.

\subsection{The even upper transfer}
\begin{proposition}[A uniformly bounded partitionable companion]\label{ex:app-rough-companion}
Fix $E$ and the finitely many prepared even sources, local types and
unit seeds used for (3.3). There are constants $A,C_E$, fixed before
$h$ and before any additional finite prime exclusion set, with the
following property. Let $h$ be positive squarefree in the prepared
odd-sign unit filter containing $1$, prime to its fixed support; in
particular $h\equiv1\pmod8$. There is a negative odd fundamental
discriminant $k'$, prime to $2Nh$ and the additional exclusions,
square at every place dividing $2Nh$, for which $hk'$ is partitionable
in the sense of Definition~\ref{ex:sec7-seed-partition},
$L_*(hk')\ne0$, and
\[
 \omega(k')\le A,\qquad
 v_2L_*(hk')\le2\omega(h)+2\omega(k')+s(hk')+C_E.
\]
Here $A$ is the fixed dummy-label count plus five. Dummy types and
their number are fixed first, but their distinct prime realizations
are selected after $h$ and the exclusions. The proof uses the
one-symbol stationary-address conclusion of
Lemma~\ref{ex:sec7-one-symbol-address}, whose proof is given in
Proposition~\ref{ex:app-simultaneous-address}.
\end{proposition}
\begin{proof}\leavevmode
For one fixed even source write (E.11) as
\[
             v_2L_*(\sigma D)=2v_2a(D)+\kappa_e.
\]
The lower bound (E.9) and a nonzero reference coefficient imply
that the nonempty set of integers
$v_2a(D)-\omega(D)$ has a finite attained minimum $\delta_e$.
A unit of the normalized symbol means precisely equality to this
minimum. It therefore gives
\[
 L_*(\sigma D)\ne0,\qquad
 v_2L_*(\sigma D)=2\omega(D)+C_{e,\rm unit},
 \quad C_{e,\rm unit}=2\delta_e+\kappa_e.                         \tag{F.1}
\]
This constant belongs to the fixed source, not to a curve later
introduced solely to obtain nonvanishing.

For the negative types used in the rough companion, the unsigned
index is $7$ modulo eight. Its total $(-1/D)$ label is nonzero, so
every unit seed is nonempty. A seed of least positive size has no
dependence on internal edge bits: a nonzero edge derivative
would be a smaller nonempty unit. Its labels therefore suffice.
If a list with the same total label contains the first $r-1$
seed anchors and a nonempty remainder, contract that remainder
into the last anchor. Its total label is forced to be the
remaining seed label. This proves partitionability regardless
of how many old $h$-primes are in the remainder.

For a partitionable $d$ already known to have Selmer corank zero,
the one-symbol address construction keeps every prime of $d$
stationary and gives arbitrarily large cubes $d_x$, $d_0=d$,
with unit symbol at each $x\ne0$. The scalar in (E.5) is common
to that cube. By (E.8) and (F.1),
\[
      v_2u_x^{\rm ev}(0)\le B_e:=C_{e,\rm unit}+C_{\rm loc}
                                      \quad(x\ne0).
\]
Choose $M-c_E>B_e$. The constant coefficient of
$2^{c_E}u^{\rm ev}\in\Lambda[G_b]$ has sign-value bits of total
Boolean degree at most $2^{M-1}-1$. For a larger cube dimension,
the Boolean comparison in Section~\ref{sec:limits} gives
$x\ne0$ with
$v_2(u_x^{\rm ev}(0)-u_0^{\rm ev}(0))\ge M-c_E$.
If the base were zero or had valuation greater than $B_e$, this
would contradict the upper bound at $x$. Therefore
\[
       v_2L_*(d)-2\omega(d)-s(d)
                       \le C_{e,\rm unit}+2C_{\rm loc}.          \tag{F.2}
\]
The required dimension may depend on $d$; this bound does not.
It precedes the odd minimum and uses no Heegner transfer.

Prepare these sources and nonempty seeds for every total type at
the fixed support. There are finitely many. A bounded supply of
dummy anchors and reserves works for all of them. For each $h$
realize that supply by distinct primes avoiding the entire fixed
support and $h$. Put every dummy and every other already fixed
forbidden prime in the exclusion set of the preceding
nonvanishing construction. It gives $k'$ with
$\omega(k')\le A$, where $A$ is the fixed dummy count plus five.
For the total type that actually occurs, $hk'$ is partitionable.
Nonvanishing gives corank zero by the forward theorem, so (F.2)
applies. Taking the maximum over the finitely many prepared
constants proves
\[
 v_2L_*(hk')\le2\omega(h)+2\omega(k')+s(hk')+C_E.                 \tag{F.3}
\]
No new even minimum is chosen after $h$ or its exclusion set.
Those excluded primes are not added as relaxed places in a
Selmer complex.
\end{proof}

\subsection{The odd minimum is chosen after its lower bound}
\begin{proposition}[The odd minimum and every finite operator budget]\label{ex:app-odd-minimum}
Fix the original odd filter, the detecting sequence of
Proposition~\ref{ex:app-detector}, and one nonprincipal ultrafilter
on its stages. The arithmetic lower bound (3.5), the native comparisons
(3.4), and Proposition~\ref{ex:app-rough-companion} give (F.4)--(F.5)
before any odd minimum is chosen. The bounded-weight visible depths
therefore have an attained finite integer minimum $\delta_o$.
For every fixed weight cutoff, the simultaneous inequality (F.6)
holds on a set in this same ultrafilter. For every finite operator
budget and $R\ge2$, one may choose a stage where all the integer
divisions and residual errors required in
Proposition~\ref{ex:app-coefficient-rules} are valid at every Fourier
index.
On a fixed squarefree-output range the unit symbol defined before
(F.7) is exhaustive, and its value $1$ implies (F.7), with the one
constant $C_o=\delta_o+C_{\rm det}$. No minimum is changed on
omitting a vertex or adding stationary tests, and no single finite
stage is claimed to serve every unbounded weight range.
\end{proposition}
\begin{proof}\leavevmode
The arithmetic result (3.5) and the two native-lattice height
comparisons (3.4) give, using (F.3),
\[
 2j_{\rm gen}(h)\ge2\omega(h)+s(h)-C_{\rm gen},\qquad
 C_{\rm gen}=2A+C_R+C_E-v_2L_*(h_*)+C_H.                         \tag{F.4}
\]
The variable $s(hk')$ cancels. This estimate precedes the theta
prime, precision, odd minimum, final field and binary dimension.

Combining (F.4) with (O.8) gives on all admissible $D$
\[
            v_2a_M(D)\ge\min\{M,\omega(D)-C_{\rm coarse}\},
 \qquad C_{\rm coarse}=C_{\rm gen}/2+C_{\rm det}.                 \tag{F.5}
\]
At visible coefficients (O.8) first gives a simple zero, allowing
(F.4); at the other odd-sign parameters the coefficient is zero
modulo $2^M$. This avoids a circular use of the future minimum.

Fix one nonprincipal ultrafilter on the actual stages with
$M=M_\theta\to\infty$. Consider the integers $r$ for which a
bounded-weight sequence of visible $D_M$ has
$v_2a_M(D_M)-\omega(D_M)=r$ on a set in the ultrafilter.
The fixed witness $D_w$ and (O.2),(O.7) make this set nonempty.
Equation (F.5) bounds it below. Let its attained minimum be
$\delta_o$. For every fixed $W$, on a set in the same ultrafilter,
\[
 v_2a_M(D)\ge\omega(D)+\delta_o
             \quad\hbox{for every visible }D,\ \omega(D)\le W.
 \tag{F.6}
\]
Otherwise choose a violating $D_M$ at each bad stage. One of the
finitely many integers between $-C_{\rm coarse}$ and
$\delta_o-1$ would occur on an ultrafilter-large set, a contradiction.
Invisible coefficients satisfy the required truncated inequality
once $M>W+\delta_o$.

For a fixed finite operator budget $|U|\le W_U$, put
$N_0=\max(0,\delta_o+W_U)$ and fix $R\ge2$.
Choose an integer $T\ge W_U+1$ with
$T-C_{\rm coarse}\ge N_0+R$, then choose a stage in the
ultrafilter-large set where (F.6) holds through weight $T$, with
$M\ge N_0+R$ and $M>T+\delta_o$.
For the index $Jrm^2$ put $D=(U\setminus J)r$.
The integral shell recurrence has an error in $2^M\mathcal O$.
If $\omega(D)>T$, (F.5) gives
\[
 v_2\bigl(2^{n_U}B_U[Jrm^2]\bigr)\ge N_0+R.
\]
If $\omega(D)\le T$, (F.6) for visible coefficients and
$v_2a_M(D)\ge M$ for invisible ones instead give
\[
 v_2\bigl(2^{n_U}B_U[Jrm^2]\bigr)
       \ge\min\{M,n_U+\omega(r)\}.
\]
Consequently $v_2B_U[Jrm^2]\ge\min\{R,\omega(r)\}$ for every
Fourier index and every $U$ in the budget at this one stage.
At $r=1$ the normalized originating coefficient
$2^{|J|-n_U}a_M(U\setminus J)$ is integral, so (O.10) proves
the unary congruence modulo the integer two.
The additional division in (O.11) leaves error valuation at least
$R-1\ge1$.  In particular $M-n_U-1\ge1$ for every reading and
first-shell omission in the budget.
It does not select one stage for infinitely many budgets.

More precisely, fix a squarefree-output cutoff $L$. On a set of
stages in the same ultrafilter, (F.6) holds through $L$ and
$M_\theta>\delta_o+L$. Define $\mathcal A_{M_\theta,L}(D)$ only
for positive squarefree $D$ with $\omega(D)\le L$: it is zero off the full coefficient filter,
one when $v_2a_{M_\theta}(D)=\delta_o+\omega(D)$ on that filter,
and zero otherwise. This is exhaustive on the tested range:
visible coefficients satisfy (F.6), while invisible ones have
valuation at least $M_\theta>\delta_o+\omega(D)$.
In particular a unit in this range is visible and satisfies
\[
 \ord_{s=1}L(E^D,s)=1,\qquad
        j_{\rm gen}(D)\le\omega(D)+C_o,\quad C_o=\delta_o+C_{\rm det}.
 \tag{F.7}
\]
For the preceding operator budget take $L\ge W_U+1$ and include
its valid-stage requirements; all fresh-vertex and omission outputs
then lie in this range. The constant $C_o$ is independent of $L$
and of the stage. No value is assigned by this definition at
unbounded weights of one finite stage. No minimum is changed when
a vertex is omitted or a stationary condition is later added.
\end{proof}

\subsection{Nonempty rational seed and a common finite stage}
\begin{proposition}[Nonempty seeds after fixing the complementary twist]\label{ex:app-nonempty-seeds}
Choose the negative even unit $k$ with the fresh ramification and
exclusions specified in Section~\ref{sec:rc}, and fix
$K=\Q(\sqrt{k})$ before the address dimension. Then $E/K$ has full
two-primary Selmer corank one. Transport the even coefficient source
to $E^k$, normalize at index $1$, and use its positive square-unit
filter. On that filter its coefficients have
$v_2a_{\rm new}(D)\ge\omega(D)$, and the source has a nonempty
$K$-split unit seed.
If the intrinsic odd limit symbol is a unit at the empty
configuration, $E$ already has a simple zero. Otherwise it too has
a nonempty unit seed. In this remaining case both seeds can be
read through the actual common alphabet, whose projections onto
the two intrinsic alphabets are separately surjective. No
surjectivity onto their product is required.
\end{proposition}
\begin{proof}\leavevmode
Choose the final negative even unit $k$ as in
Section~\ref{sec:rc}, with fresh ramification outside the old odd
label and torsion fields. Fix $K=\Q(\sqrt k)$ before the binary
dimension. The forward theorem gives corank zero for $E^k$;
the sum-and-difference isogeny, on compact Selmer modules tensored
with $\Q_2$, gives full corank one for $E/K$.
The intrinsic odd symbol omits the primes of $k$.
Later $\ell_0$ avoids them. Fresh ramification implies that every
old odd label still lifts to a $K$-split label.
Additional stationary tests can have relations; use their actual
common alphabet and its separate surjective projections.

\paragraph{Transport to the fixed complementary twist.}
Put $d_k=|k|$, so that $k=-d_k<0$. The old source has sign $-1$;
its Fourier coefficient at the positive index $d$ tests $E^{-d}$.
Choose the new source for $E^k$ with sign $+1$, normalize its
coefficient at $1$ to be $1$, and impose a support containing the
old support and all primes dividing $k$. For the following comparison
$D$ is positive and squarefree, is prime to that support, and lies
in the square unit classes of $1$ required by the new filter.
In particular $D$ is a square in every old support completion, so
$d_kD$ and $d_k$ lie in the same old filter. Waldspurger comparison
on the two sources gives
\[
 \begin{split}
  v_2 a_{\rm new}(D)
    &=v_2 a_{\rm old}(d_kD)-v_2 a_{\rm old}(d_k)\\
    &\ge \omega(D).
 \end{split}
\]
The equality is a valuation equality; the two sources may have
different root-of-unity phases. The normalized central-value ratios
agree, since both are
\[
 \sqrt D\,\frac{L(E^{kD},1)}{L(E^k,1)}.
\]
Fixed period constants cancel in each ratio. The old unit condition
at $d_k$, and $\omega(d_kD)=\omega(d_k)+\omega(D)$, prove the
inequality. The signed twist is $kD$ throughout; no Fourier
coefficient at a negative index is used. The final family already
satisfies these unit classes. In the nonempty-seed argument the
fresh prime $p'$ is chosen to match all fixed labels of $p$, so
$pp'$ satisfies the same filter, including its dyadic condition.

For a unit of this normalized source equality holds, so
$v_2L_*(h_xk)\le2\omega(h_x)+C_{e,\rm final}$ with one fixed
period-normalization constant.
Its $K$-split seed must be nonempty. Suppose otherwise. The integral
coefficient valuations then improve the preceding lower bound to
$\omega(D)+1$ on every nonempty product of allowed $K$-split primes.
For a fresh good $K$-split prime $p$ outside the support, apply the
trace identity to $f=\theta H_{\{p\}}$:
\[
 U_p^2t(\phi)=U_p^3+d(\phi)U_p,\qquad
 (t(\phi)f)[p^2]=f[p^3]+(d(\phi)f)[p].
\]
Here and in this calculation all congruences are modulo the ideal
$2\mathcal O_{\rm coef}$ before passage to the residue field.
Write $a_p=a_p(E^{(k)})$.
We check every theta convolution term in the three coefficients.
Put
\[
\begin{gathered}
 B=B_{\{p\}}=\frac{(U_p-\eta(p)V_p)g}{2},\qquad
 b_1=B[1]=\frac{a(p)}2,\quad b_p=B[p],\\
 H=H_{\{p\}}=\frac{B-b_1\Theta_1-b_p\Theta_{\{p\}}}{2}.
\end{gathered}
\]
The inherited depth is zero, and $H$ is integral by (O.11).
Also $H[0]=0$, since the cuspidal source and the unary series have
zero constant term.
The ordinary lower bound gives $b_1\in\mathcal O_{\rm coef}$.
Under the assumed absence of a nonempty $K$-split unit seed one
has the stronger assertion
\[
                         b_1\in2\mathcal O_{\rm coef}.          \tag{G.4a}
\]
Indeed, if $p$ lies in the coefficient filter, the improved bound is
$v_2a(p)\ge2$; otherwise the actual filter gives $a(p)=0$.
No condition that $p$ itself belong to that filter is imposed.

Let $\delta_B,\delta_1,\delta_p$ be the scalars by which $d(\phi)$
acts on the full weight-two blocks
$\theta B,\theta\Theta_1,\theta\Theta_{\{p\}}$, respectively.
They are integral units.  The determinant is the cyclotomic character
times the weight-two nebentypus on every cuspidal constituent and
on the Eisenstein character pair, independently of old multiplicity.
For this positive-twist source $\psi/\eta=\chi_{-1}$ and $i=0$.
The half-integral-weight characters of $B$ and $\Theta_{\{p\}}$
are therefore both
\[
                 \psi(p/\cdot)=\eta(-p/\cdot).
\]
Multiplication by the same $\theta$ preserves this equality.
Consequently $\delta_B=\delta_p$ for the actual Frobenius lift
$\phi$, including its action on the ramified character.  The ratio
of the characters of $\theta\Theta_{\{p\}}$ and
$\theta\Theta_1$ is quadratic.  Their cyclotomic determinant factors
are identical, so $\delta_p/\delta_1\in\{1,-1\}$ and
$\delta_p-\delta_1\in2\mathcal O_{\rm coef}$.  We must not simply
equate these two scalars.  On the actual characteristic-zero
module the exact identity is
\[
 d(\phi)f=\delta_p f+
       \frac{(\delta_p-\delta_1)b_1}{2}\,\theta\Theta_1.
\]
Now $H[p]=0$, since $B[p]=b_p$, $\Theta_{\{p\}}[p]=1$ and
$\Theta_1[p]=0$.  Integrality of $H$ implies $f[p]\in2\mathcal O_{\rm coef}$.
Also $\theta\Theta_1[p]=2\sum_{0<j^2<p}\Theta_1[p-j^2]$ is even.
The displayed determinant identity therefore already proves
$(d(\phi)f)[p]\in2\mathcal O_{\rm coef}$ using only the ordinary
integrality of $b_1$.  The following direct expansion checks the
individual cross terms under the no-seed assumption in force here.

For $n_j=p-j^2>0$, one has $p\nmid n_j$ and
$\Theta_{\{p\}}[n_j]=0$.  We claim that
\[
                           B[n_j]\in2\mathcal O_{\rm coef}.    \tag{G.4b}
\]
If $n_j$ meets the old support, the coprimality sieve makes this
coefficient zero.  In particular $2\in S$, so every odd $j$ gives
an even $n_j$ and vanishes in this way.  Otherwise write
$n_j=r m^2$, with $r$ squarefree.
Since $p\nmid n_j$, the $V_p$ term is zero and
$B[n_j]=a(prm^2)/2$.  The integral good-prime shell recurrences
express $a(prm^2)$ as an integral multiple, up to unit character
factors, of $a(pr)$.  If $r>1$, the ordinary coefficient lower
bound gives $v_2a(pr)\ge\omega(pr)=1+\omega(r)\ge2$, or the
coefficient is zero off the filter.  This uses no assertion that
the primes dividing $r$ split in $K$.  If $r=1$, the same shell
recurrences express $B[m^2]$ as an integral multiple of $b_1$,
and (G.4a) applies.  This proves (G.4b).  Also every coefficient
of $b_1\Theta_1$ is even by (G.4a).

Now expand $\theta=1+2\sum_{j>0}q^{j^2}$ before discarding any
term.  Linearity of the determinant operator on the three actual
blocks gives the exact expression
\[
\begin{split}
 (d(\phi)f)[p]
 ={}&\frac{\delta_B B[p]-\delta_p b_p}{2}\\
 &+\sum_{0<j^2<p}
 \bigl(\delta_B B[p-j^2]
       -\delta_1 b_1\Theta_1[p-j^2]
       -\delta_p b_p\Theta_{\{p\}}[p-j^2]\bigr).
\end{split}                                                     \tag{G.4c}
\]
There is no constant-theta $\Theta_1[p]$ term since $p$ is not
a square, and $\Theta_{\{p\}}[p]=1$.
The first line is zero because $B[p]=b_p$ and
$\delta_B=\delta_p$.  Every summand in the second line is even
by (G.4a)--(G.4b) and the vanishing of its last term.  In particular
representations $p=j^2+m^2$, which can occur when $p\equiv1\pmod4$,
are covered by the case $r=1$ and are not omitted.  Thus
$(d(\phi)f)[p]\in2\mathcal O_{\rm coef}$.

The other two readings have no determinant operator.  Since $H$
is integral, for every positive integer $N$,
\[
 f[N]=H[N]+2\sum_{0<j^2\le N}H[N-j^2]
                 \equiv H[N]\pmod{2\mathcal O_{\rm coef}}.     \tag{G.4d}
\]
This includes $N=p^3$ and every term in its theta convolution.
Choose a fresh $K$-split prime $p'$ matching $\phi$ in a sufficiently
precise finite Galois quotient for the retained integral module,
and matching all fixed rational characters.  The good Hecke formula
at $p'\ne p$ identifies $(t(\phi)f)[p^2]$ modulo two with
$f[p^2p']$, hence with $H[p^2p']$ by (G.4d).
Neither unary series has support at $p^2p'$, and the recurrence at
$p\mid pp'$ has $\kappa=0$.  Consequently
\[
 H[p^2p']=\frac{a(p^3p')-\eta(p)a(pp')}{4}
        =\frac{\eta(p)(a_p-1)a(pp')}{4}.
\]
The product $pp'$ lies in the square filter even if $p$ alone does
not.  Its two primes split in $K$, so the no-seed assumption gives
$v_2a(pp')\ge3$.  This last coefficient is therefore even.
The trace identity and (G.4c)--(G.4d) now imply
$H[p^3]\in2\mathcal O_{\rm coef}$.  The following computation
uses precisely this unary-subtracted coefficient.

\paragraph{The coefficient behind the mod-eight contradiction.}
Write $a_p=a_p(E^k)$. The transported source tests positive twists of $E^k$. Thus
$\psi/\eta=\chi_{-1}$ and, at the squarefree index $1$,
$\kappa=(\psi/\eta)(p)(-1/p)=1$. Also $\eta$ is odd, so
the unary correction uses $i=0$. With $a(1)=1$ the recurrence gives
\[
 a(p^2)=\eta(p)(a_p-1),\qquad
 a(p^4)=\eta(p)^2\{a_p(a_p-1)-p\}.
\]
The inherited depth is zero. Hence
\[
 B_{\{p\}}[p]=\frac{\eta(p)(a_p-2)}2,\qquad
 B_{\{p\}}[p^3]=\frac{\eta(p)^2((a_p-1)^2-p)}2.
\]
Only $\Theta_{\{p\}}$ contributes at $p^3$, with coefficient
$\eta(p)p$. Therefore the actual integer division gives
\[
 H_{\{p\}}[p^3]
   =\frac{\eta(p)^2(a_p-1)(a_p-p-1)}4.
\]
This expression is integral because full rational two-torsion gives
$a_p\equiv p+1\pmod4$. The positive-twist character convention is
essential in this computation. Under the assumed absence of a
nonempty unit seed, the trace identity makes this coefficient zero
in the residue field. Since its phase is a unit and $a_p-1$ is odd,
it follows that $a_p\equiv p+1\pmod8$. The contradiction is obtained
only after this conditional deduction.
Choose $p$ also with $(k/p)=1$, $p\equiv1\pmod4$ and
$(\ell/p)=-1$, outside the finite support.  These prescriptions
are compatible: $K$ has a ramified prime outside the ramification
of $\mathbf Q(i,\sqrt\ell)$, so
$K\cap\mathbf Q(i,\sqrt\ell)=\mathbf Q$.  The corresponding
quadratic residue classes are therefore nonempty and Dirichlet's
theorem supplies infinitely many such primes.  Since $k$ is a
square modulo $p$, the reductions of $E^{(k)}$ and $E_\ell$ are
isomorphic.  A half of a nonzero two-torsion point with root $e_i$
on a split cubic requires both differences $e_i-e_j$ to be
squares.  For the roots $0,\ell,-3\ell$, one of these required
differences is, respectively, $-\ell,\ell,-4\ell$.  All three are
nonsquares under the chosen conditions.  Thus none of the three
nonzero points of order two has a half over $\mathbf F_p$.
If the two-primary subgroup had order greater than four, its
structure as a finite abelian group with exactly four elements
killed by two would give an element of order four, and hence such
a half.  Its order is consequently exactly four.  This gives
$v_2(p+1-a_p)=2$, contradicting the congruence modulo eight above.
This proves the required nonempty even seed. If the intrinsic odd
symbol already has a unit at the empty configuration, its detection
formula gives a simple zero for $E$ and the converse is finished.
Otherwise take a least nonempty odd unit seed. Thus the remaining
two-symbol address argument has two nonempty seeds.
\end{proof}

\begin{proposition}[Compatible finite graph tables on one ultrafilter]\label{ex:app-odd-stabilization}
Use $\delta_o$ and the ultrafilter of
Proposition~\ref{ex:app-odd-minimum}. First fix the abstract type of
the bounded elementary label groups and their projection data on
a set in this same ultrafilter. For every finite cutoff $W$,
the valid odd stages give finitely many possible symbol tables on
the compatible typed graphs with at most $W+1$ vertices.
One table occurs on a set in that ultrafilter. These selected tables
agree on their common restrictions, and their union defines the
intrinsic odd symbol on all compatible finite typed graphs.
It retains its bounded-weight witness and the nonempty-output
coefficient rules. After fixing $k$, its stationary tests may be
included in a common alphabet with separately surjective
projections on the same ultrafilter, without changing $\delta_o$.
Every specified finite collection of tests has a common actual
stage. No stage common to all cutoffs is selected.
\end{proposition}
\begin{proof}\leavevmode
Here the finite graph argument has a definite arithmetic order.
The moving theta prime contributes only a bounded number of quadratic
label coordinates. The root characters at the old level were fixed
before $M_\theta$; growing Kummer precision introduces no growing-order
character on the variable primes. Thus the elementary label groups,
their projections and their finite addition tables have bounded size.
Choose their abstract type on the same ultrafilter as $\delta_o$.
For each fixed weight cutoff $W$, restrict first to the
ultrafilter-large set of stages where the preceding budget and
precision conditions hold for every $|U|\le W$, with $R=2$.
On those stages the coefficient reading, factorization, contraction
and isolation rules hold on all configurations of at most $W+1$
vertices, with the stated nonempty-output restriction.  They therefore define actual symbol truth tables on
the compatible typed graph data.  There are finitely many such
tables; select the one occurring on a set in the same ultrafilter.
Two selected tables agree on their common restriction by taking
the finite intersection of their valid-stage sets. Their union defines
the odd limit symbol on compatible finite typed graphs of every size;
it is not a total valuation symbol at any one finite stage.
Every seed or contraction identity uses only finitely many such graphs
and is inherited at a cutoff containing all of them.
The bounded-weight witness remains a witness in these tables.
No stage common to all cutoffs is selected.
After fixing $k$, include its tests and choose the common projection
data on the same filter, without redefining the intrinsic odd minimum.
Stabilize the finite alphabet and its two projections in this way.
\end{proof}

\begin{proposition}[Stationary and simultaneous address constructions]\label{ex:app-simultaneous-address}
Use the definitions and realization freedoms in
Definitions~\ref{ex:sec7-typed-graph}, \ref{ex:sec7-symbol},
\ref{ex:sec7-address} and Lemma~\ref{ex:sec7-realization},
and the coefficient rules of Proposition~\ref{ex:app-coefficient-rules}.
Fix $b\ge1$. Either take one symbol on a nonempty set of always
active terminals with a nonzero polynomial in their pair bits, or
take one or two symbols on terminals with nonzero linear activations,
including a terminal of every nonzero activation, such that each
symbol separately has a nonzero terminal polynomial at every
$x\ne0$. In the latter case the symbols may read different
compatible projections of the common alphabet and fixed stationary
external data. Require the active data to satisfy each prescribed
filter at every nonzero address.
Starting from fixed compatible typed terminal and stationary
data, one can adjoin finitely many primary auxiliaries with nonzero
linear activations and choose only incidences involving them so that
all the symbols equal $1$ at every $x\ne0$. All original data stay
fixed, no auxiliary is present at zero, and the required neutral
external incidences and splitting conditions are retained.
Only nonempty configurations and nonempty contraction outputs are
used. A nonzero iterated difference, hence a partitionability
certificate, suffices for the required nonzero terminal polynomial.
For the two final symbols the zero-total packets below supply these
terminal hypotheses using only separate projection surjectivity.
For the even stationary application, the original prime representatives
may be fixed at the outset. For a stabilized odd symbol, first fix the
abstract data and the finite budget, then choose an actual stage
carrying them, and only then choose prime representatives. Previously
fixed representatives may be retained when compatible with the full
type at that stage; arbitrary earlier primes are not asserted to
satisfy a later moving condition.
\end{proposition}
\begin{proof}\leavevmode
\paragraph{Total labels of the final seeds.}
Both final symbols have projected total label zero. The intrinsic
odd filter contains $1$: its fixed unit squareclasses, positive sign,
and moving condition $(D/\ell_0)=1$ are those of $1$. The even source
for $E^k$ is likewise normalized in the unit filter containing $1$.
Every nonzero coefficient, and hence every unit seed, lies in its
respective filter. Additional stationary tests at primes dividing
$k$ refine the common alphabet but are forgotten by the intrinsic
odd projection. They do not add conditions to that odd symbol.

Take repeated packets with total zero in the common elementary
label group, including sufficiently many copies of every label.
For either projection separately, keep the first $r-1$ anchors of
its unit seed and contract the nonempty remainder into the last
anchor. Since the total projected label is zero, the remainder has
the last seed label. This argument uses surjectivity of each
projection separately. It does not require that a fixed old odd
seed admit lifts whose total is zero in the common alphabet.
The whole packet already has that total, and contractions preserve
it. Thus the final twists satisfy the stationary tests as well as
both coefficient filters.
Concretely, choose a packet with an even number of copies of every
common label, and enough copies to provide the anchors of either
seed under its respective projection.  Its total label is zero.
For every nonzero activation $\lambda$, take a disjoint copy of
this whole packet and give all its vertices activation $\lambda$.
At every $x\ne0$ there are exactly $2^{b-1}$ active packets.
The active list is therefore nonempty, has total common label zero,
and contains enough anchors for each projected seed separately.
Its terminal polynomial for each symbol is nonzero by the
contraction certificate, regardless of the initially fixed bits.
It also contains a terminal of each required activation.  This
constructs the terminal input of the address lemma explicitly.
No simultaneous lifting of a prescribed pair of seed labels is
needed.  Every added network vertex is primary, so the total
common label remains zero at each address.

\paragraph{Network operators and their degree.}
We spell out the simultaneous argument, including the product over
all addresses.  Write $\partial_e$ for the finite difference in a
terminal pair bit $e$.  Products of these operators commute and
$\partial_e^2=0$.  The contraction rule identifies a forest product
with its contracted graph.  A cycle product is zero: after a
spanning path is contracted its last edge is internal and changes
neither a type nor any external incidence.  No contraction removes
the last vertex.

A network of primary auxiliary vertices acts on any of the symbols
by an operator in the $\partial_e$.  To see this directly, expand in
the bits incident to the auxiliary vertices.  Contract the chosen
forest edges.  A component with no terminal becomes an isolated
primary vertex, and its contribution vanishes by primary isolation.
A component meeting several terminals gives their forest difference.
If $Q$ is the grounded auxiliary Laplacian and $c_a$ is its attachment
column at terminal $a$, the degree-zero and degree-one terms are
\[
 \det(Q)\,1+
 \sum_{a<b}c_a^{\mathsf t}\operatorname{adj}(Q)c_b\,
                    \partial_{ab}.
\]
All other terms have difference degree at least two.  The displayed
coefficients are the rooted-forest minors, by Cauchy--Binet on the
incidence matrix.  They are polynomial identities, so apply to
singular $Q$ as well.  Distinct networks have no edges between their
auxiliary vertices.  Their operators multiply: expanding both
forests gives their union, and unions containing a terminal cycle
give zero by the preceding observation.  This also shows that the
operator is the same for both projected symbols.

Fix a target $x_*\ne0$.  Choose a basis
$\lambda_1,\ldots,\lambda_b$ of the dual address space with
$\lambda_i(x_*)=1$ for all $i$.  Such a basis is obtained from one
functional taking value one by adding it to a basis of the space
of functionals vanishing at $x_*$.  Group $i$ has three primary vertices, each activated
by $\lambda_i$.  In its physical three-dimensional coordinate space
let $e_i=(1,1,1)$ and choose $u_i,v_i$ spanning its orthogonal
complement.  The vector $e_i$ has norm one.  Prescribe the symmetric
matrix $Q$ by its bilinear matrix $I_b$ on the $e_i$ and
\[
 H(T,B)=\begin{pmatrix}0&T\\T^{\mathsf t}&B\end{pmatrix}
\]
on the remaining coordinates.  In physical coordinates each
$h_i=Qe_i=e_i$ is supported on its own group.  Attach this column to
a terminal activated by $\lambda_i$, which exists by hypothesis.
For an active set $J$ the sum of these active attachment columns is
$Q_J\mathbf1_J$.  Thus the prescribed principal matrix $Q_J$ really
is the grounded graph Laplacian on that address, including its
diagonal entries.

For a selected terminal pair $a,b$ active at $x_*$, let their
activations be $\alpha,\beta$.  Attach an extra column $c$, dual to
$u_1$, to both.  If $\alpha\ne\beta$, attach the same column to a
terminal of activation $\alpha+\beta$; it is absent at $x_*$.  The
extra columns sum to zero at every address since
$\alpha+\beta+(\alpha+\beta)=0$.  If $\alpha=\beta$, the two
columns already cancel.  Hence they do not alter $Q_J$.

For the singular network set $T=I+P$, where $P$ is the cyclic shift,
and give $B$ one diagonal entry one.  On every proper group subset,
the restricted cyclic shift is a union of paths, so $T_J$ is
invertible and $\det Q_J=\det(T_J)^2=1$.  On all groups, the kernels
of $T$ and $T^{\mathsf t}$ are generated by the all-one vector.
The one nonzero diagonal entry of $B$ forces the lower component of
a vector in $\ker H$ to vanish.  Thus $\ker Q$ is precisely the line
generated by $r=\sum_i u_i$, and
$\operatorname{adj}(Q)=rr^{\mathsf t}$.  The anchor columns pair
to zero with $r$, while $c^{\mathsf t}r=1$.  Only $a,b$ receive
active extra columns at $x_*$.  Consequently its network operator
$S_{x_*,e}$, for $e=ab$, has the expansions
\[
 S_{x_*,e}|_{x_*}=\partial_e+\text{terms of degree at least two},
 \qquad
 S_{x_*,e}|_x=1+\text{terms of positive degree}\quad(x\ne x_*).
\]
For $b=1$, $T=0$ and $B=(1)$ give exactly the same calculation.

For a varying network set $T=I+L$, with $L$ the consecutive lower
shift, and $B_\varepsilon=\varepsilon E_{bb}$ for
$\varepsilon\in\{0,1\}$.  Every $T_J$ is invertible and both
operators $V^\varepsilon_{x_*,e}|_x$ have constant coefficient one.
The upper-left block of $H^{-1}$ is
$T^{-\mathsf t}B_\varepsilon T^{-1}$.  Its change on the selected
columns is $((T_J^{-1})_{b1})^2$, interpreted as zero when an endpoint
is absent.  This is one exactly when the entire path is present,
which means $x=x_*$.  Anchor columns have zero change.  Therefore
\[
 \begin{aligned}
 (V^0_{x_*,e}+V^1_{x_*,e})|_{x_*}
       &=\partial_e+\text{terms of degree at least two},\\
 (V^0_{x_*,e}+V^1_{x_*,e})|_x
       &=\text{terms of degree at least two}\quad(x\ne x_*).
 \end{aligned}
\]

\paragraph{One choice for both symbols at every nonzero address.}
For each $x\ne0$ let $P_{1,x},P_{2,x}$ be the two nonzero terminal
polynomials, of degrees $d_{1,x},d_{2,x}$, and let $H_{s,x}$ be their
highest homogeneous parts.  Work in the ordinary polynomial ring
in the terminal pair variables for this address.  Since it is a
domain, $H_{1,x}H_{2,x}\ne0$.  Choose a coefficient-one monomial
\[
 M_x=\prod_e X_e^{n_{x,e}},\qquad n_{x,e}\in\{0,1,2\},\qquad
 \sum_e n_{x,e}=d_{1,x}+d_{2,x}.
\]
For $n_{x,e}=2$ install one singular network $S_{x,e}$; it acts on
both symbol factors.  For $n_{x,e}=1$ install one varying network
$V^{\varepsilon_{x,e}}_{x,e}$.  Use disjoint auxiliary sets for all
these networks.  Let $\Gamma_y(\varepsilon)$ be the resulting graph
at address $y$, and consider the finite sum in $\mathbf F_2$
\[
 \sum_{\varepsilon}
       \prod_{y\ne0}\prod_{s=1}^2 A_s(\Gamma_y(\varepsilon)).
                                                               \tag{G.1a}
\]
The total degree available to all terminal polynomials is
$D=\sum_{y\ne0}(d_{1,y}+d_{2,y})$.

Each singular network at $x$ has zero constant at $x$, so its action
costs at least one difference in each of the two factors there, hence
at least two in (G.1a).  For a varying network, summing over its
$\varepsilon$ cancels the all-constant term.  Its first-degree
terms are independent of $\varepsilon$ away from its target, and
their difference at the target is $\partial_e$.  Thus every
surviving term costs at least one difference, and a cost-one term
acts at its target on exactly one of the two factors.  The sum of
all these minimum costs is exactly $D$.  Any higher-degree action,
or any additional action at a wrong address, exceeds $D$ and
annihilates at least one terminal polynomial.  Hence only the
minimum-cost terms survive.

At address $x$, a variable with exponent two is therefore
differentiated once in each factor; a variable with exponent one
is differentiated in exactly one factor, summed over the two
possibilities.  The resulting scalar is precisely
$[M_x](H_{1,x}H_{2,x})=1$.  Lower-degree parts cannot contribute.
Taking the product over all $x\ne0$ proves that (G.1a) equals one.
At least one choice of the finitely many $\varepsilon$ consequently
makes every factor in that product one.  This is one graph for both
symbols at every nonzero address.  Boolean reduction of
$H_{1,x}H_{2,x}$ before selecting $M_x$ would destroy this
exponent-two accounting and is not used.
Here is the stationary and one-symbol specialization, so that the
application in Proposition~\ref{ex:app-rough-companion} uses the same
proved construction. If all terminals are always active, attach
each column $h_i=Qe_i$ to any one fixed terminal. Since $h_i$ is
supported on its own group, their restricted sum is
$Q_J\mathbf1_J$ at every address. The two additional columns $c$
at a selected terminal pair cancel directly; no balancing terminal
of activation $\alpha+\beta$ is needed. The same singular program
therefore isolates that pair difference at its target address.
For one symbol, choose a coefficient-one top monomial of each
nonzero $P_x$ and put a separate singular program at $x$ for each
variable in that monomial. Its own programs already spend
$\deg P_x$ contraction degrees at $x$; any extra degree kills the
term. Thus exactly the selected top coefficient survives, and it
is one. A constant nonzero $P_x$ needs no program. This proves the
stationary one-symbol case as well as the moving one-symbol case.
In all cases the primary realization of
Lemma~\ref{ex:sec7-realization} changes only incidences with
the newly adjoined vertices. These are finite constructions for
each $b$, with no assertion of one infinite prime array.

\end{proof}

\begin{proposition}[One actual stage for the final family]\label{ex:app-final-family}
In the nonempty-seed case of Proposition~\ref{ex:app-nonempty-seeds},
fix the final $k,K$ and the stabilized common alphabet. For every
sufficiently large $b$, the address construction is realized at one
actual theta stage by a finite set $Q$ of distinct good odd primes
split in $K$, avoiding the fixed support and the specified finite
exclusions. With nonzero activations $\lambda_q$ it gives exactly
the family in (3.9),
\[
 h_x=\prod_{q\in Q}(q^*)^{\lambda_q(x)},\qquad h_0=1,
 \qquad q^*=(-1)^{(q-1)/2}q.
\]
All fixed local types are retained, and for every $x\ne0$,
\[
 \operatorname{ord}_{s=1}L(E^{h_x},s)=1,\quad L_*(h_xk)\ne0,
 \qquad j_{\rm gen}(h_x)\le\omega(h_x)+C_o,\quad
 v_2L_*(h_xk)\le2\omega(h_x)+C_e.
\]
The constants precede $b$. The finite graph budget and then the
actual theta precision are chosen after $b$. The theta prime remains
a coefficient prime and exclusion, not a relaxed Selmer prime.
\end{proposition}
\begin{proof}\leavevmode
For each chosen $b$, this construction has a finite budget $W_b$.
Enlarge it to bound every graph product and squarefree output in its
finite auxiliary tests, as well as every operator set. Contractions
only decrease the number of vertices. Thus every final $h_x$ lies in
the tested range of (F.7); this enlargement changes no fixed constant.
Only now take one actual theta stage in the finite intersection
of the selected table sets through $W_b+1$, the valid-stage set
for the full budget $|U|\le W_b$ with $R=2$, and the required
precision and finite exclusion sets.  In particular (F.6) holds
through the threshold $T$ used for that budget.
Realize its prescribed forward residue bits by distinct primes using CRT and Dirichlet;
quadratic reciprocity supplies the reverse bits.
There is no uniform bound on $W_b$, no intersection over all
budgets, and no product-surjectivity assumption on the alphabets.
The resulting one family satisfies (3.9), with constants fixed
before $b$. The theta prime $\ell_0$ is a coefficient prime and
an exclusion, not an active relaxed Selmer prime or a moving
arithmetic prime $r_i$.
\paragraph{Finite intersection and the upper bound.}
The required stage conditions form a finite list once the address
program is fixed. Let $\mathcal U$ be the chosen ultrafilter and
let $A_1,\ldots,A_t\in\mathcal U$ be the selected truth-table
sets used by that program, including all intermediate graphs.
Let $B\in\mathcal U$ impose (F.6) through the threshold $T$
for its operator budget and impose
$M\ge\max(0,\delta_o+W_b)+2$ and $M>T+\delta_o$.
Finite prime exclusions are included in $B$.
Then $B\cap\bigcap_{i=1}^t A_i\in\mathcal U$, hence is nonempty.
Choose a stage there before realizing the prime labels. The
coefficient rules hold for every realization at this stage, not
merely for a previously selected finite list of Fourier indices.
The address construction makes each nonzero address a unit of
both sources. For the odd source this says exactly
\[
 v_2a_M(h_x)=\delta_o+\omega(h_x)<M.
\]
The strict inequality follows from the tested-range precision
condition. Applying (O.8) gives the simple zero and
$j_{\rm gen}(h_x)\le\omega(h_x)+\delta_o+C_{\rm det}$.
The normalized even source gives
$v_2L_*(h_xk)\le2\omega(h_x)+C_{e,\rm final}$.
Thus one may take $C_o=\delta_o+C_{\rm det}$ and
$C_e=C_{e,\rm final}$, both fixed before $b$. Only the number
of tests, the threshold $T$, and the chosen stage depend on the
address dimension.
\end{proof}

\subsection{One coordinate in the final contradiction}
The family in Proposition~\ref{ex:app-final-family} is now used with
the same coordinate of Lemma~\ref{ex:arith-common-coordinate}
and its central calibration in
Proposition~\ref{ex:arith-central-calibration}. The denominator and
coefficient comparisons are those of
Propositions~\ref{ex:limits-fixed-field-clearing}
and~\ref{ex:limits-coefficient-transfer}; the zero-base assertion is
Lemma~\ref{ex:arith-base-zero}.
For this fixed $K$, (3.4), (F.7), and the even bound just proved give
\[
 2j_{\rm rc}(h_x,k)\le4\omega(h_x)+C_J,\qquad
 C_J=2C_o+C_{e,\rm final}-v_2L_*(h_*)+C_{H,\rm final}
\]
at every $x\ne0$.
The finite support comparison uses the same class, duality map,
bases and coefficient stages. Its integral correction
\[
 \eta=\prod_q(1+v_q(1+g_q)),\qquad l_q/e_q^2=1+2v_q,
 \qquad \widetilde u=\eta u_{\rm rel}
\]
is formed before characters. The actual integral determinant
triangle gives $\widetilde u_x=\epsilon_xu_{{\rm prim},x}$ with
$\epsilon_x\in\Lambda^\times$, rather than independently chosen
units in $\mathscr O^\times$. Therefore the central identity is
\[
 \begin{split}
 v_2\widetilde u_x(0)={}&2j_{\rm rc}(h_x,k)-4\omega(h_x)
       -s(h_x)-s(h_xk)\\
       &-\sum_{p\mid k}t_p+\epsilon'_x,\qquad
 0\le t_p\le2,\quad |\epsilon'_x|\le C_{\rm cal}.
 \end{split}                                                    \tag{F.8}
\]
The nonnegative finite $\Sha$ terms give one upper bound
$v_2\widetilde u_x(0)\le B$ for all $x\ne0$, independent of $b$.
At the unknown base no finiteness of $\Sha$ is assumed.
The retained central class and corank-one argument instead give
$\widetilde u_0(0)=0$ if the base derivative vanishes.

The fixed-field construction of Section~\ref{sec:limits} has
\[
 \begin{gathered}
 c_{\rm det}=a+4(h_K+2c_E+2t_K),\quad
 L_T(M)=16(d_K+5)M,\\
 f_x\widetilde u_x=d_x,\qquad
 c=v_2f_0(0)\le c_{\rm det}.
 \end{gathered}
\]
Choose $M>\max(c_{\rm det},B+2c_{\rm det},0)$, then
$b>(L_T(M)+2)(2^M-1)$ large enough for the address construction,
then its finite graph budget and actual theta stage $M_\theta$.
The residual comparison tests coefficients of $f$ from
$-L_T(M)$ through zero and only the constant coefficient of $d$.
It gives $x\ne0$ with
\[
       v_2(\widetilde u_x(0)-\widetilde u_0(0))\ge M-2c>B.
\]
The estimate uses both $v_{\mathscr O}(f_x)\le c$ and
$\widetilde u_x\in\Lambda[1/2]$; neither a bound on the whole
numerator nor a finite-Artin inverse of $t$ is asserted.
This contradicts (F.8) under the supposed base zero.
The Rankin factorization and $L(E^k,1)\ne0$ give the asserted
simple zero of $E$.

This completes the proof of the restricted converse,
Theorem~\ref{hyp:RC}, from the displayed classical inputs and
the arithmetic, coefficient and determinant constructions.
The conversion from that converse to pole witnesses, the integer
model, the natural compiler and finally the ordinary rational value
is the separate argument in the main text.

% End 09-normalization-endpoints.tex



% Begin references.tex
% Original exposition of the prime-pattern argument, prepared for Version 5.
% Requires amsmath, amssymb, amsthm, hyperref and the manuscript theorem environments.
% Publication-status qualifications in this file are intentional.
\section{Simultaneous primes with coprime adjacent norms}
\label{sec:v5-sun-proof}

This appendix gives the prime-pattern argument needed for the construction.
Its strategy is due to Sun~\cite{v5sun}. We supply the local calculations,
the two ranges of the sieve, the normalization of the ray-class densities,
and the order in which the parameters are chosen. For the analytic input
we prove the special case for fixed ray class fields of $\Q(i)$ directly.
The ideal-norm method follows Frei and Tinkov\'a~\cite{v5ft}; no theorem
from that work is assumed. The proof uses the published fixed-field prime
estimate of Mitsui, the one-parameter equidistribution and factorization
results of Green--Tao, and published ordinary-prime transference results.
The field and its ray modulus are fixed before every asymptotic limit.
The exact source versions checked are identified in the bibliography
and the accompanying theorem-application inventory.

Write $\mathrm N(z)=z\bar z$ for $z\in\mathbb Z[i]$, and put
$d_-(z)=\mathrm N(z-1)$ and $d_+(z)=\mathrm N(z+1)$.
An Artin or Frobenius element in this appendix is arithmetic Frobenius:
on residue fields it sends $x$ to $x^{\mathrm N\mathfrak p}$.

\begin{theorem}[The simultaneous prime input]\label{thm:v5-simultaneous-primes}
Let $c_1,\ldots,c_k$ be distinct integers, where $k\geq2$.
Let $q\geq1$, let $M$ be an odd squarefree positive integer with $7\nmid M$,
and let $H,Y>0$. There are positive integers $u,w$ such that
\[
 q\mid w,\qquad u/w^2>H,\qquad p_j=u+c_jw^2>Y
 \quad(1\leq j\leq k),
\]
where every $p_j$ is prime. Each $p_j$ has a Gaussian prime generator
$\pi_j$ of norm $p_j$ such that
\[
 \pi_j\equiv9\pmod {16},\qquad
 \pi_j\equiv1+3i\pmod7,\qquad
 \pi_j\equiv i\pmod M.
\]
No odd prime divides both $d_-(\pi_j)$ and $d_+(\pi_j)$.
For $j\ne h$, no odd prime divides both $d_\epsilon(\pi_j)$ and
$d_\eta(\pi_h)$, for any $\epsilon,\eta\in\{-,+\}$.
Moreover,
\[
 v_2(d_-(\pi_j))=6,\quad v_2(d_+(\pi_j))=2,\quad
 \gcd(d_-(\pi_j),d_+(\pi_h))=4\quad\hbox{for all }j,h,
\]
and all these norms are coprime to $7M$.
\end{theorem}

The quantifiers concern fixed data. No estimate uniform in $q,M,k$ or the
$c_j$ is asserted. We first prove the result for $q=1$. Replacing the
coefficients by $q^2c_j$ and then replacing the parameter $v$ by $w=qv$
proves the general case, since the ratio before this substitution can be
made greater than $q^2H$.

\subsection{Local counts and the medium primes}\label{subsec:v5-local-sieve}

Fix a nonzero integer $\kappa$ and write
\[
 Q_\kappa=a^2+b^2-c^2-d^2-\kappa w^2,\quad
 G_\epsilon=(a+\epsilon)^2+b^2,\quad
 H_\eta=(c+\eta)^2+d^2,
 \qquad \epsilon,\eta\in\{1,-1\}.
\]
We use $\chi(\ell)=(-1/\ell)$ for odd primes $\ell$.

\begin{lemma}[Finite-field counts]\label{lem:v5-local-counts}
If $\ell\nmid2\kappa$, then $Q_\kappa=0$ has $\ell^4-1$ nonzero
solutions over $\mathbb F_\ell$. At most $4\ell^2$ of its solutions satisfy
$G_\epsilon=H_\eta=0$. If $\chi(\ell)=-1$, the latter set has exactly one
element.

Let $\gamma_1,\ldots,\gamma_k$ be distinct integers, including $0$ and
$-\kappa$, and put $L_j=a^2+b^2+\gamma_jw^2$. Outside a fixed set of primes
containing the divisors of $2\kappa\prod_{j<h}(\gamma_j-\gamma_h)$, the
proportion $g(\ell)$ of nonzero zeros of $Q_\kappa$ on which some $L_j$
vanishes is
\begin{equation}\label{eq:v5-local-g}
 g(\ell)=
 \frac{(\ell-1)(\ell-\chi(\ell))
 [2\nu_0+(k-2)(\ell-\chi(\ell))]+\nu_0^2-1}{\ell^4-1},
 \qquad \nu_0=\ell+(\ell-1)\chi(\ell).
\end{equation}
Consequently $g(\ell)=(k+2\chi(\ell))/\ell+O_k(\ell^{-2})$.
\end{lemma}
\begin{proof}
The norm $x^2+y^2$ has $\nu_0$ representations of zero and
$\ell-\chi(\ell)$ representations of every nonzero element. This follows
either by factoring over $\mathbb F_\ell$ when $\chi=1$, or by the norm map
from $\mathbb F_{\ell^2}$ when $\chi=-1$. Summing these fibres for
$Q_\kappa$ gives $\ell^4$ solutions including the origin.
In the split case, each of the equations $G_\epsilon=0$ and $H_\eta=0$
is a union of two lines. On each of the four pairs of lines the remaining
equation is
$-2\epsilon a+2\eta c-\kappa w^2=0$, leaving at most $\ell^2$ choices.
In the inert case these equations force
$(a,b,c,d)=(-\epsilon,0,-\eta,0)$, and then $w=0$.

For $w\ne0$, the conditions $L_j=0$ are disjoint. For the two distinguished
coefficients one norm is zero and the other is nonzero, giving
$2\nu_0(\ell-\chi)$ choices. Each other coefficient gives
$(\ell-\chi)^2$ choices. There are $\ell-1$ possible nonzero $w$.
For $w=0$, both norms are zero, giving $\nu_0^2-1$ nonzero vectors.
This proves \eqref{eq:v5-local-g}. In particular, when $k=2$ and
$\chi=-1$, its value is $2/(\ell^2+1)$, not zero.
\end{proof}

\begin{lemma}[Primitive congruence count]\label{lem:v5-primitive-count}
Let $\psi$ be a smooth nonnegative compactly supported function away from
the origin, and let $m$ be squarefree and prime to $2\kappa$.
For each $\ell\mid m$ prescribe an arbitrary subset $\Omega_\ell$ of the
nonzero zeros of $Q_\kappa$ over $\mathbb F_\ell$. The weighted number of
primitive integral vectors on $Q_\kappa=0$ with these reductions equals
\begin{equation}\label{eq:v5-primitive-count}
 C_\kappa\sigma_\infty(\psi)B^3
 \prod_{\ell\mid m}\frac{|\Omega_\ell|}{\ell^4-1}
 +O_{\kappa,\psi,\varepsilon}
       (B^{5/2+\varepsilon}m^{5/2+\varepsilon}).
\end{equation}
The constant $C_\kappa$ is positive and independent of $m$ and the subsets.
\end{lemma}
\begin{proof}
Apply Browning--Loughran~\cite{v5bl}, Theorem~4.1, to $Q_\kappa$ in five
variables. It is isotropic, since $(1,0,1,0,0)$ is a zero, and nonsingular,
since $\kappa\ne0$. Its discriminant has no prime divisor outside
$2\kappa$. The theorem permits $n=5$ and arbitrary nonzero residue subsets;
its error is uniform in the modulus and in these subsets.
Impose primitivity by $\sum_{h\mid(a,b,c,d,w)}\mu(h)$.
If $(h,m)>1$ there are no permitted vectors. Otherwise the subsets become
$h^{-1}\Omega_\ell$, of the same cardinalities, and the scale becomes $B/h$.
The sum of the error terms is convergent. Extending the main-term sum to
infinity costs $O(B)$.
At a good prime the unrestricted local density is
\[
 \sigma_\ell=\frac{1-\ell^{-4}}{1-\ell^{-3}}.
\]
Indeed, nonsingular primitive zeros lift with four free coordinates, and
the vectors divisible by $\ell$ contribute the geometric factor
$\ell^{-3}$. The primitive density is therefore $1-\ell^{-4}$.
This turns the factor $|\Omega_\ell|/\ell^4$ of the unprimitive count into
$|\Omega_\ell|/(\ell^4-1)$ in \eqref{eq:v5-primitive-count}.
The remaining finitely many local factors are positive because the form
has a nonsingular rational zero. Empty subsets give the assertion directly.
\end{proof}

In using this count, the two norms lie between fixed positive multiples of
$B^2$, and $|w|\leq\tau B$, where $0<\tau\leq\tau_0$.
A smooth majorant can be chosen with
$\sigma_\infty(\psi)\ll\tau$ uniformly in $\tau$: integrate first in the
two radial norm variables, on which the derivative of their difference is
bounded away from zero. The constants in error terms may depend on
derivatives of this majorant, and hence on $\tau$.

\begin{lemma}[Medium-prime bound]\label{lem:v5-medium-sieve}
For fixed coefficients and sufficiently large fixed $L$, the contribution
of prime patterns with a common adjacent-norm divisor
$L<\ell\leq B^{1/100}$ is at most
\begin{equation}\label{eq:v5-medium-sieve}
 \frac{C\tau B^3}{(\log B)^k L\log L}+O_\tau(B^{27/10}).
\end{equation}
The count is unweighted by logarithms, and $C$ is independent of $L,\tau$.
The bound covers both distinct positions and the two signs at one position.
\end{lemma}
\begin{proof}
For two positions, use $Q_\kappa$ with their fixed nonzero coefficient
difference and impose $G_\epsilon=H_\eta=0\pmod\ell$.
For squarefree $d$ supported on good primes other than $\ell$, impose also
that some $L_j$ vanishes at every prime dividing $d$. The preceding lemmas
give a sieve sequence
\[
 A_{\ell,d}=X_\ell g(d)+r_{\ell,d},\qquad
 X_\ell\ll\tau B^3/\ell^2,\qquad
 |r_{\ell,d}|\ll_{\tau,\varepsilon}
 B^{5/2+\varepsilon}(\ell d)^{5/2+\varepsilon}.
\]
Enlarge the fixed exceptional set so that $0<g(p)<1$ outside it.
Set $h(p)=g(p)/(1-g(p))$ and extend $h$ multiplicatively to squarefree
integers. The elementary Selberg upper sieve has denominator
\[
 \mathcal G_\ell(z)=\sum_{\substack{d<z\ \text{squarefree}\\
                         (d,\ell\prod_{p\in S}p)=1}}h(d)
 \gg (\log z)^k.
\]
Here is the dimension calculation. Williams~\cite{v5williams}, Theorem~1
with fixed modulus $4$ and residues $1,3$, gives Mertens' estimates in the
two reduced classes. Taking logarithms gives
$\sum_{p\leq y}g(p)=k\log\log y+O(1)$ and
$\sum_{p\leq y}g(p)\log p\ll\log y$.
For the latter bound only the elementary Chebyshev estimate is needed:
$\prod_{n<p\leq2n}p\mid\binom{2n}{n}<4^n$ gives
$\vartheta(x)\ll x$ by dyadic summation, and partial summation gives
$\sum_{p\leq y}(\log p)/p\ll\log y$.
In the Euler product $\prod_{p\leq y}(1+h(p))$, normalize the weights
$h(d)$ to a probability distribution. Its expected value of $\log d$ is
$\sum_{p\leq y}g(p)\log p$. Taking $y$ to be a sufficiently small fixed
power of $z$ leaves at least half of the mass on $d<z$, by Markov's
inequality. Removing one prime $\ell$ changes the estimate by a bounded
factor.

For completeness the sieve weights may be taken as
\[
 \lambda_d=\frac{\mu(d)}{\mathcal G_\ell(z)}
 \prod_{p\mid d}(1-g(p))^{-1}
 \sum_{\substack{r<z/d\ (r,d)=1\ r\ \text{squarefree}\\
                      (r,\ell\prod_{p\in S}p)=1}}h(r).
\]
They have $\lambda_1=1$, $|\lambda_d|\ll_\varepsilon d^\varepsilon$,
and direct expansion gives
$\sum_{d,e}\lambda_d\lambda_e g([d,e])=1/\mathcal G_\ell(z)$.
Thus the number with no prime divisor below $z$ in any $L_j$ is at most
\[
 \frac{C\tau B^3}{\ell^2(\log z)^k}
 +O_{\tau,\varepsilon}
  (B^{5/2+\varepsilon}\ell^{5/2+\varepsilon}z^{7+4\varepsilon}).
\]
All the prime values under consideration exceed $z$. Put $z=B^{1/100}$
and sum over $L<\ell\leq B^{1/100}$. The main term uses
$\sum_{\ell>L}\ell^{-2}\ll1/(L\log L)$, another consequence of the
same Chebyshev bound; for example
$\varepsilon=1/1000$ makes the error smaller than $B^{27/10}$.

At one position, an odd common divisor of $(a-1)^2+b^2$ and
$(a+1)^2+b^2$ forces $a=0$ and $b^2=-1\pmod\ell$.
There are at most $2\ell$ such triples $(a,b,w)$ over $\mathbb F_\ell$.
For $w\ne0$ the zero sets of the $k$ prime forms are disjoint, and the
ordinary three-dimensional local sieve density is
$(k+\chi(\ell))/\ell+O_k(\ell^{-2})$.
Counting residue classes in a smooth box gives the corresponding
three-dimensional congruence estimate. The same Selberg weights and the
same, weaker error in \eqref{eq:v5-medium-sieve} apply.
\end{proof}

\subsection{The large-prime tail}\label{subsec:v5-large-sieve}

We use Lemmas~2.1 and 4.1 of Browning--Heath-Brown~\cite{v5bhb}.
The adjacent-norm polynomials are inhomogeneous, so we give the
fixed-coordinate argument rather than invoke their homogeneous
quadric-sieve theorem directly.
Their Lemma~2.1 says that coprime polynomials in $\mathbb Z[X_1,\ldots,X_m]$
of bounded degree and coefficient height at most $H$ have at most
\begin{equation}\label{eq:v5-box-sieve}
 O\left(\frac{V\log(VH)}{L\log L}
           +\frac{V\log(VH)}{B_{\min}}\right)
\end{equation}
integer points in a box of side parameters $B_i\geq2$ at which a prime
greater than $L\geq2$ divides every polynomial; $V=\prod B_i$.
``Coprime'' here excludes a common constant prime factor as well as a
common nonconstant factor. Their Lemma~4.1 covers the zeros modulo a
squarefree $q$ of a rank-$m$ quadratic form, with $m\geq3$, by at most
$(3m)^{\omega(q)}q^{m-2}$ lattices of determinant $q^{m-1}$.
The number with largest Euclidean successive minimum exceeding $T$ is
\begin{equation}\label{eq:v5-lattice-tail}
 O_m((3m)^{\omega(q)}q^{2m-3}T^{-m}).
\end{equation}
Every prime dividing $q$ must be odd and preserve the rank of the form.
The lattices in that lemma contain $q\mathbb Z^m$; in particular their
largest successive minima are at most $q$.

\begin{lemma}[Large common divisors on the quadric]\label{lem:v5-large-sieve}
Fix the implied constant in the box of size $O(B)$.
For fixed $\kappa\ne0$ and $\epsilon,\eta\in\{1,-1\}$, uniformly for
$B,L\geq2$, the number of integral vectors in this box satisfying
$Q_\kappa=0$ and
$\ell\mid G_\epsilon,H_\eta$ for some prime $\ell>L$ is
\begin{equation}\label{eq:v5-large-sieve}
 O_{\kappa,\alpha}\left(
    \frac{B^{3+\alpha}}{L\log L}+B^{11/4+\alpha}\right)
 \qquad(\alpha>0).
\end{equation}
\end{lemma}
\begin{proof}
Sign changes of $a,c$ reduce to $\epsilon=\eta=1$.
Put $x_0=a+c$, $x_1=a-c$, $x=b-d$, $z=b+d$, and
$R=\kappa w^2-xz$. Then $x_0x_1=R$.
Parity restrictions can be discarded in an upper bound.
For a fixed nonzero value of either $x_0$ or $x_1$, the triple $(x,z,w)$
determines the other coordinate uniquely whenever it gives an integral
solution. We shall use this injection when discarding divisibility
conditions.

First suppose that $\ell\mid x_0,x_1$. Then $\ell^2\mid R$.
For each fixed integer $v=O(B^2)$, the equation $R=v$ has
$O_{\kappa,\alpha}(B^{1+\alpha})$ triples: for each $w$, count the
factorizations $xz=\kappa w^2-v$ by the divisor bound; the zero right-hand
side occurs for at most two $w$ and contributes $O(B)$ each.
If $R\ne0$, it has $O_\alpha(B^\alpha)$ factorizations $x_0x_1=R$.
Since $0<|h|\ll B^2/\ell^2$ and $\ell\ll B$, summation gives
\[
 O_{\kappa,\alpha}\left(
 B^{3+2\alpha}\sum_{\ell>L}\ell^{-2}\right)
 =O_{\kappa,\alpha}(B^{3+2\alpha}/(L\log L)).
\]
The same bound for triples with $R=0$, followed by the $O(B)$ choices
with $x_0x_1=0$, gives $O_{\kappa,\alpha}(B^{2+\alpha})$.
We may discard all these $R=0$ solutions before proceeding.

In the remaining cases choose $t=x_1$ or $t=x_0$ such that $\ell\nmid2t$.
Take $t=x_1$ first. Here
$G_1-H_1=\kappa w^2+2t$ and $4t^2G_1=K$.
Thus the two norm congruences are equivalent to divisibility of
\[
 P=\kappa w^2+2t,\qquad
 K=(\kappa w^2-xz+t^2+2t)^2+t^2(x+z)^2.
\]
The identity
\[
 K-(t^2+x^2)(t^2+z^2)=P(P+2t^2-2xz)
\]
shows that it suffices to sieve the coprime pair
$P,(t^2+x^2)(t^2+z^2)$. They have no common nonconstant factor over
$\mathbb Q$, because the first depends nontrivially only on $w$, whereas
the second is independent of $w$. Their contents have no common prime,
since the second has coefficient $1$ on $x^2z^2$.
For $t=x_0$ we still have $4t^2G_1=K$, whereas
$t(G_1-H_1)=\kappa(t+2)w^2-2xz$. The pair is therefore
\[
 P'=\kappa(t+2)w^2-2xz,\qquad K.
\]
If $t\ne-2$, $P'$ is an irreducible ternary quadratic form of rank three
and $K$ has nonzero constant term $t^2(t+2)^2$; homogeneity rules out
divisibility. If $t=-2$, neither $x$ nor $z$ divides $K$.
The coefficient $1$ of $x^2z^2$ in $K$ excludes every common content
prime. In both charts the fixed value $t$ is nonzero.

In either chart, $t\mid R$. Put
\[
 S_R=\{p:p\mid2\kappa\},\qquad
 q=\prod_{\substack{p\mid t\ p\notin S_R}}p.
\]
The matrix of $R$ has determinant $-\kappa/4$, so $R$ has rank three
at every prime in this product, and those primes are odd. Also $q\ll B$.
For a fixed $q$, every eligible nonzero $t=O(B)$ is supported on
$S_R\cup\{p:p\mid q\}$. For every $\delta>0$, Rankin's bound gives
\[
 \#\{t:q(t)=q,\ 0<|t|\ll B\}
 \ll B^\delta
   \prod_{p\in S_R\ \text{or}\ p\mid q}(1-p^{-\delta})^{-1}
 \ll_{\kappa,\delta} B^\delta q^\delta
 \ll_{\kappa,\delta} B^{2\delta}.
\]
For the middle inequality, remove the finitely many primes with
$p^{-\delta}>1/2$ and bound the remaining factors by $2^{\omega(q)}
\ll_\delta q^\delta$. Renaming the small exponent gives
$O_{\kappa,\alpha}(B^\alpha)$ choices of $t$.

Apply \eqref{eq:v5-lattice-tail} with $m=3$ to $q\mid R$.
If $\lambda_1\leq\lambda_2\leq\lambda_3$ are the successive minima,
then $\lambda_3\leq q\ll B$ and
$\lambda_1\lambda_2\lambda_3\asymp q^2$.
A basis comparable to these minima therefore gives
$O(\prod_i(1+B/\lambda_i))=O(B^3/q^2)$ points in the box.
The lattices with $\lambda_3>T$ contribute at most
\[
 B^\alpha\sum_{q\ll B}
 O\left(9^{\omega(q)}q^3T^{-3}\frac{B^3}{q^2}\right)
 \ll_{\kappa,\alpha} B^{5+2\alpha}T^{-3}.
\]
For the other lattices choose a basis with lengths comparable to the
successive minima. The inverse basis matrix bounds the new coordinates by
$O(B/\lambda_i)$, so the resulting rectangular box has volume
$O(B^3/q^2)$ and minimum side parameter $\gg B/T$.
Enlarge its side parameters to at least $2$, as required in
\eqref{eq:v5-box-sieve}; since $\lambda_i\ll B$, this changes volume
only by a fixed factor.
The substitution into the polynomial pair is invertible over $\mathbb Q$.
For $\ell\nmid2t$ its determinant, of absolute value $q^2$, is invertible
modulo $\ell$ too. Thus no new common content prime outside $2t$ appears.
Remove common content supported on $2t$, which cannot affect the primes
being counted. Coefficient heights are powers of $B$ with fixed exponents.
Formula \eqref{eq:v5-box-sieve} bounds this lattice's contribution by
\[
 O\left(\frac{B^3\log B}{q^2L\log L}
                 +\frac{B^2T\log B}{q^2}\right).
\]
There are at most $9^{\omega(q)}q$ lattices. Since
$9^{\omega(q)}\ll_\alpha q^\alpha$, summing over $q\ll B$ and the
$O_{\kappa,\alpha}(B^\alpha)$ values of $t$ gives
$O_{\kappa,\alpha}(B^{3+3\alpha}/(L\log L)+B^{2+3\alpha}T)$.
Choose $T=B^{3/4}$ and rename the small exponent to obtain
\eqref{eq:v5-large-sieve}.
\end{proof}

\begin{proposition}[Weighted obstruction tail]\label{prop:v5-obstruction-tail}
For fixed distinct $c_j$ there are $L_*\geq2$ and $\tau_0,C>0$
such that, for every fixed $0<\tau\leq\tau_0$, uniformly for $L\geq L_*$,
the logarithmically weighted number of pairs
\[
 N/2\leq u\leq N,\qquad 1\leq w\leq\tau\sqrt N,
 \qquad p_j=u+c_jw^2\text{ all prime},
\]
having any of the odd common divisors forbidden in
Theorem~\ref{thm:v5-simultaneous-primes} greater than $L$, is at most
\begin{equation}\label{eq:v5-obstruction-tail}
 \frac{C\tau N^{3/2}}{L\log L}
       +O_\tau(N^{3/2-1/800}).
\end{equation}
The weight is $\prod_j\log p_j$, and the bound allows every Gaussian
generator, so it also holds for any prescribed choice of generators.
\end{proposition}
\begin{proof}
Take $B=\sqrt N$ and $L_0=B^{1/100}$, and choose $\tau_0$ so that
every prime value in the displayed region is between fixed positive
multiples of $B^2$. Thus $\prod_j\log p_j\ll (\log B)^k$.
Choose $L_*$ above the fixed exceptional primes in the medium-prime bound.
When $L\leq L_0$, that bound gives
\[
 \frac{C\tau B^3}{L\log L}+O_\tau(B^{27/10}(\log B)^k)
\]
for $L<\ell\leq L_0$. Its proof is uniform in this interval of $L$:
the remainder sum is bounded by the sum over all $\ell\leq L_0$.
When $L>L_0$ this part is empty.

For distinct positions, apply Lemma~\ref{lem:v5-large-sieve} at $L_0$
with $\alpha=1/1000$. After multiplication by the weight, the bound is
\[
 O\left(\left(
   \frac{B^{3+1/1000}}{B^{1/100}\log L_0}
       +B^{11/4+1/1000}\right)(\log B)^k\right)
 =O(B^{3-1/200}).
\]
This bounds all the remaining primes, also when $L>L_0$.
At one position, an odd divisor of both adjacent norms divides
$a$ and $b^2+1$. The coprime pair $a,b^2+1$ in variables $(a,b,w)$,
inserted in \eqref{eq:v5-box-sieve} at $L_0$, gives
\[
 O\left(\left(
 \frac{B^3\log B}{L_0\log L_0}+B^2\log B\right)(\log B)^k\right)
 =O(B^{3-1/200}).
\]
No factor $\tau$ is needed in these power-saving errors.

For a fixed pair of positions, the vector $(a,b,c,d,w)$ determines
$u=a^2+b^2-c_jw^2$; at one position $(a,b,w)$ does so.
Thus the vector counts bound the pattern counts. There are finitely many
pairs of positions and signs. Each split rational prime has eight
Gaussian elements of that norm, so counting all generator choices
changes the bound by at most a fixed factor. Every vector used in the
medium-prime count is primitive, since a common integer divisor greater
than one would have its square divide a prime norm.
All logarithmic factors in the remainders are absorbed for sufficiently
large $B$, with thresholds depending only on the fixed data and $\tau$.
Weakening the error to $B^{3-1/400}=N^{3/2-1/800}$ proves the claim.
\end{proof}

\subsection{Gaussian ray weights and their progression main terms}
\label{subsec:gaussian-progression-mainterms}

Let $E=\Q(i)$ and let $H$ be a fixed positive integer divisible by $16$.
Let $K/E$ be the ray class field for $H\Z[i]$.  This field is Galois over
$\Q$, since complex conjugation preserves both the modulus and its
principal ray subgroup.  Put
\[
 A=\operatorname{Gal}(K/E)
   \simeq (\Z[i]/H\Z[i])^\times/\{1,-1,i,-i\},
 \qquad G=\operatorname{Gal}(K/\Q).
\]
We use arithmetic Frobenius throughout.  Let $C\subseteq A$ be invariant
under complex conjugation, so that $C$ is a union of conjugacy classes in
$G$.  Define $\Lambda'_C(p)=\log p$ for unramified rational primes whose
Frobenius lies in $C$, and put $\Lambda'_C(n)=0$ otherwise.  In particular
this function is supported on rational primes which split in $E$.

For ideals coprime to $H$ let $[\mathfrak a]\in A$ denote their ray class,
and set
\[
 T_C(n)=\sum_{\substack{\mathrm N\mathfrak a=n\\(\mathfrak a,H)=1}}
       \Lambda_E(\mathfrak a)1_C([\mathfrak a]).
\]
Here $\Lambda_E(\mathfrak p^j)=\log\mathrm N\mathfrak p$ and it is zero
on all other ideals.  Every unramified split rational prime has two
conjugate prime ideals of norm $p$, so
\begin{equation}\label{eq:gaussian-ray-prime-factor}
 T_C(p)=2\Lambda'_C(p)\qquad(p\nmid2H).
\end{equation}
The remaining difference between $T_C$ and $2\Lambda'_C$ is supported on
rational prime powers of exponent at least two and the finitely many
excluded primes.  Its total absolute value up to $X$ is
$O_H(\sqrt X\log^2(2X))$.  This bound follows by summing the at most two
prime ideals above each rational prime and their powers; no density
theorem is used in this comparison.

For $\chi\in\widehat A$ put
\[
 c_C(\chi)=|A|^{-1}\sum_{a\in C}\overline{\chi(a)},\qquad
 T_\chi(n)=\sum_{\mathrm N\mathfrak a=n}
                  \Lambda_E(\mathfrak a)\chi(\mathfrak a),
\]
where each character is extended by zero on ideals not coprime to $H$.
Finite Fourier inversion gives the exact identity
\begin{equation}\label{eq:gaussian-ray-fourier}
                 T_C(n)=\sum_{\chi\in\widehat A}c_C(\chi)T_\chi(n).
\end{equation}
The field and all these characters are fixed before $X$ tends to infinity.

\begin{lemma}[Fixed-field Hecke Siegel--Walfisz estimate]
\label{lem:gaussian-hecke-sw}
Let $E=\Q(i)$, let $\chi$ be a fixed finite-order ray character of $E$
with finite modulus $\mathfrak m$, and let $B>0$.
There is $c=c(\chi,B)>0$ such that, uniformly for $X\geq3$,
$1\leq Y\leq X$, positive integers $q\leq(\log X)^B$, and Dirichlet
characters $\rho$ modulo $q$,
\begin{equation}\label{eq:gaussian-hecke-sw-estimate}
 \sum_{\mathrm N\mathfrak a\leq Y}
   \Lambda_E(\mathfrak a)\chi(\mathfrak a)\rho(\mathrm N\mathfrak a)
 =r_{\chi,\rho}Y+
 O_{\chi,B}\bigl(X\exp(-c\sqrt{\log X})\bigr).
\end{equation}
Here $\Lambda_E(\mathfrak p^j)=\log\mathrm N\mathfrak p$ for a prime
ideal $\mathfrak p$ and $j\geq1$, and it is zero otherwise; characters
are extended by zero on ideals not coprime to their moduli.
The number $r_{\chi,\rho}$ is one if the primitive Hecke character
inducing $\chi(\rho\circ\mathrm N)$ is trivial, and zero otherwise.
There is also $c_B>0$ such that, with the same ranges of $X,Y,q,\rho$,
\begin{equation}\label{eq:rational-sw-estimate}
 \sum_{n\leq Y}\Lambda(n)\rho(n)
 =1_{\rho\ \mathrm{principal}}Y+
 O_B\bigl(X\exp(-c_B\sqrt{\log X})\bigr).
\end{equation}
In particular, either error is $O(X(\log X)^{-D})$ for every fixed
$D>0$, with the corresponding dependence on $D$ allowed.
\end{lemma}

\begin{proof}
We use the published fixed-field estimate of Mitsui
\cite[\S1, Lemma~15, p.~12]{Mitsui1956}.
In his notation, for a fixed number field $K$, a fixed $A>0$, and any
ray character $\vartheta$ whose finite modulus has norm at most
$(\log x)^A$, it states
\begin{equation}\label{eq:mitsui-weighted-primes}
 \sum_{\mathrm N\mathfrak p\leq x}
       \vartheta(\mathfrak p)\log\mathrm N\mathfrak p
 =E(\vartheta)x+
 O_{K,A}\bigl(x\exp(-c_{K,A}\sqrt{\log x})\bigr),
\end{equation}
where $E(\vartheta)=1$ for the principal character and is zero otherwise.
The sum is over prime ideals of every residue degree.
Mitsui defines these characters and $E(\vartheta)$ on p.~2, includes
imprimitive characters, and states the uniformity in the character
and its modulus immediately before Lemma~14 on p.~11.
Infinite primes may occur in the ray modulus; only its finite part
enters the norm bound.

For completeness, the elementary passage from Mitsui's proved
smoothed estimate (Lemma~14, pp.~11--12) to
\eqref{eq:mitsui-weighted-primes} also applies to complex characters.
Put
\[
 P(x,\vartheta)=\sum_{\mathrm N\mathfrak p\leq x}
  \vartheta(\mathfrak p)\log\mathrm N\mathfrak p
                         \log\frac{x}{\mathrm N\mathfrak p}.
\]
Lemma~14 gives $P(x,\vartheta)=E(\vartheta)x+
O_{K,A}(xe^{-c_1\sqrt{\log x}})$.
Subtract its values at $xe^h$ and $x$, where
$h=e^{-c_1\sqrt{\log x}/2}$, and divide by $h$.
The resulting left side is the prime sum in
\eqref{eq:mitsui-weighted-primes} plus boundary terms of absolute
value at most
\[
 \sum_{x<\mathrm N\mathfrak p\leq xe^h}\log\mathrm N\mathfrak p
       \ll_K(xh+1)\log x.
\]
Indeed, for each integer norm there are at most $[K:\Q]$ prime ideals.
The main term is $E(\vartheta)x(e^h-1)/h=E(\vartheta)x+O(xh)$,
and the error divided by $h$ is
$O_{K,A}(xe^{-c_1\sqrt{\log x}/2})$.
Reducing the positive exponent constant proves
\eqref{eq:mitsui-weighted-primes}.

Prime powers of exponent at least two contribute
$O_K(\sqrt x\log^2(2x))$: sum over exponents $j\geq2$ and use the
elementary bound $O_K(x^{1/j})$ for the number of prime ideals of norm
at most $x^{1/j}$.
Consequently \eqref{eq:mitsui-weighted-primes} holds with the prime
sum replaced by
$\sum_{\mathrm N\mathfrak a\leq x}
\Lambda_K(\mathfrak a)\vartheta(\mathfrak a)$.

The character $\rho\circ\mathrm N$ on $E$ has ray modulus $(q)$:
if $\alpha\equiv1\pmod{q\Z[i]}$, then
$\mathrm N\alpha\equiv1\pmod q$.
Thus $\chi(\rho\circ\mathrm N)$ has a modulus dividing
$\mathfrak m(q)$, whose norm is $\mathrm N\mathfrak m\,q^2$.
For $\sqrt X\leq Y\leq X$, the bounds
$q\leq(\log X)^B$ and $\log Y\geq\frac12\log X$ imply
\[
 \mathrm N\mathfrak m\,q^2\leq(\log Y)^{2B+1}
\]
for all sufficiently large $X$, with the threshold depending only
on $\mathfrak m$ and $B$.
Apply the preceding estimate with $K=E$, $x=Y$, and $A=2B+1$.
Its error is at most the error in
\eqref{eq:gaussian-hecke-sw-estimate}, after reducing $c$.
For $1\leq Y<\sqrt X$, the elementary bound
$\sum_{\mathrm N\mathfrak a\leq Y}\Lambda_E(\mathfrak a)
\ll_E Y\log^2(2Y)$ suffices, and the remaining bounded values of
$X$ are absorbed in the implied constant.

The principal character at the chosen modulus is exactly the
character induced by the trivial primitive character, giving the
stated $r_{\chi,\rho}$.
Equivalently one can work at the primitive conductor and remove
the extra Euler factors: the corresponding prime-power towers
contribute only
$O_\chi(\log(q+2)\log(Y+2))$, which is absorbed in the error.

For \eqref{eq:rational-sw-estimate}, take $K=\Q$ and use the ray
character attached to $\rho$, allowing the real infinite place
for odd characters.  Its finite modulus norm is $q$.
The same proof uses $A=B+1$ for $Y\geq\sqrt X$ and the elementary
bound for smaller $Y$.  Finally,
$e^{-c\sqrt{\log X}}\ll_{c,D}(\log X)^{-D}$ for every fixed $D>0$.
\end{proof}


The preceding lemma gives the following consequence, including $Y=0$ by the empty-sum convention.
  For every $B,D>0$, every fixed finite-order
ray character $\chi$ of $E$, and every Dirichlet character $\rho$ of
modulus $q\leq(\log X)^B$,
\begin{equation}\label{eq:gaussian-hecke-sw-input}
 \sum_{\mathrm N\mathfrak a\leq Y}
       \Lambda_E(\mathfrak a)\chi(\mathfrak a)
                         \rho(\mathrm N\mathfrak a)
  = r_{\chi,\rho}Y+O_{H,B,D}(X(\log X)^{-D})
 \quad(0\leq Y\leq X),
\end{equation}
where $r_{\chi,\rho}=1$ if the induced Hecke character
$\chi(\rho\circ\mathrm N)$ is trivial and is zero otherwise.
The omitted Euler factors only change the sum by a polylogarithmic
quantity.  The constant may be ineffective.  We require uniformity in
$\rho,q,Y$, but not in $H$.

\begin{lemma}[Progression density from the fixed-field input]
\label{lem:gaussian-ray-progression}
For $(a,q)=1$ define
$J_q=K\cap\Q(\zeta_q)$ and let $\sigma_a(\zeta_q)=\zeta_q^a$.
For fixed $B,D>0$, uniformly for $q\leq(\log X)^B$ and $0\leq Y\leq X$,
\[
 \sum_{\substack{n\leq Y\\ n\equiv a\pmod q}}\Lambda'_C(n)
 =\frac{\#\{g\in C:g|_{J_q}=\sigma_a|_{J_q}\}}
             {[K\Q(\zeta_q):\Q]}Y
       +O_{H,B,D}(X(\log X)^{-D}).
\]
\end{lemma}
\begin{proof}
Orthogonality of Dirichlet characters and
\eqref{eq:gaussian-ray-fourier} show that the coefficient of $Y$ for
$T_C$ is
\[
 \frac1{\varphi(q)}\sum_{\chi\in\widehat A}c_C(\chi)
       \sum_{\rho\bmod q}\overline{\rho(a)}r_{\chi,\rho}.
\]
The finite sum over $\chi$ has size depending only on $H$, while the
average over $\rho$ cancels its own factor $\varphi(q)$ in the error.

To evaluate this coefficient without a convention about inverse Artin
symbols, put $L=K\Q(\zeta_q)$ and $B_q=\operatorname{Gal}(L/E)$.
Restriction identifies $B_q$ with
\[
 \{(g,s)\in A\times\operatorname{Gal}(\Q(\zeta_q)/\Q):
                       g|_{J_q}=s|_{J_q}\}.
\]
The characters $\chi(g)\rho(s)$ restrict to characters of $B_q$.
Their average on $B_q$ is one precisely when their pullback by Artin
reciprocity is the trivial Hecke character; otherwise it is zero.
This is exactly $r_{\chi,\rho}$.  Expanding $1_C(g)$ and
$1_{s=\sigma_a}$ by Fourier inversion therefore identifies the displayed
coefficient with the proportion of elements of $B_q$ satisfying
$g\in C$ and $s=\sigma_a$.  Its numerator is the numerator in the
lemma and its denominator is $|B_q|=[L:E]$.
Dividing by two using \eqref{eq:gaussian-ray-prime-factor} changes the
denominator to $[L:\Q]$.  The prime-power error is smaller than the stated
error.  This proves the formula.
\end{proof}

Choose a fixed integer $W$ divisible by the conductor of $K^{\mathrm{ab}}$.
It is in particular divisible by $4$.
For a unit $b\pmod W$ put
\[
\begin{aligned}
 \mu_C(b)&=
 \frac{\#\{g\in C:g|_{K^{\mathrm{ab}}}=\sigma_b|_{K^{\mathrm{ab}}}\}}
      {[K:K^{\mathrm{ab}}]},\\
 f_{C,b,W}(n)&=\frac{\varphi(W)}W\Lambda'_C(Wn+b),\\
 \lambda_{b,W}(n)&=\frac{\varphi(W)}W\Lambda'(Wn+b).
\end{aligned}
\]
The prime-supported ordinary von Mangoldt function is denoted by
$\Lambda'$.  Replacing it by the usual $\Lambda$ costs only the same
negligible prime-power error.
The centered weight is
\begin{equation}\label{eq:gaussian-centered-weight}
       h_{C,b,W}=f_{C,b,W}-\mu_C(b)\lambda_{b,W}.
\end{equation}
Here $0\leq\mu_C(b)\leq1$, so
$|h_{C,b,W}(n)|\leq\lambda_{b,W}(n)=O_W(\log(2N))$ on $[N]$.

\begin{corollary}[Centered sums on progressions]
\label{cor:gaussian-centered-progressions}
For every $D>0$,
\[
       \sum_{n\in P}h_{C,b,W}(n)=O_{H,W,D}(N(\log N)^{-D})
\]
uniformly over arithmetic progressions $P$ contained in $[N]$.
\end{corollary}
\begin{proof}
Write $P$ as a residue class modulo $q$ intersected with an interval.
If $q>(\log N)^{D+3}$, its cardinality is at most
$N/(\log N)^{D+3}+1$, and the pointwise logarithmic bound proves the
claim directly.  Otherwise the condition $n\in P$ becomes one residue
class $a\pmod{Wq}$ and an interval for $Wn+b$.

If $(a,Wq)>1$, every prime in this residue class divides $Wq$.  There
are at most $O(\log(Wq))$ such primes, with total logarithmic weight
polylogarithmic in $N$, which is negligible here.  Suppose instead that
$(a,Wq)=1$.
Since $\Q(\zeta_{Wq})$ is abelian and already contains
$K^{\mathrm{ab}}$, one has
$K\cap\Q(\zeta_{Wq})=K^{\mathrm{ab}}$.
Also $a\equiv b\pmod W$, so the lemma's main coefficient is
\[
   \frac{\mu_C(b)}{\varphi(Wq)}.
\]
For the ordinary prime weight, the Siegel--Walfisz main coefficient
is $1/\varphi(Wq)$.  The two terms of
\eqref{eq:gaussian-centered-weight} therefore have equal main terms.
Apply the progression lemma at both interval endpoints, with an error
exponent larger than $D$ and with $X=O_W(N)$, and subtract.  This proves
the claimed bound.  All moduli used are at most a fixed power of
$\log N$, and $W$ is fixed throughout this argument.
\end{proof}


% Standalone section fragment; needs amsmath, amssymb, amsthm.
% Bibliography keys: GTOrbits2012 and GTMobius2012.
\subsection{Transfer from Gaussian ideal norms to polynomial nilsequences}
\label{sec:gn-transfer}

Throughout this subsection $E=\mathbb Q(i)$, $\mathcal O_E=\mathbb Z[i]$,
and $H_0$ is a fixed positive integer. A ray character $\chi$ modulo
$H_0\mathcal O_E$ is extended by zero to ideals not coprime to
$H_0\mathcal O_E$. The trivial character on all ideals is included by
taking $H_0=1$. We use the conventions of \cite{GTOrbits2012}: $G$ is
connected, simply connected and nilpotent, $\Gamma$ is a discrete cocompact
subgroup, and the filtration consists of closed connected rational
subgroups. Lipschitz norms include the supremum norm.

\begin{lemma}[Gaussian norm transfer]\label{lem:gn-norm-transfer}
Fix positive integers $m,d$. There are $\kappa,\sigma>0$, depending only
on $m,d$, and constants $C,X_0,\delta_0>0$, depending on $H_0,m,d$, with
the following property. Suppose
\[
 X\geq X_0,\qquad X^{-\sigma}\leq\delta\leq\delta_0,
 \qquad 1\leq Q\leq\delta^{-\kappa},\qquad
 1\leq q\leq\delta^{-\kappa},\quad q\in\mathbb Z.
\]
Let $G/\Gamma$ have dimension $m$, a filtration $G_\bullet$ of degree at
most $d$, and a $Q$-rational adapted Mal'cev basis. Suppose
$g\in\operatorname{poly}(\mathbb Z,G_\bullet)$ is totally
$\delta$-equidistributed on $[X]$. If $F:G/\Gamma\to\mathbb C$ is
Lipschitz and $\int F=0$, then for every $0\leq A<B\leq X$ and every
$b\in\mathbb Z$,
\begin{equation}\label{eq:gn-transfer}
 \left|\sum_{\substack{A<\mathrm N\mathfrak a\leq B\\
                 \mathrm N\mathfrak a\equiv b\pmod q}}
       \chi(\mathfrak a)F(g(\mathrm N\mathfrak a)\Gamma)\right|
 \leq C X\delta^\kappa\|F\|_{\mathrm{Lip}}.
\end{equation}
There is no lower bound on $B-A$ or on the number of progression terms.
Inserting $\log\mathrm N\mathfrak a$ on the left costs at most an
additional factor $O(\log(2X))$.
\end{lemma}

We prove the lemma using only the one-parameter quantitative Leibman
theorem \cite[Theorem 2.9]{GTOrbits2012}, together with elementary
polynomial and lattice arguments. No multiparameter equidistribution
theorem is used.

\begin{lemma}[Quadratic coefficient extrapolation]\label{lem:gn-quadratic-coefficients}
Let $X\geq1$, $P(t)=a_0+a_1t+a_2t^2\in\mathbb Z[t]$, $a_2\ne0$, and suppose
\[
 |a_i|\leq B_0X^{1-i/2},\qquad
 L\geq X^{1/2}/B_0,\qquad B_0\geq2.
\]
For every scalar polynomial $p:\mathbb Z\to\mathbb R/\mathbb Z$ of
degree at most $d$, there is an integer $1\leq k\leq C_dB_0^{C_d}$ such
that
\[
 \|kp\|_{C^\infty[X]}
   \leq C_dB_0^{C_d}\|p\circ P\|_{C^\infty[L]}.
\]
Here, for $p(n)=\sum_j\alpha_j\binom nj$, we use
$\|p\|_{C^\infty[X]}=\max_{1\leq j\leq d}X^j
\|\alpha_j\|_{\mathbb R/\mathbb Z}$.
\end{lemma}
\begin{proof}
Write $\binom{P(t)}j=\sum_\ell s_{j,\ell}\binom t\ell$. These
coefficients are integers since the left side is integer-valued. Its
degree is $2j$, and comparison of leading coefficients gives
\[
 s_{j,2j}=\frac{(2j)!}{j!}a_2^j\ne0,
 \qquad |s_{j,\ell}|\ll_d B_0^jX^{j-\ell/2}.
\]
For the bound, expand powers of $P$ and change between the monomial and
binomial bases. The change-of-basis matrices have fixed rational
coefficients. Descending induction in $\ell$ shows that every term
contributing to $s_{j,\ell}$ has weighted input degree at least $\ell$,
which gives the displayed estimate.

If $p=\sum_j\alpha_j\binom nj$ and
$p\circ P=\sum_\ell\beta_\ell\binom t\ell$, then
\[
 \beta_{2j}=\sum_{h=j}^d s_{h,2j}\alpha_h.
\]
Put $k_j=\prod_{h=j}^d s_{h,2h}$. Descending induction in $j$ gives
\[
 X^j\|k_j\alpha_j\|_{\mathbb R/\mathbb Z}
   \ll_d B_0^{O_d(1)}\|p\circ P\|_{C^\infty[L]}.
\]
Indeed the $\beta_{2j}$ term is bounded using
$X^jL^{-2j}\leq B_0^{2j}$, and a higher coefficient is multiplied by
$s_{h,2j}=O_d(B_0^hX^{h-j})$, canceling its $X^{-h}$ error.
All multipliers used to clear higher rows are integers of size
$B_0^{O_d(1)}$. Multiplying to the common nonzero integer $|k_1|$ proves
the assertion. The constant coefficient $\alpha_0$ is not used.
\end{proof}

\begin{lemma}[Small horizontal character obstruction]\label{lem:gn-horizontal-obstruction}
There are $c_0,\sigma_0>0$, depending only on $m,d$, such that the
following holds for sufficiently large $X$ and sufficiently small
$\delta\geq X^{-\sigma_0}$. Suppose the basis is $T$-rational and
a nontrivial horizontal character $\eta$ satisfies
\[
 |\eta|\leq T,\qquad \|\eta\circ g\|_{C^\infty[X]}\leq T,
 \qquad T\leq\delta^{-c_0}.
\]
Then $g$ is not totally $\delta$-equidistributed on $[X]$.
\end{lemma}
\begin{proof}
Choose representatives of the nonconstant binomial coefficients of
$\eta\circ g$ satisfying $|\alpha_j|\leq TX^{-j}$. On the initial
interval of length $\lfloor X/(100dT)\rfloor$ its values lie in an arc
of length less than $1/10$, translated by its constant coefficient.
Choose a Lipschitz function on $\mathbb R/\mathbb Z$ which equals one on
that arc, has mean zero and bounded Lipschitz norm. Its composition with
$\eta$ has Lipschitz norm at most $T^{C(m,d)}$, by the Mal'cev metric
estimates in \cite[Appendix A]{GTOrbits2012}. Its Haar mean is zero:
a nontrivial horizontal character maps Haar measure onto Haar measure
on the circle. Choose $c_0$ with $C(m,d)c_0<1/4$ and $c_0<1/4$.
The initial interval is longer than $\delta X$, while the discrepancy
one exceeds $\delta T^{C(m,d)}$. Taking $\sigma_0$ sufficiently small
absorbs the floors. This contradicts total equidistribution.
\end{proof}

\begin{proof}[Proof of Lemma~\ref{lem:gn-norm-transfer}]
Fix a constant $C_*$ which dominates the exponents in
Lemma~\ref{lem:gn-quadratic-coefficients}, the degree-$2d$ case of
\cite[Theorem 2.9]{GTOrbits2012}, and the preceding obstruction.
Choose
\[
 0<\kappa<\min\{(1000C_*^2)^{-1},c_0/(12C_*)\},\qquad
 0<\sigma<\min\{(1000C_*^2)^{-1},\sigma_0\},
 \qquad 2\kappa\sigma<1/4.
\]
These choices depend only on $m,d$. Decrease $\delta_0$, depending on
$H_0,m,d$, whenever a fixed multiplicative constant must be absorbed.

We first record a pullback consequence. Suppose $P$ and $L$ satisfy
Lemma~\ref{lem:gn-quadratic-coefficients} with
$B_0\leq\delta^{-5\kappa}$. Then
$(g(P(t))\Gamma)_{t\in[L]}$ is $\delta^\kappa$-equidistributed.
Otherwise the quantitative Leibman theorem gives a nontrivial $\eta$
with
\[
 |\eta|,\ \|\eta\circ g\circ P\|_{C^\infty[L]}
     \ll\delta^{-C_*\kappa}.
\]
The basis is $\delta^{-\kappa}$-rational as required. The composed
sequence has degree at most $2d$ for the stretched filtration
$\widetilde G_j=G_{\lceil j/2\rceil}$, with the same adapted basis.
This follows from the elementary one-parameter coordinate
characterization in \cite[Lemma 6.7]{GTOrbits2012}.
Extrapolation gives a nonzero integer multiple $k\eta$ with
\[
 |k\eta|,\ \|k\eta\circ g\|_{C^\infty[X]}
       \ll\delta^{-6C_*\kappa},
\]
with $C_*$ chosen sufficiently large initially. It remains nontrivial. The choice
$6C_*\kappa<c_0/2$ and a reduction of $\delta_0$ make this contradict
Lemma~\ref{lem:gn-horizontal-obstruction}. Only ordinary
equidistribution of this pullback is needed.

Every nonzero Gaussian ideal has exactly four generators. The function
$z\mapsto\chi((z))$, extended by zero when $z$ is not coprime to $H_0$,
is periodic modulo $H_0\mathbb Z[i]$ and invariant under units. Thus the
sum in \eqref{eq:gn-transfer} is exactly
\[
 \frac14\sum_{\substack{z=x+iy\ne0\\A<x^2+y^2\leq B\\
                         x^2+y^2\equiv b\pmod q}}
       \chi((z))F(g(x^2+y^2)\Gamma).
\]
The factor $1/4$ also applies to ideals with generators on an axis;
we have introduced no angular fundamental domain.

Set $\rho=\delta^\kappa$ and
$M=\lfloor\rho\sqrt X/H_0\rfloor$. On each horizontal row $y\in
\mathbb Z$, and in each residue class of $x$ modulo $H_0$, partition
the points into consecutive blocks of $M$ points with step $H_0$.
The character is constant on a block. The intersection of the annulus
with a horizontal row consists of at most two intervals, so at most
four blocks per residue class meet an endpoint. There are $O(\sqrt X)$
relevant rows. Consequently the boundary blocks contain
\[
 O(H_0\sqrt X\,M)=O(\rho X)
\]
points. Estimate their contribution trivially. Excluding $z=0$, if
necessary, costs one more bounded summand, absorbed by $\rho X$.

Every remaining contributing block lies wholly in the annulus. Split
its index modulo $q$. The norm congruence is constant on each resulting
class. Such a class has the form
\[
 x=x_0+H_0qn,\quad 1\leq n\leq L,\qquad
 P(n)=(x_0+H_0qn)^2+y^2,
 \qquad L\geq\frac{\rho\sqrt X}{4H_0q}.
\]
The hypotheses on $\delta$ and $\sigma$ ensure $M\geq2q$ and absorb
the floors for large $X$. A harmless endpoint adjustment gives
$|x_0|\leq2\sqrt X$. The coefficients of $P$ satisfy
\[
 |a_0|\ll X,\qquad |a_1|\ll H_0q\sqrt X,\qquad a_2=H_0^2q^2.
\]
Since $H_0$ is fixed and $q\leq\delta^{-\kappa}$, these bounds and the
lower bound for $L$ fit the pullback consequence with
$B_0\leq\delta^{-5\kappa}$ after decreasing $\delta_0$.
Hence the sum on the class is
$O(L\delta^\kappa\|F\|_{\mathrm{Lip}})$. The classes partition the
interior points, of which there are $O(X)$; this elementary bound follows
from the enclosing square. Summing and adding the boundary contribution
proves \eqref{eq:gn-transfer}.

The number of row endpoints is independent of $B-A$, explaining why no
lower interval length is required. Finally apply Abel summation to the
uniform bound with variable upper endpoint to insert $\log\mathrm N
\mathfrak a$ at cost $O(\log(2X))$.
\end{proof}

\begin{remark}\label{rem:gn-product-and-progression}
For Type II, apply the lemma in dimension $2m$ to
$(g_1(n),g_2(n))$ and $F_1\otimes\overline{F_2}$, when this entire
product orbit is totally equidistributed. Mean-zero factors ensure that
the product test function has mean zero; they do not imply
equidistribution of the product orbit. Use the smaller exponents for
dimensions $m$ and $2m$. In dimension zero mean-zero test functions
vanish, so the assertion is immediate.

The constraints $rn\equiv b\pmod q$ and $r'n\equiv b\pmod q$ either
are inconsistent or specify one residue class modulo a divisor of $q$.
Intersecting their interval restrictions only changes $A,B$. Thus all
Type I and Type II truncations are covered by the lemma. If
$\delta=(\log X)^{-B_1}$ and $Q,q\leq(\log X)^{C_1}$, any desired
saving $X(\log X)^{-A_1}$ follows once
$B_1\kappa>\max\{C_1,A_1\}$, with $H_0$ fixed.
\end{remark}


% Depends on the separately proved Lemma lem:gn-norm-transfer.
% No Frei--Tinkova theorem is assumed in this fragment.
% Their ideal-norm Type I/II strategy is retained as methodological attribution.
\subsection{Gaussian ideal Vaughan decomposition and Type I/II estimates}
\label{gi-section}

We give the ideal decomposition and the subsequent integer-parameter
recombination in full.  The architecture is that of
\cite[Section~3]{moebius}, with the Gaussian ideal-norm adaptation
attributed to Frei--Tinkova.  The analytic inputs used here are the
separately proved Lemma~\ref{lem:gn-norm-transfer} and the published
one-parameter results \cite[Theorem~2.9 and Lemma~4.5]{GTOrbits2012}.
In particular, no theorem of Frei--Tinkova is an assumed input.

Fix the Gaussian ray modulus $H_0$.  All ideals in this subsection are
nonzero integral ideals of $\mathbb Z[i]$.  Write $\mathrm N\mathfrak a$
for the absolute ideal norm, $\mu_E$ for ideal Mobius, and
$\Lambda_E(\mathfrak p^j)=\log\mathrm N\mathfrak p$ for $j\geq1$, with
$\Lambda_E$ zero on other ideals.  A ray character $\chi$ is extended by
zero on ideals not coprime to $H_0$.  It is completely multiplicative
and satisfies $|\chi|\leq1$.  We use a Lipschitz norm that includes the
supremum norm.  All nilmanifolds have connected, simply connected
nilpotent covering group, a cocompact lattice, and an adapted rational
Mal'cev basis, with the filtration and polynomial-sequence conventions
of \cite{GTOrbits2012}.

\begin{lemma}[Ideal multiplicities and moments]\label{gi-moments}
Let $a_E(n)=\#\{\mathfrak a:\mathrm N\mathfrak a=n\}$.  For $X\geq1$,
\begin{align}
 a_E(n)&\leq\tau(n),\label{gi-multiplicity}\\
 \sum_{n\leq X}a_E(n)^2&\leq X(1+\log X)^3,&
 \sum_{n\leq X}\frac{a_E(n)^2}{n}&\leq(1+\log X)^4.\label{gi-two-moments}
\end{align}
Moreover
\begin{equation}\label{gi-ideal-count}
 A_E(X):=\#\{\mathfrak a:\mathrm N\mathfrak a\leq X\}
       =\frac{\pi}{4}X+O(\sqrt X+1).
\end{equation}
In particular $A_E(X)\ll X$, and the number of ideals of norm in
$(B/2,B]$ is comparable to $B$ for all sufficiently large $B$.
\end{lemma}
\begin{proof}
At a split rational prime $p$ one has $a_E(p^k)=k+1$; at an inert prime
it is one for even $k$ and zero for odd $k$; and $a_E(2^k)=1$.
Multiplicativity proves \eqref{gi-multiplicity}.  Prime by prime,
$(k+1)^2\leq\binom{k+3}{3}$, so $\tau(n)^2\leq d_4(n)$.
Summing over ordered quadruples with product at most $X$, and summing
the last coordinate first, bounds $\sum_{n\leq X}d_4(n)$ by
$X(1+\log X)^3$.  Bounding each of the four harmonic sums separately
gives $\sum_{n\leq X}d_4(n)/n\leq(1+\log X)^4$.
Every ideal has exactly four Gaussian generators.  Counting lattice
points in the disk of radius $\sqrt X$, and dividing by four, proves
\eqref{gi-ideal-count}: the unit squares meeting its boundary lie
between circles of radii $\sqrt X-\sqrt2$ and $\sqrt X+\sqrt2$ and
number $O(\sqrt X+1)$.
\end{proof}

\begin{lemma}[Exact ideal Vaughan identity]\label{gi-vaughan}
Let $N\geq8$, $U=V=N^{1/3}$, $L=\log(2N)$, and let
$f:\mathbb N\longrightarrow\mathbb C$ satisfy $|f|\leq1$ and vanish
outside $(N/2,N]$.  Then
\begin{equation}\label{gi-vaughan-sum}
 \sum_{\mathfrak a}\Lambda_E(\mathfrak a)\chi(\mathfrak a)
                       f(\mathrm N\mathfrak a)=T_0-T_1+T_2,
\end{equation}
where
\begin{align*}
 T_0&=\sum_{\mathrm N\mathfrak d\leq U}\mu_E(\mathfrak d)\chi(\mathfrak d)
       \sum_{\mathfrak b}\log(\mathrm N\mathfrak b)\chi(\mathfrak b)
                  f(\mathrm N\mathfrak d\,\mathrm N\mathfrak b),\\
 T_1&=\sum_{\mathrm N\mathfrak d\leq UV}\alpha_{\mathfrak d}\chi(\mathfrak d)
       \sum_{\mathfrak b}\chi(\mathfrak b)
                  f(\mathrm N\mathfrak d\,\mathrm N\mathfrak b),\\
 T_2&=\sum_{V<\mathrm N\mathfrak d\leq N/U}\beta_{\mathfrak d}\chi(\mathfrak d)
       \sum_{\mathrm N\mathfrak b>U}\mu_E(\mathfrak b)\chi(\mathfrak b)
                  f(\mathrm N\mathfrak d\,\mathrm N\mathfrak b),
\end{align*}
with
\begin{align*}
 \alpha_{\mathfrak d}
   &=\sum_{\substack{\mathfrak b\mathfrak c=\mathfrak d\\
              \mathrm N\mathfrak b\leq U,\ \mathrm N\mathfrak c\leq V}}
                        \mu_E(\mathfrak b)\Lambda_E(\mathfrak c),&
 \beta_{\mathfrak d}
   &=\sum_{\substack{\mathfrak c\mid\mathfrak d\\\mathrm N\mathfrak c>V}}
                        \Lambda_E(\mathfrak c).
\end{align*}
Both coefficients have absolute value at most $L$.
\end{lemma}
\begin{proof}
Convolution on ideals satisfies
$\mu_E*1=\mathbf1_{\mathfrak a=(1)}$, $1*\Lambda_E=\log\mathrm N$,
and $\mu_E*\log\mathrm N=\Lambda_E$, as one checks on prime powers.
Splitting $\mu_E$ at $U$ and $\Lambda_E$ at $V$ gives the exact identity
\begin{equation}\label{gi-convolution-identity}
 \Lambda_E=\mu_{\leq U}*\log\mathrm N+\Lambda_{\leq V}
       -\mu_{\leq U}*1*\Lambda_{\leq V}
       +\mu_{>U}*1*\Lambda_{>V}.
\end{equation}
The second term vanishes on the support of $f$.  Complete multiplicativity
of $\chi$, including its zero extension, permits separation of each
convolution factor and yields \eqref{gi-vaughan-sum} with exactly the
stated cutoffs.  Finally
$\sum_{\mathfrak c\mid\mathfrak d}\Lambda_E(\mathfrak c)
=\log\mathrm N\mathfrak d$ and every summand is nonnegative.  This
proves the coefficient bounds, with no discarded endpoint or error term.
\end{proof}

\begin{proposition}[Quantitative extraction]\label{gi-extraction}
There are absolute constants $c_0>0$ and $N_0$ such that the following
holds.  Under Lemma~\ref{gi-vaughan}, suppose $N\geq N_0$,
$0<\epsilon\leq1$, and
\begin{equation}\label{gi-large-correlation}
 \left|\sum_{\mathfrak a}\Lambda_E(\mathfrak a)\chi(\mathfrak a)
                         f(\mathrm N\mathfrak a)\right|\geq\epsilon NL.
\end{equation}
Put $\theta=c_0\epsilon^4/L^{24}$.  At least one of the following holds.
\begin{enumerate}
\item Some dyadic $R\leq2N^{1/3}$ has at least $\theta R$ integers
$r\in(R/2,R]$ satisfying
\[
 \left|\sum_{\mathfrak b}\log(\mathrm N\mathfrak b)\chi(\mathfrak b)
                   f(r\mathrm N\mathfrak b)\right|\geq\theta N/r.
\]
\item Some dyadic $R\leq2N^{2/3}$ has at least $\theta R$ integers
$r\in(R/2,R]$ satisfying the same inequality without the logarithm.
\item There are dyadic $R,B$ with
\begin{equation}\label{gi-dyadic-ranges}
 N^{1/3}<R\leq2N^{2/3},\qquad N/2<RB<4N,
\end{equation}
and at least $\theta R^2$ pairs $r,r'\in(R/2,R]$ satisfying
\begin{equation}\label{gi-type-two-correlation}
 \left|\sum_{B/2<\mathrm N\mathfrak b\leq B}
       f(r\mathrm N\mathfrak b)\overline{f(r'\mathrm N\mathfrak b)}
       \right|\geq\theta B.
\end{equation}
\end{enumerate}
\end{proposition}
\begin{proof}
One of $|T_i|$ is at least $\epsilon NL/3$.  For $T_0$ put
\[
 G(r)=\frac{r}{NL}\left|\sum_{\mathfrak b}\log(\mathrm N\mathfrak b)
                   \chi(\mathfrak b)f(r\mathrm N\mathfrak b)\right|;
\]
for $T_1$ use $r/N$ times the absolute unweighted inner sum.  In either
case Lemma~\ref{gi-moments} gives $G\ll1$ and
\[
 \epsilon\ll\sum_{r\leq N^e}\frac{a_E(r)G(r)}r,
 \qquad e=1/3\text{ or }2/3.
\]
Cauchy--Schwarz and \eqref{gi-two-moments} imply
$\sum_{r\leq N^e}G(r)^2/r\gg\epsilon^2/L^4$.
There are $O(L)$ dyadic cells, so one satisfies
\[
 \sum_{R/2<r\leq R}G(r)^2\gg R\epsilon^2/L^5.
\]
Since $G\ll1$, at least $cR\epsilon^2/L^5$ terms have
$G(r)\geq c\epsilon/L^{5/2}$.  Both the count and the claimed amplitudes
follow after decreasing $c_0$.

For $T_2$, split both ideal norms into dyadic cells.  There are $O(L^2)$
rectangles, and a nonzero cell satisfies \eqref{gi-dyadic-ranges} and
$B\geq N^{1/3}/4$.  A rectangle contributes at least $c\epsilon N/L$;
we use the weaker bound $c\epsilon N/L^2$.  Grouping $\mathfrak d$ by
$r=\mathrm N\mathfrak d$ writes this contribution as
\[
 T=\sum_{R/2<r\leq R}A_r\sum_{\mathfrak b\in\mathcal B}c_{\mathfrak b}
                   f(r\mathrm N\mathfrak b),\qquad
 \mathcal B=\{\mathfrak b:B/2<\mathrm N\mathfrak b\leq B\}.
\]
Here $|A_r|\leq L a_E(r)$ and $|c_{\mathfrak b}|\leq1$; the restriction
$\mathrm N\mathfrak b>U$ belongs to $c_{\mathfrak b}$.  All product
cutoffs remain in $f$, so the domain is rectangular.  By
\eqref{gi-two-moments}, $\sum_r|A_r|^2\ll RL^5$.  Thus, with
$F_{r,\mathfrak b}=f(r\mathrm N\mathfrak b)$,
\begin{equation}\label{gi-first-cauchy}
 T_*:=\sum_r\left|\sum_{\mathfrak b\in\mathcal B}
                          c_{\mathfrak b}F_{r,\mathfrak b}\right|^2
          \gg \epsilon^2 RB^2/L^9.
\end{equation}
Put $M=|\mathcal B|\asymp B$.  Expansion and Cauchy--Schwarz in the
ordered pair $(\mathfrak b,\mathfrak b')$ give
\begin{align}
 T_*^2
 &=\left|\sum_{\mathfrak b,\mathfrak b'}
 c_{\mathfrak b}\overline{c_{\mathfrak b'}}
 \sum_r F_{r,\mathfrak b}\overline{F_{r,\mathfrak b'}}\right|^2\notag\\
 &\leq M^2\sum_{\mathfrak b,\mathfrak b'}
       \left|\sum_rF_{r,\mathfrak b}\overline{F_{r,\mathfrak b'}}\right|^2
 =M^2\sum_{r,r'}\left|\sum_{\mathfrak b}
                  F_{r,\mathfrak b}\overline{F_{r',\mathfrak b}}\right|^2.
 \label{gi-second-cauchy}
\end{align}
Consequently
\[
 \mathbb E_{r,r'}\left|\mathbb E_{\mathfrak b\in\mathcal B}
              F_{r,\mathfrak b}\overline{F_{r',\mathfrak b}}\right|^2
       \gg\epsilon^4/L^{18}.
\]
Every summand is at most one.  Thresholding gives
\eqref{gi-type-two-correlation} and its stated count, with the weaker
common parameter $\theta=c_0\epsilon^4/L^{24}$.  The character and Mobius
coefficients have disappeared.  All conjugates in
\eqref{gi-second-cauchy} are necessary.
\end{proof}

\paragraph{Progressions and interval restrictions.}\label{gi-progressions}
In the application,
\[
 f(n)=1_{N/2<n\leq N}\,1_{u<n\leq v}\,1_{n\equiv b\pmod q}
                              F(g(n)\Gamma).
\]
For Type I the congruence $rn\equiv b\pmod q$ is either insoluble or
one class modulo $q/\gcd(r,q)$.  In Type II its intersection with
$r'n\equiv b\pmod q$ is either empty or one class modulo
\[
 \operatorname{lcm}\bigl(q/\gcd(r,q),q/\gcd(r',q)\bigr),
\]
which divides $q$.  The remaining restrictions are interval intersections.
Lemma~\ref{lem:gn-norm-transfer} covers them without a lower bound on
their lengths.  On the product nilmanifold the function is
$F(x)\overline{F(y)}$, whose Haar integral is
$|\int F|^2=0$.  After normalizing $\|F\|_{\mathrm{Lip}}=1$, its
Lipschitz norm for the sum metric is at most two.  Thus all mean-zero,
modulus and interval hypotheses for product norm transfer are satisfied.

\begin{lemma}[Integer-parameter recombination]\label{gi-recombination}
For each $m,d$ there are $C_{\rm rec},K_{\rm rec}\geq2$ with the following
property.  Let $0<\eta<\eta_0$, $N\geq\eta^{-K_{\rm rec}}$, and let $g$
have degree at most $d$ on an $m$-dimensional nilmanifold with an
$\eta^{-1}$-rational adapted basis.  Each of the following implies that
$(g(n)\Gamma)_{n\leq N}$ is not totally $\eta^{C_{\rm rec}}$-equidistributed.
\begin{enumerate}
\item For at least $\eta R$ values $r\in(R/2,R]$, where
$R\leq2N^{2/3}$, the orbit $(g(rn)\Gamma)_{n\leq\lceil N/r\rceil}$
is not totally $\eta$-equidistributed.
\item For at least $\eta R^2$ pairs $r,r'\in(R/2,R]$, where
$N^{1/3}/4\leq R\leq2N^{2/3}$ and $N/2<RB<4N$, the product orbit
$(g(rn),g(r'n))_{n\leq B}$ is not totally $\eta$-equidistributed.
\end{enumerate}
\end{lemma}
\begin{proof}
All constants in this proof depend only on $m,d$.  Failure of total
equidistribution gives a progression $a+q[\ell]\subseteq[M]$ with
$\ell\geq\eta M$, $q\leq2/\eta$, and $|a|\leq2M$ on which ordinary
$\eta$-equidistribution fails.  Apply
\cite[Theorem~2.9]{GTOrbits2012}, in dimension $m$ or $2m$.  A nonzero
horizontal character has magnitude and restricted smoothness norm at
most $\eta^{-C_1}$, after shrinking $\eta_0$ to absorb fixed factors.

Here is the required removal of a possibly large shift.  Write the scalar
phase $p(a+qt)$ in monomials.  Multiplication by $d!$ clears conversion
from its binomial coefficients.  Inverting $t=(n-a)/q$ and multiplying
by $q^d$ clears every remaining denominator.  The error in the coefficient
of $n^i$ is bounded by
\begin{equation}\label{gi-shift-removal}
 C_d\eta^{-C_1}\sum_{j=i}^d q^{d-j}|a|^{j-i}\ell^{-j}
       \leq\eta^{-C_2}M^{-i},\qquad i\geq1.
\end{equation}
The integer parts disappear modulo one.  Thus a nonzero integer multiple
of the character has magnitude and $C^\infty[M]$ norm at most
$\eta^{-C_2}$.  This argument does not treat $a=O(M)$ as a bounded-height
rational parameter.  There are at most $\eta^{-C_3}$ possible resulting
characters, since their lattice has rank at most $2m$.  Pigeonholing
therefore fixes one character on a set of density at least $\eta^{C_4}$.

In Type I write
$\psi(g(n))=\sum_{j=0}^d\beta_jn^j$ modulo one.  For the selected set
$S\subseteq(R/2,R]$, coefficient conversion gives
\begin{equation}\label{gi-one-coefficients}
 \|d!\beta_jr^j\|_{\mathbb R/\mathbb Z}
       \leq\eta^{-C_5}(R/N)^j,\qquad 1\leq j\leq d.
\end{equation}
If $R\leq\eta^{-C_6}$ for a sufficiently large fixed $C_6$, choose one
$r\in S$ and multiply the character by $d!r^d$.  Multiplication of the
$j$th bound by $r^{d-j}$ gives magnitude and $C^\infty[N]$ norm bounded
by a fixed power of $\eta^{-1}$.  This treats every small dyadic scale,
including $R=1$.

Otherwise apply \cite[Lemma~4.5]{GTOrbits2012} to
$r\mapsto d!\beta_jr^j$ on $[\lceil R\rceil]$.  The selected set has
density at least a fixed power of $\eta$, and its values lie in an arc
of length at most $2\eta^{-C_5}(R/N)^j$.  Taking $K_{\rm rec}$ large,
the bound $R\leq2N^{2/3}$ makes this length at most half the density.
The choice of $C_6$ also places the recurrence length beyond the fixed
inverse density powers.  The published recurrence lemma supplies
nonzero integers $t_j$, $|t_j|\leq\eta^{-C_7}$, such that
\begin{equation}\label{gi-recurrent-coefficients}
 \|t_j(d!)^2\beta_j\|_{\mathbb R/\mathbb Z}
                    \leq\eta^{-C_8}N^{-j}.
\end{equation}
The extra factorial clears the leading-coefficient conversion.  Taking
the product of the $t_j$ and the fixed factorials gives a single nonzero
horizontal character of magnitude and $C^\infty[N]$ norm at most
$\eta^{-C_9}$.

For Type II the pigeonholed character is $\psi_1\oplus\psi_2$, not both
zero.  Write the two scalar phases with coefficients $\beta_j,\beta'_j$.
The information replacing \eqref{gi-one-coefficients} is
\begin{equation}\label{gi-two-coefficients}
 \|d!(\beta_jr^j+\beta'_jr'^j)\|_{\mathbb R/\mathbb Z}
                        \leq\eta^{-C_5}(R/N)^j.
\end{equation}
Interchange the two coordinates if necessary so that $\psi_1\ne0$.
Fix an $r'$ with at least $\eta^{C_4}R$ partners.  Equation
\eqref{gi-two-coefficients} puts $d!\beta_jr^j$ in an arc centered at
$-d!\beta'_jr'^j$.  The same published recurrence lemma permits this
arbitrary center and gives \eqref{gi-recurrent-coefficients} for
$\psi_1$.  Its large-length requirement follows from
$R\geq N^{1/3}/4$ and a sufficiently large $K_{\rm rec}$.  There is no
substitution of $(r-r_0)^j$ for $r^j-r_0^j$.

Finally let $\Psi$ denote the nonzero character just obtained.  Its
binomial expansion shows that for
$1\leq n\leq c\eta^{2C_9+4}N$ the phase $\Psi(g(n))$ stays within
$1/10$ of $\Psi(g(0))$.  Choose a mean-zero Lipschitz circle function
equal to one on that arc and compose it with $\Psi$.  Its mean is zero
because a nonzero horizontal character pushes Haar measure to circle
Haar measure.  Its Lipschitz norm is bounded by a fixed power of
$\eta^{-1}$ by \cite[Appendix~A]{GTOrbits2012}.  Choose $C_{\rm rec}$
larger than the length and Lipschitz exponents, and then enlarge
$K_{\rm rec}$ so the initial interval is nonempty.  Its average one
contradicts total $\eta^{C_{\rm rec}}$-equidistribution.  This proves
both assertions.  All exponent constants are finite combinations of
the published one-parameter constants and degree-$d$ polynomial algebra.
\end{proof}

\begin{proposition}[Power-form Gaussian ideal minor-arc bound]
\label{gi-minor-arc}
For each $m,d$ there is $c_{\rm G}>0$ with the following property.
Fix $H_0$ and a finite-order Gaussian ray character $\chi$ of modulus
$H_0$, extended by zero.  Let $G/\Gamma$ have dimension $m$, filtration
degree $d$, and a $Q$-rational adapted basis, with $Q\geq2$.  Suppose
$(g(n)\Gamma)_{n\leq N}$ is totally $\delta$-equidistributed, where
$0<\delta<1$.  If $F$ is Lipschitz and $\int F=0$, then for every
interval progression $P\subseteq[N]$ of common difference at most $Q$,
\begin{equation}\label{gi-power-bound}
 \left|\sum_{\mathrm N\mathfrak a\leq N}\Lambda_E(\mathfrak a)
       \chi(\mathfrak a)1_P(\mathrm N\mathfrak a)
       F(g(\mathrm N\mathfrak a)\Gamma)\right|
 \ll_{H_0,m,d} QN\log(2N)^7\delta^{c_{\rm G}}\|F\|_{\mathrm{Lip}}.
\end{equation}
No lower bound on the length of $P$ or on $\delta$ is required.
\end{proposition}
\begin{proof}
Normalize $\|F\|_{\mathrm{Lip}}=1$.  Take the smaller positive exponents
$\kappa,\sigma$ in Lemma~\ref{lem:gn-norm-transfer} for dimensions $m$
and $2m$, and decrease $\kappa$ to at most one.  Partial summation gives
the same transfer bound with an extra factor $\log(2M)$ for a
$\log\mathrm N\mathfrak b$ weight: the unweighted lemma holds for
every partial interval, and $\int_1^Mdt/t\leq\log(2M)$.
Let $C_{\rm rec},K_{\rm rec}$ be as in Lemma~\ref{gi-recombination},
and choose
\begin{equation}\label{gi-exponent-choice}
 a_*=\frac1{4C_{\rm rec}},\qquad
 c_* =\frac{a_*\kappa}{32},\qquad
 s_* =\min\left\{\frac1{100},\frac1{4a_*K_{\rm rec}},
                             \frac{\sigma}{12a_*}\right\}.
\end{equation}
First assume $N^{-s_*}\leq\delta\leq\delta_0$, with the fixed cutoff
$\delta_0$ sufficiently small, and restrict the sum to $(N/2,N]$.
We claim that its absolute value is
\begin{equation}\label{gi-shell-bound}
                  \ll QNL^7\delta^{c_*},\qquad L=\log(2N).
\end{equation}
The elementary bound is $O(NL)$.  Thus the claim is immediate when
$QL^6\delta^{c_*}\geq1$, and in the other case
$Q\leq\delta^{-c_*}$ and $L\leq\delta^{-c_*/6}$.
If \eqref{gi-shell-bound} fails with a sufficiently large implied
constant, Proposition~\ref{gi-extraction} applies with
$\epsilon$ a fixed large multiple of $QL^6\delta^{c_*}$; if this
exceeds the permitted range, the elementary bound already proves the
claim.  The extracted density and normalized amplitude are at least
a fixed positive constant times $\delta^{4c_*}$, hence at least
$\delta^{5c_*}$ after decreasing $\delta_0$.

Set $\eta=\delta^{a_*}$.  If an extracted orbit were totally
$\eta$-equidistributed, Lemma~\ref{lem:gn-norm-transfer}, with the
progression analysis above, would bound the extracted sum by
\[
 O(M\delta^{a_*\kappa}L^e),\qquad e=0\text{ or }1.
\]
Here $M=\lceil N/r\rceil\geq N^{1/3}/2$ in Type I and
$M=B\geq N^{1/3}/4$ in Type II.  The rationality and modulus hypotheses
hold because $c_*<a_*\kappa$.  The lower-parameter hypothesis holds
because $a_*s_*\leq\sigma/12$; increasing the fixed threshold for $N$
absorbs the constants two and four.  Since $a_*\kappa=32c_*$, the
displayed bound is smaller than $M\delta^{5c_*}$, even for $e=1$.
This is a contradiction.  Thus every extracted orbit fails total
$\eta$-equidistribution.  Its density is at least $\eta$, since
$a_*>5c_*$.  Also $N\geq\eta^{-K_{\rm rec}}$ by
$a_*s_*K_{\rm rec}\leq1/4$, and $Q\leq\eta^{-1}$.
Lemma~\ref{gi-recombination} now contradicts total
$\delta$-equidistribution, because
$\eta^{C_{\rm rec}}=\delta^{1/4}>\delta$.  This proves
\eqref{gi-shell-bound}.  Fixed values $\delta\geq\delta_0$ are covered
by the elementary bound with a larger constant.

We spell out the passage through all scales.  Use exact real endpoints
$M_j=2^{-j}N$ and retain shells with $M_j\geq2\sqrt\delta N$.
The discarded norms are at most $2\sqrt\delta N$ and their total
absolute weight is $O(\sqrt\delta NL)$ by Lemma~\ref{gi-moments}.
There is no rounding of ideal-norm cutoffs.  On $[\lfloor M_j\rfloor]$
the orbit is totally $2\sqrt\delta$-equidistributed: every progression
of the requisite length has at least $\delta N$ terms for sufficiently
large $N$, and the original discrepancy bound is stronger.
For $\delta\geq N^{-s_*/4}$ the lower-parameter condition for the shell
bound holds on every retained scale, since
$M_j\geq N^{1-s_*/8}$ up to fixed constants.  The progression is merely
intersected with the shell; its common difference remains at most $Q$.
Summing the shell estimates and using $\sum_jM_j\leq2N$ gives
\begin{equation}\label{gi-full-intermediate}
                  O(QNL^7\delta^{c_*/2}).
\end{equation}
For $1/N<\delta<N^{-s_*/4}$ replace $\delta$ by
$\delta^{s_*/8}\geq N^{-s_*/8}$.  Total equidistribution at the
original parameter implies it at this larger one.  The exponent in
\eqref{gi-full-intermediate} becomes $c_*s_*/16$.
For $\delta\leq1/N$, singleton progressions imply
$|F(g(n)\Gamma)|\leq\delta$ at each integer $n$, so the elementary
$O(NL)$ weight bound suffices.  All remaining bounded $N$ are absorbed
by the fixed constant.  Thus one may take
\begin{equation}\label{gi-final-exponent}
                         c_{\rm G}=c_*s_*/16>0.
\end{equation}
Its dependence is only on $m,d$ and the corresponding norm-transfer
exponents; $H_0$ enters constants and thresholds, not this exponent.
\end{proof}

\begin{corollary}[Ordinary prime term]\label{gi-ordinary-primes}
Under the nilmanifold hypotheses of Proposition~\ref{gi-minor-arc},
there is $c_{\mathbb Q}>0$, depending only on $m,d$, such that
\begin{equation}\label{gi-ordinary-bound}
 \left|\sum_{n\leq N}\Lambda'(n)1_P(n)F(g(n)\Gamma)\right|
 \ll_{m,d} QN\log(2N)^7\delta^{c_{\mathbb Q}}\|F\|_{\mathrm{Lip}},
\end{equation}
where $\Lambda'(p)=\log p$ on rational primes and is zero otherwise.
The same bound holds for ordinary von Mangoldt $\Lambda$.
\end{corollary}
\begin{proof}
For positive integers the convolution identities
$\mu*1=1_{n=1}$, $1*\Lambda=\log$, and $\mu*\log=\Lambda$ give
\eqref{gi-convolution-identity} with exactly the displayed norm cutoffs
replaced by integer cutoffs.  Its coefficients are bounded by
$\sum_{c\mid d}\Lambda(c)=\log d$.  Integer multiplicity is one,
so both moment bounds in Lemma~\ref{gi-moments} remain valid.  The
two Cauchy--Schwarz calculations \eqref{gi-first-cauchy} and
\eqref{gi-second-cauchy} apply to the resulting finite integer sums
and give Proposition~\ref{gi-extraction} with the same conservative
parameter $c_0\epsilon^4/L^{24}$.  Lemma~\ref{gi-recombination} concerns
only integer dilation parameters and is unchanged.

The unweighted replacement for norm transfer follows directly from
total equidistribution.  If an interval progression $J\subseteq[X]$
has $|J|\geq\zeta X$, then
\[
 \left|\sum_{n\in J}F(g(n)\Gamma)\right|
       \leq |J|\zeta\|F\|_{\mathrm{Lip}}
       \leq X\zeta\|F\|_{\mathrm{Lip}}.
\]
If $|J|<\zeta X$, its cardinality and the supremum norm give the same
bound.  This treats every modulus and every truncated interval.
Partial summation supplies the logarithmic weight; the identical
long-progression/short-progression dichotomy treats the mean-zero
product function $F\otimes\overline F$ on the product nilmanifold.
Thus the exponent choices and all shell steps in the proof of
Proposition~\ref{gi-minor-arc} apply, with any sufficiently small
fixed positive transfer exponents, to give \eqref{gi-ordinary-bound}
for $\Lambda$.  No extension of a Mobius theorem to $\Lambda$ is assumed.

Finally proper prime powers contribute at most
$O(\sqrt N\log(2N)^2\|F\|_{\mathrm{Lip}})$, by counting at most
$N^{1/k}$ possible bases for each $2\leq k\leq\log_2N$.
Shrink $c_{\mathbb Q}$ below $1/4$.  If $\delta>1/N$, this error is
absorbed because $N\delta^{c_{\mathbb Q}}\geq N^{1-c_{\mathbb Q}}$.
If $\delta\leq1/N$, singleton equidistribution and
$\sum_{n\leq N}\Lambda'(n)\leq N\log(2N)$ prove the bound directly.
This gives the assertion for $\Lambda'$ as well.
\end{proof}



% Section fragment. Requires the gn-* lemmas from norm_transfer.tex.
\subsection{Fixed-modulus affine interpolation}
\label{sec:ga-affine}

Fix integers $W\geq1$ and $0\leq b<W$. The modulus $W$ remains fixed
as the size parameter tends to infinity. We first justify the affine
interpolation needed to apply ideal-norm estimates to a $W$-tricked
sequence.

\begin{lemma}[Linear coefficient extrapolation]\label{lem:ga-linear-coefficients}
Let $X\geq1$ and let $p:\mathbb Z\to\mathbb R/\mathbb Z$ have degree at most $d$.
Suppose $P(t)=u+vt$ with $u,v\in\mathbb Z$, $v\ne0$,
$|u|\leq C X$, $|v|\leq E$, and $L\geq X/E$, where $E\geq2$.
There is an integer $1\leq k\leq C_dE^{C_d}$ such that
\[
 \|kp\|_{C^\infty[X]}
 \leq C_d(C+E)^{C_d}\|p\circ P\|_{C^\infty[L]}.
\]
\end{lemma}
\begin{proof}
Expand $\binom{u+vt}j$ in the binomial basis in $t$. All coefficients
are integers, the diagonal coefficient is $v^j$, and the coefficient in
row $\ell$ is $O_d((C+E)^jX^{j-\ell})$. Starting in degree $d$, clear
the nonzero diagonal entries successively. Their product is a nonzero
integer bounded by $E^{O_d(1)}$. At degree $j$ use
$X^j/L^j\leq E^j$; a higher-row error is multiplied by
$O_d((C+E)^hX^{h-j})$, canceling its $X^{-h}$ factor. Multiplication
to the common product of diagonal entries proves the claim, exactly
as in Lemma~\ref{lem:gn-quadratic-coefficients}.
\end{proof}

\begin{lemma}[Affine interpolation preserves quantitative distribution]
\label{lem:ga-affine-interpolation}
Fix dimension and degree bounds $m,d$. There are $a,\sigma>0$, depending
only on $m,d$, and thresholds depending on $W,m,d$, with the following
property. Suppose a nilmanifold and its rational filtration satisfy the
Green--Tao conventions, its adapted basis is $Q$-rational, and
$g\in\operatorname{poly}(\mathbb Z,G_\bullet)$. Suppose also that
\[
 N^{-\sigma}\leq\delta\leq\delta_0(W,m,d),\qquad
 Q\leq\delta^{-a},
\]
and that $g$ is totally $\delta$-equidistributed on $[N]$. There is
$\widetilde g\in\operatorname{poly}(\mathbb Z,G_\bullet)$ satisfying
\[
 \widetilde g(Wn+b)=g(n)\quad(n\in\mathbb Z)
\]
which is totally $\delta^a$-equidistributed on $[WN+b]$.
\end{lemma}
\begin{proof}
Use the one-parameter coordinate characterization in
\cite[Lemma 6.7]{GTOrbits2012}. Each second-kind Mal'cev coordinate of
$g(n)$ is a polynomial, with the prescribed degree bound from the
filtration. Evaluate those same polynomials at real arguments and apply
the inverse coordinate map, obtaining $g_{\mathbb R}(t)$. Set
\[
 \widetilde g(t)=g_{\mathbb R}((t-b)/W).
\]
Affine composition preserves the coordinate degree bounds, so the same
characterization shows that this is a polynomial sequence for the
original filtration. It gives the stated identity on integer $n$.
Equivalently, one may use the Taylor product with real powers; those
powers exist in a connected simply connected nilpotent group.

Put $X=WN+b$ and $\epsilon=\delta^a$. Suppose the asserted total
$\epsilon$-equidistribution fails. Then on a progression
\[
 u+v[L]\subseteq[X],\qquad L\geq\epsilon X,
 \quad u,v\in\mathbb Z,\quad
 1\leq v\leq2/\epsilon,\quad |u|\leq2X,
\]
the sequence is not $\epsilon$-equidistributed. By the one-parameter
quantitative Leibman theorem \cite[Theorem 2.9]{GTOrbits2012}, there is
a nontrivial horizontal character $\eta$ satisfying
\[
 |\eta|\ll\epsilon^{-C},\qquad
 \|\eta\circ\widetilde g\circ(u+vt)\|_{C^\infty[L]}
       \ll\epsilon^{-C}.
\]
The rationality hypothesis holds since $Q\leq\epsilon^{-1}$.
Lemma~\ref{lem:ga-linear-coefficients} gives a positive integer
$k\ll\epsilon^{-C}$ such that
\[
 \|k\eta\circ\widetilde g\|_{C^\infty[X]}
      \ll\epsilon^{-C}.
\]
Forward composition with the fixed integer polynomial $Wn+b$ gives
\[
 \|k\eta\circ g\|_{C^\infty[N]}
      \ll_{W,d}\epsilon^{-C}.
\]
For this last bound expand $\binom{Wn+b}j$ in the binomial basis; its
matrix is integral, and $X$ is comparable to $WN$. The character
$k\eta$ remains nontrivial.

Let $C_{\mathrm{aff}}$ bound all exponents in these estimates, including
the character size. Fix explicitly
\[
 0<a<\min\{1,c_0/(2C_{\mathrm{aff}})\},\qquad
 0<\sigma<\min\{\sigma_0,1/(2a)\},
\]
where $c_0,\sigma_0$ are from
Lemma~\ref{lem:gn-horizontal-obstruction}. Decrease $\delta_0$ to
absorb fixed $W$-dependent constants. The resulting character bounds
contradict that lemma and total $\delta$-equidistribution of $g$.
The inequality $a\sigma<1/2$ ensures $\epsilon X\gg N^{1/2}$, so all
progression parametrizations and floor estimates used above hold for
large $N$. This proves the assertion on the entire interval $[WN+b]$,
including progressions outside the residue class $b$ modulo $W$.
\end{proof}

The Gaussian input in the following lemma is furnished by
Proposition~\ref{gi-minor-arc}; its rational counterpart is
Corollary~\ref{gi-ordinary-primes}.

\begin{lemma}[$W$-tricked minor-arc interface]\label{lem:ga-wtrick-ma}
For $E=\mathbb Q(i)$ and a fixed ray character $\chi$, put
\[
 T_\chi(t)=\sum_{\mathrm N\mathfrak a=t}
             \Lambda_E(\mathfrak a)\chi(\mathfrak a).
\]
Assume the mean-zero ideal estimate
\begin{equation}\label{eq:ga-ideal-input}
 \left|\sum_{\substack{\mathrm N\mathfrak a\leq X\\
                          \mathrm N\mathfrak a\in R}}
  \Lambda_E(\mathfrak a)\chi(\mathfrak a)
             F_0(h(\mathrm N\mathfrak a)\Gamma)\right|
 \ll Q_1X\log(2X)^7\eta^{c_E}\|F_0\|_{\mathrm{Lip}}
\end{equation}
for a positive $c_E$, whenever $h$ is totally $\eta$-equidistributed,
the basis is $Q_1$-rational, and $R$ is any interval progression of
modulus at most $Q_1$. Then, under the hypotheses of
Lemma~\ref{lem:ga-affine-interpolation}, for every Lipschitz $F_0$
with Haar mean zero and every interval progression
$P\subseteq[N]$ of modulus $q\leq Q$,
\begin{equation}\label{eq:ga-wtrick-ideal}
 \left|\sum_{n\in P}T_\chi(Wn+b)F_0(g(n)\Gamma)\right|
 \ll_{H_0,W,m,d}QN\log(2N)^7\delta^{ac_E}
                     \|F_0\|_{\mathrm{Lip}}.
\end{equation}
The same implication holds for $E=\mathbb Q$ and the corresponding
ordinary von Mangoldt estimate.
\end{lemma}
\begin{proof}
Take $X=WN+b$ and the sequence $\widetilde g$ constructed above. The
conditions $t=Wn+b$ and $n\in P$ specify an interval progression $R$
of modulus $Wq$. Apply \eqref{eq:ga-ideal-input} with
$Q_1=\max\{Q,Wq,2\}\ll_WQ$ and $\eta=\delta^a$. The ideal sum
then equals the left side of \eqref{eq:ga-wtrick-ideal} exactly, while
$X\ll_WN$ and $\log(2X)\ll_W\log(2N)$. No assertion about
fractional-argument distribution has been assumed: it was proved in
Lemma~\ref{lem:ga-affine-interpolation}.
\end{proof}

\begin{corollary}[Centered ray-prime weights]\label{cor:ga-centered-primes}
Suppose a conjugation-invariant fixed Gaussian ray condition has the
checked finite-character expansion
\[
 \Lambda'_C(t)=\frac12\sum_\chi c_\chi T_\chi(t)+e(t),
 \qquad \sum_{t\leq X}|e(t)|\ll_{H_0}\sqrt X\log(2X)^2.
\]
Assume both the Gaussian estimate \eqref{eq:ga-ideal-input} and its
ordinary rational counterpart, with positive exponents $c_E,c_{\mathbb Q}$.
Let $G/\Gamma$ and $g$ satisfy the hypotheses of
Lemma~\ref{lem:ga-affine-interpolation}, and let $F_0$ be Lipschitz with
Haar mean zero.
Then any fixed linear combination of the $W$-tricked ray-prime weight
and the $W$-tricked ordinary-prime weight obeys, for every interval
progression $P\subseteq[N]$ of modulus at most $Q$,
\[
 \left|\sum_{n\in P}w_N(n)F_0(g(n)\Gamma)\right|
 \ll QN\log(2N)^7\delta^{c_W}\|F_0\|_{\mathrm{Lip}},
 \qquad c_W=a\min\{c_E,c_{\mathbb Q}\}>0,
\]
for mean-zero $F_0$, in a regime
$\delta\geq N^{-\sigma_1}$ and $Q\leq\delta^{-a}$ with fixed
$\sigma_1>0$. Constants may depend on $H_0,W$ and the fixed linear
combination. In particular this applies to the proposed centering
\[
 w_N(n)=\frac{\varphi(W)}W\Lambda'_C(Wn+b)
             -\mu_C(b)\lambda_{b,W}(n),
\]
once the normalization of $\lambda_{b,W}$, the scalar $\mu_C(b)$ and
the displayed character expansion have been fixed and checked.
\end{corollary}
\begin{proof}
Apply Lemma~\ref{lem:ga-wtrick-ma} to the finitely many characters and
to the rational von Mangoldt function. The Gaussian factor $1/2$
accounts for the two conjugate prime ideals above an eligible split
rational prime. Replacing ordinary von Mangoldt by its prime-only
version also has a proper-prime-power error
$O_W(\sqrt N\log(2N)^2)$. Thus the stated estimate holds with the
additional error
$O_{H_0,W}(\sqrt N\log(2N)^2\|F_0\|_\infty)$.
Choose $\sigma_1\leq\sigma$ so small that $\sigma_1c_W<1/4$.
Then $N\delta^{c_W}\geq N^{3/4}$ throughout the claimed regime, and
the square-root errors are absorbed. All normalization factors are
fixed constants because $W$ is fixed. This proves the minor-arc
estimate but does not assert the centered progression estimate
\eqref{eq:gf-sw}, which must be proved separately from the arithmetic
main terms.
\end{proof}

\begin{remark}\label{rem:ga-choice-order}
Fix the arithmetic moduli first, then the positive exponents in this
subsection. They do not depend on the later factorization exponent
$B$. Proposition~\ref{prop:gf-factorization} chooses $B$ large enough
to satisfy the resulting upper complexity condition. Its output
$M$ is polylogarithmic in $N$, so its equidistribution parameter also
obeys $\delta\geq N^{-\sigma_1}$ for sufficiently large $N$.
The affine interpolation argument is motivated by the corresponding
step of Frei--Tinkova, but uses no theorem from that work as an input.
\end{remark}


% Section fragment. Bibliography: GTOrbits2012, GTMobius2012.
\subsection{Quantitative factorization for centered weights}
\label{sec:gf-factorization}

We isolate the analytic implication needed after the arithmetic
progression main terms have been established. Throughout,
$L=\log(2N)$, all nilpotent groups are connected and simply connected,
and all filtrations are by closed connected rational subgroups. Constants
may depend on fixed arithmetic moduli and on the fixed dimension and
degree bounds.

\begin{proposition}[Factorization implication]\label{prop:gf-factorization}
Fix dimension and degree bounds $m,d$. Let $w_N:[N]\to\mathbb C$ satisfy
the following three hypotheses.
\begin{enumerate}
\item For a fixed $K_w\geq0$, $|w_N(n)|\ll L^{K_w}$.
\item For every $A_0,C_0>0$, uniformly over all interval progressions
$P\subseteq[N]$ of common difference at most $L^{C_0}$,
\begin{equation}\label{eq:gf-sw}
 \left|\sum_{n\in P}w_N(n)\right|
       \ll_{A_0,C_0}N L^{-A_0}.
\end{equation}
No lower bound on $|P|$ is imposed here.
\item There are fixed $c>0$, $C_1,K_1\geq0$ such that if a nilmanifold
of dimension at most $m$ and filtration degree at most $d$ has a
$Q$-rational adapted Mal'cev basis, $h$ is a polynomial sequence totally
$\delta$-equidistributed on $[N]$, and $F_0$ is Lipschitz with Haar mean
zero, then
\begin{equation}\label{eq:gf-ma}
 \left|\sum_{n\in P}w_N(n)F_0(h(n)\Gamma)\right|
 \ll N\delta^c Q^{C_1}L^{K_1}\|F_0\|_{\mathrm{Lip}}
\end{equation}
for every interval progression $P\subseteq[N]$ of modulus at most $Q$
and length at least $N/Q$. It is sufficient for this hypothesis to hold
only when
\begin{equation}\label{eq:gf-ma-regime}
 \delta\geq N^{-\sigma},\qquad Q\leq\delta^{-\kappa}
\end{equation}
for fixed positive $\sigma,\kappa$. An estimate allowing arbitrary
progression lengths is stronger than required.
\end{enumerate}
Then, for every $A,C_{\mathrm{init}}>0$, uniformly over initial
nilmanifolds of dimension at most $m$ and degree at most $d$ with an
$L^{C_{\mathrm{init}}}$-rational adapted basis, all polynomial sequences
$g$, and all Lipschitz $F$,
\begin{equation}\label{eq:gf-conclusion}
 \left|\sum_{n\leq N}w_N(n)F(g(n)\Gamma)\right|
     \ll_{A,C_{\mathrm{init}}}N L^{-A}\|F\|_{\mathrm{Lip}}.
\end{equation}
\end{proposition}

\begin{proof}
Normalize $\|F\|_{\mathrm{Lip}}=1$ and put
\[
 D=A+K_w+3,\qquad M_0=L^{\max\{1,C_{\mathrm{init}}\}}.
\]
Leave a factorization exponent $B$ unspecified for the moment. By
\cite[Theorem 1.19]{GTOrbits2012}, there is an integer $M$ with
\[
 M_0\leq M\ll_{B,m,d}M_0^{E(B,m,d)}
\]
and a factorization $g=\varepsilon g'\gamma$ into polynomial sequences.
Here $\varepsilon$ is $(M,N)$-smooth, $\gamma(n)\Gamma$ has period
$q\leq M$, and $g'$ is totally $M^{-B}$-equidistributed in a rational
subgroup quotient. The new adapted basis consists of $M$-rational
combinations of the initial one.

For each residue $j$ modulo $q$, choose a reduced rational representative
$\gamma_j$ of $\gamma(j)\Gamma$. Split that residue class into
$\lceil ML^D\rceil$ consecutive blocks of nearly equal cardinality.
For sufficiently large $N$ after $B$ is fixed, the resulting progressions
$P_{j,k}$ satisfy
\begin{equation}\label{eq:gf-cells}
 \#\{P_{j,k}\}\ll M^2L^D,\qquad
 |P_{j,k}|\gg\frac{N}{M^2L^D},\qquad
 \operatorname{diam}(P_{j,k})\ll\frac{N}{ML^D}.
\end{equation}
Their common difference is $q$. Indeed each original residue class has
cardinality comparable to $N/q$, and rounding the balanced block sizes
changes these estimates only by an absolute factor, since all powers of
$M$ are polylogarithmic in $N$ once $B$ is fixed.

Choose $n_{j,k}\in P_{j,k}$ and put $a_{j,k}=\varepsilon(n_{j,k})$.
The right-invariant ambient metric and smoothness imply
\begin{align*}
 &|F(\varepsilon(n)g'(n)\gamma_j\Gamma)
       -F(a_{j,k}g'(n)\gamma_j\Gamma)|\\
 &\hspace{1cm}\leq d(\varepsilon(n),\varepsilon(n_{j,k}))
       \ll L^{-D}\qquad(n\in P_{j,k}).
\end{align*}
There is no factor depending on the coordinates of $g'(n)$ in this
estimate: right invariance removes it. The total weighted freezing
error is therefore $O(NL^{K_w-D})=O(NL^{-A-3})$.

We use the quantitative conjugation claim proved in
\cite[p.~547]{moebius}, whose inputs are the rational-subgroup,
rational-element and metric comparisons of
\cite[Appendix A]{GTOrbits2012}. There are fixed $a>0$, $C_2\geq1$,
depending only on $m,d$, such that, with
\[
 H_j=\gamma_j^{-1}G'\gamma_j,\quad
 \Lambda_j=H_j\cap\Gamma,\quad
 h_j(n)=\gamma_j^{-1}g'(n)\gamma_j,
\]
the quotient $H_j/\Lambda_j$ has an adapted basis of rationality at most
$M^{C_2}$, $h_j$ is totally $M^{-aB+C_2}$-equidistributed, and
\[
 F_{j,k}(x\Lambda_j)=F(a_{j,k}\gamma_jx\Gamma)
 \quad\hbox{satisfies}\quad
 \|F_{j,k}\|_{\mathrm{Lip}}\leq M^{C_2},\qquad
 \|F_{j,k}\|_\infty\leq1.
\]
We have enlarged $C_2$ to absorb fixed multiplicative constants. These
exponents do not depend on $B$. In applying this published claim, one
must account for the commensurable lattice resulting from rational
conjugation; conjugation is not asserted to be an isometry of the
original quotients. Its finite-index and metric losses are included in
$a,C_2$.

Write $F_{j,k}=m_{j,k}+F_{j,k}^0$ by subtracting its Haar mean on
$H_j/\Lambda_j$. Then $|m_{j,k}|\leq1$,
$\int F_{j,k}^0=0$, and
$\|F_{j,k}^0\|_{\mathrm{Lip}}\leq2M^{C_2}$. This subtraction is needed
even if the initial function $F$ happened to have Haar mean zero.
For a dimension-zero factor the centered function simply vanishes.

Put
\[
 U=C_2+D+5,\qquad Q_*=M^U,\qquad b_*=aB-C_2.
\]
Since $M\geq L$, $Q_*$ dominates the rationality, cell modulus and
$N/|P_{j,k}|$. Applying \eqref{eq:gf-ma} to all centered cell functions
and using \eqref{eq:gf-cells} gives a total of at most
\begin{equation}\label{eq:gf-total-ma}
 C N L^{K_1}M^{-b_*c+C_1U+C_2+2+D}.
\end{equation}
Every cell is counted here. In particular, the bound $N$ in the
minor-arc hypothesis has not been replaced by an unsupported cell
cardinality.

We now fix $B$ sufficiently large that
\begin{equation}\label{eq:gf-choice-B}
 b_*>1,\qquad
 b_*c\geq C_1U+C_2+2+D+A+K_1+4,
 \qquad b_*\kappa\geq U+1
\end{equation}
when the restricted regime \eqref{eq:gf-ma-regime} is used. The last
condition is omitted for an unrestricted minor-arc estimate. These
inequalities involve only fixed exponents and not $E(B,m,d)$.
They make \eqref{eq:gf-total-ma} $O(NL^{-A-4})$. They also ensure
$Q_*\leq(M^{-b_*})^{-\kappa}$. The other regime condition,
$M^{-b_*}\geq N^{-\sigma}$, holds for all sufficiently large $N$:
after fixing $B$ the upper bound on $M$ is a fixed power of $\log N$.

Only at this point use that upper bound, writing $M\leq L^{E_*}$ after
enlarging the fixed exponent $E_*$. There are at most
$O(L^{2E_*+D})$ cells of modulus at most $L^{E_*}$. Apply
\eqref{eq:gf-sw} with
\[
 C_0=E_*+1,\qquad A_0=A+2E_*+D+4.
\]
Summing the constant terms $m_{j,k}\sum_{n\in P_{j,k}}w_N(n)$ gives
$O(NL^{-A-4})$. Combined with the centered estimates and freezing error,
this proves \eqref{eq:gf-conclusion}.
\end{proof}

\begin{remark}[Order of parameter choices]\label{rem:gf-order}
The order in the proof is
\[
 (A,K_w,C_{\mathrm{init}},c,C_1,K_1,a,C_2,\kappa,\sigma)
 \longrightarrow (D,U)\longrightarrow B\longrightarrow E_*
 \longrightarrow(A_0,C_0)\longrightarrow N_0.
\]
Thus there is no requirement that $B$ dominate the factorization upper
exponent $E(B,m,d)$. The centered progression estimate allows every
logarithmic saving, so its exponent can be chosen after $E_*$ is known.
The minor-arc input must have a fixed power form such as
\eqref{eq:gf-ma}; unspecified polylogarithmic specializations alone
would not justify this argument.
\end{remark}


\begin{corollary}[Fixed-modulus centered nilsequence estimate]
\label{cor:ga-fixed-centered-nilsequence}
Fix the Gaussian ray field and a conjugation-invariant union $C$ of ray
classes as above. Let $W$ be a fixed multiple of the conductor of
$K^{\mathrm{ab}}/\Q$, and let $b$ be a unit residue modulo $W$.
For every $A,C_{\mathrm{init}}>0$ and fixed dimension and degree bounds,
uniformly over the nilmanifolds and polynomial sequences in
Proposition~\ref{prop:gf-factorization} with an
$(\log(2N))^{C_{\mathrm{init}}}$-rational adapted basis,
\[
 \left|\sum_{n\leq N}h_{C,b,W}(n)F(g(n)\Gamma)\right|
 \ll_{H,W,A,C_{\mathrm{init}},m,d}
 N(\log(2N))^{-A}\|F\|_{\mathrm{Lip}}.
\]
The constants are uniform in $b$ and in the polynomial orbit. No
uniformity as $W$ varies is asserted.
\end{corollary}
\begin{proof}
We check the three hypotheses of
Proposition~\ref{prop:gf-factorization}. The pointwise inequality
$|h_{C,b,W}|\leq\lambda_{b,W}$ gives its first hypothesis with $K_w=1$.
Corollary~\ref{cor:gaussian-centered-progressions} gives its second
hypothesis for every $A_0,C_0$, in fact for every interval progression.

For the third hypothesis, take $H_0=H$ in
Corollary~\ref{cor:ga-centered-primes}. Its character expansion is
\eqref{eq:gaussian-ray-fourier}, with the factor $1/2$ and the
prime-power error supplied by \eqref{eq:gaussian-ray-prime-factor}
and the bound following that equation. Its two assumed minor-arc
estimates are Proposition~\ref{gi-minor-arc} and
Corollary~\ref{gi-ordinary-primes}. Apply that corollary to the fixed
linear combination defining $h_{C,b,W}$ in
\eqref{eq:gaussian-centered-weight}. It gives \eqref{eq:gf-ma} with
$c=c_W$, $C_1=1$ and $K_1=7$ in the regime
$\delta\geq N^{-\sigma_1}$, $Q\leq\delta^{-a}$.
The affine-interpolation result initially restricts $\delta$ to a
fixed sufficiently small interval. For larger $\delta$, the bound
$|h_{C,b,W}|\ll_W\log(2N)$ gives \eqref{eq:gf-ma} directly after
enlarging its fixed constant. For a dimension-zero quotient the
mean-zero test function vanishes. For all other dimensions and degrees
within the fixed bounds, take the minimum of the finitely many positive
exponents $c_W,a,\sigma_1$ and the maximum of the fixed constants.
Thus the restricted regime \eqref{eq:gf-ma-regime} holds with fixed
positive $\kappa,\sigma$, uniformly throughout the required class.

Proposition~\ref{prop:gf-factorization} now proves the displayed
estimate. There are only finitely many unit residues modulo the fixed
$W$, so their constants can be bounded by a single constant. In
particular, division by $N$ and the relabeling $W=V$, $N=L$ give the
fixed-complexity estimate \eqref{eq:gaussian-centered-nil-input}.
\end{proof}


% The analytic bridge below uses the fixed-modulus centered nilsequence
% estimate just proved for the Gaussian ray fields under consideration.
\subsection{From Gaussian nilsequence estimates to the polynomial count}
\label{subsec:gaussian-gowers-bridge}

Fix the Gaussian ray field $K/\Q(i)$ and the conjugation-invariant ray
classes used in the preceding argument. Let $D$ be divisible by the
conductor of $K^{\rm ab}/\Q$. For each relevant union $C$ of classes,
let $\Lambda'_C$ be its rational prime weight and let $\mu_C(b)$ be the
conditional density in \eqref{eq:v5-mu}. In particular
$0\leq\mu_C(b)\leq1$. We use the following conclusion of the preceding
Gaussian norm, progression and factorization arguments, as assembled in
Corollary~\ref{cor:ga-fixed-centered-nilsequence}: for every fixed
multiple $V$ of $D$, every unit $b\pmod V$, every fixed bounded-complexity
nilsequence $F(g(n)\Gamma)$, and every $B>0$,
\begin{equation}\label{eq:gaussian-centered-nil-input}
\begin{aligned}
 \left|\mathbb E_{n\leq L}h_{C,V,b}(n)F(g(n)\Gamma)\right|
 &\ll_{V,B,\mathrm{complexity}}(\log L)^{-B},\\
 h_{C,V,b}(n)&=\frac{\varphi(V)}V
 \bigl(\Lambda'_C(Vn+b)-\mu_C(b)\Lambda'(Vn+b)\bigr).
\end{aligned}
\end{equation}
The bound is uniform in the polynomial orbit and functions with the stated
complexity bound. The field, its ray modulus and $V$ are fixed in this
statement. We do not assume any uniform rate as $V$ varies.

\paragraph{Published transference inputs.}
We use Green--Tao, \emph{Linear equations in primes}~\cite{greentao},
Proposition~6.4 and Appendix~D for the ordinary prime majorant,
Proposition~7.1 and Appendix~C for its generalized von Neumann bound,
and Lemma~4.4 with its subsequent normal-form reduction.
These are statements about the ordinary prime majorant and arbitrary
functions dominated by it; they do not assert an Artin-prime theorem.
For the inverse step we use the relative inverse theorem
\cite[Proposition~10.1]{greentao}, together with the local proof of the
bounded inverse theorem in Section~\ref{v15:gi-full}.
The latter supplies $\mathrm{GI}(k)$ for every fixed $k\geq3$, with the
normalized interval norm and a fixed finite family of same-step linear
nilsequences. Its supporting modules give the required constructions;
the retained published inputs are specified there. The relative theorem
applies to a real $f:[L]\to\mathbb R$, dominated by a nonnegative measure on
$\Z/Q\Z$, where $Q\in[C_*L,2C_*L]$ is prime and $C_*\geq20$.
For target norm $U^{k+1}[L]$, it requires
$(k+2)2^{k+1}$-pseudorandomness in the sense of
\cite[Definitions~6.1--6.3]{greentao}, including both the full
linear-forms condition and the correlation condition.
A fixed positive norm lower bound then gives a fixed positive
correlation with a $1$-bounded Lipschitz linear nilsequence from a
fixed finite family of $k$-step nilmanifolds. The family, metrics,
Lipschitz bound and correlation threshold depend only on the fixed
parameters, not on $L$ or the function. Complex test functions in the
bounded inverse theorem can be replaced by a real or imaginary part,
losing only a fixed factor in the correlation.
We use the normalized zero-extension interval norm, as in
\cite[Appendix~B, (B.13)--(B.15)]{greentao} and
Section~\ref{v15:gi-full}.

\begin{lemma}[Full ordinary-prime majorant with a fixed modulus factor]
\label{lem:ga-full-prime-majorant}
Fix positive integers $E,T$, an integer $S\geq2$, and $C_*\geq20$.
Suppose $z=z(L)\to\infty$, $z\leq\log\log\log L$, and
$V=E\prod_{p\leq z}p$. For each $L$ let $b_1,\ldots,b_T$ be any
unit residues in $\{0,\ldots,V-1\}$ and let
$Q\in[C_*L,2C_*L]$ be prime. There is a nonnegative
$S$-pseudorandom measure $\nu$ on $\Z/Q\Z$ and a fixed $\alpha>0$
such that, on $[L^{3/5},L]$,
\[
 \alpha\sum_{i=1}^T\frac{\varphi(V)}V\Lambda'(Vn+b_i)\leq\nu(n).
\]
All error terms and domination constants are uniform in these unit
residues. More explicitly, $\mathbb E\nu=1+o(1)$, and for
$d,m\leq S$ and every finite-complexity system of $m$ affine-linear
forms in $d$ variables with integer linear coefficients bounded by $S$,
\[
 \mathbb E_{x\in(\Z/Q\Z)^d}\prod_{j=1}^m\nu(\psi_j(x))=1+o(1),
\]
uniformly also in their affine constant terms. For each $2\leq m\leq S$
there is $\tau_m:\Z/Q\Z\to\mathbb R_{\geq0}$ such that, for all shifts,
including repetitions,
\[
 \mathbb E_x\prod_{j=1}^m\nu(x+h_j)
 \leq\sum_{i<j}\tau_m(h_i-h_j),\qquad
 \mathbb E\tau_m^q=O_{S,T,C_*,E,q}(1)\quad(1\leq q<\infty).
\]
\end{lemma}
\begin{proof}
We spell out the fixed-factor and normalization points in the
construction of \cite[Proposition~6.4 and Appendix~D]{greentao}.
Fix the smooth even function
\[
 \chi(t)=\begin{cases}\exp(1-(1-t^2)^{-1}),&|t|<1,\\0,&|t|\geq1,
 \end{cases}\qquad
 c_\chi=\int_0^1|\chi'(t)|^2\,dt>0.
\]
We retain $c_\chi$ rather than impose a second normalization on $\chi$.
Put $\gamma=1/(32S)$, $R=L^\gamma$, and
\[
 \Lambda_{\chi,R,2}(a)
 =\log R\left(\sum_{d\mid a}\mu(d)\chi(\log d/\log R)\right)^2,
 \qquad
 \widetilde\nu(n)=\frac1{Tc_\chi}\frac{\varphi(V)}V
       \sum_{i=1}^T\Lambda_{\chi,R,2}(Vn+b_i).
\]
Set $\nu(n)=(1+\widetilde\nu(n))/2$ on $[L]\subset\Z/Q\Z$ and
$\nu(n)=1$ outside it. These functions are nonnegative. For a prime
$Vn+b_i>R$, the corresponding divisor square is $\log R$.
Since $V=L^{o(1)}$, eventually $V\leq L$ and
$\log(Vn+b_i)\leq4\log L$ for $n\leq L$.
Thus $\alpha=\gamma/(8Tc_\chi)$ proves the displayed domination.

Here is why the divisor-sum estimates are uniform with this fixed
$\gamma$. In the proof of \cite[Theorem~D.3]{greentao}, a product of
$m\leq S$ squares introduces $2m$ divisors, each at most $R$.
The residue-class lattice count has error $O_d(qL^{d-1})$, where
$q\leq R^{2m}$ is their common modulus. Summing its absolute values
gives $O_{S,\chi}(R^{4m}L^{d-1}(\log R)^m)=o(L^d)$.
This count is independent of the coefficients and affine constants
of the forms. Consequently $\gamma$ need not shrink with $V$.
The remaining Euler-product argument is
\cite[Appendix~D, (D.4)--(D.23), especially Proposition~D.4]{greentao}.
For the present forms its constants are uniform: at $p\mid V$ every
nonempty divisor event is empty, by the unit-residue condition; at
$p\nmid V$, each nonzero linear form has one-equation density $1/p$.
For the linear-forms condition, choose $z$ larger than all nonzero
two-by-two linear-coefficient minors that can occur. Every such minor
has absolute value at most $2S^2$. Then every pair of distinct linear
parts has rank two at $p>z$, so every nonvertical local event has
density $O_S(p^{-2})$. No bound on affine constants is used here.
The finitely many smaller primes are either in $V$ or absorbed in
the fixed initial threshold. The Taylor and Euler-product bounds
therefore have constants depending only on $S,\chi$; their exceptional
prime factor is $\exp(O_S(\sum_{p\leq z}p^{-1/2}))$.
Here we include every $p\leq z$ in the exceptional set, also for a
one-form product; this covers its altered vertical local events.
This times $(\log R)^{-1/20}$ tends to zero under the stated bound on $z$.

For all sufficiently large $z$, the prime divisors of $E$ are at most
$z$. Thus $\varphi(V)/V=\prod_{p\leq z}(1-1/p)$, regardless of the
exponents in $E$. The local factors at these primes are exactly
$(1-1/p)^{-m}$ and cancel this normalizing factor to the $m$th power.
The factors at $p>z$ have product $1+o(1)$. The sieve factor for each
square is $c_\chi$ by \cite[Lemma~D.2]{greentao}, and is cancelled by
our factor $c_\chi^{-m}$. Expanding the average over the $T$ residues
therefore gives the required $1+o(1)$ for every subproduct of
$\widetilde\nu$ on a convex region.
To pass to the cyclic group, represent its coordinates in $[1,Q]$ and
partition into the finitely many convex regions recording each
wraparound and whether each form lies in $[L]$. The number of regions
is bounded in terms of $S,C_*$, uniformly in affine constants.
Expand $\nu=1+\tfrac12 1_{[L]}(\widetilde\nu-1)$ on these regions.
Every nonempty centered product has mean $o(1)$ by inclusion--exclusion.
This proves the full cyclic linear-forms condition and its mass case.

For clarity we also verify the correlation part of Appendix~D.
On each of the finitely many wraparound intervals, expand the factors
of $\nu$ and then their averages over the $T$ residues. Distinct
shifted forms $V(n+h_j)+b_{i_j}$ have exceptional primes outside $z$
only among the divisors of
$V(h_j-h_{j'})+b_{i_j}-b_{i_{j'}}$.
The local upper bound and the error bound in
\cite[Appendix~D, (D.25) and the following estimates]{greentao},
followed by the arithmetic--geometric mean inequality over pairs,
bound the result by a sum of fixed multiples of
\[
 H_A(u;a,b)=\exp\left(A\sum_{\substack{p>z\\p\mid Vu+a-b}}p^{-1/2}\right),
 \qquad Vu+a-b\ne0,
\]
where $A$ is fixed in terms of $S$. Terms with fewer than two active
factors are covered by adding a fixed positive constant.
All small-prime error factors are uniformly bounded by the preceding
choice of $z$. Sum these bounds over $a,b\in\{b_1,\ldots,b_T\}$
and over the possible integer lifts $u,u-Q,u+Q$ of a cyclic difference.
This defines $\tau_m(u)$ for nonzero cyclic differences, after
increasing a fixed multiplicative constant.

These weights have every fixed moment bounded. For fixed $q$, the
exponential with $Aq$ is bounded by a constant depending on $A,q$
times $\prod_{p>z,\,p\mid Vu+a-b}(1+p^{-1/4})$.
Expanding this product reduces its average over any of the above
integer intervals to
\[
 O_{A,q,C_*}\left(
  \sum_{d\geq1}d^{-5/4}
  +\frac1L\sum_{d\ll_{C_*}VL}d^{-1/4}\right)=O_{A,q,C_*}(1).
\]
Indeed all such divisors are coprime to $V$, so the Chinese remainder
count is at most $O_{C_*}(L/d+1)$, uniformly in $a-b$.
The second term is $O(V^{3/4}L^{-1/4})=o(1)$.
The finitely many summands and lifts preserve the fixed-moment bound.
An integer equality $Vu+a-b=0$, with $0\leq a,b<V$, requires $u=0$
and $a=b$. To include repeated cyclic shifts, define instead
$\tau_m(0)=(1+\|\nu\|_\infty)^S$.
The divisor bound $d(n)=n^{o(1)}$ gives $\|\nu\|_\infty=L^{o(1)}$,
uniformly here, so this changes every fixed moment by $L^{-1+o(1)}$.
For completeness that divisor bound follows by fixing $\epsilon>0$:
for all sufficiently large primes, $a+1\leq p^{\epsilon a}$;
the finitely many smaller primes contribute a fixed constant after
slightly enlarging $\epsilon$. Apply this to the prime factorization
of $n\ll VL$ and then let $\epsilon$ decrease.
This proves the correlation condition, including repetitions, and
completes the construction.
\end{proof}

\paragraph{An auxiliary cutoff, distinct from the polynomial cutoff.}
Put
\[
 V_z=D\prod_{\substack{p\leq z\\p\nmid D}}p.
\]
Here $z$ is an auxiliary cutoff used only to prove the next linear
replacement statement. It is not the cutoff $v$ defining $W,A$ in
\eqref{eq:v5-global-gowers}.
For every fixed integer $s\geq2$, one can choose $z=z(L)\to\infty$
sufficiently slowly that
\begin{equation}\label{eq:gaussian-refined-gowers}
 \max_{C,\ (b,V_z)=1}\|h_{C,V_z,b}\|_{U^s[L]}=o(1).
\end{equation}
The maximum is over the fixed finite list of ray-class unions and the
finitely many unit residues for the current $V_z$.

Here are the normalization and limit details for the inverse step.
For a required norm order $s\geq2$, put $r=\max(s,4)$ and $k=r-1$.
Choose a fixed integer $S\geq(r+1)2^r$ at least as large as the
requirements of the intended generalized von Neumann bound.
Choose the fixed ambient constant $C_*>2^r$, $C_*\geq20$, large
enough for that bound. The number $T$ of ordinary prime weights to be
dominated at any one application is fixed by the finite system.
All these choices precede a norm threshold $\delta$.

For sufficiently large $z$ our $V_z$ has the form
$E\prod_{p\leq z}p$ with $E$ fixed. Apply
Lemma~\ref{lem:ga-full-prime-majorant}. It gives one fixed $\alpha>0$
such that $f=\alpha 1_{[L^{3/5},L]}h_{C,V_z,b}$ is dominated by the
constructed $\nu$ on $[L]$, uniformly in the relevant unit residues.
This uses $|h_{C,V,b}|\leq(\varphi(V)/V)\Lambda'(Vn+b)$.
For the inverse step one may take $T=1$ and use a separate majorant
for each residue; its constants are uniform. The fixed larger $T$
is used when several weights enter the generalized von Neumann bound.

The removed initial interval has $U^s[L]$ norm $o(1)$. The weights are
$O(\log L)$ since $V_z=L^{o(1)}$, and cubes whose initial vertex lies
in that interval form an $O(L^{-2/5})$ fraction. Thus the $2^s$th
power of the removed norm is
$O((\log L)^{2^s}L^{-2/5})$. The same interval contributes $o(1)$ to
correlations with bounded tests. If the unscaled norm is at least
$\delta>0$, eventually $\|f\|_{U^s[L]}\geq\alpha\delta/2$.

Take $Q\in[C_*L,2C_*L]$ prime and write $I=[L]\subset\Z/Q\Z$.
The common ambient choice is valid for all interval norms up to order
$r$. The finite-group Gowers--Cauchy--Schwarz inequality gives
monotonicity of the group norms and
$\|1_I\|_{U^s(\Z/Q\Z)}\geq |I|/Q$, while
$\|1_I\|_{U^r(\Z/Q\Z)}\leq1$. Hence, for the zero extension
$\widetilde f$,
\[
 \|f\|_{U^r[L]}
 =\frac{\|\widetilde f\|_{U^r(\Z/Q\Z)}}
        {\|1_I\|_{U^r(\Z/Q\Z)}}
 \geq\|f\|_{U^s[L]}\frac LQ
 \geq\frac{\alpha\delta}{4C_*}.
\]
This does not require monotonicity with constant one for normalized
interval norms. Apply \cite[Proposition~10.1]{greentao} with
$k=r-1$ and threshold $\min(\alpha\delta/(4C_*),1/2)$;
Section~\ref{v15:gi-full} supplies $\mathrm{GI}(k)$.
The measure has the required $(k+2)2^{k+1}$-pseudorandomness because
it is $S$-pseudorandom. We obtain a fixed positive correlation with a
bounded linear nilsequence from a fixed finite family, independently
of $L,z,b$ and the ray-class union. This contradicts
\eqref{eq:gaussian-centered-nil-input} along a sufficiently slow cutoff.

Here is the order of that cutoff choice. The finite set of required
norm orders, the majorant complexities, $C_*,T,\chi,\alpha$ and $E$
are fixed first. For each order enumerate the tests
$\delta=1,1/2,1/3,\ldots$. Their inverse-theorem families and positive
thresholds depend only on these fixed parameters and on $\delta$.
At stage $j$, choose a cutoff $z_j$ larger than all fixed excluded
primes and large enough that the limiting linear-forms errors for
the finitely many bounded linear-coefficient types and subproducts
are below $1/j$, uniformly in affine constants. Then fix $V_{z_j}$
and choose $L_j$ so that all remaining
majorant errors are below $1/j$, the truncation bounds hold, and the
nilsequence bounds for the first $j$ tests and all unit residues
modulo $V_{z_j}$ are smaller than their fixed correlation thresholds.
Enlarge $L_j$ to ensure $z_j\leq\log\log\log L$ for $L\geq L_j$.
The uniform correlation-moment bounds of the lemma hold throughout;
they are not inferred from the linear-forms errors. On these stages
the full pseudorandom family tends to its required limit, so the
qualitative relative inverse theorem applies to any hypothetical
sequence violating \eqref{eq:gaussian-refined-gowers}.
Take the maximum of the finitely many nilmanifold thresholds at each
test. Their fixed smooth metrics are comparable to fixed Mal'cev
metrics, and a linear orbit with any base point is a polynomial orbit.
Thus the existing nilsequence estimate covers the supplied tests.

When this construction is used at the refined scale
$L\asymp N/V_{z_j}$ below, enlarge a physical threshold $N_j$ so that
every constant multiple of $N/V_{z_j}$ used by the normal-form
reduction is at least $L_j$ for $N\geq N_j$.
The cutoff is constant on each physical-scale stage. No equation
$z=z(N/V_z)$ is assumed. This proves
\eqref{eq:gaussian-refined-gowers} for the required orders.

\begin{lemma}[Fixed-system linear replacement]
\label{lem:gaussian-linear-replacement}
Let $\Psi=(\psi_1,\ldots,\psi_t):\Z^d\to\Z^t$ have fixed nonzero integer
linear parts of finite complexity. Allow its constant terms to vary
with $N$, subject to $\|\Psi\|_N\leq L_0$ for a fixed $L_0$.
Uniformly over convex regions $\mathcal K\subseteq[-N,N]^d$ with
$\Psi(\mathcal K)\subseteq[1,L_0N]^t$,
\begin{equation}\label{eq:gaussian-linear-replacement}
 \sum_{\mathbf n\in\mathcal K\cap\Z^d}
 \left(\prod_{i=1}^t\Lambda'_{C_i}(\psi_i(\mathbf n))
 -\prod_{i=1}^t\mu_{C_i}(\psi_i(\mathbf n)\bmod D)
                         \Lambda'(\psi_i(\mathbf n))\right)=o(N^d).
\end{equation}
The same conclusion holds for a mixed product in which any selected
factors have already been replaced by their ordinary-prime models.
\end{lemma}
\begin{proof}
We can discard points where some $\psi_i\leq N^{9/10}$.
For a nonconstant fixed linear form their number is
$O_{\Psi,L_0}(N^{d-1}(N^{9/10}+1))$; the product of weights is
$O((\log N)^t)$. This is $o(N^d)$, uniformly in the constant terms.
Apply the auxiliary $V_z$-trick to the variables:
$\mathbf n=V_z\mathbf m+\mathbf r$.
If one output residue is not a unit modulo $V_z$, its prime weight
vanishes on the remaining range, since the output is larger than
every prime dividing $V_z$. For an admissible residue vector write
\[
 \psi_i(V_z\mathbf m+\mathbf r)
       =V_z\widetilde\psi_{i,\mathbf r}(\mathbf m)+b_{i,\mathbf r},
       \qquad1\leq b_{i,\mathbf r}\leq V_z.
\]
The new linear coefficients are the original ones. Its normalized
constant terms remain bounded at scale $N/V_z$. Because $D\mid V_z$,
the factor $\mu_{C_i}$ is constant on this residue vector.

Telescope the difference of the two products. After normalization by
$(\varphi(V_z)/V_z)^t$, each term has one factor $h_{C_i,V_z,b_i}$;
all other factors are nonnegative and at most
$\lambda_{V_z,b_j}$. Introduce a common ordinary-prime majorant for
these at most $t$ weights and truncate as above. Multiply every
factor by a common fixed $\alpha_0>0$ so that each has absolute value
at most that majorant, and apply the normal-form extension of
Green--Tao Lemma~4.4 and Proposition~7.1. Here $\alpha_0$ depends
only on the fixed system, the chosen complexities and the fixed
conductor, not on $z$ or the residue vector. The small factor has
the required norm by \eqref{eq:gaussian-refined-gowers}. Dividing
the resulting estimate by $\alpha_0^t$ changes only a fixed constant.
The normal-form extension adds only a fixed number of variables with
fixed coefficients; its convex-region reduction is the one following
Lemma~4.4, and divides out their box volume. The error on each residue
vector is therefore $o((N/V_z)^d)$, uniformly in the vector and convex
region. The discarded initial segment stays outside the retained
output range since $V_z=N^{o(1)}$.

One can make this reduction explicit. For each pair $i\ne j$ choose
an integer vector $v_{ij}$ with $\dot\psi_j(v_{ij})=0$ and
$\dot\psi_i(v_{ij})\ne0$; fixed nonproportional linear parts ensure
such vectors exist with fixed size. Replace $\mathbf m$ by
$\mathbf x+\sum_{i\ne j}y_{ij}v_{ij}$ and allow each new $y_{ij}$
to range over an integer interval of length comparable to $N/V_z$.
The block $y_{ij}$ for fixed $i$ certifies normal form. The region
defined by the original convex constraint on this sum and the new box
is convex. The integer change $(\mathbf x,\mathbf y)\mapsto
(\mathbf x+\sum y_{ij}v_{ij},\mathbf y)$ is bijective, so its sum
is exactly the original sum times the number of points in the new box.
Division by that number gives the asserted error. If there is only one
form, choose an integer vector $v$ with $\dot\psi_1(v)\ne0$ and replace
$\mathbf m$ by $\mathbf x+yv$. The same box and bijection argument
gives the reduction in this case.

No growing factor is lost when summing residue vectors. If
$\mathcal R_z$ is the set of admissible vectors, then
\begin{equation}\label{eq:gaussian-admissibility-mass}
 \frac{|\mathcal R_z|}{V_z^d}
       \left(\frac{V_z}{\varphi(V_z)}\right)^t
   =\prod_{p\mid V_z}\beta_p(\Psi)=O_{\Psi,L_0}(1).
\end{equation}
Exclude the finitely many primes dividing a chosen nonzero
coefficient of each linear part, and a chosen nonzero two-by-two
minor for each pair of linear parts. At every remaining prime each
output hyperplane has probability $p^{-1}$, and any two meet with
probability $p^{-2}$. Hence $\beta_p=1+O_t(p^{-2})$, uniformly in
the constant terms. This also covers $t=1$, when there are no pair
minors and the good-prime factor is exactly one.
The finitely many remaining factors are bounded by $(p/(p-1))^t$.
The Chinese remainder theorem proves the equality, also for the fixed
prime powers dividing $D$. Multiplying the per-residue error by this
bounded mass gives $o(N^d)$. The proof works unchanged for the mixed
products. All auxiliary thresholds can be chosen uniformly over bounded
constant terms and convex regions, because the published majorant and
generalized von Neumann conclusions have this uniformity.
\end{proof}

\paragraph{The exact local cube needed in the appendix.}
Now fix the appendix cutoff $v$, and put $W,A,n,R$ as in its proof.
This is a new limiting argument; the auxiliary cutoff in the proof of
Lemma~\ref{lem:gaussian-linear-replacement} has already been eliminated
from that lemma's statement. Let
\[
 f_{C,b}(x)=\frac{\varphi(W)}W\Lambda'_C(Wx+b)1_{R\leq x\leq n},
 \qquad g_{C,b}=f_{C,b}-\mu_C(b).
\]
For fixed $a\geq2$ and $C_1<C_2<C_3$, uniformly over
$q\leq A^{C_1}$ and $nA^{-C_3}\leq J\leq nA^{-C_2}$,
\begin{equation}\label{eq:gaussian-global-gowers}
 \|g_{C,b}\|_{U^a_{q[J]}}^{2^a}
       \ll_{a,C_1,C_2,C_3,\rho,D}\frac1{\log v}+o_v(1).
\end{equation}
The norm uses the difference multiset $q(h-h')$, $h,h'\in[J]$, on
$\Z/n\Z$.

To prove this, remove the cutoff $x<R$ at cost
$O((\log n)^{2^a}R/n)=o_v(1)$. Expand the centered cube and partition
its domain according to every wrap-around integer $m_\omega$.
The relevant affine forms are precisely
\[
 W\left(x+q\sum_{r=1}^a\omega_r(h_r-h'_r)-m_\omega n\right)+b,
 \qquad\omega\in\{0,1\}^a.
\]
They have pairwise nonproportional linear parts. Each region is convex;
there are $O_a(1)$ regions because $qJ/n<1$. For fixed $v$, the number
of possible $q$ is finite and $J/n$ stays in a fixed positive compact
interval. Apply Lemma~\ref{lem:gaussian-linear-replacement} to every
subproduct and region. Its error is $o_v(n^{2a+1})$ uniformly in $J$,
and division by $nJ^{2a}$ gives $o_v(1)$. Since every output is $b$
modulo $D$, all its conditional densities are the same $\mu_C(b)$.
Thus the centered Artin cube equals $\mu_C(b)^{2^a}$ times the
centered ordinary-prime cube, plus $o_v(1)$.

The ordinary cube is evaluated exactly as in Tao--Ziegler~\cite{v5tz},
Proposition~6, published pages 39--40, using the ordinary-prime linear
correlation theorem, for example the published formulation in
Bienvenu~\cite{v6bienvenu}, Theorem~1.1, pages 79--80.
The factors at $p\mid W$ cancel the normalization.
For $p\nmid Wq$, pairwise independent gradients give local factor
$1+O_a(p^{-2})$, uniformly in the wrap constants. For $p\mid q$,
$p>v$, the factor is $1+O_a(p^{-1})$. Since
$\log q=O_{C_1,\rho,D}(v)$, there are $O(v/\log v)$ such prime divisors.
Their product is $1+O(1/\log v)$, and the other large-prime factors
give $1+O(1/v)$. Centering cancels the common main term. Prime powers,
if the ordinary theorem is stated for $\Lambda$, contribute $o_v(1)$
by their $O(\sqrt n\log^2 n)$ total weight and the fixed linear-fiber
bound. This proves \eqref{eq:gaussian-global-gowers}, and hence supplies
\eqref{eq:v5-global-gowers} without an Artin linear-correlations theorem.

\paragraph{Completion by the published polynomial transfer.}
The remainder of the proof of Proposition~\ref{prop:v5-artin-patterns},
starting with its arbitrary-function dense-model step, now applies
unchanged to these fixed Gaussian ray fields. Specifically use the
published Tao--Ziegler~\cite{v5tz} Theorem~11 and Corollary~3, then Theorem~9 and
Theorem~13. The model for an arbitrary $f\leq\lambda_{W,b}$ is used;
one does not assert that an Artin weight has model one. Equation
\eqref{eq:gaussian-global-gowers} shows that its bounded model, minus
$\mu_C(b)$, has small global norms. The model is bounded by a fixed
constant depending on the chosen sieve exponent. Divide its difference
from $\mu_C(b)$ by a fixed bound for that difference before applying
Theorem~9, whose hypothesis is a bound by one. After restoring this
fixed factor, concatenation and model approximation give the required
local norm bounds. None of these scaling constants depends on $v$.

Keep the original parameter order: local box order and coefficient
bound first, then $B_0>C+2$, then the global order $a$, and finally
small $\rho$. Only then let $v$ increase. The quantities
$M_0=\lfloor T/W\rfloor$, $M_1=\lfloor M_0/A\rfloor$ and
$J=\lfloor A^{-2B_0a}M_1^2\rfloor$ satisfy the size hypotheses recorded
in the appendix when $T/\sqrt N\geq4\sqrt W/A$.
For each fixed $v$ the estimates are uniform for $T$ in that allowed
range. The final assertion $T=o(\sqrt N)$ is imposed only along the
subsequent diagonal $v=v(N)\to\infty$, not while $v$ is frozen.

After $u=Wx+b$, $w=Wm+c$, the polynomials are
$c_jWm^2+2c_jcm+(b+c_jc^2-b_j)/W$, with distinct leading coefficients.
The bounded-model transfer kills every term containing a centered
factor. Output-boundary strips have width $O(T^2)=o(N)$.
For completeness, dominate them first by strips of width $\epsilon N$.
Here is the needed localization of the ordinary majorant estimate.
This argument also covers negative coefficients $c_j$. For the true
integer count use an enlarged period $n_* = \lceil B N/W\rceil$, where
$B>2+\max_j|c_j|$ is fixed. We may assume $T\leq\sqrt N$.
Nonpositive outputs have prime weight zero, and every positive output
is at most $(1+\max_j|c_j|)N$. Discard also tuples for which some
positive output is at most $W\lceil n_*^\rho\rceil+W$. For each $w$
there are $O_k(W n_*^\rho+W)$ such $u$, so their total weight is
$O_k((W n_*^\rho+W)T(\log N)^k)=o_v(NT)$, uniformly in the
allowed $T$ range. All remaining true outputs have their ordinary
prime majorants on $[n_*]$ without wrap-around. Use the same
localization calculation on the original period $n$ to bound the
artificial wrapped terms of the cyclic count. The parameters $n_*$
and $n$ differ by a fixed factor, so both applications satisfy the
same polynomial-forms hypotheses. In the following calculation $n$
denotes either chosen ambient period.
Write $\nu_0(x)$ for the majorant factor with constant shift if there
is one, and set $\nu_0=1$ otherwise. There is at most one such factor.
Put $F(x)=\mathbb E_m\prod_i\nu_i(x+P_i(m))$, omitting that factor.
Tao--Ziegler, Theorem~15 in Appendix~A (published page~57), applied to
$\nu_0F$, $\nu_0F^2$ and $\nu_0$, gives
\[
 \mathbb E_x\nu_0(x)(F(x)-1)^2=o_{v\to\infty}(1)+o_v(1).
\]
In the square use independent variables $m,m'$; the differences
between all its nonconstant shift polynomials are nonconstant,
including those from distinct copies. Weighted Cauchy--Schwarz now
shows, for every interval $I$ in the $x$ variable,
\[
 \mathbb E_x1_I(x)\nu_0(x)F(x)
 =\mathbb E_x1_I(x)\nu_0(x)+o_{v\to\infty}(1)+o_v(1).
\]
For fixed relative interval length the one-dimensional divisor-sum
mean estimate gives
$\mathbb E_x1_I\nu_0=|I|/n+o_{v\to\infty}(1)+o_v(1)$;
it follows either directly from the ordinary majorant construction
or from its linear-forms estimate on that fixed interval.
The normalized prime weights are bounded by fixed multiples of
their sieve majorants, with constants depending on the fixed sieve
exponent. Thus a strip of width $\epsilon N$ has total weight
$[O(\epsilon)+o_{v\to\infty}(1)+o_{v,\epsilon}(1)]NT$, with that
fixed domination factor included in the implied constants.
The errors here have the same iterated order as the rest of the
polynomial transfer. Then take $\epsilon\downarrow0$ as described
below. This avoids a spurious factor $(\log N)^k$ from a pointwise
bound. The same positive majorant treats interval endpoints.

Choose increasing cutoffs $v_j$ for which the global/local errors and
$4\sqrt{W(v_j)}/A(v_j)$ tend to zero, and then choose size thresholds
valid for all the finitely many residues, step values and transfer
parameters at that cutoff.
At stage $j$ impose the boundary estimates for every
$\epsilon\in\{1,1/2,\ldots,1/j\}$, preserving all earlier tests.
Let $v(N)=v_j$ between successive thresholds and take
$\omega(N)\geq4\sqrt{W(v(N))}/A(v(N))$, tending to zero sufficiently
slowly. This gives the uniform assertion for
$\omega(N)\sqrt N\leq T=o(\sqrt N)$.
Indeed, for an arbitrary such $T$, first fix $k$; eventually
$T^2/N<1/k$ and the estimate for the strip of width $N/k$ applies.
Take the upper limit as $N\to\infty$ and then $k\to\infty$.
There is no requirement that $T^2/N<1/j(N)$ along the chosen stages.
The main term is precisely \eqref{eq:v5-artin-series}; the admissible
residue mass is bounded, so summation does not amplify errors.
The $u\in[N/2,N]$ version follows by using that convex interval directly,
or by subtraction with thresholds enlarged for $N/2$.

This bridge uses the separately established Gaussian nilsequence input
and published ordinary-prime transference. It uses no statement or
estimate from Frei--Tinkov\'a as a theorem.


\subsection{Fixed Artin conditions in quadratic patterns}
\label{subsec:v5-artin-transfer}



Here the field is fixed throughout the limiting argument. Let $K$ be the
ray class field of $\Q(i)$ for $H\Z[i]$, where $16\mid H$. Put
$G=\operatorname{Gal}(K/\Q)$, and let $C_1,\ldots,C_k$ be
conjugation-invariant subsets of $\operatorname{Gal}(K/\Q(i))$,
regarded as conjugacy-invariant subsets of $G$.
Let $D$ be divisible by the conductor of $K^{\rm ab}$.
For a unit $b\pmod D$, define
\begin{equation}\label{eq:v5-mu}
 \mu_j(b)=\frac{\#\{g\in C_j:g|_{K^{\rm ab}}=\sigma_b|_{K^{\rm ab}}\}}
                  {[K:K^{\rm ab}]},
 \qquad \sigma_b(\zeta_D)=\zeta_D^b,
\end{equation}
and set $\mu_j(b)=0$ for nonunits. Define $\Lambda'_j(p)=\log p$ for
unramified primes with arithmetic Frobenius in $C_j$, and zero otherwise.

\begin{proposition}[Fixed-field polynomial count]\label{prop:v5-artin-patterns}
There is a function $\omega(N)\to0$, depending on the fixed data, such that
whenever $\omega(N)\sqrt N\leq T=o(\sqrt N)$,
\begin{equation}\label{eq:v5-artin-patterns}
 \sum_{u\leq N}\sum_{1\leq w\leq T}
       \prod_{j=1}^k\Lambda'_j(u+c_jw^2)
       =\mathfrak S\,NT+o(NT),
\end{equation}
where
\begin{equation}\label{eq:v5-artin-series}
 \mathfrak S=
 \left(\frac D{\varphi(D)}\right)^k
 \mathbb E_{u,w\bmod D}\prod_j\mu_j(u+c_jw^2)
 \prod_{p\nmid D}\beta_p,
 \quad
 \beta_p=\mathbb E_{u,w\bmod p}
       \prod_j\frac p{p-1}1_{p\nmid u+c_jw^2}.
\end{equation}
The same statement on $N/2\leq u\leq N$ has main term
$\mathfrak SNT/2$.
\end{proposition}
\begin{proof}
The Gaussian linear replacement of
Lemma~\ref{lem:gaussian-linear-replacement} and its local-cube consequence
\eqref{eq:gaussian-global-gowers} supply the analytic input. We record
the normalization and polynomial-transfer parameters here.

Choose a small positive $\rho$ and first keep $v$ fixed. Put
\[
 W=D\prod_{\substack{p\leq v\\p\nmid D}}p,\quad
 A=W^{1/\rho},\quad n=\lfloor N/W\rfloor,\quad R=n^\rho,
\]
and, for unit $b\pmod W$, put
\[
 f_{j,b}(x)=\frac{\varphi(W)}W\Lambda'_j(Wx+b)1_{R\leq x\leq n},
 \qquad g_{j,b}=f_{j,b}-\mu_j(b).
\]
View these functions on $\mathbb Z/n\mathbb Z$, with representatives
$1,\ldots,n$. If $a\geq2$ and $C_1<C_2<C_3$ are fixed,
\eqref{eq:gaussian-global-gowers} gives
\begin{equation}\label{eq:v5-global-gowers}
 \|g_{j,b}\|_{U^a_{q[J]}}^{2^a}
       \ll_a(\log v)^{-1}+o_v(1),
 \quad q\leq A^{C_1},\quad
 nA^{-C_3}\leq J\leq nA^{-C_2}.
\end{equation}
The implicit constant may also depend on the other fixed complexities
and the fixed ray field. The norm, wrap-around regions, conditional
densities, and order of the auxiliary cutoff are all treated explicitly
in Section~\ref{subsec:gaussian-gowers-bridge}.

\phantomsection\label{q214foundation:tz-adaptation}
We now adapt the arbitrary-function dense-model argument of
Tao--Ziegler~\cite{v5tz}. For their squarefree primorial modulus, the published
Theorem~11 and Corollary~3 give a single bounded model accurate in both
global Gowers norms and the local polynomial box norms. Their Theorem~9 is
the bounded-function concatenation statement, and their Theorem~13 bounds
the polynomial average by one of those local norms. All majorants here are
the ordinary prime majorants, since $0\leq\Lambda'_j\leq\Lambda'$.
Our modulus can contain a fixed extra factor. Thus the uses of their
modulus-dependent results here are adaptations of their proofs, including
the polynomial-forms estimate in Theorem~15, rather than literal
applications of the published statements. The bounded-function
Theorem~9 does not involve this modulus. The changes in the divisor-sum
construction are as follows:
for sufficiently large $v$, $W=E\prod_{p\leq v}p$ with a fixed $E$ whose
prime factors are at most $v$. A unit residue makes every small-prime
divisibility condition identically false, including at the extra prime
powers. At the remaining primes the divisor-sum local factors satisfy the same
uniform estimates. Polynomial coefficient bounds absorb the fixed factor $E$.
Thus the polynomial-forms estimates for the majorant have the same
$1+o(1)$ conclusion.

The following choice records the size hypotheses of these results.
First choose their local box order $d'$ and a coefficient exponent $C$;
choose an integer $B_0>C+2$, then a sufficiently large global order $a$,
and finally a sufficiently small $\rho$. Set $\rho_* =\rho/B_0$ and
$\mathcal A=A^{B_0}$. Put
$M_0=\lfloor T/W\rfloor$, $M_1=\lfloor M_0/A\rfloor$.
For $T/\sqrt N\geq4\sqrt W/A$ one has
\[
 A^{-1}\sqrt n\leq M_0\leq\sqrt n,
 \qquad A^{-3}\sqrt n\leq M_1\leq\sqrt n.
\]
The model at cutoff $n^{\rho_*}$ dominates the function already truncated
at $n^\rho$. Corollary~3 supplies approximation for the degree-one local
step polynomials, whose coefficients are $O(A^C)$, and simultaneously for
$q\leq A^{B_0a}$ with
$J=\lfloor A^{-2B_0a}M_1^2\rfloor$.
In particular
\[
 nA^{-2B_0a-7}\leq J\leq nA^{-2B_0a}.
\]
This is within \eqref{eq:v5-global-gowers}. Let $f'_{j,b}$ be the
model supplied by Corollary~3 at cutoff $n^{\rho_*}$. Its bound is
$0\leq f'_{j,b}\leq B_*-1$, where $B_*\geq2$ depends on the fixed
ray field, complexities and on $\rho_*$, but not on $v,n,b$ or the polynomials
within the stated coefficient bounds. This is the bounded model of
Tao--Ziegler's Proposition~4: its construction first rescales the
prime weight by a fixed positive constant and then divides the model
by that constant. In particular, it need not be bounded by one.
Since $0\leq\mu_j(b)\leq1$, the function
\[
 F_{j,b}=\frac{f'_{j,b}-\mu_j(b)}{B_*}
\]
is bounded in magnitude by one. The triangle inequality, the global
part of Corollary~3, and \eqref{eq:v5-global-gowers} show that its
required global norms tend to zero under the diagonal limit below.
Apply Theorem~9 to $F_{j,b}$
with parameters $d=d_0=2$, $M=M_1$, $A=\mathcal A$, and global
order $a$. Its smallness threshold is independent of $\mathcal A$;
$B_*$ has already been fixed before $v$ is chosen. Multiplication
by $B_*$ and the local part of Corollary~3 therefore show that the
local norm of $g_{j,b}$ tends to zero. Theorem~13, with the original
parameter $A$ and coarse scale $M_0$, makes every polynomial average
containing a centered factor $o(1)$. These are the reasons for choosing
the complexities before $v$.

For the final $W$-trick write $u=Wx+b$ and $w=Wm+c$.
If $b_j$ is the residue of $b+c_jc^2$ modulo $W$, the resulting integer
polynomial is
\[
 c_jWm^2+2c_jcm+(b+c_jc^2-b_j)/W.
\]
Its leading coefficient differences are nonzero, and its coefficients
are $O(W)$. Expand the product into its means and centered factors and
apply the preceding conclusion. End intervals of length $O(T^2)=o(N)$
are controlled by the weighted second-moment localization and the
fixed-strip argument in Section~\ref{subsec:gaussian-gowers-bridge}. Summing residue classes does not
amplify the error: their normalized admissible mass is
$\prod_{p\mid W}\beta_p=O(1)$.
Indeed, outside the fixed coefficient-difference primes,
\[
 \beta_p=\left(\frac p{p-1}\right)^k
                \frac{(p-1)(p-k+1)}{p^2}=1+O_k(p^{-2}).
\]
This also proves convergence of the product in
\eqref{eq:v5-artin-series}.
Let $N$ tend to infinity for each fixed $v$ and then let $v$ tend to
infinity sufficiently slowly. A diagonal choice supplies $\omega(N)$
and yields \eqref{eq:v5-artin-patterns}. Subtract the corresponding
formula with upper limit $N/2$, enlarging $\omega$ if necessary, to
obtain the last assertion.
\end{proof}

\begin{remark}[Published numbering]\label{rem:v5-tz-numbering}
The labels in the published 2018 paper differ from those in arXiv
1603.07817v1. The relevant dictionary is
\[
\begin{array}{c|ccccc}
 \text{arXiv} &\text{Theorem 1.4}&\text{Theorem 2.5}&\text{Theorem 3.1}
               &\text{Corollary 3.9}&\text{Theorem 4.1}\\
 \text{published}&\text{Theorem 4}&\text{Theorem 9}&\text{Theorem 11}
               &\text{Corollary 3}&\text{Theorem 13}
\end{array}
\]
The published locations are respectively pages $4,21,28,38,41$.
The preprint's equations (3.5) and (4.1) are published equations (36)
and (57). The dense model used above is the theorem's assertion for an
arbitrary function bounded by the prime weight. The special model $1$
for the ordinary prime weight is not an assertion about Artin weights.
\end{remark}

\subsection{The ray-class lower bound}\label{subsec:v5-ray-density}

For ray class fields we use the published existence and reciprocity
statements in Kopp--Lagarias~\cite{v5kl}, Theorems~1.1, 1.3 and 7.2,
and the modulus exact sequence of Theorem~6.5, specialized to the maximal
order $\mathbb Z[i]$. There are no real places. This replaces an appeal
to unpublished course notes. In this principal ideal domain the ray group
for the finite modulus $H\mathbb Z[i]$ is
\begin{equation}\label{eq:v5-ray-group}
 (\mathbb Z[i]/H\mathbb Z[i])^\times/\{1,-1,i,-i\}.
\end{equation}
It also follows directly by sending a residue $r$ to the ideal class of
$(r)$: surjectivity is principality, and the kernel consists exactly of
unit residues. The Artin map sends this ideal class to arithmetic
Frobenius. Conjugation acts by $r\mapsto\bar r$.

Choose a fixed $L_1$ larger than all coefficient-difference primes and
larger than $11$. Let $H_0$ be divisible by $16\cdot7M$ and by all odd
primes at most $L_1$. Prescribe a unit residue $r\pmod{H_0}$ which is
$9$ modulo $16$, $1+3i$ modulo $7$, and $i$ modulo every other odd
prime dividing $H_0$. Compatible higher powers can be chosen arbitrarily.
The associated rational-prime Frobenius condition is the two-element
conjugacy class $C_0=\{h_r,h_{\bar r}\}$ in the Galois group $G_0$ of
the ray field $K_0$ over $\mathbb Q$.
The two elements are distinct: an equality of their ray classes would
require $r=\epsilon\bar r\pmod{H_0}$ for a Gaussian unit $\epsilon$;
modulo $16$ this forces $\epsilon=1$, contradicted modulo $7$.
The field is Galois over $\mathbb Q$, because conjugation preserves the
modulus and the defining ray subgroup.

Let $D_0$ be the conductor of $K_0^{\rm ab}$ and choose a unit
$b_0\pmod{D_0}$ inducing the restriction of $h_r$.
The value in \eqref{eq:v5-mu} is then
$\mu_0=|C_0|/|G_0'|>0$.
For $L\geq L_1$ let
$T_L=\prod_{L_1<\ell\leq L,\,\ell\nmid H_0}\ell$ and let $K_L$ be the
ray field for $H_0T_L\mathbb Z[i]$. For each new prime write
\[
 B_\ell=(\mathbb Z[i]/\ell\mathbb Z[i])^\times,
 \qquad N_\ell=\ker(B_\ell\xrightarrow{\mathrm N}\mathbb F_\ell^\times).
\]
The norm map is onto and $|N_\ell|=\ell-\chi(\ell)$.

\begin{lemma}[Abelian fibres and their surviving proportion]
\label{lem:v5-ray-fibres}
Put $N_L=\prod_{\ell\mid T_L}N_\ell$ and
$G_L=\operatorname{Gal}(K_L/\mathbb Q)$.
Then
\[
 N_L\subseteq G_L',\qquad
 G_L'/N_L\simeq G_0',\qquad
 K_L^{\rm ab}=K_0^{\rm ab}\mathbb Q(\zeta_{T_L}).
\]
Let $C_L$ consist of those elements above $C_0$ whose normalized Gaussian
residue $z$ satisfies $\mathrm N(z-1)\mathrm N(z+1)\not\equiv0\pmod\ell$
at every $\ell\mid T_L$. For the density $\mu_L$ of this set,
\begin{equation}\label{eq:v5-ray-density}
 \mu_L(b)\geq\mu_0
       \prod_{\ell\mid T_L}\left(1-\frac4{\ell-1}\right)
 \qquad (b\equiv b_0\pmod{D_0},\ (b,D_0T_L)=1).
\end{equation}
\end{lemma}
\begin{proof}
In the CRT presentation of \eqref{eq:v5-ray-group}, the unit group
embeds in the $H_0$ factor, because it already embeds modulo $16$.
Once the representative above $r$ or $\bar r$ is normalized, the kernel
of reduction to the base ray group is exactly $\prod B_\ell$.
There is no further quotient by the four units.

The map $z\mapsto z/\bar z$ from $B_\ell$ onto $N_\ell$ is surjective.
For split primes this is $(x,y)\mapsto(x/y,y/x)$.
For inert primes it is $z\mapsto z^{1-\ell}$ in the cyclic group
$\mathbb F_{\ell^2}^\times$.
These maps are commutators with complex conjugation, so $N_L\subseteq G_L'$.
The norm on each new residue factor gives the cyclotomic action
$\zeta_\ell\mapsto\zeta_\ell^{\mathrm N(z)}$.
This is also seen by applying arithmetic Frobenius to $\zeta_\ell$;
it is compatible with the ray map by reciprocity.
The reduction map gives $K_0\subseteq K_L^{N_L}$. The cyclotomic
action just computed is trivial on $N_L$, so also
$\mathbb Q(\zeta_{T_L})\subseteq K_L^{N_L}$.
The fields $K_0$ and $\mathbb Q(\zeta_{T_L})$ are disjoint over
$\mathbb Q$. Indeed, $K_0$ is unramified at every prime dividing
$T_L$, whereas the inertia groups at these primes generate the
Galois group of $\mathbb Q(\zeta_{T_L})/\mathbb Q$. Thus their
intersection is trivial. Moreover, the kernel computation above gives
\[
 [K_L^{N_L}:\mathbb Q]
 =[K_0:\mathbb Q]\prod_{\ell\mid T_L}\frac{|B_\ell|}{|N_\ell|}
 =[K_0:\mathbb Q]\varphi(T_L).
\]
This is the degree of $K_0\mathbb Q(\zeta_{T_L})$, proving
\[
 K_L^{N_L}=K_0\mathbb Q(\zeta_{T_L}),\qquad
 G_L/N_L\simeq G_0\times\prod_{\ell\mid T_L}\mathbb F_\ell^\times.
\]
Since $N_L\subseteq G_L'$, taking commutators in this quotient gives
$G_L'/N_L\simeq G_0'$ and
$K_L^{\rm ab}=K_0^{\rm ab}\mathbb Q(\zeta_{T_L})$.

Fix a nonzero norm $s\in\mathbb F_\ell^\times$.
For $z=a+bi$ in that norm fibre,
$\mathrm N(z\pm1)=s+1\pm2a$.
For each sign, vanishing fixes $a$ and leaves at most two values of $b$.
Thus at most four of the $\ell-\chi(\ell)$ elements are excluded.
The conditions are invariant under conjugation, so $C_L$ is a union of
rational conjugacy classes.
In every compatible abelian fibre its numerator in \eqref{eq:v5-mu} is
at least $|C_0|\prod(\ell-\chi(\ell)-4)$, and its denominator is
$|G_0'|\prod(\ell-\chi(\ell))$.
Their quotient proves \eqref{eq:v5-ray-density}.
\end{proof}

\begin{proposition}[Finite-cutoff lower bound]\label{prop:v5-cutoff-lower}
There is a constant $c>0$ independent of $L$ such that the singular series
in Proposition~\ref{prop:v5-artin-patterns} for the fixed field $K_L$ and
the set $C_L$ at every position satisfies
\[
 \mathfrak S_L\geq c(\log L)^{-4k}.
\]
\end{proposition}
\begin{proof}
At the base modulus retain just $u=b_0$, $w=0\pmod{D_0}$.
At each new prime retain all pairs on which the $k$ norm values are units.
Apply \eqref{eq:v5-ray-density} separately to each position.
The norm-fibre denominators cancel as shown in the previous lemma; a
factor inverse to the degree of $K_L$ is not left in the bound.
The normalized proportion of the retained pairs at each new prime is
$\beta_\ell=1+O_k(\ell^{-2})$. Each is positive, since $w=0$, $u\ne0$
is allowed. Their product, together with the tail Euler product, is
bounded below by a fixed positive number. Williams' theorem with modulus
$1$, or its product for the two reduced classes modulo $4$, gives
\[
 \prod_{\ell\mid T_L}(1-4/(\ell-1))^k
       \gg(\log L)^{-4k}.
\]
The base modulus contributes another fixed positive factor.
\end{proof}

\subsection{The final order of limits}\label{subsec:v5-diagonal}

\begin{proof}[Proof of Theorem~\ref{thm:v5-simultaneous-primes}]
Fix the coefficients, $M$, and the base modulus first. Choose one fixed
$L\geq\max(L_1,L_*)$ so large that
\[
 \frac C{L\log L}<\frac c{32(\log L)^{4k}},
\]
where $C,c$ are from the obstruction and lower-bound propositions.
Now fix $K_L,C_L$. Only after this choice do we use the asymptotic
polynomial theorem or let its small-prime cutoff tend to infinity.

Choose numbers $\tau_j\downarrow0$, all at most $\tau_0$.
Choose successively $N_j$ so large that for $N\geq N_j$
$\omega(N)\leq\tau_j/2$,
that $\tau_j\sqrt N\geq2$, and that the error term in
\eqref{eq:v5-obstruction-tail} is less than
$c\tau_jN^{3/2}/(32(\log L)^{4k})$.
This is possible for each fixed $\tau_j$.
For $N_j\leq N<N_{j+1}$ use $\tau(N)=\tau_j$ and define
$T=\lfloor\tau(N)\sqrt N\rfloor$.
Then $T=o(\sqrt N)$ and $T\geq\tau\sqrt N/2$.
Apply Proposition~\ref{prop:v5-artin-patterns} to this piecewise choice
of $T$. For all sufficiently large $N$ its lower count is at least
$cNT/(4(\log L)^{4k})$ and exceeds the full obstruction tail.
Hence a pattern remains with no forbidden odd
common divisor greater than $L$.
The ray conditions exclude every odd divisor at most $L$ from each
individual adjacent norm. At $7$ the adjacent norms of $1+3i$ are $2$
and $6$; at every other odd base prime the adjacent norms of $i$ are $2$.
Thus no forbidden odd common divisor remains at all.

The residue modulo $16$ gives, on writing $a=9+16A$ and $b=16B$,
\[
 (a+1)^2+b^2\equiv4\pmod8,\qquad
 (a-1)^2+b^2\equiv64\pmod{128}.
\]
These yield the exact two-adic valuations and hence the cross-sign gcd
$4$. The residue modulo $M$ gives coprimality with $M$.
Finally $u\geq N/2$ and $w\leq\tau\sqrt N$ imply
$u/w^2\geq1/(2\tau^2)\to\infty$.
For sufficiently small $\tau_0$, all $u+c_jw^2$ are between fixed positive
multiples of $N$, so they exceed $Y$ when $N$ is large.
The dilation at the start imposes $q\mid w$.
\end{proof}

% BEGIN VERSION 9 ADDITION
\section{The odd-product quartics in 523 and 540 variables}\label{sec:v9}

The following construction incorporates the November 7 contribution of
Gabriel Conant, developed with ChatGPT. We give its arithmetic argument,
including the connection to the results proved above. Variable counts
include the input. The older quartics remain as source and proof baselines.

\begin{theorem}[The Version 9 quartics]\label{thm:v9-quartics}
There are specified integer polynomials $P_{523}$ and $P_{540}$, of total
degree four in respectively 523 and 540 variables, such that
\[
 q\in\Z\ \Longleftrightarrow\ \exists v\in\Q^{522}:P_{523}(q,v)=0
 \ \Longleftrightarrow\ \exists w\in\Q^{539}:P_{540}(q,w)=0.
\]
The first has 35,163 nonzero monomials and largest coefficient magnitude
$64\cdot3^{2^{27}}$. Its coefficient table uses the fixed integer
$S_0=3^{2^{26}}$. The second has 35,213 nonzero monomials and largest
coefficient magnitude $3^{2048}+1005072314631333<2^{3247}$.
Its complete decimal table occupies 708,848 bytes. Both have content one.
\end{theorem}

The proof is the composition of the source equivalences, local and global
arguments, and exact polynomial certificates below. The intermediate
554-variable construction is included to separate the odd-product
argument from the additional formal-depth refinement. It is not a further
claimed improvement over the two polynomials in the theorem.

\subsection{Inherited inputs and the exact changes of convention}\label{subsec:v9-inputs}

Use the normalized curve and generator of
Section~\ref{subsec:v5-point-source}, with $P=P_*$, not a new generator.
The coefficient basis is $(1,\sqrt2)$ and conjugation sends
$a+b\sqrt2$ to $a-b\sqrt2$. At unramified finite places $v(p)=1$;
at ramified places the integral valuation has value group $\Z$.
Heights are absolute logarithmic Weil heights, with the local degree
weights divided by $[K:\Q]=2$. Throughout this section set
$\mathcal H=H_{19}$ of Theorem~\ref{thm:quartic-height19}.

The source is the system in Section~\ref{subsec:v5-complete-source},
with the following explicit changes. Replace the comparison targets
$[4]P_*,[4]Q_i$ by $P_*,Q_i$; replace the orders 16 and 52 by 12 and 40
before the further refinement below. Retain the sixteen-test product,
the same five contacts, the scalar norm tree with orders $(1,1,1,1,32,32)$,
the Boolean correction and the preparation-zero pole equations.
Change $\beta$ to $2X^2+3$; use the point $S-P$ only in the scalar and
physical-anchor blocks; and change the fixed coefficient $\mathsf S$
to $S_0=3^{2^{26}}$. The point-chart replacement, the three new scales,
and the final order-one/order-forty test list are given below. All
unchanged quotient guards and permitted zero values of pole witnesses
are retained. This description, the displayed equations below, and the
pinned rational straight-line program specify the source without a
choice of arithmetic witnesses.

\begin{center}\small
\begin{tabular}{p{.29\textwidth}p{.63\textwidth}}
\toprule
Input & Exact result used and translation\\\midrule
Elliptic labels and domains & Section~\ref{subsec:v5-point-source},
Lemma~\ref{lem:v5-point-domains}, and \eqref{eq:o8-indexheight}:
$2E(K)=\Z P_*$, nonzero coordinate domains, and the same lower bound
$h(x([n]P_*))\ge n^2-6$.\\
Height of rational contacts & Theorem~\ref{thm:quartic-height19},
with $\mathcal H<2^{2093966}$, applied through the fractional linear
map and contact calculation of Lemma~\ref{lem:v5-height-transfer}.
Only that calculation is reused; its exponent 24 is replaced by 19.\\
Pole obstruction and realization & Lemmas~\ref{hyp:o8-pole},
\ref{lem:original-seven-simultaneous-realization}, and
\ref{lem:v5-one-place}. These include the restricted-converse dependency
proved in the earlier arithmetic sections; it has not been discarded.\\
Scalar and anchor & Lemma~\ref{v5:scalar32-bridge}, with its proof
applied to $2X^2+3$ and $S-P$, and
Lemma~\ref{lem:v5-three-witness-anchor} with the smaller exponent
calculated just below.\\
Local comparisons & Lemma~\ref{lem:v5-unramified-log} and the
recurrence of Lemma~\ref{lem:quartic-reciprocal-log}; orders and
target changes are justified below, not inferred by renaming indices.\\
Global choices & Theorem~\ref{thm:v5-simultaneous-primes},
Lemma~\ref{lem:v5-two-reductions}, Proposition~\ref{prop:v5-artin-patterns},
Section~\ref{subsec:v5-diagonal}, and
Lemma~\ref{v5:real-progression-bridge}. All fixed data precede the
asymptotic prime choice.\\
Final contradictions & Lemma~\ref{lem:o8-lucas} and the explicit
new height inequalities below.\\\bottomrule
\end{tabular}
\end{center}

The published results behind these inputs retain the versions and
convention qualifications in the theorem-application inventory. The new
contribution does not supply an assumed replacement for the restricted
converse, the height theorem, or the Gaussian prime argument.

\begin{lemma}[Direct targets and smaller anchor]\label{lem:v9-bridge}
For the changed source above, the local comparison and completeness
arguments remain valid with direct targets. The scalar anchor on
$R=[a-1]P_*$ implies
\[
 0\ne a-1\in2\,3^{2^{24}-1}\Z,
 \qquad \log_2|a-1|>25165823.5,
 \qquad h(\beta)\le\mathcal H|a-1|^{19/32}.
\]
It permits arbitrarily divisible positive scalar indices $a-1$.
\end{lemma}
\begin{proof}
For direct targets set $V_0=H-P_*$ and $V_i=U_i-Q_i$, with
$H=[h]P_*$, $U_i=[u_i]P_*$ and $Q_i=[n_i]P_*$. The quotient guards
give finite nonidentity points. At a consumed prime the full logarithm
ratios are $u_i/h$ and $(u_i-n_i)/(h-1)$. One of $h,h-1$ is a unit.
As in Lemma~\ref{lem:v5-two-anchors}, the cross congruence first makes
both approximate ratios integral, and then makes their relative
errors absolute. Subtraction of the two integral congruences gives
$L_i\equiv n_i$, without division by four. The logarithm recurrence
and denominator bound in Lemma~\ref{lem:quartic-reciprocal-log} give
the same estimates at orders 12 and 40. The test list below strengthens
the input order to 40 in the final refinement.

For completeness, the reduced elliptic groups have all four rational
two-torsion points, since $p>76176$ and the three roots $0,\pm d$
are distinct. Thus their doubling images have order $d_\pm(p)/4$.
The point $P_*$ is itself a double, so these integers annihilate its
reductions. The prime theorem and Lemma~\ref{lem:v5-two-reductions}
therefore give, on setting $M_\pm=\operatorname{lcm}_p(d_\pm(p)/4)$,
\[
 \gcd(M_-,M_+)=1,\quad\gcd(M_-,L)=16,\quad\gcd(M_+,L)=1.
\]
Here $64\mid L$, the exact two-parts are $2^6$ and $2^2$ before
division by four, and all odd prime divisors of $L$ were avoided.
Consequently $h\equiv0\pmod{M_-}$, $h\equiv1\pmod{M_+}$ and
$h\equiv h_0\pmod L$ are compatible when $64\mid h_0$.
The real-progression argument supplies large positive choices and
avoids the finite forbidden point and quotient loci. The later
prime-separation lemma justifies adding the factor $\Pi^{19}$.

For the anchor, repeat the proof of
Lemma~\ref{lem:v5-three-witness-anchor} with $C=3^{2^{25}}$ and
$\mathsf S=C^2=S_0$. Its quaternion criterion and anisotropic
binary form are unchanged. Thus the two rational coefficients of
$f(R)/C$ belong to $\Z_{(2)}\cap\Z_{(3)}$. The exact local calculation
of Section~\ref{subsec:o8-anchor-subgroup} now gives
$a-1\in2\,3^{2^{24}-1}\Z\setminus\{0\}$, and
\[
 \log_2|a-1|\ge1+(2^{24}-1)\log_2 3
 >1+\tfrac32(2^{24}-1)=25165823.5.
\]
For the scalar proof the new label is at least three, its five-contact
product never vanishes, and all its negative valuations at the
consumed odd primes are even. These are exactly the properties of
$X^2+2$ used in Lemma~\ref{v5:scalar32-bridge}; they hold for
$2X^2+3$ by the calculation below. That proof therefore gives
$M_5(\beta)^{32}\mid a-1$. Applying
Theorem~\ref{thm:quartic-height19} gives the asserted exponent $19/32$.
Its finite correction and continuity arguments apply to the same
norm tree with the scalar point $R$; choosing $a-1$ sufficiently
divisible meets the anchor and any further prescribed divisibility.
\end{proof}


\subsection{An exact replacement of the point chart}
Put $K=\Q(\sqrt2)$, $\delta=3-\sqrt2$ and $D=\delta^2=11-6\sqrt2$. The old chart, with letters in $K$, is
\begin{equation}\label{v9chart:eq:old}
 u=t^2,\qquad x=h^2,\qquad u(x-u/4)=D,\qquad y=h(u/2-x).
\end{equation}
The new chart is
\begin{equation}\label{v9chart:eq:new}
 x=h^2,\qquad j^2=x^2-D,\qquad y=hj.
\end{equation}
An equation over $K$ means its two rational component equations. The last expression in \eqref{v9chart:eq:new} need not be assigned its own coordinates.

\begin{lemma}\label{v9chart:lem:squares}
If $U,V\in K$, $U\ne0$, and $V^2=U^3+4DU$, then $U$ is a square in $K$.
\end{lemma}
\begin{proof}
The ring $\mathcal O=\Z[\sqrt2]$ is a PID: rounding both coefficients of an element of $K$ to nearest integers leaves a remainder of absolute norm at most $1/2$, giving a Euclidean algorithm. Its units are $\pm(1+\sqrt2)^n$, and every totally positive unit is a square. For completeness, multiply a positive unit by a power of $1+\sqrt2$ to place its first real embedding in $[1,1+\sqrt2)$. Its other embedding is the positive or negative reciprocal. Writing the unit as $a+b\sqrt2$ then forces $a,b\in\{0,1\}$, and the norm condition and excluded upper endpoint leave only the unit~1. The two signs of the fundamental unit under the real embeddings give the assertion about totally positive units.

Let $m=v_{\mathfrak p}(U)$, with integral normalization. Away from $(\sqrt2)$ and $(\delta)$ the two terms $U^3$ and $4DU$ have valuations $3m$ and $m$. If $m$ is odd, their unequal minimum is odd, a contradiction. At $(\sqrt2)$ these valuations are $3m$ and $m+4$, with equality only at the even value $m=2$. Again $m$ cannot be odd. At $(\delta)$ they are $3m$ and $m+2$, with equality only at $m=1$. In that exceptional case write $U=\delta u$ with $u$ a local unit. Then
\[
 V^2=\delta^3u(u^2+4).
\]
The residue field is $\mathbf F_7$, in which $-4$ is a nonsquare. Thus the right side has valuation three, also impossible. Every finite valuation of $U$ is even.

Under each real embedding, $D>0$ and $U(U^2+4D)=V^2$, so $U>0$. By the PID property, $U=\gamma^2\eta$ for $\gamma\in K^*$ and a unit $\eta$. This unit is totally positive, hence a square.
\end{proof}

\begin{proposition}\label{v9chart:prop:chart}
The charts \eqref{v9chart:eq:old} and \eqref{v9chart:eq:new} have exactly the same projected pairs $(x,y)\in K^2$. Every solution of the new chart has $h,j,x,y\ne0$, and $x$ is positive under both real embeddings.
\end{proposition}
\begin{proof}
Given \eqref{v9chart:eq:old}, put $j=t^2/2-h^2$. Then
\[
 j^2=h^4-t^2h^2+t^4/4=h^4-D,
\]
and $y=hj$. Conversely, given \eqref{v9chart:eq:new}, set $U=2(h^2+j)$ and $V=2hU$. The identity $j^2=h^4-D$ gives
\[
 U^2-4h^2U+4D=0,
\]
hence $V^2=U^3+4DU$. Also $U\ne0$, since $U=0$ would imply $D=0$. By Lemma~\ref{v9chart:lem:squares}, choose $t\in K$ with $t^2=U$. It follows that
\[
 t^2(h^2-t^2/4)=D,\qquad h(t^2/2-h^2)=hj,
\]
as required.

If $h=0$, then $j^2=-D$, contradicting either real embedding. If $j=0$, then $h^4=D$, so $h^2=\pm\delta$. Taking norms yields $\Norm(h)^2=7$, impossible for $\Norm(h)\in\Q$. The remaining assertions follow immediately.
\end{proof}

The proof is specific to this field and this value of $D$. It does not use the generally false assertion that square abscissa characterizes doubles on every elliptic curve. It also avoids adding a Mordell--Weil rank or generator hypothesis to the chart replacement.

\subsection{Using the chart without ordinate variables}
All nine point modules use the new chart. It has six rational coordinates: the two components of each of $h,j,x$, subject to four quadratic component equations. The old module had ten coordinates, including the components of $t,h,u,x,y$. Thus the change saves four coordinates per point, provided each downstream ordinate is consumed as the expression $hj$.

For the difference of chart points $Q,R$, retain the existing inverse equations for $x_Q-x_R$ and define the slope by
\[
 (x_Q-x_R)m=h_Qj_Q+h_Rj_R.
\]
Its numerator and defining equation are quadratic in the available coordinates. Define the difference abscissa by the existing square of $m$, and its ordinate by
\[
 y_{Q-R}=m(x_Q-x_{Q-R})-h_Qj_Q.
\]
This again is a quadratic assignment. The retained inverse ensures that the slope has exactly its old value.

For a chart point,
\[
 z=-x/y=-h/j.
\]
Thus the local parameter is defined by $jz=-h$. When a logarithm value is combined with its parameter, replace $y\ell=-xG$ by $j\ell=-hG$, where $G$ is the unchanged truncated logarithm polynomial. At the scalar ratio replace
\[
 x_H\rho=-y_Hz_S\qquad\text{by}\qquad h_H\rho=-j_Hz_S.
\]
All cancellations are reversible by Proposition~\ref{v9chart:prop:chart}. Difference-point denominators and all other guards remain unchanged. The fixed target has
\[
 h=-1+\tfrac74\sqrt2,\qquad j=\tfrac{39}{8}-\tfrac92\sqrt2;
\]
these give exactly its previous abscissa and ordinate.

Together with Proposition~\ref{v9chart:prop:chart}, these identities give rational witness maps in both directions. To reconstruct the old chart, choose the square root supplied by Lemma~\ref{v9chart:lem:squares}; then evaluate all old assignments. The choice of sign does not alter $(x,y)$ or any downstream value.

\subsection{Shorter weight circuits}\label{subsec:v9-laurent}
For each of nine detectors, write $e$ for its nonzero weight and $q$ for its order-one weight. The required outputs are $q_{12}=qe^{22}$ and $q_{40}=qe^{78}$. The old construction uses seven pure powers of $e$ and two final quotients. Instead construct
\[
 a_1=e^{-1},\quad a_3=a_1^2/e,\quad a_7=a_3^2/e,\quad
 a_{14}=a_7^2,\quad a_{29}=a_{14}^2/e,\quad a_{55}=a_{29}^2/a_3,
\]
and impose
\[
 a_{29}q_{12}=qa_7,\qquad a_{55}q_{40}=q_{12}e.
\]
These are six power coordinates and two output coordinates. They give the same outputs and save one coordinate per detector.

The nonzero condition is inherited, rather than presumed from a quotient notation. In the source, $w_0$ has a retained inverse and
\[
 w_0n=(w_0^2-2w_1^2)/2,\qquad (e\Delta)b=n^2.
\]
Since $w_0\ne0$ and 2 is a nonsquare in $\Q$, the first equation gives $n\ne0$. The second gives $e\Delta\ne0$. Therefore every Laurent denominator above is nonzero. Conversely all Laurent coordinates are uniquely defined for each source solution. The changes introduce no new zero branches and exclude no source solution.

\subsection{Scalar preparation with one fewer coordinate}\label{subsec:v9-preparation}
The scalar pole uses a nonzero input $b$ and a preparation witness $t$, with
\[
 U=\frac{b(1-bt^2)}{1+bt^2}.
\]
The old circuit stores $b^{-1},t_2=t^2,bt_2,U$. Replace these four coordinates by $t_2,T,V$, with
\[
 t_2=t^2,\qquad b(T-t_2)=1,\qquad TV=2,
\]
and use the affine expression $U=V-b$. The equation $b(T-t_2)=1$ forces $b\ne0$, and $TV=2$ forces $T\ne0$. Consequently $T=t^2+b^{-1}$, $1+bt^2=bT\ne0$, and
\[
 U=2/T-b=\frac{b(1-bt^2)}{1+bt^2}.
\]
Conversely every old solution extends uniquely by these values of $T,V$. Thus the modification saves one coordinate and preserves all downstream quantities. It also supplies the nonzero condition formerly imposed by the separate coordinate $b^{-1}$.


\subsection{An exact product-height identity}\label{v9:sec:height}
Put
\[
 K=\Q(\sqrt2),\quad d=3-\sqrt2,\quad D=d^2=11-6\sqrt2,
 \qquad E:y^2=x^3-Dx,
\]
and put
\[
 R_0=(1+\sqrt2,2\sqrt2),\qquad
 P=2R_0=\left(\frac{57-28\sqrt2}{8},\frac{-660+417\sqrt2}{32}\right).
\]
The model calculation in Section~\ref{sec:degree8-original} proves that $P=P_*$ has infinite order and $2E(K)=\Z P_*$; this is the same normalized curve and generator, not the earlier model of Proposition~\ref{prop:generator}. Write $\operatorname{h}$ for the absolute logarithmic Weil height. The following lemma uses the displayed coordinates and doubling identity, together with the infinite order of $P$.

\begin{lemma}\label{v9:lem:product}
For odd integers $n_1,\ldots,n_r$,
\[
 \operatorname{h}\!\left(\prod_{i=1}^r x([n_i]P)\right)
 =\sum_{i=1}^r\operatorname{h}(x([n_i]P)).
\]
\end{lemma}
\begin{proof}
Let $T_0=(0,0)$. At a finite place $v$ with $v(D)=0$, use the finite quotient $E(K_v)/E_1(K_v)$, where $E_1$ denotes the formal neighborhood of infinity for this integral equation.
Here we use the integral formal group construction of
\cite[Chapter IV, Sections 1--3]{silverman}. Its parameter is $z=-x/y$,
and $x=z^{-2}(1+O(z))$, $y=-z^{-3}(1+O(z))$ identify this neighborhood
with the maximal ideal. It is an open subgroup; the proper curve
$E(K_v)$ is compact, so its quotient is finite. This argument does not
assume that the equation is minimal at the place above two. For a finite nonidentity point $Q$, membership in $E_1$ is equivalent to $v(x(Q))<0$. Since
\[
 x(Q+T_0)=-D/x(Q),
\]
positive valuation of $x(Q)$ is equivalent to $Q+T_0\in E_1$. The image of $T_0$ in this quotient has order two.

If an odd multiple of $P$ belongs to $E_1$, the image of $P$ has odd order. None of its multiples can then equal the image of $T_0$, so all $x([n_i]P)$ have nonpositive valuation. If no odd multiple belongs to $E_1$, all these abscissae have nonnegative valuation. Thus the valuations have a common weak sign. This argument includes the place above two.

The only place with $v(D)>0$ is $\mathfrak d=(3-\sqrt2)$ above seven. The displayed integral model is minimal there, since $v_{\mathfrak d}(-D)=2<4$. Its reduction is $y^2=x^3$ over $\F_7$. Substituting $\sqrt2=3$ shows that $P$ reduces to $(1,6)$, a nonsingular point. By \cite[Proposition VII.2.1, pp. 188--190]{silverman}, all multiples therefore belong to the subgroup $E_0(K_{\mathfrak d})$ of points with nonsingular reduction. A positive valuation of $x$ would force a positive valuation of $y$, and hence reduction to the singular point $(0,0)$. It cannot occur. The valuations again have a common weak sign.

At either real embedding, each $[n_i]P$ is a double of a non-torsion point. The duplication formula makes its abscissa a nonzero square. Its ordinate is nonzero, so the equation of $E$ gives $x([n_i]P)>d$ at that embedding. Both real values $3\pm\sqrt2$ of $d$ exceed one.

At a finite place, numbers whose valuations have the same weak sign satisfy
\[
 \log\max\left(1,\left|\prod_i x([n_i]P)\right|_v\right)
 =\sum_i\log\max(1,|x([n_i]P)|_v).
\]
At both real places the same equality holds because all the factors exceed one. Summing with normalized local weights proves the assertion.
\end{proof}

The product here is an element of $K$. It is not its norm down to $\Q$. Norms can cancel a pole at one split place against a zero at the other, whereas the proof above works separately at every place.

\begin{corollary}\label{v9:cor:scales}
Suppose the four root points $Q_i=[n_i]P$ have odd indices, and write
\[
 C=\prod_{i=1}^4x(Q_i)=C_0+C_1\sqrt2.
\]
The original nine height-detector scales may be replaced by
\[
 1,\qquad C_0^2,\qquad (1+C_1^2)^2.
\]
The root-height estimate retains the original constant $7\mathcal H$, where $\mathcal H$ is the exponent-19 height constant of the source.
\end{corollary}
\begin{proof}
Both real values of $C$ are positive, so $C_0>0$. All three scales are positive nonzero rational squares. If a nonzero integer $N$ is divisible by the $k$th powers of the three contact radicals, the unchanged detector estimate gives
\[
 \operatorname{h}(\lambda),\ \operatorname{h}(C_0^2\lambda),\
 \operatorname{h}((1+C_1^2)^2\lambda)
 \le T:=(\mathcal H+\log3)|N|^{19/k}.
\]
The elementary rational height identities yield
\[
 \operatorname{h}(C_0)\le T,\qquad \operatorname{h}(C_1)\le T/2.
\]
Indeed, height subadditivity first bounds the heights of the two scales by $2T$; also $\operatorname{h}(C_0^2)=2\operatorname{h}(C_0)$ and $\operatorname{h}((1+C_1^2)^2)\ge4\operatorname{h}(C_1)$. Hence
\[
 \operatorname{h}(C)\le\operatorname{h}(C_0)+\operatorname{h}(C_1)+\tfrac32\log2
 \le\tfrac32T+\tfrac32\log2
 <7\mathcal H |N|^{19/k}.
\]
Here $\mathcal H\ge1$. Lemma~\ref{v9:lem:product} bounds each root height by $\operatorname{h}(C)$. In particular, the inherited lower bound, with $B=\max_i|n_i|$, gives
\begin{equation}\label{v9:eq:height-transfer}
 B^2-6<7\mathcal H |N|^{19/k}.
\end{equation}
The old completeness proof uses its scales only after the roots have been fixed, and requires fixed positive rational squares. Thus its prime selection and pole construction apply to this shorter list.
\end{proof}

\subsection{Enforcing odd indices by quadratic equations}\label{v9:sec:chart}
Put $\epsilon=1+\sqrt2$. The existing chart for a point $Q$ has $K$-coordinates
\[
 x=h^2,\qquad j^2=x^2-D,\qquad Q=(x,hj).
\]
To this chart add $t,r\in K$ and the equations
\begin{equation}\label{v9:eq:odd-chart}
 t^2=2(x+j),\qquad th+x+j=\epsilon r^2.
\end{equation}
These are four rational quadratic equations in four additional rational coordinates. No extra inverse witnesses are needed.

\begin{proposition}\label{v9:prop:odd-chart}
The projected points of the augmented chart are exactly $P+4E(K)$, that is, the odd multiples of $P$.
\end{proposition}
\begin{proof}
Use the squareclass homomorphism
\[
 \alpha:E(K)\longrightarrow K^*/K^{*2},\qquad
 \alpha(O)=1,\quad \alpha(T_0)=[-D]=[-1],\quad
 \alpha((a,b))=[a]\quad(a\ne0).
\]
For clarity, its homomorphism property follows from the line-intersection identity
\[
 x(R+S)x(R)x(S)=
 \left(\frac{x(R)y(S)-x(S)y(R)}{x(R)-x(S)}\right)^2
\]
when the displayed abscissae are nonzero and distinct. Doubling follows from
$x(2R)=((x(R)^2+D)/(2y(R)))^2$; addition by $T_0$ follows from $x(R+T_0)=-D/x(R)$. Inverse pairs and identity summands are immediate. If $R+S=T_0$, then $x(R)x(S)=-D$. These cases cover the exceptional pairs.

Since $2E(K)=\Z P$ and $2R_0=P$, every point of $E(K)$ has a unique expression $mR_0+T$, where $m\in\Z$ and $T\in E[2](K)$. The two-torsion is generated by $T_0$ and $T_d=(d,0)$, so
\[
 \alpha(mR_0+aT_0+bT_d)=[\epsilon]^m[-1]^a[d]^b.
\]
The classes $[\epsilon],[-1],[d]$ are independent. A square $\epsilon^c(-1)^a d^b$, with $a,b,c\in\{0,1\}$, would have rational square norm $(-1)^c7^b$. This forces $c=b=0$, and the real embeddings then force $a=0$. Consequently
\begin{equation}\label{v9:eq:squareclass}
 \alpha(R)=[\epsilon]\quad\Longleftrightarrow\quad R\in R_0+2E(K).
\end{equation}

Given the augmented chart, set $u=t^2=2(x+j)$ and
\[
 a=th+u/2,\qquad b=at.
\]
Expansion gives
\[
 u(x-u/4)=D,\qquad a^2-ua=D,\qquad b^2=a^3-Da.
\]
In particular $t,a,b$ are nonzero. The tangent slope at $R=(a,b)$ is $h+t$; duplication gives $2R=(x,hj)=Q$. Equation~\eqref{v9:eq:odd-chart} says $a=\epsilon r^2$, so $\alpha(R)=[\epsilon]$. By~\eqref{v9:eq:squareclass}, $Q\in P+4E(K)$.

Conversely, for $Q=(2m+1)P$, take $R=(2m+1)R_0=(a,b)$. Then $a,b\ne0$, and $a=\epsilon r^2$ for some $r\in K$. Set
\[
 t=b/a,\qquad h=(a^2+D)/(2b),\qquad x=h^2,
 \qquad j=t^2/2-x.
\]
The same duplication identities give $Q=(x,hj)$, $j^2=x^2-D$, and both equations~\eqref{v9:eq:odd-chart}.
\end{proof}

Only the four root points receive this augmentation. The shared point and the four comparison points retain their old charts: their indices need not be odd. The total cost is sixteen coordinates and sixteen quadratic residuals. The projected-point formulation matters, since the added equations can select a sign representation of the old $h,j$ chart.

\subsection{Odd labels, scalar bounds, and completeness}\label{v9:sec:labels}
Write the root points as
\[
 A=[n]P,\qquad S=[a]P,\qquad T=[t]P,\qquad Z=[b]P.
\]
Here $S$ denotes a point; the fixed coefficient remains $S_0$. Replace the old input expression by
\[
 \beta=2X^2+3.
\]
For rational $X$, integrality of $\beta$ is equivalent to integrality of $X$: if $X=u/v$ in lowest terms, then $\beta\in\Z$ forces $v^2\mid2$, and hence $v=1$. For integral $X$, $\beta$ is odd.

In the scalar norm block and the physical anchor, replace $S$ by
\[
 R=S-P=[a-1]P.
\]
The comparison and product-height blocks still use $S$. The retained subtraction guard excludes $a=\pm1$, which is harmless for completeness. The original scalar contact product is
\[
 \beta(\beta-1)(\beta-2401)(32\beta+343)(2\beta-1).
\]
It never vanishes for rational $X$: $\beta\ge3$, and $\beta=2401$ would give $X^2=1199$, which is not a rational square. At an odd prime, a negative valuation of $\beta$ is twice that of $X$, as required by the scalar argument.

\subsubsection{Soundness}\label{subsec:v9-soundness}
Put $c=|a-1|$. Applying the scalar and anchor lemmas to $R$ gives
\[
 \operatorname{h}(\beta)\le\mathcal H c^{19/32},\qquad
 \log_2c>25{,}165{,}823.5,
 \qquad 1\le\mathcal H<2^{2{,}093{,}966}.
\]
The anchor is still $S_0=3^{2^{26}}$. These explicit bounds imply
\[
 c>(80\mathcal H)^{10},\qquad
 |a|\ge c-1\ge c/2>(40\mathcal H)^{10},
\]
and
\begin{equation}\label{v9:eq:margins}
 \operatorname{h}(\beta)<c/80\le |a|/40,
 \qquad\log(40\mathcal H)<c/80\le |a|/40.
\end{equation}
For example $10(7+2{,}093{,}966)<25{,}165{,}823.5$, and $80<2^7$. The first estimate in~\eqref{v9:eq:margins} follows from $c^{13/32}>80\mathcal H$; the logarithmic estimate is weaker still.

The local comparison lemmas retain their original conclusions. Let
\[
 F=t^2-a(a-2)tb+(a-2)^2b^2-1,
 \qquad B=\max(|n|,|a|,|t|,|b|).
\]
For $B>4$, the triangle inequality gives $|F|\le2B^4+6B^3+5B^2+1<5B^4$. If $F\ne0$, the three radicals divide it to order 40, so~\eqref{v9:eq:height-transfer} gives
\[
 B^{1/10}<7\mathcal H\,5^{19/40}+6<24\mathcal H.
\]
This contradicts the anchor. Thus $F=0$. Since the point $Z$ is finite and nonidentity, $b\ne0$. Lemma~\ref{lem:o8-lucas} now gives $\log B>|a|-3$.

If $n\ne\beta$, the nonzero numerator $N$ of $n-\beta$, in reduced form, satisfies $|N|\le2e^{\operatorname{h}(\beta)}B$, and the radicals divide it to order 12. Equation~\eqref{v9:eq:height-transfer} and $B>4$ give
\[
 B^{5/12}<40\mathcal H e^{(19/12)\operatorname{h}(\beta)},
\]
so
\[
 \log B<\frac{12}{5}\log(40\mathcal H)+\frac{19}{5}\operatorname{h}(\beta)
 <\frac{31}{200}|a|<|a|-3.
\]
This is the other contradiction. Therefore $n=\beta\in\Z$, and $X\in\Z$.

\subsubsection{Completeness}\label{subsec:v9-completeness}
Fix $X\in\Z$, and set $n=\beta$. Perform the old scalar preparation: choose its finite correction and then a positive shared index $h_0$ divisible by 64 and sufficiently deep at the fixed small primes. Choose a positive scalar index $s$ sufficiently divisible to satisfy the scalar leaves and physical anchor for $[s]P$. We may require $6\mid s$ and take $s$ arbitrarily large. Set $a=s+1$. Then $a$ is odd, $a\equiv1\pmod3$, and $S-P=[s]P$ is exactly the prepared scalar point.

Define
\[
 U_0=0,\qquad U_1=1,\qquad U_{j+1}=aU_j-U_{j-1},
 \qquad k=2(a-2),
\]
and take $b=U_k/(a-2)$ and $t=U_{k+1}$. Modulo $a-2$, induction gives $U_j\equiv j$, so $b\in\Z$. Modulo two the sequence has period $0,1,1$; since $k\equiv1\pmod3$, both $U_k$ and $U_{k+1}$ are odd. The positive odd denominator $a-2$ shows that $b$ is odd. The Lucas identity gives
\[
 t^2-a(a-2)tb+(a-2)^2b^2=1.
\]
Thus all four root indices $n,a,t,b$ are positive and odd. Proposition~\ref{v9:prop:odd-chart} supplies their chart witnesses.

Fix these roots and the three scales of Corollary~\ref{v9:cor:scales}. Local continuity of the scalar and comparison functions at the prepared shared point gives the same fixed modulus $L$, with $64\mid L$. Only then choose the contact primes by the inherited simultaneous-prime theorem. The direct-target construction chooses a common index $h$ satisfying the primary and opposite reduction-order congruences and $h\equiv h_0\pmod L$, with the required real bounds. Its comparison points are $H=[h]P$ and $U_i=[hn_i]P$. Their shifted points have indices $h-1$ and $(h-1)n_i$.

All quotient-domain arguments remain valid. The product scales are fixed positive squares, so the original pole witnesses apply unchanged. The scalar block is preserved by the congruence modulo $L$. This gives a rational solution for every integral $X$ and completes the arithmetic proof of the primary construction.

\subsection{Replacing five logarithms by asymmetric formal depth}\label{v9:sec:depth}
This section is used only for the 523-variable construction. Let
\[
 U_0=H=[h]P,\quad U_i=[u_i]P,\quad
 V_0=H-P,\quad V_i=U_i-Q_i,
\]
where $Q_i=[n_i]P$ are the four roots. Put $f=1/x$ and $z=-x/y$. Keep the alternate truncated logarithms $\ell_{40}(V_i)=z(V_i)G_9(Df(V_i)^2)$, but replace every primary truncated logarithm by $z(U_i)$. Thus
\[
 L_i=z(U_i)/z(H),\qquad
 L_i'=\overline{\ell_{40}(V_i)/\ell_{40}(V_0)}.
\]
The bar denotes conjugation in $K$. The sixteen tests now have the following orders:
\[
\begin{array}{c|l}
1 & f(U_1),\ldots,f(U_4),\ \overline{f(V_0)},\ldots,\overline{f(V_4)}\\[2pt]
40 & f(H),\ L_1-L_1',\ldots,L_4-L_4',\ L_1-\beta,\
       \text{the direct Lucas residual}.
\end{array}
\]
The scalar tests retain their original order 32. The stronger beta tests imply all former order-12 conclusions.

\subsubsection{Local soundness}\label{subsec:v9-local-depth}
At a consumed unramified place above $p>76176$, a formal point $Q$ with $r=v(z(Q))\ge1$ satisfies
\begin{equation}\label{v9:eq:zerror}
 \log_E(Q)/z(Q)\in1+p^{4r}\mathcal O.
\end{equation}
This is the same formal-logarithm expansion used in the source, starting its tail at the first nonconstant term. Namely $v(f(Q))=2r$, and for its coefficients $a_j$ the factorial estimate is $v(a_j)\ge-(4j+1)/(p-1)$. If $4j+1<p$, the coefficient is integral. For the remaining $j\ge19044$, the displayed bound still makes the $j$th term have valuation at least $4r$. Both $z(Q)$ and its full logarithm have valuation $r$.

The new test on $f(H)$ gives $v(z(H))\ge20$. The order-one tests put the other points on the formal branch. The full logarithm ratios at the primary and conjugate places are
\[
 r_i=u_i/h,\qquad r_i'=(u_i-n_i)/(h-1).
\]
At least one of $h,h-1$ is a $p$-unit, so at least one of these ratios is integral. Even before knowing the depth of $U_i$, the weak unit estimate~\eqref{v9:eq:zerror} and the alternate truncation estimate show that $L_i,L_i'$ have the same valuations as $r_i,r_i'$. The order-40 cross test transfers integrality to both sides. Now $L_i$ integral implies $v(z(U_i))\ge v(z(H))\ge20$. The primary logarithmic corrections are therefore $1$ modulo $p^{80}$. The alternate corrections are $1$ modulo $p^{40}$. Since the ratios are integral, these relative estimates become absolute:
\[
 L_i\equiv r_i\equiv r_i'\pmod {p^{40}}.
\]
Multiplying by $h$ and $h-1$ and subtracting gives $L_i\equiv n_i\pmod {p^{40}}$. At an inert prime conjugation preserves the local valuation, so the same proof applies. The height contradictions of Section~\ref{v9:sec:labels} remain unchanged.

\subsubsection{A stronger choice within the existing prime-pattern argument}
For completeness, the common index must acquire extra primary divisibility. The following observation ensures compatibility with the opposite reduction orders.

\begin{lemma}\label{v9:lem:separation}
The contact primes in the source completeness argument may be chosen so that no selected prime divides any opposite reduced group order.
\end{lemma}
\begin{proof}
After merging repeated contact forms and fixing all scale and congruence data, the prime-pattern proof uses distinct forms $p_j=u+c_jw^2$, with fixed coefficients $c_j$. Integer dilations and ray congruences change only fixed constants. It counts $u\in[N/2,N]$ and $1\le w\le T$, where $T=\lfloor\tau(N)\sqrt N\rfloor$ and $\tau(N)\to0$ can be chosen arbitrarily slowly in the final diagonal argument. Its weighted count, after the prescribed obstruction removal, retains a fixed positive proportion of $NT$. These are the fixed-field weighted count in Proposition~\ref{prop:v5-artin-patterns} and the final diagonal argument of Section~\ref{subsec:v5-diagonal}, including its obstruction removal.

If there are $k$ forms, impose also
\[
 \frac{\tau(N)N^{1/6}}{(\log N)^k}\longrightarrow\infty.
\]
This is compatible with the same diagonal choice: for each fixed positive value of $\tau$, its threshold can be added to the existing finite list of thresholds. The strip $w<N^{1/3}$ has at most $O(N^{4/3})$ integer pairs and weight at most $O((\log N)^k)$ per pair. Its total weight is $o(NT)$, so an admissible pattern remains outside it.

For these patterns $w^2\ge N^{2/3}$ and $w^2=o(N)$. All $p_j$ are asymptotic to the same $u$. Hasse's bound \cite[Theorem V.1.1, p. 138]{silverman} gives $d_+(p_j)=p_j+O(\sqrt N)$ for the opposite reduced group order. Every ratio $d_+(p_j)/p_i$ is eventually strictly between $1/2$ and $2$. Hence $p_i\mid d_+(p_j)$ would force equality. If $i\ne j$, this contradicts
\[
 |p_i-p_j|\ge c_*N^{2/3},\qquad
 c_*:=\min_{i\ne j}|c_i-c_j|>0.
\]
If $i=j$, equality is impossible because $4\mid d_+(p_i)$ whereas $p_i$ is odd.
\end{proof}

\subsubsection{Completeness with deeper primary points}\label{subsec:v9-deep-completeness}
Keep the choice order in Section~\ref{v9:sec:labels}: the roots, scales and continuity modulus $L$ precede the prime selection. Choose primes using Lemma~\ref{v9:lem:separation}, beyond all prime factors of $L$ and the fixed scale data. Put
\[
 M_\pm=\mathop{\mathrm{lcm}}_p(d_\pm(p)/4),\qquad \Pi=\prod_p p.
\]
The direct-target prime theorem gives
\[
 \gcd(M_-,M_+)=1,\quad
 \gcd(M_-,L)=16,\quad\gcd(M_+,L)=1.
\]
Lemma~\ref{v9:lem:separation} gives $\gcd(\Pi,M_+)=1$, and the cutoff gives $\gcd(\Pi,L)=1$. Thus
\[
 h\equiv0\pmod{\Pi^{19}M_-},\qquad
 h\equiv1\pmod{M_+},\qquad h\equiv h_0\pmod L
\]
are compatible, since $64\mid h_0$. The inherited real-progression argument supplies a sufficiently large $h$ in this progression with the same real bounds.

The point $P$ is globally a double, so multiplication by the primary doubling-image order first makes it formal, with $z$-valuation at least one. Multiplication by $p^{19}$ raises this valuation by nineteen, using the formal logarithm isometry of \cite[Theorem IV.6.4, p. 132]{silverman}; these completions are unramified, $v(p)=1$, and $1>1/(p-1)$. Consequently $H=[h]P$ and $U_i=[hn_i]P$ have primary depth at least twenty, and their parameter ratios equal $n_i$ to order eighty. At the opposite places, $V_0=[h-1]P$ and $V_i=[(h-1)n_i]P$ are formal and their retained logarithm ratios equal $n_i$ to order forty.

The pole-completeness estimates still apply to these fixed test orders. In the source notation their large-real-parameter decay is $\lambda^{-3}$ for each order $m$, and the small-prime lower bound is $3\nu+(10m-8)\eta-2c$, which still tends to infinity with the preparatory divisibility. Scalar conditions and the anchor are preserved by $h\equiv h_0\pmod L$. This proves completeness for the optional refinement.

\subsection{Polynomial construction and exact verification}\label{subsec:v9-finite}
The generators apply three changes to the earlier point-chart source: the four odd-coset charts, the product and three scales, and the shifted scalar point with $\beta=2X^2+3$. The depth variant first changes the logarithms and test orders as in Section~\ref{v9:sec:depth}. Each field equation is represented by two rational component equations. Intermediate products are assigned coordinates as needed to keep every residual quadratic.

The full-log source has 577 raw coordinates; the pinned upstream compiler reduces this to 561. Seven further certified eliminations give 554. The depth source has 540 raw coordinates; compilation gives 532 and nine certified eliminations give 523. Every elimination is justified by exact polynomial identities and reversible substitutions preserving the input. Clearing residual denominators and taking the sum of their squares gives the displayed integer quartics. Over $\Q$, the sum vanishes exactly when every residual does.

The constant-chain versions add $c_0,\ldots,c_{16}$ with
\[
 c_0=3^{1024},\qquad c_{i+1}=c_i^2\quad(0\le i<16).
\]
These witnesses have unique rational values and $c_{16}=S_0$. Substitution for $S_0$ leaves every residual quadratic. Direct expansion gives the 3,247-bit maximum coefficient bound.

The accompanying archive contains the complete polynomial tables, residual systems, coordinate maps, elimination certificates, pinned source generators, arithmetic notes and independent checkers. Each final package has a standard-library-only \texttt{verify\_all.py}. It replays all eliminations, verifies the unique constant lift, independently expands the sums of squares, and compares every coefficient. Both full packages have passed fresh independent runs. A separate source replay rebuilds the two compiled precursors and checks the match to their frozen residual systems.

These finite checks certify the counts, degrees and stated algebraic transformations. The arithmetic assertions in Sections~\ref{v9:sec:height}--\ref{v9:sec:depth} received separate adversarial AI review; they have not received independent human peer review or a complete formal proof check. The inherited arithmetic inputs are the results of this manuscript identified in Section~\ref{subsec:v9-inputs}; no theorem from the contribution archive is assumed as an additional arithmetic input.



\begin{proof}[Proof of Theorem~\ref{thm:v9-quartics}]
The chart identities and Lemma~\ref{lem:v9-bridge} identify the exact
source used by the generators with the arithmetic system described here.
Proposition~\ref{v9:prop:odd-chart} enforces the odd root indices, so
Corollary~\ref{v9:cor:scales} supplies the required height bound.
Section~\ref{subsec:v9-local-depth} verifies the local implications for
the final test orders. Section~\ref{subsec:v9-soundness} then forces
the input to be an integer. Conversely, for an integer input,
Sections~\ref{subsec:v9-completeness} and
\ref{subsec:v9-deep-completeness}, using Lemma~\ref{v9:lem:separation},
construct witnesses for every source equation and retained guard.
The certified substitutions preserve this projection in both directions.
The sum of squares vanishes over $\Q$ precisely when all its residuals
vanish. Finally the 17-coordinate constant chain has exactly one rational
solution and ends at $S_0$, so it preserves the same projection. This
proves the statements for both polynomials.
\end{proof}

The complete tables are in the files
\begin{center}\small
\nolinkurl{quartics/quartic-523-symbolic-coefficients.json}\\
\nolinkurl{quartics/quartic-540-fully-expanded.txt}.
\end{center}
Their SHA-256 digests, respectively, are
\begin{center}\small
\nolinkurl{6a3685980de4512bdc7f3a781d0c0d4b32cce5829ec83f03f318264931bdbbe3}\\
\nolinkurl{83e7a31e11f9ad2edc2f67935e50bff437d87ef4b77d37b6a4391e8117770680}.
\end{center}
The reproduction command is
\nolinkurl{python scripts/check_odd_product.py}; the complete source,
certificates and repaired, non-mutating verifiers are in
\nolinkurl{quartics/quartic-523-540/}.

% END VERSION 9 ADDITION

\section{A quartic in 498 variables}\label{sec:v12}

We now reduce the 523-variable construction of
Theorem~\ref{thm:v9-quartics} to 498 variables.  The input remains one
rational variable $X$; every other variable remains existentially
quantified over $\Q$.  The changes below use the same normalized curve,
the same integer labels, the same three height detectors, and the same
published and manuscript inputs as Section~\ref{sec:v9}.  In particular,
they do not replace the restricted converse, the exponent-nineteen
height theorem, or the simultaneous-prime argument by new assumptions.

\begin{theorem}\label{thm:v12-quartic}
There is a specified integer polynomial $P_{498}$ of total degree four
in 498 variables such that
\[
 X\in\Z\quad\Longleftrightarrow\quad
 \exists\mathbf v\in\Q^{497}:P_{498}(X,\mathbf v)=0.
\]
The polynomial has 24,669 nonzero monomials, content one, and largest
absolute coefficient $16\cdot3^{2^{26}}$.  Its symbolic coefficient
table uses the fixed integer $C=3^{2^{25}}$, not an additional variable.
There is also a specified integer quartic $P_{514}$ with
\[
 X\in\Z\quad\Longleftrightarrow\quad
 \exists\mathbf w\in\Q^{513}:P_{514}(X,\mathbf w)=0.
\]
It has 24,716 nonzero monomials and content one. Its largest absolute coefficient is
\[3^{2048}+1005072314676622<2^{3247}.\]
\end{theorem}

The proof distinguishes two sorts of changes.  Some are rational
changes of private witnesses with a rational inverse.  Others alter
the arithmetic certificate while preserving both soundness and the
existence of witnesses for every integer input.  For the latter we
prove those two assertions separately; we do not claim equality of
all intermediate witness assignments.

\subsection{Conventions and retained inputs}\label{subsec:v12-inputs}

Throughout put
\[
 K=\Q(\rho),\quad \rho^2=2,\quad
 d_*=3-\rho,\quad D=d_*^2=11-6\rho,\quad
 E:y^2=x^3-Dx.
\]
The point $P=P_*$ is the same normalized generator as in
Section~\ref{subsec:v5-point-source} and Section~\ref{sec:v9}:
\[
 P=(c_*,a_*b_*),\qquad
 a_*=-1+\tfrac74\rho,\quad
 b_*=\tfrac{39}{8}-\tfrac92\rho,\quad
 c_*=a_*^2=\tfrac{57}{8}-\tfrac72\rho.
\]
Thus $b_*^2=c_*^2-D$.  In particular $a_*,b_*,c_*,D$ are nonzero.
We use the inherited equality $2E(K)=\Z P$, the nonzero point domains
of Lemma~\ref{lem:v5-point-domains}, and the inherited estimate
$\operatorname h(x([n]P))\ge n^2-6$.  Here $\operatorname h$ is the
absolute logarithmic Weil height.  At every unramified place above a
rational prime $p$ we normalize $v(p)=1$.  The coefficient basis is
always $(1,\rho)$ and conjugation sends $\rho$ to $-\rho$.

Retain the four odd-coset root charts of
Proposition~\ref{v9:prop:odd-chart}, with labels
\[
 Q_1=[n]P,\quad Q_2=S=[a]P,\quad
 Q_3=[t]P,\quad Q_4=[b]P,
 \qquad n,a,t,b\text{ odd}.
\]
The letter $b$ in this display is an integer point label.  We use
$\mathfrak b$ below for a rational input to a pole predicate.
Put $Q_0=P$, $U_0=H=[h]P$, $U_i=[u_i]P$, and
$V_i=U_i-Q_i$.  The scalar point is
\[
 R=S-P=[s]P,\qquad s=a-1,
\]
and the input label is $\beta=2X^2+3$.

We retain the pole obstruction and pole-completeness statements
identified in Section~\ref{subsec:v9-inputs}, specifically
Lemmas~\ref{hyp:o8-pole}, \ref{lem:original-seven-simultaneous-realization},
and \ref{lem:v5-one-place}, including their existing source dependencies.
We also retain the three scales of Corollary~\ref{v9:cor:scales} and
its product-height estimate.  Write $\mathcal H=H_{19}$, so that
\begin{equation}\label{eq:v12-height-input}
 1\le\mathcal H<2^{2093966},\qquad
 B^2-6<7\mathcal H|N|^{19/k}
\end{equation}
whenever the three contact radicals divide a nonzero integer $N$
to order $k$, where \[ B=\max(|n|,|a|,|t|,|b|). \]
This is exactly \eqref{v9:eq:height-transfer}; none of its scales is
replaced by a field norm in the present construction.

The local formal-group inputs are the integral formal group in
$z=-x/y$, its multiplication series, and its logarithm.  The integral
construction is the one in \cite[IV\S1, Example~IV.2.2.3 and
Proposition~IV.2.3(a)]{silverman}.  At the unramified places used below,
$p>76176$, and the logarithm is an isometry on depth at least one by
\cite[Theorem IV.6.4, p.~132]{silverman}; its hypothesis is
$1>1/(p-1)$.  The coefficient and truncation estimates are those of
Lemma~\ref{lem:quartic-reciprocal-log} and its unramified extension in
Lemma~\ref{lem:v5-unramified-log}.  We spell out their new applications
below.  No minimal-model assertion at the prime above two is needed.

\subsection{Mixed logarithmic precision}\label{subsec:v12-mixed}

Let $f=1/x$ and $z=-x/y$, as before.  Define
\[
 a_0=1,\quad a_1=-\tfrac25,\quad
 a_j=\frac{4j-5}{4j+1}a_{j-1}\quad(j\ge2),\qquad
 G_J(T)=\sum_{j=0}^J a_jT^j.
\]
Thus the source truncations are
$\ell_{12}=zG_2(Df^2)$ and $\ell_{40}=zG_9(Df^2)$.
Keep every primary logarithm equal to the uncorrected parameter $z$.
Only the alternate numerator for $V_1$ is shortened:
\begin{equation}\label{eq:v12-ratios}
 L_i=\frac{z(U_i)}{z(H)},\qquad
 L'_1=\overline{\frac{\ell_{12}(V_1)}{\ell_{40}(V_0)}},\qquad
 L'_i=\overline{\frac{\ell_{40}(V_i)}{\ell_{40}(V_0)}}\quad(2\le i\le4).
\end{equation}
The only denominator logarithm remains $\ell_{40}(V_0)$, with its
original nonzero domain.  In particular this change introduces no
new logarithm denominator.

The sixteen field tests and their orders are now
\begin{equation}\label{eq:v12-tests}
\begin{array}{c|l}
1&f(U_1),\ldots,f(U_4),\ \overline{f(V_0)},\ldots,\overline{f(V_4)}\\[2pt]
12&L_1-L'_1,\ L_1-\beta\\[2pt]
40&f(H),\ L_2-L'_2,\ L_3-L'_3,\ L_4-L'_4,\ F(L_2,L_3,L_4),
\end{array}
\end{equation}
where
\[
 F(A,T,B)=T^2-A(A-2)TB+(A-2)^2B^2-1.
\]
This changes the two order-forty input tests in the 523-variable
source to order twelve.  The three other label comparisons and the
Lucas test remain of order forty.

\begin{lemma}\label{lem:v12-mixed-local}
At every consumed unramified place above $p>76176$, the tests
\eqref{eq:v12-tests} imply
\[
 L_1\equiv n\pmod {p^{12}},\qquad
 L_i\equiv n_i\pmod {p^{40}}\quad(2\le i\le4),
 \qquad(n_2,n_3,n_4)=(a,t,b).
\]
The primary points $U_i$ have depth at least twenty after integrality
of the ratios has been established.
\end{lemma}
\begin{proof}
The order-one tests put the indicated points in the formal group.
The test on $f(H)$ gives $v(z(H))\ge20$.  For a formal point $Q$ of
depth $r\ge1$, the unchanged logarithm expansion gives
\[
 \log_E(Q)/z(Q)\in1+p^{4r}\mathcal O.
\]
Its truncation after $G_2$ has relative error of order at least twelve,
and its truncation after $G_9$ has relative error of order at least
forty.  To check these particular applications, the first omitted
indices are respectively $j=3$ and $j=10$ in the source recurrence.
For $4j+1<p$ the coefficients are integral.  For the remaining indices
the source factorial bound
$v_p(a_j)\ge-(4j+1)/(p-1)$ is dominated by $4jr$ and gives the same
stated lower bounds.  Both truncating factors are units.

The full ratios are the rational numbers
\[
 r_i=\frac{u_i}{h},\qquad r'_i=\frac{u_i-n_i}{h-1}.
\]
At least one of $h,h-1$ is a $p$-unit, so at least one of these two
ratios is integral.  Before any improved depth is inferred for $U_i$,
the unit factors just described show
$v(L_i)=v(r_i)$ and $v(L'_i)=v(r'_i)$.  Each cross-test has positive
order, so integrality transfers to both approximate ratios and both
full ratios.  This first step is essential.

Now $L_i$ integral gives $v(z(U_i))\ge v(z(H))\ge20$.
The primary relative errors have order at least eighty.  The alternate
relative errors have order twelve for $i=1$ and forty for $i\ge2$:
in the first case the denominator still has order forty, so it does
not reduce the numerator's order twelve.  Since the ratios are
integral, these relative errors give absolute congruences.  Together
with the cross-tests they yield $L_i\equiv r_i\equiv r'_i$ to the
asserted order.  Multiply by $h$ and $h-1$ and subtract.  The result
is $L_i\equiv n_i$, without division by either integer.  At an inert
prime conjugation preserves the single local valuation, so this
argument has the same conclusion there.
\end{proof}

For clarity, the detector's product is unchanged as a mathematical
operation.  For a test $D_j$ of order $m_j$, a scale $\mathcal A$, and
$t=\mathcal A\lambda$, $\delta=t-1$, its factors are
\[
 1+\mathcal R\left(\frac{D_j}{\delta^{m_j+1}\mathcal R^{m_j}}\right)^2,
 \qquad
 \mathcal R=\mathcal A^{-5}F_5(t)/\delta^6,
 \qquad F_5(t)=\prod_{\theta\in\Theta}(t-\theta),
\]
with the same $\Theta$ as \eqref{eq:v5-detector-data}.  On writing
$e=(\mathcal A\delta)^5/F_5(t)$, the multiplier of $D_j^2$ is
$q_{m_j}=q_1e^{2m_j-2}$, where $q_1=e/\delta^3$.
For each $m\in\{1,12,40\}$, form the coefficients of
\[
 \prod_{j:m_j=m}(1+TD_j^2)
\]
and evaluate at $T=q_m$.  Their degrees are respectively $9,2,5$.
Their product is exactly the same sixteen-factor detector expression.
The smaller $G_2$ numerator and these three grouped products explain
the changed quadratic circuit.  Multiplying in $K$ means the two
rational bilinear identities in the fixed basis.  Every reassociation
or grouping here follows from distributivity, with no division or
new witness domain.  The retained coefficient tables and exact
quadratic elimination identities specify the resulting circuit.

\subsection{Five rational half-trace guards}\label{subsec:v12-guards}

For each $0\le i\le4$ write the subtraction denominator as
\[
 d_i=x(U_i)-x(Q_i)=A_i+B_i\rho.
\]
The old field inverse used two private rational coordinates.  Replace
those two equations by the single equation
\begin{equation}\label{eq:v12-half-trace}
 A_i\gamma_i=1,\qquad\gamma_i\in\Q.
\end{equation}
No other equation used the removed inverse coordinates.  This saves
one coordinate in each of five blocks.  It strengthens the domain;
there is no claim that every old trace-zero assignment is retained.

Soundness has an explicit rational map.  From \eqref{eq:v12-half-trace}
we have $A_i\ne0$, hence $d_i\ne0$ and
$A_i^2-2B_i^2\ne0$.  Restore the old field inverse as
\[
 d_i^{-1}=\frac{A_i-B_i\rho}{A_i^2-2B_i^2}.
\]
All old subtraction outputs and all their consumers are unchanged.

For completeness fix the four root points and the scalar preparation
before choosing contact primes, as in Section~\ref{subsec:v9-completeness}.
Increase the allowed finite prime cutoff to exclude every prime at
which any of $P,Q_1,\ldots,Q_4$ has nonintegral abscissa.  This is a
finite restriction on already fixed points.  Choose any one of the
split contact primes supplied by Lemma~\ref{lem:v5-two-reductions},
and denote its primary and opposite places by $v,\bar v$.
The unchanged CRT construction makes every $U_i$ formal at $v$ and
makes it reduce to $Q_i$ at $\bar v$.  All five $Q_i$ have finite
integral reduction.  Consequently
\[
 v(d_i)<0,\qquad \bar v(d_i)\ge0.
\]
Conjugation exchanges the two places, so
$v(\overline{d_i})=\bar v(d_i)\ge0$.  The two terms of
$d_i+\overline{d_i}=2A_i$ have unequal valuations; the sum is nonzero.
Thus all five inverses $\gamma_i=1/A_i$ exist for the same completeness
witnesses.  This includes $i=0$.  It imposes no unit condition on the
root labels and no condition on the permitted zero pole witnesses.

\subsection{Normalizing the three detector poles}\label{subsec:v12-poles}

Here $\mathfrak b$ is one detector's rational pole input.  In the
compiler's rational coordinates the seven relevant equations are
\begin{align*}
 p^2&=p_2,&17\mathfrak b p_2+3e^2-34e\mathfrak b-289\mathfrak b^2&=0,\\
 u_2&=\mathfrak b^2,&17e\mathfrak b-eu_2-17g&=0,\\
 s_0^2&=s_2,&q_0^2&=q_2,\\
 &&17\mathfrak b q_2+867g^2-34gs_2-s_2^2&=0.
\end{align*}
In particular the stored $g$ is the source's quantity $Bh/17$,
not $Bh$.  The following substitution respects that factor:
\[
 e=\mathfrak b E_0,\quad p=\mathfrak b P_0,\quad
 s_0=\mathfrak b T_0,\quad q_0=\mathfrak b^2Q_0,
\]
\[
 p_2=\mathfrak b^2P_2,\quad s_2=\mathfrak b^2T_2,\quad
 q_2=\mathfrak b^4Q_2,\quad
 u_2=\mathfrak b^2,\quad g=\mathfrak b^2G_0.
\]
The replacement consists of six quadratic equations:
\begin{align}\label{eq:v12-pole-normalized}
 P_0^2&=P_2,&17\mathfrak b P_2+3E_0^2-34E_0-289&=0,\notag\\
 17E_0-E_0\mathfrak b-17G_0&=0,&T_0^2&=T_2,\\
 Q_0^2&=Q_2,&17\mathfrak b Q_2+867G_0^2-34G_0T_2-T_2^2&=0.\notag
\end{align}

The division by $\mathfrak b$ in the reverse coordinate change is
legitimate in both systems.  Indeed the unchanged detector equations,
with its actual product $W=w_0+w_1\rho$, include
\[
 w_0\iota=1,\qquad
 2w_0\nu=w_0^2-2w_1^2,
\]
and an equation whose specialization at $\mathfrak b=0$ is $\nu^2=0$.
Thus $\mathfrak b=0$ would give $\nu=0$ and $w_0^2=2w_1^2$.
Rationality then forces $w_0=w_1=0$, contradicting $w_0\iota=1$.

Under the displayed substitution the old $u_2$ equation is an identity.
The old $p_2$, first-conic, $g$, and $s_2$ equations are
$\mathfrak b^2$ times their replacements; the old $q_2$ and second-conic
equations are $\mathfrak b^4$ times their replacements.  These identities
prove mutually inverse rational maps fixing every outside coordinate.
They remain valid when any of $e,p,s_0,q_0,g$ is zero.  Applying the
change to three detectors saves three coordinates.  The scalar pole
with its preparation variable is not altered by this normalization.

\subsection{One coefficient in the physical anchor}\label{subsec:v12-anchor}

Put $C=3^{2^{25}}$ and write the actual scalar value as
$f(R)=f_0+f_1\rho$.  The previous fixed parameter was $S_0=C^2$.
The old three-witness anchor encoded
$f_0/C,f_1/C\in\Z_{(2)}\cap\Z_{(3)}$.  On the present integer orbit it
is enough to impose that condition on $f_1/C$.

\begin{lemma}\label{lem:v12-anchor-coefficient}
For every nonzero even integer $s$, write
$f([s]P)=f_0+f_1\rho$.  Both coefficients are $2$-integral, and
\[
 v_3(f_1)=2+2v_3(s),\qquad
 v_3(f_0)\ge3+2v_3(s).
\]
In particular $f_1\ne0$, and
\[
 f_1/C\in\Z_{(2)}\cap\Z_{(3)}
 \quad\Longleftrightarrow\quad
 f_0/C,f_1/C\in\Z_{(2)}\cap\Z_{(3)}.
\]
\end{lemma}
\begin{proof}
At $3$ the field is unramified quadratic with residue field
$\F_9=\F_3[\rho]/(\rho^2-2)$.  The exact starting calculation in
Section~\ref{subsec:o8-anchor-subgroup} is
\[
 z(2P)=
 \frac{3396356159088-4135204317756\rho}{3287811625151},
 \qquad z(2P)/3\equiv2+\rho\pmod3.
\]
The second congruence follows by reducing the displayed rational
coefficients after division by three.  Its square is $\rho$ in
$\F_9$.  Write $s=2\cdot3^k u$ with $3\nmid u$, allowing negative $u$.
The integral formal group has inverse $-z$, and hence odd
multiplication series.  Thus
\[
 [u]z=uz+O(z^3),\qquad [3]z=3z+O(z^3)
\]
with integral coefficients.  At depth $r\ge1$, division of the second
identity by $3^{r+1}$ leaves an error of valuation at least $2r-1>0$.
Multiplication by three therefore preserves the leading residue.
Multiplication by $u$ multiplies it by the nonzero residue of $u$.
Consequently
\[
 z([s]P)/3^{k+1}\equiv u(2+\rho)\pmod3.
\]
The exact curve relation is $f=z^2(1-Df^2)$.  On the formal branch its
second factor is one modulo three, and $u^2=1$ in $\F_3$.  Therefore
\[
 f([s]P)/3^{2k+2}\equiv\rho\pmod3.
\]
Independence of $1,\rho$ over $\F_3$ proves the two asserted
coefficient valuations, including the possible zero value of $f_0$.

At the place above two normalize $v(2)=1$.  The literal point has
$v(x(P))=-3$, $v(y(P))=-9/2$, and $v(z(P))=3/2$.  It lies in the
integral formal subgroup, which contains all its integral multiples.
One may construct that subgroup directly from
$w=z^3-Dzw^2$, where $w=-1/y$, as in
Section~\ref{subsec:o8-anchor-subgroup}; this uses no assertion that
the displayed equation is minimal at two.  Thus $f([s]P)$ is integral.
The equality
\[
 v(a+b\rho)=\min\{v_2(a),v_2(b)+\tfrac12\}\quad(a,b\in\Q)
\]
shows that both rational coefficients are $2$-integral.
Since $C$ is a $2$-unit, the equivalence now follows from the
coefficient valuations at three.
\end{proof}

The three-witness criterion in the proof of
Lemma~\ref{lem:v5-three-witness-anchor} applies to every rational $t$:
\begin{equation}\label{eq:v12-quaternion-criterion}
 t\in\Z_{(2)}\cap\Z_{(3)}
 \quad\Longleftrightarrow\quad
 \exists u,v,w\in\Q:\ u^2+6v^2-2w^2-3t^2+2=0.
\end{equation}
To recall its exact input, $(2,3)_\Q$ ramifies precisely at two and
three.  The embedding criterion
\cite[Proposition 14.6.7, p.~232]{voight-quaternion} tests the nonsquareness of
$1-2/(3t^2)$ at those places when $t\ne0$.  At two, positive valuation
of $t$ gives odd negative valuation of the expression, valuation zero
gives residue three modulo eight, and negative valuation gives a
square in $1+8\Z_2$.  At three, nonnegative valuation gives odd negative
valuation of the expression, whereas negative valuation gives a square
in $1+3\Z_3$.  The expression is nonzero because $t^2=2/3$ has no
rational solution.  Multiplication by $3t^2$ gives
\eqref{eq:v12-quaternion-criterion}; for $t=0$ the witnesses are
$(0,0,1)$.  These are exactly the field, local conditions, and
exceptional case already present in that source proof.

Apply this criterion to $t=f_1/C$, multiply its witnesses by $C$,
and then replace the old $w$ by $w+C$.  The resulting anchor is
\begin{equation}\label{eq:v12-anchor}
 u^2+6v^2-2w^2-4Cw-3f_1^2=0.
\end{equation}
It uses the same three private rational witnesses and removes the
old coordinate for $f_0^2+f_0f_1-f_1^2$.  Since $a$ is odd and $R$ is
nonidentity, $s=a-1$ is a nonzero even integer.  By
Lemma~\ref{lem:v12-anchor-coefficient}, old and new anchors have the
same projection onto all their nonprivate coordinates.  In each
direction use \eqref{eq:v12-quaternion-criterion} to choose the private
witnesses; they need not be the same witnesses in both anchors.
The affine translation of $w$ is invertible.

The exact allowed subgroup is unchanged:
\begin{equation}\label{eq:v12-anchor-subgroup}
 s\in2\,3^{2^{24}-1}\Z\setminus\{0\},\qquad
 \log_2|s|\ge1+(2^{24}-1)\log_2 3>25165823.5.
\end{equation}
Indeed $v_3(f_1)\ge v_3(C)=2^{25}$ is precisely
$2+2v_3(s)\ge2^{25}$.  The subgroup permits arbitrarily divisible
positive $s$, as required by completeness.

\subsection{A shorter scalar norm tree}\label{subsec:v12-scalar}

Keep the exact contact polynomial and finite correction
\[
 r=\beta(\beta-1)(\beta-2401)(32\beta+343)(2\beta-1),\qquad
 c\in\left\{\pm\prod_{p\in I}p:
 I\subseteq\{2,3,5,7,11,13,17\}\right\}.
\]
Here $r\ne0$ for rational $X$.  The only positive contact not excluded
by $\beta\ge3$ is $\beta=2401$, which would imply $X^2=1199$,
impossible over $\Q$.  The old inverse equation for $r$ and its
power chain remain in place.

Write
\[
 f(H)=h_0+h_1\rho,\qquad z(R)/z(H)=z_0+z_1\rho,
 \qquad q=h_0^2-2h_1^2=\operatorname{Norm}_{K/\Q}(f(H)).
\]
Use $N_r(u,v)=u^2+rv^2$ and replace the old six-leaf scalar tree by
\begin{equation}\label{eq:v12-scalar-tree}
 W=N_r\left(N_r(1,q/r),\ N_r(z_0/r^{32},z_1/r^{32})\right).
\end{equation}
The scalar pole input is still exactly $\mathfrak b=crW^2$, with the
original preparation and nonzero condition.  This tree changes that
input's value, so completeness must supply new pole witnesses.

For an explicit quadratic encoding, let $i_{32}=r^{-32}$ be the
retained power-chain output and introduce eight rational coordinates
$q,A,U,V,V_2,B,B_2,W$ with equations
\begin{equation}\label{eq:v12-scalar-eight}
\begin{aligned}
 q&=h_0^2-2h_1^2,& r(A-1)&=q^2,\\
 U&=z_0i_{32},&V&=z_1i_{32},\\
 V_2&=V^2,&B&=U^2+rV_2,\\
 B_2&=B^2,&W&=A^2+rB_2.
\end{aligned}
\end{equation}
The existing $W^2$ coordinate is reattached to this $W$ with its
actual fixed rational normalization.  Thirteen old leaf and norm
coordinates disappear, giving a net saving of five at this stage.
The final elimination of $q$ is given below.

\begin{lemma}\label{lem:v12-scalar-sound}
The new scalar certificate implies
$M_5(\beta)^{32}\mid s$ and
$\operatorname h(\beta)\le\mathcal H|s|^{19/32}$.
\end{lemma}
\begin{proof}
Let $p>76176$ be a prime of the same five-contact radical as in
Lemma~\ref{v5:scalar32-bridge}, and put $\eta=v_p(r)$.
The unchanged contact calculation makes $\eta$ odd; the correction
$c$ is a $p$-unit.  The original pole obstruction, applied to the
retained nonzero input $crW^2$, gives
\[
 \eta+2v_p(W)>0.
\]
At a node $u^2+rv^2$, nonzero summands have opposite valuation parity,
so they cannot cancel.  Zero summands are harmless with valuation
$+\infty$.  If $\eta<0$, the constant-left path in
\eqref{eq:v12-scalar-tree} gives $v_p(A)\le0$ and $v_p(W)\le0$, a
contradiction.  Thus $\eta>0$ and $v_p(W)>-\eta/2$.
Descending the two levels gives
\[
 v_p(A)>-\eta/4,\quad v_p(B)>-3\eta/4,
\]
\[
 v_p(q/r)>-5\eta/8,\quad
 v_p(U)>-3\eta/8,\quad v_p(V)>-7\eta/8.
\]
Therefore
\begin{equation}\label{eq:v12-scalar-valuations}
 v_p(q)>3\eta/8>0,\qquad
 v_p(z_0)>(32-3/8)\eta,\qquad
 v_p(z_1)>(32-7/8)\eta.
\end{equation}
Since $\eta$ is a positive integer, rational integral valuations give
$v_p(z_0),v_p(z_1)\ge32$.

At a split prime $v_p(q)$ is the sum of the two valuations of $f(H)$;
at an inert prime it is twice the single valuation.  Positivity of
$v_p(q)$ therefore supplies at least one place $v$ above $p$ with
$v(f(H))>0$.  We claim only existence of one such place.  At that
place $H$ is formal and $v(z(H))\ge1$.  By the two coefficient bounds,
\[
 v(z(R))\ge v(z(H))+32>0.
\]
This need not make $R$ itself formal: it may lie near $(0,0)$.
We handle both branches explicitly.

Put $m=v(x(R))$.  The retained point domains give $x(R)y(R)\ne0$.
Since $D$ is a unit and
\[
 z(R)^2=\frac{x(R)}{x(R)^2-D},
\]
positive valuation of $z(R)$ excludes $m=0$.  For $m>0$ its valuation
is $m/2$, and for $m<0$ it is $-m/2$.  The exact duplication identity
\begin{equation}\label{eq:v12-duplication}
 x([2]R)=\frac{(x(R)^2+D)^2}{4x(R)(x(R)^2-D)}
\end{equation}
then gives, for either sign of $m$,
\[
 v(x([2]R))=-2v(z(R))<0,\qquad v(z([2]R))=v(z(R)).
\]
Its numerator and denominator are nonzero: with $m\ne0$ their
relevant terms have unequal valuations, and $p$ is odd.
Also $[2]R$ is nonidentity because $s\ne0$ and $P$ has infinite order.

Both $H$ and $[2]R$ are now in the formal subgroup where the logarithm
is a homomorphism and isometry.  Apply it to the exact global identity
\[
 [h]([2]R)=[2s]H.
\]
The points are nontorsion, so their nonzero formal logarithms permit
comparison of valuations.  Since $2$ is a $p$-unit,
\[
 v_p(h)+v(z([2]R))=v_p(s)+v(z(H)),
 \qquad v_p(s/h)\ge32.
\]
The nonzero integer $h$ has nonnegative $p$-valuation, hence
$v_p(s)\ge32$.  This proves the claimed divisibility at every prime
of $M_5(\beta)$.  The same exponent-nineteen height theorem as in
Lemma~\ref{lem:v9-bridge} gives the asserted height bound.
\end{proof}

For completeness use the old scalar preparation with its stronger
six-leaf conditions.  For integral $X$, $r$ is a nonzero integer.
Choose $c$, a sufficiently divisible positive $h_0$, and then a still
deeper positive $s$ satisfying the anchor, exactly as in the proof of
Lemma~\ref{v5:scalar32-bridge}.  At primes dividing $r$, both
coefficients of $f(H)$ are divisible by $r$ locally, so $q/r$ is
integral; the two retained ratio leaves are integral.  At primes
dividing $c$ but not $r$, the same leaves are divisible by the prime.
The constant-left tree is consequently one modulo every prime on
the support of $cr$.  It is nonzero and a unit at those primes.
Thus $crW^2$ is positive, has square reduced denominator, and has the
same local squareclasses as $cr$ at two, three, and seventeen.
The original pole-completeness lemma supplies fresh scalar pole
witnesses.  The new norm and ratio functions are continuous at the
chosen points, so the same finite-neighborhood argument provides a
congruence class of $h_0$ preserving all required conditions.

\subsection{A reciprocal chart for the scalar difference}
\label{subsec:v12-reciprocal}

This replacement is used only for $R=S-P$.  The five comparison
subtractions retain their strengthened guards
\eqref{eq:v12-half-trace}.  Write the incoming scalar point as
$S=(x,hj)$ with $x=h^2$ and $j^2=x^2-D$.  Put
\[
 d=c_*x+D,\quad N=a_*j-b_*h,\quad
 M=-2a_*Dh(x+c_*)-b_*j(c_*x-D).
\]
Both real embeddings give $c_*>0$, $D>0$, and $x\ge0$; hence $d\ne0$.
Introduce $g,k\in K$, retain the existing outputs $f,z\in K$, and
introduce a rational guard $\sigma$.  The final equations are
\begin{equation}\label{eq:v12-reciprocal}
 dg=N,\qquad
 a_*dk+2Dh+a_*b_*j-2Db_*g=0,\qquad
 f=g^2,\qquad f_0\sigma=1,\qquad kz=-g.
\end{equation}
Each field equation means its two rational component equations.
Every component equation is quadratic.  Since
$f_0=g_0^2+2g_1^2$, the guard is precisely $g\ne0$ over rational
coordinates, and imposes no unmentioned sign choice.

We give the identities and all exceptional cases needed to identify
these outputs with the actual group difference.  If
\[
 G=dg-N,\quad K_0=d^2k-M,\quad
 L=a_*dk+2Dh+a_*b_*j-2Db_*g,
\]
then expansion using $a_*^2=c_*$ and $b_*^2=c_*^2-D$ gives
\begin{equation}\label{eq:v12-dcancel}
 dL=a_*K_0-2Db_*G.
\end{equation}
Thus the first two equations of \eqref{eq:v12-reciprocal} are
exactly equivalent to $g=N/d$, $k=M/d^2$, because $a_*d\ne0$.
The latter expressions satisfy
\begin{equation}\label{eq:v12-reciprocal-identities}
 k^2+Dg^4=1,\qquad
 (a_*j-b_*h)(a_*j+b_*h)=(x-c_*)(c_*x+D).
\end{equation}
For example the first is the polynomial identity
$M^2+DN^4=d^4$ modulo $h^2=x$, $j^2=x^2-D$ and the two fixed-constant
identities.  This also specifies an exact denominator-cleared check.

The guard gives $g\ne0$.  If $k=0$, the first identity in
\eqref{eq:v12-reciprocal-identities} gives $d_*g^2=1$ or $-1$.
Taking rational norms gives $7\operatorname{Norm}(g)^2=1$, impossible
for a rational norm.  Thus $k\ne0$, and all displayed outputs are
well defined.  If $x\ne c_*$, use the ordinary subtraction slope
$m=(hj+a_*b_*)/(x-c_*)$.  Substitution and clearing denominators give
\[
 m^2-x-c_*=1/g^2,\qquad
 m(x-1/g^2)-hj=k/g^3.
\]
These are the abscissa and ordinate of $S-P$.  Therefore the retained
outputs are exactly $f(R)=g^2$ and $z(R)=-g/k$.

If $x=c_*$, the incoming chart forces
$h=\pm a_*$ and $j=\pm b_*$.  Equal choices of signs give $S=P$
and $g=0$, excluded by the guard.  The other choices give $S=-P$.
For $h=a_*,j=-b_*$ the explicit values are
\[
 g=\frac{-2a_*b_*}{c_*^2+D},\qquad
 k=\frac{c_*^4-6c_*^2D+D^2}{(c_*^2+D)^2}.
\]
They give the abscissa and ordinate of $-2P$ by the ordinary duplication
formula.  Simultaneously negating $h,j$ negates $g,k$ and leaves
$1/g^2$, $k/g^3$, and $-g/k$ unchanged.  Thus this branch also represents
the true difference, not an arbitrary value of an old zero-over-zero
slope.

In the complete present system that extra branch is impossible.
It would give $R=-2P$ and $v_3(f_1(R))=2$ by
Lemma~\ref{lem:v12-anchor-coefficient}, whereas
\eqref{eq:v12-anchor} requires $v_3(f_1(R))\ge2^{25}$.
It follows that every final solution has $x\ne c_*$.  Restore the old
inverse of $x-c_*$, slope, slope square, and ordinate products by their
ordinary rational formulas.  All retained outputs coincide.  Conversely
every old solution extends to \eqref{eq:v12-reciprocal} by the displayed
formulas, with $\sigma=1/f_0$; its actual nonzero $f(R)=g^2$ ensures
$f_0>0$.  This is an equivalence on all nonprivate coordinates of the
complete system, with the anchor explicitly used to exclude the
additional branch.

For the coordinate count one may first use $d_2=d^2$ and $d_2k=M$.
That intermediate reciprocal chart replaces eight old scalar
difference coordinates by the six components of $g,d_2,k$ and one
rational guard, saving one.  Identity \eqref{eq:v12-dcancel} then removes
the two components of $d_2$.  Forward, divide $dL=0$ by the already
nonzero $d$; backward, reinsert the unique $d_2=d^2$ and use
$a_*K_0=0$.  This second step is a direct rational equivalence requiring
no new exceptional-point argument.  The combined saving is three.

\subsection{A shorter weight circuit with both outputs retained}
\label{subsec:v12-toric}

In each detector let $e$ be its weight $e$, $q=q_1$, and keep both
outputs $q_{12},q_{40}$.  The grouped evaluations already contain
$Q=q^3$ and $T=q_{12}^2$.  Replace the six old intermediate powers of
$e$ by five rational coordinates $\alpha,\gamma,u,v,\tau$ and impose
\begin{equation}\label{eq:v12-toric}
 e\alpha=1,\quad e\gamma=\alpha^2,\quad
 \gamma u=eq,\quad \gamma v=u^2,\quad
 Qq_{12}=v^2,\quad q\tau=eT,\quad vq_{40}=T\tau.
\end{equation}
These are seven quadratic equations for seven coordinates when the
two retained outputs are counted.  The old circuit had eight.

The first equation gives $e\ne0$.  The unchanged definition of $q$
is $q\delta^3=e$, with the powers of $\delta$ represented by their
retained coordinates, so $q\ne0$ as well.  The equations now determine,
in order,
\[
 \alpha=e^{-1},\quad\gamma=e^{-3},\quad
 u=qe^4,\quad v=q^2e^{11},\quad q_{12}=qe^{22}.
\]
The division by $Q=q^3$ is legitimate.  The retained equation for
$T$ then gives $T=q^2e^{44}$, and the last two equations give
\[
 \tau=qe^{45},\qquad q_{40}=qe^{78}.
\]
In particular $v\ne0$ whenever it is used as a divisor.  This is a
successive computation, not a circular specification involving $T$.

The forward map from the old circuit assigns these five Laurent
monomials and retains $q_{12},q_{40}$, $Q,T$, and every other coordinate.
For the reverse map restore the old six powers as
$e^{-1},e^{-3},e^{-7},e^{-14},e^{-29},e^{-55}$.  The old output equations
then give the same $qe^{22}$ and $qe^{78}$.  The replaced internal
powers have no outside consumers.  Thus these maps preserve every
outside coordinate, not merely the input.  Applying this equivalence
to all three detectors saves three coordinates.

\subsection{Removing the scalar norm coordinate}\label{subsec:v12-norm-elim}

One further elimination uses coordinates already present in the
order-forty test $f(H)$.  With $f(H)=h_0+h_1\rho$, its square is stored
through
\[
 t=h_0^2+2h_1^2,\qquad u=h_0h_1.
\]
Here $u$ is half the $\rho$-coefficient, a normalization that matters.
The scalar coordinate $q=h_0^2-2h_1^2$ from
\eqref{eq:v12-scalar-eight} has only one other consumer, the equation
$r(A-1)=q^2$.  Use the identity
\begin{equation}\label{eq:v12-norm-square}
 (h_0^2-2h_1^2)^2=(h_0^2+2h_1^2)^2-8(h_0h_1)^2=t^2-8u^2.
\end{equation}
Delete $q$ and its defining row and replace the consumer by
\[
 r(A-1)=t^2-8u^2.
\]
This is still quadratic.  Forward, forget $q$.  Backward, reinsert
$q=h_0^2-2h_1^2$.  Identity \eqref{eq:v12-norm-square} proves both
rows in the old system and makes these maps mutually inverse.
No division, nonzero hypothesis, or extra arithmetic input is used.
The proof of Lemma~\ref{lem:v12-scalar-sound} applies by this unique
reinsertion; it does not require retaining $q$ as a quantified variable.

\subsection{Soundness of the combined construction}
\label{subsec:v12-soundness}

Let a rational solution of the final system be given.  The inherited
odd charts identify the four root points with odd integer labels.
The reciprocal chart supplies the actual $R=S-P$, is nonidentity,
and, together with the anchor, excludes $S=-P$.  The explicit maps
above therefore restore all unchanged point domains and quotient
guards needed by the earlier arguments.  The removed scalar norm
coordinate also has its unique rational value.

Put $\kappa=|s|=|a-1|$.  By
\eqref{eq:v12-anchor-subgroup} and Lemma~\ref{lem:v12-scalar-sound},
\[
 \log_2\kappa>25165823.5,\qquad
 \operatorname h(\beta)\le\mathcal H\kappa^{19/32}.
\]
These are the exact bounds consumed in the Version 9 height argument.
For completeness we repeat the numerical deductions.  Since
\[
 10(7+2093966)<25165823.5,\qquad80<2^7,
\]
we have
\[
 \kappa>(80\mathcal H)^{10},\qquad
 |a|\ge\kappa-1\ge\kappa/2>(40\mathcal H)^{10},
\]
and
\begin{equation}\label{eq:v12-margins}
 \operatorname h(\beta)<\kappa/80\le|a|/40,\qquad
 \log(40\mathcal H)<\kappa/80\le|a|/40.
\end{equation}
The first inequality uses
$\kappa^{13/32}>80\mathcal H$; the logarithmic inequality is weaker.
In particular all labels are in the large range needed below.

Set
\[
 F=t^2-a(a-2)tb+(a-2)^2b^2-1,\qquad
 B=\max(|n|,|a|,|t|,|b|)>4.
\]
The triangle inequality gives $|F|<5B^4$.  If $F\ne0$,
Lemma~\ref{lem:v12-mixed-local} and the unchanged detector implication
make the three contact radicals divide $F$ to order forty.
Equation~\eqref{eq:v12-height-input} then gives
\[
 B^{1/10}<7\mathcal H\,5^{19/40}+6<24\mathcal H,
\]
contradicting $B\ge|a|>(40\mathcal H)^{10}$.  Hence $F=0$.
The point $Q_4$ is nonidentity, so its integer label $b\ne0$.
Lemma~\ref{lem:o8-lucas} therefore gives $\log B>|a|-3$.

If $n\ne\beta$, let $N$ be the nonzero reduced numerator of
$n-\beta$.  Then $|N|\le2e^{\operatorname h(\beta)}B$, and the same
three radicals divide $N$ to order twelve.  Using
$B^2-6>(5/8)B^2$ for $B>4$, equation~\eqref{eq:v12-height-input}
gives
\[
 B^{5/12}<40\mathcal H e^{(19/12)\operatorname h(\beta)}.
\]
Consequently
\[
 \log B<\frac{12}{5}\log(40\mathcal H)
            +\frac{19}{5}\operatorname h(\beta)
       <\frac{31}{200}|a|<|a|-3,
\]
a contradiction.  Thus $n=\beta\in\Z$.  If $X=u/v$ in lowest terms
with $v>0$, integrality of $2X^2+3$ forces $v^2\mid2$, so $v=1$.
Every rational solution therefore projects to an integer input.

\subsection{Completeness with all replacements at once}
\label{subsec:v12-completeness}

Fix an integer $X$ and put $n=\beta=2X^2+3$, a positive odd integer.
First perform the scalar preparation with the original six-leaf
bounds, which, as shown above, suffice for the new scalar tree.
Choose its finite correction and a sufficiently divisible positive
shared index $h_0$, with $64\mid h_0$.  Choose a positive scalar index
$s$ still more divisible, large enough for the physical anchor, and
with $6\mid s$.  Put $a=s+1$.  The actual prepared scalar point is
$R=[s]P=S-P$.  The new anchor has exactly the old allowed subgroup,
and the new scalar tree has the required unit root.  Fresh scalar
pole witnesses therefore exist.

To choose the two other root labels, define
\[
 U_0=0,\quad U_1=1,\quad U_{j+1}=aU_j-U_{j-1},\qquad
 k=2(a-2),\quad b=U_k/(a-2),\quad t=U_{k+1}.
\]
Induction gives $U_j\equiv j\pmod{a-2}$, so $b$ is an integer.
Modulo two the sequence has period $0,1,1$.  Since $a\equiv1\pmod3$,
$k\equiv1\pmod3$, and both $U_k,U_{k+1}$ are odd.  The positive odd
denominator $a-2$ makes $b$ odd.  The Lucas identity gives
$F(a,t,b)=0$.  Thus all four root indices are positive and odd, and
Proposition~\ref{v9:prop:odd-chart} supplies their chart witnesses.

Fix these root points and the three positive rational square scales
of Corollary~\ref{v9:cor:scales}.  The new scalar functions and the
mixed alternate numerator are rational functions defined at the
prepared points.  Their required finite local conditions persist in
a sufficiently small neighborhood.  As in the original argument,
choose a common continuity modulus $L$, with $64\mid L$, so that
$h\equiv h_0\pmod L$ preserves them.  All these choices precede the
contact-prime selection.  Increase the finite prime cutoff to include
the fixed root denominators needed for the five half-trace guards.

Choose the contact primes using the same simultaneous-prime theorem
and Lemma~\ref{v9:lem:separation}.  Put
\[
 M_\pm=\operatorname{lcm}_p(d_\pm(p)/4),\qquad\Pi=\prod_p p.
\]
The retained direct-target construction supplies
\[
 \gcd(M_-,M_+)=1,\quad\gcd(M_-,L)=16,\quad\gcd(M_+,L)=1,
 \quad\gcd(\Pi,M_+)=\gcd(\Pi,L)=1.
\]
Thus the same three congruences are compatible:
\[
 h\equiv0\pmod{\Pi^{19}M_-},\qquad
 h\equiv1\pmod{M_+},\qquad h\equiv h_0\pmod L.
\]
Use the inherited real-progression argument to choose a sufficiently
large positive $h$ in this progression, avoiding the finitely many
forbidden point and quotient loci.  Set
$U_i=[hn_i]P$ and $V_i=[(h-1)n_i]P$.

At every primary place the points $H,U_i$ have depth at least twenty:
the doubling-image order first makes them formal, and multiplication
by $p^{19}$ raises depth by nineteen under the same formal logarithm
isometry.  At opposite places the points $V_i$ are formal.  Hence
$L_i$ agrees with its integer label to order at least eighty; the
alternate ratio agrees to order twelve for $i=1$ and forty for the
other three labels.  All tests \eqref{eq:v12-tests} hold.  The same
primary/opposite reductions prove that all five half-traces are
nonzero, by the argument following \eqref{eq:v12-half-trace}.
The chosen scalar label has $S\ne\pm P$, so its reciprocal chart
has the ordinary rational lift described above.

The pole-realization estimates are valid for these fixed orders
$1,12,40$.  In the source notation their large-real-parameter decay
is still $\lambda^{-3}$; the small-prime lower bound is
$3\nu+(10m-8)\eta-2c$, which tends to infinity under the same
preparatory divisibility.  No new prime-pattern assertion is used.
Obtain the detector pole witnesses and apply the rational
normalization \eqref{eq:v12-pole-normalized}.  Apply the weight-circuit
maps \eqref{eq:v12-toric}, the reciprocal denominator cancellation,
and the norm-square elimination.  These exact maps preserve all
outside coordinates.  They therefore give one simultaneous rational
solution of the final system for the chosen integer $X$.

This proves completeness in the required order of choices.  In
particular, it does not combine separately chosen primes or
incompatible continuity moduli from different optimizations.

\subsection{Exact polynomial, counts, and coefficient lift}
\label{subsec:v12-finite}

Each equation above is represented by its rational component rows;
intermediate products in the retained straight-line program are
represented by their specified coordinates.  Normalize each nonzero
rational residual by clearing denominators and dividing its integer
content, and let $P_{498}$ be the sum of the squares of those 467
normalized residuals.  This is an integer polynomial.  A sum of
rational squares vanishes exactly when every residual vanishes.
The soundness and completeness just proved therefore establish
Theorem~\ref{thm:v12-quartic}.

The following count records the actual transported, fully reduced
systems; it does not count unreduced formulas in isolation.
\begin{center}\small
\begin{tabular}{lrr}
\toprule
System or replacement & Total variables & Quadratic residuals\\
\midrule
The retained Version 9 baseline &523&492\\
Mixed input precision and its checked compilation &519&488\\
Three normalized detector poles &516&485\\
Five half-trace guards &511&480\\
Single-coefficient translated anchor &510&479\\
Shorter scalar norm tree &505&474\\
Scalar reciprocal chart with $d^2$ retained &504&473\\
Three weight circuits retaining $q_{12},q_{40}$ &501&470\\
Cancellation of the scalar $d^2$ coordinates &499&468\\
Norm-square coordinate elimination &498&467\\
\bottomrule
\end{tabular}
\end{center}

The mixed-depth compilation uses only exact quadratic identities,
constant-pivot substitutions, and invertible shears in shared
coefficient coordinates.  A constant-pivot substitution inserts its
unique polynomial value in the reverse direction; a shear inserts
its inverse linear transformation.  Both maps fix the input.
The final norm-square elimination is the explicit identity
\eqref{eq:v12-norm-square}.  Thus no compiler reduction introduces a
variable denominator or alters an arithmetic guard.  The exact
certificates also check that private coordinates deleted in the
other replacements have no undeclared consumers.  In particular,
the half-trace witnesses occur only in their guard equations; the
removed weight powers occur only in their old weight circuit; the
removed scalar $d^2$ coordinates occur only in their definitions and
$k$ equations; and the removed norm occurs only in its definition and
constant-node equation.

The frozen residual file is \texttt{candidate498-ir.json}, with
SHA-256
\begin{center}\small
\nolinkurl{7682a1229149c75e27322e86dec424449b55321b06c48e468f90fa78bcebb61a}.
\end{center}
The complete sparse coefficient table \texttt{polynomial498.json}
has SHA-256
\begin{center}\small
\nolinkurl{817fec936dee3b358fc8b5b0b9a5570694adf009b134432467ef0be914aeedc4}.
\end{center}
Its exact expansion has 24,669 monomials and degree four; all 498
variables occur.  A coefficient of absolute value one proves content
one.  The only quadratic-in-$C$ coefficient is $16C^2$, at the square
of the translated anchor witness $w$, with no lower-$C$ term at that
monomial.  All other coefficients are independent of $C$ or linear
in $C$ and are smaller in absolute value.  The coefficient comparison
is carried out symbolically; no expansion of $3^{2^{25}}$ is needed.
It gives the exact maximum $16C^2=16\cdot3^{2^{26}}$.

There is also an ordinary-coefficient lift with sixteen additional
coordinates:
\[
 c_0=3^{1024},\qquad c_{j+1}=c_j^2\quad(0\le j<15),\qquad C=c_{15}.
\]
The recurrence has exactly one rational solution and ends at
$3^{1024\cdot2^{15}}=3^{2^{25}}$.  Replace $-4Cw$ in
\eqref{eq:v12-anchor} by $-4c_{15}w$.  All residuals remain quadratic.
This gives a specified quartic in 514 variables with the same input
projection and 513 existential witnesses.  It has 483 quadratic
residuals and 24,716 monomials, with content one and maximum coefficient
$3^{2048}+1005072314676622<2^{3247}$.  Its coefficient table,
\texttt{polynomial514.json}, has SHA-256
\begin{center}\small
\nolinkurl{d247a27709afda5f2be626390cffa1a175c3dfcbee1a5c832a1c457b32629bd5}.
\end{center}
The lift is unique, so it requires no additional arithmetic argument.

The exact checks establish the residual identities, the witness
counts, and every coefficient of the displayed sums of squares.
They are separate from the arithmetic arguments above.  The
inherited source qualifications of Section~\ref{subsec:v9-inputs}
remain in force; a finite polynomial check is not a new proof of the
restricted converse, the height theorem, or the simultaneous-prime
input.

\section{A quartic presentation with 394 variables}
\label{sec:v13}

The transformations in this section give the following specified
polynomials. The arithmetic inputs and their scope are those retained
in Section~\ref{subsec:v12-inputs}.
\begin{theorem}\label{thm:v13-quartic}
There is a specified integer polynomial $P_{394}$ of total degree four
in 394 variables such that
\[
 X\in\mathbb Z\quad\Longleftrightarrow\quad
 \exists\mathbf v\in\mathbb Q^{393}:P_{394}(X,\mathbf v)=0.
\]
It has 6,338 nonzero monomials, content one, and largest absolute
coefficient $16\cdot3^{2^{26}}$. Its symbolic coefficient table uses
$C=3^{2^{25}}$ as a fixed integer, not an additional variable.
There is also a specified integer quartic $P_{411}$ with
\[
 X\in\mathbb Z\quad\Longleftrightarrow\quad
 \exists\mathbf w\in\mathbb Q^{410}:P_{411}(X,\mathbf w)=0.
\]
It has 6,388 nonzero monomials, content one, and all coefficients have
absolute value less than $2^{3247}$. Its largest coefficient has exactly
3,247 bits; its exact value is included in the coefficient artifact.
\end{theorem}

Some transformations below are inverse rational changes of private
witnesses, fixing every outside coordinate. Others change the arithmetic
certificate or enlarge a point domain; for those we prove soundness and
existence of witnesses for every integer input separately.

\subsection{Conventions and retained arithmetic inputs}
Put $K=\mathbb Q(\rho)$, $\rho^2=2$, $d_0=3-\rho$, $D=d_0^2$, and
$E:y^2=x^3-Dx$. Retain $2E(K)=\mathbb ZP$, with
\[
 P=(c_*,a_*b_*),\quad a_*=-1+7\rho/4,\quad
 b_*=39/8-9\rho/2,\quad c_*=a_*^2.
\]
The roots $Q_i=[n_i]P$ have nonzero odd indices, with
$(n_1,n_2,n_3,n_4)=(n,a,t,b)$. In the intermediate source the auxiliary point is
$H=[h]P\ne O$, and $U_i=[u_i]P$, $V_0=H-P$, $V_i=U_i-Q_i$.
We will allow $U_i$ and $V_i$ to be $O$ for $i>0$. The proof below
still excludes $V_0=O$. The scalar point is $R=Q_2-P=[s]P$, where
$s=a-1$ is nonzero and even.

The retained root-height implication is
\begin{equation}\label{eq:v13-height}
 B^2-6<7\mathcal H|N|^{19/k},\qquad
 B=\max_i|n_i|,\quad 1\le\mathcal H<2^{2093966},
\end{equation}
whenever the $k$th powers of the three contact radicals divide the
nonzero integer $N$. Their scales are $1,C_0^2,(1+C_1^2)^2$, where
$\prod_i x(Q_i)=C_0+C_1\rho$. The retained scalar anchor gives
\begin{equation}\label{eq:v13-anchor}
 \log_2|s|>25165823.5.
\end{equation}
These are precisely the root-height and anchor statements proved earlier;
neither is inferred from a finite coefficient calculation.

\subsection{Rational logarithms and the two order-ten tests}
Write $f=1/x$, $z=-x/y$, and extend these regular functions by
$f(O)=z(O)=0$. Define
\begin{align*}
 Q(w)&=1-\frac{1526}{1073}w+\frac{513}{1073}w^2,\\
 N(w)&=1-\frac{9776}{5365}w+\frac{2941}{3219}w^2
 -\frac{1024}{13949}w^3-\frac{1920}{237133}w^4
 -\frac{1024}{711399}w^5-\frac{1408}{5928325}w^6
 +\frac{2432}{103152855}w^8,\\
 \mathcal R(w)&=N(w)/Q(w).
\end{align*}
The exact identity
\[
 QG_9-N=\frac{2970726}{254443709}w^{10}
       -\frac{41697666}{6361092725}w^{11}
\]
proves that this replacement has the same order-forty local accuracy.
Indeed, on $[0,1]$ the denominator is at least $60/1073$ and the
Bernstein coefficients of the numerator are
\[
 1,\ \frac{4143}{5365},\ \frac{260069}{450660},\
 \frac{27829}{67340},\ \frac{4636683}{16599310},\
 \frac{82391}{474266},\ \frac{2245591}{23713300},\
 \frac{955437}{23713300},\ \frac{83904}{10113025}.
\]
They are all positive and give $N\ge83904/10113025$. Every nonidentity point of $2E(K)$ has
$0<Df^2<1$ at both real embeddings; at $O$ its value is zero.
Thus it is defined and nonzero on these doubled-point domains. For
the later enlarged H domain and its difference H-P, we instead use
the global numerator and denominator nonvanishing certificates below.

At an unramified place above $p>76176$, a nonidentity formal point of
depth $r=v(z)\ge1$ has $v(Df^2)=4r$. All coefficient denominators
above are supported at primes at most37. Consequently
$v(\mathcal R-G_9)\ge40r$. The earlier full-tail estimate gives
$v(\log_E(Q)-z\mathcal R(Df^2))\ge41r$.
The point $O$ has exact zero error. We retain the $G_2$ numerator at
$V_1$, of relative precision12, and use $\mathcal R$ at $V_0,V_2,V_3,V_4$.

Set $\ell_0=z(V_0)\mathcal R(Df(V_0)^2)$ and define
\[
 L_i=z(U_i)/z(H),\quad
 L'_1=\overline{z(V_1)G_2(Df(V_1)^2)/\ell_0},\quad
 L'_i=\overline{z(V_i)\mathcal R(Df(V_i)^2)/\ell_0}\quad(i>1).
\]
The nonzero denominator is proved below, before these quotients are used.
The tests have orders1 on $f(U_i)$ and $\overline{f(V_i)}$, order40
on $f(H)$, order10 on $L_1-L'_1,L_1-\beta$, and order40 on the other
three cross-tests and
\[
 F(L_2,L_3,L_4),\qquad F(A,T,B)=T^2-A(A-2)TB+(A-2)^2B^2-1.
\]
Conjugation and the place conventions are unchanged.

For completeness, spell out the zero-safe local argument. Let $A(Q)$ be
the analytic unit defined by $\log_E(Q)=z(Q)A(Q)$, with $A(O)=1$.
Define $T_2=G_2/A$ and $T_{40}=\mathcal R/A$ by their power series.
These units equal1 at $O$. Their errors at nonidentity formal points
have the stated orders. The primary ratio is exactly
$(u_i/h)A(H)/A(U_i)$, also when $u_i=0$. The alternate ratio is
$(u_i-n_i)/(h-1)$ times the corresponding conjugated unit ratio.
At least one of $h,h-1$ is a $p$-unit. Therefore one full ratio is
integral, and the positive-order cross-test transfers integrality to
both approximate and full ratios. Only now infer depth at least20 for
a nonzero $U_i$; if $U_i=O$, its error is zero. The primary error has
order80. Multiplication by $h,h-1$ and subtraction give
\[
 L_1\equiv n\pmod{p^{10}},\qquad L_i\equiv n_i\pmod{p^{40}} (i>1).
\]
No quotient $\log_E(O)/z(O)$ is used.

\subsection{Odd roots, shared reciprocals and identity-inclusive charts}
The retained finite chart for a root is
$x=h_Q^2$, $j_Q^2=x^2-D$, $Q=(x,h_Qj_Q)$.
It has $h_Qj_Q\ne0$ by the real embeddings and $\operatorname{Norm}(d_0)=7$.
Replace its old odd-index witnesses by $A\in K,u\in\mathbb Q$ with
\[
 A^2-2(x+j_Q)A-D=0,\qquad \operatorname{Norm}(A)=-u^2.
\]
Here is the reverse implication, including the squareclass choice.
Put $q=x+j_Q$ and $t=(A-q)/h_Q$. The equation gives
$t^2=2q$. The point $(A,At)$ lies on $E$ and its tangent slope is
$h_Q+t$, so its double is exactly $(x,h_Qj_Q)$.
The inherited squareclass decomposition writes its abscissa class as
$[\epsilon]^m[-1]^e[d_0]^f$, where $\epsilon=1+\rho$ and $e,f\in\{0,1\}$.
The norm condition forces $m$ odd and $f=0$. Thus $[A]=[\epsilon]$
or $[-\epsilon]$. In the latter case replace $A,t$ by $2q-A,-t$;
the product of the two abscissae is $-D$, so the new class is
$[\epsilon]$. Choosing the corresponding square root recovers the old
odd witnesses without changing the root. Conversely the old witnesses
$t,r$ give $A=th_Q+q=\epsilon r^2$, $u=\operatorname{Norm}(r)$.
This is equality of exterior projections, not a unique private inverse.

For the primary ratios keep one shared field value
$S=j_H/h_H=-1/z(H)$. The finite-point equations become
\[
 h_H S=j_H,\qquad j_iL_i=h_iS.
\]
These remove the four private primary parameters $z(U_i)$ without
changing their ratios. The scalar equation can be fused in the same way.
If its old equations are $Tz_R+a_*N=0$ and
$h_H\varrho+j_Hz_R=0$, replace them by
\[
 T\varrho-a_*NS=0.
\]
The inverse is $z_R=-\varrho/S$; $S\ne0$ follows from the retained
finite H chart.

For $i>0$ an identity-inclusive primary chart instead uses
\[
 f_i=r_i^2,\qquad J_i^2+Df_i^2=1,\qquad J_iL_i=r_iS.
\]
When $r_i\ne0$, set $h_i=1/r_i,j_i=J_i/f_i$ to recover the old
finite chart. When $r_i=0$, interpret $U_i=O$. Both choices $J_i=\pm1$
give that same point. The factor $J_i$ cannot vanish: it would give
$D r_i^4=1$, and taking real signs and then norms gives
$7\operatorname{Norm}(r_i)^2=1$, impossible over $\mathbb Q$.
Thus $z(U_i)=-r_i/J_i$ and the primary ratio is correct also at $O$.

Write the corresponding fixed root as $Q_i=(c,ab)$ with
$c=a^2$, $b^2=c^2-D$. Define
\[
 d=c+Df_i,\quad dg=aJ_i-br_i,\quad u=cr_i+bg,
 \quad dk=bJ_i-2au,\quad f(V_i)=g^2,\quad kz(V_i)=-g.
\]
The factor $d$ is positive at both real embeddings. Reduction by the
primary conic gives $k^2+Dg^4=1$, so $k\ne0$ by the same norm argument.
On the ordinary subtraction domain these expressions give
$V_i=(1/g^2,k/g^3)=U_i-Q_i$. They have the following exact boundary values.
If $r_i=0$, then $g=J_i/a$, $k=-bJ_i/c$, and $V_i=-Q_i$.
If $g=0$, substitution gives
$(cr_i^2-1)(c+Dr_i^2)=0$, hence $U_i=Q_i$ and $V_i=O$.
The remaining equal-abscissa case is the regular doubling limit for
$U_i=-Q_i$. Equality with the group-law morphism follows from its equality
on the dense ordinary domain and regularity of these formulas.

For H and the denominator difference retain the finite chart, or its
following exact rescaling. With $x=h_H^2$, put
\[
 d=c_*x+D,\quad N=a_*j_H-b_*h_H,\quad dg=N,
 \quad T=-2Dh_H-a_*b_*j_H+2Db_*g.
\]
The scalar reciprocal elimination is the identity
$L=a_*dZ-zK-a_*G$, where
$G=dg-N$, $K=a_*dk-T$, $Z=kz+g$, $L=Tz+a_*N$.
Since $a_*d\ne0$, reinserting $k=T/(a_*d)$ recovers the old equations.

For V0 one may also eliminate $j_H$ by writing $j_H=h_HS$ and $g=h_Hv$.
The resulting exact equations are
\[
 S^2=x-Df(H),\quad (c_*x+D)v=a_*S-b_*,\quad
 f(V_0)f(H)=v^2,
\]
\[
 (-2D-a_*b_*S+2Db_*v)\ell_0
       +a_*(a_*S-b_*)\mathcal R(Df(V_0)^2)=0.
\]
Both maps divide only by the already nonzero $h_H$ and $x$.
These exact maps precede the separately justified enlargement of H below.

\paragraph{The denominator cannot be the identity.}
The old argument using $V_1\ne O$ is no longer available. Instead suppose
$V_0=O$. The reciprocal chart gives $H=P$, and
\[
 f(H)=\frac{8(57+28\rho)}{1681}
\]
is a unit at every consumed place over $p>76176$. The one-place detector
lemma would force its order-forty test at any contact prime, a contradiction.
Thus all three contact radicals equal1. Equation~\eqref{eq:v13-height}
with $N=1$ gives $B^2-6<7\mathcal H$, contrary to
$B\ge|a|>(40\mathcal H)^{10}$ from the unchanged anchor.
This argument precedes every interpretation of an LP quotient.
Now $V_0\ne O$, its parameter is nonzero, and the established positivity
of $\mathcal R$ gives $\ell_0\ne0$. The local ratio proof above therefore
applies to all new solutions. It proves soundness directly, without
asserting that identity-branch witnesses lift to the old affine source.

\subsection{Scaled biquadratic comparison charts}
For one root write $Q=(c,ab)$, $c=a^2$ and $b^2=c^2-D$.
The reciprocal comparison chart above uses $r,J,f,g,k,u,F$, with
$f=r^2$, $F=g^2$. Its retained squared outputs are
$Y_U=f^2$ and $Y_V=F^2$, as elements of $K$, not separate component
squares. Replace the seven field values by five:
\[
 V=ar,\quad Q'=ag,\quad I=aJ,\quad K'=ak,\quad W=rg.
\]
The new equations are
\begin{align*}
 VQ'&=cW,& V^2+(Q')^2+DW^2+2bW&=1,\\
 aI&=bV+cQ'+DVW,& aK'&=-bQ'-cV-DQ'W,\\
 IL&=VS,& K'z&=-Q',\\
 DcY_U+I^2&=c,& DcY_V+(K')^2&=c.
\end{align*}
Every displayed equation is quadratic in retained coordinates. In
particular c is an existing coordinate, so cY is a quadratic product.

Here is an algebraic proof including all zero branches. Put
$d=c+Dr^2$, $e=c+Dg^2$, and define
\[
 T=c(r^2+g^2)+Dr^2g^2+2brg-1,
\]
\[
 G=dg-aJ+br,\quad E'=ak+bg+er,\quad
 H'=dk-bJ+2a(cr+bg).
\]
Modulo only $c=a^2$ and $b^2=c^2-D$, direct expansion gives
\begin{align*}
 a^2(J^2+Dr^4-1)&=dT-G(dg+br+aJ),\\
 a^2(k^2+Dg^4-1)&=eT+E'(ak-bg-er),\\
 dE'&=aH'+DrT-bG.
\end{align*}
The factors a,c,D are nonzero, and d,e are positive at both real
embeddings. An old solution has G=H'=0, hence T=E'=0, proving
all new equations and the second reciprocal conic. Conversely set
$r=V/a$, $g=Q'/a$, $J=I/a$, $k=K'/a$; the product equation
forces W=rg. The new equations give T=G=E'=0, so the three
identities recover both conics and $H'=0$. Set $f=r^2$, $F=g^2$ and
$u=cr+bg$. Comparing each new Y equation with its recovered conic
gives $Y_U=r^4$ and $Y_V=g^4$. Dividing the ratio equations by a
restores the old ratios. Thus these are inverse rational maps,
fixing every outside coordinate.

There is no division by r,g,J,k,W,L,z or either test value. When
r=0 the first point is O and its ratio is zero; when g=0 the
difference is O and its parameter is zero. The biquadratic equation
prevents both from vanishing. These are exactly the already accepted
identity branches. This transformation replaces the intermediate inverse-addition change;
it is not composed with a second copy of that elimination.

\subsection{Enlarging the auxiliary H point domain}
After the exact H.j rescaling, the two coordinates H.h occur only
in the two component equations x=H.h$^2$. Delete them and keep
\[
                  xf=1,\qquad S^2=x-Df.
\]
They define the finite point H=(x,xS) of E(K), with x nonzero,
but no longer require H to be a double. The rank and torsion
calculation in the generator proof gives
\[
 H=[m]R_0+T,\quad P=2R_0,\quad m\in\mathbb Z,\quad T\in E[2](K).
\]
Write h=m/2 only for its rational free index. It is not notation
for multiplication by a half-integer on the elliptic curve.

We first exclude all finite points in $E[2]$, $P+E[2]$ and
$-P+E[2]$ before interpreting ratios. The points O and (0,0) are
excluded by the finite chart. Put $d_0=3-\rho$, $c_*=x(P)$.
The exact norm and translation identities are
\[
 \operatorname{Norm}(c_*)=41^2/64,\quad
 \operatorname{Norm}(c_*-d_0)=17^2/64,\quad
 \operatorname{Norm}(c_*+d_0)=63^2/64,
\]
\[
 x(P+(0,0))=-D/c_*,\quad
 x(P+(d_0,0))=d_0(c_*+d_0)/(c_*-d_0),
\]
\[
 x(P+(-d_0,0))=-d_0(c_*-d_0)/(c_*+d_0).
\]
Each of these abscissae and its inverse is integral outside
$\{2,3,7,17,41\}$. The same is true for the nonzero torsion
abscissae, and negating P changes none of these abscissae. Thus
f(H) is a unit at every consumed place. The retained order40 f(H)
test would then rule out every contact prime. All three radicals
would equal1, contradicting the root-height bound and scalar anchor
exactly as in the preceding V0 exclusion. This proof needs no
ratio involving H. Consequently m is neither0 nor2, H and H-P are
nontorsion, S is nonzero, and H is neither P nor -P.

In the rescaled difference equations put
\[
 d=c_*x+D,\quad N=a_*S-b_*,\quad dv=N,\quad
 T'=-2D-a_*b_*S+2Db_*v.
\]
The former positivity proof for d no longer applies. If d=0, the
curve equation gives $x=-D/c_*$, $S=\pm b_*/a_*$, so H is one of
the excluded points $\pm P+(0,0)$. If N=0, it gives
$(x-c_*)(x+D/c_*)=0$, also excluded. Thus d,N,v are nonzero.
Direct subtraction now gives
\[
 x(H-P)=1/(xv^2),\qquad y(H-P)=T'/(a_*dxv^3),
 \qquad z(H-P)=-a_*dv/T'.
\]
There is no unresolved equal-abscissa case, since $x=c_*$ means
H=P or -P. Since H-P is nontorsion, its ordinate and T' are
nonzero. In particular the retained equation $f(V_0)f(H)=v^2$ has
its intended interpretation, even if either real abscissa is negative.

The old positivity interval for the rational logarithm is no longer
enough. A stronger certificate is available: N(w) above has no root
modulo71, and Q(w) has no root modulo7. At the respective prime all
coefficients are integral and the leading coefficient is a unit.
These assertions are finite substitutions through all residues. Both
primes split in K, with square roots of2 equal to12,59 modulo71
and3,4 modulo7. A K-root would give a local root, necessarily integral,
and then a forbidden residue root. Hence N and Q never vanish in K.
The actual equation $T'\ell_0+a_*N\mathcal R(Df(V_0)^2)=0$
therefore implies $\ell_0\ne0$ and gives the correct denominator
logarithm. No sign assertion for this logarithm is needed.

The formal ratio argument changes by using integral group identities:
\[
 [m]U_i=[2u_i]H,\qquad
 [m-2]V_i=[2(u_i-n_i)]V_0.
\]
They kill T. Applying the formal logarithm only to formal points
gives ratios $u_i/h$ and $(u_i-n_i)/(h-1)$. At every consumed odd
prime h and h-1 are integral and differ by1. The earlier transfer
and subtraction proof therefore applies, including zero numerators.
At opposite split places use the same conjugation as before; at an
inert place it preserves the valuation.

The scalar transfer uses the same integral-index convention. For $R=[s]P$ use
\[
                    [m](2R)=[4s]H.
\]
If the scalar coefficient test has depth $d_s$, formal logarithm
isometry gives $v_p(s/(m/2))\ge d_s$. Since m/2 is p-integral,
$v_p(s)\ge d_s$. This includes the selected depth24. The proof
that 2R is formal when R is near O or (0,0) is unchanged. Thus
both the comparison and scalar height arguments survive the larger
domain. Completeness uses the original H in $2E(K)$ and forgets its
two square witnesses. It need not realize any newly admitted half-index
or torsion translate.

\subsection{A positive scalar input and the depth-twenty-four certificate}
Use the selected label and contact polynomial
\[
 \tau=2X+1,\qquad \beta=51\tau^2+2504,
 \qquad r=\beta(\beta-1)(\beta-2401)(32\beta+343)(2\beta-1).
\]
For every rational $X$, $\beta\ge2504$, hence $r>0$.
For integral $X$, $\tau$ is odd and $\beta$ is a positive odd integer.
Conversely, the final conclusion must be that $\beta$ is an \emph{odd}
integer. If $\tau=u/v$ in lowest terms, integrality gives $v^2\mid51$;
as $51$ is squarefree, $v=1$. Oddness of $\beta$ then makes $\tau$ odd
and $X=(\tau-1)/2$ integral. Mere integrality would not suffice:
$X=-1/2$ gives the even integer $\beta=2504$.
The unchanged odd chart for $Q_1=[n]P$ supplies the required oddness
once $n=\beta$ has been proved.

At a consumed prime $p>76176$, $51$ is a unit. A negative valuation
of $\beta$ is $2v_p(\tau)$, hence even. The five contact factors then
have the same valuation, so an odd valuation of $r$ must be positive.
This argument makes no assertion at the exceptional primes3 and17.

For integral $X$ the following simultaneous local calculation holds.
Since $\tau$ is odd, $\beta\equiv3\pmod8$; the two even contact
factors each have valuation one and equal odd parts modulo8. The
other factors have residues3,7,5. Therefore $v_2(r)=2$ and
$r/4\equiv1\pmod8$. At3, $\beta\equiv2$ and the first four factors
have product1 modulo3. Furthermore
\[
 (2\beta-1)/3=34\tau^2+1669\equiv1+\tau^2\pmod3,
\]
which is nonzero. Thus $v_3(r)=1$. At17, the contact residues are
5,4,1,10,9 with product15, a nonzero square. Thus $v_{17}(r)=0$.

An intermediate one-choice construction used $c=33033$ when
$3\mid\tau$ and $c=33$ otherwise. Both are square units at2 and17,
and their unit parts at3 are respectively1 and2, correcting the two
possible classes of $r$. The selected construction removes even
this choice and sets
\[
                 c=r+4896,\qquad4896=2^5\cdot3^2\cdot17.
\]
At2 the quotient $c/r=1+4896/r$ belongs to $1+8\mathbb Z_2$;
at3 and17 it belongs to $1+p\mathbb Z_p$. Hence $cr$ is a square
in all three local fields. Also $c,r$ are positive integers.
At a consumed scalar contact prime, positivity of the odd valuation
$v_p(r)$ was proved before using $c$; since4896 is a unit there,
$c$ is a unit there too. We do not assert that $c$ is a unit at
\emph{every} large prime. Its additional finite prime support is
included in the preparation before the continuity modulus is chosen.

The physical input square $Y=X^2$ is retained. In each other consumer,
the old occurrence of $Y$ is replaced by $102Y+102X+76$, since
$2(102Y+102X+76)+2403=51(2X+1)^2+2504$. This substitution must
include every scalar test and height-input consumer, not the defining
row $Y=X^2$ itself.

For the deterministic correction the two old physical rows, in their
fixed normalization, are
\[
 cs_0-20619812F=0,\qquad
 (7895415c+20619812F)V-20619812b_0=0.
\]
The retained contact expression has $r=-4(s_0+7895415)$.
Substitution and multiplication by $-4$ give
$rcV+82479248b_0=0$. With the already retained coordinate $r_2=r^2$,
the selected row is
\[
                    (r_2+4896r)V+82479248b_0=0.
\]
This is quadratic and preserves the physical sign of the scalar input.
The free correction coordinate and its product coordinate have no
remaining consumers. In the reverse arithmetic interpretation the
correction is defined to be $r+4896$; the removed finite-choice equation
is not reimposed.

Write $q=\operatorname{Norm}_{K/\mathbb Q}(f(H))$ and
$z(R)/z(H)=z_0+z_1\rho$. The selected scalar root is
\begin{equation}\label{eq:v13-scalar24}
 B_s=\frac{z_0^2+r z_1^2}{r^{32}},\qquad
 W=\frac{B_s^2+r q^2}{r^{32}},\qquad \mathfrak b=crW^2.
\end{equation}
Here $q\ne0$ and $r>0$, so $W>0$. Forward powers compute $r^{32}$;
there is no need to retain an inverse-power tower or $r^{64}$.
The square $q^2$ is computed without a separate norm coordinate:
if $f(H)=f_0+f_1\rho$, use
\[
 q^2=(f_0^2+2f_1^2)^2-8(f_0f_1)^2.
\]
The factor $8$ is essential: the second stored quantity is a product,
not the full $\rho$-coefficient.

Here is the changed local application. At a consumed prime let
$\eta=v_p(r)>0$ be odd. The correction is a unit there.
The pole obstruction gives $\eta+2v_p(W)>0$.
The two terms in each numerator of~\eqref{eq:v13-scalar24} have
opposite valuation parity, so they cannot cancel; assign valuation
$+\infty$ to a zero term. Successively this gives
\[
 v_p(B_s)>63\eta/4,\quad v_p(q)>61\eta/4,\quad
 v_p(z_0)>191\eta/8,\quad v_p(z_1)>187\eta/8.
\]
Thus both rational coefficient valuations are at least24.
Positive valuation of the norm supplies a place with $v(f(H))>0$;
there $H$ is formal and $v(z(R)/z(H))\ge24$.
Small $z(R)$ alone does not imply that $R$ is formal. If $e=v(x(R))$,
then $e\ne0$, and the duplication identity gives, for either sign of $e$,
\[
 v(x(2R))=-2v(z(R)),\qquad v(z(2R))=v(z(R)).
\]
For the enlarged H domain apply the formal logarithm to
$[m](2R)=[4s]H$, where $h=m/2$, to obtain
$v_p(s/(m/2))\ge24$. Since $m/2$ is p-integral at this odd prime,
$v_p(s)\ge24$. This uses integer group multiples only.
The same height theorem therefore gives
\[
 \operatorname h(\beta)\le\mathcal H|s|^{19/24}.
\]
The exact inequality
$(5/24)\,25165823.5>2093966+7$ and~\eqref{eq:v13-anchor} imply
$\operatorname h(\beta)<|s|/80\le|a|/40$.
The retained anchor also gives $\log(48\mathcal H)<|a|/40$.

For completeness choose $c$ first. The existing preparation permits
arbitrary finite coefficient-depth conditions over the support of $cr$.
At a prime dividing $r$, impose depth $8v_p(r)$ on the coefficients of
$f(H)$ and depth $24v_p(r)$ on $z_0,z_1$. Then $W$ is integral there.
At primes dividing $c$ but not $r$, sufficiently positive depths give
integrality as well. The actual pole input $crW^2$ is positive, has
square reduced denominator and the prescribed local squareclasses.
These are the retained pole-completeness hypotheses; $W$ need not be
a unit. Fresh pole witnesses are supplied for this new input.

The scalar coefficient guard may be imposed on $f(R)_1$ instead of
$f(R)_0$. In the reciprocal chart $f(R)=g^2$ one has
$f(R)_0=g_0^2+2g_1^2$ and $f(R)_1=2g_0g_1$. A nonzero second
coefficient therefore implies a positive first coefficient and recovers
the old guard. Completeness uses the established exact formula
$v_3(f([s]P)_1)=2+2v_3(s)$ for nonzero even $s$, so this second
coefficient is nonzero on every prepared scalar orbit point.
This is an arithmetic replacement of the guard, not an equivalence
for arbitrary pairs $g_0,g_1$.

\subsection{The two final height contradictions}
The local order-forty tests force the contact radicals to divide
$F=t^2-a(a-2)tb+(a-2)^2b^2-1$ to order40.
If $F\ne0$, then $|F|<5B^4$ and~\eqref{eq:v13-height} contradicts
the retained lower bound $B\ge|a|>(40\mathcal H)^{10}$.
Thus $F=0$. The nonzero odd root label $b$ and the retained Lucas lemma
give $\log B>|a|-3$.
If $n\ne\beta$, its nonzero reduced numerator $N$ satisfies
$|N|\le2e^{\operatorname h(\beta)}B$ and is divisible by the three
contact radicals to order10. Since $B^2-6>(5/8)B^2$ for $B>4$,
\[
 B^{1/10}<48\mathcal H e^{(19/10)\operatorname h(\beta)},\qquad
 \log B<10\log(48\mathcal H)+19\operatorname h(\beta)
          <29|a|/40<|a|-3.
\]
This contradiction proves $n=\beta$. The unchanged root chart makes $n$ odd, so the parity-sensitive input argument above gives $X\in\mathbb Z$.
The first contradiction uses order40; lowering only the scalar depth
does not weaken that test.

\subsection{Polynomial circuit reductions}
For retained field values $Y_i$, each order block evaluates
$\prod_i(1+qY_i)$. In order1 group the nine factors as
\[
 [V_0,U_1,U_2],\quad[V_1,U_3,U_4],\quad[V_2,V_3,V_4],
\]
where labels denote the squared test values with their existing
conjugations. Order10 uses its pair; order40 retains its pair and
triple. For a triple $a,b,c$ the polynomial is
\[
 1+q(a+b+c)+q^2(ab+ac+bc)+q^3abc.
\]
Its coefficients can be shared among the three detectors. If an old
private coordinate is $d=bc$ and $e_2=a(b+c)+d$, eliminate it by
\[
 d=e_2-a(b+c),\qquad e_3=ae_2-a^2(b+c).
\]
For the chosen first members the square $a^2$ already has a retained
quadratic expression from the logarithm block. Substitution therefore
keeps the new residual quadratic. The reverse map reinserts the displayed
d. Every other use is substituted as well, including affine-normalized
W coordinates. Such a coordinate has the explicit inverse
$W_{\rm old}=(W_{\rm new}-r)/c$ with a fixed nonzero rational c.
These are polynomial reconstructions apart from division by constants.
They do not alter the field test values or pole conditions.

For the order weights let E and q be the retained nonzero scale and
order-one weight. Use
\[
 Ea=q,\quad Eb=a^2,\quad qc=b^2,\quad cd=q^2q^3,
 \quad cq_{10}=d^2,\quad cq_{40}=q_{10}^4.
\]
The last fourth power is evaluated through its retained square.
These equations give $c=q^3/E^6$, $d=q^2E^6$,
$q_{10}=qE^{18}$ and $q_{40}=qE^{78}$. All divisions are licensed by
the unchanged norm/trace guards. The expression $q^5=q^2q^3$
removes the former fifth-power coordinate without increasing degree.

An exact elimination is certified by a residual identity
\[
 F_{\rm new}=F_{\rm old}|_{x=q}
              +\sum_j A_jG_j,
\]
where the retained $G_j$ vanish and the old defining equation has a
nonzero constant pivot in x. Forgetting x proves the forward map;
reinserting x=q proves the reverse map. Checking all consumers is
part of this argument, not an inference from a lower variable count.

\subsection{Factored poles and the preparation conic}
For a normalized detector, the old second-conic rows include
\[
 17E-Eb-17F=0,\quad t_2=t^2,\quad q_2=q^2,\quad
 17bq_2+867F^2-34Ft_2-t_2^2=0.
\]
Replace $F=t v$, $q=t Q$, $q_2=t^2Q_2$ and use
\[
 Q_2=Q^2,\qquad17bQ_2+867v^2-34tv-t^2=0,
 \qquad17E-Eb-17tv=0.
\]
The old conic is $t^2$ times the new conic, so every new solution
reconstructs an old one, including t=0. For completeness the retained
construction supplies positive semantic b. If t!=0 divide by t;
if t=0 the old equation forces F=q=0, so the zero lift works.
The same argument applies to the scalar pole after its exact fixed
normalization. Its stored input has the opposite sign from the semantic
positive input; positivity must not be inferred from the stored symbol.

For scalar preparation, only the square N of the prepared value is
consumed outside its private block. The old equations are
\[
 t_2=t^2,\quad\mathfrak b(T-t_2)=1,\quad TW=2,
 \quad N=(W-\mathfrak b)^2.
\]
Replace them by
\[
 U^2=N,\qquad \mathfrak bJ=1,\qquad JN+Y^2=\mathfrak b.
\]
Forward take $U=W-\mathfrak b$, $J=T-t_2$, $Y=tW$. Conversely, if
$d=\mathfrak b+U\ne0$, take $W=d$, $t=Y/d$, $t_2=t^2$, $T=t_2+J$.
The conic identity gives TW=2. If d=0, then Y=0 and
$N=\mathfrak b^2$; choose $W=2\mathfrak b$, $t=t_2=0$, $T=J$.
This changes only the sign of a private unsquared value and preserves N.
The inverse guard is retained, and no rational point on an unrelated
conic is being assumed.

\subsection{Specializing the scalar preparation to zero}
Write $q=\operatorname{Norm}(f(H))$, and write the actual scalar ratio as
$z(R)/z(H)=z_0+z_1\rho$. The retained expressions are
\[
 B=(z_0^2+rz_1^2)/r^{32},\qquad
 W=(B^2+rq^2)/r^{32},\qquad \mathfrak b=crW^2.
\]
Here $r,c$ are positive and $q\ne0$, so $W,\mathfrak b>0$ at every
rational solution, independently of a private preparation inverse.
The earlier preparation-zero lemma applies when $\mathfrak b<17$, its
reduced denominator is a square, it is locally square at $2,3,17$, and
its valuations there are at least $3,1,2$, respectively. We now arrange
these conditions at the same points used in the final construction.

Let $G$ be the closure of the simultaneous real orbit of $P$ and its
conjugate, let $G^0$ be its identity component, and put $m_0=|G/G^0|$.
The real-progression lemma says that a progression with both residue
and modulus divisible by $m_0$ has dense positive tail in $G^0$.
The functions $f,z$ extend continuously to the identity with value zero.
Choose a multiple $M$ of $64m_0$ whose multiples give all desired formal
depths at primes in $S=\operatorname{supp}(cr)\cup\{2,3,17\}$.
In particular require coefficient depth $8\eta$ for $f(H)$ when
$\eta=v_p(r)>0$, and positive coefficient depth at17. Density permits
$H_0=[h_0]P$, $h_0>0$, $M\mid h_0$, as close to the two real identities
as needed to make $q(H_0)^2/r^{31}$ arbitrarily small.
Its real $z$ values are nonzero.

With $H_0$ fixed, choose a multiple $N_0$ of $6m_0$ imposing the anchor
and all required scalar finite depths. At $p\mid r$ require coefficient
depth $24\eta$ for $z(R)/z(H_0)$; impose positive depths also at17 and
at primes dividing $c$ but not $r$. The formal subgroups have finite
index, so arbitrary further depths can be imposed this way. Density of
the positive progression $s\equiv0\pmod{N_0}$ permits $s$ arbitrarily
large and $R=[s]P$ arbitrarily close to both real identities. The two
real ratios, hence their rational coefficients $z_0,z_1$, become small.
Choose them so that $\mathfrak b(H_0,R)<17/4$. This also permits the
size part of the anchor; set $a=s+1$ as before.

At $p\mid r$ these depths give $v_p(q)\ge16\eta$,
$v_p(B)\ge16\eta$ and $v_p(W)\ge0$. At $p\mid c$, $p\nmid r$,
the imposed positive depths again give $v_p(W)\ge0$. Outside this
support $v_p(\mathfrak b)=2v_p(W)$, so the reduced denominator of
$\mathfrak b$ is a square. Its squareclasses at $2,3,17$ are those of
$cr$. Since $v_2(cr)=4$, $v_3(cr)=2$, the first two required bounds
hold. At17 both $r,c$ are units; the separately imposed depths give
$v_{17}(q),v_{17}(B)\ge2$, hence $v_{17}(W)\ge4$ and
$v_{17}(\mathfrak b)\ge8$. Inclusion of17 does not depend on it
belonging to the support of $cr$.

Keep $R$ fixed. A nonempty real neighborhood $U\subset G^0$ of $H_0$
has $0<\mathfrak b(H,R)<17$. A sufficiently divisible continuity
modulus $L$, divisible by $64m_0$, preserves all the finite conditions
under $h\equiv h_0\pmod L$. Choose the remaining labels, root points
and height scales before the contact primes, and enlarge $L$ for their
finite continuity conditions. The finitely many comparison pole sets
have measure zero in $G^0$ by the real-progression lemma. Choose a
nonempty open $U'$, with compact closure in $U$, avoiding them. All
comparison functions are now uniformly bounded on that closure, by
bounds fixed before any contact prime is chosen.

The original simultaneous-prime and separation argument is valid for
this enlarged fixed $L$. Its compatible congruences are
\[
 h\equiv0\pmod{\Pi^{19}M_-},\qquad
 h\equiv1\pmod{M_+},\qquad h\equiv h_0\pmod L.
\]
Their solution progression has residue and modulus divisible by $m_0$.
Its positive tail is therefore dense in $G^0$ and meets $U'$ at
arbitrarily large indices. Choose the final $H$ there. It has every
required primary and opposite-place depth, all scalar finite depths,
the real comparison bounds and $0<\mathfrak b<17$ simultaneously.
The preparation-zero lemma supplies fresh pole witnesses with prepared
value exactly $\mathfrak b$. A zero factor in the factored pole is
handled by the zero lift already proved above.

For clarity, the physical elimination is also exact. With its fixed
positive normalization $K$, the semantic input is
$\mathfrak b=-Kx$ and the consumed square is $K^2n$. The old rows are
\[
 U^2=K^2n,\qquad1+KxJ=0,\qquad Kx+K^2Jn+V^2=0.
\]
Replace all three by $n=x^2$ and forget $U,J,V$. At any new solution
$\mathfrak b>0$ as proved independently above; restore
$U=\mathfrak b$, $J=1/\mathfrak b$, $V=0$. Each old row is then zero,
and every outside value is unchanged. The opposite direction for the
integer input projection is the simultaneous construction, not a claim
that an arbitrary old preparation witness can be specialized.

\subsection{Normalizing the remaining scalar pole}
After the preparation-zero replacement, let $x$ again denote the
physical scalar input coordinate, so $\mathfrak b=-Kx>0$ and $x\ne0$.
This nonvanishing was proved from the unchanged point and scalar root
rows; it does not depend on the preparation square or on the private
pole equations. Assign weight one to the five private slots
$(e,h,p,q,v)$ and weight two to their square slots $(p_2,q_2)$.
Replace each old slot of weight $j$ by $x^j$ times its new slot,
and discard the preparation-square slot $n$, setting it equal to $x^2$.
Every one of the five remaining pole residuals is then $x^2$ times
its new residual. For example the first conic and the h equation have
respectively their $ex$ and $x^2$ terms replaced by $e$ and1, and
\[
 17ex+Ken-17hv=0
       \quad\longmapsto\quad 17e+Kxe-17hv=0.
\]
The square rows remain $p_2=p^2$, $q_2=q^2$, and the factored
second-conic coefficient list is unchanged. In the physical indexing
of the reviewed component, the weight-one slots are
$51,52,53,54,532$, the weight-two slots are $527,534$, and $531$ is n.
These are implementation indices, not mathematical constants.

From an old solution divide each private slot by its indicated power
of x. From a new solution multiply by that power and restore $n=x^2$.
These maps are inverse and fix every outside coordinate. The sign of x
is immaterial, and no division uses a private pole value. In particular
all zero branches of p,q,v or h are preserved. These seven slots occur only in the five displayed pole rows, and n
has only its former definition and these consumers. Thus there is no
untransformed outside use of a rescaled slot.

\subsection{Rescaling the scalar anchor}
The odd scalar root S has the unchanged chart
$S=(x,hj)$, $x=h^2$, $j^2=x^2-D$, hence $S=[a]P$ for an odd integer a.
Write the retained reciprocal difference equation as
\[
 (c_*x+D)g=a_*j-b_*h,\qquad g=g_0+g_1\rho.
\]
Its denominator is positive in both real embeddings. We replace the
old inverse of the irrational coefficient of f(S-P) by
$g_0\tau=1$. This implies $g\ne0$ before using the anchor.
The subtraction chart identifies $g^2=f(S-P)$: at $S=P$ the quotient
would be zero, which is excluded, and at $S=-P$ the explicit duplication
formula gives the same value $f(-2P)$. Thus $S-P=[s]P$ has a nonzero
even index $s=a-1$ without any prior anchor divisibility assertion.

The earlier scalar coefficient lemma applies to every such s and gives
\[
 v_3(f_1([s]P))=2+2v_3(s).
\]
Its hypotheses contain no anchor equation. Indeed
$z(2P)/3\equiv2+\rho\pmod3$; multiplication by units and then by
powers of3 in the integral formal group preserve the stated leading
term. The curve equation makes the leading residue of
$f/3^{2+2v_3(s)}$ a nonzero rational multiple of $\rho$.
Independence of1 and $\rho$ in the residue field proves the displayed
coefficient valuation, also for negative s. Consequently
$F=f_1([s]P)=2g_0g_1\ne0$, in particular $g_1\ne0$.
This is the noncircular justification for restoring the old inverse.

The old three-row subsystem is
\[
 F=2g_0g_1,\quad F\sigma=1,\quad
 u^2+6v^2-2w^2-3F^2-4Cw=0.
\]
The new subsystem is
\[
 g_0\tau=1,\qquad
 U^2+6V^2-2W^2-12g_1^2-4C\tau W=0.
\]
Forward put $(U,V,W)=(u,v,w)/g_0$, $\tau=1/g_0$ and forget F.
Backward put $(u,v,w)=g_0(U,V,W)$, restore $F=2g_0g_1$ and
$\sigma=1/F$. If $q_{\rm old},q_{\rm new}$ denote the anchor residuals,
the exact identity is
\[
 q_{\rm old}|_{(u,v,w)=g_0(U,V,W),\ F=2g_0g_1}
       -g_0^2q_{\rm new}
       =4Cg_0W(g_0\tau-1).
\]
These inverse maps fix the input and all other coordinates. They permit
zero values of U,V,W. They are equivalences for the full point system,
not for an isolated anchor with arbitrary g: that isolated equation
would admit $g_1=0$, which the coefficient lemma just excluded.

When C remains a fixed integer coefficient, both new residuals are
quadratic and their squared sum contributes $16C^2\tau^2W^2$.
For the ordinary companion compute C by the quadratic chain
$c_0=3^{1024}$, $c_j=c_{j-1}^2$ for $1\le j\le15$, so that
$c_{15}=3^{2^{25}}=C$. Introduce a separate multiplication
slot $T_C$ and equation
\[
                 T_C=C\tau.
\]
Use $-4T_CW$ in the anchor. Treating a computed C as a coefficient would
instead make $C\tau W$ cubic, so this extra slot is necessary for that
companion. Both directions restore $T_C$ uniquely; no nonvanishing or
point-domain argument is changed. This slot and its gate, together with the sixteen constant-chain slots,
are included in the ordinary companion.

\subsection{Completeness and finite certification}
Fix an integer X and compute the positive deterministic scalar correction.
Choose the preliminary shared point and then the deep nonzero scalar
point by the simultaneous real and finite-place construction just proved.
Choose the odd Lucas labels, then their actual root points and positive
root scales. Fix a continuity modulus for all the new rational functions.
Only after these choices select the contact primes and solve the retained
primary/opposite congruences. Take a sufficiently large common index h in the prescribed real open
set, avoiding the finitely many old exceptional point loci. This gives
$U_i=[hn_i]P$ and $V_i=[(h-1)n_i]P$, all nonidentity in this completeness
construction. The new rational logarithms have the required local
precision, and the original real pole estimates and arbitrary finite
preparation depths give positive semantic pole inputs. Obtain fresh
pole witnesses with scalar preparation equal to its input, apply the
factored-conic maps, and reconstruct each
finite private circuit. These are simultaneous choices; no claim is
made that an arbitrary old pole witness survives a changed logarithm
or changed scalar root.

\subsection{The specified polynomials and coefficient tables}
Clear denominators and divide each nonzero quadratic residual by its
integer content. Define $P_{394}$ to be the sum of the squares of the
364 resulting integer residuals. Over the rational numbers this sum
vanishes exactly when every residual vanishes. The soundness and
simultaneous completeness arguments above therefore prove its stated
integer projection. The sixteen constant-chain equations and the
$T_C=C\tau$ equation give the ordinary companion: its 381 normalized
residual squares define $P_{411}$, with the same integer projection.

Expansion of the specified residual lists gives respectively 6,338 and
6,388 distinct nonzero monomials. Both polynomials have total degree
four and content one. For the symbolic polynomial the largest absolute
coefficient is $16C^2=16\cdot3^{2^{26}}$; for the ordinary polynomial it
has 3,247 bits and is strictly less than $2^{3247}$. These counts and
bounds are computed from the complete coefficient lists, not by adding
component savings. Every listed coordinate occurs.

The coefficient artifacts use consecutive coordinate indices, with
index one for X. The symbolic JSON coefficient table has 973,221 bytes,
and the ordinary JSON coefficient table has 982,235 bytes. Their plain
text tables have 212,067 and 214,957 bytes, respectively. The tables,
residual lists, coordinate maps and exact maximum coefficients specify
the two polynomials independently of their factored construction.

% Isolated Version 14 author draft; assembled independent review is required.
% This file does not modify the original polynomial or Version 13.
\section{A quartic presentation with 299 variables}\label{sec:v14}

We retain the arithmetic inputs, qualifications and source editions of
Sections~\ref{subsec:v12-inputs} and~\ref{sec:v13}. The polynomial of
the main construction is unchanged. The following are smaller
presentations of the same rational input projection.
\begin{theorem}\label{thm:v14-quartic}
There is a specified integer polynomial $P_{299}$ of total degree four
in 299 variables such that
\[
 X\in\mathbb Z\quad\Longleftrightarrow\quad
 \exists\mathbf v\in\mathbb Q^{298}:P_{299}(X,\mathbf v)=0.
\]
It has 6,630 nonzero monomials, content one, and largest absolute
coefficient $186624\cdot3^{2^{26}}$. Its coefficient artifact writes
$C=3^{2^{25}}$ as a fixed integer, not a quantified variable.
There is also a specified integer quartic $P_{314}$ with
\[
 X\in\mathbb Z\quad\Longleftrightarrow\quad
 \exists\mathbf w\in\mathbb Q^{313}:P_{314}(X,\mathbf w)=0.
\]
It has 6,672 nonzero monomials and content one. Every coefficient has
absolute value strictly less than $2^{6496}$; its largest coefficient
has exactly 6,496 bits. The compact ordinary decimal coefficient table
occupies 264,430 bytes in its specified ASCII format. The separate table
retaining source coordinate identifiers occupies 222,890 bytes.
\end{theorem}
The ordinary coefficient bound here is different from the bound for
$P_{411}$. We prove the new bound for the new artifact. Some changes
below are inverse changes of private witnesses; others change an
arithmetic certificate or point domain. For the latter we prove
soundness and completeness separately.

\subsection{Conventions and retained interfaces}\label{subsec:v14-inputs}
Put $K=\mathbb Q(\rho)$, $\rho^2=2$, $d=3-\rho$, $D=d^2=11-6\rho$,
and $E:y^2=x^3-Dx$. Retain
$E(K)=\mathbb ZR_0+E[2](K)$ and $2E(K)=\mathbb ZP$, where
\[
 R_0=(1+\rho,2\rho),\quad P=(c_*,a_*b_*),\quad
 a_*=-1+7\rho/4,\quad b_*=39/8-9\rho/2,\quad c_*=a_*^2.
\]
Thus $b_*^2=c_*^2-D$. Write $f(Q)=1/x(Q)$ and
$z(Q)=-x(Q)/y(Q)$, with both functions extended regularly by zero
at $O$. The roots are $Q_i=[n_i]P$ with labels
$(n_1,n_2,n_3,n_4)=(n,a,t,b)$. Their nonzero odd integral labels
will be restored before height or logarithm arguments. The scalar
point is $Q_2-P=[s]P$, where $s=a-1$.

The retained inputs include the squareclass calculation, formal
expansions and logarithm estimates, simultaneous contact-prime
construction with its real-density condition, preparation-zero pole
criterion, root-height implication, Lucas argument and scalar-depth-24
argument. In particular the three positive rational square scales are
\[
 A_0=1,\quad A_1=C_0^2,\quad A_2=(1+C_1^2)^2,\qquad
 \prod_i x(Q_i)=C_0+C_1\rho .
\]
The same root-height conclusion is
\[
 B^2-6<7\mathcal H|N|^{19/k},\qquad
 B=\max_i|n_i|,\quad 1\leq\mathcal H<2^{2093966},
\]
when the $k$th powers of the three required contact radicals divide
the nonzero integer $N$. The independently restored scalar anchor
gives $B>(40\mathcal H)^{10}$, excluding all three radicals equal
to one without using H or auxiliary logarithms.
The retained preparation-zero equations exclude negative odd
valuations of a pole input at required large primes. Their
completeness interface uses the same small-prime squareclasses and
depths, square reduced denominator, and real interval $(0,17)$.
We verify the changed consumers of these inputs below.

\subsection{The fixed-difference chart for H}\label{subsec:v14-H}
First replace the old $x(H)$ pair by the equations
\begin{equation}\label{eq:v14-H-raw}
 c_*(f+F)+Dv^2+2b_*v=1,\qquad v=ft,\qquad F=vt.
\end{equation}
The common primary normalizer is the fixed affine expression
\begin{equation}\label{eq:v14-normalizer}
 \mathcal S=(b_*+c_*t+Dv)/a_*.
\end{equation}
Its substitution into every primary, logarithm and scalar comparison
row preserves quadratic degree. If $f=0$, the last two equations
force $v=F=0$, contradicting the first. Set $x=1/f$. They give
$v^2=fF$ and $a_*f\mathcal S-b_*f-(c_*+Df)v=0$.
Here is an unconditional certificate for the recovered curve.
Put $d_1=c_*+Df$, $B_1=b_*f+d_1v$,
$A_1=a_*f\mathcal S-B_1$, and
$K_1=c_*(f+F)+Dv^2+2b_*v-1$. Then
\[
 c_*f(f\mathcal S^2+Df^2-1)
 =fd_1K_1+c_*d_1(v^2-fF)
 +(b_*^2-c_*^2+D)f^2+A_1(a_*f\mathcal S+B_1).
\]
Thus $\mathcal S^2=x-Df$ and the old fixed-difference equations
hold. This direction divides only by $c_*f$. Conversely, the old
full chart gives $t=v/f$ and the same affine expression for
$\mathcal S$; its previously established exceptional-point exclusion
supplies the forward implication. All outside values are preserved.

Next use the already needed outputs $U=f(H)^2$, $V=f(H-P)^2$.
Store $W,v,t,U,V$, put $Q=(1-DW-2b_*v)/c_*$, and impose
\begin{equation}\label{eq:v14-H-squared}
 W=v^2,\qquad UV=W^2,\qquad Q^2=U+V+2W,\qquad t(U+W)=Qv.
\end{equation}
On $UQ\ne0$, set $f=(U+W)/Q$, $F=(V+W)/Q$. If
$E_1=UV-W^2$ and $E_2=Q^2-U-V-2W$, then
\begin{align*}
 (U+W)^2-UQ^2&=-UE_2-E_1,\\
 (V+W)^2-VQ^2&=-VE_2-E_1,\\
 (U+W)(V+W)-WQ^2&=-WE_2+E_1.
\end{align*}
These give $f^2=U$, $F^2=V$, $fF=W$, $f+F=Q$, $ft=v$,
and $F=vt$, restoring~\eqref{eq:v14-H-raw}.

The extra branches are classified before this reconstruction.
For $U=0$ one has $W=v=0$, $Q=1/c_*$, $V=1/c_*^2$.
For $Q=0$ one has $U=V=-v^2$ and
\[
 Dv^2+2b_*v-1=0,\qquad
 v=\frac{3+2\rho}{4}\quad\hbox{or}\quad
 v=\frac{-36+16\rho}{49}.
\]
These are units at every place above $p>76176$:
$c_*+b_*=12-8\rho$ has norm16 and
$c_*-b_*=(9+4\rho)/4$ has norm$49/16$.
Section~\ref{subsec:v14-groups} excludes these branches by their
unchanged $V_0$ unit leaf and the pre-H height/anchor contradiction.
No division by $UQ$ is used to prove its own legitimacy.

\subsection{Root parity and the shifted anchor}\label{subsec:v14-parity}
The root chart $c_i=a_i^2$, $b_i^2=c_i^2-D$ represents
$(c_i,a_ib_i)$. Its abscissa is nonzero; for the final presentation
this follows from the retained product-height inverse as proved in
Section~\ref{subsec:v14-reciprocal}.
The squareclass map has independent values
\[
 \alpha(R_0)=[1+\rho],\qquad \alpha(T_0)=[-1],\qquad
 \alpha(T_d)=[d].
\]
Indeed norms in a putative square relation force the exponents
of $1+\rho$ and $d$ even, and the real embeddings exclude the
remaining class $[-1]$. Thus the square abscissa already puts
each root in $2E(K)=\mathbb ZP$, with nonzero integral label,
before an odd-chart condition.

Keep the scalar reciprocal equations
\[
 (c_*c_2+D)g=a_*b_2-b_*a_2,\qquad
 g=g_0+g_1\rho,\qquad g_0\tau=1.
\]
The denominator is nonzero at both real embeddings. The subtraction
identity gives $f(Q_2-P)=g^2$. At equal abscissae, $Q_2=P$
contradicts the guard and $Q_2=-P$ gives the exact duplication
value for $-2P$. Hence $s=n_2-1\ne0$ is an integer before height
or logarithms. Write $F_s=2g_0g_1$.

Let $X_n$ be the $\rho$ coefficient of $x([n]P)$. For $n\ne0$,
\begin{equation}\label{eq:v14-X}
 v_3(X_n)=
 \begin{cases}0&n\text{ odd},\\-2-2v_3(n)&n\text{ even},\end{cases}
 \qquad v_2(X_n)=-1\quad(n\text{ odd}).
\end{equation}
At3 the reduction of $P$ is $(\rho,0)$ over $\mathbb F_9$,
which proves the odd case. For $n=2\cdot3^ku$, $u$ a3-unit,
the formal expansion gives
$z([n]P)/3^{k+1}\equiv u(2+\rho)\pmod3$.
Since $x(z)=z^{-2}+O(z^2)$ and
$(2+\rho)^{-2}=3/2-\rho\equiv2\rho\pmod3$, the scaled
abscissa has nonzero imaginary coefficient modulo3. This proves
the even case and $X_n\ne0$.
At2, $z(P)=(232+86\rho)/357$ has component valuations3 and1.
The integral odd multiplication series preserves these values:
its error has field valuation at least$9/2$ when $v(2)=1$.
For $z=A+B\rho$ with those valuations the imaginary coefficient
of $z^{-2}$ is $-2AB/(A^2-2B^2)^2$, of valuation$-1$.
The Laurent correction has larger valuation. This uses integral
formal coefficients, not good reduction at2.

The retained scalar-coefficient lemma gives $v_3(F_n)=0$ for odd
$n$, $2+2v_3(n)$ for even nonzero $n$, and dyadic integrality.
It is applied with both parities. Put
\begin{equation}\label{eq:v14-q}
 q=12(X_A^2+X_T^2)+108X_Z^2,\qquad T=F_s+Cq,
\end{equation}
and use the private anchor
\begin{equation}\label{eq:v14-anchor}
 u^2+6v_a^2-2w_a^2-4Cw_a-3T^2=0.
\end{equation}
The inherited quaternion criterion, including parameter zero,
supplies rational witnesses exactly when $T/C$ is integral at2
and3. Its3-integrality forces all three summands in $F_s/C+q$
to be3-integral. Their possible finite valuations lie in the
disjoint classes $2\mathbb Z$, $1+4\mathbb Z$, $3+4\mathbb Z$;
the binary sum uses anisotropy modulo3 and~\eqref{eq:v14-X}.
A negative minimum cannot cancel. Thus $X_A,X_T,X_Z$ are3-units,
so the three root labels are odd. Also $F_s/C$ is3-integral
and is dyadically integral by the scalar lemma. The old unshifted
anchor can therefore be restored with new private witnesses
before H or height. An odd $s$ would make $F_s$ a3-unit,
so $s$ is even and the S-root label $s+1$ is odd.

For an odd root, its half point $[n_i]R_0=(A,B)$ has
$A=(1+\rho)r^2$. Duplication gives
\[
 a_{\rm half}=(A^2+D)/(2B),\qquad
 b_{\rm half}=(B/A)^2/2-c_i,
\]
with $a_{\rm half}^2=c_i$ and the correct product ordinate.
After changing both root signs if necessary, the old half-chart
equations are
$A^2-2(c_i+b_{\rm half})A-D=0$ and
$\operatorname N(A)+\operatorname N(r)^2=0$.
In the final A,T,Z modules the sign freedom is $a_i\mapsto-a_i$,
with reconstructed $b_i\mapsto-b_i$. For S negate
$a_2,b_2,g_0,g_1,\tau$ together. These actions fix the actual
points, $F_s$, all detector values and the private unscaled anchor
witnesses. The scalar rows change at most by an overall sign.

Conversely the old odd roots make $q$ integral at2 and3 by
\eqref{eq:v14-X}, so the same quaternion criterion supplies the
shifted anchor. No prime or H choice changes. The exact private
unscaling used here multiplies the old anchor witnesses and its
parameter $B$ by $g_0$, with inverse multiplication by $\tau$.
It gives $T=g_0B=F_s+Cq$ and~\eqref{eq:v14-anchor}.
This exact step assumes no private anchor witness nonzero.

More explicitly, let the old private residuals be
$R_Q=Q-q$, $R_B=B-2g_1-CQ\tau$, and
$R_A=U^2+6V^2-2W^2-4C\tau W-3B^2$.
Let $N_T=T-2g_0g_1-Cq$ and let $N_A$ be the new
anchor residual. With $G=g_0\tau-1$, pullback by the
forward map gives the unconditional identities
\[
 N_T=g_0R_B+CR_Q+CQG,\qquad
 N_A=g_0^2R_A+4Cg_0WG.
\]
The reverse map is
$Q=q$, $B=\tau T$, $U=\tau u$, $V=\tau v_a$,
$W=\tau w_a$. Its pullbacks are
\[
 R_Q=0,\qquad R_B=\tau N_T+2g_1G,\qquad R_A=\tau^2N_A.
\]
Thus every private anchor consumer is restored with its
literal normalization, and all nonprivate coordinates are fixed.

\subsection{Contact poles and their conditional domains}\label{subsec:v14-contact}
For each scale put
\[
 R_j=F_5(A_j\lambda),\quad E_jR_j=1,\qquad
 F_5(t)=t(t+1/2)(t+1)(t-1/2399)(t+32/407).
\]
The retained contact gates make $R_j$ affine in existing coordinates.
With $W_j=G_{j,0}G_{j,1}G_{j,2}$ retain only
\begin{equation}\label{eq:v14-pole}
 \operatorname N(W_j)=2\operatorname{Re}(W_j)z_j,\qquad
 B_j(R_j+E_j+2)=z_j^2.
\end{equation}
There is no quantified inverse of $\operatorname{Re}(W_j)$.
At a required contact prime, $\alpha=v_p(R_j)>0$ is odd
and $1+R_j$ is a unit. The groups below give $W_j\ne0$.
The norm row then forces its real component and $z_j$ nonzero:
$w_0^2-2w_1^2=0$ over $\mathbb Q$ forces both components zero.
Thus
\[
 T_j^{\rm tr}=\operatorname{Tr}(W_j^{-1})=1/z_j,\qquad
 B_j=\frac{R_j}{(T_j^{\rm tr})^2(1+R_j)^2}.
\]
These domains are needed only at required contact primes.
A detector with radical one has empty local obligations;
no global trace inverse is asserted for it. The case $R_j=-1$
likewise cannot occur at a required contact.

In completeness the ratio between $R_j$ and the earlier contact
rational function is $A_j^5(t-1)^6$, a rational square and a
unit at selected contacts. Hence $R_j$ has valuation one there
and even valuation outside the prescribed odd support.
Both denominators in $B_j$ change valuations by even integers.
This includes additional primes introduced by later witnesses.

\subsection{Common-scale groups and the successful place}\label{subsec:v14-groups}
Use the physical tests
\[
 H_0=f(H),\quad H_i=L_{i+1}-\overline{LP_{i+1}}\ (1\leq i\leq3),
 \quad H_4=\mathcal L(L_2,L_3,L_4),
\]
\[
 M_0=L_1-\overline{LP_1},\quad M_1=L_1-\beta(X),\quad
 P_i=f(U_i)^2,\quad Z_1=\overline{f(V_1)^4},\quad
 T_i=\overline{f(V_i)^8}.
\]
Here $\beta(X)=204(X^2+X)+2555$ is the retained comparison
polynomial. Conjugation is part of the physical tests.
The Lucas polynomial is
\[
 \mathcal L(A,T,Z)=T^2-A(A-2)TZ+(A-2)^2Z^2-1.
\]
The affine recentering $C_L=(A-2)Z$, $V_L=AT$ gives
$\mathcal L+1=T^2+C_L^2-C_LV_L$ without changing its value.

For a nonzero common rational $\gamma$, store
\[
 ZH_i=H_i^4/\gamma^{160},\quad ZM_i=M_i^4/\gamma^{40},\quad
 ZP_2=P_2^2/\gamma^{160},\quad
 ZP_3=P_3^2/\gamma^{80},\quad ZP_4=P_4^2/\gamma^{80},
\]
and $ZP_1=P_1^2$. These definitions are quadratic using the
retained square inputs and common powers. For each detector,
\begin{align}
 G_0&=1+RZH_1+R^2ZH_2+EU+E^4Z_1+\gamma^{16}E^8,\label{eq:v14-g0}\\
 G_1&=1+RZH_3+R^2ZH_4+R^3ZP_2+E^4ZP_4+EZP_1+E^8T_0,\label{eq:v14-g1}\\
 G_2&=1+RZM_0+R^2ZM_1+R^3ZP_3+E^8T_2+E^4T_3+ET_4.\label{eq:v14-g2}
\end{align}
Here $U=H_0^2$ is the actual squared-H output even before
H is reconstructed. The product, trace and pole use these
same final groups.

At a consumed place set $\alpha=v(R)>0$ odd and
$c_\gamma=v(\gamma)\in\mathbb Z$. For a reciprocal point
$h^2=f(1-Df^2)$, a nonzero $f$ has even valuation: the leading
valuation is $v(f)$ when positive and $3v(f)$ when negative.
Zero leaves are omitted. Thus the relevant powers have
periods4,8,16 before formality or logarithms. In the generic
squared-H branch its algebraic reconstruction gives the same
curve fact; its exceptional branches have the fixed signed
values already classified. Other high and middle tests are
literal fourth powers, independently of their later ratio meanings.

In term order, the valuation periods and coefficients of
$\alpha$ are
\[
 \begin{array}{c|l}
 G_0&(4,1),(4,2),(4,-1),(8,-4),(16,-8)\\
 G_1&(4,1),(4,2),(8,3),(8,-4),(8,-1),(16,-8)\\
 G_2&(4,1),(4,2),(8,3),(16,-8),(16,-4),(16,-1).
 \end{array}
\]
All gamma exponents are divisible by the relevant periods.
Within each row the classes modulo16 are disjoint and exclude
zero for every odd $\alpha$. Every group therefore has a unique
minimum valuation, is nonzero, and has valuation at most zero.
This covers arbitrary $c_\gamma$ and every zero-leaf pattern.

Let $\ell_p$ be the maximum valuation of $W$ over places above
$p$. Trace integrality gives $v_p(T^{\rm tr})\geq-\ell_p$.
The pole obstruction and~\eqref{eq:v14-pole} give
\[
 0\leq v_p(B)=\alpha-2v_p(T^{\rm tr})\leq\alpha+2\ell_p.
\]
Its odd parity makes $\ell_p>-\alpha/2$. At a place attaining
the maximum, every group, and then every individual term,
has valuation greater than $-\alpha/2$. This is one place for
all tests, also at an inert unramified prime.

The guard gives $c_\gamma>15\alpha/32$. The high and middle
terms give
\[
 4v(H_i)>160c_\gamma-(r+1/2)\alpha,\qquad
 4v(M_i)>40c_\gamma-(r+1/2)\alpha\quad(r=1\text{ or }2).
\]
For $\alpha=1$ the worst bounds are157.5 and37.5; for
$\alpha=3$ they are312.5 and72.5; for $\alpha\geq5$ they
exceed$72.5\alpha$ and$16.25\alpha$. Thus $v(H_i)\geq40$
and $v(M_i)\geq10$. The primary bounds are
\[
 \begin{split}
 4v(f(U_2))&>160c_\gamma-3.5\alpha,\\
 4v(f(U_3))&>80c_\gamma-3.5\alpha,\\
 4v(f(U_4))&>80c_\gamma+3.5\alpha.
 \end{split}
\]
Using even valuations, these give respectively40,20,22 for the
reciprocal depths. The other low terms force positive valuation
of every required nonzero primary or alternate reciprocal value.

On either exceptional squared-H branch $T_0$ is a unit at every
consumed place. Its $E^8T_0$ term has valuation$-8\alpha$,
contradicting the common strict bound. Either branch would
therefore make all radicals one, contrary to the pre-H
root-height/anchor bounds. The unchanged fixed-point exclusions
from Section~\ref{sec:v13} then restore the full H/V0 domain
before logarithms. Here their required hypothesis is only
positive valuation of $f(H)$, rather than the former stronger
order40 test. A nonidentity torsion value of H, or a value
in either of the fixed translates $\pm P+E[2](K)$, has the
same fixed unit obstruction used there at every consumed
large prime. The new term $EU$, where $U=f(H)^2$ after
reconstruction, gives $2v(f(H))>\alpha/2>0$ and excludes
those fixed values. The identity is already outside the
reconstructed affine H chart. Thus weakening the old H test
does not weaken these domain exclusions; no high-depth
logarithm estimate is used in them.

\subsection{Reciprocal point charts and their orientation}
\label{subsec:v14-reciprocal}
For each of the four root indices use reciprocal coordinates
\[
 f=1/x(U),\quad h=-y(U)/x(U)^2,\qquad
 F=1/x(V),\quad k=-y(V)/x(V)^2.
\]
The pair $(0,0)$ denotes $O$. Retain the square caches $Y=f^2$,
$B=F^2$. In particular the sign of $h$ and $k$ is negative.
The reduced chart uses only $f,h,F,s,Y,B$ and imposes
\begin{align}
 h^2&=f-DfY,\label{eq:v14-primary-curve}\\
 e&=s^2-c+D(f+F)=0,\label{eq:v14-slope-square}\\
 q&=2sh-c(f-F)-1+D(2Y+fF)=0.\label{eq:v14-reduced-slope}
\end{align}
Here $s$ is a private slope, unrelated to the scalar integer
$a-1$ when that integer is denoted by $s$ elsewhere.
All displayed equations have degree at most two in the quantified
coordinates. Products such as $fY$ use the retained cache $Y$.
Define, without adding a coordinate,
\[
 k=s(f-F)-h.
\]
Put $r_1=s(f-F)-h-k$, $r_2=s(h-k)-1+D(Y+fF+B)$ and
$r_4=h^2-f+Df^3$. Exact expansion, using the two caches, gives
\begin{align}
 r_1&=0,\qquad r_2=q-(f-F)e,\label{eq:v14-slope-restoration}\\
 k^2-F+DF^3&=r_4-(f-F)q+(f-F)^2e.\label{eq:v14-alt-restoration}
\end{align}
Thus both reciprocal cubics and the three slope equations hold
before a ratio, valuation or logarithm is interpreted.

These cubics describe actual points of $E(K)$. For $f\ne0$ the
inverse is $x=1/f$, $y=-h/f^2$; for $f=0$ the cubic gives $h=0$
and the point is $O$. This chart omits $(0,0)\in E(K)$ but includes
the other two finite points of order two. Those latter points will
be excluded by the low tests before division by a ratio denominator.
Translation by $(0,0)\in E(K)$ sends the represented point to
$(-Df,-Dh)$, including the case in which the represented point
is $O$.

The oriented difference is
\begin{equation}\label{eq:v14-root-ordinate}
 U-V=(c,y_R),\qquad y_R=Dh-s(Df+c).
\end{equation}
For completeness, this assertion includes every branch of the
equations. If $f\ne F$, the line with slope $s$ joins
$(-Df,-Dh)$ and $(-DF,Dk)$, so the usual secant formula gives
the stated difference. If $f=F$, the first slope equation gives
$k=-h$. The case $h=0$ is impossible: the cubic gives either
$f=0$ or $Df^2=1$, whereas the second slope equation would give
$3Df^2=1$. Hence $h\ne0$, and that equation gives
$2sh=1-3Df^2$, exactly the tangent slope at $(-Df,-Dh)$.
The tangent formula again gives~\eqref{eq:v14-root-ordinate}.
If $U=O$ the equations give $Fc=1$ and $R=-V$; if $V=O$ they
give $fc=1$ and $R=U$. Both points cannot be $O$.

There is also a division-free check on the root curve:
\begin{equation}\label{eq:v14-root-identity}
 y_R^2-c^3+Dc
 =D^2r_4-D(Df+c)(r_2-sr_1)
       +e(c+Df)(c+DF).
\end{equation}
Consequently the retained square equation $a_i^2=c_i$ recovers
the old root chart once $c_i\ne0$: put $b_i=y_R/a_i$, obtaining
$b_i^2=c_i^2-D$.

The nonzero assertion uses no auxiliary ratio. In the literal
system the successive root products are
\[
 c_1c_2=v_{337}+2\rho v_{338},\quad
 c_1c_2c_3=2(v_{339}+2\rho v_{340}),\quad
 c_1c_2c_3c_4=2(v_{341}+2\rho v_{342}).
\]
If any $c_i$ vanished, both last components would vanish.
The unchanged height-product chain would then give
$v_{343}=v_{391}=v_{392}=v_{393}=v_{394}=0$, contradicting its
retained inverse equation. Thus every $c_i$ and $a_i$ is nonzero.
The squareclass calculation now places every root in $2E(K)$.
This supplies the root chart used in the parity and pre-H
arguments above; it does not require the auxiliary points to
lie in $2E(K)$.

For the scalar root retain both old coordinates $a_2,b_2$ and
the old equation $b_2^2=c_2^2-D$. Add the quadratic orientation
equation
\begin{equation}\label{eq:v14-S-orientation}
 a_2b_2=Dh-s(Df+c_2).
\end{equation}
The scalar subsystem consumes this particular ordinate. A root
curve equation alone would allow its negative and would not
identify its scalar point with $U_2-V_2$.
The physical caches for this chart are
$Y=v_{257}+2\rho v_{258}$ and $B=v_{196}+2\rho v_{197}$.
Every old scalar consumer of $a_2,b_2,c_2$ is retained.

Set $N_0=\mathcal S$, the fixed affine expression
in~\eqref{eq:v14-normalizer}. After the H domain is restored
its meaning is $-1/z(H)$. It is not a new quotient or a new
coordinate in the pre-H equations. The primary equations are
\begin{equation}\label{eq:v14-primary-ratios}
 hL+fN_0=0,\qquad (1-DY)L+hN_0=0.
\end{equation}
The retained alternate equation is
\begin{equation}\label{eq:v14-regular-ratio}
 (1-DB)z=s(f-F)-h=k.
\end{equation}
It remains quadratic after substituting for $k$.
At $O$ these equations force $L=0$ or $z=0$, respectively,
rather than allowing a free ratio. The raw alternate equation
$kz=F$ is not imposed. Its recovery requires the global order of
the proof. With $A=1-DB$ and $g=Az-k$,
\begin{equation}\label{eq:v14-raw-ratio}
 A(kz-F)=kg+(k^2-F+DFB).
\end{equation}
The right side vanishes by the curve and regular ratio equations.
One may cancel $A$ only after the low tests have been applied.
At a consumed place they give $v(F)>0$ for a nonidentity point,
so $A$ is a local unit. At the identity $F=0$ and $A=1$.
Thus $A\ne0$ in $K$, and the raw ratio follows.

The order is therefore the same for all four root indices.
First restore the cubics, oriented group differences and nonzero
root charts. Next restore root parity and the scalar-anchor
contradiction to empty contact support. Apply the same-place low
tests using the actual cubics and literal caches. Only then
cancel $1-DB$ and interpret the ratios. In particular the scalar
auxiliary pair is not passed to an older argument which assumed
an unavailable square root of $f$ or $F$.

\subsection{Auxiliary point domains and the H-depth bootstrap}
\label{subsec:v14-logarithms}
We give the domain argument for all four pairs simultaneously.
At a good consumed place $D$ is a unit. On
$h^2=f-Df^3$, a nonzero $f$ has even valuation. If $v(f)>0$,
$2v(h)=v(f)$; if $v(f)<0$, $2v(h)=3v(f)$; and if $v(f)=0$
there is nothing to prove. The identity has $f=0$ and every
positive-depth condition is automatic. Consequently the
valuations of $f^2,f^4,f^8$ have periods $4,8,16$ on all of
$E(K)$ represented by this chart. This proves the periods used
in the group calculation before H has been restored.
The low tests exclude the finite points of order two, whose
reciprocal abscissae are units at these places.

After the fixed H/V0 domain and the raw ratios have been restored,
write
\[
 H=mR_0+T_H,\quad U_i=k_iR_0+T_i,\quad
 V_i=(k_i-2n_i)R_0+T_i,\quad
 V_0=(m-2)R_0+T_H,
\]
where all the $T$'s are in $E[2](K)$. The correct identities,
which kill torsion without imposing a doubled-point domain, are
\begin{equation}\label{eq:v14-doubled-identities}
 [2m]U_i=[2k_i]H,\qquad
 [2(m-2)]V_i=[2(k_i-2n_i)]V_0.
\end{equation}
The restored H exclusions make $H,V_0$ nontorsion.
Taking formal logarithms gives full ratios
\[
 \frac{u_i}{h_0},\qquad \frac{u_i-n_i}{h_0-1},
 \qquad u_i=k_i/2,\quad h_0=m/2.
\]
The half-indices are integral at each consumed odd prime.
At the identity the formal unit
$\log(Q)/z(Q)$ means its regular extension, whose value is one.
No quotient $0/0$ is taken.

One of $h_0,h_0-1$ is a unit. The corresponding full ratio is
integral. The depth-one formal unit estimates and the positive
cross-test transfer integrality to both approximate ratios and
both full ratios. In particular $L_2,L_3,L_4$ are integral at
the common place. Their retained Lucas test has positive
valuation and its value at $(0,0,0)$ is $-1$. At least one of
these three $L_i$ is therefore a unit. The primary bounds proved
above give formal depth at least10 for its point $U_i$.
The exact identity $z(U_i)=L_i z(H)$ then gives depth at least10
for $H$. All integral primary $L_i$ give the same depth for
every nonidentity primary point.

This bootstrap precedes the order40 estimate. At depth10 the
relative primary logarithm error has valuation at least40.
The alternate estimates have order40 for the degree-nine
approximants and order12 for the middle approximants. Their
denominators are units at the consumed place. Since the full
ratios are integral, these relative errors are also absolute
errors. Multiplying the two ratios by $h_0$ and $h_0-1$ and
subtracting gives the retained congruences to the integer root
labels $n_i$, with their required orders10 or40.
No assertion that the global auxiliary labels $u_i$ are integers
is used.

The larger auxiliary domain also requires global nonvanishing
of the rational approximant numerator and denominator.
In its variable $w=Df^2$ their coefficient lists are
\[
\begin{split}
N(w):\;&1,-9776/5365,2941/3219,-1024/13949,\\
 &-1920/237133,-1024/711399,-1408/5928325,0,
 2432/103152855,\\
Q(w):\;&1,-1526/1073,513/1073 .
\end{split}
\]
The unchanged certificates modulo71 for $N$ and modulo7 for $Q$
apply to every element of $K$. Both primes split in $K$; the
respective coefficients are integral, with unit leading
coefficient, and their reductions have no root. A root with
negative valuation is excluded by its unique leading term,
and an integral root is excluded by reduction. Thus the old
nonzero products and inverses remain justified. A positivity
interval restricted to doubled real points is not used here.

The retained root-height, Lucas and scalar-depth arguments can
now be applied with their original inputs. The roots are the
same odd integer multiples, the congruences have the required
precision, the scalar point has the orientation
in~\eqref{eq:v14-S-orientation}, and the scalar depth24 transfer
is unchanged. They conclude that the distinguished input is
an integer. This proves soundness of the changed arithmetic
certificates.

\subsection{Simultaneous completeness of the common scale}
\label{subsec:v14-completeness}
We check that a single choice of witnesses meets all the new
tests. Start with the original simultaneous construction for
an integer input, with the same positive progression and the
same contact-prime restrictions. Use its depth19 preparation.
The resulting primary formal depths are at least20, the high
tests have valuation at least40, and the middle tests have
valuation at least12. The fixed H/V0 exclusions remain part of
this construction.

At each contact prime $\alpha=v(R)=1$. Prescribe $v(\gamma)=1$.
In the three groups the high-term valuations are at least1
or2 and the middle-term valuations at least9 or10. The primary
terms have lower bounds159,3,83,76 in root order; the guard has
lower bound8. The alternate low terms are unchanged and have
positive valuation. Every group is therefore one plus an
element of the maximal ideal at the successful place.
At the other place above a split contact prime, the inverse
of each group has residue zero or one. Thus the trace of the
product of inverses has residue one or two and is a unit.
The pole value has positive odd valuation one.

At each preparation prime $p=2,3,17$, put
$\eta_j=v_p(A_j)$, prescribe $v_p(\gamma)=v_p(\lambda)=-\nu_p$,
and enlarge $\nu_p$ before choosing the contact primes.
Since $v_p(R)=5(\eta_j-\nu_p)$, the coefficients multiplying
the bounded test values have valuations
\[
\begin{array}{c|c}
 \gamma^{-160}R&155\nu_p+5\eta_j\\
 \gamma^{-160}R^2&150\nu_p+10\eta_j\\
 \gamma^{-160}R^3&145\nu_p+15\eta_j\\
 \gamma^{-40}R&35\nu_p+5\eta_j\\
 \gamma^{-40}R^2&30\nu_p+10\eta_j\\
 \gamma^{-80}R^3&65\nu_p+15\eta_j\\
 \gamma^{-80}E^4&100\nu_p-20\eta_j\\
 \gamma^{16}E^8&24\nu_p-40\eta_j .
\end{array}
\]
Each tends to $+\infty$ with $\nu_p$. These are all the new
coefficient forms. The previous finite preparation therefore
works after increasing its bounds.
At the real places impose $\lambda<\gamma<2\lambda$.
The corresponding powers of $\lambda$ are
$-155,-150,-145,-35,-30,-65,-100,-24$.
They all tend to zero uniformly in this interval. The groups
tend to one, their trace to two, and the pole value tends to
$1/(4R)$ through positive values. This gives the real conditions
required by the unchanged local-global pole certificates.

Away from the contact support, $R$ has even valuation. Since
the other factors in its pole value are squares, new primes in
the denominator of $\gamma$ create no odd pole support.
This also explains why $\gamma$ may be chosen on the anchor
quadric below rather than independently. The argument requires
its specified local valuations and real interval, not an
absence of other denominator primes.

The choices have a fixed order. First prepare the scalar and
root witnesses and all finite and real bounds. Choose the
contact primes and then $\lambda$. Choose the anchor point
and hence $\gamma$ with the required local conditions.
Finally choose H in the original progression, deep enough at
the finitely many selected places and in the retained real
neighborhoods. The bounds above are uniform in the anchor
choice. Thus this last H choice meets all three groups
simultaneously.

\subsection{Coupling the scale to the quaternion anchor}
\label{subsec:v14-anchor}
Let $T$ be the shifted scalar anchor already constructed.
Writing the old quaternion coordinates as $u,v,w$, put
$z=w+C$. Its equation is
\begin{equation}\label{eq:v14-anchor-quadric}
 u^2+6v^2-2z^2=N,\qquad N=3T^2-2C^2.
\end{equation}
Here $N\ne0$, since $N=0$ would make the rational square
$(T/C)^2$ equal $2/3$.
Use $u$ to define $\gamma=1/u$. The independent common-scale
chain is replaced by the eight coordinates specified by
\[
 a_2=u^2,\quad a_4=a_2^2,\quad a_8=a_4^2,\quad
 a_{16}=a_8^2,\quad a_{16}J=1,\quad
 g_{40}a_8=J^2,\quad g_{80}=g_{40}^2,\quad
 g_{160}=g_{80}^2.
\]
These equations force $u\ne0$, $J=\gamma^{16}$ and
$g_j=\gamma^j$ for $j=40,80,160$. They recover the former
common-scale coordinates uniquely and prove soundness of
this replacement.

We prove that the extra choice of $u$ does not obstruct
completeness. In the original scalar preparation require also
the allowed condition $F_s\in8\mathbb Z_2$. The same positive
progression and further formal depths remain available.
For the three odd roots $v_2(X_i)=-1$, and
\[
 q=3(2X_A)^2+3(2X_T)^2+27(2X_Z)^2\equiv1\pmod8.
\]
The shifted anchor therefore has $t=T/C$ odd at2 and integral
at3. The retained quaternion criterion supplies one rational
point $P$ on~\eqref{eq:v14-anchor-quadric}.
Fix it before choosing the contact primes.

The required local sets on this quadric are nonempty. Choose
$\nu_2,\nu_{17}\geq1$ and
$\nu_3>v_3(C)=2^{25}$, enlarged to meet all earlier preparation
bounds. At2 take
\[
 u=2^{\nu_2},\qquad v=T/2,\qquad
 z=(C/2)\sqrt{4-3t^2+2(u/C)^2}.
\]
The radicand is one modulo8. At3 take
\[
 u=3^{\nu_3},\quad v=0,\quad
 z=C\sqrt{1-3t^2/2+(u/C)^2/2};
\]
the radicand is one modulo3. At17 take
\[
 u=17^{\nu_{17}},\quad v=T/\sqrt2,\quad
 z=C\sqrt{1+u^2/(2C^2)}.
\]
Here2 is a square in $\mathbb Q_{17}$ and the last radicand
is one modulo17. Each construction solves the quadric exactly.

At a contact prime $p>76176$ exclude the finitely many
denominator primes of $T$, and prescribe $u=p^{-1}$.
Setting $V=pv$, $Z=pz$ leaves
$6V^2-2Z^2=p^2N-1$.
Modulo $p$ the right side is $-1$. The binary form
$Z^2-3V^2$ represents every nonzero residue: it is a split
product if3 is a square and a surjective finite-field norm
otherwise. The resulting solution is nonsingular, so Hensel
lifting gives the desired local point. At the real place take
$1/(2\lambda)<u<1/\lambda$, choose $v$ sufficiently large,
and solve for a positive real $z$. Hence $\lambda<\gamma<2\lambda$.

For clarity, weak approximation on this particular quadric
can be proved without adding quantified witnesses.
Let $Q(u,v,z)=u^2+6v^2-2z^2$ and let $B$ be its associated
symmetric bilinear form, so $B(X,X)=Q(X)$.
For a direction $d$ with $Q(d)\ne0$ put
\[
 \Phi(d)=P-\frac{2B(P,d)}{Q(d)}d .
\]
Then $Q(\Phi(d))=N$. Every point $X$ on the quadric with
$Q(X-P)\ne0$ is obtained by taking $d=X-P$.
This image meets every nonempty local open set. Indeed the
excluded tangent section has equation $B(P,X)=N$.
If $P$ and $X$ are not proportional, choose
$h\perp X$ with $B(P,h)\ne0$ and use the local curve
\[
 X(s)=\sqrt{1-s^2Q(h)/N}\,X+sh .
\]
It stays on the quadric and leaves that section to first
order. If $P$ is proportional to $X$ on the section, then
$P=X$. The tangent plane is nondegenerate because $N\ne0$;
choose $h$ there with $Q(h)\ne0$. The same curve leaves the
section to second order. The square root exists for
sufficiently small $s$, also over $\mathbb Q_2$.

Finally approximate the three direction coordinates
simultaneously by rationals. After clearing the local
denominators with an integer $D_0$, impose the required finite
congruences modulo an integer $M$. For integers
$L\equiv1\pmod M$ tending to infinity, coordinates of the form
$(A+Mk)/(D_0L)$ retain those congruences and have real mesh
$M/(D_0L)$ tending to zero. The Chinese remainder theorem
chooses $A$ for each coordinate. Continuity of $\Phi$ on the
chosen local neighborhoods now supplies a rational point in
all the required open sets. Its $u$ gives precisely the
valuations and real interval required for $\gamma$.
This proves simultaneous completeness of the coupled anchor.

\subsection{Exact private changes and the finite polynomial}
\label{subsec:v14-finite}
We relate the preceding arithmetic argument to the literal
quadratic presentation. Coordinates represent elements of
$K$ by their rational components; each displayed field
equation contributes its two component equations.
All rational coefficients are cleared and each equation is
divided by its integer content before taking the sum of
squares. These operations do not change its zero set.

For the reciprocal replacement, an old root chart with
scaled private coordinates $V_{\rm old}$, $I_{\rm old}$,
$Q_{\rm old}$, $K_{\rm old}$ has $c=a_i^2\ne0$.
Its actual reciprocal coordinates are
\[
 f=V_{\rm old}^2/c,\quad
 h=-I_{\rm old}V_{\rm old}/c,\quad
 F=Q_{\rm old}^2/c,\quad
 k=-K_{\rm old}Q_{\rm old}/c.
\]
The old identities $I_{\rm old}^2=c(1-DY)$ and
$K_{\rm old}^2=c(1-DB)$, together with
$I_{\rm old}L=V_{\rm old}N_0$ and
$K_{\rm old}z=-Q_{\rm old}$, give the new curve and both
regular ratio equations, including the identity cases.
The actual equality $U-V=Q_i$ supplies a rational secant
slope or its rational tangent limit. For the scalar pair it
also supplies~\eqref{eq:v14-S-orientation}.
No new local neighborhood, root sign, anchor choice or H
choice is needed for this change.
Conversely~\eqref{eq:v14-root-identity} restores the deleted
root ordinates and the preceding soundness proof applies to
the enlarged auxiliary domain. We assert equality of input
projections, not a bijection to all the old private
square-root fibers.

The subsequent deletion of each alternate ordinate is an
exact change on the full zero set. Its inverse is
$k=s(f-F)-h$, and~\eqref{eq:v14-alt-restoration} and the
globally justified~\eqref{eq:v14-raw-ratio} recover every
deleted equation. This accounts for six rational coordinates
in the A, T and Z pairs. The scalar pair saves two further
coordinates while retaining its root ordinate and adding
the orientation equation.

The remaining private replacements are exact power chains
or polynomial changes. For a local inverse $ER=1$, a shorter
chain may specify
\[
\begin{gathered}
 EJ=R^2,\quad E_4J=E,\quad E_8=E_4^2,\quad E_{19}J=E_8^2,\\
 RE_{39}=E_{19}^2,\quad RE_{79}=E_{39}^2,\quad
 E_{158}=E_{79}^2,\quad E_{157}=RE_{158}.
\end{gathered}
\]
It forces $J=R^3$ and exactly the indicated powers of $E$.
There is no zero or sign branch because $ER=1$.
The final common-scale presentation needs only the local
coordinates for $E,R^2,R^3,E^4,E^8$, with shared powers
specified above. Its equivalence as an arithmetic
certificate was proved in the common-place and completeness
arguments, rather than inferred from coordinate counting.

The sparse Lucas change replaces its private coordinate
$C_{\rm old}=AZ$ by $C_{\rm new}=C_{\rm old}-2Z$.
The inverse is $C_{\rm old}=C_{\rm new}+2Z$, and the defining
polynomial becomes
\[
 F+1=T^2+C_{\rm new}^2-C_{\rm new}(AT).
\]
Both substitutions are affine in the private coordinates.
The cached product $AT$ makes the displayed equation
quadratic. Every external Lucas value is unchanged.
This change reduces terms, not the number of witnesses.

For reference the successive total-variable counts of the
selected construction are
\[
394,\ 392,\ 390,\ 372,\ 369,\ 366,\ 326,\ 319,\
318,\ 315,\ 314,\ 313,\ 307,\ 301,\ 299 .
\]
The early exact eliminations, affine H chart and squared-H
replacement are justified above. The count326 replaces
separate pole blocks by the three groups; the final group
argument proves their required arithmetic properties
directly. The count315 precedes common scaling.
Common scaling removes twelve local coordinates, adds nine
shared coordinates and two scaled test coordinates, giving314.
The coupled quaternion anchor replaces those nine by eight,
giving313. Reciprocal A, T and Z root charts give307,
alternate-ordinate elimination gives301, and the oriented
scalar auxiliary chart gives299.
Intermediate alternatives with different depth requirements
are not needed for this selected proof.

The literal final system has279 quadratic component
equations in299 active rational variables, including the
input. Its polynomial is their unweighted sum of squares.
Over $\mathbb Q$ it vanishes exactly when every component
equation vanishes. Exact expansion, collection by exponent
vector and removal of zero coefficients give6,630 monomials,
total degree four, content one, and maximum absolute
coefficient $186624C^2=186624\cdot3^{2^{26}}$.
Every coordinate is active. The coefficient artifacts specify
the exponent-vector order and the integer coefficient
representation, so these finite assertions refer to a
single polynomial, not just to its zero set.

\subsection{The ordinary coefficient presentation}
\label{subsec:v14-ordinary}
The large fixed coefficient can be represented with fourteen
constant coordinates. Put
\[
 (a_0,\ldots,a_{13})=(2049,
 \underbrace{2048,\ldots,2048}_{11\ {\rm entries}},0,0)
\]
and impose
$c_0=3^{a_0}$ and $c_j=3^{a_j}c_{j-1}^2$ for $1\leq j\leq13$.
If $c_j=3^{e_j}$ then
$e_0=2049$ and $e_j=2e_{j-1}+a_j$.
After the eleven2048 steps $e_{11}=2^{23}$, and the two
last squares give $e_{13}=2^{25}$. Thus $c_{13}=C$.
Each coordinate is uniquely determined; no square-root sign
is introduced.

Add one more coordinate
\[
 Q=48v_{83}^2+48v_{87}^2+432v_{89}^2.
\]
Only two rows of the symbolic system contain $C$. Replace
them by
\[
\begin{split}
 v_{55}^2+6v_{56}^2-2v_{57}^2
       -4c_{13}v_{57}-3v_{903}^2&=0,\\
 v_{903}-2v_{543}v_{544}-c_{13}Q&=0 .
\end{split}
\]
Both are quadratic. Substitution of the unique constants
and $Q$ gives the exact symbolic rows. Projection and unique
extension therefore give inverse maps on full rational
solution sets and fix the input.
There are $299+14+1=314$ variables and $279+14+1=294$
quadratic equations.

Apply the same primitive integer normalization and
unweighted sum of squares. The ordinary polynomial has
6,672 nonzero monomials and content one.
Let $K_0$ denote the constant coefficient of the symbolic
sum of squares, an explicitly tabulated positive integer
of1,514 bits. The exact largest absolute coefficient of
the ordinary expansion is
\[
 3^{4098}+K_0,\qquad
 2^{6495}\leq3^{4098}+K_0<2^{6496}.
\]
This proves the asserted ordinary coefficient bound.
The source-coordinate table retains the original coordinate
identifiers and has222,890 bytes. A compact table renames
the active variables consecutively by an explicit bijection
fixing the input; its byte count is a separate artifact
property. That renaming changes neither the equations nor
the coefficient statistics. The proof and the coefficient
tables above establish Theorem~\ref{thm:v14-quartic}.

\section{Unary theta series with imprimitive sieves}\label{v15:unary-poisson}
We give the unary-series proof used in the theta construction. Poisson summation is the Schwartz--Gaussian specialization of Weil, I.19, Theorem 4 and I.20 \cite{v15-weil-poisson}; the multiplier convention is Ramsey equation (1.1), p.~600 \cite{ramsey}. Let $S$ contain $2$, the conductors of $\eta$ and the old characters, and all late sieve primes, including $\ell_0$. Choose $i$ so that $\lambda_0=\eta\chi_{-4}^i$ is odd. Choose $C$ divisible by its conductor and all primes in $S$, with no prime factors outside $S$, and induce $\lambda_0$ to a character $\lambda$ modulo $C$, extended by zero off the units. It need not be primitive. Put $e(t)=\exp(2\pi i t)$ and
\[F_\lambda(z)=\sum_{m>0}\lambda(m)m e(m^2z).\]
Then, in the theta-multiplier convention used above,
\[F_\lambda\in S_{3/2}(\Gamma_0(4C^2),\lambda\chi_{-4}).\]
For squarefree positive $J$ prime to $S$, dilation gives
\[F_\lambda(Jz)\in S_{3/2}(\Gamma_0(4C^2J),\lambda\chi_{-4}(J/\cdot)).\]
This is an available level, not a minimal-level assertion. In particular each prime of $J$ occurs to exponent one, and no fresh prime outside $J$ occurs.

Write $\Theta(z,w)=\sum_{n\in\Z}e(n^2z+nw)$ and, for $\gamma=\left(\begin{smallmatrix}a&b\\c&d\end{smallmatrix}\right)\in\Gamma_0(4)$, put $j(\gamma,z)=\epsilon_d^{-1}(c/d)(cz+d)^{1/2}$, with the fixed theta branch. Then $j^2=\chi_{-4}(d)(cz+d)$. Poisson summation for a Gaussian gives
\[\vartheta(-1/\tau,u/\tau)=(-i\tau)^{1/2}\exp(\pi i u^2/\tau)\vartheta(\tau,u),\]
where $\vartheta(\tau,u)=\sum_n\exp(\pi i n^2\tau+2\pi i nu)$, and direct substitution gives invariance under $\tau\mapsto\tau+2$. With $\tau=2z$, the associated matrix is $\left(\begin{smallmatrix}a&2b\\c/2&d\end{smallmatrix}\right)\in\Gamma(2)$. The Euclidean algorithm with even row additions gives the generators $T^2,ST^2S^{-1},-I$. Composition of the two identities, with the scalar fixed at $w=0$, yields
\begin{equation}\label{v15:unary-jacobi}
\Theta(\gamma z,w/(cz+d))=j(\gamma,z)\exp\!\left(\frac{\pi i c w^2}{2(cz+d)}\right)\Theta(z,w).
\end{equation}
For $\gamma\in\Gamma_0(4C^2)$ set $D=cz+d$. Differentiate in $w$ and then set $w=Dv$:
\[\Theta_w(\gamma z,v)=D j(\gamma,z)e^{\pi i cDv^2/2}\bigl(\Theta_w(z,Dv)+\pi i cv\Theta(z,Dv)\bigr).\]
Directly shifting the summation variable proves
\begin{align*}
\Theta(z,w+2mz)&=e(-m^2z-mw)\Theta(z,w),\\
\Theta_w(z,w+2mz)&=e(-m^2z-mw)\bigl(\Theta_w(z,w)-2\pi im\Theta(z,w)\bigr).
\end{align*}
Take $v=r/C$ and $m=cr/(2C)\in\Z$. Since $Dv=2mz+dr/C$, the derivative corrections cancel. The remaining exponential is
\[e^{\pi i cDr^2/(2C^2)}e(-m^2z-mdr/C)=e(-cdr^2/(4C^2))=1,\]
including the full dyadic divisor $4C^2\mid c$. Hence
\[\Theta_w(\gamma z,r/C)=D j(\gamma,z)\Theta_w(z,dr/C).\]
For $\widehat\lambda(r)=C^{-1}\sum_{a\bmod C}\lambda(a)e(-ar/C)$, Fourier inversion and oddness give
\[F_\lambda(z)=(4\pi i)^{-1}\sum_{r\bmod C}\widehat\lambda(r)\Theta_w(z,r/C).\]
The transformed coefficient is $\lambda(dn)=\lambda(d)\lambda(n)$, because $d$ is a unit modulo $C$. Thus
\[F_\lambda(\gamma z)=\lambda(d)D j(\gamma,z)F_\lambda(z)=\lambda(d)\chi_{-4}(d)j(\gamma,z)^3F_\lambda(z).\]
This gives the character, not its inverse, also for imprimitive $\lambda$. The finite Fourier transform is only a complex-analytic proof device: the original coefficients remain the integral algebraic numbers $\lambda(m)m$.

For cusp vanishing take any $\gamma\in\mathrm{SL}_2(\Z)$ with $c>0$. Put $L=cC$, $A=2i/(cD)$, so $\Re A>0$. With $F_{\rm full}=2F_\lambda$,
\[F_{\rm full}(\gamma z)=\sum_{r\bmod L}\lambda(r)e(ar^2/c)\sum_{n\equiv r\ (L)}n e^{-\pi An^2}.\]
The outer coefficient is $L$-periodic. The Fourier transform of $x e^{-\pi Ax^2}$ is $-i\xi A^{-3/2}e^{-\pi\xi^2/A}$. Poisson summation on each coset therefore gives, with a fixed nonzero scalar $\beta$,
\[D^{-3/2}F_{\rm full}(\gamma z)=\beta\sum_{k\in\Z}kS(k)e(cd k^2/(4L^2))e(c^2k^2z/(4L^2)),\]
where $S(k)=\sum_{r\bmod L}\lambda(r)e(ar^2/c+kr/L)$ is bounded. The exponents are $k^2/(4C^2)$; the zero coefficient vanishes. Normal convergence and decay prove cusp vanishing. Infinity follows from the original series; $c<0$ follows on replacing $\gamma$ by $-\gamma$.

For dilation, use $\gamma_J=\left(\begin{smallmatrix}a&bJ\\c/J&d\end{smallmatrix}\right)$ and $J\gamma z=\gamma_J(Jz)$. The ratio $j(\gamma_J,Jz)/j(\gamma,z)=(J/d)$ proves the stated character and level. Dilation sends rational cusps to rational cusps. Multiplication by ordinary theta gives weight two since $j^4=(cz+d)^2$. The support of $C$ also proves the fresh-prime omissions required by the application.

\section{Coefficient conjugation and finite theta twists}\label{v15:theta-transport}
Let $4\mid L$, let $g$ be an algebraic-coefficient cusp form of weight $3/2$ and character $\psi$ at level $L$, and let $\tau\in\operatorname{Aut}(\overline\Q)$. In Ramsey's geometric theta convention (Definition 3.1 \cite{ramsey}), $g/\theta^3$ is a section on $X_1(L)$ with its prescribed divisor and cusp conditions. These conditions commute with field extension and conjugation; its diamond character becomes $\tau\psi$. The chosen Tate cusp may not be rational. If $\tau(\zeta_L)=\zeta_L^a$, its conjugate marking differs by diamond $a$, so the two expansions differ by a nonzero constant. Rescaling by this constant gives exactly the coefficientwise conjugate at the fixed marking, where theta has rational coefficients. This preserves all cusp conditions. Only flat field extension is used, not a weight-at-least-five vanishing theorem.

Ramsey's Theorem 6.2 \cite{ramsey}, at weight $3/2$, gives the good-prime Hecke formula
\[a(p^2n)+\psi(p)(-n/p)a(n)+p\psi(p^2)a(n/p^2),\]
with a nonintegral-index coefficient interpreted as zero. Its non-character factors are rational. Conjugation therefore sends eigenvalues $\eta(p)a_p(E)$ to $\tau(\eta(p))a_p(E)$.

For a primitive character $\nu$ of conductor $m$, set
\[h(z)=G(\overline\nu)^{-1}\sum_{t\bmod m}\overline\nu(t)g(z+t/m).\]
The primitive Gauss identity, including its zero at indices not prime to $m$, gives $h=\sum_n\nu(n)a(n)q^n$; for $m=1$ this is $g$. For $\gamma\in\Gamma_0(Lm^2)$ choose $t'\equiv d^2t\pmod m$. Matrix multiplication gives
\[n(t/m)\gamma=\gamma_t n(t'/m),\qquad
\gamma_t=\begin{pmatrix}a+tc/m&b+(td-at')/m-ctt'/m^2\\c&d-ct'/m\end{pmatrix}\in\Gamma_0(L).
\]
Put $d'=d-ct'/m$. The theta multiplier is unchanged: $d'\equiv d\pmod L$ and modulo four, and $(c/d')=(c/d)$. Indeed for $c\ne0$ write $c=sr^2$ with signed squarefree $s$. The relevant quadratic character has conductor dividing $4|s|$, which divides $c/m$: at odd primes in $s$, $v_p(c)-v_p(m)\ge1$; at two, $v_2(c)\ge2+2v_2(m)$ gives the required exponent two, or three when $v_2(c)$ is odd. The signed Kronecker convention handles negative $c,d$. For $c=0$, periodicity and the original $-I$ law give the same assertion. The square-root factors agree because $c(z+t'/m)+d'=cz+d$. Reindexing multiplies the weights by $\nu(d)^2$, so $h$ has character $\psi\nu^2$.

Rational translations preserve cusp vanishing: after cusp scaling their residual upper-triangular matrix rescales the imaginary part by a positive rational number. Choose a common cusp width for the finitely many translations. Thus $h$ is cuspidal at level supported on $Lm$. For an imprimitive presentation in the application, use the inducing primitive character; the prescribed unit support of $g$ makes the coefficient product identical, and the conductor is supported on the old set $S_0$. No new fresh-prime factor is introduced.

Direct substitution in the Hecke formula gives $T(p^2)h=\nu(p)^2\lambda_p h$ for $p\nmid Lm$. For $\nu=\eta^{(e-1)/2}$, $e$ odd, and $\psi/\eta$ quadratic, one has $\psi\nu^2=\psi^e$ and $\nu(p)^2\eta(p)=\eta(p)^e$. Apply the full good-Hecke-eigenspace proportionality already proved in the theta section, at a common level supported on $S_0$. With its squareclass conditions, $D_0$ prime to $S_0$ and $a(D_0)=1$, the reference denominator is the nonzero $\nu(D_0)$; consequently
\[\tau(a(D))=\nu(D/D_0)a(D).\]
The full-eigenspace proportionality and unit support remain explicit inputs from the earlier theta construction. No rationality of an $L$-value quotient is assumed.

\part{The polynomial inverse theorem and its local joins}\label{v15:gi-full}
The following internal modules prove the polynomial inverse theorem used above, including the degree--rank induction and its current selection joins. All finite data and standardness conventions are retained. Module names below identify internal sections, not external theorem assumptions.
\section{Assembly of the precise inverse-theorem input}\label{v15:gi:ASSEMBLY}
\newcommand{\GIAZ}{\mathbb Z}
\newcommand{\GIAN}{\mathbb N}
\newcommand{\GIAC}{\mathbb C}
\newcommand{\GIAR}{\mathbb R}
\newcommand{\GIAE}{\mathbb E}
\newcommand{\GIAGI}{\mathrm{GI}}



The following sections assemble the inverse theorem from the internal local results listed below and their stated published inputs.

\subsection{Exact target and conventions}

The target, denoted $\GIAGI(k)$, is the following unchanged linear formulation.
For each fixed integer $k\geq3$ and each $0<\delta\leq1$, there are a finite
family $\mathcal M_{k,\delta}$ of nilmanifolds $G/\Gamma$, each with connected
simply connected nilpotent group of step at most $k$ and a discrete cocompact
lattice, fixed smooth metrics, and constants $C,c>0$, such that every
$1$-bounded $f:[N]\to\GIAC$ with $\|f\|_{U^{k+1}[N]}\geq\delta$ correlates with
a $1$-bounded test $F$ of Lipschitz constant at most $C$ along a linear orbit:
\[
 \left|\GIAE_{n\in[N]}f(n)F(g^nx)\right|\geq c.
\]
The family and constants depend on $k,\delta$, not on $N,f$ or the orbit
parameters. Strict source thresholds can be used with $\delta/2$ to give
this weak-inequality version. No effective numerical bounds are asserted.

For $d\geq2$, the interval norm is the norm of the zero extension to a
cyclic group of sufficiently large order divided by the norm of the
interval indicator, as in the published GTZ definition on p.~1232.
We use the same normalization at every inverse step. On a finite abelian
group $A$, write
\[
 \|v\|_{U^d(A)}^{2^d}
 =\GIAE_{x,t_1,\ldots,t_d\in A}
   \prod_{\omega\in\{0,1\}^d}\mathcal C^{|\omega|}
     v(x+\omega\cdot t),
\]
where $\mathcal C$ is complex conjugation. Our correlation convention is
$|\GIAE f\psi|$; conjugating the test function converts the other common
convention without changing its degree or bounds.

All limit nilcharacters in the assembly are internal in the explicitly
specified sense of the local Lemma 13.2 and tautification units: the
vertical identity holds for all limit group elements and all limit points,
with the natural extension of a standard vertical character. Every
constructor used below must supply that domain, rather than only
standard-point equivariance. This is an explicit application check, not a
consequence of bounded Lipschitz continuity alone.

The manuscript currently requires $\GIAGI(20)$ and $\GIAGI(39)$. The proposed
relative inverse replacement also requires $\GIAGI(r-1)$ for each fixed
$r=\max(s,4)$ chosen before the norm threshold and asymptotic variable.
A polynomial orbit without a same-step linearization or an
$N$-dependent family does not satisfy this target.

\subsection{Intermediate interfaces that must be supplied}

The intermediate conclusions and their hypotheses are listed below. Each is proved in the indicated internal section or supplied by the explicitly cited published input.

\begin{enumerate}
\item The compatible bounded polynomial inverse theorem of order $2$,
proved in Section~\ref{v15:gi:GI-BASE2} relative to
its retained published forward bracket theorem and Bohr-set inputs.
It gives the fixed finite family and normalized interval $U^3$ norm.
The needed trivial-subgroup specialization of the false quadratic
extension is proved locally, and the periodic connected polynomial
construction replaces the affected skew-iterate route.
\item The nilcharacter approximation/selection interface: from positive
correlation with a bounded degree-$d$ nilsequence, obtain a degree-$d$
unit-vector nilcharacter and a scalar component giving positive correlation.
For a dense family of parameters, make the choices in a limit fashion and,
after passage to a dense subfamily if necessary, keep the finite dimensions,
complexities, bounds and positive correlation standard and uniform.
Section 4 below gives the finite-net and unit-vector construction used
at this interface using the previous finite inverse family. Later selection applications use Section~\ref{v15:gi-selection}.
\item Linearization at inverse order $s$: for a dense family of degree
$s-1$ nilcharacters $\chi_h$ correlated with $\Delta_hf$, obtain one
multidegree $(1,s-1)$ nilcharacter $\chi(h,n)$ with which $\Delta_h f$
has degree-at-most-$s-2$ correlation for many $h$.
This is the precise Theorem 7.1 interface. The modified Proposition 7.3
is supplied by \texttt{local-proof-sol-lp02-03-r3}; the quotient and
rank-lower comparison by \texttt{lp04-001}; the final orbit by the
Section~\ref{v15:gi:GI-SECTION12}; and the degree conversions, unit-vector
recovery, finite rank iteration and rank-zero endpoint by
\texttt{rank-iteration-001}. The frozen Sections 9--11 reconstruction
proves the one-step frequency premise at
every newly constructed rank family, conditional on its precise
one-variable D and F contracts. Its theorem takes the same uniform
unit-vector degree--rank family and accumulated multidegree factor as
the rank iteration. It reconstructs all core/linear/independent data
anew; none is inherited. Its fixed unit padding is removed by projecting
the first factor of the entire LP04 matrix, so the matrix is exactly
$\chi_h\otimes\overline{\zeta_h}$ required by the next rank step.
Section~\ref{v15:gi:GI-SECTION12} supplies the fixed coordinate isometry, with the zero-index groups explicit. All these
steps retain the full-limit vertical identities and uniform bounds.


For clarity, the invariant concerns represented families: ordinary
unit-vector tests with exact global vertical identities, fixed standard
bounds and finitely many standard models and characters. The initial
positive-degree selection is made among these representations, by the
parameter-level selection argument for the preceding inverse order;
it does not upgrade an arbitrary weak nilcharacter or preserve a
prescribed symbol. Degree-to-rank conversion retains the represented
function and restricts its character, so this invariant starts at
$(m,m)$, where $m=s-1$, with accumulated factor $A=1$.

Distinguish the fixed inclusion $J$ that pads the incoming $\chi_h$
in Sections 9--11 from the final Section 12 inclusion $I$, for which
$Z(h,n)=I\zeta_h(n)$ on the retained shift set. The full ordinary-top
character in Section 12 is chosen before the LP04 comparison and
restricts to its cancelling rank-top character. The comparison thus
supplies the entire vector
\[
 W_h=(J\chi_h)\otimes\overline{\zeta_h}
\]
on one common model of degree-rank at most $(m,R-1)$. If $QJ=\mathrm{id}$
is the coordinate projection on the first factor, then
\[
 (Q\otimes\mathrm{id})W_h
   =\chi_h\otimes\overline{\zeta_h}=:V_h.
\]
This fixed bounded linear map preserves the common model, internal
parameter dependence and standard Lipschitz bound. The projection
recovering $\zeta_h$ from $Z$ is instead the one associated with $I$.
In particular, this conclusion concerns every matrix entry, not a
separately selected scalar entry.

The identity
$\chi_{h,a}=\sum_b\zeta_{h,b}(\chi_{h,a}\overline{\zeta_{h,b}})$
uses the unit norm of the whole vector $\zeta_h$ and therefore uses
exactly $V_h$ in the correlation recovery. The finite ordinary restart
dictionary preserves the global unit and vertical identities after
internal selection; the ordinary tests chosen in the proof of E.5 are
extended only after their finite selection. Tensoring the accumulated
factor with $Z$ preserves its unit norm and multidegree $(1,m)$.
Sections 9--11 then reconstruct all frequency data for this new family.
No ordinary degree drop is claimed before the rank endpoint: at rank
zero the degree-rank axiom makes the remaining factor ordinary degree
at most $m-1$, which is then absorbed into the lower-degree test.
The geometric, analytic and source premises of all these component
constructions remain those of their stated contracts.

The centered fractional-part convention in this chain is
$(-1/2,1/2]$, a fixed fundamental cube permitted by Section 12.
The reconstruction uses the affine-offset coefficient transfer supplied by the local multifactor lemma, as connected below.
\item Symmetrization at inverse order $s$: from that input obtain one
degree-$s$ nilcharacter $\Theta$ and a bounded nilsequence $\Psi(h,n)$
supported on
\begin{equation}\label{v15:gi:0:eq:J}
 J=\{(i,j):i+j\leq s-1\}\ \cup\ \{(i,s-i):2\leq i\leq s\},
\end{equation}
such that the relevant $\chi(h,n)$ is a bounded finite linear combination
of components of
$\Theta(n+h)\otimes\overline{\Theta(n)}\otimes\Psi(h,n)$.
This is the required Theorem 7.4 output, supplied by the locally authored
Section~\ref{v15:gi:GI-SECTION13}. It takes the
degree-at-most-$s-2$ correlation actually supplied above; the printed
strict hypothesis is not used. Its proof binds the corrected
multilinearization, literal block symmetry, local Lemma 13.2 and later
shift/sign calculations. It also retains the precise stronger
Proposition 8.4, E.11 and E.13 applications. The precise plain
Proposition 8.4 input is supplied by the accepted local
\texttt{plain-quadruples-001}: select the same scalar coordinate of
$\chi(h,\cdot)\otimes\phi_h$ and retain that entire lower factor.
For $H\subset[-N,N]$ of density $\alpha$, correlation $\delta$ and
global scalar bound $B$, that proof gives standard positive density
and correlation depending only on $\alpha,\delta,B$. Its output has
signs $++--$ and shifts in the second and fourth factors by
$h_1-h_4$, exactly as used in Section 13. The functions have global
polynomial representations; the interval cutoff belongs to $f$.
The proof removes its remaining bounded cutoff by an inequality,
not by an asserted identity. Finite component selection and the
internal version preserve standard uniform constants.
This is an explicit substitution of the supplying proof for the
published Proposition 8.4 invocation in the frozen Section 13 proof,
lines 88--125. That earlier file remains unchanged and does not itself
refer to the later local proof. Its intervals $I_1=[-N,N]$ and $J=[N]$,
scalar component, global test domain, and quadruple count on the
$N^3$ scale are exactly the local theorem's data. The reparametrization
$(h_1,h_2,h_3,h_4)=(h_0+a,h_0+b,h_0+a+b,h_0)$ is therefore unchanged;
the subsequent pigeonhole has $h_0\in I_1$ and $a,b\in[-2N,2N]$.
The E.11 and E.13 inputs require the
actual multivariable Appendix D.5 factorization and all-box reconstruction;
a one-variable D.4 proof does not supply that premise. This obligation
is supplied by the local multifactor theorem, as connected below.
\item Splitting for a union of two finite downsets: a fixed bounded
nilsequence on their union is uniformly approximable, to any fixed positive
error, by a finite sum of products of nilsequences supported on the two
downsets, with fixed finite complexity and bounds. This is the E.4
application used explicitly in the next section.
\item The converse correlation bound: a bounded scalar function with
positive correlation against a bounded degree-$d$ nilsequence has positive
$U^{d+1}[N]$ norm, with a bound for each fixed complexity and correlation.
The local \texttt{converse-001} proves the required implication using
the published compact cube constraint, an explicitly normalized smooth
average, a uniform Fourier bound and Gowers--Cauchy--Schwarz. It uses
accepted exact same-step linearization for polynomial tests. It does not
import the disputed averaging-domain transition in Green--Tao 11.2 or
invoke the inverse theorem being proved. Its exact source application
is the one stated in Section~\ref{v15:gi:GI-CONVERSE}.
\item Countability, up to filtered isomorphism, of the fixed rational
filtered nilmanifolds of bounded degree, and equivalence of fixed smooth
metrics on a compact nilmanifold. The published countability assertion is
on p.~1250 with its Mal'cev-coordinate justification. Its exact scope must
include the taut rational objects constructed by this project.
\item The local increment linearization theorem: an exact polynomial
nilsequence of degree $s$ has a same-step linear representation on fixed
geometry with a fixed Lipschitz comparison. This is supplied by the
internal $P_1\rtimes\GIAR$ construction,
rather than by the defective full-$P_0$ suspension in Appendix C.
\end{enumerate}

The increment unit supplies the exact sequence and finite-family join
used here. It does not assert the stronger clause in the original
Proposition C.2 about one universal projection into the original,
unnormalized nilmanifold for every polynomial: bounded normalization is
absorbed in the test function $F_q$. That distinction changes no part of
the target in section 1.

The affine-transfer and multivariable-factorization connections are as follows.

First, the one-variable weak factorization uses a progression
$A+b[M]$ with $|A|=O(N)$, standard bounded nonzero integer $b$, and
$M\asymp N$. Apply the local inverse-affine lemma with one variable
to a real lift of the horizontal polynomial. Its input is the positive
binomial coefficient bound on $m\mapsto f(A+bm)$. In the inverse
substitution, the coefficient from input degree $j$ to positive output
degree $i$ has size $O(N^{j-i})$ and a standard denominator independent
of $A$. One standard integer multiplier clears the integer portions;
the residual $O(N^{-j})$ becomes $O(N^{-i})$. Thus the same original
horizontal character, multiplied by a standard positive integer,
has the required ambient coefficient bound. It remains nonzero.
No bound on the constant coefficient is asserted or needed. This is
the exact proof-level transfer at DF002 lines 132--145, distinct from
its right-coset periodicity argument.

Second, at Section 13 lines 327--343 take the local multifactor bias
theorem with dimension $s+1$, degree $s$, and its scale equal to $8N$.
The input is the actual reflected unit nilcharacter
$\sigma_\varepsilon$ and reflected convex polytope obtained by the
orthant selection. Every positive progression is the full $[8N]$;
the auxiliary lower-degree test is constant. Its conclusion supplies
a globally represented lower-degree nilsequence agreeing with
$\sigma_\varepsilon$ on all of $[-8N,8N]^{s+1}$.
The inverse reflection is an integer linear substitution, so the
same full-box conclusion holds for $\sigma$. This does not extend
an identity from only the selected positive orthant by assumption.

Third, after the displayed diagonal pullback in Section 13, the
actual unit tensor-product representative $\xi(h,n)$ of the
antisymmetric symbol satisfies $s![\xi]=0$ on $[-8N,8N]^2$.
By the definition of the zero symbol, the coordinates of
$\xi^{\otimes s!}$ are lower-degree nilsequences on that box;
their finite product of models is one common lower-degree model.
Apply the local torsion theorem in dimension two, degree $s$, scale
$8N$, and positive standard tensor power $s!$. It gives a globally
represented degree-at-most-$(s-1)$ nilsequence equal to $\xi$ on
that same symmetric box. Only afterward restrict to
$[-N,N]\times[N]$, as in Section 13 lines 345--363.

There is also a distinct converse use at Section 13 lines 298--305.
For $s\geq4$, put $d=s-3\geq1$ in the accepted local converse and
apply it to the uniformly selected lower-degree test correlated
with the chosen coordinate of $\alpha'$. This gives the positive
$U^{s-2}[N]$ bound there, uniformly after the stated selections.
For $s=3$, a degree-zero nilsequence is constant: a fixed positive
correlation with a bounded constant directly implies a fixed
positive $|\mathbb E_{n\in[N]}\alpha'(n)|$: if the correlation is
at least $c>0$ and the constant has modulus at most $M$, this mean
has modulus at least $c/M$. Zero extension has group mean
$(N/|Q|)\mathbb E_{[N]}\alpha'$, and the interval indicator has
$U^1$ quantity $N/|Q|$, so normalization gives exactly that interval
mean. No degree-one test is inserted into this case.

Closure of fixed-degree nilsequences under conjugation, scalar
multiplication and finite products uses the actual product nilmanifold and
product filtration. The top degree is the maximum of the input degree
bounds. This elementary object identification is also required when it is
used in the induction below.

\subsection{Correct splitting of the symmetrization remainder}

We first give the lattice detail needed by the splitting application.
It uses the integral Hall-coordinate lemma in the internal LP04 section, not a general assertion that disjoint rational Lie subgroups
automatically separate points after quotienting by a lattice.

\begin{lemma}[Splitting on two finite multidegree downsets]
Let $J_1,J_2\subset\GIAN^a$ be finite downsets. Every bounded Lipschitz
nilsequence of multidegree supported on $J_1\cup J_2$ is uniformly
approximable, to each fixed positive error, by a finite sum of products
of nilsequences supported on $J_1$ and $J_2$. The construction uses
fixed geometries and bounds before the polynomial coefficients or the
interval length vary, for each fixed input geometry and test bound.
\end{lemma}
\begin{proof}
Use bounded tautification to set the polynomial's value at zero to the
identity and restrict the ambient group to the group generated by its
positive coordinate groups. For each positive $j\in J_1\cup J_2$ choose
a finite generating set of $\Gamma\cap G_j$. Its real generated group
is $G_j$, by the rational subgroup assertion. Introduce one formal
letter of weight $j$ for each such generator. Take the free real Lie
algebra on these letters and discard every individual-letter multidegree
whose total vector weight is outside $J_1\cup J_2$. This is an ideal:
the complement of a downset is upward closed. Every letter has positive
total weight, so the surviving length is bounded by
$\max_{j\in J_1\cup J_2}|j|$.

Let $U$ be the exponential group and $\Lambda$ the group generated by
the exponential letters. The Hall-coordinate lemma proves that
$\Lambda$ is a lattice and that all weight filtration groups are
rational. Sending each letter to its chosen original lattice generator
gives a filtered Lie homomorphism $\phi:U\to G$ with
$\phi(\Lambda)\subseteq\Gamma$. It maps $U_j$ onto $G_j$: the letters
of weight exactly $j$ already generate $G_j$ as a real group. Lift each
Taylor coefficient of the normalized original polynomial through this
surjection. The multivariate Taylor theorem, B.9, makes their ordered
binomial product an adapted polynomial $\widetilde g$ with
$\phi\widetilde g=g$. Pulling the test back along
$U/\Lambda\to G/\Gamma$ therefore preserves the sequence exactly.
Its Lipschitz bound increases by a fixed geometric factor only.

For $\nu=1,2$ let $I_\nu$ be the connected ideal subgroup spanned by
the surviving Hall words whose weight is outside $J_\nu$. The quotient
$U/I_\nu$ has degree support in $J_\nu$. The integral Hall lemma says
that its quotient lattice is obtained by deleting precisely those
coordinates. No surviving Hall word is deleted by both quotients, since
its weight belongs to $J_1\cup J_2$. Consequently the map
\[
 U/\Lambda\longrightarrow
 U/(I_1\Lambda)\ \times\ U/(I_2\Lambda)
\]
is injective. In detail, if two cosets have equal images, their difference
has integral coordinates in each projected Hall system. The union of
the retained coordinate sets is the full Hall system, so every coordinate
of the difference is integral. The difference is in $\Lambda$.
This checks the lattice condition as well as the Lie algebra condition.

The pullbacks of the two algebras of Lipschitz functions on these compact
quotients contain the constants, are closed under conjugation, and
together separate points by that injection. The complex
Stone--Weierstrass theorem approximates the pulled-back test uniformly
by the algebra they generate. Each element of that algebra is a finite
sum of products with one factor from each quotient, since each of the
two individual algebras is closed under multiplication. Composing with
$\widetilde g$ gives the asserted splitting. For an internal bounded
Lipschitz test, first approximate by one standard Lipschitz test on the
fixed compact manifold, using Definition 5.1, and then perform this
standard finite approximation. For a fixed uniform Lipschitz and norm
bound, compactness of that class in the uniform norm supplies a finite
net and thus fixed finite bounds. Only this fixed-bound form is used
below. If the positive groups are trivial, the sequence is constant and
the assertion follows with one constant product.
\end{proof}

The coordinate order is $(h,n)$. Set
\[
 A=\{(0,j):0\leq j\leq s-1\},\qquad
 B=\{(i,j):j\leq s-2,\ i+j\leq s\}.
\]
Both are downsets and $A\cup B=J$. Indeed the only point of $J$ with
$j=s-1$ is $(0,s-1)$, covered by $A$. Every other point of $J$ has
$j\leq s-2$ and lies in $B$. Conversely a point of $B$ of total degree
$s$ has $i\geq2$, so belongs to the second part of \eqref{v15:gi:0:eq:J}; all its
other points lie in the first part. This proves equality in both directions.

Apply interface 5 to the actual $\Psi$ with these two downsets. A sequence
supported on $A$ is independent of $h$: its polynomial map has every
$h$-direction derivative in the trivial first-coordinate group.
It is a degree-at-most-$s-1$ function $\psi(n)$ of $n$ alone.
For a factor supported on $B$, fixing $h$ gives a polynomial nilsequence
in $n$ of degree at most $s-2$, since all higher $n$-direction groups are
trivial. If normalization of this section at fixed $h$ is needed, the
accepted bounded tautification theorem supplies it with uniform geometry
and parameter dependence.

Here is the needed approximation-to-correlation connection. After uniform
selection, suppose the preceding correlation is at least a standard
$\epsilon_0>0$ on a dense set of $h$, and the lower-degree test factors
have magnitude at most a standard $M$. Choose the uniform splitting error
less than $\epsilon_0/(2M)$, also absorbing the fixed finite component and
coefficient bounds into $M$. The error in each normalized correlation is
then at most $\epsilon_0/2$. A finite sum remains, so one summand has a
positive standard correlation on a dense subfamily by finite pigeonhole.
The same finite choice of components of $\Theta$ can be made on a dense
subfamily. The result is two fixed bounded scalar nilsequences
$\theta,\theta'$ of degree at most $s$ such that
\begin{equation}\label{v15:gi:0:eq:factorcorr}
 n\longmapsto f(n+h)\overline{f(n)}\theta(n+h)\theta'(n)
\end{equation}
is degree-at-most-$s-2$ biased for a dense set of $h$.
The $B$-factor is absorbed into that lower-degree test. The $A$-factor
$\psi(n)$ is multiplied into the unshifted component forming $\theta'$.
No scalar component is divided by, or assumed to have modulus one.

The printed source on p.~1272 uses the opposite coordinate restrictions
while claiming independence of $h$. The corrected downsets and the
approximation estimate above are the explicit replacement for that step;
its historical printed form is preserved in the findings record.

\subsection{Uniform nilcharacter selection and fixed geometries}

Here is the selection needed at the start of each induction step. Use
the finite-family inverse statement already obtained at the previous
order; the finite-family extraction below is performed after each order,
before proceeding to the next. This avoids assuming that an arbitrary
external assignment of nilmanifolds to parameters is an internal family.

On one fixed taut degree-$d$ nilmanifold, the top group is central and
its quotient by its lattice is a compact torus $T$. For each integral
character $\eta$ of $T$, a unit-vector smooth function with that vertical
character exists as follows. Cover the compact quotient by $T$ by finitely
many trivializing charts. Choose smooth real functions $\rho_j$ supported
inside the charts with $\sum_j\rho_j^2=1$, using a smooth partition of
unity and normalizing its vector by its positive Euclidean norm. In the
$j$th trivialization define the $j$th component by
$\rho_j(b)e(\eta(t))$, and extend it by zero off that chart. These
components are smooth, have the exact character, and their squared
magnitudes sum to one. The central rational quotient and its lattice
are justified by the accepted quotient geometry in the local Lemma 13.2
unit. This is the explicit construction used in published Lemma 6.4.

Components of such unit-vector functions, including frequency zero,
form a conjugation-stable multiplicative family: use tensor products
for multiplication and conjugation for conjugation. They contain
constants and separate points. Distinct points on the same $T$-fibre
are separated by an integral character, with a nonzero component of
the unit vector. Points on different fibres are separated by a smooth
frequency-zero scalar unit function pulled back from the base.
Stone--Weierstrass therefore gives uniform approximation of every
continuous test by a finite linear combination of these components.

For the fixed norm and Lipschitz bound in the previous inverse theorem,
Arzela--Ascoli gives a finite uniform net of tests. Apply the preceding
approximation to each element of that net with error smaller than a
fixed fraction of the positive inverse correlation. There are finitely
many models and net elements, so the total list of unit-vector test
maps, dimensions, frequencies, coefficient bounds and Lipschitz bounds
is fixed. Pigeonholing a scalar component of an approximating sum
preserves a fixed positive correlation. No component is assumed to
have modulus one; the entire vector has norm one.

For each finite interval and each of its good shifts, ordinary choice
selects one polynomial orbit and one item from this fixed finite list.
Taking the ultralimit of these finite assignments gives an internal
family. A finite pigeonhole step restricts to a dense set of shifts with
one common fixed model and fixed test map if required. The unit and
vertical identities hold at every representative and hence for all
limit arguments, with the natural extension of the chosen standard
character. This establishes the selection interface needed here from
the preceding finite inverse theorem. It does not replace separate
selection applications later in Sections 8--13, which remain with their
actual consumers.

For the countability interface, use the rational Mal'cev-coordinate
description of a lattice in a connected simply connected nilpotent Lie
group, as invoked in published Section 3 and p.~1250. In each finite
dimension the rational structure constants have countably many
possibilities. The rational subspaces defining the finitely many
filtration groups are specified by finite matrices with rational
entries, again a countable choice. In the rational structure determined
by the lattice, a finite list of rational logarithmic generators
specifies the lattice. Thus there are at most countably many such
filtered lattice data up to isomorphism. All local constructions here
have finite rational filtrations and discrete cocompact lattices, so
belong to this class. On each compact quotient any two fixed smooth
Riemannian metrics are bi-Lipschitz equivalent: compare their positive
quadratic forms on the compact unit tangent bundle. These are the exact
countability and metric comparisons used in the extraction below.

\subsection{The finite-group inequality needed for the last induction step}

\begin{lemma}\label{v15:gi:0:lem:CS}
For bounded complex functions $a,b$ on a finite abelian group $Q$ and an
integer $d\geq1$,
\[
 \GIAE_{h\in Q}\|a(\cdot+h)b(\cdot)\|_{U^d(Q)}^{2^d}
 \leq
 \|a\|_{U^{d+1}(Q)}^{2^d}
 \|b\|_{U^{d+1}(Q)}^{2^d}.
\]
\end{lemma}
\begin{proof}
For $t=(t_1,\ldots,t_d)$ write
$C_a(x;t)=\prod_\omega\mathcal C^{|\omega|}a(x+\omega\cdot t)$,
and define $C_b$ likewise. Expanding the left side and using the bijection
$(x,h)\mapsto(x,y=x+h)$ gives
\[
 \GIAE_t(\GIAE_y C_a(y;t))(\GIAE_x C_b(x;t)).
\]
Its absolute value is at most
$\bigl(\GIAE_t|\GIAE_yC_a(y;t)|^2\bigr)^{1/2}
 \bigl(\GIAE_t|\GIAE_xC_b(x;t)|^2\bigr)^{1/2}$.
Each squared expression is the $(d+1)$-cube average: expand the square
and substitute its second base point as the first plus a new increment.
It therefore equals the corresponding $U^{d+1}$ norm to the power
$2^{d+1}$. This proves the inequality. The original left side is
nonnegative because it is an average of powers of Gowers seminorms.
\end{proof}

Choose a cyclic group $Q$ with size between $C_sN$ and $2C_sN$ for a fixed
sufficiently large $C_s$, and zero-extend all functions supported on $[N]$.
The normalized interval and group norms differ by the fixed indicator
factor. For each fixed order that factor is bounded above by one and below
by $N/|Q|$, using the finite-group Gowers--Cauchy--Schwarz inequality.
Thus a fixed positive interval norm gives a fixed positive group norm,
and conversely; the constants depend only on the fixed ambient ratio and
order. Taking $C_s$ larger than $2^{s+1}$ ensures the same ambient group
is valid for all interval orders used here. The zero-extended product in
\eqref{v15:gi:0:eq:factorcorr} is exactly the group product for $|h|\leq N$;
there is no wrapped overlap since the ambient size is larger than $4N$.

By interface 6, \eqref{v15:gi:0:eq:factorcorr} has a uniformly positive
$U^{s-1}[N]$ norm for a dense set of $h$. Let $a=f\theta$ and
$b=\overline f\theta'$, both zero-extended. The average over all group
shifts of $\|a(\cdot+h)b\|_{U^{s-1}(Q)}^{2^{s-1}}$ is bounded below by a
fixed positive constant: the contributing set of interval shifts has
positive density also in $Q$, and the other summands are nonnegative.
Lemma~\ref{v15:gi:0:lem:CS} and the fixed bound on $b$ imply
$\|a\|_{U^s(Q)}$ is positive by a standard amount, hence so is
$\|f\theta\|_{U^s[N]}$. The inverse theorem at order $s-1$ now makes
$f\theta$ correlate with a degree-at-most-$s-1$ nilsequence $\psi$.
The actual product $\theta\psi$ is a degree-at-most-$s$ nilsequence on
the product filtered nilmanifold. Thus $f$ is $s$-biased.

\subsection{Starting and completing the induction}

The bounded induction base is Section~\ref{v15:gi:GI-BASE2}. It starts from the exact forward bracket theorem,
proves the affected trivial-subgroup specialization, constructs a
periodic connected polynomial test and transfers its correlation
directly to the normalized interval.

For completeness the initial derivative step has the same normalization
control. The cube identity on the fixed ambient cyclic group gives
\[
 \GIAE_h\|\Delta_h f\|_{U^s(Q)}^{2^s}
 =\|f\|_{U^{s+1}(Q)}^{2^{s+1}}.
\]
If the interval norm is positive by a standard amount, so is the right
side. Since $f$ is bounded, all summands have a fixed upper bound. The
nonzero differences occur only on the shifts represented by $[-N,N]$,
so a dense subset of those shifts has a uniformly positive $U^s[N]$
norm. Apply the inverse theorem of order $s-1$, then interface 2, to
obtain the family required by interface 3. Interfaces 3--5 and the preceding
section complete the induction step. The bounded limit-$N$ case is included:
a bounded interval has standard finite size, and any function on it is
represented by a Lipschitz function along a rational circle rotation,
with a standard finite Lipschitz bound. Thus the induction proves the
limit polynomial inverse conclusion for every limit $N$, conditional on
the listed intermediate interfaces and base.

We spell out the finite-family extraction rather than replace it by
pointwise compactness. Enumerate the rational filtered nilmanifolds of
degree at most $s$ up to isomorphism as $M_1,M_2,\ldots$, with one fixed
smooth metric on each, using interface 7. If the finite polynomial inverse
statement failed at some fixed $\delta>0$, for each $j$ choose a
$1$-bounded $f_j:[N_j]\to\GIAC$ with norm at least $\delta$, but correlation
at most $1/j$ against every polynomial nilsequence from the first $j$
manifolds with $1$-bounded test and Lipschitz norm at most $j$.
Its ultralimit has positive norm, so the limit inverse conclusion gives a
fixed standard filtered nilmanifold, bounded Lipschitz limit test and
positive standard correlation. Normalize the test by a fixed bound to
make it $1$-bounded. It is isomorphic to some fixed $M_{j_0}$; fixed
metric equivalence transfers its Lipschitz bound to a standard constant.
Representatives of the polynomial and test satisfy the same bounds with
one slightly larger standard constant on an ultrafilter-large set.
For $j$ larger than that constant and $j_0$, the selected correlation
upper bound tends to zero, contradicting the positive lower bound.
This yields the finite-family polynomial statement with all parameters
fixed independently of $N$.

Finally interface 8 converts every member of the finite polynomial family
to a fixed nilmanifold of step at most $s$, with a fixed Lipschitz bound
and exactly equal sequence. A finite maximum of those bounds and the
unchanged positive correlation give the exact linear target in section 1.
This final join depends on the reviewed increment construction; it does
not import the failed full-polynomial suspension claim.

\subsection{Conclusion and published inputs}

The preceding argument assembles the local results and their stated hypotheses. The affine coefficient transfer, multivariable factorization, bias and torsion conclusions are used through the explicit substitutions above.

The published forward bracket and Bohr theorems, quantitative equal-side equidistribution, rational-coordinate results and compact cube constraint retain their stated scope. The conclusion has the normalized norm, quantifier order, fixed finite family and same-step linear output specified at the start.


\section{The polynomial GI(2) base from the published bracket theorem}\label{v15:gi:GI-BASE2}
\newcommand{\GIBR}{\mathbb R}
\newcommand{\GIBZ}{\mathbb Z}
\newcommand{\GIBC}{\mathbb C}
\newcommand{\GIBE}{\mathbb E}
\newcommand{\GIBee}{\mathrm e}


This section derives the exact polynomial interval base from the forward
bracket-quadratic theorem and Bohr regularity in Green--Tao,
\emph{An inverse theorem for the Gowers $U^3(G)$ norm}, Proceedings of
the Edinburgh Mathematical Society 51 (2008), 73--153. It replaces the
needed use of the false quadratic-extension assertion by an explicit
finite-box proof. It also constructs the bracket tests directly on
connected polynomial models, so neither Theorem 12.9 nor its erroneous
skew-shift iterate parameters are a theorem premise. The forward
Bohr/bracket inverse theorem and its unchanged proof inputs remain
published inputs, not newly proved by this unit. The complete inverse theorem is assembled in Section~\ref{v15:gi:ASSEMBLY}.

\subsection{Exact target and published input}

The required base is this statement. For every $\delta>0$ there are
finitely many fixed rational filtered nilmanifolds of degree at most two,
with connected simply connected groups, and constants $L<\infty,c>0$
such that every complex $1$-bounded $f:[N]\to\GIBC$ with
$\|f\|_{U^3[N]}\geq\delta$ correlates by at least $c$ with
$F(g(n)\Gamma)$ on one of them. The polynomial $g$ is adapted,
$\|F\|_\infty\leq1$, and $\operatorname{Lip}(F)\leq L$.
All geometry and constants are fixed before $N$ and $f$.
The interval norm is zero extension to a sufficiently large cyclic
group divided by the norm of the interval indicator.

The published forward input is Theorem 10.11, pp.~124--125, especially
(10.7) and (10.9). For each fixed sufficiently small $\eta>0$, and
each prime $P>2$, a function $u:\GIBZ/P\GIBZ\to\GIBC$ with
$|u|\leq1$ and $\|u\|_{U^3}\geq\eta$ yields a frequency set
$S$ of size at most $D(\eta)$, a proper progression $Q$ containing
a Bohr set $B(S,\rho_0)$ with $\rho_0\geq\rho_*(\eta)>0$,
and one bracket quadratic
\[
 \phi(n)=c+\sum_j a_j\{\alpha_jn\}
                  +\sum_{j,k}a_{jk}\{\alpha_jn\}\{\alpha_kn\}
       \pmod1 .
\]
Here $\alpha_j=\xi_j/P$ represents a frequency in $S$,
$\{t\}\in[-1/2,1/2)$, and coefficients are unrestricted real numbers.
For every nonempty $A\subseteq Q$ there is a shift $h_A$ such that
\[
 \left|\GIBE_{n\in A}u(n+h_A)\GIBee(-\phi(n))\right|
 \geq b_*(\eta)>0.
\]
The actual powers and exponential constants printed in 10.11 are
stronger than needed; only the stated positive/finiteness dependence
on $\eta$ is consumed. Its proof uses the forward Theorem 2.7,
Corollary 10.6 and Proposition 10.10. The affected special case in the
latter is supplied next.

\subsection{The exact specialization replacing quadratic extension}

The final PDF's Lemma 10.7, pp.~122--123, invokes quadratic extension
from a subgroup. The author notice of December 2023 identifies
Proposition 3.2, p.~88, as false. The later use needed here is
Proposition 10.10, where the progression is proper and $H=\{0\}$.
The notice numbers these items 10.6 and 10.8 in its preprint version;
we bind to the final PDF numbering and argument.

\begin{lemma}\label{v15:gi:1:lem:box}
On a finite box $B=\prod_{j=1}^d\{0,\ldots,L_j-1\}$, any
$v:B\to\GIBR/\GIBZ$ whose third additive differences vanish whenever
the whole cube lies in $B$ has the form
\[
 v(l)=c+\sum_j b_jl_j+\sum_j d_j\binom{l_j}{2}
                       +\sum_{j<k}e_{jk}l_jl_k.
\]
Consequently a locally quadratic phase on a proper progression with
trivial subgroup has the polynomial coordinate representation needed
in Proposition 10.10, without extension to the ambient group.
\end{lemma}
\begin{proof}
Coordinates of length one can be omitted. Each mixed second difference
is constant on its rectangular domain: every coordinate difference of
it is a vanishing third difference, and the domain is connected by
coordinate steps. The same holds for each available diagonal second
difference when $L_j\geq3$. Choose their constants as $e_{jk}$ and
$d_j$, setting $d_j=0$ when $L_j=2$. Subtract these quadratic terms.
Every remaining first difference is constant because its differences
in every available coordinate vanish. When $L_j=2$ its domain has only
one slice in direction $j$, so mixed differences suffice. Subtract the
resulting linear function; the remainder is constant by connectedness.
The coordinate map of a proper progression is injective and pulls
back the local quadratic identities to the box. This proves the
claim. Division by two in $\GIBR/\GIBZ$, if one wants the symmetric
coefficient convention of (10.5), is possible; no uniqueness is used.
\end{proof}

In Proposition 10.10 the independent vectors supplied by Corollary 10.6
express each progression coordinate as a real linear combination of
$\{\xi n/P\}$ on that progression. Substituting the lemma's polynomial
gives precisely its bracket quadratic. The later localization and
shift in the proof of 10.11 use this specialization, not the extension
of a phase from a nontrivial subgroup. The remaining published forward
inverse and Bohr-set applications are retained as explicit inputs.
This does not claim to prove a corrected general version of 3.2.

\subsection{A smooth periodic bracket test with positive correlation}

Use Lemmas 8.1--8.2 to choose a regular Bohr set
$B=B(S,\rho)\subseteq Q$, where $\rho$ has a positive lower bound
depending only on $\eta$ and may be taken at most $1/8$.
Its density in $\GIBZ/P\GIBZ$ is at least $b_0(\eta)>0$.
Apply the displayed (10.9) input with $A=B$ and write its shift as $h$.

For $d=|S|>0$, choose a fixed $0<\epsilon<1/(100d)$ sufficiently
small in terms of $b_*(\eta)$ and $d$. Let $\chi:\GIBR/\GIBZ\to[0,1]$
be one on $|x|\leq\rho(1-\epsilon)$, zero on
$|x|\geq\rho(1+\epsilon)$, and linear on the intervening strips.
Its Lipschitz bound is at most $(2\epsilon\rho)^{-1}$, with an
inessential larger fixed bound permitted. In particular it vanishes
on a fixed neighbourhood of the fractional-part seam at $1/2$.
Put
\[
 W(n)=\prod_{j=1}^d\chi(\alpha_j n)^{1+4d},\qquad
 \Psi(n)=W(n)\GIBee(-\phi(n)).
\]
The difference $W-1_B$ vanishes outside the regularity boundary shell
$B(S,\rho(1+\epsilon))\setminus B(S,\rho(1-\epsilon))$ and
is bounded by one. Bohr regularity bounds the shell by
$C d\epsilon |B|$ with an absolute fixed $C$. Choose $\epsilon$
so that $Cd\epsilon\leq b_*(\eta)/2$. It follows that
\begin{equation}\label{v15:gi:1:eq:bias}
 \left|\GIBE_{n\in\GIBZ/P\GIBZ}u(n+h)\Psi(n)\right|
 \geq \tfrac12 b_*(\eta)b_0(\eta)=:c_0(\eta)>0.
\end{equation}
When $S$ is empty the phase is constant and this conclusion already
uses a constant test; that case needs no cutoff construction.

Every factor in $\Psi$ is periodic with period $P$, because
$P\alpha_j$ is integral. We now represent this same function for
all integers, not merely signed representatives of $\GIBZ/P\GIBZ$.

\subsection{Connected polynomial models for every bracket factor}

Write the coefficients below as $a=q+s$, with $q\in\GIBZ$ and
$|s|\leq1/2$. A linear bracket factor is
\[
 \chi(\alpha n)\GIBee(a\{\alpha n\})
 =\chi(\alpha n)\GIBee(s\{\alpha n\})\GIBee(q\alpha n).
\]
It is the test $F(x,y)=\chi(x)\GIBee(sx+y)$ along the linear
polynomial $(\alpha n,q\alpha n)$ on $\GIBR^2/\GIBZ^2$.
The seam is harmless because $\chi$ vanishes there. This test is
$1$-bounded and has Lipschitz bound $O(1+\operatorname{Lip}\chi)$,
uniformly in $a,\alpha$.

A quadratic bracket factor satisfies the exact identity
\begin{align*}
 &\chi(\alpha n)^2\chi(\gamma n)^2
                    \GIBee(a\{\alpha n\}\{\gamma n\})\\
 &=\bigl[\chi(\alpha n)\chi(\gamma n)
             \GIBee(s\{\alpha n\}\{\gamma n\})\bigr]
   \GIBee(q\alpha\gamma n^2)
   \bigl[\chi(\gamma n)\GIBee(-q\alpha n[\gamma n])\bigr]
   \bigl[\chi(\alpha n)\GIBee(-q\gamma n[\alpha n])\bigr].
\end{align*}
Indeed the omitted term in expanding the integer part is
$q[\alpha n][\gamma n]\in\GIBZ$, where $[t]=t-\{t\}$.
The first bracket is a test on a two-torus with uniformly bounded
Lipschitz constant. The pure quadratic phase is a test $\GIBee(z)$
along $z=q\alpha\gamma n^2$ on $\GIBR/\GIBZ$, using the degree-two
filtration $G_0=G_1=G_2=\GIBR$, $G_3=0$.

For the remaining bracket use the connected simply connected Heisenberg
group with coordinates $(u,v,w)$ and law
\[
 (u,v,w)(u',v',w')=(u+u',v+v',w+w'+uv').
\]
Its lattice is $\GIBZ^3$ in these coordinates; its degree-two filtration
has $G_0=G_1=G$ and $G_2=\{(0,0,w)\}$.
On a right lattice quotient, reduce $v$ first, then $u$. The resulting
coordinates are
\[
 (\{u\},\{v\},w-u[v]\pmod1).
\]
Define a test on this fundamental cube by
$F_H(\{u\},\{v\},z)=\chi(\{v\})\GIBee(z)$.
It is well-defined and globally Lipschitz. Across the $u$ faces the
central coordinate does not change; the central faces are periodic;
across the $v$ faces the cutoff is identically zero in a neighbourhood.
Within the remaining finitely many fixed charts its Lipschitz constant
is bounded by a fixed geometric multiple of
$1+\operatorname{Lip}\chi$.

Evaluate this test on the polynomial
$p(n)=(\beta n,\alpha n,0)$. The exact value is
\[
 F_H(p(n)\Gamma)=\chi(\alpha n)\GIBee(-\beta n[\alpha n]).
\]
For completeness its left first difference is
$(\beta h,\alpha h,-\alpha\beta hn)$, its next difference
is $(0,0,-\alpha\beta hk)$, and its third difference is the
identity. These satisfy the degree-two filtration for every integer
shift, so this is an adapted polynomial without any appeal to an
incorrect formula for linear skew iterates. Take $\beta=q\gamma$
and interchange $\alpha,\gamma$ for the two remaining factors.

Apply the linear construction to each coefficient of $-\phi$ and
the quadratic construction to all ordered pairs, inserting zero
coefficients if needed. Each frequency receives one cutoff from the
linear factor and $4d$ from the quadratic factors, exactly the power
in $W$. The constant phase is harmless. Products give a fixed
connected rational polynomial model of degree at most two, of dimension
at most $2d+9d^2$. The product test is $1$-bounded and has a finite
Lipschitz bound depending only on $d,\rho,\epsilon$. Arbitrarily
large coefficients occur only in polynomial orbit parameters.
There are finitely many geometric product types for $d\leq D(\eta)$.

Thus the model represents the periodic function $\Psi(n)$ exactly
for all $n\in\GIBZ$. Individual polynomial factors need not be
$P$-periodic; their product is periodic by the exact bracket identity.
That distinction is used in the next change of variables.

\subsection{The normalized interval transfer}

Fix an integer $C>16$ and choose a prime $P$ with $CN<P<2CN$.
Zero-extend the given interval function to $u$ on $\GIBZ/P\GIBZ$.
The interval-indicator Gowers norm is at least its mean $N/P$, so
$\|u\|_{U^3}\geq\delta/(2C)$. Choose once and for all a positive
$\eta$ at most this bound and within the smallness range of the
published forward theorem. All preceding data now depend only on
$\delta$.

Using the exact periodicity of $\Psi$ in \eqref{v15:gi:1:eq:bias}, change
variables to obtain
\[
 \left|\GIBE_{n\in\GIBZ/P\GIBZ}u(n)\Psi(n-h)\right|\geq c_0(\eta).
\]
The zero extension gives
\[
 \left|\GIBE_{n\in[N]}f(n)\Psi(n-h)\right|
 \geq (P/N)c_0(\eta)>0.
\]
The translated sequence is represented by the same test on the same
fixed model, with polynomial orbit $g(n-h)$. Translation preserves
adapted polynomiality, either by the derivative criterion or B.5.
No signed-representative split, discontinuous interval indicator or
additional $N$-dependent geometry is introduced. Norm, Lipschitz
and correlation bounds have precisely the required fixed dependence.
Bertrand's postulate applies to the integer $CN>1$, so the prime
choice also covers every positive $N$.

\subsection{Source comparison and bounded related-use findings}

The unpublished author notice identifies the false general quadratic
extension, but is used here only to locate the affected proof. The
needed trivial-subgroup specialization is proved in Section 2.
Its published source locator is Proposition 3.2 on p.~88, correcting
the p.~89 metadata typo in assembly001's RELATED-USES record.

The actual publisher image of Example 12.3 on p.~136 gives
$T(x,y)=(x+\alpha,y+\beta+mx)$ but prints an iterate omitting
$mnx$ and using $mn(n+1)\alpha/2$. Direct iteration gives
\[
 T^n(x,y)=(x+n\alpha,
 y+n\beta+mnx+m\alpha\tbinom n2).
\]
Already $n=1$ distinguishes these formulas. The later pure quadratic
choice in Lemma 12.11 inherits the printed iterate convention.
Rather than silently change its parameters, Section 4 constructs
all needed factors directly on connected polynomial models.
This preserves the qualitative base theorem's strength and avoids
those printed iterates entirely.

The cutoff's derivative bound also depends on its transition width.
This proof retains the finite factors $(\epsilon\rho)^{-1}$ instead
of importing the universal numerical bounds stated for those tests
in Lemma 12.11. The required application needs only finite bounds
after $\eta$ is fixed.

The base uses the forward chain
$2.7\to10.6\to10.10\to10.11$, the affected trivial-subgroup use
of 10.7, Bohr regularity 8.1--8.2, and the elementary construction
12.2--12.11 as consumed by the interval base. The chain notation records its published dependencies.
The converse half of 2.7, the finite-field-only extension use and
unrelated applications are not premises of this unit. The base uses the precise bracket theorem and the local construction, with the converse proved separately.


\section{Tautification and the shifted nilcharacter construction}\label{v15:gi:GTZ-LP-TAUTIFICATION}
\newcommand{\GICZ}{\mathbb Z}
\newcommand{\GICN}{\mathbb N}
\newcommand{\GICR}{\mathbb R}
\newcommand{\GICpoly}{\operatorname{poly}}
\newcommand{\GICNil}{\operatorname{Nil}}


The geometric input is the rational-subgroup and quotient argument in Section~\ref{v15:gi:GTZ-LP13-MULTILINEARITY}.

\subsection{The exact change of filtration}

All Lie groups and rationality conventions are those of
\cite[Sections 3--6]{v15-gi-2-GTZ}. A rational subgroup is connected, and its intersection
with the ambient lattice is cocompact. The three domain filtrations used here
are the degree filtration on $\GICZ^a$, the multidegree filtration on $\GICZ^a$,
and the degree-rank filtration on $\GICZ$ from \cite[Definition 6.17]{v15-gi-2-GTZ}.
For limit objects, defining relations are interpreted by their natural
ultrapower extension as in \cite[Appendix A.2--A.3, pp.~1324--1325]{v15-gi-2-GTZ}.
Thus vertical equivariance quantifies over every limit point of the
nilmanifold and every limit element of the top group, using the extension
of the standard vertical character. We do not infer this exact internal
identity from an identity asserted only at standard points.
For these three cases respectively put
\[
 H=G_1,\qquad H=\langle G_{e_1},\ldots,G_{e_a}\rangle_{\GICR},
 \qquad H=G_{(1,0)}.
\]
In the middle case the connected real generated subgroup equals the group
generated by the connected subgroups: the latter is the connected analytic
subgroup with the generated Lie algebra, and such subgroups in a simply
connected nilpotent group are closed. The rational-subgroup lemma in
Section~\ref{v15:gi:GTZ-LP13-MULTILINEARITY} shows that $H$ is rational.

\begin{theorem}\label{v15:gi:2:thm:taut}
Suppose $G/\Gamma$ is a rational filtered nilmanifold with finite downset of
degrees $J$ for one of these three domain filtrations. Every polynomial
nilsequence $F(g(n)\Gamma)$ admits an exactly equal representation on
$H/(H\cap\Gamma)$ with filtration
\[
 G'_0=H,\qquad G'_i=G_i\quad(i>0).
\]
For degree-rank, $0$ denotes $(0,0)$. This filtration is taut and has the
same positive-degree groups and no larger degree support, hence the same
allowed degree bound. Its test-function Lipschitz
bound depends only on the original fixed filtered nilmanifold and original
Lipschitz bound, not on the polynomial coefficients. For a nilcharacter
of positive degree $d$, the same standard vertical character on $G_d$ is
preserved. These assertions hold uniformly for limit maps and families of
polynomial coefficients with fixed geometric data.
\end{theorem}
\begin{proof}
As in Section~\ref{v15:gi:GTZ-LP13-MULTILINEARITY}, first restrict to $G_0$ and its lattice if
necessary. Choose $g(0)=q\gamma$ with $q$ in a fixed compact representative
set, and set
\[
 h(n)=q^{-1}g(n)\gamma^{-1},\qquad F_q(x\Gamma)=F(qx\Gamma).
\]
Then $h(0)=1$, $F_q(h(n)\Gamma)=F(g(n)\Gamma)$, and $h$ is polynomial
for the old filtration. Indeed every $G_i$ is normal in $G_0$; use the
constant map $g(0)^{-1}$, product closure in \cite[Corollary B.4]{v15-gi-2-GTZ},
and filtered conjugation by $q$. The Lipschitz bound is uniform because
$q$ lies in a fixed compact set.

We show that $h(n)$ belongs to $H$. In the degree case the difference
$h(n)h(0)^{-1}$ is in $G_1$ for every $n\in\GICZ^a$, by the derivative
condition. In the degree-rank case the same statement holds with
$G_{(1,0)}$. In the multidegree case connect $0$ to $n$ by a finite path
of positive or negative coordinate steps. At a step in the $j$th coordinate,
$h(v\pm e_j)h(v)^{-1}\in G_{e_j}$. Multiplying these increments along the
path and using $h(0)=1$ proves the assertion. This argument applies to every
integer input and puts no bound on its size.

Every original positive group is contained in $H$. For degree and
degree-rank this follows from nesting and the least positive index.
For multidegree, any positive index dominates some $e_j$.
Consequently the displayed new groups are well-defined subgroups of $H$,
are descending, and satisfy the commutator axioms inherited from $G$,
including those with index zero because $H\subseteq G_0$.
Their intersections with $H\cap\Gamma$ equal their intersections with
$\Gamma$ and are lattices. In degree-rank all the required equalities
$G_{(d,0)}=G_{(d,1)}$ for $d>0$ remain unchanged. The new group at zero is
exactly the required positive group or group generated by the minimal
positive groups, so the new filtration is taut.

All positive iterated derivatives of $h$ have the same required group
membership as before, and the zeroth-order membership was just checked.
Thus $h$ is polynomial for the new filtration. Restrict $F_q$ through the
fixed smooth inclusion $H/(H\cap\Gamma)\longrightarrow G/\Gamma$.
The inclusion is an embedding of compact manifolds, so this costs only a
fixed Lipschitz comparison factor and gives the asserted exact sequence.

For a nilcharacter let $\eta:G_d\longrightarrow\GICR$ be its standard
vertical character. The proof in Section~\ref{v15:gi:GTZ-LP13-MULTILINEARITY} applies verbatim
with $G_d$ in place of its top cube group: for each $t\in G_d\cap\Gamma$,
the continuous map $q\mapsto\eta(qtq^{-1})$ on connected $G_0$ is
integer-valued by the unit-norm vertical identity at $q\Gamma$.
It is therefore constant. The lattice spans $G_d$ as a connected real
group, so $\eta(qtq^{-1})=\eta(t)$ for all $t\in G_d$ and $q\in G_0$.
Hence $F_q$ and its restriction have the same standard vertical character.

For families, choose one bounded standard representative section of the
fixed compact quotient. Such a section can be formed from finitely many
local smooth sections, taking the first available chart. Apply this one
section to each constant coefficient. The resulting choices are limit maps
when the coefficients are limit maps. Discontinuities in the parameter do
not enter a Lipschitz bound in the orbit variable. All metric constants and
groups remain fixed. The exact identities can be imposed on representatives
on an ultrafilter-large set by the witness argument of
Section~\ref{v15:gi:GTZ-LP13-MULTILINEARITY}; the ordinary proof then passes to the limit.
\end{proof}

This also proves equality of the classes of nilsequences under the old and
taut conventions, at each finite degree support: one inclusion is the theorem
and the reverse inclusion is immediate. It therefore preserves the
equivalence relation defined by a lower-degree nilsequence. It is not merely
an assertion that the top-degree function has a similar replacement.

\begin{corollary}\label{v15:gi:2:cor:central}
In a taut degree or degree-rank filtration of degree at most $(d,r)$ or $d$,
the ambient group has nilpotency class at most $d$. In a taut multidegree
filtration supported on a finite downset $J$, its class is at most
$D=\max_{a\in J}|a|$. If the support is contained in the downset below one
positive index $d$, the group $G_d$ is central.
\end{corollary}
\begin{proof}
For degree and degree-rank, brackets of $j$ elements of the ambient first
positive group belong respectively to $G_j$ or $G_{(j,0)}$, which vanish
when $j>d$. For multidegree, expand an iterated Lie bracket using the
generating Lie algebras $\mathfrak g_{e_i}$. A bracket of length $j$ lies in
some $\mathfrak g_a$ with $|a|=j$, hence vanishes for $j>D$.
The generated Lie algebra has the asserted nilpotency class, and so does
its simply connected group. Finally $[G_d,G_{e_i}]\subseteq G_{d+e_i}=1$
in the multidegree case. The same argument uses the first positive group
in the other two cases. Since these groups generate the ambient group,
$G_d$ is central.
\end{proof}

\subsection{A complete shifted nilcharacter construction}

\begin{theorem}\label{v15:gi:2:thm:shift}
For any of the three domain filtrations above, let $\chi$ be a nilcharacter
of positive degree $d$. Let $i>0$ and let $u$ belong to the domain group at
index $i$. Then
\[
 n\longmapsto\chi(n+u)\otimes\overline{\chi(n)}
\]
is a nilsequence with support strictly below $d$. The representing geometric
data and Lipschitz bound can be fixed independently of $u$ within this
domain group. For limit $u$ the representation and test function can be
chosen in a limit fashion. This supplies the shift assertion required in
\cite[Lemma E.8(iv)]{v15-gi-2-GTZ} for the three domains used in the proof.
\end{theorem}
\begin{proof}
By Theorem~\ref{v15:gi:2:thm:taut} represent the same $\chi$ on a taut nilmanifold
$G/\Gamma$, with the same degree and vertical character. Write its
polynomial map as $g$. In the degree-rank case use $i=(1,0)$: any
nonzero element of a positive domain group requires this index, and if the
given domain group is trivial then $u=0$ also belongs to this first group.
Put $U(v)=(v,1)$ and $D(v)=(v,v)$ in $G^2$.
For each index $j$ let $L_j$ be the connected subgroup generated by
\[
 U(G_{j+i})\quad\hbox{and}\quad D(G_j).
\]
These groups are rational by Section~\ref{v15:gi:GTZ-LP13-MULTILINEARITY}, and are descending.
The generator commutators have the following memberships:
\begin{align*}
 [D(G_j),D(G_k)]&\subseteq D(G_{j+k}),\\
 [D(G_j),U(G_{k+i})]&\subseteq U(G_{j+k+i}),\\
 [U(G_{j+i}),U(G_{k+i})]&\subseteq U(G_{j+k+2i})
                                      \subseteq U(G_{j+k+i}).
\end{align*}
Passing to generated Lie algebras and using BCH, as in
Section~\ref{v15:gi:GTZ-LP13-MULTILINEARITY}, proves $[L_j,L_k]\subseteq L_{j+k}$, including
index zero. Thus $L_0/\Delta$, with $\Delta=L_0\cap\Gamma^2$, is a rational
filtered nilmanifold of degree at most $d$. In the degree-rank case its
additional equalities hold: $L_{(a,0)}=L_{(a,1)}$ follows from
$G_{(a,0)}=G_{(a,1)}$ and $G_{(a+1,0)}=G_{(a+1,1)}$ for $a>0$.

The difference $a_u(n)=g(n+u)g(n)^{-1}$ is polynomial for the shifted
filtration $(G_{j+i})_j$. Its zeroth-order value belongs to $G_i$.
Every further iterated derivative is an iterated derivative of $g$ with
the derivative in direction $u$ innermost, so the defining index sum has
the additional term $i$. Therefore $U(a_u(n))$ and $D(g(n))$ are polynomial
for $(L_j)$. By \cite[Corollary B.4]{v15-gi-2-GTZ}, so is their product
\[
 P_u(n)=(g(n+u),g(n)).
\]
This verifies the actual orbit, not only its coordinate projections.

As $i>0$, $G_{d+i}=1$ and $L_d=D(G_d)$. Corollary~\ref{v15:gi:2:cor:central}
shows that $G_d$ is central in $G_0$, hence $L_d$ is central in $L_0$.
The function
\[
 T(x\Gamma,y\Gamma)=F(x\Gamma)\otimes\overline{F(y\Gamma)}
\]
restricted to $L_0/\Delta$ is invariant under $D(G_d)$, since its two
vertical characters cancel. The rational central-quotient argument in
Section~\ref{v15:gi:GTZ-LP13-MULTILINEARITY} therefore gives a nilsequence on
$L_0/L_d$ with the exact value $T(P_u(n)\Delta)$. Its filtration group at
$d$ is trivial and all groups outside the downset below $d$ were already
trivial. Its support is consequently strictly below $d$.
The same argument gives a fixed Lipschitz bound on the quotient.

The quotient may not yet be taut. Apply Theorem~\ref{v15:gi:2:thm:taut} to this
lower-degree nilsequence, preserving its exact values without enlarging its
degree support.
The groups $L_j$ and their quotient depend on the index $i$ but not on the
element $u$. The map $u\mapsto P_u(0)$ is a limit map when $g$ and $u$ are
limit data. The one fixed bounded section used in Theorem~\ref{v15:gi:2:thm:taut}
then gives a limit choice of the final normalized representation, with
uniform Lipschitz bounds independent of $u$. This proves all the assertions.
\end{proof}

For the erratum's needed case the domain is $\GICZ^2$ with multidegree,
$d=(1,s-1)$, $i=(1,0)$ and $u=(k,0)$. The resulting support is contained in
\[
 \{a:a\leq(0,s-1)\}\ \cup\ \{a:a\leq(1,s-2)\},
\]
uniformly in $k$. This is exactly the support needed before applying the
published splitting and vertical Fourier lemmas in the proof of the new
difference decomposition. Those applications, and the subsequent two
Cauchy--Schwarz steps, are separate proof obligations.

\subsection{Scope of the correction}

The tautification theorem supplies the general representation repair
suggested, but not fully proved, in Section 2 of the author erratum.
It also proves the top centrality and nilpotency-class assertions for the
taut conventions. The shifted construction above supplies the needed
instance of E.8(iv) without leaving its auxiliary rationality or central
quotient implicit. This unit does not claim to have audited every filtration
introduced in the published proof. In particular, E.8(v), E.9, Lemma 12.1,
the Section 12 quotient, and the final modified Proposition 7.3 argument
still need their own exact applications or replacements.


\section{A local replacement for the filtration argument in GTZ Lemma 13.2}\label{v15:gi:GTZ-LP13-MULTILINEARITY}
\newcommand{\GIDR}{\mathbb R}
\newcommand{\GIDZ}{\mathbb Z}
\newcommand{\GIDN}{\mathbb N}
\newcommand{\GIDone}{\mathbf 1}
\newcommand{\GIDee}{\mathrm e}
\newcommand{\GIDLip}{\operatorname{Lip}}
\newcommand{\GIDpoly}{\operatorname{poly}}
\newcommand{\GIDSymb}{\operatorname{Symb}}



This section proves the multilinearity lemma and its auxiliary filtration construction.

\subsection{Statement and conventions}

We use the nilmanifold, polynomial-map and nilcharacter conventions of
\cite[Sections 3, 4 and 6]{v15-gi-3-GTZ}. In particular all Lie groups below are connected,
simply connected and nilpotent; filtration groups are connected rational Lie
subgroups. A rational subgroup $J$ means that $J\cap\Gamma$ is a lattice in $J$.
The multidegree ordering is coordinatewise. Write
$\GIDone=(1,\ldots,1)\in\GIDN^s$ and $|a|=\sum_j a_j$.
The group $G_a$ is trivial when $a\not\leq\GIDone$.
We use $\GIDee(t)=\exp(2\pi i t)$ and complex conjugation componentwise in tensor
products. A vector-valued nilcharacter has a unit-vector-valued Lipschitz
test function and a real vertical character, as in \cite[Definition 6.1]{v15-gi-3-GTZ}.
For limit objects, every defining relation is read in its natural ultrapower
extension, using \cite[Appendix A, Definition A.2 and the following discussion,
and Definition A.3, pp.~1324--1325]{v15-gi-3-GTZ}. In particular the vertical identity
is required for every limit point $x\in{}^*(G/\Gamma)$ and every
$t\in{}^*G_{\GIDone}$, with the standard homomorphism $\eta$ extended to
$^*\eta$. This is the internal nilcharacter convention used in this proof.
We do not infer this exact identity from an identity only at standard points;
bounded Lipschitz continuity alone would not justify that inference.

\begin{theorem}[The unchanged multilinearity statement]\label{v15:gi:3:thm:main}
Let $s,m\geq1$ be fixed standard integers. Let $\chi$ be a nilcharacter of
multidegree at most $\GIDone$ on $\GIDZ^s$, or a limit nilcharacter with that
description in the sense of \cite[Sections 5--6]{v15-gi-3-GTZ}. For fixed homomorphisms
$L_1,L'_1,L_2,\ldots,L_s:\GIDZ^m\longrightarrow\GIDZ$, the three pullbacks below
are nilcharacters of degree at most $s$ and satisfy
\begin{align*}
 [\chi(L_1+L'_1,L_2,\ldots,L_s)]_{\GIDSymb^s}
 ={}&[\chi(L_1,L_2,\ldots,L_s)]_{\GIDSymb^s}\\
 &+[\chi(L'_1,L_2,\ldots,L_s)]_{\GIDSymb^s}.
\end{align*}
The same assertion holds in each coordinate. More explicitly, each component
of
\begin{equation}\label{v15:gi:3:eq:triple}
 \chi(x+x',y_2,\ldots,y_s)\otimes
 \overline{\chi(x,y_2,\ldots,y_s)}\otimes
 \overline{\chi(x',y_2,\ldots,y_s)}
\end{equation}
is a polynomial nilsequence of degree at most $s-1$ on $\GIDZ^{s+1}$.
The representing nilmanifold is fixed by the original filtered nilmanifold;
its test-function Lipschitz bound is bounded in terms of that fixed data and
the original Lipschitz bound, independently of the polynomial coefficients.
\end{theorem}

The symbol notation abbreviates the degree-$s$ equivalence relation of
\cite[Definitions 6.6 and 6.22]{v15-gi-3-GTZ}. Thus the explicit assertion about
\eqref{v15:gi:3:eq:triple} supplies the asserted equality, with no approximation or
loss in degree. It is enough to prove the ordinary statement with uniform
fixed-data bounds and then take limit sequences. We supply that uniformity
below. We do not assume $G_0$ equals the group generated by its positive
multidegree groups.

\subsection{Subgroups and lattices used in the construction}

We use the following basic nilpotent Lie-group facts recorded in
\cite[Section 3, pp.~1242--1243]{v15-gi-3-GTZ}: the exponential is a homeomorphism;
the connected subgroup generated by a set is the exponential of the generated
Lie algebra; and, for a finite set $A$, its discrete generated group is
cocompact in its connected real generated group. The last assertion is applied
only when $A$ lies in an already specified discrete lattice. Here are the
compactness details needed to apply these facts.

\begin{lemma}\label{v15:gi:3:lem:rational}
Let $M$ be a connected simply connected nilpotent Lie group with lattice
$\Lambda$. Let $J_1,\ldots,J_t$ be connected rational subgroups. Then the
connected subgroup generated by the $J_i$ is closed, simply connected, and
rational in $M$. This also applies to diagonal or coordinate embeddings of
such subgroups into a finite product of $M$.
\end{lemma}
\begin{proof}
First a lattice $\Lambda_J$ in a connected Lie group $J$ is finitely generated.
Choose a proper left-invariant path metric, and $R$ such that every point is
within distance $R$ of some lattice point; compactness of $J/\Lambda_J$,
followed by inversion, gives such an $R$. Follow a path from the identity to a
lattice point in steps of length at most one and choose nearby lattice points,
using the prescribed endpoints. Consecutive chosen lattice points differ by
an element in the compact ball of radius $2R+1$. Its intersection with the
discrete lattice is finite and generates the lattice.

The connected real subgroup generated by this finite generating set is all
of $J$. To justify this last assertion, let $U$ be that subgroup. It contains
$\Lambda_J$, so $J/U$ is compact. A connected subgroup of a simply connected
nilpotent group is closed, and its homogeneous space is a Euclidean space.
One can see the latter directly by a chain of Lie subalgebras from
$\mathfrak u$ to $\mathfrak j$, each of codimension one and an ideal in the
next: the normalizer of a proper subalgebra of a nilpotent Lie algebra strictly
contains it, so such a chain is obtained by choosing one normalizing vector at
each step. Successive exponential coordinates give the Euclidean coordinates
on $J/U$. For completeness, the normalizer assertion follows by taking the
first member of the ascending central series not contained in the subalgebra;
the preceding member is contained in it, so any vector in that member but
outside the subalgebra normalizes it. Hence compactness forces $U=J$.

Choose a finite generating set for each $J_i\cap\Lambda$. Their union $A$
lies in $\Lambda$ and generates all the $J_i$ as a connected real group.
Its abstract generated group is discrete because it lies in $\Lambda$.
The finite-set fact from \cite[p.~1243]{v15-gi-3-GTZ} therefore makes it a lattice in
$\langle A\rangle_{\GIDR}$. Consequently the larger group
$\Lambda\cap\langle A\rangle_{\GIDR}$ is also a lattice there. A Lie subalgebra
of a nilpotent Lie algebra exponentiates to a closed simply connected subgroup,
which proves the other assertions. For a diagonal or coordinate embedding,
apply the same argument to the images of the finite generating sets, which
lie in the product lattice.
\end{proof}

We will also use the elementary central-quotient consequence explicitly.
If $J$ is such a group, $\Delta$ its lattice, and $C$ a connected central
rational subgroup, then $J/C$ is connected simply connected nilpotent and the
image of $\Delta$ is a lattice. Cocompactness follows from the continuous image
of $J/\Delta$. For discreteness, if lattice cosets tend to the identity coset,
choose local lifts tending to the identity and write the lattice points as
those lifts times elements of $C$. Reduce those elements modulo the lattice
$C\cap\Delta$ into a fixed compact set. The adjusted lattice points now lie
in a compact subset of $J$, so only finitely many occur. Thus their cosets
cannot accumulate at the identity. Simple connectedness follows either from
exponential coordinates or the homotopy sequence with connected kernel $C$.

\subsection{Bounded normalization at zero}

Write $\chi(n)=F(g(n)\Gamma)$, with the data in
Theorem~\ref{v15:gi:3:thm:main}. Initially restrict to the fixed subgroup $G_0$ and its
lattice $\Gamma_0=G_0\cap\Gamma$, restricting $F$ through the smooth inclusion
$G_0/\Gamma_0\longrightarrow G/\Gamma$. This only changes a fixed metric
comparison constant. In the following notation $G=G_0$ and $\Gamma=\Gamma_0$.

Choose a fixed compact set of representatives for $G/\Gamma$. Write
$g(0)=q\gamma$, with $q$ in that compact set and $\gamma\in\Gamma$. Set
\begin{equation}\label{v15:gi:3:eq:normalize}
 h(n)=q^{-1}g(n)\gamma^{-1}
     =q^{-1}\bigl(g(n)g(0)^{-1}\bigr)q,
 \qquad F_q(z\Gamma)=F(qz\Gamma).
\end{equation}
Then $h(0)=1$ and $F_q(h(n)\Gamma)=F(g(n)\Gamma)$ exactly.
Each $G_a$ is normal in $G_0$, because $[G_0,G_a]\subseteq G_a$.
The constant map with value $g(0)^{-1}$ is polynomial into the original
filtration. Product closure of polynomial maps
\cite[Corollary B.4]{v15-gi-3-GTZ}, followed by the filtered conjugation by $q$, shows
that $h$ is polynomial into the same multidegree filtration. No conjugation
of the lattice is made.

The maps $z\Gamma\mapsto qz\Gamma$ are smooth with uniformly bounded
derivatives for $q$ in the fixed compact set. Thus $F_q$ has a uniform
Lipschitz bound. If $F(tz)=\GIDee(\eta(t))F(z)$ for $t\in G_{\GIDone}$, then
\begin{equation}\label{v15:gi:3:eq:vertical}
 F_q(tz)=\GIDee(\eta_q(t))F_q(z),\qquad
 \eta_q(t)=\eta(qtq^{-1}).
\end{equation}
This is a homomorphism on $G_{\GIDone}$. Its restriction to
$G_{\GIDone}\cap\Gamma$ is integer-valued: put $z=\Gamma$ in
\eqref{v15:gi:3:eq:vertical} and use both $F_q(t\Gamma)=F_q(\Gamma)$ and the fact that
$F_q(\Gamma)$ is a unit vector. In fact $\eta_q=\eta$. For each fixed
$t\in G_{\GIDone}\cap\Gamma$, the continuous map
$q\mapsto\eta(qtq^{-1})$ on connected $G_0$ has integer values, so it equals
its value $\eta(t)$ at the identity. The lattice generates $G_{\GIDone}$ as a
connected real group, by the argument in Lemma~\ref{v15:gi:3:lem:rational}.
Two continuous real homomorphisms agreeing on that lattice therefore agree
on $G_{\GIDone}$. This proves the assertion for every $q\in G_0$.
In particular the normalization preserves the same standard vertical
character, including for limit choices of $q$.

\subsection{The descending auxiliary filtration}

For $a\in\{0,1\}^s$ with $a\ne0$, define the following homomorphisms on
$G_a$ whenever the indicated first coordinate occurs:
\[
 A_a(u)=(u,u,u)\quad(a_1=0),\qquad
 B_a(u)=(u,u,1),\quad C_a(u)=(u,1,u)\quad(a_1=1).
\]
For $1\leq i\leq s$ let $H_i$ be the connected subgroup of $G^3$ generated
by these images for all $a$ with $|a|\geq i$. Set $H_{s+1}=\{1\}$ and
\begin{equation}\label{v15:gi:3:eq:K}
 K=H_1,\qquad K_0=K_1=H_1,\qquad K_i=H_i\ (i\geq1).
\end{equation}
Every group just defined is fixed by the original filtration, not by $g$.
Lemma~\ref{v15:gi:3:lem:rational} makes it a connected closed simply connected rational
subgroup of $G^3$. Put $\Delta=K\cap\Gamma^3$.
The sequence $(K_i)$ is descending by its definition.

We check the commutator law. For generators of types $A,A$, $A,B$, and
$A,C$, the coordinatewise commutator has respectively type $A,B,C$ and
belongs to the group with index $a+b$. The weight adds.
If two generators both have first index one, every possibly nontrivial
coordinate commutator belongs to $G_{a+b}=\{1\}$ because
$(a+b)_1=2$. Thus the $B,B$, $B,C$, and $C,C$ commutators are trivial.
These assertions hold for their Lie algebras as well. Bracketing generated
Lie algebras and using the Jacobi identity gives
\[
 [\mathfrak h_i,\mathfrak h_j]\subseteq\mathfrak h_{i+j}
 \quad (i,j\geq1).
\]
Here iterated brackets remain in the indicated space: each bracket is a sum
of iterated brackets of the original generators, and its total original
weight is at least $i+j$. BCH, which terminates for a nilpotent group,
then gives $[H_i,H_j]\subseteq H_{i+j}$. The index-zero assertions also
hold since $K_0=H_1$: for example
$[K_0,K_j]\subseteq H_{j+1}\subseteq H_j=K_j$.
Hence $(K_i)$ is a rational degree filtration, is taut at zero, and has
degree at most $s$.

The top group has a particularly simple exact description:
\begin{equation}\label{v15:gi:3:eq:top}
 K_s=H_s=\{(uv,u,v):u,v\in G_{\GIDone}\}.
\end{equation}
Indeed $a=\GIDone$ is the only index of weight $s$ and has first coordinate
one. The group $G_{\GIDone}$ is abelian since
$[G_{\GIDone},G_{\GIDone}]\subseteq G_{2\GIDone}=\{1\}$. Thus its $B$ and $C$
images commute and their product is exactly the displayed connected group.
It is central in $K$, also following from the commutator law.

\subsection{The exact polynomial orbit}

Choose once and for all a total order of $\{0,1\}^s$ extending its partial
order. By \cite[Lemma B.9]{v15-gi-3-GTZ}, the normalized map has the Taylor expansion
\begin{equation}\label{v15:gi:3:eq:Taylor}
 h(n)=\prod_{a\in\{0,1\}^s\setminus\{0\}} u_a^{n^a},
 \qquad u_a\in G_a,\qquad n^a=\prod_{j=1}^s n_j^{a_j}.
\end{equation}
The absent constant coefficient equals $h(0)=1$. Define
\[
 P(x,x',y)=
 \bigl(h(x+x',y),h(x,y),h(x',y)\bigr),\qquad
 y=(y_2,\ldots,y_s).
\]
This is exactly the ordered product of the following factors, in the order
of \eqref{v15:gi:3:eq:Taylor}. If $a_1=0$, its factor is
\[
 A_a(u_a)^{\prod_{j=2}^s y_j^{a_j}}.
\]
If $a_1=1$, its factor is
\[
 B_a(u_a)^{x\prod_{j=2}^s y_j^{a_j}}
 C_a(u_a)^{x'\prod_{j=2}^s y_j^{a_j}}.
\]
The latter identity is exact because the powers of the same element $u_a$
commute. In either case the integer exponent has degree $|a|$ and the fixed
group element lies in $K_{|a|}$. A map $v^{Q(n)}$ of this form is polynomial
into $(K_i)$: its $j$-fold derivative is $v^{\Delta^j Q(n)}$, which lies in
$K_j$ for $j\leq\deg Q$, and is the identity for $j>\deg Q$.
Here derivatives commute within the cyclic one-parameter subgroup generated
by $v$, and $K_{\deg Q}\subseteq K_j$ is used in the first case.
The values also lie in $K_0$. Product closure
\cite[Corollary B.4]{v15-gi-3-GTZ} now proves that
\begin{equation}\label{v15:gi:3:eq:poly}
 P\in\GIDpoly(\GIDZ^{s+1}_{\GIDN}\longrightarrow K_{\GIDN}).
\end{equation}
This argument proves the assertion for all integer inputs, not just a box.
In particular it verifies that the orbit lies in the normalized ambient
group $K$, which the exact-weight construction at index zero did not ensure.

\subsection{The lower-degree quotient and its test function}

Let $\overline K=K/K_s$, let $\pi:K\longrightarrow\overline K$ be the
quotient map, and let $\overline\Delta=\pi(\Delta)$.
The central-quotient argument after Lemma~\ref{v15:gi:3:lem:rational} shows that
$\overline K/\overline\Delta$ is a compact nilmanifold with connected simply
connected nilpotent group. Give it the filtration
$\overline K_i=\pi(K_i)$. It is taut and has degree at most $s-1$.
For each $i\leq s$, $K_s\subseteq K_i$, so
\[
 \overline\Delta\cap\pi(K_i)=\pi(\Delta\cap K_i).
\]
The right side is cocompact in $\pi(K_i)$, by the compact quotient of
$K_i/(\Delta\cap K_i)$; it is discrete as a subgroup of
$\overline\Delta$. Thus the quotient filtration is rational at every index.
Equation~\eqref{v15:gi:3:eq:poly} remains polynomial after applying $\pi$, directly
from the derivative definition.

Define on $K/\Delta$
\[
 T_q(z_1,z_2,z_3)=
 F_q(z_1\Gamma)\otimes\overline{F_q(z_2\Gamma)}
                    \otimes\overline{F_q(z_3\Gamma)}.
\]
It is well-defined and Lipschitz, by restriction from $(G/\Gamma)^3$.
For $(uv,u,v)\in K_s$, equation~\eqref{v15:gi:3:eq:vertical} multiplies this tensor
by
\[
 \GIDee\bigl(\eta_q(uv)-\eta_q(u)-\eta_q(v)\bigr)=1.
\]
Therefore $T_q$ is invariant under $K_s$ and descends to a function
$\overline T_q$ on $\overline K/\overline\Delta$.

For explicit control of its regularity, fix a smooth metric on $K/\Delta$
and average it over the compact central torus $K_s/(K_s\cap\Delta)$.
The averaged metric and the original fixed metric are uniformly equivalent.
The quotient metric makes the quotient a Riemannian submersion. A function
constant on its fibers and Lipschitz upstairs is Lipschitz downstairs with
the same bound for these metrics, by horizontal path lifting. All the
metric comparisons are fixed by the original filtered nilmanifold.
The bound for $T_q$ is at most a fixed multiple of the bound for $F_q$,
since its three factors have norm one. The earlier compact normalization
therefore makes the Lipschitz bound for $\overline T_q$ uniform in $g$.

Finally the expression \eqref{v15:gi:3:eq:triple} is exactly
\[
 \overline T_q\bigl(\pi(P(x,x',y))\overline\Delta\bigr).
\]
This proves the claimed degree at most $s-1$ representation.

For clarity, the individual factors also have valid degree-$s$
representations. In one copy of $G$, take the connected group $J_i$ generated
by $G_a$ for $|a|\geq i\geq1$, and set $J_0=J_1$.
The same rationality and commutator arguments give a taut rational degree-$s$
filtration on $J_1$, with $J_s=G_{\GIDone}$. The Taylor factors of $h$ and of
each of its linear pullbacks are polynomial for this filtration by the same
integer-polynomial argument. Restricting $F_q$ gives a unit-vector-valued
Lipschitz test function with vertical character $\eta_q$.
Hence those factors are nilcharacters of degree at most $s$.

Substitute the homomorphism
$n\mapsto(L_1(n),L'_1(n),L_2(n),\ldots,L_s(n))$ into the lower-degree
representation. Derivatives, or the explicit polynomial factors above, show
that degree does not increase. The definition of symbol equivalence proves
Theorem~\ref{v15:gi:3:thm:main}. Relabeling coordinates gives the other coordinate
identities.

\subsection{Limit data and the actual published consumers}

For a limit nilcharacter the filtered nilmanifold and its metric data are
standard fixed data, and the polynomial and Lipschitz test functions are
limit sequences with bounded Lipschitz norm. Choose the bounded representatives
$q$ separately in each member of those sequences; this gives a limit map
with values in the same compact set. All groups $K_i$, the lattice,
the quotient and the metric constants constructed above are unchanged along
the sequence. The same standard vertical character is retained by the preceding
argument. Representatives for the test functions may be taken to satisfy
their defining unit-norm and vertical-character identities on an
ultrafilter-large set of indices: otherwise, choosing a violating point,
and for the second identity a violating group element, at every bad index
would contradict the exact identity of the limit function at the resulting
limit points. This uses the all-limit quantifiers just specified and the
evaluation rule in \cite[Definition A.3]{v15-gi-3-GTZ}; it would be invalid with only
standard-point quantifiers. Discard those bad indices. The formulas now hold at every
remaining member, and the test-function bounds are uniform. Their limit is
therefore a limit polynomial nilsequence of
degree at most $s-1$ in precisely the required sense. No compactness argument
selects a nilmanifold depending on an unbounded polynomial coefficient.

The first published use, immediately after Lemma 13.2 on p.~1318, expands
$[\widetilde\chi(n+h,\ldots,n+h)]$ by applying the identity in each
coordinate. Before any use of symmetry the exact expansion is
\[
 [\widetilde\chi(n+h,\ldots,n+h)]
 =\sum_{\omega\in\{0,1\}^s}
 [\widetilde\chi(z_{\omega_1},\ldots,z_{\omega_s})],
 \qquad z_0=n,\ z_1=h.
\]
Every map used here is a fixed integer homomorphism from $\GIDZ^2$.
Thus the coefficient counting and subsequent symmetry step in that argument
receive exactly the same multilinearity statement; no stronger inverse
theorem is asserted by this local repair.

The further explicit appeals in Lemmas 13.3, 13.4 and 13.5 on
pp.~1319--1322 likewise use addition, subtraction and repeated copies of
their variables. These are integer homomorphisms from the displayed finite
Cartesian powers of $\GIDZ$, so the theorem applies without a range or
normalization change. Constant shifts in those arguments still require the
separate symbol-shift assertion in E.8; this proof does not discharge it.
The repaired lemma supplies the multilinearity input to Theorem 7.4,
then to the published inverse-theorem argument. Other corrections, including
the false original Proposition 8.3 and its replacement, remain separate
obligations. Source and current-manuscript consumer spans are recorded in
the accompanying contracts.


\section{Difference decomposition and the additive-quadruple estimate}\label{v15:gi:GTZ-LP02-03}
\newcommand{\GIEE}{\mathbb E}
\newcommand{\GIEZ}{\mathbb Z}



\subsection{Inputs and conventions}
We use the natural-ultrapower convention of GTZ Appendix A and the
shift theorem in Section~\ref{v15:gi:GTZ-LP-TAUTIFICATION}. A bound $\gg1$ is positive and standard. Write $N$ for an
unbounded limit integer, $J=[N]$ and $I=[-N,N]\cap{}^*\GIEZ$. All changes
of variables are integer identities at the representative indices.
The author erratum guides the proof. Its Proposition 3.1 and modified
Proposition 7.3 are the conclusions here, rather than assumed premises.

The shift theorem gives a lower-support sequence with fixed geometry
and uniform Lipschitz bound. It gives neither a separated finite
decomposition nor a correlation conclusion. Products, conjugates,
and translates use product manifolds, conjugated tests, and translated
polynomial orbits, respectively. A bounded internal scalar depending
only on parameters is degree zero in the averaging variable.

\subsection{Difference decomposition}
\begin{proposition}[LP02]
For standard $s\ge3$, let $\chi$ be a multidegree $(1,s-1)$
nilcharacter. For every standard $\epsilon>0$ each component of
$D_k(h,n)=\chi(h+k,n)\otimes\overline{\chi(h,n)}$
is approximated uniformly in $k,h,n$ to error $\epsilon$ by
\[
 \sum_{j=1}^{L} \Psi_{j,k,a_j}(n) B_{j,k}(h,n).
\]
Here $L$ is standard, $\Psi_{j,k,a_j}$ is a component of a
unit-vector degree-$(s-1)$ nilcharacter, and $B_{j,k}$ is supported
on the downset below $(1,s-2)$. The families are internal with
standard finite geometry lists and uniform bounds.
\end{proposition}
\begin{proof}
Apply the project shift theorem, Theorem 3, with positive index
$(1,0)$ and shift $(k,0)$. The support is in the union of the downsets
$J_0=\downarrow(0,s-1)$ and $J_1=\downarrow(1,s-2)$.
We now give the fixed geometric split explicitly. Write $G/\Gamma$
for the taut lower-support representation from the shift theorem and
$J=J_0\cup J_1$. For each positive $a\in J$ take a rational basis of
$\mathfrak g_a$ for the rational structure of $\Gamma$. This exists
because $G_a$ is rational. Take a free rational Lie algebra with one
generator of weight $a$ for each basis vector. Kill all homogeneous
words whose total weight is outside $J$. The excluded weights form
an upward-closed set, so these words generate a homogeneous ideal.
All generator weights are positive. Thus word length is bounded by
$\max_{a\in J}|a|=s-1$ in the quotient, which is finite-dimensional
and nilpotent. Let $\mathfrak u=\bigoplus_{a\in J\setminus\{0\}}
\mathfrak u[a]$ denote it and let $U$ be its simply connected real
group. Set $\mathfrak u_b=\bigoplus_{a\ge b}\mathfrak u[a]$.
This is a descending rational multidegree filtration with the
required bracket rule.

Sending generators to their basis vectors defines a rational
homomorphism $\rho:U\to G$. The killed words map to zero by the
filtration rule in $G$. It maps each positive layer onto $G_b$:
the generators tagged $b$ span $\mathfrak g_b$, and all weights
at least $b$ map into that layer. At zero it is onto because
$G$ is taut and generated by its coordinate groups $G_{(1,0)}$ and $G_{(0,1)}$.
Thus the polynomial lifts exist by published GTZ Corollary B.10,
applied to this layerwise-surjective filtered map, and can be
selected internally at the representative indices.

Choose a homogeneous rational basis of $\mathfrak u$ and let
$L_0$ be its integer span. There is a positive integer $D$ such
that $\Lambda=\exp(D L_0)$ is a lattice and $\rho(\Lambda)\subset
\Gamma$. Here are the scaling details. In this basis BCH is a
finite rational polynomial with linear part addition. For two
inputs in $D L_0$, terms of degree at least two lie in $D L_0$
once $D$ clears all the coefficient denominators. Inversion
negates the logarithm. Thus this set is a subgroup, discrete
in exponential coordinates. It is cocompact: order the basis
by increasing total weight and successively reduce each weight's
logarithm coordinates modulo $D$. Right multiplication by an
element of the corresponding homogeneous lattice leaves lower
weights unchanged, and at the current weight is ordinary
addition. The resulting bounded box of logarithms is a compact
set of representatives. The same argument applies to every
filtration layer and every homogeneous quotient. To impose
$\rho(\Lambda)\subset\Gamma$, use rational Mal'cev coordinates
of the second kind adapted to $\Gamma$, as in GTZ Section 3.
The coordinates of $\rho(\exp x)$ are finite rational
polynomials in the chosen coordinates of $x$, with zero constant
term. Taking $D$ divisible by their denominators makes the
coordinates of $\rho(\exp(D x))$ integral for $x\in L_0$.
This gives the required lattice inclusion. Increasing $D$ preserves
the BCH and compactness arguments.

For $i=0,1$, quotient $U$ by the connected subgroup with
Lie algebra the sum of homogeneous weights outside $J_i$.
These sums are ideals by the downset property. The image
lattices are the exponential lattices in the surviving
homogeneous bases. The pair of quotient maps separates points
of $U/\Lambda$: if two points have the same image in both,
the logarithm of their group difference lies in the projected
lattice in each quotient. Each homogeneous weight survives
in at least one quotient. Hence every coordinate of that
logarithm belongs to $D\mathbb Z$, so the difference lies in
$\Lambda$. This proves separation on the quotient spaces,
rather than just intersection-triviality of the two ideals.

The pullbacks of Lipschitz tests from the two compact quotients
generate a conjugation-closed algebra containing constants
and separating points. Stone--Weierstrass gives a uniform
approximation by finite sums of products of these tests.
The lifted tests of the shift family have a common Lipschitz
and sup bound on this fixed compact manifold, since the
map to $G/\Gamma$ is fixed and smooth. A standard finite
$\epsilon/4$ net for this bounded equicontinuous class gives
finitely many standard product lists, each chosen to tolerance
$\epsilon/4$. Select the first net entry within $\epsilon/4$
using an internal sup-norm comparison at representative indices.
Use a smaller net radius if needed. Thus the splitting error
is at most $\epsilon/2$, uniformly in $k,h,n$. Pad the lists
with zero factors to a standard common length.

In the $J_0$ quotient all surviving generators have zero first
weight and strictly positive second weight. The restricted
one-variable degree filtration is therefore taut: its degree-zero
and degree-one groups coincide. The highest layer is central by
the bracket law and the degree bound. The orbit is independent
of $h$, since the first-variable derivative lies in the trivial
first-coordinate group. Restrict to $h=0$ and apply published
Lemma E.5 with this taut one-variable geometry of degree $s-1$.
Use the ordinary functions constructed in its proof on printed p.~1353.
For each of the finitely many standard $J_0$ tests, fix ordinary Lipschitz
unit-vector tests with global vertical frequency and ordinary top-invariant
coefficient functions. Extend these functions to the ultrapower after making
these finite choices. The resulting tests have exact unit norm and
equivariance at all limit points and all limit top-group elements, and
internal selection from the finite lists preserves both identities.
Its test decomposition gives linear functionals of unit-vector
nilcharacters times degree-at-most-$(s-2)$ tests. Expand the
functionals in coordinates and absorb their coefficients and
lower-degree factors into the $J_1$ factors. The product still
has support $J_1$. Choose individual Fourier tolerances after
the finite lists and their sup bounds so that the added error
is at most $\epsilon/2$. On fixing $h,k$, the lower factor has
degree at most $s-2$ in $n$ by the $(0,1)$ derivative index.
All geometry lists are fixed and finite, and all orbit data
and choices are internal.
\end{proof}

\subsection{The two averages}
\begin{proposition}[LP03]
Under the hypotheses of GTZ Theorem 7.2, for standard $s\ge3$,
there are $\gg N^3$ additive quadruples in $H$ for which
\[
 \chi_{h_1}(n)\otimes\chi_{h_2}(n+h_1-h_4)
 \otimes\overline{\chi_{h_3}(n)}
 \otimes\overline{\chi_{h_4}(n+h_1-h_4)}
\]
is $(s-2)$-biased on $J$. The witnessing tests have standard
bounded complexity and positive standard correlation.
\end{proposition}
\begin{proof}
Extend $f$ by zero outside $J$. GTZ A.11 and A.12 select the
lower-degree correlation tests internally on the dense internal set
$H\subset I$. The relation includes a selected positive lower bound
$1/j$ with standard integer $j$. Select global representatives of the lower-degree tests and of the
restricted nilcharacters by the same internal selection, storing
the representations as part of the witness. Only their original
values on $J$ enter the unsmoothed identities. The bounded-complexity descriptions
are a countable union of internal sets. The boundedness principle
A.5 gives uniform standard complexity and a common positive lower
bound. Expand the bounded-dimensional tensor contraction in
coordinates. Retain a positive-density subset of $H$ with a fixed
witnessing coordinate pair by the standard finite pigeonhole.
Write $U_h,V_h$ for the conjugates of those coordinates of
$\chi(h,\cdot)$ and $\chi_h$. Then there are scalar lower-degree
tests $c_h$ such that
\[
 \left|\GIEE_{n\in J}f(n+h)\overline{f(n)}U_h(n)V_h(n)c_h(n)\right|
 \ge\eta>0
\]
on the retained set, with standard uniform bounds. Align the phase
by multiplying $c_h$ by a unit scalar. Set $c_h=0$ outside that
set and extend $V_h$ boundedly there. This is an internal choice.
The real average over $h\in I$ is $\gg1$.

Apply Cauchy--Schwarz in $n$ to eliminate $\overline{f(n)}$.
It gives
\[
 \GIEE_{n\in J}\left|\GIEE_{h\in I}
       f(n+h)U_h(n)V_h(n)c_h(n)\right|^2\gg1.
\]
Expand the square and put $h'=h+k$, with
$K=[-2N,2N]\cap{}^*\GIEZ$. The exact new parameter cutoff is
$1_{h+k\in I}$, degree zero in $n$. On $I\times K$ the new
normalized average is $|I|/|K|$ times the old one. With
$c_{h,k}=1_{h+k\in I}c_h\overline{c_{h+k}}$ this gives
\[
 \left|\GIEE_{h\in I,k\in K,n\in J}
 f(n+h)\overline{f(n+h+k)}
 U_h(n)\overline{U_{h+k}(n)}
 V_h(n)\overline{V_{h+k}(n)}c_{h,k}(n)\right|\gg1.
\]
The $U$ product is a component of $D_k$. Take LP02 with error
less than half the correlation divided by a common bound for the
other factors. A term of the standard finite sum retains a positive
standard correlation. Absorb its lower factor into $c_{h,k}$.

For its surviving coordinate $\Psi_{k,a}(n)$ use the identity
\[
 \Psi_{k,a}(n)=\sum_b
 [\Psi_{k,a}(n)\overline{\Psi_{k,b}(n+h)}]\Psi_{k,b}(n+h).
\]
The identity uses the unit vector norm of $\Psi_k(n+h)$.
Individual coordinates need not have modulus one. The bracket
is a component of the conjugate transpose of the shift tensor.
The one-variable degree case of the project shift theorem gives
degree at most $s-2$, uniformly and internally in $h,k$, using the
finite geometries of LP02. Select a surviving term by the same
finite averaging argument. Absorb the bracket into $d_{h,k}$ and
write the remaining coordinate as $P_k(n+h)$. The lower test is
still exactly zero unless both $h,h+k$ are retained shifts.

Set $m=n+h$ and $M=[-2N,2N]\cap{}^*\GIEZ$. The change to the
full $M$ average, with the exact cutoff, multiplies the
normalized correlation by $N/|M|$. We have
\[
 \left|\GIEE_{m\in M,k\in K}
 f(m)\overline{f(m+k)}P_k(m)
 \GIEE_{h\in I} 1_J(m-h)V_h(m-h)
 \overline{V_{h+k}(m-h)}d_{h,k}(m-h)\right|\gg1 .
\]
The outer factor depends only on $(m,k)$ and is bounded.
The second Cauchy--Schwarz in $(m,k)$ gives $\gg1$ for the
mean squared modulus of the inner average. Expand with independent
$h,h'\in I$. Put
\[
 (h_1,h_2,h_3,h_4)=(h,h'+k,h+k,h'),\qquad n=m-h.
\]
The shift substitution is injective, with inverse
$(h,h',k)=(h_1,h_4,h_3-h_1)$. It has $h_1+h_2=h_3+h_4$.
The first cutoff restricts $n$ to $J$, where $m=n+h\in M$
automatically. The second is $1_J(n+h_1-h_4)$. The four factors
in the new average are
\[
 V_{h_1}(n)V_{h_2}(n+h_1-h_4)
 \overline{V_{h_3}(n)}\overline{V_{h_4}(n+h_1-h_4)} .
\]
The lower factor is
$d_{h,k}(n)\overline{d_{h',k}(n+h_1-h_4)}$
times the remaining cutoff. The ratio $N/|M|$ is comparable
to one. The parameter vanishing is exact unless all four shifts
are retained. This proof keeps zero extension throughout.

To absorb the cutoff, choose fixed $0<\lambda<1/4$.
On the circle $\mathbb R/16\GIEZ$ take a piecewise linear test
equal to one on $[\lambda,1-\lambda]$, zero off
$[-\lambda,1+\lambda]$ in $[-8,8)$, between zero and one,
with Lipschitz bound $O(1/\lambda)$. Evaluate it at
$(n+h_1-h_4)/N$. This is a one-step nilsequence in $n$,
with fixed geometry and internal orbit. Its argument lies in
$[-2,3]$ and does not wrap. The normalized mean absolute error
against the cutoff is $O(\lambda+1/N)$ uniformly in all shifts:
only two endpoint bands of total $O(\lambda N+1)$ integers
differ. Multiplying by bounded factors changes only the constant.
Choose $\lambda$ after the positive correlation bound so the
error loses at most half. Since $s\ge3$, the one-step test is
absorbed into degree at most $s-2$. Parameter vanishing persists.

If the outer correlation has modulus at least $\beta>0$ and
each inner average has modulus at most $B$, at least a fraction
$\beta/(2B)$ of the inner averages have modulus at least
$\beta/2$. The shift triple box has size comparable to $N^3$.
By injectivity, these yield $\gg N^3$ additive quadruples in $H$.
Conjugating the scalar correlations undoes the initial choice of
$V_h$ and preserves lower degree. The resulting scalar correlations
are coordinates of the claimed tensor, giving bias on the
original interval $J$.
\end{proof}

\subsection{Scope}
The published proof inputs used here are GTZ Appendix A, Corollary
B.10, E.5, the correlation conventions, and the stated
hypotheses of Theorem 7.2. The project shift input and these exact
applications await separate nonauthor review. This unit supplies the
additive-quadruple result, not the subsequent frequency analysis,
the other affected symbol rules, or the final GI theorem.

\section{One-variable factorization and approximate additive linearization}\label{v15:gi:GI-DF}
\newcommand{\GIFT}{\mathbb T}
\newcommand{\GIFZ}{\mathbb Z}
\newcommand{\GIFR}{\mathbb R}
\newcommand{\GIFE}{\mathbb E}
\newcommand{\GIFpoly}{\operatorname{poly}}
\newcommand{\GIFLie}{\operatorname{Lie}}
\newcommand{\GIFLip}{\operatorname{Lip}}


This section proves the one-variable factorization connections, using the final published sources [GTZ], [GT12], [GT14], and [GT08] identified in its references. Only one-variable factorization is asserted here.

\subsection{Objects, source differences, and the required outputs}
All nilpotent Lie groups below are standard, connected and simply connected;
all lattices and filtrations are standard and rational. Subgroups described as
rational are connected, and their intersection with the ambient lattice is
cocompact in that subgroup. We use the natural ultrapowers of these objects,
with the internal-function convention of [GTZ, Appendix A]. A bounded quantity
is bounded by some standard constant. In the finite combinatorial proof all
constants depend only on the stated positive density and error constants.
Passing that finite proof to an ultrapower therefore preserves uniformity and
internal choices. Braces denote the representative in $(-1/2,1/2]$.

An $\mathbb N$-filtration $G_\bullet$ may have $G_1\ne G_0=G$.
Its positive groups satisfy $[G_i,G_j]\subseteq G_{i+j}$ and are rational.
The polynomial convention is the ordinary degree filtration on the source.
In fixed adapted Mal'cev coordinates, write the coordinates of a polynomial
map as a finite polynomial in $n$ or $\binom n j$. Coordinate degree is bounded
in terms of the fixed filtered group. For a smooth factor, ``small Taylor
coefficients'' means that its constant coefficient is bounded and its
coefficient of index $j\geq1$ is $O(N^{-j})$, in the group Taylor convention
of [GTZ, Lemma B.9]. Equivalent coordinate bounds are justified below.

The statement of [GTZ, D.2] cannot be used as a literal equivalence: on
$\GIFR/\GIFZ$, $g(n)=n/N$ is equidistributed on $[N]$ but is not equidistributed
on $[\lfloor N/2\rfloor]$. For example its first Fourier average on that
half interval tends to $2i/\pi$, whereas its full-interval average tends to
zero. Its nontrivial horizontal character is $1/N$-Lipschitz. We instead use
the one-variable weak factorization [GT12, Proposition 9.2], whose proof uses
the quantitative theorem on an actual dense progression. The published
correction [GT14, pp.1175--1176] explicitly retains the one-variable results.
Its minor corrections on pp.1182--1183 are part of this source choice.

The properness requirement printed in the second bullet of [GTZ, D.4] is
removed. Only a weak-factorization step that starts with failed total
equidistribution forces a proper subgroup. Also, a rational periodic factor
is not asserted to take values in $\Gamma$.

\begin{theorem}[One-variable factorization, DF-D4]\label{v15:gi:5:thm:D}
Let $N$ be unbounded and let $g\in{}^*\GIFpoly(\GIFZ,G_\bullet)$.
There are a standard connected rational subgroup $P\leq G$ and internal
polynomial maps $\epsilon,p,\gamma\in{}^*\GIFpoly(\GIFZ,G_\bullet)$ such that
\[
             g(n)=\epsilon(n)p(n)\gamma(n)\qquad(n\in{}^*\GIFZ).
\]
The map $p$ takes values in ${}^*P$, with its induced filtration. On $[N]$,
$\epsilon$ is bounded, is $O(1/N)$-Lipschitz, and has the small Taylor
coefficient bounds above. There is a standard discrete subgroup
$\Lambda\geq\Gamma$ with $\gamma(n)\in{}^*\Lambda$ for every $n$, and a
standard integer $q>0$ with
$\gamma(n+q)\Gamma=\gamma(n)\Gamma$ for every $n$.
For every standard finite-index subgroup $\Delta\leq P\cap\Gamma$,
$p(n)\Delta$ is totally equidistributed in $P/\Delta$ on $[N]$.
The subgroup $P$ is allowed to equal $G$ or to be trivial.
\end{theorem}

\begin{theorem}[Approximate four-function linearization, DF-F1]\label{v15:gi:5:thm:F}
Let $N$ be unbounded, let $j\geq1$ be standard, and put $e=N^{-j}$.
Let $H\subseteq[-N,N]\cap{}^*\GIFZ$ be internal, and let
$\xi_i:H\to{}^*\GIFT$, $1\leq i\leq4$, be internal functions.
Suppose that for at least $cN^3$ ordered quadruples in $H^4$ satisfying
$h_1+h_2=h_3+h_4$ one has
\[
 \|\xi_1(h_1)+\xi_2(h_2)+\xi_3(h_3)+\xi_4(h_4)-a\|_\GIFT\leq C e,
\]
where $c>0$ and $C<\infty$ are standard and $a\in{}^*\GIFT$.
For any specified index $i$ there are standard $K<\infty$, $c'>0$,
an internal $H'\subseteq H$ with $|H'|\geq c'N$, internal
$\alpha_k\in{}^*\GIFT$, $\beta_k\in{}^*\GIFR$, and $\delta\in{}^*\GIFT$ such that
\[
 \xi_i(h)=\delta+\sum_{k=1}^K\{\alpha_kh\}\beta_k+O(e)
                    \quad(h\in H').                                      \tag{F}
\]
The implicit constant is standard and uniform over $H'$.
The bracket sum restricted to $H'$ is an exact Freiman homomorphism:
it preserves every relation $h_1+h_2=h_3+h_4$ with all four points in $H'$.
No initial Freiman property, boundedness of the real $\beta_k$, or regularity
of the input functions is assumed. All choices may be made internally.
\end{theorem}

\subsection{Weak factorization and coefficient bounds}
\begin{lemma}[Weak step, DF-D3]\label{v15:gi:5:lem:weak}
Suppose first that $G_1=G$. If $g(n)\Gamma$ is not totally equidistributed
on $[N]$, it admits the factorization of Theorem~\ref{v15:gi:5:thm:D} with $p$ taking
values in a standard proper connected rational subgroup, without the final
equidistribution requirement. All other bounds and the discrete-overgroup
and period assertions hold.
\end{lemma}
\begin{proof}
Failed total equidistribution supplies an internal dense progression and a
standard Lipschitz test with noninfinitesimal discrepancy. Thus for some
standard $0<\delta<1/2$ it fails total $\delta$-equidistribution in the
quantitative convention of [GT12]. Reduce $\delta$ further to include the
rationality constant of the fixed Mal'cev basis. Apply the one-variable case
of [GT12, Proposition 9.2, pp.522--523] with this $\delta$, the actual degree
of $G_\bullet$, and the actual dimension of $G$. Its bounds are standard.
The same argument works for an internal Lipschitz test with a standard norm
bound: a finite uniform net in the standard Lipschitz ball first supplies a
standard test with positive discrepancy. This uses compactness of $G/\Gamma$.

Here are the stronger bounds extracted from that source's construction,
including the distinctions needed for later use. Write $g(0)=c\lambda$
with $c$ in a fixed bounded fundamental set and $\lambda\in{}^*\Gamma$.
Then $\widetilde g=c^{-1}g\lambda^{-1}$ has value $1$ at zero. Left
translation by bounded $c$ changes a Lipschitz norm by a standard factor,
and right multiplication by $\lambda^{-1}$ changes no point of $G/\Gamma$.
Thus the same failed-total-equidistribution argument applies to
$\widetilde g$, with a possibly smaller standard $\delta$.

The source's horizontal-character step gives a nonzero standard integer
vector $k$ with
\[
 \psi(\widetilde g(n))=\sum_{r\geq1}t_r\binom n r,
 \qquad \|k\cdot t_r\|_\GIFT=O(N^{-r}).                         \tag{D1}
\]
Here $k\cdot\psi$ is the real lift of the horizontal homomorphism; in
particular it kills all commutators. Standardness of $k$ follows because the
quantitative bound leaves finitely many integer vectors. The source obtains
(D1) by applying its Theorem 2.9 to the witnessing progression and undoing
that progression by its Lemma 8.4. This use does not assert the false
single-scale equivalence described above.

Choose integers $a_r$ nearest to $k\cdot t_r$.
The coordinate support condition in [GT12, Lemma 6.7] says that $t_r$ lies
in the coordinate subspace for $G_r$. If $k$ vanishes on that subspace,
take $a_r=0$ and $v_r=w_r=0$. Otherwise fix a rational vector $u_r$ in
that subspace with $k\cdot u_r=1$, and put
\[
 v_r=a_ru_r,\qquad w_r=(k\cdot t_r-a_r)u_r.
\]
Define $\epsilon_0$ and $\gamma_0$ by Mal'cev coordinates
\[
 \psi(\epsilon_0(n))=\sum_r w_r\binom n r,
 \qquad \psi(\gamma_0(n))=\sum_r v_r\binom n r,
 \qquad p=\epsilon_0^{-1}\widetilde g\gamma_0^{-1}.             \tag{D2}
\]
The same coordinate-support lemma and closure of polynomial maps under
products show that these are polynomial for the original filtration.
The real horizontal lift vanishes identically on $p$, since its coefficient
at every index is $k\cdot t_r-(k\cdot t_r-a_r)-a_r=0$.
Consequently $p$ lies in the connected kernel of that real homomorphism.
This is a standard proper rational subgroup. A kernel of the circle-valued
homomorphism alone would not establish connectedness; (D2) supplies the real
zero required here.

Each $w_r$ is $O(N^{-r})$. The vectors $v_r$ have a common standard
denominator, but their numerators need not be bounded. This denominator
statement is all that is required. By [GT12, A.11--A.12], $\gamma_0$ is
uniformly rational and periodic modulo $\Gamma$ with a standard period.
The more explicit common discrete overgroup argument below also applies.
Set $\epsilon=c\epsilon_0$ and $\gamma=\gamma_0\lambda$.
These are the correct ordered factors of the original $g$.
Multiplication by bounded $c$ preserves the stated small-coordinate bounds,
and the added $\lambda$ lies in the original lattice. The next lemma verifies
all closure and Taylor assertions used here and in the descent.
\end{proof}

\begin{lemma}[Closure of the controlled factors]\label{v15:gi:5:lem:closure}
In a fixed rational nilpotent group, finite products of bounded polynomial
maps with coefficient bounds $O(N^{-r})$ retain those bounds and are
$O(1/N)$-Lipschitz on $[N]$. These bounds imply the group Taylor bounds in
Theorem~\ref{v15:gi:5:thm:D}. A finite collection of rational lattices commensurable
with $\Gamma$, together with sets of points whose rational coordinates have
a fixed common denominator, is contained in a standard discrete subgroup
$\Lambda\geq\Gamma$. Any fixed-degree polynomial map into this group is
periodic modulo $\Gamma$ with a standard period.
\end{lemma}
\begin{proof}
Coordinate products, inverses, and changes between first-kind and Mal'cev
coordinates are fixed polynomials by [GT12, A.2--A.3]. In a coordinate
polynomial $\sum a_r n^r$, the bounds $a_r=O(N^{-r})$ are equivalent to
bounded coefficients after substituting $n=Nt$. Products of such polynomials
retain that property. Converting back to the binomial basis gives at index
$r$ a fixed linear combination of coefficients of monomial indices $\ell\geq r$,
hence $O(N^{-r})$. Constant coefficients stay bounded. On the fixed compact
coordinate range for $0\leq t\leq1$, the group-coordinate maps have bounded
derivatives; differentiation in $n$ contributes $N^{-1}$. This proves the
Lipschitz assertion, with any fixed standard metric.

For clarity about the group Taylor convention, one may assign formal weight
$r$ to each coefficient of $n^r$ in the logarithm of a coordinate polynomial.
The recursive Taylor extraction by $g(0)$, then $g_1^{\binom n1}$, then
$g_2^{\binom n2}$, and so on uses only fixed BCH polynomials. Its coefficient
$\log g_r$ contains no monomial of formal weight less than $r$.
Indeed after removing indices below $r$, each surviving polynomial term of
formal weight $w$ has ordinary degree at most $w$. The residual group map
is the identity at $0,\ldots,r-1$, as a formal identity in its input
coefficients. Its logarithmic terms of weight $w<r$ are therefore polynomials
of degree less than $r$ with $r$ zeros, and vanish. BCH products add weights,
so this assertion is preserved at each extraction step.
The remaining finitely many monomials have weight at least $r$, and are
$O(N^{-r})$. This proves the group Taylor assertion, including under products
and fixed rational embeddings of subgroups.

For the discrete-overgroup assertion, use a rational basis adapted to the
lower central series and let its weights be $w_i\geq1$. All points of each
given lattice have logarithmic coordinates with bounded denominators, as do
the specified coordinate sets, by the fixed polynomial coordinate changes.
Choose a standard integer $D$ divisible by these denominators and by the
denominators of all coefficients in the finite BCH formula in this basis.
Increase $D$ for the finitely many basis conversions if necessary. The set
\[
 L=\sum_i D^{-(2w_i-1)}\GIFZ X_i
\]
is closed under BCH. For a nonlinear term with $t\geq2$ inputs of weights
$w_{i_1},\ldots,w_{i_t}$ contributing to output weight $w_k$, one has
$w_k\geq\sum w_{i_a}$. Therefore the difference of denominator exponents is
\[
 (2w_k-1)-\sum_a(2w_{i_a}-1)\geq t-1\geq1.
\]
The extra power of $D$ clears that BCH coefficient. Linear terms and negation
also preserve $L$. Thus $\exp L$ is a standard discrete subgroup containing
all the required points and $\Gamma$. It is cocompact, and
$[\exp L:\Gamma]$ is finite. This argument also proves commensurability
claims needed when a lattice of a rational subgroup is enlarged inside $G$.

For the period in the required right-coset convention, write the Mal'cev
coordinate vector of $\gamma(n)$ as a polynomial of fixed degree $D_0$.
Its values have a common standard coordinate denominator, since
$\gamma(n)\in{}^*\Lambda$. Interpolating at $0,\ldots,D_0$ shows that
its coefficients have a common standard denominator too. Their numerators
may be arbitrary internal integers. By the rational polynomial formulas
for multiplication and inversion, the vector
\[
             R(n,r)=\psi\bigl(\gamma(n)^{-1}\gamma(n+r)\bigr)
\]
is a polynomial in the two integer variables $n,r$ with a common standard
coefficient denominator $A$. The identity $R(n,0)=0$ is a polynomial
identity, so every monomial has a positive power of $r$.
For the standard integer $q=A$, every coordinate of $R(n,q)$ is an integer
for every internal integer $n$. Adapted Mal'cev coordinates identify
$\Gamma$ with the integer-coordinate points. Hence
$\gamma(n)^{-1}\gamma(n+q)\in{}^*\Gamma$, exactly the assertion
$\gamma(n+q)\Gamma=\gamma(n)\Gamma$.
This proves the conclusion of [GT12, A.12] directly; its printed left
quotient $\gamma(n+q)\gamma(n)^{-1}$ is not the quotient used here.
\end{proof}

\subsection{Minimal subgroup descent, with all finite covers retained}
\begin{proof}[Proof of Theorem~\ref{v15:gi:5:thm:D}]
First address the filtration convention. If $G_1$ is proper, write
$g(0)=c\lambda$ as above and set $p=c^{-1}g\lambda^{-1}$.
Then $p(0)=1$ and
\[
 p(n)=\lambda\bigl(g(0)^{-1}g(n)\bigr)\lambda^{-1}\in G_1.
\]
The groups $G_r$ are normal in $G_0$, so conjugation preserves them.
The induced filtration on $G_1$ has zeroth and first group both $G_1$,
and the positive groups remain $G_r$. Hence the one-variable result with
the convention $G_0=G_1$ applies there. This also shows directly that a
controlled representation exists: $\epsilon=c$, $\gamma=\lambda$,
and the represented subgroup is $G_1$. If $G_1=G$, the identity-factor
representation is already enough. Trivial groups cause no exception.

Call a subgroup $P$ a representing subgroup if a factorization with all
conclusions of Theorem~\ref{v15:gi:5:thm:D} except total equidistribution exists.
Choose one with least dimension among the nonempty collection of standard
rational representing subgroups. This uses a least element of the finite
set $\{0,\ldots,\dim G\}$, not an infinite iteration or a uniform choice
over all finite covers.

Suppose that for a standard finite-index $\Delta\leq P\cap\Gamma$ the
orbit $p\Delta$ is not totally equidistributed. The induced filtration of
$P$ is rational for $\Delta$, since passing to a finite-index lattice
does not change its rational Lie algebra. If $P\cap G_1$ is proper in $P$,
the normalization just given already represents $p$ in that proper subgroup,
with a bounded constant on the left and a $\Delta$-element on the right.
Otherwise apply Lemma~\ref{v15:gi:5:lem:weak} in $P/\Delta$ to obtain
\[
                 p=\epsilon_2p_2\gamma_2
\]
with $p_2$ in a standard proper rational $P_2<P$.
Rationality relative to $\Delta$ is rationality relative to $P\cap\Gamma$,
so $P_2$ is rational in $G$. The actual new factorization is
\[
                 g=(\epsilon\epsilon_2)p_2(\gamma_2\gamma).    \tag{D3}
\]
The order of both products is necessary in a noncommutative group.
Lemma~\ref{v15:gi:5:lem:closure} proves the small-coefficient assertions for the left
product, in the ambient coordinates. For the right product, the rational
lattice coordinates of $\gamma_2$ relative to $\Delta$ have common standard
denominators in the original rational structure. The same lemma supplies a
single discrete overgroup containing its values, the old overgroup, and
$\Gamma$. It also supplies the standard period modulo $\Gamma$.
Thus (D3) is a representation by the lower-dimensional subgroup $P_2$,
contrary to minimality.

No choice of $\Delta$ was imposed in that minimality argument. Consequently
the chosen $p$ is totally equidistributed for every standard finite-index
$\Delta\leq P\cap\Gamma$, with the same factors. This is the required
quantifier order. The argument only produces standard constants for this
particular internal orbit and fixed standard geometry. It does not assert
a prescribed effective bound uniform over all internal orbits.
\end{proof}

\subsection{A finite combinatorial proof of the linearization contract}
We prove a finite assertion with constants depending only on $c,C$ and then
transfer. Take $N\geq1$, $0<e\leq1$, and suppose the hypotheses of
Theorem~\ref{v15:gi:5:thm:F} hold in finite sets. Put $M=\lceil1/e\rceil$ and
$Q=100N+1$. Round each $\xi_i(h)$ in the circle to $y_i(h)\in M^{-1}\GIFZ/\GIFZ$,
with error at most $1/(2M)\leq e/2$. Since $Me\leq2$, the approximate sum
can take only a bounded number, depending on $C$, of values in this grid.
For at least $c_1N^3$ of the original quadruples its rounded value is a
fixed $a_0$ in that grid. Here $c_1>0$ depends only on $c,C$.

Embed $[-N,N]$ in $B=\GIFZ/Q\GIFZ$ and identify the target grid with $A=\GIFZ/M\GIFZ$.
In $B\times A$ form the four sets
\[
 A_1=\{(h,y_1(h))\},\quad A_2=\{(h,y_2(h))\},\quad
 A_3=\{(-h,y_3(h))\},\quad A_4=\{(-h,y_4(h))\}.             \tag{F1}
\]
The signed first coordinates are essential. The selected quadruples now
give $c_1N^3$ representations of $(0,a_0)$ in $A_1+A_2+A_3+A_4$.
The initial additive relation has no wraparound, since $4N<Q$.
Fourier inversion and Holder's inequality in this finite abelian group give
\[
 r_{A_1+A_2+A_3+A_4}(0,a_0)
          \leq\prod_{i=1}^4 E(A_i)^{1/4},\qquad
 E(A_i)=|\{x_1+x_2=x_3+x_4:x_a\in A_i\}|.
\]
For completeness, with unnormalized Fourier transforms the inverse formula
has factor $|B\times A|^{-1}$; Holder bounds the sum of four transforms
by the fourth moments, each of which equals $|B\times A|E(A_i)$.
Those factors cancel. Each $E(A_i)\leq |H|^3\leq27N^3$.
It follows that $E(A_1)\geq c_2N^3$ for a positive $c_2=c_2(c,C)$.
The same reasoning can isolate any $A_i$; reflecting its first coordinate
back if necessary gives the graph of the specified $y_i$.
We henceforth call that function $y$ and its graph $A_1$.

Apply [GT08, Theorems 5.2 and 5.3, p.95] to $A_1$ in the actual finite
group $B\times A$. Their hypotheses follow from the energy bound and
$|A_1|\leq3N$. They give a subset $\Gamma'$ of its graph and a constant
$K_0=K_0(c,C)$ such that
\[
            |\Gamma'|\geq N/K_0,
            \qquad |9\Gamma'-8\Gamma'|\leq K_0N.             \tag{F2}
\]
We use these published theorems themselves, not the special correlation
premise in [GT08, Proposition 5.4].

\subsection*{Graph refinement with the actual cyclic target}
Let $V=\{a\in A:(0,a)\in8\Gamma'-8\Gamma'\}$.
The map $\Gamma'\times V\to9\Gamma'-8\Gamma'$ given by addition is
injective: equality determines the first graph coordinate, then its unique
second coordinate, and finally $a$. Hence $|V|\leq K_0^2$.
There is a set $S_A$ of at most $O(1+\log K_0)$ characters of $A$ for which
\[
        V\cap\{a:\|\chi(a)\|_\GIFT<1/4\ \forall\chi\in S_A\}=\{0\}.
\]
Indeed, for any nonzero $a$ a uniformly random character takes values
uniformly among the roots of unity of its order; at least one third of those
values have distance at least $1/4$ from zero. Choose enough independent
characters that the union bound over $V\setminus\{0\}$ is less than one.
This proves the required separation without identifying $A$ with the dual
of $B$, an identification that is unavailable here.

Partition each coordinate of the character vector on $A$ into intervals of
length at most $1/100$. A cell contains a subset $\Gamma''\subseteq\Gamma'$
of size $c_3N$. In any element of $8\Gamma''-8\Gamma''$, every selected
character of the second coordinate has a real representative of absolute
value less than $16/100<1/4$ when its first coordinate is zero.
It follows that the only element of this difference set with first
coordinate zero is $(0,0)$. Therefore $4\Gamma''-4\Gamma''$ is a graph.
Write $H''$ for the projection of $\Gamma''$, and let
\[
 \zeta:2H''-2H''\longrightarrow A
\]
be the function whose graph is $2\Gamma''-2\Gamma''$.
It is an exact Freiman homomorphism on this domain: two sums of two of its
graph points lie in $4\Gamma''-4\Gamma''$, where the first coordinate
determines the second. Also $\zeta(0)=0$. Thus
$\zeta(x+y)=\zeta(x)+\zeta(y)$ whenever $x,y,x+y$ lie in its domain.
This proves precisely the generalization of [GT08, Lemma 9.2] needed here.

\subsection*{A proper progression and an explicit linear bracket formula}
The set $H''$ has density at least $c_3/101$ in $B$.
By [GT08, Lemma 6.4, pp.97--98] there is a bounded set of characters $S$
with $B(S,1/4)\subseteq2H''-2H''$.
Add the faithful character $x\mapsto x/Q$ to $S$, and shrink the radius
to $\rho=1/16$. Write $d=|S|$, a bounded integer.
[GT08, Lemma 10.4, p.121] supplies a proper centered progression
\[
 P=\Bigl\{\sum_{a=1}^r l_av_a: |l_a|<L_a\Bigr\}\subseteq B(S,\rho),
              \qquad |P|\geq c_4 Q,
\]
with $r\leq d$, and linearly independent vectors
$\Phi(v_a)=(\{\chi(v_a)\})_{\chi\in S}$.
The coset subgroup in that lemma is zero because $S$ contains a faithful
character. Its proof also gives
\[
 \Phi\left(\sum_a l_av_a\right)=\sum_a l_a\Phi(v_a)\quad\hbox{on }P. \tag{F3}
\]
This is the lift of a lattice progression inside the real cube, not an
unjustified inference from reduction modulo one. Delete inactive generators
with no nonzero integer coefficient allowed; all remaining $v_a$ lie in $P$.
Every partial sum along a coordinate path to a point of this centered box
stays in $P$. The local additivity of $\zeta$ consequently gives
\[
          \zeta\left(\sum_a l_av_a\right)=\sum_a l_a\zeta(v_a).
\]
Choose real lifts $b_a$ of $\zeta(v_a)$ and real dual vectors $u_a$ to the
independent vectors $\Phi(v_a)$ on their span. Using (F3), define
$b_\chi=\sum_a b_a(u_a)_\chi$. Then, for every $p\in P$,
\[
                 \zeta(p)=\sum_{\chi\in S}b_\chi\{\chi(p)\}\pmod1. \tag{F4}
\]
This is a linear calculation. It uses neither the quadratic-extension
theorem nor the quadratic classification lemma in [GT08]. The author notice
concerning that extension theorem is therefore not a new assumption here.

\subsection*{Translation, fixed carries, and the Freiman restriction}
Averaging translates of $P$ in $B$ gives an $x\in B$ with
$|H''\cap(x+P)|\geq |H''||P|/Q\geq c_5N$.
For $h=x+p$ and $h_0=x+p_0$ in this intersection, the graph relation gives
$y(h)-y(h_0)=\zeta(p-p_0)$, since $\Gamma''-\Gamma''$ is contained in
$2\Gamma''-2\Gamma''$. Both $p,p_0$ lie in the domain of $\zeta$.
Its Freiman identity applied to $(p-p_0)+p_0=p+0$ gives
\[
                 y(h)=\delta_x+\zeta(p).
\]
For each $\chi\in S$, write $\chi(h)=\alpha_\chi h$ on the cyclic group.
The difference
\[
 \{\chi(h-x)\}-\{\chi(h)\}+\{\chi(x)\}
\]
is an integer in a fixed set of at most three values. Restrict to a cell on
which all these carries are fixed. This costs a factor at most $3^d$ and,
by (F4), gives
\[
                y(h)=\delta+\sum_{\chi\in S}b_\chi\{\alpha_\chi h\}.
\]
Finally partition each $\{\alpha_\chi h\}$ into a fixed number of intervals
of length at most $1/10$, and retain a largest cell $H'$.
For any integer relation $h_1+h_2=h_3+h_4$ in $H'$, the corresponding sum
of four fractional parts, with signs $+,+,-,-$, is an integer of absolute
value at most $1/5<1$. It is therefore zero in $\GIFR$ for each character.
The real weighted sum is an exact Freiman homomorphism on $H'$.
Undoing the rounding changes the displayed value by at most $e/2$.
All density losses depend only on $c,C$ and all sets and choices are finite.
This proves the finite assertion, and transfer with $e=N^{-j}$ proves
Theorem~\ref{v15:gi:5:thm:F}. A bounded rank can be made a fixed standard rank by
passing to its ultrapower value or by adding zero terms.

\subsection{Consumer boundary and related-use scope}
Theorem~\ref{v15:gi:5:thm:D} supplies exactly the one-variable factorization consumed
by the Sections 9--11 repair. That repair still must derive its nonzero coset
integral and its degree-rank horizontal inclusion for the actual Taylor
quotients. In particular, it may multiply a character by a standard
denominator to kill the rational periodic contribution; it may not silently
erase that contribution. No such application is accepted by this packet.

Theorem~\ref{v15:gi:5:thm:F} applies separately at each of the finitely many required
indices $j$, including $j=1$, and to any chosen one of the four functions.
The number of applications is standard. Dense refinements and all losses
remain standard only when that finite iteration is recorded by the consumer.
The theorem supplies no global Freiman extension beyond $H'$.

The related-use audit is bounded by uses of the single-scale equivalence,
properness in D.4, composition of factorization factors, deletion of rational
Taylor contributions, and the graph/Bohr/bracket joins in F.1. The
multi-variable D.5 and its consumers in Appendix E require their own exact
application review. They are not inferred from the one-variable theorem here.
Likewise, this packet does not accept Gaussian's pending Sections 9--11
argument or any downstream assembled inverse theorem.

\section{Equal-scale factorization and the bias and torsion joins}\label{v15:gi:GI-MULTIFACTOR}
\newcommand{\GIGZ}{\mathbb Z}
\newcommand{\GIGR}{\mathbb R}
\newcommand{\GIGC}{\mathbb C}
\newcommand{\GIGT}{\mathbb T}
\newcommand{\GIGE}{\mathbb E}
\newcommand{\GIGpoly}{\operatorname{poly}}
\newcommand{\GIGLip}{\operatorname{Lip}}

This section proves the multivariable factorization needed for E.11 and E.13. The only quantitative equidistribution input is the final published equal-side theorem in [GT14, pp.1176--1177], with the published minor corrections to [GT12]. No unpublished arbitrary-side extension is imported.

\subsection{Conventions and exact statements}
Fix standard positive integers $k,d$ and an unbounded integer $N$.
All groups, rational structures, lattices, coordinate bases and finite
dimensions are standard, with natural ultrapower extensions. A nilpotent
Lie group is connected and simply connected. The rational degree filtration
$G_\bullet$ has finite length; the source $\GIGZ^k$ has its total-degree
filtration. Multiindices have nonnegative integer entries, and
$|\alpha|=\sum_i\alpha_i$. All vertical identities for a limit nilcharacter
hold at every limit group element and point, as in the accepted
tautification-003 convention. We do not substitute equivariance only at
standard points.

In fixed Mal'cev coordinates, small polynomial coefficients mean bounded
constant coefficient and coefficient $O(N^{-|\alpha|})$ at positive index
$\alpha$ in either the monomial or product-binomial basis. The constants
are standard and may depend on the actual fixed geometric data and orbit.
Total equidistribution on $[N]^k$ means equidistribution on every product
of dense arithmetic progressions in $[N]$, with standard Lipschitz tests.
It also applies to internal tests with a common standard Lipschitz bound:
finite uniform nets in a compact Lipschitz ball reduce those tests to
standard ones, with the desired accuracy chosen before taking the limit.

\begin{theorem}[Equal-scale factorization, MF-D4]\label{v15:gi:6:thm:factor}
For $g\in{}^*\GIGpoly(\GIGZ^k,G_\bullet)$ there are a standard connected rational
subgroup $P\leq G$ and internal polynomial maps $\epsilon,p,\gamma$ such that
\[
                       g(n)=\epsilon(n)p(n)\gamma(n)
                         \qquad(n\in{}^*\GIGZ^k).                 \tag{1}
\]
The map $p$ lies in $P$ with its induced filtration, and is totally
equidistributed on $[N]^k$ in $P/\Delta$ for every standard finite-index
$\Delta\leq P\cap\Gamma$, with the same $p$ and $P$.
The coefficients of $\epsilon$ are small as specified above; it is bounded
and $O(1/N)$-Lipschitz on every fixed standard dilate $[-AN,AN]^k$.
Its group Taylor coefficient at positive index $\alpha$ is
$O(N^{-|\alpha|})$ as well. There is a standard discrete subgroup
$\Lambda\geq\Gamma$ containing all values of $\gamma$, and a standard
$q>0$ such that
\[
           \gamma(n+qe_i)\Gamma=\gamma(n)\Gamma
                  \qquad(n\in{}^*\GIGZ^k, 1\leq i\leq k).       \tag{2}
\]
There is no properness requirement on $P$, and no claim that $\gamma$ is
$\Gamma$-valued. All factors in (1) are defined on the full integer lattice.
\end{theorem}

Write $I=[-N,N]\cap{}^*\GIGZ$.
\begin{theorem}[Actual bias and torsion outputs, MF-E11/E13]\label{v15:gi:6:thm:outputs}
Let $d\geq2$ be standard, and let $\chi$ be a globally represented unit
vector nilcharacter of degree $d$ on ${}^*\GIGZ^k$.
If a scalar coordinate of $\chi$ has nonzero correlation with a degree
at most $d-1$ nilsequence after multiplication by the indicator of a
convex polytope in $I^k$ and a product of dense progressions in $[N]$,
then $\chi$ agrees on all of $I^k$ with a globally represented
degree-at-most-$(d-1)$ vector nilsequence. The same conclusion holds if
$\chi^{\otimes a}$ agrees on $I^k$ with a degree-at-most-$(d-1)$ vector
nilsequence for a standard integer $a\geq1$.

These assertions give exactly the lower-degree conclusions of [GTZ, E.11]
and [GTZ, E.13] on their required symmetric box. They include finite
scalar-coordinate reductions for vector correlations. They do not require
an extrapolation from an unrelated subprogression or the deletion of a
rational Taylor coefficient.
\end{theorem}

\subsection{A weighted inverse-affine estimate}
The following estimate also supplies the large-offset transition behind
[GT12, Proposition 9.2]. Merely inserting an offset of size $N$ into the
height bound of its Lemma 8.4 would not give the required standard bound.

\begin{lemma}[Large offsets, bounded steps, MF-AFFINE]\label{v15:gi:6:lem:affine}
Let $f(n)$ be a real polynomial of fixed total degree at most $D$ in $k$
variables. Let $A_i$ be internal integers with $|A_i|=O(N)$ and let $b_i$
be nonzero integers of bounded absolute value. Put
$\widetilde f(m)=f(A_1+b_1m_1,\ldots,A_k+b_km_k)$.
Suppose every positive product-binomial coefficient $\widetilde f_\beta$
satisfies $\|\widetilde f_\beta\|_\GIGT=O(N^{-|\beta|})$.
There is a standard positive integer $r$, depending only on $D,k$ and the
step bounds, such that every positive coefficient of $rf$ has the same
bound $O(N^{-|\alpha|})$ modulo one.
\end{lemma}
\begin{proof}
Invert the substitution over $\GIGR$ and expand
\[
 \prod_i\binom{(n_i-A_i)/b_i}{\beta_i}
        =\sum_{\alpha\leq\beta}c_{\alpha\beta}\binom n\alpha.
                                                               \tag{3}
\]
Expanding each falling factorial first in monomials shows that the
denominators divide a fixed integer formed from $\beta_i!$ and
$b_i^{\beta_i}$. Conversion of ordinary powers to binomial coefficients
uses integer Stirling coefficients. Thus one standard $r$ clears all
denominators, independently of the integer numerators $A_i$.
The coefficient of $n^\nu$ before the last conversion is a polynomial in
the $A_i$ of total degree at most $|\beta|-|\nu|$. Consequently
\[
 |c_{\alpha\beta}|=O(N^{|\beta|-|\alpha|}),
 \qquad rc_{\alpha\beta}\in{}^*\GIGZ.                             \tag{4}
\]
The case $\beta=0$ contributes only to $\alpha=0$.
For each $\beta>0$ write $\widetilde f_\beta=z_\beta+e_\beta$, where
$z_\beta$ is an integer and $e_\beta=O(N^{-|\beta|})$.
In the coefficient of index $\alpha>0$ in $rf$, the contributions
$rc_{\alpha\beta}z_\beta$ are integers, while
\[
 \sum_{\beta\geq\alpha}rc_{\alpha\beta}e_\beta
                         =O(N^{-|\alpha|}).                    \tag{5}
\]
There are only a fixed number of indices. This proves the claim. Negative
steps cause no change after taking a positive common denominator. The
uncontrolled constant coefficient causes no positive-index error.
\end{proof}

\subsection{From a dense rectangle to the published equal-side theorem}
\begin{lemma}[Weak factorization in finitely many parameters]\label{v15:gi:6:lem:weak}
Suppose $G_0=G_1=G$. If $g\Gamma$ is not totally equidistributed on $[N]^k$,
it has a factorization (1) with $p$ in a standard proper connected rational
subgroup, retaining all conclusions of Theorem~\ref{v15:gi:6:thm:factor} except
the final equidistribution condition.
\end{lemma}
\begin{proof}
Failure supplies a product of progressions
$P_i=\{a_i+b_it:1\leq t\leq L_i\}\subseteq[N]$ with $L_i\geq\delta N$,
standard $\delta>0$, and a bounded Lipschitz test with noninfinitesimal
discrepancy. Their integer steps obey $1\leq b_i\leq2/\delta$ and their
offsets are $O(N)$. Subtract the Haar mean, normalize the test and rotate a
constant complex phase so that the real mean discrepancy is positive.

In the index rectangle $\prod[ L_i]$, partition each axis into blocks of a
common length $L=\lfloor\eta N\rfloor$, leaving a final remainder of length
less than $L$. Choose the standard $\eta>0$ small compared with the
discrepancy, $\delta$, $k$ and the test norm. The discarded proportion is
at most $\sum_iL/L_i\leq k\eta/\delta+o(1)$, so its contribution is less
than one quarter of that discrepancy. At least one full cube of side $L$
has a positive standard discrepancy. On this cube the polynomial is
$\widetilde g(m)=g(A_1+b_1m_1,\ldots,A_k+b_km_k)$ with $|A_i|=O(N)$.

Apply the exact equal-side theorem of [GT14, pp.1176--1177] at length $L$,
using a standard error parameter below this discrepancy and below the
reciprocal rationality bound of the fixed adapted basis. The result gives
a nonzero bounded horizontal integer character $\xi$ for which all positive
coefficients of the real horizontal lift of $\xi\widetilde g$ have circle
norm $O(L^{-|\alpha|})=O(N^{-|\alpha|})$.
Only finitely many such characters are possible, so the selected character
is standard. Lemma~\ref{v15:gi:6:lem:affine} replaces it by a nonzero standard
multiple and gives the same coefficient bounds for $\xi g$ on the
original variables. This is the precise rectangle-to-equal-cube join;
the original unequal-side theorem is not used.

Normalize $g(0)=c\lambda$ with bounded $c$ and $\lambda\in{}^*\Gamma$
and apply the preceding argument to $c^{-1}g\lambda^{-1}$, whose failure
is preserved by bounded left translation. Its coordinates are
$\sum_{\alpha>0}t_\alpha\binom n\alpha$, with $t_\alpha$ supported in
$G_{|\alpha|}$ by [GT12, Lemma 6.7]. Write the real horizontal lift as
$u\cdot\psi$, with $u$ a fixed integer vector.
Choose integers $a_\alpha$ nearest to $u\cdot t_\alpha$.
On a support where $u$ is nonzero choose a fixed rational vector
$v_\alpha$ of that support with $u\cdot v_\alpha=1$, and define
\[
 w_\alpha=(u\cdot t_\alpha-a_\alpha)v_\alpha,
 \quad z_\alpha=a_\alpha v_\alpha,
 \quad \psi(\epsilon_0(n))=\sum w_\alpha\binom n\alpha,
 \quad \psi(\gamma_0(n))=\sum z_\alpha\binom n\alpha.           \tag{6}
\]
If $u$ vanishes on a support, set all these vectors and its integer to zero.
Coordinate support and polynomial group closure make these maps polynomial
for the original filtration. The real lift is identically zero on
$p=\epsilon_0^{-1}c^{-1}g\lambda^{-1}\gamma_0^{-1}$.
Thus $p$ lies in its standard proper connected rational kernel.
The original factors are $\epsilon=c\epsilon_0$ and
$\gamma=\gamma_0\lambda$. The $w_\alpha$ have precisely the required
small bounds; the $z_\alpha$ have a common standard denominator, without
a numerator bound. Controlled products and periods are checked next.
\end{proof}

\subsection{Controlled products, periods, and descent}
Scale all variables by $n_i=Nt_i$. Fixed coordinate polynomials with small
coefficients become polynomials in $t$ with bounded coefficients. The
fixed BCH and coordinate-change polynomials preserve this property under
finite products, inverses and fixed rational subgroup embeddings.
Conversion between monomial and product-binomial bases is triangular in
componentwise indices, so the coefficient of $\binom n\alpha$ is
$O(N^{-|\alpha|})$. Differentiating the scaled coordinate maps on the
compact box $[-A,A]^k$ gives the claimed $O(1/N)$ Lipschitz bound on every
fixed dilate. A real-coordinate polynomial extension of $\epsilon$ has the
same bounds, which will be used below.

These estimates also imply the group Taylor bounds. Assign formal weight
$|\alpha|$ to each logarithmic coefficient of $n^\alpha$ and weight zero to
the bounded constant term. Group Taylor extraction in any fixed admissible
order uses fixed BCH polynomials. Inductively every term of formal weight
$w$ has coordinate polynomial total degree at most $w$; removing a Taylor
factor at index $\beta$ preserves this property because its binomial power
has degree $|\beta|$. The coefficient at index $\alpha$ cannot contain a
term of formal weight less than $|\alpha|$: such a term has total degree
less than $|\alpha|$ and has zero coefficient at that index. This is the
same triangular Taylor extraction as [GTZ, B.9], now with product-binomial
indices. Each surviving BCH monomial contributes at least
$N^{-|\alpha|}$ decay. In particular no estimate is inferred solely from
separate degree in individual variables.

The common discrete overgroup argument is independent of the number of
parameters. In a rational lower-central basis of weights $w_i$, choose a
standard $D$ clearing every coordinate denominator and every fixed BCH
coefficient. As in DF002, the additive lattice
$\sum_iD^{-(2w_i-1)}\GIGZ X_i$ is BCH-closed: a nonlinear term with $t$
inputs has output denominator exponent exceeding the input sum by at least
$t-1$. Its exponential is a standard discrete overgroup $\Lambda\geq\Gamma$
containing the rational factors and any finite product of them.

Here is the full multivariable period argument, in the right-coset direction.
Each coordinate of $\gamma(n)$ is a polynomial of fixed degree whose values
have a common standard denominator. Tensor-product interpolation on a fixed
integer grid gives a standard common denominator for all its coefficients.
For every $i$ the vector
\[
 R_i(n,r)=\psi\bigl(\gamma(n)^{-1}\gamma(n+re_i)\bigr)
\]
is a rational coefficient polynomial with a common standard denominator.
It is identically zero when $r=0$. Choose a standard positive $q$ divisible
by all these denominators. Each monomial then has a positive power of $r$,
so $R_i(n,q)$ is integral for every internal integer vector $n$.
This is exactly (2). This proof does not use the wrongly oriented quotient
in the printed proof of [GT12, A.12]. The corresponding one-variable
sublemma of DF002 has a separate independent root review.

To finish Theorem~\ref{v15:gi:6:thm:factor}, first normalize the constant term and
restrict to $G_1$ if $G_1\ne G$. The identity
\[
 c^{-1}g(n)\lambda^{-1}
          =\lambda\bigl(g(0)^{-1}g(n)\bigr)\lambda^{-1}\in G_1
\]
holds in any number of variables, since first derivatives lie in $G_1$ and
the positive groups are normal in $G_0$. Thus a controlled representation
exists with a taut positive filtration.
Choose a standard rational representing subgroup of least dimension.
For any standard finite-index lattice $\Delta\leq P\cap\Gamma$, failure
of total equidistribution in $P/\Delta$ permits Lemma~\ref{v15:gi:6:lem:weak}, after
the same normalization if its induced first group is proper.
It supplies $p=\epsilon_2p_2\gamma_2$ in a proper rational $P_2<P$.
The correct new identity is
\[
                 g=(\epsilon\epsilon_2)p_2(\gamma_2\gamma).
\]
The preceding controlled-product and discrete-overgroup arguments verify
every required bound and period. Rationality relative to $\Delta$ is the
same rational structure as relative to $P\cap\Gamma$.
This contradicts minimal dimension. It proves total equidistribution for
every standard finite-index cover with the same factors. Only a minimum
among $0,\ldots,\dim G$ was selected; no infinite sequence of refinements
or dependence of the chosen factors on the later cover is hidden here.

\subsection{One quotient for all periodic residue classes}
We now work in a taut ambient degree-$d$ group $G$, so $G_d$ is central.
Tautification-003, Theorem 1, supplies this exact representation for each
nilsequence used below with the same positive-degree groups and the same
standard vertical character. Products of these representations remain taut.

Apply Theorem~\ref{v15:gi:6:thm:factor}, and choose a standard representative
$\lambda_\rho\in\Lambda$ for each value of the right coset of $\gamma(n)$
on a residue class $n\equiv\rho\pmod q$, $\rho\in(\GIGZ/q\GIGZ)^k$.
Such representatives can be standard: $\Lambda/\Gamma$ is a fixed finite
set, and periodicity makes the coset constant on each class. Let
\[
 \Gamma_0=\bigcap_{\lambda\in\Lambda}\lambda\Gamma\lambda^{-1},
 \qquad \Delta=P\cap\Gamma_0.                                \tag{7}
\]
The first intersection is the kernel of the action on the finite coset set,
so it is a standard normal finite-index subgroup of $\Lambda$ contained in
$\Gamma$. Thus $\Delta$ is finite index in $P\cap\Gamma$, and $p$ is
totally equidistributed on $P/\Delta$. For every $\rho$ the function
\[
                      x\longmapsto F(zx\lambda_\rho\Gamma)   \tag{8}
\]
is well-defined on this \emph{same} quotient: for $\delta\in\Delta$,
$\lambda_\rho^{-1}\delta\lambda_\rho\in\Gamma$.
No deletion of a rational horizontal coefficient is involved.

\begin{lemma}[A nonzero integral from the actual cutoff]\label{v15:gi:6:lem:integral}
Let $F$ be an internal bounded Lipschitz scalar test of standard norm on
$G/\Gamma$. If
\[
 \left|\GIGE_{n\in I^k}1_\Omega(n)1_{P_1\times\cdots\times P_k}(n)
                           F(g(n)\Gamma)\right|>c>0           \tag{9}
\]
for a convex polytope $\Omega\subseteq I^k$ and dense progressions
$P_i\subseteq[N]$, then for some residue $\rho$ and bounded
$z=\epsilon(n_0)$ one has
\[
          \left|\int_{P/\Delta}F(zx\lambda_\rho\Gamma)
                                           \,d\mu(x)\right|>c'>0. \tag{10}
\]
Here $c'$ is standard and positive. The factors and quotient are those
just constructed for the same actual $g$.
\end{lemma}
\begin{proof}
Partition the progression product into its finitely many nonempty residue
classes modulo $q$. Each intersection is again a product of dense
progressions, by the elementary congruence calculation for each axis.
Their steps are bounded by the old step times $q$. Rotate a common complex
phase in (9); one class retains a real sum at least $c_0N^k$ for a
standard $c_0>0$. Fix this lower bound before selecting a mesh.

Use index coordinates on that progression product. The inverse image of
$\Omega$ intersected with the index rectangle is convex and lies in a box
of side $O(N)$. Partition it by cubes of side $\lfloor\eta N\rfloor$.
The cubes meeting its boundary have at most $O_k(\eta N^k)+O_k(N^{k-1})$
integer points, with constants also depending on the fixed step bounds.
Indeed place disjoint unit cubes at those integer points and enclose them
in the $O_k(\eta N+1)$ neighborhood of the convex boundary; apply the
published boundary-volume estimate [GT10, Corollary A.2, p.1810].
For a convex set with empty interior, first thicken by an arbitrarily small
ball and then take the decreasing limit of its boundary neighborhoods.
Choose standard $\eta$ small enough that the discarded boundary
contribution is below $c_0N^k/4$. The total number of points in the full
cubes is at most $N^k$. Some cube contained in the convex set therefore
gives a dense progression product $R$ whose phase-adjusted average is at
least $3c_0/4$.

The choice of $\eta$ can also ensure that
$|F(\epsilon(n)x\lambda_\rho\Gamma)-F(\epsilon(n_0)x\lambda_\rho\Gamma)|$
is below $c_0/4$ for all $n,n_0\in R$ and all
$x\in P/\Delta$. This follows from the $O(1/N)$ variation of $\epsilon$,
its bounded image, the fixed standard representative $\lambda_\rho$, and
compactness of $P/\Delta$. Choose the mesh only after these standard
constants and the original positive discrepancy are known.
For $n\in R$, (1) and the fixed right coset give
$F(g(n)\Gamma)=F(\epsilon(n)p(n)\lambda_\rho\Gamma)$.
Freeze $\epsilon$ at $n_0$. Total equidistribution of the same $p$ in
$P/\Delta$ on $R$ identifies the resulting average with the integral in
(10), up to an infinitesimal error. Internal-test uniformity follows from
the finite Lipschitz net specified at the start. This proves (10).
\end{proof}

\subsection{Exact reconstruction on the symmetric box}
Let $G=G^{\chi}\times G^0$ be the taut product of a degree-$d$ nilcharacter
representation and any degree-at-most-$(d-1)$ auxiliary representation.
The latter has trivial degree-$d$ group. Write $F_\chi$ for the original
unit vector test, and let $\eta$ be its real standard vertical character.
Set
\[
                K=P_d=P\cap G_d.
\]
We justify the rational intersection with the original lattice explicitly.
Fix a rational Mal'cev basis for $G/\Gamma$. A connected subgroup $H$
with compact $H/(H\cap\Gamma)$ has rational Lie algebra in this basis.
Indeed, [GT12, Lemma A.2] makes the coordinates of $\log\gamma$ rational
for $\gamma\in\Gamma$. The Lie algebra $\mathfrak l$ generated by
$\log(H\cap\Gamma)$ therefore has a rational basis, by selecting finitely
many independent iterated brackets. If $\mathfrak l$ were proper in
$\mathfrak h$, put it in a maximal proper subalgebra $\mathfrak m$.
The last lower central term not contained in $\mathfrak m$ supplies a
vector outside $\mathfrak m$ normalizing it. Thus maximality makes
$\mathfrak m$ an ideal, and $\mathfrak h/\mathfrak m$ has dimension one,
since otherwise a line would be a proper nonzero subalgebra. The quotient
functional exponentiates by BCH to a surjection $H\to\mathbb R$ killing
$H\cap\Gamma$, contradicting compactness. Hence $\mathfrak l=\mathfrak h$.
Apply this argument to $P$ and $G_d$. Their Lie-algebra intersection has
a rational basis by rational linear algebra. The global exponential
homeomorphism and BCH identify $K$ with its exponential, so $K$ is closed
and connected. Equip $K$ with its rational lower central filtration.
Choose a standard $Q$ bounding the finitely many rational coordinates.
[GT12, Definition 2.5 and Proposition A.10] then give a Mal'cev basis for
$K/(K\cap\Gamma)$, so this actual intersection is a lattice. The trivial
subgroup case is immediate. Since $G_d$ is central in the taut ambient
group, $K$ is central too.
The coset map from $(K\cap\Gamma)/(K\cap\Delta)$ to
$(P\cap\Gamma)/\Delta$ is injective. Thus $K\cap\Delta$ has finite index
in $K\cap\Gamma$ and is a lattice in $K$, without requiring normality
of $\Delta$. The central-quotient argument in LP13-002, Section 2, now
applies with its rationality premise established: the image
$\overline\Delta$ of $\Delta$ is a lattice in $\overline P=P/K$.
The quotient filtration kills $P_d=K$ and hence has degree at most $d-1$.
All these objects depend only on the fixed standard geometry.

Suppose that
\[
                  \eta(\operatorname{pr}_\chi K)=0.            \tag{11}
\]
For each $\rho$ and bounded real-coordinate parameter $t\in[-1/5,1/5]^k$,
define, initially on $P/\Delta$,
\[
 H_\rho(x,t)=F_\chi\bigl(
   \widetilde\epsilon_\chi(10Nt)\,x_\chi\,(\lambda_\rho)_\chi
                                                    \Gamma_\chi\bigr), \tag{12}
\]
where $\widetilde\epsilon$ is the real-coordinate polynomial extension.
This is well-defined by (7). It is invariant under $K$ by the all-limit
vertical identity and (11); centrality permits its element to pass across
$\widetilde\epsilon$ and $\lambda_\rho$. Hence (12) descends to
$\overline P/\overline\Delta$.
The scaled-coefficient estimates imply a standard Lipschitz bound in $t$,
uniformly in $N$. Smooth multiplication on fixed compact ranges gives the
same conclusion in $x$. Invariance descends this bound to the quotient
metric; any fixed standard Riemannian metric is equivalent on this compact
nilmanifold.

Choose a fixed Lipschitz function $B(t)$ on $\GIGT^k$ equal to one on
$[-1/10,1/10]^k$ and zero outside $(-1/5,1/5)^k$.
Using centered representatives inside that chart, define
$\widetilde H_\rho(x,t)=B(t)H_\rho(x,t)$ there and zero outside.
It is a globally bounded Lipschitz vector function. The zero collar makes
the definition independent of the chart boundary. The functions are limit
functions with a standard bound, since their coordinate-level construction
has uniform bounds before taking the ultrapower.

For each $\rho\in(\GIGZ/q\GIGZ)^k$ choose a standard Lipschitz function
$b_\rho$ on $\GIGT^k$ with
$b_\rho(\sigma/q)=1_{\rho=\sigma}$. For example use products of triangular
bumps of radius $1/q$; for $q=1$ use the constant function one.
The vector function
\[
 \mathcal H(x,t,z)=\sum_\rho b_\rho(z)\widetilde H_\rho(x,t)    \tag{13}
\]
is globally bounded Lipschitz with a standard bound. Its orbit is
\[
          n\longmapsto\bigl(\overline p(n)\overline\Delta,
                                      n/(10N),\,n/q\bigr).    \tag{14}
\]
This orbit has degree at most $d-1$: its first component has that degree
and both tori have degree one, with $1\leq d-1$.
For every $n\in I^k$, exactly the correct residue value is sampled in
(13), and $B(n/(10N))=1$. The identities (1), (2), and (12) give
\[
                        \chi(n)=\mathcal H((14)).              \tag{15}
\]
This is an equality of the actual vector functions on the entire symmetric
box. It is not an extrapolation of an equality on just one progression.
No unit condition is imposed on $\mathcal H$ away from the orbit; a
lower-degree vector nilsequence does not require one.

\subsection{Bias and torsion applications on the same objects}
\begin{proof}[Proof of Theorem~\ref{v15:gi:6:thm:outputs}]
For the bias assertion, choose the scalar coordinate and the lower-degree
scalar test occurring in the actual nonzero correlation. If a vector
correlation is given, a fixed finite coordinate choice supplies one such
scalar nonzero average; its density loss is standard. Form the taut product
representation $G^\chi\times G^0$ and the product test
$F= (F_\chi)_i\overline{F_0}$, using the sign in that correlation.
It has vertical character $\eta$ on the degree-$d$ group because $G^0_d$
is trivial. Apply Theorem~\ref{v15:gi:6:thm:factor} and Lemma~\ref{v15:gi:6:lem:integral} to
this actual product orbit and test, obtaining a nonzero integral (10).

Translation of Haar measure on $P/\Delta$ by any $u\in K$ preserves that
integral. The integrand is multiplied by
$e(\eta(\operatorname{pr}_\chi u))$, since $K$ is central in the ambient
taut group. The integral is nonzero, so this phase is one for every
standard $u\in K$. The real homomorphism $\eta$ is continuous and $K$ is
connected; its image is therefore a connected subset of $\GIGZ$, hence zero.
Its natural ultrapower extension vanishes on every limit element too.
This is (11). Equations (12)--(15), with the \emph{original} vector test
$F_\chi$, now prove the asserted all-box lower-degree representation.

For the torsion assertion, let $\Psi$ be the given degree-at-most-$(d-1)$
vector nilsequence equal to $\chi^{\otimes a}$ on $I^k$.
Put its finitely many coordinates into one common product nilmanifold of
that degree. Since $\chi$ is a unit vector, pointwise on $[N]^k$ one has
\[
              \sum_j (\chi^{\otimes a})_j(n)\overline{\Psi_j(n)}=1.
\]
At least one scalar term has a nonzero standard average, with lower bound
at least the reciprocal of the fixed tensor dimension in absolute value.
Use that term as the test on the product of the original $\chi$ geometry
and the auxiliary geometry of $\Psi$. Its top vertical character is
$a\eta$. Lemma~\ref{v15:gi:6:lem:integral} applies with the full positive cube and
no nontrivial cutoff. The same Haar argument gives
$a\eta(\operatorname{pr}_\chi K)=0$ as a real homomorphism.
As $a\geq1$, (11) follows. The same reconstruction (12)--(15) uses the
original $F_\chi$, not its tensor power, and proves the required conclusion.

In the symbol formulation of [GTZ, E.13],
$a[\chi]=0$ means exactly that the tensor power is equivalent to the
constant nilcharacter. By the definition of equivalence, its coordinates
are lower-degree nilsequences on the specified box. This is precisely the
$\Psi$ premise just used, with a finite common representation.
No torsion-free conclusion in degree one is asserted; the time and residue
torus factors explain the requirement $d\geq2$.
\end{proof}

\subsection{Scope of the factorization}
This proof supplies the finite-parameter factorization actually needed by
E.11/E.13 and their all-symmetric-box reconstruction. It does not assert
the false horizontal Taylor closeness printed in D.5. The one-variable D.6 argument must retain its denominator correction in the
degree-rank horizontal quotient.
The inverse-affine Lemma~\ref{v15:gi:6:lem:affine} also validates the large-offset
coefficient transition in the one-variable weak-factorization source.
The affine-transfer application and the E.11/E.13 substitutions are made explicitly in Section~\ref{v15:gi:ASSEMBLY}.


\section{Sections 9--11 at every rank step}\label{v15:gi:GI-SECTION9-11}
\newcommand{\GIHR}{\mathbb R}
\newcommand{\GIHZ}{\mathbb Z}
\newcommand{\GIHQ}{\mathbb Q}
\newcommand{\GIHC}{\mathbb C}
\newcommand{\GIHT}{\mathbb T}
\newcommand{\GIHN}{\mathbb N}
\newcommand{\GIHee}{\mathrm e}
\newcommand{\GIHDR}{\mathrm{DR}}
\newcommand{\GIHHor}{\operatorname{Horiz}}
\newcommand{\GIHTay}{\operatorname{Taylor}}


This section separates two analytic hypotheses, \textup{F} and \textup{D}, from the algebraic, selection and representation arguments proved here. Their supplying results are proved in Sections~\ref{v15:gi:GI-DF} and \ref{v15:gi:GI-MULTIFACTOR}. The universal group uses the explicit fixed-rank weighted presentation below. No final inverse theorem is used as a premise.

\subsection{The family at an arbitrary rank}
Fix a standard $s\geq3$, $m=s-1$ and $1\leq R\leq m$, and an unbounded
limit integer $N$. All groups, lattices, metrics and finite dimensions below
are standard unless a natural ultrapower is displayed or understood.
Braces use $(-1/2,1/2]$ throughout this packet. The Section 12 constructor
allows either fixed half-open fundamental cube; use this same choice there.
Let $H\subseteq[-N,N]\cap{}^*\GIHZ$ be dense and internal. The input consists of
a bounded function $f$ on $[N]$, a global unit-vector multidegree $(1,m)$
nilcharacter $A(h,n)$, and an internal family of unit-vector degree--rank
$(m,R)$ nilcharacters $\chi_h$. The latter have a fixed finite list of rational
models and standard uniform Lipschitz bounds. For every $h\in H$,
$\Delta_h f$ has uniformly positive degree-at-most-$m-1$ correlation with
$A(h,\cdot)\otimes\chi_h$, with bounded internal witnessing tests.
Every vertical identity is on the full natural extension of its group and
compact quotient, not merely their standard points.

This is exactly the new input in rank-iteration001, Section 3. In particular
no core/petal representation or independence statement is inherited from
the preceding rank. They will be reconstructed. The corrected LP03 gives,
after any dense internal restriction of $H$, many additive quadruples
$h_1+h_2=h_3+h_4$ for which
\begin{equation}\label{v15:gi:7:eq:quad}
 \chi_{h_1}(n)\otimes\chi_{h_2}(n+h_1-h_4)
 \otimes\overline{\chi_{h_3}(n)}
 \otimes\overline{\chi_{h_4}(n+h_1-h_4)}
\end{equation}
is $(m-1)$-biased, with standard uniform witnesses. This hereditary statement
means reapplying LP03 to the restricted correlation hypothesis. It does not
assert that a previously chosen dense set of quadruples survives restriction.

\paragraph{Analytic contract F.}
For internal functions $X_1,\ldots,X_4:H\to{}^*\GIHT$ and $\epsilon=N^{-j}$,
if $X_1(h_1)+\cdots+X_4(h_4)=O(\epsilon)$ on many additive quadruples,
there are a dense internal $H'$, a standard finite $K$, and internal
$\alpha_k\in{}^*\GIHT$, $\beta_k\in{}^*\GIHR$, $\delta\in{}^*\GIHT$ such that
\[
 X_4(h)=\sum_{k=1}^K\{\alpha_k h\}\beta_k+\delta+O(\epsilon)
 \quad(h\in H').
\]
The errors have a common standard bound. A permutation of the four positions
changes which function is recovered. This is the precise Appendix F.1 input;
the signed-graph and no-wrap argument is a separate analytic supplier.

\paragraph{Analytic contract D.}
For a polynomial sequence in one variable in a fixed rational filtered
nilmanifold there is a factorization $g=\varepsilon g'\gamma$ with
$\varepsilon$ bounded, Lipschitz $O(1/N)$ and positive-index Taylor
coefficients $O(N^{-i})$; $g'$ takes values in a connected rational subgroup
$Q$ and is totally equidistributed on every standard finite cover of
$Q/(Q\cap\Gamma)$; and $\gamma$ takes values in a fixed discrete overgroup
$\widehat\Gamma\supseteq\Gamma$ and is periodic modulo $\Gamma$ with standard
period. The subgroup $Q$ is allowed to equal $G$. All factors remain adapted
to the ordinary degree filtration. Only this one-variable contract is needed.
Its quantitative published supplier and corrected proof are separate.

\begin{theorem}\label{v15:gi:7:thm:output}
Subject to these two analytic contracts and the stated LP03 input, a dense
internal subfamily has the following representation. There are fixed standard
dimension blocks $C,L,B$, fixed filtered lattice-compatible maps from a
weighted universal group $U$, and one integral top character $\eta$ on
$U_{(m,R)}$. Up to a fixed coordinate isometry $\iota$,
\[
 \iota\chi_h(n)=F_h(g_h(n)\Gamma,g_{0,h}(n)\Gamma_0)
 \qquad(n\in[N]).
\]
The lower model has degree--rank at most $(m,R-1)$. The functions $F_h$ are
unit vectors with uniform standard bounds and exact top-character law;
they need not be equal as functions. The horizontal coefficients are fixed
core frequencies, linear petals $\{\alpha_a h\}\beta_a$ for $a\in L$, and
independent petals for $a\in B$. The linear petals are Freiman homomorphisms
on the retained set. For almost every additive quadruple, the core, all four
independent-petal blocks and any three linear-petal blocks are independent
modulo $O(N^{-i})$ separately at each weight $i$.

For every weight-$m$, length-$R$ iterated commutator $w$, the real number
$\eta(w)$ is zero if $w$ contains two petals, or contains one letter from
$B$. These are precisely the annihilations required by LP04. They hold for
every new rank family satisfying the initial hypotheses, without inherited
frequency information. The construction of the final two-variable test is
the separate Section 12 join.
\end{theorem}

\subsection{Selection and the corrected universal representation}
We use the following elementary internal-choice rule, which also makes the
countability uses precise. For an increasing internal sequence of witness
sets $S_k$, suppose each member of an internal parameter set has a witness
in some standard $S_k$. The least successful index is an internal integer.
It is uniformly bounded: otherwise the internal nonempty sets on which it
exceeds $1,2,\ldots$ have a common element by countable saturation, contrary
to pointwise standard boundedness. At representative indices choose the
first witness in a fixed choice order. This gives internal choices with one
standard bound. Countable saturation here follows by diagonalizing nested
finite intersections: at representative index $v$, choose a witness for the
largest $k\leq v$ whose first $k$ conditions hold there. This $k$ tends to
infinity along the ultrafilter. The unsupported diagonal $x_{v,v}$ alone
would not establish this conclusion.

We also need the ``almost all'' form. For nested internal bad sets $B_k$
in a hyperfinite set $X$, if every $B_k$ has infinitesimal relative size,
the same diagonal construction chooses an unbounded $K$ for which $B_K$
has infinitesimal relative size and contains every standard $B_k$.
Consequently, if their countable union is not contained in a sparse internal
set, some $B_k$ is dense. Applying finite pigeonhole inside that $B_k$ fixes
an integer coefficient vector and a standard error bound. This is the exact
quantifier step used for failed frequency independence below.

First restrict to one of the fixed models for $\chi_h$. Bounded
tautification and normalization use their separately proved project
interfaces. In a fixed model, only finitely many integral vertical characters
can occur under a common Lipschitz bound. Indeed, along a fixed unit-speed
top-torus direction $X$, the identity
$F(\exp(tX)x)=\GIHee(t\eta(X))F(x)$ and $\|F(x)\|=1$ give
$2\pi|\eta(X)|\leq C_X\operatorname{Lip}(F)$ by taking small $t$.
On an integral basis this bounds all integer character coordinates. Finite
pigeonhole therefore fixes the character.

For each weight $i$, choose finitely many elements of $\Gamma\cap G_{(i,0)}$
whose logarithms generate its rational Lie algebra. Such elements are obtained
by clearing denominators in an adapted rational Mal'cev basis; their real
generated subgroup is $G_{(i,0)}$. Use one formal letter of weight $i$ for
each. Define $U$ by killing all Lie words of weight greater than $m$ and all
weight-$m$ words of length greater than $R$. Set
\[
 U_{(d,r)}=\exp\operatorname{span}\{\text{surviving words of weight}>d,
 \text{ or weight}=d\text{ and length}\geq\max(r,1)\}
\]
for positive $d$, with $U_{(0,0)}=U$. This is the exact LP04 weighted
presentation before the additional petal relations. Its Hall-coordinate
lemma supplies the lattice generated by the letters and its filtered
intersections. The upward-closed word conditions prove every bracket
inclusion. In particular $U_{(d,0)}=U_{(d,1)}$.

Mapping each letter to its chosen lattice element gives a filtered
homomorphism $\phi:U\to G$. The killed words vanish by the degree--rank
bracket law in $G$; the map sends the generated lattice into $\Gamma$.
It is onto each $G_{(i,0)}$. In
$\GIHHor_i(G)=G_{(i,1)}/G_{(i,2)}$ the images of the weight-$i$ letters span
the real vector space and generate a full sublattice. Thus the induced
map of compact horizontal tori is onto. Select internal preimages of
$\GIHTay_i(g_h\Gamma)$ to get frequencies $\xi_{h,i,a}\in{}^*\GIHT$.
Pull the fixed top character back by $\phi$ and initially take $G_0$ trivial.
This proves the exact Definition 9.11 representation, including lattice
compatibility and all-limit equivariance, for every new family.

Here and below horizontal Taylor coefficients are homomorphisms on the
polynomial-sequence group, are unchanged by translating $n$, and for
$c\in G_{(i,1)}$ the sequence $c^{\binom ni}$ has only the $i$th horizontal
coefficient $c$ and zero others. These facts use the iterated-difference
commutator rule B.7: all noncommutative correction terms have rank at least
two and disappear modulo $G_{(i,2)}$; an extra difference has weight above
$i$. Binomial differences give the last assertion. Polynomial products use
B.4. No claim that arbitrary Taylor factors commute is used.

If $\xi_{i,a}=\sum_b c_{i,a,b}\xi'_{i,b}$ with standard integers $c_{i,a,b}$,
map the new weight-$i$ letter $e'_{i,b}$ to
$\prod_a e_{i,a}^{c_{i,a,b}}$ in a fixed order. The universal relations and
lattice inclusion hold, giving $\psi:U'\to U$. In the abelian horizontal
quotient this is exactly the indicated integer matrix. Replace $\phi$ by
$\phi\psi$ and $\eta$ by $\eta\psi$; the original orbit and test are unchanged.
When these matrices depend on $h$, the preceding bounded-choice rule and
finite pigeonhole fix them on a dense internal subfamily before this step.

\subsection{Frequency regularization with the actual quantifiers}
Fix one weight and put $\epsilon=N^{-i}$. A partial description consists of
fixed core frequencies, independent-candidate petals $b_h$, bracket-linear
petals $l_h$, fixed rational frequencies, and $O(\epsilon)$ errors, with all
original frequencies standard integer combinations of them. Every list is
standard finite, and every family is internal. Dense restriction fixes the
integer matrices and error bounds. Put all original frequencies in $b_h$
to obtain an initial partial description.

Minimize lexicographically the numbers of $b$ petals, then $l$ petals,
then core frequencies, among all such partial descriptions on dense
subsets. Bracket-linear petals are required to be individually Freiman
homomorphisms. This costs only a finite dense restriction: place each
$\{\alpha h\}$ in one of a fixed finite collection of intervals of width
less than $1/4$. For an additive quadruple the difference of the four
fractional parts is an integer of absolute value less than one, hence zero.
Multiplication by any real $\beta$ preserves that identity. The intervals
may be anywhere in the fundamental cube; none is required to surround zero.

If independence of the core, four $b$ blocks and the first three $l$ blocks
fails on more than a sparse set of quadruples, the bad-set rule above fixes
one nonzero standard integer relation with error $O(\epsilon)$ on many
quadruples. There are three exhaustive cases.

If it contains only core coefficients, solve for a core with nonzero
coefficient $a$. Choose internal $a$th roots of the other core frequencies.
Division introduces a residue in the finite group $a^{-1}\GIHZ/\GIHZ$ and an
$O(\epsilon)$ error. Move them to the rational and error lists. This reduces
the core count without changing either petal count.

If the relation involves $l$ but no $b$, it is a relation on three free
parameters. Fix two by averaging over the many triples. On a dense set of
the remaining parameter, solve for one $l$ coefficient. Divide the remaining
bracket-linear coefficients as real numbers $\beta/a$, and root the core
coefficients. Freeze the finite rational carry and add the values at the
two fixed parameters to the core. A further fixed-box restriction restores
the Freiman laws for the divided bracket-linear terms. This reduces the
$l$ count without changing the $b$ count.

If a $b$ coefficient occurs, use a symmetry of $h_1+h_2=h_3+h_4$ to place
a nonzero one in position four. The Freiman laws allow elimination of the
fourth $l$ block after this permutation. Group the relation as
$X_1(h_1)+\cdots+X_4(h_4)=O(\epsilon)$, with $X_4$ containing that $b$
coefficient and fixed core terms. Contract F expresses $X_4$ as finitely many
bracket-linear terms plus a constant and error. Solve for the chosen $b$
coefficient, using roots of the other $b$ and core frequencies and dividing
the real bracket-linear coefficients. Freeze the finite rational carry,
enlarge the other lists, and refine into fixed boxes as above. The original
frequencies are still integer combinations of the new lists, because the
old frequencies are integer multiples of their chosen roots. The $b$ count
has decreased. Each case contradicts the lexicographic minimum.

Thus the required independence holds for almost all quadruples. Apply this
argument successively for the finitely many weights $1,\ldots,m$. Dense
restrictions preserve previously obtained independence: their bad sets
remain $o(N^3)$, while the retained parameter set has standard positive
density. Standard integer relations, rational carries, error constants,
finite dimensions and fixed model choices are selected internally with
the explicit bounded-choice rule. This proves the regularization step of
10.5--10.7 conditional only on contract F, with its needed family uniformity.

\phantomsection\label{q214foundation:gi-normalization}
\subsection{Removing rational and small coordinates without losing unit norm}
Apply the integer change-of-basis construction to the regularized frequency
lists. Before removing any coordinate, make the horizontal identities exact.
Fix internal real lifts $x_{h,i,a}$ of the regularized frequencies: the small
lift for a small coordinate, a rational lift for a rational coordinate, and
$\{\alpha h\}\beta$ for a linear petal. For each $i$, the difference between
the $i$th horizontal Taylor coefficient of $g_h$ and
$\sum_a x_{h,i,a}\phi(e_{i,a})$ lies in the horizontal lattice of $\Gamma$.
Choose $\gamma_{h,i}\in{}^*(\Gamma\cap G_{(i,1)})$ mapping to this difference,
and replace $g_h(n)$ by
\[
 g_h(n)\prod_i\gamma_{h,i}^{-\binom ni}.
\]
The product is taken in a fixed order. It takes values in $\Gamma$, so this
right multiplication does not change the orbit modulo $\Gamma$. The Taylor
criterion B.9 and the polynomial-product rule B.4 make it polynomial. The
horizontal Taylor map is a homomorphism, so the identities now hold in the
horizontal vector groups themselves. These are internal choices, as in the
lattice normalization used below. In the two removals that follow, use these
fixed lifts as the exponents. Each removal then subtracts the removed term
exactly, preserving the identities for the remaining coordinates.

We now remove each rational or small coordinate in turn. The
remaining top character and lattice map must be updated; forgetting them
would not give a new representation.

For a small coordinate $\xi=O(N^{-i})$, choose its small real lift and put
$c=\phi(e_{i,a})^{\xi}\in G_{(i,1)}$. Set
$g'_h(n)=c^{-\binom ni}g_h(n)$. B.4 proves polynomiality and horizontal
homomorphism shows that exactly the selected coordinate has been removed.
On $[N]$, $c^{\binom ni}$ is a bounded smooth function of $u=n/(2N)$.
Indeed its logarithm is $\binom{2Nu}{i}\xi\log\phi(e_{i,a})$; its coefficients
and all derivatives on a fixed compact interval are bounded because
$N^i\xi=O(1)$. Extend this logarithm to the unit circle by using its real
formula on the lift $(-1/4,3/4)$ and multiplying by a fixed smooth cutoff
equal to one on $[0,1/2]$ and zero near the lift's endpoints. Let $c_h(u)$
be its exponential. Replace the test by
\[
 F'_h(x,x_0,u)=F_h(c_h(u)x,x_0).
\]
It is globally a unit vector and uniformly Lipschitz. The top image
$\phi(U_{(m,R)})$ is central, so the same exact vertical law holds everywhere.
Add the circle orbit $n/(2N)$ to the lower model. It has degree one, which
is at most $(m,R-1)$ because $m\geq2$. Thus this is an exact representation,
not an approximation to the original sequence.

For a rational coordinate $q\in\GIHQ/\GIHZ$, put $\gamma=\phi(e_{i,a})^q$ and
$g'_h(n)=g_h(n)\gamma^{-\binom ni}$. Again only that horizontal coordinate
is removed. The element $\gamma$ is rational relative to $\Gamma$.
The subgroup $\widehat\Gamma=\langle\Gamma,\gamma\rangle$ is discrete:
in rational Mal'cev coordinates the finite nilpotent collection formulas
express every word as a polynomial in integer exponents of a finite set
of generators and their finitely many commutators; clearing their fixed
rational denominators gives a discrete coordinate grid. Since $\Gamma$
is cocompact, its index in $\widehat\Gamma$ is finite. Its normal core
$\Gamma'=\bigcap_{\delta\in\widehat\Gamma}\delta\Gamma\delta^{-1}$ is
therefore a finite-index subgroup, normalized by $\gamma$.

Choose a standard $b$ with $\gamma^b\in\Gamma$ and a standard period $q_0$
such that $\binom{n+q_0}{i}\equiv\binom ni\pmod b$. For example a sufficiently
divisible multiple of $b\,i!$ works, by the integer binomial difference
formula. Put $\gamma_r=\gamma^{\binom ri}$ for $r\pmod{q_0}$. Then
\[
 F_h(g_h(n)\Gamma,x_0)=F_h(g'_h(n)\gamma_r\Gamma,x_0)
 \quad(n\equiv r\pmod{q_0}).
\]
The test $W_{r,h}(x\Gamma',x_0)=F_h(x\gamma_r\Gamma,x_0)$ is well-defined,
unit, and has the same top-frequency law. Also
$W_{0,h}(x\Gamma',x_0)=F_h(x\Gamma,x_0)$ is such a global unit test.

To interpolate the finitely many residue tests, choose disjoint small arcs
around $r/q_0$ in $\GIHT$. In the $r$th arc choose a smooth angle $\theta_r$
which is zero at $r/q_0$ and $\pi/2$ near its boundary. Set $a_r=\cos\theta_r$
there and zero outside; set $b=\sin\theta_r$ there and $b=1$ outside all
arcs. The supports of the $a_r$ are disjoint and $\sum_r a_r^2+b^2=1$.
The test
\[
 \left(\sum_r a_r(u)W_{r,h}(x,x_0),\ b(u)W_{0,h}(x,x_0)\right)
\]
is globally Lipschitz, unit and has the same vertical character. At the orbit
$u=n/q_0$ it is exactly $(\chi_h(n),0)$. This fixed zero padding is explicit;
we do not assert that an arbitrary scalar interpolation stays unit or that
the old dimension has been retained. Correlation and subsequent matrix
comparisons recover the original family by the fixed coordinate projection.

\phantomsection\label{q214foundation:gi-cover}
The map from the remaining universal group must now land in the covering
lattice $\Gamma'$. Choose a standard $M$ for which every remaining lattice
generator's $M$th power lies in $\Gamma'$. Send each new letter to that
$M$th power, and divide its horizontal frequency by $M$. Divide the fixed
real lift: since
\[
 \sum_a (x_{h,i,a}/M)\,M\phi(e_{i,a})
       =\sum_a x_{h,i,a}\phi(e_{i,a}),
\]
the vector identities are unchanged, and hence hold modulo the horizontal
lattice of $\Gamma'$. Pull back the character by the same map. For linear petals choose the division as
$\{\alpha h\}(\beta/M)$; for core coordinates choose fixed roots and for
independent coordinates choose internal roots. Repeating this finite process
produces a common standard denominator. It is not silently dropped, as in
the displayed conclusion of the printed proof of 10.10.

Independence is preserved: multiply an alleged standard integer relation
between divided frequencies by the common denominator to contradict the
old independence. After dividing linear petals, refine again into the fixed
short boxes to make their real fractional-part sums exactly additive, hence
their divided values Freiman. This is valid at arbitrary box locations.
All rational covers, maps, dimensions, and periods are fixed on a dense
internal subfamily before use. Small-coordinate curve bounds are uniform
by the internal bounded-choice rule. A finite composition of paddings is
one fixed coordinate isometry. We have proved the precise representation
part of Theorem~\ref{v15:gi:7:thm:output}, without replacing the varying $F_h$ by a
single equal function through an approximation argument.

\subsection{The horizontal consequence of factorization, with denominators}
Here we derive the needed D.6 conclusion from contract D. View a degree--rank
group as ordinarily filtered by $G_i=G_{(i,0)}$. We retain the finer horizontal
spaces $\GIHHor_i(G)=G_{(i,1)}/G_{(i,2)}$. For a polynomial orbit $g\Gamma$ put
\[
 \Xi_i=\{\lambda\in\widehat{\GIHHor_i(G/\Gamma)}:
                 \lambda(\GIHTay_i(g\Gamma))=O(N^{-i})\pmod1\}.
\]
Characters here have standard integer coordinates. Their real linear lifts
define $\Xi_i^\perp$ inside the real horizontal vector space.

Suppose $F$ is a bounded Lipschitz test with nonzero standard average.
Take one scalar component if needed, and factor $g=\varepsilon g'\gamma$.
On a progression of the standard periodic modulus, choose a bounded standard
representative $\gamma_P\in\widehat\Gamma$ of the constant coset of $\gamma$.
Set $Q_P=\gamma_P^{-1}Q\gamma_P$ and
$g_P=\gamma_P^{-1}g'\gamma_P$. The finite-cover equidistribution clause in
D, applied to the common finite-index lattice, makes $g_P$ totally
equidistributed in $Q_P/(Q_P\cap\Gamma)$ on every retained dense progression.
Rationality of $\gamma_P$ gives rationality of $Q_P$.

Partition into those residue progressions and then into intervals of length
$\delta N$ for a sufficiently small standard $\delta$. Averaging selects a
piece with an average bounded below by the original fixed threshold.
On that piece freeze $\varepsilon(n)\gamma_P$ at one point. The uniform
Lipschitz error is $O(\delta)$, so equidistribution gives
\begin{equation}\label{v15:gi:7:eq:coset}
 \left|\int_{P/(P\cap\Gamma)}F(\varepsilon_0x)\,d\mu(x)\right|>c>0
\end{equation}
for the connected rational $P=Q_P$ and bounded $\varepsilon_0$.
This also holds for an internal test with a standard Lipschitz bound:
compactness gives a finite standard uniform net in that Lipschitz ball.
Apply total equidistribution to each net entry and then let the fixed net
error tend to zero. Thus no unjustified standardness of the selected $F_h$
or of the frozen bounded translate is being used.

We claim $\pi_i(P\cap G_i)\supseteq\Xi_i^\perp$. Let $\lambda$ be an
integer character annihilating the rational space $\pi_i(P\cap G_i)$.
Then $\lambda(\GIHTay_i(g_P))=0$. Conjugation changes no horizontal
coefficient since $[G,G_i]\subseteq G_{(i,2)}$. Hence the horizontal
homomorphism identity gives
\[
 \lambda(\GIHTay_i(g))=
 \lambda(\GIHTay_i(\varepsilon))+
 \lambda(\GIHTay_i(\gamma)) \pmod1.
\]
The first term is $O(N^{-i})$. The second is rational of bounded standard
denominator: derivatives of $\gamma$ lie in $\widehat\Gamma\cap G_i$, whose
horizontal image is a finite-index over-lattice of that of $\Gamma\cap G_i$.
Thus for some standard positive $q$, $q\lambda\in\Xi_i$. Every annihilator
of the rational image has this property. Taking real annihilators proves the
claim. In particular no false assertion that $\gamma$ takes values in
$\Gamma$ or has zero horizontal coefficient is needed.

\subsection{Applying the coset conclusion to the actual quadruple}
After all the preceding dense restrictions, reapply hereditary LP03.
There are $\gg N^3$ biased quadruples. The exceptional sets for frequency
independence have size $o(N^3)$. Permutations preserving the additive
quadruple equation make the omitted linear-petal position any of the four;
the union of the four exceptional sets is still sparse. Fix a biased
quadruple with all four required independence statements.

Use the representations of its four nilcharacters and the lower bias test
to form a product group
\[
 \mathcal G=G^{(1)}\times G^{(2)}\times G^{(3)}\times G^{(4)}\times G^{(0)}.
\]
The first four groups have degree--rank at most $(m,R)$; the last has degree
at most $(m,R-1)$, including all four lower representation models and the
ordinary degree-at-most-$m-1$ bias witness. The stronger strict inequality
printed for this last group is unnecessary and is not assumed.
Translations in $n$ preserve horizontal Taylor coefficients by their
homomorphism/difference identity. Zero extensions in \eqref{v15:gi:7:eq:quad}
introduce only an interval-intersection mask. Approximate this mask in
normalized $L^1$ to error $\delta$ by a bounded Lipschitz circle test in
$n/(cN)$, with fixed large $c$ and Lipschitz constant $O(1/\delta)$.
Its two transition strips have total length $O(\delta N)$; multiplying by
the uniformly bounded remaining tests controls the average error. Choosing
$\delta$ below the fixed bias threshold preserves positive bias. This circle
has degree one and belongs in $G^{(0)}$ because $m\geq2$.

The resulting product test has a nonzero scalar average and exact vertical law
\[
 F(\phi_1(t_1)x_1,\ldots,\phi_4(t_4)x_4,x_0)
 =\GIHee(\eta(t_1)+\eta(t_2)-\eta(t_3)-\eta(t_4))F(x_1,\ldots,x_4,x_0).
\]
All factors, including the cutoff, are uniformly Lipschitz on a fixed
standard product geometry. Apply the preceding factorization consequence
to obtain a connected rational $P\subseteq\mathcal G$ and the nonzero
coset integral \eqref{v15:gi:7:eq:coset}.

\subsection{Rational duality and the two commutator conclusions}
For each weight $i$, let $V_{123,i}$ be the real space spanned by the
three-coordinate diagonal images of core letters and by each individual
coordinate image of every petal letter in positions $1,2,3$. Define
$V_{124,i}$ likewise. The projection of $\pi_i(P\cap\mathcal G_i)$ to
these three factors contains the indicated $V$.

Here is the exact duality argument. If containment failed, rationality gives
a rational linear functional that vanishes on the projected space but not
on $V$. Extend it by zero in the omitted factors and clear denominators to
make it an integer character $\lambda$. Because the full space contains
$\Xi_i^\perp$, this character is in the rational span of $\Xi_i$. After
multiplication by another standard integer it belongs to $\Xi_i$ itself.
Evaluating it on the Taylor coefficient gives a relation
\[
 \sum_{a\in C}(c_{1a}+c_{2a}+c_{3a})\xi_{C,a}
       +\sum_{k=1}^3\sum_{a\in L\cup B}c_{ka}\xi_{h_k,a}
       =O(N^{-i}).
\]
The selected independence forces every petal coefficient and each displayed
core sum to vanish. That says the character vanishes on $V$, a contradiction.
The $124$ proof is identical. This uses rational real subspaces, not an
unjustified identification of an arbitrary real annihilator with one
integral member of $\Xi_i$.

Similarly the projection onto all four main factors contains the space
$V_{B,i}$ generated by diagonal core images and by the individual-coordinate
images of every letter in $B$. In the analogous integer relation eliminate
the fourth $L$ coefficient using
$\xi_{h_4,a}=\xi_{h_1,a}+\xi_{h_2,a}-\xi_{h_3,a}$ for $a\in L$.
Independence makes the four $B$ coefficients vanish separately and the sum
of core coefficients vanish. This is exactly annihilation of $V_{B,i}$.
No numerical range for $B$ is borrowed from the conflicting printed order
of the $L$ and $B$ blocks.

Let $w$ be any bracketing of $R$ letters of weights summing to $m$.
First suppose it has two petal occurrences. For one occurrence lift its
first-coordinate image, with zero second and third coordinates, from
$V_{123,i}$ to $P\cap\mathcal G_i$. For the second lift its first-coordinate
image with zero second and fourth coordinates from $V_{124,i}$. Lift the
remaining letters so that their first horizontal coordinate is the required
one. Take the same iterated commutator of these lifts.

In the first factor it equals $\phi_1(w)$ exactly. Replacing one lift by its
horizontal error inserts an element of $G_{(i,2)}$ in place of
$G_{(i,1)}$, raising the total rank to at least $R+1$ at weight $m$;
every such term vanishes. Higher BCH terms have larger weight and vanish
as well. In factor two, either selected petal lift has such an error;
factor three has the first, and factor four the second. Therefore those
coordinates are the identity. In factor zero, an $R$-fold commutator of
the prescribed weights lies in its trivial $(m,R)$ group. We obtain
\[
                 (\phi_1(w),1,1,1,1)\in P.                 \tag{*}
\]
This reasoning applies to each position of a repeated letter and to every
bracketing, not only a single left-nested commutator.

If $w$ has a $B$ letter, use $V_{B,i}$ to lift that occurrence with all
other main coordinates zero horizontally. Lift the remaining occurrences
by their first coordinates and repeat the same rank-error calculation.
This proves (*) as well. It works for $R=1$; the other top coordinates then
already vanish, and the lower model has trivial weight-$m$ group.

Finally (*) gives a real-zero conclusion. The element is central in
$\mathcal G$. Since $P$ is a connected subgroup of a simply connected
nilpotent group, its entire one-parameter subgroup belongs to $P$.
Left translation by its $t$th power leaves Haar measure on $P/(P\cap\Gamma)$
invariant and commutes with the bounded coset translate. The vertical law
multiplies the nonzero integral by $\GIHee(t\eta(w))$. Hence
$\GIHee(t\eta(w))=1$ for every real $t$, so $\eta(w)=0$ as a real number.
Using only $t=1$ would give an integer condition and would be insufficient.
This proves both annihilations in Theorem~\ref{v15:gi:7:thm:output}.

\subsection{The LP04 and rank-family joins}
The output has the exact LP04 hypotheses: the corrected universal weighted
group, an integral top character with all required real-zero annihilations,
fixed lattice-compatible $\phi$, matched horizontal Taylor tori, and a unit
representation with a lower degree--rank factor. Its petal quotient kills
the whole ideal generated by the bad words at every weight; the vanishing
just proved concerns the top restriction needed to descend the character.
These distinct roles are not identified.

The Section 12 join chooses its stronger integral extension to the ordinary
top before invoking LP04 comparison, and gives one global multidegree
$(1,m)$ nilcharacter equal to a fixed coordinate embedding of $\zeta_h$.
LP04 then supplies the whole matrix
$(\iota\chi_h)\otimes\overline{\zeta_h}$ with strictly lower degree--rank.
Projecting the first tensor factor through the fixed coordinate projection
gives $\chi_h\otimes\overline{\zeta_h}$ with the same bound and exact values.
Thus the unit padding used in rational removal does not weaken the actual
matrix input required by rank-iteration001.

All selections, roots, covers, maps and tests were made internally on
representatives; each retained dimension, character and bound is standard.
There are finitely many weights and coordinates at any one rank, so all
density and correlation losses remain positive standard numbers. The next
rank family from rank-iteration001 again has exactly the hypotheses of
Section 1. Repeating the present reconstruction at that new rank uses its
own new frequency data. No independence is transported through a
Stone--Weierstrass approximation or assumed for the newly selected test.
There are $m$ rank steps, so the same finite-loss conclusion holds across
the full iteration. This is the required source-to-consumer compatibility,
conditional on the two separately supplied analytic contracts and the
stated hypotheses of the internal proof units used here.
Representation equality is required on $[N]$. Whenever a global sequence
is used, take the displayed polynomial model as its extension. Small-frequency
removal need not preserve a previous extension outside $[N]$; the interval
mask in the quadruple argument is precisely what makes that change harmless.

\subsection{Bounded correction and consistency record}
The source scope is final GTZ Sections 9--11, the immediate pp.~1307--1309
consumer, the D.4--D.6 interface, B.4/B.7, and the selection statements in
Appendix A. The attached record pins exact passages and the current project
consumers. Concrete changes are:
\begin{enumerate}
\item State the fixed rank $R$ explicitly in Definition 9.1's second
truncation relation, where the printed formula writes $r$ without the fixed
rank subscript. The printed weight cutoff $\sum i_j\geq s$ is already
equivalent to $\sum i_j>m$ and is not an error. Rendering the source rather
than relying on text extraction establishes this distinction.
\item Retain the common division from 10.9 in 10.10 and restore Freiman laws
after real coefficient division. Fixed arbitrary-location boxes suffice.
\item Rational removal uses a finite cover and explicit same-character unit
completion with a fixed isometry, rather than unproved unit interpolation.
\item Keep varying uniformly Lipschitz $F_h$; total boundedness does not make
them literally equal as claimed in the immediate p.~1308 narrative.
\item Use named core/linear/independent blocks, because p.~1302 reverses the
order used in the statement of 10.10 and Theorem 11.1.
\item Treat the repeated clauses (iii),(iv) of Theorem 11.1 as duplicates;
there are two distinct annihilation conclusions.
\item Use degree--rank at most $(m,R-1)$ for the lower product model, not an
unneeded strict inequality below that index.
\item Retain the rational Taylor term in D.5 and clear its denominator in
the annihilator argument. The factorization subgroup may equal the ambient
group. The separate analytic supplier addresses the D.2 total-equidistribution
issue and F.1 signed-graph/no-wrap details.
\item The real-zero character argument uses all real powers in the connected
subgroup. One integral phase identity alone is insufficient.
\item Reapply LP03 after each dense restriction; intersect many biased
quadruples with almost-all independence, not two unrelated dense sets.
\end{enumerate}
These are authored repairs and scoped dependencies, not a declaration that
all historical versions or the entire inverse theorem have been checked.
The original published source and all earlier frozen packets are preserved.

\section{GTZ-LP-04: the petal quotient and degree--rank construction}\label{v15:gi:GTZ-LP04}
\newcommand{\GIIR}{\mathbb R}
\newcommand{\GIIQ}{\mathbb Q}
\newcommand{\GIIZ}{\mathbb Z}
\newcommand{\GIIN}{\mathbb N}
\newcommand{\GIIDR}{\mathrm{DR}}
\newcommand{\GIIee}{\mathrm e}
\newcommand{\GIIlie}{\mathfrak u}
\newcommand{\GIIid}{\mathrm{id}}
\newcommand{\GIIHH}{\operatorname{Horiz}}
\newcommand{\GIIpoly}{\operatorname{poly}}

This section proves the algebraic quotient, its exact lattice and character statements, and the comparison needed by \cite[Lemma 12.1]{v15-gi-8-GTZ}, conditional on the explicitly stated frequency input. That frequency input is supplied in Section~\ref{v15:gi:GI-SECTION9-11}. The erratum \cite{v15-gi-8-Err} motivated the construction but is not a theorem premise. The rational-subgroup and tautification arguments are in Sections~\ref{v15:gi:GTZ-LP13-MULTILINEARITY} and \ref{v15:gi:GTZ-LP-TAUTIFICATION}.

\subsection{Indices, free groups, and the actual objects}
Fix integers $m\geq1$, $1\leq R\leq m$, and finite disjoint sets of letters
$C_i,L_i,B_i$ for $1\leq i\leq m$. Their union is $\mathcal A$.
The letters in $C_i$ are core, those in $L_i$ retained linear petals, and
those in $B_i$ discarded petals. Every letter in these three sets has weight
$i$. A bracket word has length $\ell$ equal to its number of entries, weight
$w$ equal to the sum of their weights, and a multidegree recording the number
of occurrences of each individual letter. Repeated entries count repeatedly.
All bracketings are allowed, including length-one words. Write $p$ for the
number of petal entries, and $b$ for the number of discarded entries.
We use $[x,y]=x^{-1}y^{-1}xy$, real powers
$x^t=\exp(t\log x)$, and $\GIIee(t)=\exp(2\pi\mathrm i t)$.

Let $\mathcal F$ be the free real Lie algebra on symbols $X_a$, $a\in\mathcal A$.
Its usual length grading and its individual-letter multigrading are retained.
Take the homogeneous ideal spanned by all words for which
\begin{equation}\label{v15:gi:8:eq:truncate}
 w>m\quad\hbox{or}\quad(w=m\hbox{ and }\ell>R),
\end{equation}
and let $\GIIlie$ be the quotient. These forbidden multidegrees are closed under
addition of any nonzero multidegree, so their span is indeed an ideal.
All surviving words have length at most $m$. Thus $\GIIlie$ is finite-dimensional
and nilpotent. Let $U=\exp\GIIlie$, with multiplication given by the finite BCH
polynomial, and set $e_a=\exp X_a$. The group is connected and simply connected.
These are the weighted universal objects intended in
\cite[Definition 9.1, p.~1279]{v15-gi-8-GTZ}, with the top rank in the relation fixed
to be $R$, not a running rank index.

Use $\GIIDR=\{(d,r):d,r\in\GIIN,\ 0\leq r\leq d\}$ with lexicographic order.
Put $U_{(0,0)}=U$. For $d>0$ let
\begin{equation}\label{v15:gi:8:eq:filtration}
 \GIIlie_{(d,r)}=
 \operatorname{span}_{\GIIR}\{W:w(W)>d\ \hbox{or}
 [w(W)=d,\ \ell(W)\geq\max(r,1)]\},
 \qquad U_{(d,r)}=\exp\GIIlie_{(d,r)}.
\end{equation}
In particular the first alternative has \emph{no} length restriction.
For ranks larger than $d$ use the notation $U_{(d,r)}=U_{(d+1,0)}$ only
as the next-index convention; do not put such pairs into the index monoid.
Write $U_d=U_{(d,0)}$. The top group is
\[
 T=U_{(m,R)}=\exp\operatorname{span}_{\GIIR}
       \{W:w(W)=m,\ \ell(W)=R\}.
\]

\begin{lemma}\label{v15:gi:8:lem:filtration}
The groups in \eqref{v15:gi:8:eq:filtration} form a taut degree--rank filtration of
degree at most $(m,R)$. Each is connected and normal in $U$, and $T$ is central.
\end{lemma}
\begin{proof}
Bracketing words adds weights, lengths, and individual-letter multidegrees.
Each displayed span is an ideal, since bracketing with a nonzero word increases
weight. It is a Lie subalgebra and BCH therefore identifies its exponential
with a connected subgroup. Nesting follows directly from lexicographic order:
passing to a larger degree removes all lower-weight words, and at fixed degree
passing to a larger rank increases the length cutoff.
For positive indices $(d,r),(e,t)$, a bracket of two allowed words either has
weight greater than $d+e$, or has weight $d+e$ and length at least
$\max(r,1)+\max(t,1)\geq\max(r+t,1)$. Hence the Lie bracket lies in
$\GIIlie_{(d+e,r+t)}$. The BCH commutator formula has this bracket and further
brackets, all in the same ideal, which proves the group commutator condition.
At index zero use normality. Every generator has weight at least one, so
$U_{(1,0)}=U$; also $U_{(d,0)}=U_{(d,1)}$ for all $d>0$.
The truncation gives the claimed degree support. Bracketing a top word with
any generator raises its weight above $m$, proving centrality of $T$.
\end{proof}

\subsection{An integral coordinate lemma}
Here is the integral premise used, with its scope separated from the subsequent
quotient argument. Hall collection says that in a free nilpotent group on a
finite ordered alphabet, every element has a unique collected expression in
the basic commutators with integer exponents; the length-$j$ basic commutators
give a basis of its $j$th lower-central quotient. This published statement is
recorded in \cite[Definition 5.4.3, Theorem 5.4.5 and the final paragraph
of \S5.4, pp.~96--98]{v15-gi-8-Tointon}. The weight in that reference is bracket
\emph{length}, not the additional weight $w$ used here.

The corresponding basic Lie brackets form a homogeneous basis of the free
nilpotent Lie algebra. To spell out the connection, antisymmetry and Jacobi
reorder a bracket $[c_i,c_j]$ until $i>j$, and when
$c_i=[c_s,c_t]$ until $j\geq t$. Jacobi replaces the offending bracket by
brackets with the smaller inner entry moved outward. Induction on total
length and then the finite Hall ordering gives spanning by basic brackets.
Independence follows by mapping the free Lie algebra to the real tensor of
the lower-central graded Lie ring of the free group: the bracket is induced
by the group commutator, and Hall's basis assertion makes the images of those
basic brackets independent in each length. This also respects the individual
letter multigrading because the reordering identities preserve every count.

\begin{lemma}[Upward-closed monomial quotients]\label{v15:gi:8:lem:hall}
In a free nilpotent Lie algebra of class at most $m$, discard an upward-closed
set of nonzero individual-letter multidegrees. Let $V$ be the exponential
of the quotient and $\Lambda$ the discrete group generated by the surviving
$\exp X_a$. The surviving Hall commutators give unique real coordinates
\[
 (t_1,\ldots,t_v)\longmapsto c_1^{t_1}\cdots c_v^{t_v}
\]
on $V$, and $\Lambda$ is exactly the set with all $t_i\in\GIIZ$.
It is a lattice. A further upward-closed set of multidegrees defines a
rational connected normal subgroup $V_S$. The quotient map $q$ between
two such monomial quotients obeys, for each such $S$,
\begin{equation}\label{v15:gi:8:eq:saturation}
 q(\Lambda\cap V_S)=q(\Lambda)\cap q(V_S).
\end{equation}
\end{lemma}
\begin{proof}
Order basic commutators by increasing length. For a basic group commutator
$c$, BCH gives $\log c=W_c+$ terms of strictly greater length, where $W_c$
is the corresponding basic Lie bracket. Each correction uses all the letters
of $c$ and possibly more occurrences of them. A discarded basic word
therefore vanishes in the quotient, including every BCH correction.
The surviving leading brackets are a basis, since the discarded Lie ideal
is the direct sum of precisely the discarded multigraded components.
Consequently the logarithm of the displayed coordinate map is triangular
by length with diagonal entries one. Solving successively in each length
gives a polynomial inverse and proves unique real coordinates.

Hall collection expresses every word in the original $\exp X_a$ and their
inverses as a product of basic group commutators with integer exponents.
After quotienting, discarded factors are identities, so it gives the claimed
integer coordinates. Conversely every surviving basic commutator is itself
a word in the generators, which proves the other inclusion. In particular
the subgroup is discrete. Suffix spans in the length-ordered basis are ideals,
since every nonzero bracket increases length. Right multiplication by
$c_i^{-n}$ changes coordinate $i$ by $-n$, leaves all earlier coordinates
unchanged, and can change only later coordinates. Reduce the coordinates
successively modulo integers. Every right coset has a representative with
coordinates in $[0,1]^v$, a compact set, proving cocompactness.

For $V_S$, quotient once more by its multigraded ideal. The same triangular
coordinate argument says that an element belongs to $V_S$ exactly when all
coordinates outside $S$ vanish. Thus its lattice consists of the integer
coordinates on the basic commutators in $S$, and the same compact-coordinate
argument proves rationality. A quotient deletes discarded basic factors and
keeps the remaining ones in their original order. Membership in $q(V_S)$
forces all retained coordinates outside $S$ to vanish. An element of
$q(\Lambda)\cap q(V_S)$ therefore lifts as the identical integer product
of retained basic commutators in $S$. This proves \eqref{v15:gi:8:eq:saturation}.
\end{proof}

The truncation \eqref{v15:gi:8:eq:truncate} is upward-closed, as are all the thresholds
\eqref{v15:gi:8:eq:filtration}. The lemma proves that
$\Gamma=\langle e_a:a\in\mathcal A\rangle_{\GIIZ}$ is a lattice in $U$, and
every $U_{(d,r)}$ is rational. It also proves the lattice statement below
without treating it as a general property of arbitrary rational quotients.

\subsection{The corrected kernel and the descended character}
Call a word bad when $p\geq2$ or $b\geq1$. Let $I\subseteq\GIIlie$ be the
span of all bad Lie words of every weight. Let $G^*$ be the connected real
subgroup generated by all bad iterated \emph{group} commutators in the $e_a$,
again of every weight and including length one. Equivalently, the real
generated subgroup contains every real power of each such commutator.
This convention is essential: a discrete generated subgroup would not be
the kernel of a map of simply connected Lie groups.

Let $\widetilde{\GIIlie}$ be the free Lie algebra on $C_i\cup L_i$ with relations
\eqref{v15:gi:8:eq:truncate} and the vanishing of all words with at least two petal
entries. Give $\widetilde G=\exp\widetilde{\GIIlie}$ the filtration
\eqref{v15:gi:8:eq:filtration} on its surviving words. Put
$\widetilde\Gamma=\langle\widetilde e_a:a\in C_i\cup L_i\rangle_{\GIIZ}$.
Define $q:U\to\widetilde G$ by $e_a\mapsto\widetilde e_a$ on retained
letters and $e_a\mapsto1$ on discarded letters.

\begin{theorem}\label{v15:gi:8:thm:quotient}
The homomorphism $q$ is well-defined and surjective, and
\[
 \ker q=G^*=\exp I.
\]
Thus $G^*$ is connected, rational and normal; it need not be central. For
every degree--rank index $a$, including zero,
\begin{align}
 q(U_a)&=\widetilde G_a\simeq U_a/(G^*\cap U_a),\label{v15:gi:8:eq:filteredquot}\\
 q(\Gamma\cap U_a)&=\widetilde\Gamma\cap\widetilde G_a.\label{v15:gi:8:eq:exactlattice}
\end{align}
Suppose $\eta:T\to\GIIR$ is a continuous homomorphism integral on
$T\cap\Gamma$, and for every bad group commutator $c$ of weight $m$ and
length $R$ one has $\eta(c)=0$. Then there is a unique continuous
$\widetilde\eta:\widetilde G_{(m,R)}\to\GIIR$ satisfying
\[
 \widetilde\eta(q(t))=\eta(t)\quad(t\in T),\qquad
 \widetilde\eta(\widetilde\Gamma\cap\widetilde G_{(m,R)})\subseteq\GIIZ.
\]
\end{theorem}
\begin{proof}
Bad multidegrees are upward-closed, so $I$ is an ideal. The universal property
of the free Lie algebra identifies $\GIIlie/I$ with $\widetilde{\GIIlie}$:
killing discarded letters kills every word involving one, and killing
all words with two petals gives exactly the other defining relations.
There are no additional relations because both presentations impose the
same homogeneous ideal on the same free Lie algebra. Exponentiation of this
surjective Lie map gives $q$, and its kernel is $\exp I$ by the global
exponential diffeomorphisms.

BCH puts the logarithm of every bad group commutator in $I$. Conversely
take the bad Hall commutators. Their logarithms have the bad basic Lie
brackets as leading terms, with only higher-length bad terms. The triangular
matrix on the finite-dimensional bad subspace is invertible. Their logarithms
therefore span $I$, and the connected real subgroup they generate is $\exp I$.
This proves the asserted description of $G^*$, not just its normal closure.
All these spaces are spanned by rational basic words. Lemma~\ref{v15:gi:8:lem:hall}
proves rationality and $q(\Gamma)=\widetilde\Gamma$.

The Lie map deletes bad multigraded components and preserves weight and
length of every surviving word. Applying \eqref{v15:gi:8:eq:filtration} before or
after the map therefore gives identical spaces. Exponentiating proves
\eqref{v15:gi:8:eq:filteredquot}, and Lemma~\ref{v15:gi:8:lem:hall} proves
\eqref{v15:gi:8:eq:exactlattice}.

The direct multigrading gives
\[
 I\cap\log T=\operatorname{span}_{\GIIR}
    \{W:W\hbox{ bad},\ w(W)=m,\ \ell(W)=R\}.
\]
There is no cancellation between a bad multidegree and a good one. At
weight $m$ the BCH correction terms in a group commutator have larger
weight and vanish. Thus the bad group commutators of weight $m$ and
length $R$ are precisely the exponentials of their Lie brackets.
The assumption on $\eta$, and real linearity of a continuous homomorphism
on the vector group $T$, prove that $\eta$ vanishes on $G^*\cap T$.
It consequently descends through \eqref{v15:gi:8:eq:filteredquot} at the top index.
Finally \eqref{v15:gi:8:eq:exactlattice} lifts every top lattice point downstairs to
one upstairs, where $\eta$ is integral. This proves the last assertion.
\end{proof}

\paragraph{Why the omitted lower weights matter.}
With one core letter $x$ and one discarded letter $b$, both of weight one,
and $m=3,R=3$, the quotient kills $b$ itself. A subgroup generated only by
bad weight-three words cannot contain $b$. Moreover $[b,x]$ survives in $U$,
so the actual kernel is not central. This example tests both false shortcuts.

\paragraph{The precise analytic input.}
In the notation of \cite{v15-gi-8-GTZ}, take $m=s-1$, $R=r_*$,
$|C_i|=D_{*,i}$, $|L_i|=D'_{\mathrm{lin},i}$ and
$|B_i|=D'_{\mathrm{ind},i}$. Use the \emph{final output} of Lemma 10.10,
after rational-core and small frequencies have been removed by Lemma 10.9.
Then the two hypotheses on $\eta$ in Theorem~\ref{v15:gi:8:thm:quotient} are exactly
Theorem 11.1(i),(ii): all weight-$m$, length-$R$ words with two petals or
one beyond-linear petal are annihilated. The repeated (iii),(iv) add nothing.
For a bracketing convention restricted to left-normed words, Jacobi expresses
any other bracket of these same entries as a linear combination of such words;
their petal counts and total weight are unchanged. At the central top layer
this is also an exact identity of logarithms, so the vanishing extends to
all bracketings used here.

Theorem 11.1 requires the hypotheses of Theorem 7.2, the limit family of
degree--rank $(m,R)$ nilcharacters, and the sunflower representation with
its common standard dimension vector and vertical frequency. Its proof uses
the additive-quadruple bias conclusion of Proposition 7.3 and the stated
independence from Lemma 10.10. In the repaired project, that bias conclusion
must come from LP-03, not the false original Proposition 8.3. No use of
Theorem 1.3 is made to supply it. This unit assumes these inputs; merely
partitioning generators does not prove the required vanishing.

\subsection{Vertical test functions and petal actions}
\begin{lemma}\label{v15:gi:8:lem:vertical}
For the descended integral character there is a finite-dimensional
unit-vector-valued smooth function on $\widetilde G/\widetilde\Gamma$ of
vertical frequency $\widetilde\eta$. Its dimension and Lipschitz bound depend
only on the fixed filtered group, lattice, metric and character.
\end{lemma}
\begin{proof}
The central torus $\widetilde T/(\widetilde T\cap\widetilde\Gamma)$ acts
freely. Choose finitely many local smooth sections of the resulting principal
torus bundle and smooth real functions $a_j$ subordinate to their domains,
with $\sum_j a_j^2>0$ everywhere on the base. Normalize them so
$\sum_j a_j^2=1$. In a local trivialization $(z,t)$ put the $j$th component
equal to $a_j(z)\GIIee(\widetilde\eta(t))$, extending by zero outside its
support. The support stays inside the chart, so the extension is smooth.
Transition functions change only the component's phase. The vector has
norm one, and all components have the same vertical frequency.
Compactness gives a finite Lipschitz bound. If the torus or character is
trivial the same construction, or the constant function, applies.
\end{proof}

The quotient Lie algebra is a direct sum $\mathfrak c\oplus P$, where
$\mathfrak c$ consists of words with no petals and $P$ of words with one.
Here $P$ is an abelian ideal: a bracket of two such words has two petals
and is zero. Thus $\widetilde G_{\mathrm{petal}}=\exp P$ is abelian, normal
and rational by Lemma~\ref{v15:gi:8:lem:hall}. Decompose $P=\bigoplus_{a\in L}P_a$
by the unique petal letter. Core conjugation preserves each $P_a$.
For a tuple $t=(t_a)\in\GIIR^L$, the map fixing core generators and sending
$X_a$ to $t_aX_a$ on linear petals is a Lie endomorphism $M_t$:
it preserves all defining ideals. It scales $P_a$ by $t_a$, so on the petal
group, written multiplicatively,
\[
 M_tM_u=M_{tu},\qquad M_t(p)M_u(p)=M_{t+u}(p),\qquad
 M_t(gpg^{-1})=gM_t(p)g^{-1}.
\]
The last identity follows because $M_t$ commutes with every core adjoint
action and conjugation by $P$ acts trivially on $P$. This proves the precise
identities used in (12.8)--(12.10), without treating $M_t$ as an automorphism
when some $t_a=0$. For integral tuples $t$, $M_t$ preserves
$\widetilde\Gamma$, because it sends each generator to an integer power.
It also preserves the petal lattice by its subgroup invariance. These are
the lattice facts needed by the subsequent shear semidirect products.

\subsection{The comparison filtration and corrected Lemma 12.1}
We record sufficient input explicitly. Let $G/\Lambda$, $G^0/\Lambda^0$
be fixed taut rational degree--rank nilmanifolds of degrees at most $(m,R)$
and strictly below $(m,R)$, respectively. Let $\phi:U\to G$ be filtered
and send $\Gamma$ into $\Lambda$. Suppose $a_h,c_h$ are polynomial sequences
and $F_h$ are unit-vector-valued uniformly Lipschitz functions on
$G/\Lambda\times G^0/\Lambda^0$ such that
\begin{equation}\label{v15:gi:8:eq:Fchar}
 F_h(\phi(t)x,z)=\GIIee(\eta(t))F_h(x,z)\quad(t\in T).
\end{equation}
Let $E_{h,i}=\prod_{a\in C_i\cup L_i\cup B_i}e_a^{\xi_{h,a}}$ in a fixed
order, and suppose the horizontal Taylor coefficient of $a_h\Lambda$ is
that of $\phi(E_{h,i})$ in $G_{(i,1)}/G_{(i,2)}$ modulo the horizontal
lattice. Let
\[
 b_h(n)=\prod_{i,a\in C_i}\widetilde e_a^{\xi_{h,a}\binom ni}
           \prod_{i,a\in L_i}\widetilde e_a^{\xi_{h,a}\binom ni}.
\]
These are exactly the core and retained-petal factors in (12.3)--(12.5).
Using the function from Lemma~\ref{v15:gi:8:lem:vertical}, set
\[
 \chi_h(n)=F_h(a_h(n)\Lambda,c_h(n)\Lambda^0),\qquad
 \widetilde\chi_h(n)=\widetilde F(b_h(n)\widetilde\Gamma).
\]

\begin{proposition}\label{v15:gi:8:prop:comparison}
Under these hypotheses, every component of
$\chi_h\otimes\overline{\widetilde\chi_h}$ is a nilsequence of
degree--rank strictly below $(m,R)$. The representing filtered nilmanifold
is fixed independently of $h$ and the coefficients; its Lipschitz bounds
depend only on the fixed data and original bounds.
\end{proposition}
\begin{proof}
For $d>0$ put $\bar r=\max(r,1)$ and let $\sigma(d,\bar r)$ denote the
successor of $(d,\bar r)$ in the positive degree--rank list. In the product
$G\times G^0\times\widetilde G$, let $H_{(d,r)}$ be the connected subgroup
generated by
\begin{align*}
 J_{(d,r)}&=G_{\sigma(d,\bar r)}\times G^0_{(d,\bar r)}
                        \times\widetilde G_{\sigma(d,\bar r)},\\
 D(U_{(d,\bar r)})&=\{(\phi(u),1,q(u)):u\in U_{(d,\bar r)}\}.
\end{align*}
Set $H_0=H_{(1,0)}$. This specifies the zero group; it is not an arbitrary
full product. All generator subgroups are rational. For the diagonal image,
the image of the appropriate lattice generates its real subgroup by the
Hall coordinates and $\phi(\Gamma)\subseteq\Lambda$. The rational-generated
subgroup lemma of Section~\ref{v15:gi:GTZ-LP13-MULTILINEARITY} applies.

Nesting is direct. For positive indices $i,j$, commutators of two diagonal
generators lie in $D(U_{i+j})$. A commutator involving a $J_i$ generator
has its first and third coordinates in the successor group at $i+j$,
since $\sigma(i)+j\succeq\sigma(i+j)$, with the convention replacing
rank zero by one. Its middle coordinate is in $G^0_{i+j}$. The same
argument works for two $J$ generators. Taking generated Lie algebras and
BCH gives $[H_i,H_j]\subseteq H_{i+j}$. At zero use
$H_0=H_{(1,0)}=H_{(1,1)}$ and nesting. Thus this is a taut rational
degree--rank filtration. Its top group is exactly
\begin{equation}\label{v15:gi:8:eq:Htop}
 H_{(m,R)}=D(T),
\end{equation}
because every group in $J_{(m,R)}$ is trivial.

We next check actual polynomial membership, including constant coefficients
and the distinction between a horizontal vector and its torus class.
Write $a_h(0)=v_h\lambda_h$ with $v_h$ in a fixed compact representative
set and $\lambda_h\in\Lambda$. Replace $a_h(n)$ by
$v_h^{-1}a_h(n)\lambda_h^{-1}$, and replace $F_h(x,z)$ by $F_h(v_hx,z)$.
Do the same in the second coordinate. Values of $\chi_h$ are unchanged,
the new sequences start at the identity, and Lipschitz bounds stay uniform.
Since $G$ is taut, conjugation by $v_h$ acts trivially on every
$G_{(i,1)}/G_{(i,2)}$: its commutator is in
$G_{(i+1,1)}\subseteq G_{(i,2)}$. The horizontal torus conditions remain
unchanged. Equation \eqref{v15:gi:8:eq:Fchar} remains unchanged since the top group
of the taut group is central. Tautify the original targets before constructing
the total-frequency representation, using Section~\ref{v15:gi:GTZ-LP-TAUTIFICATION} and then
\cite[Lemma 9.12]{v15-gi-8-GTZ}; no conjugated, potentially irrational replacement
of $\phi$ is being silently used here.

For each $i$, choose $\lambda_{h,i}\in\Lambda\cap G_i$ lifting the
horizontal lattice difference between the new $a_h$ and $\phi(E_{h,i})$.
Such a lift exists by the definition of the horizontal lattice. Multiply
$a_h(n)$ on the right by the inverse of
$\prod_i\lambda_{h,i}^{\binom ni}$. This does not change its orbit.
The factor is polynomial by the Taylor criterion
\cite[Lemma B.9]{v15-gi-8-GTZ}. The horizontal Taylor map is a homomorphism
\cite[Definition 9.6 and the argument following it]{v15-gi-8-GTZ}, so now the
horizontal equalities hold in the vector groups themselves.

Taylor-expand $a_h(n)=\prod_{i\geq1}a_{h,i}^{\binom ni}$. Then
$a_{h,i}\phi(E_{h,i})^{-1}\in G_{(i,2)}$. Hence
$(a_{h,i},1,q(E_{h,i}))\in H_{(i,1)}$. The Taylor criterion and product
closure \cite[Lemma B.9 and Corollary B.4]{v15-gi-8-GTZ} show that
\[
 n\mapsto\left(a_h(n),1,\prod_iq(E_{h,i})^{\binom ni}\right)
\]
is polynomial for $H$. Both its last coordinate and $b_h$ have the same
horizontal Taylor coefficients, since multiplication of factors is addition
in each horizontal quotient and discarded generators have image one.
Their ratio is polynomial and all its horizontal Taylor coefficients vanish.
Its Taylor coefficient at index $i$ therefore lies in
$\widetilde G_{(i,2)}$. These groups, with zeroth group
$\widetilde G_{(1,2)}$, form a degree filtration by the same commutator
check. The ratio starts at one; Lemma B.9 now shows that it is polynomial
for this smaller filtration. Embedded in the third coordinate it is
polynomial in $H$. The second-coordinate map $(1,c_h(n),1)$ is also
polynomial in $H$. Product closure proves that the exact triple
$(a_h(n),c_h(n),b_h(n))$ is polynomial in $H$.

Restrict $F_h\otimes\overline{\widetilde F}$ to
$H/(H\cap(\Lambda\times\Lambda^0\times\widetilde\Gamma))$.
On $D(t)$ its vertical multiplier is
$\GIIee(\eta(t)-\widetilde\eta(q(t)))=1$ for every $t\in T$.
It therefore descends through the central rational subgroup
\eqref{v15:gi:8:eq:Htop}. The central quotient and lattice argument in
Section~\ref{v15:gi:GTZ-LP13-MULTILINEARITY} applies here: this kernel is contained in
every nontrivial indexed subgroup, unlike the earlier kernel $G^*$.
The top group becomes trivial, so the quotient has degree--rank strictly
below $(m,R)$. Restriction, quotient charts, and the compact translations
cost only constants from the fixed geometry. This proves the proposition.
\end{proof}

When $R=1$, strictly below $(m,1)$ means ordinary degree at most $m-1$.
For $R>1$ it means degree--rank at most $(m,R-1)$, not necessarily degree
at most $m-1$. The first sentence of the published proof of Lemma 12.1
claims the latter unnecessarily strong conclusion. The proved former
conclusion is exactly what the following rank-induction step consumes.

\subsection{Uniformity, limits, and the remaining Section 12 interfaces}
Every group, ideal, Hall list, lattice and quotient above depends only on
$m,R$ and the fixed finite labelled dimension data. All formulas are finite
polynomials, Lie-group maps, or choices from fixed lattice and compact
representative sections. A fixed ordering chooses the first lift of a lattice
class; a fixed finite atlas chooses a bounded representative. Applying these
standard choices to representative sequences gives limit maps. The derivative
and vertical identities are required on their full ultrapower domains,
as in Section~\ref{v15:gi:GTZ-LP-TAUTIFICATION}; identities only at standard points are insufficient.
Thus no bound depends on $N$, on $h$, or on the real polynomial coefficients.
The character is fixed before choosing a vertical test function. No assertion
of a bound uniform over unbounded characters is made. Taking the dense
subfamily with fixed geometric data and character is the separate selection
step in Lemma 10.10 and Section 12, using Appendix A.

The ordinary degree filtration used to build the final semidirect product
is $\widetilde G_d=\widetilde G_{(d,0)}$. The petal ideal and scalar-action
identities proved above give precisely the normality, commutativity, and
integral stability used on pp.~1312--1314. In particular a bracket with two
petal entries is zero at \emph{every} weight. This assertion would fail if
one quotiented only by the old top-weight kernel.

There is a separate unit-norm issue at the cutoff in (12.11) and the paragraph
following it. A compactly supported scalar cutoff times a unit-vector function
is not itself unit-vector-valued. Here is an exact completion for that use.
On the final compact nilmanifold choose a smooth vertical-equivariant local
unit vector $V$ on the torus-coordinate chart, agreeing with
$\widetilde F(g\widetilde\Gamma)$ there. Choose a smooth function
$\theta$ of the torus coordinates equal to zero on the closed inner box and
equal to $\pi/2$ outside a slightly larger box still inside the chart.
These coordinates are unchanged by the top vertical action. Construct a
global unit-vector function $V_0$ of the same vertical character by
Lemma~\ref{v15:gi:8:lem:vertical}, applied to the final top torus. Then
\[
 \bigl(\cos\theta\,V,\ \sin\theta\,V_0\bigr)
\]
is smooth, unit-vector-valued, has the required vertical character, and equals
$(V,0)$ on the inner box. The first component extends by zero, since its
support stays inside the chart. Choose $\theta$ constant near both boundaries,
so both extensions are smooth. This preserves the required values for the
dense parameter set after shrinking to the inner box. Its geometry and bounds
are fixed before $N$. The final semidirect-product filtration and polynomial
orbit must still be checked at their displayed definitions; they are not
consequences of the character descent alone.

\subsection{Related applications of the quotient construction}
The failure mechanism is restriction of a generated ideal or filtration to
too few weights, followed by an unjustified quotient or lattice identification.
The audit scope is the universal degree--rank definitions and the consumers
of their commutator/Taylor quotients in Sections 9--12. It includes the
zero-index and repeated normalization assertions in those passages.
\begin{enumerate}
\item Definition 9.1, p.~1279: use the fixed $R=r_*$ in the top relation,
and count the $\ell$ actual entries. Equations \eqref{v15:gi:8:eq:truncate} and
\eqref{v15:gi:8:eq:filtration}, including all higher weights, give the intended
object. Lemmas~\ref{v15:gi:8:lem:filtration} and \ref{v15:gi:8:lem:hall} prove its filtration
and lattice, rather than inheriting typographical indices.
\item Definition 9.6, pp.~1281--1282: use
$\GIIHH_i=G_{(i,1)}/G_{(i,2)}$, and $G_{(1,2)}=G_{(2,0)}$.
The vector/torus distinction is discharged in
Proposition~\ref{v15:gi:8:prop:comparison}, not by replacing a congruence by equality.
\item Definition 9.11 and Lemma 9.12, pp.~1284--1287: their $G_i$ means
$G_{(i,0)}$. First tautify, then choose lattice generators and construct
the filtered map. This order preserves a genuine lattice map; the proof
does not assume arbitrary real conjugation preserves the lattice.
\item Lemma 9.14 and Corollary 9.15, pp.~1287--1288, are explicitly stated
there not to be used later. Their comparison-filtration argument has the
same zero-index and horizontal-lattice issues. The construction in
Proposition~\ref{v15:gi:8:prop:comparison} supplies that mechanism with $q$ trivial
when the character is zero. Their broader reuse requires this version,
not the unqualified original paragraph.
\item Lemma 10.8, p.~1297: the vertical domain is
$U_{(m,R)}$, not the larger ordinary degree group $U_m$.
This unit consistently defines vertical frequencies only on that top group.
The source argument's other linear-algebra steps are not newly reproved here.
\item Lemma 10.10 and Theorem 11.1, pp.~1299--1307: use the final
core/linear/independent decomposition and the exact two vanishing assertions.
The bias and independence premises remain named upstream inputs.
The repeated (iii),(iv) of Theorem 11.1 are redundant. Reordering linear
and independent petals requires relabelling every associated index, not
changing which generators are annihilated.
\item Quotient construction, p.~1309: Theorem~\ref{v15:gi:8:thm:quotient} proves
the full kernel, every filtered quotient, and the top-lattice saturation
needed for integral vertical descent. The kernel is not asserted central.
\item Lemma 12.1, pp.~1309--1311: Proposition~\ref{v15:gi:8:prop:comparison}
proves the rank drop with actual constant normalization, horizontal lattice
lifts, exact polynomial orbit and conjugate vertical cancellation.
The degree-only strengthening is not used or claimed.
\item Petal actions, pp.~1311--1314: the proved abelian ideal and its
individual-petal splitting verify (12.8)--(12.10), including zero and
negative scaling parameters and integral lattice stability. The displayed
final multidegree filtration, chart gluing, and polynomial orbit remain
explicit downstream construction checks for LP-00/LP-07. The unit-norm
cutoff completion above removes the identified local norm obstruction.
\end{enumerate}
The list concerns this quotient mechanism and its stated applications; it does not assert a classification of all possible errors in the cited source.


\section{The Section 12 semidirect-product join}\label{v15:gi:GI-SECTION12}
\newcommand{\GIJR}{\mathbb R}
\newcommand{\GIJZ}{\mathbb Z}
\newcommand{\GIJC}{\mathbb C}
\newcommand{\GIJN}{\mathbb N}
\newcommand{\GIJee}{\mathrm e}
\newcommand{\GIJLie}{\operatorname{Lie}}
\newcommand{\GIJid}{\mathrm{id}}
\newcommand{\GIJAd}{\operatorname{Ad}}


This section supplies the final geometric
construction on published GTZ pp.~1313--1315 from the explicit petal quotient
and comparison interface of LP04. The analytic frequency input, rank iteration and other inverse
theorem interfaces remain separate. We neither invoke the final inverse
theorem to prove an intermediate construction nor import the unpublished
erratum as a theorem. All quotients below are right lattice quotients.

\subsection{Exact input and output}
Fix the data of LP04: integers $m\geq1$, $1\leq R\leq m$, fixed finite
sets of core and retained petal letters with positive integer weights,
and the weighted free nilpotent quotient
\[
 G=\exp\mathfrak g,\qquad
 \mathfrak g=\mathfrak c\oplus\mathfrak p.
\]
The homogeneous relations discard all words of weight greater than $m$,
all weight-$m$ words of length greater than $R$, and all words containing
at least two petals. The last relation holds at every weight.
Thus $\mathfrak p$ is an abelian ideal. Write $P=\exp\mathfrak p$,
$G_j=G_{(j,0)}$, $P_j=P\cap G_j$, $G_0=G_1=G$, and $P_0=P_1=P$.
Let $\Gamma$ be the lattice generated by the exponentials of the letters,
and $\Lambda=P\cap\Gamma$. LP04's Hall-coordinate lemma supplies the
exact integer lattice coordinates and rationality of $G_j,P_j$.
Put $T=G_{(m,R)}$, and let $\eta:T\to\GIJR$ be the descended continuous
character integral on $T\cap\Gamma$.

There are $\ell$ retained petal letters. For $t\in V=\GIJR^\ell$, let
$M_t:P\to P$ scale the summand with unique petal letter $a$ by $t_a$.
We use the following proved LP04 identities, for all real parameters:
\begin{equation}\label{v15:gi:9:eq:actions}
 M_t(p)M_u(p)=M_{t+u}(p),\qquad
 M_t(gpg^{-1})=gM_t(p)g^{-1},\qquad M_t(P_j)\subseteq P_j.
\end{equation}
For integral $z$, $M_z(\Lambda)\subseteq\Lambda$. These are
endomorphisms; $M_0$ is not asserted to be invertible.

The family to be represented has the form
\begin{equation}\label{v15:gi:9:eq:inputfamily}
 b_h(n)=g_0(n)M_{\{\alpha h\}}(g_1(n)),\qquad
 \widetilde\chi_h(n)=\widetilde F(b_h(n)\Gamma).
\end{equation}
Here $g_0:\GIJZ\to G$ and $g_1:\GIJZ\to P$ are polynomial for the ordinary
weight filtrations; $\alpha\in\GIJR^\ell$ is unrestricted. Braces denote
one fixed half-open fundamental cube of $V/\GIJZ^\ell$, for example
$[-1/2,1/2)^\ell$. The factors used in GTZ are ordered products
$e_a^{\xi_a\binom ni}$ for core letters of weight $i$, and the commuting
product $\prod_{a\ \mathrm{petal}}e_a^{\beta_a\binom n{w(a)}}$ for $g_1$.
Core coefficients are independent of $h$. Identifying an arbitrary input
family with this specific factorization is an upstream obligation; it is
not inferred from the names of the factors.
For the actual coefficients in (12.3)--(12.4), the algebraic identification
is exact: $M_t(\exp(uX_a))=\exp(t_a uX_a)$, and $M_t$ is a
homomorphism on the abelian petal group. Hence applying
$M_{\{\alpha h\}}$ to the displayed $g_1$ gives precisely the product
in (12.4). No reordering past a core factor is involved. The hypotheses
that produce these coefficient data are still the separate analytic input.

The function $\widetilde F$ is a choice made below, before applying the
LP04 comparison. It is not an arbitrary previously fixed function having
only $T$-equivariance. This order of choices is necessary.

\begin{theorem}\label{v15:gi:9:thm:join}
One can choose $\widetilde F$ on the fixed $G/\Gamma$, with unit vector
norm and the required $T$-frequency $\eta$, such that the following holds.
There is a fixed connected simply connected nilpotent group $H$, a lattice
$\Delta$, and a rational taut $\GIJN^2$-filtration supported on the downset
below $(1,m)$. The nilpotency class is at most $m+1$. There is a finite
list of fixed smooth unit-vector functions on $H/\Delta$, all with the
same top vertical character, with fixed finite Lipschitz bounds.

For every $\alpha,g_0,g_1$ as above and every finite nonempty parameter set
$S\subseteq\GIJZ$, some $S'\subseteq S$ has $|S'|\geq |S|/J$, where the
fixed integer $J$ depends only on $\ell$, and one function $F'$ in this
list satisfies
\begin{equation}\label{v15:gi:9:eq:exactoutput}
 F'(O(h,n))=(\widetilde\chi_h(n),0)
 \quad(h\in S',\ n\in\GIJZ).
\end{equation}
Here $O$ is a polynomial orbit of multidegree at most $(1,m)$.
The zero padding is a fixed isometric coordinate inclusion, so projection
onto the original coordinates gives literal equality with
$\widetilde\chi_h$. No scalar component is divided by.

All constants, manifolds and functions are independent of $\alpha$,
the polynomial coefficients, $S$ and $N$. The conclusion transfers to the
full internal ultrapower domains and preserves a dense set of limit
parameters. It supplies this construction and its LP04 comparison join;
it does not by itself establish the other premises of Theorem 7.1 or GI.
\end{theorem}

\subsection{Extend the character before choosing the test}
\begin{lemma}\label{v15:gi:9:lem:extension}
The character $\eta$ extends to a continuous character
$\bar\eta:G_m\to\GIJR$ integral on $G_m\cap\Gamma$.
\end{lemma}
\begin{proof}
The weight-$m$ group is central. Its surviving Hall commutators have
weight $m$ and lengths at most $R$. At this weight the logarithm of each
group commutator is exactly its Lie bracket: every correction has greater
weight and vanishes. By LP04's exact Hall coordinates, these logarithms
form a real basis and their integer span is precisely the logarithm of
$G_m\cap\Gamma$. The subset of length $R$ is a basis of $T$, and its
integer span is $T\cap\Gamma$. Assign to those basis elements the values
of $\eta$, and assign zero to all remaining weight-$m$ basis elements.
These assigned values are integers. Linear extension on this central
vector group proves the assertion. The construction depends only on the
fixed Hall ordering and $\eta$.
\end{proof}

The central torus $G_m/(G_m\cap\Gamma)$ acts freely on $G/\Gamma$.
For any integral character of this torus there is a smooth equivariant
unit-vector function of finite dimension. Here is the construction we use
twice. Choose finitely many local smooth sections of the torus quotient
and real smooth functions $a_j$ supported in their domains with
$\sum a_j^2>0$ everywhere. Such functions can be obtained by smooth
Euclidean bumps in a finite collection of relatively compact local charts
whose smaller charts cover the compact base. Replace $a_j$ by
$a_j/(\sum a_i^2)^{1/2}$. In the $j$th local trivialization $(x,z)$ define
the $j$th component to be $a_j(x)\GIJee(\bar\eta(z))$ and extend it by zero.
The support lies inside its chart. Thus each component is globally smooth,
has the stated vertical transformation law, and the squared vector norm
is $\sum a_j^2=1$. Local sections can be seen directly: choose a linear
complement to the central Lie subspace. Centrality gives
$\exp(X+Z)=\exp X\exp Z$, so the complement gives smooth local lifts
from the group quotient. The image lattice is discrete by the Hall
coordinates with the central weight-$m$ coordinates deleted. Small
coordinate neighborhoods with disjoint lattice translates give the same
local lifts on the compact quotient.

Choose this function as $\widetilde F$. It has the original frequency
$\eta$ on $T$ and now also the required frequency on all of $G_m$.
Every cancellation in LP04's comparison uses only the restriction to $T$,
so this choice preserves that proof. The comparison does not require a
specific previously selected $\widetilde F$.

\paragraph{Adverse test.}
$T$ can be strictly smaller than $P_m$. For example, with $m=R=2$
and just one petal generator of weight two, $T$ is trivial whereas
$G_m=P\simeq\GIJR$. The unit-vector function
$(\cos 2\pi x,\sin 2\pi x)$ on $\GIJR/\GIJZ$ satisfies the vacuous $T$-law
but has no scalar vertical character for all translations. Consequently
$T$-equivariance alone cannot be cited as the final multidegree frequency.
The integral extension and new choice above are the repair, not an
assertion that every old test already had the stronger property.

\subsection{Both semidirect products and their signs}
Set $K=G\ltimes P$ with multiplication
\begin{equation}\label{v15:gi:9:eq:Klaw}
 (g,p)(g',p')=(gg',\,(g')^{-1}pg'\,p').
\end{equation}
Associativity follows either directly from conjugation or from the
injective homomorphism $(g,p)\mapsto(g,gp)$ into $G\times G$.
Its image is the closed subgroup of pairs agreeing modulo $P$.
The inverse is $(g^{-1},gp^{-1}g^{-1})$. Its underlying manifold is
$G\times P$, so $K$ is connected and simply connected.

Define $\rho_t(g,p)=(gM_t(p),p)$. To check multiplication, applying
$\rho_t$ to \eqref{v15:gi:9:eq:Klaw} gives first component
$gM_t(p)g'M_t(p')$, by \eqref{v15:gi:9:eq:actions}. Multiplying the two images
has the same first component. Its second is
\[
 M_t(p')^{-1}(g')^{-1}pg'M_t(p')p'=(g')^{-1}pg'p',
\]
since all terms surrounding $p'$ are in the abelian group $P$.
Thus $\rho_t$ is a homomorphism. Further,
$\rho_t\rho_u=\rho_{t+u}$ and $\rho_0=\GIJid$, so $\rho_t$ is an
automorphism with inverse $\rho_{-t}$. This conclusion uses the additive
shear action, not invertibility of $M_t$.

Set $H=V\ltimes_\rho K$ with the right-action convention
\begin{equation}\label{v15:gi:9:eq:Hlaw}
 (t,k)(t',k')=(t+t',\rho_{t'}(k)k').
\end{equation}
In three coordinates this reads
\[
 (t,g,p)(t',g',p')
 =(t+t',\,gM_{t'}(p)g',\,(g')^{-1}pg'p').
\]
The action identity proves associativity, and the inverse of $(t,k)$ is
$(-t,\rho_{-t}(k^{-1}))$. Smoothness follows from the polynomial BCH
formulas and the linear maps on $\mathfrak p$. The manifold
$V\times G\times P$ is connected and simply connected.

For a complete nilpotence and filtration check, use Lie coordinates
$(t,X,Y)\in V\oplus\mathfrak g\oplus\mathfrak p$. The bracket is
\begin{equation}\label{v15:gi:9:eq:bracket}
 [(t,X,Y),(t',X',Y')]
 =\bigl(0,\,[X,X']+M_{t'}Y-M_tY',\,[X,Y']+[Y,X']\bigr).
\end{equation}
Indeed the embedding of $K$ has derivative $(X,Y)\mapsto(X,X+Y)$,
giving the two bracket terms of $K$. The derivative of $\rho_t$ is
$(X,Y)\mapsto(X+M_tY,Y)$. Conjugation by $(t,1)$ on $K$ is
$\rho_{-t}$ in convention \eqref{v15:gi:9:eq:Hlaw}, giving the signs displayed.

Assign bidegree $(0,i)$ to core words of weight $i$ in the $X$ copy,
$(1,i)$ to petal words of weight $i$ in the $X$ copy, $(0,i)$ to petal
words of weight $i$ in the $Y$ copy, and $(1,0)$ to $V$.
Every nonzero term of \eqref{v15:gi:9:eq:bracket} has the sum of the input
bidegrees. For example, $[V,Y_i]$ lies in the first petal copy of degree
$(1,i)$, and brackets involving two petal words vanish. All weights are
at most $m$ and the first degree is at most one. Each generator has
positive total degree, so every bracket of length greater than $m+1$
vanishes. This proves nilpotence of $H$ of class at most $m+1$.

\subsection{The complete filtration and the lattice}
Define $H_{a,j}=1$ for every $a>1$ and every $j\geq0$.
The remaining positive-$j$ formulas, together with the zero-index groups, are
\begin{align*}
 H_{0,0}&=H, & H_{a,j}&=1\quad(a>1),\\
 H_{1,0}&=\{(t,p,1):t\in V,p\in P\},&
 H_{1,j}&=\{(0,p,1):p\in P_j\},\\
 H_{0,j}&=\{(0,g,p):g\in G_j,p\in P_j\}.
\end{align*}
For $j>m$ these groups are trivial. These definitions are also the sums
of the homogeneous Lie spaces whose bidegrees dominate the index.
Thus the groups are connected, descending and normal; brackets lie in
the sum-index subgroup by \eqref{v15:gi:9:eq:bracket}, and the BCH commutator
terms do also. To spell out the boundary cases, two $H_{1,0}$ elements
commute, $[H_{1,0},H_{0,j}]\subseteq H_{1,j}$, and
$[H_{1,i},H_{0,j}]\subseteq H_{1,i+j}$. Brackets with index zero are
covered by normality. The groups $H_{1,0}$ and $H_{0,1}=K$ generate
$H$, which proves tautness. In particular
\begin{equation}\label{v15:gi:9:eq:top}
 Z_*:=H_{1,m}=\{(0,p,1):p\in P_m\}
\end{equation}
is central. Importantly $H_{0,1}$ contains the full $G_1=G$ in its first
group coordinate, not merely $P$.

Let $\Lambda_K=\Gamma\ltimes\Lambda$ with law \eqref{v15:gi:9:eq:Klaw}.
Since $P$ is normal, $\Gamma$ normalizes $\Lambda$. Therefore this is
a subgroup, including inverses. For $z\in\GIJZ^\ell$, both $\rho_z$
and $\rho_{-z}$ preserve $\Lambda_K$, by the integral LP04 action.
Consequently
\[
 \Delta=\GIJZ^\ell\ltimes_\rho\Lambda_K
       =\{(z,\gamma,\lambda):z\in\GIJZ^\ell,
                   \gamma\in\Gamma,\lambda\in\Lambda\}
\]
is a subgroup of $H$. It is discrete in the product manifold topology.

Here is explicit cocompactness. Right-multiply $(t,k)$ by $(-z,1)$,
where $z$ reduces $t$ into $[0,1]^\ell$. This gives
$(t-z,\rho_{-z}k)$. For an arbitrary $(g,p)\in K$, choose
$\gamma\in\Gamma$ with $g\gamma^{-1}$ in a fixed compact
representative set for $G/\Gamma$. Right multiplication by
$(\gamma^{-1},1)$ gives $(g\gamma^{-1},\gamma p\gamma^{-1})$.
Now reduce the second coordinate into a fixed compact representative
set for $P/\Lambda$ by right multiplication by $(1,\lambda^{-1})$.
The resulting compact product set represents every right coset of
$\Lambda_K$. Applying this after the $t$ reduction gives a compact
product set representing every $\Delta$ coset. No bound on the
intermediate conjugate or on $z$ is needed.

The same argument inside $G_j\ltimes P_j$, using the already proved
lattices $\Gamma\cap G_j$ and $\Lambda\cap P_j$, proves rationality
of $H_{0,j}$. The intersections for $H_{1,j}$ are exactly
$\Lambda\cap P_j$, and for $H_{1,0}$ exactly $\GIJZ^\ell\times\Lambda$.
These are cocompact by LP04 and the product construction. This checks
all filtered lattice intersections, including zero and trivial groups.
The topology, filtration, lattice and smooth metric on the compact
quotient are fixed before any orbit coefficients are chosen.

\subsection{Polynomial orbit and exact reduction}
Define
\begin{equation}\label{v15:gi:9:eq:orbit}
 O(h,n)=(0,g_0(n),g_1(n))(\alpha h,1,1)\Delta.
\end{equation}
We verify polynomiality in the full group before taking the quotient.
The source filtration on $D=\GIJZ^2$ is
$D_{0,0}=\GIJZ^2$, $D_{1,0}=\GIJZ\times\{0\}$,
$D_{0,1}=\{0\}\times\GIJZ$, and $D_{a,j}=\{0\}$ at every other
nonzero index. It is an $\GIJN^2$-filtered abelian group.
The embeddings $g\mapsto(0,g,1)$ and $p\mapsto(0,1,p)$ are
homomorphisms taking the weight-$j$ groups to $H_{0,j}$.
Each independent-of-$h$ factor is therefore polynomial on $\GIJZ^2$ with
its coordinate multidegree filtration: $h$ differences vanish, and
$n$ differences have the asserted $H_{0,j}$ membership.
The map $(h,n)\mapsto(\alpha h,1,1)$ is a homomorphism in $h$ with
values in $H_{1,0}$, has zero second $h$ difference, and is constant
in $n$. Thus it too is polynomial for this same filtration.
Published Corollary B.4, GTZ p.~1334, says that polynomial maps between
the same index-filtered groups are closed under pointwise multiplication.
Here both source and target are the precise $\GIJN^2$-filtered groups
just checked. Applying it to these three factors proves polynomiality
of \eqref{v15:gi:9:eq:orbit}, including all mixed derivatives and all integer inputs.
This is the only external polynomial-product theorem used in this join.

For the explicit inputs preceding Theorem~\ref{v15:gi:9:thm:join}, their univariate
polynomiality can also be checked before this use. An $r$-fold difference
of $e_a^{c\binom ni}$ is $e_a^{c\Delta^r\binom ni}$ because all its
values lie in one abelian one-parameter subgroup. It vanishes for $r>i$;
for $r\leq i$ it lies in $G_i\subseteq G_r$. This holds for arbitrary
integer increments, including negative ones. Products are polynomial by
the same Corollary B.4; for $g_1$ all factors already commute in $P$.
No Taylor expansion with a potentially unjustified centrality hypothesis
is imported for this check.

Write $z_h=\alpha h-\{\alpha h\}\in\GIJZ^\ell$. Formula
\eqref{v15:gi:9:eq:Hlaw} gives a representative of \eqref{v15:gi:9:eq:orbit} as
\[
 (\alpha h,g_0(n)M_{\alpha h}(g_1(n)),g_1(n)).
\]
Right multiplication by $(-z_h,1,1)\in\Delta$ changes this to
\begin{equation}\label{v15:gi:9:eq:reduced}
 (\{\alpha h\},g_0(n)M_{\{\alpha h\}}(g_1(n)),g_1(n)).
\end{equation}
The equality follows from the additive identity in \eqref{v15:gi:9:eq:actions}.
It is exact for every $h,n$; there is no illegal commutation of a
fractional-part operation with a nonabelian product.

\subsection{Fixed charts, unit completion and uniform bounds}
Projection onto $t$ gives a smooth map
$\pi:H/\Delta\to V/\GIJZ^\ell$. If a box $U\subset V$ has width less
than one in each coordinate, projection is injective on $U$. On the
corresponding open torus box there is an exact smooth trivialization
\begin{equation}\label{v15:gi:9:eq:chart}
 U\times K/\Lambda_K\longrightarrow\pi^{-1}(U\bmod\GIJZ^\ell),
 \qquad(t,k\Lambda_K)\longmapsto(t,k)\Delta.
\end{equation}
Indeed two representatives in $U$ differ by an element of $\Delta$
only if their integer $t$ displacement is zero; their $k$ coordinates
then differ precisely by right multiplication in $\Lambda_K$.
Surjectivity follows by reducing the $t$ coordinate to its unique lift
in $U$. The formulas are smooth on this chart.

Choose finitely many fixed closed boxes $I_j$ covering the closed
fundamental cube $[-1/2,1/2]^\ell$, with fixed slightly larger boxes
$U_j$ of widths less than one. They may extend beyond the cube's
boundary; each $U_j$ still projects injectively. Choose a smooth
function $\theta_j$ on the torus which is zero on $I_j$ modulo integers
and $\pi/2$ outside a smaller open box with closure contained in $U_j$.
It can be constructed by a fixed Euclidean bump in $U_j$, equal to one
on $I_j$, followed by multiplication of its complement by $\pi/2$.
The chart margins are fixed positive real numbers, not functions of $N$.
At least one box contains $\{\alpha h\}$ for at least $|S|/J$ of
the parameters, where $J$ is the fixed number of boxes. Fix that box
and call this subfamily $S'$.

On its chart define
\[
 W(t,(g,p)\Lambda_K)=\widetilde F(g\Gamma).
\]
The projection $K/\Lambda_K\to G/\Gamma$ is well-defined and smooth,
since the first coordinate of \eqref{v15:gi:9:eq:Klaw} is a homomorphism and
maps $\Lambda_K$ into $\Gamma$. Thus $W$ is a smooth local unit
vector. Define a character on the actual group \eqref{v15:gi:9:eq:top} by
\[
 \eta'(0,p,1)=\bar\eta(p)\qquad(p\in P_m).
\]
It is integral on $Z_*\cap\Delta$ by Lemma~\ref{v15:gi:9:lem:extension}.
Left multiplication by $(0,p_*,1)\in Z_*$ changes the chart's first
group coordinate to $p_*g$, keeps $t$ fixed, and hence multiplies
$W$ by $\GIJee(\bar\eta(p_*))$. This uses equivariance on $P_m$,
not merely on $T\cap P$.

The torus $Z_* /(Z_*\cap\Delta)$ acts freely on $H/\Delta$:
centrality makes the stabilizer at every point exactly
$Z_*\cap\Delta$. It is compact by the rationality already proved.
The quotient by $Z_*$ has coordinates $V\times(G/P_m)\times P$,
and the image of $\Delta$ has coordinates
$\GIJZ^\ell\times q(\Gamma)\times\Lambda$. Here $q(\Gamma)$ is
discrete by LP04's Hall-coordinate quotient, so the image lattice is
discrete and the quotient base is compact. Central complements and small
coordinate neighborhoods give local lifts exactly as in Section 2.
The finite local-section construction from that section therefore gives a
global smooth unit vector $W_0$ with frequency $\eta'$. Set
\begin{equation}\label{v15:gi:9:eq:completion}
 F'_j=(\cos(\theta_j\circ\pi)W,
                 \sin(\theta_j\circ\pi)W_0).
\end{equation}
The first component is extended by zero outside the chart. It vanishes
on a neighborhood of the chart boundary, so this extension is smooth.
Outside the chart the second coefficient is one. Inside, the squared
norm is $\cos^2\theta_j+\sin^2\theta_j=1$. Both components have the
same vertical frequency because $Z_*$ fixes $\pi$. Therefore $F'_j$
is a genuine unit-vector nilcharacter test, including outside the box.

For $h\in S'$, the representative \eqref{v15:gi:9:eq:reduced} uses the unique
lift $\{\alpha h\}\in I_j$. Its cutoff angle is zero, so
\eqref{v15:gi:9:eq:completion} gives \eqref{v15:gi:9:eq:exactoutput} for every integer $n$.
The second component is exactly zero there, not an approximation.

Every chart used has compact closure in a fixed larger chart, and its
fiber $K/\Lambda_K$ is compact. The displayed projections and their
derivatives are bounded on these sets. The finite bump functions,
$\widetilde F$, $W_0$, and the fixed metrics have bounded derivatives.
Taking the maximum over the finite list proves a finite Lipschitz bound
depending only on the fixed groups, character and charts. Unbounded
polynomial coefficients change only the orbit; they do not enter these
derivative bounds. This also proves fixed finite dimension. If $\ell=0$
the torus is a point, use one chart and $\theta=0$; all arguments still
apply. Trivial top groups or zero characters likewise cause no exception.
This completes the ordinary proof of Theorem~\ref{v15:gi:9:thm:join}.

\subsection{Internal domains and the LP04 consumer}
All groups, charts, test functions and characters above are standard
objects after fixing the selected dimension data and character. All
algebraic identities hold for every real parameter and every group point.
Taking representative sequences gives internal polynomial maps, including
for unbounded limit coefficients. The standard tests extend to the whole
ultrapower of the compact quotient. Their vertical laws then hold for
every limit element of $Z_*$ and every limit point, by transfer of the
actual universal identity. We are not inferring all-limit equivariance
from an identity checked only on standard points.

Finite pigeonhole over the fixed list of boxes preserves positive standard
density of a limit parameter family. Equivalently apply the ordinary
cardinality inequality to representatives and pass to the ultralimit.
After choosing the resulting fixed box, the nilcharacter
$\chi(h,n)=F'_j(O(h,n))$ is defined on all of $\!{}^*\GIJZ^2$, not only
on the dense subfamily. Its exact restriction is the fixed isometric
inclusion of $\widetilde\chi_h$ for all limit integers $n$ there.

Choose $\widetilde F$ by Section 2 before invoking LP04's comparison
proposition. That proposition's top cancellation is unchanged because
$\bar\eta|_T=\eta$. It produces the same strictly lower degree--rank
comparison for $\chi_h\otimes\overline{\widetilde\chi_h}$.
For clarity, unit norm gives the exact component recovery identity
\[
 (\chi_h)_a=\sum_b (\widetilde\chi_h)_b
             \bigl((\chi_h)_a\overline{(\widetilde\chi_h)_b}\bigr).
\]
On the selected subfamily, replace $(\widetilde\chi_h)_b$ by the
corresponding component of the single $\chi(h,n)$. The terms in parentheses
are exactly the lower degree--rank sequences supplied by LP04. The sum
is finite and its coefficients are fixed; no division or nonvanishing of
an individual component is used. The zero padding in
\eqref{v15:gi:9:eq:exactoutput} changes none of these terms.

At $R=1$, lower degree--rank means degree at most $m-1$. At $R>1$
the subsequent rank iteration remains an upstream/downstream assembly
obligation; we do not incorrectly replace it by an immediate ordinary
degree drop. Likewise the frequency decomposition leading to
\eqref{v15:gi:9:eq:inputfamily}, the initial correlation, and the remaining
Theorem 7.1 induction premises are not established by this geometric join.

\subsection{Bounded related-use and consistency audit}
The search scope is the Section 12 construction and its immediate LP04
comparison and Theorem 7.1 consumers, published pp.~1308--1315. We also
check the polynomial-product supplier B.4 at p.~1334. This is an authored
bounded audit, not a claim of independent source-first coverage of GI.
\begin{enumerate}
\item The smaller degree--rank character domain versus the ordinary
top group is repaired by the integral extension and the new test choice.
The repeated final frequency on p.~1315 is explicitly restricted to
$P_m$ and uses $\bar\eta$. A favorable quotient-character review alone
does not supply this extension.
\item The illustrative $H_{0,1}$ on p.~1312 restricts its first group
coordinate to the petal subgroup, although the orbit contains the core
$e_2^{\beta n}$. For nonzero generic $\beta$ its $n$ difference is outside
that alleged group. The general p.~1314 definition uses full $G_j$;
we use and prove that definition, not the defective illustration.
\item Zero or negative scaling parameters do not make every $M_t$
invertible. The actual invertible automorphisms are the shears $\rho_t$,
whose inverse and integer lattice preservation are proved here.
\item The lattice in p.~1314 is a subgroup of $H$, not of the original
$G$. Every filtered lattice intersection and the exact right-coset
reduction are checked above; neither discreteness alone nor a real
endomorphism establishes the claim.
\item The cutoff on pp.~1312 and 1315 cannot alone be a unit-vector
test or guarantee equality on the selected box. Fixed inner boxes,
positive chart margins and \eqref{v15:gi:9:eq:completion} supply unit norm,
literal equality and coefficient-independent Lipschitz bounds.
\item In particular a variable box of arbitrarily small width would
not supply a fixed Lipschitz bound. The predetermined finite cover
removes this possible parameter dependence at a controlled density loss.
\item The p.~1308 claim that pigeonhole puts the fractional parts in a
short interval around zero is not valid for an arbitrary dense family.
For $\alpha=1/2$ and odd $h$, our fractional parts are all $-1/2$.
Our fixed cover allows a box at any location, including across the chosen
fundamental-cube boundary, and does not need a near-zero assertion.
\item Formula (12.11) is checked as right multiplication by an actual
lattice element; the displayed identities do not silently change the
object or introduce a different representative in the local test.
\item A lower rank comparison is not automatically a lower ordinary
degree comparison. The consumer paragraph above preserves the exact
LP04 conclusion and leaves the rank iteration to its actual proof.
\end{enumerate}
Earlier printed statements and earlier frozen project packets are retained;
this proposal does not relabel them as having contained the new argument.
Any other current copies of the smaller-domain frequency claim or the
illustrative filtration must be corrected or marked historical when this
proof is incorporated. The attached occurrence record identifies these
specific copies and the new canonical assertions for review.


\section{Degree conversion and the finite rank iteration}\label{v15:gi:GI-RANK}
\newcommand{\GIKR}{\mathbb R}
\newcommand{\GIKQ}{\mathbb Q}
\newcommand{\GIKZ}{\mathbb Z}
\newcommand{\GIKN}{\mathbb N}
\newcommand{\GIKC}{\mathbb C}
\newcommand{\GIKDR}{\mathrm{DR}}


This section proves the degree conversions and the assembly of the rank steps from the precisely stated one-step geometric/frequency interface. The analytic sunflower and commutator hypotheses are supplied in Section~\ref{v15:gi:GI-SECTION9-11}. Every internal nilcharacter below has exact unit norm and vertical equivariance for all limit elements under the natural ultrapower interpretation.

\subsection{The actual degree-to-rank conversion}

Let $G/\Gamma$ be a connected simply connected nilpotent Lie group with
lattice and rational degree filtration $(G_i)_{i\geq0}$ supported on
$0\leq i\leq m$. Apply accepted bounded tautification first, so
$G_0=G_1=G$ and polynomial orbits may be normalized at zero without
changing the represented sequence. Write $\mathfrak g_i=\log G_i$.
All the vector spaces here are rational for the rational structure
determined by $\Gamma$.

For $(d,r)\in\GIKDR=\{(d,r)\in\GIKN^2:0\leq r\leq d\}$ with $d>0$ define
$\mathfrak h_{(d,r)}$ to be the real span of all iterated Lie brackets
of length $\ell\geq1$ with entries $X_j\in\mathfrak g_{i_j}$,
$i_j\geq1$, for which
\begin{equation}\label{v15:gi:10:eq:DR}
 w:=\sum_{j=1}^{\ell}i_j>d,
 \qquad\hbox{or}\qquad
 w=d\ \hbox{and}\ \ell\geq\max(r,1).
\end{equation}
Take every bracketing and set
$\mathfrak h_{(0,0)}=\mathfrak g$, $H_{(d,r)}=\exp\mathfrak h_{(d,r)}$.
Only finitely many nonzero brackets occur, since their weight exceeds
$m$ when they can no longer contribute and every entry has positive
weight.

\begin{lemma}\label{v15:gi:10:lem:DR}
These are connected rational groups forming a taut degree--rank
filtration supported on $(m,m)$. For each $d>0$,
$H_{(d,0)}=H_{(d,1)}=G_d$. Hence the original degree-polynomial orbits
are exactly the degree--rank polynomial orbits for this refinement.
A degree-$m$ nilcharacter is a degree--rank $(m,m)$ nilcharacter
with the same function and with the top character restricted.
\end{lemma}
\begin{proof}
All bracket spans in \eqref{v15:gi:10:eq:DR} are rational, because brackets of
rational vectors are rational. Bracketing with $\mathfrak g_1$ raises
weight by at least one, so each span is an ideal of $\mathfrak g$.
It is therefore a Lie subalgebra, and its exponential is a connected
normal subgroup, which is closed by the global exponential homeomorphism.
Choose a rational Mal'cev basis for $G/\Gamma$ compatible with its rational
structure. Each bracket span has a rational basis in these coordinates.
Its lower central series is rational, since brackets of rational vectors
are rational. Choose one standard $Q$ bounding the finitely many rational
structure constants and basis coefficients. Green--Tao, \emph{Annals of
Mathematics} 175 (2012), Definition 2.5 and Proposition A.10, applied to
this subgroup with its lower central filtration, supplies a Mal'cev basis
relative to its actual intersection with $\Gamma$. That intersection is
therefore discrete and cocompact. This application uses Lie-algebra
rationality in Definition 2.5; it does not assume the lattice conclusion.

The lexicographic nesting is immediate at a fixed $d$ as $r$ increases.
If $d$ increases, every new allowed weight was strictly above the old
$d$, so its span is contained in the old one as well. For two positive
indices $(d,r),(e,t)$, brackets either have weight greater than $d+e$,
or weight exactly $d+e$ and length at least
$\max(r,1)+\max(t,1)\geq\max(r+t,1)$. Thus the Lie bracket is in
$\mathfrak h_{(d+e,r+t)}$. The BCH commutator and its higher terms
lie there too, using the ideal property. At index zero use normality.
This verifies every pair of indices, not only equal indices.

For $r=0$ or $1$, every allowed bracket has weight at least $d$ and
thus lies in $\mathfrak g_d$. Conversely the length-one brackets with
$i_1=d$ include all of $\mathfrak g_d$. This proves
$H_{(d,0)}=H_{(d,1)}=G_d$. In particular $H_{(1,0)}=G$.
If $d>m$ all brackets vanish. At degree $m$ every nonzero bracket has
length at most its weight $m$, which gives the asserted support.

The domain filtration on $\GIKZ$ has $\GIKZ$ at indices $(0,0),(1,0)$
and zero at all other indices, including $(1,1)$, as in Definition 6.17.
The derivative criterion for an adapted
map therefore asks for its $a$th differences to lie in $H_{(a,0)}$;
these groups equal $G_a$. This proves the equality of the polynomial
classes directly, matching published (6.8). The new top
$H_{(m,m)}$ is a rational subgroup of $G_m$. Restricting the original
vertical character is integral on its lattice and preserves all exact
vertical identities and unit norm. No enlargement of its character
domain is asserted.
\end{proof}

\subsection{The total-degree conversion used in related applications}

For a rational multidegree filtration $(G_a)_{a\in\GIKN^q}$ supported
on a finite downset $J$, first use bounded tautification to restrict
to $H=\langle G_{e_1},\ldots,G_{e_q}\rangle$. Let
\[
 K_0=H,\qquad
 \mathfrak k_i=\sum_{a\in J:\ |a|\geq i}\mathfrak g_a,
 \qquad K_i=\exp\mathfrak k_i\quad(i\geq1).
\]
Each $\mathfrak g_a$ is an ideal in $\mathfrak g_0$, so the sum is
a rational ideal. It equals the Lie algebra of the connected generated
group. The brackets satisfy
$[\mathfrak k_i,\mathfrak k_j]\subseteq\mathfrak k_{i+j}$,
because $[\mathfrak g_a,\mathfrak g_b]\subseteq\mathfrak g_{a+b}$
and $|a+b|=|a|+|b|$. The index-zero condition follows from normality.
Also $K_1=H$ and $K_i$ is trivial above $D=\max_{a\in J}|a|$.
The rational subgroup lemma verifies every lattice intersection.

For the diagonal restriction $n\mapsto g(n,\ldots,n)$ of a
multivariable polynomial orbit, use its Taylor expansion. A factor
with coefficient $u_a\in G_a$ has scalar exponent
$\prod_j\binom n{a_j}$ of ordinary degree $|a|$. Such a factor
is adapted to $K_\bullet$ since $u_a\in K_{|a|}$; its successive
differences have the same one-parameter coefficient and the required
scalar polynomial degree. The product is polynomial by B.4. This
checks the actual diagonal map rather than treating the domain
filtrations as identical. For a rectangle $J=\{a:a\leq d\}$,
the only index of total weight $|d|$ is $d$, so $K_{|d|}=G_d$.
A multidegree-$d$ nilcharacter therefore restricts diagonally to a
degree-$|d|$ nilcharacter with exactly the same top character.
For a nonrectangular downset only the nilsequence degree assertion
has been proved, and no unwarranted common top character is claimed.

\subsection{The precise one-step interface}

Fix $s\geq3$, write $m=s-1$, and let $1\leq R\leq m$.
Suppose $f$ is bounded on $[N]$, $H$ is a dense internal set of
shifts, $A(h,n)$ is a single unit multidegree $(1,m)$ nilcharacter,
and $(\chi_h)$ is an internal uniformly bounded family of unit
degree--rank $(m,R)$ nilcharacters. Suppose that for every retained
$h$, $\Delta_h f$ has positive degree-at-most-$m-1$ correlation with
$A(h,\cdot)\otimes\chi_h$.

The one-step premise used here is the following exact output, after
restricting to a dense internal subset and fixing uniform standard
bounds. It must be supplied by the corrected additive-quadruple result,
the actual Section 10 frequency representation, the two conclusions of
Theorem 11.1, accepted LP04 and the final Section 12 join:
\begin{enumerate}
\item There is a unit multidegree $(1,m)$ nilcharacter $Z(h,n)$,
defined globally, and a family of unit-vector sequences $\zeta_h(n)$
such that $Z(h,n)=\iota\zeta_h(n)$ for retained $h$, with one fixed
coordinate isometry $\iota$. Zero padding is allowed.
\item The entire matrix-valued function
$V_h=\chi_h\otimes\overline{\zeta_h}$ is a nilsequence of
degree--rank at most $(m,R-1)$, with standard uniform geometry and
Lipschitz bounds and internal parameter dependence.
\end{enumerate}
The second clause states a rank drop; for $R>1$ it does not state
an ordinary-degree drop. The final Section 12 constructor must choose
its test equivariant on the full ordinary top before invoking LP04's
comparison. Equivariance only on $G_{(m,R)}$ does not by itself
justify its larger multidegree top. Those obligations remain in the
separately authored geometric join and are not hidden in this premise.

\subsection{Recovering the next family without dividing by components}

The unit-vector identity gives, for every component $a$ and all relevant
integers,
\begin{equation}\label{v15:gi:10:eq:recovery}
 (\chi_h)_a(n)=\sum_b (\zeta_h)_b(n)
       \bigl((\chi_h)_a(n)\overline{(\zeta_h)_b(n)}\bigr).
\end{equation}
The terms in parentheses are components of $V_h$. Replace the first
factor by its corresponding coordinate in $Z(h,n)$ and form the unit
nilcharacter $A'(h,n)=A(h,n)\otimes Z(h,n)$ on the product model.
The product filtration has the same multidegree bound $(1,m)$.

Choose a scalar component and a degree-at-most-$m-1$ test witnessing
the incoming correlation, with fixed bounds and a fixed positive
threshold after the usual finite selection. Inserting
\eqref{v15:gi:10:eq:recovery} expresses it as a fixed finite sum of correlations
of $\Delta_h f$ against components of $A'(h,\cdot)$, the same
ordinary lower-degree test, and a component of $V_h$. At least one
summand has a positive fixed correlation. Finite pigeonhole restricts
to a dense subset with fixed component choices if necessary. No scalar
component is inverted, or assumed nonzero or of modulus one.

On each fixed rational degree--rank model for $V_h$, apply the
unit-character approximation construction to its top group at
$(m,R-1)$. That group is central after tautification, since bracketing
with $G_{(1,0)}$ raises the first coordinate above $m$. Its lattice
is rational. Smooth local sections and a finite unit partition on the
central torus produce exact-character unit-vector tests. The complex linear span of their components
is a unital conjugation-stable algebra separating points; Stone--Weierstrass
and a finite uniform net of the bounded Lipschitz test class give a
finite approximation by components of such tests. This is the same
explicit compact construction in the assembly's selection section,
now on the stated degree--rank model. It uses centrality and rationality,
not an identification of this top with the larger ordinary-degree top.

Choose the approximation error so that its correlation with the other
fixed bounded factors is less than half the surviving threshold.
Another finite selection supplies a unit degree--rank $(m,R-1)$
nilcharacter $\chi'_h$ for which $\Delta_h f$ has positive
degree-at-most-$m-1$ correlation with
$A'(h,\cdot)\otimes\chi'_h$. Finite choices and orbit choices are
made on representatives and then passed to the ultralimit, retaining
internal dependence and uniform bounds. Every test has its exact
standard vertical identity before extension, so the all-limit
convention is preserved. This is the output required by the next
rank step.

\subsection{Initialization, termination and the actual endpoint}

Start with the degree-$m$ family obtained from the preceding inverse
order. Lemma~\ref{v15:gi:10:lem:DR} makes it a degree--rank $(m,m)$ family
with the same functions. Take $A=1$, viewed as a constant unit
multidegree $(1,m)$ nilcharacter. The original positive correlation
is degree-at-most-$m-1$ correlation with the lower test equal to one.

Apply the one-step premise and the preceding recovery argument at
$R=m,m-1,\ldots,1$. There are exactly $m$ steps, a fixed standard
number. A finite sequence of positive standard density losses still
leaves a positive standard density; the same holds for positive
correlation and finite complexity bounds. Products of the successive
multidegree $(1,m)$ nilcharacters retain that bound on their fixed
product model.

At the last step the residual family has support at most $(m,0)$.
Its extra degree--rank axiom says $G_{(m,0)}=G_{(m,1)}$.
But $(m,1)>(m,0)$, so the latter group is trivial. All groups of
larger degree are also trivial. Thus the ordinary filtration
$(G_{(i,0)})_i$ has degree at most $m-1$. The residual nilcharacter,
viewed merely as a bounded nilsequence, can now be multiplied into
the already present ordinary degree-at-most-$m-1$ correlation test.
The product degree is still at most $m-1$.

We obtain one internal unit multidegree $(1,m)=(1,s-1)$ nilcharacter
$A_{\mathrm{final}}(h,n)$ such that $\Delta_h f$ has
degree-at-most-$s-2$ correlation with it on a dense set of shifts.
This is the exact Theorem 7.1 endpoint needed by the separately
authored symmetry proof. It is conditional on the one-step premise
being supplied at every input rank and every constructed family.
Acceptance of the finite iteration alone does not establish those
analytic premises or the full inverse theorem.

\subsection{Source findings and bounded related-use scope}

The actual final printed p.~1262 was inspected visually. In Example 6.11
the multidegree-to-degree group is printed with $|a|\leq i$, which
includes $G_0$ at every $i$ and does not give bounded degree. For
example $G_0=G_1=\GIKR$, $G_2=0$ would give the group $\GIKR$ at every
index. The required rule is $|a|\geq i$, proved in Section 2 above.

The same example's degree-to-rank rule prints the equality of summed
weights to the fixed bound $d$, rather than to the running index $d'$.
The desired identity $H_{(d',0)}=G_{d'}$ then need not hold. For
instance the linear filtration $G_0=G_1=\GIKR$, $G_2=0$ with allowed
degree bound $d=2$ yields zero at $(1,0)$ under that printed rule,
whereas $G_1=\GIKR$. A nonconstant linear polynomial is then lost.
Section 1 supplies the current-index rule and proves every filtration,
lattice and polynomial assertion used here. The printed length
condition is $\ell\geq\max(r',1)$; that sign is not itself an error.

The bounded related-use scope is Example 6.11, the polynomial identity
(6.8), Lemma 6.21's degree conversions, the initialization and rank-zero
termination of Theorems 7.1--7.2, and the diagonal/total-degree conversions
used in E.9 and the Section 13 assembly. Separate authors of the latter
consumers have been notified. No claim is made here to have read every
unrelated use of a filtration in the source or every historical correction.
These formula errors are not evidence that the intended inverse theorem
is false. Earlier source bytes and decisions are preserved.


\section{Scaling symbols and symmetric multilinearisation with the factorial retained}\label{v15:gi:GTZ-LP05}
\newcommand{\GILN}{\mathbb N}
\newcommand{\GILZ}{\mathbb Z}
\newcommand{\GILR}{\mathbb R}
\newcommand{\GILQ}{\mathbb Q}
\newcommand{\GILC}{\mathbb C}
\newcommand{\GILg}{\mathfrak g}
\newcommand{\GILSym}{\operatorname{Sym}}
\newcommand{\GILpoly}{\operatorname{poly}}

This section proves the multilinearization and scaling assertion E.8(v), including the required divisibility. Its tautification input is proved in Section~\ref{v15:gi:GTZ-LP-TAUTIFICATION}. The unpublished erratum motivates the factorial correction but supplies no theorem premise.

\subsection{Conventions and preliminary constructions}\label{v15:gi:11:sec:setup}
We use the rational, connected, simply connected nilpotent groups,
polynomial maps, finite-dimensional unit-sphere-valued nilcharacters and
symbol equivalence of \cite[Sections 3--6 and Appendix E]{v15-gi-11-GTZ}.
For a vector-valued function $u$, $u^{\otimes r}$ means its ordinary tensor
power if $r\geq0$, with $u^{\otimes0}=1$, and
$\overline u^{\otimes(-r)}$ if $r<0$.
Conjugation is always written explicitly. Equivalence at multidegree $d$
means that every component of $u\otimes\overline v$ is a nilsequence with
support in $\{a:a\leq d,\ a\ne d\}$; all inequalities on multiindices
are coordinatewise. In the degree case this is degree at most $d-1$.
We use the group law on symbols proved in published E.7 and E.8(i)--(iii).
Their proofs use tensor products, conjugation, contraction of a unit vector
with its conjugate, and a rational top-group quotient. With tautification,
the top group is central, so those proofs apply to the present objects.

Initially all data are ordinary, and all integers $d,k,q$ are fixed standard
integers. The last section states exactly how the constructions pass to
limits and parameter families. A multidegree nilcharacter on a limit subset
comes with a representing global polynomial orbit and test function; we
use that representation to extend it, and only then restrict back.

Apply Section~\ref{v15:gi:GTZ-LP-TAUTIFICATION} to the given representation.
We can work on a taut group and with an exactly equal represented sequence
whose polynomial map $g$ satisfies $g(0)=1$. Positive groups and the
vertical character are unchanged. For a degree filtration $G_0=G_1$;
for multidegree $G_0$ is generated by the $G_{e_i}$. In particular $G_d$
is central when the support is contained below one positive $d$.
Write $\GILg_a=\log G_a$ with its rational structure induced by $\Gamma$.

\begin{lemma}[Monomial coordinates]\label{v15:gi:11:lem:monomial}
For multidegree support $a\leq d$, the normalised map has an expression
\[
 g(n)=\prod_{0<a\leq d}\exp(X_a n^a),\qquad X_a\in\GILg_a,
 \quad n^a=\prod_i n_i^{a_i}.                 \tag{1}
\]
Fix any order refining total weight. For a degree-$D$ map on $\GILZ^k$,
replace $a\leq d$ by $0<|a|\leq D$ and require $X_a\in\GILg_{|a|}$.
Every such ordered product is polynomial for the indicated filtration.
\end{lemma}
\begin{proof}
For a degree-$D$ input, use the multidegree target $\GILg_a=\GILg_{|a|}$:
coordinate derivatives are particular degree derivatives, so the input
is polynomial for this target. Published Lemma B.9 gives ordered
binomial monomials with coefficient
at index $a$ in $G_a$. Expand the finite BCH formula. A product of scalar
binomial polynomials with labels $a_1,\ldots,a_r$ has monomials $n^b$
only for $b\leq a_1+\cdots+a_r$. Its Lie-bracket coefficient belongs to
$\GILg_{a_1+\cdots+a_r}\subseteq\GILg_b$. Thus
$\log g(n)=\sum_{0<b\leq d}Y_b n^b$, $Y_b\in\GILg_b$; terms outside the
support vanish and the constant is zero. To obtain (1), process increasing
total monomial weight: choose $X_b$ to match the coefficient at $n^b$ in
the BCH logarithm of the partial ordered product. All commutators of
previous factors have larger total monomial weight, and their coefficients
at $b$ belong to $\GILg_b$. Consequently the adjustment is in $\GILg_b$ and
does not change any previously matched weight. The finite support makes
this algorithm terminate. The argument with $\GILg_{|b|}$ is identical.
For the converse each factor is a polynomial map: in its one-parameter
abelian subgroup an iterated difference of $n^a$ vanishes if its coordinate
order exceeds $a$, and otherwise the coefficient stays in
$\GILg_a\subseteq\GILg_b$ at derivative index $b$. Apply the product theorem
\cite[Corollary B.4]{v15-gi-11-GTZ}. This use does not invoke E.8, E.9 or GI.
\end{proof}

We will use rational connected subgroups and rational central quotients.
Their generated subgroups and central quotients are the geometric
construction proved in Section~\ref{v15:gi:GTZ-LP13-MULTILINEARITY}. We spell out the additional
rational-image check used below. For each rational $G_a$, choose finitely
many elements $\gamma_{a,j}$ of $\Gamma\cap G_a$ whose logarithms span
$\GILg_a$. Such a choice exists: the real span of the logarithms of a
nilpotent subgroup is a Lie algebra. Indeed conjugation by integer powers
of any of its elements gives a polynomial in that integer, and finite
differences extract its iterated brackets; BCH then gives the assertion.
Its connected subgroup contains the lattice, so cocompactness forces it
to be $G_a$, by the homogeneous-space argument of Section~\ref{v15:gi:GTZ-LP13-MULTILINEARITY}.
The same polynomial/BCH argument gives the rational Lie structure
$\operatorname{span}_{\GILQ}\log\Gamma$; all filtration spaces are rational
for this structure.

For the scaling image below, the vectors
$(q^{|a|}\log\gamma_{a,j},\log\gamma_{a,j})$ exponentiate to the actual
lattice elements $(\gamma_{a,j}^{q^{|a|}},\gamma_{a,j})$. For the
multilinear image, multiply each $\log\gamma_{a,j}$ by an integer $N$
so that its rational second component in (8) lies in the lattice $L$
constructed next. Its first component then exponentiates to
$\gamma_{a,j}^N$. These finitely many lattice elements have logarithms
spanning each image Lie algebra $\mathfrak k_b$. The finite-generator
fact in \cite[p.~1243]{v15-gi-11-GTZ} makes their generated discrete subgroup
cocompact in $K_b$. Thus $K_b$ is rational for the actual product lattice,
not merely for an abstract rational form. This also checks $b=0$.

The following explicit lattice choice records
the extra finite-symmetry requirement needed below. In a rational
nilpotent Lie algebra with a finite group $P$ of rational automorphisms,
choose a full rational lattice $L_0$ and replace it by
$\sum_{\sigma\in P}\sigma L_0$. Clear rational bracket denominators and
the finitely many BCH denominators by replacing $L_0$ by $L=M L_0$.
Then BCH$(L,L)\subseteq L$ and $-L=L$. The group $\exp L$ is discrete.
It is cocompact: an adapted rational central-series basis, after integer
scaling, lies in $L$; successive rounding of its coordinates gives a
compact set of coset representatives. The same construction makes all
specified rational filtration subspaces rational. If a rational linear
functional on a specified rational central subspace is prescribed,
increase $M$ to make it integral on $L$ intersected with that subspace.
This gives a $P$-invariant lattice, not just a rational structure invariant
under $P$. Each choice uses only fixed geometric data.

\subsection{Scaling and divisibility on the actual orbit}\label{v15:gi:11:sec:scale}
\begin{theorem}[Corrected E.8(v)]\label{v15:gi:11:thm:scale}
Let $d\ne0$, and let $\chi(n)=F(g(n)\Gamma)$ have multidegree at most $d$.
For every standard integer $q$,
\[
 [\chi(qn)]=q^{|d|}[\chi(n)] .                \tag{2}
\]
For degree $D>0$ on $\GILZ^k$, the exponent is $D$. The equivalence is on
the whole domain, hence on every indicated restriction. Its geometric
and Lipschitz bounds are independent of the polynomial coefficients.
\end{theorem}
\begin{proof}
Use (1). Form the finite positive Rees Lie algebra
\[
 \mathcal R=\bigoplus_{0<a\leq d}\GILg_a t^a,
 \qquad[Xt^a,Yt^b]=[X,Y]t^{a+b},
\]
where a component outside the downset is zero. This is a Lie algebra:
it is a subalgebra of $\GILg\otimes\GILR[t_1,\ldots,t_k]$ modulo the ideal
spanned by monomials outside the downset. Define the rational Lie map
\[
 A_q(Xt^a)=(q^{|a|}X,X)\in\GILg\oplus\GILg.
\]
Its bracket identity is exact because
$q^{|a|}q^{|b|}=q^{|a+b|}$, including negative and zero $q$.
Let $\mathfrak k_b$ be the image of
$\bigoplus_{a\geq b,\ a\ne0}\GILg_a t^a$, and put $K_b=\exp\mathfrak k_b$.
These are rational, descending, and
$[\mathfrak k_b,\mathfrak k_c]\subseteq\mathfrak k_{b+c}$.
Also $\mathfrak k_0=\sum_i\mathfrak k_{e_i}$, so the filtration is taut.
The ambient group is $K_0$, with lattice $K_0\cap\Gamma^2$.
Every factor of
\[
 (g(qn),g(n))=\prod_a\exp\bigl(n^a A_q(X_a t^a)\bigr)
\]
is polynomial at index $a$, and hence so is their product. The top group
has Lie algebra
$\mathfrak k_d=\{(q^{|d|}X,X):X\in\GILg_d\}$.
For $c=q^{|d|}$, the restricted function
\[
 (x,y)\longmapsto F(x)\otimes\overline{F(y)^{\otimes c}}
                                                        \tag{3}
\]
is invariant under $K_d$: the two vertical phases are $c\eta(X)$ and
$-c\eta(X)$. This calculation uses $F^{\otimes c}$, with the signed
convention above, on the same second coordinate. The rational central
quotient $K_0/K_d$ has the quotient filtration with its top group zero.
Thus (3) along the displayed orbit is a lower-degree nilsequence, which
is precisely (2). Restriction and descent preserve a fixed Lipschitz
bound since the geometric maps and compact quotient are fixed.
For degree $D$ use
$\bigoplus_{0<|a|\leq D}\GILg_{|a|}t^a$ and the filtration by $|a|\geq j$.
The top image is the same graph with exponent $D$; all other steps agree.
At $q=0$ the first orbit coordinate is $1$, and (3) is the constant
$F(1)$; no $0^0$ or division is used.
\end{proof}

\begin{corollary}[Divisibility, with the representative specified]\label{v15:gi:11:cor:divide}
For positive degree and every standard $Q\geq1$, there is a nilcharacter
$\psi$ on the same domain with $[\chi]=Q[\psi]$.
\end{corollary}
\begin{proof}
In (1) replace each $X_a$ by $Q^{-|a|}X_a$ and call the resulting map
$g_Q$. It is polynomial on the same group and lattice. With the same
$F$ put $\xi(n)=F(g_Q(n)\Gamma)$. Then $g_Q(Qn)=g(n)$ exactly, factor
by factor in the same order. Theorem~\ref{v15:gi:11:thm:scale} gives
$[\chi]=Q^D[\xi]$, where $D=|d|$ or the degree bound, and
$\psi=\xi^{\otimes Q^{D-1}}$ gives the asserted equation.
The test function still has unit norm; its frequency and finite target
dimension change only by the displayed standard tensor power. In
particular this is divisibility of the actual symbol, not division of
an arbitrarily selected multilinear representative after construction.
\end{proof}

\subsection{Raw multilinearisation and the factorial graph}\label{v15:gi:11:sec:raw}
Let $d=(d_1,\ldots,d_k)\ne0$, $D=|d|$, and divide $\{1,\ldots,D\}$
into ordered blocks $B_i$ of size $d_i$. For a subset $J$ put
$w(J)=(|J\cap B_i|)_i$ and $a!=\prod_i a_i!$; set $C=d!$.
Construct
\[
 \widetilde\GILg=\bigoplus_{\varnothing\ne J\subseteq[D]}\GILg_{w(J)}\epsilon_J,
 \qquad
 [X\epsilon_J,Y\epsilon_L]=
 \begin{cases}[X,Y]\epsilon_{J\cup L},&J\cap L=\varnothing,\\
 0,&J\cap L\ne\varnothing.
 \end{cases}                                             \tag{4}
\]
This is a subalgebra of $\GILg\otimes
\GILR[\epsilon_1,\ldots,\epsilon_D]/(\epsilon_1^2,\ldots,\epsilon_D^2)$;
Jacobi therefore holds. There is no empty component: normalisation
$g(0)=1$ is what permits its removal. Put
\[
 \widetilde\GILg_b=
 \bigoplus_{1_J\geq b}\GILg_{w(J)}\epsilon_J.                 \tag{5}
\]
These are descending rational Lie subalgebras, their bracket has the
required index sum, and they vanish when any $b_i>1$. The zero group
is generated by the first-coordinate groups since each nonempty $J$
contains a coordinate. Thus (5) is taut. Its top group is central.
The block permutation group $P=\prod_i S_{d_i}$ preserves (4), permutes
the groups in (5) by the same permutation of their indices, and
acts trivially on the top copy $\GILg_d\epsilon_{[D]}$.
Choose a $P$-invariant lattice $\widetilde\Gamma$ as in Section
\ref{v15:gi:11:sec:setup}, making
\[
 \widetilde\eta(\exp(X\epsilon_{[D]}))=\eta(\exp X)
                                                        \tag{6}
\]
integral on its top lattice. This is possible since $\eta$ is integral
on $\Gamma\cap G_d$ and hence rational on the rational top Lie algebra.

Choose a Lipschitz unit-sphere-valued function $\widetilde F$ with frequency
$\widetilde\eta$. Here is the needed existence construction, with no
invariance under $P$ assumed. The central torus acts freely on
$\widetilde G/\widetilde\Gamma$: its stabilizer is precisely the top
lattice. On finitely many local trivialisations of this torus bundle,
take the character (6) times a smooth base bump. Choose the bumps so
that their supports cover the compact base and they vanish near the
chart boundaries. The vector of these functions has no common zero;
divide by its norm. This gives a globally Lipschitz unit vector of the
required frequency, with a fixed finite dimension. It is the
partition-of-unity construction underlying published (6.3).

For the coefficients of (1), define
\[
 \widetilde g(z)=\prod_{0<a\leq d}
 \exp\left(a!\sum_{w(J)=a}z_J X_a\epsilon_J\right),
 \qquad z_J=\prod_{j\in J}z_j,\qquad
 \widetilde\chi(z)=\widetilde F(\widetilde g(z)\widetilde\Gamma).
                                                        \tag{7}
\]
Each factor is polynomial for (5). To see this without commuting distinct
coefficients, fix $a$; the summands in its exponent commute because their
Lie coefficients are the same $X_a$, so (7)'s factor is a product of the
individual monomial factors. Then use product closure. Thus
$\widetilde\chi$ has multidegree at most $(1,\ldots,1)$.
Let $\Delta n$ repeat $n_i$ throughout $B_i$.

\begin{lemma}[Raw diagonal identity]\label{v15:gi:11:lem:raw}
The diagonal $\widetilde\chi(\Delta n)$ is equivalent at multidegree $d$
to $\chi(n)^{\otimes C}$. In particular for $k=1$ the multiplier is $d!$.
\end{lemma}
\begin{proof}
On the positive Rees algebra $\mathcal R$ of Section~\ref{v15:gi:11:sec:scale}, set
\[
 T(Xt^a)=\left(X,a!\sum_{w(J)=a}X\epsilon_J\right)
                       \in\GILg\oplus\widetilde\GILg.           \tag{8}
\]
This is a rational Lie map. Indeed for each $U$ with $w(U)=a+b$ the
number of ordered disjoint decompositions $U=J\sqcup L$ with
$w(J)=a,w(L)=b$ is $\prod_i\binom{a_i+b_i}{a_i}$.
Multiplying by $a!b!$ gives $(a+b)!$. Thus
$[T(Xt^a),T(Yt^b)]=T([X,Y]t^{a+b})$ exactly. If $a+b$ is outside
the support, both sides vanish by $[\GILg_a,\GILg_b]\subseteq\GILg_{a+b}=0$.

Define $\mathfrak k_b=T(\bigoplus_{a\geq b,a\ne0}\GILg_a t^a)$,
$K_b=\exp\mathfrak k_b$. As before these give a descending, rational,
taut filtration on $K_0\subseteq G\times\widetilde G$, with lattice
$\Lambda=K_0\cap(\Gamma\times\widetilde\Gamma)$. They give the actual
comparison orbit
\[
 (g(n),\widetilde g(\Delta n))
       =\prod_a\exp\left(n^aT(X_at^a)\right).             \tag{9}
\]
No asserted equality of top weights is being used to generate a
non-descending filtration: every defining sum includes all $a\geq b$.
The top is exactly
\[
 K_d=\{(\exp X,\exp(CX\epsilon_{[D]})):X\in\GILg_d\}.       \tag{10}
\]
It is rational and central. The function on the comparison nilmanifold is
\[
 \mathcal F(x,y)=F(x)^{\otimes C}\otimes\overline{\widetilde F(y)}.
                                                        \tag{11}
\]
Its two vertical phases on (10) are $C\eta(\exp X)$ and
$-\widetilde\eta(\exp(CX\epsilon_{[D]}))=-C\eta(\exp X)$.
Therefore (11) descends to a Lipschitz function $F'$ on
$K_0/K_d$ with the projected lattice $\Lambda'$ and quotient filtration.
Writing $k_a'$ for the projected $\exp T(X_a t^a)$ gives the corrected
identity
\[
 \chi(n)^{\otimes C}\otimes\overline{\widetilde\chi(\Delta n)}
       =F'\left(\prod_a(k_a')^{n^a}\Lambda'\right).        \tag{12}
\]
The right side has top group zero, so it has support strictly below $d$.
For $k=1$, (11) contains $F^{\otimes d!}$ and (12) contains
$\chi^{\otimes d!}$. Dropping either factor would leave the uncancelled
vertical phase $(1-d!)\eta$ and would not prove the claimed equivalence.

The diagonal itself is a degree-$d$ nilcharacter: coarsen (5) by taking
the direct sum of components with $w(J)\geq a$ at index $a$. This is
descending, rational and taut, its top is still the $J=[D]$ component,
and (7) on the diagonal is polynomial for it. This verifies the exact
symbol domain used in (12), without assuming that arbitrary linear
maps preserve arbitrary multidegree filtrations.
\end{proof}

\subsection{Literal symmetry and division before construction}\label{v15:gi:11:sec:symmetry}
The map $\widetilde g$ is equivariant under block permutations:
$\widetilde g(\sigma z)=\sigma\widetilde g(z)$ after fixing the compatible
action convention. An arbitrary $\widetilde F$ need not be invariant,
so equivariance alone does not prove that $\widetilde\chi$ is symmetric.
We supply the missing transition explicitly.

Let $V$ be the finite-dimensional Hermitian target of $\widetilde F$ and
$m=|P|=C$. Regard $\GILSym^m V$ as the symmetric subspace of $V^{\otimes m}$
with its induced norm, and write $\odot$ for symmetrised tensor product.
On the compact nilmanifold put
\[
 B(x)=\mathop{\bigodot}_{\sigma\in P}\widetilde F(\sigma x),
 \qquad F_{\rm sym}(x)=B(x)/\|B(x)\|.                    \tag{13}
\]
The product of nonzero vectors in the symmetric algebra is nonzero:
after a choice of basis this algebra is a polynomial ring over $\GILC$,
an integral domain. Hence $\|B(x)\|>0$ everywhere. Continuity and
compactness give a positive minimum, so (13) is Lipschitz. Reindexing
the finite product proves exact $P$-invariance. Since $P$ fixes the top
group, (13) has vertical frequency $m\widetilde\eta$.
Consequently
\[
 \chi_{\rm sym}(z)=F_{\rm sym}(\widetilde g(z)\widetilde\Gamma)
                                                        \tag{14}
\]
is a literally block-symmetric multidegree-$(1,\ldots,1)$ nilcharacter.
On a block diagonal, $\sigma\widetilde g(\Delta n)=\widetilde g(\Delta n)$.
All factors in (13) then equal the same unit vector $v$, and
$v^{\odot m}=v^{\otimes m}$ has norm one. Thus the diagonal of (14)
is exactly $\widetilde\chi(\Delta n)^{\otimes m}$ in the symmetric
subspace, not merely a proportional vector with an uncontrolled scalar.
Lemma~\ref{v15:gi:11:lem:raw} therefore gives
\[
 [\chi_{\rm sym}(\Delta n)]=Cm[\chi(n)].                 \tag{15}
\]

\begin{theorem}[E.9 and E.10 with a symmetric output]\label{v15:gi:11:thm:multi}
For a multidegree-$d$ nilcharacter $\chi$ there is a literally
block-symmetric multidegree-$(1,\ldots,1)$ nilcharacter $M(\chi)$
whose block diagonal is equivalent to $\chi$. The one-variable case
is the asserted symmetric E.9 output. The assertion holds on any
given limit subset after restriction.
\end{theorem}
\begin{proof}
If $d\ne0$, first use Corollary~\ref{v15:gi:11:cor:divide} with $Q=Cm=C^2$
to choose the specified representative $\rho$ with $[\chi]=Cm[\rho]$.
Apply (4)--(14) to this $\rho$, including its actual polynomial
coefficients, lattice, function and vertical frequency. Equation (15)
now identifies its output diagonal with $[\chi]$. We do not divide
the already symmetrised function and assume that its symmetry survives.
If $C=1$, all constructions and division remain valid. Coordinates
with $d_i=0$ have empty blocks; the normalised polynomial is independent
of them. If all $d_i=0$, the original sequence is a constant unit vector
and that vector on the zero-variable domain gives the assertion exactly.
\end{proof}

\subsection{The exact input to Section 13}\label{v15:gi:11:sec:join}
Let $s\geq2$ be fixed and let $\chi(h,n)$ be the multidegree-$(1,s-1)$
nilcharacter occurring immediately before published (13.4).
First choose $\rho$ with $[\chi]=s[\rho]$ by
Corollary~\ref{v15:gi:11:cor:divide}, in that multidegree symbol group.
Then set $\widetilde\chi=M(\rho)$ from Theorem~\ref{v15:gi:11:thm:multi}.
It has $s$ variables, multidegree $(1,\ldots,1)$ and literal symmetry
in its last $s-1$ variables. Its actual diagonal gives
\[
 [\chi(h,n)]_{\mathrm{Symb}_s({}^*\GILZ^2)}
   =s[\widetilde\chi(h,n,\ldots,n)]_{\mathrm{Symb}_s({}^*\GILZ^2)}.
                                                        \tag{16}
\]
Here is the comparison of symbol groups in (16). For support $a\leq(1,s-1)$,
give the same group the degree filtration generated by the $G_a$ with
$|a|\geq j$. Commutators satisfy the index law, rationality follows
from the rational generated-subgroup lemma, and the orbit is polynomial
by its monomial expansion. At total weight $s$ the only possible index
is $(1,s-1)$, so deleting that top index leaves degree at most $s-1$.
Thus both multidegree equivalences used above become degree-$s$
equivalences, with the same functions and the same input $(h,n)$.
This proves precisely (13.4), including its multiplier $s$.

The sequence $\Theta(n)=\widetilde\chi(n,\ldots,n)$ is a degree-at-most-$s$
nilcharacter by the same coarsening with the binary cube. Fixed linear
pullbacks $n\mapsto n+h$ preserve that degree filtration and preserve
lower-degree equivalence; one may either use the monomial expansion
or published Corollary B.5 and E.8(vi) in its degree-only form.
No full symmetry in all $s$ variables is asserted here. Section 13's
later symmetrisation argument is exactly what must address the first
variable. Its uses of multilinearity require the separately proved
replacement for Lemma 13.2; that result is not inferred from (16).
Likewise the subsequent bias and density arguments remain separate.

\subsection{Uniformity and limit interpretation}\label{v15:gi:11:sec:limits}
All nilpotent groups, rational subspaces, finite tensor dimensions,
integer powers and permutation groups above are fixed by the original
geometric data, $d,k$, and the specified standard divisibility/scaling
integers. They do not depend on orbit coefficients or interval length.
The bounded-basepoint normalisation supplied by Section~\ref{v15:gi:GTZ-LP-TAUTIFICATION} has a uniform
Lipschitz constant. Subgroup restrictions, rational quotients and metrics
are fixed compact geometric maps. The lattice and the test function in
(6)--(7) are selected once from those fixed data. The positive minimum
in (13) is for that fixed test function on that fixed compact space,
not an infimum over a varying family of independently chosen functions.
For a varying original test function with bounded Lipschitz norm, the
raw comparison (11) varies, but the auxiliary $\widetilde F$ and (13)
can still be fixed because their frequency and group are fixed.

Every coefficient change is a finite BCH expression or a rational scalar
multiple in a fixed Lie algebra. Apply it coordinatewise to a representing
sequence to obtain a limit map. The finite formulas, exact identities,
positive fixed lower bound in (13), and uniform Lipschitz estimates pass
to the natural ultrapowers used in published Appendix A. Where geometric
data range over a standard finite list, perform this construction for
each member of that list. No selection depending on an unbounded interval
length or uncontrolled pointwise choice is introduced.
All vertical identities here have that natural-ultrapower meaning: they
hold for every limit point and every limit top-group element, by transfer
of the exact ordinary identity. An identity tested only on standard
points would not suffice; no such restriction is used here.

The preceding constructions supply the applications E.8(v),(vii), E.9, E.10 and (13.4) through the definition of $\Theta$, including literal symmetry in the last $s-1$ variables. The scope is these applications, not a classification of all corrections to GTZ.


\section{The plain additive-quadruple implication on a finite interval}\label{v15:gi:GI-PLAIN-QUADRUPLES}
\newcommand{\GIMZ}{\mathbb Z}
\newcommand{\GIMC}{\mathbb C}
\newcommand{\GIME}{\mathbb E}


\noindent\textbf{Scope.} This local proof supplies the exact plain
correlation statement of GTZ Proposition 8.4, p.~1277, needed by the
Section 13 join. It does not use Proposition 8.3 or infer plain bias
from lower-degree bias. The argument is the elementary two-step
Cauchy--Schwarz pattern; the finite domains and constants are written
explicitly. No source error in Proposition 8.4 is asserted.

\subsection{Precise finite assertion}
Put $J=\{1,\ldots,N\}$, $I=\{-N,\ldots,N\}$, and suppose
$H\subseteq I$ has $|H|\geq\alpha |I|$ for fixed $\alpha>0$.
Let $f:J\to\GIMC$ be $1$-bounded and extend it by zero to $\GIMZ$.
For each $h\in H$ let $\chi_h:\GIMZ\to\GIMC$ be globally bounded
by a fixed $B>0$, and assume
\[
 \left|\GIME_{n\in J} f(n+h)\overline{f(n)}\chi_h(n)\right|
 \geq\delta>0.
\]
The values outside $J$ are part of the data and are not truncated.

\begin{theorem}\label{v15:gi:12:thm:plain}
Write $a=\alpha\delta/3$ and $\eta=a^4/4$.
There are at least
\[
 \frac{\eta}{2B^4}(4N+1)(2N+1)^2
\]
distinct additive quadruples $(h_1,h_2,h_3,h_4)\in H^4$,
$h_1+h_2=h_3+h_4$, for which
\[
 \left|\GIME_{n\in J}
  \chi_{h_1}(n)\chi_{h_2}(n+h_1-h_4)
  \overline{\chi_{h_3}(n)}
  \overline{\chi_{h_4}(n+h_1-h_4)}\right|
 \geq\eta/2.
\]
In particular their number is a positive fixed multiple of $N^3$.
\end{theorem}

\subsection{Two exact Cauchy--Schwarz steps}
For each $h\in H$ choose a scalar $c_h$ of modulus one so the
given average multiplied by $c_h$ is real and at least $\delta$.
Set $u_h=c_h\chi_h$ on $H$, and $u_h=0$ for every $h\notin H$.
The functions remain defined on all integers. Then
\[
 \GIME_{h\in I}\GIME_{n\in J}
 f(n+h)\overline{f(n)}u_h(n)\geq\alpha\delta.
\]
Let $M=J+I=\{1-N,\ldots,2N\}$, so $|M|=3N$.
Set $m=n+h$. The preceding average is exactly
\[
 3\GIME_{m\in M}f(m)
       \GIME_{h\in I}\overline{f(m-h)}u_h(m-h).
\]
Terms with $m-h\notin J$ vanish because $f$ was zero-extended.
Cauchy--Schwarz in $m$, using $|f|\leq1$, therefore yields
\begin{equation}\label{v15:gi:12:eq:first}
 \GIME_{m\in M}\left|\GIME_{h\in I}
          \overline{f(m-h)}u_h(m-h)\right|^2\geq a^2.
\end{equation}

Expand the square and put $h'=h+k$, $n=m-h$.
Let $K=\{-2N,\ldots,2N\}$. The expanded left side equals
\[
 \frac{N}{|M|}\frac{|K|}{|I|}
 \GIME_{k\in K}\GIME_{n\in J}
 \overline{f(n)}f(n-k)T(k,n),
 \qquad
 T(k,n)=\GIME_{h\in I}u_h(n)\overline{u_{h+k}(n-k)}.
\]
This is an equality with the following domain checks. The factor
$\overline{f(n)}$ forces $n\in J$, and then $m=n+h\in M$
automatically for $h\in I$. The condition $h+k\in I$ is supplied
by $u_{h+k}=0$ outside $H\subseteq I$. Thus no additional cutoff
in $n-k$ is inserted into $T$. Its values there are legitimate because
the original functions are globally defined. The remaining factor
$f(n-k)$ stays in the bounded weight.

The prefactor is positive and at most two. By \eqref{v15:gi:12:eq:first},
the absolute value of the displayed average without prefactor is
at least $a^2/2$. Cauchy--Schwarz in $(k,n)$, dropping only the
bounded weight $\overline{f(n)}f(n-k)$ by an inequality, gives
\[
 \GIME_{k\in K}\GIME_{n\in J}|T(k,n)|^2\geq a^4/4=\eta.
\]
Expand this square with independent $h,h'\in I$. It is the average
over $(k,h,h')\in K\times I^2$ of
\[
 C(k,h,h')=\GIME_{n\in J}
 u_h(n)u_{h'+k}(n-k)
       \overline{u_{h'}(n)}\overline{u_{h+k}(n-k)}.
\]
The average is real and at least $\eta$, while $|C|\leq B^4$.
Consequently at least a fraction $\eta/(2B^4)$ of triples have
$|C|\geq\eta/2$. Indeed the complementary contribution to the
average of $|C|$ is at most $\eta/2$, and the contribution of the
selected fraction is at most that fraction times $B^4$.

Such a triple has $h,h',h+k,h'+k\in H$, since otherwise $C=0$.
Set
\[
 (h_1,h_2,h_3,h_4)=(h,h'+k,h',h+k).
\]
It is an additive quadruple and $h_1-h_4=-k$. The map is injective,
with inverse $h=h_1$, $h'=h_3$, $k=h_4-h_1$.
The factors $c_h$ contribute only a constant phase of modulus one
to $C$, so they do not change its absolute value. This proves the
theorem, including the exact unweighted average on $J$.

\subsection{Vector families, internal limits and actual consumer}
For vector-valued $\chi_h$ of fixed finite dimension $D$, a positive
Euclidean correlation implies that some component has magnitude at
least $\delta/\sqrt D$. Finite pigeonhole retains a component on
a subfamily of density at least $\alpha/D$. Apply the scalar theorem
there. The scalar output is a component of the full tensor product,
so it implies plain vector bias. A finite list of dimensions can first
be padded into one fixed dimension. For a bounded $f$ with a fixed
bound other than one, divide it by that bound, changing $\delta$
by a fixed factor only.

Apply the finite assertion to representatives of an internal family,
with uniform standard bounds, density and correlation. Choices of
components and phases can be made for each finite family and then
passed to the ultralimit. All constants above are positive standard
numbers, so the selected quadruples have positive standard density
on the $N^3$ scale and the correlation is positive by a standard
amount. This is exactly the plain Proposition 8.4 output.

In Section 13 the input to this theorem is the entire bounded product
that includes the selected lower nilcharacter. Keeping it inside
$\chi_h$ throughout the calculation preserves that lower information.
The theorem requires no degree bound and introduces no lower-degree
test factor. It therefore supplies the stronger plain statement that
cannot be replaced merely by the degree-biased LP03 output.

The bounded application audit checks the original $J$, the enlarged
$M=J+I$, the parameter range $K$, the exact inverse reindexing and
the global domain of $\chi_h$. A possible artificial cutoff
$1_J(n-k)$ is unnecessary: $f(n-k)$ is removed by the second
Cauchy--Schwarz inequality, not by claiming an equality after erasing
its support. This resolves the domain question without classifying
the published Proposition 8.4 as false or defective. The modified
Proposition 7.3, its nilsequence-specific factorization and the
remaining Section 13 hypotheses are outside this theorem's scope.


\section{The Section 13 correlation and symmetry join}\label{v15:gi:GI-SECTION13}
\newcommand{\GINZ}{\mathbb Z}
\newcommand{\GINN}{\mathbb N}
\newcommand{\GINC}{\mathbb C}
\newcommand{\GINE}{\mathbb E}
\newcommand{\GINSymb}{\operatorname{Symb}}
\newcommand{\GINNil}{\operatorname{Nil}}
\newcommand{\GINone}{\mathbf1}

This section proves the symmetrization conclusion from the internal multilinearization theorem and its stated published inputs. The signed identities retain the actual unconjugated $\alpha$ in published (13.10).

\subsection{Target and exact inputs}
Fix a standard integer $s\geq3$ and an unbounded limit integer $N$.
Write $I_t=[-tN,tN]\cap{}^*\GINZ$, $J=[N]$, $r=s-1$ and $m=s-2$.
All geometric data and finite dimensions are standard. Vertical identities
hold for every limit point and every limit top-group element under the
natural ultrapower extension, as in the project tautification and LP13
units. Tensor products have their Hilbert tensor norms. Every nilcharacter
is unit-vector-valued; a scalar coordinate need not have modulus one.

\begin{theorem}[Theorem 7.4 on the consumed domain]\label{v15:gi:13:thm:target}
Let $f$ be bounded on $J$, extended by zero outside $J$, and let
$H\subseteq I_1$ be a dense limit subset. Suppose that $\chi(h,n)$ is
a globally represented multidegree-$(1,s-1)$ nilcharacter and
$\Delta_h f$ has degree-at-most-$(s-2)$ correlation with $\chi(h,\cdot)$
on $J$ for every $h\in H$, with GTZ's correlation convention.
Then there are a globally represented degree-at-most-$s$ unit nilcharacter
$\Theta$ and a bounded globally represented nilsequence $\Psi(h,n)$
supported on
\[
 \mathcal J=\{(i,j):i+j\leq s-1\}\cup\{(i,s-i):2\leq i\leq s\}
                                                        \tag{1}
\]
such that, on $I_1\times J$, each component of $\chi(h,n)$ is a bounded
finite linear combination of components of
\[
 \Theta(n+h)\otimes\overline{\Theta(n)}\otimes\Psi(h,n).  \tag{2}
\]
All families and bounds are standard in the limit sense, independently of
the interval length. No effective numerical constants are asserted.
\end{theorem}

We use the following exact premises.
The project LP05-001 supplies a unit nilcharacter $\widetilde\chi$ of
cube multidegree $(1,\ldots,1)$, literally symmetric in its last $r$
variables, with
\[
 [\chi(h,n)]_{\GINSymb_s({}^*\GINZ^2)}
       =s[\widetilde\chi(h,n,\ldots,n)]_{\GINSymb_s({}^*\GINZ^2)}. \tag{3}
\]
The project LP13 theorem supplies multilinearity of the symbol of
$\widetilde\chi$ under fixed integer linear substitutions, in each slot.
The project tautification-003 Theorem 3 supplies the uniform shift tensor
of a positive-degree nilcharacter, lowering degree by one in the ordinary
degree case. It applies to limit shifts and limit families.

The unchanged published inputs are the definitions and elementary symbol
arithmetic E.3, E.7 and E.8(i)--(iii),(vi); nilcharacter approximation and
selection E.5/E.6 and Appendix A; Proposition 8.4 with its plain-correlation
hypothesis; the converse Proposition 5.6; and E.11/E.13. Their exact source
applications are recorded separately. In particular Proposition 8.4 is a
distinct input. The weaker degree-$(s-2)$ biased-quadruple conclusion of
Proposition 7.3 alone would not retain the lower factors needed here.
No corrected assertion is supplied by citing the unpublished erratum.

For an ordered list of linear forms define
\[
 T(L_1;L_2,\ldots,L_s)=[\widetilde\chi(L_1,L_2,\ldots,L_s)]_{\GINSymb_s}.
\]
The semicolon separates the distinguished first slot. We use the notation
$n^j$ in a slot list for $j$ repeated entries, not a power of $n$.
All formulas between such symbols are on the full indicated limit lattice
until a box is explicitly named.

\subsection{Plain correlation and retention of the lower factors}
By E.5/E.6 in the ordinary one-dimensional degree ordering and Appendix A
selection, after restricting to a still dense
internal subset of $H$, choose degree-at-most-$(s-2)$ unit nilcharacters
$\varphi_h$ so that $\Delta_hf$ plainly correlates with
$\chi(h,\cdot)\otimes\varphi_h$. The family is internal, with a standard
finite geometry list, bounded dimensions and common Lipschitz bounds.
Here the selection is into the countable increasing complexity presentation
of the nilcharacter class. An internal image contained in that union is
contained in a standard stage by Appendix A compactness; otherwise nested
nonempty internal sets of parameters outside successive stages give an
image outside the union. Positive correlations can likewise be made
uniform by A.6. Pad targets by zero coordinates to one fixed dimension.
Outside the retained set use a fixed constant unit nilcharacter, and keep
all parameter indicators separately. This preserves the unit norm needed
for the identities below.
Conjugate the selected lower nilcharacter if needed to match the published
correlation convention. A finite pigeonhole also fixes a scalar coordinate
of $\chi(h,\cdot)\otimes\varphi_h$ retaining plain correlation on a dense
set. Thus the scalar version of Proposition 8.4 suffices; the full unit
vectors remain available behind these coordinates.

Apply published Proposition 8.4 to this actual family. Its hypotheses are
plain correlation on $J$, internal dependence, bounded functions and a
dense subset of $I_1$; all hold. Its conclusion gives bias for many additive
quadruples $(h_1,h_2,h_3,h_4)$ in that subset, with tensor signs $+,+,-,-$
and the shift $n+h_1-h_4$ on the second and fourth factors. Parameterise
\[
 (h_1,h_2,h_3,h_4)=(h_0+a,h_0+b,h_0+a+b,h_0).
\]
This is an injective correspondence with the retained quadruples.
$h_0\in I_1$ and $a,b\in I_2$. Pigeonhole in $h_0$ to obtain one fixed
$h_0$ and a dense set of pairs $(a,b)\in I_2^2$.
Define the global unit nilcharacter
\[
 \tau(a,b,n)=\chi(h_0+a,n)\otimes\chi(h_0+b,n+a)
 \otimes\overline{\chi(h_0+a+b,n)}
 \otimes\overline{\chi(h_0,n+a)}.                         \tag{4}
\]
It has total degree at most $s$, not degree three for general $s$.

For any coordinate of a unit nilcharacter $v$, the exact identity
\[
 v_j(n+a)=\sum_k\bigl(v_j(n+a)\overline{v_k(n)}\bigr)v_k(n) \tag{5}
\]
follows from $\sum_k|v_k(n)|^2=1$. The bracket is a coordinate of the
shift tensor. Apply the project shift theorem to $\varphi_{h_0+b}$ and
$\varphi_{h_0}$, with shift $a$. These brackets have degree at most
$s-3$, uniformly in all parameters. Conjugate (5) when the shifted factor
is conjugated. Thus the original bias implies degree-at-most-$(s-3)$ bias
of
\[
 \tau(a,b,n)\otimes\varphi_{h_0+a}(n)\otimes\varphi_{h_0+b}(n)
 \otimes\overline{\varphi_{h_0+a+b}(n)}
 \otimes\overline{\varphi_{h_0}(n)}.                      \tag{6}
\]
Selection gives internal scalar degree-at-most-$(s-3)$ tests with common
bounds and finite geometries. A standard finite choice of coordinates
retains a positive density of the pairs. Phase-align their scalar averages
by a bounded parameter scalar, and set the test equal to zero on other
pairs. These factors are constant in $n$, so they have degree zero.
We obtain fixed scalar coordinates $t$ of $\tau$ and $u,v,w,z$ of the
displayed $\varphi$ factors with
\[
 \left|\GINE_{a,b\in I_2,n\in J}
 t(a,b,n)u_a(n)v_b(n)\overline{w_{a+b}(n)}\overline{z(n)}
 p_{a,b}(n)\right|\gg1.                                 \tag{7}
\]
The scalar $p_{a,b}$ is a uniformly bounded internal lower-degree family.
Finite coordinate selection, rather than division by a coordinate, justifies
this passage even if individual coordinates vanish.

\subsection{The three Cauchy--Schwarz eliminations}
Every average in this section has $n\in J$. All bounds are standard.
Applying Cauchy--Schwarz in $(a,n)$ to remove the bounded factor
$u_a\overline z$ in (7), and expanding the squared inner average, gives
\[
 \left|\GINE_{a,b,b'\in I_2,n\in J}
 t(a,b,n)\overline{t(a,b',n)}v_b(n)\overline{v_{b'}(n)}
 \overline{w_{a+b}(n)}w_{a+b'}(n)
 p_{a,b}(n)\overline{p_{a,b'}(n)}\right|\gg1.             \tag{8}
\]
Put $c=a+b+b'$. Keep $b,b'\in I_2$, take $c\in I_6$, and retain
$\GINone_{c-b-b'\in I_2}$. This exact indicator is independent of $n$.
The normalized average changes by the factor $|I_2|/|I_6|$, comparable
to one. No three-term sum of elements of $I_2$ is asserted to lie in $I_3$.
Pigeonhole in $c$ to fix $c_0\in I_6$. Set
\[
 \alpha(b,b',n)=\tau(c_0-b-b',b,n)
                    \otimes\overline{\tau(c_0-b-b',b',n)},\qquad
 V_b(n)=v_b(n)w_{c_0-b}(n).                              \tag{9}
\]
In particular the second factor of $V_b$ is not conjugated: its sign
follows from the last two $w$ factors in (8). Absorb the parameter cutoff
into the scalar lower test. Equation (8) becomes, for a fixed coordinate
$a_0$ of the actual $\alpha$,
\[
 \left|\GINE_{b,b'\in I_2,n\in J}
 a_0(b,b',n)V_b(n)\overline{V_{b'}(n)}q_{b,b'}(n)\right|\gg1. \tag{10}
\]
The functions $V_b$ need only be bounded; no inverse or unit norm for them
is used. Cauchy--Schwarz in $(b',n)$ removes $\overline{V_{b'}}$ and
produces the product
$a_0(b_1,b',n)\overline{a_0(b_2,b',n)}V_{b_1}\overline{V_{b_2}}$
with the product of the corresponding $q$ tests. Cauchy--Schwarz in
$(b_1,b_2,n)$ removes the remaining $V$ factors. Its expanded four-term
product is
\[
 a_0(b_1,b'_1,n)\overline{a_0(b_2,b'_1,n)}
 \overline{a_0(b_1,b'_2,n)}a_0(b_2,b'_2,n).               \tag{11}
\]
Its scalar test is the product of four $q$ tests with the same signs,
still degree at most $s-3$, internal and uniformly bounded.
Pigeonhole in $(b_2,b'_2)$, calling the fixed values $(b_0,b'_0)\in I_2^2$.
The resulting expression is a coordinate of the global unit nilcharacter
\[
 \alpha'(b,b',n)=\alpha(b,b',n)\otimes\overline{\alpha(b_0,b',n)}
 \otimes\overline{\alpha(b,b'_0,n)}\otimes\alpha(b_0,b'_0,n). \tag{12}
\]
Hence $\alpha'(b,b',\cdot)$ is degree-at-most-$(s-3)$ biased on $J$
for many $(b,b')\in I_2^2$. To see the last assertion quantitatively,
if the final average has modulus at least $\delta>0$ and each inner
average is at most $B$, a proportion at least $\delta/(2B)$ has modulus
at least $\delta/2$. All selections and parameter cutoffs above depend
internally on the parameters. Products and conjugates do not increase the
ordinary nilsequence degree or destroy the fixed finite complexity bounds.

\subsection{Exact symbols and the lower-slot remainder}\label{v15:gi:13:sec:algebra}
Constants such as $h_0,c_0,b_0,b'_0$ are limit integers. Give the original
multidegree group its total-degree filtration generated by all components
whose total weight is \emph{at least} the index, as verified in LP05.
Apply the uniform project ordinary-degree shift theorem on the full slot
domain to a constant translation in the relevant coordinate. Its difference
has degree at most $s-1$, and stays lower-degree under any fixed integer
linear pullback. Thus a constant added to a slot can be removed in the
degree-$s$ symbol. A slot consisting only of a constant becomes zero.
LP13 additivity gives $T(\ldots,0,\ldots)=2T(\ldots,0,\ldots)$, so that
symbol is zero. This argument does not choose an implicit subgroup
basepoint for a fixed-slot orbit. The coarsening uses $|a|\geq j$;
the opposite inequality printed in Example 6.11 is not used.

From (3),(4), linearity in the first slot, and last-block symmetry,
\begin{align}
 [\tau(a,b,n)]
 &=s\bigl(T(b;(n+a)^r)-T(b;n^r)\bigr)\notag\\
 &=s\sum_{j=1}^r {r\choose j}T(b;a^j,n^{r-j}).           \tag{13}
\end{align}
In the first line $(n+a)^r$ denotes repeated slots.
The leading term is $srT(b;a,n^m)$. Let
\[
 R_\tau(a,b,n)=s\sum_{j=2}^r {r\choose j}T(b;a^j,n^{r-j}). \tag{14}
\]
Every summand has at most $s-3$ slots equal to $n$.
Substitute $a=c_0-b-b'$ in (9), remove $c_0$ as just justified, and
expand the $j$ copies of $-b-b'$ by LP13. The leading symbol of $\alpha$ is
\[
 -sr\,T(b-b';b+b',n^m).                                \tag{15}
\]
This has the order of the two distinguished entries dictated by (13).
Interchanging them would require an unproved first-slot symmetry and is
not permitted. The remainder is the difference of the two substitutions
in (14), each fully expanded into entries in $\{b,b',n\}$ and still with
at most $s-3$ copies of $n$.

Apply the mixed difference (12). At the level of degree-$s$ symbols,
constant slots vanish; a monomial with no $b$ or no $b'$ cancels.
In (15) the pure terms $T(b;b,n^m)$ and $T(b';b',n^m)$ therefore disappear.
There is a globally represented unit nilcharacter $\beta$ such that
\[
 [\alpha'(b,b',n)]=-sr\bigl(T(b;b',n^m)-T(b';b,n^m)\bigr)+[\beta(b,b',n)].
                                                               \tag{16}
\]
Here $[\beta]$ is an explicit finite integer linear combination of symbols
$T(v_1;\ldots,v_s)$, with every $v_i\in\{b,b',n\}$ and at most $s-3$
slots equal to $n$. More precisely, expand (14) after the two substitutions
$(a,b)=(-b-b',b)$ and $(-b-b',b')$, subtract the second, then delete the
monomials missing $b$ or missing $b'$. Tensor the corresponding actual
pullbacks of $\widetilde\chi$ with their signed integer multiplicities.
This constructs $\beta$ and proves (16). Thus the needed remainder
assertion is stronger than merely a separate bound on its degree in $n$.

Define the actual unit nilcharacter
\[
 \sigma(b,b',h_1,\ldots,h_m,n)=
 \bigotimes_{\omega\in\{0,1\}^m}\mathcal C^{|\omega|}
   \alpha'\left(b,b',n+\sum_{j=1}^m\omega_jh_j\right),    \tag{17}
\]
where $\mathcal C$ is conjugation. Every shift has its own $h_j$.
It has total degree at most $s$ by linear pullback and tensor arithmetic.
Apply the alternating sum to (16). In a term with fewer than $m$ slots
equal to $n$, expand each such slot into $n$ and the available $h_j$.
If a subset $S$ of the shift labels is used, its coefficient contains
\[
 \sum_{\omega:\,\omega_j=1\ (j\in S)}(-1)^{|\omega|}=0
 \quad\hbox{unless }S=[m].
\]
There are too few slots to use all $m$ labels, so every $\beta$ term
vanishes. For the leading term there are exactly $m!$ surviving bijections
of labels to slots, each with coefficient $(-1)^m$. Last-block symmetry
identifies their values. With the minus sign in (16), the exact result is
\[
 [\sigma]=(-1)^{s-1}s!
 \bigl(T(b;b',h_1,\ldots,h_m)-T(b';b,h_1,\ldots,h_m)\bigr). \tag{18}
\]
The unused coordinate $n$ on the right is still a coordinate of the same
domain. Formula (18) keeps the actual definitions (9),(12),(17). If one
instead conjugates $\alpha$ and $\alpha'$ explicitly, its coefficient
changes sign; silently using the other object is not an identification.

\subsection{From correlation to a zero symbol on the correct box}
By Proposition 5.6, the positive degree-at-most-$(s-3)$ bias obtained
after (12) implies a positive $U^m[J]$ norm for a scalar coordinate of
$\alpha'$, for many $(b,b')$. Appendix A uniformity gives a standard
positive lower bound on a dense internal set. Equivalently the average
of the vector $U^m[J]^{2^m}$ quantities is $\gg1$. Since the dimension is
standard finite, at least one fixed scalar coordinate contributes $\gg1$.
For $s=3$, this is the converse at degree zero and $U^1$, so no extra
one-step test has been inserted.

Expand the interval Gowers norm using its zero-extension definition and
the same normalization as GTZ. One can choose the auxiliary cyclic group
of order a fixed multiple of $N$ greater than $2(m+1)N$. For a nonzero
cube all vertices lie in $J$, so its differences have unique lifts
$h_j\in I_1$, and no lifted sum can wrap around that cyclic group.
The normalizing denominator is a positive standard constant multiple of
the count of such cubes. Thus the expansion gives a coordinate of (17)
with a positive standard average, up to a fixed positive normalizing factor,
on the convex polytope
\[
 \Omega=\{(b,b',h,n): b,b'\in I_2,\quad
       1\leq n+\omega\cdot h\leq N\text{ for all }\omega\in\{0,1\}^m\}.
                                                               \tag{19}
\]
In particular $\Omega\subseteq I_8^{s+1}$, and
$\left|\GINE_{x\in I_8^{s+1}}\GINone_\Omega(x)\sigma(x)\right|\gg1$.
All defining inequalities in (19) are affine linear; their number depends
only on $s$. The selection of biased pairs is used only for a lower bound
on the nonnegative norm average, not as an extra arbitrary cutoff in (19).

Published E.11 asks for dense subprogressions of a positive interval.
Partition $I_8^{s+1}$ into its finitely many sign orthants and coordinate-zero
hyperplanes. Since $\sigma$ is bounded, the latter have normalized mass
$O_s(1/N)$ and contribute $o(1)$. By the triangle inequality one open
orthant has positive standard weighted average. For its signs
$\varepsilon_i\in\{-1,1\}$, reflect $x_i=\varepsilon_i y_i$.
The reflected polytope is convex and the reflected nilcharacter
$\sigma_\varepsilon(y)=\sigma(\varepsilon y)$ has degree at most $s$.
Take every progression in E.11 to be the full positive interval $[8N]$.
The hypothesis is plain bias, hence bias of degree strictly below $s$.
E.11, with dimension $s+1$, degree $s\geq3$ and scale $8N$, gives
$\sigma_\varepsilon\in\GINNil_{<s}(I_8^{s+1})$. Reflecting back gives
\[
 [\sigma]_{\GINSymb_s(I_8^{s+1})}=0.                        \tag{20}
\]
Only one orthant was used to obtain bias; E.11 itself supplies the conclusion
on the full symmetric box. No extrapolation from positive points is implicit.

Pull (18),(20) back along the fixed integer linear map
\[
 (h,n)\longmapsto(b,b',h_1,\ldots,h_m,n_{\rm aux})
                 =(h,n,n,\ldots,n,0).
\]
It maps $I_8^2$ into $I_8^{s+1}$. Lower-degree nilsequences pull back with
the same degree, so
\[
 s!\bigl(T(h;n^r)-T(n;h,n^{r-1})\bigr)=0
                    \quad\hbox{in }\GINSymb_s(I_8^2).       \tag{21}
\]
The sign in (18) can be removed because both sides equal zero. Apply E.13
on this symmetric box, with dimension two, degree $s\geq3$ and the fixed
positive integer $s!$. Its torsion-free conclusion gives
\[
 T(h;n^r)=T(n;h,n^{r-1})\quad\hbox{in }\GINSymb_s(I_8^2).  \tag{22}
\]
We restrict to $I_1\times J$ only after this use. E.13 is not being applied
directly to an asymmetric rectangle.

\subsection{Exact finite-component factorization}
Set $\Theta(n)=\widetilde\chi(n,\ldots,n)$, the same object as in LP05.
It is a unit degree-at-most-$s$ nilcharacter. Expand
$[\Theta(n+h)]-[\Theta(n)]$ by LP13. Terms with exactly one $h$ equal
\[
 T(h;n^r)+rT(n;h,n^{r-1})=sT(h;n^r)=[\chi(h,n)]          \tag{23}
\]
on $I_1\times J$, by (22),(3). Let $R(h,n)$ be the tensor product of
the actual pullbacks of $\widetilde\chi$ for all subsets of the $s$ slots
of cardinality at least two, putting $h$ in those slots and $n$ elsewhere.
There is no permutation of the distinguished first slot in this definition.
Every factor has multidegree at most $(i,s-i)$ for its subset size $i$:
coarsen the cube filtration by counting the coordinates assigned to $h$
and $n$, using all weights at least the index. The polynomial orbit,
rationality and top frequency are as in the LP05 block coarsening.
Their product is therefore supported on the union of these downsets,
contained in $\mathcal J$.

By the symbol equality just proved, the unit-vector expression
\[
 W(h,n)=\chi(h,n)\otimes\overline{\Theta(n+h)}\otimes\Theta(n)\otimes R(h,n)
                                                               \tag{24}
\]
is a degree-at-most-$(s-1)$ nilsequence on $I_1\times J$.
By the definition of the restricted nilsequence class, it has a global
representation whose restriction is (24); no global equality beyond this
rectangle is claimed. Such a degree filtration gives multidegree support
$i+j\leq s-1$ by assigning the group at $(i,j)$ its total-degree group.
Choose that representation and set $\Psi=\overline R\otimes W$.
It has support in (1) by the product filtration and is uniformly bounded.

Finally tensor (24) with
$\Theta(n+h)\otimes\overline{\Theta(n)}\otimes\overline R$.
Contract the matching pairs using their unit norms. Each contraction is
a fixed linear map between fixed finite tensor spaces. The result is exactly
$\chi$, so (2) holds component by component with fixed bounded coefficients.
This does not divide by any scalar coordinate or replace symbol equality
by an unproved scalar equality. It proves Theorem~\ref{v15:gi:13:thm:target}.

\subsection{Bounds and source comparison}
All finite coordinate selections, tensor dimensions, permutation counts
and parameter boxes depend only on fixed degree and geometric data.
Appendix A selection/uniformity supplies the finite family of lower tests
before the averaging manipulations. The three Cauchy--Schwarz steps use
common standard bounds, and each pigeonhole retains a standard positive
amount. Fixed parameter values may be unbounded limit integers; the project
shift theorem supplies uniform constructions for exactly that case.
All algebraic identities are finite and transfer coordinatewise, with the
all-limit vertical convention retained. E.11 and E.13 have the exact degree,
dimension, progression and symmetric-box domains stated above.

The proof retains, rather than erases, the published factorial history.
Theorem 7.4 on both pp.~1271 and 1315 prints the stricter hypothesis
``$<s-2$-correlates'', whereas the preceding Theorem 7.1 supplies
degree-at-most-$(s-2)$ correlation. The target here deliberately assumes
the latter, and Section 2 proves the required selection and shift step
under that weaker hypothesis. The printed Theorem 7.4 is not a premise.
In addition to the repeated shift index and alternating-derivative sign in
(13.14)--(13.17), the first two slots in printed Lemma 13.4 are interchanged
relative to (13.3),(13.10). This changes the sign in (13.13). Formula (18)
is the resulting coefficient for the actual unconjugated objects. The
published simplification in Lemma 13.3 also retains an $a$ where cancellation
requires $b$; its tail has degree at most $s-3$ in $n$. The scalar lower-factor
bookkeeping and the $I_6$ substitution above give their own exact formulas.
These source discrepancies are not counterexamples to the final zero relation.
Their required related uses and current representations must be checked
independently before a complete correction claim.

This packet proves a local composition relative to its stated project and
published premises. It does not reprove E.11/E.13's equidistribution inputs,
Proposition 5.6's published converse sources, Proposition 8.4, the frequency
analysis, or the final finite-family inverse-theorem assembly. Those are
direct dependencies or separate consumers, not silently discharged here.

\section{Same-step linearization using the increment polynomial group}\label{v15:gi:GTZ-LP-INCREMENT-LINEARIZATION}
\newcommand{\GIOR}{\mathbb R}
\newcommand{\GIOZ}{\mathbb Z}
\newcommand{\GION}{\mathbb N}
\newcommand{\GIOpoly}{\operatorname{poly}}
\newcommand{\GIOfg}{\mathfrak g}
\newcommand{\GIOid}{\mathrm{id}}



This section proves the same-step linearization construction. It uses the bounded normalization and tautification proved in Section~\ref{v15:gi:GTZ-LP-TAUTIFICATION}, and the published algebraic polynomial-map results specified below. It does not assume GI($s$).

\subsection{The required conclusion}

Let $G/\Gamma$ be a connected simply connected nilpotent Lie group modulo a
lattice, with a rational degree filtration $G_\bullet$ of degree at most
$s\geq1$. Every $G_i$ is connected and rational. Let a scalar polynomial
nilsequence be $F(g(n)\Gamma)$, with $F$ complex-valued, $1$-bounded and
Lipschitz. The polynomial $g$ is adapted to $G_\bullet$ in the right-difference
convention $\partial_u g(n)=g(n+u)g(n)^{-1}$.

\begin{theorem}\label{v15:gi:14:thm:linear}
For this fixed filtered nilmanifold there is a fixed $s$-step nilmanifold
$\widetilde G/\widetilde\Gamma$ with a fixed smooth metric, and a constant
$C$ depending only on the original fixed geometric data, with the following
property. For every polynomial $g$ and every such $F$, there are
$a\in\widetilde G$ and a $1$-bounded Lipschitz function
$\widetilde F:\widetilde G/\widetilde\Gamma\longrightarrow\mathbb C$ such that
\[
 F(g(n)\Gamma)=\widetilde F(a^n\widetilde\Gamma)
 \quad(n\in\GIOZ),\qquad
 \operatorname{Lip}(\widetilde F)\leq C\operatorname{Lip}(F).
\]
Fixed smooth metric comparisons are included in $C$.
The same construction applies to limit polynomial maps and bounded
Lipschitz limit functions, with fixed nilmanifold and uniform constants.
Consequently a finite-family polynomial inverse theorem of degree $s$
implies the finite-family linear inverse theorem of step $s$, without
changing the input Gowers norm or losing correlation.
\end{theorem}

If $\operatorname{Lip}(F)=0$, the constant test function has a constant
extension and the displayed bound is understood literally. A point
nilmanifold can be used for a constant sequence. Step at most $s$ is the
meaning of $s$-step throughout.

By Section~\ref{v15:gi:GTZ-LP-TAUTIFICATION}, restrict to a fixed taut representation with
$G_0=G_1=G$ and replace the polynomial by $h$, where $h(0)=1$.
The exact identity is
\[
 g(0)=q\gamma,\qquad h(n)=q^{-1}g(n)\gamma^{-1},\qquad
 F_q(x\Gamma)=F(qx\Gamma),\qquad
 F_q(h(n)\Gamma)=F(g(n)\Gamma).
\]
Here $q$ lies in a fixed compact set and
$\operatorname{Lip}(F_q)$ is bounded by a fixed multiple of
$\operatorname{Lip}(F)$. The underlying taut filtered nilmanifold is fixed
independently of $g$. It therefore suffices to prove the theorem for this
normalized $h$ and $F_q$. In what follows all geometric data refer to this
fixed taut representation.

\subsection{Polynomial groups and their Lie algebras}

For $i\geq0$ put
\[
 P_i=\GIOpoly(\GIOZ_\GION\longrightarrow G_{\bullet+i}).
\]
The zeroth group of the shifted filtration is $G_i$; it need not be taut.
The group operation is pointwise multiplication.
By \cite[Corollary B.4 and Proposition B.6]{v15-gi-14-GTZ}, these are groups with
\begin{equation}\label{v15:gi:14:eq:polyfiltered}
 [P_i,P_j]\subseteq P_{i+j},\qquad P_i\triangleleft P_0.
\end{equation}
Those results concern abstract filtered groups and require no equality of
the first two groups. Their conditions hold because the shifted groups
inherit nesting and the commutator law from $G_\bullet$.

Here is an explicit finite-dimensional description, including the smooth
real shift needed below. Let $\GIOfg_i=\log G_i$ and define
\[
 V_i=\left\{X(t)=\sum_{j=0}^{s-i} X_jt^j:X_j\in\GIOfg_{j+i}\right\}
 \quad(0\leq i\leq s),\qquad V_i=0\quad(i>s).
\]
Use pointwise Lie bracket. Bracketing coefficients gives
\begin{equation}\label{v15:gi:14:eq:Vbracket}
 [V_i,V_j]\subseteq V_{i+j}.
\end{equation}
For instance, a bracket of a degree-$a$ coefficient in $\GIOfg_{a+i}$ and a
degree-$b$ coefficient in $\GIOfg_{b+j}$ belongs to $\GIOfg_{a+b+i+j}$ and has
polynomial degree $a+b$. Thus $V_i$ is a finite-dimensional nilpotent Lie
algebra for every $i$. Its simply connected group is obtained by the
terminating BCH formula, or equivalently by pointwise exponentiation.

\begin{lemma}\label{v15:gi:14:lem:Pi}
The pointwise exponential identifies that group with $P_i$. In particular
$P_i$ is connected and simply connected, and its smooth topology can be
described either by the coefficients of $\log p(t)$ or by the Taylor
coefficients in \cite[Lemma B.9]{v15-gi-14-GTZ}.
\end{lemma}
\begin{proof}
If $i>s$, the zeroth group of the shifted filtration is trivial, so both
$P_i$ and $V_i$ are trivial. Fix $0\leq i\leq s$ and put $m=s-i$.
We use the connectedness of every $G_j$ stated in the setup.

First recall explicitly why the global exponential coordinates used here
are smooth. For a connected nilpotent Lie group $H$, let $Q_a(x)=x^a$,
where $a$ is a positive integer. In left tangent trivializations,
differentiating the $a$ factors gives
\[
 DQ_a(x)=\sum_{k=0}^{a-1}\operatorname{Ad}(x)^{-k}.
\]
The infinitesimal adjoint action lowers the lower central series.
An exponential neighborhood generates $H$, so $\operatorname{Ad}(x)$
induces the identity on each successive quotient of that series.
Consequently $DQ_a(x)$ induces multiplication by $a$ on each quotient
and is invertible. The usual local Lie exponential is smooth with
derivative the identity at zero. Given $X\in\mathfrak h$, choose $a$ so
large that $X/a$ lies in this local chart. The identity
\[
 \exp_H(Y)=Q_a(\exp_H(Y/a))
\]
then proves that the exponential has invertible derivative near $X$.
For simply connected nilpotent $H$, its global bijectivity is the
exponential homeomorphism recalled in \cite[Section 3]{v15-gi-14-GTZ}.
Its local smooth inverses therefore agree to give a smooth global inverse.
This argument uses the local Lie exponential, differentiation of group
operations and the inverse function theorem. It does not infer
smoothness from the word \emph{homeomorphism}.

By \cite[Proposition B.6]{v15-gi-14-GTZ}, $P_i$ is a group under pointwise
multiplication. By \cite[Lemma B.9]{v15-gi-14-GTZ}, its elements have the unique
increasing-order Taylor expansion
\[
 p(n)=\prod_{j=0}^{m}u_j^{\binom nj},\qquad u_j\in G_{i+j},
 \qquad n\in\GIOZ.
\]
The coefficients are recovered by the triangular formulas
\[
 u_0=p(0),\qquad
 u_j=\left(\prod_{k<j}u_k^{\binom jk}\right)^{-1}p(j).
\]
The order of the displayed inverse product is part of the formula.
Transfer the product manifold structure of $\prod_{j=0}^m G_{i+j}$
to $P_i$ using these coordinates. Each evaluation at an integer is
smooth. Evaluation of a product or inverse followed by the recovery
formulas proves smoothness of the group operations. The recovered
coefficients lie in the prescribed embedded subgroups; smoothness into
those subgroups follows from their local embedded charts.
The same formulas identify this topology with the pointwise topology.

Evaluation at $0,\ldots,m$ is an injective group homomorphism
$P_i\longrightarrow G^{m+1}$. Hence $P_i$ is nilpotent, with class at most
the class of $G$. Each closed connected subgroup of a simply connected
nilpotent group is simply connected. In exponential coordinates it is
the exponential of its Lie subalgebra: the finite nilpotent BCH formula
makes that exponential image a connected subgroup, and uniqueness of
the connected subgroup with that Lie algebra gives the assertion.
Thus the product manifold above, and therefore $P_i$, is connected and
simply connected. This step does not assert that $P_0$ has class at most $s$.

At the identity, differentiating the Taylor parametrization gives
\[
 (X_0,\ldots,X_m)\longmapsto
 \left(n\longmapsto\sum_{j=0}^m\binom nj X_j\right),
 \qquad X_j\in\GIOfg_{i+j}.
\]
The binomial-to-monomial matrix is triangular with nonzero diagonal.
A contribution to degree $b$ comes only from $j\geq b$, and
$\GIOfg_{i+j}\subseteq\GIOfg_{i+b}$. Descending triangular inversion preserves
these same coefficient spaces. This identifies the tangent space
with exactly $V_i$.
Evaluation maps are smooth homomorphisms, so their derivatives preserve
Lie brackets. The finitely many evaluations just used determine every
tangent vector. They therefore identify the bracket with the pointwise
bracket on $V_i$.

Uniqueness of one-parameter subgroups now gives
\[
 \operatorname{ev}_n(\exp_{P_i}X)=\exp_G(X(n))\qquad(n\in\GIOZ).
\]
Apply the smooth-global-exponential argument above to $P_i$.
Its inverse assigns to $p$ exactly the Lie polynomial whose values
are $\log_G(p(n))$. Conversely every $X\in V_i$ gives an element of $P_i$.
This proves the pointwise exponential identification and the smooth
logarithmic coefficient coordinates. The group law is determined by
multiplication in $G$ at the integer arguments. Under the ordinary
nilpotent BCH identification it is the pointwise BCH law.
The coordinate spaces and finite evaluation maps depend only on the
fixed filtration and $i$.
\end{proof}

For real $u$, define $T^u$ on $V_i$ by $T^uX(t)=X(t+u)$.
Coefficient expansion and nesting show that this preserves $V_i$.
It is a smooth one-parameter group of Lie algebra automorphisms, since
the bracket is pointwise. Exponentiating defines a smooth automorphism
group on $P_i$, which at integer $u$ is the ordinary shift of polynomial
sequences. Its infinitesimal action is ordinary differentiation, and
\begin{equation}\label{v15:gi:14:eq:derivative}
 \frac{d}{dt}V_i\subseteq V_{i+1}.
\end{equation}
Indeed the coefficient of $t^j$ in $X'(t)$ belongs to
$\GIOfg_{j+1+i}$.

\subsection{The lattice and the same-step semidirect product}

Define
\[
 \Lambda_i=\{p\in P_i:p(n)\in\Gamma\ \hbox{for every }n\in\GIOZ\}.
\]
It is a subgroup by pointwise multiplication. We give the needed lattice
argument to avoid relying on the disputed step estimate in Appendix C.

\begin{lemma}\label{v15:gi:14:lem:lattice}
$\Lambda_i$ is a discrete cocompact subgroup of $P_i$. Integer shifts
preserve $\Lambda_i$.
\end{lemma}
\begin{proof}
It is enough to consider $i\leq s$ and put $d=s-i$; for larger $i$ both
groups are trivial. Discreteness follows from evaluation at
$0,\ldots,d$. A lattice-valued polynomial close enough to the identity
at these finitely many points equals the identity at all those points,
since $\Gamma$ is discrete. Taylor uniqueness then makes it the identity
polynomial.

For cocompactness, use induction to write
\[
 p=\gamma_j h_j c_j,
\]
where $\gamma_j\in\Lambda_i$, $c_j$ ranges over a fixed compact subset
of $P_i$, and $h_j$ is the identity at $0,\ldots,j-1$.
For $j=0$ take $\gamma_0=c_0=1$ and $h_0=p$.
Taylor expansion shows $h_j(j)\in G_{j+i}$.
Choose a fixed compact left fundamental set for the lattice
$G_{j+i}\cap\Gamma$ in $G_{j+i}$ and write
$h_j(j)=\lambda_j v_j$, where $\lambda_j$ is in that lattice and $v_j$
is in the compact set. Define
\[
 \ell_j(t)=\lambda_j^{\binom tj},\qquad
 v_j^*(t)=v_j^{\binom tj}.
\]
Both are in $P_i$, by the same one-parameter difference check as above;
$\ell_j\in\Lambda_i$ because $\binom nj$ is an integer for every
$n\in\GIOZ$. Now put
\[
 \gamma_{j+1}=\gamma_j\ell_j,\qquad
 h_{j+1}=\ell_j^{-1}h_j(v_j^*)^{-1},\qquad
 c_{j+1}=v_j^*c_j.
\]
These formulas preserve the factorization and make $h_{j+1}$ the identity
at $0,\ldots,j$. The possible $v_j^*$ form a compact subset of $P_i$,
so the possible $c_{j+1}$ remain in a fixed compact set.
At $j=d+1$, Taylor uniqueness gives $h_{d+1}=1$.
Thus $p=\gamma c$ with $c$ in a fixed compact set, which proves left
cocompactness; inversion gives right cocompactness. Finally an integer
shift preserves both polynomiality and all integer lattice values, and
its inverse does also.
\end{proof}

Now use the \emph{increment} group $P_1$, not $P_0$, and set
\begin{equation}\label{v15:gi:14:eq:semi}
 \widetilde G=P_1\rtimes_T\GIOR,
 \qquad (p,u)(q,v)=(p\,T^u q,u+v),
 \qquad \widetilde\Gamma=\Lambda_1\rtimes_T\GIOZ.
\end{equation}
The real action just constructed makes $\widetilde G$ a connected simply
connected Lie group, diffeomorphic to $P_1\times\GIOR$.
Its Lie algebra is $V_1\rtimes\GIOR\tau$, where $[\tau,X]=X'$.
Give this Lie algebra the filtration
\[
 \mathfrak l_0=\mathfrak l_1=V_1\oplus\GIOR\tau,
 \qquad \mathfrak l_j=V_j\quad(j\geq2).
\]
Equations \eqref{v15:gi:14:eq:Vbracket} and \eqref{v15:gi:14:eq:derivative} prove
$[\mathfrak l_i,\mathfrak l_j]\subseteq\mathfrak l_{i+j}$ for all indices.
In particular $\mathfrak l_{s+1}=0$. Thus $\widetilde G$ is nilpotent of
step at most $s$. This explicit shifted-filtration calculation is the
same-step bound used in the theorem.

The subgroup $\widetilde\Gamma$ is discrete, because its clock coordinate
is integral and $\Lambda_1$ is discrete. It is cocompact as well.
Given $(p,u)$, left-multiply by $(1,-\lfloor u\rfloor)$, then left-reduce
the resulting polynomial coordinate modulo $\Lambda_1$ using Lemma
\ref{v15:gi:14:lem:lattice}. The resulting representative lies in the product of a
fixed compact set in $P_1$ and $[0,1]$. This proves left cocompactness;
inversion proves right cocompactness. The operations used for left reduction
belong to $\widetilde\Gamma$ by the integer-shift assertion of the lemma.

\subsection{The orbit identity and fixed compact fiber}

Let $\theta:\widetilde G/\widetilde\Gamma\longrightarrow\GIOR/\GIOZ$ be the
map induced by the clock coordinate, and set $K=\theta^{-1}(0)$.
It is a fixed compact submanifold canonically identified with
$P_1/\Lambda_1$: every element of the fiber can be written
$(p,0)\widetilde\Gamma$, and two such representatives agree precisely when
their polynomial coordinates differ on the right by an element of
$\Lambda_1$.
Evaluation at zero gives a smooth map
\begin{equation}\label{v15:gi:14:eq:pi}
 \pi:K\longrightarrow G/\Gamma,\qquad
 \pi((p,0)\widetilde\Gamma)=p(0)\Gamma.
\end{equation}
It is well-defined because $\lambda(0)\in\Gamma$ for
$\lambda\in\Lambda_1$. Since $K$ is compact, $\pi$ is Lipschitz for
fixed smooth metrics.

The normalized polynomial $h\in P_0$ has $h(0)=1$. Define its right
increment by
\[
 b(t)=h(t)^{-1}h(t+1).
\]
The left increment $h(t+1)h(t)^{-1}$ lies in $P_1$ directly from the
iterated-difference definition. Proposition B.6 gives $P_1\triangleleft
P_0$, so conjugating that left increment by $h(t)^{-1}$ shows
$b\in P_1$. This verifies the orientation of the increment actually
used, rather than substituting a left increment into a right-telescoping
product.

Put $a=(b,1)\in\widetilde G$. For every integer $n$ define
$b_n(t)=h(t)^{-1}h(t+n)$. The exact cocycle identity
\[
 b_{n+m}(t)=b_n(t)\,b_m(t+n)
\]
and $b_0=1$ prove, for positive and negative integers alike,
\[
 a^n=(b_n,n).
\]
Modulo $\widetilde\Gamma$, right multiplication by $(1,-n)$ gives
$a^n\widetilde\Gamma=(b_n,0)\widetilde\Gamma\in K$.
Consequently
\begin{equation}\label{v15:gi:14:eq:orbit}
 \pi(a^n\widetilde\Gamma)=b_n(0)\Gamma=h(n)\Gamma,
 \qquad n\in\GIOZ.
\end{equation}
The last equality uses the normalization $h(0)=1$.

\subsection{Test functions, uniformity and the inverse-theorem join}

The function $F_q\circ\pi$ is $1$-bounded on $K$ and has a Lipschitz
constant bounded by a fixed multiple $L$ of $\operatorname{Lip}(F)$.
Extend its real and imaginary parts separately by the explicit McShane
formula
\[
 E_u(x)=\inf_{y\in K}\bigl(u(y)+L\,d(x,y)\bigr).
\]
This equals $u$ on $K$ and is $L$-Lipschitz: the triangle inequality
proves the Lipschitz assertion, and the original Lipschitz inequality
proves equality at points of $K$. Clip each extension to $[-1,1]$,
combine them as a complex function, and apply metric projection onto
the closed unit disk. This projection $Q$ is $1$-Lipschitz: the
nearest-point inequalities for the disk give
\[
 \langle x-Qx,Qy-Qx\rangle\leq0,\qquad
 \langle y-Qy,Qx-Qy\rangle\leq0,
\]
and therefore $\|Qx-Qy\|^2\leq\langle Qx-Qy,x-y\rangle
\leq\|Qx-Qy\|\|x-y\|$. It also fixes the original values on $K$.
Thus the result $\widetilde F$ is $1$-bounded, agrees with
$F_q\circ\pi$ on $K$, and is at most $\sqrt2 L$-Lipschitz.
Equation \eqref{v15:gi:14:eq:orbit} proves the exact equality in Theorem
\ref{v15:gi:14:thm:linear}.

The nilmanifold, lattice, compact fiber, evaluation map and all metric
constants depend only on the fixed original filtered nilmanifold and $s$.
The possibly unbounded polynomial coefficients occur only in the allowed
translation element $a$ and in the bounded representative $q$ affecting the
test function. They do not alter the geometry or its Lipschitz bound.
For limit data the same algebraic formulas and one fixed bounded
representative section give limit choices of $b$ and $a$. Apply the explicit
extension formula to each representative test function. Its uniform bound
and Lipschitz constant give a bounded Lipschitz limit function on the same
standard nilmanifold. The ordinary orbit identities hold at every member
of the defining sequences and therefore at every limit integer.

Suppose now that a polynomial inverse theorem supplies, for fixed $s$ and
$\delta$, a finite family of rational degree-$s$ filtered nilmanifolds,
a fixed Lipschitz bound, and a positive correlation lower bound against
$1$-bounded test functions. Apply the theorem to each member of that finite
family, taking the maximum of its finitely many fixed metric constants.
There is still a finite family, now of step at most $s$ with linear orbits,
and the correlation is unchanged because the represented sequences are
equal pointwise. No comparison of Gowers norms is involved in this join.
This proves the final assertion of Theorem~\ref{v15:gi:14:thm:linear}.

The construction also covers a point input group: $P_1$ is trivial and the
clock group yields a circle with a constant test function, or it may be
replaced by a point. For $s=1$ the increment group consists of constant
polynomials in an abelian group and the shift action is trivial, so
$\widetilde G$ is abelian, as the required bound predicts.

\subsection{Comparison with the failed full-group construction}

The original Appendix C used $P_0\rtimes\GIOR$. For the taut abelian
degree-$s$ filtration on $\GIOR$, the commutator of its clock with
$p(t)=\tfrac12\binom ts$, repeated $s$ times, is the nonzero constant
$1/2$. That is a commutator of length $s+1$, so that group is not
$s$-step. This adverse example is recorded in the separate finding packet.
The present construction uses $P_1$ and its increment $b$, reducing
the polynomial degree available in the clock action by one. Its Lie
filtration above verifies the required step bound for all inputs, not
only this example. It changes the construction, while retaining the
same-step linear conclusion needed by GI($s$).

The finite-family polynomial inverse theorem is still a separate premise.
This unit does not prove its rank reduction, symmetrization, or
ultralimit-to-finite-family extraction. Those dependencies remain explicit
in the theorem assembly.


\section{A normalized local converse for the GI assembly}\label{v15:gi:GI-CONVERSE}
\newcommand{\GIPR}{\mathbb R}
\newcommand{\GIPZ}{\mathbb Z}
\newcommand{\GIPC}{\mathbb C}
\newcommand{\GIPE}{\mathbb E}


This section derives the precise correlation-to-Gowers-norm implication used by the
assembly. It imports the published compact cube constraint, Green--Tao,
\emph{Linear equations in primes}, Proposition 11.5, p.~1799, with its
definitions in 11.4 and proof in Appendix E. It uses the internal increment linearization for polynomial tests. It does not import
Proposition 11.2's averaging-domain transition, assume an inverse
theorem, or claim to reprove the cube constraint.

\subsection{Available geometric assertion and needed conclusion}

Fix an integer $d\geq1$ and a fixed $d$-step nilmanifold $M=G/\Gamma$
with a smooth metric. The available cube constraint says there is a
compact subset
$\Sigma\subseteq M^{\{0,1\}^{d+1}\setminus\{0\}}$ and a continuous
map $P:\Sigma\to M$ such that, for every $g\in G$, $x\in M$,
$n\in\GIPZ$ and $h\in\GIPZ^{d+1}$, the nonzero vertices
$(g^{n+\omega\cdot h}x)_{\omega\ne0}$ lie in $\Sigma$ and
\[
 g^nx=P((g^{n+\omega\cdot h}x)_{\omega\ne0}).
\]
No bound on $g$ or the integers is imposed. In our application $G$ is
connected simply connected and at most $d$-step; these satisfy the
published nilmanifold hypothesis.

\begin{theorem}\label{v15:gi:15:thm:converse}
For fixed $d,M,L$ and every $\kappa>0$ there is $c>0$ such that
every $1$-bounded $v:[N]\to\GIPC$ and $1$-bounded $L$-Lipschitz test
$F:M\to\GIPC$ satisfying
\[
 |\GIPE_{n\in[N]}v(n)F(g^nx)|\geq\kappa
\]
obey $\|v\|_{U^{d+1}[N]}\geq c$. The bound is uniform in $N,g,x,F$.
The same conclusion holds for a fixed finite family of rational
degree-at-most-$d$ polynomial models, using the exact same-step
increment linearization and taking a finite maximum of its bounds.
\end{theorem}

\subsection{A finite tensor approximation uniform in the test}

The algebra generated by coordinate functions that are Lipschitz on
$M$ separates points on the compact set $\Sigma$, contains constants
and is closed under conjugation. Complex Stone--Weierstrass gives,
for every fixed continuous $F$ and $\epsilon>0$, an approximation
\[
 F(P(z))=\sum_{a=1}^A c_a\prod_{\omega\ne0}H_{\omega,a}(z_\omega)
                 +R(z),\qquad |R(z)|\leq\epsilon,
\]
where the factors are $1$-bounded Lipschitz functions. Scaling each
factor into the unit disk is absorbed into the finite coefficients
$c_a$; no coefficient is discarded. By Arzela--Ascoli the class
$\|F\|_\infty\leq1$, $\operatorname{Lip}(F)\leq L$ is compact
in the uniform norm. A finite uniform net, approximated as above,
makes $A$, $\sum_a|c_a|$ and all factor bounds depend only on
$M,L,\epsilon$. Choose a net radius $\epsilon/2$ and approximate
its members to $\epsilon/2$ to retain the displayed total error.

For each orbit, put
$b_{\omega,a}(t)=H_{\omega,a}(g^tx)$. The cube constraint gives
\begin{equation}\label{v15:gi:15:eq:pointwise}
 F(g^nx)=\sum_a c_a\prod_{\omega\ne0}
 b_{\omega,a}(n+\omega\cdot h)+O(\epsilon)
\end{equation}
for every integer $n,h$, uniformly. The factors are $1$-bounded;
their dependence on the orbit will no longer matter in the analytic
estimate.

\subsection{Normalized smooth averaging and a uniform Fourier bound}

Choose once and for all a smooth $\sigma:\GIPR\to[0,1]$ supported on
$[-1,2]$ and equal to one on $[0,1]$. Define
\[
 S_N=\sum_{h\in\GIPZ}\sigma(h/N),\qquad
 \mu_N(h)=\sigma(h/N)/S_N.
\]
Then $N\leq S_N\leq3N+1\leq4N$, and $\mu_N$ is a probability
measure with finite support. Average \eqref{v15:gi:15:eq:pointwise} against
the product probability measure $\prod_j\mu_N(h_j)$ on all
integer $h$. The error remains at most $\epsilon$, with no missing
normalizing scalar. Thus
\[
 F(g^nx)=\sum_a c_a B_a(n)+O(\epsilon),\quad
 B_a(n)=S_N^{-(d+1)}\sum_{h\in\GIPZ^{d+1}}
   \prod_j\sigma(h_j/N)\prod_{\omega\ne0}
                  b_{\omega,a}(n+\omega\cdot h).
\]

Take $Q=C_dN$ with a fixed integer $C_d>10(d+2)$, and identify
representatives of $\GIPZ/Q\GIPZ$ with a centered interval. The sampled
cutoff $w_N(h)=\sigma(h/N)$ has support strictly inside that interval.
Write its discrete Fourier expansion as
$w_N(h)=\sum_r\widehat w_N(r)e(rh/Q)$ with normalized coefficients.
There is a bound
\begin{equation}\label{v15:gi:15:eq:l1}
 \sum_{r\in\GIPZ/Q\GIPZ}|\widehat w_N(r)|\leq C'_d
\end{equation}
independent of $N$. To prove it, the zero coefficient and every
coefficient are at most $S_N/Q=O_d(1)$. For $r\ne0$, summation by
parts twice gives
\[
 |1-e(r/Q)|^2 |\widehat w_N(r)|
 \leq Q^{-1}\sum_h|\Delta^2w_N(h)|\leq C/(NQ).
\]
The final bound follows from the second derivative of the fixed smooth
cutoff, with only $O(N)$ possibly nonzero terms. For the distance
$r_*\in[1,Q/2]$ to zero modulo $Q$,
$|1-e(r/Q)|\geq4r_*/Q$. Thus
$|\widehat w_N(r)|\leq C_d/r_*^2$ and summation proves
\eqref{v15:gi:15:eq:l1}. Smooth zero extension of $\sigma$ ensures that the
same finite-difference bound holds at its support boundaries.

\subsection{The exact Gowers estimate}

Zero-extend $v$ to $\GIPZ/Q\GIPZ$. Extend each bounded integer function
$b_{\omega,a}$ to the cyclic group so that it agrees on all integer
arguments that occur when $n\in[N]$ and $h_j\in[-N,2N]$.
Those arguments lie between $-(d+1)N$ and $(2d+3)N$; the choice
of $C_d$ places the whole interval strictly inside the centered
representative range. There is therefore no conflicting cyclic value.
The normalized correlation with $B_a$ is exactly
\[
 \frac{Q}{N}\left(\frac{Q}{S_N}\right)^{d+1}
 \GIPE_{n,h\in\GIPZ/Q\GIPZ}
 v(n)\prod_j w_N(h_j)
          \prod_{\omega\ne0}b_{\omega,a}(n+\omega\cdot h).
\]
The prefactor is bounded above by a constant depending only on $d$.
This identity comes from an all-integer smooth average; no hard box
average is replaced by a larger summation domain.

Expand the product of cutoffs by Fourier series. The sum of absolute
values of its coefficients is at most $(C'_d)^{d+1}$. In each term,
the factor $e(\sum_jr_jh_j/Q)$ can be absorbed into the nonzero
vertex functions. For each $j$, choose an index $k\ne j$, which
exists since $d+1\geq2$, and use
$h_j=(n+h_j+h_k)-(n+h_k)$. Multiply the functions at vertices
$e_j+e_k$ and $e_k$ by the corresponding opposite linear phases.
Their absolute bounds remain one. The Gowers--Cauchy--Schwarz
inequality for the full $(d+1)$-cube gives the absolute bound
$\|v\|_{U^{d+1}(\GIPZ/Q\GIPZ)}$ for every such term. Complex
conjugations are absorbed into the vertex functions, leaving the
norm and absolute bounds unchanged. Hence
\[
 |\GIPE_{[N]}v B_a|\leq K_d\|v\|_{U^{d+1}(\GIPZ/Q\GIPZ)}
 \leq K_d\|v\|_{U^{d+1}[N]}.
\]
The last inequality uses the prescribed normalized interval norm:
the cyclic norm is the interval norm multiplied by the cyclic norm
of the interval indicator, which is at most one. If necessary enlarge
$C_d$ also to exceed $2^{d+1}$ so this single ambient group realizes
that interval normalization.

Combining the finite tensor approximation with this estimate yields
\[
 |\GIPE_{[N]}v(n)F(g^nx)|
 \leq K_d\Bigl(\sum_a|c_a|\Bigr)
                     \|v\|_{U^{d+1}[N]}+\epsilon.
\]
Choose $\epsilon<\kappa/2$. This proves
Theorem~\ref{v15:gi:15:thm:converse} with a positive bound depending only on
$d,M,L,\kappa$. A fixed bound other than one for $v$ or $F$ is
handled by dividing by that bound first. This proof is complex-valued
throughout; it does not need a real-valued dual-norm convention.

\subsection{Polynomial input and source comparison}

For the actual assembly input, a scalar polynomial nilsequence has
degree $d=s-2\geq1$ on fixed rational geometry. The accepted
increment linearization001 represents it exactly on a fixed group
of step at most $d$, with a fixed Lipschitz multiplier and unchanged
norm bound. Applying the theorem there and taking finite maxima
over a finite family gives precisely the positive normalized
$U^{s-1}[N]$ norm required by the last induction step. This is a
converse implication and has no inverse-theorem premise.

The literal published proof of 11.2 on p.~1801 defines an average
over $h\in[N]^{d+1}$ with $\sigma(h_j/N)$, while $\sigma=1$ on
$[0,1]$. Its later passage from (11.6) to (11.7) widens the shift
domain to a cyclic group with the same cutoff supported on $[-1,2]$.
That changes the set of summands. The local proof above instead
starts with the normalized all-integer smooth average and tracks
its exact cyclic prefactor. Whether an earlier source correction
already supplied this interpretation has not been established.
The observation is a bounded proof-transition issue, not a
counterexample to Proposition 11.2 or the converse theorem.

The applications are the normalized average, its dual-norm estimate, Corollary 11.6 and the present converse. Published Proposition 11.5 remains the compact cube-constraint input. The full inverse theorem is assembled in Section~\ref{v15:gi:ASSEMBLY}.


\section{Uniform selection of represented lower tests}\label{v15:gi-selection}
Fix standard $s\ge3$, $d=s-2\ge1$, a nonempty internal interval $J=[N]$, and an internal parameter set $H$. The actual represented unit families in the inverse-theorem construction are retained throughout. Suppose the lower-test correlation premise holds for every $h$. In its tensor convention, a bounded scalar component $u_h$ has positive standard-size correlation with a scalar nilsequence of degree at most $d$. Apply the ordinary-test approximation in GTZ Lemma 6.4, with error smaller than half this correlation divided by $\max(1,\|u_h\|_\infty)$. Its finite component sum, with finite coefficient sum, contains a component of an ordinary globally equivariant unit-vector test retaining positive standard-size correlation. Conjugation preserves that class. Compatible polynomial Taylor coefficients may be limit coefficients. This proves pointwise existence before uniformity.

Enumerate the countable rational filtered model types, equip each with one fixed smooth metric, and enumerate their integral vertical characters and finite target dimensions. Let the ordinary set $T_M$ contain the first $M$ tags, all ordinary global unit tests with Lipschitz and supremum bounds at most $M$, all compatible Taylor coefficients without any bound on them, the finite coordinate and conjugation choices, and reciprocal correlation thresholds with positive integer denominator at most $M$. Choose these sets increasing in an ordinary disjoint union $T$. The natural extension of the single ordinary sequence $M\mapsto T_M$ supplies an internal exhaustion. This is stronger than an arbitrary countable family of internal sets. Each pointwise witness belongs to ${}^*T_M$ for some standard $M$.

The relation $R(h,t)$ says that the modulus of the normalized finite correlation sum is at least the encoded threshold. Evaluation and finite sums are internal; norm one and global equivariance are already conditions on the ordinary parameter set and transfer. No external predicate of positive standard size enters $R$. The least internal stage supporting a witness is pointwise bounded by a standard integer, hence bounded by one standard $M_0$ on $H$. Indeed, otherwise the internal sets $E_M$ of points requiring a stage greater than $M$ would be decreasing and nonempty. Choose decreasing representatives $E_{M,i}$ and put $S_M=\{i:E_{M,i}\ne\varnothing\}$. The $S_M$ are decreasing and ultrafilter-large. Put
\[k_i=\max\bigl(\{k\le i:i\in S_k\}\cup\{0\}\bigr)\]
and choose a point of $E_{k_i,i}$ when $k_i>0$. For each fixed $M$, on $S_M\cap\{i\ge M\}$ it belongs to $E_{M,i}$. Its limit belongs to every $E_M$, contradicting pointwise standard-stage existence. Internal choice inside ${}^*T_{M_0}$ now supplies the witness map. This is the bounded-stage argument of GTZ Appendix A, Definitions A.5 and A.10 and Lemma A.11; the countability of model tags is as on p.~1250. No compactness of the unbounded orbit coefficients is used.

Pad the finitely many target dimensions with zeros. If $H$ has positive density, finite pigeonhole fixes the coordinate and convention choices on a positive-density internal subset. Outside it use a constant unit representative and keep the indicator separate. This is the later rank-step selection. For the earlier two-family application, keep both original strong represented families $\chi(h,n)$ and $\chi_h(n)$ unchanged and apply the same argument to
\[u_{h,a,b}(n)=\Delta_hf(n)\,\overline{\chi(h,n)_a}\,\overline{\chi_h(n)_b},\]
with conjugations as required by the tensor convention. No polynomial dependence on $h$ is needed. Finite pigeonhole fixes the coordinate pair. Multiplication of the selected test by the internal unit phase of its nonzero correlation aligns its average; extend this scalar test by zero outside the retained subset. The phase has degree zero. The resulting lower tests and correlation threshold have common standard bounds. This selects a lower correlation witness, not a replacement of either incoming family, and does not upgrade arbitrary weak nilcharacters to strong ones. The following shift and whole-matrix reconstruction arguments are the ones proved in the preceding inverse-theorem sections.

\section{Integral endomorphisms and polarized automorphisms}\label{v15:tate-polarized}
We supply the integral-lattice step in the marked-surface construction. The inputs are Milne's published chapter \cite{milne1986-abelian}, especially \S8 (torsion and Tate modules), \S12 (decomposition, degree polynomial and divisibility), and \S17 (Rosati positivity). Work over an algebraically closed field of characteristic different from $3$.

For an abelian surface $A$, $T_3A$ is free of rank four and $T_3A/3T_3A=A[3]$: successively lift torsion points for surjectivity, and shift a compatible sequence for the kernel. A unital action of the division quaternion algebra over $\mathbf Q_3$ on $V_3A$ is faithful and has rank one. An endomorphism acting trivially on the Tate module is zero. Indeed, after isogeny decomposition, a nonzero map between simple isogenous factors is an isogeny and cannot kill torsion of unbounded order.

Finite rational rank alone does not prove finite generation of integral endomorphisms. Let $S$ be simple, let $M$ be generated freely by $e_1,\ldots,e_r$ in $\operatorname{End}(S)$, put $V=\mathbf QM$ and $D=V\cap\operatorname{End}(S)$. The degree polynomial is a positive integer on every nonzero point of $D$. A neighborhood of zero where its absolute value is less than one therefore contains no such point. Applying this to differences makes $D$ uniformly discrete. A compact fundamental parallelepiped for $M$ meets $D$ finitely, so $D/M$ is finite and $N_MD\subset M$ for some positive integer $N_M$.

The matrices $T_3(e_i)$ are independent over $\mathbf Q_3$. Scale a purported relation to coefficients $a_i\in\mathbf Z_3$, not all zero, and choose integers $n_i\equiv a_i\pmod{3^m}$. The endomorphism $\sum_i n_ie_i$ kills $S[3^m]$. The divisibility lemma gives
\[3^{-m}\sum_i n_ie_i\in D,\qquad N_Mn_i/3^m\in\mathbf Z.\]
For nonzero $a_j$, taking $m>v_3(N_M)+v_3(a_j)$ contradicts the latter inclusion. Consequently $r\le(2\dim S)^2$. All finite independent sets satisfy this bound, so the rational endomorphism space is finite dimensional. Choose $M$ spanning it and repeat the lattice argument to prove that the entire integral ring is finitely generated. For a general abelian variety, decomposition embeds its endomorphisms into a finite matrix of Hom groups between simple factors. Nonisogenous factors contribute zero; composing with a fixed isogeny embeds each remaining Hom group into an endomorphism group of a simple factor. Thus the integral endomorphisms again form a lattice.

A pointed polarization-preserving automorphism $u$ satisfies $u^\dagger u=1$. Rosati positivity makes $q(u)=\operatorname{Tr}(u^\dagger u)$ positive definite. Such automorphisms lie on the compact sphere $q(u)=\operatorname{Tr}(1)=4$, whose intersection with the lattice is finite. This proves the required finiteness; translations are excluded by the pointed-morphism convention.

\section{Descent of the marked abelian scheme}\label{v15:marking-descent}
Set $R=\mathbf Z[1/30]$. The earlier moduli construction supplies the smooth projective full-five scheme $X_5$, its universal polarized surface and order action, with all required determinant components. The geometric module calculation supplies the three-markings. The order-linear Isom sheaf is finite etale: after an etale trivialization of torsion, it is the finite subset of bijections preserving addition and finitely many order generators. Stacks Tag 03RV represents it, and the geometric calculation proves surjectivity. Hence $X_{35}\to X_5$ is finite etale. A finite map is projective (0B3I); composition is projective here over the affine base (01W7, Lemma 29.44.14). Thus $X_{35}$ is smooth projective and regular.

Let $G_5=(\mathcal O/5\mathcal O)^\times$ change the five-marking. For every nonidentity $g$, its equalizer with the identity is closed and has no geometric point by the earlier automorphism calculation, hence is empty. A section of the constant group is locally a fixed $g$. The action is therefore free on arbitrary test schemes, including nonreduced ones, without inverting $|G_5|$. Tag 07S7 represents $Y=X_{35}/G_5$ and makes $q:X_{35}\to Y$ a torsor. Descent of finite and etale morphisms (02LA, 02VN) makes it finite etale surjective.

Smoothness descends (036U). The pullback of the diagonal of $Y/R$ is the disjoint union of the closed graphs of $G_5$, so 02L6 makes $Y/R$ separated. The proper cover makes it universally closed (03GN) and quasi-compact, hence proper. The norm of an ample line bundle descends existence of an ample line bundle (0BD4); properness gives projectivity (0B45). This norm proves projectivity of the quotient, not a replacement of its polarization by a multiple.

Before duality, use the regular-base projectivity assertion in BLR, \S8.4, proof of Corollary 6, p.~234, for these regular noetherian bases of finite type over $R$. The universal abelian scheme over $X_{35}$ is projective. Its relabeling descent datum is effective by Stacks 0CCJ. Full faithfulness (023Q) descends zero, addition, inverse and their identities. Properness and smoothness descend (02L1, 02VL), and geometric connectedness follows after choosing a point of the nonempty covering fibre. The result is an abelian scheme $A/Y$, again projective by the regular-base assertion.

BLR \S8.2, Theorem 1 and Proposition 4, pp.~210--211, and \S8.4, Theorem 5, p.~234, now apply to these projective flat finitely presented families with reduced irreducible geometric fibres and zero sections. They give the Picard functor, normalized Poincare bundle and dual. The fibrewise definitions of $\operatorname{Pic}^0$ and $\operatorname{Pic}^\tau$ commute with base change, so canonically
\[(A^\vee)\times_YX_{35}\simeq(A_{35})^\vee.\]
The actual principal polarization is invariant under relabeling and descends by 023Q. It remains an isomorphism (02L4) and is a polarization on each geometric fibre. Every order endomorphism and the three-marking descend similarly, with identities checked after $q$.

An automorphism of a full-three object trivial on geometric fibres acts trivially on finite etale five-torsion, since maps of finite constant sets are detected at geometric points. After adding a five-marking it is the identity by full-five rigidity, and therefore is the identity before the cover. For an arbitrary full-three object, its five-marking sheaf is a $G_5$-torsor; its map to $X_{35}$ descends uniquely to $Y$. Conversely pull back the descended universal object. These constructions are inverse after a faithfully flat cover and hence inverse. This identifies the full-three moduli problem, without an extra five-marking factor in its degree.

\section{The finite forgetful map and its Picard norm}\label{v15:picard-norm}
Diamond--Im, \S9.3, pp.~77--79 \cite{diamond-im}, supplies smooth proper models of $X_1(N)$ and $X_0(N)$ over $\mathbf Z[1/N]$ and the diamond-group quotient; for $N\le4$ the coarse curves coincide. Thus the actual forgetful map $u:X_1(N)\to X_0(N)$ is finite surjective over $\mathbf Q$. On quotient affine charts, the monic orbit polynomial makes elements integral over the invariants; finite type gives finiteness and lying-over gives surjectivity. The ineffective element $-1$ is harmless. Compose with the finite surjective modular map to the elliptic curve, obtaining $f:X_1(N)\to A$. At a closed point of the smooth target, its local ring is a DVR and the finite source module is torsion-free, hence free. Therefore $f$ is finite locally free of positive constant degree $d$.

Its norm is
\[N_f(\mathcal L)=\det(f_*\mathcal L)\otimes\det(f_*\mathcal O)^{-1}.\]
Locally trivialize $\mathcal L$ over the finite source. Changing a generator by a unit $u$ changes its determinant trivialization by $\det(m_u)$. These are exactly the function norms of Stacks 0BD2 and the transition functions of 0BCY. They prove multiplicativity; no multiplicativity of determinant for arbitrary tensor products of vector bundles is asserted. Restriction of scalars and exterior powers here commute with arbitrary base change. Moreover
\[N_{f_T}(\mathcal L\otimes\pi_X^*\mathcal M)
=N_{f_T}(\mathcal L)\otimes\pi_A^*(\mathcal M^{\otimes d}).\]
The norms therefore pass to the quotient by line bundles from $T$ and after fppf sheafification define a natural transformation of relative Picard sheaves. Representability and Yoneda give a rational group morphism. After algebraic closure it sends the connected identity component into the connected identity component, so restricts to $\operatorname{Pic}^0$. No rational point of $X_1(N)$ or universal rational line bundle is required. On divisor line bundles it is pushforward, by the valuation formula for the determinant norm (Silverman II.3.6(c),(d)). It is thus exactly the morphism used in the earlier point-difference proof of surjectivity.

\section{Labelled CM points and their common-place reduction}\label{v15:cm-reduction}
We give the CM and reduction arguments used in the conductor-one sum. Let $F$ be imaginary quadratic with maximal order $\mathcal O$, fix a complex CM type, let $q$ be an odd inert unramified prime, and let $N$ be prime to $q\operatorname{disc}(F)$ with all its prime divisors split in $F$. In the conductor-one application $\mathcal O^\times=\{\pm1\}$. For each $p^e\Vert N$ choose one prime $\mathfrak p$ over $p$, and set $\mathfrak n=\prod\mathfrak p^e$. Then $\mathcal O/\mathfrak n\simeq\mathbf Z/N\mathbf Z$. A labelled CM point consists of an elliptic curve, its actual $\mathcal O$-action of the fixed type, and the cyclic subgroup $A[\mathfrak n]$. Isomorphisms preserve these labels.

\subsection{Ideal action and reciprocity}
Complex uniformization and its CM refinement (Silverman VI.4.1, VI.4.1.1, VI.5.1.1 and VI.5.3, with the corrected formulation) identify such curves with $\mathbf C/\Lambda$, where $\Lambda$ is a fractional $\mathcal O$-ideal. Labelled isomorphisms correspond to multiplication by $F^\times$. The ideal action of an integral ideal $\mathfrak a$ is
\[\mathbf C/\Lambda\longrightarrow\mathbf C/\mathfrak a^{-1}\Lambda;\]
its kernel is $A[\mathfrak a]$ and its degree is $\mathrm N\mathfrak a$. For $\mathfrak a$ prime to $N$, it transports the specified level subgroup. Every ideal class has a representative prime to any prescribed finite set: choose an integral ideal $I$ in the inverse class and, by the ideal Chinese remainder theorem, an $x\in I$ with $v_{\mathfrak p}(x)=v_{\mathfrak p}(I)$ at that set. Then $(x)I^{-1}$ is integral and prime to the set.

The order class polynomial is monic integral (Fowler, \S2.1, pp.~612--613; Silverman C.11.2(a)), so all these $j$-invariants are algebraic integers. Silverman C.11.2(c),(d) gives the Hilbert class field $H=F(j(\Lambda))$. The Frobenius formula in C.11.4, at primes where the chosen models have good reduction, is
\[j(\Lambda)^{\operatorname{Frob}_{\mathfrak p}}=j(\Lambda\mathfrak p^{-1}).\]
To use it for all ideal classes, choose $H$-models for each of the finitely many $j$-values (Silverman III.1.4(c)) and exclude the ideal primes below every bad prime of these models, together with the primes dividing $Nq$. Choose class representatives prime to this enlarged set and factor them into prime ideals. Iterating the formula gives the class action. This does not assert good reduction at $q$ yet. An equality of $j$-invariants lifts to an isomorphism of the curves (Silverman III.1.4(b)); because the Galois action fixes $F$ and hence the CM type, faithfulness of the tangent character in complex uniformization makes it an $\mathcal O$-linear isomorphism. The selected level subgroup is then preserved. This asserts invariance of labelled isomorphism classes, not that arbitrary choices of isomorphisms are defined over $H$.

In the convention of Cai--Shu--Tian, the point indexed by $\mathfrak b$ is the isogeny $\mathbf C/\mathfrak b\to\mathbf C/(\mathfrak b\mathfrak n^{-1})$. Thus it is $[\mathfrak b^{-1}]$ times the reference point in the preceding convention. Their Artin compatibility (pp.~2524, 2527 and 2531) agrees after this one global inversion of class indices. Quadratic characters are unchanged by that inversion. All reductions below take place at one chosen place. The forward CM reduction criterion of Gross--Zagier, pp.~198--199 \cite{gz}, gives supersingularity when $(\operatorname{disc}(F)/q)=-1$, including $q=3$; no converse lifting theorem is needed.

\subsection{A simultaneous good model}
Consider the finite diagram of curves, order actions, selected level groups and ideal isogenies needed in the class sum, with isogeny degrees prime to $qN$. The integral $j$-invariants and Silverman VII.5.5 give potentially good reduction for each curve. A common finite extension gives all good models over one complete DVR $R$ of residue characteristic $q$, and defines the finitely many maps and order generators. These elliptic schemes are Neron models (BLR \S1.4, Proposition 2, p.~19). The Neron mapping property, BLR \S1.2, Definition 1 and Proposition 8, pp.~12,15, extends every generic morphism. Uniqueness extends all composition and order-action identities.

The special-fibre $\mathcal O$-action remains injective. For nonzero $\alpha$, its dual composite is multiplication by a positive integer (Silverman III.6.1 and III.6.2(a)); the dual is defined over the same field by uniqueness. After extension this identity still holds, and multiplication by that integer is nonzero on the special elliptic curve, of degree its square (III.6.2(d)), even if $q$ divides the integer.

For $n$ prime to $q$, the differential of $[n]$ is a unit. The etale criterion for maps between smooth curves (EGA IV$_4$, Corollary 17.11.2(c'), p.~84) makes it etale. It is proper by the graph and projection argument, and quasi-finite; hence finite (EGA III$_1$, Corollary 4.4.11). Therefore $A[n]$ is finite etale. Over the strict henselization it is constant, and specialization of geometric points is a bijection compatible with all endomorphisms, by EGA IV$_4$, Propositions 18.5.14--15. Simultaneous kernels are thus preserved as subsets of a constant group scheme.

Let $f:A\to B$ be one of the isogenies of degree $d$ prime to $q$. Its dual extends, and their composites are $[d]$. Their differentials are therefore units. Thus $f$ is etale, proper and quasi-finite, hence finite etale. Its degree is constantly $d$ and it is surjective: finite etale rank is locally constant (EGA IV$_4$, Proposition 18.2.3 and Remark 18.2.7(i)), and the elliptic target over the local base is connected. If $f$ is the ideal quotient for $\mathfrak a$, then $d\in\mathfrak a$, so its kernel lies in $A[d]$. Equality there with the common kernel of finitely many generators of $\mathfrak a$ holds generically and hence on the constant components; the special kernel is exactly $\overline A[\mathfrak a]$. The map
\[\overline A\times\ker\overline f\longrightarrow
 \overline A\times_{\overline B}\overline A,
 \qquad(x,t)\longmapsto(x,x+t)\]
is an isomorphism, with inverse $(x,y)\mapsto(x,y-x)$. Consequently $\overline B$ represents the quotient by this subgroup.

Choose $u,v$ with $ud+vN=1$. On $N$-torsion, $[u]\widehat f$ is the inverse of $f$, preserving all order eigensummands and therefore every chosen prime-power level subgroup, including powers of $2$. On $q^r$-torsion use compatible inverses of $d$ modulo $q^r$; this gives an isomorphism of $q$-divisible groups without an etaleness claim for them. Finally the tangent identity for the $\mathcal O$-linear map $f$, whose differential is nonzero, shows that the residue character $\mathcal O/q\to\overline{\mathbf F}_q$ is the same on both ends. It is not replaced by its conjugate.

\subsection{The oriented conductor-one correspondence}\label{v15:cm-oriented}
Use the enhanced Eichler class-set description and rational Skolem--Noether calculation proved in the quaternionic sections above. The source set of labelled points of the fixed CM type and level choices is a torsor under $\operatorname{Pic}(\mathcal O)$. We now verify the same description of the oriented supersingular target without assuming surjectivity of reduction.

Away from $Nq$, an $\mathcal O_p$-stable lattice is a principal fractional ideal locally, also at primes ramifying in $F$. At $p^e\Vert N$, the two integral idempotents of $\mathcal O_p=\mathbf Z_p\times\mathbf Z_p$ split the lattice chain. A cyclic quotient of order $p^e$ forces the two exponents to be $(e,0)$ or $(0,e)$; a mixed pair would not give a cyclic quotient. This holds at $p=2$ as well. At $q$, write the division algebra as
\[B_q=F_q+F_q\Pi,\qquad \Pi a=\bar a\Pi,\]
with $\Pi^2$ a central uniformizer times a unit. Conjugation by a maximal-order unit acts trivially on the residue field $\mathbf F_{q^2}$, whereas conjugation by $\Pi$ is its nontrivial automorphism. These are the two local embedding classes, distinguished by the residue orientation and parity of the norm valuation; they must not be identified by a normalizer quotient.

Rational Skolem--Noether fixes a labelled embedding of $F$. Its stabilizer is the centralizer $F^\times$, not the normalizer. With a fixed reference representative $z_0$, the locally compatible representatives are $\widehat F^\times z_0\widehat R^\times$. The intersection with $\widehat F^\times$ is $\widehat{\mathcal O}^\times$, since the order is maximal. Thus the fixed-orientation target is
\[F^\times\backslash\widehat F^\times/\widehat{\mathcal O}^\times
 =\operatorname{Pic}(\mathcal O).\]
It is a torsor independently of reduction. Reduce an actual reference CM point to obtain $y_0$. The injective action is optimal because $\mathcal O$ is maximal, and the residue character gives the chosen orientation. The simultaneous good-model argument proves equivariance of reduction: prime-to-$qN$ ideal quotients multiply the Tate lattices by $\mathfrak a^{-1}$, preserve the $q$-divisible group and level choices, and preserve the tangent orientation. A nonempty equivariant map between these two torsors is a bijection, with multiplicity one.

All choices at primes dividing $N$ are obtained by the prime-power Atkin--Lehner quotients and duals. Their degrees are prime to $q$, so the same characteristic-zero constructions reduce compatibly, including powers of $2$. The opposite residue orientation at $q$ is the Frobenius image of the geometric reference point, or equivalently is obtained from the conjugate CM type. It is not an extra image of the same fixed-type torsor. Frobenius and inverse Frobenius have the same action on the underlying $\mathbf F_{q^2}$-defined unlabelled pair, but this does not identify their labelled embeddings.

Consequently the class-character sum on this fixed component is the reduction, followed by the same additive functional, of the characteristic-zero sum:
\[\sum_{[\mathfrak a]}\chi(\mathfrak a)P_M(\mathfrak a z_0)
 =\lambda_M\!\left(\overline{L\sum_{[\mathfrak a]}
       \chi(\mathfrak a)\pi(x_D^{[\mathfrak a]})}\right).\]
The same place, the same map $\pi$ and the same additive functional are used throughout. Global inversion of Artin indices does not change a quadratic character. No stabilizer division is hidden here: for a fixed underlying pair the embedding stabilizer is $R_z^\times\cap F^\times=\mathcal O^\times=\{\pm1\}$. Its orbit has size $|R_z^\times|/2$, and weighting by $1/|R_z^\times|$ gives $1/2$, also at the extra-automorphism curves. The embeddings $x$ and $-x$ have opposite residue orientations, so are not conjugate by a unit. For the other components, the already proved Schwartz and parametrization factors $w_{p^e}\chi_{h^*}(p)^e$ occur twice and multiply to one. In the actual application $q=\ell_0\equiv3\pmod4$ and $\chi_{h^*}(q)=-1$; the residue weight is $(-1/q)=-1$, the norm character is $-1$, and the functional's Frobenius factor is $+1$. This is the factor $c_0 2^{\omega(N)}$ in the conductor-one identity above.

The actual application has $F=\mathbf Q(\sqrt{Dh^*})$ with fundamental discriminant and maximal order, $q=\ell_0$, every prime dividing $2N$ split, and no extra units. These are exactly the hypotheses used here. No corresponding assertion for nonmaximal orders or arbitrary conductors has been inserted.

\part{The 222-variable family and its predecessors}
\section{Integral evaluation and the effective contact height}\label{v15:height}
\subsection{Conventions and the evaluation estimate}

If $K$ is a number field, $\OK$ denotes its ring of integers.
For a nonzero ideal $\mathfrak a\subset\OK$, write
$\mathrm N\mathfrak a=\#(\OK/\mathfrak a)$.
If $E$ is a projective $\OK$-module with Hermitian metrics at
its complex embeddings, $\Ardeg(E)$ means the arithmetic degree
of its determinant. In rank one, for a nonzero integral section $s$,
\[
 \Ardeg(E)=\log\#(E/\OK s)
                  -\sum_{\sigma:K\hookrightarrow\C}\log\|s\|_\sigma.
\]
The definition is independent of the choice of $s$ by the product
formula. Both embeddings in a complex conjugate pair are included.

Throughout, $\hF$ is the original stable Faltings height.
For a curve $X/K$ with semistable reduction, its Jacobian height is
\begin{equation}\label{v15:height:eq:height-convention}
 \hF(\Jac X)=\frac{\Ardeg H^0(\mathcal X,
                                  \omega_{\mathcal X/\OK})}{[K:\Q]},
\end{equation}
where $\mathcal X$ is a minimal regular semistable model. The metric
on differentials at a complex embedding is
\[
 \langle\omega,\eta\rangle
       =\frac i2\int_{X(\C)}\omega\wedge\overline\eta.
\]
This is the determinant metric corresponding to the original Faltings
normalization; see \cite[Sections~2.1 and~2.3]{v15-height-Javanpeykar} and
\cite[Section~2.1]{v15-height-vonKaenel}. Passing to a finite extension on which
the curve is semistable does not change the stable height.

A \emph{Belyi map} is a finite morphism from a smooth projective
geometrically connected curve to $\mathbb P^1$, unramified outside
$0,1,\infty$. Let
\[
 \pi:X\longrightarrow\mathbb P^1
\]
be such a map, of degree $d$, where $X$ has genus $g\ge1$.
Write $t$ for the coordinate on the target.
Choose a number field $K$ over which $X$ and $\pi$ are defined,
$X$ is semistable, and every point of $\pi^{-1}(2)$ is rational.
Let $\mathcal Y$ be the normalization of $\mathbb P^1_{\OK}$ in
$K(X)$. It is finite and flat over $\mathbb P^1_{\OK}$.
On the affine $t$-chart write its algebra as $B$, a finite projective
$\OK[t]$-algebra of rank $d$, and put
\begin{equation}\label{v15:height:eq:order}
 C=B/(t-2)B\subset K^d,
 \qquad O=\OK^d.
\end{equation}
The inclusion in $K^d$ uses the $d$ distinct points of the split
fibre. Every element of $C$ is integral over $\OK$, so $C\subset O$.
Let $\mathfrak D(C)$ be the discriminant ideal of its trace pairing.

\begin{theorem}\label{v15:height:thm:evaluation}
With this notation,
\begin{equation}\label{v15:height:eq:evaluation-bound}
 \hF(\Jac X)
 \le\frac{\log\mathrm N\mathfrak D(C)}{2[K:\Q]}
       +g\log\frac2{\sqrt\pi}.
\end{equation}
\end{theorem}

The auxiliary field $K$ is used to split a finite fibre and to
compute a stable height. Its degree will cancel from the estimates.
No upper bound for $[K:\Q]$ is assumed in this theorem.

\subsection{Trace-dual lattices and regular differentials}

\begin{lemma}\label{v15:height:lem:trace-lattice}
Let $C\subset O=\OK^d$ be an order with $C\otimes K=K^d$, and
let
\[
 C^\vee=\{x\in K^d:\operatorname{Tr}_{K^d/K}(xC)\subset\OK\}.
\]
For any set $I$ of $g$ coordinate positions, set
$F=\operatorname{pr}_I(C^\vee)\subset K^g$.
Then $F$ is a projective $\OK$-module containing $\OK^g$.
With the standard Euclidean metric on $K^g$ at every embedding,
\[
 \Ardeg(F)\le\tfrac12\log\mathrm N\mathfrak D(C).
\]
\end{lemma}
\begin{proof}
Work locally at a finite prime $\mathfrak p$. The trace pairing
on $K^d$ is the ordinary dot product, and $O$ is self-dual.
Thus $O\subset C^\vee$, and dualizing the inclusion $C\subset O$
shows
\[
 \operatorname{length}(C^\vee/O)
 =\operatorname{length}(O/C)
 =\tfrac12v_{\mathfrak p}(\mathfrak D(C)).
\]
For example, if a basis matrix for $C$ relative to $O$ is $U$,
then the trace Gram matrix is $U^{\mathsf T}U$, which proves the
last equality. These equalities are valid locally even when
$C$ is not globally free.

Projection gives a surjection
$C^\vee/O\twoheadrightarrow F/\OK^g$. The latter quotient
therefore has no larger length at any prime.
A finite torsion-free module over the Dedekind ring $\OK$ is
projective. Finally, the standard basis of $\OK^g$ has determinant
norm one, so
\[
 \Ardeg(F)=\log\#(F/\OK^g)
 \le\tfrac12\log\mathrm N\mathfrak D(C).
\]
\end{proof}

\begin{lemma}\label{v15:height:lem:integral-evaluation}
Let $\mathcal X$ be a minimal regular semistable model of $X$.
The evaluations of a differential in
$E=H^0(\mathcal X,\omega_{\mathcal X/\OK})$ at the points of
$\pi^{-1}(2)$, expressed in the coordinate $t$, lie in $C^\vee$.
Equivalently, there is an integral map
\[
 E\longrightarrow C^\vee,
 \qquad \omega\longmapsto
       \left(\frac{\omega}{dt}(P)\right)_{P\in\pi^{-1}(2)}.
\]
\end{lemma}
\begin{proof}
First, normalization is finite because the base is excellent.
A normal surface is Cohen--Macaulay, and a finite morphism from
such a surface to a regular surface is flat by miracle flatness;
see \cite[Tag~00R4]{v15-height-Stacks}. Consequently $B$ is finite projective
over $\OK[t]$.

Take a common regular model $\mathcal Z$ dominating
$\mathcal X$ and $\mathcal Y$. Pullback under a blowup of regular
arithmetic surfaces takes regular relative differentials to regular
relative differentials. One can also read this directly from the
canonical-sheaf formula for a blowup; see
\cite[Lemma~2.3.1]{v15-height-Javanpeykar}.
Hence a differential in $E$ is regular on $\mathcal Z$.
The birational morphism $\mathcal Z\to\mathcal Y$ is an
isomorphism at every codimension-one point of the normal surface
$\mathcal Y$.

The dualizing module of a Cohen--Macaulay normal surface is
reflexive: it is torsion-free and satisfies Serre's condition
$(S_2)$. Thus sections defined at all codimension-one points extend
across the codimension-two complement; see
\cite[Tags~0AWQ and~0AVT, Lemmas~31.13.12--31.13.14]{v15-height-Stacks}.
The differential therefore defines a section of the relative
dualizing sheaf of $\mathcal Y$.

Finite duality identifies this sheaf on the affine chart with
\[
 \operatorname{Hom}_{\OK[t]}(B,\OK[t])\,dt;
\]
see \cite[Tags~0BUL and~0C0I]{v15-height-Stacks}.
This statement uses the dualizing \emph{module}; it does not require
that the normal surface have an invertible dualizing sheaf.
Because $B$ is finite projective, specializing the dual at $t=2$
gives $\operatorname{Hom}_{\OK}(C,\OK)\,dt$.
On the generic fibre the identification sends a differential
$f\,dt$ to the functional $b\mapsto\operatorname{Tr}(fb)$.
After splitting the fibre, its coordinates are exactly the values
$f(P)=\omega/dt(P)$. This proves the assertion.
\end{proof}

\begin{proof}[Proof of Theorem~\ref{v15:height:thm:evaluation}]
Riemann--Hurwitz gives
\[
 2g-2=-2d+\deg\operatorname{Ram}(\pi)\le d-3,
\]
since each of the three branch fibres contributes at most $d-1$.
In particular $d>2g-2$. A nonzero holomorphic differential cannot
vanish at all $d$ distinct points of $\pi^{-1}(2)$. Therefore the
evaluation map has rank $g$. Choose $g$ coordinate positions $I$
for which projection remains injective. By
Lemmas~\ref{v15:height:lem:trace-lattice} and~\ref{v15:height:lem:integral-evaluation},
we obtain an integral, generically invertible map
\[
 \phi:E\longrightarrow F=\operatorname{pr}_I(C^\vee),
 \qquad\Ardeg(F)\le\tfrac12\log\mathrm N\mathfrak D(C).
\]

Fix a complex embedding of $K$. The disc
$D=\{t:|t-2|<1/2\}$ is simply connected and avoids the branch
locus. Its inverse image is a disjoint union of $d$ analytic
sheets, each mapped isomorphically to $D$. On one such sheet write
$\omega=f(t)\,dt$. The mean-value inequality implies
\[
 |f(2)|^2\le\frac4\pi\int_D|f(t)|^2\,dx\,dy
       \le\frac4\pi\|\omega\|_{L^2}^2.
\]
Hence the norm of each evaluation functional is at most
$2/\sqrt\pi$. In an orthonormal basis of $E\otimes_\sigma\C$,
each row of the evaluation matrix has this same bound on its
Euclidean norm. Hadamard's determinant inequality gives
\[
 \|\det\phi\|_\sigma\le(2/\sqrt\pi)^g.
\]
The determinant is nonzero and integral. Taking its arithmetic
degree yields
\[
 \Ardeg(E)\le\Ardeg(F)
       +\sum_\sigma\log\|\det\phi\|_\sigma
 \le\tfrac12\log\mathrm N\mathfrak D(C)
       +[K:\Q]g\log(2/\sqrt\pi).
\]
Now apply \eqref{v15:height:eq:height-convention}.
\end{proof}

\subsection{The discriminant of the fibre}

Let $S$ be a finite set of rational primes, and put
\[
 R_S=\prod_{p\in S}p,\qquad \omega(S)=\#S.
\]
Assume that the normalization defining the Belyi cover is finite
\'etale over
\begin{equation}\label{v15:height:eq:good-open}
 \operatorname{Spec}\mathcal O_{K_0}
       [S^{-1},t,t^{-1},(t-1)^{-1}],
\end{equation}
where $K_0$ is a field of definition.
This property persists after any finite field extension:
finite \'etale base change is finite \'etale, and a finite \'etale
algebra over a normal base is normal. Thus it agrees with the
normalization of the new generic cover on this open set.
In particular the auxiliary splitting field in Section~1 introduces
no new vertical branch primes outside $S$.

\begin{lemma}\label{v15:height:lem:discriminant}
Under \eqref{v15:height:eq:good-open}, the order $C$ in \eqref{v15:height:eq:order}
satisfies
\[
 \frac{\log\mathrm N\mathfrak D(C)}{[K:\Q]}
 \le d\left(\log2+\log R_S+\omega(S)\log d\right).
\]
\end{lemma}
\begin{proof}
The determinant of the trace pairing of $B/\OK[t]$ defines an
effective discriminant divisor on the base. Formation of the trace
pairing commutes with specialization. Its intersection with the
section $t=2$ is therefore $\log\mathrm N\mathfrak D(C)$.
The coefficient of this divisor at a codimension-one point is
$\sum_j f_jr_j$, where $e_j,f_j,r_j$ are the ramification index,
residue degree, and different exponent of the corresponding
extensions of discrete valuation rings. They satisfy
$\sum_j e_jf_j=d$. These are the usual local
trace-discriminant and different identities; see
\cite[Chapter~III, Sections~3 and~6]{v15-height-Serre} and
\cite[Section~4.2]{v15-height-Javanpeykar}.

At a horizontal branch divisor the residue characteristic is zero.
The different exponent is $r_j=e_j-1$, so the discriminant
coefficient is at most $d$. The section $t=2$ meets $t=0$
with normalized intersection $\log2$, and meets neither $t=1$
nor $t=\infty$. Thus the horizontal contribution is at most
$d\log2$.

Let $\mathfrak p\mid p$ be a prime of $K$, and write
$e_{\mathfrak p}=v_{\mathfrak p}(p)$.
Lenstra's different bound, in the form
\cite[Proposition~4.1.3]{v15-height-Javanpeykar}, gives
\[
 r_j\le e_j-1+e_j e_{\mathfrak p}
                \left\lfloor\frac{\log d}{\log p}\right\rfloor.
\]
This bound also applies to the possibly imperfect residue field
at the generic point of a vertical fibre.
Its discriminant coefficient is consequently at most
\[
 d\left(1+e_{\mathfrak p}
                \left\lfloor\frac{\log d}{\log p}\right\rfloor\right).
\]
Only $p\in S$ can contribute. The intersection of the vertical
fibre over $\mathfrak p$ with the section has value
$f_{\mathfrak p}\log p$. Using
\[
 \sum_{\mathfrak p\mid p}e_{\mathfrak p}f_{\mathfrak p}=[K:\Q],
 \qquad \sum_{\mathfrak p\mid p}f_{\mathfrak p}\le[K:\Q],
\]
and
$\lfloor\log d/\log p\rfloor\log p\le\log d$
proves the estimate.
\end{proof}

\begin{theorem}\label{v15:height:thm:coverheight}
For a Belyi cover satisfying \eqref{v15:height:eq:good-open},
\begin{equation}\label{v15:height:eq:coverheight}
 \hF(\Jac X)
 \le\frac d2\left(\log2+\log R_S+\omega(S)\log d\right)
       +g\log\frac2{\sqrt\pi}.
\end{equation}
\end{theorem}
\begin{proof}
Combine Theorem~\ref{v15:height:thm:evaluation} and
Lemma~\ref{v15:height:lem:discriminant}.
\end{proof}

\subsection{The integral modular cover}

Let $N$ be a positive integer and let
$H\subset(\Z/N\Z)^\times$ contain $-1$. Write
\[
 X_H=X_1(N)/H,\qquad t_H=[(\Z/N\Z)^\times:H],
 \qquad J_H=\Jac X_H.
\]
The action is the diamond action. Assume that $X_H$ has positive
genus, and use the rational map $t=j/1728$ to the projective line.

\begin{lemma}\label{v15:height:lem:modular-model}
The map $X_H\to\mathbb P^1_t$ has degree
\[
 d_H=Nt_H\prod_{p\mid N}(1+p^{-1}).
\]
It is unramified outside $0,1,\infty$ and admits a finite \'etale
model over
\[
 U=\operatorname{Spec}\Z[1/(6N),t,t^{-1},(t-1)^{-1}].
\]
\end{lemma}
\begin{proof}
Over $U$ consider
\begin{equation}\label{v15:height:eq:universal-curve}
 E_t:\quad y^2=x^3-\frac{3t}{t-1}x-\frac{2t}{t-1}.
\end{equation}
Its discriminant and $j$-invariant are
\[
 \Delta(E_t)=\frac{1728t^2}{(t-1)^3},\qquad j(E_t)=1728t.
\]
Thus it is an elliptic curve over $U$. Since $N$ is invertible,
its scheme of points of exact order $N$ is finite \'etale.
The constant group $H$ acts freely on this scheme: if $aP=P$
and $P$ has exact order $N$, then $a=1$ in $(\Z/N\Z)^\times$.
The quotient is finite \'etale. This follows also from the
classification of finite \'etale covers by locally constant finite
sets; freeness is enough, even when a residue characteristic divides
$|H|$.

To identify its generic fibre, fix an elliptic curve with
$j$-invariant different from $0$ and $1728$ over an algebraic
closure. Its automorphism group is $\{1,-1\}$. The geometric
fibre of $X_1(N)$ is the set of primitive torsion points modulo
this automorphism group. Since $-1\in H$, the fibre of $X_H$
is precisely the set of primitive torsion points modulo $H$.
Different isomorphisms to \eqref{v15:height:eq:universal-curve} differ by
$\pm1$ and hence give the same $H$-orbit. In particular quadratic
twisting introduces no extra double cover. This identifies the
finite \'etale model with the normalization of $X_H$ over $U$.

There are $N^2\prod_{p\mid N}(1-p^{-2})$ primitive points.
Division by $|H|=\varphi(N)/t_H$ gives the asserted degree.
Equivalently, the congruence subgroup contains $-I$, so its
$\operatorname{SL}_2(\Z)$ index is already the degree of the
coarse map to $X(1)$; see \cite[Section~4]{v15-height-vonKaenel}.
The only possible branch points of that map are the two elliptic
points and the cusp, which become $0,1,\infty$ in the $t$ coordinate.
\end{proof}

\begin{lemma}\label{v15:height:lem:torsion-different}
Let $V$ be an excellent discrete valuation ring of characteristic zero,
let $E/V$ be an elliptic scheme, and let $n$ be a positive integer.
Write $E[n]=\operatorname{Spec} A$, and let $\widetilde A$ be its
normalization in its generic finite \'etale algebra. Under the generic
trace pairing, the complementary modules satisfy
\[
 A^\vee=n^{-1}A,\qquad
 \widetilde A^\vee\subseteq n^{-1}\widetilde A.
\]
The different of every field component of $\widetilde A/V$ therefore
contains $n$. The same assertion holds for every intermediate field
of such a component.
\end{lemma}
\begin{proof}
Multiplication $[n]:E\to E$ is finite flat and a local complete
intersection morphism. Its derivative in invariant differential bases
is multiplication by $n$. Consequently its relative differential
module has zeroth Fitting ideal $(n)$. For a finite flat local complete
intersection, this K\"ahler different is the different; see
\cite[Tag~0BWB, Lemmas~49.12.3 and~49.12.6]{v15-height-Stacks}.

Base change along the zero section gives $E[n]/V$. The relative
differential module and its zeroth Fitting ideal commute with base
change, while the morphism remains finite flat and a local complete
intersection. Its different is therefore $(n)$. Its relative
dualizing module is invertible. Under the trace pairing this is the
complementary fractional ideal $A^\vee$, whose inverse is the different;
see \cite[Tag~0BW0, Lemmas~49.8.2--49.8.3, and Tag~0BWB,
Lemma~49.12.4]{v15-height-Stacks}. Thus $A^\vee=n^{-1}A$ as fractional ideals,
not merely as abstract modules.

Normalization is finite, and inclusion reverses trace duals:
\[
 \widetilde A^\vee\subseteq A^\vee=n^{-1}A
                              \subseteq n^{-1}\widetilde A.
\]
Hence $n\widetilde A^\vee\subseteq\widetilde A$, which says exactly
that the different contains $n$. Finally, for normal discrete
valuation rings $V\subseteq W\subseteq U$, the different formula
\[
 \mathfrak D_{U/V}=\mathfrak D_{U/W}\mathfrak D_{W/V}U
\]
shows that $v_U(\mathfrak D_{U/V})\le v_U(n)$ implies
$v_W(\mathfrak D_{W/V})\le v_W(n)$, because the different exponent
of $U/W$ is nonnegative. Apply this in each field component and
at each place. These finite-duality and different arguments do not
require a perfect residue field.
\end{proof}

\begin{corollary}\label{v15:height:cor:modularheight}
Let $g_H$ be the genus of $X_H$. Then
\begin{equation}\label{v15:height:eq:modularheight}
 \hF(J_H)\le
 \frac{d_H}{2}\left(\log12+\log N+2\log d_H\right)
                  +g_H\log\frac2{\sqrt\pi}.
\end{equation}
\end{corollary}
\begin{proof}
Use the auxiliary field $K$ from Theorem~\ref{v15:height:thm:evaluation}.
At a prime $\mathfrak p$ of $K$ over $p>3$, localize the $t$-line
at the generic point of the special fibre. The family
\eqref{v15:height:eq:universal-curve} is an elliptic scheme over this discrete
valuation ring: its discriminant $1728t^2/(t-1)^3$ is a unit there.
By Lemma~\ref{v15:height:lem:modular-model}, the generic cover for $X_H$ is
the primitive $N$-torsion cover modulo $H$. Each of its field
components is an intermediate field of a field generated by a single
$N$-torsion point. Lemma~\ref{v15:height:lem:torsion-different} therefore bounds
every different exponent above $\mathfrak p$ by the valuation of $N$.
The coefficient of the vertical discriminant divisor is at most
\[
                    d_H v_{\mathfrak p}(N).
\]
After summing over these primes and dividing by $[K:\Q]$, the
contribution is at most $d_H\log N$.

At $p=2,3$, the estimate in Lemma~\ref{v15:height:lem:discriminant} gives at
most $d_H(\log6+2\log d_H)$. The horizontal intersection at
$t=2$ contributes at most $d_H\log2$. Thus the normalized
logarithmic discriminant of the fibre order is at most
\[
                 d_H(\log12+\log N+2\log d_H).
\]
Apply Theorem~\ref{v15:height:thm:evaluation}. In particular no degree of a
full torsion splitting field enters this estimate.
\end{proof}

\phantomsection\label{q214foundation:cmh-surface}
\subsection{An abelian surface occurring as an isogeny factor}

\begin{lemma}\label{v15:height:lem:surface}
Let $J$ be an abelian variety of dimension $g\ge2$ over a number
field $L'$, and let $A$ be an abelian surface which is an isogeny
factor of $J$ over $L'$. Put $d'=[L':\Q]$ and
\[
 H_J=\max(1,\hF(J)+3g/2).
\]
Then
\begin{equation}\label{v15:height:eq:surface}
 \hF(A)\le H_J+
 8\left(32\log14+\log d'+
                \log\max(1,\log d',H_J)\right).
\end{equation}
\end{lemma}
\begin{proof}
Apply the construction in the published proof of
\cite[Lemma~16.2(1), pp.~97--98]{v15-height-GaudronRemond}
over the fixed field $L'$ to $J$. It gives an $L'$-isogenous
abelian variety $J_0$, a product over $L'$ of simple abelian varieties,
with $\hF(J_0)\le\hF(J)$. Since $A$ is an $L'$-isogeny factor
of $J_0$, uniqueness of simple factors up to $L'$-isogeny, with
multiplicities, selects a subproduct $B$ that is $L'$-isogenous to $A$.
Write $J_0=B\times E$; then $\dim B=2$ and $\dim E=g-2$.
Additivity of stable Faltings height on products and Bost's lower bound
in the original normalization give
\[
 \hF(B)=\hF(J_0)-\hF(E)
 \le\hF(J)+3(g-2)/2\le H_J.
\]
Here the lower bound is
$\hF(E)\ge-\log(\pi\sqrt2)\dim E>-3\dim E/2$
when $\dim E>0$, and $\hF(E)=0$ when $\dim E=0$;
see \cite[\S2.3 and Corollary~8.4]{gr-periods2014}.
The direct supplier for the product construction is the published
2023 proof just cited; that proof in turn uses R\'emond's 2022
height results \cite{v15-height-Remond} and his 2017 results.
This use does not remove those underlying dependencies.

Apply \cite[Theorem~1.9(1)]{v15-height-GaudronRemond} with source $B$
and target $A$, both over $L'$. In dimension two it supplies an
isogeny of degree at most
\[
 \left(14^{32}d'\max(1,\log d',\hF(B))\right)^{16}.
\]
Stable Faltings height increases under an isogeny of degree $D$
by at most $\frac12\log D$; see \cite[\S2.3]{gr-periods2014}.
Replacing $\hF(B)$ by $H_J$ proves
\eqref{v15:height:eq:surface}.
\end{proof}

The field $L'$ in this lemma is the field on which the surface
factor is realized. It need not split a fibre of a modular curve,
and it need not be the auxiliary field $K$ of
Theorem~\ref{v15:height:thm:evaluation}. The evaluation estimate is a stable
height bound valid after arbitrary extension; the degree entering
\eqref{v15:height:eq:surface} is only $[L':\Q]$.

\subsection{The five-contact height estimate}

For $a=u/v\in\Q$ in lowest terms, with $v>0$, put
$h(a)=\log\max(|u|,v)$. Define
\[
 S_*=\{p\text{ prime}:p\le76176\},\quad
 B_*=\prod_{p\in S_*}p,\quad
 \mathcal D=\{0,1,2401,-343/32\}.
\]
For $a\notin\mathcal D$, let $M_5(a)$ be the product of the
primes $p\notin S_*$ for which at least one of
\[
 v_p(v),\quad v_p(u),\quad v_p(u-v),\quad
 v_p(u-2401v),\quad v_p(32u+343v)
\]
is positive and odd. Every prime occurs only once.
The first valuation is the contact at infinity.

We use the following consequence of the internal arithmetic construction in
Sections~\ref{sec:main}.

\begin{lemma}[Modular realization]\label{v15:height:def:input}
For each $a\in\Q\setminus\mathcal D$, put $R=B_*M_5(a)$.
There is an abelian surface $A$ such that
\begin{equation}\label{v15:height:eq:geometric}
 h(a)\le2\max(0,\hF(A))+128.
\end{equation}
Either $\hF(A)<854R$, or there are a normal number field $L'/\Q$
and an even nebentypus $\varepsilon$ of a primitive weight-two
newform, of level $N$, such that
\begin{equation}\label{v15:height:eq:input}
 \begin{gathered}
 [L':\Q]\le2304\varphi(R),\qquad
 N\le210^{300}R^2,\qquad \rad(N)\mid R,\\
 \varepsilon^2|_{G_{L'}}=1,
 \qquad
 A_{L'}\text{ is an isogeny factor of }
       \Jac(X_1(N)/\ker\varepsilon)_{L'}.
 \end{gathered}
\end{equation}
\end{lemma}
\begin{proof}
The geometric comparison and CM alternative are proved in Theorem~\ref{thm:quartic-height19} and its proof. Retain the sharper scalar-character order from (LR2): $m\le2\prod_{p\mid R,\ p\text{ odd}}(p-1)=2\varphi(R)$. The construction (LR3) gives $[L':\Q]\le2m[L:\Q]$ and $[L:\Q]\le576$, hence $[L':\Q]\le2304\varphi(R)$. Equations (LR4)--(LR5) give the all-prime level bound and its support. The determinant and diamond-quotient argument in the proof of Theorem~\ref{thm:quartic-height19} gives the same even character, its square triviality over $L'$, and the indicated isogeny factor over that field. No coefficient-field degree or extra cyclotomic extension is used.
\end{proof}

Here $R$ is squarefree and divisible by $6$.
The degree bound in \eqref{v15:height:eq:input} retains the totient from the
scalar character: its order satisfies
$m\le2\prod_{p\mid R,\ p\text{ odd}}(p-1)=2\varphi(R)$,
and the construction gives $[L':\Q]\le1152m$.
The square of the nebentypus is trivial over $L'$ by the determinant
identity for the compatible representations. The newform quotient
factors through the diamond quotient because every element of
$\ker\varepsilon$ fixes all conjugates of its differentials.
No bound on the degree of the newform's coefficient field is assumed.

\begin{theorem}\label{v15:height:thm:contactheight}
Use the modular realization just proved. Put
\[
 D_0=4608\cdot210^{300},\qquad
 C_3=19D_0B_*^3\log B_*.
\]
Then $C_3<2^{330770}$ and
\begin{equation}\label{v15:height:eq:contact3}
 h(a)\le C_3 M_5(a)^3(1+\log M_5(a))
 \quad(a\in\Q\setminus\mathcal D).
\end{equation}
In particular,
\begin{equation}\label{v15:height:eq:contact-seven-halves}
 h(a)<2^{330789}M_5(a)^{7/2}.
\end{equation}
For every fixed $\alpha>3$ there is an effective constant
$C_\alpha$ such that $h(a)\le C_\alpha M_5(a)^\alpha$.
\end{theorem}
\begin{proof}
Set $M=M_5(a)$ and $R=B_*M$. In the non-CM alternative, let
$H=\ker\varepsilon$ and $t_H=\operatorname{ord}\varepsilon$.
Since $\varepsilon^2$ is trivial on $G_{L'}$, the image of
$\varepsilon^2$ has order dividing $[L':\Q]$. Therefore
$\operatorname{ord}(\varepsilon)\le2[L':\Q]$, and hence
\[
 t_H\le2[L':\Q]\le4608\varphi(R).
\]
The exact index identity in Lemma~\ref{v15:height:lem:modular-model} gives
\begin{equation}\label{v15:height:eq:cubic-degree}
 \begin{split}
 d_H
 &\le4608\cdot210^{300}R^2\varphi(R)
                    \prod_{p\mid R}(1+p^{-1})\\
 &=D_0R^3\prod_{p\mid R}(1-p^{-2})<D_0R^3.
 \end{split}
\end{equation}
Also $2\le g_H\le d_H$, because the Jacobian contains a surface
factor and the cover is a Belyi map.

Write $\ell=\log R$. The exact numerical verification in Section~\ref{v15:numerical-height} gives
\[
 D_0^2<B_*,\qquad144<B_*,\qquad \ell\ge1.
\]
Consequently
\[
 \log N\le\tfrac52\ell,\qquad
 \log d_H\le\tfrac72\ell,\qquad
 \log12\le\tfrac12\ell.
\]
Using $\log(2/\sqrt\pi)<1$ in
Corollary~\ref{v15:height:cor:modularheight}, we obtain
\begin{equation}\label{v15:height:eq:jacobian-numeric}
 \hF(J_H)\le5d_H\ell+d_H\le6d_H\ell.
\end{equation}
Thus the number $H_J$ in Lemma~\ref{v15:height:lem:surface} is at most
\[
 H_*=8D_0R^3\ell.
\]
Let $d'=[L':\Q]$. Since $d'\le2304R$, $2304<e^8$,
$\ell\ge1$, and $\log D_0\le\ell$,
\[
 \log d'<9\ell,\qquad
 \log H_*\le9\ell.
\]
Moreover $H_*>\max(1,\log d')$.
Lemma~\ref{v15:height:lem:surface}, with $\log14<3$, now gives
\[
 \begin{split}
 \hF(A)
 &\le H_*+8(96+18\ell)\\
 &\le H_*+912\ell
 \le9D_0R^3\ell.
 \end{split}
\]
The CM bound $854R$ is smaller than the same expression.
Combining with \eqref{v15:height:eq:geometric}, and using
$D_0R^3\ell>128$, gives
\[
 h(a)\le19D_0R^3\ell.
\]
Finally $\log R=\log B_*+\log M\le
(\log B_*)(1+\log M)$, since $\log B_*>1$ and $M\ge1$.
This proves \eqref{v15:height:eq:contact3}.

The explicit numerical bound uses
\[
 19<2^5,\qquad D_0<2^{2413},\qquad
 B_*<2^{109445},\qquad \log B_*<2^{17}.
\]
Consequently
\[
 C_3<2^{5+2413+3\cdot109445+17}=2^{330770}.
\]
For $u\ge0$, the maximum of $(1+u)e^{-u/2}$ occurs at
$u=1$ and is less than two. Taking $u=\log M$ gives a
constant smaller than $2^{330771}$, and in particular proves
the convenient bound~\eqref{v15:height:eq:contact-seven-halves}.
More generally, for each $\eta>0$,
$(1+u)e^{-\eta u}$ is bounded on $[0,\infty)$ by an explicitly
computable number. Apply this with $\eta=\alpha-3$.
\end{proof}

\subsection{Disjoint sets of contacts}

\begin{corollary}\label{v15:height:cor:amplification}
Fix $r\ge1$ and put
\[
 \Theta_r=\bigcup_{j=0}^{r-1}
 \left(3j+\left\{0,-\tfrac12,-1,
                     \tfrac1{2399},-\tfrac{32}{407}\right\}\right).
\]
There is an explicitly determined finite exceptional set $\Sigma_r$
such that, if $M_r(t)$ is the product of primes outside $\Sigma_r$
at which some $t-\theta$, $\theta\in\Theta_r$, has positive odd
valuation, then
\[
 h(t)\ll_r M_r(t)^{3/r}(1+\log M_r(t))
       \ll_r M_r(t)^{7/(2r)}
       \qquad(t\in\Q\setminus\Theta_r).
\]
All implied constants are effective under the stated modular
realization input.
\end{corollary}
\begin{proof}
Apply Theorem~\ref{v15:height:thm:contactheight} to
$\phi_j(t)=2+1/(t-3j)$ for $0\le j<r$.
Their matrices have determinant $-1$, so
$h(t)\le h(\phi_j(t))+O_r(1)$.
In addition to $S_*$, exclude the primes dividing contact
coefficients and the nonzero determinants of pairs of primitive
contact forms. The $5r$ rational contacts are distinct.
Outside the resulting fixed set, different contact forms cannot
vanish at the same residue point. Thus the radicals $M_j(t)$
for the different blocks have disjoint prime support, and
$\prod_jM_j(t)=M_r(t)$. Removing the exceptional primes from
a block changes its height inequality by at most an effective
fixed factor.

Multiplying the $r$ inequalities gives
\[
 h(t)^r\ll_r M_r(t)^3
                  \prod_{j=0}^{r-1}(1+\log M_j(t)).
\]
Each logarithmic factor is at most $1+\log M_r(t)$.
Taking $r$th roots proves the first bound. Applying
\eqref{v15:height:eq:contact-seven-halves} in each block proves the second.
The primitive contact forms, their pairwise determinants, and the
fixed exceptions can be constructed directly as in
Sections~\ref{v15:amplification}.
\end{proof}

Four translated sets of five contacts therefore give height exponent
$7/8$ with twenty distinct finite contacts. This statement concerns
logarithmic height; replacing it by an estimate for multiplicative
height would be a different assertion.


\subsection{Numerical constants in the contact-height bound}\label{v15:numerical-height}
Here $B_*$ is the squarefree prime product defined above and $D_0=4608\cdot210^{300}$. The integer comparisons used in the preceding estimates are
\[D_0^2<B_*,\qquad144<B_*,\qquad D_0<2^{2413},\qquad B_*<2^{109445}.\]
These are finite comparisons of the displayed integers. They can be checked by multiplying the primes in the definition of $B_*$; primality up to its stated cutoff is decided by trial division up to the square root. The constant estimates follow by substituting these four comparisons into the preceding displayed bounds. No numerical comparison determines the ineffective aggregate-height constant or a defining value of $N$.

\section{Disjoint contacts and height amplification}\label{v15:amplification}
\subsection{The height hypothesis and the conclusion}

For $t=m/n\in\Q$ in lowest terms with $n>0$, put
$h(t)=\log\max(|m|,n)$. Valuations are normalized by $v_p(p)=1$.
Let $S_0$ be the set of primes at most $76176$, and put
\[
 H=2^{110}210^{1800}\left(\prod_{p\in S_0}p\right)^{19}+128.
\]
The arithmetic premise used in this paper is the following five-contact
height bound from Sections~\ref{sec:main}:
for $s=u/v\in\Q\setminus\{0,1,2401,-343/32\}$ in lowest terms,
\begin{equation}\label{v15:amplification:eq:base}
 h(s)\le H M_5(s)^{19},
\end{equation}
where $M_5(s)$ is the product of primes outside $S_0$ for which at
least one of
\[
 v_p(v),\quad v_p(u),\quad v_p(u-v),\quad
 v_p(u-2401v),\quad v_p(32u+343v)
\]
is positive and odd. The contact at infinity is represented by $v$.
All height conclusions below are implications of \eqref{v15:amplification:eq:base}.
No stronger height premise is assumed.

For $r\ge1$, define
\[
 \Theta=\left\{0,-\frac12,-1,\frac1{2399},-\frac{32}{407}\right\},
 \qquad \Theta_r=\bigcup_{j=0}^{r-1}(3j+\Theta).
\]
Introduce the following $5r$ primitive integral linear forms:
\begin{align*}
 L_{j,0}(m,n)&=m-3jn,\\
 L_{j,1}(m,n)&=2m+(1-6j)n,\\
 L_{j,2}(m,n)&=m+(1-3j)n,\\
 L_{j,3}(m,n)&=2399m-(7197j+1)n,\\
 L_{j,4}(m,n)&=407m+(32-1221j)n.
\end{align*}
Their zeros $m/n$ are exactly $\Theta_r$.
Write these forms as $L_i=a_i m+b_i n$, using a single index when
the pair $(j,k)$ is immaterial.

\begin{definition}\label{v15:amplification:def:exceptional}
The finite set $\Sigma_r$ consists of $S_0$ and all prime divisors of
the nonzero integers
\[
 a_i,\qquad a_i b_{i'}-a_{i'}b_i\ (i\ne i'),\qquad a_i+b_i.
\]
Put
\[
 E_r=\prod_{p\in\Sigma_r\setminus S_0}p,\qquad
 C_r=\log\max(3,9r-10),\qquad H_r=C_r+HE_r^{19}.
\]
For $t\in\Q\setminus\Theta_r$, let $M_r(t)$ be the product of
primes $p\notin\Sigma_r$ for which
$v_p(t-\theta)$ is positive and odd for some $\theta\in\Theta_r$.
\end{definition}

Every integer used to define $\Sigma_r$ is nonzero.
Indeed, the five elements of $\Theta$ are distinct and lie in
$[-1,1/2399]$, so their translates by $3j$ are disjoint.
Also $1\notin\Theta_r$: for $j=0$ this is immediate, and for $j\ge1$
the least element is $3j-1\ge2$. Every $a_i$ is positive.
In particular, all quantities in Definition~\ref{v15:amplification:def:exceptional}
are explicitly computable from $r$.

\begin{theorem}\label{v15:amplification:thm:amplification}
Assume \eqref{v15:amplification:eq:base}. For every positive integer $r$ and every
$t\in\Q\setminus\Theta_r$,
\begin{equation}\label{v15:amplification:eq:amplified}
 h(t)\le C_r+HE_r^{19}M_r(t)^{19/r}
       \le H_r M_r(t)^{19/r}.
\end{equation}
Moreover, $M_r(1)=1$.
\end{theorem}

\subsection{Disjoint radicals and the proof}

\begin{lemma}\label{v15:amplification:lem:disjoint}
Let $t=m/n$ be in lowest terms. At a prime $p\notin\Sigma_r$,
at most one of the integers $L_i(m,n)$ is divisible by $p$.
If $p\mid n$, none of them is divisible by $p$.
Consequently $M_r(t)$ is the product of the pairwise coprime
radicals $M_{r,j}(t)$ obtained by using only the five forms with
first index $j$.
\end{lemma}
\begin{proof}
If $p\mid n$, then $m$ and each $a_i$ are units modulo $p$.
Thus $L_i(m,n)$ is a unit. If $p\nmid n$ and two forms vanish
modulo $p$, their determinant vanishes modulo $p$, contrary to
the definition of $\Sigma_r$. Since $a_i$ and $n$ are units in
the latter case,
\[
 v_p\!\left(t+\frac{b_i}{a_i}\right)=v_p(L_i(m,n)).
\]
The assertions about the radicals follow.
\end{proof}

\begin{proof}[Proof of Theorem~\ref{v15:amplification:thm:amplification}]
For $0\le j<r$, put
\[
 \phi_j(t)=2+\frac1{t-3j}
          =\frac{2m+(1-6j)n}{m-3jn}.
\]
The numerator and denominator are relatively prime: the matrix
\[
 A_j=\begin{pmatrix}2&1-6j\\1&-3j\end{pmatrix}
 \quad\hbox{has determinant }-1.
\]
Its inverse is
\[
 A_j^{-1}=\begin{pmatrix}3j&1-6j\\1&-2\end{pmatrix}.
\]
The operator norm for the maximum norm is
$\max(3,9j-1)$ for $j\ge0$, interpreting the formula at $j=0$
as $3$. Hence
\begin{equation}\label{v15:amplification:eq:height-change}
 h(t)\le h(\phi_j(t))+C_r.
\end{equation}
The five projective contact forms for $\phi_j(t)$ are exactly
$L_{j,0},\ldots,L_{j,4}$, up to signs. In particular, the hypothesis
$t\notin\Theta_r$ excludes every forbidden value and pole of
$\phi_j$. Primes removed by enlarging $S_0$ to $\Sigma_r$ occur
at most once in the radical, regardless of their valuations.
Therefore
\[
 M_5(\phi_j(t))\le E_r M_{r,j}(t).
\]
If $h(t)\le C_r$, the result is immediate. Otherwise
\eqref{v15:amplification:eq:base} and \eqref{v15:amplification:eq:height-change} give, for every $j$,
\[
 0<h(t)-C_r\le HE_r^{19}M_{r,j}(t)^{19}.
\]
Multiplying these $r$ inequalities and applying
Lemma~\ref{v15:amplification:lem:disjoint} proves the first inequality in
\eqref{v15:amplification:eq:amplified}. The second follows from $M_r(t)\ge1$.
Finally, every prime dividing a contact value $L_i(1,1)$ belongs
to $\Sigma_r$, so $M_r(1)=1$.
\end{proof}

\begin{remark}
This argument uses the product of disjoint odd-valuation radicals.
It does not replace an odd-valuation radical by an ordinary radical.
At the removed primes, an arbitrary positive valuation contributes
at most one prime factor to the comparison with $M_5$.
\end{remark}

\subsection{A compatible rational normalization}

The following identities describe how the larger contact set can
enter a local valuation construction. They are algebraic identities;
they do not assert that an arbitrary circuit using them has the
projection property of a defining polynomial.

Let $d=5r$ and set $\epsilon=1$ for odd $r$, $\epsilon=2$ for even
$r$. Thus $d+\epsilon$ is even and $1\le\epsilon\le2$.
Put
\[
 F_r(t)=\prod_{\theta\in\Theta_r}(t-\theta).
\]
Fix a positive rational square $a$, take $t=a\lambda$, and write
$\delta=t-1$. When $\delta F_r(t)\ne0$, define
\begin{equation}\label{v15:amplification:eq:normalization}
 R=\frac{F_r(t)}{a^d\delta^{d+\epsilon}},\qquad
 e=\frac{(a\delta)^d}{F_r(t)},\qquad
 q=\frac1{R\delta^4}.
\end{equation}
The corresponding cleared identities are
\begin{equation}\label{v15:amplification:eq:cleared}
 F_r(t)e=(a\delta)^d,\qquad q\delta^{4-\epsilon}=e,
 \qquad e\delta^\epsilon=R^{-1}.
\end{equation}

\begin{proposition}\label{v15:amplification:prop:normalization}
The following assertions hold.
\begin{enumerate}
\item Every solution of the first identity in \eqref{v15:amplification:eq:cleared}
has $F_r(t)\ne0$. If $q=0$ in the first two identities, then
$\delta=0$ and $M_r(t)=1$.
\item If $p\mid M_r(t)$, then $\delta$ is a $p$-adic unit and
$v_p(R)$ is odd. In particular, $v_p(q)=-v_p(R)$.
\item For $\lambda\to+\infty$,
\[
 R\sim a^{-(d+\epsilon)}\lambda^{-\epsilon}>0,
 \qquad q\sim a^{d+\epsilon-4}\lambda^{\epsilon-4}\longrightarrow0.
\]
\item At a fixed prime where $v_p(t)<0$ and is smaller than every
valuation of a nonzero contact, one has
\[
 v_p(q)=d\,v_p(a)-(4-\epsilon)v_p(t).
\]
For fixed $a$, this tends to $+\infty$ as $v_p(\lambda)$ tends
to $-\infty$.
\end{enumerate}
\end{proposition}
\begin{proof}
If $F_r(t)=0$, then $t\in\Theta_r$, whence $\delta\ne0$ and
$(a\delta)^d\ne0$, a contradiction. If $q=0$, the second identity
gives $e=0$, and the first gives $\delta=0$; Theorem~\ref{v15:amplification:thm:amplification}
gives $M_r(t)=1$.

At a consumed prime, exactly one contact difference has positive odd
valuation, all other contact differences are units, and $t-1$ is a
unit. These facts follow from the determinant and $L_i(1,1)$
conditions in $\Sigma_r$. Since $v_p(a)$ is even, the valuation
$v_p(F_r(t))-d\,v_p(a)$ is odd. This proves the second assertion.
The third follows by comparing the leading terms in
\eqref{v15:amplification:eq:normalization}. For the fourth, every contact difference
and $\delta$ has valuation $v_p(t)$. Thus $v_p(e)=d\,v_p(a)$,
and the second identity in \eqref{v15:amplification:eq:cleared} proves the formula.
\end{proof}

\begin{proposition}\label{v15:amplification:prop:prime-pattern}
Suppose $\lambda=u/v$ with $v$ a positive integral square and
$u\in\Z$, and suppose the $d$ numbers
\[
 p_\theta=u-\frac{\theta}{a}v\qquad(\theta\in\Theta_r)
\]
are positive primes. Assume they avoid $\Sigma_r$ and every prime
dividing the numerator or denominator of $a$. If they are distinct,
then $R$ is positive for $\lambda$ sufficiently large and its
reduced denominator is a square. Its primes of odd valuation
are exactly the $p_\theta$, with valuation one. If in addition each
$p_\theta$ is a square in $\Q_2$, $\Q_3$, and $\Q_{17}$, then
so is $R$.
\end{proposition}
\begin{proof}
Direct substitution gives
\[
 R=\frac{\prod_{\theta\in\Theta_r}p_\theta}
          {v^d\delta^{d+\epsilon}}.
\]
The denominator is a rational square. We check that it is a unit at
each $p_\theta$. Such a prime cannot divide $v$: otherwise it also
divides $u$ and hence all the distinct primes $p_{\theta'}$, since
all contact and scale denominators are units there. Thus $a\lambda$
reduces to $\theta$ at this prime, and $\delta$ reduces to
$\theta-1\ne0$. Its denominator is therefore a unit, and the
valuation of $R$ is one. At every other prime the valuation is even.
This proves the assertion about the reduced denominator and the odd
valuations. Positivity follows from $d+\epsilon$ even and the
positive numerator. The final assertion follows because a product
of local squares divided by a nonzero square is a square.
\end{proof}

The prime-pattern hypothesis in Proposition~\ref{v15:amplification:prop:prime-pattern}
is the same finite-pattern requirement used in the arithmetic
construction: the contacts and scales are fixed first, rational
coefficients are cleared by a fixed square denominator, and a
finite set of primes is avoided by Section~\ref{v15:dependencies} and Theorem~\ref{thm:v5-simultaneous-primes}.
Increasing $r$ changes this finite list and the associated constants.
An explicit numerical cutoff must consequently be recalculated.

\subsection{The divisibility exponents required by height arguments}

\begin{lemma}\label{v15:amplification:lem:two-scales}
Let $k\in\mathbb Z_{>0}$, $\alpha>0$, $H_*>0$, and suppose a radical $M$ satisfies
$h(s)\le H_*M(s)^\alpha$ at two nonzero rational values
$s=\lambda$ and $s=T^2\lambda$, where $T\ne0$.
If $N\ne0$ is an integer and
\[
 M(\lambda)^k\mid N,\qquad M(T^2\lambda)^k\mid N,
\]
then
\[
 h(T)\le H_*|N|^{\alpha/k}.
\]
The two radicals need not be relatively prime.
\end{lemma}
\begin{proof}
The elementary height inequalities give
\[
 2h(T)=h(T^2)\le h(T^2\lambda)+h(\lambda).
\]
Each radical is at most $|N|^{1/k}$, proving the result.
\end{proof}

\begin{corollary}\label{v15:amplification:cor:residual}
Assume \eqref{v15:amplification:eq:base}. Suppose $B\ge1$, $A>0$, $J\ge0$,
$T,\lambda\in\Q^\times$, and
\[
 B^2\le A h(T)+J.
\]
Assume both $\lambda$ and $T^2\lambda$ avoid $\Theta_r$.
Let an integer residual $\Delta$ satisfy
\[
 |\Delta|\le5B^4,\qquad
 M_r(\lambda)^k\mid\Delta,\qquad
 M_r(T^2\lambda)^k\mid\Delta.
\]
If $rk>38$, then all solutions with $B$ sufficiently large in terms
of $r,k,A,J,H_r$ satisfy $\Delta=0$.
\end{corollary}
\begin{proof}
If $\Delta\ne0$, Theorem~\ref{v15:amplification:thm:amplification} and
Lemma~\ref{v15:amplification:lem:two-scales} imply
\[
 B^2\le AH_r 5^{19/(rk)}B^{76/(rk)}+J.
\]
Since $76/(rk)<2$, this bounds $B$ above by a constant depending
only on the displayed data.
\end{proof}

\begin{corollary}\label{v15:amplification:cor:input}
Under the same hypotheses, suppose an integer $N\ne0$ instead
satisfies
\[
 |N|\le2\exp(h(\beta))B,\qquad
 M_r(\lambda)^\mu\mid N,\qquad
 M_r(T^2\lambda)^\mu\mid N.
\]
If $2r\mu>19$, then
\[
 \log B\le \frac{19}{2r\mu-19}\,h(\beta)+C,
\]
where $C$ depends only on $r,\mu,A,J,H_r$.
\end{corollary}
\begin{proof}
Put $\gamma=19/(r\mu)<2$. The preceding argument gives
\[
 B^2\le AH_r2^\gamma\exp(\gamma h(\beta))B^\gamma+J.
\]
If $B^2\le2J$, the assertion follows by increasing $C$.
Otherwise absorb $J$ into $B^2/2$, take logarithms, and divide by
$2-\gamma$. The coefficient of $h(\beta)$ becomes
$\gamma/(2-\gamma)=19/(2r\mu-19)$.
\end{proof}

The least positive integers meeting these two strict inequalities are
\[
 k_r=\left\lfloor\frac{38}{r}\right\rfloor+1,
 \qquad \mu_r=\left\lfloor\frac{19}{2r}\right\rfloor+1.
\]
Some useful values are:
\begin{center}
\begin{tabular}{rrrrrrrrrrrr}\toprule
$r$&1&2&3&4&5&6&7&10&13&20&39\\\midrule
$5r$&5&10&15&20&25&30&35&50&65&100&195\\
$k_r$&39&20&13&10&8&7&6&4&3&2&1\\
$\mu_r$&10&5&4&3&2&2&2&1&1&1&1\\\bottomrule
\end{tabular}
\end{center}
These are requirements on divisibility of integer residuals. A formal
approximation formula must separately be proved accurate enough to
produce those divisibilities. Lower exponents can shorten such a
formula, but the enlarged product $F_r$ requires additional arithmetic
operations. Neither this table nor Theorem~\ref{v15:amplification:thm:amplification}
is a coordinate count for a polynomial definition of the integers.

\subsection{Exact finite verification}

The accompanying standard-library program
\nolinkurl{check_contact_height_amplification.py} reconstructs the
linear forms for the values of $r$ in the table, verifies their
primitivity and pairwise distinct zeros, and factors every integer
defining $\Sigma_r$. It records the exceptional primes larger than
$76176$, the exact product $E_r$, and the height-transformation norm.
It also checks the integer threshold inequalities and the normalized
valuation identities. This calculation verifies the finite algebra;
the arithmetic premise \eqref{v15:amplification:eq:base} remains explicit.


\section{Uniform aggregate trace height}\label{v15:aggregate}
\subsection{Functions with one torsion zero and one pole}

Throughout the geometric argument the base field is algebraically
closed of characteristic zero. Let $E$ be an elliptic curve, let
$O$ be its identity, and write $\mathcal R=\End(E)$. This ring is
commutative and has no zero divisors. Let $T\in E$ have finite
order $N>1$. Suppose $f\in k(E)$ satisfies
\[
 \divisor(f)=m[T]-m[O]\qquad(m>0).
\]
Thus $f$ has its only zero at $T$ and its only pole at $O$.
Such a function exists when $m=N$, because $N(T-O)$ is a
principal divisor on an elliptic curve.

\begin{lemma}\label{v15:aggregate:lem:no-internal-cancellation}
Let $P\in E$ have infinite order. Let $n_1,\ldots,n_r$ be
nonzero integers prime to $N$, and let
$\alpha_1,\ldots,\alpha_r\in\mathcal R$. At every point $U$
of $E$, a pole of one function
$f([n_i]P+\alpha_iU)$ cannot coincide with a zero of another.
Consequently, if some $\alpha_i\ne0$, the product of these
functions has a pole and is nonconstant.
\end{lemma}
\begin{proof}
Suppose
\[
 [n_i]P+\alpha_iU=O,\qquad
 [n_j]P+\alpha_jU=T.
\]
The first equation requires $\alpha_i\ne0$; a constant factor
cannot be a zero or a pole because its argument is nontorsion.
Applying the two endomorphisms and subtracting gives
\[
 (n_i\alpha_j-n_j\alpha_i)P=-\alpha_iT.
\]
The right side is torsion. A nonzero endomorphism is an isogeny
and cannot take a nontorsion point to a torsion point. Hence
$n_i\alpha_j=n_j\alpha_i$.

Put $g=\gcd(|n_i|,|n_j|)$, $u=n_i/g$, and $v=n_j/g$.
There are integers $a,b$ with $au+bv=1$. Set
$\gamma=a\alpha_i+b\alpha_j$; the preceding equality gives
$\alpha_i=u\gamma$ and $\alpha_j=v\gamma$. Therefore
$R=[g]P+\gamma U$ satisfies $[u]R=O$ and $[v]R=T$.
The order of $R$ divides $|u|$, which is prime to $N$, whereas
$T$ has order $N>1$. This is impossible. Finally, every nonzero
$\alpha_i$ is surjective, so the corresponding factor has a
pole, and the absence of a zero at that point keeps it a pole
of the product.
\end{proof}

Two nontorsion points $P_1,P_2$ will be called independent over
$\mathcal R$ modulo torsion if
$\alpha P_1+\beta P_2$ torsion implies $\alpha=\beta=0$.

\begin{lemma}\label{v15:aggregate:lem:cross-cancellation}
Suppose $P_1,P_2$ are independent over $\mathcal R$ modulo
torsion. If $n,m$ are nonzero integers and $\alpha,\beta$ are
nonzero endomorphisms, the functions
$f([n]P_1+\alpha U)$ and $f([m]P_2+\beta U)$ have no common
pole.
\end{lemma}
\begin{proof}
A common pole would give $\alpha U=-[n]P_1$ and
$\beta U=-[m]P_2$. Commutativity then gives
$n\beta P_1-m\alpha P_2=O$, contradicting independence.
\end{proof}

\subsection{Excluding abelian translates in an aggregate fiber}

Fix nonzero scalars $c_1,c_2$ and an integer $r\ge1$. On
$E^{2r}$ consider the rational function
\[
 F(Q_1,\ldots,Q_r,Q'_1,\ldots,Q'_r)
 =c_1\prod_{i=1}^r f(Q_i)+c_2\prod_{i=1}^r f(Q'_i).
\]

\begin{theorem}\label{v15:aggregate:thm:no-translate}
Let $P_1,P_2$ be independent over $\mathcal R$ modulo torsion.
Suppose every $n_i,m_i$ is nonzero and prime to $N$. Put
\[
 q=([n_1]P_1,\ldots,[n_r]P_1,
       [m_1]P_2,\ldots,[m_r]P_2).
\]
No positive-dimensional abelian translate through $q$ is
contained in a level set of $F$.
\end{theorem}
\begin{proof}
Suppose $q+B$ is such a translate, with $B$ a nonzero abelian
subvariety of $E^{2r}$. There is a nonzero homomorphism
$\phi:E\to B$. For example, a nonzero coordinate projection
$B\to E$ is surjective; dualizing and composing with an inverse
to a polarization up to multiplication produces such a map.
Write its coordinates as
$(\alpha_1,\ldots,\alpha_r,\beta_1,\ldots,\beta_r)$,
not all zero. The restriction of $F$ to $q+\phi(E)$ is
\[
 c_1\prod_i f([n_i]P_1+\alpha_iU)
 +c_2\prod_i f([m_i]P_2+\beta_iU).
\]
Within either product a pole survives by
Lemma~\ref{v15:aggregate:lem:no-internal-cancellation}. A pole of one product
is not a pole of the other by
Lemma~\ref{v15:aggregate:lem:cross-cancellation}. Since at least one
endomorphism is nonzero, at least one product has a pole, and
it remains a pole of the sum. This contradicts constancy.
\end{proof}

The same proof permits any finite number of blocks based on
points that are pairwise independent over $\mathcal R$ modulo
torsion. The integers in each block need only be prime to the
order of the unique torsion zero of its chosen function.

\subsection{Sums of abscissas without a parity restriction}

Give $E$ a Weierstrass equation and write $x$ for its abscissa.
This function has its only pole at $O$, of order two. Let
$w_1,\ldots,w_r,w'_1,\ldots,w'_r$ be positive rational numbers
and let $c_1,c_2$ be nonzero algebraic numbers. Define
\[
 F_+(Q_1,\ldots,Q_r,Q'_1,\ldots,Q'_r)
   =c_1\sum_i w_i x(Q_i)+c_2\sum_i w'_i x(Q'_i).
\]

\begin{lemma}\label{v15:aggregate:lem:sum-poles}
Let $P\in E$ have infinite order, let all integers $n_i$ be
nonzero, and let $\alpha_i\in\mathcal R$. If some
$\alpha_i\ne0$, the function
\[
 \sum_i w_i x([n_i]P+\alpha_iU)
\]
has a pole and is nonconstant.
\end{lemma}
\begin{proof}
Choose a pole $U_0$ of a nonconstant summand, say the $i$th.
If another summand, the $j$th, has a pole there, then
$[n_i]P+\alpha_iU_0=[n_j]P+\alpha_jU_0=O$.
Applying the endomorphisms and subtracting gives
$(n_i\alpha_j-n_j\alpha_i)P=O$ and hence
$n_i\alpha_j=n_j\alpha_i$.

Let $z$ be a uniformizer at $O$ with
$x=z^{-2}+O(z^{-1})$, and use $z(U-U_0)$ locally at $U_0$.
Let $\lambda_j$ be the nonzero tangent multiplier of
$\alpha_j$ at $O$. The leading coefficient of the $j$th
summand at this pole is $w_j\lambda_j^{-2}$, and the
endomorphism identity gives
$\lambda_j=(n_j/n_i)\lambda_i$. Consequently the coefficient
of $z^{-2}$ in the full sum is
\[
 \lambda_i^{-2}
       \sum_{j:\,[n_j]P+\alpha_jU_0=O}
                    w_j(n_i/n_j)^2.
\]
The sum is a nonzero positive rational number. The pole
therefore cannot cancel.
\end{proof}

\begin{theorem}\label{v15:aggregate:thm:sum-no-translate}
Suppose $P_1,P_2$ are independent over $\mathcal R$ modulo
torsion, and all $n_i,m_i$ are nonzero. No
positive-dimensional abelian translate through
\[
 q=([n_i]P_1)_{i=1}^r\times([m_i]P_2)_{i=1}^r
\]
is contained in a level set of $F_+$.
\end{theorem}
\begin{proof}
Restrict to an elliptic subgroup direction exactly as in
Theorem~\ref{v15:aggregate:thm:no-translate}. Every nonconstant block sum
has a pole by Lemma~\ref{v15:aggregate:lem:sum-poles}. A pole in one block
cannot occur in the other block, because the same argument
as Lemma~\ref{v15:aggregate:lem:cross-cancellation} would give a nontrivial
endomorphism relation between $P_1$ and $P_2$. At least one
block is nonconstant, so their weighted sum is nonconstant.
\end{proof}

This result uses no torsion zero of $x$ and imposes no parity
condition. Its positivity assumption concerns the fixed
rational weights within each block, not the values of the
abscissas.

\begin{theorem}[Several blocks]\label{v15:aggregate:thm:many-blocks}
Let $P_1,\ldots,P_s$ be nontorsion points on $E$ such that
each pair is independent over $\mathcal R$ modulo torsion.
Fix positive integers $r_j$, positive rational weights
$w_{ji}$, and nonzero algebraic scalars $c_j$. On
$E^{\sum_jr_j}$ put
\[
 F_{\mathrm{sum}}((Q_{ji}))
      =\sum_{j=1}^s c_j\sum_{i=1}^{r_j}w_{ji}x(Q_{ji}).
\]
For any tuple $q=([n_{ji}]P_j)$ with all $n_{ji}\ne0$, no
positive-dimensional abelian translate through $q$ is
contained in a fiber of $F_{\mathrm{sum}}$.

There is a product version. For each $j$ suppose
$\divisor(f_j)=m_j[T_j]-m_j[O]$, where $T_j$ has order
$N_j>1$. Replace $F_{\mathrm{sum}}$ by
\[
 F_{\mathrm{product}}((Q_{ji}))
      =\sum_{j=1}^s c_j\prod_{i=1}^{r_j}f_j(Q_{ji}).
\]
The same conclusion holds if every $n_{ji}$ is prime to $N_j$.
\end{theorem}
\begin{proof}
Restrict to a nonzero homomorphism $E\to B$ in any proposed
positive-dimensional subgroup $B$. Some block has a
nonconstant coordinate. Within that block a pole survives
by Lemma~\ref{v15:aggregate:lem:sum-poles}, or by
Lemma~\ref{v15:aggregate:lem:no-internal-cancellation} in the product case.
A pole cannot occur in two distinct blocks: their two pole
equations would contradict the pairwise independence of
their base points. The pole therefore survives in the full
sum, contradicting constancy.
\end{proof}

The scalars $c_j$ need not be real or positive. Pairwise
independence suffices; no independence assumption on the
entire collection is required.

\subsection{A linear height comparison}

We use the following consequence of Von Buhren's family height
theorem \cite[Theorem~1.1]{v15-aggregate-VonBuhren}. Let $A$ be an abelian
variety, $S$ a projective variety, and $V\subseteq A\times S$
an integral closed subvariety. For a fixed subgroup
$\Gamma\subseteq A(\overline\Q)$ of finite rank, there are
constants $\mu,\nu\ge0$ such that
\begin{equation}\label{v15:aggregate:eq:VB}
 \widehat h_A(a)\le\mu h_S(s)+\nu
\end{equation}
whenever $a\in\Gamma$ lies in the fiber $V_s$ and does not lie
on a positive-dimensional abelian translate contained in that
fiber. The slope $\mu$ can be computed from projective
geometric data. The theorem does not provide an effective
value for $\nu$.

\begin{corollary}\label{v15:aggregate:cor:height}
Assume the data of Theorem~\ref{v15:aggregate:thm:no-translate} are algebraic
numbers. There are constants $M>0$ and $M_0\ge0$, independent
of all the integers $n_i,m_i$, such that
\[
 \sum_{i=1}^r\bigl(\widehat h_E([n_i]P_1)
                 +\widehat h_E([m_i]P_2)\bigr)
 \le M h(F(q))+M_0
\]
for every tuple in that theorem. Here $h$ is the absolute
logarithmic height on $\Pj^1$. One may choose an effective
geometric slope $M$ after fixing compatible height
normalizations. No effective $M_0$ is asserted.
The same conclusion holds with $F_+$ and the hypotheses of
Theorem~\ref{v15:aggregate:thm:sum-no-translate}, allowing all nonzero labels.
\end{corollary}
\begin{proof}
Let $A=E^{2r}$ and take for $V\subseteq A\times\Pj^1$ the
closure of the graph of $F$. This is integral and projective.
Let $\Gamma$ be the subgroup generated by $P_1$ in each of
the first $r$ coordinates and by $P_2$ in each of the last
$r$ coordinates. It is finitely generated.

Every indicated $q$ is in the domain of $F$, since no
$[n_i]P_1$ or $[m_i]P_2$ is the identity. If an abelian
translate through $q$ were contained in the projective fiber
$V_{F(q)}$, it would meet the domain of $F$ at $q$. Its dense
open intersection with that domain would make $F$ constant
on the translate. Theorem~\ref{v15:aggregate:thm:no-translate} rules this
out. This argument also handles possible boundary components
introduced by closing the graph; none through $q$ changes
the conclusion.

Apply~\eqref{v15:aggregate:eq:VB}, using a symmetric ample product line
bundle on $A$. Its canonical height is a fixed positive
multiple of the sum of the coordinate canonical heights.
Rescaling gives the stated inequality. For $F_+$ the same
proof uses Theorem~\ref{v15:aggregate:thm:sum-no-translate} to exclude the
exceptional translates.
\end{proof}

For either function of Theorem~\ref{v15:aggregate:thm:many-blocks} the same
argument gives a uniform inequality
\[
 \sum_{j,i}\widehat h_E([n_{ji}]P_j)
       \le M h(F(q))+M_0.
\]
The subgroup is generated by the finitely many base points
placed in their respective coordinates. The constants are
fixed before any of the integer labels are chosen.

\subsection{The quadratic-field elliptic curve}

Let
\[
 K=\Q(\sqrt2),\quad d=3-\sqrt2,\quad d'=3+\sqrt2,
 \quad E:\ y^2=x^3-d^2x,
\]
and let $\sigma$ be the nonidentity automorphism of $K/\Q$.
The point used in the rational source is
\[
 R_0=(1+\sqrt2,2\sqrt2),\qquad
 P=2R_0=\left(\frac{57}{8}-\frac72\sqrt2,
             -\frac{165}{8}+\frac{417}{32}\sqrt2\right).
\]
For this section its infinite order is the only group
property needed; no assertion about the full Mordell--Weil
group is required.

\begin{lemma}
The displayed point $P$ has infinite order.
\end{lemma}
\begin{proof}
Substitution verifies $R_0\in E(K)$, and the usual tangent
formula gives the displayed double. At the prime
$\mathfrak p=(\sqrt2)$ normalize $v_{\mathfrak p}(\sqrt2)=1$,
so $v_{\mathfrak p}(2)=2$. The displayed coordinates give
\[
 v_{\mathfrak p}(x(P))=-6,\qquad
 v_{\mathfrak p}(y(P))=-9,\qquad
 v_{\mathfrak p}(-x(P)/y(P))=3.
\]
The parameter $z=-x/y$ places $P$ in the formal group of
this integral Weierstrass equation. Its multiplication series
have integral coefficients and linear term $[m](z)=mz+O(z^2)$.
For odd $m$ and $v(z)\ge3$, the linear term uniquely has
least valuation, so $v([m]z)=v(z)$. For multiplication by
two, the linear term has valuation $v(z)+2<2v(z)$,
and therefore $v([2]z)=v(z)+2$. No successive multiplication
by an odd integer and by a power of two can give zero.
Every nonzero integer has this form, proving the claim.
\end{proof}

\begin{lemma}\label{v15:aggregate:lem:independent}
Over a field containing $\sqrt d$ and $\sqrt{d'}$, identify
$E$ and $E^\sigma$ with $E_0:Y^2=X^3-X$ by
\[
 (X,Y)=\left(\frac{x}{d},\frac{y}{d\sqrt d}\right),
 \qquad
 (X',Y')=\left(\frac{x'}{d'},\frac{y'}{d'\sqrt{d'}}\right).
\]
The images $P_1,P_2$ of $P,P^\sigma$ are independent over
$\End(E_0)=\Z[i]$ modulo torsion.
\end{lemma}
\begin{proof}
The squareclasses of $d,d'$ are independent in $K(i)$.
Indeed, a positive element of the real field $K$ that is a
square in $K(i)$ is already a square in $K$: writing a square
root as $a+ib$ gives $ab=0$, and the purely imaginary case
would be negative at a real embedding. Neither $d$ nor $d'$
is a square in $K$, since its norm to $\Q$ is $7$. Their
product $7$ is not a square in $K$: expanding
$(a+b\sqrt2)^2=7$ for rational $a,b$ gives $ab=0$ and then
requires either $7$ or $7/2$ to be a rational square.

It follows that $K(i,\sqrt d,\sqrt{d'})/K(i)$ has an
automorphism fixing $\sqrt d$ and negating $\sqrt{d'}$.
It fixes $P_1$, negates $P_2$, and fixes every endomorphism
in $\Z[i]$. If $\alpha P_1+\beta P_2$ were torsion, applying
the automorphism and subtracting would make $2\beta P_2$
torsion. Infinite order forces $\beta=0$, and then
$\alpha=0$.
\end{proof}

\begin{theorem}\label{v15:aggregate:thm:aggregate}
For any fixed $r\ge1$, there are constants $M_+>0$ and
$M_{+,0}\ge0$ such that, for all nonzero integers
$n_1,\ldots,n_r$,
\begin{equation}\label{v15:aggregate:eq:sum-height}
 \sum_{i=1}^r\widehat h_E([n_i]P)
 \le M_+ h\!\left(\frac12\Tr_{K/\Q}
                 \sum_{i=1}^r x([n_i]P)\right)+M_{+,0}.
\end{equation}
There are also constants $M_\times>0$, $M_{\times,0}\ge0$
such that, when all the labels are odd,
\begin{equation}\label{v15:aggregate:eq:aggregate-height}
 \sum_{i=1}^r\widehat h_E([n_i]P)
 \le M_\times h\!\left(\frac12\Tr_{K/\Q}
                 \prod_{i=1}^r x([n_i]P)\right)+M_{\times,0}.
\end{equation}
Both slopes can be chosen effectively from the fixed geometry;
the additive constants are not determined effectively here.
\end{theorem}
\begin{proof}
On $E_0$, the coordinate $X$ has divisor
$2[(0,0)]-2[O]$. Its zero is a nonzero point of order two,
so the index condition in Theorem~\ref{v15:aggregate:thm:no-translate} is
exactly oddness. Under the isomorphisms of
Lemma~\ref{v15:aggregate:lem:independent}, the rational function in the
height on the right of~\eqref{v15:aggregate:eq:aggregate-height} is
\[
 F=\frac{d^r}{2}\prod_{i=1}^r X_i
       +\frac{(d')^r}{2}\prod_{i=1}^r X'_i.
\]
For the sum, the corresponding function is
$F_+=(d/2)\sum_iX_i+(d'/2)\sum_iX'_i$ and
Theorem~\ref{v15:aggregate:thm:sum-no-translate} applies without parity
restrictions. Apply Corollary~\ref{v15:aggregate:cor:height} to both
functions with $m_i=n_i$.
Canonical heights are unchanged by these pointed
isomorphisms, and conjugation preserves absolute canonical
height. The coordinate height sum there is twice the sum
in~\eqref{v15:aggregate:eq:aggregate-height}; divide by two.
\end{proof}

Notice that the inequality is linear. Finiteness of each
individual fiber alone would not give this conclusion. The
additive constant is the remaining effectiveness issue; it is
not an additional mathematical conjecture.

\subsection{A trace-height consequence over number fields}

\begin{corollary}\label{v15:aggregate:cor:number-fields}
Let $K/F$ be a finite extension of number fields, let $E/K$
be an elliptic curve, and let $P\in E(K)$ have infinite order.
For every embedding $\sigma:K\hookrightarrow\overline F$
over $F$, suppose there is a pointed isomorphism
$\phi_\sigma:E^\sigma\to E_0$ to a fixed elliptic curve over
$\overline F$. Assume that the points
$P_\sigma=\phi_\sigma(P^\sigma)$ are pairwise independent
over $\End(E_0)$ modulo torsion.

For every fixed positive integer $r$ there are constants
$M>0$, $M_0\ge0$ such that, for all nonzero integer labels,
\[
 \sum_{i=1}^r\widehat h_E([n_i]P)
 \le M h\!\left(\Tr_{K/F}
                    \sum_{i=1}^r x([n_i]P)\right)+M_0.
\]
The slope can be chosen effectively from the fixed geometric
data; the additive constant has the same ineffective status
as in Corollary~\ref{v15:aggregate:cor:height}.
\end{corollary}
\begin{proof}
Under a pointed isomorphism of Weierstrass models the
abscissa transforms as
$x^\sigma=a_\sigma x_0+b_\sigma$, with $a_\sigma\ne0$.
The function in the height is therefore
\[
 \sum_\sigma a_\sigma\sum_i x_0([n_i]P_\sigma)
                  +r\sum_\sigma b_\sigma.
\]
Theorem~\ref{v15:aggregate:thm:many-blocks} excludes constant translates
through the indicated points for the nonconstant part.
Adding the fixed scalar does not change that conclusion.
The graph-family argument gives a linear bound for the sum
of the canonical heights in all the embedding blocks.
Absolute canonical height is unchanged by conjugation and
pointed isomorphism, so this sum is $[K:F]$ times the left
side of the displayed inequality. Divide by $[K:F]$.
\end{proof}

This corollary is a height comparison under explicit
independence hypotheses. It does not assert a Diophantine
definition of the ring of integers of $K$.


\section{Plain formal parameters and rational approximations}\label{v15:logs}
\subsection{Uniform logarithm truncations of arbitrary order}
\label{v15:logs:sec:uniform-logarithm}

Let $E:y^2=x^3-Dx$ over a field of characteristic zero, with $D\ne0$.
Write $f=1/x$ and $z=-x/y$ near the identity. The equation gives
\[
                         f=z^2(1-Df^2).
\]
Define rational numbers
\[
 c_0=1,\qquad c_1=-\frac25,\qquad
 c_j=c_{j-1}\frac{4j-5}{4j+1}\quad(j\ge2),
\]
and put
\[
 Q_J(T)=\sum_{j=0}^{J-1}c_jT^j,
 \qquad L_J(Q)=z(Q)Q_J(Df(Q)^2)\quad(J\ge1).
\]

\begin{lemma}\label{v15:logs:lem:logarithm-series}
The formal series
\[
                  L(Q)=z(Q)\sum_{j\ge0}c_j(Df(Q)^2)^j
\]
is the normalized formal logarithm of the displayed equation.
For every odd prime $p$ and $j\ge1$,
\[
                  v_p(c_j)\ge-\lfloor\log_p(4j+1)\rfloor.
\]
\end{lemma}
\begin{proof}
The invariant differential is
$\omega=dx/(2y)=z\,df/(2f)$. Put $T=Df^2$ and $S(T)=\sum c_jT^j$.
From $z^2=f/(1-T)$, differentiation shows that $d(zS(T))=\omega$
is equivalent to
\[
                 4T(1-T)S'(T)+(1+T)S(T)=1-T.
\]
The constant coefficient is $c_0=1$, the linear coefficient is
$5c_1=-2$, and every subsequent coefficient is precisely the stated
recurrence. The resulting series is $z+O(z^5)$ and hence is the
normalized formal logarithm.

For $j\ge1$ its coefficient is exactly
\[
                 c_j=2\prod_{i=1}^j\frac{4i-5}{4i+1}.
\]
For an odd prime power $p^r$, each of the two arithmetic progressions
in this quotient has length $j$ and invertible step four. The counts
of terms divisible by $p^r$ therefore differ by at most one.
No term has absolute value exceeding $4j+1$. Summing the possible
loss of one valuation over the prime powers at most $4j+1$ proves
the claimed lower bound.
\end{proof}

\begin{theorem}\label{v15:logs:thm:uniform-logarithm-tail}
Let $F/\Q_p$ be unramified, where $p$ is odd, normalize $v(p)=1$,
and suppose $D$ is integral. For any nonzero point $Q$ in the formal
group, so that $v(z(Q))\ge1$, and every integer $J\ge1$,
\[
 \begin{aligned}
 v(L_J(Q))&=v(z(Q)),\\
 v(L(Q)-L_J(Q))&\ge v(z(Q))+3J.
 \end{aligned}
\]
In particular, the conclusion is uniform in $J$ and does not require
$p>J$ or $p>4J$.
\end{theorem}
\begin{proof}
The relation between $f$ and $z$ gives $v(f(Q))=2v(z(Q))$.
For every $j\ge1$ and odd $p$ one has
\[
                 \lfloor\log_p(4j+1)\rfloor\le j.
\]
Indeed $3^{j+1}>4j+1$ for $j\ge1$, by induction. Consequently
Lemma~\ref{v15:logs:lem:logarithm-series} gives
\[
 \begin{split}
 v\bigl(z(Q)c_j(Df(Q)^2)^j\bigr)
 &\ge v(z(Q))+4jv(z(Q))-\lfloor\log_p(4j+1)\rfloor\\
 &\ge v(z(Q))+3j.
 \end{split}
\]
All nonconstant terms are strictly smaller in absolute value than
the leading term $z(Q)$, proving the first equality. The same estimates
show that the infinite series converges, and that its tail beginning
at $j=J$ has the stated valuation.
\end{proof}

\begin{corollary}\label{v15:logs:cor:uniform-logarithm-ratio}
In the setting of Theorem~\ref{v15:logs:thm:uniform-logarithm-tail}, suppose
$Q=[n]P$ and $R=[m]P$ are nonzero formal points, where $m,n$ are
nonzero integers. For truncation lengths $J_1,J_2\ge1$, put
$k=3\min(J_1,J_2)$. Then
\[
 v\left(\frac{L_{J_1}(Q)}{L_{J_2}(R)}-\frac nm\right)
                  \ge v(n/m)+k.
\]
If $n/m$ is integral, the right side is at least $k$.
\end{corollary}
\begin{proof}
Formal additivity applied to $[m]Q=[n]R$ gives $L(Q)/L(R)=n/m$.
The leading-term estimate in the theorem makes both logarithms
nonzero. Write $L_{J_1}(Q)=L(Q)(1+\varepsilon_Q)$ and
$L_{J_2}(R)=L(R)(1+\varepsilon_R)$; their two errors have valuations
at least $k$. Subtracting the exact ratio proves the assertion.
\end{proof}

For the field $\Q(\sqrt2)$, every odd completion is unramified and
$D=11-6\sqrt2$ is integral. Thus these statements apply at every
contact prime $p>76176$, independently of the chosen truncation
length. If $M$ is a power of two with $M\ge2$, the choice $J=M/2$
has relative precision at least $3M/2$, which is stronger than
precision $M$. Arbitrarily long truncations therefore leave the fixed
contact-height constant and its exceptional-prime set unchanged.

\subsection{Nonvanishing at the real embeddings}

\begin{lemma}\label{v15:logs:lem:positive-logarithm-polynomial}
For every integer $J\ge1$, one has $Q_J(T)>0$ for $0\le T\le1$.
\end{lemma}
\begin{proof}
All $c_j$ for $j\ge1$ are negative. For $r\ge1$, the recurrence
telescopes to
\[
                    \sum_{j=0}^{r}c_j=-\frac{4r-1}{2}c_r>0.
\]
The identity holds at $r=1$ and its successive differences are $c_r$.
Because $0\le T^j\le1$ and $c_j<0$, it follows that
$Q_J(T)\ge Q_J(1)>0$; $J=1$ is immediate.
\end{proof}

\begin{corollary}\label{v15:logs:cor:nonzero-logarithm-denominator}
For $K=\Q(\sqrt2)$ and $D=(3-\sqrt2)^2$, let $Q\in2E(K)$
be nonzero. Then $L_J(Q)\ne0$ for every $J\ge1$.
\end{corollary}
\begin{proof}
The duplication identity
\[
                x(2R)=\left(\frac{x(R)^2+D}{2y(R)}\right)^2
\]
shows that the abscissa of a nonzero double is positive at both real
embeddings. Its ordinate cannot be zero: the two nonzero
ordinate-zero abscissas are $\pm(3-\sqrt2)$, neither a square in $K$
since their norm is seven, while abscissa zero is excluded by the
duplication formula. At either embedding the curve equation now
implies $x(Q)>\sqrt D$. Thus $0<Df(Q)^2<1$ and $z(Q)\ne0$.
Lemma~\ref{v15:logs:lem:positive-logarithm-polynomial} proves the assertion.
\end{proof}

Consequently any rational system whose logarithm denominators are
evaluated at nonzero doubles may use these truncations at arbitrary
finite length. Increasing the length introduces no new
nonvanishing condition at those points.

\subsection{A finite integer coefficient specification}

For a polynomial $G$, write $\|G\|_1$ for the sum of the absolute
values of its coefficients. In a finite free algebra over $\mathbb Q$
or an integer polynomial ring, choose a basis and let $M_A$ be the
matrix of multiplication by $A$. Use the maximum column sum of the
coefficient $\ell^1$ norms as its matrix norm. This norm is
submultiplicative, has $\|I\|=1$, and satisfies $\|cI\|=|c|$ for
rational scalars.

Compatible weights mean nonnegative weights on base variables and
positive weights on the algebra generators for which multiplication
and reduction by the monic defining relations do not increase the
sum of the input weights. The weighted degree of the zero element is
$-\infty$; all other weighted degrees are nonnegative.

Define
\[
 D_J=\prod_{i=1}^{J-1}(4i+1),\qquad
 G_J(T)=D_J Q_J(T)\in\Z[T].
\]
Empty products are one. Its coefficients are explicitly
\[
 [T^0]G_J=D_J,\qquad
 [T^j]G_J=
 2\left(\prod_{i=1}^{j}(4i-5)\right)
  \left(\prod_{i=j+1}^{J-1}(4i+1)\right)
 \quad(1\le j<J).
\]
This finite product-and-sum definition specifies the ordinary integer
polynomial without requiring a list of all its coefficients.

\begin{proposition}\label{v15:logs:prop:indexed-logarithm-coefficients}
For $J\ge2$, the degree of $G_J$ is $J-1$ and
\[
 D_J\le(4J)^{J-1},\qquad
              \|G_J\|_1\le2J(4J)^{J-1}.
\]
If $J=2^{b-1}$, with $b\ge2$, these imply
\[
                    D_J<2^{(b+1)J},\qquad
                    \|G_J\|_1<2^{(b+1)J}.
\]
If $A$ is an element of a finite free algebra with
$\|M_A\|\le2^B$, $B\ge0$, then
\[
                  \|M_{G_J(A)}\|
                  \le\|G_J\|_1\,2^{(J-1)B}.
\]
If the algebra has the multiplicative weighted filtration associated
with compatible positive generator weights, then
$\operatorname{wdeg}(G_J(A))\le(J-1)\max(0,\operatorname{wdeg}(A))$.
\end{proposition}
\begin{proof}
Every nonconstant coefficient is a product of $J-1$ integers of
absolute value at most $4J$, times two. The same bound without two
holds for the constant coefficient. Summing the coefficient bounds
gives the displayed estimate. Its binary exponent when
$J=2^{b-1}$ is
\[
       1+(b-1)+(J-1)(b+1)=(b+1)J-1.
\]
The operator estimate follows from the triangle inequality and
submultiplicativity, and the weight estimate follows term by term.
\end{proof}

For a rational argument $A/B$, its homogeneous numerator is the
explicit polynomial
\[
 G_J^{\rm hom}(A,B)=\sum_{j=0}^{J-1}[T^j]G_J\,A^jB^{J-1-j}.
\]
Thus $Q_J(A/B)=G_J^{\rm hom}(A,B)/(D_JB^{J-1})$ on $B\ne0$.
This form is convenient for a finite algebra over an integer polynomial
ring. Its degree and operator-norm estimates follow from the same
proof with the larger of the two argument bounds.


\section{Prime selection and simultaneous denominator preparation}\label{v15:selection}
\subsection{The precise analytic input}
Fix distinct integers $c_1,\ldots,c_k$, with $k\ge2$, and put
\[
 p_j=u+c_jw^2.
\]
The fixed congruences, Gaussian ray conditions and reduction-order
avoidance conditions in Theorem~\ref{thm:v5-simultaneous-primes} are imposed throughout.
A fixed dilation of $w$ can be absorbed into the integers $c_j$.
All such data are chosen before the parameter $X$ tends to infinity.

We derive the following quantitative statement internally. Proposition~\ref{prop:v5-artin-patterns} supplies the fixed Gaussian ray-class count. Proposition~\ref{prop:v5-cutoff-lower} bounds the finite-cutoff singular series below by a positive multiple of $(\log L)^{-4k}$; Proposition~\ref{prop:v5-obstruction-tail} bounds the full forbidden common-divisor tail by $C\tau X^{3/2}/(L\log L)+O_\tau(X^{3/2-1/800})$. Choose $L$ first so that this tail is smaller than the positive cutoff mass, then fix the field and ray data, then choose the diagonal thresholds as in Section~\ref{subsec:v5-diagonal}. Each surviving weight is a product of $k$ prime logarithms. This proves the following statement for the fixed compatible Gaussian data used here; no uniformity in growing fields or moduli and no general-Galois preprint estimate is invoked. There are constants
$c,C>0$ and a sufficiently slowly decreasing function $\tau(X)>0$,
with $\tau(X)\to0$, such that, for
\[
 T=\lfloor\tau(X)\sqrt X\rfloor,
\]
the admissible pairs in $X/2\le u\le X$, $1\le w\le T$ carry
nonnegative weights $\omega_X(u,w)$ satisfying
\begin{equation}\label{v15:selection:eq:count}
 \sum_{u,w}\omega_X(u,w)\ge cXT,
 \qquad 0\le\omega_X(u,w)\le C(\log X)^k.
\end{equation}
The weights vanish outside admissible pairs. The lower bound is after
the removal of all self-order and cross-order obstructions, not merely
after a finite list of congruences.

For clarity, the order of limits in the cited argument is essential.
Its finite cutoff, denoted by $L$ there, is fixed first. The surviving
weight is at least
\[
 \frac{c_1XT}{8(\log L)^{4k}}
\]
for a positive constant $c_1$ and all sufficiently large $X$.
This gives the first inequality in \eqref{v15:selection:eq:count} with a fixed
positive constant. Each surviving weight is a product of $k$ prime
logarithms, up to fixed normalizations, giving the second inequality.
The diagonal construction of $\tau$ allows additional lower thresholds
which hold eventually for every fixed positive value of $\tau$.

All assertions below that select prime patterns depend on these cited
estimates. The elementary statements about weights, integers and
valuations are proved here.

\subsection{Removing a thin strip}
\begin{lemma}\label{v15:selection:lem:strip}
Assume \eqref{v15:selection:eq:count}. Let $0<\alpha<1/2$ be fixed.
The function $\tau$ may be chosen sufficiently slowly that admissible
pairs with $w\ge X^\alpha$ still have total weight at least $cXT/2$
for every sufficiently large $X$.
\end{lemma}
\begin{proof}
The strip $w<X^\alpha$ contains at most $X^{1+\alpha}+X$
integer pairs. Its weight is at most
\[
 C(X^{1+\alpha}+X)(\log X)^k.
\]
Choose the diagonal thresholds so that
\[
 \frac{\tau(X)X^{1/2-\alpha}}{(\log X)^k}\longrightarrow\infty.
\]
This is compatible with every prior threshold: for a fixed positive
value $\tau_j$, the displayed quotient tends to infinity, so the start
of the $j$th diagonal interval can be increased to meet this additional
condition. Also require $\tau(X)\sqrt X\ge2$, so
$T\ge\tau(X)\sqrt X/2$. The strip weight divided by $XT$ tends
to zero. Subtracting it from \eqref{v15:selection:eq:count} proves the assertion.
\end{proof}

\begin{theorem}\label{v15:selection:thm:separation}
For each selected prime $p_j$, let $d_j$ be a positive integer such that
\[
 d_j=p_j+O(\sqrt{p_j}),\qquad 4\mid d_j,
\]
where the implied constant is uniform in $j$ and the pattern.
Under the cited analytic input, the pattern can be chosen so that
\[
                       p_i\nmid d_j\qquad(1\le i,j\le k).
\]
The fixed congruence and avoidance conditions remain satisfied, all
selected primes tend to infinity, and $u/w^2$ tends to infinity.
\end{theorem}
\begin{proof}
Apply Lemma~\ref{v15:selection:lem:strip} with $1/4<\alpha<1/2$.
Since $w^2\le T^2=o(X)$, every $p_j/u$ tends to one, uniformly
over the selected patterns. The same is true of $d_j/u$.
Consequently $1/2<d_j/p_i<2$ for sufficiently large $X$.
Divisibility would therefore imply $d_j=p_i$.

For $i\ne j$, let $c_*:=\min_{i\ne j}|c_i-c_j|>0$. Then
\[
 |p_i-p_j|\ge c_*w^2\ge c_*X^{2\alpha},
\]
whereas $|d_j-p_j|=O(\sqrt X)$. Because $2\alpha>1/2$,
equality is impossible for sufficiently large $X$.
For $i=j$, the divisibility $4\mid d_i$ excludes $d_i=p_i$,
since $p_i$ is an odd prime. The other conditions follow from the
support of the weights and from $u\ge X/2$, $w^2=o(X)$.
\end{proof}

For the elliptic curves used in the defining-polynomial constructions,
Hasse's bound supplies the displayed estimate for $d_j$, and the fixed
Gaussian congruences make the opposite order divisible by four.
Thus Theorem~\ref{v15:selection:thm:separation} gives precisely the separation used
to impose additional formal depth on the primary points.
The application of this observation to those constructions is already
present in Sections~\ref{sec:main}; the
argument above states its quantitative source and permitted limits.

\subsection{Simultaneous denominator preparation}
Let $S$ be a finite set of rational primes, and let
$a_1,\ldots,a_s\in\Q^\times$ be fixed.

\begin{lemma}\label{v15:selection:lem:denominator}
There is a positive integer $D$ such that, whenever $u$ is a prime
outside $S$ and $w$ is a positive integer, the numbers
\[
 \lambda=\frac{u}{D^2w^2},\qquad t_i=a_i\lambda,
 \qquad\delta_i=t_i-1
\]
satisfy $v_p(\delta_i)=v_p(t_i)<0$ for every $p\in S$ and every $i$.
The integer $D$ can also satisfy any prescribed finite set of
divisibility conditions.
\end{lemma}
\begin{proof}
Choose $D$ with $2v_p(D)>\max_i v_p(a_i)$ for every $p\in S$,
and enlarge it to satisfy the other fixed divisibilities. Since
$v_p(u)=0$ and $v_p(w)\ge0$,
\[
 v_p(t_i)=v_p(a_i)-2v_p(D)-2v_p(w)<0.
\]
The two terms of $t_i-1$ have distinct valuations, so its valuation
equals $v_p(t_i)$. All requirements concern a finite set of primes.
\end{proof}

This choice can be made after fixing the scales and the finite
preparation set, but before selecting primes. Multiplying the
denominator dilation only changes fixed quadratic coefficients of the
contact forms. The coefficient of the zero contact is zero, so $u$
itself is one of the selected primes; the lower cutoff keeps it outside
$S$. This is within the fixed-dilation scope of the cited theorem.

\subsection{Powers of one local weight}
For nonzero rational $R,\delta$, define
\[
 q_m=R^{1-2m}\delta^{-2m-2},\qquad
 q=q_1=R^{-1}\delta^{-4},\qquad m\ge1.
\]
\begin{proposition}\label{v15:selection:prop:powers}
The alternative weights $q'_m=q^{2m-1}$ satisfy
\[
                 q'_m=q_m\delta^{-6(m-1)}.
\]
At every valuation where $\delta$ is a unit, $q'_m$ and $q_m$ have
the same valuation. At every valuation where $v(\delta)\le0$,
$v(q'_m)\ge v(q_m)$. In particular, a lower bound for
$v(q_mD^2)$ remains true with $q'_m$ in place of $q_m$.
\end{proposition}
\begin{proof}
The exponent of $R$ in $q^{2m-1}$ is $1-2m$ and its exponent of
$\delta$ is $4-8m$. Subtracting the exponent $-2m-2$ in $q_m$
gives $-6(m-1)$. Taking valuations proves the assertions.
\end{proof}

\begin{corollary}\label{v15:selection:cor:preservation}
Consider factors
\[
 1+\sum_{j=0}^{r-1}q_{m_j}^{2^j}D_j^{2^{j+1}}
\]
in a local test. At contact places where $v(R)$ is odd and $\delta$
is a unit, replacing $q_{m_j}$ by $q^{2m_j-1}$ preserves every
valuation and the distinct-valuation argument for its summands.
At the finite places prepared by Lemma~\ref{v15:selection:lem:denominator}, any
prescribed positive lower bounds on the summand valuations are
preserved. If, as $\lambda\to+\infty$, one has
$R\asymp\lambda^{-1}$, $\delta\asymp\lambda$, and all $D_j$
bounded at the real embeddings, the modified factors tend to one.
\end{corollary}
\begin{proof}
The finite-place conclusions follow term by term from
Proposition~\ref{v15:selection:prop:powers}. At a contact place, the valuation of
the $j$th nonconstant summand is $2^j$ times an odd integer; these
valuations and zero are pairwise distinct. At infinity,
$q=O(\lambda^{-3})$, hence
$q^{2m_j-1}=O(\lambda^{-3(2m_j-1)})$. Each nonconstant summand
therefore tends to zero.
\end{proof}

The corollary is a statement about the specified local tests and their
simultaneous preparation. A complete defining-polynomial construction
must still prove the pole realization and index-height steps, specify
its equations, and check their possible zero strata. Those separate
arithmetic requirements are unaffected by this normalization.


\section{Finite avoidance and simultaneous scalar preparation}\label{v15:dependencies}
\subsection{Finite avoidance and the common congruence progression}
\label{v15:dependencies:sec:avoidance}

For a Gaussian integer $\pi$, set
\[
 d_-(\pi)=N_{\Q(i)/\Q}(\pi-1),\qquad
 d_+(\pi)=N_{\Q(i)/\Q}(\pi+1).
\]
Theorem~\ref{thm:v5-simultaneous-primes} permits any fixed finite list of distinct integral
coefficients, any fixed divisibility of the square-denominator parameter,
and any fixed lower bounds for the prime values and for the slope.
Its avoidance modulus is odd, squarefree, and prime to seven.
The fixed Gaussian residues are $9$ modulo $16$ and $1+3i$ modulo seven;
at every prime dividing the avoidance modulus the residue is $i$.
The theorem excludes every odd common divisor between the two order
norms at the same position and between any two order norms at distinct
positions.

\begin{lemma}\label{v15:dependencies:lem:arbitrary-avoidance}
Let $L\ge1$ be divisible by $64$, and fix distinct integers
$c_1,\ldots,c_k$, with $k\ge2$. Prescribe $q\ge1$ and real numbers
$Y,H_0>0$. Theorem~\ref{thm:v5-simultaneous-primes} supplies positive integers $u,w$ such that
\[
 q\mid w,\qquad u/w^2>H_0,\qquad p_j=u+c_jw^2>Y
\]
are primes, with Gaussian generators whose order norms satisfy
\[
 \begin{gathered}
 v_2(d_-(\pi_j))=6,\qquad v_2(d_+(\pi_j))=2,\\
 \gcd(d_-(\pi_j),d_+(\pi_h))=4\qquad(1\le j,h\le k),
 \end{gathered}
\]
and are all prime to the odd part of $L$. The primes may simultaneously
be required to satisfy $p_j\equiv1\pmod{8\cdot3\cdot17}$.
\end{lemma}
\begin{proof}
Take the avoidance modulus to be the product of the odd primes dividing
$3\cdot17\cdot L$, with seven omitted. It is squarefree and prime
to seven, as required by Theorem~\ref{thm:v5-simultaneous-primes}. At each of its prime factors,
the two adjacent norms of $i$ are both two. They are therefore units.
At seven the fixed residue $1+3i$ instead gives
\[
       d_-(1+3i)=9\equiv2\pmod7,\qquad
       d_+(1+3i)=13\equiv6\pmod7.
\]
Thus seven is avoided even when it divides $L$.

Write $\pi=9+16a+16bi$. Direct calculation gives
\[
 d_+(\pi)\equiv4\pmod8,
 \qquad d_-(\pi)\equiv64\pmod{128}.
\]
The exact powers of two in the statement follow. The exclusion of odd
common factors supplied by Theorem~\ref{thm:v5-simultaneous-primes} then gives every stated
cross-sign gcd, including a position paired with itself.
Finally the Gaussian norm is one modulo eight, and its residues at
three and seventeen are the norm of $i$, namely one. These are precisely
the required congruences on $p_j$.
\end{proof}

\begin{corollary}\label{v15:dependencies:cor:compatible-progressions}
In the situation of Lemma~\ref{v15:dependencies:lem:arbitrary-avoidance}, put
\[
 U=\operatorname{lcm}_j\bigl(d_-(\pi_j)/4\bigr),\qquad
 V=\operatorname{lcm}_j\bigl(d_+(\pi_j)/4\bigr).
\]
Then
\[
              \gcd(U,V)=1,\qquad\gcd(U,L)=16,
              \qquad\gcd(V,L)=1.
\]
For every integer $h_*$ divisible by $64$, the congruences
\[
                 h\equiv0\pmod U,\qquad
                 h\equiv1\pmod V,\qquad
                 h\equiv h_*\pmod L
\]
have an arithmetic progression of solutions.
\end{corollary}
\begin{proof}
Division by four makes every minus order have two-adic valuation four
and every plus order odd. The cross-sign gcds make their least common
multiples coprime. Lemma~\ref{v15:dependencies:lem:arbitrary-avoidance} excludes all odd
factors of $L$ from both least common multiples. The displayed gcds
follow. The only nontrivial pairwise compatibility condition is
$h_*\equiv0\pmod{16}$, which holds. The generalized Chinese remainder
theorem now supplies the progression.
\end{proof}

The reduction orders at the two places of $K$ above $p_j$ are the
unordered pair $d_-(\pi_j),d_+(\pi_j)$, by the source lemma
\nolinkurl{lem:v5-two-reductions}. Both reductions have full rational
two-torsion. For a finite group of elliptic-curve points of order $d$,
full two-torsion implies that the subgroup of doubles has order $d/4$.
Every double is therefore killed by $d/4$. This justifies the divisions
by four in Corollary~\ref{v15:dependencies:cor:compatible-progressions}; the source point
$P$ is itself a double.

\begin{lemma}\label{v15:dependencies:lem:rational-contacts}
The preceding statements also apply to any fixed finite list of
rational contact coefficients, with a common square denominator for
the slope and with arbitrarily deep fixed denominator divisibility.
\end{lemma}
\begin{proof}
For coefficients $\gamma_j\in\Q$, choose a positive integer $D_0$
so that $D_0^2\gamma_j$ are integers, and remove repeated coefficients.
Apply Lemma~\ref{v15:dependencies:lem:arbitrary-avoidance} with
$c_j=D_0^2\gamma_j$ and put $v=D_0^2w^2$.
Then $u+c_jw^2=u+\gamma_jv$ and $u/v>D_0^{-2}H_0$.
Enlarge the prescribed lower bound before applying the lemma to obtain
any desired slope bound. By choosing $q$ divisible by suitable powers
of a fixed finite set of primes, one obtains any prescribed lower
bounds on their valuations in $v$. Enlarge $Y$ to exclude any fixed
finite list of primes. None of these choices changes the quantified
form of Theorem~\ref{thm:v5-simultaneous-primes}.
\end{proof}

\begin{lemma}\label{v15:dependencies:lem:real-progression}
Let $G$ be the closure of a cyclic subgroup generated by a point $Q$
in a compact real Lie group, and let $G^0$ be its identity component.
Suppose $L$ is divisible by the order of $G/G^0$. Any arithmetic
progression $h\equiv h_*\pmod L$, with an additional compatible
finite collection of congruences, is dense after multiplication by
$Q$ in the component $h_*Q+G^0$. Arbitrarily large positive indices
can be used.
\end{lemma}
\begin{proof}
The compatible congruences give one coset modulo an integer $L'$
divisible by $L$. A compact connected abelian Lie group is a torus,
and multiplication by any nonzero integer is onto it. It follows that
$[L']G=G^0$. Thus the two-sided progression has the stated closure.
In a compact group, the closure of a positive cyclic orbit is also
a group: positive return times approach the identity and consequently
approximate the negative multiples. Removing a finite initial segment
does not change the closure. This proves the positive-index assertion.
\end{proof}

For the construction, apply this lemma to the closure of
$(P,P^\sigma)$ in the product of the two real elliptic groups.
The scalar point, the four integer root labels, the real compact
neighborhoods, all contact coefficients, and the modulus $L$ are fixed
before the prime-pattern theorem is used. Corollary~\ref{v15:dependencies:cor:compatible-progressions}
then supplies the final shared index without disturbing those
neighborhoods. No estimate uniform in the coefficients, $q$, $L$, or
the number of contact positions is required.

\subsection{Pole realization with zero preparation}
\label{v15:dependencies:sec:poles}

For $b\in\Q^\times$, put $B=17b$. The zero-preparation predicate is
\[
 \begin{aligned}
 p^2&=B(1-e)(1+3e),\\
 h_0&=e(B-b^2),\\
 Bq^2&=(s^2-Bh_0)(s^2+3Bh_0),
 \end{aligned}
\]
with rational $e,p,q,s$; $h_0$ is the displayed expression.
No nonzero condition is imposed on these four witnesses.

Lemma~\ref{v5:preparation-zero-bridge} states exactly
the following two implications. A solution excludes negative odd
valuations of $b$ outside $\{2,3,17\}$. Conversely a solution exists if
the reduced denominator of $b$ is a square and
\[
 \begin{gathered}
 0<b<17,\qquad b\in\Q_2^{\times2}\cap\Q_3^{\times2}
                              \cap\Q_{17}^{\times2},\\
 v_2(b)\ge3,\qquad v_3(b)\ge1,\qquad v_{17}(b)\ge2.
 \end{gathered}
\]
These are precisely the hypotheses verified for each of the four
normalized pole inputs in the finite-algebra construction.

The special role of zero preparation can also be checked directly in
the source argument. With preparation parameter zero, the auxiliary
quantity there is
\[
                    B(B-b^2)=B^2(1-b/17).
\]
It is positive and is a square in the three specified completions,
by the displayed bounds. If $\mu,\theta$ are the positive squarefree
representatives of $B$ and this quantity, their supports are disjoint,
$\mu\equiv17\pmod{24}$, $\theta\equiv1\pmod{24}$, and each is a
local square at primes dividing the other. To see the last condition,
an odd prime in the support of $\mu$ has positive odd valuation on $b$,
so $1-b/17$ is a square unit; a prime in the support of $\theta$ outside
$2,3,17$ must have $v_p(b)=0$ and $b\equiv17\pmod p$, making
$B\equiv17^2\pmod p$. The conditions at $2,3,17$ were already imposed.
The simultaneous-realization argument in Lemma~\ref{v5:preparation-zero-bridge} therefore applies without
an additional hypothesis. Its deep restricted-converse and rank
dependencies are retained in this use.

\begin{lemma}\label{v15:dependencies:lem:conic-complete}
For $B\in\Q^\times$, the first conic is covered by the two formulas
\[
 e=\frac{\tau^2+(4\varepsilon-1)B}{\tau^2+3B},\qquad
 p=\frac{-4B\tau(1-\varepsilon)}{\tau^2+3B},\qquad
 \varepsilon\in\{0,1\},\quad\tau^2+3B\ne0.
\]
Thus its parametrization imposes no additional restriction on the
pole-realization theorem.
\end{lemma}
\begin{proof}
Substitution verifies the conic equation. If $e\ne1$, take
$\varepsilon=0$ and $\tau=p/(e-1)$. The conic equation gives
$\tau^2+3B=4B/(1-e)\ne0$ and recovers the displayed point.
If $e=1$, then $p=0$, and $\varepsilon=1,\tau=0$ works.
The case $e=-1/3,p=0$ is included in the first chart.
\end{proof}

The two structural coordinates attached to different pole inputs are
independent. Once all four inputs satisfy the source hypotheses, choose
each rational certificate separately and use Lemma~\ref{v15:dependencies:lem:conic-complete}.
Zero second-square roots are included by the explicit finite support
split in the finite-algebra construction. No common structural witness
for distinct pole inputs is asserted or needed.


\paragraph{Role of this baseline.} This section proves the intermediate 248-variable construction with its original shifted scalar index. It is not the final variable count. Section~\ref{v15:reductions} subsequently replaces the shifted scalar argument and gives the fixed-input transformations to 238 variables. References below to arithmetic inputs refer to the internal sections just established, not external manuscript assumptions.
\section{The 248-variable arithmetic baseline}\label{v15:base248}
\subsection{The family and its arithmetic scope}
All variable counts include the distinguished input $X$. For a fixed
integer $N\ge1$ let $A_N=149^{2N+1}$. The complete tables
\begin{center}\small
\nolinkurl{quadratic-family248-v17.json}\\
\nolinkurl{polynomial-family248-v17.json}
\end{center}
specify 222 primitive integer
quadratic polynomials and their sum of squares
\[
 P_N(X,Y_1,\ldots,Y_{247})=\sum_{j=1}^{222}Q_{N,j}^2.
\]
The parameter $A_N$ is a coefficient, not a variable.

\begin{theorem}\label{v15:base248:thm:main}
Using the internally proved arithmetic consequences stated below, there is a fixed integer
$N_0\ge53$ such that every $N\ge N_0$ satisfies
\[
 x\in\Z\quad\Longleftrightarrow\quad
 \exists\mathbf y\in\Q^{247}\ P_N(x,\mathbf y)=0
 \qquad(x\in\Q).
\]
For every $N\ge1$ the polynomial $P_N$ has degree exactly four,
content one, and 3189 nonzero coordinate monomials. For $N\ge53$,
\[
                       H(P_N)=149^{4N+2}+1.
\]
The construction specifies every member exactly. The proof does not
determine a sufficient numerical value of $N_0$.
\end{theorem}

Set
\[
 K=\Q(\rho),\quad \rho^2=2,\quad D=(3-\rho)^2=11-6\rho,
 \qquad E:y^2=x^3-Dx,
\]
\[
 P=2(1+\rho,2\rho)
 =\left(\frac{57-28\rho}{8},\frac{-660+417\rho}{32}\right).
\]
We use the following previously proved consequences, with their precise domains. The normalized group and index-height bound are established in the elliptic-model sections; Sections~\ref{v15:dependencies} and \ref{v15:selection} give the simultaneous finite and prime preparation:
\begin{enumerate}
\item $2E(K)=\Z P$ and $h(x([n]P))\ge n^2-6$ for every nonzero
integer $n$, together with the specified point charts and group domains.
\item The modular-realization lemma and contact-height theorem of Section~\ref{v15:height} prove
\begin{equation}\label{v15:base248:eq:baseheight}
 h(u)\le H M_5(u)^{7/2},\qquad H=2^{330789}.
\end{equation}
Here $M_5(u)$ is the odd contact radical outside $76176$ for
$\infty,0,1,2401,-343/32$, and $u$ avoids the four finite contacts.
The same theorem also proves the sharper estimate
\begin{equation}\label{v15:base248:eq:baseheightsharp}
 h(u)\le C_3M_5(u)^3(1+\log M_5(u)),\qquad C_3=2^{330770}.
\end{equation}
Both bounds are consequences of the modular input, rather than an
additional $abc$ assumption.
\item The pole obstruction outside $2,3,17$, and its simultaneous
realization for positive inputs less than $17$ with square reduced
denominator and the stated local squareclasses and valuation bounds
at those three primes. Its normalized quadratic equations are retained.
\item The simultaneous prime-pattern, conjugate reduction-order,
finite congruence, and real approximation assertions used by the source.
The weighted selection and fixed denominator dilation are justified
from the proved counting estimates in Section~\ref{v15:selection}.
\end{enumerate}
The aggregate trace-height theorem used below is the proved consequence
of the family height theorem described in Sections~\ref{v15:aggregate}; its additive
constant is ineffective. The contact amplification below is an elementary consequence of~\eqref{v15:base248:eq:baseheight},
not an additional arithmetic assumption. The finite coefficient calculation is separate from these arithmetic arguments.

\subsection{Two contact blocks and the comparison scales}
Put
\[
 \Theta=\{0,-1/2,-1,1/2399,-32/407\},\qquad
 \Theta_2=\Theta\cup(2-\Theta),\qquad
 F_2(t)=\prod_{\theta\in\Theta_2}(t-\theta).
\]
The ten contacts are distinct, omit one, and are symmetric about one.
Let $\Sigma$ consist of the primes at most $76176$ and the prime
$2029147$. Outside this set, all contact differences and all numbers
$1-\theta$ are units. Define $M_2^\circ(t)$ as the product of primes
$p\notin\Sigma$ for which $v_p(t-\theta)$ is positive and odd for
some $\theta\in\Theta_2$. Exact factorizations in the independent
certificate verify this exceptional set. In particular
\[
 F_2(1)=-\frac{82429855236}{7878870169},\qquad M_2^\circ(1)=1.
\]

Set $E_2=2029147$ and $H_2=\log7+HE_2^{7/2}$. Then
\begin{equation}\label{v15:base248:eq:contactheight}
                   h(t)\le H_2M_2^\circ(t)^{7/4}
                   \qquad(t\notin\Theta_2).
\end{equation}
Indeed, apply~\eqref{v15:base248:eq:baseheight} to $2+1/t$ and $2+1/(2-t)$.
The inverse projective matrices have maximum absolute row sums three
and seven, respectively, so the inverse height change is at most
$\log7$. If $M_+,M_-$ are the corresponding block radicals with
$\Sigma$ omitted, the original radicals are at most $E_2M_+$ and
$E_2M_-$. Outside $\Sigma$ the two radicals are coprime and their
product is $M_2^\circ(t)$. Multiplying the two inequalities
\[
              h(t)-\log7\le HE_2^{7/2}M_\pm^{7/2}
\]
when their left side is positive proves~\eqref{v15:base248:eq:contactheight};
if it is nonpositive, the conclusion is immediate.

The retained root points are
\[
 Q_1=[n]P,\quad Q_2=[a]P,\quad Q_3=[t]P,\quad Q_4=[b]P,
 \qquad B=\max(|n|,|a|,|t|,|b|),
\]
with nonzero integer labels. Set
\[
 T=\frac14\Tr_{K/\Q}\left(\sum_{i=1}^4x(Q_i)\right).
\]
It is positive, and the aggregate trace theorem gives fixed constants
$M_1>0$, $M_2\ge1$ such that
\begin{equation}\label{v15:base248:eq:aggregate}
                         B^2\le M_1h(T)+M_2.
\end{equation}
Use the two positive rational scales $a_*=1,T$. No square condition
on a scale is needed. Replacing the half-trace in the aggregate
height theorem by its half changes its logarithmic height by at
most $\log2$; this is absorbed into $M_2$.
For each scale put $t_*=a_*\lambda$ and $\delta=t_*-1$.
On the nonzero domain use the balanced contact and the local weight
\begin{equation}\label{v15:base248:eq:detector}
 R=\frac{F_2(t_*)}{\delta^{10}},\qquad
 q=\frac{\delta^{10}}{F_2(t_*)}-1=R^{-1}-1,
 \qquad \gamma=\frac{408}{409}.
\end{equation}
The contact circuit stores $r=\delta^2$ and its powers $r^2,r^3$.
Its actual quadratic normalization is
\begin{equation}\label{v15:base248:eq:clearing}
                         F_2(t_*)(1+q)=r^2r^3.
\end{equation}
Thus $R^{-1}=1+q$ when $\delta\ne0$, and no additional coordinate
for this reciprocal is needed. No contact value can occur:
$F_2(t_*)=0$ would force $r=0$ and $t_*=1$, contradicting $F_2(1)\ne0$.
There are two distinct exceptional strata. If $\delta=0$, then $q=-1$
and $M_2^\circ(t_*)=1$; no reciprocal is used. If $q=0$, then $\delta\ne0$
and $F_2(t_*)=\delta^{10}$, so again $M_2^\circ(t_*)=1$. Indeed a consumed
contact prime would have odd valuation on the left and valuation zero
on the right, since $\delta$ is a unit there.

If both radicals divide a nonzero integer $J$ to order $k$, then
$h(T)\le h(\lambda)+h(T\lambda)$ and
\eqref{v15:base248:eq:contactheight}--\eqref{v15:base248:eq:aggregate} imply
\begin{equation}\label{v15:base248:eq:heighttransfer}
 B^2\le C|J|^{7/(4k)}+M_2,\qquad C=2M_1H_2.
\end{equation}
The two radicals need not be coprime for this implication.

\subsection{The scalar certificate and the point domains}\label{v15:base248:sec:scalar}
Put $\beta=51X^2+2504$, let $S=Q_2=[a]P$, and write
$C_S=S-P=[a-1]P$. The scalar normalizer is
\begin{equation}\label{v15:base248:eq:scalarnormalizer}
 \begin{split}
 \ell&=149,\qquad A_N=\ell^{2N+1},\\
 r_s&=\beta(\beta-1)(\beta-2401)(32\beta+343)(2\beta-1),\\
 R_s&=\frac{A_Nr_s}{1+\ell^2r_s}.
 \end{split}
\end{equation}
All five factors of $r_s$ are positive, since $\beta\ge2504$, so
$R_s>0$. The retained square-abscissa chart for $S$ is
\[
 x_S=h_S^2,\qquad j_S^2=x_S^2-D,
 \qquad S=(x_S,h_Sj_S).
\]
It has $h_Sx_Sj_S\ne0$. The fixed-subtraction coordinate $g_S\in K$
satisfies
\[
 8(x_Px_S+D)g_S+(39-36\rho)h_S+(8-14\rho)j_S=0.
\]
The rational component of $g_S$ has a retained inverse witness.
At $S=P$ the coordinate $g_S$ is zero, so this guard gives $a\ne1$.
All nonzero points of $\Z P$ are nontorsion. Thus $C_S$ is a nonzero
nontorsion point, and the usual fixed-subtraction formulas have no
remaining zero denominator. Explicitly, define
\[
 \begin{split}
 N_S&=(-660+417\rho)h_S+(-228+112\rho)j_S,\\
 D_S&=(-704+384\rho)h_S+(660-417\rho)j_S
                         +(6888-5040\rho)g_S.
 \end{split}
\]
Then $D_S\ne0$ and $z(C_S)=N_S/D_S$, where $z=-x/y$.
These identities are obtained from the chord-and-tangent formulas;
they are the same fixed-subtraction identities independently checked
in the exact function field by the supplied certificate.

The common point chart retains $\chi,f,S_H\in K$ with
\begin{equation}\label{v15:base248:eq:Hchart}
                  \chi f=1,\qquad S_H^2=\chi-Df,
                  \qquad H=(\chi,\chi S_H).
\end{equation}
Write $H=[m]R_0+T_0$, where $P=2R_0$ and $T_0\in E[2](K)$,
and put $h=m/2$. Before any logarithm is used, the scalar ratio
coordinate $Z=Z_0+\rho Z_1$ is determined by
\[
                         ZD_S=-N_SS_H.
\]
When $H$ is nontorsion this is $Z=z(C_S)/z(H)$.
Set $F=\Nm_{K/\Q}(f)\ne0$. The reduced scalar norm equations are
\begin{equation}\label{v15:base248:eq:scalarnorm}
 A=\frac{Z_0^2+R_sZ_1^2}{R_s^4},\qquad
 W_s=\frac{A^2+R_sF^2}{R_s^4},\qquad
 b_s=R_s(R_s+9792)W_s^2.
\end{equation}
The actual quadratic source stores $R_s^2,R_s^4,Z_1^2,F^2,A,W_s$
and a scaled square of $W_s$; thus these equations are quadratic after
clearing. The normalized scalar pole has exactly the ordinary pole
input $b_s$ in~\eqref{v15:base248:eq:scalarnorm}. To verify its scaling literally,
let $M$ be the coefficient of the stored $W_s$-square coordinate in
source row491, let $K$ be the coefficient of $e\bar b$ in row499,
and let $L$ be the coefficient of $\bar b p^2$ in row494. The exact
coefficient identities are
\[
                  K=82479248M,\qquad 3L=17K.
\]
The rational substitutions
\[
 b_s=-K\bar b,\quad e'=-e/K,\quad p'=p/K,\quad
 u'=u/K,\quad v'=-v,\quad w'=w/K
\]
convert those rows to the retained normalized pole equations. The
independent checker verifies these identities. Since $R_s>0$ and
$F\ne0$, one has $W_s>0$ and $b_s>0$ at every rational source solution.

\begin{lemma}\label{v15:base248:lem:scalarextract}
Let $p>76176$ or $p=149$, and suppose $\nu=v_p(R_s)>0$ is odd. Then
\[
 v_p(F)>5\nu/4,\qquad
 v_p(Z_0)>23\nu/8,\qquad v_p(Z_1)>19\nu/8.
\]
Once the point interpretations hold, these imply
\[
                v_p(a-1)>19\nu/8.
\]
\end{lemma}
\begin{proof}
The number $9792=2^6\cdot3^2\cdot17$ is a unit at $p$.
The valuation of the nonzero pole input is
$\nu+2v_p(W_s)$, so the pole obstruction makes it positive and
$v_p(W_s)>-\nu/2$.
The two nonzero terms of a norm $u^2+R_sv^2$ have opposite
valuation parity. A zero term can be omitted. Hence the second
norm in~\eqref{v15:base248:eq:scalarnorm} gives
\[
 v_p(A)>7\nu/4,\qquad v_p(F)>5\nu/4.
\]
The first norm gives the two asserted bounds on $Z_0,Z_1$.

At some place above $p$ the value $f(H)$ has positive valuation,
because the valuation of its norm is positive. Thus $H$ is in the
formal group there. The displayed bounds give $v(Z)>19\nu/8>0$,
so $v(z(C_S))>0$. On $y^2=x^3-Dx$, at a good odd prime a point
with positive $z$-valuation is either formal or reduces to $(0,0)$.
Indeed $z^2=x/(x^2-D)$, and $D$ is a unit. Translation by $(0,0)$
sends $z$ to $-z$ and exchanges these two cases. Choose
$C'=C_S+\epsilon(0,0)$ to be formal. The integral identity
\[
                       [2m]C'=[4(a-1)]H
\]
kills both two-torsion terms. The formal logarithm is an isometry,
so its ratio has valuation $v_p((a-1)/h)=v(Z)$. Since $h=m/2$ is
$p$-integral, $v_p(a-1)\ge v(Z)>19\nu/8$.
\end{proof}

The integer $\ell=149$ is prime. For the five contact roots
$0,1,2401,-343/32,1/2$, the residues of $(\theta-2504)/51$ are
respectively $59,97,109,108,78$ modulo $149$. Their $74$th powers
are all $148$, so they are nonsquares. Trial division through $12$
certifies primality. The preceding lemma applies at $149$: the curve
has good reduction there, $9792$ is a unit, and the pole obstruction
excludes only $2,3,17$. The generic contact cutoff is unnecessary
at this separately verified prime. Thus $r_s$ is an $\ell$-adic unit when $X$ is integral
at $\ell$. If $v_\ell(X)<0$, then $v_\ell(r_s)=10v_\ell(X)$.
It follows that
\begin{equation}\label{v15:base248:eq:forcedvaluation}
 v_\ell(R_s)=
 \begin{cases}2N+1,&v_\ell(X)\ge0,\\2N-1,&v_\ell(X)<0.
 \end{cases}
\end{equation}
This is always positive and odd. These finite facts are included
in the independent exact certificate.

Lemma~\ref{v15:base248:lem:scalarextract}, using only its norm inequalities at
this stage, gives $v_\ell(F)>0$. Every finite allowed point in
\[
                E[2](K)\ \cup\ (P+E[2](K))\ \cup\ (-P+E[2](K))
\]
has reciprocal-abscissa norm an $\ell$-adic unit. To check this,
use the two-torsion translations and
\[
 \Nm(x_P)=41^2/64,\quad
 \Nm(x_P-(3-\rho))=17^2/64,\quad
 \Nm(x_P+(3-\rho))=63^2/64.
\]
The remaining torsion abscissas have norm seven. Thus none of
these points can equal $H$. In particular $h\notin\{-1,0,1\}$,
$H$ and $H-P$ are nontorsion, and all required denominators and
formal-logarithm interpretations are now justified. This exclusion
precedes every use of a comparison logarithm.

Let $c=|a-1|>0$. Outside the fixed cutoff, every prime of
$M_5(\beta)$ is a positive odd contact of $r_s$ and $R_s$.
The contact at infinity contributes none, since a negative valuation
of $\beta=51X^2+2504$ is even there. Lemma~\ref{v15:base248:lem:scalarextract}
therefore implies $M_5(\beta)^3\mid c$. The prime $\ell=149$ is below the defining cutoff of
$M_5$, so it is automatically absent from $M_5(\beta)$.
The forced-prime estimate is consequently independent of these radical
factors. We obtain
\begin{equation}\label{v15:base248:eq:scalardivisibility}
 \ell^{L_N}M_5(\beta)^3\mid c,\qquad
 L_N=\left\lfloor\frac{19(2N-1)}8\right\rfloor+1.
\end{equation}
Together with~\eqref{v15:base248:eq:baseheightsharp}, this gives
\begin{equation}\label{v15:base248:eq:scalar}
 h(\beta)\le C_3\ell^{-L_N}c\left(1+\frac{\log c}{3}\right),
 \qquad c\ge\ell^{L_N}.
\end{equation}
Indeed $M_5(\beta)^3\le c/\ell^{L_N}$ and
$\log M_5(\beta)\le(\log c)/3$.

\subsection{Plain parameter ratios and the local tests}
For $f=1/x$ and $z=-x/y$, extend their regular values at the
identity by zero. The retained point relations are
\[
 U_0=H,\quad V_0=H-P,\quad U_i=[u_i]P,\quad V_i=U_i-Q_i
 \qquad(1\le i\le4).
\]
The charts permit $U_i=O$ or $V_i=O$ for $i>0$.
The fixed subtraction for $H$ stores $g,F_0\in K$ with
\[
 8(x_P\chi+D)g+(39-36\rho)+(8-14\rho)S_H=0,
 \qquad fF_0=g^2.
\]
Its affine numerator and denominator, after multiplying by the
common unit $3+2\rho$, are
\[
 \begin{split}
 N_H&=(-312-69\rho)+(-236-120\rho)S_H,\\
 D_H'&=(-576-256\rho)+(312+69\rho)S_H+(504-1344\rho)g.
 \end{split}
\]
On the established domain, $F_0=f(H-P)$ and $z(H-P)=N_H/D_H'$.
The exact function-field certificate checks these identities against
the addition law, including the exceptional point
$-P+(0,0)$ already excluded above.
For each $i$, retain the primary ratio $L_i$ and an unconjugated
alternate ratio $L_{i,0}'$ using
\[
             L_{i,0}'N_H=z(V_i)D_H',\qquad L_i'=\sigma L_{i,0}'.
\]
These are quadratic equations and mean precisely
\[
       L_i=\frac{z(U_i)}{z(H)},\qquad
       L_i'=\sigma\!\left(\frac{z(V_i)}{z(H-P)}\right).
\]
No logarithm polynomial or coordinate for $z(H-P)$ is retained.
Put
\[
 D_H=f(H),\quad A_i=f(U_i),\quad B_i=\sigma f(V_i),\quad
 C_i=L_i-L_i',\quad I=L_1-\beta,
\]
\[
 G=L_3^2-L_2(L_2-2)L_3L_4+(L_2-2)^2L_4^2-1.
\]
The required orders are two for $D_H$, four for $C_2,C_3,C_4,G$,
one for $C_1$, two for $I$, and one for the nine $A_i,B_i$.
The final input comparison uses order one.

\subsection{Five factors and extraction at one place}
At a contact prime $p\notin\Sigma$, $\delta$ is a unit and
$\nu=v_p(R)=v_p(F_2(t_*))$ is positive and odd, and $v_p(q)=-\nu$. This follows
from the unique contact factor with positive odd valuation and
does not require any condition on the valuation of the scale. All valuations
of $K$ are unramified and normalized by $v(p)=1$.
For a point on $E$, every finite value of $v(f)$ is even: if
$v(x)<0$ then $2v(y)=3v(x)$, while if $v(x)>0$ then
$2v(y)=v(x)$; the remaining case is valuation zero.
This also holds for conjugated reciprocal abscissas. Before excluding
the exceptional $H$ classes, the stored value for $f(H-P)$ still has even
valuation: its literal equation is $f(H)F_0=g^2$, and $f(H)$ has even
valuation. Thus the parity argument and the preceding exclusion do not
assume the validity of a division in the fixed subtraction chart.

Use
\begin{align*}
 T_1&=1+q^7C_2^2+q^2D_H^2+q^4B_1^4,\\
 T_2&=1+q^7G^2+q^2A_3^2+q^4B_0^4,\\
 T_3&=1+q^7I^4+qA_2^2+q^2A_4^2+q^4B_3^4,\\
 T_4&=1+q^7C_3^2+q^2C_1^4+q^4B_4^4,\\
 T_5&=1+q^7C_4^2+q^2A_1^2+q^4B_2^4,
 \qquad W=T_1T_2T_3T_4T_5.
\end{align*}

\begin{lemma}\label{v15:base248:lem:residues}
At a contact prime, every factor is nonzero and has valuation at most
zero. A factor of valuation zero is congruent to one in the residue field.
\end{lemma}
\begin{proof}
First suppose every displayed test in a factor is nonzero. In each
factor other than $T_3$, the three nonconstant valuations belong,
respectively, to the odd integers, to $2+4\Z$, and to $4+8\Z$.
For example $T_1$ has valuations
\[
 2v(C_2)-7\nu,\qquad 2v(D_H)-2\nu,\qquad4v(B_1)-4\nu.
\]
The second has exact factor two because $v(D_H)$ is even and $\nu$ is
odd. The third has exact factor four because $v(B_1)$ is even.
The term $q^2C_1^4$ in $T_4$ has valuation
$2(2v(C_1)-\nu)$, again twice an odd integer.

In $T_3$ the four nonconstant valuations are
\[
 4v(I)-7\nu,\quad 2v(A_2)-\nu,\quad
 2v(A_4)-2\nu,\quad4v(B_3)-4\nu.
\]
The first three have residues $\nu,-\nu,2$ modulo four. These are
pairwise distinct and nonzero because $\nu$ is odd and the $A_i$
have even valuations. The fourth has residue four modulo eight;
it therefore differs from each of the first three and from zero.
Thus, in every factor, all finite term valuations, including zero
for the constant term, are distinct. Zero tests simply remove terms
from this argument. The minimum is attained exactly once. This proves
nonvanishing and the asserted bound. If the minimum is zero, every
nonconstant term has positive valuation, so the residue is one.
\end{proof}

\begin{proposition}\label{v15:base248:prop:extract}
At each contact prime, the retained pole system and the cleared trace
equations imply $\nu>0$ and give all sixteen required test orders at
one place above that prime.
\end{proposition}
\begin{proof}
Write $W=w_0+w_1\rho$. The literal cleared trace equations are
\[
 w_0^2-2w_0\tau-2w_1^2=0,\qquad (1+q)\mathfrak b=\gamma^2\tau^2.
\]
By Lemma~\ref{v15:base248:lem:residues}, $W\ne0$. The first equation therefore
gives $w_0\ne0$ and $\tau=1/\Tr(W^{-1})\ne0$; indeed $w_0=0$
would give $w_1=0$. Hence
$\mathfrak b=\gamma^2R/\Tr(W^{-1})^2$. The factor $\gamma$ is a unit at every consumed contact prime. If
$\ell=\max_{v\mid p}v(W)\le0$, the trace inequality gives
$v_p(\mathfrak b)\le\nu+2\ell$.
This valuation is odd; the retained pole obstruction forces it positive.
Consequently $\nu>0$ and $\ell>-\nu/2$.
At a place attaining $\ell$, the factor valuations are nonpositive
and add to $\ell$. Each factor valuation is therefore at least $\ell$,
and each of its term valuations is strictly greater than $-\nu/2$.

The resulting lower bounds, before integer rounding, are
\[
\begin{array}{c|c|c}
\text{test}&\text{strict lower bound}&\text{integer order}\\\hline
C_2,C_3,C_4,G&13\nu/4&4\\
I&13\nu/8&2\\
C_1&3\nu/8&1\\
D_H&3\nu/4&2\\
A_2&\nu/4&1\\
A_1,A_3,A_4&3\nu/4&1\\
B_0,B_1,B_2,B_3,B_4&7\nu/8&1
\end{array}
\]
Here $\nu$ is a positive odd integer. The order two for $D_H$ uses
its even valuation. In fact each of the nine reciprocal-abscissa
tests has order at least two; order one suffices for its subsequent
formal-group use. Zero tests satisfy all the stated inequalities.
\end{proof}

\subsection{Integrality before the formal congruences}
At a contact place selected by Proposition~\ref{v15:base248:prop:extract}, the
points $H,U_i$ are formal in the primary completion, and $H-P,V_i$
are formal in the conjugate completion. The reciprocal-abscissa
tests have positive even valuation, so every nonidentity point has
formal depth at least one. In the parameter $z=-x/y$, write
\[
 \log_E(Q)=z(Q)\mathcal A(Q),\qquad
 \mathcal A(Q)\equiv1\pmod{p^{4v(z(Q))}}.
\]
The value of the unit is taken to be one at the identity. This is
the constant-term formal logarithm estimate proved in~Sections~\ref{v15:logs}.

The integral group identities
\[
                  [m]U_i=[2u_i]H,\qquad
                  [m-2]V_i=[2(u_i-n_i)](H-P)
\]
kill the torsion in $H=[m]R_0+T_0$. Their full logarithmic ratios
are $u_i/h$ and $(u_i-n_i)/(h-1)$, where
$(n_1,n_2,n_3,n_4)=(n,a,t,b)$. At least one of $h,h-1$ is a
$p$-adic unit, so at least one full ratio is integral. The formal
units preserve the valuation of its plain parameter ratio.
The positive-order test $C_i$ transfers integrality to the other
plain ratio, and then to the other full ratio.

Only after this transfer may the relative errors of order four
be used as absolute congruences. They give
\[
 L_i\equiv u_i/h\pmod{p^4},\qquad
 L_i'\equiv(u_i-n_i)/(h-1)\pmod{p^4}.
\]
Since $n_i=h(u_i/h)-(h-1)((u_i-n_i)/(h-1))$, the high cross tests
give $L_i\equiv n_i\pmod{p^4}$ for $i>1$. The first cross test
gives $L_1\equiv n\pmod p$, while $I$ gives $L_1\equiv\beta\pmod p$.
These arguments include identity numerator points, half-integral $h$,
and inert primes, where conjugation preserves the unique valuation.

\subsection{Global soundness for sufficiently large parameters}\label{v15:base248:sec:soundness}
The scalar bounds~\eqref{v15:base248:eq:scalar} hold independently of the
comparison logarithms. In particular $c\to\infty$ uniformly with $N$,
and $B\ge c-1$.

Let
$\Delta=t^2-a(a-2)tb+(a-2)^2b^2-1$.
The local comparison shows that each fourth power of a comparison
radical divides $\Delta$. If $\Delta\ne0$, then
$|\Delta|\le14B^4$, and \eqref{v15:base248:eq:heighttransfer} gives
\[
                  B^2\le C14^{7/16}B^{7/4}+M_2.
\]
Since $7/4<2$, this bounds $B$ by a fixed constant. Equation
\eqref{v15:base248:eq:scalar} gives $B\ge c-1\to\infty$ uniformly as
$N\to\infty$, a contradiction. Thus $\Delta=0$ eventually.
A stronger form of the integral Lucas descent gives
\begin{equation}\label{v15:base248:eq:stronglucas}
              \log B\ge(c-5)\log(c-2)\qquad(c\ge8).
\end{equation}
Here is the complete elementary argument. Put $A_0=|a|$ and
$w=\operatorname{sgn}(a)(a-2)b$. Then
$t^2-A_0tw+w^2=1$. For
$\eta=(A_0+\sqrt{A_0^2-4})/2$, every positive norm-one unit of
$\Z[\eta]$ is a power of $\eta$. Indeed multiply a unit by a
power of $\eta$ to place it in $[1,\eta)$; its coefficient of
$\eta$ is $(v-v^{-1})/(\eta-\eta^{-1})$, an integer in $[0,1)$,
so it is zero and $v=1$.
Let $U_0=0$, $U_1=1$, and $U_{j+1}=A_0U_j-U_{j-1}$ for $j\ge1$.
The unit identities show that $|w|=U_j$ and
$|t|\in\{U_{j-1},U_{j+1}\}$ for some $j\ge1$.
Modulo $A_0-2$ one has $U_j\equiv j$; modulo $A_0+2$ one has
$U_j\equiv(-1)^{j-1}j$. Since $w$ is a nonzero multiple of
$|a-2|$, it follows that $j\ge A_0-2$ if $a>0$ and
$j\ge A_0+2$ if $a<0$. Induction gives
$U_j\ge(A_0-1)^{j-1}$ for $j\ge1$. Therefore
$|t|\ge(A_0-1)^{A_0-4}$. Finally $A_0\ge c-1$, which proves
\eqref{v15:base248:eq:stronglucas}.

If $n-\beta\ne0$, let $J$ be its nonzero reduced numerator.
The local input comparisons make both radicals divide $J$,
while $|J|\le2B\exp(h(\beta))$. For $B^2\ge2M_2$,
\eqref{v15:base248:eq:heighttransfer} gives
\[
 \log B\le4\log(2C)+7\log2+7h(\beta)
 \le C'+7C_3\ell^{-L_N}c\left(1+\frac{\log c}{3}\right).
\]
Choose $N$ sufficiently large that $7C_3\ell^{-L_N}\le1$.
The right side is then at most $C'+c+(c\log c)/3$, which is
strictly smaller than $(c-5)\log(c-2)$ for all sufficiently large
$c$. Since $c\ge\ell^{L_N}\to\infty$ uniformly with $N$,
this contradicts~\eqref{v15:base248:eq:stronglucas}.
Hence $n=\beta\in\Z$. If $X=u/v$ is reduced with $v>0$, then
$v^2\mid51$, and squarefreeness of $51$ gives $v=1$.
All bounds depend only on fixed constants. Since $L_N$ is increasing,
a single $N_0$ works for every larger $N$. Its dependence on the
aggregate additive constant is ineffective.

\subsection{Simultaneous completeness}
Fix an integer input. The scalar construction can be made uniformly
complete for all sufficiently large $N$, as follows. First choose a fixed
real neighborhood of a nonzero reference point $[h_*]P$, with
$64\mid h_*$, on which both real values of $|z(H)|$ are bounded below.
Choose a fixed neighborhood of the identity for $C_S$ on which its two
$z$ values are bounded. Then $|Z_0|,|Z_1|$ are bounded by a constant
independent of the input and $N$. Also $|\Nm(f(H))|\le1/7$ when
$H\in\Z P$ is nonzero, by the positive-square abscissa bounds.
For $R_s\ge1$, equations~\eqref{v15:base248:eq:scalarnorm} consequently give
\[
                 W_s=O(R_s^{-3}),\qquad b_s=O(R_s^{-4}),
\]
with fixed constants. Since $r_s\ge1$,
$R_s\ge A_N/(1+\ell^2)\to\infty$ uniformly in the input.
Hence $0<b_s<17$ for all sufficiently large $N$ on these neighborhoods.

The finite preparation uses no new arithmetic theorem. For integer
input, $r_s$ is a positive integer and
\[
 R_s(R_s+9792)=
 \frac{A_Nr_s\bigl(A_Nr_s+9792(1+\ell^2r_s)\bigr)}
                    {(1+\ell^2r_s)^2}.
\]
The numerator is coprime to the displayed denominator, so its reduced
denominator is square. At the three preparation primes the elementary
input congruences give
\[
 v_2(R_s)\in\{2,3\},\qquad v_3(R_s)=1,
                       \qquad v_{17}(R_s)=0.
\]
In particular $r_s\equiv15\pmod{17}$ and
$1+149^2r_s\equiv3\pmod{17}$, so the normalizer denominator is
also a unit at $17$. Thus
$R_s(R_s+9792)=R_s^2(1+9792/R_s)$ is a local square at all three primes.
Choose the shared index $h_0$ sufficiently divisible at the finite
support of this rational number, and then choose the index
$c=a-1$ still more divisible. The points $H=[h_0]P$ and $C_S=[c]P$
are then formal at the required finite places, and all components of
$f(H)$ and $Z$ have the needed arbitrarily large valuations. Make
$W_s$ integral at every prime where the displayed numerator has odd
valuation, and deep enough at $2,3,17$ to supply the pole bounds.
At every other prime the valuation of $b_s$ is even. These conditions
give a square reduced denominator and all pole preparation hypotheses.
The progression of Section~\ref{v15:dependencies} and the compact-real-group recurrence supplies the
chosen real neighborhoods at the same time, together with the retained
nonzero chart guards. It also permits $c$ to be arbitrarily large and
to satisfy the retained divisibility conditions for the Lucas construction.

After $S$ is fixed, continuity preserves this finite scalar preparation
and its guards for all $h$ in one fixed congruence class modulo a
sufficiently divisible integer $L$. Retain the source's integral Lucas
construction, choose the four root labels, and fix both positive scales
and compact neighborhoods for all comparison tests before selecting
contact primes. The additional contact primes avoid the finite scalar
preparation support. Hence these later congruences do not conflict with
any scalar congruence.

Choose a fixed positive integer $D_0$ clearing all contact coefficients
$D_0^2\theta/a_*$ and sufficiently divisible at the prepared primes,
including $\Sigma$, and the prime divisors of the scales. In the prime
pattern write
\[
                         \lambda=\frac{u}{D_0^2w^2}.
\]
The zero contact makes $u$ one of the selected primes, outside that
finite set. Thus $v_p(\lambda)\le-2v_p(D_0)$ there. Make this upper
bound low enough that $t_*=a_*\lambda$ has valuation below zero and
all nonzero contacts. Then $v_p(\delta)=v_p(t_*)$, and
\begin{equation}\label{v15:base248:eq:preparedsquare}
 R=\prod_{\theta\in\Theta_2}
                 \left(1+\frac{1-\theta}{\delta}\right)
       \longrightarrow1
\end{equation}
in every prepared completion. The weight $q=R^{-1}-1$ tends to zero.
Increasing the fixed denominator divisibility makes its valuation
uniformly arbitrarily large. Thus all weighted test terms are small
at every prepared place, and $\Tr(W^{-1})$ is arbitrarily close to two.

The factors in~\eqref{v15:base248:eq:preparedsquare} are arbitrarily close to one.
Requiring congruence to one modulo eight at $2$, and modulo $p$ at an odd
prepared prime $p$, makes their product a local square. The pole input
is $\gamma^2R/\Tr(W^{-1})^2$, so it is also a square at $2,3,17$.
The fixed factor has valuations
\[
 v_2(\gamma)=3,\qquad v_3(\gamma)=v_{17}(\gamma)=1.
\]
Since the reciprocal trace is close to $1/2$, the pole input has
valuations $4,2,2$ at those three primes. These satisfy all retained
preparation inequalities.

This preparation is valid for any fixed nonzero rational scales.
The prime-pattern input in Sections~\ref{v15:selection}, \ref{v15:dependencies} applies to
arbitrary fixed distinct integer coefficients. After clearing
$D_0^2\theta/a_*$ and removing repeated coefficients, it applies
here without a square restriction on $a_*$. All denominator
conditions are fixed before prime selection.

Select the contact primes with the inherited opposite-order separation.
Take $H=[h]P$, $U_i=[hn_i]P$ and $V_i=[(h-1)n_i]P$.
The common-index congruences preserve the scalar progression and make
$H,U_i$ formal at one place and $H-P,V_i$ formal at its conjugate.
Only natural depth one is needed on either branch; the source's stronger
primary preparation is also permitted. Both full logarithmic ratios
are exactly $n_i$. Plain-parameter errors therefore have order at least
four. All $C_i$, and $I$ since $n_1=\beta$, have valuation at least four,
and the Lucas equation for the chosen labels gives $v(G)\ge4$.
Every reciprocal-abscissa test has valuation at least two.
At the selected place $v_p(q)=-1$, so every nonconstant term of every
packed factor has positive valuation. Thus $W\equiv1$ there. At the
opposite place, each factor has valuation at most zero, with residue one
when its valuation is zero. The trace of $W^{-1}$ is a unit with residue
one or two. The comparison pole input has valuation one at every
selected contact prime.

The exact prime-pattern identity is
\[
 R=\frac{\displaystyle a_*^{10}\prod_{\theta\in\Theta_2}
          (u-(\theta/a_*)D_0^2w^2)}
        {(D_0^2w^2)^{10}\delta^{10}}.
\]
Every factor outside the selected prime numerators is a rational
square, including $a_*^{10}$ even when $a_*$ is not itself a square.
Thus at every prime outside those numerators the valuation is even. Multiplication by $\gamma^2$ and division by the squared trace preserve this parity. This argument includes primes dividing $\delta$;
no unit assertion about $\delta$ is made away from contacts.
Equation~\eqref{v15:base248:eq:preparedsquare} and the preparatory divisibility
give the required local squareclasses and valuation bounds at $2,3,17$.

For the real preparation, write $r=\delta^2$ and let
$\alpha_1,\ldots,\alpha_5>0$ be the five squared distances of the
contact pairs from one. Then
\[
 R=\prod_{j=1}^{5}(1-\alpha_j/r),\qquad
 q\sim\frac{\alpha_1+\cdots+\alpha_5}r>0
 \quad\text{as }r\longrightarrow+\infty.
\]
On the fixed compact real neighborhoods the tests remain bounded, so
$W\to1$ and $\Tr(W^{-1})\to2$. Consequently
\[
              \mathfrak b\longrightarrow\frac{\gamma^2}{4}
                              =\frac{41616}{167281}\in(0,17).
\]
All factor values are positive at both real embeddings for sufficiently
large $\lambda$, because $q>0$ and all test powers are even.
The simultaneous selection supplies the required real size after all
fixed data have been chosen. The pole-realization input supplies its
witnesses, and the exact power identities fill the quadratic circuit.
This proves completeness for every integer input and all sufficiently large $N$.

\subsection{Exact quadratic specification}
The local weights are specified by three scalar quadratic equations:
\[
                     q_2=q^2,\qquad q_4=q_2^2,\qquad qq_7=q_4^2.
\]
At every contact $q\ne0$, so $q_7=q^7$. If $q=0$, the contact
radical is one; its possibly free $q_7$ value is not used for extraction.
No inverse coordinate is required.

The four contact outputs are
\[
 r=(t_*-1)^2,\qquad r_2=r^2,\qquad r_3=rr_2,\qquad
 F_2(t_*)=P_5(r)=\prod_{\theta\in\Theta}(r-(\theta-1)^2).
\]
Every integer from zero through five is a sum of two elements of
$\{0,1,2,3\}$, so the last expression is quadratic in the retained
outputs. Its normalization is $F_2(t_*)(1+q)=r_2r_3$.
The scalar normalizer is the quadratic equation
$R_s+22201r_sR_s-A_Nr_s=0$.
The scalar chain is $R_2=R_s^2$, $R_4=R_2^2$.
The two scalar norm rows use $R_4$, as in~\eqref{v15:base248:eq:scalarnorm}.

The retained fourth powers of all five $B_i$ are needed by the
local factors; none is a logarithm correction. The three separately
computed shared fourth powers are $C_1^4,I^4,B_1^4$, each with two
rational coordinates. All other fourth-power coordinates and ordinary
point equations are fixed by the literal coefficient tables.
The exact count is
\[
\begin{array}{lr}\toprule
\text{point, scalar, contact, ratio, and weight coordinates}&210\\
\text{three shared fourth powers}&6\\
\text{five factors and intermediate products for two detectors}&32\\\midrule
\text{total coordinates, including the input}&248\\\midrule
\text{point, scalar, contact, ratio, and weight quadratic rows}&180\\
\text{shared fourth-power rows}&6\\
\text{factor and product rows}&36\\\midrule
\text{total quadratic rows}&222\\\bottomrule
\end{array}
\]

The complete tables fix every coordinate, equation, coefficient, and
normalization; the input is coordinate one. The producer clears
denominators and primitive-normalizes each row before squaring.
The independent checkers import no producer. The contact certificate recovers $P_5$ by
Newton finite differences at six exact arguments, checks all five roots,
and independently verifies $P_5((t-1)^2)=F_2(t)$ as a degree-ten
polynomial identity. It reconstructs every modified quadratic equation,
compares the other equations with their pinned coefficient source, and
checks the contact exceptional set and both height matrices. It also verifies the scalar pole scaling, the forced prime and its nonsquare tests, and the fourth-power extraction bounds.
The balanced normalization, its two exceptional strata, the fixed factor $\gamma$, and the linear trace scale are verified directly.
The projection proof uses these literal equations and the preceding
arithmetic argument.
It expands the entire sum of squares independently by unordered
pairs of terms. The supplementary local certificate exhausts all
odd residues of $\nu$ modulo eight and all possible valuation residues
of every test. These complete residue classes verify the
noncancellation assertion for all integer valuations, together with
the exact extraction bounds and sufficient completeness depths.
The infinite arithmetic arguments are the proofs above.

The expanded quartic has 3190 terms before powers of $A_N$ belonging
to the same coordinate monomial are collected, and 3189 monomials
after collection. Its unique coefficient involving $A_N^2$ is
$A_N^2+1$, at the square of the actual scalar contact-product coordinate.
Every other coefficient is a nonzero constant or a nonzero constant
multiple of $A_N$. Consequently no support cancellation occurs for
positive $A_N$. The largest constant coefficient $C_0$ has 1520 bits,
the largest linear multiplier is $44402$, and exact comparison
gives
\[
 C_0>44402^2,\qquad
 149^{210}+1<C_0<149^{214}+1.
\]
Thus $H(P_N)=\max(C_0,A_N^2+1)$ for every $N\ge1$ and the asserted
exact height holds for $N\ge53$. This is a coefficient statement;
it does not make the defining threshold effective.


\section{From the baseline to the 238-variable predecessor}\label{v15:reductions}
\subsection{The arithmetic change: 248 to 245}
Use the source field $K=\Q(\sqrt2)$ and retain its elliptic curve, generator $P$, and all local conventions. The source scalar comparison uses $S-P$, where $S=[a]P$. Replace it by $S$ itself. In the retained charts
\[
 S=(h_S^2,h_Sj_S),\qquad H=(\chi,\chi S_H),\qquad
 z(S)=-h_S/j_S,\quad z(H)=-1/S_H,
\]
so the new comparison is the two rational components of
\[ Zj_S=h_SS_H. \]
The finite root chart excludes $O$ and represents a nonzero point of the retained torsion-free subgroup; consequently $a\ne0$ and $j_S\ne0$. This eliminates the two shifted-point coordinates and their inverse guard: original 248-coordinate IDs 201,202,203, with original rows 135,136,177. Original rows 124,125 are replaced by the displayed two-component equation. The other 217 rows are unchanged up to compaction. There are now 245 coordinates and 219 residuals.

This step is not claimed to be a bijection between the old and new witnesses at fixed root tuple. Its soundness uses the unchanged two scalar norm tests. At a consumed odd prime with positive odd valuation $\nu$ of the scalar normalizer they give
\[
 v(F)>5\nu/4,\qquad v(Z_0)>23\nu/8,\qquad v(Z_1)>19\nu/8.
\]
The retained forced-149 exclusion is applied before taking logarithms. At a suitable place, $H$ is formal and $S$ is formal after the permitted two-torsion translation. Writing $H=[m]Q+T$ and $P=2Q$, the relation
\[ [2m]S'=[4a]H \]
kills the torsion. The same formal-logarithm isometry gives the required divisibility with the nonzero index $c=|a|$, and the retained height argument gives
\[
 h(\beta)\le C_3\,149^{-L_N}c(1+\log(c)/3).
\]
For sufficiently large $c$, the source's integral Lucas calculation in $|a|$ supplies the stronger lower bound $(c-4)\log(c-1)$ for $\log B$. Since $c\ge149^{L_N}$ and $L_N\to\infty$, increasing the single ineffective $N_0$ ensures this size condition uniformly for every $N\ge N_0$.

For completeness, choose $H$ with the retained finite depths and real bounds, then choose a positive sufficiently divisible $a$ with $[a]P$ near $O$ in the real compact group. Recurrence in the required finite-index subgroup gives this while retaining finite local conditions. The ratio is bounded near $O$. For every such $a\ge3$, put $U_0=0,U_1=1,U_{j+1}=aU_j-U_{j-1}$. Choose positive $j$ divisible by $a-2$, set $b=U_j/(a-2)$ and $t=U_{j+1}$. Then $U_j\equiv j\pmod{a-2}$ and
\[t^2-a(a-2)tb+(a-2)^2b^2=1.\]
The remaining contact selection is the retained source construction. In particular no condition $a\equiv1$ is introduced.

\subsection{Root-product charts: 245 to 240}
Put $d=3-\sqrt2$ and $D=d^2$. For each of the four source root blocks the old field coordinates satisfy $x=h^2$, $j^2=x^2-D$ and
\begin{align*}
hI&=jV+DVW+xQ,&hK_1&=-jQ-xV-DQW,\\
VQ&=xW,&V^2+Q^2+2jW+DW^2&=1.
\end{align*}
Here $F_U$ and $F_V$ are the retained reciprocal-abscissa-square tests of the associated points $U,V$: the source ratio equations give $I^2=x(1-DF_U)$ and $K_1^2=x(1-DF_V)$. The retained $z_V$ is the plain formal coordinate $-x(V)/y(V)$, with the source zero convention at the identity. The primary comparison $L$ satisfies $V S_H=I L$, so it represents the original primary formal-coordinate ratio wherever the denominator is nonzero, with the same cleared equation at zero. These are existing consumers, not newly imposed nonvanishing assumptions. Replace these eight field coordinates by seven $A,B,C,E,i,k,u$, imposing
\begin{gather*}
 A^2+B^2+2CE=0,\quad AC=d(i-k),\quad BE=d(i+k),\\
 (A+B)^2=4Du,\quad F_U=u^2,\quad i^2+DF_U=1,\quad k^2+DF_V=1,\\
 A-B+2dkz_V=0,\qquad (A+B)S_H=2diL.
\end{gather*}
The last two equations replace the old reciprocal and primary ratio equations. Their number is unchanged. Take rational components with the explicit conventions of Section~\ref{v15:component-typing}.

The forward map is
\begin{gather*}
 A=d(V+Q)/h,\quad B=d(V-Q)/h,\quad C=(x+j+DW)/h,\\
 E=C-2h,\quad i=I/h,\quad k=K_1/h,\quad u=(A+B)^2/(4D).
\end{gather*}
Here $h\ne0$ follows from $j^2=h^4-D$ at a real embedding. The inverse is
\begin{gather*}
 h=(C-E)/2,\quad x=h^2,\quad j=(B^2+C^2-A^2-E^2)/4,\\
 V=h(A+B)/(2d),\quad Q=h(A-B)/(2d),\quad
 I=hi,\quad K_1=hk,\quad W=(A^2-B^2)/(4D).
\end{gather*}
For verification put $\alpha=(B^2+C^2)/2$ and $\beta_1=(A^2+E^2)/2$. The new equations give $\alpha+\beta_1=2h^2$ and $\alpha\beta_1=D$, using
\[(B^2+C^2)(A^2+E^2)=(AB-CE)^2+(AC+BE)^2.\]
Thus $j^2=h^4-D$ and $h\ne0$. Also $j\ne0$, since $h^4=d^2$ would force $d$ to be a square in $K$, contrary to $\operatorname{Norm}(d)=7$. Substitution proves both compositions and all retained equations. No division by $A\pm B,i,k,C,E$ is made, so their zero strata survive.

Four blocks save eight rational coordinates. Retaining the two components of $j_S$ and adding one rational trace output $T=\tfrac14\operatorname{Tr}_{K/\Q}\sum h_\ell^2$ gives $245-8+2+1=240$, with 214 residuals. The scalar ratio uses $h_S=(C_S-E_S)/2$ and the retained quadratic expression for $j_S$.

\subsection{Reciprocal scale: 240 to 239}
The remaining scale equations are $T=\tfrac14\operatorname{Tr}\sum h_\ell^2$ and $t_1=\lambda T$. Replace $\lambda$ by $\rho=1/\lambda$, delete $T$, and impose $\rho t_1=\tfrac14\operatorname{Tr}\sum h_\ell^2$. The right side is positive, so $\rho\ne0$ without a new guard.

For the exact source contact polynomial $F$, set $G(\rho)=\rho^{10}F(1/\rho)$. The cleared weight equation becomes
\[ G(\rho)(1+q)=(1-\rho)^{10}. \]
Three existing power coordinates now represent $\rho^2,\rho^4,\rho^5$. Degrees zero through ten have quadratic representatives in $1,\rho,\rho^2,\rho^4,\rho^5$, so both displayed polynomials remain quadratic in the expanded coordinates. The exact coefficients are the pinned 239 IR, not an independently chosen contact polynomial. The old equation excludes $\lambda=0$ since $F(0)=0$. Multiplication or division by $\rho^{10}$ proves the inverse maps. Values $\rho=1$ and $q=0$ are not discarded. The count is 239 coordinates and 213 residuals.

\subsection{Global scalar rescaling: 239 to 238}
This section uses the 239 IR indices, before deleting coordinate 140. Write $R=x_{143}$; coordinates 139 and 140 represent $R^2,R^4$. For every other index impose $x_i=R^{w_i}y_{\mu(i)}$, where $\mu(i)=i$ for $i<140$ and $\mu(i)=i-1$ for $i>140$. All unspecified weights are zero. The nonzero ordinary-coordinate weights are
\begin{align*}
w_i=-1:&\quad i\in[27,32]\cup[55,66]\cup[68,71]\cup[184,195]\\
 &\qquad\cup[198,205]\cup[212,219]\cup[226,233],\\
w_i=-2:&\quad i\in\{2,3,110,111,112,130,138,145\},\\
w_{113}=-4,&\qquad w_{128}=w_{129}=1,\qquad w_{135}=2.
\end{align*}
Intervals here include all integer indices. Retain $y_{139}=R^2$ with the explicit equation $y_{139}-y_{142}^2=0$. Delete old coordinate 140 and only its defining row 87. This equation must not be replaced by a tautology when powers are collected.

For each remaining row, the row-multiplier table below specifies the integer $k$ for which its substituted old residual equals $R^k$ times the new residual. Powers are represented using $1,R,R^2,R\cdot R^2,(R^2)^2$, subject to total degree two. The complete 212-row output is the specified final IR.

The zero-domain check is independent of the rescaled charts. Old rows 25,72--74,92 remain unchanged. If $z=X^2\ge0$, their scalar contact product is
\begin{align*}
r={}&260105411931107592+149961348584209227z\\
 &+11129085312521652z^2+328736168366319z^3\\
 &+4385243879406z^4+22081616064z^5>0.
\end{align*}
The normalizer equation is $(1+22201r)R=A_Nr$, so $R>0$. The forward map divides by the indicated powers of $R$; the inverse multiplies by them and restores $R^4$. Every division is therefore legitimate on both solution sets. Together with the row identities this proves a fixed-input, fixed-parameter rational witness bijection. The count is 238 coordinates and 212 residuals.

\subsection{Sparse presentation of the 238-variable predecessor}
The preceding section describes the raw 238-coordinate presentation. The predecessor presentation is the independently accepted composition in \texttt{composition238-007}. Apply its affine coordinate stages to the raw system: the Lucas $P/E$ basis, four constant shifts, five non-Lucas shears, and the four root sum/difference bases. In each root block the old slots become
\[
 A=(X_1+Y_1)/2,\quad B=(X_1-Y_1)/2,\quad
 C=(J_1+H_1)/2,\quad E=(J_1-H_1)/2.
\]
Here $X_1$ denotes a block coordinate, not the distinguished input. Every consumer is transformed, including the scalar and aggregate equations. The intermediate sequence also makes 19 invertible row additions, replaces the two root component rows by their sum and difference, and applies recorded nonzero constant primitive normalizations. No variable-dependent divisor is introduced. These are applied in the exact order in \path{EXACT-MAPS.json}, rather than presumed to commute with the preceding $R$ rescaling.

For each coordinate stage, that file states old coordinates as affine polynomials in new ones; unspecified coordinates are fixed. A raw solution is sent to the final system by applying the inverse coordinate stages in forward stage order. Recover the raw solution by applying the forward stages in reverse order. \path{INVERSE-MAPS.json} records these inverses. Row additions are undone by subtraction, sum/difference by half-sum/half-difference, and primitive normalization by its recorded nonzero rational multiplier. Thus simultaneous vanishing is equivalent in both directions with the input and $A_N$ fixed. All 238 variables occur.

The exact final definitions are the 212 equations in Section~\ref{v15:equations}; the coordinate and row transformations are specified below. The companion files are reproducibility aids only. All 212 explicit quadratic equations are listed in the appendix, in the IR order. Its table specifies every large integer coefficient. Squaring and collecting these final equations yields 2342 coordinate monomials (2343 parameter-separated terms). The sparse composition changes the polynomial's coefficients, but not its rational zero-set projection onto the input. Its independent review is identified by the prefix \texttt{07f3b18c}; the complete receipt is included in the delivery evidence.


\subsection{Field components and scalar gate meanings}\label{v15:component-typing}
Here is an explicit interpretation of the raw 239-coordinate dictionary, before the signed $R$ rescaling and the subsequent sparse affine changes. Write its coordinates as $x_i$ and put $\epsilon=\sqrt2$. The following formulas specify the component conventions; a suffix in a historical coordinate name is not itself a convention.
\begin{align*}
j_S&=x_2+x_3\epsilon,\qquad\rho=x_4,\qquad X=x_1,\\
x(H)&=x_{17}+2x_{18}\epsilon,\\
f(H)&=x_{19}+x_{20}\epsilon,\\
S_H=-1/z_H&=x_{21}+x_{22}\epsilon,\\
f(V_0)&=x_{23}+x_{24}\epsilon,\\
z_{V_n}&=x_{23+2n}+x_{24+2n}\epsilon\quad(1\le n\le4),\\
L_n&=x_{47+4n}+x_{48+4n}\epsilon,\\
M'_n&=x_{49+4n}+x_{50+4n}\epsilon\quad(1\le n\le4),\\
Z&=x_{128}+x_{129}\epsilon,\\
J_0&=x_{141}+x_{142}\epsilon.
\end{align*}
Here $J_0$ is the retained reciprocal auxiliary, with $J_0^2=f(H)f(V_0)$; its two cleared component equations are retained in the printed system. The raw pair $M'_n$ is the unconjugated second ratio; the baseline notation is $L'_n=\overline{M'_n}$. These primes do not denote differentiation. With $D=(3-\epsilon)^2=11-6\epsilon$, the stored square components are
\begin{align*}
f(V_0)^2&=x_{33}+2x_{34}\epsilon,&
f(V_1)^2&=x_{37}+2x_{38}\epsilon,\\
f(V_2)^2&=x_{39}+2x_{40}\epsilon,&
f(V_3)^2&=x_{43}+2x_{44}\epsilon,&
f(V_4)^2&=x_{47}+2x_{48}\epsilon.
\end{align*}
For $(a,b)=(33,35),(39,41),(43,45),(47,49)$, the relation is
\[D^2(x_a+2x_{a+1}\epsilon)^2=193x_b-132x_{b+1}\epsilon.\]
This follows by expanding $D^2=193-132\epsilon$, and records both exceptional multipliers. Likewise, for $n=1,2,3,4$,
\[f(U_n)^2=x_{72+2n}+2x_{73+2n}\epsilon.\]
In particular these are the $F_U$ inputs to the root equations; the $F_V$ inputs are the corresponding four squares of $V_1,\ldots,V_4$ just displayed. The root variables $A,B,C,E,i,k,u$ for the $n$th root block have ordinary components starting at $184+14(n-1)$: the successive pairs have offsets $0,2,4,6,8,10,12$. Thus $h=(C-E)/2$, the specified root expression for $j$, $f(U)=u$, and every root equation are literal equations in $\mathbf Q(\epsilon)$ with no further component factors.

The remaining special square conventions are
\begin{align*}
(L_1-\overline{M'_1})^2&=x_{82}+2x_{83}\epsilon,\\
(L_1-(51X^2+2501))^2&=x_{84}+2x_{85}\epsilon,\\
f(H)^2&=x_{86}+2x_{87}\epsilon,\\
(L_n-\overline{M'_n})^2&=x_{84+2n}+2x_{85+2n}\epsilon\quad(2\le n\le4).
\end{align*}
The Lucas products have ordinary components $L_2L_4=x_{68}+x_{69}\epsilon$ and $L_2L_3=x_{70}+x_{71}\epsilon$. The raw Lucas quantity is $F=G+1=x_{72}+2x_{73}\epsilon$, where $G$ is the baseline Lucas defect, and hence $F^2=x_{94}+4x_{95}\epsilon$; the defect square is recovered as $G^2=F^2-2F+1$. Put $\beta=51X^2+2504$ and $I=L_1-\beta$. Thus the square at $84$ is $(I+3)^2$, not $I^2$. The fourth-power pairs at $146,148$ are ordinary components of $(L_1-\overline{M'_1})^4$ and $I^4$, respectively; indeed $I^2=(I+3)^2-6(I+3)+9$. The pair at $150$ is instead the ordinary pair of $\overline{f(V_1)^4}$, as is seen from $x_{150}=x_{37}^2+8x_{38}^2$ and $x_{151}=-4x_{37}x_{38}$. All detector-factor and detector-product pairs, and the two actual-$W$ pairs at $103$ and $118$, use ordinary components. Their defining equations in the final system, transported backwards by the displayed maps, specify their signs. No imaginary multiplier is inferred from their names.

For completeness the scalar singletons that occur in the rescaling argument have the following literal meanings. Put $R=x_{143}$, $s=X^2=x_{67}$ and $K=(2^2\cdot3\cdot5\cdot11\cdot109\cdot439)^{14}$. Then
\begin{align*}
x_{125}&=(51s+2501)^2/4,\\
1525940x_{126}&=2386189056485+97317586470s+992242485s^2\\
&\quad+15531948s x_{125}+761674548x_{125}-512x_{125}^2,\\
x_{127}&=-71677797244068-2921180304792s-29762681184s^2\\
&\quad-429432240s x_{125}-155645880s x_{126}\\
&\quad-21080070840x_{125}-7640381580x_{126},\\
x_{139}&=R^2,\qquad x_{140}=R^4,\qquad x_{135}=x_{129}^2,\\
x_{136}R^4&=x_{128}^2+Rx_{135},\\
x_{137}&=x_{86}^2-8x_{87}^2=\operatorname{Norm}(f(H))^2,\\
x_{138}R^4&=x_{136}^2+Rx_{137},\\
Kx_{130}&=x_{138}^2,\\
82479248x_{131}&=-x_{130}(R^2+9792R),\\
(22201x_{127}+1)R&=A_Nx_{127}.
\end{align*}
These are polynomial identities even at a zero input; division by $R$ is used only after the positivity argument. The other scalar square auxiliaries are the ordinary squares indicated by their defining rows. Thus the backward reconstruction from the final 212 equations consists of the explicit inverse sparse affine maps, reinsertion of $R^4$ and the signed-power inverse map, followed by the component formulas above. It requires neither an external coordinate file nor an assumption that every stored pair has identical scaling.

\subsection{Index dictionary for the rescaling proof}\label{v15:dictionary}
The following raw 239-coordinate dictionary fixes the indices in the weight table. The field interpretation and every exceptional stored-component multiplier are specified in Section~\ref{v15:component-typing}. The final sparse stage subsequently changes these positions by its explicit affine formulas.
\begin{longtable}{r p{0.78\textwidth}}
Index&Coordinate label\\\hline
1&\path{input}\\
2&\path{S.chart.j0}\\
3&\path{S.chart.j1}\\
4&\path{reciprocal.rho}\\
5&\path{detector0.normalized.scaled_e}\\
6&\path{detector0.normalized.s}\\
7&\path{detector0.normalized.p}\\
8&\path{detector0.factored.q}\\
9&\path{detector1.normalized.scaled_e}\\
10&\path{detector1.normalized.s}\\
11&\path{detector1.normalized.p}\\
12&\path{detector1.factored.q}\\
13&\path{normalized_scalar.scaled_e}\\
14&\path{normalized_scalar.s}\\
15&\path{normalized_scalar.p}\\
16&\path{normalized_scalar.factored.q}\\
17&\path{H.chart.x.0.quadratic}\\
18&\path{H.chart.x.1.quadratic}\\
19&\path{U0.f.0}\\
20&\path{U0.f.1}\\
21&\path{primary.negativeReciprocalZ0.0}\\
22&\path{primary.negativeReciprocalZ0.1}\\
23&\path{V0.f.0}\\
24&\path{V0.f.1}\\
25&\path{V1.z.0}\\
26&\path{V1.z.1}\\
27&\path{V2.z.0}\\
28&\path{V2.z.1}\\
29&\path{V3.z.0}\\
30&\path{V3.z.1}\\
31&\path{V4.z.0}\\
32&\path{V4.z.1}\\
33&\path{logV0.fSquare.0.quadratic}\\
34&\path{logV0.fSquare.1.quadratic}\\
35&\path{logV0.wPower2.0.quadratic}\\
36&\path{logV0.wPower2.1.quadratic}\\
37&\path{logV1.fSquare.0.quadratic}\\
38&\path{logV1.fSquare.1.quadratic}\\
39&\path{logV2.fSquare.0.quadratic}\\
40&\path{logV2.fSquare.1.quadratic}\\
41&\path{logV2.wPower2.0.quadratic}\\
42&\path{logV2.wPower2.1.quadratic}\\
43&\path{logV3.fSquare.0.quadratic}\\
44&\path{logV3.fSquare.1.quadratic}\\
45&\path{logV3.wPower2.0.quadratic}\\
46&\path{logV3.wPower2.1.quadratic}\\
47&\path{logV4.fSquare.0.quadratic}\\
48&\path{logV4.fSquare.1.quadratic}\\
49&\path{logV4.wPower2.0.quadratic}\\
50&\path{logV4.wPower2.1.quadratic}\\
51&\path{L1.0}\\
52&\path{L1.1}\\
53&\path{LP1.0}\\
54&\path{LP1.1}\\
55&\path{L2.0}\\
56&\path{L2.1}\\
57&\path{LP2.0}\\
58&\path{LP2.1}\\
59&\path{L3.0}\\
60&\path{L3.1}\\
61&\path{LP3.0}\\
62&\path{LP3.1}\\
63&\path{L4.0}\\
64&\path{L4.1}\\
65&\path{LP4.0}\\
66&\path{LP4.1}\\
67&\path{inputSquare.quadratic}\\
68&\path{Lucas.C.0.quadratic}\\
69&\path{Lucas.C.1.quadratic}\\
70&\path{Lucas.AT.0.quadratic}\\
71&\path{Lucas.AT.1.quadratic}\\
72&\path{Lucas.F0.quadratic}\\
73&\path{Lucas.F1.quadratic}\\
74&\path{testSquare.1.0.0.quadratic}\\
75&\path{testSquare.1.0.1.quadratic}\\
76&\path{testSquare.1.1.0.quadratic}\\
77&\path{testSquare.1.1.1.quadratic}\\
78&\path{testSquare.1.2.0.quadratic}\\
79&\path{testSquare.1.2.1.quadratic}\\
80&\path{testSquare.1.3.0.quadratic}\\
81&\path{testSquare.1.3.1.quadratic}\\
82&\path{testSquare.12.0.0.quadratic}\\
83&\path{testSquare.12.0.1.quadratic}\\
84&\path{testSquare.12.1.0.quadratic}\\
85&\path{testSquare.12.1.1.quadratic}\\
86&\path{testSquare.40.0.0.quadratic}\\
87&\path{testSquare.40.0.1.quadratic}\\
88&\path{testSquare.40.1.0.quadratic}\\
89&\path{testSquare.40.1.1.quadratic}\\
90&\path{testSquare.40.2.0.quadratic}\\
91&\path{testSquare.40.2.1.quadratic}\\
92&\path{testSquare.40.3.0.quadratic}\\
93&\path{testSquare.40.3.1.quadratic}\\
94&\path{testSquare.40.4.0.quadratic}\\
95&\path{testSquare.40.4.1.quadratic}\\
96&\path{reciprocal.rhoSquared}\\
97&\path{reciprocal.rhoFourth}\\
98&\path{reciprocal.contactPolynomial}\\
99&\path{detector0.q1}\\
100&\path{reflected.detector0.weightPower7}\\
101&\path{detector0.E1.power2.quadratic}\\
102&\path{detector0.E1.power4.quadratic}\\
103&\path{detector0.actualW.0}\\
104&\path{detector0.actualW.1}\\
105&\path{detector0.normOver2w0}\\
106&\path{detector0.B}\\
107&\path{detector0.normalized.p2.quadratic}\\
108&\path{detector0.factored.v}\\
109&\path{detector0.factored.q2}\\
110&\path{detector1.t.quadratic}\\
111&\path{symmetric.detector1.deltaSquared}\\
112&\path{symmetric.detector1.deltaFourthPower}\\
113&\path{symmetric.detector1.contactPolynomial}\\
114&\path{detector1.q1}\\
115&\path{reflected.detector1.weightPower7}\\
116&\path{detector1.E1.power2.quadratic}\\
117&\path{detector1.E1.power4.quadratic}\\
118&\path{detector1.actualW.0}\\
119&\path{detector1.actualW.1}\\
120&\path{detector1.normOver2w0}\\
121&\path{detector1.B}\\
122&\path{detector1.normalized.p2.quadratic}\\
123&\path{detector1.factored.v}\\
124&\path{detector1.factored.q2}\\
125&\path{scalar.rPair1.quadratic}\\
126&\path{scalar.rFour.quadratic}\\
127&\path{scalar.contactProduct}\\
128&\path{scalar.zRatio.0}\\
129&\path{scalar.zRatio.1}\\
130&\path{scalar.W2.quadratic}\\
131&\path{scalar.b.quadratic}\\
132&\path{normalized_scalar.p2.quadratic}\\
133&\path{normalized_scalar.factored.v}\\
134&\path{normalized_scalar.factored.q2}\\
135&\path{scalar.ratio1_unscaled_square}\\
136&\path{reduced_scalar.ratio_node}\\
137&\path{reduced_scalar.H_norm_squared}\\
138&\path{reduced_scalar.root}\\
139&\path{scalar.r_power_2}\\
140&\path{scalar.r_power_4}\\
141&\path{V0.reciprocal.gOverH.0}\\
142&\path{V0.reciprocal.gOverH.1}\\
143&\path{scalar.forced_positive_parameter}\\
144&\path{reciprocal.rhoFifth}\\
145&\path{reflected.detector1.deltaSixthPower}\\
146&\path{amplified.C1.power4.0}\\
147&\path{amplified.C1.power4.1}\\
148&\path{residue.I.power4.0}\\
149&\path{residue.I.power4.1}\\
150&\path{residue.V1.power4.0}\\
151&\path{residue.V1.power4.1}\\
152&\path{amplified.det0.factor0.0}\\
153&\path{amplified.det0.factor0.1}\\
154&\path{amplified.det0.factor1.0}\\
155&\path{amplified.det0.factor1.1}\\
156&\path{amplified.det0.factor2.0}\\
157&\path{amplified.det0.factor2.1}\\
158&\path{amplified.det0.factor3.0}\\
159&\path{amplified.det0.factor3.1}\\
160&\path{amplified.det0.factor4.0}\\
161&\path{amplified.det0.factor4.1}\\
162&\path{amplified.det0.product1.0}\\
163&\path{amplified.det0.product1.1}\\
164&\path{amplified.det0.product2.0}\\
165&\path{amplified.det0.product2.1}\\
166&\path{amplified.det0.product3.0}\\
167&\path{amplified.det0.product3.1}\\
168&\path{amplified.det1.factor0.0}\\
169&\path{amplified.det1.factor0.1}\\
170&\path{amplified.det1.factor1.0}\\
171&\path{amplified.det1.factor1.1}\\
172&\path{amplified.det1.factor2.0}\\
173&\path{amplified.det1.factor2.1}\\
174&\path{amplified.det1.factor3.0}\\
175&\path{amplified.det1.factor3.1}\\
176&\path{amplified.det1.factor4.0}\\
177&\path{amplified.det1.factor4.1}\\
178&\path{amplified.det1.product1.0}\\
179&\path{amplified.det1.product1.1}\\
180&\path{amplified.det1.product2.0}\\
181&\path{amplified.det1.product2.1}\\
182&\path{amplified.det1.product3.0}\\
183&\path{amplified.det1.product3.1}\\
184&\path{rootProduct.1.A.0}\\
185&\path{rootProduct.1.A.1}\\
186&\path{rootProduct.1.B.0}\\
187&\path{rootProduct.1.B.1}\\
188&\path{rootProduct.1.C.0}\\
189&\path{rootProduct.1.C.1}\\
190&\path{rootProduct.1.E.0}\\
191&\path{rootProduct.1.E.1}\\
192&\path{rootProduct.1.i.0}\\
193&\path{rootProduct.1.i.1}\\
194&\path{rootProduct.1.k.0}\\
195&\path{rootProduct.1.k.1}\\
196&\path{rootProduct.1.u.0}\\
197&\path{rootProduct.1.u.1}\\
198&\path{rootProduct.2.A.0}\\
199&\path{rootProduct.2.A.1}\\
200&\path{rootProduct.2.B.0}\\
201&\path{rootProduct.2.B.1}\\
202&\path{rootProduct.2.C.0}\\
203&\path{rootProduct.2.C.1}\\
204&\path{rootProduct.2.E.0}\\
205&\path{rootProduct.2.E.1}\\
206&\path{rootProduct.2.i.0}\\
207&\path{rootProduct.2.i.1}\\
208&\path{rootProduct.2.k.0}\\
209&\path{rootProduct.2.k.1}\\
210&\path{rootProduct.2.u.0}\\
211&\path{rootProduct.2.u.1}\\
212&\path{rootProduct.3.A.0}\\
213&\path{rootProduct.3.A.1}\\
214&\path{rootProduct.3.B.0}\\
215&\path{rootProduct.3.B.1}\\
216&\path{rootProduct.3.C.0}\\
217&\path{rootProduct.3.C.1}\\
218&\path{rootProduct.3.E.0}\\
219&\path{rootProduct.3.E.1}\\
220&\path{rootProduct.3.i.0}\\
221&\path{rootProduct.3.i.1}\\
222&\path{rootProduct.3.k.0}\\
223&\path{rootProduct.3.k.1}\\
224&\path{rootProduct.3.u.0}\\
225&\path{rootProduct.3.u.1}\\
226&\path{rootProduct.4.A.0}\\
227&\path{rootProduct.4.A.1}\\
228&\path{rootProduct.4.B.0}\\
229&\path{rootProduct.4.B.1}\\
230&\path{rootProduct.4.C.0}\\
231&\path{rootProduct.4.C.1}\\
232&\path{rootProduct.4.E.0}\\
233&\path{rootProduct.4.E.1}\\
234&\path{rootProduct.4.i.0}\\
235&\path{rootProduct.4.i.1}\\
236&\path{rootProduct.4.k.0}\\
237&\path{rootProduct.4.k.1}\\
238&\path{rootProduct.4.u.0}\\
239&\path{rootProduct.4.u.1}\\
\end{longtable}
\paragraph{Row multipliers for global rescaling.} The omitted raw row is87. For every other raw row $i$, the substituted residual is $R^{k_i}$ times new row $j$. Zero powers are omitted from this table; row order is unchanged apart from deletion87.
\begin{longtable}{rrr}
$i$ & $j$ & $k_i$\\\hline
15&15&-1\\
16&16&-1\\
19&19&-1\\
20&20&-1\\
23&23&-1\\
24&24&-1\\
26&26&-2\\
27&27&-2\\
28&28&-2\\
29&29&-2\\
30&30&-2\\
31&31&-2\\
38&38&-2\\
39&39&-2\\
40&40&-2\\
41&41&-2\\
42&42&-2\\
43&43&-2\\
59&59&-4\\
60&60&-4\\
61&61&-4\\
62&62&-4\\
75&75&-4\\
76&76&-1\\
88&87&2\\
89&88&2\\
94&93&-4\\
137&136&-2\\
138&137&-2\\
139&138&-2\\
140&139&-2\\
141&140&-2\\
142&141&-2\\
143&142&-2\\
144&143&-2\\
147&146&-2\\
148&147&-2\\
149&148&-2\\
150&149&-2\\
151&150&-1\\
152&151&-1\\
153&152&-1\\
154&153&-1\\
155&154&-2\\
156&155&-2\\
157&156&-2\\
158&157&-2\\
159&158&-2\\
160&159&-2\\
161&160&-2\\
162&161&-2\\
169&168&-1\\
170&169&-1\\
171&170&-1\\
172&171&-1\\
173&172&-2\\
174&173&-2\\
175&174&-2\\
176&175&-2\\
177&176&-2\\
178&177&-2\\
179&178&-2\\
180&179&-2\\
187&186&-1\\
188&187&-1\\
189&188&-1\\
190&189&-1\\
191&190&-2\\
192&191&-2\\
193&192&-2\\
194&193&-2\\
195&194&-2\\
196&195&-2\\
197&196&-2\\
198&197&-2\\
205&204&-1\\
206&205&-1\\
207&206&-1\\
208&207&-1\\
209&208&-2\\
210&209&-2\\
211&210&-1\\
212&211&-1\\
213&212&-2\\
\end{longtable}

\subsection{Explicit ordered sparse transformations}\label{v15:maps}
At each coordinate stage $z$ on the left denotes the old coordinates and $y$ on the right the new coordinates; unlisted coordinates are fixed. After the stage rename $y$ to $z$. Row indices below are one-based. Apply stages in order.
\paragraph{Stage 1.} Lucas P/E basis.
\begin{gather*}
z_{68}=(2)y_{63}+(1)y_{68}\\
z_{69}=(2)y_{64}+(1)y_{69}\\
z_{70}=(1)y_{68}+(1)y_{70}\\
z_{71}=(1)y_{69}+(1)y_{71}
\end{gather*}
\paragraph{Stage 2.} $Q_{96}\leftarrow Q_{96}+(-36)Q_{34}$.
\paragraph{Stage 3.} $Q_{11}\leftarrow Q_{11}+(-23/104)Q_{12}$.
\paragraph{Stage 4.} $Q_{12}\leftarrow Q_{12}+(-2392/4879)Q_{11}$.
\paragraph{Stage 5.} $Q_{15}\leftarrow Q_{15}+(-23/104)Q_{16}$.
\paragraph{Stage 6.} $Q_{16}\leftarrow Q_{16}+(-2392/4879)Q_{15}$.
\paragraph{Stage 7.} $Q_{19}\leftarrow Q_{19}+(-23/104)Q_{20}$.
\paragraph{Stage 8.} $Q_{20}\leftarrow Q_{20}+(-2392/4879)Q_{19}$.
\paragraph{Stage 9.} $Q_{23}\leftarrow Q_{23}+(-52/23)Q_{24}$.
\paragraph{Stage 10.} $Q_{24}\leftarrow Q_{24}+(2392/4879)Q_{23}$.
\paragraph{Stage 11.} $Q_{73}\leftarrow Q_{73}+(-381485)Q_{72}$.
\paragraph{Stage 12.} $Q_{89}\leftarrow Q_{89}+(-56/57)Q_{90}$.
\paragraph{Stage 13.} $Q_{90}\leftarrow Q_{90}+(-1596/1681)Q_{89}$.
\paragraph{Stage 14.} $Q_{97}\leftarrow Q_{97}+(-72)Q_{35}$.
\paragraph{Stage 15.} $Q_{9}\leftarrow Q_{9}+(-193/33)Q_{10}$.
\paragraph{Stage 16.} $Q_{13}\leftarrow Q_{13}+(-193/33)Q_{14}$.
\paragraph{Stage 17.} $Q_{17}\leftarrow Q_{17}+(-193/33)Q_{18}$.
\paragraph{Stage 18.} $Q_{21}\leftarrow Q_{21}+(-193/33)Q_{22}$.
\paragraph{Stage 19.} $Q_{48}\leftarrow Q_{48}+(3813373161796)Q_{49}$.
\paragraph{Stage 20.} $Q_{61}\leftarrow Q_{61}+(3813373161796)Q_{62}$.
\paragraph{Stage 21.} four constant shifts.
\begin{gather*}
z_{140}=(-1672/18921)+(1)y_{140}\\
z_{67}=(-2501/51)+(1)y_{67}\\
z_{84}=(-9)+(1)y_{84}\\
z_{94}=(-1)+(1)y_{94}
\end{gather*}
\paragraph{Stage 22.} non-Lucas shear.
\begin{gather*}
z_{35}=(1)y_{35}+(1)y_{36}
\end{gather*}
\paragraph{Stage 23.} non-Lucas shear.
\begin{gather*}
z_{41}=(1)y_{41}+(1)y_{42}
\end{gather*}
\paragraph{Stage 24.} non-Lucas shear.
\begin{gather*}
z_{45}=(1)y_{45}+(1)y_{46}
\end{gather*}
\paragraph{Stage 25.} non-Lucas shear.
\begin{gather*}
z_{49}=(1)y_{49}+(1)y_{50}
\end{gather*}
\paragraph{Stage 26.} non-Lucas shear.
\begin{gather*}
z_{126}=(1)y_{125}+(1)y_{126}
\end{gather*}
\paragraph{Stage 27.} four root sum/difference bases.
\begin{gather*}
z_{183}=(1/2)y_{183}+(1/2)y_{185}\\
z_{185}=(1/2)y_{183}+(-1/2)y_{185}\\
z_{187}=(1/2)y_{187}+(1/2)y_{189}\\
z_{189}=(1/2)y_{187}+(-1/2)y_{189}\\
z_{184}=(1/2)y_{184}+(1/2)y_{186}\\
z_{186}=(1/2)y_{184}+(-1/2)y_{186}\\
z_{188}=(1/2)y_{188}+(1/2)y_{190}\\
z_{190}=(1/2)y_{188}+(-1/2)y_{190}\\
z_{197}=(1/2)y_{197}+(1/2)y_{199}\\
z_{199}=(1/2)y_{197}+(-1/2)y_{199}\\
z_{201}=(1/2)y_{201}+(1/2)y_{203}\\
z_{203}=(1/2)y_{201}+(-1/2)y_{203}\\
z_{198}=(1/2)y_{198}+(1/2)y_{200}\\
z_{200}=(1/2)y_{198}+(-1/2)y_{200}\\
z_{202}=(1/2)y_{202}+(1/2)y_{204}\\
z_{204}=(1/2)y_{202}+(-1/2)y_{204}\\
z_{211}=(1/2)y_{211}+(1/2)y_{213}\\
z_{213}=(1/2)y_{211}+(-1/2)y_{213}\\
z_{215}=(1/2)y_{215}+(1/2)y_{217}\\
z_{217}=(1/2)y_{215}+(-1/2)y_{217}\\
z_{212}=(1/2)y_{212}+(1/2)y_{214}\\
z_{214}=(1/2)y_{212}+(-1/2)y_{214}\\
z_{216}=(1/2)y_{216}+(1/2)y_{218}\\
z_{218}=(1/2)y_{216}+(-1/2)y_{218}\\
z_{225}=(1/2)y_{225}+(1/2)y_{227}\\
z_{227}=(1/2)y_{225}+(-1/2)y_{227}\\
z_{229}=(1/2)y_{229}+(1/2)y_{231}\\
z_{231}=(1/2)y_{229}+(-1/2)y_{231}\\
z_{226}=(1/2)y_{226}+(1/2)y_{228}\\
z_{228}=(1/2)y_{226}+(-1/2)y_{228}\\
z_{230}=(1/2)y_{230}+(1/2)y_{232}\\
z_{232}=(1/2)y_{230}+(-1/2)y_{232}
\end{gather*}
\paragraph{Stage 28.} $ (Q_{138},Q_{140})\leftarrow(Q_{138}+Q_{140},Q_{138}-Q_{140})$ simultaneously.
\paragraph{Stage 29.} $ (Q_{139},Q_{141})\leftarrow(Q_{139}+Q_{141},Q_{139}-Q_{141})$ simultaneously.
\paragraph{Stage 30.} $ (Q_{156},Q_{158})\leftarrow(Q_{156}+Q_{158},Q_{156}-Q_{158})$ simultaneously.
\paragraph{Stage 31.} $ (Q_{157},Q_{159})\leftarrow(Q_{157}+Q_{159},Q_{157}-Q_{159})$ simultaneously.
\paragraph{Stage 32.} $ (Q_{174},Q_{176})\leftarrow(Q_{174}+Q_{176},Q_{174}-Q_{176})$ simultaneously.
\paragraph{Stage 33.} $ (Q_{175},Q_{177})\leftarrow(Q_{175}+Q_{177},Q_{175}-Q_{177})$ simultaneously.
\paragraph{Stage 34.} $ (Q_{192},Q_{194})\leftarrow(Q_{192}+Q_{194},Q_{192}-Q_{194})$ simultaneously.
\paragraph{Stage 35.} $ (Q_{193},Q_{195})\leftarrow(Q_{193}+Q_{195},Q_{193}-Q_{195})$ simultaneously.
\paragraph{Primitive normalization.} Finally multiply row $i$ by the following factor $f_i$; unlisted factors are one. These nonzero constants do not change simultaneous vanishing.
\begin{longtable}{rr}
$i$&$f_i$\\\hline
5&$358004241$\\
6&$18921$\\
9&$33$\\
11&$5512$\\
12&$258587/104$\\
13&$33$\\
15&$5512$\\
16&$258587/104$\\
17&$33$\\
19&$5512$\\
20&$258587/104$\\
21&$33$\\
23&$1219$\\
24&$258587/23$\\
25&$51$\\
73&$1/4$\\
89&$1078497$\\
90&$10602067/19$\\
136&$2$\\
137&$2$\\
138&$2$\\
139&$2$\\
140&$2$\\
141&$2$\\
154&$2$\\
155&$2$\\
156&$2$\\
157&$2$\\
158&$2$\\
159&$2$\\
172&$2$\\
173&$2$\\
174&$2$\\
175&$2$\\
176&$2$\\
177&$2$\\
190&$2$\\
191&$2$\\
192&$2$\\
193&$2$\\
194&$2$\\
195&$2$\\
209&$2$\\
\end{longtable}

% Generated from the frozen candidate238-ir.json; all 212 normalized equations.
% Requires amsmath. Variables z_1,...,z_238; A_N=149^(2N+1).
\subsection{Completion of the 238-variable predecessor}
The soundness and completeness proof for the 248-coordinate baseline is established in Section~\ref{v15:base248}. The scalar-S argument in Section~\ref{v15:reductions} changes its nonzero integer index to $|a|$ and preserves both directions. The three subsequent coordinate changes and the displayed affine and row operations are rational witness bijections, with the only nonconstant inverse powers justified by $R>0$. Their composition therefore proves the defining equivalence for this 238-variable predecessor. The further reductions pass through the 233-variable predecessor in Section~\ref{v15x:reductions} and culminate in the 222-variable family in Sections~\ref{v222:construction}--\ref{v222:height}. The same sufficiently large $N_0$ works for all larger $N$: the fixed aggregate bound first excludes the nonzero defect, and $149^{-L_N}\to0$ then gives the Lucas contradiction uniformly. The explicit equations below specify the resulting family. Their integer square expansion has coefficient one on $z_4^4$, degree four, and the stated coefficient profile; this last step is finite integer arithmetic and is independently reproducible from the complete displayed system. In particular it does not make $N_0$ effective.

\section{The explicit 238-variable predecessor}\label{v15:equations}
Put $A_N=149^{2N+1}$. The distinguished input is $z_1$. Define $Q_1,\ldots,Q_{212}$ below and $P_N=\sum_{r=1}^{212}Q_r^2$.
% H.chart.equation.0
\begin{equation*}\begin{aligned}
Q_{1}={}&
z_{17} -11 z_{19} +12 z_{20}\\
&-z_{21}^{2} -2 z_{22}^{2}
\end{aligned}\end{equation*}
% H.chart.equation.1
\begin{equation*}\begin{aligned}
Q_{2}={}&
2 z_{18} +6 z_{19} -11 z_{20}\\
&-2 z_{21} z_{22}
\end{aligned}\end{equation*}
% U0.f.definition.0
\begin{equation*}\begin{aligned}
Q_{3}={}&
1 -z_{17} z_{19} -4 z_{18} z_{20}
\end{aligned}\end{equation*}
% U0.f.definition.1
\begin{equation*}\begin{aligned}
Q_{4}={}&
z_{17} z_{20} +2 z_{18} z_{19}
\end{aligned}\end{equation*}
% V0.f.definition.0
\begin{equation*}\begin{aligned}
Q_{5}={}&
-2795584 +358004241 z_{19} z_{23} +716008482 z_{20} z_{24}\\
&+63271824 z_{140} -358004241 z_{140}^{2} -716008482 z_{141}^{2}
\end{aligned}\end{equation*}
% V0.f.definition.1
\begin{equation*}\begin{aligned}
Q_{6}={}&
18921 z_{19} z_{24} +18921 z_{20} z_{23} -37842 z_{140} z_{141}\\
&+3344 z_{141}
\end{aligned}\end{equation*}
% logV0.fSquare.0.definition
\begin{equation*}\begin{aligned}
Q_{7}={}&
z_{23}^{2} +2 z_{24}^{2} -z_{33}
\end{aligned}\end{equation*}
% logV0.fSquare.1.definition
\begin{equation*}\begin{aligned}
Q_{8}={}&
z_{23} z_{24} -z_{34}
\end{aligned}\end{equation*}
% logV0.wPower2.0.definition
\begin{equation*}\begin{aligned}
Q_{9}={}&
2401 z_{33} z_{34} -6369 z_{35}
\end{aligned}\end{equation*}
% logV0.wPower2.1.definition
\begin{equation*}\begin{aligned}
Q_{10}={}&
33 z_{33}^{2} -193 z_{33} z_{34} +264 z_{34}^{2}\\
&-33 z_{36}
\end{aligned}\end{equation*}
% plain.alternateRatio1.cleared.0
\begin{equation*}\begin{aligned}
Q_{11}={}&
775761 z_{21} z_{25} +504136 z_{21} z_{53} +373756 z_{21} z_{54}\\
&+1551522 z_{22} z_{26} +373756 z_{22} z_{53} +1008272 z_{22} z_{54}\\
&-1542912 z_{25} +3027360 z_{25} z_{140} -8022504 z_{25} z_{141}\\
&-8022504 z_{26} z_{140} +6054720 z_{26} z_{141} +775761 z_{53}
\end{aligned}\end{equation*}
% plain.alternateRatio1.cleared.1
\begin{equation*}\begin{aligned}
Q_{12}={}&
775761 z_{21} z_{26} +186878 z_{21} z_{53} +504136 z_{21} z_{54}\\
&+775761 z_{22} z_{25} +504136 z_{22} z_{53} +373756 z_{22} z_{54}\\
&-4011252 z_{25} z_{140} +3027360 z_{25} z_{141} -1542912 z_{26}\\
&+3027360 z_{26} z_{140} -8022504 z_{26} z_{141} +775761 z_{54}
\end{aligned}\end{equation*}
% logV2.wPower2.0.definition
\begin{equation*}\begin{aligned}
Q_{13}={}&
2401 z_{39} z_{40} -6369 z_{41}
\end{aligned}\end{equation*}
% logV2.wPower2.1.definition
\begin{equation*}\begin{aligned}
Q_{14}={}&
33 z_{39}^{2} -193 z_{39} z_{40} +264 z_{40}^{2}\\
&-33 z_{42}
\end{aligned}\end{equation*}
% plain.alternateRatio2.cleared.0
\begin{equation*}\begin{aligned}
Q_{15}={}&
775761 z_{21} z_{27} +504136 z_{21} z_{57} +373756 z_{21} z_{58}\\
&+1551522 z_{22} z_{28} +373756 z_{22} z_{57} +1008272 z_{22} z_{58}\\
&-1542912 z_{27} +3027360 z_{27} z_{140} -8022504 z_{27} z_{141}\\
&-8022504 z_{28} z_{140} +6054720 z_{28} z_{141} +775761 z_{57}
\end{aligned}\end{equation*}
% plain.alternateRatio2.cleared.1
\begin{equation*}\begin{aligned}
Q_{16}={}&
775761 z_{21} z_{28} +186878 z_{21} z_{57} +504136 z_{21} z_{58}\\
&+775761 z_{22} z_{27} +504136 z_{22} z_{57} +373756 z_{22} z_{58}\\
&-4011252 z_{27} z_{140} +3027360 z_{27} z_{141} -1542912 z_{28}\\
&+3027360 z_{28} z_{140} -8022504 z_{28} z_{141} +775761 z_{58}
\end{aligned}\end{equation*}
% logV3.wPower2.0.definition
\begin{equation*}\begin{aligned}
Q_{17}={}&
2401 z_{43} z_{44} -6369 z_{45}
\end{aligned}\end{equation*}
% logV3.wPower2.1.definition
\begin{equation*}\begin{aligned}
Q_{18}={}&
33 z_{43}^{2} -193 z_{43} z_{44} +264 z_{44}^{2}\\
&-33 z_{46}
\end{aligned}\end{equation*}
% plain.alternateRatio3.cleared.0
\begin{equation*}\begin{aligned}
Q_{19}={}&
775761 z_{21} z_{29} +504136 z_{21} z_{61} +373756 z_{21} z_{62}\\
&+1551522 z_{22} z_{30} +373756 z_{22} z_{61} +1008272 z_{22} z_{62}\\
&-1542912 z_{29} +3027360 z_{29} z_{140} -8022504 z_{29} z_{141}\\
&-8022504 z_{30} z_{140} +6054720 z_{30} z_{141} +775761 z_{61}
\end{aligned}\end{equation*}
% plain.alternateRatio3.cleared.1
\begin{equation*}\begin{aligned}
Q_{20}={}&
775761 z_{21} z_{30} +186878 z_{21} z_{61} +504136 z_{21} z_{62}\\
&+775761 z_{22} z_{29} +504136 z_{22} z_{61} +373756 z_{22} z_{62}\\
&-4011252 z_{29} z_{140} +3027360 z_{29} z_{141} -1542912 z_{30}\\
&+3027360 z_{30} z_{140} -8022504 z_{30} z_{141} +775761 z_{62}
\end{aligned}\end{equation*}
% logV4.wPower2.0.definition
\begin{equation*}\begin{aligned}
Q_{21}={}&
2401 z_{47} z_{48} -6369 z_{49}
\end{aligned}\end{equation*}
% logV4.wPower2.1.definition
\begin{equation*}\begin{aligned}
Q_{22}={}&
33 z_{47}^{2} -193 z_{47} z_{48} +264 z_{48}^{2}\\
&-33 z_{50}
\end{aligned}\end{equation*}
% plain.alternateRatio4.cleared.0
\begin{equation*}\begin{aligned}
Q_{23}={}&
-775761 z_{21} z_{32} -186878 z_{21} z_{65} -504136 z_{21} z_{66}\\
&-775761 z_{22} z_{31} -504136 z_{22} z_{65} -373756 z_{22} z_{66}\\
&+4011252 z_{31} z_{140} -3027360 z_{31} z_{141} +1542912 z_{32}\\
&-3027360 z_{32} z_{140} +8022504 z_{32} z_{141} -775761 z_{66}
\end{aligned}\end{equation*}
% plain.alternateRatio4.cleared.1
\begin{equation*}\begin{aligned}
Q_{24}={}&
775761 z_{21} z_{31} +504136 z_{21} z_{65} +373756 z_{21} z_{66}\\
&+1551522 z_{22} z_{32} +373756 z_{22} z_{65} +1008272 z_{22} z_{66}\\
&-1542912 z_{31} +3027360 z_{31} z_{140} -8022504 z_{31} z_{141}\\
&-8022504 z_{32} z_{140} +6054720 z_{32} z_{141} +775761 z_{65}
\end{aligned}\end{equation*}
% inputSquare.definition
\begin{equation*}\begin{aligned}
Q_{25}={}&
2501 +51 z_{1}^{2} -51 z_{67}
\end{aligned}\end{equation*}
% Lucas.C.0.definition
\begin{equation*}\begin{aligned}
Q_{26}={}&
z_{55} z_{63} +2 z_{56} z_{64} -2 z_{63} z_{142}\\
&-z_{68} z_{142}
\end{aligned}\end{equation*}
% Lucas.C.1.definition
\begin{equation*}\begin{aligned}
Q_{27}={}&
z_{55} z_{64} +z_{56} z_{63} -2 z_{64} z_{142}\\
&-z_{69} z_{142}
\end{aligned}\end{equation*}
% Lucas.AT.0.definition
\begin{equation*}\begin{aligned}
Q_{28}={}&
z_{55} z_{59} +2 z_{56} z_{60} -z_{68} z_{142}\\
&-z_{70} z_{142}
\end{aligned}\end{equation*}
% Lucas.AT.1.definition
\begin{equation*}\begin{aligned}
Q_{29}={}&
z_{55} z_{60} +z_{56} z_{59} -z_{69} z_{142}\\
&-z_{71} z_{142}
\end{aligned}\end{equation*}
% Lucas.F0.definition
\begin{equation*}\begin{aligned}
Q_{30}={}&
z_{59}^{2} +2 z_{60}^{2} -z_{68} z_{70}\\
&-2 z_{69} z_{71} -z_{72} z_{139}
\end{aligned}\end{equation*}
% Lucas.F1.definition
\begin{equation*}\begin{aligned}
Q_{31}={}&
2 z_{59} z_{60} -z_{68} z_{71} -z_{69} z_{70}\\
&-2 z_{73} z_{139}
\end{aligned}\end{equation*}
% testSquare.12.0.0.definition
\begin{equation*}\begin{aligned}
Q_{32}={}&
z_{51}^{2} -2 z_{51} z_{53} +2 z_{52}^{2}\\
&+4 z_{52} z_{54} +z_{53}^{2} +2 z_{54}^{2}\\
&-z_{82}
\end{aligned}\end{equation*}
% testSquare.12.0.1.definition
\begin{equation*}\begin{aligned}
Q_{33}={}&
z_{51} z_{52} +z_{51} z_{54} -z_{52} z_{53}\\
&-z_{53} z_{54} -z_{83}
\end{aligned}\end{equation*}
% testSquare.12.1.0.definition
\begin{equation*}\begin{aligned}
Q_{34}={}&
9 +z_{51}^{2} -102 z_{51} z_{67}\\
&+2 z_{52}^{2} +2601 z_{67}^{2} -z_{84}
\end{aligned}\end{equation*}
% testSquare.12.1.1.definition
\begin{equation*}\begin{aligned}
Q_{35}={}&
z_{51} z_{52} -51 z_{52} z_{67} -z_{85}
\end{aligned}\end{equation*}
% testSquare.40.0.0.definition
\begin{equation*}\begin{aligned}
Q_{36}={}&
z_{19}^{2} +2 z_{20}^{2} -z_{86}
\end{aligned}\end{equation*}
% testSquare.40.0.1.definition
\begin{equation*}\begin{aligned}
Q_{37}={}&
z_{19} z_{20} -z_{87}
\end{aligned}\end{equation*}
% testSquare.40.1.0.definition
\begin{equation*}\begin{aligned}
Q_{38}={}&
z_{55}^{2} -2 z_{55} z_{57} +2 z_{56}^{2}\\
&+4 z_{56} z_{58} +z_{57}^{2} +2 z_{58}^{2}\\
&-z_{88} z_{139}
\end{aligned}\end{equation*}
% testSquare.40.1.1.definition
\begin{equation*}\begin{aligned}
Q_{39}={}&
z_{55} z_{56} +z_{55} z_{58} -z_{56} z_{57}\\
&-z_{57} z_{58} -z_{89} z_{139}
\end{aligned}\end{equation*}
% testSquare.40.2.0.definition
\begin{equation*}\begin{aligned}
Q_{40}={}&
z_{59}^{2} -2 z_{59} z_{61} +2 z_{60}^{2}\\
&+4 z_{60} z_{62} +z_{61}^{2} +2 z_{62}^{2}\\
&-z_{90} z_{139}
\end{aligned}\end{equation*}
% testSquare.40.2.1.definition
\begin{equation*}\begin{aligned}
Q_{41}={}&
z_{59} z_{60} +z_{59} z_{62} -z_{60} z_{61}\\
&-z_{61} z_{62} -z_{91} z_{139}
\end{aligned}\end{equation*}
% testSquare.40.3.0.definition
\begin{equation*}\begin{aligned}
Q_{42}={}&
z_{63}^{2} -2 z_{63} z_{65} +2 z_{64}^{2}\\
&+4 z_{64} z_{66} +z_{65}^{2} +2 z_{66}^{2}\\
&-z_{92} z_{139}
\end{aligned}\end{equation*}
% testSquare.40.3.1.definition
\begin{equation*}\begin{aligned}
Q_{43}={}&
z_{63} z_{64} +z_{63} z_{66} -z_{64} z_{65}\\
&-z_{65} z_{66} -z_{93} z_{139}
\end{aligned}\end{equation*}
% testSquare.40.4.0.definition
\begin{equation*}\begin{aligned}
Q_{44}={}&
1 +z_{72}^{2} +8 z_{73}^{2}\\
&-z_{94}
\end{aligned}\end{equation*}
% testSquare.40.4.1.definition
\begin{equation*}\begin{aligned}
Q_{45}={}&
z_{72} z_{73} -z_{95}
\end{aligned}\end{equation*}
% reciprocal.symmetric.detector0.deltaSquared.definition
\begin{equation*}\begin{aligned}
Q_{46}={}&
z_{4}^{2} -z_{96}
\end{aligned}\end{equation*}
% reciprocal.symmetric.detector0.deltaFourthPower.definition
\begin{equation*}\begin{aligned}
Q_{47}={}&
z_{96}^{2} -z_{97}
\end{aligned}\end{equation*}
% reciprocal.reflected.detector0.contactPolynomial.definition
\begin{equation*}\begin{aligned}
Q_{48}={}&
7626746323592 -76267463235920 z_{4} -628059618995128 z_{4} z_{96}\\
&+732473028362323 z_{4} z_{143} +307309842082151 z_{96} -583791423311868 z_{96} z_{143}\\
&+718967979842546 z_{97} +162308064100032 z_{97}^{2} -38129835686440 z_{97} z_{143}\\
&-7626746323592 z_{98} -3813373161796 z_{98} z_{99} -646146742577308 z_{143}\\
&+3813373161796 z_{143}^{2}
\end{aligned}\end{equation*}
% reciprocal.reflected.detector0.weight.definition
\begin{equation*}\begin{aligned}
Q_{49}={}&
1 -10 z_{4} -120 z_{4} z_{96}\\
&+210 z_{4} z_{143} +45 z_{96} -120 z_{96} z_{143}\\
&+210 z_{97} +45 z_{97}^{2} -10 z_{97} z_{143}\\
&-z_{98} -z_{98} z_{99} -252 z_{143}\\
&+z_{143}^{2}
\end{aligned}\end{equation*}
% detector0.E1.power2.definition
\begin{equation*}\begin{aligned}
Q_{50}={}&
z_{99}^{2} -z_{101}
\end{aligned}\end{equation*}
% detector0.E1.power4.definition
\begin{equation*}\begin{aligned}
Q_{51}={}&
z_{101}^{2} -z_{102}
\end{aligned}\end{equation*}
% detector0.normOver2w0.definition
\begin{equation*}\begin{aligned}
Q_{52}={}&
z_{103}^{2} -2 z_{103} z_{105} -2 z_{104}^{2}
\end{aligned}\end{equation*}
% balanced.detector0.poleInput.definition
\begin{equation*}\begin{aligned}
Q_{53}={}&
167281 z_{99} z_{106} -166464 z_{105}^{2} +167281 z_{106}
\end{aligned}\end{equation*}
% detector0.p2.definition
\begin{equation*}\begin{aligned}
Q_{54}={}&
z_{7}^{2} -z_{107}
\end{aligned}\end{equation*}
% detector0.firstConic
\begin{equation*}\begin{aligned}
Q_{55}={}&
289 +34 z_{5} -3 z_{5}^{2}\\
&-17 z_{106} z_{107}
\end{aligned}\end{equation*}
% detector0.Bh.definition
\begin{equation*}\begin{aligned}
Q_{56}={}&
17 z_{5} -z_{5} z_{106} -17 z_{6} z_{108}
\end{aligned}\end{equation*}
% detector0.q2.definition
\begin{equation*}\begin{aligned}
Q_{57}={}&
z_{8}^{2} -z_{109}
\end{aligned}\end{equation*}
% detector0.secondConic
\begin{equation*}\begin{aligned}
Q_{58}={}&
z_{6}^{2} +34 z_{6} z_{108} -17 z_{106} z_{109}\\
&-867 z_{108}^{2}
\end{aligned}\end{equation*}
% symmetric.detector1.deltaSquared.definition
\begin{equation*}\begin{aligned}
Q_{59}={}&
z_{110}^{2} -2 z_{110} z_{139} -z_{111} z_{139}\\
&+z_{139}^{2}
\end{aligned}\end{equation*}
% symmetric.detector1.deltaFourthPower.definition
\begin{equation*}\begin{aligned}
Q_{60}={}&
z_{111}^{2} -z_{112} z_{139}
\end{aligned}\end{equation*}
% reflected.detector1.contactPolynomial.definition
\begin{equation*}\begin{aligned}
Q_{61}={}&
-141822723674536 z_{111} z_{139} +35893742479489 z_{111} z_{144} +192222345736537 z_{112} z_{139}\\
&-3813373161796 z_{113} z_{114} +39896049934224 z_{139}^{2} -122376041313918 z_{139} z_{144}
\end{aligned}\end{equation*}
% reflected.detector1.weight.definition
\begin{equation*}\begin{aligned}
Q_{62}={}&
z_{112} z_{144} -z_{113} -z_{113} z_{114}
\end{aligned}\end{equation*}
% detector1.E1.power2.definition
\begin{equation*}\begin{aligned}
Q_{63}={}&
z_{114}^{2} -z_{116}
\end{aligned}\end{equation*}
% detector1.E1.power4.definition
\begin{equation*}\begin{aligned}
Q_{64}={}&
z_{116}^{2} -z_{117}
\end{aligned}\end{equation*}
% detector1.normOver2w0.definition
\begin{equation*}\begin{aligned}
Q_{65}={}&
z_{118}^{2} -2 z_{118} z_{120} -2 z_{119}^{2}
\end{aligned}\end{equation*}
% balanced.detector1.poleInput.definition
\begin{equation*}\begin{aligned}
Q_{66}={}&
167281 z_{114} z_{121} -166464 z_{120}^{2} +167281 z_{121}
\end{aligned}\end{equation*}
% detector1.p2.definition
\begin{equation*}\begin{aligned}
Q_{67}={}&
z_{11}^{2} -z_{122}
\end{aligned}\end{equation*}
% detector1.firstConic
\begin{equation*}\begin{aligned}
Q_{68}={}&
289 +34 z_{9} -3 z_{9}^{2}\\
&-17 z_{121} z_{122}
\end{aligned}\end{equation*}
% detector1.Bh.definition
\begin{equation*}\begin{aligned}
Q_{69}={}&
17 z_{9} -z_{9} z_{121} -17 z_{10} z_{123}
\end{aligned}\end{equation*}
% detector1.q2.definition
\begin{equation*}\begin{aligned}
Q_{70}={}&
z_{12}^{2} -z_{124}
\end{aligned}\end{equation*}
% detector1.secondConic
\begin{equation*}\begin{aligned}
Q_{71}={}&
z_{10}^{2} +34 z_{10} z_{123} -17 z_{121} z_{124}\\
&-867 z_{123}^{2}
\end{aligned}\end{equation*}
% scalar.rPair1.definition
\begin{equation*}\begin{aligned}
Q_{72}={}&
2601 z_{67}^{2} -4 z_{125}
\end{aligned}\end{equation*}
% scalar.rFour.definition
\begin{equation*}\begin{aligned}
Q_{73}={}&
3882987 z_{67} z_{125} -128 z_{125}^{2} -381485 z_{126}
\end{aligned}\end{equation*}
% scalar.contact_product.definition
\begin{equation*}\begin{aligned}
Q_{74}={}&
31581660 +2103220824 z_{67} +29762681184 z_{67}^{2}\\
&+585078120 z_{67} z_{125} +155645880 z_{67} z_{126} +28680300 z_{125}\\
&+7629700 z_{126} +z_{127}
\end{aligned}\end{equation*}
% scalar.W2.definition
\begin{equation*}\begin{aligned}
Q_{75}={}&
K_{1} z_{130} z_{139} -z_{138}^{2}
\end{aligned}\end{equation*}
% scalar.b.definition
\begin{equation*}\begin{aligned}
Q_{76}={}&
9792 z_{130} +z_{130} z_{142} +82479248 z_{131} z_{142}
\end{aligned}\end{equation*}
% scalar.p2.definition
\begin{equation*}\begin{aligned}
Q_{77}={}&
z_{15}^{2} -z_{132}
\end{aligned}\end{equation*}
% scalar.firstConic
\begin{equation*}\begin{aligned}
Q_{78}={}&
K_{7} -K_{4} z_{13} -z_{13}^{2}\\
&+K_{3} z_{131} z_{132}
\end{aligned}\end{equation*}
% scalar.Bh.definition
\begin{equation*}\begin{aligned}
Q_{79}={}&
17 z_{13} +K_{2} z_{13} z_{131} -17 z_{14} z_{133}
\end{aligned}\end{equation*}
% scalar.q2.definition
\begin{equation*}\begin{aligned}
Q_{80}={}&
z_{16}^{2} -z_{134}
\end{aligned}\end{equation*}
% scalar.secondConic
\begin{equation*}\begin{aligned}
Q_{81}={}&
z_{14}^{2} -K_{6} z_{14} z_{133} +K_{5} z_{131} z_{134}\\
&-K_{8} z_{133}^{2}
\end{aligned}\end{equation*}
% reduced_scalar.ratio_node_squared
\begin{equation*}\begin{aligned}
Q_{82}={}&
z_{86}^{2} -8 z_{87}^{2} -z_{137}
\end{aligned}\end{equation*}
% reduced_scalar.root.power4
\begin{equation*}\begin{aligned}
Q_{83}={}&
z_{136}^{2} +z_{137} z_{142} -z_{138} z_{139}
\end{aligned}\end{equation*}
% reflected.detector0.weightPower7.definition
\begin{equation*}\begin{aligned}
Q_{84}={}&
z_{99} z_{100} -z_{102}^{2}
\end{aligned}\end{equation*}
% reflected.detector1.weightPower7.definition
\begin{equation*}\begin{aligned}
Q_{85}={}&
z_{114} z_{115} -z_{117}^{2}
\end{aligned}\end{equation*}
% scalar.r_power_2.definition
\begin{equation*}\begin{aligned}
Q_{86}={}&
z_{139} -z_{142}^{2}
\end{aligned}\end{equation*}
% scalar.ratio1_unscaled_square.definition
\begin{equation*}\begin{aligned}
Q_{87}={}&
z_{129}^{2} -z_{135}
\end{aligned}\end{equation*}
% scalar.ratio_node.consolidated_denominator.power4
\begin{equation*}\begin{aligned}
Q_{88}={}&
z_{128}^{2} +z_{135} z_{142} -z_{136} z_{139}
\end{aligned}\end{equation*}
% V0.reciprocal.g.definition.0
\begin{equation*}\begin{aligned}
Q_{89}={}&
24231 -2810632 z_{17} +31806201 z_{17} z_{140}\\
&+127224804 z_{18} z_{141} -6206088 z_{21} -21721308 z_{22}\\
&+44048088 z_{140} -10293024 z_{141}
\end{aligned}\end{equation*}
% V0.reciprocal.g.definition.1
\begin{equation*}\begin{aligned}
Q_{90}={}&
17709376 -31806201 z_{17} z_{141} +5621264 z_{18}\\
&-63612402 z_{18} z_{140} +10860654 z_{21} +6206088 z_{22}\\
&+5146512 z_{140} -44048088 z_{141}
\end{aligned}\end{equation*}
% scalar.forced_prime.definition
\begin{equation*}\begin{aligned}
Q_{91}={}&
22201 z_{127} z_{142} +z_{142} -A_N z_{127}
\end{aligned}\end{equation*}
% reciprocal.reflected.detector0.deltaSixthPower.definition
\begin{equation*}\begin{aligned}
Q_{92}={}&
z_{4} z_{97} -z_{143}
\end{aligned}\end{equation*}
% reflected.detector1.deltaSixthPower.definition
\begin{equation*}\begin{aligned}
Q_{93}={}&
z_{111} z_{112} -z_{139} z_{144}
\end{aligned}\end{equation*}
% amplified.C1.power4.0.definition
\begin{equation*}\begin{aligned}
Q_{94}={}&
z_{82}^{2} +8 z_{83}^{2} -z_{145}
\end{aligned}\end{equation*}
% amplified.C1.power4.1.definition
\begin{equation*}\begin{aligned}
Q_{95}={}&
4 z_{82} z_{83} -z_{146}
\end{aligned}\end{equation*}
% residue.I.power4.0.definition
\begin{equation*}\begin{aligned}
Q_{96}={}&
-324 -12 z_{51} z_{84} -48 z_{52} z_{85}\\
&+612 z_{67} z_{84} +36 z_{84} +z_{84}^{2}\\
&+8 z_{85}^{2} -z_{147}
\end{aligned}\end{equation*}
% residue.I.power4.1.definition
\begin{equation*}\begin{aligned}
Q_{97}={}&
-24 z_{51} z_{85} -12 z_{52} z_{84} +1224 z_{67} z_{85}\\
&+4 z_{84} z_{85} +72 z_{85} -z_{148}
\end{aligned}\end{equation*}
% residue.V1.power4.0.definition
\begin{equation*}\begin{aligned}
Q_{98}={}&
z_{37}^{2} +8 z_{38}^{2} -z_{149}
\end{aligned}\end{equation*}
% residue.V1.power4.1.definition
\begin{equation*}\begin{aligned}
Q_{99}={}&
4 z_{37} z_{38} +z_{150}
\end{aligned}\end{equation*}
% residue.det0.factor0.0.definition
\begin{equation*}\begin{aligned}
Q_{100}={}&
1 +z_{86} z_{101} +z_{88} z_{100}\\
&+z_{102} z_{149} -z_{151}
\end{aligned}\end{equation*}
% residue.det0.factor0.1.definition
\begin{equation*}\begin{aligned}
Q_{101}={}&
2 z_{87} z_{101} +2 z_{89} z_{100} +z_{102} z_{150}\\
&-z_{152}
\end{aligned}\end{equation*}
% residue.det0.factor1.0.definition
\begin{equation*}\begin{aligned}
Q_{102}={}&
2401 +37249 z_{35} z_{102} +2401 z_{36} z_{102}\\
&-4802 z_{72} z_{100} +2401 z_{78} z_{101} +2401 z_{94} z_{100}\\
&-2401 z_{153}
\end{aligned}\end{equation*}
% residue.det0.factor1.1.definition
\begin{equation*}\begin{aligned}
Q_{103}={}&
25476 z_{35} z_{102} +9604 z_{73} z_{100} -4802 z_{79} z_{101}\\
&-9604 z_{95} z_{100} +2401 z_{154}
\end{aligned}\end{equation*}
% residue.det0.factor2.0.definition
\begin{equation*}\begin{aligned}
Q_{104}={}&
2401 +37249 z_{45} z_{102} +2401 z_{46} z_{102}\\
&+2401 z_{76} z_{99} +2401 z_{80} z_{101} +2401 z_{100} z_{147}\\
&-2401 z_{155}
\end{aligned}\end{equation*}
% residue.det0.factor2.1.definition
\begin{equation*}\begin{aligned}
Q_{105}={}&
25476 z_{45} z_{102} -4802 z_{77} z_{99} -4802 z_{81} z_{101}\\
&-2401 z_{100} z_{148} +2401 z_{156}
\end{aligned}\end{equation*}
% residue.det0.factor3.0.definition
\begin{equation*}\begin{aligned}
Q_{106}={}&
2401 +37249 z_{49} z_{102} +2401 z_{50} z_{102}\\
&+2401 z_{90} z_{100} +2401 z_{101} z_{145} -2401 z_{157}
\end{aligned}\end{equation*}
% residue.det0.factor3.1.definition
\begin{equation*}\begin{aligned}
Q_{107}={}&
25476 z_{49} z_{102} -4802 z_{91} z_{100} -2401 z_{101} z_{146}\\
&+2401 z_{158}
\end{aligned}\end{equation*}
% residue.det0.factor4.0.definition
\begin{equation*}\begin{aligned}
Q_{108}={}&
2401 +37249 z_{41} z_{102} +2401 z_{42} z_{102}\\
&+2401 z_{74} z_{101} +2401 z_{92} z_{100} -2401 z_{159}
\end{aligned}\end{equation*}
% residue.det0.factor4.1.definition
\begin{equation*}\begin{aligned}
Q_{109}={}&
25476 z_{41} z_{102} -4802 z_{75} z_{101} -4802 z_{93} z_{100}\\
&+2401 z_{160}
\end{aligned}\end{equation*}
% amplified.det0.product1.0.definition
\begin{equation*}\begin{aligned}
Q_{110}={}&
z_{151} z_{153} +2 z_{152} z_{154} -z_{161}
\end{aligned}\end{equation*}
% amplified.det0.product1.1.definition
\begin{equation*}\begin{aligned}
Q_{111}={}&
z_{151} z_{154} +z_{152} z_{153} -z_{162}
\end{aligned}\end{equation*}
% amplified.det0.product2.0.definition
\begin{equation*}\begin{aligned}
Q_{112}={}&
z_{155} z_{161} +2 z_{156} z_{162} -z_{163}
\end{aligned}\end{equation*}
% amplified.det0.product2.1.definition
\begin{equation*}\begin{aligned}
Q_{113}={}&
z_{155} z_{162} +z_{156} z_{161} -z_{164}
\end{aligned}\end{equation*}
% amplified.det0.product3.0.definition
\begin{equation*}\begin{aligned}
Q_{114}={}&
z_{157} z_{163} +2 z_{158} z_{164} -z_{165}
\end{aligned}\end{equation*}
% amplified.det0.product3.1.definition
\begin{equation*}\begin{aligned}
Q_{115}={}&
z_{157} z_{164} +z_{158} z_{163} -z_{166}
\end{aligned}\end{equation*}
% amplified.det0.product4.0.definition
\begin{equation*}\begin{aligned}
Q_{116}={}&
z_{103} -z_{159} z_{165} -2 z_{160} z_{166}
\end{aligned}\end{equation*}
% amplified.det0.product4.1.definition
\begin{equation*}\begin{aligned}
Q_{117}={}&
z_{104} -z_{159} z_{166} -z_{160} z_{165}
\end{aligned}\end{equation*}
% residue.det1.factor0.0.definition
\begin{equation*}\begin{aligned}
Q_{118}={}&
1 +z_{86} z_{116} +z_{88} z_{115}\\
&+z_{117} z_{149} -z_{167}
\end{aligned}\end{equation*}
% residue.det1.factor0.1.definition
\begin{equation*}\begin{aligned}
Q_{119}={}&
2 z_{87} z_{116} +2 z_{89} z_{115} +z_{117} z_{150}\\
&-z_{168}
\end{aligned}\end{equation*}
% residue.det1.factor1.0.definition
\begin{equation*}\begin{aligned}
Q_{120}={}&
2401 +37249 z_{35} z_{117} +2401 z_{36} z_{117}\\
&-4802 z_{72} z_{115} +2401 z_{78} z_{116} +2401 z_{94} z_{115}\\
&-2401 z_{169}
\end{aligned}\end{equation*}
% residue.det1.factor1.1.definition
\begin{equation*}\begin{aligned}
Q_{121}={}&
25476 z_{35} z_{117} +9604 z_{73} z_{115} -4802 z_{79} z_{116}\\
&-9604 z_{95} z_{115} +2401 z_{170}
\end{aligned}\end{equation*}
% residue.det1.factor2.0.definition
\begin{equation*}\begin{aligned}
Q_{122}={}&
2401 +37249 z_{45} z_{117} +2401 z_{46} z_{117}\\
&+2401 z_{76} z_{114} +2401 z_{80} z_{116} +2401 z_{115} z_{147}\\
&-2401 z_{171}
\end{aligned}\end{equation*}
% residue.det1.factor2.1.definition
\begin{equation*}\begin{aligned}
Q_{123}={}&
25476 z_{45} z_{117} -4802 z_{77} z_{114} -4802 z_{81} z_{116}\\
&-2401 z_{115} z_{148} +2401 z_{172}
\end{aligned}\end{equation*}
% residue.det1.factor3.0.definition
\begin{equation*}\begin{aligned}
Q_{124}={}&
2401 +37249 z_{49} z_{117} +2401 z_{50} z_{117}\\
&+2401 z_{90} z_{115} +2401 z_{116} z_{145} -2401 z_{173}
\end{aligned}\end{equation*}
% residue.det1.factor3.1.definition
\begin{equation*}\begin{aligned}
Q_{125}={}&
25476 z_{49} z_{117} -4802 z_{91} z_{115} -2401 z_{116} z_{146}\\
&+2401 z_{174}
\end{aligned}\end{equation*}
% residue.det1.factor4.0.definition
\begin{equation*}\begin{aligned}
Q_{126}={}&
2401 +37249 z_{41} z_{117} +2401 z_{42} z_{117}\\
&+2401 z_{74} z_{116} +2401 z_{92} z_{115} -2401 z_{175}
\end{aligned}\end{equation*}
% residue.det1.factor4.1.definition
\begin{equation*}\begin{aligned}
Q_{127}={}&
25476 z_{41} z_{117} -4802 z_{75} z_{116} -4802 z_{93} z_{115}\\
&+2401 z_{176}
\end{aligned}\end{equation*}
% amplified.det1.product1.0.definition
\begin{equation*}\begin{aligned}
Q_{128}={}&
z_{167} z_{169} +2 z_{168} z_{170} -z_{177}
\end{aligned}\end{equation*}
% amplified.det1.product1.1.definition
\begin{equation*}\begin{aligned}
Q_{129}={}&
z_{167} z_{170} +z_{168} z_{169} -z_{178}
\end{aligned}\end{equation*}
% amplified.det1.product2.0.definition
\begin{equation*}\begin{aligned}
Q_{130}={}&
z_{171} z_{177} +2 z_{172} z_{178} -z_{179}
\end{aligned}\end{equation*}
% amplified.det1.product2.1.definition
\begin{equation*}\begin{aligned}
Q_{131}={}&
z_{171} z_{178} +z_{172} z_{177} -z_{180}
\end{aligned}\end{equation*}
% amplified.det1.product3.0.definition
\begin{equation*}\begin{aligned}
Q_{132}={}&
z_{173} z_{179} +2 z_{174} z_{180} -z_{181}
\end{aligned}\end{equation*}
% amplified.det1.product3.1.definition
\begin{equation*}\begin{aligned}
Q_{133}={}&
z_{173} z_{180} +z_{174} z_{179} -z_{182}
\end{aligned}\end{equation*}
% amplified.det1.product4.0.definition
\begin{equation*}\begin{aligned}
Q_{134}={}&
z_{118} -z_{175} z_{181} -2 z_{176} z_{182}
\end{aligned}\end{equation*}
% amplified.det1.product4.1.definition
\begin{equation*}\begin{aligned}
Q_{135}={}&
z_{119} -z_{175} z_{182} -z_{176} z_{181}
\end{aligned}\end{equation*}
% rootProduct.1.equation1.0
\begin{equation*}\begin{aligned}
Q_{136}={}&
z_{183}^{2} +2 z_{184}^{2} +z_{185}^{2}\\
&+2 z_{186}^{2} +z_{187}^{2} +2 z_{188}^{2}\\
&-z_{189}^{2} -2 z_{190}^{2}
\end{aligned}\end{equation*}
% rootProduct.1.equation1.1
\begin{equation*}\begin{aligned}
Q_{137}={}&
z_{183} z_{184} +z_{185} z_{186} +z_{187} z_{188}\\
&-z_{189} z_{190}
\end{aligned}\end{equation*}
% rootProduct.1.equation2.0
\begin{equation*}\begin{aligned}
Q_{138}={}&
-12 z_{142} z_{191} +8 z_{142} z_{192} +z_{183} z_{187}\\
&+2 z_{184} z_{188} +z_{185} z_{189} +2 z_{186} z_{190}
\end{aligned}\end{equation*}
% rootProduct.1.equation2.1
\begin{equation*}\begin{aligned}
Q_{139}={}&
4 z_{142} z_{191} -12 z_{142} z_{192} +z_{183} z_{188}\\
&+z_{184} z_{187} +z_{185} z_{190} +z_{186} z_{189}
\end{aligned}\end{equation*}
% rootProduct.1.equation3.0
\begin{equation*}\begin{aligned}
Q_{140}={}&
12 z_{142} z_{193} -8 z_{142} z_{194} +z_{183} z_{189}\\
&+2 z_{184} z_{190} +z_{185} z_{187} +2 z_{186} z_{188}
\end{aligned}\end{equation*}
% rootProduct.1.equation3.1
\begin{equation*}\begin{aligned}
Q_{141}={}&
-4 z_{142} z_{193} +12 z_{142} z_{194} +z_{183} z_{190}\\
&+z_{184} z_{189} +z_{185} z_{188} +z_{186} z_{187}
\end{aligned}\end{equation*}
% rootProduct.1.equation4.0
\begin{equation*}\begin{aligned}
Q_{142}={}&
-44 z_{139} z_{195} +48 z_{139} z_{196} +z_{183}^{2}\\
&+2 z_{184}^{2}
\end{aligned}\end{equation*}
% rootProduct.1.equation4.1
\begin{equation*}\begin{aligned}
Q_{143}={}&
12 z_{139} z_{195} -22 z_{139} z_{196} +z_{183} z_{184}
\end{aligned}\end{equation*}
% rootProduct.1.equation5.0
\begin{equation*}\begin{aligned}
Q_{144}={}&
z_{74} -z_{195}^{2} -2 z_{196}^{2}
\end{aligned}\end{equation*}
% rootProduct.1.equation5.1
\begin{equation*}\begin{aligned}
Q_{145}={}&
z_{75} -z_{195} z_{196}
\end{aligned}\end{equation*}
% rootProduct.1.equation6.0
\begin{equation*}\begin{aligned}
Q_{146}={}&
11 z_{74} z_{139} -24 z_{75} z_{139} -z_{139}\\
&+z_{191}^{2} +2 z_{192}^{2}
\end{aligned}\end{equation*}
% rootProduct.1.equation6.1
\begin{equation*}\begin{aligned}
Q_{147}={}&
-3 z_{74} z_{139} +11 z_{75} z_{139} +z_{191} z_{192}
\end{aligned}\end{equation*}
% rootProduct.1.equation7.0
\begin{equation*}\begin{aligned}
Q_{148}={}&
11 z_{37} z_{139} -24 z_{38} z_{139} -z_{139}\\
&+z_{193}^{2} +2 z_{194}^{2}
\end{aligned}\end{equation*}
% rootProduct.1.equation7.1
\begin{equation*}\begin{aligned}
Q_{149}={}&
-3 z_{37} z_{139} +11 z_{38} z_{139} +z_{193} z_{194}
\end{aligned}\end{equation*}
% rootProduct.1.equation8.0
\begin{equation*}\begin{aligned}
Q_{150}={}&
6 z_{25} z_{193} -4 z_{25} z_{194} -4 z_{26} z_{193}\\
&+12 z_{26} z_{194} +z_{185}
\end{aligned}\end{equation*}
% rootProduct.1.equation8.1
\begin{equation*}\begin{aligned}
Q_{151}={}&
-2 z_{25} z_{193} +6 z_{25} z_{194} +6 z_{26} z_{193}\\
&-4 z_{26} z_{194} +z_{186}
\end{aligned}\end{equation*}
% rootProduct.1.equation9.0
\begin{equation*}\begin{aligned}
Q_{152}={}&
z_{21} z_{183} +2 z_{22} z_{184} -6 z_{51} z_{191}\\
&+4 z_{51} z_{192} +4 z_{52} z_{191} -12 z_{52} z_{192}
\end{aligned}\end{equation*}
% rootProduct.1.equation9.1
\begin{equation*}\begin{aligned}
Q_{153}={}&
z_{21} z_{184} +z_{22} z_{183} +2 z_{51} z_{191}\\
&-6 z_{51} z_{192} -6 z_{52} z_{191} +4 z_{52} z_{192}
\end{aligned}\end{equation*}
% rootProduct.2.equation1.0
\begin{equation*}\begin{aligned}
Q_{154}={}&
z_{197}^{2} +2 z_{198}^{2} +z_{199}^{2}\\
&+2 z_{200}^{2} +z_{201}^{2} +2 z_{202}^{2}\\
&-z_{203}^{2} -2 z_{204}^{2}
\end{aligned}\end{equation*}
% rootProduct.2.equation1.1
\begin{equation*}\begin{aligned}
Q_{155}={}&
z_{197} z_{198} +z_{199} z_{200} +z_{201} z_{202}\\
&-z_{203} z_{204}
\end{aligned}\end{equation*}
% rootProduct.2.equation2.0
\begin{equation*}\begin{aligned}
Q_{156}={}&
-12 z_{139} z_{205} +8 z_{139} z_{206} +z_{197} z_{201}\\
&+2 z_{198} z_{202} +z_{199} z_{203} +2 z_{200} z_{204}
\end{aligned}\end{equation*}
% rootProduct.2.equation2.1
\begin{equation*}\begin{aligned}
Q_{157}={}&
4 z_{139} z_{205} -12 z_{139} z_{206} +z_{197} z_{202}\\
&+z_{198} z_{201} +z_{199} z_{204} +z_{200} z_{203}
\end{aligned}\end{equation*}
% rootProduct.2.equation3.0
\begin{equation*}\begin{aligned}
Q_{158}={}&
12 z_{139} z_{207} -8 z_{139} z_{208} +z_{197} z_{203}\\
&+2 z_{198} z_{204} +z_{199} z_{201} +2 z_{200} z_{202}
\end{aligned}\end{equation*}
% rootProduct.2.equation3.1
\begin{equation*}\begin{aligned}
Q_{159}={}&
-4 z_{139} z_{207} +12 z_{139} z_{208} +z_{197} z_{204}\\
&+z_{198} z_{203} +z_{199} z_{202} +z_{200} z_{201}
\end{aligned}\end{equation*}
% rootProduct.2.equation4.0
\begin{equation*}\begin{aligned}
Q_{160}={}&
-44 z_{139} z_{209} +48 z_{139} z_{210} +z_{197}^{2}\\
&+2 z_{198}^{2}
\end{aligned}\end{equation*}
% rootProduct.2.equation4.1
\begin{equation*}\begin{aligned}
Q_{161}={}&
12 z_{139} z_{209} -22 z_{139} z_{210} +z_{197} z_{198}
\end{aligned}\end{equation*}
% rootProduct.2.equation5.0
\begin{equation*}\begin{aligned}
Q_{162}={}&
z_{76} -z_{209}^{2} -2 z_{210}^{2}
\end{aligned}\end{equation*}
% rootProduct.2.equation5.1
\begin{equation*}\begin{aligned}
Q_{163}={}&
z_{77} -z_{209} z_{210}
\end{aligned}\end{equation*}
% rootProduct.2.equation6.0
\begin{equation*}\begin{aligned}
Q_{164}={}&
-1 +11 z_{76} -24 z_{77}\\
&+z_{205}^{2} +2 z_{206}^{2}
\end{aligned}\end{equation*}
% rootProduct.2.equation6.1
\begin{equation*}\begin{aligned}
Q_{165}={}&
-3 z_{76} +11 z_{77} +z_{205} z_{206}
\end{aligned}\end{equation*}
% rootProduct.2.equation7.0
\begin{equation*}\begin{aligned}
Q_{166}={}&
-1 +11 z_{39} -24 z_{40}\\
&+z_{207}^{2} +2 z_{208}^{2}
\end{aligned}\end{equation*}
% rootProduct.2.equation7.1
\begin{equation*}\begin{aligned}
Q_{167}={}&
-3 z_{39} +11 z_{40} +z_{207} z_{208}
\end{aligned}\end{equation*}
% rootProduct.2.equation8.0
\begin{equation*}\begin{aligned}
Q_{168}={}&
6 z_{27} z_{207} -4 z_{27} z_{208} -4 z_{28} z_{207}\\
&+12 z_{28} z_{208} +z_{199}
\end{aligned}\end{equation*}
% rootProduct.2.equation8.1
\begin{equation*}\begin{aligned}
Q_{169}={}&
-2 z_{27} z_{207} +6 z_{27} z_{208} +6 z_{28} z_{207}\\
&-4 z_{28} z_{208} +z_{200}
\end{aligned}\end{equation*}
% rootProduct.2.equation9.0
\begin{equation*}\begin{aligned}
Q_{170}={}&
z_{21} z_{197} +2 z_{22} z_{198} -6 z_{55} z_{205}\\
&+4 z_{55} z_{206} +4 z_{56} z_{205} -12 z_{56} z_{206}
\end{aligned}\end{equation*}
% rootProduct.2.equation9.1
\begin{equation*}\begin{aligned}
Q_{171}={}&
z_{21} z_{198} +z_{22} z_{197} +2 z_{55} z_{205}\\
&-6 z_{55} z_{206} -6 z_{56} z_{205} +4 z_{56} z_{206}
\end{aligned}\end{equation*}
% rootProduct.3.equation1.0
\begin{equation*}\begin{aligned}
Q_{172}={}&
z_{211}^{2} +2 z_{212}^{2} +z_{213}^{2}\\
&+2 z_{214}^{2} +z_{215}^{2} +2 z_{216}^{2}\\
&-z_{217}^{2} -2 z_{218}^{2}
\end{aligned}\end{equation*}
% rootProduct.3.equation1.1
\begin{equation*}\begin{aligned}
Q_{173}={}&
z_{211} z_{212} +z_{213} z_{214} +z_{215} z_{216}\\
&-z_{217} z_{218}
\end{aligned}\end{equation*}
% rootProduct.3.equation2.0
\begin{equation*}\begin{aligned}
Q_{174}={}&
-12 z_{139} z_{219} +8 z_{139} z_{220} +z_{211} z_{215}\\
&+2 z_{212} z_{216} +z_{213} z_{217} +2 z_{214} z_{218}
\end{aligned}\end{equation*}
% rootProduct.3.equation2.1
\begin{equation*}\begin{aligned}
Q_{175}={}&
4 z_{139} z_{219} -12 z_{139} z_{220} +z_{211} z_{216}\\
&+z_{212} z_{215} +z_{213} z_{218} +z_{214} z_{217}
\end{aligned}\end{equation*}
% rootProduct.3.equation3.0
\begin{equation*}\begin{aligned}
Q_{176}={}&
12 z_{139} z_{221} -8 z_{139} z_{222} +z_{211} z_{217}\\
&+2 z_{212} z_{218} +z_{213} z_{215} +2 z_{214} z_{216}
\end{aligned}\end{equation*}
% rootProduct.3.equation3.1
\begin{equation*}\begin{aligned}
Q_{177}={}&
-4 z_{139} z_{221} +12 z_{139} z_{222} +z_{211} z_{218}\\
&+z_{212} z_{217} +z_{213} z_{216} +z_{214} z_{215}
\end{aligned}\end{equation*}
% rootProduct.3.equation4.0
\begin{equation*}\begin{aligned}
Q_{178}={}&
-44 z_{139} z_{223} +48 z_{139} z_{224} +z_{211}^{2}\\
&+2 z_{212}^{2}
\end{aligned}\end{equation*}
% rootProduct.3.equation4.1
\begin{equation*}\begin{aligned}
Q_{179}={}&
12 z_{139} z_{223} -22 z_{139} z_{224} +z_{211} z_{212}
\end{aligned}\end{equation*}
% rootProduct.3.equation5.0
\begin{equation*}\begin{aligned}
Q_{180}={}&
z_{78} -z_{223}^{2} -2 z_{224}^{2}
\end{aligned}\end{equation*}
% rootProduct.3.equation5.1
\begin{equation*}\begin{aligned}
Q_{181}={}&
z_{79} -z_{223} z_{224}
\end{aligned}\end{equation*}
% rootProduct.3.equation6.0
\begin{equation*}\begin{aligned}
Q_{182}={}&
-1 +11 z_{78} -24 z_{79}\\
&+z_{219}^{2} +2 z_{220}^{2}
\end{aligned}\end{equation*}
% rootProduct.3.equation6.1
\begin{equation*}\begin{aligned}
Q_{183}={}&
-3 z_{78} +11 z_{79} +z_{219} z_{220}
\end{aligned}\end{equation*}
% rootProduct.3.equation7.0
\begin{equation*}\begin{aligned}
Q_{184}={}&
-1 +11 z_{43} -24 z_{44}\\
&+z_{221}^{2} +2 z_{222}^{2}
\end{aligned}\end{equation*}
% rootProduct.3.equation7.1
\begin{equation*}\begin{aligned}
Q_{185}={}&
-3 z_{43} +11 z_{44} +z_{221} z_{222}
\end{aligned}\end{equation*}
% rootProduct.3.equation8.0
\begin{equation*}\begin{aligned}
Q_{186}={}&
6 z_{29} z_{221} -4 z_{29} z_{222} -4 z_{30} z_{221}\\
&+12 z_{30} z_{222} +z_{213}
\end{aligned}\end{equation*}
% rootProduct.3.equation8.1
\begin{equation*}\begin{aligned}
Q_{187}={}&
-2 z_{29} z_{221} +6 z_{29} z_{222} +6 z_{30} z_{221}\\
&-4 z_{30} z_{222} +z_{214}
\end{aligned}\end{equation*}
% rootProduct.3.equation9.0
\begin{equation*}\begin{aligned}
Q_{188}={}&
z_{21} z_{211} +2 z_{22} z_{212} -6 z_{59} z_{219}\\
&+4 z_{59} z_{220} +4 z_{60} z_{219} -12 z_{60} z_{220}
\end{aligned}\end{equation*}
% rootProduct.3.equation9.1
\begin{equation*}\begin{aligned}
Q_{189}={}&
z_{21} z_{212} +z_{22} z_{211} +2 z_{59} z_{219}\\
&-6 z_{59} z_{220} -6 z_{60} z_{219} +4 z_{60} z_{220}
\end{aligned}\end{equation*}
% rootProduct.4.equation1.0
\begin{equation*}\begin{aligned}
Q_{190}={}&
z_{225}^{2} +2 z_{226}^{2} +z_{227}^{2}\\
&+2 z_{228}^{2} +z_{229}^{2} +2 z_{230}^{2}\\
&-z_{231}^{2} -2 z_{232}^{2}
\end{aligned}\end{equation*}
% rootProduct.4.equation1.1
\begin{equation*}\begin{aligned}
Q_{191}={}&
z_{225} z_{226} +z_{227} z_{228} +z_{229} z_{230}\\
&-z_{231} z_{232}
\end{aligned}\end{equation*}
% rootProduct.4.equation2.0
\begin{equation*}\begin{aligned}
Q_{192}={}&
-12 z_{139} z_{233} +8 z_{139} z_{234} +z_{225} z_{229}\\
&+2 z_{226} z_{230} +z_{227} z_{231} +2 z_{228} z_{232}
\end{aligned}\end{equation*}
% rootProduct.4.equation2.1
\begin{equation*}\begin{aligned}
Q_{193}={}&
4 z_{139} z_{233} -12 z_{139} z_{234} +z_{225} z_{230}\\
&+z_{226} z_{229} +z_{227} z_{232} +z_{228} z_{231}
\end{aligned}\end{equation*}
% rootProduct.4.equation3.0
\begin{equation*}\begin{aligned}
Q_{194}={}&
12 z_{139} z_{235} -8 z_{139} z_{236} +z_{225} z_{231}\\
&+2 z_{226} z_{232} +z_{227} z_{229} +2 z_{228} z_{230}
\end{aligned}\end{equation*}
% rootProduct.4.equation3.1
\begin{equation*}\begin{aligned}
Q_{195}={}&
-4 z_{139} z_{235} +12 z_{139} z_{236} +z_{225} z_{232}\\
&+z_{226} z_{231} +z_{227} z_{230} +z_{228} z_{229}
\end{aligned}\end{equation*}
% rootProduct.4.equation4.0
\begin{equation*}\begin{aligned}
Q_{196}={}&
-44 z_{139} z_{237} +48 z_{139} z_{238} +z_{225}^{2}\\
&+2 z_{226}^{2}
\end{aligned}\end{equation*}
% rootProduct.4.equation4.1
\begin{equation*}\begin{aligned}
Q_{197}={}&
12 z_{139} z_{237} -22 z_{139} z_{238} +z_{225} z_{226}
\end{aligned}\end{equation*}
% rootProduct.4.equation5.0
\begin{equation*}\begin{aligned}
Q_{198}={}&
z_{80} -z_{237}^{2} -2 z_{238}^{2}
\end{aligned}\end{equation*}
% rootProduct.4.equation5.1
\begin{equation*}\begin{aligned}
Q_{199}={}&
z_{81} -z_{237} z_{238}
\end{aligned}\end{equation*}
% rootProduct.4.equation6.0
\begin{equation*}\begin{aligned}
Q_{200}={}&
-1 +11 z_{80} -24 z_{81}\\
&+z_{233}^{2} +2 z_{234}^{2}
\end{aligned}\end{equation*}
% rootProduct.4.equation6.1
\begin{equation*}\begin{aligned}
Q_{201}={}&
-3 z_{80} +11 z_{81} +z_{233} z_{234}
\end{aligned}\end{equation*}
% rootProduct.4.equation7.0
\begin{equation*}\begin{aligned}
Q_{202}={}&
-1 +11 z_{47} -24 z_{48}\\
&+z_{235}^{2} +2 z_{236}^{2}
\end{aligned}\end{equation*}
% rootProduct.4.equation7.1
\begin{equation*}\begin{aligned}
Q_{203}={}&
-3 z_{47} +11 z_{48} +z_{235} z_{236}
\end{aligned}\end{equation*}
% rootProduct.4.equation8.0
\begin{equation*}\begin{aligned}
Q_{204}={}&
6 z_{31} z_{235} -4 z_{31} z_{236} -4 z_{32} z_{235}\\
&+12 z_{32} z_{236} +z_{227}
\end{aligned}\end{equation*}
% rootProduct.4.equation8.1
\begin{equation*}\begin{aligned}
Q_{205}={}&
-2 z_{31} z_{235} +6 z_{31} z_{236} +6 z_{32} z_{235}\\
&-4 z_{32} z_{236} +z_{228}
\end{aligned}\end{equation*}
% rootProduct.4.equation9.0
\begin{equation*}\begin{aligned}
Q_{206}={}&
z_{21} z_{225} +2 z_{22} z_{226} -6 z_{63} z_{233}\\
&+4 z_{63} z_{234} +4 z_{64} z_{233} -12 z_{64} z_{234}
\end{aligned}\end{equation*}
% rootProduct.4.equation9.1
\begin{equation*}\begin{aligned}
Q_{207}={}&
z_{21} z_{226} +z_{22} z_{225} +2 z_{63} z_{233}\\
&-6 z_{63} z_{234} -6 z_{64} z_{233} +4 z_{64} z_{234}
\end{aligned}\end{equation*}
% rootProduct.S.retained_j.0
\begin{equation*}\begin{aligned}
Q_{208}={}&
4 z_{2} +z_{197} z_{199} +2 z_{198} z_{200}\\
&-z_{201} z_{203} -2 z_{202} z_{204}
\end{aligned}\end{equation*}
% rootProduct.S.retained_j.1
\begin{equation*}\begin{aligned}
Q_{209}={}&
4 z_{3} +z_{197} z_{200} +z_{198} z_{199}\\
&-z_{201} z_{204} -z_{202} z_{203}
\end{aligned}\end{equation*}
% scalar.S.direct_ratio.0
\begin{equation*}\begin{aligned}
Q_{210}={}&
2 z_{2} z_{128} +4 z_{3} z_{129} -z_{21} z_{203}\\
&-2 z_{22} z_{204}
\end{aligned}\end{equation*}
% scalar.S.direct_ratio.1
\begin{equation*}\begin{aligned}
Q_{211}={}&
2 z_{2} z_{129} +2 z_{3} z_{128} -z_{21} z_{204}\\
&-z_{22} z_{203}
\end{aligned}\end{equation*}
% reciprocal.aggregate.rootTrace.definition
\begin{equation*}\begin{aligned}
Q_{212}={}&
8 z_{4} z_{110} -z_{189}^{2} -2 z_{190}^{2}\\
&-z_{203}^{2} -2 z_{204}^{2} -z_{217}^{2}\\
&-2 z_{218}^{2} -z_{231}^{2} -2 z_{232}^{2}
\end{aligned}\end{equation*}

\section*{Large integer coefficients}
The decimal digits on consecutive lines in each row are concatenated. Thus every $K_j$ below denotes one fixed positive integer; it is not a variable or a parameter.
\small
\begin{longtable}{r p{0.78\textwidth}}
$K_{1}$ & \texttt{9819499495210116462795099322880135176040}\newline \texttt{9466991314003358996422433259542652766869}\newline \texttt{1070368153600000000000000} \\[3pt]
$K_{2}$ & \texttt{8099049341013100078437597702364627434582}\newline \texttt{0490095244015289194989569855772668241364}\newline \texttt{832628280392076492800000000000000} \\[3pt]
$K_{3}$ & \texttt{4589461293240756711114638698006622212929}\newline \texttt{8277720638275330543827422918271178670106}\newline \texttt{7384893588888433459200000000000000} \\[3pt]
$K_{4}$ & \texttt{9178922586481513422229277396013244425859}\newline \texttt{6555441276550661087654845836542357340213}\newline \texttt{4769787177776866918400000000000000} \\[3pt]
$K_{5}$ & \texttt{1376838387972227013334391609401986663878}\newline \texttt{9483316191482599163148226875481353601032}\newline \texttt{02154680766665300377600000000000000} \\[3pt]
$K_{6}$ & \texttt{2753676775944454026668783218803973327757}\newline \texttt{8966632382965198326296453750962707202064}\newline \texttt{04309361533330600755200000000000000} \\[3pt]
$K_{7}$ & \texttt{6318946488646535718731759452292806620184}\newline \texttt{4497924729301762684931640143801547296154}\newline \texttt{1724103664894830836146098454811858898268}\newline \texttt{6734559746745916595796472604244988128011}\newline \texttt{1493975224936910108363237969187823419392}\newline \texttt{0000000000000000000000000000} \\[3pt]
$K_{8}$ & \texttt{5687051839781882146858583507063525958166}\newline \texttt{0048132256371586416438476129421392566538}\newline \texttt{7551693298405347752531488609330673008441}\newline \texttt{8061103772071324936216825343820489315210}\newline \texttt{0344577702443219097526914172269041077452}\newline \texttt{80000000000000000000000000000} \\[3pt]
\end{longtable}
\normalsize


\clearpage
\section{The 233-variable family: arithmetic and exact reductions}\label{v15x:reductions}
The 238-variable system above is retained as a predecessor, with its own
coordinates and polynomial. This section constructs the 233-variable predecessor used in
Sections~\ref{v222:construction}--\ref{v222:height}. Throughout this section the coefficient is
\(A_N=149^{2N+1}\), with \(N\in\N\). Statements about the defining property
are for all sufficiently large \(N\); statements about coefficient height
for \(N\ge53\) have a different, purely finite meaning.

We first work before the optional sparse presentation of the 238-variable
system. The preceding coordinate maps recover that raw presentation
without changing the input. The first step below changes the scalar
criterion and is proved arithmetically in both directions. The next two
steps are rational witness bijections on the admitted family. Finally,
invertible affine coordinate maps and constant row operations give the
207 displayed equations in Section~\ref{v15x:equations}.

\subsection{A three-residue scalar criterion: 238 to 237}\label{v15x:parity}
Use the raw 238-coordinate indices of the global-rescaling construction.
Write \(R=x_{142}>0\), \(R_2=x_{139}=R^2\), \(u=x_{128}\),
\(v=x_{129}\), \(D_1=x_{135}=v^2\), and \(w=x_{138}\).
The physical formal-coordinate ratio is
\(Z=R(u+\sqrt2v)=Z_0+\sqrt2Z_1\), and the physical scalar pole variable
is \(W=w/R^2\). Put \(f=f(H)\) and \(F=\Norm_{K/\Q}(f)\), where
\(K=\Q(\sqrt2)\). The retained coordinates satisfy
\[
 f^2=x_{86}+2\sqrt2x_{87},\qquad F^2=x_{86}^2-8x_{87}^2.
\]
Delete coordinate 137 and its defining row 82. Repurpose coordinate 136
as \(C\). Replace raw row 88 and raw row 83, respectively, by
\[
 u^2-RC=0,\qquad C^2+RD_1+x_{86}^2-8x_{87}^2-R_2w=0.
\]
All other rows are retained. Compact the indices above the deleted
coordinate by one; compact the row indices after row 82 by one.
Thus there are 237 coordinates and 211 quadratic residuals. These
instructions refer to the raw predecessor, not its affine sparse basis.
In physical variables the new equations say
\begin{equation}\label{v15x:W}
 W=\frac{Z_0^4}{R^{10}}+\frac{Z_1^2}{R^5}+\frac{F^2}{R^4},
 \qquad b=R(R+9792)W^2.
\end{equation}
The definitions of the positive normalizer, the pole certificate, and
all non-scalar interfaces remain unchanged.

\begin{lemma}\label{v15x:scalar-local}
At each consumed scalar contact prime, with \(\nu=v_p(R)>0\) odd,
the new criterion implies
\[
 v_p(Z_0)>19\nu/8,\qquad v_p(Z_1)>9\nu/4,
 \qquad v_p(F)>7\nu/4.
\]
The convention \(v_p(0)=+\infty\) is allowed for either component of \(Z\).
\end{lemma}
\begin{proof}
The primes in question are \(149\) or primes greater than \(76176\),
so \(9792\) is a unit and the curve has good unramified reduction.
First, \(W>0\): \(R>0\) and \(F\ne0\), so the last summand
in~\eqref{v15x:W} is strictly positive. Thus \(b\ne0\).
The unchanged pole obstruction rules out a negative odd valuation of
\(b\). Since \(v_p(b)=\nu+2v_p(W)\) is odd, it follows that
\(v_p(W)>-\nu/2\).

For the actual point on \(y^2=x^3-Dx\), with \(D=11-6\sqrt2\),
the finite valuation of \(x\), hence that of \(f=1/x\), is even:
if \(v(x)>0\), then \(2v(y)=v(x)\); if \(v(x)<0\), then
\(2v(y)=3v(x)\). If \(v(x)=0\) the conclusion is immediate.
Consequently \(v_p(F)\) is even, both at split primes (sum of the two
even valuations) and at inert primes (residue degree two).
The three summands of \(R^4W\) have valuations
\[
 4v_p(Z_0)-6\nu\in2+4\Z,\quad
 2v_p(Z_1)-\nu\in1+2\Z,\quad
 2v_p(F)\in4\Z.
\]
Their finite valuations are distinct, so their minimum cannot cancel.
Omit a zero summand. The last summand is nonzero. Since
\(v_p(R^4W)>7\nu/2\), each finite valuation exceeds \(7\nu/2\).
Solving the three inequalities gives the assertion.
\end{proof}

The retained formal-group argument now applies with the same physical
ratio \(Z\). At a suitable place, \(H\) is formal, and the allowed
two-torsion translation makes the comparison point formal. With
\(H=[m]Q+T\), \(P=2Q\), and comparison index \(a\ne0\), the identity
\([2m]S'=[4a]H\) and the formal-logarithm isometry give
\(v_p(a)>9\nu/4\). This uses no extra divisibility hypothesis on \(m\).
In particular every selected contact contributes exponent at least three.
The forced normalizer still has
\(v_{149}(R)=2N+1\) or \(2N-1\). The bound on \(F\) excludes the same
finite exceptional reference points before their ratios are interpreted.
Writing \(c_0=|a|\) and
\[
 L'_N=\lfloor9(2N-1)/4\rfloor+1,
\]
the unchanged contact and aggregate argument gives, for sufficiently
large positive \(N\),
\[
 149^{L'_N}M_5(\beta)^3\mid c_0,\qquad
 h(\beta)\le C_3\,149^{-L'_N}c_0(1+\log(c_0)/3).
\]
Here \(M_5\), \(\beta\), and \(C_3\) retain their meanings in the
preceding height proof. All five-factor tests and both aggregate
contacts are unchanged. The first rationality and Lucas-defect
arguments therefore still apply. In the last contradiction replace
\(L_N\) by \(L'_N\): the upper bound
\[
 \log B\le C'+7C_3\,149^{-L'_N}c_0(1+\log(c_0)/3)
\]
is incompatible with \(\log B\ge(c_0-4)\log(c_0-1)\) for all
sufficiently large \(N\), uniformly in the input, because
\(c_0\ge149^{L'_N}\) and \(L'_N\to\infty\).
This proves soundness with one, still ineffective, threshold.

For completeness, retain the compact real neighborhoods for the
reference and comparison points used in the scalar-S construction.
On these neighborhoods \(Z_0,Z_1,F\) are uniformly bounded.
The unchanged normalizer has
\(R\ge A_N/(1+149^2)\), so~\eqref{v15x:W} gives
\(W=O(R^{-4})\) and \(b=O(R^{-6})\). In particular \(0<b<17\)
for sufficiently large \(N\), uniformly in the integer input.
Its squareclass is still \(R(R+9792)\); its reduced denominator
is a square. At the finite support of this squareclass, first choose
\(H\) sufficiently formal and then the comparison index sufficiently
divisible. This makes \(F\) and both components of \(Z\) arbitrarily
deep. Thus the three terms of~\eqref{v15x:W} can be made integral
at every odd-valuation numerator prime and sufficiently deep at
\(2,3,17\), despite the fixed powers of \(R\) in their denominators.
At all other primes \(b\) has even valuation. These are precisely
the retained pole-realization hypotheses.

Compact-group recurrence meets these finite conditions and the real
neighborhoods simultaneously. After fixing the comparison point,
continuity preserves the conditions for \(H\) in a fixed progression.
The later contact selection avoids the finite support already fixed;
its ordering, nonzero Lucas labels, opposite-order conditions, and
four pairs of root points are unchanged. Finally set
\(C=u^2/R\), \(D_1=v^2\), and \(w=R^2W\), and fill the existing
scaled pole witnesses. These are rational and satisfy the changed
rows. The positive-normalizer proof also holds when the distinguished
input is zero. This proves completeness for all integers with the
same type of uniform, ineffective threshold. It does not assert a
bijection of the old and new scalar witnesses.

\subsection{Reference identities and removal of two inverse coordinates}
\label{v15x:235}
Now use raw 237-coordinate indices, and put
\[
 f=x_{19}+\sqrt2x_{20},\quad S=x_{21}+\sqrt2x_{22},\quad
 v=x_{23}+\sqrt2x_{24},\quad g=x_{139}+\sqrt2x_{140},
 \quad \chi=x_{17}+2\sqrt2x_{18}.
\]
Set
\[
 a=57-28\sqrt2,\quad b_0=8-14\sqrt2,\quad
 c=88-48\sqrt2,\quad d=39-36\sqrt2,\quad A=a+cf.
\]
The notation \(b_0\) distinguishes this fixed field constant from the
physical pole input \(b\) above; \(A\) is not the coefficient \(A_N\).
The old reference rows are
\[
 \chi f=1,\qquad S^2=\chi-Df,\qquad
 G:=g^2-fv=0,\qquad (a\chi+c)g+b_0S+d=0.
\]
Replace the inverse and reference rows by the two field equations
\begin{align}
 \mathcal F&=Ag+b_0fS+df=0,\label{v15x:F}\\
 \mathcal E&=(b_0S-d)g-af+8=0.\label{v15x:E}
\end{align}
Retain \(G=0\). Delete coordinates 17 and 18, and raw rows
1--4 and 88--89. Insert the four rational component rows of
\eqref{v15x:F}--\eqref{v15x:E}; every other row is unchanged under
\(i\mapsto i\) for \(i<17\) and \(i\mapsto i-2\) for \(i>18\).
The result has 235 coordinates and 209 rows.

Direct expansion in \(\Q(\sqrt2)\) gives
\begin{align}
 A(af-8)&=b_0^2Df^2+d^2f-b_0^2,\label{v15x:constant-identity}\\
 \mathcal F(b_0S-d)-A\mathcal E
 &=b_0^2(fS^2-1+Df^2).\label{v15x:saturation}
\end{align}
For a new solution, the second identity proves the old curve equation
and also proves \(f\ne0\), since its right-hand factor at \(f=0\)
is \(-1\). Define \(\chi=1/f\). Division of~\eqref{v15x:F} by
\(f\) recovers the old reference equation, and the recovered curve
equation is \(S^2=\chi-Df\). This direction requires no division
by \(A,S,g\), or \(v\).

For the forward direction, an old solution gives \(\mathcal F=0\)
and \(A\mathcal E=0\). We must exclude \(A=0\), rather than discard
that stratum without proof. Its only possible value is
\[
 f_*=-a/c=-291/392-(17/196)\sqrt2,\qquad
 \Norm(f_*)=(41/56)^2.
\]
This norm is a 149-adic unit. The scalar criterion of
Section~\ref{v15x:parity}, at the forced positive odd contact, gives
\(v_{149}(\Norm f)>7\nu/4>0\). Therefore \(f\ne f_*\), and
\(\mathcal E=0\). This proves both directions on the admitted
\(A_N\) family in the range used for the defining theorem.

Writing \(f=f_0+\sqrt2f_1\), the inverse coordinate map is explicitly
\[
 x_{17}=\frac{f_0}{f_0^2-2f_1^2},\qquad
 x_{18}=\frac{-f_1}{2(f_0^2-2f_1^2)}.
\]
The denominator is nonzero because \(f\ne0\) in the field.
Either rational component may vanish. Every retained coordinate is
transported by the inverse index map. Thus these maps are inverse
on the full solution sets under consideration, not just on prepared
completeness witnesses.

\subsection{Keeping the square, eliminating the reference value: 235 to 233}
\label{v15x:233}
In raw 235 coordinates, the retained fields are
\[
 f=x_{17}+\sqrt2x_{18},\quad S=x_{19}+\sqrt2x_{20},\quad
 v=x_{21}+\sqrt2x_{22},\quad g=x_{137}+\sqrt2x_{138},
\]
with \(B=x_{84}+2\sqrt2x_{85}=f^2\) and
\(w=x_{31}+2\sqrt2x_{32}=v^2\).
All occurrences of \(v\) are in the two component rows of \(G=0\)
and the two component rows defining \(w=v^2\). Put
\[
 T=8-af-2dg.
\]
Delete coordinates 21 and 22 and rows 5--8, and replace the four rows
by the two rational components of
\begin{equation}\label{v15x:square}
 (a^2+2acf+c^2B)w-T^2=0.
\end{equation}
This is quadratic in the retained coordinates: \(B\) is an existing
coordinate pair, not an instruction to substitute a fresh square into
the displayed residual. All 205 other rows are retained. Put the two
new rows first; the old rows 1--4 follow, then the old rows 9--209.
Compact coordinate indices above 22 by two. This specifies the
corrected raw 233-coordinate, 207-row presentation.

Here is a two-sided verification, including the nonvanishing needed
for the inverse. Define
\[
 Q=Av-T,\qquad R_1=Av+b_0Sg+dg.
\]
Then \(R_1=Q+\mathcal E\), and direct multiplication gives
\begin{equation}\label{v15x:Gidentity}
 \mathcal Fg-AG=fR_1.
\end{equation}
On an old solution, \(f\ne0\), \(\mathcal F=\mathcal E=G=0\),
so \(R_1=Q=0\). Together with \(w=v^2\), this proves
\eqref{v15x:square}. Conversely, the retained
\(\mathcal F=\mathcal E=0\) first recovers
\(fS^2-1+Df^2=0\) by~\eqref{v15x:saturation}, hence \(f\ne0\).
The unchanged scalar argument at 149 can now be applied to that
actual reference point, before interpreting \(v\) or \(w\).
It excludes \(A=0\) by the same norm calculation as above.
Set \(v=T/A\). Since \(B=f^2\), equation~\eqref{v15x:square}
gives \(w=v^2\). Now \(Q=0\), hence \(R_1=0\), and
\eqref{v15x:Gidentity} gives \(G=0\). All deleted rows are restored.

For explicit rational formulas in the corrected raw 233 coordinates,
write \(f=f_0+\sqrt2f_1\), \(g=g_0+\sqrt2g_1\), where
\((f_0,f_1)=(x_{17},x_{18})\) and
\((g_0,g_1)=(x_{135},x_{136})\). Then
\begin{align*}
 A_0&=57+88f_0-96f_1,& A_1&=-28+88f_1-48f_0,\\
 T_0&=8-57f_0+56f_1-78g_0+144g_1,&
 T_1&=28f_0-57f_1+72g_0-78g_1,\\
 v_0&=(T_0A_0-2T_1A_1)/(A_0^2-2A_1^2),&
 v_1&=(T_1A_0-T_0A_1)/(A_0^2-2A_1^2).
\end{align*}
Neither component of \(A\) is assumed nonzero separately. Only its
field norm is inverted. No division by \(T,g,S,w\), or a component
of \(f\) occurs. In particular the maps cover their zero strata.
The input coordinate and \(A_N\) are fixed. The two inverse removals
introduce no new arithmetic assumption and do not make the threshold
effective.

\subsection{Exact sparse presentation and its inverse}\label{v15x:maps}
All indices in this subsection refer to the corrected raw 233 system.
The input is coordinate 1 and is fixed by every stage. A coordinate stage
means substitute the displayed old coordinates as affine expressions in
new coordinates, simultaneously within that stage; all unspecified
coordinates are fixed. Stages are applied in the order printed, without
commuting coordinate changes past row operations. A row operation acts on
the current list of 207 residuals; its indices below are one-based.
These operations preserve simultaneous vanishing over the rationals.
\paragraph{Stage 1.}
Lucas P/E basis.
\[x_{64}=2 y_{59} +y_{64}.\]
\[x_{65}=2 y_{60} +y_{65}.\]
\[x_{66}=y_{64} +y_{66}.\]
\[x_{67}=y_{65} +y_{67}.\]
\paragraph{Stage 2.}
\[R_{91}\longleftarrow R_{91} + (-36)R_{32}.\]
\paragraph{Stage 3.}
\[R_{9}\longleftarrow R_{9} + (\frac{-23}{104})R_{10}.\]
\paragraph{Stage 4.}
\[R_{10}\longleftarrow R_{10} + (\frac{-2392}{4879})R_{9}.\]
\paragraph{Stage 5.}
\[R_{13}\longleftarrow R_{13} + (\frac{-23}{104})R_{14}.\]
\paragraph{Stage 6.}
\[R_{14}\longleftarrow R_{14} + (\frac{-2392}{4879})R_{13}.\]
\paragraph{Stage 7.}
\[R_{17}\longleftarrow R_{17} + (\frac{-23}{104})R_{18}.\]
\paragraph{Stage 8.}
\[R_{18}\longleftarrow R_{18} + (\frac{-2392}{4879})R_{17}.\]
\paragraph{Stage 9.}
\[R_{21}\longleftarrow R_{21} + (\frac{-52}{23})R_{22}.\]
\paragraph{Stage 10.}
\[R_{22}\longleftarrow R_{22} + (\frac{2392}{4879})R_{21}.\]
\paragraph{Stage 11.}
\[R_{71}\longleftarrow R_{71} + (-381485)R_{70}.\]
\paragraph{Stage 12.}
\[R_{92}\longleftarrow R_{92} + (-72)R_{33}.\]
\paragraph{Stage 13.}
\[R_{7}\longleftarrow R_{7} + (\frac{-193}{33})R_{8}.\]
\paragraph{Stage 14.}
\[R_{11}\longleftarrow R_{11} + (\frac{-193}{33})R_{12}.\]
\paragraph{Stage 15.}
\[R_{15}\longleftarrow R_{15} + (\frac{-193}{33})R_{16}.\]
\paragraph{Stage 16.}
\[R_{19}\longleftarrow R_{19} + (\frac{-193}{33})R_{20}.\]
\paragraph{Stage 17.}
\[R_{46}\longleftarrow R_{46} + (3813373161796)R_{47}.\]
\paragraph{Stage 18.}
\[R_{59}\longleftarrow R_{59} + (3813373161796)R_{60}.\]
\paragraph{Stage 19.}
four constant shifts.
\[x_{135}=-\frac{1672}{18921} +y_{135}.\]
\[x_{63}=-\frac{2501}{51} +y_{63}.\]
\[x_{80}=-9 +y_{80}.\]
\[x_{90}=-1 +y_{90}.\]
\paragraph{Stage 20.}
non-Lucas shear.
\[x_{31}=y_{31} +y_{32}.\]
\paragraph{Stage 21.}
non-Lucas shear.
\[x_{37}=y_{37} +y_{38}.\]
\paragraph{Stage 22.}
non-Lucas shear.
\[x_{41}=y_{41} +y_{42}.\]
\paragraph{Stage 23.}
non-Lucas shear.
\[x_{45}=y_{45} +y_{46}.\]
\paragraph{Stage 24.}
non-Lucas shear.
\[x_{122}=y_{121} +y_{122}.\]
\paragraph{Stage 25.}
four root sum/difference bases.
\[x_{178}=\frac{1}{2} y_{178} +\frac{1}{2} y_{180}.\]
\[x_{180}=\frac{1}{2} y_{178} -\frac{1}{2} y_{180}.\]
\[x_{182}=\frac{1}{2} y_{182} +\frac{1}{2} y_{184}.\]
\[x_{184}=\frac{1}{2} y_{182} -\frac{1}{2} y_{184}.\]
\[x_{179}=\frac{1}{2} y_{179} +\frac{1}{2} y_{181}.\]
\[x_{181}=\frac{1}{2} y_{179} -\frac{1}{2} y_{181}.\]
\[x_{183}=\frac{1}{2} y_{183} +\frac{1}{2} y_{185}.\]
\[x_{185}=\frac{1}{2} y_{183} -\frac{1}{2} y_{185}.\]
\[x_{192}=\frac{1}{2} y_{192} +\frac{1}{2} y_{194}.\]
\[x_{194}=\frac{1}{2} y_{192} -\frac{1}{2} y_{194}.\]
\[x_{196}=\frac{1}{2} y_{196} +\frac{1}{2} y_{198}.\]
\[x_{198}=\frac{1}{2} y_{196} -\frac{1}{2} y_{198}.\]
\[x_{193}=\frac{1}{2} y_{193} +\frac{1}{2} y_{195}.\]
\[x_{195}=\frac{1}{2} y_{193} -\frac{1}{2} y_{195}.\]
\[x_{197}=\frac{1}{2} y_{197} +\frac{1}{2} y_{199}.\]
\[x_{199}=\frac{1}{2} y_{197} -\frac{1}{2} y_{199}.\]
\[x_{206}=\frac{1}{2} y_{206} +\frac{1}{2} y_{208}.\]
\[x_{208}=\frac{1}{2} y_{206} -\frac{1}{2} y_{208}.\]
\[x_{210}=\frac{1}{2} y_{210} +\frac{1}{2} y_{212}.\]
\[x_{212}=\frac{1}{2} y_{210} -\frac{1}{2} y_{212}.\]
\[x_{207}=\frac{1}{2} y_{207} +\frac{1}{2} y_{209}.\]
\[x_{209}=\frac{1}{2} y_{207} -\frac{1}{2} y_{209}.\]
\[x_{211}=\frac{1}{2} y_{211} +\frac{1}{2} y_{213}.\]
\[x_{213}=\frac{1}{2} y_{211} -\frac{1}{2} y_{213}.\]
\[x_{220}=\frac{1}{2} y_{220} +\frac{1}{2} y_{222}.\]
\[x_{222}=\frac{1}{2} y_{220} -\frac{1}{2} y_{222}.\]
\[x_{224}=\frac{1}{2} y_{224} +\frac{1}{2} y_{226}.\]
\[x_{226}=\frac{1}{2} y_{224} -\frac{1}{2} y_{226}.\]
\[x_{221}=\frac{1}{2} y_{221} +\frac{1}{2} y_{223}.\]
\[x_{223}=\frac{1}{2} y_{221} -\frac{1}{2} y_{223}.\]
\[x_{225}=\frac{1}{2} y_{225} +\frac{1}{2} y_{227}.\]
\[x_{227}=\frac{1}{2} y_{225} -\frac{1}{2} y_{227}.\]
\paragraph{Stage 26.}
\[(R_{133},R_{135})\longleftarrow(R_{133}+R_{135},R_{133}-R_{135}).\]
\paragraph{Stage 27.}
\[(R_{134},R_{136})\longleftarrow(R_{134}+R_{136},R_{134}-R_{136}).\]
\paragraph{Stage 28.}
\[(R_{151},R_{153})\longleftarrow(R_{151}+R_{153},R_{151}-R_{153}).\]
\paragraph{Stage 29.}
\[(R_{152},R_{154})\longleftarrow(R_{152}+R_{154},R_{152}-R_{154}).\]
\paragraph{Stage 30.}
\[(R_{169},R_{171})\longleftarrow(R_{169}+R_{171},R_{169}-R_{171}).\]
\paragraph{Stage 31.}
\[(R_{170},R_{172})\longleftarrow(R_{170}+R_{172},R_{170}-R_{172}).\]
\paragraph{Stage 32.}
\[(R_{187},R_{189})\longleftarrow(R_{187}+R_{189},R_{187}-R_{189}).\]
\paragraph{Stage 33.}
\[(R_{188},R_{190})\longleftarrow(R_{188}+R_{190},R_{188}-R_{190}).\]
\paragraph{Stage 34.}
centered Lucas defect and square.
\[x_{68}=1 +y_{68}.\]
\[x_{90}=2 +2 y_{68} +y_{90}.\]
\[x_{91}=y_{69} +y_{91}.\]
After all stages, multiply each row by the following nonzero factor;
unlisted factors are 1. This clears denominators and makes each row
primitive. The sign convention is the one in the displayed equations.
\begin{longtable}{rr}Row&Multiplier\\\hline
1&$39778249$\\
2&$39778249$\\
3&$18921$\\
4&$18921$\\
5&$18921$\\
6&$18921$\\
7&$33$\\
9&$5512$\\
10&$\frac{258587}{104}$\\
11&$33$\\
13&$5512$\\
14&$\frac{258587}{104}$\\
15&$33$\\
17&$5512$\\
18&$\frac{258587}{104}$\\
19&$33$\\
21&$1219$\\
22&$\frac{258587}{23}$\\
23&$51$\\
71&$\frac{1}{4}$\\
131&$2$\\
132&$2$\\
133&$2$\\
134&$2$\\
135&$2$\\
136&$2$\\
149&$2$\\
150&$2$\\
151&$2$\\
152&$2$\\
153&$2$\\
154&$2$\\
167&$2$\\
168&$2$\\
169&$2$\\
170&$2$\\
171&$2$\\
172&$2$\\
185&$2$\\
186&$2$\\
187&$2$\\
188&$2$\\
189&$2$\\
190&$2$\\
204&$2$\\
\end{longtable}
For completeness the inverse coordinate substitutions are printed next,
in reverse stage order. Here $y$ denotes the new coordinate and $x$ the
old coordinate for that individual stage. They are simultaneous within
a stage. Thus one obtains final coordinates from a raw witness by solving
the forward stages in their original order; one recovers the raw witness
from a final witness by substituting old-in-new formulas in reverse order.
Row addition is inverted by subtraction, sum/difference by half-sum and
half-difference, and row normalization by its reciprocal. These inverses
introduce only fixed rational denominators.
\paragraph{Inverse coordinate stage 1.}
\[y_{68}=-1 +x_{68}.\]
\[y_{90}=-2 x_{68} +x_{90}.\]
\[y_{91}=-x_{69} +x_{91}.\]
\paragraph{Inverse coordinate stage 2.}
\[y_{178}=x_{178} +x_{180}.\]
\[y_{179}=x_{179} +x_{181}.\]
\[y_{180}=x_{178} -x_{180}.\]
\[y_{181}=x_{179} -x_{181}.\]
\[y_{182}=x_{182} +x_{184}.\]
\[y_{183}=x_{183} +x_{185}.\]
\[y_{184}=x_{182} -x_{184}.\]
\[y_{185}=x_{183} -x_{185}.\]
\[y_{192}=x_{192} +x_{194}.\]
\[y_{193}=x_{193} +x_{195}.\]
\[y_{194}=x_{192} -x_{194}.\]
\[y_{195}=x_{193} -x_{195}.\]
\[y_{196}=x_{196} +x_{198}.\]
\[y_{197}=x_{197} +x_{199}.\]
\[y_{198}=x_{196} -x_{198}.\]
\[y_{199}=x_{197} -x_{199}.\]
\[y_{206}=x_{206} +x_{208}.\]
\[y_{207}=x_{207} +x_{209}.\]
\[y_{208}=x_{206} -x_{208}.\]
\[y_{209}=x_{207} -x_{209}.\]
\[y_{210}=x_{210} +x_{212}.\]
\[y_{211}=x_{211} +x_{213}.\]
\[y_{212}=x_{210} -x_{212}.\]
\[y_{213}=x_{211} -x_{213}.\]
\[y_{220}=x_{220} +x_{222}.\]
\[y_{221}=x_{221} +x_{223}.\]
\[y_{222}=x_{220} -x_{222}.\]
\[y_{223}=x_{221} -x_{223}.\]
\[y_{224}=x_{224} +x_{226}.\]
\[y_{225}=x_{225} +x_{227}.\]
\[y_{226}=x_{224} -x_{226}.\]
\[y_{227}=x_{225} -x_{227}.\]
\paragraph{Inverse coordinate stage 3.}
\[y_{122}=-x_{121} +x_{122}.\]
\paragraph{Inverse coordinate stage 4.}
\[y_{45}=x_{45} -x_{46}.\]
\paragraph{Inverse coordinate stage 5.}
\[y_{41}=x_{41} -x_{42}.\]
\paragraph{Inverse coordinate stage 6.}
\[y_{37}=x_{37} -x_{38}.\]
\paragraph{Inverse coordinate stage 7.}
\[y_{31}=x_{31} -x_{32}.\]
\paragraph{Inverse coordinate stage 8.}
\[y_{63}=\frac{2501}{51} +x_{63}.\]
\[y_{80}=9 +x_{80}.\]
\[y_{90}=1 +x_{90}.\]
\[y_{135}=\frac{1672}{18921} +x_{135}.\]
\paragraph{Inverse coordinate stage 9.}
\[y_{64}=-2 x_{59} +x_{64}.\]
\[y_{65}=-2 x_{60} +x_{65}.\]
\[y_{66}=2 x_{59} -x_{64} +x_{66}.\]
\[y_{67}=2 x_{60} -x_{65} +x_{67}.\]
\subsection{Coordinate dictionary and index bridge}\label{v15x:dictionary}
The dictionary below names the raw 233 coordinates before the preceding
sparse stages. A name records its inherited role; after a coordinate
stage its value is determined by that stage, not by reimposing its old
definition. This convention is essential for shifted and sheared slots.
The raw 238 to raw 237 map deletes 137; the raw 237 to raw 235 map deletes
17 and 18; the raw 235 to raw 233 map deletes 21 and 22. At each deletion
all higher indices decrease by the number of deleted indices below them.
The exceptional repurposed scalar slot is described in
Section~\ref{v15x:parity}. These index maps together with the full
predecessor dictionary specify every retained arithmetic variable.
\begin{longtable}{r p{0.79\textwidth}}Coordinate&Raw inherited role\\\hline
1&\nolinkurl{input}\\
2&\nolinkurl{S.chart.j0.rescaled}\\
3&\nolinkurl{S.chart.j1.rescaled}\\
4&\nolinkurl{reciprocal.rho}\\
5&\nolinkurl{detector0.normalized.scaled_e}\\
6&\nolinkurl{detector0.normalized.s}\\
7&\nolinkurl{detector0.normalized.p}\\
8&\nolinkurl{detector0.factored.q}\\
9&\nolinkurl{detector1.normalized.scaled_e}\\
10&\nolinkurl{detector1.normalized.s}\\
11&\nolinkurl{detector1.normalized.p}\\
12&\nolinkurl{detector1.factored.q}\\
13&\nolinkurl{normalized_scalar.scaled_e}\\
14&\nolinkurl{normalized_scalar.s}\\
15&\nolinkurl{normalized_scalar.p}\\
16&\nolinkurl{normalized_scalar.factored.q}\\
17&\nolinkurl{U0.f.0}\\
18&\nolinkurl{U0.f.1}\\
19&\nolinkurl{primary.negativeReciprocalZ0.0}\\
20&\nolinkurl{primary.negativeReciprocalZ0.1}\\
21&\nolinkurl{V1.z.0}\\
22&\nolinkurl{V1.z.1}\\
23&\nolinkurl{V2.z.0.rescaled}\\
24&\nolinkurl{V2.z.1.rescaled}\\
25&\nolinkurl{V3.z.0.rescaled}\\
26&\nolinkurl{V3.z.1.rescaled}\\
27&\nolinkurl{V4.z.0.rescaled}\\
28&\nolinkurl{V4.z.1.rescaled}\\
29&\nolinkurl{logV0.fSquare.0.quadratic}\\
30&\nolinkurl{logV0.fSquare.1.quadratic}\\
31&\nolinkurl{logV0.wPower2.0.quadratic}\\
32&\nolinkurl{logV0.wPower2.1.quadratic}\\
33&\nolinkurl{logV1.fSquare.0.quadratic}\\
34&\nolinkurl{logV1.fSquare.1.quadratic}\\
35&\nolinkurl{logV2.fSquare.0.quadratic}\\
36&\nolinkurl{logV2.fSquare.1.quadratic}\\
37&\nolinkurl{logV2.wPower2.0.quadratic}\\
38&\nolinkurl{logV2.wPower2.1.quadratic}\\
39&\nolinkurl{logV3.fSquare.0.quadratic}\\
40&\nolinkurl{logV3.fSquare.1.quadratic}\\
41&\nolinkurl{logV3.wPower2.0.quadratic}\\
42&\nolinkurl{logV3.wPower2.1.quadratic}\\
43&\nolinkurl{logV4.fSquare.0.quadratic}\\
44&\nolinkurl{logV4.fSquare.1.quadratic}\\
45&\nolinkurl{logV4.wPower2.0.quadratic}\\
46&\nolinkurl{logV4.wPower2.1.quadratic}\\
47&\nolinkurl{L1.0}\\
48&\nolinkurl{L1.1}\\
49&\nolinkurl{LP1.0}\\
50&\nolinkurl{LP1.1}\\
51&\nolinkurl{L2.0.rescaled}\\
52&\nolinkurl{L2.1.rescaled}\\
53&\nolinkurl{LP2.0.rescaled}\\
54&\nolinkurl{LP2.1.rescaled}\\
55&\nolinkurl{L3.0.rescaled}\\
56&\nolinkurl{L3.1.rescaled}\\
57&\nolinkurl{LP3.0.rescaled}\\
58&\nolinkurl{LP3.1.rescaled}\\
59&\nolinkurl{L4.0.rescaled}\\
60&\nolinkurl{L4.1.rescaled}\\
61&\nolinkurl{LP4.0.rescaled}\\
62&\nolinkurl{LP4.1.rescaled}\\
63&\nolinkurl{inputSquare.quadratic}\\
64&\nolinkurl{Lucas.C.0.quadratic.rescaled}\\
65&\nolinkurl{Lucas.C.1.quadratic.rescaled}\\
66&\nolinkurl{Lucas.AT.0.quadratic.rescaled}\\
67&\nolinkurl{Lucas.AT.1.quadratic.rescaled}\\
68&\nolinkurl{Lucas.F0.quadratic}\\
69&\nolinkurl{Lucas.F1.quadratic}\\
70&\nolinkurl{testSquare.1.0.0.quadratic}\\
71&\nolinkurl{testSquare.1.0.1.quadratic}\\
72&\nolinkurl{testSquare.1.1.0.quadratic}\\
73&\nolinkurl{testSquare.1.1.1.quadratic}\\
74&\nolinkurl{testSquare.1.2.0.quadratic}\\
75&\nolinkurl{testSquare.1.2.1.quadratic}\\
76&\nolinkurl{testSquare.1.3.0.quadratic}\\
77&\nolinkurl{testSquare.1.3.1.quadratic}\\
78&\nolinkurl{testSquare.12.0.0.quadratic}\\
79&\nolinkurl{testSquare.12.0.1.quadratic}\\
80&\nolinkurl{testSquare.12.1.0.quadratic}\\
81&\nolinkurl{testSquare.12.1.1.quadratic}\\
82&\nolinkurl{testSquare.40.0.0.quadratic}\\
83&\nolinkurl{testSquare.40.0.1.quadratic}\\
84&\nolinkurl{testSquare.40.1.0.quadratic}\\
85&\nolinkurl{testSquare.40.1.1.quadratic}\\
86&\nolinkurl{testSquare.40.2.0.quadratic}\\
87&\nolinkurl{testSquare.40.2.1.quadratic}\\
88&\nolinkurl{testSquare.40.3.0.quadratic}\\
89&\nolinkurl{testSquare.40.3.1.quadratic}\\
90&\nolinkurl{testSquare.40.4.0.quadratic}\\
91&\nolinkurl{testSquare.40.4.1.quadratic}\\
92&\nolinkurl{reciprocal.rhoSquared}\\
93&\nolinkurl{reciprocal.rhoFourth}\\
94&\nolinkurl{reciprocal.contactPolynomial}\\
95&\nolinkurl{detector0.q1}\\
96&\nolinkurl{reflected.detector0.weightPower7}\\
97&\nolinkurl{detector0.E1.power2.quadratic}\\
98&\nolinkurl{detector0.E1.power4.quadratic}\\
99&\nolinkurl{detector0.actualW.0}\\
100&\nolinkurl{detector0.actualW.1}\\
101&\nolinkurl{detector0.normOver2w0}\\
102&\nolinkurl{detector0.B}\\
103&\nolinkurl{detector0.normalized.p2.quadratic}\\
104&\nolinkurl{detector0.factored.v}\\
105&\nolinkurl{detector0.factored.q2}\\
106&\nolinkurl{detector1.t.quadratic.rescaled}\\
107&\nolinkurl{symmetric.detector1.deltaSquared.rescaled}\\
108&\nolinkurl{symmetric.detector1.deltaFourthPower.rescaled}\\
109&\nolinkurl{symmetric.detector1.contactPolynomial.rescaled}\\
110&\nolinkurl{detector1.q1}\\
111&\nolinkurl{reflected.detector1.weightPower7}\\
112&\nolinkurl{detector1.E1.power2.quadratic}\\
113&\nolinkurl{detector1.E1.power4.quadratic}\\
114&\nolinkurl{detector1.actualW.0}\\
115&\nolinkurl{detector1.actualW.1}\\
116&\nolinkurl{detector1.normOver2w0}\\
117&\nolinkurl{detector1.B}\\
118&\nolinkurl{detector1.normalized.p2.quadratic}\\
119&\nolinkurl{detector1.factored.v}\\
120&\nolinkurl{detector1.factored.q2}\\
121&\nolinkurl{scalar.rPair1.quadratic}\\
122&\nolinkurl{scalar.rFour.quadratic}\\
123&\nolinkurl{scalar.contactProduct}\\
124&\nolinkurl{scalar.zRatio.0.rescaled}\\
125&\nolinkurl{scalar.zRatio.1.rescaled}\\
126&\nolinkurl{scalar.W2.quadratic.rescaled}\\
127&\nolinkurl{scalar.b.quadratic}\\
128&\nolinkurl{normalized_scalar.p2.quadratic}\\
129&\nolinkurl{normalized_scalar.factored.v}\\
130&\nolinkurl{normalized_scalar.factored.q2}\\
131&\nolinkurl{scalar.ratio1_unscaled_square.rescaled}\\
132&\nolinkurl{parity_scalar.C_Z0square_over_R}\\
133&\nolinkurl{reduced_scalar.root.rescaled}\\
134&\nolinkurl{scalar.r_power_2}\\
135&\nolinkurl{V0.reciprocal.gOverH.0}\\
136&\nolinkurl{V0.reciprocal.gOverH.1}\\
137&\nolinkurl{scalar.forced_positive_parameter}\\
138&\nolinkurl{reciprocal.rhoFifth}\\
139&\nolinkurl{reflected.detector1.deltaSixthPower.rescaled}\\
140&\nolinkurl{amplified.C1.power4.0}\\
141&\nolinkurl{amplified.C1.power4.1}\\
142&\nolinkurl{residue.I.power4.0}\\
143&\nolinkurl{residue.I.power4.1}\\
144&\nolinkurl{residue.V1.power4.0}\\
145&\nolinkurl{residue.V1.power4.1}\\
146&\nolinkurl{amplified.det0.factor0.0}\\
147&\nolinkurl{amplified.det0.factor0.1}\\
148&\nolinkurl{amplified.det0.factor1.0}\\
149&\nolinkurl{amplified.det0.factor1.1}\\
150&\nolinkurl{amplified.det0.factor2.0}\\
151&\nolinkurl{amplified.det0.factor2.1}\\
152&\nolinkurl{amplified.det0.factor3.0}\\
153&\nolinkurl{amplified.det0.factor3.1}\\
154&\nolinkurl{amplified.det0.factor4.0}\\
155&\nolinkurl{amplified.det0.factor4.1}\\
156&\nolinkurl{amplified.det0.product1.0}\\
157&\nolinkurl{amplified.det0.product1.1}\\
158&\nolinkurl{amplified.det0.product2.0}\\
159&\nolinkurl{amplified.det0.product2.1}\\
160&\nolinkurl{amplified.det0.product3.0}\\
161&\nolinkurl{amplified.det0.product3.1}\\
162&\nolinkurl{amplified.det1.factor0.0}\\
163&\nolinkurl{amplified.det1.factor0.1}\\
164&\nolinkurl{amplified.det1.factor1.0}\\
165&\nolinkurl{amplified.det1.factor1.1}\\
166&\nolinkurl{amplified.det1.factor2.0}\\
167&\nolinkurl{amplified.det1.factor2.1}\\
168&\nolinkurl{amplified.det1.factor3.0}\\
169&\nolinkurl{amplified.det1.factor3.1}\\
170&\nolinkurl{amplified.det1.factor4.0}\\
171&\nolinkurl{amplified.det1.factor4.1}\\
172&\nolinkurl{amplified.det1.product1.0}\\
173&\nolinkurl{amplified.det1.product1.1}\\
174&\nolinkurl{amplified.det1.product2.0}\\
175&\nolinkurl{amplified.det1.product2.1}\\
176&\nolinkurl{amplified.det1.product3.0}\\
177&\nolinkurl{amplified.det1.product3.1}\\
178&\nolinkurl{rootProduct.1.A.0.rescaled}\\
179&\nolinkurl{rootProduct.1.A.1.rescaled}\\
180&\nolinkurl{rootProduct.1.B.0.rescaled}\\
181&\nolinkurl{rootProduct.1.B.1.rescaled}\\
182&\nolinkurl{rootProduct.1.C.0.rescaled}\\
183&\nolinkurl{rootProduct.1.C.1.rescaled}\\
184&\nolinkurl{rootProduct.1.E.0.rescaled}\\
185&\nolinkurl{rootProduct.1.E.1.rescaled}\\
186&\nolinkurl{rootProduct.1.i.0.rescaled}\\
187&\nolinkurl{rootProduct.1.i.1.rescaled}\\
188&\nolinkurl{rootProduct.1.k.0.rescaled}\\
189&\nolinkurl{rootProduct.1.k.1.rescaled}\\
190&\nolinkurl{rootProduct.1.u.0}\\
191&\nolinkurl{rootProduct.1.u.1}\\
192&\nolinkurl{rootProduct.2.A.0.rescaled}\\
193&\nolinkurl{rootProduct.2.A.1.rescaled}\\
194&\nolinkurl{rootProduct.2.B.0.rescaled}\\
195&\nolinkurl{rootProduct.2.B.1.rescaled}\\
196&\nolinkurl{rootProduct.2.C.0.rescaled}\\
197&\nolinkurl{rootProduct.2.C.1.rescaled}\\
198&\nolinkurl{rootProduct.2.E.0.rescaled}\\
199&\nolinkurl{rootProduct.2.E.1.rescaled}\\
200&\nolinkurl{rootProduct.2.i.0}\\
201&\nolinkurl{rootProduct.2.i.1}\\
202&\nolinkurl{rootProduct.2.k.0}\\
203&\nolinkurl{rootProduct.2.k.1}\\
204&\nolinkurl{rootProduct.2.u.0}\\
205&\nolinkurl{rootProduct.2.u.1}\\
206&\nolinkurl{rootProduct.3.A.0.rescaled}\\
207&\nolinkurl{rootProduct.3.A.1.rescaled}\\
208&\nolinkurl{rootProduct.3.B.0.rescaled}\\
209&\nolinkurl{rootProduct.3.B.1.rescaled}\\
210&\nolinkurl{rootProduct.3.C.0.rescaled}\\
211&\nolinkurl{rootProduct.3.C.1.rescaled}\\
212&\nolinkurl{rootProduct.3.E.0.rescaled}\\
213&\nolinkurl{rootProduct.3.E.1.rescaled}\\
214&\nolinkurl{rootProduct.3.i.0}\\
215&\nolinkurl{rootProduct.3.i.1}\\
216&\nolinkurl{rootProduct.3.k.0}\\
217&\nolinkurl{rootProduct.3.k.1}\\
218&\nolinkurl{rootProduct.3.u.0}\\
219&\nolinkurl{rootProduct.3.u.1}\\
220&\nolinkurl{rootProduct.4.A.0.rescaled}\\
221&\nolinkurl{rootProduct.4.A.1.rescaled}\\
222&\nolinkurl{rootProduct.4.B.0.rescaled}\\
223&\nolinkurl{rootProduct.4.B.1.rescaled}\\
224&\nolinkurl{rootProduct.4.C.0.rescaled}\\
225&\nolinkurl{rootProduct.4.C.1.rescaled}\\
226&\nolinkurl{rootProduct.4.E.0.rescaled}\\
227&\nolinkurl{rootProduct.4.E.1.rescaled}\\
228&\nolinkurl{rootProduct.4.i.0}\\
229&\nolinkurl{rootProduct.4.i.1}\\
230&\nolinkurl{rootProduct.4.k.0}\\
231&\nolinkurl{rootProduct.4.k.1}\\
232&\nolinkurl{rootProduct.4.u.0}\\
233&\nolinkurl{rootProduct.4.u.1}\\
\end{longtable}
\subsection{Completion of the 233-variable theorem}\label{v15x:completion}
The full arithmetic proof of the 248-variable predecessor is retained
in Section~\ref{v15:base248}. The scalar-S change and root, reciprocal,
and global-scaling changes prove the 238-variable predecessor as above.
Section~\ref{v15x:parity} proves both arithmetic directions for the new
237-variable scalar criterion, with one sufficiently large ineffective
threshold uniform in the input. Sections~\ref{v15x:235} and
\ref{v15x:233} prove two-sided maps to 235 and 233 variables on this
family. All exceptional denominators have been justified before their
use. The constant affine and row maps of Section~\ref{v15x:maps}
then preserve rational solubility with the input fixed.

The final residuals have integer coefficients and degree at most two.
Their sum of squares vanishes over $\Q$ exactly when each residual does.
Consequently they prove the defining equivalence for this 233-variable
predecessor. The additional reduction to Theorem~\ref{thm:v15-main} is
given in Sections~\ref{v222:construction}--\ref{v222:height}. The coefficient calculation below proves
quartic degree, content one, and the stated height. This finite
calculation does not supply a numerical value of the arithmetic
threshold $N_0$.

The following233-variable system is retained as a historical proof intermediate; the current family is in Section~\ref{v222:equations}.
\section{The explicit 233-variable system}\label{v15x:equations}
In this section only, $z_1,\ldots,z_{233}$ are the final sparse coordinates,
with distinguished input $z_1$. Put $A_N=149^{2N+1}$ and
$P_N=\sum_{r=1}^{207}Q_{r,N}^2$. We suppress the subscript $N$ on
individual residuals. Each $K_j$ in this section is the fixed positive
integer specified in the following table; these symbols are local to
this section and do not refer to the predecessor's table.
% BEGIN V15X FINAL 207 EQUATIONS -- IR SHA256 10fe99b5ee41b8ed7b602fd9e8bd759f8640cd18b37326ce12ca62832320d702

% V0.reference.square_eliminated.0
\begin{equation*}\begin{aligned}
Q_{1}={}&
-12042981952 +95879821296 z_{17} -191611825433 z_{17}^{2}\\
&+507888683232 z_{17} z_{18} +612903260592 z_{17} z_{29} -1654775158400 z_{17} z_{30}\\
&-674479990044 z_{17} z_{135} +1000502518848 z_{17} z_{136} -124053139840 z_{18}\\
&-383223650866 z_{18}^{2} -827387579200 z_{18} z_{29} +2451613042368 z_{18} z_{30}\\
&+1000502518848 z_{18} z_{135} -1348959980088 z_{18} z_{136} +191611825433 z_{29}\\
&+491340931648 z_{29} z_{82} -1344186590208 z_{29} z_{83} -507888683232 z_{30}\\
&-1344186590208 z_{30} z_{82} +3930727453184 z_{30} z_{83} +165304149024 z_{135}\\
&-654431752548 z_{135}^{2} +1787157171072 z_{135} z_{136} -249575558400 z_{136}\\
&-1308863505096 z_{136}^{2}
\end{aligned}\end{equation*}
% V0.reference.square_eliminated.1
\begin{equation*}\begin{aligned}
Q_{2}={}&
3769142784 -31013284960 z_{17} +63486085404 z_{17}^{2}\\
&-191611825433 z_{17} z_{18} -206846894800 z_{17} z_{29} +612903260592 z_{17} z_{30}\\
&+250125629712 z_{17} z_{135} -337239995022 z_{17} z_{136} +47939910648 z_{18}\\
&+126972170808 z_{18}^{2} +306451630296 z_{18} z_{29} -827387579200 z_{18} z_{30}\\
&-337239995022 z_{18} z_{135} +500251259424 z_{18} z_{136} -63486085404 z_{29}\\
&-168023323776 z_{29} z_{82} +491340931648 z_{29} z_{83} +191611825433 z_{30}\\
&+491340931648 z_{30} z_{82} -1344186590208 z_{30} z_{83} -62393889600 z_{135}\\
&+223394646384 z_{135}^{2} -654431752548 z_{135} z_{136} +82652074512 z_{136}\\
&+446789292768 z_{136}^{2}
\end{aligned}\end{equation*}
% H.reference.linear.0
\begin{equation*}\begin{aligned}
Q_{3}={}&
-95304 +590783 z_{17} +151368 z_{17} z_{19}\\
&-529788 z_{17} z_{20} +1665048 z_{17} z_{135} -1816416 z_{17} z_{136}\\
&-1201800 z_{18} -529788 z_{18} z_{19} +302736 z_{18} z_{20}\\
&-1816416 z_{18} z_{135} +3330096 z_{18} z_{136} +1078497 z_{135}\\
&-1059576 z_{136}
\end{aligned}\end{equation*}
% H.reference.linear.1
\begin{equation*}\begin{aligned}
Q_{4}={}&
46816 -600900 z_{17} -264894 z_{17} z_{19}\\
&+151368 z_{17} z_{20} -908208 z_{17} z_{135} +1665048 z_{17} z_{136}\\
&+590783 z_{18} +151368 z_{18} z_{19} -529788 z_{18} z_{20}\\
&+1665048 z_{18} z_{135} -1816416 z_{18} z_{136} -529788 z_{135}\\
&+1078497 z_{136}
\end{aligned}\end{equation*}
% H.reference.nondegenerate.0
\begin{equation*}\begin{aligned}
Q_{5}={}&
216576 -1078497 z_{17} +1059576 z_{18}\\
&-13376 z_{19} +151368 z_{19} z_{135} -529788 z_{19} z_{136}\\
&+46816 z_{20} -529788 z_{20} z_{135} +302736 z_{20} z_{136}\\
&-737919 z_{135} +1362312 z_{136}
\end{aligned}\end{equation*}
% H.reference.nondegenerate.1
\begin{equation*}\begin{aligned}
Q_{6}={}&
-60192 +529788 z_{17} -1078497 z_{18}\\
&+23408 z_{19} -264894 z_{19} z_{135} +151368 z_{19} z_{136}\\
&-13376 z_{20} +151368 z_{20} z_{135} -529788 z_{20} z_{136}\\
&+681156 z_{135} -737919 z_{136}
\end{aligned}\end{equation*}
% logV0.wPower2.0.definition
\begin{equation*}\begin{aligned}
Q_{7}={}&
2401 z_{29} z_{30} -6369 z_{31}
\end{aligned}\end{equation*}
% logV0.wPower2.1.definition
\begin{equation*}\begin{aligned}
Q_{8}={}&
33 z_{29}^{2} -193 z_{29} z_{30} +264 z_{30}^{2}\\
&-33 z_{32}
\end{aligned}\end{equation*}
% plain.alternateRatio1.cleared.0
\begin{equation*}\begin{aligned}
Q_{9}={}&
775761 z_{19} z_{21} +504136 z_{19} z_{49} +373756 z_{19} z_{50}\\
&+1551522 z_{20} z_{22} +373756 z_{20} z_{49} +1008272 z_{20} z_{50}\\
&-1542912 z_{21} +3027360 z_{21} z_{135} -8022504 z_{21} z_{136}\\
&-8022504 z_{22} z_{135} +6054720 z_{22} z_{136} +775761 z_{49}
\end{aligned}\end{equation*}
% plain.alternateRatio1.cleared.1
\begin{equation*}\begin{aligned}
Q_{10}={}&
775761 z_{19} z_{22} +186878 z_{19} z_{49} +504136 z_{19} z_{50}\\
&+775761 z_{20} z_{21} +504136 z_{20} z_{49} +373756 z_{20} z_{50}\\
&-4011252 z_{21} z_{135} +3027360 z_{21} z_{136} -1542912 z_{22}\\
&+3027360 z_{22} z_{135} -8022504 z_{22} z_{136} +775761 z_{50}
\end{aligned}\end{equation*}
% logV2.wPower2.0.definition
\begin{equation*}\begin{aligned}
Q_{11}={}&
2401 z_{35} z_{36} -6369 z_{37}
\end{aligned}\end{equation*}
% logV2.wPower2.1.definition
\begin{equation*}\begin{aligned}
Q_{12}={}&
33 z_{35}^{2} -193 z_{35} z_{36} +264 z_{36}^{2}\\
&-33 z_{38}
\end{aligned}\end{equation*}
% plain.alternateRatio2.cleared.0
\begin{equation*}\begin{aligned}
Q_{13}={}&
775761 z_{19} z_{23} +504136 z_{19} z_{53} +373756 z_{19} z_{54}\\
&+1551522 z_{20} z_{24} +373756 z_{20} z_{53} +1008272 z_{20} z_{54}\\
&-1542912 z_{23} +3027360 z_{23} z_{135} -8022504 z_{23} z_{136}\\
&-8022504 z_{24} z_{135} +6054720 z_{24} z_{136} +775761 z_{53}
\end{aligned}\end{equation*}
% plain.alternateRatio2.cleared.1
\begin{equation*}\begin{aligned}
Q_{14}={}&
775761 z_{19} z_{24} +186878 z_{19} z_{53} +504136 z_{19} z_{54}\\
&+775761 z_{20} z_{23} +504136 z_{20} z_{53} +373756 z_{20} z_{54}\\
&-4011252 z_{23} z_{135} +3027360 z_{23} z_{136} -1542912 z_{24}\\
&+3027360 z_{24} z_{135} -8022504 z_{24} z_{136} +775761 z_{54}
\end{aligned}\end{equation*}
% logV3.wPower2.0.definition
\begin{equation*}\begin{aligned}
Q_{15}={}&
2401 z_{39} z_{40} -6369 z_{41}
\end{aligned}\end{equation*}
% logV3.wPower2.1.definition
\begin{equation*}\begin{aligned}
Q_{16}={}&
33 z_{39}^{2} -193 z_{39} z_{40} +264 z_{40}^{2}\\
&-33 z_{42}
\end{aligned}\end{equation*}
% plain.alternateRatio3.cleared.0
\begin{equation*}\begin{aligned}
Q_{17}={}&
775761 z_{19} z_{25} +504136 z_{19} z_{57} +373756 z_{19} z_{58}\\
&+1551522 z_{20} z_{26} +373756 z_{20} z_{57} +1008272 z_{20} z_{58}\\
&-1542912 z_{25} +3027360 z_{25} z_{135} -8022504 z_{25} z_{136}\\
&-8022504 z_{26} z_{135} +6054720 z_{26} z_{136} +775761 z_{57}
\end{aligned}\end{equation*}
% plain.alternateRatio3.cleared.1
\begin{equation*}\begin{aligned}
Q_{18}={}&
775761 z_{19} z_{26} +186878 z_{19} z_{57} +504136 z_{19} z_{58}\\
&+775761 z_{20} z_{25} +504136 z_{20} z_{57} +373756 z_{20} z_{58}\\
&-4011252 z_{25} z_{135} +3027360 z_{25} z_{136} -1542912 z_{26}\\
&+3027360 z_{26} z_{135} -8022504 z_{26} z_{136} +775761 z_{58}
\end{aligned}\end{equation*}
% logV4.wPower2.0.definition
\begin{equation*}\begin{aligned}
Q_{19}={}&
2401 z_{43} z_{44} -6369 z_{45}
\end{aligned}\end{equation*}
% logV4.wPower2.1.definition
\begin{equation*}\begin{aligned}
Q_{20}={}&
33 z_{43}^{2} -193 z_{43} z_{44} +264 z_{44}^{2}\\
&-33 z_{46}
\end{aligned}\end{equation*}
% plain.alternateRatio4.cleared.0
\begin{equation*}\begin{aligned}
Q_{21}={}&
-775761 z_{19} z_{28} -186878 z_{19} z_{61} -504136 z_{19} z_{62}\\
&-775761 z_{20} z_{27} -504136 z_{20} z_{61} -373756 z_{20} z_{62}\\
&+4011252 z_{27} z_{135} -3027360 z_{27} z_{136} +1542912 z_{28}\\
&-3027360 z_{28} z_{135} +8022504 z_{28} z_{136} -775761 z_{62}
\end{aligned}\end{equation*}
% plain.alternateRatio4.cleared.1
\begin{equation*}\begin{aligned}
Q_{22}={}&
775761 z_{19} z_{27} +504136 z_{19} z_{61} +373756 z_{19} z_{62}\\
&+1551522 z_{20} z_{28} +373756 z_{20} z_{61} +1008272 z_{20} z_{62}\\
&-1542912 z_{27} +3027360 z_{27} z_{135} -8022504 z_{27} z_{136}\\
&-8022504 z_{28} z_{135} +6054720 z_{28} z_{136} +775761 z_{61}
\end{aligned}\end{equation*}
% inputSquare.definition
\begin{equation*}\begin{aligned}
Q_{23}={}&
2501 +51 z_{1}^{2} -51 z_{63}
\end{aligned}\end{equation*}
% Lucas.C.0.definition
\begin{equation*}\begin{aligned}
Q_{24}={}&
z_{51} z_{59} +2 z_{52} z_{60} -2 z_{59} z_{137}\\
&-z_{64} z_{137}
\end{aligned}\end{equation*}
% Lucas.C.1.definition
\begin{equation*}\begin{aligned}
Q_{25}={}&
z_{51} z_{60} +z_{52} z_{59} -2 z_{60} z_{137}\\
&-z_{65} z_{137}
\end{aligned}\end{equation*}
% Lucas.AT.0.definition
\begin{equation*}\begin{aligned}
Q_{26}={}&
z_{51} z_{55} +2 z_{52} z_{56} -z_{64} z_{137}\\
&-z_{66} z_{137}
\end{aligned}\end{equation*}
% Lucas.AT.1.definition
\begin{equation*}\begin{aligned}
Q_{27}={}&
z_{51} z_{56} +z_{52} z_{55} -z_{65} z_{137}\\
&-z_{67} z_{137}
\end{aligned}\end{equation*}
% Lucas.F0.definition
\begin{equation*}\begin{aligned}
Q_{28}={}&
z_{55}^{2} +2 z_{56}^{2} -z_{64} z_{66}\\
&-2 z_{65} z_{67} -z_{68} z_{134} -z_{134}
\end{aligned}\end{equation*}
% Lucas.F1.definition
\begin{equation*}\begin{aligned}
Q_{29}={}&
2 z_{55} z_{56} -z_{64} z_{67} -z_{65} z_{66}\\
&-2 z_{69} z_{134}
\end{aligned}\end{equation*}
% testSquare.12.0.0.definition
\begin{equation*}\begin{aligned}
Q_{30}={}&
z_{47}^{2} -2 z_{47} z_{49} +2 z_{48}^{2}\\
&+4 z_{48} z_{50} +z_{49}^{2} +2 z_{50}^{2}\\
&-z_{78}
\end{aligned}\end{equation*}
% testSquare.12.0.1.definition
\begin{equation*}\begin{aligned}
Q_{31}={}&
z_{47} z_{48} +z_{47} z_{50} -z_{48} z_{49}\\
&-z_{49} z_{50} -z_{79}
\end{aligned}\end{equation*}
% testSquare.12.1.0.definition
\begin{equation*}\begin{aligned}
Q_{32}={}&
9 +z_{47}^{2} -102 z_{47} z_{63}\\
&+2 z_{48}^{2} +2601 z_{63}^{2} -z_{80}
\end{aligned}\end{equation*}
% testSquare.12.1.1.definition
\begin{equation*}\begin{aligned}
Q_{33}={}&
z_{47} z_{48} -51 z_{48} z_{63} -z_{81}
\end{aligned}\end{equation*}
% testSquare.40.0.0.definition
\begin{equation*}\begin{aligned}
Q_{34}={}&
z_{17}^{2} +2 z_{18}^{2} -z_{82}
\end{aligned}\end{equation*}
% testSquare.40.0.1.definition
\begin{equation*}\begin{aligned}
Q_{35}={}&
z_{17} z_{18} -z_{83}
\end{aligned}\end{equation*}
% testSquare.40.1.0.definition
\begin{equation*}\begin{aligned}
Q_{36}={}&
z_{51}^{2} -2 z_{51} z_{53} +2 z_{52}^{2}\\
&+4 z_{52} z_{54} +z_{53}^{2} +2 z_{54}^{2}\\
&-z_{84} z_{134}
\end{aligned}\end{equation*}
% testSquare.40.1.1.definition
\begin{equation*}\begin{aligned}
Q_{37}={}&
z_{51} z_{52} +z_{51} z_{54} -z_{52} z_{53}\\
&-z_{53} z_{54} -z_{85} z_{134}
\end{aligned}\end{equation*}
% testSquare.40.2.0.definition
\begin{equation*}\begin{aligned}
Q_{38}={}&
z_{55}^{2} -2 z_{55} z_{57} +2 z_{56}^{2}\\
&+4 z_{56} z_{58} +z_{57}^{2} +2 z_{58}^{2}\\
&-z_{86} z_{134}
\end{aligned}\end{equation*}
% testSquare.40.2.1.definition
\begin{equation*}\begin{aligned}
Q_{39}={}&
z_{55} z_{56} +z_{55} z_{58} -z_{56} z_{57}\\
&-z_{57} z_{58} -z_{87} z_{134}
\end{aligned}\end{equation*}
% testSquare.40.3.0.definition
\begin{equation*}\begin{aligned}
Q_{40}={}&
z_{59}^{2} -2 z_{59} z_{61} +2 z_{60}^{2}\\
&+4 z_{60} z_{62} +z_{61}^{2} +2 z_{62}^{2}\\
&-z_{88} z_{134}
\end{aligned}\end{equation*}
% testSquare.40.3.1.definition
\begin{equation*}\begin{aligned}
Q_{41}={}&
z_{59} z_{60} +z_{59} z_{62} -z_{60} z_{61}\\
&-z_{61} z_{62} -z_{89} z_{134}
\end{aligned}\end{equation*}
% testSquare.40.4.0.definition
\begin{equation*}\begin{aligned}
Q_{42}={}&
z_{68}^{2} +8 z_{69}^{2} -z_{90}
\end{aligned}\end{equation*}
% testSquare.40.4.1.definition
\begin{equation*}\begin{aligned}
Q_{43}={}&
z_{68} z_{69} -z_{91}
\end{aligned}\end{equation*}
% reciprocal.symmetric.detector0.deltaSquared.definition
\begin{equation*}\begin{aligned}
Q_{44}={}&
z_{4}^{2} -z_{92}
\end{aligned}\end{equation*}
% reciprocal.symmetric.detector0.deltaFourthPower.definition
\begin{equation*}\begin{aligned}
Q_{45}={}&
z_{92}^{2} -z_{93}
\end{aligned}\end{equation*}
% reciprocal.reflected.detector0.contactPolynomial.definition
\begin{equation*}\begin{aligned}
Q_{46}={}&
7626746323592 -76267463235920 z_{4} -628059618995128 z_{4} z_{92}\\
&+732473028362323 z_{4} z_{138} +307309842082151 z_{92} -583791423311868 z_{92} z_{138}\\
&+718967979842546 z_{93} +162308064100032 z_{93}^{2} -38129835686440 z_{93} z_{138}\\
&-7626746323592 z_{94} -3813373161796 z_{94} z_{95} -646146742577308 z_{138}\\
&+3813373161796 z_{138}^{2}
\end{aligned}\end{equation*}
% reciprocal.reflected.detector0.weight.definition
\begin{equation*}\begin{aligned}
Q_{47}={}&
1 -10 z_{4} -120 z_{4} z_{92}\\
&+210 z_{4} z_{138} +45 z_{92} -120 z_{92} z_{138}\\
&+210 z_{93} +45 z_{93}^{2} -10 z_{93} z_{138}\\
&-z_{94} -z_{94} z_{95} -252 z_{138}\\
&+z_{138}^{2}
\end{aligned}\end{equation*}
% detector0.E1.power2.definition
\begin{equation*}\begin{aligned}
Q_{48}={}&
z_{95}^{2} -z_{97}
\end{aligned}\end{equation*}
% detector0.E1.power4.definition
\begin{equation*}\begin{aligned}
Q_{49}={}&
z_{97}^{2} -z_{98}
\end{aligned}\end{equation*}
% detector0.normOver2w0.definition
\begin{equation*}\begin{aligned}
Q_{50}={}&
z_{99}^{2} -2 z_{99} z_{101} -2 z_{100}^{2}
\end{aligned}\end{equation*}
% balanced.detector0.poleInput.definition
\begin{equation*}\begin{aligned}
Q_{51}={}&
167281 z_{95} z_{102} -166464 z_{101}^{2} +167281 z_{102}
\end{aligned}\end{equation*}
% detector0.p2.definition
\begin{equation*}\begin{aligned}
Q_{52}={}&
z_{7}^{2} -z_{103}
\end{aligned}\end{equation*}
% detector0.firstConic
\begin{equation*}\begin{aligned}
Q_{53}={}&
289 +34 z_{5} -3 z_{5}^{2}\\
&-17 z_{102} z_{103}
\end{aligned}\end{equation*}
% detector0.Bh.definition
\begin{equation*}\begin{aligned}
Q_{54}={}&
17 z_{5} -z_{5} z_{102} -17 z_{6} z_{104}
\end{aligned}\end{equation*}
% detector0.q2.definition
\begin{equation*}\begin{aligned}
Q_{55}={}&
z_{8}^{2} -z_{105}
\end{aligned}\end{equation*}
% detector0.secondConic
\begin{equation*}\begin{aligned}
Q_{56}={}&
z_{6}^{2} +34 z_{6} z_{104} -17 z_{102} z_{105}\\
&-867 z_{104}^{2}
\end{aligned}\end{equation*}
% symmetric.detector1.deltaSquared.definition
\begin{equation*}\begin{aligned}
Q_{57}={}&
z_{106}^{2} -2 z_{106} z_{134} -z_{107} z_{134}\\
&+z_{134}^{2}
\end{aligned}\end{equation*}
% symmetric.detector1.deltaFourthPower.definition
\begin{equation*}\begin{aligned}
Q_{58}={}&
z_{107}^{2} -z_{108} z_{134}
\end{aligned}\end{equation*}
% reflected.detector1.contactPolynomial.definition
\begin{equation*}\begin{aligned}
Q_{59}={}&
-141822723674536 z_{107} z_{134} +35893742479489 z_{107} z_{139} +192222345736537 z_{108} z_{134}\\
&-3813373161796 z_{109} z_{110} +39896049934224 z_{134}^{2} -122376041313918 z_{134} z_{139}
\end{aligned}\end{equation*}
% reflected.detector1.weight.definition
\begin{equation*}\begin{aligned}
Q_{60}={}&
z_{108} z_{139} -z_{109} -z_{109} z_{110}
\end{aligned}\end{equation*}
% detector1.E1.power2.definition
\begin{equation*}\begin{aligned}
Q_{61}={}&
z_{110}^{2} -z_{112}
\end{aligned}\end{equation*}
% detector1.E1.power4.definition
\begin{equation*}\begin{aligned}
Q_{62}={}&
z_{112}^{2} -z_{113}
\end{aligned}\end{equation*}
% detector1.normOver2w0.definition
\begin{equation*}\begin{aligned}
Q_{63}={}&
z_{114}^{2} -2 z_{114} z_{116} -2 z_{115}^{2}
\end{aligned}\end{equation*}
% balanced.detector1.poleInput.definition
\begin{equation*}\begin{aligned}
Q_{64}={}&
167281 z_{110} z_{117} -166464 z_{116}^{2} +167281 z_{117}
\end{aligned}\end{equation*}
% detector1.p2.definition
\begin{equation*}\begin{aligned}
Q_{65}={}&
z_{11}^{2} -z_{118}
\end{aligned}\end{equation*}
% detector1.firstConic
\begin{equation*}\begin{aligned}
Q_{66}={}&
289 +34 z_{9} -3 z_{9}^{2}\\
&-17 z_{117} z_{118}
\end{aligned}\end{equation*}
% detector1.Bh.definition
\begin{equation*}\begin{aligned}
Q_{67}={}&
17 z_{9} -z_{9} z_{117} -17 z_{10} z_{119}
\end{aligned}\end{equation*}
% detector1.q2.definition
\begin{equation*}\begin{aligned}
Q_{68}={}&
z_{12}^{2} -z_{120}
\end{aligned}\end{equation*}
% detector1.secondConic
\begin{equation*}\begin{aligned}
Q_{69}={}&
z_{10}^{2} +34 z_{10} z_{119} -17 z_{117} z_{120}\\
&-867 z_{119}^{2}
\end{aligned}\end{equation*}
% scalar.rPair1.definition
\begin{equation*}\begin{aligned}
Q_{70}={}&
2601 z_{63}^{2} -4 z_{121}
\end{aligned}\end{equation*}
% scalar.rFour.definition
\begin{equation*}\begin{aligned}
Q_{71}={}&
3882987 z_{63} z_{121} -128 z_{121}^{2} -381485 z_{122}
\end{aligned}\end{equation*}
% scalar.contact_product.definition
\begin{equation*}\begin{aligned}
Q_{72}={}&
31581660 +2103220824 z_{63} +29762681184 z_{63}^{2}\\
&+585078120 z_{63} z_{121} +155645880 z_{63} z_{122} +28680300 z_{121}\\
&+7629700 z_{122} +z_{123}
\end{aligned}\end{equation*}
% scalar.W2.definition
\begin{equation*}\begin{aligned}
Q_{73}={}&
K_{1} z_{126} z_{134} -z_{133}^{2}
\end{aligned}\end{equation*}
% scalar.b.definition
\begin{equation*}\begin{aligned}
Q_{74}={}&
9792 z_{126} +z_{126} z_{137} +82479248 z_{127} z_{137}
\end{aligned}\end{equation*}
% scalar.p2.definition
\begin{equation*}\begin{aligned}
Q_{75}={}&
z_{15}^{2} -z_{128}
\end{aligned}\end{equation*}
% scalar.firstConic
\begin{equation*}\begin{aligned}
Q_{76}={}&
K_{7} -K_{4} z_{13} -z_{13}^{2}\\
&+K_{3} z_{127} z_{128}
\end{aligned}\end{equation*}
% scalar.Bh.definition
\begin{equation*}\begin{aligned}
Q_{77}={}&
17 z_{13} +K_{2} z_{13} z_{127} -17 z_{14} z_{129}
\end{aligned}\end{equation*}
% scalar.q2.definition
\begin{equation*}\begin{aligned}
Q_{78}={}&
z_{16}^{2} -z_{130}
\end{aligned}\end{equation*}
% scalar.secondConic
\begin{equation*}\begin{aligned}
Q_{79}={}&
z_{14}^{2} -K_{6} z_{14} z_{129} +K_{5} z_{127} z_{130}\\
&-K_{8} z_{129}^{2}
\end{aligned}\end{equation*}
% parity_scalar.reduced_scalar.root.power4
\begin{equation*}\begin{aligned}
Q_{80}={}&
z_{82}^{2} -8 z_{83}^{2} +z_{131} z_{137}\\
&+z_{132}^{2} -z_{133} z_{134}
\end{aligned}\end{equation*}
% reflected.detector0.weightPower7.definition
\begin{equation*}\begin{aligned}
Q_{81}={}&
z_{95} z_{96} -z_{98}^{2}
\end{aligned}\end{equation*}
% reflected.detector1.weightPower7.definition
\begin{equation*}\begin{aligned}
Q_{82}={}&
z_{110} z_{111} -z_{113}^{2}
\end{aligned}\end{equation*}
% scalar.r_power_2.definition
\begin{equation*}\begin{aligned}
Q_{83}={}&
z_{134} -z_{137}^{2}
\end{aligned}\end{equation*}
% scalar.ratio1_unscaled_square.definition
\begin{equation*}\begin{aligned}
Q_{84}={}&
z_{125}^{2} -z_{131}
\end{aligned}\end{equation*}
% parity_scalar.scalar.ratio_node.consolidated_denominator.power4
\begin{equation*}\begin{aligned}
Q_{85}={}&
z_{124}^{2} -z_{132} z_{137}
\end{aligned}\end{equation*}
% scalar.forced_prime.definition
\begin{equation*}\begin{aligned}
Q_{86}={}&
22201 z_{123} z_{137} +z_{137} -A_N z_{123}
\end{aligned}\end{equation*}
% reciprocal.reflected.detector0.deltaSixthPower.definition
\begin{equation*}\begin{aligned}
Q_{87}={}&
z_{4} z_{93} -z_{138}
\end{aligned}\end{equation*}
% reflected.detector1.deltaSixthPower.definition
\begin{equation*}\begin{aligned}
Q_{88}={}&
z_{107} z_{108} -z_{134} z_{139}
\end{aligned}\end{equation*}
% amplified.C1.power4.0.definition
\begin{equation*}\begin{aligned}
Q_{89}={}&
z_{78}^{2} +8 z_{79}^{2} -z_{140}
\end{aligned}\end{equation*}
% amplified.C1.power4.1.definition
\begin{equation*}\begin{aligned}
Q_{90}={}&
4 z_{78} z_{79} -z_{141}
\end{aligned}\end{equation*}
% residue.I.power4.0.definition
\begin{equation*}\begin{aligned}
Q_{91}={}&
-324 -12 z_{47} z_{80} -48 z_{48} z_{81}\\
&+612 z_{63} z_{80} +36 z_{80} +z_{80}^{2}\\
&+8 z_{81}^{2} -z_{142}
\end{aligned}\end{equation*}
% residue.I.power4.1.definition
\begin{equation*}\begin{aligned}
Q_{92}={}&
-24 z_{47} z_{81} -12 z_{48} z_{80} +1224 z_{63} z_{81}\\
&+4 z_{80} z_{81} +72 z_{81} -z_{143}
\end{aligned}\end{equation*}
% residue.V1.power4.0.definition
\begin{equation*}\begin{aligned}
Q_{93}={}&
z_{33}^{2} +8 z_{34}^{2} -z_{144}
\end{aligned}\end{equation*}
% residue.V1.power4.1.definition
\begin{equation*}\begin{aligned}
Q_{94}={}&
4 z_{33} z_{34} +z_{145}
\end{aligned}\end{equation*}
% residue.det0.factor0.0.definition
\begin{equation*}\begin{aligned}
Q_{95}={}&
1 +z_{82} z_{97} +z_{84} z_{96}\\
&+z_{98} z_{144} -z_{146}
\end{aligned}\end{equation*}
% residue.det0.factor0.1.definition
\begin{equation*}\begin{aligned}
Q_{96}={}&
2 z_{83} z_{97} +2 z_{85} z_{96} +z_{98} z_{145}\\
&-z_{147}
\end{aligned}\end{equation*}
% residue.det0.factor1.0.definition
\begin{equation*}\begin{aligned}
Q_{97}={}&
2401 +37249 z_{31} z_{98} +2401 z_{32} z_{98}\\
&+2401 z_{74} z_{97} +2401 z_{90} z_{96} -2401 z_{148}
\end{aligned}\end{equation*}
% residue.det0.factor1.1.definition
\begin{equation*}\begin{aligned}
Q_{98}={}&
25476 z_{31} z_{98} -4802 z_{75} z_{97} -9604 z_{91} z_{96}\\
&+2401 z_{149}
\end{aligned}\end{equation*}
% residue.det0.factor2.0.definition
\begin{equation*}\begin{aligned}
Q_{99}={}&
2401 +37249 z_{41} z_{98} +2401 z_{42} z_{98}\\
&+2401 z_{72} z_{95} +2401 z_{76} z_{97} +2401 z_{96} z_{142}\\
&-2401 z_{150}
\end{aligned}\end{equation*}
% residue.det0.factor2.1.definition
\begin{equation*}\begin{aligned}
Q_{100}={}&
25476 z_{41} z_{98} -4802 z_{73} z_{95} -4802 z_{77} z_{97}\\
&-2401 z_{96} z_{143} +2401 z_{151}
\end{aligned}\end{equation*}
% residue.det0.factor3.0.definition
\begin{equation*}\begin{aligned}
Q_{101}={}&
2401 +37249 z_{45} z_{98} +2401 z_{46} z_{98}\\
&+2401 z_{86} z_{96} +2401 z_{97} z_{140} -2401 z_{152}
\end{aligned}\end{equation*}
% residue.det0.factor3.1.definition
\begin{equation*}\begin{aligned}
Q_{102}={}&
25476 z_{45} z_{98} -4802 z_{87} z_{96} -2401 z_{97} z_{141}\\
&+2401 z_{153}
\end{aligned}\end{equation*}
% residue.det0.factor4.0.definition
\begin{equation*}\begin{aligned}
Q_{103}={}&
2401 +37249 z_{37} z_{98} +2401 z_{38} z_{98}\\
&+2401 z_{70} z_{97} +2401 z_{88} z_{96} -2401 z_{154}
\end{aligned}\end{equation*}
% residue.det0.factor4.1.definition
\begin{equation*}\begin{aligned}
Q_{104}={}&
25476 z_{37} z_{98} -4802 z_{71} z_{97} -4802 z_{89} z_{96}\\
&+2401 z_{155}
\end{aligned}\end{equation*}
% amplified.det0.product1.0.definition
\begin{equation*}\begin{aligned}
Q_{105}={}&
z_{146} z_{148} +2 z_{147} z_{149} -z_{156}
\end{aligned}\end{equation*}
% amplified.det0.product1.1.definition
\begin{equation*}\begin{aligned}
Q_{106}={}&
z_{146} z_{149} +z_{147} z_{148} -z_{157}
\end{aligned}\end{equation*}
% amplified.det0.product2.0.definition
\begin{equation*}\begin{aligned}
Q_{107}={}&
z_{150} z_{156} +2 z_{151} z_{157} -z_{158}
\end{aligned}\end{equation*}
% amplified.det0.product2.1.definition
\begin{equation*}\begin{aligned}
Q_{108}={}&
z_{150} z_{157} +z_{151} z_{156} -z_{159}
\end{aligned}\end{equation*}
% amplified.det0.product3.0.definition
\begin{equation*}\begin{aligned}
Q_{109}={}&
z_{152} z_{158} +2 z_{153} z_{159} -z_{160}
\end{aligned}\end{equation*}
% amplified.det0.product3.1.definition
\begin{equation*}\begin{aligned}
Q_{110}={}&
z_{152} z_{159} +z_{153} z_{158} -z_{161}
\end{aligned}\end{equation*}
% amplified.det0.product4.0.definition
\begin{equation*}\begin{aligned}
Q_{111}={}&
z_{99} -z_{154} z_{160} -2 z_{155} z_{161}
\end{aligned}\end{equation*}
% amplified.det0.product4.1.definition
\begin{equation*}\begin{aligned}
Q_{112}={}&
z_{100} -z_{154} z_{161} -z_{155} z_{160}
\end{aligned}\end{equation*}
% residue.det1.factor0.0.definition
\begin{equation*}\begin{aligned}
Q_{113}={}&
1 +z_{82} z_{112} +z_{84} z_{111}\\
&+z_{113} z_{144} -z_{162}
\end{aligned}\end{equation*}
% residue.det1.factor0.1.definition
\begin{equation*}\begin{aligned}
Q_{114}={}&
2 z_{83} z_{112} +2 z_{85} z_{111} +z_{113} z_{145}\\
&-z_{163}
\end{aligned}\end{equation*}
% residue.det1.factor1.0.definition
\begin{equation*}\begin{aligned}
Q_{115}={}&
2401 +37249 z_{31} z_{113} +2401 z_{32} z_{113}\\
&+2401 z_{74} z_{112} +2401 z_{90} z_{111} -2401 z_{164}
\end{aligned}\end{equation*}
% residue.det1.factor1.1.definition
\begin{equation*}\begin{aligned}
Q_{116}={}&
25476 z_{31} z_{113} -4802 z_{75} z_{112} -9604 z_{91} z_{111}\\
&+2401 z_{165}
\end{aligned}\end{equation*}
% residue.det1.factor2.0.definition
\begin{equation*}\begin{aligned}
Q_{117}={}&
2401 +37249 z_{41} z_{113} +2401 z_{42} z_{113}\\
&+2401 z_{72} z_{110} +2401 z_{76} z_{112} +2401 z_{111} z_{142}\\
&-2401 z_{166}
\end{aligned}\end{equation*}
% residue.det1.factor2.1.definition
\begin{equation*}\begin{aligned}
Q_{118}={}&
25476 z_{41} z_{113} -4802 z_{73} z_{110} -4802 z_{77} z_{112}\\
&-2401 z_{111} z_{143} +2401 z_{167}
\end{aligned}\end{equation*}
% residue.det1.factor3.0.definition
\begin{equation*}\begin{aligned}
Q_{119}={}&
2401 +37249 z_{45} z_{113} +2401 z_{46} z_{113}\\
&+2401 z_{86} z_{111} +2401 z_{112} z_{140} -2401 z_{168}
\end{aligned}\end{equation*}
% residue.det1.factor3.1.definition
\begin{equation*}\begin{aligned}
Q_{120}={}&
25476 z_{45} z_{113} -4802 z_{87} z_{111} -2401 z_{112} z_{141}\\
&+2401 z_{169}
\end{aligned}\end{equation*}
% residue.det1.factor4.0.definition
\begin{equation*}\begin{aligned}
Q_{121}={}&
2401 +37249 z_{37} z_{113} +2401 z_{38} z_{113}\\
&+2401 z_{70} z_{112} +2401 z_{88} z_{111} -2401 z_{170}
\end{aligned}\end{equation*}
% residue.det1.factor4.1.definition
\begin{equation*}\begin{aligned}
Q_{122}={}&
25476 z_{37} z_{113} -4802 z_{71} z_{112} -4802 z_{89} z_{111}\\
&+2401 z_{171}
\end{aligned}\end{equation*}
% amplified.det1.product1.0.definition
\begin{equation*}\begin{aligned}
Q_{123}={}&
z_{162} z_{164} +2 z_{163} z_{165} -z_{172}
\end{aligned}\end{equation*}
% amplified.det1.product1.1.definition
\begin{equation*}\begin{aligned}
Q_{124}={}&
z_{162} z_{165} +z_{163} z_{164} -z_{173}
\end{aligned}\end{equation*}
% amplified.det1.product2.0.definition
\begin{equation*}\begin{aligned}
Q_{125}={}&
z_{166} z_{172} +2 z_{167} z_{173} -z_{174}
\end{aligned}\end{equation*}
% amplified.det1.product2.1.definition
\begin{equation*}\begin{aligned}
Q_{126}={}&
z_{166} z_{173} +z_{167} z_{172} -z_{175}
\end{aligned}\end{equation*}
% amplified.det1.product3.0.definition
\begin{equation*}\begin{aligned}
Q_{127}={}&
z_{168} z_{174} +2 z_{169} z_{175} -z_{176}
\end{aligned}\end{equation*}
% amplified.det1.product3.1.definition
\begin{equation*}\begin{aligned}
Q_{128}={}&
z_{168} z_{175} +z_{169} z_{174} -z_{177}
\end{aligned}\end{equation*}
% amplified.det1.product4.0.definition
\begin{equation*}\begin{aligned}
Q_{129}={}&
z_{114} -z_{170} z_{176} -2 z_{171} z_{177}
\end{aligned}\end{equation*}
% amplified.det1.product4.1.definition
\begin{equation*}\begin{aligned}
Q_{130}={}&
z_{115} -z_{170} z_{177} -z_{171} z_{176}
\end{aligned}\end{equation*}
% rootProduct.1.equation1.0
\begin{equation*}\begin{aligned}
Q_{131}={}&
z_{178}^{2} +2 z_{179}^{2} +z_{180}^{2}\\
&+2 z_{181}^{2} +z_{182}^{2} +2 z_{183}^{2}\\
&-z_{184}^{2} -2 z_{185}^{2}
\end{aligned}\end{equation*}
% rootProduct.1.equation1.1
\begin{equation*}\begin{aligned}
Q_{132}={}&
z_{178} z_{179} +z_{180} z_{181} +z_{182} z_{183}\\
&-z_{184} z_{185}
\end{aligned}\end{equation*}
% rootProduct.1.equation2.0
\begin{equation*}\begin{aligned}
Q_{133}={}&
-12 z_{137} z_{186} +8 z_{137} z_{187} +z_{178} z_{182}\\
&+2 z_{179} z_{183} +z_{180} z_{184} +2 z_{181} z_{185}
\end{aligned}\end{equation*}
% rootProduct.1.equation2.1
\begin{equation*}\begin{aligned}
Q_{134}={}&
4 z_{137} z_{186} -12 z_{137} z_{187} +z_{178} z_{183}\\
&+z_{179} z_{182} +z_{180} z_{185} +z_{181} z_{184}
\end{aligned}\end{equation*}
% rootProduct.1.equation3.0
\begin{equation*}\begin{aligned}
Q_{135}={}&
12 z_{137} z_{188} -8 z_{137} z_{189} +z_{178} z_{184}\\
&+2 z_{179} z_{185} +z_{180} z_{182} +2 z_{181} z_{183}
\end{aligned}\end{equation*}
% rootProduct.1.equation3.1
\begin{equation*}\begin{aligned}
Q_{136}={}&
-4 z_{137} z_{188} +12 z_{137} z_{189} +z_{178} z_{185}\\
&+z_{179} z_{184} +z_{180} z_{183} +z_{181} z_{182}
\end{aligned}\end{equation*}
% rootProduct.1.equation4.0
\begin{equation*}\begin{aligned}
Q_{137}={}&
-44 z_{134} z_{190} +48 z_{134} z_{191} +z_{178}^{2}\\
&+2 z_{179}^{2}
\end{aligned}\end{equation*}
% rootProduct.1.equation4.1
\begin{equation*}\begin{aligned}
Q_{138}={}&
12 z_{134} z_{190} -22 z_{134} z_{191} +z_{178} z_{179}
\end{aligned}\end{equation*}
% rootProduct.1.equation5.0
\begin{equation*}\begin{aligned}
Q_{139}={}&
z_{70} -z_{190}^{2} -2 z_{191}^{2}
\end{aligned}\end{equation*}
% rootProduct.1.equation5.1
\begin{equation*}\begin{aligned}
Q_{140}={}&
z_{71} -z_{190} z_{191}
\end{aligned}\end{equation*}
% rootProduct.1.equation6.0
\begin{equation*}\begin{aligned}
Q_{141}={}&
11 z_{70} z_{134} -24 z_{71} z_{134} -z_{134}\\
&+z_{186}^{2} +2 z_{187}^{2}
\end{aligned}\end{equation*}
% rootProduct.1.equation6.1
\begin{equation*}\begin{aligned}
Q_{142}={}&
-3 z_{70} z_{134} +11 z_{71} z_{134} +z_{186} z_{187}
\end{aligned}\end{equation*}
% rootProduct.1.equation7.0
\begin{equation*}\begin{aligned}
Q_{143}={}&
11 z_{33} z_{134} -24 z_{34} z_{134} -z_{134}\\
&+z_{188}^{2} +2 z_{189}^{2}
\end{aligned}\end{equation*}
% rootProduct.1.equation7.1
\begin{equation*}\begin{aligned}
Q_{144}={}&
-3 z_{33} z_{134} +11 z_{34} z_{134} +z_{188} z_{189}
\end{aligned}\end{equation*}
% rootProduct.1.equation8.0
\begin{equation*}\begin{aligned}
Q_{145}={}&
6 z_{21} z_{188} -4 z_{21} z_{189} -4 z_{22} z_{188}\\
&+12 z_{22} z_{189} +z_{180}
\end{aligned}\end{equation*}
% rootProduct.1.equation8.1
\begin{equation*}\begin{aligned}
Q_{146}={}&
-2 z_{21} z_{188} +6 z_{21} z_{189} +6 z_{22} z_{188}\\
&-4 z_{22} z_{189} +z_{181}
\end{aligned}\end{equation*}
% rootProduct.1.equation9.0
\begin{equation*}\begin{aligned}
Q_{147}={}&
z_{19} z_{178} +2 z_{20} z_{179} -6 z_{47} z_{186}\\
&+4 z_{47} z_{187} +4 z_{48} z_{186} -12 z_{48} z_{187}
\end{aligned}\end{equation*}
% rootProduct.1.equation9.1
\begin{equation*}\begin{aligned}
Q_{148}={}&
z_{19} z_{179} +z_{20} z_{178} +2 z_{47} z_{186}\\
&-6 z_{47} z_{187} -6 z_{48} z_{186} +4 z_{48} z_{187}
\end{aligned}\end{equation*}
% rootProduct.2.equation1.0
\begin{equation*}\begin{aligned}
Q_{149}={}&
z_{192}^{2} +2 z_{193}^{2} +z_{194}^{2}\\
&+2 z_{195}^{2} +z_{196}^{2} +2 z_{197}^{2}\\
&-z_{198}^{2} -2 z_{199}^{2}
\end{aligned}\end{equation*}
% rootProduct.2.equation1.1
\begin{equation*}\begin{aligned}
Q_{150}={}&
z_{192} z_{193} +z_{194} z_{195} +z_{196} z_{197}\\
&-z_{198} z_{199}
\end{aligned}\end{equation*}
% rootProduct.2.equation2.0
\begin{equation*}\begin{aligned}
Q_{151}={}&
-12 z_{134} z_{200} +8 z_{134} z_{201} +z_{192} z_{196}\\
&+2 z_{193} z_{197} +z_{194} z_{198} +2 z_{195} z_{199}
\end{aligned}\end{equation*}
% rootProduct.2.equation2.1
\begin{equation*}\begin{aligned}
Q_{152}={}&
4 z_{134} z_{200} -12 z_{134} z_{201} +z_{192} z_{197}\\
&+z_{193} z_{196} +z_{194} z_{199} +z_{195} z_{198}
\end{aligned}\end{equation*}
% rootProduct.2.equation3.0
\begin{equation*}\begin{aligned}
Q_{153}={}&
12 z_{134} z_{202} -8 z_{134} z_{203} +z_{192} z_{198}\\
&+2 z_{193} z_{199} +z_{194} z_{196} +2 z_{195} z_{197}
\end{aligned}\end{equation*}
% rootProduct.2.equation3.1
\begin{equation*}\begin{aligned}
Q_{154}={}&
-4 z_{134} z_{202} +12 z_{134} z_{203} +z_{192} z_{199}\\
&+z_{193} z_{198} +z_{194} z_{197} +z_{195} z_{196}
\end{aligned}\end{equation*}
% rootProduct.2.equation4.0
\begin{equation*}\begin{aligned}
Q_{155}={}&
-44 z_{134} z_{204} +48 z_{134} z_{205} +z_{192}^{2}\\
&+2 z_{193}^{2}
\end{aligned}\end{equation*}
% rootProduct.2.equation4.1
\begin{equation*}\begin{aligned}
Q_{156}={}&
12 z_{134} z_{204} -22 z_{134} z_{205} +z_{192} z_{193}
\end{aligned}\end{equation*}
% rootProduct.2.equation5.0
\begin{equation*}\begin{aligned}
Q_{157}={}&
z_{72} -z_{204}^{2} -2 z_{205}^{2}
\end{aligned}\end{equation*}
% rootProduct.2.equation5.1
\begin{equation*}\begin{aligned}
Q_{158}={}&
z_{73} -z_{204} z_{205}
\end{aligned}\end{equation*}
% rootProduct.2.equation6.0
\begin{equation*}\begin{aligned}
Q_{159}={}&
-1 +11 z_{72} -24 z_{73}\\
&+z_{200}^{2} +2 z_{201}^{2}
\end{aligned}\end{equation*}
% rootProduct.2.equation6.1
\begin{equation*}\begin{aligned}
Q_{160}={}&
-3 z_{72} +11 z_{73} +z_{200} z_{201}
\end{aligned}\end{equation*}
% rootProduct.2.equation7.0
\begin{equation*}\begin{aligned}
Q_{161}={}&
-1 +11 z_{35} -24 z_{36}\\
&+z_{202}^{2} +2 z_{203}^{2}
\end{aligned}\end{equation*}
% rootProduct.2.equation7.1
\begin{equation*}\begin{aligned}
Q_{162}={}&
-3 z_{35} +11 z_{36} +z_{202} z_{203}
\end{aligned}\end{equation*}
% rootProduct.2.equation8.0
\begin{equation*}\begin{aligned}
Q_{163}={}&
6 z_{23} z_{202} -4 z_{23} z_{203} -4 z_{24} z_{202}\\
&+12 z_{24} z_{203} +z_{194}
\end{aligned}\end{equation*}
% rootProduct.2.equation8.1
\begin{equation*}\begin{aligned}
Q_{164}={}&
-2 z_{23} z_{202} +6 z_{23} z_{203} +6 z_{24} z_{202}\\
&-4 z_{24} z_{203} +z_{195}
\end{aligned}\end{equation*}
% rootProduct.2.equation9.0
\begin{equation*}\begin{aligned}
Q_{165}={}&
z_{19} z_{192} +2 z_{20} z_{193} -6 z_{51} z_{200}\\
&+4 z_{51} z_{201} +4 z_{52} z_{200} -12 z_{52} z_{201}
\end{aligned}\end{equation*}
% rootProduct.2.equation9.1
\begin{equation*}\begin{aligned}
Q_{166}={}&
z_{19} z_{193} +z_{20} z_{192} +2 z_{51} z_{200}\\
&-6 z_{51} z_{201} -6 z_{52} z_{200} +4 z_{52} z_{201}
\end{aligned}\end{equation*}
% rootProduct.3.equation1.0
\begin{equation*}\begin{aligned}
Q_{167}={}&
z_{206}^{2} +2 z_{207}^{2} +z_{208}^{2}\\
&+2 z_{209}^{2} +z_{210}^{2} +2 z_{211}^{2}\\
&-z_{212}^{2} -2 z_{213}^{2}
\end{aligned}\end{equation*}
% rootProduct.3.equation1.1
\begin{equation*}\begin{aligned}
Q_{168}={}&
z_{206} z_{207} +z_{208} z_{209} +z_{210} z_{211}\\
&-z_{212} z_{213}
\end{aligned}\end{equation*}
% rootProduct.3.equation2.0
\begin{equation*}\begin{aligned}
Q_{169}={}&
-12 z_{134} z_{214} +8 z_{134} z_{215} +z_{206} z_{210}\\
&+2 z_{207} z_{211} +z_{208} z_{212} +2 z_{209} z_{213}
\end{aligned}\end{equation*}
% rootProduct.3.equation2.1
\begin{equation*}\begin{aligned}
Q_{170}={}&
4 z_{134} z_{214} -12 z_{134} z_{215} +z_{206} z_{211}\\
&+z_{207} z_{210} +z_{208} z_{213} +z_{209} z_{212}
\end{aligned}\end{equation*}
% rootProduct.3.equation3.0
\begin{equation*}\begin{aligned}
Q_{171}={}&
12 z_{134} z_{216} -8 z_{134} z_{217} +z_{206} z_{212}\\
&+2 z_{207} z_{213} +z_{208} z_{210} +2 z_{209} z_{211}
\end{aligned}\end{equation*}
% rootProduct.3.equation3.1
\begin{equation*}\begin{aligned}
Q_{172}={}&
-4 z_{134} z_{216} +12 z_{134} z_{217} +z_{206} z_{213}\\
&+z_{207} z_{212} +z_{208} z_{211} +z_{209} z_{210}
\end{aligned}\end{equation*}
% rootProduct.3.equation4.0
\begin{equation*}\begin{aligned}
Q_{173}={}&
-44 z_{134} z_{218} +48 z_{134} z_{219} +z_{206}^{2}\\
&+2 z_{207}^{2}
\end{aligned}\end{equation*}
% rootProduct.3.equation4.1
\begin{equation*}\begin{aligned}
Q_{174}={}&
12 z_{134} z_{218} -22 z_{134} z_{219} +z_{206} z_{207}
\end{aligned}\end{equation*}
% rootProduct.3.equation5.0
\begin{equation*}\begin{aligned}
Q_{175}={}&
z_{74} -z_{218}^{2} -2 z_{219}^{2}
\end{aligned}\end{equation*}
% rootProduct.3.equation5.1
\begin{equation*}\begin{aligned}
Q_{176}={}&
z_{75} -z_{218} z_{219}
\end{aligned}\end{equation*}
% rootProduct.3.equation6.0
\begin{equation*}\begin{aligned}
Q_{177}={}&
-1 +11 z_{74} -24 z_{75}\\
&+z_{214}^{2} +2 z_{215}^{2}
\end{aligned}\end{equation*}
% rootProduct.3.equation6.1
\begin{equation*}\begin{aligned}
Q_{178}={}&
-3 z_{74} +11 z_{75} +z_{214} z_{215}
\end{aligned}\end{equation*}
% rootProduct.3.equation7.0
\begin{equation*}\begin{aligned}
Q_{179}={}&
-1 +11 z_{39} -24 z_{40}\\
&+z_{216}^{2} +2 z_{217}^{2}
\end{aligned}\end{equation*}
% rootProduct.3.equation7.1
\begin{equation*}\begin{aligned}
Q_{180}={}&
-3 z_{39} +11 z_{40} +z_{216} z_{217}
\end{aligned}\end{equation*}
% rootProduct.3.equation8.0
\begin{equation*}\begin{aligned}
Q_{181}={}&
6 z_{25} z_{216} -4 z_{25} z_{217} -4 z_{26} z_{216}\\
&+12 z_{26} z_{217} +z_{208}
\end{aligned}\end{equation*}
% rootProduct.3.equation8.1
\begin{equation*}\begin{aligned}
Q_{182}={}&
-2 z_{25} z_{216} +6 z_{25} z_{217} +6 z_{26} z_{216}\\
&-4 z_{26} z_{217} +z_{209}
\end{aligned}\end{equation*}
% rootProduct.3.equation9.0
\begin{equation*}\begin{aligned}
Q_{183}={}&
z_{19} z_{206} +2 z_{20} z_{207} -6 z_{55} z_{214}\\
&+4 z_{55} z_{215} +4 z_{56} z_{214} -12 z_{56} z_{215}
\end{aligned}\end{equation*}
% rootProduct.3.equation9.1
\begin{equation*}\begin{aligned}
Q_{184}={}&
z_{19} z_{207} +z_{20} z_{206} +2 z_{55} z_{214}\\
&-6 z_{55} z_{215} -6 z_{56} z_{214} +4 z_{56} z_{215}
\end{aligned}\end{equation*}
% rootProduct.4.equation1.0
\begin{equation*}\begin{aligned}
Q_{185}={}&
z_{220}^{2} +2 z_{221}^{2} +z_{222}^{2}\\
&+2 z_{223}^{2} +z_{224}^{2} +2 z_{225}^{2}\\
&-z_{226}^{2} -2 z_{227}^{2}
\end{aligned}\end{equation*}
% rootProduct.4.equation1.1
\begin{equation*}\begin{aligned}
Q_{186}={}&
z_{220} z_{221} +z_{222} z_{223} +z_{224} z_{225}\\
&-z_{226} z_{227}
\end{aligned}\end{equation*}
% rootProduct.4.equation2.0
\begin{equation*}\begin{aligned}
Q_{187}={}&
-12 z_{134} z_{228} +8 z_{134} z_{229} +z_{220} z_{224}\\
&+2 z_{221} z_{225} +z_{222} z_{226} +2 z_{223} z_{227}
\end{aligned}\end{equation*}
% rootProduct.4.equation2.1
\begin{equation*}\begin{aligned}
Q_{188}={}&
4 z_{134} z_{228} -12 z_{134} z_{229} +z_{220} z_{225}\\
&+z_{221} z_{224} +z_{222} z_{227} +z_{223} z_{226}
\end{aligned}\end{equation*}
% rootProduct.4.equation3.0
\begin{equation*}\begin{aligned}
Q_{189}={}&
12 z_{134} z_{230} -8 z_{134} z_{231} +z_{220} z_{226}\\
&+2 z_{221} z_{227} +z_{222} z_{224} +2 z_{223} z_{225}
\end{aligned}\end{equation*}
% rootProduct.4.equation3.1
\begin{equation*}\begin{aligned}
Q_{190}={}&
-4 z_{134} z_{230} +12 z_{134} z_{231} +z_{220} z_{227}\\
&+z_{221} z_{226} +z_{222} z_{225} +z_{223} z_{224}
\end{aligned}\end{equation*}
% rootProduct.4.equation4.0
\begin{equation*}\begin{aligned}
Q_{191}={}&
-44 z_{134} z_{232} +48 z_{134} z_{233} +z_{220}^{2}\\
&+2 z_{221}^{2}
\end{aligned}\end{equation*}
% rootProduct.4.equation4.1
\begin{equation*}\begin{aligned}
Q_{192}={}&
12 z_{134} z_{232} -22 z_{134} z_{233} +z_{220} z_{221}
\end{aligned}\end{equation*}
% rootProduct.4.equation5.0
\begin{equation*}\begin{aligned}
Q_{193}={}&
z_{76} -z_{232}^{2} -2 z_{233}^{2}
\end{aligned}\end{equation*}
% rootProduct.4.equation5.1
\begin{equation*}\begin{aligned}
Q_{194}={}&
z_{77} -z_{232} z_{233}
\end{aligned}\end{equation*}
% rootProduct.4.equation6.0
\begin{equation*}\begin{aligned}
Q_{195}={}&
-1 +11 z_{76} -24 z_{77}\\
&+z_{228}^{2} +2 z_{229}^{2}
\end{aligned}\end{equation*}
% rootProduct.4.equation6.1
\begin{equation*}\begin{aligned}
Q_{196}={}&
-3 z_{76} +11 z_{77} +z_{228} z_{229}
\end{aligned}\end{equation*}
% rootProduct.4.equation7.0
\begin{equation*}\begin{aligned}
Q_{197}={}&
-1 +11 z_{43} -24 z_{44}\\
&+z_{230}^{2} +2 z_{231}^{2}
\end{aligned}\end{equation*}
% rootProduct.4.equation7.1
\begin{equation*}\begin{aligned}
Q_{198}={}&
-3 z_{43} +11 z_{44} +z_{230} z_{231}
\end{aligned}\end{equation*}
% rootProduct.4.equation8.0
\begin{equation*}\begin{aligned}
Q_{199}={}&
6 z_{27} z_{230} -4 z_{27} z_{231} -4 z_{28} z_{230}\\
&+12 z_{28} z_{231} +z_{222}
\end{aligned}\end{equation*}
% rootProduct.4.equation8.1
\begin{equation*}\begin{aligned}
Q_{200}={}&
-2 z_{27} z_{230} +6 z_{27} z_{231} +6 z_{28} z_{230}\\
&-4 z_{28} z_{231} +z_{223}
\end{aligned}\end{equation*}
% rootProduct.4.equation9.0
\begin{equation*}\begin{aligned}
Q_{201}={}&
z_{19} z_{220} +2 z_{20} z_{221} -6 z_{59} z_{228}\\
&+4 z_{59} z_{229} +4 z_{60} z_{228} -12 z_{60} z_{229}
\end{aligned}\end{equation*}
% rootProduct.4.equation9.1
\begin{equation*}\begin{aligned}
Q_{202}={}&
z_{19} z_{221} +z_{20} z_{220} +2 z_{59} z_{228}\\
&-6 z_{59} z_{229} -6 z_{60} z_{228} +4 z_{60} z_{229}
\end{aligned}\end{equation*}
% rootProduct.S.retained_j.0
\begin{equation*}\begin{aligned}
Q_{203}={}&
4 z_{2} +z_{192} z_{194} +2 z_{193} z_{195}\\
&-z_{196} z_{198} -2 z_{197} z_{199}
\end{aligned}\end{equation*}
% rootProduct.S.retained_j.1
\begin{equation*}\begin{aligned}
Q_{204}={}&
4 z_{3} +z_{192} z_{195} +z_{193} z_{194}\\
&-z_{196} z_{199} -z_{197} z_{198}
\end{aligned}\end{equation*}
% scalar.S.direct_ratio.0
\begin{equation*}\begin{aligned}
Q_{205}={}&
2 z_{2} z_{124} +4 z_{3} z_{125} -z_{19} z_{198}\\
&-2 z_{20} z_{199}
\end{aligned}\end{equation*}
% scalar.S.direct_ratio.1
\begin{equation*}\begin{aligned}
Q_{206}={}&
2 z_{2} z_{125} +2 z_{3} z_{124} -z_{19} z_{199}\\
&-z_{20} z_{198}
\end{aligned}\end{equation*}
% reciprocal.aggregate.rootTrace.definition
\begin{equation*}\begin{aligned}
Q_{207}={}&
8 z_{4} z_{106} -z_{184}^{2} -2 z_{185}^{2}\\
&-z_{198}^{2} -2 z_{199}^{2} -z_{212}^{2}\\
&-2 z_{213}^{2} -z_{226}^{2} -2 z_{227}^{2}
\end{aligned}\end{equation*}
% END V15X FINAL 207 EQUATIONS
\subsection{Fixed integer constants for the final system}
The following decimal strings are concatenated across line breaks.
\begin{longtable}{r p{0.82\textwidth}}Constant&Exact positive integer\\\hline
$K_{1}$&{\ttfamily 98194994952101164627950993228801351760\newline 40946699131400335899642243325954265276\newline 68691070368153600000000000000}\\[4pt]
$K_{2}$&{\ttfamily 80990493410131000784375977023646274345\newline 82049009524401528919498956985577266824\newline 1364832628280392076492800000000000000}\\[4pt]
$K_{3}$&{\ttfamily 45894612932407567111146386980066222129\newline 29827772063827533054382742291827117867\newline 01067384893588888433459200000000000000}\\[4pt]
$K_{4}$&{\ttfamily 91789225864815134222292773960132444258\newline 59655544127655066108765484583654235734\newline 02134769787177776866918400000000000000}\\[4pt]
$K_{5}$&{\ttfamily 13768383879722270133343916094019866638\newline 78948331619148259916314822687548135360\newline 10320215468076666530037760000000000000\newline 0}\\[4pt]
$K_{6}$&{\ttfamily 27536767759444540266687832188039733277\newline 57896663238296519832629645375096270720\newline 20640430936153333060075520000000000000\newline 0}\\[4pt]
$K_{7}$&{\ttfamily 63189464886465357187317594522928066201\newline 84449792472930176268493164014380154729\newline 61541724103664894830836146098454811858\newline 89826867345597467459165957964726042449\newline 88128011149397522493691010836323796918\newline 78234193920000000000000000000000000000}\\[4pt]
$K_{8}$&{\ttfamily 56870518397818821468585835070635259581\newline 66004813225637158641643847612942139256\newline 65387551693298405347752531488609330673\newline 00844180611037720713249362168253438204\newline 89315210034457770244321909752691417226\newline 90410774528000000000000000000000000000\newline 0}\\[4pt]
\end{longtable}
\subsection{Exact coefficient profile and its scope}\label{v15x:height}
Expanding these 207 integer squares and collecting coordinate monomials
gives 2514 nonzero monomials (2515 terms if different powers of $A_N$
are kept separate). The coefficient of $z_4^4$ is 1, so the degree is
four and the content is one. All 233 coordinates occur. The only
parameter-dependent residual is
\[
 Q_{86}=22201z_{123}z_{137}+z_{137}-A_Nz_{123}.
\]
Thus every parameter-dependent coefficient can be checked directly
from its square; all other contributions are constant. The largest
absolute coefficient independent of $A_N$ is the integer $C_*$ below,
and the largest absolute coefficient of $A_N$ is 44402.
\begin{center}\small\ttfamily
3234255862836649052395809050513757712720766678827201051\par 4121214447529782641164116187119030726093150720667902069\par 1146333036638193446287119593493174002525964438855275658\par 9725679097130070419240959939043673118145265566835699850\par 6023069549107151770397486097251277767085767051395915412\par 3470651877805122130468753046467052593626986399988193307\par 1819781041661355765844830798721391874678242849435643063\par 5201635285362278400000000000000000000000000000000000000\par 000000000000000000
\end{center}
The coefficient of $z_{123}^2$ is $A_N^2+1$. The remaining
parameter-dependent coefficients have absolute value at most
$44402A_N$. Direct integer comparison gives
$C_*<149^{214}+1$. For $N\ge53$, $A_N\ge149^{107}>44402$;
hence both $C_*$ and $44402A_N$ are at most $A_N^2+1$.
Therefore the exact height is
\[
 H(P_N)=A_N^2+1=149^{4N+2}+1\qquad(N\ge53).
\]
The displayed equations, constants and finite coefficient calculation
specify the family without choosing a numerical defining member.
In particular $N=53$ is a coefficient threshold only. The assertion
that $P_N$ defines $\Z$ requires $N\ge N_0$, with $N_0$ still ineffective.



\clearpage
\allowdisplaybreaks
\section{The 222-variable family}\label{v222:construction}
The 233-variable construction above is a predecessor. We now give the
additional arithmetic change and exact algebraic maps leading to the
222-variable family in Theorem~\ref{thm:v15-main}. The distinguished
input is fixed throughout. Write $A=A_N=149^{2N+1}$ and
$K=\Q(\rho)$, $\rho^2=2$, $d=3-\rho$, $D_c=d^2=11-6\rho$.
The numerical condition $N\ge53$ below concerns coefficient size only;
the sufficient defining threshold remains ineffective.

\subsection{Four finite contacts and 232 coordinates}\label{v222:contacts}
This first step is a change of arithmetic construction, not a bijection
between all predecessor witnesses. Put
\[
 \beta=2504+51x^2,\qquad
 r_4=\beta(\beta-1)(\beta-2401)(32\beta+343),\qquad
 R={A r_4\over1+149^2r_4}.
\]
The finite contacts of the retained $M_5(\beta)$ are precisely
$0,1,2401,-343/32$; its fifth contact is infinity. Thus the extra
predecessor factor $2\beta-1$ is unnecessary. In the corrected raw233
indexing of Section~\ref{v15x:dictionary}, retain
$t=x_{63}=x^2$ and $u=x_{121}=(51t+2501)^2/4$, delete $x_{122}$ and
its defining equation, and replace the contact-product equation by
\[
 x_{123}+14315517724+291790992t+15531948tu
       +767410608u-512u^2=0.                 \tag{C232}
\]
This is exactly $x_{123}=r_4$. Compact indices after the deleted slot.
The other205 equations retain their meanings, giving232 coordinates
and206 quadratic equations.

For every rational $x$, all four factors of $r_4$ are positive and
$r_4\ge1$. Consequently $R>0$ and
$A/(1+149^2)\le R<A/149^2$. At149 the four residues
$(\theta-2504)/51$, for the four finite contacts $\theta$, are
$59,97,109,108$, each a nonsquare. If $v_{149}(x)\ge0$, $r_4$ is a
unit; otherwise $v_{149}(r_4)=8v_{149}(x)$ and
$v_{149}(1+149^2r_4)=2+8v_{149}(x)$. Thus for $N\ge1$, $v_{149}(R)$ is respectively
$2N+1$ or $2N-1$, positive and odd. The unchanged scalar argument of
Section~\ref{v15x:parity} therefore still gives the forced149 depth
and the reference exceptional-value exclusions.

For $p>76176$ contributing to $M_5(\beta)$, exactly one finite
contact factor has positive odd valuation. The constant differences
and32 are units at such a prime, so $v_p(R)=v_p(r_4)>0$ is odd.
If $\beta$ is nonintegral then $v_p(\beta)=2v_p(x)$ is even and the
infinite contact supplies no odd factor. The same local extraction
therefore gives, for the nonzero scalar index $a$,
\[
149^{\lfloor9(2N-1)/4\rfloor+1}M_5(\beta)^3\mid |a|.
\]
This is the divisibility input in the two retained height/Lucas
contradictions. Those arguments, and the squarefree51 denominator
argument recovering $x\in\Z$, apply unchanged.

For completeness, fix an integer $x$. A direct parity calculation gives
$v_2(r_4)\in\{2,3\}$. Modulo3, $r_4=1$ and $1+149^2r_4=2$;
modulo17, $r_4=13$ and $1+149^2r_4=5$. Hence
$v_3(R)=v_{17}(R)=0$ and $v_2(R)\in\{2,3\}$.
Since $9792=2^6 3^2 17$, the element $1+9792/R$ is a square in each
of $\Q_2,\Q_3,\Q_{17}$. In particular the changed value $v_3(R)=0$
meets the actual local-square requirement. If $D_0=1+149^2r_4$,
then $\gcd(Ar_4,D_0)=\gcd(Ar_4+9792D_0,D_0)=1$, so the reduced
denominator of $R(R+9792)$ is the square $D_0^2$.
The same compact real neighborhoods and finite-prime recurrence used
above now give
\[
 W=Z_0^4/R^{10}+Z_1^2/R^5+F^2/R^4=O(R^{-4}),\qquad
 b=R(R+9792)W^2=O(R^{-6})>0.
\]
The bound is uniform in the integer input, so $0<b<17$ for a common
sufficiently large $N$. First choose $H$ sufficiently formal at the
finite support of $R(R+9792)$ and then choose the scalar index deeper;
all terms of $W$ meet the same local integrality and square conditions.
The retained recurrence/progression argument meets these conditions
and the real neighborhoods simultaneously. The later contact primes
are chosen outside this finite support, in the original order.
This constructs new witnesses for every integer, without imposing any
new exact-square condition or claiming a bijection of old pole witnesses.

\subsection{One common alternate ratio}\label{v222:ratio}
The useful direct-ratio idea is due to Jos\'e Brox's
\emph{Fixed-difference charts}, version18, public revision
\texttt{34454098637ba9100b19071da55a38199589f5e0}, equations for formal-parameter
ratios. Here we retain the preceding reference chart and prove the
needed domain from the scalar149 argument; no new arithmetic theorem
from that revision is assumed.

In raw232 coordinates write $f=x_{17}+\rho x_{18}$,
$S=x_{19}+\rho x_{20}$ and $g=x_{134}+\rho x_{135}$. Put
\[
 \alpha=156+\tfrac{69}{2}\rho,\quad \beta_0=118+60\rho,
 \quad\gamma=-288-128\rho,\quad\delta=252-672\rho.
\]
\[
N_r=\alpha S+\gamma+\delta g,\quad D_r=\beta_0S+\alpha.
\]
For $i=0,1,2,3$, the actual alternate equations have the form
$N_r z_i+D_rL_i'=0$ and $E_i+K_i z_i=0$, where
\[
\begin{split}
z_i&=x_{21+2i}+\rho x_{22+2i},\\
L_i'&=x_{49+4i}+\rho x_{50+4i},\\
K_i&=2d(x_{187+14i}+\rho x_{188+14i}),\\
E_i&=(x_{177+14i}+\rho x_{178+14i})
 -(x_{179+14i}+\rho x_{180+14i}).
\end{split}
\]
These are all the consumers of the eight $z_i$ components. Replace them
by a single field coordinate $T$ and the five field equations
\[
D_rT=N_r,\qquad E_iT=K_iL_i'\quad(0\le i\le3).
\]
At the inert prime149 the unchanged reference and scalar equations,
without any alternate-ratio equation, give
$v(f)=m>0$, $v(S)=-m/2$ and $v(g)=m/2$.
The norms $\Norm(\alpha)=43911/2$ and $\Norm(\beta_0)=6724$ are units.
Thus $N_r,D_r$ have the same finite valuation $-m/2$, and $T$ is a
nonzero unit. The maps are $T=N_r/D_r$ and $z_i=-L_i'/T$.
They never divide by a numerator, $K_i$, or a rational component;
all their zero strata are included. The identity
\[
D_r(E_iT-K_iL_i')=D_rT(E_i+K_iz_i)
 -K_i(N_rz_i+D_rL_i')-K_iz_i(D_rT-N_r)
\]
checks the forward implication directly. Delete raw232 slots21--28,
compact the survivors and append the two components of $T$.
This gives226 coordinates and200 equations.

\subsection{Polynomial reference elimination}\label{v222:reference}
Write $a=57-28\rho$, $b=8-14\rho$, $c=88-48\rho$, $e=39-36\rho$.
The reference equations are
\[
\mathcal F=(a+cf)g+bfS+ef=0,\quad
\mathcal E=bSg-eg-af+8=0,\quad
(a+cf)^2w=(8-af-2eg)^2.
\]
Here $B=f^2$, $w=v^2$ and $W=w^2$ are the retained square outputs.
Eliminate $g$ by the polynomial formula
\[
g=((\beta_0S+\alpha)T-\alpha S-\gamma)/\delta.       \tag{G}
\]
Only a nonzero constant is divided by. The following six explicitly
specified quadrics in $f,S,B,w,W,T$ are denoted $Q_0,\ldots,Q_5$.
Their zero numbering is local to this subsection.

{\scriptsize
\begin{align*}
Q_0&=\left(-1+0\rho\right) B\\
&+\left(1+0\rho\right) f^{2} .
\end{align*}
\begin{align*}
Q_1&=\left(\frac{21373536}{4905889}-\frac{4638784}{4905889}\rho\right) \\
&+\left(\frac{76828374}{4905889}-\frac{48988044}{4905889}\rho\right) f\\
&+\left(-\frac{87}{41}+\frac{129}{164}\rho\right) S\\
&+\left(\frac{279649719}{201141449}-\frac{615250593}{804565796}\rho\right) T\\
&+\left(\frac{26325600}{4905889}-\frac{13327692}{4905889}\rho\right) f T\\
&+\left(1+0\rho\right) S T .
\end{align*}
\begin{align*}
Q_2&=\left(-\frac{287384773949078688806717371}{430762914076126579926624016}-\frac{81557625814705269604141077}{215381457038063289963312008}\rho\right) \\
&+\left(\frac{1256455567493584412665635}{2197769969776156020033796}+\frac{214486582192171922617190}{549442492444039005008449}\rho\right) f\\
&+\left(-\frac{4766844858863040}{111996519375803041}-\frac{3149175549566448}{111996519375803041}\rho\right) S\\
&+\left(-\frac{86993}{153664}-\frac{4947}{38416}\rho\right) B\\
&+\left(\frac{3317607896141091533}{7167777240051394624}+\frac{677690583653850567}{1791944310012848656}\rho\right) w\\
&+\left(\frac{9191218419936655190467815}{26922682129757911245414001}+\frac{47475418018975282967719917}{215381457038063289963312008}\rho\right) T\\
&+\left(\frac{171357}{235298}+\frac{2519877}{7529536}\rho\right) f S\\
&+\left(\frac{23335979380411305}{9142573010269636}-\frac{256685364214285}{4571286505134818}\rho\right) f w\\
&+\left(-\frac{47892880950890591270656509}{53845364259515822490828002}-\frac{647265901015780366495244709}{1723051656304506319706496064}\rho\right) f T\\
&+\left(\frac{44649296511110307}{111996519375803041}-\frac{2886885451827729}{447986077503212164}\rho\right) S w\\
&+\left(1+0\rho\right) B w .
\end{align*}
\begin{align*}
Q_3&=\left(\frac{570}{2009}+\frac{669}{8036}\rho\right) \\
&+\left(\frac{87}{41}-\frac{129}{164}\rho\right) B\\
&+\left(-\frac{11}{49}-\frac{6}{49}\rho\right) T\\
&+\left(-\frac{667923}{658952}-\frac{72481}{329476}\rho\right) f S\\
&+\left(\frac{310419}{658952}-\frac{15327}{329476}\rho\right) f T\\
&+\left(1+0\rho\right) B T .
\end{align*}
\begin{align*}
Q_4&=\left(-1+0\rho\right) W\\
&+\left(1+0\rho\right) w^{2} .
\end{align*}
\begin{align*}
Q_5&=\left(-\frac{118216066002595658088988042821122451634498060}{217387430662605583226669899282540047335778051}-\frac{50186624466822596896077228706581721089715127}{434774861325211166453339798565080094671556102}\rho\right) \\
&+\left(\frac{1200047839547464829478234769214860335942940}{31055347237515083318095699897505721047968293}+\frac{3580415723464424486437662870330074238560465}{31055347237515083318095699897505721047968293}\rho\right) f\\
&+\left(\frac{668475976652228193529571247574784}{17731703300321543710787069886827041}+\frac{537773276227497287456656778348544}{17731703300321543710787069886827041}\rho\right) S\\
&+\left(-\frac{239794457122086328917179}{542833069831184086820628}+\frac{221507061335393600426173}{2171332279324736347282512}\rho\right) B\\
&+\left(\frac{439546863040207888071363077708623743}{2907999341252733168569079461439634724}+\frac{4872456519496576590629658909470099559}{11631997365010932674276317845758538896}\rho\right) w\\
&+\left(\frac{16558750559000902978384}{135708267457796021705157}+\frac{39160642255580373322972}{135708267457796021705157}\rho\right) W\\
&+\left(\frac{1461967417941685926050299573323630374958742}{4262498640443246729934703907500785241878001}+\frac{760475724129912204044824303873051497411852}{4262498640443246729934703907500785241878001}\rho\right) T\\
&+\left(\frac{64581164203274943033219}{74506499780750757014596}+\frac{2334048428878964391793}{37253249890375378507298}\rho\right) f S\\
&+\left(\frac{69171671661784723704468037462721535400}{6330218078214791104750983949597253637}-\frac{29830129278260699152917591268341094453}{6330218078214791104750983949597253637}\rho\right) f w\\
&+\left(\frac{798843646842276177602560}{135708267457796021705157}-\frac{438773961053992851161920}{135708267457796021705157}\rho\right) f W\\
&+\left(-\frac{20372751135700507179178457331278023758792291}{17049994561772986919738815630003140967512004}-\frac{603007143259447425286649316925571213407697}{8524997280886493459869407815001570483756002}\rho\right) f T\\
&+\left(\frac{2843253671326083223981015767800324}{17731703300321543710787069886827041}-\frac{7340159548393698461511516883879680}{17731703300321543710787069886827041}\rho\right) S w\\
&+\left(\frac{347346545590041266304}{380135202963014066401}-\frac{189974428707723093504}{380135202963014066401}\rho\right) S W\\
&+\left(\frac{12533845232345798175800320}{5564038965769636889911437}-\frac{6772286266558612146808064}{5564038965769636889911437}\rho\right) B W\\
&+\left(1+0\rho\right) w T .
\end{align*}
}
The following polynomial identities certify reference recovery after (G):
{\scriptsize
\begin{align*}
\mathcal F&=\left(-\frac{92}{119}-\frac{1280}{357}\rho\right) f Q_1\\
&+\left(-\frac{75218560}{4905889}+\frac{84084944}{4905889}\rho\right) T Q_0\\
&+\left(-\frac{291890088}{4905889}+\frac{237589968}{4905889}\rho\right)  Q_0\\
&+\left(-\frac{75218560}{4905889}+\frac{84084944}{4905889}\rho\right)  Q_3\\
&+\left(-\frac{5977}{4998}-\frac{2273}{833}\rho\right)  Q_1 .
\end{align*}
\begin{align*}
\mathcal E&=\left(\frac{5240}{2499}+\frac{2351}{2499}\rho\right) S Q_1\\
&+\left(-\frac{30123784}{4905889}+\frac{3179520}{4905889}\rho\right) f Q_1\\
&+\left(\frac{877778014606080}{24067746880321}-\frac{485183286738528}{24067746880321}\rho\right) T Q_0\\
&+\left(\frac{2625878274764976}{24067746880321}-\frac{1719982607738976}{24067746880321}\rho\right)  Q_0\\
&+\left(-\frac{709839747}{480777122}+\frac{71029071}{240388561}\rho\right)  Q_1\\
&+\left(\frac{877778014606080}{24067746880321}-\frac{485183286738528}{24067746880321}\rho\right)  Q_3 .
\end{align*}
}

The recovery identities use only $Q_0,Q_1,Q_3$. If $f=0$, they force
$B=0$, $T=87/41-129\rho/164$, and then the nonzero constant
$114441/13448-12311\rho/3362=0$, whose norm is $4905889/107584$.
Furthermore
\[(bS-e)\mathcal F-(a+cf)\mathcal E=8a(fS^2+D_cf^2-1).\]
Thus the finite $H$ curve and $f\ne0$ hold before interpreting $w$.
The unchanged scalar149 argument gives $v(f)=m>0$, $v(S)=-m/2$,
and $v(g)=m/2$, independently of $w$. In $Q_2=Mw+L$ the coefficient is
\begin{align*}
M={}&B+\left({23335979380411305\over9142573010269636}
 -{256685364214285\over4571286505134818}\rho\right)f\\
&+\left({44649296511110307\over111996519375803041}
 -{2886885451827729\over447986077503212164}\rho\right)S\\
&+{3317607896141091533\over7167777240051394624}
 +{677690583653850567\over1791944310012848656}\rho .
\end{align*}
All coefficients are149-integral, and the norm of the $S$ coefficient
is $414947904808977/2612163717219896$, a149-unit. Thus its $S$ term is
uniquely lowest and $M\ne0$. Substitute $g=-f(bS+e)/(a+cf)$,
$w_*=g^4/f^2$ and $T=N_r/D_r$ into the listed quadrics and reduce by
$S^2=f^{-1}-D_cf$. After multiplying by $(a+cf)^{16}D_r^2$ each
remainder is zero. Both factors are nonzero by the preceding valuation
argument. In particular $Q_2=0$ forces $w=w_*$; $Q_4=0$ gives $W=w^2$,
and $Q_5$ follows. Retain only $Q_0,Q_1,Q_2,Q_3,Q_4$.
Delete raw226 slots126,127 and compact. This gives224 coordinates and198
quadrics with inverse (G), including all zero numerator and component
branches. In raw224 the actual output bases are
\[
 B=x_{74}+2\rho x_{75},\quad w=x_{21}+2\rho x_{22},\quad
 W={37249x_{23}-34848x_{24}\over2401}
 +\rho{25476(x_{23}-x_{24})\over2401}.
\]

\subsection{Paired factors and three root normalizations}\label{v222:roots}
Until the next subsection indices refer to raw224. Set
$q_0=x_{87}$, $q=x_{102}$, $Q_0=x_{89}=q_0^2$,
$Q_1=x_{104}=q^2$, $h_0=x_{88}=q_0^7$, $h_1=x_{103}=q^7$,
$j_0=x_{90}=q_0^4$, $j_1=x_{105}=q^4$, $R=x_{126}$, $R^2=x_{125}$.
The retained contact equations imply $q\ne0$ without using a root or
trace equation. After division by the nonzero $R^2$ they give
$(1+q)P_5(r)=r^5$, $r=(t-1)^2$. Modulo5,
$P_5(r)=(r-1)^2(r-4)^3$. If $q=0$ and $t$ is5-integral, the square
residues $r=0,1,4$ contradict this equation. If $v_5(t)<0$,
$P_5(r)-r^5$ has degree4 and unit leading coefficient, so its leading
term cannot cancel. This argument preserves $q=-1$.

Delete slots62,63,66,67,68,69. For each tuple
\[
(F_0,F_1,U_0,U_1)=(62,63,179,180),(66,67,207,208),(68,69,221,222)
\]
replace $U$ by $qU$ and recover the deleted output by
\[
 x_{F_0}^{\rm old}={(U_0^{\rm new})^2+2(U_1^{\rm new})^2\over q^2},
 \qquad x_{F_1}^{\rm old}={U_0^{\rm new}U_1^{\rm new}\over q^2}.
\]
The other changes, with unlisted coordinates fixed, are
\[
\begin{array}{c|c}
\text{raw224 slots}&\text{new value}\\\hline
167\!:\!178,\ 195\!:\!202,\ 209\!:\!216 &(q/R)\,\text{old value}\\
179,180,203\!:\!208,217\!:\!222&q\,\text{old value}\\
47\!:\!54,56\!:\!59&\text{old value}/R
\end{array}
\]
The inverse multiplies by the reciprocal indicated factor. In particular
the second root and physical scalar ratio stay fixed. The differing
first-root $i,k$ normalization is included. Literal substitution in
all product, comparison and Lucas equations gives quadrics; the
ordinary $R^2$ factors cancel against the rescaled label/product
coordinates. The six deleted square definitions become identities.

Introduce
\[
H=Q_1h_0-Q_0h_1,\quad J=Q_1j_0-Q_0j_1,\quad K_*=Q_1q_0-Q_0q.
\]
For matching factor expressions $F_j=1+h_jA+Q_jB+j_jC+q_jS_*$,
where $S_*$ occurs only in factor2,
\[
Q_1F_0-Q_0F_1=Q_1-Q_0+HA+JC+K_*S_*.
\]
Apply this to factors1,2,4, retaining the detector1 equation and replacing
the detector0 equation by the combination. The deleted $q^2F_U$ terms
in detector1 become the quadratic expressions in $U$ above. The special
$qF_U$ output of root2 is retained. Since $Q_1\ne0$, this row operation
is reversible. All actual factor values, product trees, norm outputs
and pole inputs are unchanged.

Introduce also
\[
G=8x_4x_{98}-[(x_{185}-x_{187})^2+2(x_{186}-x_{188})^2]
\]
using the unchanged second root, and let $Z$ be the sum of the three
new trace squares for pairs $(171,173),(199,201),(213,215)$ and their
imaginary pairs with coefficient2. The new trace is $Q_1G-R^2Z=0$,
which is $q^2$ times the old trace after substitution. Delete the six
output slots, compact in increasing order, and append $H,J,K_*,G,Z$.
There are223 coordinates and197 quadrics. No component, label or root
numerator has been divided by.

\subsection{The final common quotient}\label{v222:quotient}
In raw223 the only consumers of $G=x_{222},Z=x_{223}$ are
\[
Q_1G-R^2Z=0,\quad G=P-T_2,\quad Z=T_1+T_3+T_4,
\]
where $Q_1=x_{98}$, $R^2=x_{119}$ and $T_i$ are the trace squares.
Keep one coordinate $U_*=Z/Q_1$ and replace these equations by
\[
R^2U_*-P+T_2=0,\qquad Q_1U_*-T_1-T_3-T_4=0.
\]
They are quadratic. The inverse is polynomial, $G=R^2U_*$,
$Z=Q_1U_*$, and the old trace becomes identically zero.
The forward quotient exists by the unchanged contact proof of $q\ne0$.
No trace value or rational component is inverted. Keep slots1 through221
and number $U_*$ as222. This proves222 variables and196 quadrics with
all zero branches preserved. The input, arithmetic hypotheses and
ineffective sufficient threshold remain unchanged.

\subsection{Exact final affine and row presentation}\label{v222:maps}
The following ordered affine substitutions use raw223 indices and give
old coordinates in terms of new ones. They do not touch slots222,223,98
or119, so commute with the final common quotient. Unlisted coordinates
are fixed. Each has a rational linear inverse.

\paragraph{Affine stage 1.}
\begin{align*}
x_{56}^{\rm old}&=2 x_{51}^{\rm new}+x_{56}^{\rm new}\\
x_{57}^{\rm old}&=2 x_{52}^{\rm new}+x_{57}^{\rm new}\\
x_{58}^{\rm old}&=x_{56}^{\rm new}+x_{58}^{\rm new}\\
x_{59}^{\rm old}&=x_{57}^{\rm new}+x_{59}^{\rm new}
\end{align*}
Replace residual 76 by itself plus $-36$ times residual 18.
Replace residual 77 by itself plus $-72$ times residual 19.
Replace residual 1 by itself plus $-\frac{193}{33}$ times residual 2.
Replace residual 3 by itself plus $-\frac{193}{33}$ times residual 4.
Replace residual 5 by itself plus $-\frac{193}{33}$ times residual 6.
Replace residual 7 by itself plus $-\frac{193}{33}$ times residual 8.
Replace residual 32 by itself plus $3813373161796$ times residual 33.
Replace residual 45 by itself plus $3813373161796$ times residual 46.

\paragraph{Affine stage 2.}
\begin{align*}
x_{55}^{\rm old}&=-\frac{2501}{51} +x_{55}^{\rm new}\\
x_{66}^{\rm old}&=-9 +x_{66}^{\rm new}\\
x_{76}^{\rm old}&=-1 +x_{76}^{\rm new}
\end{align*}

\paragraph{Affine stage 3.}
\begin{align*}
x_{23}^{\rm old}&=x_{23}^{\rm new}+x_{24}^{\rm new}
\end{align*}

\paragraph{Affine stage 4.}
\begin{align*}
x_{29}^{\rm old}&=x_{29}^{\rm new}+x_{30}^{\rm new}
\end{align*}

\paragraph{Affine stage 5.}
\begin{align*}
x_{33}^{\rm old}&=x_{33}^{\rm new}+x_{34}^{\rm new}
\end{align*}

\paragraph{Affine stage 6.}
\begin{align*}
x_{37}^{\rm old}&=x_{37}^{\rm new}+x_{38}^{\rm new}
\end{align*}

\paragraph{Affine stage 7.}
\begin{align*}
x_{161}^{\rm old}&=\frac{1}{2} x_{161}^{\rm new}+\frac{1}{2} x_{163}^{\rm new}\\
x_{162}^{\rm old}&=\frac{1}{2} x_{162}^{\rm new}+\frac{1}{2} x_{164}^{\rm new}\\
x_{163}^{\rm old}&=\frac{1}{2} x_{161}^{\rm new}-\frac{1}{2} x_{163}^{\rm new}\\
x_{164}^{\rm old}&=\frac{1}{2} x_{162}^{\rm new}-\frac{1}{2} x_{164}^{\rm new}\\
x_{165}^{\rm old}&=\frac{1}{2} x_{165}^{\rm new}+\frac{1}{2} x_{167}^{\rm new}\\
x_{166}^{\rm old}&=\frac{1}{2} x_{166}^{\rm new}+\frac{1}{2} x_{168}^{\rm new}\\
x_{167}^{\rm old}&=\frac{1}{2} x_{165}^{\rm new}-\frac{1}{2} x_{167}^{\rm new}\\
x_{168}^{\rm old}&=\frac{1}{2} x_{166}^{\rm new}-\frac{1}{2} x_{168}^{\rm new}\\
x_{175}^{\rm old}&=\frac{1}{2} x_{175}^{\rm new}+\frac{1}{2} x_{177}^{\rm new}\\
x_{176}^{\rm old}&=\frac{1}{2} x_{176}^{\rm new}+\frac{1}{2} x_{178}^{\rm new}\\
x_{177}^{\rm old}&=\frac{1}{2} x_{175}^{\rm new}-\frac{1}{2} x_{177}^{\rm new}\\
x_{178}^{\rm old}&=\frac{1}{2} x_{176}^{\rm new}-\frac{1}{2} x_{178}^{\rm new}\\
x_{179}^{\rm old}&=\frac{1}{2} x_{179}^{\rm new}+\frac{1}{2} x_{181}^{\rm new}\\
x_{180}^{\rm old}&=\frac{1}{2} x_{180}^{\rm new}+\frac{1}{2} x_{182}^{\rm new}\\
x_{181}^{\rm old}&=\frac{1}{2} x_{179}^{\rm new}-\frac{1}{2} x_{181}^{\rm new}\\
x_{182}^{\rm old}&=\frac{1}{2} x_{180}^{\rm new}-\frac{1}{2} x_{182}^{\rm new}\\
x_{189}^{\rm old}&=\frac{1}{2} x_{189}^{\rm new}+\frac{1}{2} x_{191}^{\rm new}\\
x_{190}^{\rm old}&=\frac{1}{2} x_{190}^{\rm new}+\frac{1}{2} x_{192}^{\rm new}\\
x_{191}^{\rm old}&=\frac{1}{2} x_{189}^{\rm new}-\frac{1}{2} x_{191}^{\rm new}\\
x_{192}^{\rm old}&=\frac{1}{2} x_{190}^{\rm new}-\frac{1}{2} x_{192}^{\rm new}\\
x_{193}^{\rm old}&=\frac{1}{2} x_{193}^{\rm new}+\frac{1}{2} x_{195}^{\rm new}\\
x_{194}^{\rm old}&=\frac{1}{2} x_{194}^{\rm new}+\frac{1}{2} x_{196}^{\rm new}\\
x_{195}^{\rm old}&=\frac{1}{2} x_{193}^{\rm new}-\frac{1}{2} x_{195}^{\rm new}\\
x_{196}^{\rm old}&=\frac{1}{2} x_{194}^{\rm new}-\frac{1}{2} x_{196}^{\rm new}\\
x_{203}^{\rm old}&=\frac{1}{2} x_{203}^{\rm new}+\frac{1}{2} x_{205}^{\rm new}\\
x_{204}^{\rm old}&=\frac{1}{2} x_{204}^{\rm new}+\frac{1}{2} x_{206}^{\rm new}\\
x_{205}^{\rm old}&=\frac{1}{2} x_{203}^{\rm new}-\frac{1}{2} x_{205}^{\rm new}\\
x_{206}^{\rm old}&=\frac{1}{2} x_{204}^{\rm new}-\frac{1}{2} x_{206}^{\rm new}\\
x_{207}^{\rm old}&=\frac{1}{2} x_{207}^{\rm new}+\frac{1}{2} x_{209}^{\rm new}\\
x_{208}^{\rm old}&=\frac{1}{2} x_{208}^{\rm new}+\frac{1}{2} x_{210}^{\rm new}\\
x_{209}^{\rm old}&=\frac{1}{2} x_{207}^{\rm new}-\frac{1}{2} x_{209}^{\rm new}\\
x_{210}^{\rm old}&=\frac{1}{2} x_{208}^{\rm new}-\frac{1}{2} x_{210}^{\rm new}
\end{align*}
Replace residuals 118,120 by their sum and difference, respectively.
Replace residuals 119,121 by their sum and difference, respectively.
Replace residuals 134,136 by their sum and difference, respectively.
Replace residuals 135,137 by their sum and difference, respectively.
Replace residuals 152,154 by their sum and difference, respectively.
Replace residuals 153,155 by their sum and difference, respectively.
Replace residuals 168,170 by their sum and difference, respectively.
Replace residuals 169,171 by their sum and difference, respectively.

\paragraph{Affine stage 8.}
\begin{align*}
x_{60}^{\rm old}&=1 +x_{60}^{\rm new}\\
x_{76}^{\rm old}&=2 +2 x_{60}^{\rm new}+x_{76}^{\rm new}\\
x_{77}^{\rm old}&=x_{61}^{\rm new}+x_{77}^{\rm new}
\end{align*}

Finally clear denominators and divide each row by the positive gcd of
its integer coefficients. These are nonzero constant row scalings.
Then replace row124 by $3Q_{124}+11Q_{125}$, leaving row125 fixed;
the inverse subtracts $11Q_{125}$ and divides by3. The common quotient
removes former row186, and all subsequent row indices decrease by one.
The explicit list below fixes every sign and normalization. All these
row operations are invertible on simultaneous zero sets.

\section{The explicit 222-variable quartic}\label{v222:equations}
Here $z_1,\ldots,z_{222}$ are the final coordinates and $A_N=149^{2N+1}$.
For the integer quadratic residuals below put
\[
 P_N(z)=\sum_{i=1}^{196}w_iQ_{i,N}(z)^2.
\]
Every weight is a positive integer. Unlisted weights are2; the exceptions are
\[
\begin{array}{c|l}w_i&i\\\hline
4&40,53,119,121,153,155,169,171\\
1&57,71,116,132,150,166\\
8&117,151,167\\
6&133,183\\
\end{array}
\]
The weights sum to428. Thus this is also a sum of428 repeated squares,
not196 unit-weight squares. The positive weights preserve the common zero set.

The following exact integer constants abbreviate long coefficients.
{\scriptsize
\begin{align*}\kappa_{1}={}&568705183978188214685858350706\cdot10^{199}\\
&+352595816600481322563715864164\cdot10^{169}\\
&+384761294213925665387551693298\cdot10^{139}\\
&+405347752531488609330673008441\cdot10^{109}\\
&+806110377207132493621682534382\cdot10^{79}\\
&+048931521003445777024432190975\cdot10^{49}\\
&+269141722690410774528000000000\cdot10^{19}\\
&+0000000000000000000\end{align*}
\begin{align*}\kappa_{2}={}&631894648864653571873175945229\cdot10^{198}\\
&+280662018444979247293017626849\cdot10^{168}\\
&+316401438015472961541724103664\cdot10^{138}\\
&+894830836146098454811858898268\cdot10^{108}\\
&+673455974674591659579647260424\cdot10^{78}\\
&+498812801114939752249369101083\cdot10^{48}\\
&+632379691878234193920000000000\cdot10^{18}\\
&+000000000000000000\end{align*}
\begin{align*}\kappa_{3}={}&275367677594445402666878321880\cdot10^{85}\\
&+397332775789666323829651983262\cdot10^{55}\\
&+964537509627072020640430936153\cdot10^{25}\\
&+3330600755200000000000000\end{align*}
\begin{align*}\kappa_{4}={}&137683838797222701333439160940\cdot10^{85}\\
&+198666387894833161914825991631\cdot10^{55}\\
&+482268754813536010320215468076\cdot10^{25}\\
&+6665300377600000000000000\end{align*}
\begin{align*}\kappa_{5}={}&917892258648151342222927739601\cdot10^{84}\\
&+324442585965554412765506610876\cdot10^{54}\\
&+548458365423573402134769787177\cdot10^{24}\\
&+776866918400000000000000\end{align*}
\begin{align*}\kappa_{6}={}&458946129324075671111463869800\cdot10^{84}\\
&+662221292982777206382753305438\cdot10^{54}\\
&+274229182711786701067384893588\cdot10^{24}\\
&+888433459200000000000000\end{align*}
\begin{align*}\kappa_{7}={}&809904934101310007843759770236\cdot10^{83}\\
&+462743458204900952440152891949\cdot10^{53}\\
&+895698557726682413648326282803\cdot10^{23}\\
&+92076492800000000000000\end{align*}
\begin{align*}\kappa_{8}={}&981949949521011646279509932288\cdot10^{75}\\
&+013517604094669913140033589964\cdot10^{45}\\
&+224332595426527668691070368153\cdot10^{15}\\
&+600000000000000\end{align*}
}
\allowdisplaybreaks

{\scriptsize\begin{equation}
\begin{aligned}
Q_{1,N}={}&+2401 z_{21} z_{22}-6369 z_{23}
\end{aligned}\tag{Q1}\label{v222:q1}
\end{equation}}

{\scriptsize\begin{equation}
\begin{aligned}
Q_{2,N}={}&+33 z_{21}^{2}-193 z_{21} z_{22}\\
&+264 z_{22}^{2}-33 z_{24}
\end{aligned}\tag{Q2}\label{v222:q2}
\end{equation}}

{\scriptsize\begin{equation}
\begin{aligned}
Q_{3,N}={}&+2401 z_{27} z_{28}-6369 z_{29}
\end{aligned}\tag{Q3}\label{v222:q3}
\end{equation}}

{\scriptsize\begin{equation}
\begin{aligned}
Q_{4,N}={}&+33 z_{27}^{2}-193 z_{27} z_{28}\\
&+264 z_{28}^{2}-33 z_{30}
\end{aligned}\tag{Q4}\label{v222:q4}
\end{equation}}

{\scriptsize\begin{equation}
\begin{aligned}
Q_{5,N}={}&+2401 z_{31} z_{32}-6369 z_{33}
\end{aligned}\tag{Q5}\label{v222:q5}
\end{equation}}

{\scriptsize\begin{equation}
\begin{aligned}
Q_{6,N}={}&+33 z_{31}^{2}-193 z_{31} z_{32}\\
&+264 z_{32}^{2}-33 z_{34}
\end{aligned}\tag{Q6}\label{v222:q6}
\end{equation}}

{\scriptsize\begin{equation}
\begin{aligned}
Q_{7,N}={}&+2401 z_{35} z_{36}-6369 z_{37}
\end{aligned}\tag{Q7}\label{v222:q7}
\end{equation}}

{\scriptsize\begin{equation}
\begin{aligned}
Q_{8,N}={}&+33 z_{35}^{2}-193 z_{35} z_{36}\\
&+264 z_{36}^{2}-33 z_{38}
\end{aligned}\tag{Q8}\label{v222:q8}
\end{equation}}

{\scriptsize\begin{equation}
\begin{aligned}
Q_{9,N}={}&+2501 +51 z_{1}^{2}\\
&-51 z_{55}
\end{aligned}\tag{Q9}\label{v222:q9}
\end{equation}}

{\scriptsize\begin{equation}
\begin{aligned}
Q_{10,N}={}&+z_{43} z_{51}+2 z_{44} z_{52}\\
&-2 z_{51} z_{120}-z_{56} z_{120}
\end{aligned}\tag{Q10}\label{v222:q10}
\end{equation}}

{\scriptsize\begin{equation}
\begin{aligned}
Q_{11,N}={}&+z_{43} z_{52}+z_{44} z_{51}\\
&-2 z_{52} z_{120}-z_{57} z_{120}
\end{aligned}\tag{Q11}\label{v222:q11}
\end{equation}}

{\scriptsize\begin{equation}
\begin{aligned}
Q_{12,N}={}&+z_{43} z_{47}+2 z_{44} z_{48}\\
&-z_{56} z_{120}-z_{58} z_{120}
\end{aligned}\tag{Q12}\label{v222:q12}
\end{equation}}

{\scriptsize\begin{equation}
\begin{aligned}
Q_{13,N}={}&+z_{43} z_{48}+z_{44} z_{47}\\
&-z_{57} z_{120}-z_{59} z_{120}
\end{aligned}\tag{Q13}\label{v222:q13}
\end{equation}}

{\scriptsize\begin{equation}
\begin{aligned}
Q_{14,N}={}&-1 +z_{47}^{2}\\
&+2 z_{48}^{2}-z_{56} z_{58}\\
&-2 z_{57} z_{59}-z_{60}
\end{aligned}\tag{Q14}\label{v222:q14}
\end{equation}}

{\scriptsize\begin{equation}
\begin{aligned}
Q_{15,N}={}&+2 z_{47} z_{48}-z_{56} z_{59}\\
&-z_{57} z_{58}-2 z_{61}
\end{aligned}\tag{Q15}\label{v222:q15}
\end{equation}}

{\scriptsize\begin{equation}
\begin{aligned}
Q_{16,N}={}&+z_{39}^{2}-2 z_{39} z_{41}\\
&+2 z_{40}^{2}+4 z_{40} z_{42}\\
&+z_{41}^{2}+2 z_{42}^{2}\\
&-z_{64}
\end{aligned}\tag{Q16}\label{v222:q16}
\end{equation}}

{\scriptsize\begin{equation}
\begin{aligned}
Q_{17,N}={}&+z_{39} z_{40}+z_{39} z_{42}\\
&-z_{40} z_{41}-z_{41} z_{42}\\
&-z_{65}
\end{aligned}\tag{Q17}\label{v222:q17}
\end{equation}}

{\scriptsize\begin{equation}
\begin{aligned}
Q_{18,N}={}&+9 +z_{39}^{2}\\
&-102 z_{39} z_{55}+2 z_{40}^{2}\\
&+2601 z_{55}^{2}-z_{66}
\end{aligned}\tag{Q18}\label{v222:q18}
\end{equation}}

{\scriptsize\begin{equation}
\begin{aligned}
Q_{19,N}={}&+z_{39} z_{40}-51 z_{40} z_{55}\\
&-z_{67}
\end{aligned}\tag{Q19}\label{v222:q19}
\end{equation}}

{\scriptsize\begin{equation}
\begin{aligned}
Q_{20,N}={}&+z_{17}^{2}+2 z_{18}^{2}\\
&-z_{68}
\end{aligned}\tag{Q20}\label{v222:q20}
\end{equation}}

{\scriptsize\begin{equation}
\begin{aligned}
Q_{21,N}={}&+z_{17} z_{18}-z_{69}
\end{aligned}\tag{Q21}\label{v222:q21}
\end{equation}}

{\scriptsize\begin{equation}
\begin{aligned}
Q_{22,N}={}&+z_{43}^{2}-2 z_{43} z_{45}\\
&+2 z_{44}^{2}+4 z_{44} z_{46}\\
&+z_{45}^{2}+2 z_{46}^{2}\\
&-z_{70} z_{119}
\end{aligned}\tag{Q22}\label{v222:q22}
\end{equation}}

{\scriptsize\begin{equation}
\begin{aligned}
Q_{23,N}={}&+z_{43} z_{44}+z_{43} z_{46}\\
&-z_{44} z_{45}-z_{45} z_{46}\\
&-z_{71} z_{119}
\end{aligned}\tag{Q23}\label{v222:q23}
\end{equation}}

{\scriptsize\begin{equation}
\begin{aligned}
Q_{24,N}={}&+z_{47}^{2}-2 z_{47} z_{49}\\
&+2 z_{48}^{2}+4 z_{48} z_{50}\\
&+z_{49}^{2}+2 z_{50}^{2}\\
&-z_{72}
\end{aligned}\tag{Q24}\label{v222:q24}
\end{equation}}

{\scriptsize\begin{equation}
\begin{aligned}
Q_{25,N}={}&+z_{47} z_{48}+z_{47} z_{50}\\
&-z_{48} z_{49}-z_{49} z_{50}\\
&-z_{73}
\end{aligned}\tag{Q25}\label{v222:q25}
\end{equation}}

{\scriptsize\begin{equation}
\begin{aligned}
Q_{26,N}={}&+z_{51}^{2}-2 z_{51} z_{53}\\
&+2 z_{52}^{2}+4 z_{52} z_{54}\\
&+z_{53}^{2}+2 z_{54}^{2}\\
&-z_{74}
\end{aligned}\tag{Q26}\label{v222:q26}
\end{equation}}

{\scriptsize\begin{equation}
\begin{aligned}
Q_{27,N}={}&+z_{51} z_{52}+z_{51} z_{54}\\
&-z_{52} z_{53}-z_{53} z_{54}\\
&-z_{75}
\end{aligned}\tag{Q27}\label{v222:q27}
\end{equation}}

{\scriptsize\begin{equation}
\begin{aligned}
Q_{28,N}={}&+z_{60}^{2}+8 z_{61}^{2}\\
&-z_{76}
\end{aligned}\tag{Q28}\label{v222:q28}
\end{equation}}

{\scriptsize\begin{equation}
\begin{aligned}
Q_{29,N}={}&+z_{60} z_{61}-z_{77}
\end{aligned}\tag{Q29}\label{v222:q29}
\end{equation}}

{\scriptsize\begin{equation}
\begin{aligned}
Q_{30,N}={}&+z_{4}^{2}-z_{78}
\end{aligned}\tag{Q30}\label{v222:q30}
\end{equation}}

{\scriptsize\begin{equation}
\begin{aligned}
Q_{31,N}={}&+z_{78}^{2}-z_{79}
\end{aligned}\tag{Q31}\label{v222:q31}
\end{equation}}

{\scriptsize\begin{equation}
\begin{aligned}
Q_{32,N}={}&+7626746323592 -76267463235920 z_{4}\\
&-628059618995128 z_{4} z_{78}+732473028362323 z_{4} z_{121}\\
&+307309842082151 z_{78}-583791423311868 z_{78} z_{121}\\
&+718967979842546 z_{79}+162308064100032 z_{79}^{2}\\
&-38129835686440 z_{79} z_{121}-7626746323592 z_{80}\\
&-3813373161796 z_{80} z_{81}-646146742577308 z_{121}\\
&+3813373161796 z_{121}^{2}
\end{aligned}\tag{Q32}\label{v222:q32}
\end{equation}}

{\scriptsize\begin{equation}
\begin{aligned}
Q_{33,N}={}&+1 -10 z_{4}\\
&-120 z_{4} z_{78}+210 z_{4} z_{121}\\
&+45 z_{78}-120 z_{78} z_{121}\\
&+210 z_{79}+45 z_{79}^{2}\\
&-10 z_{79} z_{121}-z_{80}\\
&-z_{80} z_{81}-252 z_{121}\\
&+z_{121}^{2}
\end{aligned}\tag{Q33}\label{v222:q33}
\end{equation}}

{\scriptsize\begin{equation}
\begin{aligned}
Q_{34,N}={}&+z_{81}^{2}-z_{83}
\end{aligned}\tag{Q34}\label{v222:q34}
\end{equation}}

{\scriptsize\begin{equation}
\begin{aligned}
Q_{35,N}={}&+z_{83}^{2}-z_{84}
\end{aligned}\tag{Q35}\label{v222:q35}
\end{equation}}

{\scriptsize\begin{equation}
\begin{aligned}
Q_{36,N}={}&+z_{85}^{2}-2 z_{85} z_{87}\\
&-2 z_{86}^{2}
\end{aligned}\tag{Q36}\label{v222:q36}
\end{equation}}

{\scriptsize\begin{equation}
\begin{aligned}
Q_{37,N}={}&+167281 z_{81} z_{88}-166464 z_{87}^{2}\\
&+167281 z_{88}
\end{aligned}\tag{Q37}\label{v222:q37}
\end{equation}}

{\scriptsize\begin{equation}
\begin{aligned}
Q_{38,N}={}&+z_{7}^{2}-z_{89}
\end{aligned}\tag{Q38}\label{v222:q38}
\end{equation}}

{\scriptsize\begin{equation}
\begin{aligned}
Q_{39,N}={}&+289 +34 z_{5}\\
&-3 z_{5}^{2}-17 z_{88} z_{89}
\end{aligned}\tag{Q39}\label{v222:q39}
\end{equation}}

{\scriptsize\begin{equation}
\begin{aligned}
Q_{40,N}={}&+17 z_{5}-z_{5} z_{88}\\
&-17 z_{6} z_{90}
\end{aligned}\tag{Q40}\label{v222:q40}
\end{equation}}

{\scriptsize\begin{equation}
\begin{aligned}
Q_{41,N}={}&+z_{8}^{2}-z_{91}
\end{aligned}\tag{Q41}\label{v222:q41}
\end{equation}}

{\scriptsize\begin{equation}
\begin{aligned}
Q_{42,N}={}&+z_{6}^{2}+34 z_{6} z_{90}\\
&-17 z_{88} z_{91}-867 z_{90}^{2}
\end{aligned}\tag{Q42}\label{v222:q42}
\end{equation}}

{\scriptsize\begin{equation}
\begin{aligned}
Q_{43,N}={}&+z_{92}^{2}-2 z_{92} z_{119}\\
&-z_{93} z_{119}+z_{119}^{2}
\end{aligned}\tag{Q43}\label{v222:q43}
\end{equation}}

{\scriptsize\begin{equation}
\begin{aligned}
Q_{44,N}={}&+z_{93}^{2}-z_{94} z_{119}
\end{aligned}\tag{Q44}\label{v222:q44}
\end{equation}}

{\scriptsize\begin{equation}
\begin{aligned}
Q_{45,N}={}&-141822723674536 z_{93} z_{119}+35893742479489 z_{93} z_{122}\\
&+192222345736537 z_{94} z_{119}-3813373161796 z_{95} z_{96}\\
&+39896049934224 z_{119}^{2}-122376041313918 z_{119} z_{122}
\end{aligned}\tag{Q45}\label{v222:q45}
\end{equation}}

{\scriptsize\begin{equation}
\begin{aligned}
Q_{46,N}={}&+z_{94} z_{122}-z_{95}\\
&-z_{95} z_{96}
\end{aligned}\tag{Q46}\label{v222:q46}
\end{equation}}

{\scriptsize\begin{equation}
\begin{aligned}
Q_{47,N}={}&+z_{96}^{2}-z_{98}
\end{aligned}\tag{Q47}\label{v222:q47}
\end{equation}}

{\scriptsize\begin{equation}
\begin{aligned}
Q_{48,N}={}&+z_{98}^{2}-z_{99}
\end{aligned}\tag{Q48}\label{v222:q48}
\end{equation}}

{\scriptsize\begin{equation}
\begin{aligned}
Q_{49,N}={}&+z_{100}^{2}-2 z_{100} z_{102}\\
&-2 z_{101}^{2}
\end{aligned}\tag{Q49}\label{v222:q49}
\end{equation}}

{\scriptsize\begin{equation}
\begin{aligned}
Q_{50,N}={}&+167281 z_{96} z_{103}-166464 z_{102}^{2}\\
&+167281 z_{103}
\end{aligned}\tag{Q50}\label{v222:q50}
\end{equation}}

{\scriptsize\begin{equation}
\begin{aligned}
Q_{51,N}={}&+z_{11}^{2}-z_{104}
\end{aligned}\tag{Q51}\label{v222:q51}
\end{equation}}

{\scriptsize\begin{equation}
\begin{aligned}
Q_{52,N}={}&+289 +34 z_{9}\\
&-3 z_{9}^{2}-17 z_{103} z_{104}
\end{aligned}\tag{Q52}\label{v222:q52}
\end{equation}}

{\scriptsize\begin{equation}
\begin{aligned}
Q_{53,N}={}&+17 z_{9}-z_{9} z_{103}\\
&-17 z_{10} z_{105}
\end{aligned}\tag{Q53}\label{v222:q53}
\end{equation}}

{\scriptsize\begin{equation}
\begin{aligned}
Q_{54,N}={}&+z_{12}^{2}-z_{106}
\end{aligned}\tag{Q54}\label{v222:q54}
\end{equation}}

{\scriptsize\begin{equation}
\begin{aligned}
Q_{55,N}={}&+z_{10}^{2}+34 z_{10} z_{105}\\
&-17 z_{103} z_{106}-867 z_{105}^{2}
\end{aligned}\tag{Q55}\label{v222:q55}
\end{equation}}

{\scriptsize\begin{equation}
\begin{aligned}
Q_{56,N}={}&+2601 z_{55}^{2}-4 z_{107}
\end{aligned}\tag{Q56}\label{v222:q56}
\end{equation}}

{\scriptsize\begin{equation}
\begin{aligned}
Q_{57,N}={}&+6316332 +291790992 z_{55}\\
&+15531948 z_{55} z_{107}+5736060 z_{107}\\
&-512 z_{107}^{2}+z_{108}
\end{aligned}\tag{Q57}\label{v222:q57}
\end{equation}}

{\scriptsize\begin{equation}
\begin{aligned}
Q_{58,N}={}&+\kappa_{8} z_{111} z_{119}\\
&-z_{118}^{2}
\end{aligned}\tag{Q58}\label{v222:q58}
\end{equation}}

{\scriptsize\begin{equation}
\begin{aligned}
Q_{59,N}={}&+9792 z_{111}+z_{111} z_{120}\\
&+82479248 z_{112} z_{120}
\end{aligned}\tag{Q59}\label{v222:q59}
\end{equation}}

{\scriptsize\begin{equation}
\begin{aligned}
Q_{60,N}={}&+z_{15}^{2}-z_{113}
\end{aligned}\tag{Q60}\label{v222:q60}
\end{equation}}

{\scriptsize\begin{equation}
\begin{aligned}
Q_{61,N}={}&+\kappa_{2} \\
&-\kappa_{5} z_{13}\\
&-z_{13}^{2}\\
&+\kappa_{6} z_{112} z_{113}
\end{aligned}\tag{Q61}\label{v222:q61}
\end{equation}}

{\scriptsize\begin{equation}
\begin{aligned}
Q_{62,N}={}&+17 z_{13}\\
&+\kappa_{7} z_{13} z_{112}\\
&-17 z_{14} z_{114}
\end{aligned}\tag{Q62}\label{v222:q62}
\end{equation}}

{\scriptsize\begin{equation}
\begin{aligned}
Q_{63,N}={}&+z_{16}^{2}-z_{115}
\end{aligned}\tag{Q63}\label{v222:q63}
\end{equation}}

{\scriptsize\begin{equation}
\begin{aligned}
Q_{64,N}={}&+z_{14}^{2}\\
&-\kappa_{3} z_{14} z_{114}\\
&+\kappa_{4} z_{112} z_{115}\\
&-\kappa_{1} z_{114}^{2}
\end{aligned}\tag{Q64}\label{v222:q64}
\end{equation}}

{\scriptsize\begin{equation}
\begin{aligned}
Q_{65,N}={}&+z_{68}^{2}-8 z_{69}^{2}\\
&+z_{116} z_{120}+z_{117}^{2}\\
&-z_{118} z_{119}
\end{aligned}\tag{Q65}\label{v222:q65}
\end{equation}}

{\scriptsize\begin{equation}
\begin{aligned}
Q_{66,N}={}&+z_{81} z_{82}-z_{84}^{2}
\end{aligned}\tag{Q66}\label{v222:q66}
\end{equation}}

{\scriptsize\begin{equation}
\begin{aligned}
Q_{67,N}={}&+z_{96} z_{97}-z_{99}^{2}
\end{aligned}\tag{Q67}\label{v222:q67}
\end{equation}}

{\scriptsize\begin{equation}
\begin{aligned}
Q_{68,N}={}&+z_{119}-z_{120}^{2}
\end{aligned}\tag{Q68}\label{v222:q68}
\end{equation}}

{\scriptsize\begin{equation}
\begin{aligned}
Q_{69,N}={}&+z_{110}^{2}-z_{116}
\end{aligned}\tag{Q69}\label{v222:q69}
\end{equation}}

{\scriptsize\begin{equation}
\begin{aligned}
Q_{70,N}={}&+z_{109}^{2}-z_{117} z_{120}
\end{aligned}\tag{Q70}\label{v222:q70}
\end{equation}}

{\scriptsize\begin{equation}
\begin{aligned}
Q_{71,N}={}&+22201 z_{108} z_{120}+z_{120}\\
&-A_N z_{108}
\end{aligned}\tag{Q71}\label{v222:q71}
\end{equation}}

{\scriptsize\begin{equation}
\begin{aligned}
Q_{72,N}={}&+z_{4} z_{79}-z_{121}
\end{aligned}\tag{Q72}\label{v222:q72}
\end{equation}}

{\scriptsize\begin{equation}
\begin{aligned}
Q_{73,N}={}&+z_{93} z_{94}-z_{119} z_{122}
\end{aligned}\tag{Q73}\label{v222:q73}
\end{equation}}

{\scriptsize\begin{equation}
\begin{aligned}
Q_{74,N}={}&+z_{64}^{2}+8 z_{65}^{2}\\
&-z_{123}
\end{aligned}\tag{Q74}\label{v222:q74}
\end{equation}}

{\scriptsize\begin{equation}
\begin{aligned}
Q_{75,N}={}&+4 z_{64} z_{65}-z_{124}
\end{aligned}\tag{Q75}\label{v222:q75}
\end{equation}}

{\scriptsize\begin{equation}
\begin{aligned}
Q_{76,N}={}&-324 -12 z_{39} z_{66}\\
&-48 z_{40} z_{67}+612 z_{55} z_{66}\\
&+36 z_{66}+z_{66}^{2}\\
&+8 z_{67}^{2}-z_{125}
\end{aligned}\tag{Q76}\label{v222:q76}
\end{equation}}

{\scriptsize\begin{equation}
\begin{aligned}
Q_{77,N}={}&-24 z_{39} z_{67}-12 z_{40} z_{66}\\
&+1224 z_{55} z_{67}+4 z_{66} z_{67}\\
&+72 z_{67}-z_{126}
\end{aligned}\tag{Q77}\label{v222:q77}
\end{equation}}

{\scriptsize\begin{equation}
\begin{aligned}
Q_{78,N}={}&+z_{25}^{2}+8 z_{26}^{2}\\
&-z_{127}
\end{aligned}\tag{Q78}\label{v222:q78}
\end{equation}}

{\scriptsize\begin{equation}
\begin{aligned}
Q_{79,N}={}&+4 z_{25} z_{26}+z_{128}
\end{aligned}\tag{Q79}\label{v222:q79}
\end{equation}}

{\scriptsize\begin{equation}
\begin{aligned}
Q_{80,N}={}&+1 +z_{68} z_{83}\\
&+z_{70} z_{82}+z_{84} z_{127}\\
&-z_{129}
\end{aligned}\tag{Q80}\label{v222:q80}
\end{equation}}

{\scriptsize\begin{equation}
\begin{aligned}
Q_{81,N}={}&+2 z_{69} z_{83}+2 z_{71} z_{82}\\
&+z_{84} z_{128}-z_{130}
\end{aligned}\tag{Q81}\label{v222:q81}
\end{equation}}

{\scriptsize\begin{equation}
\begin{aligned}
Q_{82,N}={}&-37249 z_{23} z_{220}-2401 z_{24} z_{220}\\
&-2401 z_{76} z_{219}+2401 z_{83}\\
&-2401 z_{83} z_{147}-2401 z_{98}\\
&+2401 z_{98} z_{131}
\end{aligned}\tag{Q82}\label{v222:q82}
\end{equation}}

{\scriptsize\begin{equation}
\begin{aligned}
Q_{83,N}={}&+25476 z_{23} z_{220}-9604 z_{77} z_{219}\\
&-2401 z_{83} z_{148}+2401 z_{98} z_{132}
\end{aligned}\tag{Q83}\label{v222:q83}
\end{equation}}

{\scriptsize\begin{equation}
\begin{aligned}
Q_{84,N}={}&-37249 z_{33} z_{220}-2401 z_{34} z_{220}\\
&-2401 z_{62} z_{221}+2401 z_{83}\\
&-2401 z_{83} z_{149}-2401 z_{98}\\
&+2401 z_{98} z_{133}-2401 z_{125} z_{219}
\end{aligned}\tag{Q84}\label{v222:q84}
\end{equation}}

{\scriptsize\begin{equation}
\begin{aligned}
Q_{85,N}={}&+25476 z_{33} z_{220}-4802 z_{63} z_{221}\\
&-2401 z_{83} z_{150}+2401 z_{98} z_{134}\\
&-2401 z_{126} z_{219}
\end{aligned}\tag{Q85}\label{v222:q85}
\end{equation}}

{\scriptsize\begin{equation}
\begin{aligned}
Q_{86,N}={}&+2401 +37249 z_{37} z_{84}\\
&+2401 z_{38} z_{84}+2401 z_{72} z_{82}\\
&+2401 z_{83} z_{123}-2401 z_{135}
\end{aligned}\tag{Q86}\label{v222:q86}
\end{equation}}

{\scriptsize\begin{equation}
\begin{aligned}
Q_{87,N}={}&+25476 z_{37} z_{84}-4802 z_{73} z_{82}\\
&-2401 z_{83} z_{124}+2401 z_{136}
\end{aligned}\tag{Q87}\label{v222:q87}
\end{equation}}

{\scriptsize\begin{equation}
\begin{aligned}
Q_{88,N}={}&-37249 z_{29} z_{220}-2401 z_{30} z_{220}\\
&-2401 z_{74} z_{219}+2401 z_{83}\\
&-2401 z_{83} z_{153}-2401 z_{98}\\
&+2401 z_{98} z_{137}
\end{aligned}\tag{Q88}\label{v222:q88}
\end{equation}}

{\scriptsize\begin{equation}
\begin{aligned}
Q_{89,N}={}&+25476 z_{29} z_{220}-4802 z_{75} z_{219}\\
&-2401 z_{83} z_{154}+2401 z_{98} z_{138}
\end{aligned}\tag{Q89}\label{v222:q89}
\end{equation}}

{\scriptsize\begin{equation}
\begin{aligned}
Q_{90,N}={}&+z_{129} z_{131}+2 z_{130} z_{132}\\
&-z_{139}
\end{aligned}\tag{Q90}\label{v222:q90}
\end{equation}}

{\scriptsize\begin{equation}
\begin{aligned}
Q_{91,N}={}&+z_{129} z_{132}+z_{130} z_{131}\\
&-z_{140}
\end{aligned}\tag{Q91}\label{v222:q91}
\end{equation}}

{\scriptsize\begin{equation}
\begin{aligned}
Q_{92,N}={}&+z_{133} z_{139}+2 z_{134} z_{140}\\
&-z_{141}
\end{aligned}\tag{Q92}\label{v222:q92}
\end{equation}}

{\scriptsize\begin{equation}
\begin{aligned}
Q_{93,N}={}&+z_{133} z_{140}+z_{134} z_{139}\\
&-z_{142}
\end{aligned}\tag{Q93}\label{v222:q93}
\end{equation}}

{\scriptsize\begin{equation}
\begin{aligned}
Q_{94,N}={}&+z_{135} z_{141}+2 z_{136} z_{142}\\
&-z_{143}
\end{aligned}\tag{Q94}\label{v222:q94}
\end{equation}}

{\scriptsize\begin{equation}
\begin{aligned}
Q_{95,N}={}&+z_{135} z_{142}+z_{136} z_{141}\\
&-z_{144}
\end{aligned}\tag{Q95}\label{v222:q95}
\end{equation}}

{\scriptsize\begin{equation}
\begin{aligned}
Q_{96,N}={}&+z_{85}-z_{137} z_{143}\\
&-2 z_{138} z_{144}
\end{aligned}\tag{Q96}\label{v222:q96}
\end{equation}}

{\scriptsize\begin{equation}
\begin{aligned}
Q_{97,N}={}&+z_{86}-z_{137} z_{144}\\
&-z_{138} z_{143}
\end{aligned}\tag{Q97}\label{v222:q97}
\end{equation}}

{\scriptsize\begin{equation}
\begin{aligned}
Q_{98,N}={}&+1 +z_{68} z_{98}\\
&+z_{70} z_{97}+z_{99} z_{127}\\
&-z_{145}
\end{aligned}\tag{Q98}\label{v222:q98}
\end{equation}}

{\scriptsize\begin{equation}
\begin{aligned}
Q_{99,N}={}&+2 z_{69} z_{98}+2 z_{71} z_{97}\\
&+z_{99} z_{128}-z_{146}
\end{aligned}\tag{Q99}\label{v222:q99}
\end{equation}}

{\scriptsize\begin{equation}
\begin{aligned}
Q_{100,N}={}&-2401 -37249 z_{23} z_{99}\\
&-2401 z_{24} z_{99}-2401 z_{76} z_{97}\\
&+2401 z_{147}-2401 z_{201}^{2}\\
&-4802 z_{202}^{2}
\end{aligned}\tag{Q100}\label{v222:q100}
\end{equation}}

{\scriptsize\begin{equation}
\begin{aligned}
Q_{101,N}={}&+25476 z_{23} z_{99}-9604 z_{77} z_{97}\\
&+2401 z_{148}-4802 z_{201} z_{202}
\end{aligned}\tag{Q101}\label{v222:q101}
\end{equation}}

{\scriptsize\begin{equation}
\begin{aligned}
Q_{102,N}={}&-2401 -37249 z_{33} z_{99}\\
&-2401 z_{34} z_{99}-2401 z_{62} z_{96}\\
&-2401 z_{97} z_{125}+2401 z_{149}\\
&-2401 z_{215}^{2}-4802 z_{216}^{2}
\end{aligned}\tag{Q102}\label{v222:q102}
\end{equation}}

{\scriptsize\begin{equation}
\begin{aligned}
Q_{103,N}={}&+25476 z_{33} z_{99}-4802 z_{63} z_{96}\\
&-2401 z_{97} z_{126}+2401 z_{150}\\
&-4802 z_{215} z_{216}
\end{aligned}\tag{Q103}\label{v222:q103}
\end{equation}}

{\scriptsize\begin{equation}
\begin{aligned}
Q_{104,N}={}&+2401 +37249 z_{37} z_{99}\\
&+2401 z_{38} z_{99}+2401 z_{72} z_{97}\\
&+2401 z_{98} z_{123}-2401 z_{151}
\end{aligned}\tag{Q104}\label{v222:q104}
\end{equation}}

{\scriptsize\begin{equation}
\begin{aligned}
Q_{105,N}={}&+25476 z_{37} z_{99}-4802 z_{73} z_{97}\\
&-2401 z_{98} z_{124}+2401 z_{152}
\end{aligned}\tag{Q105}\label{v222:q105}
\end{equation}}

{\scriptsize\begin{equation}
\begin{aligned}
Q_{106,N}={}&-2401 -37249 z_{29} z_{99}\\
&-2401 z_{30} z_{99}-2401 z_{74} z_{97}\\
&+2401 z_{153}-2401 z_{173}^{2}\\
&-4802 z_{174}^{2}
\end{aligned}\tag{Q106}\label{v222:q106}
\end{equation}}

{\scriptsize\begin{equation}
\begin{aligned}
Q_{107,N}={}&+25476 z_{29} z_{99}-4802 z_{75} z_{97}\\
&+2401 z_{154}-4802 z_{173} z_{174}
\end{aligned}\tag{Q107}\label{v222:q107}
\end{equation}}

{\scriptsize\begin{equation}
\begin{aligned}
Q_{108,N}={}&+z_{145} z_{147}+2 z_{146} z_{148}\\
&-z_{155}
\end{aligned}\tag{Q108}\label{v222:q108}
\end{equation}}

{\scriptsize\begin{equation}
\begin{aligned}
Q_{109,N}={}&+z_{145} z_{148}+z_{146} z_{147}\\
&-z_{156}
\end{aligned}\tag{Q109}\label{v222:q109}
\end{equation}}

{\scriptsize\begin{equation}
\begin{aligned}
Q_{110,N}={}&+z_{149} z_{155}+2 z_{150} z_{156}\\
&-z_{157}
\end{aligned}\tag{Q110}\label{v222:q110}
\end{equation}}

{\scriptsize\begin{equation}
\begin{aligned}
Q_{111,N}={}&+z_{149} z_{156}+z_{150} z_{155}\\
&-z_{158}
\end{aligned}\tag{Q111}\label{v222:q111}
\end{equation}}

{\scriptsize\begin{equation}
\begin{aligned}
Q_{112,N}={}&+z_{151} z_{157}+2 z_{152} z_{158}\\
&-z_{159}
\end{aligned}\tag{Q112}\label{v222:q112}
\end{equation}}

{\scriptsize\begin{equation}
\begin{aligned}
Q_{113,N}={}&+z_{151} z_{158}+z_{152} z_{157}\\
&-z_{160}
\end{aligned}\tag{Q113}\label{v222:q113}
\end{equation}}

{\scriptsize\begin{equation}
\begin{aligned}
Q_{114,N}={}&+z_{100}-z_{153} z_{159}\\
&-2 z_{154} z_{160}
\end{aligned}\tag{Q114}\label{v222:q114}
\end{equation}}

{\scriptsize\begin{equation}
\begin{aligned}
Q_{115,N}={}&+z_{101}-z_{153} z_{160}\\
&-z_{154} z_{159}
\end{aligned}\tag{Q115}\label{v222:q115}
\end{equation}}

{\scriptsize\begin{equation}
\begin{aligned}
Q_{116,N}={}&+z_{161}^{2}+2 z_{162}^{2}\\
&+z_{163}^{2}+2 z_{164}^{2}\\
&+z_{165}^{2}+2 z_{166}^{2}\\
&-z_{167}^{2}-2 z_{168}^{2}
\end{aligned}\tag{Q116}\label{v222:q116}
\end{equation}}

{\scriptsize\begin{equation}
\begin{aligned}
Q_{117,N}={}&+z_{161} z_{162}+z_{163} z_{164}\\
&+z_{165} z_{166}-z_{167} z_{168}
\end{aligned}\tag{Q117}\label{v222:q117}
\end{equation}}

{\scriptsize\begin{equation}
\begin{aligned}
Q_{118,N}={}&-12 z_{96} z_{169}+8 z_{96} z_{170}\\
&+z_{161} z_{165}+2 z_{162} z_{166}\\
&+z_{163} z_{167}+2 z_{164} z_{168}
\end{aligned}\tag{Q118}\label{v222:q118}
\end{equation}}

{\scriptsize\begin{equation}
\begin{aligned}
Q_{119,N}={}&+4 z_{96} z_{169}-12 z_{96} z_{170}\\
&+z_{161} z_{166}+z_{162} z_{165}\\
&+z_{163} z_{168}+z_{164} z_{167}
\end{aligned}\tag{Q119}\label{v222:q119}
\end{equation}}

{\scriptsize\begin{equation}
\begin{aligned}
Q_{120,N}={}&+12 z_{96} z_{171}-8 z_{96} z_{172}\\
&+z_{161} z_{167}+2 z_{162} z_{168}\\
&+z_{163} z_{165}+2 z_{164} z_{166}
\end{aligned}\tag{Q120}\label{v222:q120}
\end{equation}}

{\scriptsize\begin{equation}
\begin{aligned}
Q_{121,N}={}&-4 z_{96} z_{171}+12 z_{96} z_{172}\\
&+z_{161} z_{168}+z_{162} z_{167}\\
&+z_{163} z_{166}+z_{164} z_{165}
\end{aligned}\tag{Q121}\label{v222:q121}
\end{equation}}

{\scriptsize\begin{equation}
\begin{aligned}
Q_{122,N}={}&-44 z_{96} z_{173}+48 z_{96} z_{174}\\
&+z_{161}^{2}+2 z_{162}^{2}
\end{aligned}\tag{Q122}\label{v222:q122}
\end{equation}}

{\scriptsize\begin{equation}
\begin{aligned}
Q_{123,N}={}&+12 z_{96} z_{173}-22 z_{96} z_{174}\\
&+z_{161} z_{162}
\end{aligned}\tag{Q123}\label{v222:q123}
\end{equation}}

{\scriptsize\begin{equation}
\begin{aligned}
Q_{124,N}={}&-3 z_{98}+3 z_{169}^{2}\\
&+11 z_{169} z_{170}+6 z_{170}^{2}\\
&+49 z_{173} z_{174}
\end{aligned}\tag{Q124}\label{v222:q124}
\end{equation}}

{\scriptsize\begin{equation}
\begin{aligned}
Q_{125,N}={}&+z_{169} z_{170}-3 z_{173}^{2}\\
&+11 z_{173} z_{174}-6 z_{174}^{2}
\end{aligned}\tag{Q125}\label{v222:q125}
\end{equation}}

{\scriptsize\begin{equation}
\begin{aligned}
Q_{126,N}={}&+11 z_{25} z_{98}-24 z_{26} z_{98}\\
&-z_{98}+z_{171}^{2}\\
&+2 z_{172}^{2}
\end{aligned}\tag{Q126}\label{v222:q126}
\end{equation}}

{\scriptsize\begin{equation}
\begin{aligned}
Q_{127,N}={}&-3 z_{25} z_{98}+11 z_{26} z_{98}\\
&+z_{171} z_{172}
\end{aligned}\tag{Q127}\label{v222:q127}
\end{equation}}

{\scriptsize\begin{equation}
\begin{aligned}
Q_{128,N}={}&-6 z_{41} z_{171}+4 z_{41} z_{172}\\
&+4 z_{42} z_{171}-12 z_{42} z_{172}\\
&+z_{163} z_{217}+2 z_{164} z_{218}
\end{aligned}\tag{Q128}\label{v222:q128}
\end{equation}}

{\scriptsize\begin{equation}
\begin{aligned}
Q_{129,N}={}&+2 z_{41} z_{171}-6 z_{41} z_{172}\\
&-6 z_{42} z_{171}+4 z_{42} z_{172}\\
&+z_{163} z_{218}+z_{164} z_{217}
\end{aligned}\tag{Q129}\label{v222:q129}
\end{equation}}

{\scriptsize\begin{equation}
\begin{aligned}
Q_{130,N}={}&+z_{19} z_{161}+2 z_{20} z_{162}\\
&-6 z_{39} z_{169}+4 z_{39} z_{170}\\
&+4 z_{40} z_{169}-12 z_{40} z_{170}
\end{aligned}\tag{Q130}\label{v222:q130}
\end{equation}}

{\scriptsize\begin{equation}
\begin{aligned}
Q_{131,N}={}&+z_{19} z_{162}+z_{20} z_{161}\\
&+2 z_{39} z_{169}-6 z_{39} z_{170}\\
&-6 z_{40} z_{169}+4 z_{40} z_{170}
\end{aligned}\tag{Q131}\label{v222:q131}
\end{equation}}

{\scriptsize\begin{equation}
\begin{aligned}
Q_{132,N}={}&+z_{175}^{2}+2 z_{176}^{2}\\
&+z_{177}^{2}+2 z_{178}^{2}\\
&+z_{179}^{2}+2 z_{180}^{2}\\
&-z_{181}^{2}-2 z_{182}^{2}
\end{aligned}\tag{Q132}\label{v222:q132}
\end{equation}}

{\scriptsize\begin{equation}
\begin{aligned}
Q_{133,N}={}&+z_{175} z_{176}+z_{177} z_{178}\\
&+z_{179} z_{180}-z_{181} z_{182}
\end{aligned}\tag{Q133}\label{v222:q133}
\end{equation}}

{\scriptsize\begin{equation}
\begin{aligned}
Q_{134,N}={}&-12 z_{119} z_{183}+8 z_{119} z_{184}\\
&+z_{175} z_{179}+2 z_{176} z_{180}\\
&+z_{177} z_{181}+2 z_{178} z_{182}
\end{aligned}\tag{Q134}\label{v222:q134}
\end{equation}}

{\scriptsize\begin{equation}
\begin{aligned}
Q_{135,N}={}&+4 z_{119} z_{183}-12 z_{119} z_{184}\\
&+z_{175} z_{180}+z_{176} z_{179}\\
&+z_{177} z_{182}+z_{178} z_{181}
\end{aligned}\tag{Q135}\label{v222:q135}
\end{equation}}

{\scriptsize\begin{equation}
\begin{aligned}
Q_{136,N}={}&+12 z_{119} z_{185}-8 z_{119} z_{186}\\
&+z_{175} z_{181}+2 z_{176} z_{182}\\
&+z_{177} z_{179}+2 z_{178} z_{180}
\end{aligned}\tag{Q136}\label{v222:q136}
\end{equation}}

{\scriptsize\begin{equation}
\begin{aligned}
Q_{137,N}={}&-4 z_{119} z_{185}+12 z_{119} z_{186}\\
&+z_{175} z_{182}+z_{176} z_{181}\\
&+z_{177} z_{180}+z_{178} z_{179}
\end{aligned}\tag{Q137}\label{v222:q137}
\end{equation}}

{\scriptsize\begin{equation}
\begin{aligned}
Q_{138,N}={}&-44 z_{119} z_{187}+48 z_{119} z_{188}\\
&+z_{175}^{2}+2 z_{176}^{2}
\end{aligned}\tag{Q138}\label{v222:q138}
\end{equation}}

{\scriptsize\begin{equation}
\begin{aligned}
Q_{139,N}={}&+12 z_{119} z_{187}-22 z_{119} z_{188}\\
&+z_{175} z_{176}
\end{aligned}\tag{Q139}\label{v222:q139}
\end{equation}}

{\scriptsize\begin{equation}
\begin{aligned}
Q_{140,N}={}&+z_{62}-z_{187}^{2}\\
&-2 z_{188}^{2}
\end{aligned}\tag{Q140}\label{v222:q140}
\end{equation}}

{\scriptsize\begin{equation}
\begin{aligned}
Q_{141,N}={}&+z_{63}-z_{187} z_{188}
\end{aligned}\tag{Q141}\label{v222:q141}
\end{equation}}

{\scriptsize\begin{equation}
\begin{aligned}
Q_{142,N}={}&-1 +11 z_{62}\\
&-24 z_{63}+z_{183}^{2}\\
&+2 z_{184}^{2}
\end{aligned}\tag{Q142}\label{v222:q142}
\end{equation}}

{\scriptsize\begin{equation}
\begin{aligned}
Q_{143,N}={}&-3 z_{62}+11 z_{63}\\
&+z_{183} z_{184}
\end{aligned}\tag{Q143}\label{v222:q143}
\end{equation}}

{\scriptsize\begin{equation}
\begin{aligned}
Q_{144,N}={}&-1 +11 z_{27}\\
&-24 z_{28}+z_{185}^{2}\\
&+2 z_{186}^{2}
\end{aligned}\tag{Q144}\label{v222:q144}
\end{equation}}

{\scriptsize\begin{equation}
\begin{aligned}
Q_{145,N}={}&-3 z_{27}+11 z_{28}\\
&+z_{185} z_{186}
\end{aligned}\tag{Q145}\label{v222:q145}
\end{equation}}

{\scriptsize\begin{equation}
\begin{aligned}
Q_{146,N}={}&-6 z_{45} z_{185}+4 z_{45} z_{186}\\
&+4 z_{46} z_{185}-12 z_{46} z_{186}\\
&+z_{177} z_{217}+2 z_{178} z_{218}
\end{aligned}\tag{Q146}\label{v222:q146}
\end{equation}}

{\scriptsize\begin{equation}
\begin{aligned}
Q_{147,N}={}&+2 z_{45} z_{185}-6 z_{45} z_{186}\\
&-6 z_{46} z_{185}+4 z_{46} z_{186}\\
&+z_{177} z_{218}+z_{178} z_{217}
\end{aligned}\tag{Q147}\label{v222:q147}
\end{equation}}

{\scriptsize\begin{equation}
\begin{aligned}
Q_{148,N}={}&+z_{19} z_{175}+2 z_{20} z_{176}\\
&-6 z_{43} z_{183}+4 z_{43} z_{184}\\
&+4 z_{44} z_{183}-12 z_{44} z_{184}
\end{aligned}\tag{Q148}\label{v222:q148}
\end{equation}}

{\scriptsize\begin{equation}
\begin{aligned}
Q_{149,N}={}&+z_{19} z_{176}+z_{20} z_{175}\\
&+2 z_{43} z_{183}-6 z_{43} z_{184}\\
&-6 z_{44} z_{183}+4 z_{44} z_{184}
\end{aligned}\tag{Q149}\label{v222:q149}
\end{equation}}

{\scriptsize\begin{equation}
\begin{aligned}
Q_{150,N}={}&+z_{189}^{2}+2 z_{190}^{2}\\
&+z_{191}^{2}+2 z_{192}^{2}\\
&+z_{193}^{2}+2 z_{194}^{2}\\
&-z_{195}^{2}-2 z_{196}^{2}
\end{aligned}\tag{Q150}\label{v222:q150}
\end{equation}}

{\scriptsize\begin{equation}
\begin{aligned}
Q_{151,N}={}&+z_{189} z_{190}+z_{191} z_{192}\\
&+z_{193} z_{194}-z_{195} z_{196}
\end{aligned}\tag{Q151}\label{v222:q151}
\end{equation}}

{\scriptsize\begin{equation}
\begin{aligned}
Q_{152,N}={}&-12 z_{96} z_{197}+8 z_{96} z_{198}\\
&+z_{189} z_{193}+2 z_{190} z_{194}\\
&+z_{191} z_{195}+2 z_{192} z_{196}
\end{aligned}\tag{Q152}\label{v222:q152}
\end{equation}}

{\scriptsize\begin{equation}
\begin{aligned}
Q_{153,N}={}&+4 z_{96} z_{197}-12 z_{96} z_{198}\\
&+z_{189} z_{194}+z_{190} z_{193}\\
&+z_{191} z_{196}+z_{192} z_{195}
\end{aligned}\tag{Q153}\label{v222:q153}
\end{equation}}

{\scriptsize\begin{equation}
\begin{aligned}
Q_{154,N}={}&+12 z_{96} z_{199}-8 z_{96} z_{200}\\
&+z_{189} z_{195}+2 z_{190} z_{196}\\
&+z_{191} z_{193}+2 z_{192} z_{194}
\end{aligned}\tag{Q154}\label{v222:q154}
\end{equation}}

{\scriptsize\begin{equation}
\begin{aligned}
Q_{155,N}={}&-4 z_{96} z_{199}+12 z_{96} z_{200}\\
&+z_{189} z_{196}+z_{190} z_{195}\\
&+z_{191} z_{194}+z_{192} z_{193}
\end{aligned}\tag{Q155}\label{v222:q155}
\end{equation}}

{\scriptsize\begin{equation}
\begin{aligned}
Q_{156,N}={}&-44 z_{96} z_{201}+48 z_{96} z_{202}\\
&+z_{189}^{2}+2 z_{190}^{2}
\end{aligned}\tag{Q156}\label{v222:q156}
\end{equation}}

{\scriptsize\begin{equation}
\begin{aligned}
Q_{157,N}={}&+12 z_{96} z_{201}-22 z_{96} z_{202}\\
&+z_{189} z_{190}
\end{aligned}\tag{Q157}\label{v222:q157}
\end{equation}}

{\scriptsize\begin{equation}
\begin{aligned}
Q_{158,N}={}&-z_{98}+z_{197}^{2}\\
&+2 z_{198}^{2}+11 z_{201}^{2}\\
&-24 z_{201} z_{202}+22 z_{202}^{2}
\end{aligned}\tag{Q158}\label{v222:q158}
\end{equation}}

{\scriptsize\begin{equation}
\begin{aligned}
Q_{159,N}={}&+z_{197} z_{198}-3 z_{201}^{2}\\
&+11 z_{201} z_{202}-6 z_{202}^{2}
\end{aligned}\tag{Q159}\label{v222:q159}
\end{equation}}

{\scriptsize\begin{equation}
\begin{aligned}
Q_{160,N}={}&+11 z_{31} z_{98}-24 z_{32} z_{98}\\
&-z_{98}+z_{199}^{2}\\
&+2 z_{200}^{2}
\end{aligned}\tag{Q160}\label{v222:q160}
\end{equation}}

{\scriptsize\begin{equation}
\begin{aligned}
Q_{161,N}={}&-3 z_{31} z_{98}+11 z_{32} z_{98}\\
&+z_{199} z_{200}
\end{aligned}\tag{Q161}\label{v222:q161}
\end{equation}}

{\scriptsize\begin{equation}
\begin{aligned}
Q_{162,N}={}&-6 z_{49} z_{199}+4 z_{49} z_{200}\\
&+4 z_{50} z_{199}-12 z_{50} z_{200}\\
&+z_{191} z_{217}+2 z_{192} z_{218}
\end{aligned}\tag{Q162}\label{v222:q162}
\end{equation}}

{\scriptsize\begin{equation}
\begin{aligned}
Q_{163,N}={}&+2 z_{49} z_{199}-6 z_{49} z_{200}\\
&-6 z_{50} z_{199}+4 z_{50} z_{200}\\
&+z_{191} z_{218}+z_{192} z_{217}
\end{aligned}\tag{Q163}\label{v222:q163}
\end{equation}}

{\scriptsize\begin{equation}
\begin{aligned}
Q_{164,N}={}&+z_{19} z_{189}+2 z_{20} z_{190}\\
&-6 z_{47} z_{197}+4 z_{47} z_{198}\\
&+4 z_{48} z_{197}-12 z_{48} z_{198}
\end{aligned}\tag{Q164}\label{v222:q164}
\end{equation}}

{\scriptsize\begin{equation}
\begin{aligned}
Q_{165,N}={}&+z_{19} z_{190}+z_{20} z_{189}\\
&+2 z_{47} z_{197}-6 z_{47} z_{198}\\
&-6 z_{48} z_{197}+4 z_{48} z_{198}
\end{aligned}\tag{Q165}\label{v222:q165}
\end{equation}}

{\scriptsize\begin{equation}
\begin{aligned}
Q_{166,N}={}&+z_{203}^{2}+2 z_{204}^{2}\\
&+z_{205}^{2}+2 z_{206}^{2}\\
&+z_{207}^{2}+2 z_{208}^{2}\\
&-z_{209}^{2}-2 z_{210}^{2}
\end{aligned}\tag{Q166}\label{v222:q166}
\end{equation}}

{\scriptsize\begin{equation}
\begin{aligned}
Q_{167,N}={}&+z_{203} z_{204}+z_{205} z_{206}\\
&+z_{207} z_{208}-z_{209} z_{210}
\end{aligned}\tag{Q167}\label{v222:q167}
\end{equation}}

{\scriptsize\begin{equation}
\begin{aligned}
Q_{168,N}={}&-12 z_{96} z_{211}+8 z_{96} z_{212}\\
&+z_{203} z_{207}+2 z_{204} z_{208}\\
&+z_{205} z_{209}+2 z_{206} z_{210}
\end{aligned}\tag{Q168}\label{v222:q168}
\end{equation}}

{\scriptsize\begin{equation}
\begin{aligned}
Q_{169,N}={}&+4 z_{96} z_{211}-12 z_{96} z_{212}\\
&+z_{203} z_{208}+z_{204} z_{207}\\
&+z_{205} z_{210}+z_{206} z_{209}
\end{aligned}\tag{Q169}\label{v222:q169}
\end{equation}}

{\scriptsize\begin{equation}
\begin{aligned}
Q_{170,N}={}&+12 z_{96} z_{213}-8 z_{96} z_{214}\\
&+z_{203} z_{209}+2 z_{204} z_{210}\\
&+z_{205} z_{207}+2 z_{206} z_{208}
\end{aligned}\tag{Q170}\label{v222:q170}
\end{equation}}

{\scriptsize\begin{equation}
\begin{aligned}
Q_{171,N}={}&-4 z_{96} z_{213}+12 z_{96} z_{214}\\
&+z_{203} z_{210}+z_{204} z_{209}\\
&+z_{205} z_{208}+z_{206} z_{207}
\end{aligned}\tag{Q171}\label{v222:q171}
\end{equation}}

{\scriptsize\begin{equation}
\begin{aligned}
Q_{172,N}={}&-44 z_{96} z_{215}+48 z_{96} z_{216}\\
&+z_{203}^{2}+2 z_{204}^{2}
\end{aligned}\tag{Q172}\label{v222:q172}
\end{equation}}

{\scriptsize\begin{equation}
\begin{aligned}
Q_{173,N}={}&+12 z_{96} z_{215}-22 z_{96} z_{216}\\
&+z_{203} z_{204}
\end{aligned}\tag{Q173}\label{v222:q173}
\end{equation}}

{\scriptsize\begin{equation}
\begin{aligned}
Q_{174,N}={}&-z_{98}+z_{211}^{2}\\
&+2 z_{212}^{2}+11 z_{215}^{2}\\
&-24 z_{215} z_{216}+22 z_{216}^{2}
\end{aligned}\tag{Q174}\label{v222:q174}
\end{equation}}

{\scriptsize\begin{equation}
\begin{aligned}
Q_{175,N}={}&+z_{211} z_{212}-3 z_{215}^{2}\\
&+11 z_{215} z_{216}-6 z_{216}^{2}
\end{aligned}\tag{Q175}\label{v222:q175}
\end{equation}}

{\scriptsize\begin{equation}
\begin{aligned}
Q_{176,N}={}&+11 z_{35} z_{98}-24 z_{36} z_{98}\\
&-z_{98}+z_{213}^{2}\\
&+2 z_{214}^{2}
\end{aligned}\tag{Q176}\label{v222:q176}
\end{equation}}

{\scriptsize\begin{equation}
\begin{aligned}
Q_{177,N}={}&-3 z_{35} z_{98}+11 z_{36} z_{98}\\
&+z_{213} z_{214}
\end{aligned}\tag{Q177}\label{v222:q177}
\end{equation}}

{\scriptsize\begin{equation}
\begin{aligned}
Q_{178,N}={}&-6 z_{53} z_{213}+4 z_{53} z_{214}\\
&+4 z_{54} z_{213}-12 z_{54} z_{214}\\
&+z_{205} z_{217}+2 z_{206} z_{218}
\end{aligned}\tag{Q178}\label{v222:q178}
\end{equation}}

{\scriptsize\begin{equation}
\begin{aligned}
Q_{179,N}={}&+2 z_{53} z_{213}-6 z_{53} z_{214}\\
&-6 z_{54} z_{213}+4 z_{54} z_{214}\\
&+z_{205} z_{218}+z_{206} z_{217}
\end{aligned}\tag{Q179}\label{v222:q179}
\end{equation}}

{\scriptsize\begin{equation}
\begin{aligned}
Q_{180,N}={}&+z_{19} z_{203}+2 z_{20} z_{204}\\
&-6 z_{51} z_{211}+4 z_{51} z_{212}\\
&+4 z_{52} z_{211}-12 z_{52} z_{212}
\end{aligned}\tag{Q180}\label{v222:q180}
\end{equation}}

{\scriptsize\begin{equation}
\begin{aligned}
Q_{181,N}={}&+z_{19} z_{204}+z_{20} z_{203}\\
&+2 z_{51} z_{211}-6 z_{51} z_{212}\\
&-6 z_{52} z_{211}+4 z_{52} z_{212}
\end{aligned}\tag{Q181}\label{v222:q181}
\end{equation}}

{\scriptsize\begin{equation}
\begin{aligned}
Q_{182,N}={}&+4 z_{2}+z_{175} z_{177}\\
&+2 z_{176} z_{178}-z_{179} z_{181}\\
&-2 z_{180} z_{182}
\end{aligned}\tag{Q182}\label{v222:q182}
\end{equation}}

{\scriptsize\begin{equation}
\begin{aligned}
Q_{183,N}={}&+4 z_{3}+z_{175} z_{178}\\
&+z_{176} z_{177}-z_{179} z_{182}\\
&-z_{180} z_{181}
\end{aligned}\tag{Q183}\label{v222:q183}
\end{equation}}

{\scriptsize\begin{equation}
\begin{aligned}
Q_{184,N}={}&+2 z_{2} z_{109}+4 z_{3} z_{110}\\
&-z_{19} z_{181}-2 z_{20} z_{182}
\end{aligned}\tag{Q184}\label{v222:q184}
\end{equation}}

{\scriptsize\begin{equation}
\begin{aligned}
Q_{185,N}={}&+2 z_{2} z_{110}+2 z_{3} z_{109}\\
&-z_{19} z_{182}-z_{20} z_{181}
\end{aligned}\tag{Q185}\label{v222:q185}
\end{equation}}

{\scriptsize\begin{equation}
\begin{aligned}
Q_{186,N}={}&+1752629952 +6299926668 z_{17}\\
&+2158699200 z_{17} z_{217}-2185741488 z_{17} z_{218}\\
&-8034039216 z_{18}-2185741488 z_{18} z_{217}\\
&+4317398400 z_{18} z_{218}-853624686 z_{19}\\
&+402282898 z_{19} z_{217}+632859681 z_{20}\\
&+804565796 z_{20} z_{218}+559299438 z_{217}\\
&-615250593 z_{218}
\end{aligned}\tag{Q186}\label{v222:q186}
\end{equation}}

{\scriptsize\begin{equation}
\begin{aligned}
Q_{187,N}={}&-760760576 -8034039216 z_{17}\\
&-2185741488 z_{17} z_{217}+4317398400 z_{17} z_{218}\\
&+12599853336 z_{18}+4317398400 z_{18} z_{217}\\
&-4371482976 z_{18} z_{218}+632859681 z_{19}\\
&+804565796 z_{19} z_{218}-1707249372 z_{20}\\
&+804565796 z_{20} z_{217}-615250593 z_{217}\\
&+1118598876 z_{218}
\end{aligned}\tag{Q187}\label{v222:q187}
\end{equation}}

{\scriptsize\begin{equation}
\begin{aligned}
Q_{188,N}={}&-1149539095796314755226869484 +985061164914970179529857840 z_{17}\\
&+1254821386791945912952706976 z_{17} z_{19}+1153292377786262120636131146 z_{17} z_{20}\\
&+4398006762181672234512370320 z_{17} z_{21}-387008901290036540531246720 z_{17} z_{22}\\
&-1532572190428498920661008288 z_{17} z_{217}-1294531802031560732990489418 z_{17} z_{218}\\
&+1345259843509302298655015680 z_{18}+1153292377786262120636131146 z_{18} z_{19}\\
&+2509642773583891825905413952 z_{18} z_{20}-193504450645018270265623360 z_{18} z_{21}\\
&+17592027048726688938049481280 z_{18} z_{22}-1294531802031560732990489418 z_{18} z_{217}\\
&-3065144380856997841322016576 z_{18} z_{218}-73337278472469394027868160 z_{19}\\
&+686923528829960141567886528 z_{19} z_{21}-44414351330348966029310016 z_{19} z_{22}\\
&-96899299673172813900969984 z_{20}-22207175665174483014655008 z_{20} z_{21}\\
&+2747694115319840566271546112 z_{20} z_{22}+797514988115594446587154013 z_{21}\\
&+1723051656304506319706496064 z_{21} z_{68}+2606545027324788158562625392 z_{22}\\
&+13784413250436050557651968512 z_{22} z_{69}-975462260105801738014285793 z_{68}\\
&-887540248202664802539362352 z_{69}+588237978875945932189940160 z_{217}\\
&+759606688303604527483518672 z_{218}
\end{aligned}\tag{Q188}\label{v222:q188}
\end{equation}}

{\scriptsize\begin{equation}
\begin{aligned}
Q_{189,N}={}&-652461006517642156833128616 +672629921754651149327507840 z_{17}\\
&+576646188893131060318065573 z_{17} z_{19}+1254821386791945912952706976 z_{17} z_{20}\\
&-96752225322509135132811680 z_{17} z_{21}+8796013524363344469024740640 z_{17} z_{22}\\
&-647265901015780366495244709 z_{17} z_{217}-1532572190428498920661008288 z_{17} z_{218}\\
&+985061164914970179529857840 z_{18}+1254821386791945912952706976 z_{18} z_{19}\\
&+1153292377786262120636131146 z_{18} z_{20}+4398006762181672234512370320 z_{18} z_{21}\\
&-387008901290036540531246720 z_{18} z_{22}-1532572190428498920661008288 z_{18} z_{217}\\
&-1294531802031560732990489418 z_{18} z_{218}-48449649836586406950484992 z_{19}\\
&-11103587832587241507327504 z_{19} z_{21}+1373847057659920283135773056 z_{19} z_{22}\\
&-73337278472469394027868160 z_{20}+686923528829960141567886528 z_{20} z_{21}\\
&-44414351330348966029310016 z_{20} z_{22}+651636256831197039640656348 z_{21}\\
&+3446103312609012639412992128 z_{21} z_{69}+1595029976231188893174308026 z_{22}\\
&+3446103312609012639412992128 z_{22} z_{68}-221885062050666200634840588 z_{68}\\
&-1950924520211603476028571586 z_{69}+379803344151802263741759336 z_{217}\\
&+588237978875945932189940160 z_{218}
\end{aligned}\tag{Q189}\label{v222:q189}
\end{equation}}

{\scriptsize\begin{equation}
\begin{aligned}
Q_{190,N}={}&+186960 -667923 z_{17} z_{19}\\
&-289924 z_{17} z_{20}+310419 z_{17} z_{217}\\
&-61308 z_{17} z_{218}-289924 z_{18} z_{19}\\
&-1335846 z_{18} z_{20}-61308 z_{18} z_{217}\\
&+620838 z_{18} z_{218}+1398264 z_{68}\\
&+658952 z_{68} z_{217}-2073288 z_{69}\\
&+2635808 z_{69} z_{218}-147928 z_{217}\\
&-161376 z_{218}
\end{aligned}\tag{Q190}\label{v222:q190}
\end{equation}}

{\scriptsize\begin{equation}
\begin{aligned}
Q_{191,N}={}&+54858 -144962 z_{17} z_{19}\\
&-667923 z_{17} z_{20}-30654 z_{17} z_{217}\\
&+310419 z_{17} z_{218}-667923 z_{18} z_{19}\\
&-289924 z_{18} z_{20}+310419 z_{18} z_{217}\\
&-61308 z_{18} z_{218}-518322 z_{68}\\
&+658952 z_{68} z_{218}+2796528 z_{69}\\
&+1317904 z_{69} z_{217}-80688 z_{217}\\
&-147928 z_{218}
\end{aligned}\tag{Q191}\label{v222:q191}
\end{equation}}

{\scriptsize\begin{equation}
\begin{aligned}
Q_{192,N}={}&-z_{82} z_{98}+z_{83} z_{97}\\
&+z_{219}
\end{aligned}\tag{Q192}\label{v222:q192}
\end{equation}}

{\scriptsize\begin{equation}
\begin{aligned}
Q_{193,N}={}&+z_{83} z_{99}-z_{84} z_{98}\\
&+z_{220}
\end{aligned}\tag{Q193}\label{v222:q193}
\end{equation}}

{\scriptsize\begin{equation}
\begin{aligned}
Q_{194,N}={}&-z_{81} z_{98}+z_{83} z_{96}\\
&+z_{221}
\end{aligned}\tag{Q194}\label{v222:q194}
\end{equation}}

{\scriptsize\begin{equation}
\begin{aligned}
Q_{195,N}={}&-8 z_{4} z_{92}+z_{119} z_{222}\\
&+z_{181}^{2}+2 z_{182}^{2}
\end{aligned}\tag{Q195}\label{v222:q195}
\end{equation}}

{\scriptsize\begin{equation}
\begin{aligned}
Q_{196,N}={}&+z_{98} z_{222}-z_{167}^{2}\\
&-2 z_{168}^{2}-z_{195}^{2}\\
&-2 z_{196}^{2}-z_{209}^{2}\\
&-2 z_{210}^{2}
\end{aligned}\tag{Q196}\label{v222:q196}
\end{equation}}

\subsection{Completion and coefficient profile}\label{v222:height}
The four-contact construction proves both arithmetic directions from
Section~\ref{v15x:completion}, with a possibly increased ineffective
uniform threshold. Every subsequent transformation above has a two-sided
map on the admitted-family solution set, or removes a relation proved
from the retained ones in the stated order. The last affine and row
operations preserve that zero set, and positive weights preserve it
again. This proves Theorem~\ref{thm:v15-main} for the displayed polynomial.

Expanding the196 weighted squares gives2277 coordinate monomials,
or2278 terms if the coefficient parameter is separated. The polynomial
has degree four and content one, and all222 coordinates occur. The only
quadratic dependence on $A_N$ occurs in the coefficient of $z_{108}^2$,
which is $A_N^2+1$. Linear parameter coefficients have absolute value at
most44402; every parameter-free coefficient has magnitude less than
$149^{214}$. Consequently for $N\ge53$ the exact largest absolute
coefficient is $A_N^2+1=149^{4N+2}+1$. These are finite coefficient facts,
not an effective bound for the defining threshold. No particular
numerical $N$ is certified to define the integers solely by this calculation.

The full ordinary decimal314-variable example and the original huge
polynomial in the earlier sections are retained independently of this
family. The222 construction does not replace their distinct finite
coefficient specifications or strengthen any inherited arithmetic premise.

\begin{thebibliography}{99}
\bibitem{direct-p5-modularity-project} Version 14 project proof supplement,
\emph{Direct prime-five modularity for the rational scalar descent},
version 1. Exact source applications, finite-twist connection and historical
route map at
\path{research/version14/proofs/direct-p5-modularity-v1/}.

\bibitem{Mitsui1956}
T.~Mitsui,
\emph{Generalized Prime Number Theorem},
Japanese Journal of Mathematics \textbf{26} (1956), 1--42.
\url{https://doi.org/10.4099/jjm1924.26.0_1}.

\bibitem{GTOrbits2012}
B. Green and T. Tao, \emph{The quantitative behaviour of polynomial
orbits on nilmanifolds}, Annals of Mathematics \textbf{175} (2012), 465--540.
\url{https://doi.org/10.4007/annals.2012.175.2.2}.
\bibitem{v6bienvenu}
P.-Y. Bienvenu, \emph{A higher-dimensional Siegel--Walfisz theorem},
Acta Arithmetica \textbf{179} (2017), 79--100.
\url{https://doi.org/10.4064/aa8600-10-2016}.

\bibitem{CTS1989Cubic} J.-L. Colliot-Th\'el\`ene and P. Salberger,
\emph{Arithmetic on some singular cubic hypersurfaces},
Proc. London Math. Soc. (3) \textbf{58} (1989), 519--549,
\S1, Proposition~1.3, pp.~520--521.
\bibitem{CTSSD1987} J.-L. Colliot-Th\'el\`ene, J.-J. Sansuc and
P. Swinnerton-Dyer, \emph{Intersections of two quadrics and Ch\^atelet
surfaces. I}, J. reine angew. Math. \textbf{373} (1987), 37--107,
Remark~2.3.1, p.~56.
\bibitem{cubic-source-project} Project source connection,
\emph{The published characteristic-zero cubic input}, version~1,
\href{https://github.com/elliotglazer/diophantine-z-q/blob/main/research/version14/proofs/cubic-unirationality-v1/PROOF.md}{exact hypotheses and quotient application}.

\bibitem{Muller2002HIT} P. M\"uller,
\emph{Finiteness results for Hilbert's irreducibility theorem},
Ann. Inst. Fourier \textbf{52} (2002), no.~4, 983--1015.
\href{https://doi.org/10.5802/aif.1907}{doi:10.5802/aif.1907}.

\bibitem{agentic-improvements} Agentic notes,
\emph{Conditional polynomial constructions, lower bounds, and number-field
obstructions for rational Diophantine integrality}, working version 3,
6 November 2026. Public source at Git revision
\nolinkurl{f4d1b5e6c23d4adeb3ae1f986ab0bf87884bdff6}.
\href{https://github.com/JoseBrox/diophantine-z-q_improvements/blob/f4d1b5e6c23d4adeb3ae1f986ab0bf87884bdff6/paper/diophantine_optimization_20261106_v3.tex}{Public source at the pinned revision}.
% Insert these items inside the manuscript's existing thebibliography.
\bibitem{v5sun}
Z.-W. Sun, \emph{86 unknowns for Hilbert's tenth problem over the rationals},
arXiv:2610.20691v2, 2026. Theorem~4.1, Sections~4--6 and Appendix~A.
\url{https://arxiv.org/abs/2610.20691v2}.

\bibitem{v5bl}
T. Browning and D. Loughran, \emph{Sieving rational points on varieties},
Trans. Amer. Math. Soc. \textbf{371} (2019), 5757--5785.
\url{https://doi.org/10.1090/tran/7514}.
Theorem~4.1 and its adjacent conventions were compared with the final
published text, pp.~5772--5773, including uniformity in the modulus and
the allowed local subsets.

\bibitem{v5bhb}
T. D. Browning and D. R. Heath-Brown, \emph{The geometric sieve for quadrics},
Forum Math. \textbf{33} (2021), no.~1, 147--165.
\url{https://doi.org/10.1515/forum-2020-0074}.
Lemmas~2.1 and 4.1 and the Section~4 lattice-basis estimates were compared
with the final published text, pp.~150 and 155--158. The local notation
slips and their proof-based interpretations are recorded in the
\href{https://github.com/elliotglazer/diophantine-z-q/blob/main/research/version14/evidence/final-source-bhb/BHB-FINAL-STATEMENT-DISPOSITION.md}{source audit}.

\bibitem{v5ft}
C. Frei and M. Tinkov\'a, with an appendix by C. Frei, J. K\"onig and
M. Tinkov\'a, \emph{Linear constellations in primes with arithmetic restrictions},
arXiv:2408.17441v3, 5 November 2025.
\url{https://arxiv.org/abs/2408.17441v3}.
Theorem~1.4. Listed as \emph{Mathematika}, to appear, on
\url{https://www.math.tugraz.at/~frei/}, checked 6 November 2026.
No final published version was verified.

\bibitem{v5tz}
T. Tao and T. Ziegler, \emph{Polynomial patterns in the primes},
Forum Math. Pi \textbf{6} (2018), e1, 60 pp.
\url{https://doi.org/10.1017/fmp.2017.3}.
Published Theorems~9, 11, 13 and Corollary~3.

\bibitem{v5kl}
G. S. Kopp and J. C. Lagarias,
\emph{Ray class groups and ray class fields for orders of number fields},
Essential Number Theory \textbf{4} (2025), no.~1, 1--65.
\url{https://doi.org/10.2140/ent.2025.4.1}.
Theorems~1.1, 1.3, 6.5, 7.2 and Remark~7.3.

\bibitem{v5williams}
K. S. Williams, \emph{Mertens' theorem for arithmetic progressions},
J. Number Theory \textbf{6} (1974), 353--359, Theorem~1, p.~356.
\url{https://doi.org/10.1016/0022-314X(74)90032-8}.

\bibitem{scholze2013} P. Scholze,
\emph{$p$-adic Hodge theory for rigid-analytic varieties},
Forum of Mathematics, Pi 1 (2013), e1, 1--77,
Corollary 1.8 and Theorems 9.1, 9.3.
\url{https://doi.org/10.1017/fmp.2013.1}.
\bibitem{scholze-erratum} P. Scholze,
\emph{$p$-adic Hodge theory for rigid-analytic varieties---corrigendum},
Forum of Mathematics, Pi 4 (2016), e6, 1--4.
\url{https://doi.org/10.1017/fmp.2016.4}.
\bibitem{fontaine1994-representations} J.-M. Fontaine,
\emph{Repr\'esentations $p$-adiques semi-stables},
in \emph{P\'eriodes $p$-adiques}, Ast\'erisque 223 (1994),
113--184, Theorem 3.8, p.~142; definitions in Sections 3.4--3.7.
\url{https://www.numdam.org/item/AST_1994__223__113_0/}.
\bibitem{brinon-conrad} O. Brinon and B. Conrad,
\emph{CMI summer school notes on $p$-adic Hodge theory},
preliminary version, June 24, 2009,
Theorem 5.2.1 and Propositions 6.3.3, 6.3.8.
\url{https://math.stanford.edu/~conrad/papers/notes.pdf}.
\bibitem{milne-shimura2005} J. S. Milne, \emph{Introduction to Shimura varieties},
Clay Mathematics Proceedings 4 (2005), 265--378,
Theorems 6.7--6.8 and Corollary 6.9, p.~319; for H5,
Theorems 7.4 and 14.8, pp.~321--322 and 352,
13.6--13.7, p.~349, and the type-A/C criterion, pp.~354--355.
The full Hodge-orbit condition of 7.4 is used; the general implication
in the original 8.14/8.17 is not used without its correction.
\url{https://www.jmilne.org/math/xnotes/svi2005.pdf}.
\bibitem{milne-av} J. S. Milne, \emph{Abelian varieties},
version 2.00 (2008), Part I, Theorem 12.1(a).
\url{https://www.jmilne.org/math/CourseNotes/AV.pdf}.
\bibitem{babagranath} S. Baba and H. Granath,
\emph{Genus 2 curves with quaternionic multiplication},
Canadian Journal of Mathematics 60 (2008), 734--757,
\S2.1, Theorem 2.1, Lemma 3.1 and the fixed-point table on p.~740.
\url{https://doi.org/10.4153/CJM-2008-033-7}.
The printed principal-form normalization and the local index slip in
Lemma 3.1 are treated in \cite{a001-moduli-project}; the explicit
invariant normalization is supplied separately by
\cite{ly-d6-normalization-project}.
\bibitem{a001-moduli-project} Project proof,
\emph{The fixed principal datum and its discriminant-six coarse quotient},
\href{https://github.com/elliotglazer/diophantine-z-q/blob/main/research/version14/proofs/a001-moduli-current-v1/PROOF.md}{Fixed-datum and coarse-moduli proof}.
The adjacent Baba--Granath product-map finding is preserved separately.
\bibitem{ly-d6-normalization-project} Project proof,
\emph{Discriminant-six normalization from Lin--Yang Example 30},
\href{https://github.com/elliotglazer/diophantine-z-q/blob/main/research/version14/proofs/ly-d6-normalization-v1/PROOF.md}{Normalization proof},
with its current-consumer connection and exact source inputs.
\bibitem{b11-integral-model-project} Project proof,
\emph{The full-three integral model and its normalizer degree},
\href{https://github.com/elliotglazer/diophantine-z-q/blob/main/research/version14/proofs/b11-integral-model-v1/ASSEMBLY.md}{Integral-model assembly}.
Includes the local-order, normalizer, five-marking, Tate and
scheme-descent arguments with their separately stated inputs.

\bibitem{Deligne1971} P. Deligne, \emph{Travaux de Shimura},
Seminaire Bourbaki, expose 389 (1970--1971),
Proposition 1.15, Theorems 4.7 and 4.21,
\S\S4.9--4.17, Corollaries 5.5 and 5.7, and Example 5.8.
\bibitem{stacks-h5-fine} The Stacks Project,
Tags 03YR, 0ADD, 0A18, 0E83, 0D3A, 0D32 and 0B45,
official statements retrieved November 8, 2026.
\url{https://stacks.math.columbia.edu/}.
These supply the stated trivial-inertia, curve, Stein-factorization,
genus and ampleness clauses. Exact versions and the full H5 construction
are recorded in \path{research/version14/proofs/h5-full-three-v1/}.
\bibitem{buzzard} K. Buzzard, \emph{Integral models of certain Shimura
curves}, Duke Mathematical Journal \textbf{87} (1997), 591--612,
\S1, Theorem 2.1 and Corollary 2.3.
\url{https://doi.org/10.1215/S0012-7094-97-08719-6}.
\bibitem{linyang} Y.-H. Lin and Y. Yang, \emph{Quaternionic loci in
Siegel's modular threefold}, Mathematische Zeitschrift \textbf{295} (2020),
775--819. \url{https://doi.org/10.1007/s00209-019-02372-z}.
\bibitem{berger2002} L. Berger, \emph{Repr\'esentations $p$-adiques et
\'equations diff\'erentielles}, Inventiones Mathematicae 148 (2002),
219--284, $p$-adic monodromy theorem.
\url{https://doi.org/10.1007/s002220100202}.
\bibitem{pan2022} L. Pan, \emph{The Fontaine--Mazur conjecture in the
residually reducible case}, arXiv:1901.07166v2,
Theorem 1.0.4, p.~5, the historical author-proof version. Published in Journal of the
American Mathematical Society 35 (2022), 1031--1169.
\url{https://arxiv.org/abs/1901.07166v2}.
\bibitem{schneider-teitelbaum2002} P. Schneider and J. Teitelbaum,
\emph{Banach space representations and Iwasawa theory},
Israel Journal of Mathematics 127 (2002), 359--380.
\url{https://doi.org/10.1007/BF02784538}.
\bibitem{pan-scalar-supplement} Project proof supplement,
\emph{A published replacement for the surjectivity step in Pan's Lemma 3.5.9},
version 1. The scalar coefficient case, with upstream admissibility
and universal completion retained as premises.
\url{https://github.com/elliotglazer/diophantine-z-q/blob/main/docs/research-review/proof-sources/pan-scalar-closed-image/proof-v1.md}.
\bibitem{serre-ladic-published} J.-P. Serre,
\emph{Abelian $\ell$-adic representations and elliptic curves},
Addison--Wesley, 1989, reprint of the 1968 edition,
Chapter III, \S1.2, p.~III-7, and Appendix A.7,
Theorem 3, pp.~III-52--53.
\url{https://www.math.mcgill.ca/darmon/courses/18-19/gs/serre-mcgill.pdf}.
\bibitem{li-newforms} W.-C. W. Li,
\emph{Newforms and functional equations}, Mathematische Annalen
212 (1975), 285--315, p.~296.
\bibitem{shimura-factors} G. Shimura,
\emph{On the factors of the jacobian variety of a modular function field},
Journal of the Mathematical Society of Japan 25 (1973), 523--544,
Theorem 1.
\url{https://doi.org/10.2969/jmsj/02530523}.
\bibitem{javanpeykar} A. Javanpeykar, with an appendix by P. Bruin,
\emph{Polynomial bounds for Arakelov invariants of Belyi curves},
Algebra \& Number Theory 8 (2014), 89--140,
Theorem 1.1.1 and Section 2.3.
\url{https://doi.org/10.2140/ant.2014.8.89}.
\bibitem{gr-periods2014} \'E. Gaudron and G. R\'emond,
\emph{Th\'eor\`eme des p\'eriodes et degr\'es minimaux d'isog\'enies},
Commentarii Mathematici Helvetici 89 (2014), no.~2, 343--403,
\S2.3 and Corollaire~8.4.
\url{https://doi.org/10.4171/CMH/322}.
\bibitem{gr2014} \'E. Gaudron and G. R\'emond,
\emph{Polarisations et isog\'enies}, Duke Mathematical Journal
163 (2014), 2057--2108, Theorem 1.4 and Section 9.
\url{https://doi.org/10.1215/00127094-2782528}.
\bibitem{fowler} G. Fowler, \emph{Some uniform effective results on
Andr\'e--Oort for sums of powers in $\mathbb C^n$}, Mathematical
Proceedings of the Cambridge Philosophical Society 180 (2026),
607--641, \S2.1 and Proposition 2.7.
\url{https://doi.org/10.1017/S0305004125101825}.
\bibitem{cmheight} A. Barquero-Sanchez, L. Cadwallader, O. Cannon,
T. Genao and R. Masri, \emph{Faltings heights of CM elliptic curves
and special Gamma values}, Research in Number Theory 3 (2017),
article 13, Section 6, equation (6.2).
\url{https://doi.org/10.1007/s40993-017-0077-7}.
\bibitem{silverbergzarhin} A. Silverberg and Yu. G. Zarhin,
\emph{Semistable reduction and torsion subgroups of abelian varieties},
Annales de l'Institut Fourier 45 (1995), 403--420, Theorem 3.5.
\url{https://www.numdam.org/article/AIF_1995__45_2_403_0.pdf}.
\bibitem{serrelocal} J.-P. Serre, \emph{Local class field theory},
in J. W. S. Cassels and A. Fr\"ohlich (eds.), \emph{Algebraic Number Theory},
Academic Press, 1967, Chapter VI, 128--161, Sections 4.1--4.3.
\bibitem{faltings1983} G. Faltings, \emph{Endlichkeitss\"atze f\"ur
abelsche Variet\"aten \"uber Zahlk\"orpern}, Inventiones Mathematicae
73 (1983), 349--366; erratum, 75 (1984), 381.
\bibitem{faltings1988} G. Faltings, \emph{$p$-adic Hodge theory},
Journal of the American Mathematical Society 1 (1988), 255--299.

\bibitem{igusa1967} J.-I. Igusa, \emph{Modular forms and projective invariants},
American Journal of Mathematics \textbf{89} (1967), 817--855,
pp.~833--834 and 848. \url{https://doi.org/10.2307/2373243}.
\bibitem{igusa-theta-project} Project source connection,
\emph{Igusa's theta identities at the actual genus-two point}, version 1,
\href{https://github.com/elliotglazer/diophantine-z-q/blob/main/research/version14/proofs/igusa-theta-v1/PROOF.md}{definitions, coordinate comparison and exact substitution}.
\bibitem{pazuki} F. Pazuki, \emph{Theta height and Faltings height},
Bulletin de la Soci\'et\'e Math\'ematique de France, 2012,
Sections 2.2--2.3 and Corollary 1.3.
\url{https://www.numdam.org/item/10.24033/bsmf.2623.pdf}.
\bibitem{poonenstoll} B. Poonen and M. Stoll, \emph{The Cassels--Tate
pairing on polarized abelian varieties}, Annals of Mathematics 150
(1999), 1109--1149, Corollaries 4, 7 and Theorem 8.
\url{https://ftp.gwdg.de/pub/misc/EMIS/journals/Annals/150_3/10.html}.
\bibitem{greentao} B. Green and T. Tao, \emph{Linear equations in primes},
Annals of Mathematics 171 (2010), 1753--1850.
\url{https://annals.math.princeton.edu/2010/171-3/p08}.
\bibitem{moebius} B. Green and T. Tao, \emph{The M\"obius function is
strongly orthogonal to nilsequences}, Annals of Mathematics 175 (2012).
\url{https://annals.math.princeton.edu/2012/175-2/p03}.
\bibitem{local-gi} Project proof supplement,
\emph{A local proof supplement for the required Gowers inverse theorem},
version~1, 8 November 2026, accompanying this manuscript as
\texttt{quartic-local-gi-supplement.pdf}.
Assembly, Sections~1--6, with the supporting proof modules and
the exact published inputs listed in its source register.
\bibitem{bmor} M. A. Bennett, G. Martin, K. O'Bryant and A. Rechnitzer,
\emph{Explicit bounds for primes in arithmetic progressions},
Illinois Journal of Mathematics 62 (2018), 427--532,
arXiv:1802.00085v3, Theorem 1.2.
\url{https://arxiv.org/abs/1802.00085v3}.
\bibitem{milne} J. S. Milne, \emph{Arithmetic duality theorems},
second edition, 2006, Chapter I.
\url{https://www.jmilne.org/math/Books/ADTnot.pdf}.
\bibitem{stackscohom} The Stacks Project,
\emph{Stalks of higher direct images} and \emph{The Leray spectral sequence},
Theorem 59.53.1, Tag 03Q9, and Proposition 59.54.2, Section Tag 03QA,
HTML snapshots of November 5, 2026.
\url{https://stacks.math.columbia.edu/tag/03Q9};
\url{https://stacks.math.columbia.edu/tag/03QA}.
\bibitem{nekovar} J. Nekov\'a\v r, \emph{Selmer complexes},
Ast\'erisque 310, 2006, Proposition 5.2.4, Theorems 6.3.4 and 7.8.8, and Sections 5.7 and 6.4.
\url{https://www.numdam.org/item/AST_2006__310__1_0/}.
\bibitem{afp} Archive of Formal Proofs, \emph{DPRM\_Theorem},
\texttt{Alpha\_Sequence}, \texttt{Exponential\_Relation}, and
\texttt{Exponentiation}, directory
\texttt{thys/DPRM\_Theorem/Diophantine/}, AFP development git mirror
revision \texttt{b4f3b892d0ae25cea27412ba4007cdb63e44b0eb}
(May 10, 2026). The three source files match this revision byte for byte.
\url{https://isa-afp.org/entries/DPRM_Theorem.html}.
\bibitem{edixhoven}
B. Edixhoven, \emph{N\'eron models and tame ramification},
Compositio Mathematica 81 (1992), 291--306,
Theorem~5.3 and Remark~5.4.1.
\url{https://www.numdam.org/item/CM_1992__81_3_291_0/}.
\bibitem{sutherland} A. V. Sutherland,
\emph{Isogeny kernels and division polynomials}, MIT 18.783,
Fall 2025, Lecture 5, September 18, 2025, \S\S5.5--5.6.
\url{https://ocw.mit.edu/courses/18-783-elliptic-curves-fall-2025/mit18_783_f25_lec05.pdf}.
\bibitem{silverman} J. H. Silverman, \emph{The arithmetic of elliptic
curves}, second edition, Springer, 2009, Chapters IV and VII.
\bibitem{ddselmer} T. Dokchitser and V. Dokchitser,
\emph{On the Birch--Swinnerton-Dyer quotients modulo squares},
Annals of Mathematics 172 (2010), 567--596,
Theorem 1.4 and Conjecture 1.2, p.~568; Theorem 4.19, p.~593.
\url{https://doi.org/10.4007/annals.2010.172.567}.
\bibitem{bcdt} C. Breuil, B. Conrad, F. Diamond and R. Taylor,
\emph{On the modularity of elliptic curves over $\Q$: wild 3-adic exercises},
author manuscript, Theorem A and Introduction, conditions (1)--(6).
\url{https://virtualmath1.stanford.edu/~conrad/papers/tswfinal.pdf}.
Published in Journal of the American Mathematical Society 14 (2001), 843--939.
\bibitem{amirhong}
M. Amir and L. Hong, \emph{On L-functions of modular elliptic curves
and certain K3 surfaces}, Ramanujan Journal 57 (2022), 1001--1019,
Theorem 3.1. \url{https://doi.org/10.1007/s11139-021-00388-w}.

\bibitem{milnemf} J. S. Milne, \emph{Modular functions and modular forms},
v1.31, March 22, 2017, Remark 11.18 and Theorem 11.19.
\url{https://www.jmilne.org/math/CourseNotes/MF.pdf}.
\bibitem{elkik} R. Elkik, \emph{Le th\'eor\`eme de Manin--Drinfeld},
Ast\'erisque 183 (1990), 59--67, Theorem, p.~59.
\url{https://www.numdam.org/item/AST_1990__183__59_0.pdf}.
\bibitem{gz} B. Gross and D. Zagier, \emph{Heegner points and derivatives
of $L$-series}, Inventiones Mathematicae 84 (1986), 225--320.
\bibitem{kolyvagin} V. A. Kolyvagin,
\emph{On the Mordell--Weil group and the Shafarevich--Tate group
of modular elliptic curves}, Proceedings of the International Congress
of Mathematicians, Kyoto 1990, Vol.~I, Mathematical Society of Japan,
1991, pp.~429--436, Theorems 1--5, pp.~430--431.
\url{https://www.mathunion.org/fileadmin/ICM/Proceedings/ICM1990.1/ICM1990.1.ocr.pdf}.
\bibitem{deligne-serre1974} P. Deligne and J.-P. Serre,
\emph{Formes modulaires de poids 1},
Annales scientifiques de l'\'Ecole Normale Sup\'erieure (4)
\textbf{7} (1974), no.~4, 507--530.
\url{https://www.numdam.org/item/ASENS_1974_4_7_4_507_0/}

\bibitem{kato} K. Kato, \emph{$p$-adic Hodge theory and values of zeta
functions of modular forms}, Ast\'erisque 295 (2004), 117--290.
\url{https://www.numdam.org/article/AST_2004__295__117_0.pdf}.
\bibitem{bfh} D. Bump, S. Friedberg and J. Hoffstein,
\emph{Nonvanishing theorems for $L$-functions of modular forms and their
derivatives}, Inventiones Mathematicae 102 (1990), 543--618.
\url{https://wstein.org/papers/bib/bump-friedberg-hoffstein-nonvanishing.pdf}.
\bibitem{hl} J. Hoffstein and W. Luo, \emph{Nonvanishing of $L$-series and
the combinatorial sieve}, Mathematical Research Letters 4 (1997), 435--444.
\url{https://intlpress.com/site/pub/files/_fulltext/journals/mrl/1997/0004/0003/MRL-1997-0004-0003-a012.pdf}.
\bibitem{desjardins} J. Desjardins, \emph{Root number of twists of an
elliptic curve}, Journal de Th\'eorie des Nombres de Bordeaux
32 (2020), 73--101. \url{https://jtnb.centre-mersenne.org/item/10.5802/jtnb.1112.pdf}.
\bibitem{cst} L. Cai, J. Shu and Y. Tian, \emph{Explicit Gross--Zagier
and Waldspurger formulae}, Algebra \& Number Theory 8 (2014).
\url{https://arxiv.org/abs/1408.1733}.
\bibitem{ccl} L. Cai, Y. Chen and Y. Liu, \emph{Heegner points on modular
curves}, Transactions of the American Mathematical Society 370 (2018),
3721--3743. \url{https://archive.ymsc.tsinghua.edu.cn/pacm_download/632/12220-Heegner_points_on_modular_curves.pdf}.
\bibitem{ribet-geometry-1990} K. A. Ribet,
\emph{On modular representations of $\operatorname{Gal}(\overline{\mathbb Q}/\mathbb Q)$
arising from modular forms}, Inventiones Mathematicae 100 (1990),
431--476, Section 3, Proposition 3.3, Remarks 3.5(a)--(b), and Proposition 3.6.
\url{https://math.berkeley.edu/~ribet/Articles/invent_100.pdf}.
\bibitem{jetchevkane} D. Jetchev and B. Kane,
\emph{Equidistribution of Heegner Points and Ternary Quadratic Forms},
arXiv:0908.3905 (2009), Sections 1.1--1.2 and 2.5.
\url{https://arxiv.org/abs/0908.3905}.
\bibitem{cj-conductor-one-project}
Versioned project proof, \emph{A direct conductor-one reduction argument},
version 1. Exact proof, explicit retained inputs and bounded review at
\href{https://github.com/elliotglazer/diophantine-z-q/blob/main/research/version14/proofs/cj-conductor-one-v1/README.md}
{\texttt{research/version14/proofs/cj-conductor-one-v1/}}.
\bibitem{cornutjetchev} C. Cornut and D. Jetchev,
\emph{Liftings of reduction maps for quaternion algebras}, 2010 preprint.
\url{https://webusers.imj-prg.fr/~christophe.cornut/papers/liftings.pdf}.
\bibitem{waldspurger} J.-L. Waldspurger, \emph{Sur les coefficients de
Fourier des formes modulaires de poids demi-entier},
Journal de Math\'ematiques Pures et Appliqu\'ees 60 (1981), 375--484,
Proposition 18 and Corollary 2.
\bibitem{ramsey} N. Ramsey, \emph{Geometric and $p$-adic modular forms
of half-integral weight}, Annales de l'Institut Fourier 56 (2006),
599--624. \url{https://www.numdam.org/item/10.5802/aif.2195.pdf}.
\bibitem{casselman}
W. Casselman, \emph{On some results of Atkin and Lehner},
Mathematische Annalen 201 (1973), 301--314,
Theorem~1 and its proof, pp.~302--306, and Theorem~4, p.~313.
\bibitem{carayol} H. Carayol, \emph{Sur les repr\'esentations
$\ell$-adiques associ\'ees aux formes modulaires de Hilbert},
Annales Scientifiques de l'\'Ecole Normale Sup\'erieure 19 (1986), 409--468,
Theorem A. \url{https://www.numdam.org/item/ASENS_1986_4_19_3_409_0/}.
\bibitem{waldspurgerICM} J.-L. Waldspurger,
\emph{Correspondances de Shimura}, Proceedings of the International
Congress of Mathematicians, 1983, 525--531.
\url{https://www.imj-prg.fr/wp-content/uploads/2020/prix/waldspurger1983.pdf}.
\bibitem{weil} A. Weil, \emph{Sur certains groupes d'op\'erateurs unitaires},
Acta Mathematica 111 (1964), 143--211, I.13, II.25, II.28, III, Theorem 6.
\url{https://webhomes.maths.ed.ac.uk/~v1ranick/papers/weil2.pdf}.
\bibitem{serre} J.-P. Serre, \emph{Propri\'et\'es galoisiennes des points
d'ordre fini des courbes elliptiques}, Inventiones Mathematicae 15
(1972), 259--331.
\bibitem{mazur1992} B. Mazur,
\emph{The topology of rational points}, Experimental Mathematics 1
(1992), 35--45.
\url{https://doi.org/10.1080/10586458.1992.10504244}.
\bibitem{robinson1949} J. Robinson,
\emph{Definability and decision problems in arithmetic},
Journal of Symbolic Logic 14 (1949), 98--114, Theorem 3.1.
\url{https://doi.org/10.2307/2266510}.
\bibitem{poonen2009} B. Poonen,
\emph{Characterizing integers among rational numbers with a
universal-existential formula}, American Journal of Mathematics 131
(2009), 675--682, Theorems 3.1 and 4.1. Author version with the April 2019
correction to Lemma 2.3.
\url{https://math.mit.edu/~poonen/papers/ae.pdf}.
\bibitem{kato1999} K. Kato,
\emph{Euler systems, Iwasawa theory, and Selmer groups},
Kodai Mathematical Journal 22 (1999), 313--372,
Sections 0.1 and 0.4--0.7, Example 0.6.2 and Theorem 0.8,
pp.~314--318.
\url{https://doi.org/10.2996/kmj/1138044090}.
\bibitem{openaipointwise} OpenAI,
\emph{A pointwise 2-converse for elliptic curves with rational
two-torsion}, preprint, 7 October 2026, Theorem 1.1; source revision
\nolinkurl{4eaee87548250b2682affda7d05d947c05554bae}.
\href{https://github.com/openai/math/blob/4eaee87548250b2682affda7d05d947c05554bae/preprints/A-pointwise-2-converse-for-elliptic-curves-with-rational-two-torsion-October-7-2026/paper.pdf}{Pinned source PDF}.
\bibitem{openaigoldfeld} OpenAI,
\emph{Goldfeld's analytic density conjecture and the 2-converse
for elliptic curves}, preprint, 7 October 2026, Theorem 1.1;
same frozen source revision.
\href{https://github.com/openai/math/blob/4eaee87548250b2682affda7d05d947c05554bae/preprints/Goldfelds-analytic-density-conjecture-and-the-2-converse-for-elliptic-curves-October-7-2026/paper.pdf}{Pinned source PDF}.
\bibitem{mazur1995} B. Mazur,
\emph{Speculations about the topology of rational points: an up-date},
Ast\'erisque 228 (1995), 165--181, Conjecture 3, p.~165.
\url{https://www.numdam.org/article/AST_1995__228__165_0.pdf}.
\bibitem{matiyasevich1971} Yu. V. Matiyasevich,
\emph{Diophantine representation of enumerable predicates},
Mathematics of the USSR-Izvestiya 5 (1971), 1--28,
Basic Theorem, p.~3; integer-witness conversion, p.~2.
\url{https://doi.org/10.1070/IM1971v005n01ABEH001004}.
\bibitem{openaihtp} OpenAI,
\emph{Hilbert's tenth problem over the rational numbers},
OpenAI Math Release preprint, 6 October 2026.
\url{https://github.com/openai/math/blob/main/preprints/Hilberts-tenth-problem-over-the-rational-numbers-October-6-2026/main.pdf}.
\bibitem{howard2012} B. Howard, \emph{The Heegner point Kolyvagin system},
Compositio Mathematica 140 (2004), 1439--1472; arXiv:1202.6340v1.
\bibitem{lang1956} S. Lang, \emph{Algebraic groups over finite fields},
American Journal of Mathematics 78 (1956), 555--563.
\bibitem{williams-mertens}
K. S. Williams, \emph{Mertens' theorem for arithmetic progressions},
J. Number Theory \textbf{6} (1974), 353--359, Theorem 1.
\url{https://people.math.carleton.ca/~williams/papers/pdf/057.pdf}.

\bibitem{voight-quaternion}
J. Voight, \emph{Quaternion algebras}, Graduate Texts in Mathematics
\textbf{288}, Springer, 2021, Main Theorem 14.6.1 and
Proposition 14.6.7, pp.~231--232.
The H5 division-order calculation also uses 13.3.4 and 13.3.11,
pp.~206--207, with valuation $\operatorname{ord}_3(\operatorname{nrd})/2$.
\url{https://doi.org/10.1007/978-3-030-56694-4_14}.


% Insert inside the existing thebibliography environment.
% This is the only new bibliography key used by assembled-original-extension.tex.
\bibitem{koenigsmann-defining} J. Koenigsmann,
\emph{Defining $\mathbb Z$ in $\mathbb Q$},
Annals of Mathematics \textbf{183} (2016), 73--93.
\href{https://doi.org/10.4007/annals.2016.183.1.2}{doi:10.4007/annals.2016.183.1.2}.

% Proposed additions only. No existing bibliography entry is replaced.
\bibitem{GonzalezRotger2006} J. Gonz\'alez and V. Rotger,
\emph{Non-elliptic Shimura curves of genus one},
Journal of the Mathematical Society of Japan \textbf{58} (2006), 927--948,
Theorem 4.2 and Table 2, pp.~938--939; Proposition 5.2, pp.~940--941.
\url{https://doi.org/10.2969/jmsj/1179759530}.

\bibitem{Yuan2024} X. Yuan,
\emph{Explicit Kodaira--Spencer maps over Shimura curves},
Acta Mathematica Sinica, Chinese Series \textbf{67} (2024), 227--249.
\url{https://doi.org/10.12386/A20220154}.
The published locators used here are Section 2.1, pp.~230--231;
Theorem 1.1, p.~229; and the metric definitions, pp.~228--229.
The English translation dated April 30, 2024 is retained for comparison:
\url{https://faculty.bicmr.pku.edu.cn/~yxy/preprints/explicit_KS.pdf}.
Its proof of Lemma 2.1 is longer than the published proof; that integral
coarse-model lemma is not an input to our direct stack calculation.
The source comparison and exact consumer distinction are recorded in
\texttt{research/version14/proofs/yuan-final-source-v1/PROOF.md}.

\bibitem{Deligne1979} P. Deligne,
\emph{Vari\'et\'es de Shimura: interpr\'etation modulaire, et techniques
de construction de mod\`eles canoniques},
in \emph{Automorphic Forms, Representations and $L$-Functions},
Proceedings of Symposia in Pure Mathematics \textbf{33}, Part 2 (1979),
247--289, Sections 2.2.5 and 2.3.1, pp.~269--270.
The cited statements were compared with the original French published scan:
\url{https://publications.ias.edu/sites/default/files/34_VarietesdeShimura.pdf}.
English translation by J. S. Milne, corrected September 6, 2020, p.~29:
\url{https://www.jmilne.org/math/Documents/DeligneSV.pdf}.

\bibitem{GaudronRemond2023} \'E. Gaudron and G. R\'emond,
\emph{Nouveaux th\'eor\`emes d'isog\'enie},
M\'emoires de la Soci\'et\'e Math\'ematique de France \textbf{176} (2023),
vi+123 pp., Section 1.5 and Theorem 1.9(1), pp.~7--8.
\url{https://doi.org/10.24033/msmf.484}.

\bibitem{vonKaenel2021} R. von K\"anel,
\emph{The effective Shafarevich conjecture for abelian varieties of
$\mathrm{GL}_2$-type}, Forum of Mathematics, Sigma \textbf{9} (2021), e39,
1--29, equations (2.2) and (5.9), pp.~5 and 13.
\url{https://doi.org/10.1017/fms.2021.29}.

\bibitem{voight2021} J. Voight, \emph{Quaternion algebras},
Graduate Texts in Mathematics \textbf{288}, Springer, 2021,
Chapter 7, Corollary 7.1.5 and its proof, p.~97, with the first-edition
errata of August 5, 2026, p.~2, item~19;
Chapter 43, equations (43.1.2)--(43.1.3), p.~800, and
Main Theorem 43.6.14, p.~819; proof pp.~822--823.
\url{https://doi.org/10.1007/978-3-030-56694-4_7};
\url{https://doi.org/10.1007/978-3-030-56694-4_43}.
Errata: \url{https://jvoight.github.io/quat-errata.pdf}.

\bibitem{calegari-icm2022} F. Calegari,
\emph{Reciprocity in the Langlands program since Fermat's Last Theorem},
in \emph{Proceedings of the International Congress of Mathematicians 2022},
Vol.~2, EMS Press, 2023, 610--651, Section 7, p.~624;
definition of geometric representations, pp.~611--612.
\url{https://doi.org/10.4171/ICM2022/101}.

\bibitem{diamond-im} F. Diamond and J. Im,
\emph{Modular forms and modular curves}, in Seminar on Fermat's Last Theorem,
CMS Conference Proceedings 17 (1995), 39--133, \S12.5.
\bibitem{baker-et-al-finiteness} M. H. Baker, E. Gonz\'alez-Jim\'enez,
J. Gonz\'alez and B. Poonen, \emph{Finiteness results for modular curves
of genus at least 2}, American Journal of Mathematics 127 (2005),
1325--1387, \S\S3.1, 3.2.1.
\bibitem{mascot-moduli} N. Mascot, \emph{Moduli-friendly Eisenstein series
over the p-adics and the computation of modular Galois representations},
Research in Number Theory 8 (2022), article 48,
\href{https://doi.org/10.1007/s40993-022-00329-6}{doi:10.1007/s40993-022-00329-6},
\S2.1, (2.1.8).
\bibitem{pheidas2000} T. Pheidas,
\emph{An effort to prove that the existential theory of $\Q$ is undecidable},
in \emph{Hilbert's Tenth Problem: Relations with Arithmetic and Algebraic Geometry},
J. Denef, L. Lipshitz, T. Pheidas and J. Van Geel, eds.,
Contemporary Mathematics \textbf{270}, American Mathematical Society,
Providence, RI, 2000, 237--252.
\url{https://doi.org/10.1090/conm/270/04376}.
\bibitem{pan-p5-sw99-project}
Versioned project proof, \emph{Prime-five auxiliary cohomology and element
selection}, version 1. Authored proof, supporting sources and bounded
review records at
\href{https://github.com/elliotglazer/diophantine-z-q/blob/main/research/version14/proofs/pan-p5-sw99-v1/README.md}
{\texttt{research/version14/proofs/pan-p5-sw99-v1/}}.

\bibitem{pan-local-rings-project}
Versioned project proof, \emph{Local matrix connections for Pan at prime five},
version 1. Exact source qualifications, integral proof and bounded
independent review at
\href{https://github.com/elliotglazer/diophantine-z-q/blob/main/research/version14/proofs/pan-local-rings-v1/README.md}
{\texttt{research/version14/proofs/pan-local-rings-v1/}}.

\bibitem{pan-utility-project}
Versioned project proof, \emph{Four utility applications in Pan},
version 1. Proof, exact source qualifications and bounded review at
\href{https://github.com/elliotglazer/diophantine-z-q/blob/main/research/version14/proofs/pan-utility-p5-v1/README.md}
{\texttt{research/version14/proofs/pan-utility-p5-v1/}}.

\bibitem{cj-cm-inputs-project}
Project proof supplement, \emph{Retained CM inputs for the conductor-one comparison},
version 1, \href{https://github.com/elliotglazer/diophantine-z-q/blob/main/research/version14/proofs/cj-cm-inputs-v1/README.md}{proof and source contracts}.
\bibitem{cj-good-reduction-project}
Project proof supplement, \emph{Good-reduction specialization}, version 1,
\href{https://github.com/elliotglazer/diophantine-z-q/blob/main/research/version14/proofs/cj-good-reduction-v1/README.md}{proof and independent scoped review}.

\bibitem{milne1986-abelian} J. S. Milne, \emph{Abelian varieties},
in \emph{Arithmetic Geometry}, Springer, 1986, 103--150,
Proposition 12.1 and Lemma 12.2, p.~122;
Theorem 8.2 and Remark 8.4, pp.~115--116.
\url{https://www.jmilne.org/math/articles/1986b.pdf}.

\bibitem{faltings1986-complements} G. Faltings,
\emph{Complements to Mordell}, in \emph{Rational Points},
second edition, Vieweg, 1986, Theorem 1(a)--(b), p.~204.
\url{https://doi.org/10.1007/978-3-663-06812-9}.

\bibitem{washington1997}
L.~C. Washington,
\emph{Introduction to Cyclotomic Fields}, second edition,
Graduate Texts in Mathematics 83, Springer, New York, 1997.
Chapter 14, Theorem 14.1, p.~321.
\url{https://doi.org/10.1007/978-1-4612-1934-7_14}.

\bibitem{serre-tate-good-reduction}
J.-P. Serre and J. Tate,
\emph{Good reduction of abelian varieties},
Ann. of Math. (2) \textbf{88} (1968), 492--517.
\url{https://doi.org/10.2307/1970722}.

\bibitem{carayol1987}
H. Carayol,
\emph{Formes modulaires et repr\'esentations $\ell$-adiques},
Ast\'erisque \textbf{147--148} (1987), 33--47.
\url{https://www.numdam.org/item/AST_1987__147-148__33_0/}.

\bibitem{v15-height-Stacks}
The Stacks Project Authors, \emph{The Stacks Project},
Tags~00R4, 0AWQ, 0AVT, 0BUL, 0C0I, 0BW0, and 0BWB,
\url{https://stacks.math.columbia.edu}.


\bibitem{v15-height-Javanpeykar}
A.~Javanpeykar, with an appendix by P.~Bruin,
\emph{Polynomial bounds for Arakelov invariants of Belyi curves},
Algebra \& Number Theory \textbf{8} (2014), 89--140,
\url{https://doi.org/10.2140/ant.2014.8.89}.


\bibitem{v15-height-Serre}
J.-P.~Serre, \emph{Local Fields}, Graduate Texts in Mathematics
\textbf{67}, Springer, 1979, Chapter~III,
\url{https://doi.org/10.1007/978-1-4757-5673-9}.


\bibitem{v15-height-vonKaenel}
R.~von K\"anel, \emph{The effective Shafarevich conjecture for abelian
varieties of $\mathrm{GL}_2$-type}, Forum of Mathematics, Sigma
\textbf{9} (2021), e39,
\url{https://doi.org/10.1017/fms.2021.29}.


\bibitem{v15-height-Remond}
G.~R\'emond, \emph{Propri\'et\'es de la hauteur de Faltings},
Annali della Scuola Normale Superiore di Pisa, Classe di Scienze
\textbf{23} (2022), 1589--1596,
\url{https://doi.org/10.2422/2036-2145.202010_062}.


\bibitem{v15-height-GaudronRemond}
\'E.~Gaudron and G.~R\'emond,
\emph{Nouveaux th\'eor\`emes d'isog\'enie},
M\'emoires de la Soci\'et\'e Math\'ematique de France
\textbf{176} (2023), Theorem~1.9(1),
\url{https://doi.org/10.24033/msmf.484}.


\bibitem{v15-aggregate-VonBuhren}
J.~Von Buhren, \emph{Borne de hauteur semi-effective pour le
probl\`eme de Mordell--Lang dans une vari\'et\'e ab\'elienne},
Journal de Th\'eorie des Nombres de Bordeaux
\textbf{29} (2017), no.~1, 307--320.
\href{https://doi.org/10.5802/jtnb.981}{doi:10.5802/jtnb.981}.
\href{https://www.numdam.org/item/10.5802/jtnb.981.pdf}{Full text}.


\bibitem{v15-amplification-Sun} Z.-W. Sun,
\emph{86 unknowns for Hilbert's tenth problem over the rationals},
arXiv:2610.20691v2, Theorem 4.1,
\url{https://arxiv.org/abs/2610.20691v2}.

\bibitem{v15-selection-Sun} Z.-W. Sun, \emph{86 unknowns for Hilbert's tenth problem
over the rationals}, arXiv:2610.20691v2 (2026), Theorem 4.1,
Propositions 5.1 and 6.2, Sections 6.1--8.3.
\url{https://arxiv.org/html/2610.20691v2}.

\bibitem{v15-gi-2-GTZ} B. Green, T. Tao and T. Ziegler,
\emph{An inverse theorem for the Gowers $U^{s+1}[N]$-norm},
Annals of Mathematics \textbf{176} (2012), 1231--1372.

\bibitem{v15-gi-2-LP13} Project proof, \emph{A local replacement for the filtration
argument in GTZ Lemma 13.2}, draft 002, 7 November 2026.
Manifest SHA-256:
\par\noindent{\small\texttt{95771803bed9d9135430152e2e20ad450c46aaae3ef0eeba0cab33bee59fd06c}}
\par\noindent
This is an explicitly linked author proof awaiting independent acceptance.

\bibitem{v15-gi-3-GTZ}
B. Green, T. Tao and T. Ziegler,
\emph{An inverse theorem for the Gowers $U^{s+1}[N]$-norm},
Annals of Mathematics \textbf{176} (2012), 1231--1372.
The frozen final published PDF is the source for the numbered results used
here. The arXiv v5 correction motivates the auxiliary generator types but is
not used as an accepted theorem in this proof.

\bibitem{v15-gi-5-GTZ}[GTZ] B. Green, T. Tao and T. Ziegler,
An inverse theorem for the Gowers $U^{s+1}[N]$-norm,
\emph{Annals of Mathematics} 176 (2012), 1231--1372.
DOI: 10.4007/annals.2012.176.2.11.

\bibitem{v15-gi-5-GT12}[GT12] B. Green and T. Tao,
The quantitative behaviour of polynomial orbits on nilmanifolds,
\emph{Annals of Mathematics} 175 (2012), 465--540.
DOI: 10.4007/annals.2012.175.2.2.

\bibitem{v15-gi-5-GT14}[GT14] B. Green and T. Tao,
On the quantitative distribution of polynomial nilsequences---erratum,
\emph{Annals of Mathematics} 179 (2014), 1175--1183.
DOI: 10.4007/annals.2014.179.3.8.

\bibitem{v15-gi-5-GT08}[GT08] B. Green and T. Tao,
An inverse theorem for the Gowers $U^3(G)$ norm,
\emph{Proceedings of the Edinburgh Mathematical Society} 51 (2008), 73--153.
DOI: 10.1017/S0013091505000325.

\bibitem{v15-gi-6-GT14}[GT14] B. Green and T. Tao, On the quantitative distribution
of polynomial nilsequences---erratum, Ann. of Math. 179 (2014), 1175--1183.
DOI: 10.4007/annals.2014.179.3.8.

\bibitem{v15-gi-6-GT12}[GT12] B. Green and T. Tao, The quantitative behaviour of
polynomial orbits on nilmanifolds, Ann. of Math. 175 (2012), 465--540.
DOI: 10.4007/annals.2012.175.2.2.

\bibitem{v15-gi-6-GT10}[GT10] B. Green and T. Tao, Linear equations in primes,
Ann. of Math. 171 (2010), 1753--1850. DOI: 10.4007/annals.2010.171.1753.

\bibitem{v15-gi-6-GTZ}[GTZ] B. Green, T. Tao and T. Ziegler, An inverse theorem for
the Gowers $U^{s+1}[N]$-norm, Ann. of Math. 176 (2012), 1231--1372.
DOI: 10.4007/annals.2012.176.2.11.

\bibitem{v15-gi-8-GTZ} B. Green, T. Tao and T. Ziegler,
An inverse theorem for the Gowers $U^{s+1}[N]$-norm,
Annals of Mathematics 176 (2012), 1231--1372.
\url{https://doi.org/10.4007/annals.2012.176.2.11}.

\bibitem{v15-gi-8-Tointon} M. C. H. Tointon,
Introduction to Approximate Groups, LMS Student Texts 94,
Cambridge University Press, 2020, \S5.4, pp.~96--98.
\url{https://doi.org/10.1017/9781108652865}.
The author-posted published-book PDF is pinned in the source manifest;
its author page lists corrections, none affecting this interface.

\bibitem{v15-gi-8-LP13} Project proof by /root/astra\_v14\_baseline\_theta,
Lemma 13.2 repair, version lemma13-2-002, 7 November 2026.
Exact manifest and proof hashes are in the accompanying source-use record.

\bibitem{v15-gi-8-TAUT} Same author, Tautification and the shifted nilcharacter
construction, version tautification-003, 7 November 2026.

\bibitem{v15-gi-8-Err} B. Green, T. Tao and T. Ziegler,
Erratum for the inverse theorem, six-page author document,
pp.~5--6. Attribution and problem identification only; no correction theorem
is imported from this unpublished document.

\bibitem{v15-gi-9-GTZ} B.~Green, T.~Tao and T.~Ziegler,
\emph{An inverse theorem for the Gowers $U^{s+1}[N]$-norm},
Annals of Mathematics 176 (2012), 1231--1372,
Section 12 and Corollary B.4.
\url{https://doi.org/10.4007/annals.2012.176.2.11}.

\bibitem{v15-gi-9-LP04} Project proof by /root/astra\_v14\_assertions\_arithmetic,
\emph{GTZ-LP-04: the petal quotient and degree--rank construction},
lp04-001. Its exact manifest and proof hashes are recorded in the
accompanying source-use record. Its dependencies and independent-review
status remain separate.

\bibitem{v15-gi-11-GTZ} B. Green, T. Tao and T. Ziegler,
An inverse theorem for the Gowers $U^{s+1}[N]$-norm,
Annals of Mathematics 176 (2012), 1231--1372.
DOI: 10.4007/annals.2012.176.2.11. Authoritative final PDF pinned in contracts.

\bibitem{v15-gi-11-TAUT} Project proof, Tautification and the shifted nilcharacter
construction, project author theta,
tautification-003, Theorem 1 and Corollary 2. Pending independent review.

\bibitem{v15-gi-11-LP13} Project proof, local Lemma 13.2, author
theta, lemma13-2-002, Section 2.
Pending independent review. Exact version pinned in contracts.

\bibitem{v15-gi-14-GTZ} B. Green, T. Tao and T. Ziegler,
\emph{An inverse theorem for the Gowers $U^{s+1}[N]$-norm},
Annals of Mathematics \textbf{176} (2012), 1231--1372.
Only the specified algebraic polynomial-map and Lie-group facts are used;
the Appendix C step claim is replaced here.

\bibitem{v15-gi-14-Taut} Project proof, \emph{Tautification and the shifted
nilcharacter construction}, draft 003, 7 November 2026.
The exact author manifest and independent scope decision are pinned in
the accompanying application record.

\bibitem{v15-weil-poisson} A. Weil, \emph{Sur certains groupes d'op'erateurs unitaires}, Acta Mathematica \textbf{111} (1964), 143--211, I.19, Theorem 4 and I.20.
\end{thebibliography}

% End references.tex

\clearpage
\section*{Version 15 provenance}
The historical 233-variable system uses the three-residue 237-variable scalar criterion, the two-reference-identity reduction to 235 variables, and the corrected reference-square elimination to 233 variables. The full sparse presentation has 207 quadratic residuals and 2514 monomials. Its exact source IR has SHA256
\begin{center}\small\ttfamily 10fe99b5ee41b8ed7b602fd9e8bd759f8640cd18b37326ce12ca62832320d702\end{center}
These reductions and finite transformations were developed and internally reviewed by separate Astra agents. The present manuscript integrates those arguments; its companion audit records source versions, equation generation, and subsequent review separately. Internal model review is not independent human peer review or formal verification. The earlier complete development follows with its provenance preserved.
\section*{The 238-variable predecessor: provenance}
The 238-variable development begins with Jos\'e Brox's 9 November 2026 version-17 sources, pinned to repository commit \texttt{d25e0aa6046662d9a878755aca8ed742821d3e61}. Its height, aggregate and local-construction arguments have been incorporated into this manuscript with internal references. Astra agents authored and independently reviewed the scalar-S replacement, root-product charts, reciprocal scale, global scalar rescaling, sparse transformations, exact expansions and source matching, under Elliot Glazer's direction. These are internal model reviews, not independent human peer review or formal verification. The companion dependency map preserves exact source hashes and transformation provenance. Earlier provenance follows unchanged.
The Version 9 odd-product and asymmetric-depth improvements were developed
by Gabriel Conant with ChatGPT and incorporated with additional source,
convention, and finite-certificate checks.
\section*{AI usage statement}
This manuscript was developed primarily by OpenAI Codex agents (Astra 6 and Sol 6.1)
and Anthropic Claude agents (Opus 5.5), with auditing from Fable 5.1 and 5.2.
Their contributions included mathematical
exploration, construction and analysis of the polynomial, proof drafting,
literature searches, computational checks, internal review, and drafting
and revision of the exposition. Reviews by multiple AI agents formed part
of this process; they do not constitute independent human peer review.
\end{document}
